\documentclass[default]{subfiles}

\addbibresource{cite.bib}

\begin{document}

{%
\selectlanguage{english}

\graphicspath{{image/}}

\udc{004.729, 519.872}

\pacs{%
  89.20.Ff, %Computer science and technology
  02.60.Pn, %Numerical optimization
  07.05.Tp %Computer modeling and simulation
}

\edn{xxxxxx}

\documentType{Research article}

\title{Comparative Analysis of Active Queue Management Algorithms: RED, ARED, Gentle~RED}

\author[rudn]{Abuli M. Palagashvili}

\address[rudn]{RUDN University, 6 Miklukho-Maklaya St, Moscow, 117198, Russian Federation}

\begin{authordescription}
  \authorfull[english]{Palagashvili, Abuli M.}
  \authorpost[english]{PhD student of Department of Applied Probability and Informatics of Peoples' Friendship University of Russia (RUDN University)}
  \email{1042250185@rudn.ru}
  \orcid{0009-0007-7009-7038}
\end{authordescription}

\abstracts{%
  This paper presents a systematic comparative analysis of three
  Active Queue Management (AQM) algorithms from the RED family:
  RED (Random Early Detection), ARED (Adaptive RED), and Gentle~RED.
  A dumbbell-topology network model is implemented in the ns-3
  simulator with ten TCP bulk-send sources creating sustained
  congestion at a bottleneck link. All three algorithms are evaluated
  under identical conditions (same queue thresholds, load profile,
  and topology). An extended set of metrics is proposed, including
  EWMA zone distribution and temporal drop structure. Results show
  that while all algorithms achieve the same average throughput
  ($\sim$90\% link utilization), they exhibit fundamentally different
  buffer management dynamics: RED provides the most predictable
  throughput (CV~=~18.7\%) at the cost of highest drop count;
  ARED minimizes drops ($-19\%$ vs.\ RED) but produces the largest
  queue oscillations (p99~=~49~packets); Gentle~RED achieves the
  lowest average queue length (6.40~packets) with the highest
  throughput variability (CV~=~24.7\%). Practical recommendations
  for algorithm selection based on scenario priorities are formulated.
}

\keywords{active queue management, AQM, RED, Adaptive RED, Gentle RED,
  congestion control, ns-3, network simulation, bufferbloat}

\alttitle{Сравнительный анализ алгоритмов активного управления очередями: RED, ARED, Gentle~RED}

\altauthor[rudn]{А. М. Палагашвили}

\altaddress[rudn]{Российский университет дружбы народов,
  ул. Миклухо-Маклая, д.6, Москва, 117198, Российская Федерация}

\altabstracts{%
  В данной работе представлен систематический сравнительный анализ
  трёх алгоритмов активного управления очередями (AQM) семейства
  RED: RED (Random Early Detection), ARED (Adaptive RED) и
  Gentle~RED. Реализована симуляционная модель сети с топологией
  <<гантель>> в среде ns-3 с десятью TCP-источниками, создающими
  устойчивую перегрузку узкого канала. Все три алгоритма
  оцениваются в строго идентичных условиях (единые пороги очереди,
  профиль нагрузки и топология). Предложен расширенный набор метрик
  оценки, включающий зональное распределение EWMA-средней и
  временну́ю структуру сбросов. Результаты показывают, что при
  одинаковой средней пропускной способности ($\sim$90\% утилизации
  канала) алгоритмы принципиально различаются по динамике
  управления буфером: RED обеспечивает наиболее предсказуемую
  пропускную способность (CV~=~18,7\%) ценой наибольшего числа
  сбросов; ARED минимизирует сбросы ($-19\%$ к RED), но порождает
  наибольшие осцилляции очереди (p99~=~49~пакетов); Gentle~RED
  достигает наименьшей средней очереди (6,40~пакета) при наибольшей
  вариабельности пропускной способности (CV~=~24,7\%).
  Сформулированы практические рекомендации по выбору алгоритма в
  зависимости от приоритетных требований сценария.}

\altkeywords{активное управление очередями, AQM, RED, Adaptive RED,
  Gentle RED, управление перегрузкой, ns-3, сетевое моделирование,
  bufferbloat}


\maketitle

\section{Introduction}
\label{sec:intro}

The continuous growth of Internet traffic, driven by the development of cloud computing, video streaming, and 5G/6G networks, imposes increasingly stringent requirements on congestion management mechanisms in network nodes. The traditional approach---tail-drop, in which packets are discarded only when the buffer is completely full---leads to global synchronization of TCP flows, sharp oscillations of throughput, and the bufferbloat effect~\cite{gettys2012}. The bufferbloat problem---systematic buffer overflow in routers causing end-to-end delay growth by orders of magnitude---has been identified as one of the key threats to quality of service~\cite{rfc8289}.

Active Queue Management (AQM) algorithms, recommended by IETF in RFC~2309~\cite{rfc2309}, solve this problem by preemptive dropping or marking of packets before buffer overflow. The RED (Random Early Detection) family of algorithms remains the most widely deployed class of AQM mechanisms: RED is implemented in all major network operating systems (Cisco IOS, Linux tc, FreeBSD ALTQ), and its variants---ARED and Gentle RED---are shipped as standard modules in the ns-3 simulator and the Linux kernel. Despite the emergence of new algorithms (CoDel~\cite{nichols2012}, PIE~\cite{pan2013}), the choice of a specific RED variant for specific operating conditions still requires experimental justification, since analytical models do not account for all aspects of interaction with heterogeneous TCP traffic.

\subsection{Literature review}

The RED algorithm was proposed by Floyd and Jacobson in 1993~\cite{floyd1993} and became the first practical AQM mechanism. The authors demonstrated that probabilistic early dropping of packets based on the exponentially weighted moving average (EWMA) of queue length avoids global synchronization of TCP flows and reduces average buffer delay compared to tail-drop.

Practical deployment of RED revealed its main drawback---high sensitivity to parameter tuning. Christiansen et~al.~\cite{christiansen2001} showed that RED performance significantly depends on the ratio of $\mathrm{min}_{\mathrm{Th}}$, $\mathrm{max}_{\mathrm{Th}}$, and $\mathrm{maxP}$, with optimal values changing as the number of flows and load level change. To overcome this limitation, Floyd et~al.~\cite{floyd2001} proposed ARED (Adaptive RED), in which the $\mathrm{maxP}$ parameter automatically adapts to the current load based on the deviation of the average queue from the target range.

In parallel, Floyd~\cite{floyd2000} proposed the Gentle RED modification, which eliminates the sharp transition of drop probability from $\mathrm{maxP}$ to~$1$ at $\bar{q} = \mathrm{max}_{\mathrm{Th}}$. Gentle RED introduces an additional linear zone $[\mathrm{max}_{\mathrm{Th}},\, 2\cdot\mathrm{max}_{\mathrm{Th}})$, in which the probability smoothly increases to unity, reducing TCP flow synchronization and decreasing delay variance during short-term bursts.

A significant contribution to mathematical modeling of TCP and RED interaction was made by Korolkova~A.V.~\cite{korolkova2010}, who constructed a mathematical model of the traffic transmission process with dynamic flow intensity regulated by a RED-type algorithm. In her work, using the apparatus of stochastic differential equations with a Poisson process, an autonomous system of three ordinary differential equations was formalized for the first time, describing the joint dynamics of TCP window size, instantaneous queue length, and EWMA average. A method was proposed for determining the region of self-oscillation occurrence, and qualitative analysis of the model was performed using RED, GRED, DSRED, and ARED algorithms in the NS-2 simulator.

The development of this direction was continued by Velieva~T.R.~\cite{velieva2019}, who investigated nonlinear phenomena and boundaries of the self-oscillation region in continuous-discrete dynamical models. The key contribution of this work was the adaptation of the Krylov--Bogolyubov method and harmonic linearization method to the continuous-discrete TCP/RED model, which made it possible to analytically determine the amplitude-frequency characteristics of self-oscillations and the boundaries of their occurrence region in the parameter space.

Among alternative approaches to AQM, the BLUE algorithm~\cite{feng2002}, which controls drop probability based on loss events and link idle events, and the family of control-theory-based algorithms---the PI controller of Hollot et~al.~\cite{hollot2001} and its development PIE~\cite{pan2013}---should be noted. The CoDel algorithm~\cite{nichols2012} uses a fundamentally different approach, tracking the sojourn time of a packet in the queue instead of queue length.

Thus, despite extensive research on individual aspects of the RED family algorithms---both analytical (stability, self-oscillations) and experimental---a systematic three-way comparison of RED, ARED, and Gentle~RED under strictly identical conditions with an extended set of metrics (including temporal drop structure and zone distribution of the queue) is insufficiently represented in the literature.

\subsection{Research objectives}

The \textbf{goal} of this work is to conduct a comparative analysis of three algorithms from the RED family (RED, ARED, Gentle~RED) under conditions of sustained congestion of a network node and to identify quantitative differences in buffer management dynamics, throughput stability, and the nature of congestion signals.

To achieve this goal, the following \textbf{tasks} are set:
\begin{enumerate}
  \item Implement a simulation model of a network with dumbbell topology in the ns-3 environment with the ability to switch the AQM algorithm while maintaining other conditions.
  \item Conduct a series of experiments with identical queue parameters and load for each of the three algorithms.
  \item Compare algorithms by a set of metrics: mean and tail queue length, coefficient of variation of throughput, number and temporal structure of packet drops.
  \item Formulate practical recommendations for AQM algorithm selection depending on the priority requirements of the scenario.
\end{enumerate}

\subsection{Scientific novelty}

The scientific novelty of this work consists of the following:
\begin{enumerate}
  \item Unlike the works of Floyd and Jacobson~\cite{floyd1993}, Floyd~\cite{floyd2000, floyd2001}, and Christiansen et~al.~\cite{christiansen2001}, where each algorithm was analyzed in isolation or in pairwise comparison with tail-drop, this work conducts a systematic three-way comparison of RED, ARED, and Gentle~RED under strictly identical conditions (unified topology, identical thresholds, identical load profile), which eliminates the influence of methodological differences on results.
  \item Unlike the analytical approach of Korolkova~\cite{korolkova2010}, who studied self-oscillation modes using ODE systems, and Velieva~\cite{velieva2019}, who applied the harmonic linearization method to determine stability boundaries, this work proposes an extended set of \textit{experimental} evaluation metrics, including zone distribution of EWMA average operating time and temporal structure of drops (median interval, clustering), which allows characterizing not only static indicators but also the dynamics of control actions in the simulation model.
  \item Unlike the works of Hollot et~al.~\cite{hollot2001} and Feng et~al.~\cite{feng2002}, focused on comparing RED with fundamentally different AQM approaches (PI controller, BLUE), this work focuses on \textit{intra-family} comparison of RED variants and for the first time quantitatively describes the fundamental trade-off between throughput predictability (RED, CV~$= 18.7\%$), drop minimization (ARED, $-19\%$), and buffer delay minimization (Gentle~RED, average queue 6.40~pkt).
\end{enumerate}

\subsection{RED family algorithms}

Active queue management algorithms begin dropping or marking packets \textit{preemptively}---before the buffer overflows---based on average queue length statistics. TCP connections receive congestion signals earlier and reduce their sending rate more smoothly, preventing collapse.

\subsubsection{RED (Random Early Detection)}

RED is the base AQM algorithm proposed by Floyd and Jacobson in 1993~\cite{floyd1993}. The algorithm computes the exponentially weighted moving average of queue length (EWMA) and makes drop decisions based on it.

\textbf{Average queue update} on each arriving packet:
\begin{equation}
\bar{q} = (1 - w_q)\,\bar{q} + w_q\,q
\end{equation}
where $q$ is the instantaneous queue length, $w_q \in (0,1)$ is the filter weight.

\textbf{Drop probability} is determined by three zones:
\begin{equation}
p(\bar{q}) = \begin{cases}
0 & \bar{q} < \mathrm{min}_{\mathrm{Th}} \\
\mathrm{maxP} \cdot \dfrac{\bar{q} - \mathrm{min}_{\mathrm{Th}}}{\mathrm{max}_{\mathrm{Th}} - \mathrm{min}_{\mathrm{Th}}} & \mathrm{min}_{\mathrm{Th}} \le \bar{q} < \mathrm{max}_{\mathrm{Th}} \\
1 & \bar{q} \ge \mathrm{max}_{\mathrm{Th}}
\end{cases}
\end{equation}

The hard forced drop above $\mathrm{max}_{\mathrm{Th}}$ ensures fast reaction but synchronizes TCP flows and creates sawtooth queue oscillations.

\subsubsection{ARED (Adaptive RED)}

ARED extends RED with an adaptive mechanism~\cite{floyd2001}: the $\mathrm{maxP}$ parameter is periodically recomputed every $T_{\mathrm{interval}}$ depending on the deviation of the average queue from the target range:
\begin{equation}
\mathrm{maxP} = \begin{cases}
\mathrm{maxP} + \alpha & \bar{q} > \dfrac{\mathrm{min}_{\mathrm{Th}} + \mathrm{max}_{\mathrm{Th}}}{2} \\[6pt]
\mathrm{maxP} - \beta & \bar{q} < \dfrac{\mathrm{min}_{\mathrm{Th}} + \mathrm{max}_{\mathrm{Th}}}{2}
\end{cases}
\end{equation}
where $\alpha$ and $\beta$ are increment and decrement coefficients ($\alpha > \beta$). The value of $\mathrm{maxP}$ is bounded within an admissible range. This allows the algorithm to adapt to changing load without manual tuning; however, the adaptation introduces additional inertia and may lead to wide queue oscillations.

\subsubsection{Gentle RED}

Gentle RED~\cite{floyd2000} modifies RED behavior in the zone above $\mathrm{max}_{\mathrm{Th}}$: instead of immediate forced drop, the algorithm continues the probabilistic approach up to $2 \cdot \mathrm{max}_{\mathrm{Th}}$:
\begin{equation}
p(\bar{q}) = \begin{cases}
0 & \bar{q} < \mathrm{min}_{\mathrm{Th}} \\[4pt]
\mathrm{maxP} \cdot \dfrac{\bar{q} - \mathrm{min}_{\mathrm{Th}}}{\mathrm{max}_{\mathrm{Th}} - \mathrm{min}_{\mathrm{Th}}}
  & \mathrm{min}_{\mathrm{Th}} \le \bar{q} < \mathrm{max}_{\mathrm{Th}} \\[8pt]
\mathrm{maxP} + (1 - \mathrm{maxP}) \cdot \dfrac{\bar{q} - \mathrm{max}_{\mathrm{Th}}}{\mathrm{max}_{\mathrm{Th}}}
  & \mathrm{max}_{\mathrm{Th}} \le \bar{q} < 2\,\mathrm{max}_{\mathrm{Th}} \\[4pt]
1 & \bar{q} \ge 2\,\mathrm{max}_{\mathrm{Th}}
\end{cases}
\end{equation}

In the zone $[\mathrm{max}_{\mathrm{Th}},\, 2\cdot\mathrm{max}_{\mathrm{Th}})$ the probability linearly grows from $\mathrm{maxP}$ to~$1$, and forced drop occurs only at $\bar{q} \ge 2\cdot\mathrm{max}_{\mathrm{Th}}$. This softens the reaction to short-term bursts, reduces TCP flow synchronization, and decreases the average queue length---at the cost of somewhat greater throughput variability.

\section{Methods}
\label{sec:methods}

\subsection{Experimental topology}

The experiment uses a dumbbell topology: multiple senders are connected to a common router through fast links, and the router is connected to a single receiver through a narrow bottleneck link with limited bandwidth. The narrow link is the only possible congestion point, and the AQM queue discipline is installed on its outgoing interface.

\begin{verbatim}
  Sender 0  --\
  Sender 1  --|
  Sender 2  --|  1 Mbit/s         2 Mbit/s
      ...     |----------> [Router] ----------> Receiver
  Sender 8  --|  delay 2 ms      delay 10 ms
  Sender 9  --/                ^
                          AQM queue
                       (RED / ARED / Gentle)
\end{verbatim}

Each of the ten senders establishes one TCP connection with the receiver and continuously transmits data (BulkSend). The total incoming load on the router potentially reaches $10 \times 1$~Mbit/s~$= 10$~Mbit/s, which is 5~times the bottleneck link capacity. This ensures sustained congestion and allows evaluation of AQM algorithm behavior under real buffer saturation conditions.

Topology parameters are listed in Table~\ref{tab:topology}.

\begin{table}[H]
\centering
\caption{Topology and simulation parameters}
\label{tab:topology}
\begin{tabular}{lll}
\toprule
Parameter & Value & Description \\
\midrule
\texttt{nSenders} & 10 & Number of TCP senders \\
\texttt{fastRate} & 1 Mbit/s & Sender $\to$ router link speed (delay 2~ms) \\
\texttt{bottleneckRate} & 2 Mbit/s & Bottleneck link speed (router $\to$ receiver) \\
\texttt{bottleneckDelay} & 10 ms & Bottleneck link delay \\
\texttt{simTime} & 60.0 & Simulation duration (s) \\
Transport & TCP BulkSend & Segment size 512 bytes \\
\bottomrule
\end{tabular}
\end{table}

\subsection{AQM configuration}

All three algorithms are launched with identical threshold parameters (Table~\ref{tab:aqm-config}).

\begin{table}[H]
\centering
\caption{Unified threshold parameters for RED, ARED, and Gentle RED}
\label{tab:aqm-config}
\begin{tabular}{lll}
\toprule
Parameter & Value & Description \\
\midrule
\texttt{minTh} & 5 packets & Lower average queue threshold---no drops below \\
\texttt{maxTh} & 10 packets & Upper threshold; above it RED/ARED drop all packets \\
\texttt{MaxSize} & 100 packets & Hard buffer limit \\
\bottomrule
\end{tabular}
\end{table}

\subsection{Instrumentation}

The discrete-event network simulator \textbf{ns-3} version~3.41 was used. Simulation scripts are placed in the \texttt{scratch/} directory. For this experiment, a subdirectory was created with the following structure:

\begin{verbatim}
ns-3.41/
+-- scratch/
    +-- red-experiments/
        +-- red.cc
        +-- configs/
        |   +-- red.txt
        |   +-- ared.txt
        |   +-- gentle.txt
        +-- analysis/
            +-- vis.py
            +-- ...
\end{verbatim}

Configuration files specify the AQM type and queue parameters. Example \texttt{configs/red.txt}:

\begin{verbatim}
aqmType=red
minTh=5
maxTh=10
MaxSize=100
\end{verbatim}

Analogous files for ARED (\texttt{aqmType=ared}) and Gentle RED (\texttt{aqmType=gentle}) differ only in the \texttt{aqmType} field value.

Simulations were launched with:
\begin{verbatim}
./ns3 run "scratch/red-experiments/red.cc --config=configs/<algorithm>.txt"
\end{verbatim}

Upon completion of each simulation, three CSV files were generated:
\begin{itemize}
  \item \texttt{<algorithm>\_queue.csv}---queue length sampled every 100~ms;
  \item \texttt{<algorithm>\_drops.csv}---timestamps of each packet drop;
  \item \texttt{<algorithm>\_throughput.csv}---throughput sampled every 100~ms.
\end{itemize}

For plotting, the script \texttt{vis.py} with moving average smoothing (window of 30~samples) was used:
\begin{verbatim}
python vis.py --algorithms red ared gentle --window 30
\end{verbatim}

\subsection{Evaluation metrics}

\begin{table}[H]
\centering
\caption{Algorithm evaluation metrics}
\label{tab:metrics}
\begin{tabular}{lll}
\toprule
Metric & Data source & What it characterizes \\
\midrule
Mean and median queue length & \texttt{*\_queue.csv} & Average buffer delay \\
p90 / p99 queue length & \texttt{*\_queue.csv} & Tail delays \\
Fraction of time with empty queue & \texttt{*\_queue.csv} & Link underutilization \\
Mean throughput & \texttt{*\_throughput.csv} & Link utilization \\
Coefficient of variation (CV) & \texttt{*\_throughput.csv} & Flow stability \\
Total drop count & \texttt{*\_drops.csv} & Load on TCP congestion mechanism \\
Median inter-drop interval & \texttt{*\_drops.csv} & Nature of control signals \\
\bottomrule
\end{tabular}
\end{table}

\section{Results}
\label{sec:results}

The experiment consisted of three sequential simulations---one for each algorithm. Each algorithm was launched with identical queue parameters (Table~\ref{tab:aqm-config}) for 60~seconds (the first second is warm-up, excluded from analysis).

\subsection{Queue dynamics}

Figure~\ref{fig:queue} shows the queue length dynamics for the three algorithms over the entire simulation time.

\begin{figure}[H]
\centering
\resizebox{\textwidth}{!}{%% Creator: Matplotlib, PGF backend
%%
%% To include the figure in your LaTeX document, write
%%   \input{<filename>.pgf}
%%
%% Make sure the required packages are loaded in your preamble
%%   \usepackage{pgf}
%%
%% Also ensure that all the required font packages are loaded; for instance,
%% the lmodern package is sometimes necessary when using math font.
%%   \usepackage{lmodern}
%%
%% Figures using additional raster images can only be included by \input if
%% they are in the same directory as the main LaTeX file. For loading figures
%% from other directories you can use the `import` package
%%   \usepackage{import}
%%
%% and then include the figures with
%%   \import{<path to file>}{<filename>.pgf}
%%
%% Matplotlib used the following preamble
%%   \def\mathdefault#1{#1}
%%   \everymath=\expandafter{\the\everymath\displaystyle}
%%   \IfFileExists{scrextend.sty}{
%%     \usepackage[fontsize=10.000000pt]{scrextend}
%%   }{
%%     \renewcommand{\normalsize}{\fontsize{10.000000}{12.000000}\selectfont}
%%     \normalsize
%%   }
%%   \usepackage{fontspec}
%%   \setmainfont{TeX Gyre Termes}
%%   \setsansfont{TeX Gyre Heros}
%%   \setmonofont{TeX Gyre Cursor}
%%   \makeatletter\@ifpackageloaded{underscore}{}{\usepackage[strings]{underscore}}\makeatother
%%
\begingroup%
\makeatletter%
\begin{pgfpicture}%
\pgfpathrectangle{\pgfpointorigin}{\pgfqpoint{5.491112in}{3.494944in}}%
\pgfusepath{use as bounding box, clip}%
\begin{pgfscope}%
\pgfsetbuttcap%
\pgfsetmiterjoin%
\definecolor{currentfill}{rgb}{1.000000,1.000000,1.000000}%
\pgfsetfillcolor{currentfill}%
\pgfsetlinewidth{0.000000pt}%
\definecolor{currentstroke}{rgb}{1.000000,1.000000,1.000000}%
\pgfsetstrokecolor{currentstroke}%
\pgfsetdash{}{0pt}%
\pgfpathmoveto{\pgfqpoint{0.000000in}{0.000000in}}%
\pgfpathlineto{\pgfqpoint{5.491112in}{0.000000in}}%
\pgfpathlineto{\pgfqpoint{5.491112in}{3.494944in}}%
\pgfpathlineto{\pgfqpoint{0.000000in}{3.494944in}}%
\pgfpathlineto{\pgfqpoint{0.000000in}{0.000000in}}%
\pgfpathclose%
\pgfusepath{fill}%
\end{pgfscope}%
\begin{pgfscope}%
\pgfsetbuttcap%
\pgfsetmiterjoin%
\definecolor{currentfill}{rgb}{1.000000,1.000000,1.000000}%
\pgfsetfillcolor{currentfill}%
\pgfsetlinewidth{0.000000pt}%
\definecolor{currentstroke}{rgb}{0.000000,0.000000,0.000000}%
\pgfsetstrokecolor{currentstroke}%
\pgfsetstrokeopacity{0.000000}%
\pgfsetdash{}{0pt}%
\pgfpathmoveto{\pgfqpoint{0.586112in}{0.502778in}}%
\pgfpathlineto{\pgfqpoint{5.391112in}{0.502778in}}%
\pgfpathlineto{\pgfqpoint{5.391112in}{3.197778in}}%
\pgfpathlineto{\pgfqpoint{0.586112in}{3.197778in}}%
\pgfpathlineto{\pgfqpoint{0.586112in}{0.502778in}}%
\pgfpathclose%
\pgfusepath{fill}%
\end{pgfscope}%
\begin{pgfscope}%
\pgfpathrectangle{\pgfqpoint{0.586112in}{0.502778in}}{\pgfqpoint{4.805000in}{2.695000in}}%
\pgfusepath{clip}%
\pgfsetrectcap%
\pgfsetroundjoin%
\pgfsetlinewidth{0.803000pt}%
\definecolor{currentstroke}{rgb}{0.690196,0.690196,0.690196}%
\pgfsetstrokecolor{currentstroke}%
\pgfsetdash{}{0pt}%
\pgfpathmoveto{\pgfqpoint{0.730297in}{0.502778in}}%
\pgfpathlineto{\pgfqpoint{0.730297in}{3.197778in}}%
\pgfusepath{stroke}%
\end{pgfscope}%
\begin{pgfscope}%
\pgfsetbuttcap%
\pgfsetroundjoin%
\definecolor{currentfill}{rgb}{0.000000,0.000000,0.000000}%
\pgfsetfillcolor{currentfill}%
\pgfsetlinewidth{0.803000pt}%
\definecolor{currentstroke}{rgb}{0.000000,0.000000,0.000000}%
\pgfsetstrokecolor{currentstroke}%
\pgfsetdash{}{0pt}%
\pgfsys@defobject{currentmarker}{\pgfqpoint{0.000000in}{-0.048611in}}{\pgfqpoint{0.000000in}{0.000000in}}{%
\pgfpathmoveto{\pgfqpoint{0.000000in}{0.000000in}}%
\pgfpathlineto{\pgfqpoint{0.000000in}{-0.048611in}}%
\pgfusepath{stroke,fill}%
}%
\begin{pgfscope}%
\pgfsys@transformshift{0.730297in}{0.502778in}%
\pgfsys@useobject{currentmarker}{}%
\end{pgfscope}%
\end{pgfscope}%
\begin{pgfscope}%
\definecolor{textcolor}{rgb}{0.000000,0.000000,0.000000}%
\pgfsetstrokecolor{textcolor}%
\pgfsetfillcolor{textcolor}%
\pgftext[x=0.730297in,y=0.405556in,,top]{\color{textcolor}{\rmfamily\fontsize{10.000000}{12.000000}\selectfont\catcode`\^=\active\def^{\ifmmode\sp\else\^{}\fi}\catcode`\%=\active\def%{\%}$\mathdefault{0}$}}%
\end{pgfscope}%
\begin{pgfscope}%
\pgfpathrectangle{\pgfqpoint{0.586112in}{0.502778in}}{\pgfqpoint{4.805000in}{2.695000in}}%
\pgfusepath{clip}%
\pgfsetrectcap%
\pgfsetroundjoin%
\pgfsetlinewidth{0.803000pt}%
\definecolor{currentstroke}{rgb}{0.690196,0.690196,0.690196}%
\pgfsetstrokecolor{currentstroke}%
\pgfsetdash{}{0pt}%
\pgfpathmoveto{\pgfqpoint{1.470716in}{0.502778in}}%
\pgfpathlineto{\pgfqpoint{1.470716in}{3.197778in}}%
\pgfusepath{stroke}%
\end{pgfscope}%
\begin{pgfscope}%
\pgfsetbuttcap%
\pgfsetroundjoin%
\definecolor{currentfill}{rgb}{0.000000,0.000000,0.000000}%
\pgfsetfillcolor{currentfill}%
\pgfsetlinewidth{0.803000pt}%
\definecolor{currentstroke}{rgb}{0.000000,0.000000,0.000000}%
\pgfsetstrokecolor{currentstroke}%
\pgfsetdash{}{0pt}%
\pgfsys@defobject{currentmarker}{\pgfqpoint{0.000000in}{-0.048611in}}{\pgfqpoint{0.000000in}{0.000000in}}{%
\pgfpathmoveto{\pgfqpoint{0.000000in}{0.000000in}}%
\pgfpathlineto{\pgfqpoint{0.000000in}{-0.048611in}}%
\pgfusepath{stroke,fill}%
}%
\begin{pgfscope}%
\pgfsys@transformshift{1.470716in}{0.502778in}%
\pgfsys@useobject{currentmarker}{}%
\end{pgfscope}%
\end{pgfscope}%
\begin{pgfscope}%
\definecolor{textcolor}{rgb}{0.000000,0.000000,0.000000}%
\pgfsetstrokecolor{textcolor}%
\pgfsetfillcolor{textcolor}%
\pgftext[x=1.470716in,y=0.405556in,,top]{\color{textcolor}{\rmfamily\fontsize{10.000000}{12.000000}\selectfont\catcode`\^=\active\def^{\ifmmode\sp\else\^{}\fi}\catcode`\%=\active\def%{\%}$\mathdefault{10}$}}%
\end{pgfscope}%
\begin{pgfscope}%
\pgfpathrectangle{\pgfqpoint{0.586112in}{0.502778in}}{\pgfqpoint{4.805000in}{2.695000in}}%
\pgfusepath{clip}%
\pgfsetrectcap%
\pgfsetroundjoin%
\pgfsetlinewidth{0.803000pt}%
\definecolor{currentstroke}{rgb}{0.690196,0.690196,0.690196}%
\pgfsetstrokecolor{currentstroke}%
\pgfsetdash{}{0pt}%
\pgfpathmoveto{\pgfqpoint{2.211136in}{0.502778in}}%
\pgfpathlineto{\pgfqpoint{2.211136in}{3.197778in}}%
\pgfusepath{stroke}%
\end{pgfscope}%
\begin{pgfscope}%
\pgfsetbuttcap%
\pgfsetroundjoin%
\definecolor{currentfill}{rgb}{0.000000,0.000000,0.000000}%
\pgfsetfillcolor{currentfill}%
\pgfsetlinewidth{0.803000pt}%
\definecolor{currentstroke}{rgb}{0.000000,0.000000,0.000000}%
\pgfsetstrokecolor{currentstroke}%
\pgfsetdash{}{0pt}%
\pgfsys@defobject{currentmarker}{\pgfqpoint{0.000000in}{-0.048611in}}{\pgfqpoint{0.000000in}{0.000000in}}{%
\pgfpathmoveto{\pgfqpoint{0.000000in}{0.000000in}}%
\pgfpathlineto{\pgfqpoint{0.000000in}{-0.048611in}}%
\pgfusepath{stroke,fill}%
}%
\begin{pgfscope}%
\pgfsys@transformshift{2.211136in}{0.502778in}%
\pgfsys@useobject{currentmarker}{}%
\end{pgfscope}%
\end{pgfscope}%
\begin{pgfscope}%
\definecolor{textcolor}{rgb}{0.000000,0.000000,0.000000}%
\pgfsetstrokecolor{textcolor}%
\pgfsetfillcolor{textcolor}%
\pgftext[x=2.211136in,y=0.405556in,,top]{\color{textcolor}{\rmfamily\fontsize{10.000000}{12.000000}\selectfont\catcode`\^=\active\def^{\ifmmode\sp\else\^{}\fi}\catcode`\%=\active\def%{\%}$\mathdefault{20}$}}%
\end{pgfscope}%
\begin{pgfscope}%
\pgfpathrectangle{\pgfqpoint{0.586112in}{0.502778in}}{\pgfqpoint{4.805000in}{2.695000in}}%
\pgfusepath{clip}%
\pgfsetrectcap%
\pgfsetroundjoin%
\pgfsetlinewidth{0.803000pt}%
\definecolor{currentstroke}{rgb}{0.690196,0.690196,0.690196}%
\pgfsetstrokecolor{currentstroke}%
\pgfsetdash{}{0pt}%
\pgfpathmoveto{\pgfqpoint{2.951555in}{0.502778in}}%
\pgfpathlineto{\pgfqpoint{2.951555in}{3.197778in}}%
\pgfusepath{stroke}%
\end{pgfscope}%
\begin{pgfscope}%
\pgfsetbuttcap%
\pgfsetroundjoin%
\definecolor{currentfill}{rgb}{0.000000,0.000000,0.000000}%
\pgfsetfillcolor{currentfill}%
\pgfsetlinewidth{0.803000pt}%
\definecolor{currentstroke}{rgb}{0.000000,0.000000,0.000000}%
\pgfsetstrokecolor{currentstroke}%
\pgfsetdash{}{0pt}%
\pgfsys@defobject{currentmarker}{\pgfqpoint{0.000000in}{-0.048611in}}{\pgfqpoint{0.000000in}{0.000000in}}{%
\pgfpathmoveto{\pgfqpoint{0.000000in}{0.000000in}}%
\pgfpathlineto{\pgfqpoint{0.000000in}{-0.048611in}}%
\pgfusepath{stroke,fill}%
}%
\begin{pgfscope}%
\pgfsys@transformshift{2.951555in}{0.502778in}%
\pgfsys@useobject{currentmarker}{}%
\end{pgfscope}%
\end{pgfscope}%
\begin{pgfscope}%
\definecolor{textcolor}{rgb}{0.000000,0.000000,0.000000}%
\pgfsetstrokecolor{textcolor}%
\pgfsetfillcolor{textcolor}%
\pgftext[x=2.951555in,y=0.405556in,,top]{\color{textcolor}{\rmfamily\fontsize{10.000000}{12.000000}\selectfont\catcode`\^=\active\def^{\ifmmode\sp\else\^{}\fi}\catcode`\%=\active\def%{\%}$\mathdefault{30}$}}%
\end{pgfscope}%
\begin{pgfscope}%
\pgfpathrectangle{\pgfqpoint{0.586112in}{0.502778in}}{\pgfqpoint{4.805000in}{2.695000in}}%
\pgfusepath{clip}%
\pgfsetrectcap%
\pgfsetroundjoin%
\pgfsetlinewidth{0.803000pt}%
\definecolor{currentstroke}{rgb}{0.690196,0.690196,0.690196}%
\pgfsetstrokecolor{currentstroke}%
\pgfsetdash{}{0pt}%
\pgfpathmoveto{\pgfqpoint{3.691975in}{0.502778in}}%
\pgfpathlineto{\pgfqpoint{3.691975in}{3.197778in}}%
\pgfusepath{stroke}%
\end{pgfscope}%
\begin{pgfscope}%
\pgfsetbuttcap%
\pgfsetroundjoin%
\definecolor{currentfill}{rgb}{0.000000,0.000000,0.000000}%
\pgfsetfillcolor{currentfill}%
\pgfsetlinewidth{0.803000pt}%
\definecolor{currentstroke}{rgb}{0.000000,0.000000,0.000000}%
\pgfsetstrokecolor{currentstroke}%
\pgfsetdash{}{0pt}%
\pgfsys@defobject{currentmarker}{\pgfqpoint{0.000000in}{-0.048611in}}{\pgfqpoint{0.000000in}{0.000000in}}{%
\pgfpathmoveto{\pgfqpoint{0.000000in}{0.000000in}}%
\pgfpathlineto{\pgfqpoint{0.000000in}{-0.048611in}}%
\pgfusepath{stroke,fill}%
}%
\begin{pgfscope}%
\pgfsys@transformshift{3.691975in}{0.502778in}%
\pgfsys@useobject{currentmarker}{}%
\end{pgfscope}%
\end{pgfscope}%
\begin{pgfscope}%
\definecolor{textcolor}{rgb}{0.000000,0.000000,0.000000}%
\pgfsetstrokecolor{textcolor}%
\pgfsetfillcolor{textcolor}%
\pgftext[x=3.691975in,y=0.405556in,,top]{\color{textcolor}{\rmfamily\fontsize{10.000000}{12.000000}\selectfont\catcode`\^=\active\def^{\ifmmode\sp\else\^{}\fi}\catcode`\%=\active\def%{\%}$\mathdefault{40}$}}%
\end{pgfscope}%
\begin{pgfscope}%
\pgfpathrectangle{\pgfqpoint{0.586112in}{0.502778in}}{\pgfqpoint{4.805000in}{2.695000in}}%
\pgfusepath{clip}%
\pgfsetrectcap%
\pgfsetroundjoin%
\pgfsetlinewidth{0.803000pt}%
\definecolor{currentstroke}{rgb}{0.690196,0.690196,0.690196}%
\pgfsetstrokecolor{currentstroke}%
\pgfsetdash{}{0pt}%
\pgfpathmoveto{\pgfqpoint{4.432394in}{0.502778in}}%
\pgfpathlineto{\pgfqpoint{4.432394in}{3.197778in}}%
\pgfusepath{stroke}%
\end{pgfscope}%
\begin{pgfscope}%
\pgfsetbuttcap%
\pgfsetroundjoin%
\definecolor{currentfill}{rgb}{0.000000,0.000000,0.000000}%
\pgfsetfillcolor{currentfill}%
\pgfsetlinewidth{0.803000pt}%
\definecolor{currentstroke}{rgb}{0.000000,0.000000,0.000000}%
\pgfsetstrokecolor{currentstroke}%
\pgfsetdash{}{0pt}%
\pgfsys@defobject{currentmarker}{\pgfqpoint{0.000000in}{-0.048611in}}{\pgfqpoint{0.000000in}{0.000000in}}{%
\pgfpathmoveto{\pgfqpoint{0.000000in}{0.000000in}}%
\pgfpathlineto{\pgfqpoint{0.000000in}{-0.048611in}}%
\pgfusepath{stroke,fill}%
}%
\begin{pgfscope}%
\pgfsys@transformshift{4.432394in}{0.502778in}%
\pgfsys@useobject{currentmarker}{}%
\end{pgfscope}%
\end{pgfscope}%
\begin{pgfscope}%
\definecolor{textcolor}{rgb}{0.000000,0.000000,0.000000}%
\pgfsetstrokecolor{textcolor}%
\pgfsetfillcolor{textcolor}%
\pgftext[x=4.432394in,y=0.405556in,,top]{\color{textcolor}{\rmfamily\fontsize{10.000000}{12.000000}\selectfont\catcode`\^=\active\def^{\ifmmode\sp\else\^{}\fi}\catcode`\%=\active\def%{\%}$\mathdefault{50}$}}%
\end{pgfscope}%
\begin{pgfscope}%
\pgfpathrectangle{\pgfqpoint{0.586112in}{0.502778in}}{\pgfqpoint{4.805000in}{2.695000in}}%
\pgfusepath{clip}%
\pgfsetrectcap%
\pgfsetroundjoin%
\pgfsetlinewidth{0.803000pt}%
\definecolor{currentstroke}{rgb}{0.690196,0.690196,0.690196}%
\pgfsetstrokecolor{currentstroke}%
\pgfsetdash{}{0pt}%
\pgfpathmoveto{\pgfqpoint{5.172814in}{0.502778in}}%
\pgfpathlineto{\pgfqpoint{5.172814in}{3.197778in}}%
\pgfusepath{stroke}%
\end{pgfscope}%
\begin{pgfscope}%
\pgfsetbuttcap%
\pgfsetroundjoin%
\definecolor{currentfill}{rgb}{0.000000,0.000000,0.000000}%
\pgfsetfillcolor{currentfill}%
\pgfsetlinewidth{0.803000pt}%
\definecolor{currentstroke}{rgb}{0.000000,0.000000,0.000000}%
\pgfsetstrokecolor{currentstroke}%
\pgfsetdash{}{0pt}%
\pgfsys@defobject{currentmarker}{\pgfqpoint{0.000000in}{-0.048611in}}{\pgfqpoint{0.000000in}{0.000000in}}{%
\pgfpathmoveto{\pgfqpoint{0.000000in}{0.000000in}}%
\pgfpathlineto{\pgfqpoint{0.000000in}{-0.048611in}}%
\pgfusepath{stroke,fill}%
}%
\begin{pgfscope}%
\pgfsys@transformshift{5.172814in}{0.502778in}%
\pgfsys@useobject{currentmarker}{}%
\end{pgfscope}%
\end{pgfscope}%
\begin{pgfscope}%
\definecolor{textcolor}{rgb}{0.000000,0.000000,0.000000}%
\pgfsetstrokecolor{textcolor}%
\pgfsetfillcolor{textcolor}%
\pgftext[x=5.172814in,y=0.405556in,,top]{\color{textcolor}{\rmfamily\fontsize{10.000000}{12.000000}\selectfont\catcode`\^=\active\def^{\ifmmode\sp\else\^{}\fi}\catcode`\%=\active\def%{\%}$\mathdefault{60}$}}%
\end{pgfscope}%
\begin{pgfscope}%
\definecolor{textcolor}{rgb}{0.000000,0.000000,0.000000}%
\pgfsetstrokecolor{textcolor}%
\pgfsetfillcolor{textcolor}%
\pgftext[x=2.988612in,y=0.225000in,,top]{\color{textcolor}{\rmfamily\fontsize{10.000000}{12.000000}\selectfont\catcode`\^=\active\def^{\ifmmode\sp\else\^{}\fi}\catcode`\%=\active\def%{\%}Time (s)}}%
\end{pgfscope}%
\begin{pgfscope}%
\pgfpathrectangle{\pgfqpoint{0.586112in}{0.502778in}}{\pgfqpoint{4.805000in}{2.695000in}}%
\pgfusepath{clip}%
\pgfsetrectcap%
\pgfsetroundjoin%
\pgfsetlinewidth{0.803000pt}%
\definecolor{currentstroke}{rgb}{0.690196,0.690196,0.690196}%
\pgfsetstrokecolor{currentstroke}%
\pgfsetdash{}{0pt}%
\pgfpathmoveto{\pgfqpoint{0.586112in}{0.617029in}}%
\pgfpathlineto{\pgfqpoint{5.391112in}{0.617029in}}%
\pgfusepath{stroke}%
\end{pgfscope}%
\begin{pgfscope}%
\pgfsetbuttcap%
\pgfsetroundjoin%
\definecolor{currentfill}{rgb}{0.000000,0.000000,0.000000}%
\pgfsetfillcolor{currentfill}%
\pgfsetlinewidth{0.803000pt}%
\definecolor{currentstroke}{rgb}{0.000000,0.000000,0.000000}%
\pgfsetstrokecolor{currentstroke}%
\pgfsetdash{}{0pt}%
\pgfsys@defobject{currentmarker}{\pgfqpoint{-0.048611in}{0.000000in}}{\pgfqpoint{-0.000000in}{0.000000in}}{%
\pgfpathmoveto{\pgfqpoint{-0.000000in}{0.000000in}}%
\pgfpathlineto{\pgfqpoint{-0.048611in}{0.000000in}}%
\pgfusepath{stroke,fill}%
}%
\begin{pgfscope}%
\pgfsys@transformshift{0.586112in}{0.617029in}%
\pgfsys@useobject{currentmarker}{}%
\end{pgfscope}%
\end{pgfscope}%
\begin{pgfscope}%
\definecolor{textcolor}{rgb}{0.000000,0.000000,0.000000}%
\pgfsetstrokecolor{textcolor}%
\pgfsetfillcolor{textcolor}%
\pgftext[x=0.419445in, y=0.569598in, left, base]{\color{textcolor}{\rmfamily\fontsize{10.000000}{12.000000}\selectfont\catcode`\^=\active\def^{\ifmmode\sp\else\^{}\fi}\catcode`\%=\active\def%{\%}$\mathdefault{0}$}}%
\end{pgfscope}%
\begin{pgfscope}%
\pgfpathrectangle{\pgfqpoint{0.586112in}{0.502778in}}{\pgfqpoint{4.805000in}{2.695000in}}%
\pgfusepath{clip}%
\pgfsetrectcap%
\pgfsetroundjoin%
\pgfsetlinewidth{0.803000pt}%
\definecolor{currentstroke}{rgb}{0.690196,0.690196,0.690196}%
\pgfsetstrokecolor{currentstroke}%
\pgfsetdash{}{0pt}%
\pgfpathmoveto{\pgfqpoint{0.586112in}{1.111978in}}%
\pgfpathlineto{\pgfqpoint{5.391112in}{1.111978in}}%
\pgfusepath{stroke}%
\end{pgfscope}%
\begin{pgfscope}%
\pgfsetbuttcap%
\pgfsetroundjoin%
\definecolor{currentfill}{rgb}{0.000000,0.000000,0.000000}%
\pgfsetfillcolor{currentfill}%
\pgfsetlinewidth{0.803000pt}%
\definecolor{currentstroke}{rgb}{0.000000,0.000000,0.000000}%
\pgfsetstrokecolor{currentstroke}%
\pgfsetdash{}{0pt}%
\pgfsys@defobject{currentmarker}{\pgfqpoint{-0.048611in}{0.000000in}}{\pgfqpoint{-0.000000in}{0.000000in}}{%
\pgfpathmoveto{\pgfqpoint{-0.000000in}{0.000000in}}%
\pgfpathlineto{\pgfqpoint{-0.048611in}{0.000000in}}%
\pgfusepath{stroke,fill}%
}%
\begin{pgfscope}%
\pgfsys@transformshift{0.586112in}{1.111978in}%
\pgfsys@useobject{currentmarker}{}%
\end{pgfscope}%
\end{pgfscope}%
\begin{pgfscope}%
\definecolor{textcolor}{rgb}{0.000000,0.000000,0.000000}%
\pgfsetstrokecolor{textcolor}%
\pgfsetfillcolor{textcolor}%
\pgftext[x=0.350000in, y=1.064548in, left, base]{\color{textcolor}{\rmfamily\fontsize{10.000000}{12.000000}\selectfont\catcode`\^=\active\def^{\ifmmode\sp\else\^{}\fi}\catcode`\%=\active\def%{\%}$\mathdefault{20}$}}%
\end{pgfscope}%
\begin{pgfscope}%
\pgfpathrectangle{\pgfqpoint{0.586112in}{0.502778in}}{\pgfqpoint{4.805000in}{2.695000in}}%
\pgfusepath{clip}%
\pgfsetrectcap%
\pgfsetroundjoin%
\pgfsetlinewidth{0.803000pt}%
\definecolor{currentstroke}{rgb}{0.690196,0.690196,0.690196}%
\pgfsetstrokecolor{currentstroke}%
\pgfsetdash{}{0pt}%
\pgfpathmoveto{\pgfqpoint{0.586112in}{1.606928in}}%
\pgfpathlineto{\pgfqpoint{5.391112in}{1.606928in}}%
\pgfusepath{stroke}%
\end{pgfscope}%
\begin{pgfscope}%
\pgfsetbuttcap%
\pgfsetroundjoin%
\definecolor{currentfill}{rgb}{0.000000,0.000000,0.000000}%
\pgfsetfillcolor{currentfill}%
\pgfsetlinewidth{0.803000pt}%
\definecolor{currentstroke}{rgb}{0.000000,0.000000,0.000000}%
\pgfsetstrokecolor{currentstroke}%
\pgfsetdash{}{0pt}%
\pgfsys@defobject{currentmarker}{\pgfqpoint{-0.048611in}{0.000000in}}{\pgfqpoint{-0.000000in}{0.000000in}}{%
\pgfpathmoveto{\pgfqpoint{-0.000000in}{0.000000in}}%
\pgfpathlineto{\pgfqpoint{-0.048611in}{0.000000in}}%
\pgfusepath{stroke,fill}%
}%
\begin{pgfscope}%
\pgfsys@transformshift{0.586112in}{1.606928in}%
\pgfsys@useobject{currentmarker}{}%
\end{pgfscope}%
\end{pgfscope}%
\begin{pgfscope}%
\definecolor{textcolor}{rgb}{0.000000,0.000000,0.000000}%
\pgfsetstrokecolor{textcolor}%
\pgfsetfillcolor{textcolor}%
\pgftext[x=0.350000in, y=1.559497in, left, base]{\color{textcolor}{\rmfamily\fontsize{10.000000}{12.000000}\selectfont\catcode`\^=\active\def^{\ifmmode\sp\else\^{}\fi}\catcode`\%=\active\def%{\%}$\mathdefault{40}$}}%
\end{pgfscope}%
\begin{pgfscope}%
\pgfpathrectangle{\pgfqpoint{0.586112in}{0.502778in}}{\pgfqpoint{4.805000in}{2.695000in}}%
\pgfusepath{clip}%
\pgfsetrectcap%
\pgfsetroundjoin%
\pgfsetlinewidth{0.803000pt}%
\definecolor{currentstroke}{rgb}{0.690196,0.690196,0.690196}%
\pgfsetstrokecolor{currentstroke}%
\pgfsetdash{}{0pt}%
\pgfpathmoveto{\pgfqpoint{0.586112in}{2.101877in}}%
\pgfpathlineto{\pgfqpoint{5.391112in}{2.101877in}}%
\pgfusepath{stroke}%
\end{pgfscope}%
\begin{pgfscope}%
\pgfsetbuttcap%
\pgfsetroundjoin%
\definecolor{currentfill}{rgb}{0.000000,0.000000,0.000000}%
\pgfsetfillcolor{currentfill}%
\pgfsetlinewidth{0.803000pt}%
\definecolor{currentstroke}{rgb}{0.000000,0.000000,0.000000}%
\pgfsetstrokecolor{currentstroke}%
\pgfsetdash{}{0pt}%
\pgfsys@defobject{currentmarker}{\pgfqpoint{-0.048611in}{0.000000in}}{\pgfqpoint{-0.000000in}{0.000000in}}{%
\pgfpathmoveto{\pgfqpoint{-0.000000in}{0.000000in}}%
\pgfpathlineto{\pgfqpoint{-0.048611in}{0.000000in}}%
\pgfusepath{stroke,fill}%
}%
\begin{pgfscope}%
\pgfsys@transformshift{0.586112in}{2.101877in}%
\pgfsys@useobject{currentmarker}{}%
\end{pgfscope}%
\end{pgfscope}%
\begin{pgfscope}%
\definecolor{textcolor}{rgb}{0.000000,0.000000,0.000000}%
\pgfsetstrokecolor{textcolor}%
\pgfsetfillcolor{textcolor}%
\pgftext[x=0.350000in, y=2.054447in, left, base]{\color{textcolor}{\rmfamily\fontsize{10.000000}{12.000000}\selectfont\catcode`\^=\active\def^{\ifmmode\sp\else\^{}\fi}\catcode`\%=\active\def%{\%}$\mathdefault{60}$}}%
\end{pgfscope}%
\begin{pgfscope}%
\pgfpathrectangle{\pgfqpoint{0.586112in}{0.502778in}}{\pgfqpoint{4.805000in}{2.695000in}}%
\pgfusepath{clip}%
\pgfsetrectcap%
\pgfsetroundjoin%
\pgfsetlinewidth{0.803000pt}%
\definecolor{currentstroke}{rgb}{0.690196,0.690196,0.690196}%
\pgfsetstrokecolor{currentstroke}%
\pgfsetdash{}{0pt}%
\pgfpathmoveto{\pgfqpoint{0.586112in}{2.596827in}}%
\pgfpathlineto{\pgfqpoint{5.391112in}{2.596827in}}%
\pgfusepath{stroke}%
\end{pgfscope}%
\begin{pgfscope}%
\pgfsetbuttcap%
\pgfsetroundjoin%
\definecolor{currentfill}{rgb}{0.000000,0.000000,0.000000}%
\pgfsetfillcolor{currentfill}%
\pgfsetlinewidth{0.803000pt}%
\definecolor{currentstroke}{rgb}{0.000000,0.000000,0.000000}%
\pgfsetstrokecolor{currentstroke}%
\pgfsetdash{}{0pt}%
\pgfsys@defobject{currentmarker}{\pgfqpoint{-0.048611in}{0.000000in}}{\pgfqpoint{-0.000000in}{0.000000in}}{%
\pgfpathmoveto{\pgfqpoint{-0.000000in}{0.000000in}}%
\pgfpathlineto{\pgfqpoint{-0.048611in}{0.000000in}}%
\pgfusepath{stroke,fill}%
}%
\begin{pgfscope}%
\pgfsys@transformshift{0.586112in}{2.596827in}%
\pgfsys@useobject{currentmarker}{}%
\end{pgfscope}%
\end{pgfscope}%
\begin{pgfscope}%
\definecolor{textcolor}{rgb}{0.000000,0.000000,0.000000}%
\pgfsetstrokecolor{textcolor}%
\pgfsetfillcolor{textcolor}%
\pgftext[x=0.350000in, y=2.549396in, left, base]{\color{textcolor}{\rmfamily\fontsize{10.000000}{12.000000}\selectfont\catcode`\^=\active\def^{\ifmmode\sp\else\^{}\fi}\catcode`\%=\active\def%{\%}$\mathdefault{80}$}}%
\end{pgfscope}%
\begin{pgfscope}%
\pgfpathrectangle{\pgfqpoint{0.586112in}{0.502778in}}{\pgfqpoint{4.805000in}{2.695000in}}%
\pgfusepath{clip}%
\pgfsetrectcap%
\pgfsetroundjoin%
\pgfsetlinewidth{0.803000pt}%
\definecolor{currentstroke}{rgb}{0.690196,0.690196,0.690196}%
\pgfsetstrokecolor{currentstroke}%
\pgfsetdash{}{0pt}%
\pgfpathmoveto{\pgfqpoint{0.586112in}{3.091776in}}%
\pgfpathlineto{\pgfqpoint{5.391112in}{3.091776in}}%
\pgfusepath{stroke}%
\end{pgfscope}%
\begin{pgfscope}%
\pgfsetbuttcap%
\pgfsetroundjoin%
\definecolor{currentfill}{rgb}{0.000000,0.000000,0.000000}%
\pgfsetfillcolor{currentfill}%
\pgfsetlinewidth{0.803000pt}%
\definecolor{currentstroke}{rgb}{0.000000,0.000000,0.000000}%
\pgfsetstrokecolor{currentstroke}%
\pgfsetdash{}{0pt}%
\pgfsys@defobject{currentmarker}{\pgfqpoint{-0.048611in}{0.000000in}}{\pgfqpoint{-0.000000in}{0.000000in}}{%
\pgfpathmoveto{\pgfqpoint{-0.000000in}{0.000000in}}%
\pgfpathlineto{\pgfqpoint{-0.048611in}{0.000000in}}%
\pgfusepath{stroke,fill}%
}%
\begin{pgfscope}%
\pgfsys@transformshift{0.586112in}{3.091776in}%
\pgfsys@useobject{currentmarker}{}%
\end{pgfscope}%
\end{pgfscope}%
\begin{pgfscope}%
\definecolor{textcolor}{rgb}{0.000000,0.000000,0.000000}%
\pgfsetstrokecolor{textcolor}%
\pgfsetfillcolor{textcolor}%
\pgftext[x=0.280556in, y=3.044346in, left, base]{\color{textcolor}{\rmfamily\fontsize{10.000000}{12.000000}\selectfont\catcode`\^=\active\def^{\ifmmode\sp\else\^{}\fi}\catcode`\%=\active\def%{\%}$\mathdefault{100}$}}%
\end{pgfscope}%
\begin{pgfscope}%
\definecolor{textcolor}{rgb}{0.000000,0.000000,0.000000}%
\pgfsetstrokecolor{textcolor}%
\pgfsetfillcolor{textcolor}%
\pgftext[x=0.225000in,y=1.850278in,,bottom,rotate=90.000000]{\color{textcolor}{\rmfamily\fontsize{10.000000}{12.000000}\selectfont\catcode`\^=\active\def^{\ifmmode\sp\else\^{}\fi}\catcode`\%=\active\def%{\%}Queue size (packets)}}%
\end{pgfscope}%
\begin{pgfscope}%
\pgfpathrectangle{\pgfqpoint{0.586112in}{0.502778in}}{\pgfqpoint{4.805000in}{2.695000in}}%
\pgfusepath{clip}%
\pgfsetrectcap%
\pgfsetroundjoin%
\pgfsetlinewidth{1.505625pt}%
\definecolor{currentstroke}{rgb}{0.000000,0.000000,0.000000}%
\pgfsetstrokecolor{currentstroke}%
\pgfsetdash{}{0pt}%
\pgfpathmoveto{\pgfqpoint{0.804521in}{0.633527in}}%
\pgfpathlineto{\pgfqpoint{0.804521in}{0.625278in}}%
\pgfpathlineto{\pgfqpoint{0.804521in}{0.625278in}}%
\pgfpathlineto{\pgfqpoint{0.806399in}{0.633527in}}%
\pgfpathlineto{\pgfqpoint{0.806416in}{0.625278in}}%
\pgfpathlineto{\pgfqpoint{0.807508in}{0.625278in}}%
\pgfpathlineto{\pgfqpoint{0.808565in}{0.633527in}}%
\pgfpathlineto{\pgfqpoint{0.808565in}{0.625278in}}%
\pgfpathlineto{\pgfqpoint{0.809751in}{0.633527in}}%
\pgfpathlineto{\pgfqpoint{0.809769in}{0.625278in}}%
\pgfpathlineto{\pgfqpoint{0.810738in}{0.625278in}}%
\pgfpathlineto{\pgfqpoint{0.810906in}{0.633527in}}%
\pgfpathlineto{\pgfqpoint{0.810906in}{0.625278in}}%
\pgfpathlineto{\pgfqpoint{0.812415in}{0.633527in}}%
\pgfpathlineto{\pgfqpoint{0.812583in}{0.625278in}}%
\pgfpathlineto{\pgfqpoint{0.813421in}{0.625278in}}%
\pgfpathlineto{\pgfqpoint{0.813923in}{0.691271in}}%
\pgfpathlineto{\pgfqpoint{0.814762in}{0.881002in}}%
\pgfpathlineto{\pgfqpoint{0.814948in}{0.872753in}}%
\pgfpathlineto{\pgfqpoint{0.815767in}{0.782012in}}%
\pgfpathlineto{\pgfqpoint{0.815954in}{0.773763in}}%
\pgfpathlineto{\pgfqpoint{0.816605in}{0.889251in}}%
\pgfpathlineto{\pgfqpoint{0.817630in}{1.318207in}}%
\pgfpathlineto{\pgfqpoint{0.818282in}{1.375951in}}%
\pgfpathlineto{\pgfqpoint{0.818636in}{1.334705in}}%
\pgfpathlineto{\pgfqpoint{0.819306in}{1.268712in}}%
\pgfpathlineto{\pgfqpoint{0.819791in}{1.326456in}}%
\pgfpathlineto{\pgfqpoint{0.820815in}{1.763662in}}%
\pgfpathlineto{\pgfqpoint{0.821299in}{1.846153in}}%
\pgfpathlineto{\pgfqpoint{0.821467in}{1.870901in}}%
\pgfpathlineto{\pgfqpoint{0.822324in}{1.780160in}}%
\pgfpathlineto{\pgfqpoint{0.822659in}{1.763662in}}%
\pgfpathlineto{\pgfqpoint{0.822975in}{1.804907in}}%
\pgfpathlineto{\pgfqpoint{0.824987in}{2.365850in}}%
\pgfpathlineto{\pgfqpoint{0.825173in}{2.349352in}}%
\pgfpathlineto{\pgfqpoint{0.826012in}{2.258611in}}%
\pgfpathlineto{\pgfqpoint{0.826496in}{2.316355in}}%
\pgfpathlineto{\pgfqpoint{0.826682in}{2.407096in}}%
\pgfpathlineto{\pgfqpoint{0.827521in}{2.299857in}}%
\pgfpathlineto{\pgfqpoint{0.839738in}{0.633527in}}%
\pgfpathlineto{\pgfqpoint{0.839757in}{0.650025in}}%
\pgfpathlineto{\pgfqpoint{0.841080in}{0.633527in}}%
\pgfpathlineto{\pgfqpoint{0.841247in}{0.658274in}}%
\pgfpathlineto{\pgfqpoint{0.842105in}{0.650025in}}%
\pgfpathlineto{\pgfqpoint{0.843259in}{0.658274in}}%
\pgfpathlineto{\pgfqpoint{0.844284in}{0.650025in}}%
\pgfpathlineto{\pgfqpoint{0.845103in}{0.658274in}}%
\pgfpathlineto{\pgfqpoint{0.845122in}{0.650025in}}%
\pgfpathlineto{\pgfqpoint{0.846443in}{0.633527in}}%
\pgfpathlineto{\pgfqpoint{0.846463in}{0.650025in}}%
\pgfpathlineto{\pgfqpoint{0.847468in}{0.650025in}}%
\pgfpathlineto{\pgfqpoint{0.847785in}{0.633527in}}%
\pgfpathlineto{\pgfqpoint{0.847952in}{0.658274in}}%
\pgfpathlineto{\pgfqpoint{0.848306in}{0.650025in}}%
\pgfpathlineto{\pgfqpoint{0.849461in}{0.658274in}}%
\pgfpathlineto{\pgfqpoint{0.850486in}{0.650025in}}%
\pgfpathlineto{\pgfqpoint{0.851640in}{0.658274in}}%
\pgfpathlineto{\pgfqpoint{0.852479in}{0.633527in}}%
\pgfpathlineto{\pgfqpoint{0.852665in}{0.641776in}}%
\pgfpathlineto{\pgfqpoint{0.852833in}{0.650025in}}%
\pgfpathlineto{\pgfqpoint{0.853820in}{0.633527in}}%
\pgfpathlineto{\pgfqpoint{0.854658in}{0.658274in}}%
\pgfpathlineto{\pgfqpoint{0.854844in}{0.650025in}}%
\pgfpathlineto{\pgfqpoint{0.855999in}{0.658274in}}%
\pgfpathlineto{\pgfqpoint{0.856837in}{0.633527in}}%
\pgfpathlineto{\pgfqpoint{0.857023in}{0.650025in}}%
\pgfpathlineto{\pgfqpoint{0.857843in}{0.658274in}}%
\pgfpathlineto{\pgfqpoint{0.857861in}{0.650025in}}%
\pgfpathlineto{\pgfqpoint{0.859023in}{0.633527in}}%
\pgfpathlineto{\pgfqpoint{0.858352in}{0.658274in}}%
\pgfpathlineto{\pgfqpoint{0.859023in}{0.641776in}}%
\pgfpathlineto{\pgfqpoint{0.859359in}{0.633527in}}%
\pgfpathlineto{\pgfqpoint{0.860041in}{0.650025in}}%
\pgfpathlineto{\pgfqpoint{0.860364in}{0.633527in}}%
\pgfpathlineto{\pgfqpoint{0.860532in}{0.658274in}}%
\pgfpathlineto{\pgfqpoint{0.860879in}{0.650025in}}%
\pgfpathlineto{\pgfqpoint{0.862048in}{0.658274in}}%
\pgfpathlineto{\pgfqpoint{0.863058in}{0.650025in}}%
\pgfpathlineto{\pgfqpoint{0.865561in}{0.839756in}}%
\pgfpathlineto{\pgfqpoint{0.866914in}{1.021237in}}%
\pgfpathlineto{\pgfqpoint{0.867919in}{1.095480in}}%
\pgfpathlineto{\pgfqpoint{0.868255in}{1.062483in}}%
\pgfpathlineto{\pgfqpoint{0.871440in}{0.625278in}}%
\pgfpathlineto{\pgfqpoint{0.871778in}{0.691271in}}%
\pgfpathlineto{\pgfqpoint{0.872112in}{0.683022in}}%
\pgfpathlineto{\pgfqpoint{0.872783in}{0.790261in}}%
\pgfpathlineto{\pgfqpoint{0.873789in}{0.889251in}}%
\pgfpathlineto{\pgfqpoint{0.874795in}{0.988241in}}%
\pgfpathlineto{\pgfqpoint{0.875800in}{1.087231in}}%
\pgfpathlineto{\pgfqpoint{0.875968in}{1.103729in}}%
\pgfpathlineto{\pgfqpoint{0.876301in}{1.087231in}}%
\pgfpathlineto{\pgfqpoint{0.879318in}{0.650025in}}%
\pgfpathlineto{\pgfqpoint{0.879656in}{0.740766in}}%
\pgfpathlineto{\pgfqpoint{0.881821in}{0.806759in}}%
\pgfpathlineto{\pgfqpoint{0.881833in}{0.798510in}}%
\pgfpathlineto{\pgfqpoint{0.882674in}{0.782012in}}%
\pgfpathlineto{\pgfqpoint{0.882674in}{0.790261in}}%
\pgfpathlineto{\pgfqpoint{0.884347in}{0.897500in}}%
\pgfpathlineto{\pgfqpoint{0.884517in}{0.881002in}}%
\pgfpathlineto{\pgfqpoint{0.884517in}{0.913998in}}%
\pgfpathlineto{\pgfqpoint{0.885523in}{1.004739in}}%
\pgfpathlineto{\pgfqpoint{0.885858in}{0.979992in}}%
\pgfpathlineto{\pgfqpoint{0.886861in}{0.946995in}}%
\pgfpathlineto{\pgfqpoint{0.886864in}{0.963493in}}%
\pgfpathlineto{\pgfqpoint{0.887367in}{1.004739in}}%
\pgfpathlineto{\pgfqpoint{0.887868in}{0.971742in}}%
\pgfpathlineto{\pgfqpoint{0.889376in}{0.946995in}}%
\pgfpathlineto{\pgfqpoint{0.889546in}{0.955244in}}%
\pgfpathlineto{\pgfqpoint{0.890382in}{0.922247in}}%
\pgfpathlineto{\pgfqpoint{0.892058in}{0.790261in}}%
\pgfpathlineto{\pgfqpoint{0.894072in}{0.625278in}}%
\pgfpathlineto{\pgfqpoint{0.894911in}{0.633527in}}%
\pgfpathlineto{\pgfqpoint{0.894911in}{0.625278in}}%
\pgfpathlineto{\pgfqpoint{0.896244in}{0.633527in}}%
\pgfpathlineto{\pgfqpoint{0.896579in}{0.625278in}}%
\pgfpathlineto{\pgfqpoint{0.897250in}{0.625278in}}%
\pgfpathlineto{\pgfqpoint{0.898431in}{0.633527in}}%
\pgfpathlineto{\pgfqpoint{0.898599in}{0.625278in}}%
\pgfpathlineto{\pgfqpoint{0.899437in}{0.625278in}}%
\pgfpathlineto{\pgfqpoint{0.900610in}{0.633527in}}%
\pgfpathlineto{\pgfqpoint{0.900945in}{0.625278in}}%
\pgfpathlineto{\pgfqpoint{0.901616in}{0.625278in}}%
\pgfpathlineto{\pgfqpoint{0.902789in}{0.633527in}}%
\pgfpathlineto{\pgfqpoint{0.903124in}{0.625278in}}%
\pgfpathlineto{\pgfqpoint{0.903795in}{0.625278in}}%
\pgfpathlineto{\pgfqpoint{0.905136in}{0.633527in}}%
\pgfpathlineto{\pgfqpoint{0.905471in}{0.625278in}}%
\pgfpathlineto{\pgfqpoint{0.906142in}{0.625278in}}%
\pgfpathlineto{\pgfqpoint{0.907650in}{0.633527in}}%
\pgfpathlineto{\pgfqpoint{0.907818in}{0.625278in}}%
\pgfpathlineto{\pgfqpoint{0.908656in}{0.625278in}}%
\pgfpathlineto{\pgfqpoint{0.909829in}{0.633527in}}%
\pgfpathlineto{\pgfqpoint{0.909997in}{0.625278in}}%
\pgfpathlineto{\pgfqpoint{0.910836in}{0.625278in}}%
\pgfpathlineto{\pgfqpoint{0.911338in}{0.633527in}}%
\pgfpathlineto{\pgfqpoint{0.911338in}{0.625278in}}%
\pgfpathlineto{\pgfqpoint{0.912512in}{0.633527in}}%
\pgfpathlineto{\pgfqpoint{0.912847in}{0.625278in}}%
\pgfpathlineto{\pgfqpoint{0.913517in}{0.625278in}}%
\pgfpathlineto{\pgfqpoint{0.914852in}{0.633527in}}%
\pgfpathlineto{\pgfqpoint{0.915194in}{0.625278in}}%
\pgfpathlineto{\pgfqpoint{0.915864in}{0.625278in}}%
\pgfpathlineto{\pgfqpoint{0.917038in}{0.633527in}}%
\pgfpathlineto{\pgfqpoint{0.917198in}{0.625278in}}%
\pgfpathlineto{\pgfqpoint{0.918036in}{0.625278in}}%
\pgfpathlineto{\pgfqpoint{0.919203in}{0.633527in}}%
\pgfpathlineto{\pgfqpoint{0.919538in}{0.625278in}}%
\pgfpathlineto{\pgfqpoint{0.920208in}{0.625278in}}%
\pgfpathlineto{\pgfqpoint{0.921229in}{0.633527in}}%
\pgfpathlineto{\pgfqpoint{0.921229in}{0.625278in}}%
\pgfpathlineto{\pgfqpoint{0.922570in}{0.633527in}}%
\pgfpathlineto{\pgfqpoint{0.922905in}{0.625278in}}%
\pgfpathlineto{\pgfqpoint{0.923575in}{0.625278in}}%
\pgfpathlineto{\pgfqpoint{0.924735in}{0.633527in}}%
\pgfpathlineto{\pgfqpoint{0.924916in}{0.625278in}}%
\pgfpathlineto{\pgfqpoint{0.925754in}{0.625278in}}%
\pgfpathlineto{\pgfqpoint{0.926914in}{0.633527in}}%
\pgfpathlineto{\pgfqpoint{0.927249in}{0.625278in}}%
\pgfpathlineto{\pgfqpoint{0.927919in}{0.625278in}}%
\pgfpathlineto{\pgfqpoint{0.929202in}{0.633527in}}%
\pgfpathlineto{\pgfqpoint{0.929261in}{0.625278in}}%
\pgfpathlineto{\pgfqpoint{0.930266in}{0.625278in}}%
\pgfpathlineto{\pgfqpoint{0.931440in}{0.633527in}}%
\pgfpathlineto{\pgfqpoint{0.931775in}{0.625278in}}%
\pgfpathlineto{\pgfqpoint{0.932445in}{0.625278in}}%
\pgfpathlineto{\pgfqpoint{0.933452in}{0.633527in}}%
\pgfpathlineto{\pgfqpoint{0.933452in}{0.625278in}}%
\pgfpathlineto{\pgfqpoint{0.934793in}{0.633527in}}%
\pgfpathlineto{\pgfqpoint{0.934793in}{0.625278in}}%
\pgfpathlineto{\pgfqpoint{0.935798in}{0.625278in}}%
\pgfpathlineto{\pgfqpoint{0.936972in}{0.633527in}}%
\pgfpathlineto{\pgfqpoint{0.937139in}{0.625278in}}%
\pgfpathlineto{\pgfqpoint{0.937810in}{0.625278in}}%
\pgfpathlineto{\pgfqpoint{0.938997in}{0.633527in}}%
\pgfpathlineto{\pgfqpoint{0.939151in}{0.625278in}}%
\pgfpathlineto{\pgfqpoint{0.939989in}{0.625278in}}%
\pgfpathlineto{\pgfqpoint{0.941176in}{0.633527in}}%
\pgfpathlineto{\pgfqpoint{0.941330in}{0.625278in}}%
\pgfpathlineto{\pgfqpoint{0.942183in}{0.625278in}}%
\pgfpathlineto{\pgfqpoint{0.943355in}{0.633527in}}%
\pgfpathlineto{\pgfqpoint{0.943355in}{0.625278in}}%
\pgfpathlineto{\pgfqpoint{0.944362in}{0.625278in}}%
\pgfpathlineto{\pgfqpoint{0.945520in}{0.633527in}}%
\pgfpathlineto{\pgfqpoint{0.945703in}{0.625278in}}%
\pgfpathlineto{\pgfqpoint{0.946527in}{0.625278in}}%
\pgfpathlineto{\pgfqpoint{0.947700in}{0.633527in}}%
\pgfpathlineto{\pgfqpoint{0.947882in}{0.625278in}}%
\pgfpathlineto{\pgfqpoint{0.948706in}{0.625278in}}%
\pgfpathlineto{\pgfqpoint{0.949893in}{0.633527in}}%
\pgfpathlineto{\pgfqpoint{0.949893in}{0.625278in}}%
\pgfpathlineto{\pgfqpoint{0.950899in}{0.625278in}}%
\pgfpathlineto{\pgfqpoint{0.952058in}{0.633527in}}%
\pgfpathlineto{\pgfqpoint{0.952226in}{0.625278in}}%
\pgfpathlineto{\pgfqpoint{0.953079in}{0.625278in}}%
\pgfpathlineto{\pgfqpoint{0.954084in}{0.633527in}}%
\pgfpathlineto{\pgfqpoint{0.954084in}{0.625278in}}%
\pgfpathlineto{\pgfqpoint{0.955244in}{0.633527in}}%
\pgfpathlineto{\pgfqpoint{0.955425in}{0.625278in}}%
\pgfpathlineto{\pgfqpoint{0.956074in}{0.625278in}}%
\pgfpathlineto{\pgfqpoint{0.957269in}{0.633527in}}%
\pgfpathlineto{\pgfqpoint{0.957415in}{0.625278in}}%
\pgfpathlineto{\pgfqpoint{0.958086in}{0.625278in}}%
\pgfpathlineto{\pgfqpoint{0.959427in}{0.633527in}}%
\pgfpathlineto{\pgfqpoint{0.959448in}{0.625278in}}%
\pgfpathlineto{\pgfqpoint{0.960454in}{0.625278in}}%
\pgfpathlineto{\pgfqpoint{0.961606in}{0.633527in}}%
\pgfpathlineto{\pgfqpoint{0.961795in}{0.625278in}}%
\pgfpathlineto{\pgfqpoint{0.962612in}{0.625278in}}%
\pgfpathlineto{\pgfqpoint{0.963618in}{0.633527in}}%
\pgfpathlineto{\pgfqpoint{0.963618in}{0.625278in}}%
\pgfpathlineto{\pgfqpoint{0.964791in}{0.633527in}}%
\pgfpathlineto{\pgfqpoint{0.964959in}{0.625278in}}%
\pgfpathlineto{\pgfqpoint{0.965797in}{0.625278in}}%
\pgfpathlineto{\pgfqpoint{0.966970in}{0.633527in}}%
\pgfpathlineto{\pgfqpoint{0.967138in}{0.625278in}}%
\pgfpathlineto{\pgfqpoint{0.967976in}{0.625278in}}%
\pgfpathlineto{\pgfqpoint{0.968982in}{0.633527in}}%
\pgfpathlineto{\pgfqpoint{0.968982in}{0.625278in}}%
\pgfpathlineto{\pgfqpoint{0.970156in}{0.633527in}}%
\pgfpathlineto{\pgfqpoint{0.970323in}{0.625278in}}%
\pgfpathlineto{\pgfqpoint{0.971183in}{0.625278in}}%
\pgfpathlineto{\pgfqpoint{0.972335in}{0.633527in}}%
\pgfpathlineto{\pgfqpoint{0.972335in}{0.625278in}}%
\pgfpathlineto{\pgfqpoint{0.973340in}{0.625278in}}%
\pgfpathlineto{\pgfqpoint{0.974514in}{0.633527in}}%
\pgfpathlineto{\pgfqpoint{0.974681in}{0.625278in}}%
\pgfpathlineto{\pgfqpoint{0.975519in}{0.625278in}}%
\pgfpathlineto{\pgfqpoint{0.976693in}{0.633527in}}%
\pgfpathlineto{\pgfqpoint{0.976861in}{0.625278in}}%
\pgfpathlineto{\pgfqpoint{0.977699in}{0.625278in}}%
\pgfpathlineto{\pgfqpoint{0.978872in}{0.633527in}}%
\pgfpathlineto{\pgfqpoint{0.978893in}{0.625278in}}%
\pgfpathlineto{\pgfqpoint{0.979878in}{0.625278in}}%
\pgfpathlineto{\pgfqpoint{0.981052in}{0.633527in}}%
\pgfpathlineto{\pgfqpoint{0.981052in}{0.625278in}}%
\pgfpathlineto{\pgfqpoint{0.982057in}{0.625278in}}%
\pgfpathlineto{\pgfqpoint{0.983252in}{0.633527in}}%
\pgfpathlineto{\pgfqpoint{0.983398in}{0.625278in}}%
\pgfpathlineto{\pgfqpoint{0.984069in}{0.625278in}}%
\pgfpathlineto{\pgfqpoint{0.985242in}{0.633527in}}%
\pgfpathlineto{\pgfqpoint{0.985410in}{0.625278in}}%
\pgfpathlineto{\pgfqpoint{0.986248in}{0.625278in}}%
\pgfpathlineto{\pgfqpoint{0.987422in}{0.633527in}}%
\pgfpathlineto{\pgfqpoint{0.987589in}{0.625278in}}%
\pgfpathlineto{\pgfqpoint{0.988427in}{0.625278in}}%
\pgfpathlineto{\pgfqpoint{0.989601in}{0.633527in}}%
\pgfpathlineto{\pgfqpoint{0.989775in}{0.625278in}}%
\pgfpathlineto{\pgfqpoint{0.990606in}{0.625278in}}%
\pgfpathlineto{\pgfqpoint{0.991634in}{0.633527in}}%
\pgfpathlineto{\pgfqpoint{0.991634in}{0.625278in}}%
\pgfpathlineto{\pgfqpoint{0.992807in}{0.633527in}}%
\pgfpathlineto{\pgfqpoint{0.992975in}{0.625278in}}%
\pgfpathlineto{\pgfqpoint{0.993813in}{0.625278in}}%
\pgfpathlineto{\pgfqpoint{0.994965in}{0.633527in}}%
\pgfpathlineto{\pgfqpoint{0.995970in}{0.625278in}}%
\pgfpathlineto{\pgfqpoint{0.997144in}{0.633527in}}%
\pgfpathlineto{\pgfqpoint{0.997311in}{0.625278in}}%
\pgfpathlineto{\pgfqpoint{0.998149in}{0.625278in}}%
\pgfpathlineto{\pgfqpoint{0.999323in}{0.633527in}}%
\pgfpathlineto{\pgfqpoint{0.999491in}{0.625278in}}%
\pgfpathlineto{\pgfqpoint{1.000329in}{0.625278in}}%
\pgfpathlineto{\pgfqpoint{1.001502in}{0.633527in}}%
\pgfpathlineto{\pgfqpoint{1.001838in}{0.658274in}}%
\pgfpathlineto{\pgfqpoint{1.002527in}{0.625278in}}%
\pgfpathlineto{\pgfqpoint{1.003011in}{0.658274in}}%
\pgfpathlineto{\pgfqpoint{1.003682in}{0.633527in}}%
\pgfpathlineto{\pgfqpoint{1.003700in}{0.625278in}}%
\pgfpathlineto{\pgfqpoint{1.004352in}{0.683022in}}%
\pgfpathlineto{\pgfqpoint{1.004539in}{0.674773in}}%
\pgfpathlineto{\pgfqpoint{1.007705in}{0.905749in}}%
\pgfpathlineto{\pgfqpoint{1.008711in}{0.955244in}}%
\pgfpathlineto{\pgfqpoint{1.008732in}{0.955244in}}%
\pgfpathlineto{\pgfqpoint{1.013404in}{1.276961in}}%
\pgfpathlineto{\pgfqpoint{1.014467in}{1.425446in}}%
\pgfpathlineto{\pgfqpoint{1.016608in}{1.565682in}}%
\pgfpathlineto{\pgfqpoint{1.017111in}{1.540934in}}%
\pgfpathlineto{\pgfqpoint{1.017427in}{1.524436in}}%
\pgfpathlineto{\pgfqpoint{1.017595in}{1.549184in}}%
\pgfpathlineto{\pgfqpoint{1.017949in}{1.540934in}}%
\pgfpathlineto{\pgfqpoint{1.019961in}{1.615177in}}%
\pgfpathlineto{\pgfqpoint{1.020128in}{1.606928in}}%
\pgfpathlineto{\pgfqpoint{1.020445in}{1.623426in}}%
\pgfpathlineto{\pgfqpoint{1.021302in}{1.590429in}}%
\pgfpathlineto{\pgfqpoint{1.021786in}{1.598678in}}%
\pgfpathlineto{\pgfqpoint{1.021972in}{1.582180in}}%
\pgfpathlineto{\pgfqpoint{1.022121in}{1.573931in}}%
\pgfpathlineto{\pgfqpoint{1.022289in}{1.598678in}}%
\pgfpathlineto{\pgfqpoint{1.022643in}{1.590429in}}%
\pgfpathlineto{\pgfqpoint{1.025325in}{1.714167in}}%
\pgfpathlineto{\pgfqpoint{1.025641in}{1.697668in}}%
\pgfpathlineto{\pgfqpoint{1.027001in}{1.582180in}}%
\pgfpathlineto{\pgfqpoint{1.034377in}{0.650025in}}%
\pgfpathlineto{\pgfqpoint{1.034526in}{0.658274in}}%
\pgfpathlineto{\pgfqpoint{1.034880in}{0.625278in}}%
\pgfpathlineto{\pgfqpoint{1.035028in}{0.633527in}}%
\pgfpathlineto{\pgfqpoint{1.035215in}{0.625278in}}%
\pgfpathlineto{\pgfqpoint{1.035532in}{0.658274in}}%
\pgfpathlineto{\pgfqpoint{1.035886in}{0.641776in}}%
\pgfpathlineto{\pgfqpoint{1.036891in}{0.625278in}}%
\pgfpathlineto{\pgfqpoint{1.037040in}{0.633527in}}%
\pgfpathlineto{\pgfqpoint{1.037059in}{0.650025in}}%
\pgfpathlineto{\pgfqpoint{1.038065in}{0.625278in}}%
\pgfpathlineto{\pgfqpoint{1.039052in}{0.658274in}}%
\pgfpathlineto{\pgfqpoint{1.039238in}{0.641776in}}%
\pgfpathlineto{\pgfqpoint{1.039387in}{0.633527in}}%
\pgfpathlineto{\pgfqpoint{1.039554in}{0.658274in}}%
\pgfpathlineto{\pgfqpoint{1.039909in}{0.650025in}}%
\pgfpathlineto{\pgfqpoint{1.040057in}{0.658274in}}%
\pgfpathlineto{\pgfqpoint{1.040412in}{0.625278in}}%
\pgfpathlineto{\pgfqpoint{1.040561in}{0.633527in}}%
\pgfpathlineto{\pgfqpoint{1.041585in}{0.625278in}}%
\pgfpathlineto{\pgfqpoint{1.042424in}{0.650025in}}%
\pgfpathlineto{\pgfqpoint{1.042740in}{0.633527in}}%
\pgfpathlineto{\pgfqpoint{1.043765in}{0.625278in}}%
\pgfpathlineto{\pgfqpoint{1.044919in}{0.633527in}}%
\pgfpathlineto{\pgfqpoint{1.044938in}{0.625278in}}%
\pgfpathlineto{\pgfqpoint{1.045422in}{0.658274in}}%
\pgfpathlineto{\pgfqpoint{1.045944in}{0.650025in}}%
\pgfpathlineto{\pgfqpoint{1.047787in}{0.749015in}}%
\pgfpathlineto{\pgfqpoint{1.048439in}{0.732517in}}%
\pgfpathlineto{\pgfqpoint{1.049464in}{0.699520in}}%
\pgfpathlineto{\pgfqpoint{1.051959in}{0.782012in}}%
\pgfpathlineto{\pgfqpoint{1.051978in}{0.773763in}}%
\pgfpathlineto{\pgfqpoint{1.052295in}{0.806759in}}%
\pgfpathlineto{\pgfqpoint{1.052984in}{0.798510in}}%
\pgfpathlineto{\pgfqpoint{1.057666in}{1.004739in}}%
\pgfpathlineto{\pgfqpoint{1.057848in}{0.979992in}}%
\pgfpathlineto{\pgfqpoint{1.058683in}{1.045985in}}%
\pgfpathlineto{\pgfqpoint{1.059670in}{1.054234in}}%
\pgfpathlineto{\pgfqpoint{1.059689in}{1.045985in}}%
\pgfpathlineto{\pgfqpoint{1.060528in}{1.021237in}}%
\pgfpathlineto{\pgfqpoint{1.060695in}{1.045985in}}%
\pgfpathlineto{\pgfqpoint{1.063545in}{1.169722in}}%
\pgfpathlineto{\pgfqpoint{1.063712in}{1.161473in}}%
\pgfpathlineto{\pgfqpoint{1.063868in}{1.153224in}}%
\pgfpathlineto{\pgfqpoint{1.064204in}{1.186221in}}%
\pgfpathlineto{\pgfqpoint{1.064364in}{1.227466in}}%
\pgfpathlineto{\pgfqpoint{1.065221in}{1.111978in}}%
\pgfpathlineto{\pgfqpoint{1.068729in}{0.633527in}}%
\pgfpathlineto{\pgfqpoint{1.069412in}{0.650025in}}%
\pgfpathlineto{\pgfqpoint{1.071587in}{0.732517in}}%
\pgfpathlineto{\pgfqpoint{1.071759in}{0.716019in}}%
\pgfpathlineto{\pgfqpoint{1.072257in}{0.707769in}}%
\pgfpathlineto{\pgfqpoint{1.072417in}{0.732517in}}%
\pgfpathlineto{\pgfqpoint{1.072429in}{0.724268in}}%
\pgfpathlineto{\pgfqpoint{1.072592in}{0.732517in}}%
\pgfpathlineto{\pgfqpoint{1.072932in}{0.699520in}}%
\pgfpathlineto{\pgfqpoint{1.072935in}{0.707769in}}%
\pgfpathlineto{\pgfqpoint{1.073267in}{0.699520in}}%
\pgfpathlineto{\pgfqpoint{1.073435in}{0.724268in}}%
\pgfpathlineto{\pgfqpoint{1.074604in}{0.732517in}}%
\pgfpathlineto{\pgfqpoint{1.074946in}{0.707769in}}%
\pgfpathlineto{\pgfqpoint{1.075112in}{0.749015in}}%
\pgfpathlineto{\pgfqpoint{1.075617in}{0.732517in}}%
\pgfpathlineto{\pgfqpoint{1.078452in}{0.831507in}}%
\pgfpathlineto{\pgfqpoint{1.079470in}{0.798510in}}%
\pgfpathlineto{\pgfqpoint{1.080631in}{0.806759in}}%
\pgfpathlineto{\pgfqpoint{1.081817in}{0.740766in}}%
\pgfpathlineto{\pgfqpoint{1.084499in}{0.650025in}}%
\pgfpathlineto{\pgfqpoint{1.084654in}{0.658274in}}%
\pgfpathlineto{\pgfqpoint{1.085001in}{0.625278in}}%
\pgfpathlineto{\pgfqpoint{1.085996in}{0.633527in}}%
\pgfpathlineto{\pgfqpoint{1.085996in}{0.625278in}}%
\pgfpathlineto{\pgfqpoint{1.087162in}{0.633527in}}%
\pgfpathlineto{\pgfqpoint{1.087337in}{0.625278in}}%
\pgfpathlineto{\pgfqpoint{1.088175in}{0.625278in}}%
\pgfpathlineto{\pgfqpoint{1.089355in}{0.633527in}}%
\pgfpathlineto{\pgfqpoint{1.089516in}{0.625278in}}%
\pgfpathlineto{\pgfqpoint{1.090354in}{0.625278in}}%
\pgfpathlineto{\pgfqpoint{1.091527in}{0.633527in}}%
\pgfpathlineto{\pgfqpoint{1.091695in}{0.625278in}}%
\pgfpathlineto{\pgfqpoint{1.092533in}{0.625278in}}%
\pgfpathlineto{\pgfqpoint{1.095208in}{0.707769in}}%
\pgfpathlineto{\pgfqpoint{1.096066in}{0.699520in}}%
\pgfpathlineto{\pgfqpoint{1.096233in}{0.724268in}}%
\pgfpathlineto{\pgfqpoint{1.096382in}{0.732517in}}%
\pgfpathlineto{\pgfqpoint{1.096736in}{0.699520in}}%
\pgfpathlineto{\pgfqpoint{1.096884in}{0.707769in}}%
\pgfpathlineto{\pgfqpoint{1.097071in}{0.699520in}}%
\pgfpathlineto{\pgfqpoint{1.097742in}{0.773763in}}%
\pgfpathlineto{\pgfqpoint{1.098393in}{0.806759in}}%
\pgfpathlineto{\pgfqpoint{1.098903in}{0.782012in}}%
\pgfpathlineto{\pgfqpoint{1.099734in}{0.732517in}}%
\pgfpathlineto{\pgfqpoint{1.099921in}{0.749015in}}%
\pgfpathlineto{\pgfqpoint{1.100926in}{0.773763in}}%
\pgfpathlineto{\pgfqpoint{1.101094in}{0.765513in}}%
\pgfpathlineto{\pgfqpoint{1.101250in}{0.757264in}}%
\pgfpathlineto{\pgfqpoint{1.101411in}{0.782012in}}%
\pgfpathlineto{\pgfqpoint{1.101746in}{0.782012in}}%
\pgfpathlineto{\pgfqpoint{1.102759in}{0.806759in}}%
\pgfpathlineto{\pgfqpoint{1.102771in}{0.798510in}}%
\pgfpathlineto{\pgfqpoint{1.104435in}{0.831507in}}%
\pgfpathlineto{\pgfqpoint{1.104447in}{0.823258in}}%
\pgfpathlineto{\pgfqpoint{1.104614in}{0.848005in}}%
\pgfpathlineto{\pgfqpoint{1.105453in}{0.839756in}}%
\pgfpathlineto{\pgfqpoint{1.106271in}{0.806759in}}%
\pgfpathlineto{\pgfqpoint{1.105620in}{0.848005in}}%
\pgfpathlineto{\pgfqpoint{1.106626in}{0.823258in}}%
\pgfpathlineto{\pgfqpoint{1.107459in}{0.856254in}}%
\pgfpathlineto{\pgfqpoint{1.107795in}{0.831507in}}%
\pgfpathlineto{\pgfqpoint{1.108451in}{0.856254in}}%
\pgfpathlineto{\pgfqpoint{1.108805in}{0.823258in}}%
\pgfpathlineto{\pgfqpoint{1.109959in}{0.831507in}}%
\pgfpathlineto{\pgfqpoint{1.110314in}{0.848005in}}%
\pgfpathlineto{\pgfqpoint{1.110984in}{0.823258in}}%
\pgfpathlineto{\pgfqpoint{1.111636in}{0.856254in}}%
\pgfpathlineto{\pgfqpoint{1.112158in}{0.839756in}}%
\pgfpathlineto{\pgfqpoint{1.112493in}{0.823258in}}%
\pgfpathlineto{\pgfqpoint{1.112817in}{0.856254in}}%
\pgfpathlineto{\pgfqpoint{1.112829in}{0.848005in}}%
\pgfpathlineto{\pgfqpoint{1.113151in}{0.856254in}}%
\pgfpathlineto{\pgfqpoint{1.113487in}{0.831507in}}%
\pgfpathlineto{\pgfqpoint{1.114505in}{0.798510in}}%
\pgfpathlineto{\pgfqpoint{1.115659in}{0.806759in}}%
\pgfpathlineto{\pgfqpoint{1.116013in}{0.823258in}}%
\pgfpathlineto{\pgfqpoint{1.116851in}{0.765513in}}%
\pgfpathlineto{\pgfqpoint{1.117671in}{0.757264in}}%
\pgfpathlineto{\pgfqpoint{1.117019in}{0.773763in}}%
\pgfpathlineto{\pgfqpoint{1.117689in}{0.773763in}}%
\pgfpathlineto{\pgfqpoint{1.118025in}{0.798510in}}%
\pgfpathlineto{\pgfqpoint{1.118844in}{0.782012in}}%
\pgfpathlineto{\pgfqpoint{1.118863in}{0.773763in}}%
\pgfpathlineto{\pgfqpoint{1.119179in}{0.806759in}}%
\pgfpathlineto{\pgfqpoint{1.119869in}{0.798510in}}%
\pgfpathlineto{\pgfqpoint{1.120185in}{0.782012in}}%
\pgfpathlineto{\pgfqpoint{1.120204in}{0.823258in}}%
\pgfpathlineto{\pgfqpoint{1.120520in}{0.815008in}}%
\pgfpathlineto{\pgfqpoint{1.121701in}{0.881002in}}%
\pgfpathlineto{\pgfqpoint{1.120856in}{0.806759in}}%
\pgfpathlineto{\pgfqpoint{1.121713in}{0.872753in}}%
\pgfpathlineto{\pgfqpoint{1.122383in}{0.922247in}}%
\pgfpathlineto{\pgfqpoint{1.123377in}{0.905749in}}%
\pgfpathlineto{\pgfqpoint{1.124563in}{0.823258in}}%
\pgfpathlineto{\pgfqpoint{1.124879in}{0.831507in}}%
\pgfpathlineto{\pgfqpoint{1.125221in}{0.806759in}}%
\pgfpathlineto{\pgfqpoint{1.126220in}{0.757264in}}%
\pgfpathlineto{\pgfqpoint{1.126239in}{0.773763in}}%
\pgfpathlineto{\pgfqpoint{1.126891in}{0.757264in}}%
\pgfpathlineto{\pgfqpoint{1.126905in}{0.790261in}}%
\pgfpathlineto{\pgfqpoint{1.127225in}{0.782012in}}%
\pgfpathlineto{\pgfqpoint{1.127915in}{0.823258in}}%
\pgfpathlineto{\pgfqpoint{1.128754in}{0.773763in}}%
\pgfpathlineto{\pgfqpoint{1.129237in}{0.806759in}}%
\pgfpathlineto{\pgfqpoint{1.129405in}{0.831507in}}%
\pgfpathlineto{\pgfqpoint{1.129759in}{0.798510in}}%
\pgfpathlineto{\pgfqpoint{1.130262in}{0.798510in}}%
\pgfpathlineto{\pgfqpoint{1.131268in}{0.823258in}}%
\pgfpathlineto{\pgfqpoint{1.131435in}{0.815008in}}%
\pgfpathlineto{\pgfqpoint{1.131584in}{0.806759in}}%
\pgfpathlineto{\pgfqpoint{1.131752in}{0.831507in}}%
\pgfpathlineto{\pgfqpoint{1.132106in}{0.823258in}}%
\pgfpathlineto{\pgfqpoint{1.133763in}{0.856254in}}%
\pgfpathlineto{\pgfqpoint{1.134621in}{0.848005in}}%
\pgfpathlineto{\pgfqpoint{1.134788in}{0.872753in}}%
\pgfpathlineto{\pgfqpoint{1.134937in}{0.881002in}}%
\pgfpathlineto{\pgfqpoint{1.135291in}{0.848005in}}%
\pgfpathlineto{\pgfqpoint{1.135440in}{0.856254in}}%
\pgfpathlineto{\pgfqpoint{1.136297in}{0.848005in}}%
\pgfpathlineto{\pgfqpoint{1.135775in}{0.881002in}}%
\pgfpathlineto{\pgfqpoint{1.136464in}{0.872753in}}%
\pgfpathlineto{\pgfqpoint{1.136613in}{0.881002in}}%
\pgfpathlineto{\pgfqpoint{1.136967in}{0.848005in}}%
\pgfpathlineto{\pgfqpoint{1.137283in}{0.831507in}}%
\pgfpathlineto{\pgfqpoint{1.137458in}{0.856254in}}%
\pgfpathlineto{\pgfqpoint{1.137805in}{0.848005in}}%
\pgfpathlineto{\pgfqpoint{1.138792in}{0.881002in}}%
\pgfpathlineto{\pgfqpoint{1.138979in}{0.864503in}}%
\pgfpathlineto{\pgfqpoint{1.140468in}{0.806759in}}%
\pgfpathlineto{\pgfqpoint{1.140643in}{0.831507in}}%
\pgfpathlineto{\pgfqpoint{1.141493in}{0.815008in}}%
\pgfpathlineto{\pgfqpoint{1.141977in}{0.806759in}}%
\pgfpathlineto{\pgfqpoint{1.141996in}{0.848005in}}%
\pgfpathlineto{\pgfqpoint{1.142648in}{0.831507in}}%
\pgfpathlineto{\pgfqpoint{1.142655in}{0.864503in}}%
\pgfpathlineto{\pgfqpoint{1.143505in}{0.897500in}}%
\pgfpathlineto{\pgfqpoint{1.143672in}{0.889251in}}%
\pgfpathlineto{\pgfqpoint{1.143840in}{0.897500in}}%
\pgfpathlineto{\pgfqpoint{1.144827in}{0.881002in}}%
\pgfpathlineto{\pgfqpoint{1.144846in}{0.897500in}}%
\pgfpathlineto{\pgfqpoint{1.145517in}{0.872753in}}%
\pgfpathlineto{\pgfqpoint{1.145851in}{0.897500in}}%
\pgfpathlineto{\pgfqpoint{1.147006in}{0.955244in}}%
\pgfpathlineto{\pgfqpoint{1.147509in}{0.930497in}}%
\pgfpathlineto{\pgfqpoint{1.147528in}{0.922247in}}%
\pgfpathlineto{\pgfqpoint{1.147845in}{0.955244in}}%
\pgfpathlineto{\pgfqpoint{1.148534in}{0.946995in}}%
\pgfpathlineto{\pgfqpoint{1.148689in}{0.955244in}}%
\pgfpathlineto{\pgfqpoint{1.149017in}{0.930497in}}%
\pgfpathlineto{\pgfqpoint{1.149204in}{0.938746in}}%
\pgfpathlineto{\pgfqpoint{1.149353in}{0.930497in}}%
\pgfpathlineto{\pgfqpoint{1.149521in}{0.955244in}}%
\pgfpathlineto{\pgfqpoint{1.149875in}{0.946995in}}%
\pgfpathlineto{\pgfqpoint{1.150042in}{0.971742in}}%
\pgfpathlineto{\pgfqpoint{1.150713in}{0.938746in}}%
\pgfpathlineto{\pgfqpoint{1.151551in}{0.946995in}}%
\pgfpathlineto{\pgfqpoint{1.151886in}{0.922247in}}%
\pgfpathlineto{\pgfqpoint{1.152370in}{0.905749in}}%
\pgfpathlineto{\pgfqpoint{1.153041in}{0.930497in}}%
\pgfpathlineto{\pgfqpoint{1.154066in}{0.897500in}}%
\pgfpathlineto{\pgfqpoint{1.155388in}{0.930497in}}%
\pgfpathlineto{\pgfqpoint{1.155406in}{0.922247in}}%
\pgfpathlineto{\pgfqpoint{1.156226in}{0.930497in}}%
\pgfpathlineto{\pgfqpoint{1.156747in}{0.897500in}}%
\pgfpathlineto{\pgfqpoint{1.157064in}{0.905749in}}%
\pgfpathlineto{\pgfqpoint{1.157251in}{0.889251in}}%
\pgfpathlineto{\pgfqpoint{1.157791in}{0.831507in}}%
\pgfpathlineto{\pgfqpoint{1.158592in}{0.848005in}}%
\pgfpathlineto{\pgfqpoint{1.159075in}{0.856254in}}%
\pgfpathlineto{\pgfqpoint{1.159262in}{0.839756in}}%
\pgfpathlineto{\pgfqpoint{1.160920in}{0.757264in}}%
\pgfpathlineto{\pgfqpoint{1.160938in}{0.773763in}}%
\pgfpathlineto{\pgfqpoint{1.161944in}{0.724268in}}%
\pgfpathlineto{\pgfqpoint{1.163769in}{0.658274in}}%
\pgfpathlineto{\pgfqpoint{1.163788in}{0.674773in}}%
\pgfpathlineto{\pgfqpoint{1.164440in}{0.707769in}}%
\pgfpathlineto{\pgfqpoint{1.164962in}{0.691271in}}%
\pgfpathlineto{\pgfqpoint{1.165129in}{0.699520in}}%
\pgfpathlineto{\pgfqpoint{1.166116in}{0.683022in}}%
\pgfpathlineto{\pgfqpoint{1.167141in}{0.749015in}}%
\pgfpathlineto{\pgfqpoint{1.167457in}{0.757264in}}%
\pgfpathlineto{\pgfqpoint{1.167792in}{0.732517in}}%
\pgfpathlineto{\pgfqpoint{1.168817in}{0.699520in}}%
\pgfpathlineto{\pgfqpoint{1.170475in}{0.732517in}}%
\pgfpathlineto{\pgfqpoint{1.171331in}{0.724268in}}%
\pgfpathlineto{\pgfqpoint{1.171499in}{0.749015in}}%
\pgfpathlineto{\pgfqpoint{1.172989in}{0.806759in}}%
\pgfpathlineto{\pgfqpoint{1.173176in}{0.790261in}}%
\pgfpathlineto{\pgfqpoint{1.173324in}{0.782012in}}%
\pgfpathlineto{\pgfqpoint{1.173679in}{0.823258in}}%
\pgfpathlineto{\pgfqpoint{1.173995in}{0.831507in}}%
\pgfpathlineto{\pgfqpoint{1.174330in}{0.806759in}}%
\pgfpathlineto{\pgfqpoint{1.174665in}{0.782012in}}%
\pgfpathlineto{\pgfqpoint{1.175355in}{0.790261in}}%
\pgfpathlineto{\pgfqpoint{1.175504in}{0.782012in}}%
\pgfpathlineto{\pgfqpoint{1.175671in}{0.806759in}}%
\pgfpathlineto{\pgfqpoint{1.176025in}{0.798510in}}%
\pgfpathlineto{\pgfqpoint{1.177534in}{0.848005in}}%
\pgfpathlineto{\pgfqpoint{1.177701in}{0.839756in}}%
\pgfpathlineto{\pgfqpoint{1.178856in}{0.806759in}}%
\pgfpathlineto{\pgfqpoint{1.177869in}{0.848005in}}%
\pgfpathlineto{\pgfqpoint{1.178875in}{0.823258in}}%
\pgfpathlineto{\pgfqpoint{1.179534in}{0.856254in}}%
\pgfpathlineto{\pgfqpoint{1.180048in}{0.839756in}}%
\pgfpathlineto{\pgfqpoint{1.180216in}{0.848005in}}%
\pgfpathlineto{\pgfqpoint{1.181725in}{0.798510in}}%
\pgfpathlineto{\pgfqpoint{1.183550in}{0.856254in}}%
\pgfpathlineto{\pgfqpoint{1.182209in}{0.782012in}}%
\pgfpathlineto{\pgfqpoint{1.183568in}{0.848005in}}%
\pgfpathlineto{\pgfqpoint{1.184239in}{0.823258in}}%
\pgfpathlineto{\pgfqpoint{1.184723in}{0.839756in}}%
\pgfpathlineto{\pgfqpoint{1.185058in}{0.831507in}}%
\pgfpathlineto{\pgfqpoint{1.185748in}{0.848005in}}%
\pgfpathlineto{\pgfqpoint{1.186064in}{0.831507in}}%
\pgfpathlineto{\pgfqpoint{1.186232in}{0.856254in}}%
\pgfpathlineto{\pgfqpoint{1.186567in}{0.856254in}}%
\pgfpathlineto{\pgfqpoint{1.186902in}{0.905749in}}%
\pgfpathlineto{\pgfqpoint{1.187592in}{0.864503in}}%
\pgfpathlineto{\pgfqpoint{1.190925in}{0.683022in}}%
\pgfpathlineto{\pgfqpoint{1.187759in}{0.872753in}}%
\pgfpathlineto{\pgfqpoint{1.191447in}{0.699520in}}%
\pgfpathlineto{\pgfqpoint{1.192621in}{0.749015in}}%
\pgfpathlineto{\pgfqpoint{1.193105in}{0.732517in}}%
\pgfpathlineto{\pgfqpoint{1.194130in}{0.724268in}}%
\pgfpathlineto{\pgfqpoint{1.194446in}{0.732517in}}%
\pgfpathlineto{\pgfqpoint{1.194781in}{0.707769in}}%
\pgfpathlineto{\pgfqpoint{1.195806in}{0.674773in}}%
\pgfpathlineto{\pgfqpoint{1.196458in}{0.707769in}}%
\pgfpathlineto{\pgfqpoint{1.196644in}{0.691271in}}%
\pgfpathlineto{\pgfqpoint{1.198152in}{0.625278in}}%
\pgfpathlineto{\pgfqpoint{1.198804in}{0.658274in}}%
\pgfpathlineto{\pgfqpoint{1.199159in}{0.625278in}}%
\pgfpathlineto{\pgfqpoint{1.200145in}{0.658274in}}%
\pgfpathlineto{\pgfqpoint{1.200480in}{0.633527in}}%
\pgfpathlineto{\pgfqpoint{1.201505in}{0.625278in}}%
\pgfpathlineto{\pgfqpoint{1.203014in}{0.699520in}}%
\pgfpathlineto{\pgfqpoint{1.203349in}{0.666524in}}%
\pgfpathlineto{\pgfqpoint{1.204504in}{0.633527in}}%
\pgfpathlineto{\pgfqpoint{1.204522in}{0.650025in}}%
\pgfpathlineto{\pgfqpoint{1.206515in}{0.707769in}}%
\pgfpathlineto{\pgfqpoint{1.207037in}{0.699520in}}%
\pgfpathlineto{\pgfqpoint{1.207372in}{0.749015in}}%
\pgfpathlineto{\pgfqpoint{1.207688in}{0.757264in}}%
\pgfpathlineto{\pgfqpoint{1.208024in}{0.732517in}}%
\pgfpathlineto{\pgfqpoint{1.209048in}{0.699520in}}%
\pgfpathlineto{\pgfqpoint{1.209868in}{0.707769in}}%
\pgfpathlineto{\pgfqpoint{1.209887in}{0.699520in}}%
\pgfpathlineto{\pgfqpoint{1.210725in}{0.650025in}}%
\pgfpathlineto{\pgfqpoint{1.211209in}{0.683022in}}%
\pgfpathlineto{\pgfqpoint{1.211731in}{0.749015in}}%
\pgfpathlineto{\pgfqpoint{1.212234in}{0.724268in}}%
\pgfpathlineto{\pgfqpoint{1.212885in}{0.782012in}}%
\pgfpathlineto{\pgfqpoint{1.213891in}{0.757264in}}%
\pgfpathlineto{\pgfqpoint{1.214915in}{0.749015in}}%
\pgfpathlineto{\pgfqpoint{1.216089in}{0.773763in}}%
\pgfpathlineto{\pgfqpoint{1.216256in}{0.765513in}}%
\pgfpathlineto{\pgfqpoint{1.216741in}{0.757264in}}%
\pgfpathlineto{\pgfqpoint{1.216908in}{0.782012in}}%
\pgfpathlineto{\pgfqpoint{1.216927in}{0.773763in}}%
\pgfpathlineto{\pgfqpoint{1.217411in}{0.806759in}}%
\pgfpathlineto{\pgfqpoint{1.218081in}{0.782012in}}%
\pgfpathlineto{\pgfqpoint{1.218752in}{0.757264in}}%
\pgfpathlineto{\pgfqpoint{1.219106in}{0.798510in}}%
\pgfpathlineto{\pgfqpoint{1.219422in}{0.806759in}}%
\pgfpathlineto{\pgfqpoint{1.219758in}{0.782012in}}%
\pgfpathlineto{\pgfqpoint{1.219944in}{0.790261in}}%
\pgfpathlineto{\pgfqpoint{1.220112in}{0.798510in}}%
\pgfpathlineto{\pgfqpoint{1.221099in}{0.782012in}}%
\pgfpathlineto{\pgfqpoint{1.221602in}{0.806759in}}%
\pgfpathlineto{\pgfqpoint{1.222123in}{0.798510in}}%
\pgfpathlineto{\pgfqpoint{1.222775in}{0.831507in}}%
\pgfpathlineto{\pgfqpoint{1.223278in}{0.806759in}}%
\pgfpathlineto{\pgfqpoint{1.224303in}{0.773763in}}%
\pgfpathlineto{\pgfqpoint{1.224787in}{0.806759in}}%
\pgfpathlineto{\pgfqpoint{1.225458in}{0.782012in}}%
\pgfpathlineto{\pgfqpoint{1.226482in}{0.773763in}}%
\pgfpathlineto{\pgfqpoint{1.227320in}{0.872753in}}%
\pgfpathlineto{\pgfqpoint{1.228661in}{0.839756in}}%
\pgfpathlineto{\pgfqpoint{1.230338in}{0.773763in}}%
\pgfpathlineto{\pgfqpoint{1.230505in}{0.798510in}}%
\pgfpathlineto{\pgfqpoint{1.230822in}{0.806759in}}%
\pgfpathlineto{\pgfqpoint{1.231164in}{0.782012in}}%
\pgfpathlineto{\pgfqpoint{1.231511in}{0.773763in}}%
\pgfpathlineto{\pgfqpoint{1.231827in}{0.806759in}}%
\pgfpathlineto{\pgfqpoint{1.232181in}{0.798510in}}%
\pgfpathlineto{\pgfqpoint{1.236521in}{1.004739in}}%
\pgfpathlineto{\pgfqpoint{1.236707in}{0.988241in}}%
\pgfpathlineto{\pgfqpoint{1.237210in}{0.971742in}}%
\pgfpathlineto{\pgfqpoint{1.237378in}{0.996490in}}%
\pgfpathlineto{\pgfqpoint{1.237694in}{0.988241in}}%
\pgfpathlineto{\pgfqpoint{1.238552in}{0.996490in}}%
\pgfpathlineto{\pgfqpoint{1.238384in}{0.971742in}}%
\pgfpathlineto{\pgfqpoint{1.238719in}{0.988241in}}%
\pgfpathlineto{\pgfqpoint{1.239222in}{0.971742in}}%
\pgfpathlineto{\pgfqpoint{1.239390in}{0.996490in}}%
\pgfpathlineto{\pgfqpoint{1.239873in}{1.004739in}}%
\pgfpathlineto{\pgfqpoint{1.240060in}{0.988241in}}%
\pgfpathlineto{\pgfqpoint{1.241736in}{0.946995in}}%
\pgfpathlineto{\pgfqpoint{1.243077in}{0.996490in}}%
\pgfpathlineto{\pgfqpoint{1.243394in}{0.979992in}}%
\pgfpathlineto{\pgfqpoint{1.244251in}{0.971742in}}%
\pgfpathlineto{\pgfqpoint{1.244419in}{0.996490in}}%
\pgfpathlineto{\pgfqpoint{1.246076in}{1.029487in}}%
\pgfpathlineto{\pgfqpoint{1.247101in}{1.021237in}}%
\pgfpathlineto{\pgfqpoint{1.247603in}{0.996490in}}%
\pgfpathlineto{\pgfqpoint{1.247920in}{1.029487in}}%
\pgfpathlineto{\pgfqpoint{1.247939in}{1.021237in}}%
\pgfpathlineto{\pgfqpoint{1.249112in}{1.070732in}}%
\pgfpathlineto{\pgfqpoint{1.249448in}{1.045985in}}%
\pgfpathlineto{\pgfqpoint{1.250789in}{1.021237in}}%
\pgfpathlineto{\pgfqpoint{1.250938in}{1.029487in}}%
\pgfpathlineto{\pgfqpoint{1.251272in}{1.004739in}}%
\pgfpathlineto{\pgfqpoint{1.251459in}{1.012988in}}%
\pgfpathlineto{\pgfqpoint{1.251943in}{1.004739in}}%
\pgfpathlineto{\pgfqpoint{1.252110in}{1.029487in}}%
\pgfpathlineto{\pgfqpoint{1.252130in}{1.021237in}}%
\pgfpathlineto{\pgfqpoint{1.253117in}{1.078981in}}%
\pgfpathlineto{\pgfqpoint{1.253471in}{1.037736in}}%
\pgfpathlineto{\pgfqpoint{1.253806in}{1.021237in}}%
\pgfpathlineto{\pgfqpoint{1.254122in}{1.062483in}}%
\pgfpathlineto{\pgfqpoint{1.254141in}{1.070732in}}%
\pgfpathlineto{\pgfqpoint{1.254811in}{1.021237in}}%
\pgfpathlineto{\pgfqpoint{1.255128in}{1.029487in}}%
\pgfpathlineto{\pgfqpoint{1.255650in}{1.021237in}}%
\pgfpathlineto{\pgfqpoint{1.255967in}{1.054234in}}%
\pgfpathlineto{\pgfqpoint{1.256153in}{1.045985in}}%
\pgfpathlineto{\pgfqpoint{1.256823in}{1.120227in}}%
\pgfpathlineto{\pgfqpoint{1.258146in}{1.103729in}}%
\pgfpathlineto{\pgfqpoint{1.258816in}{1.078981in}}%
\pgfpathlineto{\pgfqpoint{1.259170in}{1.087231in}}%
\pgfpathlineto{\pgfqpoint{1.259338in}{1.095480in}}%
\pgfpathlineto{\pgfqpoint{1.260827in}{1.054234in}}%
\pgfpathlineto{\pgfqpoint{1.260995in}{1.078981in}}%
\pgfpathlineto{\pgfqpoint{1.261852in}{1.070732in}}%
\pgfpathlineto{\pgfqpoint{1.263174in}{1.054234in}}%
\pgfpathlineto{\pgfqpoint{1.263193in}{1.095480in}}%
\pgfpathlineto{\pgfqpoint{1.264199in}{1.070732in}}%
\pgfpathlineto{\pgfqpoint{1.265354in}{1.078981in}}%
\pgfpathlineto{\pgfqpoint{1.265876in}{1.070732in}}%
\pgfpathlineto{\pgfqpoint{1.266192in}{1.103729in}}%
\pgfpathlineto{\pgfqpoint{1.266378in}{1.095480in}}%
\pgfpathlineto{\pgfqpoint{1.267197in}{1.153224in}}%
\pgfpathlineto{\pgfqpoint{1.267552in}{1.111978in}}%
\pgfpathlineto{\pgfqpoint{1.271407in}{0.650025in}}%
\pgfpathlineto{\pgfqpoint{1.271575in}{0.674773in}}%
\pgfpathlineto{\pgfqpoint{1.272916in}{0.650025in}}%
\pgfpathlineto{\pgfqpoint{1.273064in}{0.658274in}}%
\pgfpathlineto{\pgfqpoint{1.273400in}{0.633527in}}%
\pgfpathlineto{\pgfqpoint{1.273586in}{0.641776in}}%
\pgfpathlineto{\pgfqpoint{1.273735in}{0.633527in}}%
\pgfpathlineto{\pgfqpoint{1.274238in}{0.666524in}}%
\pgfpathlineto{\pgfqpoint{1.275244in}{0.683022in}}%
\pgfpathlineto{\pgfqpoint{1.275263in}{0.674773in}}%
\pgfpathlineto{\pgfqpoint{1.275933in}{0.625278in}}%
\pgfpathlineto{\pgfqpoint{1.276269in}{0.674773in}}%
\pgfpathlineto{\pgfqpoint{1.276603in}{0.724268in}}%
\pgfpathlineto{\pgfqpoint{1.277423in}{0.683022in}}%
\pgfpathlineto{\pgfqpoint{1.277777in}{0.674773in}}%
\pgfpathlineto{\pgfqpoint{1.278093in}{0.707769in}}%
\pgfpathlineto{\pgfqpoint{1.278448in}{0.699520in}}%
\pgfpathlineto{\pgfqpoint{1.280105in}{0.732517in}}%
\pgfpathlineto{\pgfqpoint{1.280776in}{0.707769in}}%
\pgfpathlineto{\pgfqpoint{1.280794in}{0.749015in}}%
\pgfpathlineto{\pgfqpoint{1.281130in}{0.749015in}}%
\pgfpathlineto{\pgfqpoint{1.281781in}{0.732517in}}%
\pgfpathlineto{\pgfqpoint{1.281949in}{0.757264in}}%
\pgfpathlineto{\pgfqpoint{1.281968in}{0.749015in}}%
\pgfpathlineto{\pgfqpoint{1.282955in}{0.757264in}}%
\pgfpathlineto{\pgfqpoint{1.282974in}{0.749015in}}%
\pgfpathlineto{\pgfqpoint{1.285302in}{0.683022in}}%
\pgfpathlineto{\pgfqpoint{1.285805in}{0.707769in}}%
\pgfpathlineto{\pgfqpoint{1.286327in}{0.666524in}}%
\pgfpathlineto{\pgfqpoint{1.287984in}{0.625278in}}%
\pgfpathlineto{\pgfqpoint{1.289157in}{0.633527in}}%
\pgfpathlineto{\pgfqpoint{1.289325in}{0.625278in}}%
\pgfpathlineto{\pgfqpoint{1.290182in}{0.625278in}}%
\pgfpathlineto{\pgfqpoint{1.291020in}{0.650025in}}%
\pgfpathlineto{\pgfqpoint{1.291336in}{0.633527in}}%
\pgfpathlineto{\pgfqpoint{1.292361in}{0.625278in}}%
\pgfpathlineto{\pgfqpoint{1.293529in}{0.633527in}}%
\pgfpathlineto{\pgfqpoint{1.294205in}{0.625278in}}%
\pgfpathlineto{\pgfqpoint{1.294528in}{0.658274in}}%
\pgfpathlineto{\pgfqpoint{1.294540in}{0.650025in}}%
\pgfpathlineto{\pgfqpoint{1.296379in}{0.732517in}}%
\pgfpathlineto{\pgfqpoint{1.296715in}{0.707769in}}%
\pgfpathlineto{\pgfqpoint{1.297714in}{0.683022in}}%
\pgfpathlineto{\pgfqpoint{1.297725in}{0.699520in}}%
\pgfpathlineto{\pgfqpoint{1.299893in}{0.806759in}}%
\pgfpathlineto{\pgfqpoint{1.300407in}{0.790261in}}%
\pgfpathlineto{\pgfqpoint{1.300563in}{0.782012in}}%
\pgfpathlineto{\pgfqpoint{1.300898in}{0.815008in}}%
\pgfpathlineto{\pgfqpoint{1.301916in}{0.848005in}}%
\pgfpathlineto{\pgfqpoint{1.302232in}{0.831507in}}%
\pgfpathlineto{\pgfqpoint{1.302574in}{0.864503in}}%
\pgfpathlineto{\pgfqpoint{1.303071in}{0.881002in}}%
\pgfpathlineto{\pgfqpoint{1.303593in}{0.839756in}}%
\pgfpathlineto{\pgfqpoint{1.303760in}{0.848005in}}%
\pgfpathlineto{\pgfqpoint{1.304747in}{0.831507in}}%
\pgfpathlineto{\pgfqpoint{1.305772in}{0.897500in}}%
\pgfpathlineto{\pgfqpoint{1.306423in}{0.881002in}}%
\pgfpathlineto{\pgfqpoint{1.306598in}{0.905749in}}%
\pgfpathlineto{\pgfqpoint{1.306610in}{0.897500in}}%
\pgfpathlineto{\pgfqpoint{1.307094in}{0.930497in}}%
\pgfpathlineto{\pgfqpoint{1.307764in}{0.905749in}}%
\pgfpathlineto{\pgfqpoint{1.308789in}{0.872753in}}%
\pgfpathlineto{\pgfqpoint{1.309112in}{0.905749in}}%
\pgfpathlineto{\pgfqpoint{1.309950in}{0.881002in}}%
\pgfpathlineto{\pgfqpoint{1.310968in}{0.848005in}}%
\pgfpathlineto{\pgfqpoint{1.311452in}{0.856254in}}%
\pgfpathlineto{\pgfqpoint{1.311639in}{0.839756in}}%
\pgfpathlineto{\pgfqpoint{1.312812in}{0.823258in}}%
\pgfpathlineto{\pgfqpoint{1.313818in}{0.897500in}}%
\pgfpathlineto{\pgfqpoint{1.314469in}{0.856254in}}%
\pgfpathlineto{\pgfqpoint{1.315494in}{0.823258in}}%
\pgfpathlineto{\pgfqpoint{1.317327in}{0.881002in}}%
\pgfpathlineto{\pgfqpoint{1.317338in}{0.872753in}}%
\pgfpathlineto{\pgfqpoint{1.317990in}{0.831507in}}%
\pgfpathlineto{\pgfqpoint{1.318344in}{0.872753in}}%
\pgfpathlineto{\pgfqpoint{1.318660in}{0.881002in}}%
\pgfpathlineto{\pgfqpoint{1.318995in}{0.856254in}}%
\pgfpathlineto{\pgfqpoint{1.319517in}{0.823258in}}%
\pgfpathlineto{\pgfqpoint{1.320020in}{0.848005in}}%
\pgfpathlineto{\pgfqpoint{1.320169in}{0.856254in}}%
\pgfpathlineto{\pgfqpoint{1.320504in}{0.831507in}}%
\pgfpathlineto{\pgfqpoint{1.320691in}{0.839756in}}%
\pgfpathlineto{\pgfqpoint{1.323205in}{0.724268in}}%
\pgfpathlineto{\pgfqpoint{1.321361in}{0.848005in}}%
\pgfpathlineto{\pgfqpoint{1.323373in}{0.749015in}}%
\pgfpathlineto{\pgfqpoint{1.323857in}{0.757264in}}%
\pgfpathlineto{\pgfqpoint{1.324081in}{0.732517in}}%
\pgfpathlineto{\pgfqpoint{1.324211in}{0.749015in}}%
\pgfpathlineto{\pgfqpoint{1.325422in}{0.732517in}}%
\pgfpathlineto{\pgfqpoint{1.326706in}{0.757264in}}%
\pgfpathlineto{\pgfqpoint{1.326725in}{0.749015in}}%
\pgfpathlineto{\pgfqpoint{1.327880in}{0.757264in}}%
\pgfpathlineto{\pgfqpoint{1.328234in}{0.749015in}}%
\pgfpathlineto{\pgfqpoint{1.328886in}{0.790261in}}%
\pgfpathlineto{\pgfqpoint{1.329240in}{0.773763in}}%
\pgfpathlineto{\pgfqpoint{1.329911in}{0.798510in}}%
\pgfpathlineto{\pgfqpoint{1.331065in}{0.806759in}}%
\pgfpathlineto{\pgfqpoint{1.331084in}{0.798510in}}%
\pgfpathlineto{\pgfqpoint{1.331735in}{0.856254in}}%
\pgfpathlineto{\pgfqpoint{1.331922in}{0.848005in}}%
\pgfpathlineto{\pgfqpoint{1.332406in}{0.856254in}}%
\pgfpathlineto{\pgfqpoint{1.332593in}{0.839756in}}%
\pgfpathlineto{\pgfqpoint{1.332760in}{0.848005in}}%
\pgfpathlineto{\pgfqpoint{1.333747in}{0.831507in}}%
\pgfpathlineto{\pgfqpoint{1.333766in}{0.848005in}}%
\pgfpathlineto{\pgfqpoint{1.334604in}{0.798510in}}%
\pgfpathlineto{\pgfqpoint{1.334772in}{0.798510in}}%
\pgfpathlineto{\pgfqpoint{1.335423in}{0.782012in}}%
\pgfpathlineto{\pgfqpoint{1.335598in}{0.806759in}}%
\pgfpathlineto{\pgfqpoint{1.335610in}{0.798510in}}%
\pgfpathlineto{\pgfqpoint{1.336783in}{0.848005in}}%
\pgfpathlineto{\pgfqpoint{1.337268in}{0.831507in}}%
\pgfpathlineto{\pgfqpoint{1.337286in}{0.823258in}}%
\pgfpathlineto{\pgfqpoint{1.337777in}{0.881002in}}%
\pgfpathlineto{\pgfqpoint{1.338292in}{0.864503in}}%
\pgfpathlineto{\pgfqpoint{1.338776in}{0.856254in}}%
\pgfpathlineto{\pgfqpoint{1.338795in}{0.897500in}}%
\pgfpathlineto{\pgfqpoint{1.339782in}{0.881002in}}%
\pgfpathlineto{\pgfqpoint{1.339782in}{0.889251in}}%
\pgfpathlineto{\pgfqpoint{1.340285in}{0.955244in}}%
\pgfpathlineto{\pgfqpoint{1.340807in}{0.922247in}}%
\pgfpathlineto{\pgfqpoint{1.341793in}{0.881002in}}%
\pgfpathlineto{\pgfqpoint{1.342148in}{0.897500in}}%
\pgfpathlineto{\pgfqpoint{1.343153in}{0.946995in}}%
\pgfpathlineto{\pgfqpoint{1.343805in}{0.930497in}}%
\pgfpathlineto{\pgfqpoint{1.343824in}{0.922247in}}%
\pgfpathlineto{\pgfqpoint{1.344140in}{0.955244in}}%
\pgfpathlineto{\pgfqpoint{1.344829in}{0.946995in}}%
\pgfpathlineto{\pgfqpoint{1.345984in}{0.955244in}}%
\pgfpathlineto{\pgfqpoint{1.347008in}{0.946995in}}%
\pgfpathlineto{\pgfqpoint{1.348164in}{0.955244in}}%
\pgfpathlineto{\pgfqpoint{1.348517in}{0.946995in}}%
\pgfpathlineto{\pgfqpoint{1.348834in}{0.979992in}}%
\pgfpathlineto{\pgfqpoint{1.349188in}{0.971742in}}%
\pgfpathlineto{\pgfqpoint{1.350845in}{1.004739in}}%
\pgfpathlineto{\pgfqpoint{1.351703in}{0.996490in}}%
\pgfpathlineto{\pgfqpoint{1.351870in}{1.021237in}}%
\pgfpathlineto{\pgfqpoint{1.353024in}{1.078981in}}%
\pgfpathlineto{\pgfqpoint{1.353379in}{1.045985in}}%
\pgfpathlineto{\pgfqpoint{1.353920in}{1.029487in}}%
\pgfpathlineto{\pgfqpoint{1.354031in}{1.062483in}}%
\pgfpathlineto{\pgfqpoint{1.354255in}{1.054234in}}%
\pgfpathlineto{\pgfqpoint{1.355203in}{1.103729in}}%
\pgfpathlineto{\pgfqpoint{1.355223in}{1.095480in}}%
\pgfpathlineto{\pgfqpoint{1.355725in}{1.144975in}}%
\pgfpathlineto{\pgfqpoint{1.356061in}{1.144975in}}%
\pgfpathlineto{\pgfqpoint{1.356899in}{1.169722in}}%
\pgfpathlineto{\pgfqpoint{1.357215in}{1.153224in}}%
\pgfpathlineto{\pgfqpoint{1.358575in}{1.045985in}}%
\pgfpathlineto{\pgfqpoint{1.359581in}{1.070732in}}%
\pgfpathlineto{\pgfqpoint{1.359749in}{1.062483in}}%
\pgfpathlineto{\pgfqpoint{1.360922in}{1.045985in}}%
\pgfpathlineto{\pgfqpoint{1.361239in}{1.054234in}}%
\pgfpathlineto{\pgfqpoint{1.361574in}{1.029487in}}%
\pgfpathlineto{\pgfqpoint{1.362599in}{0.996490in}}%
\pgfpathlineto{\pgfqpoint{1.365261in}{1.103729in}}%
\pgfpathlineto{\pgfqpoint{1.365281in}{1.095480in}}%
\pgfpathlineto{\pgfqpoint{1.366957in}{0.963493in}}%
\pgfpathlineto{\pgfqpoint{1.370142in}{0.625278in}}%
\pgfpathlineto{\pgfqpoint{1.370290in}{0.633527in}}%
\pgfpathlineto{\pgfqpoint{1.371128in}{0.683022in}}%
\pgfpathlineto{\pgfqpoint{1.371315in}{0.674773in}}%
\pgfpathlineto{\pgfqpoint{1.371986in}{0.699520in}}%
\pgfpathlineto{\pgfqpoint{1.372470in}{0.683022in}}%
\pgfpathlineto{\pgfqpoint{1.373308in}{0.658274in}}%
\pgfpathlineto{\pgfqpoint{1.373495in}{0.666524in}}%
\pgfpathlineto{\pgfqpoint{1.373643in}{0.658274in}}%
\pgfpathlineto{\pgfqpoint{1.373662in}{0.699520in}}%
\pgfpathlineto{\pgfqpoint{1.373998in}{0.699520in}}%
\pgfpathlineto{\pgfqpoint{1.374500in}{0.674773in}}%
\pgfpathlineto{\pgfqpoint{1.374816in}{0.707769in}}%
\pgfpathlineto{\pgfqpoint{1.374836in}{0.699520in}}%
\pgfpathlineto{\pgfqpoint{1.375152in}{0.707769in}}%
\pgfpathlineto{\pgfqpoint{1.375487in}{0.683022in}}%
\pgfpathlineto{\pgfqpoint{1.376157in}{0.658274in}}%
\pgfpathlineto{\pgfqpoint{1.376493in}{0.716019in}}%
\pgfpathlineto{\pgfqpoint{1.376512in}{0.724268in}}%
\pgfpathlineto{\pgfqpoint{1.377182in}{0.674773in}}%
\pgfpathlineto{\pgfqpoint{1.377517in}{0.699520in}}%
\pgfpathlineto{\pgfqpoint{1.377834in}{0.707769in}}%
\pgfpathlineto{\pgfqpoint{1.378169in}{0.683022in}}%
\pgfpathlineto{\pgfqpoint{1.378524in}{0.674773in}}%
\pgfpathlineto{\pgfqpoint{1.378840in}{0.707769in}}%
\pgfpathlineto{\pgfqpoint{1.379194in}{0.699520in}}%
\pgfpathlineto{\pgfqpoint{1.379697in}{0.674773in}}%
\pgfpathlineto{\pgfqpoint{1.380013in}{0.707769in}}%
\pgfpathlineto{\pgfqpoint{1.380032in}{0.699520in}}%
\pgfpathlineto{\pgfqpoint{1.380181in}{0.707769in}}%
\pgfpathlineto{\pgfqpoint{1.380535in}{0.674773in}}%
\pgfpathlineto{\pgfqpoint{1.380684in}{0.683022in}}%
\pgfpathlineto{\pgfqpoint{1.380703in}{0.674773in}}%
\pgfpathlineto{\pgfqpoint{1.381690in}{0.716019in}}%
\pgfpathlineto{\pgfqpoint{1.382546in}{0.773763in}}%
\pgfpathlineto{\pgfqpoint{1.382870in}{0.757264in}}%
\pgfpathlineto{\pgfqpoint{1.383217in}{0.749015in}}%
\pgfpathlineto{\pgfqpoint{1.383875in}{0.806759in}}%
\pgfpathlineto{\pgfqpoint{1.383887in}{0.798510in}}%
\pgfpathlineto{\pgfqpoint{1.384055in}{0.823258in}}%
\pgfpathlineto{\pgfqpoint{1.384558in}{0.790261in}}%
\pgfpathlineto{\pgfqpoint{1.386105in}{0.757264in}}%
\pgfpathlineto{\pgfqpoint{1.386719in}{0.806759in}}%
\pgfpathlineto{\pgfqpoint{1.387073in}{0.765513in}}%
\pgfpathlineto{\pgfqpoint{1.387575in}{0.749015in}}%
\pgfpathlineto{\pgfqpoint{1.387743in}{0.773763in}}%
\pgfpathlineto{\pgfqpoint{1.388904in}{0.782012in}}%
\pgfpathlineto{\pgfqpoint{1.389420in}{0.773763in}}%
\pgfpathlineto{\pgfqpoint{1.389587in}{0.798510in}}%
\pgfpathlineto{\pgfqpoint{1.389922in}{0.798510in}}%
\pgfpathlineto{\pgfqpoint{1.390928in}{0.848005in}}%
\pgfpathlineto{\pgfqpoint{1.391587in}{0.831507in}}%
\pgfpathlineto{\pgfqpoint{1.391599in}{0.823258in}}%
\pgfpathlineto{\pgfqpoint{1.391922in}{0.856254in}}%
\pgfpathlineto{\pgfqpoint{1.392604in}{0.848005in}}%
\pgfpathlineto{\pgfqpoint{1.393778in}{0.897500in}}%
\pgfpathlineto{\pgfqpoint{1.394113in}{0.872753in}}%
\pgfpathlineto{\pgfqpoint{1.394429in}{0.856254in}}%
\pgfpathlineto{\pgfqpoint{1.394765in}{0.889251in}}%
\pgfpathlineto{\pgfqpoint{1.395945in}{0.930497in}}%
\pgfpathlineto{\pgfqpoint{1.397111in}{0.856254in}}%
\pgfpathlineto{\pgfqpoint{1.397130in}{0.872753in}}%
\pgfpathlineto{\pgfqpoint{1.397615in}{0.905749in}}%
\pgfpathlineto{\pgfqpoint{1.398285in}{0.881002in}}%
\pgfpathlineto{\pgfqpoint{1.398620in}{0.905749in}}%
\pgfpathlineto{\pgfqpoint{1.399310in}{0.872753in}}%
\pgfpathlineto{\pgfqpoint{1.399626in}{0.856254in}}%
\pgfpathlineto{\pgfqpoint{1.399961in}{0.889251in}}%
\pgfpathlineto{\pgfqpoint{1.400967in}{0.905749in}}%
\pgfpathlineto{\pgfqpoint{1.400483in}{0.872753in}}%
\pgfpathlineto{\pgfqpoint{1.400986in}{0.897500in}}%
\pgfpathlineto{\pgfqpoint{1.402148in}{0.905749in}}%
\pgfpathlineto{\pgfqpoint{1.403333in}{0.848005in}}%
\pgfpathlineto{\pgfqpoint{1.404171in}{0.823258in}}%
\pgfpathlineto{\pgfqpoint{1.404338in}{0.848005in}}%
\pgfpathlineto{\pgfqpoint{1.404997in}{0.856254in}}%
\pgfpathlineto{\pgfqpoint{1.405177in}{0.839756in}}%
\pgfpathlineto{\pgfqpoint{1.406003in}{0.806759in}}%
\pgfpathlineto{\pgfqpoint{1.406350in}{0.823258in}}%
\pgfpathlineto{\pgfqpoint{1.408015in}{0.856254in}}%
\pgfpathlineto{\pgfqpoint{1.409032in}{0.823258in}}%
\pgfpathlineto{\pgfqpoint{1.409516in}{0.831507in}}%
\pgfpathlineto{\pgfqpoint{1.409703in}{0.815008in}}%
\pgfpathlineto{\pgfqpoint{1.409852in}{0.806759in}}%
\pgfpathlineto{\pgfqpoint{1.410362in}{0.839756in}}%
\pgfpathlineto{\pgfqpoint{1.410374in}{0.848005in}}%
\pgfpathlineto{\pgfqpoint{1.410708in}{0.823258in}}%
\pgfpathlineto{\pgfqpoint{1.411379in}{0.823258in}}%
\pgfpathlineto{\pgfqpoint{1.412198in}{0.732517in}}%
\pgfpathlineto{\pgfqpoint{1.412385in}{0.724268in}}%
\pgfpathlineto{\pgfqpoint{1.412701in}{0.757264in}}%
\pgfpathlineto{\pgfqpoint{1.413223in}{0.749015in}}%
\pgfpathlineto{\pgfqpoint{1.414210in}{0.757264in}}%
\pgfpathlineto{\pgfqpoint{1.414229in}{0.749015in}}%
\pgfpathlineto{\pgfqpoint{1.414899in}{0.724268in}}%
\pgfpathlineto{\pgfqpoint{1.415067in}{0.749015in}}%
\pgfpathlineto{\pgfqpoint{1.416724in}{0.782012in}}%
\pgfpathlineto{\pgfqpoint{1.416743in}{0.773763in}}%
\pgfpathlineto{\pgfqpoint{1.417395in}{0.831507in}}%
\pgfpathlineto{\pgfqpoint{1.417749in}{0.798510in}}%
\pgfpathlineto{\pgfqpoint{1.419425in}{0.848005in}}%
\pgfpathlineto{\pgfqpoint{1.419593in}{0.839756in}}%
\pgfpathlineto{\pgfqpoint{1.420766in}{0.798510in}}%
\pgfpathlineto{\pgfqpoint{1.420934in}{0.823258in}}%
\pgfpathlineto{\pgfqpoint{1.421270in}{0.848005in}}%
\pgfpathlineto{\pgfqpoint{1.422089in}{0.831507in}}%
\pgfpathlineto{\pgfqpoint{1.422778in}{0.823258in}}%
\pgfpathlineto{\pgfqpoint{1.423101in}{0.856254in}}%
\pgfpathlineto{\pgfqpoint{1.423113in}{0.848005in}}%
\pgfpathlineto{\pgfqpoint{1.423765in}{0.856254in}}%
\pgfpathlineto{\pgfqpoint{1.423951in}{0.839756in}}%
\pgfpathlineto{\pgfqpoint{1.424957in}{0.823258in}}%
\pgfpathlineto{\pgfqpoint{1.424275in}{0.856254in}}%
\pgfpathlineto{\pgfqpoint{1.425106in}{0.831507in}}%
\pgfpathlineto{\pgfqpoint{1.425616in}{0.881002in}}%
\pgfpathlineto{\pgfqpoint{1.426130in}{0.848005in}}%
\pgfpathlineto{\pgfqpoint{1.426447in}{0.831507in}}%
\pgfpathlineto{\pgfqpoint{1.426782in}{0.864503in}}%
\pgfpathlineto{\pgfqpoint{1.427956in}{0.955244in}}%
\pgfpathlineto{\pgfqpoint{1.427975in}{0.946995in}}%
\pgfpathlineto{\pgfqpoint{1.428626in}{0.930497in}}%
\pgfpathlineto{\pgfqpoint{1.429129in}{0.955244in}}%
\pgfpathlineto{\pgfqpoint{1.430154in}{0.946995in}}%
\pgfpathlineto{\pgfqpoint{1.432314in}{1.054234in}}%
\pgfpathlineto{\pgfqpoint{1.432649in}{1.029487in}}%
\pgfpathlineto{\pgfqpoint{1.433004in}{1.021237in}}%
\pgfpathlineto{\pgfqpoint{1.433320in}{1.054234in}}%
\pgfpathlineto{\pgfqpoint{1.433674in}{1.045985in}}%
\pgfpathlineto{\pgfqpoint{1.434325in}{1.054234in}}%
\pgfpathlineto{\pgfqpoint{1.434512in}{1.037736in}}%
\pgfpathlineto{\pgfqpoint{1.435015in}{1.021237in}}%
\pgfpathlineto{\pgfqpoint{1.435183in}{1.045985in}}%
\pgfpathlineto{\pgfqpoint{1.435499in}{1.037736in}}%
\pgfpathlineto{\pgfqpoint{1.435518in}{1.045985in}}%
\pgfpathlineto{\pgfqpoint{1.436524in}{1.012988in}}%
\pgfpathlineto{\pgfqpoint{1.437865in}{0.971742in}}%
\pgfpathlineto{\pgfqpoint{1.438852in}{0.955244in}}%
\pgfpathlineto{\pgfqpoint{1.438871in}{0.971742in}}%
\pgfpathlineto{\pgfqpoint{1.440025in}{0.979992in}}%
\pgfpathlineto{\pgfqpoint{1.440212in}{0.971742in}}%
\pgfpathlineto{\pgfqpoint{1.440882in}{1.021237in}}%
\pgfpathlineto{\pgfqpoint{1.441050in}{1.012988in}}%
\pgfpathlineto{\pgfqpoint{1.441199in}{1.004739in}}%
\pgfpathlineto{\pgfqpoint{1.441373in}{1.029487in}}%
\pgfpathlineto{\pgfqpoint{1.441720in}{1.021237in}}%
\pgfpathlineto{\pgfqpoint{1.442223in}{1.070732in}}%
\pgfpathlineto{\pgfqpoint{1.442894in}{1.037736in}}%
\pgfpathlineto{\pgfqpoint{1.443396in}{1.021237in}}%
\pgfpathlineto{\pgfqpoint{1.443564in}{1.045985in}}%
\pgfpathlineto{\pgfqpoint{1.445054in}{1.103729in}}%
\pgfpathlineto{\pgfqpoint{1.445241in}{1.087231in}}%
\pgfpathlineto{\pgfqpoint{1.445557in}{1.103729in}}%
\pgfpathlineto{\pgfqpoint{1.446917in}{0.938746in}}%
\pgfpathlineto{\pgfqpoint{1.449915in}{0.658274in}}%
\pgfpathlineto{\pgfqpoint{1.449934in}{0.674773in}}%
\pgfpathlineto{\pgfqpoint{1.450270in}{0.650025in}}%
\pgfpathlineto{\pgfqpoint{1.450940in}{0.650025in}}%
\pgfpathlineto{\pgfqpoint{1.451089in}{0.658274in}}%
\pgfpathlineto{\pgfqpoint{1.451443in}{0.625278in}}%
\pgfpathlineto{\pgfqpoint{1.451591in}{0.633527in}}%
\pgfpathlineto{\pgfqpoint{1.451778in}{0.625278in}}%
\pgfpathlineto{\pgfqpoint{1.452095in}{0.658274in}}%
\pgfpathlineto{\pgfqpoint{1.452616in}{0.650025in}}%
\pgfpathlineto{\pgfqpoint{1.453287in}{0.699520in}}%
\pgfpathlineto{\pgfqpoint{1.454274in}{0.683022in}}%
\pgfpathlineto{\pgfqpoint{1.454963in}{0.650025in}}%
\pgfpathlineto{\pgfqpoint{1.455299in}{0.666524in}}%
\pgfpathlineto{\pgfqpoint{1.456956in}{0.633527in}}%
\pgfpathlineto{\pgfqpoint{1.457962in}{0.683022in}}%
\pgfpathlineto{\pgfqpoint{1.457980in}{0.674773in}}%
\pgfpathlineto{\pgfqpoint{1.458967in}{0.707769in}}%
\pgfpathlineto{\pgfqpoint{1.459154in}{0.691271in}}%
\pgfpathlineto{\pgfqpoint{1.459973in}{0.658274in}}%
\pgfpathlineto{\pgfqpoint{1.459992in}{0.699520in}}%
\pgfpathlineto{\pgfqpoint{1.460308in}{0.683022in}}%
\pgfpathlineto{\pgfqpoint{1.460327in}{0.699520in}}%
\pgfpathlineto{\pgfqpoint{1.460663in}{0.674773in}}%
\pgfpathlineto{\pgfqpoint{1.461333in}{0.674773in}}%
\pgfpathlineto{\pgfqpoint{1.462339in}{0.699520in}}%
\pgfpathlineto{\pgfqpoint{1.462507in}{0.691271in}}%
\pgfpathlineto{\pgfqpoint{1.462991in}{0.683022in}}%
\pgfpathlineto{\pgfqpoint{1.463166in}{0.707769in}}%
\pgfpathlineto{\pgfqpoint{1.463177in}{0.699520in}}%
\pgfpathlineto{\pgfqpoint{1.464686in}{0.773763in}}%
\pgfpathlineto{\pgfqpoint{1.465021in}{0.749015in}}%
\pgfpathlineto{\pgfqpoint{1.466008in}{0.732517in}}%
\pgfpathlineto{\pgfqpoint{1.466026in}{0.749015in}}%
\pgfpathlineto{\pgfqpoint{1.467368in}{0.848005in}}%
\pgfpathlineto{\pgfqpoint{1.468020in}{0.806759in}}%
\pgfpathlineto{\pgfqpoint{1.468541in}{0.798510in}}%
\pgfpathlineto{\pgfqpoint{1.469033in}{0.839756in}}%
\pgfpathlineto{\pgfqpoint{1.469715in}{0.897500in}}%
\pgfpathlineto{\pgfqpoint{1.470050in}{0.897500in}}%
\pgfpathlineto{\pgfqpoint{1.470701in}{0.881002in}}%
\pgfpathlineto{\pgfqpoint{1.470709in}{0.913998in}}%
\pgfpathlineto{\pgfqpoint{1.471716in}{0.979992in}}%
\pgfpathlineto{\pgfqpoint{1.471035in}{0.905749in}}%
\pgfpathlineto{\pgfqpoint{1.471894in}{0.963493in}}%
\pgfpathlineto{\pgfqpoint{1.472042in}{0.955244in}}%
\pgfpathlineto{\pgfqpoint{1.472397in}{0.996490in}}%
\pgfpathlineto{\pgfqpoint{1.473737in}{0.971742in}}%
\pgfpathlineto{\pgfqpoint{1.475736in}{1.054234in}}%
\pgfpathlineto{\pgfqpoint{1.475914in}{1.037736in}}%
\pgfpathlineto{\pgfqpoint{1.476736in}{1.004739in}}%
\pgfpathlineto{\pgfqpoint{1.476758in}{1.045985in}}%
\pgfpathlineto{\pgfqpoint{1.477077in}{1.029487in}}%
\pgfpathlineto{\pgfqpoint{1.477928in}{1.070732in}}%
\pgfpathlineto{\pgfqpoint{1.478098in}{1.062483in}}%
\pgfpathlineto{\pgfqpoint{1.481600in}{0.881002in}}%
\pgfpathlineto{\pgfqpoint{1.482119in}{0.897500in}}%
\pgfpathlineto{\pgfqpoint{1.482793in}{0.922247in}}%
\pgfpathlineto{\pgfqpoint{1.483274in}{0.905749in}}%
\pgfpathlineto{\pgfqpoint{1.483792in}{0.922247in}}%
\pgfpathlineto{\pgfqpoint{1.484296in}{0.897500in}}%
\pgfpathlineto{\pgfqpoint{1.484621in}{0.905749in}}%
\pgfpathlineto{\pgfqpoint{1.484947in}{0.881002in}}%
\pgfpathlineto{\pgfqpoint{1.485288in}{0.856254in}}%
\pgfpathlineto{\pgfqpoint{1.485303in}{0.897500in}}%
\pgfpathlineto{\pgfqpoint{1.485976in}{0.897500in}}%
\pgfpathlineto{\pgfqpoint{1.486643in}{0.946995in}}%
\pgfpathlineto{\pgfqpoint{1.487642in}{0.930497in}}%
\pgfpathlineto{\pgfqpoint{1.488479in}{0.905749in}}%
\pgfpathlineto{\pgfqpoint{1.488657in}{0.913998in}}%
\pgfpathlineto{\pgfqpoint{1.489145in}{0.905749in}}%
\pgfpathlineto{\pgfqpoint{1.489316in}{0.930497in}}%
\pgfpathlineto{\pgfqpoint{1.489330in}{0.922247in}}%
\pgfpathlineto{\pgfqpoint{1.489486in}{0.930497in}}%
\pgfpathlineto{\pgfqpoint{1.489834in}{0.897500in}}%
\pgfpathlineto{\pgfqpoint{1.489989in}{0.905749in}}%
\pgfpathlineto{\pgfqpoint{1.491004in}{0.872753in}}%
\pgfpathlineto{\pgfqpoint{1.493854in}{0.691271in}}%
\pgfpathlineto{\pgfqpoint{1.495365in}{0.625278in}}%
\pgfpathlineto{\pgfqpoint{1.496520in}{0.633527in}}%
\pgfpathlineto{\pgfqpoint{1.496520in}{0.625278in}}%
\pgfpathlineto{\pgfqpoint{1.497519in}{0.625278in}}%
\pgfpathlineto{\pgfqpoint{1.498526in}{0.633527in}}%
\pgfpathlineto{\pgfqpoint{1.498526in}{0.625278in}}%
\pgfpathlineto{\pgfqpoint{1.499704in}{0.633527in}}%
\pgfpathlineto{\pgfqpoint{1.499867in}{0.625278in}}%
\pgfpathlineto{\pgfqpoint{1.500718in}{0.625278in}}%
\pgfpathlineto{\pgfqpoint{1.501881in}{0.633527in}}%
\pgfpathlineto{\pgfqpoint{1.501881in}{0.625278in}}%
\pgfpathlineto{\pgfqpoint{1.502717in}{0.625278in}}%
\pgfpathlineto{\pgfqpoint{1.503894in}{0.633527in}}%
\pgfpathlineto{\pgfqpoint{1.504057in}{0.625278in}}%
\pgfpathlineto{\pgfqpoint{1.504909in}{0.625278in}}%
\pgfpathlineto{\pgfqpoint{1.506071in}{0.633527in}}%
\pgfpathlineto{\pgfqpoint{1.506079in}{0.625278in}}%
\pgfpathlineto{\pgfqpoint{1.507078in}{0.625278in}}%
\pgfpathlineto{\pgfqpoint{1.508418in}{0.633527in}}%
\pgfpathlineto{\pgfqpoint{1.508426in}{0.625278in}}%
\pgfpathlineto{\pgfqpoint{1.509425in}{0.625278in}}%
\pgfpathlineto{\pgfqpoint{1.510595in}{0.633527in}}%
\pgfpathlineto{\pgfqpoint{1.510773in}{0.625278in}}%
\pgfpathlineto{\pgfqpoint{1.511602in}{0.625278in}}%
\pgfpathlineto{\pgfqpoint{1.512780in}{0.633527in}}%
\pgfpathlineto{\pgfqpoint{1.512942in}{0.625278in}}%
\pgfpathlineto{\pgfqpoint{1.513801in}{0.625278in}}%
\pgfpathlineto{\pgfqpoint{1.514623in}{0.658274in}}%
\pgfpathlineto{\pgfqpoint{1.514956in}{0.633527in}}%
\pgfpathlineto{\pgfqpoint{1.515475in}{0.650025in}}%
\pgfpathlineto{\pgfqpoint{1.515815in}{0.625278in}}%
\pgfpathlineto{\pgfqpoint{1.517155in}{0.650025in}}%
\pgfpathlineto{\pgfqpoint{1.517318in}{0.641776in}}%
\pgfpathlineto{\pgfqpoint{1.518140in}{0.633527in}}%
\pgfpathlineto{\pgfqpoint{1.517489in}{0.650025in}}%
\pgfpathlineto{\pgfqpoint{1.518162in}{0.650025in}}%
\pgfpathlineto{\pgfqpoint{1.518666in}{0.674773in}}%
\pgfpathlineto{\pgfqpoint{1.519317in}{0.658274in}}%
\pgfpathlineto{\pgfqpoint{1.519821in}{0.633527in}}%
\pgfpathlineto{\pgfqpoint{1.520317in}{0.707769in}}%
\pgfpathlineto{\pgfqpoint{1.520339in}{0.699520in}}%
\pgfpathlineto{\pgfqpoint{1.521324in}{0.732517in}}%
\pgfpathlineto{\pgfqpoint{1.521509in}{0.716019in}}%
\pgfpathlineto{\pgfqpoint{1.522686in}{0.699520in}}%
\pgfpathlineto{\pgfqpoint{1.523020in}{0.724268in}}%
\pgfpathlineto{\pgfqpoint{1.523523in}{0.699520in}}%
\pgfpathlineto{\pgfqpoint{1.524508in}{0.658274in}}%
\pgfpathlineto{\pgfqpoint{1.524863in}{0.674773in}}%
\pgfpathlineto{\pgfqpoint{1.526522in}{0.707769in}}%
\pgfpathlineto{\pgfqpoint{1.527047in}{0.724268in}}%
\pgfpathlineto{\pgfqpoint{1.527551in}{0.699520in}}%
\pgfpathlineto{\pgfqpoint{1.528047in}{0.724268in}}%
\pgfpathlineto{\pgfqpoint{1.528699in}{0.707769in}}%
\pgfpathlineto{\pgfqpoint{1.529706in}{0.683022in}}%
\pgfpathlineto{\pgfqpoint{1.529728in}{0.699520in}}%
\pgfpathlineto{\pgfqpoint{1.530231in}{0.724268in}}%
\pgfpathlineto{\pgfqpoint{1.530564in}{0.699520in}}%
\pgfpathlineto{\pgfqpoint{1.531216in}{0.658274in}}%
\pgfpathlineto{\pgfqpoint{1.531549in}{0.691271in}}%
\pgfpathlineto{\pgfqpoint{1.531719in}{0.732517in}}%
\pgfpathlineto{\pgfqpoint{1.532408in}{0.674773in}}%
\pgfpathlineto{\pgfqpoint{1.532578in}{0.674773in}}%
\pgfpathlineto{\pgfqpoint{1.533060in}{0.707769in}}%
\pgfpathlineto{\pgfqpoint{1.533733in}{0.683022in}}%
\pgfpathlineto{\pgfqpoint{1.534570in}{0.658274in}}%
\pgfpathlineto{\pgfqpoint{1.534755in}{0.674773in}}%
\pgfpathlineto{\pgfqpoint{1.536081in}{0.707769in}}%
\pgfpathlineto{\pgfqpoint{1.536095in}{0.699520in}}%
\pgfpathlineto{\pgfqpoint{1.537421in}{0.683022in}}%
\pgfpathlineto{\pgfqpoint{1.537584in}{0.707769in}}%
\pgfpathlineto{\pgfqpoint{1.538442in}{0.666524in}}%
\pgfpathlineto{\pgfqpoint{1.540101in}{0.625278in}}%
\pgfpathlineto{\pgfqpoint{1.541108in}{0.633527in}}%
\pgfpathlineto{\pgfqpoint{1.541108in}{0.625278in}}%
\pgfpathlineto{\pgfqpoint{1.542278in}{0.633527in}}%
\pgfpathlineto{\pgfqpoint{1.542448in}{0.625278in}}%
\pgfpathlineto{\pgfqpoint{1.543285in}{0.625278in}}%
\pgfpathlineto{\pgfqpoint{1.544462in}{0.633527in}}%
\pgfpathlineto{\pgfqpoint{1.544625in}{0.625278in}}%
\pgfpathlineto{\pgfqpoint{1.544980in}{0.650025in}}%
\pgfpathlineto{\pgfqpoint{1.545484in}{0.641776in}}%
\pgfpathlineto{\pgfqpoint{1.545965in}{0.633527in}}%
\pgfpathlineto{\pgfqpoint{1.545987in}{0.674773in}}%
\pgfpathlineto{\pgfqpoint{1.547328in}{0.650025in}}%
\pgfpathlineto{\pgfqpoint{1.548490in}{0.658274in}}%
\pgfpathlineto{\pgfqpoint{1.549208in}{0.633527in}}%
\pgfpathlineto{\pgfqpoint{1.549497in}{0.683022in}}%
\pgfpathlineto{\pgfqpoint{1.549504in}{0.674773in}}%
\pgfpathlineto{\pgfqpoint{1.549675in}{0.699520in}}%
\pgfpathlineto{\pgfqpoint{1.550348in}{0.666524in}}%
\pgfpathlineto{\pgfqpoint{1.551163in}{0.633527in}}%
\pgfpathlineto{\pgfqpoint{1.551185in}{0.674773in}}%
\pgfpathlineto{\pgfqpoint{1.551837in}{0.683022in}}%
\pgfpathlineto{\pgfqpoint{1.552022in}{0.666524in}}%
\pgfpathlineto{\pgfqpoint{1.552059in}{0.658274in}}%
\pgfpathlineto{\pgfqpoint{1.552858in}{0.699520in}}%
\pgfpathlineto{\pgfqpoint{1.553347in}{0.707769in}}%
\pgfpathlineto{\pgfqpoint{1.553532in}{0.691271in}}%
\pgfpathlineto{\pgfqpoint{1.554013in}{0.658274in}}%
\pgfpathlineto{\pgfqpoint{1.554702in}{0.674773in}}%
\pgfpathlineto{\pgfqpoint{1.555206in}{0.699520in}}%
\pgfpathlineto{\pgfqpoint{1.555709in}{0.674773in}}%
\pgfpathlineto{\pgfqpoint{1.557034in}{0.633527in}}%
\pgfpathlineto{\pgfqpoint{1.557049in}{0.650025in}}%
\pgfpathlineto{\pgfqpoint{1.557701in}{0.658274in}}%
\pgfpathlineto{\pgfqpoint{1.557886in}{0.641776in}}%
\pgfpathlineto{\pgfqpoint{1.558560in}{0.625278in}}%
\pgfpathlineto{\pgfqpoint{1.558730in}{0.650025in}}%
\pgfpathlineto{\pgfqpoint{1.559878in}{0.658274in}}%
\pgfpathlineto{\pgfqpoint{1.560403in}{0.674773in}}%
\pgfpathlineto{\pgfqpoint{1.560907in}{0.650025in}}%
\pgfpathlineto{\pgfqpoint{1.562062in}{0.658274in}}%
\pgfpathlineto{\pgfqpoint{1.563084in}{0.641776in}}%
\pgfpathlineto{\pgfqpoint{1.563587in}{0.625278in}}%
\pgfpathlineto{\pgfqpoint{1.563757in}{0.650025in}}%
\pgfpathlineto{\pgfqpoint{1.564912in}{0.658274in}}%
\pgfpathlineto{\pgfqpoint{1.565268in}{0.650025in}}%
\pgfpathlineto{\pgfqpoint{1.565579in}{0.683022in}}%
\pgfpathlineto{\pgfqpoint{1.565934in}{0.674773in}}%
\pgfpathlineto{\pgfqpoint{1.566416in}{0.683022in}}%
\pgfpathlineto{\pgfqpoint{1.566608in}{0.666524in}}%
\pgfpathlineto{\pgfqpoint{1.566771in}{0.674773in}}%
\pgfpathlineto{\pgfqpoint{1.568259in}{0.625278in}}%
\pgfpathlineto{\pgfqpoint{1.569436in}{0.633527in}}%
\pgfpathlineto{\pgfqpoint{1.569770in}{0.625278in}}%
\pgfpathlineto{\pgfqpoint{1.570443in}{0.625278in}}%
\pgfpathlineto{\pgfqpoint{1.571613in}{0.633527in}}%
\pgfpathlineto{\pgfqpoint{1.571635in}{0.625278in}}%
\pgfpathlineto{\pgfqpoint{1.572642in}{0.625278in}}%
\pgfpathlineto{\pgfqpoint{1.572805in}{0.650025in}}%
\pgfpathlineto{\pgfqpoint{1.573798in}{0.633527in}}%
\pgfpathlineto{\pgfqpoint{1.573812in}{0.625278in}}%
\pgfpathlineto{\pgfqpoint{1.574819in}{0.674773in}}%
\pgfpathlineto{\pgfqpoint{1.576308in}{0.633527in}}%
\pgfpathlineto{\pgfqpoint{1.576330in}{0.650025in}}%
\pgfpathlineto{\pgfqpoint{1.576981in}{0.683022in}}%
\pgfpathlineto{\pgfqpoint{1.577485in}{0.658274in}}%
\pgfpathlineto{\pgfqpoint{1.577500in}{0.650025in}}%
\pgfpathlineto{\pgfqpoint{1.578003in}{0.699520in}}%
\pgfpathlineto{\pgfqpoint{1.578507in}{0.691271in}}%
\pgfpathlineto{\pgfqpoint{1.578655in}{0.683022in}}%
\pgfpathlineto{\pgfqpoint{1.578825in}{0.707769in}}%
\pgfpathlineto{\pgfqpoint{1.579180in}{0.699520in}}%
\pgfpathlineto{\pgfqpoint{1.580831in}{0.732517in}}%
\pgfpathlineto{\pgfqpoint{1.581861in}{0.724268in}}%
\pgfpathlineto{\pgfqpoint{1.582179in}{0.732517in}}%
\pgfpathlineto{\pgfqpoint{1.582527in}{0.699520in}}%
\pgfpathlineto{\pgfqpoint{1.582845in}{0.683022in}}%
\pgfpathlineto{\pgfqpoint{1.583179in}{0.716019in}}%
\pgfpathlineto{\pgfqpoint{1.583682in}{0.732517in}}%
\pgfpathlineto{\pgfqpoint{1.584038in}{0.699520in}}%
\pgfpathlineto{\pgfqpoint{1.584208in}{0.699520in}}%
\pgfpathlineto{\pgfqpoint{1.585866in}{0.732517in}}%
\pgfpathlineto{\pgfqpoint{1.586888in}{0.699520in}}%
\pgfpathlineto{\pgfqpoint{1.587540in}{0.658274in}}%
\pgfpathlineto{\pgfqpoint{1.587895in}{0.699520in}}%
\pgfpathlineto{\pgfqpoint{1.588547in}{0.683022in}}%
\pgfpathlineto{\pgfqpoint{1.589554in}{0.732517in}}%
\pgfpathlineto{\pgfqpoint{1.590072in}{0.724268in}}%
\pgfpathlineto{\pgfqpoint{1.590390in}{0.757264in}}%
\pgfpathlineto{\pgfqpoint{1.590575in}{0.749015in}}%
\pgfpathlineto{\pgfqpoint{1.591227in}{0.707769in}}%
\pgfpathlineto{\pgfqpoint{1.591582in}{0.749015in}}%
\pgfpathlineto{\pgfqpoint{1.591901in}{0.757264in}}%
\pgfpathlineto{\pgfqpoint{1.592234in}{0.732517in}}%
\pgfpathlineto{\pgfqpoint{1.593426in}{0.699520in}}%
\pgfpathlineto{\pgfqpoint{1.594603in}{0.749015in}}%
\pgfpathlineto{\pgfqpoint{1.595085in}{0.732517in}}%
\pgfpathlineto{\pgfqpoint{1.595943in}{0.724268in}}%
\pgfpathlineto{\pgfqpoint{1.596106in}{0.749015in}}%
\pgfpathlineto{\pgfqpoint{1.596425in}{0.757264in}}%
\pgfpathlineto{\pgfqpoint{1.596758in}{0.732517in}}%
\pgfpathlineto{\pgfqpoint{1.596943in}{0.740766in}}%
\pgfpathlineto{\pgfqpoint{1.597099in}{0.732517in}}%
\pgfpathlineto{\pgfqpoint{1.597113in}{0.773763in}}%
\pgfpathlineto{\pgfqpoint{1.597432in}{0.765513in}}%
\pgfpathlineto{\pgfqpoint{1.597787in}{0.749015in}}%
\pgfpathlineto{\pgfqpoint{1.598602in}{0.806759in}}%
\pgfpathlineto{\pgfqpoint{1.599460in}{0.798510in}}%
\pgfpathlineto{\pgfqpoint{1.599631in}{0.823258in}}%
\pgfpathlineto{\pgfqpoint{1.600971in}{0.798510in}}%
\pgfpathlineto{\pgfqpoint{1.601289in}{0.806759in}}%
\pgfpathlineto{\pgfqpoint{1.601637in}{0.773763in}}%
\pgfpathlineto{\pgfqpoint{1.601785in}{0.782012in}}%
\pgfpathlineto{\pgfqpoint{1.602963in}{0.732517in}}%
\pgfpathlineto{\pgfqpoint{1.602985in}{0.749015in}}%
\pgfpathlineto{\pgfqpoint{1.603636in}{0.782012in}}%
\pgfpathlineto{\pgfqpoint{1.604132in}{0.757264in}}%
\pgfpathlineto{\pgfqpoint{1.605162in}{0.749015in}}%
\pgfpathlineto{\pgfqpoint{1.605976in}{0.757264in}}%
\pgfpathlineto{\pgfqpoint{1.605998in}{0.749015in}}%
\pgfpathlineto{\pgfqpoint{1.606650in}{0.732517in}}%
\pgfpathlineto{\pgfqpoint{1.606650in}{0.765513in}}%
\pgfpathlineto{\pgfqpoint{1.606983in}{0.757264in}}%
\pgfpathlineto{\pgfqpoint{1.607672in}{0.773763in}}%
\pgfpathlineto{\pgfqpoint{1.607990in}{0.757264in}}%
\pgfpathlineto{\pgfqpoint{1.608323in}{0.790261in}}%
\pgfpathlineto{\pgfqpoint{1.608664in}{0.806759in}}%
\pgfpathlineto{\pgfqpoint{1.609330in}{0.757264in}}%
\pgfpathlineto{\pgfqpoint{1.609686in}{0.749015in}}%
\pgfpathlineto{\pgfqpoint{1.610189in}{0.798510in}}%
\pgfpathlineto{\pgfqpoint{1.611862in}{0.848005in}}%
\pgfpathlineto{\pgfqpoint{1.610841in}{0.782012in}}%
\pgfpathlineto{\pgfqpoint{1.612033in}{0.839756in}}%
\pgfpathlineto{\pgfqpoint{1.613188in}{0.806759in}}%
\pgfpathlineto{\pgfqpoint{1.613210in}{0.823258in}}%
\pgfpathlineto{\pgfqpoint{1.614698in}{0.856254in}}%
\pgfpathlineto{\pgfqpoint{1.614713in}{0.848005in}}%
\pgfpathlineto{\pgfqpoint{1.615387in}{0.823258in}}%
\pgfpathlineto{\pgfqpoint{1.615557in}{0.848005in}}%
\pgfpathlineto{\pgfqpoint{1.615890in}{0.872753in}}%
\pgfpathlineto{\pgfqpoint{1.616705in}{0.856254in}}%
\pgfpathlineto{\pgfqpoint{1.617379in}{0.831507in}}%
\pgfpathlineto{\pgfqpoint{1.617401in}{0.872753in}}%
\pgfpathlineto{\pgfqpoint{1.617734in}{0.872753in}}%
\pgfpathlineto{\pgfqpoint{1.618067in}{0.848005in}}%
\pgfpathlineto{\pgfqpoint{1.618548in}{0.864503in}}%
\pgfpathlineto{\pgfqpoint{1.619393in}{0.979992in}}%
\pgfpathlineto{\pgfqpoint{1.619578in}{0.971742in}}%
\pgfpathlineto{\pgfqpoint{1.619896in}{0.979992in}}%
\pgfpathlineto{\pgfqpoint{1.620229in}{0.955244in}}%
\pgfpathlineto{\pgfqpoint{1.620414in}{0.963493in}}%
\pgfpathlineto{\pgfqpoint{1.620748in}{0.946995in}}%
\pgfpathlineto{\pgfqpoint{1.621066in}{0.979992in}}%
\pgfpathlineto{\pgfqpoint{1.621088in}{0.971742in}}%
\pgfpathlineto{\pgfqpoint{1.621903in}{0.979992in}}%
\pgfpathlineto{\pgfqpoint{1.621925in}{0.971742in}}%
\pgfpathlineto{\pgfqpoint{1.623243in}{0.955244in}}%
\pgfpathlineto{\pgfqpoint{1.624087in}{0.979992in}}%
\pgfpathlineto{\pgfqpoint{1.624272in}{0.971742in}}%
\pgfpathlineto{\pgfqpoint{1.625427in}{1.004739in}}%
\pgfpathlineto{\pgfqpoint{1.624923in}{0.955244in}}%
\pgfpathlineto{\pgfqpoint{1.625612in}{0.988241in}}%
\pgfpathlineto{\pgfqpoint{1.626597in}{0.955244in}}%
\pgfpathlineto{\pgfqpoint{1.626782in}{0.971742in}}%
\pgfpathlineto{\pgfqpoint{1.628781in}{1.029487in}}%
\pgfpathlineto{\pgfqpoint{1.629803in}{1.021237in}}%
\pgfpathlineto{\pgfqpoint{1.630121in}{1.029487in}}%
\pgfpathlineto{\pgfqpoint{1.630469in}{0.996490in}}%
\pgfpathlineto{\pgfqpoint{1.631121in}{1.029487in}}%
\pgfpathlineto{\pgfqpoint{1.631624in}{1.004739in}}%
\pgfpathlineto{\pgfqpoint{1.631646in}{0.996490in}}%
\pgfpathlineto{\pgfqpoint{1.631965in}{1.029487in}}%
\pgfpathlineto{\pgfqpoint{1.632653in}{1.021237in}}%
\pgfpathlineto{\pgfqpoint{1.633638in}{1.029487in}}%
\pgfpathlineto{\pgfqpoint{1.633660in}{1.021237in}}%
\pgfpathlineto{\pgfqpoint{1.633971in}{1.004739in}}%
\pgfpathlineto{\pgfqpoint{1.634142in}{1.029487in}}%
\pgfpathlineto{\pgfqpoint{1.634497in}{1.021237in}}%
\pgfpathlineto{\pgfqpoint{1.635504in}{1.045985in}}%
\pgfpathlineto{\pgfqpoint{1.634978in}{1.004739in}}%
\pgfpathlineto{\pgfqpoint{1.635667in}{1.037736in}}%
\pgfpathlineto{\pgfqpoint{1.636659in}{1.004739in}}%
\pgfpathlineto{\pgfqpoint{1.636844in}{1.021237in}}%
\pgfpathlineto{\pgfqpoint{1.637325in}{1.029487in}}%
\pgfpathlineto{\pgfqpoint{1.637659in}{1.004739in}}%
\pgfpathlineto{\pgfqpoint{1.638332in}{0.979992in}}%
\pgfpathlineto{\pgfqpoint{1.638666in}{1.037736in}}%
\pgfpathlineto{\pgfqpoint{1.638688in}{1.045985in}}%
\pgfpathlineto{\pgfqpoint{1.639021in}{1.021237in}}%
\pgfpathlineto{\pgfqpoint{1.639695in}{1.021237in}}%
\pgfpathlineto{\pgfqpoint{1.641353in}{1.054234in}}%
\pgfpathlineto{\pgfqpoint{1.642375in}{1.045985in}}%
\pgfpathlineto{\pgfqpoint{1.642694in}{1.029487in}}%
\pgfpathlineto{\pgfqpoint{1.643027in}{1.062483in}}%
\pgfpathlineto{\pgfqpoint{1.643249in}{1.054234in}}%
\pgfpathlineto{\pgfqpoint{1.644049in}{1.070732in}}%
\pgfpathlineto{\pgfqpoint{1.645041in}{1.029487in}}%
\pgfpathlineto{\pgfqpoint{1.645374in}{1.054234in}}%
\pgfpathlineto{\pgfqpoint{1.646040in}{1.078981in}}%
\pgfpathlineto{\pgfqpoint{1.645729in}{1.045985in}}%
\pgfpathlineto{\pgfqpoint{1.646396in}{1.070732in}}%
\pgfpathlineto{\pgfqpoint{1.646899in}{1.045985in}}%
\pgfpathlineto{\pgfqpoint{1.647217in}{1.078981in}}%
\pgfpathlineto{\pgfqpoint{1.647232in}{1.070732in}}%
\pgfpathlineto{\pgfqpoint{1.648054in}{1.078981in}}%
\pgfpathlineto{\pgfqpoint{1.648076in}{1.070732in}}%
\pgfpathlineto{\pgfqpoint{1.648580in}{1.045985in}}%
\pgfpathlineto{\pgfqpoint{1.648891in}{1.078981in}}%
\pgfpathlineto{\pgfqpoint{1.648913in}{1.070732in}}%
\pgfpathlineto{\pgfqpoint{1.649417in}{1.095480in}}%
\pgfpathlineto{\pgfqpoint{1.649920in}{1.070732in}}%
\pgfpathlineto{\pgfqpoint{1.650231in}{1.054234in}}%
\pgfpathlineto{\pgfqpoint{1.650401in}{1.078981in}}%
\pgfpathlineto{\pgfqpoint{1.650734in}{1.078981in}}%
\pgfpathlineto{\pgfqpoint{1.650757in}{1.095480in}}%
\pgfpathlineto{\pgfqpoint{1.651593in}{1.045985in}}%
\pgfpathlineto{\pgfqpoint{1.651764in}{1.045985in}}%
\pgfpathlineto{\pgfqpoint{1.655096in}{1.177971in}}%
\pgfpathlineto{\pgfqpoint{1.656103in}{1.153224in}}%
\pgfpathlineto{\pgfqpoint{1.656117in}{1.169722in}}%
\pgfpathlineto{\pgfqpoint{1.656769in}{1.202719in}}%
\pgfpathlineto{\pgfqpoint{1.657272in}{1.177971in}}%
\pgfpathlineto{\pgfqpoint{1.658279in}{1.153224in}}%
\pgfpathlineto{\pgfqpoint{1.657613in}{1.202719in}}%
\pgfpathlineto{\pgfqpoint{1.658302in}{1.169722in}}%
\pgfpathlineto{\pgfqpoint{1.658613in}{1.153224in}}%
\pgfpathlineto{\pgfqpoint{1.658783in}{1.177971in}}%
\pgfpathlineto{\pgfqpoint{1.659138in}{1.169722in}}%
\pgfpathlineto{\pgfqpoint{1.660797in}{1.227466in}}%
\pgfpathlineto{\pgfqpoint{1.660982in}{1.210968in}}%
\pgfpathlineto{\pgfqpoint{1.661300in}{1.227466in}}%
\pgfpathlineto{\pgfqpoint{1.662996in}{1.037736in}}%
\pgfpathlineto{\pgfqpoint{1.666513in}{0.625278in}}%
\pgfpathlineto{\pgfqpoint{1.667520in}{0.650025in}}%
\pgfpathlineto{\pgfqpoint{1.667838in}{0.633527in}}%
\pgfpathlineto{\pgfqpoint{1.668838in}{0.625278in}}%
\pgfpathlineto{\pgfqpoint{1.670015in}{0.633527in}}%
\pgfpathlineto{\pgfqpoint{1.671022in}{0.625278in}}%
\pgfpathlineto{\pgfqpoint{1.672192in}{0.633527in}}%
\pgfpathlineto{\pgfqpoint{1.673051in}{0.625278in}}%
\pgfpathlineto{\pgfqpoint{1.673221in}{0.650025in}}%
\pgfpathlineto{\pgfqpoint{1.674539in}{0.683022in}}%
\pgfpathlineto{\pgfqpoint{1.674561in}{0.674773in}}%
\pgfpathlineto{\pgfqpoint{1.675228in}{0.650025in}}%
\pgfpathlineto{\pgfqpoint{1.675398in}{0.674773in}}%
\pgfpathlineto{\pgfqpoint{1.676553in}{0.683022in}}%
\pgfpathlineto{\pgfqpoint{1.677071in}{0.699520in}}%
\pgfpathlineto{\pgfqpoint{1.677575in}{0.674773in}}%
\pgfpathlineto{\pgfqpoint{1.678737in}{0.683022in}}%
\pgfpathlineto{\pgfqpoint{1.678915in}{0.674773in}}%
\pgfpathlineto{\pgfqpoint{1.679418in}{0.724268in}}%
\pgfpathlineto{\pgfqpoint{1.679796in}{0.707769in}}%
\pgfpathlineto{\pgfqpoint{1.679922in}{0.699520in}}%
\pgfpathlineto{\pgfqpoint{1.680240in}{0.732517in}}%
\pgfpathlineto{\pgfqpoint{1.680411in}{0.732517in}}%
\pgfpathlineto{\pgfqpoint{1.681432in}{0.773763in}}%
\pgfpathlineto{\pgfqpoint{1.681751in}{0.757264in}}%
\pgfpathlineto{\pgfqpoint{1.681914in}{0.782012in}}%
\pgfpathlineto{\pgfqpoint{1.682306in}{0.782012in}}%
\pgfpathlineto{\pgfqpoint{1.682935in}{0.831507in}}%
\pgfpathlineto{\pgfqpoint{1.683446in}{0.790261in}}%
\pgfpathlineto{\pgfqpoint{1.683772in}{0.806759in}}%
\pgfpathlineto{\pgfqpoint{1.684616in}{0.773763in}}%
\pgfpathlineto{\pgfqpoint{1.685453in}{0.823258in}}%
\pgfpathlineto{\pgfqpoint{1.686282in}{0.806759in}}%
\pgfpathlineto{\pgfqpoint{1.686460in}{0.823258in}}%
\pgfpathlineto{\pgfqpoint{1.687296in}{0.798510in}}%
\pgfpathlineto{\pgfqpoint{1.688459in}{0.806759in}}%
\pgfpathlineto{\pgfqpoint{1.688474in}{0.798510in}}%
\pgfpathlineto{\pgfqpoint{1.689133in}{0.856254in}}%
\pgfpathlineto{\pgfqpoint{1.689481in}{0.848005in}}%
\pgfpathlineto{\pgfqpoint{1.692316in}{1.062483in}}%
\pgfpathlineto{\pgfqpoint{1.692487in}{1.078981in}}%
\pgfpathlineto{\pgfqpoint{1.693331in}{1.037736in}}%
\pgfpathlineto{\pgfqpoint{1.693671in}{1.021237in}}%
\pgfpathlineto{\pgfqpoint{1.693990in}{1.054234in}}%
\pgfpathlineto{\pgfqpoint{1.694005in}{1.045985in}}%
\pgfpathlineto{\pgfqpoint{1.694175in}{1.070732in}}%
\pgfpathlineto{\pgfqpoint{1.695167in}{1.054234in}}%
\pgfpathlineto{\pgfqpoint{1.695848in}{0.971742in}}%
\pgfpathlineto{\pgfqpoint{1.696181in}{1.021237in}}%
\pgfpathlineto{\pgfqpoint{1.697351in}{1.029487in}}%
\pgfpathlineto{\pgfqpoint{1.698173in}{1.004739in}}%
\pgfpathlineto{\pgfqpoint{1.698366in}{1.012988in}}%
\pgfpathlineto{\pgfqpoint{1.700039in}{0.971742in}}%
\pgfpathlineto{\pgfqpoint{1.700876in}{0.996490in}}%
\pgfpathlineto{\pgfqpoint{1.701194in}{0.979992in}}%
\pgfpathlineto{\pgfqpoint{1.701868in}{0.955244in}}%
\pgfpathlineto{\pgfqpoint{1.702216in}{0.971742in}}%
\pgfpathlineto{\pgfqpoint{1.702371in}{0.979992in}}%
\pgfpathlineto{\pgfqpoint{1.702705in}{0.955244in}}%
\pgfpathlineto{\pgfqpoint{1.702890in}{0.963493in}}%
\pgfpathlineto{\pgfqpoint{1.703882in}{0.955244in}}%
\pgfpathlineto{\pgfqpoint{1.703215in}{0.979992in}}%
\pgfpathlineto{\pgfqpoint{1.703882in}{0.963493in}}%
\pgfpathlineto{\pgfqpoint{1.704889in}{0.979992in}}%
\pgfpathlineto{\pgfqpoint{1.704718in}{0.955244in}}%
\pgfpathlineto{\pgfqpoint{1.704904in}{0.971742in}}%
\pgfpathlineto{\pgfqpoint{1.705392in}{0.979992in}}%
\pgfpathlineto{\pgfqpoint{1.705570in}{0.963493in}}%
\pgfpathlineto{\pgfqpoint{1.708250in}{0.765513in}}%
\pgfpathlineto{\pgfqpoint{1.710264in}{0.625278in}}%
\pgfpathlineto{\pgfqpoint{1.711419in}{0.633527in}}%
\pgfpathlineto{\pgfqpoint{1.711441in}{0.625278in}}%
\pgfpathlineto{\pgfqpoint{1.711923in}{0.658274in}}%
\pgfpathlineto{\pgfqpoint{1.712441in}{0.625278in}}%
\pgfpathlineto{\pgfqpoint{1.713596in}{0.633527in}}%
\pgfpathlineto{\pgfqpoint{1.713618in}{0.625278in}}%
\pgfpathlineto{\pgfqpoint{1.714270in}{0.683022in}}%
\pgfpathlineto{\pgfqpoint{1.714625in}{0.650025in}}%
\pgfpathlineto{\pgfqpoint{1.714936in}{0.658274in}}%
\pgfpathlineto{\pgfqpoint{1.715277in}{0.633527in}}%
\pgfpathlineto{\pgfqpoint{1.715965in}{0.650025in}}%
\pgfpathlineto{\pgfqpoint{1.716299in}{0.625278in}}%
\pgfpathlineto{\pgfqpoint{1.717461in}{0.633527in}}%
\pgfpathlineto{\pgfqpoint{1.717631in}{0.625278in}}%
\pgfpathlineto{\pgfqpoint{1.718475in}{0.674773in}}%
\pgfpathlineto{\pgfqpoint{1.719631in}{0.683022in}}%
\pgfpathlineto{\pgfqpoint{1.720660in}{0.674773in}}%
\pgfpathlineto{\pgfqpoint{1.722651in}{0.732517in}}%
\pgfpathlineto{\pgfqpoint{1.723318in}{0.707769in}}%
\pgfpathlineto{\pgfqpoint{1.723673in}{0.724268in}}%
\pgfpathlineto{\pgfqpoint{1.724325in}{0.732517in}}%
\pgfpathlineto{\pgfqpoint{1.724517in}{0.716019in}}%
\pgfpathlineto{\pgfqpoint{1.725184in}{0.699520in}}%
\pgfpathlineto{\pgfqpoint{1.725354in}{0.724268in}}%
\pgfpathlineto{\pgfqpoint{1.727679in}{0.831507in}}%
\pgfpathlineto{\pgfqpoint{1.727864in}{0.815008in}}%
\pgfpathlineto{\pgfqpoint{1.728205in}{0.798510in}}%
\pgfpathlineto{\pgfqpoint{1.728523in}{0.839756in}}%
\pgfpathlineto{\pgfqpoint{1.729523in}{0.856254in}}%
\pgfpathlineto{\pgfqpoint{1.729360in}{0.831507in}}%
\pgfpathlineto{\pgfqpoint{1.729545in}{0.848005in}}%
\pgfpathlineto{\pgfqpoint{1.730707in}{0.856254in}}%
\pgfpathlineto{\pgfqpoint{1.731048in}{0.848005in}}%
\pgfpathlineto{\pgfqpoint{1.731722in}{0.897500in}}%
\pgfpathlineto{\pgfqpoint{1.733380in}{0.856254in}}%
\pgfpathlineto{\pgfqpoint{1.733380in}{0.864503in}}%
\pgfpathlineto{\pgfqpoint{1.734046in}{0.930497in}}%
\pgfpathlineto{\pgfqpoint{1.734550in}{0.905749in}}%
\pgfpathlineto{\pgfqpoint{1.734572in}{0.897500in}}%
\pgfpathlineto{\pgfqpoint{1.735239in}{0.971742in}}%
\pgfpathlineto{\pgfqpoint{1.735409in}{0.963493in}}%
\pgfpathlineto{\pgfqpoint{1.736231in}{0.955244in}}%
\pgfpathlineto{\pgfqpoint{1.735579in}{0.971742in}}%
\pgfpathlineto{\pgfqpoint{1.736246in}{0.971742in}}%
\pgfpathlineto{\pgfqpoint{1.737238in}{1.004739in}}%
\pgfpathlineto{\pgfqpoint{1.737423in}{0.988241in}}%
\pgfpathlineto{\pgfqpoint{1.738911in}{0.955244in}}%
\pgfpathlineto{\pgfqpoint{1.738259in}{0.996490in}}%
\pgfpathlineto{\pgfqpoint{1.738911in}{0.963493in}}%
\pgfpathlineto{\pgfqpoint{1.739244in}{0.955244in}}%
\pgfpathlineto{\pgfqpoint{1.739933in}{0.971742in}}%
\pgfpathlineto{\pgfqpoint{1.740584in}{0.955244in}}%
\pgfpathlineto{\pgfqpoint{1.740584in}{0.988241in}}%
\pgfpathlineto{\pgfqpoint{1.740940in}{1.021237in}}%
\pgfpathlineto{\pgfqpoint{1.740925in}{0.979992in}}%
\pgfpathlineto{\pgfqpoint{1.741614in}{0.996490in}}%
\pgfpathlineto{\pgfqpoint{1.743620in}{1.070732in}}%
\pgfpathlineto{\pgfqpoint{1.743790in}{1.062483in}}%
\pgfpathlineto{\pgfqpoint{1.744272in}{1.054234in}}%
\pgfpathlineto{\pgfqpoint{1.744294in}{1.095480in}}%
\pgfpathlineto{\pgfqpoint{1.744945in}{1.103729in}}%
\pgfpathlineto{\pgfqpoint{1.745131in}{1.087231in}}%
\pgfpathlineto{\pgfqpoint{1.745619in}{1.078981in}}%
\pgfpathlineto{\pgfqpoint{1.745782in}{1.103729in}}%
\pgfpathlineto{\pgfqpoint{1.745804in}{1.095480in}}%
\pgfpathlineto{\pgfqpoint{1.747463in}{1.177971in}}%
\pgfpathlineto{\pgfqpoint{1.747648in}{1.161473in}}%
\pgfpathlineto{\pgfqpoint{1.748803in}{1.103729in}}%
\pgfpathlineto{\pgfqpoint{1.748988in}{1.120227in}}%
\pgfpathlineto{\pgfqpoint{1.749307in}{1.177971in}}%
\pgfpathlineto{\pgfqpoint{1.749995in}{1.120227in}}%
\pgfpathlineto{\pgfqpoint{1.752505in}{1.021237in}}%
\pgfpathlineto{\pgfqpoint{1.752653in}{1.037736in}}%
\pgfpathlineto{\pgfqpoint{1.753660in}{1.054234in}}%
\pgfpathlineto{\pgfqpoint{1.753497in}{1.029487in}}%
\pgfpathlineto{\pgfqpoint{1.753682in}{1.045985in}}%
\pgfpathlineto{\pgfqpoint{1.754001in}{1.054234in}}%
\pgfpathlineto{\pgfqpoint{1.754186in}{1.037736in}}%
\pgfpathlineto{\pgfqpoint{1.754852in}{0.996490in}}%
\pgfpathlineto{\pgfqpoint{1.755356in}{1.021237in}}%
\pgfpathlineto{\pgfqpoint{1.756178in}{1.078981in}}%
\pgfpathlineto{\pgfqpoint{1.756533in}{1.037736in}}%
\pgfpathlineto{\pgfqpoint{1.759695in}{0.683022in}}%
\pgfpathlineto{\pgfqpoint{1.759717in}{0.699520in}}%
\pgfpathlineto{\pgfqpoint{1.760887in}{0.724268in}}%
\pgfpathlineto{\pgfqpoint{1.761057in}{0.716019in}}%
\pgfpathlineto{\pgfqpoint{1.761538in}{0.707769in}}%
\pgfpathlineto{\pgfqpoint{1.761709in}{0.732517in}}%
\pgfpathlineto{\pgfqpoint{1.761731in}{0.724268in}}%
\pgfpathlineto{\pgfqpoint{1.763574in}{0.848005in}}%
\pgfpathlineto{\pgfqpoint{1.764411in}{0.823258in}}%
\pgfpathlineto{\pgfqpoint{1.766255in}{0.773763in}}%
\pgfpathlineto{\pgfqpoint{1.767239in}{0.806759in}}%
\pgfpathlineto{\pgfqpoint{1.767425in}{0.790261in}}%
\pgfpathlineto{\pgfqpoint{1.767595in}{0.798510in}}%
\pgfpathlineto{\pgfqpoint{1.768602in}{0.773763in}}%
\pgfpathlineto{\pgfqpoint{1.769609in}{0.823258in}}%
\pgfpathlineto{\pgfqpoint{1.769942in}{0.790261in}}%
\pgfpathlineto{\pgfqpoint{1.770275in}{0.773763in}}%
\pgfpathlineto{\pgfqpoint{1.770594in}{0.815008in}}%
\pgfpathlineto{\pgfqpoint{1.771786in}{0.872753in}}%
\pgfpathlineto{\pgfqpoint{1.772104in}{0.881002in}}%
\pgfpathlineto{\pgfqpoint{1.772437in}{0.856254in}}%
\pgfpathlineto{\pgfqpoint{1.772793in}{0.848005in}}%
\pgfpathlineto{\pgfqpoint{1.772956in}{0.872753in}}%
\pgfpathlineto{\pgfqpoint{1.773459in}{0.864503in}}%
\pgfpathlineto{\pgfqpoint{1.773963in}{0.848005in}}%
\pgfpathlineto{\pgfqpoint{1.774133in}{0.872753in}}%
\pgfpathlineto{\pgfqpoint{1.775458in}{0.905749in}}%
\pgfpathlineto{\pgfqpoint{1.775473in}{0.897500in}}%
\pgfpathlineto{\pgfqpoint{1.776147in}{0.872753in}}%
\pgfpathlineto{\pgfqpoint{1.776310in}{0.897500in}}%
\pgfpathlineto{\pgfqpoint{1.777798in}{0.955244in}}%
\pgfpathlineto{\pgfqpoint{1.777990in}{0.938746in}}%
\pgfpathlineto{\pgfqpoint{1.778975in}{0.881002in}}%
\pgfpathlineto{\pgfqpoint{1.779316in}{0.905749in}}%
\pgfpathlineto{\pgfqpoint{1.780152in}{0.955244in}}%
\pgfpathlineto{\pgfqpoint{1.780338in}{0.946995in}}%
\pgfpathlineto{\pgfqpoint{1.781996in}{0.979992in}}%
\pgfpathlineto{\pgfqpoint{1.783188in}{0.922247in}}%
\pgfpathlineto{\pgfqpoint{1.785528in}{1.095480in}}%
\pgfpathlineto{\pgfqpoint{1.786372in}{1.062483in}}%
\pgfpathlineto{\pgfqpoint{1.789385in}{0.674773in}}%
\pgfpathlineto{\pgfqpoint{1.789719in}{0.699520in}}%
\pgfpathlineto{\pgfqpoint{1.790222in}{0.724268in}}%
\pgfpathlineto{\pgfqpoint{1.790874in}{0.707769in}}%
\pgfpathlineto{\pgfqpoint{1.790896in}{0.699520in}}%
\pgfpathlineto{\pgfqpoint{1.791392in}{0.757264in}}%
\pgfpathlineto{\pgfqpoint{1.791733in}{0.749015in}}%
\pgfpathlineto{\pgfqpoint{1.792384in}{0.732517in}}%
\pgfpathlineto{\pgfqpoint{1.792554in}{0.757264in}}%
\pgfpathlineto{\pgfqpoint{1.792569in}{0.749015in}}%
\pgfpathlineto{\pgfqpoint{1.793732in}{0.757264in}}%
\pgfpathlineto{\pgfqpoint{1.794916in}{0.691271in}}%
\pgfpathlineto{\pgfqpoint{1.796597in}{0.650025in}}%
\pgfpathlineto{\pgfqpoint{1.797604in}{0.674773in}}%
\pgfpathlineto{\pgfqpoint{1.797767in}{0.666524in}}%
\pgfpathlineto{\pgfqpoint{1.799107in}{0.625278in}}%
\pgfpathlineto{\pgfqpoint{1.797937in}{0.674773in}}%
\pgfpathlineto{\pgfqpoint{1.799263in}{0.641776in}}%
\pgfpathlineto{\pgfqpoint{1.799951in}{0.674773in}}%
\pgfpathlineto{\pgfqpoint{1.799596in}{0.633527in}}%
\pgfpathlineto{\pgfqpoint{1.800284in}{0.650025in}}%
\pgfpathlineto{\pgfqpoint{1.800788in}{0.674773in}}%
\pgfpathlineto{\pgfqpoint{1.801121in}{0.650025in}}%
\pgfpathlineto{\pgfqpoint{1.801328in}{0.633527in}}%
\pgfpathlineto{\pgfqpoint{1.801610in}{0.658274in}}%
\pgfpathlineto{\pgfqpoint{1.801958in}{0.650025in}}%
\pgfpathlineto{\pgfqpoint{1.802446in}{0.658274in}}%
\pgfpathlineto{\pgfqpoint{1.802632in}{0.641776in}}%
\pgfpathlineto{\pgfqpoint{1.803453in}{0.625278in}}%
\pgfpathlineto{\pgfqpoint{1.804794in}{0.633527in}}%
\pgfpathlineto{\pgfqpoint{1.804956in}{0.625278in}}%
\pgfpathlineto{\pgfqpoint{1.805630in}{0.625278in}}%
\pgfpathlineto{\pgfqpoint{1.806970in}{0.633527in}}%
\pgfpathlineto{\pgfqpoint{1.807133in}{0.625278in}}%
\pgfpathlineto{\pgfqpoint{1.807977in}{0.625278in}}%
\pgfpathlineto{\pgfqpoint{1.808821in}{0.633527in}}%
\pgfpathlineto{\pgfqpoint{1.808821in}{0.625278in}}%
\pgfpathlineto{\pgfqpoint{1.809999in}{0.633527in}}%
\pgfpathlineto{\pgfqpoint{1.810154in}{0.625278in}}%
\pgfpathlineto{\pgfqpoint{1.810835in}{0.625278in}}%
\pgfpathlineto{\pgfqpoint{1.812013in}{0.633527in}}%
\pgfpathlineto{\pgfqpoint{1.812176in}{0.625278in}}%
\pgfpathlineto{\pgfqpoint{1.813020in}{0.625278in}}%
\pgfpathlineto{\pgfqpoint{1.814189in}{0.633527in}}%
\pgfpathlineto{\pgfqpoint{1.814352in}{0.625278in}}%
\pgfpathlineto{\pgfqpoint{1.815189in}{0.625278in}}%
\pgfpathlineto{\pgfqpoint{1.816366in}{0.633527in}}%
\pgfpathlineto{\pgfqpoint{1.816529in}{0.625278in}}%
\pgfpathlineto{\pgfqpoint{1.817366in}{0.625278in}}%
\pgfpathlineto{\pgfqpoint{1.818536in}{0.633527in}}%
\pgfpathlineto{\pgfqpoint{1.818706in}{0.625278in}}%
\pgfpathlineto{\pgfqpoint{1.819543in}{0.625278in}}%
\pgfpathlineto{\pgfqpoint{1.820720in}{0.633527in}}%
\pgfpathlineto{\pgfqpoint{1.820883in}{0.625278in}}%
\pgfpathlineto{\pgfqpoint{1.821727in}{0.625278in}}%
\pgfpathlineto{\pgfqpoint{1.822904in}{0.633527in}}%
\pgfpathlineto{\pgfqpoint{1.823060in}{0.625278in}}%
\pgfpathlineto{\pgfqpoint{1.823919in}{0.625278in}}%
\pgfpathlineto{\pgfqpoint{1.826088in}{0.658274in}}%
\pgfpathlineto{\pgfqpoint{1.826095in}{0.650025in}}%
\pgfpathlineto{\pgfqpoint{1.826421in}{0.683022in}}%
\pgfpathlineto{\pgfqpoint{1.827102in}{0.666524in}}%
\pgfpathlineto{\pgfqpoint{1.827591in}{0.633527in}}%
\pgfpathlineto{\pgfqpoint{1.827939in}{0.674773in}}%
\pgfpathlineto{\pgfqpoint{1.828265in}{0.683022in}}%
\pgfpathlineto{\pgfqpoint{1.828591in}{0.658274in}}%
\pgfpathlineto{\pgfqpoint{1.829768in}{0.625278in}}%
\pgfpathlineto{\pgfqpoint{1.830612in}{0.633527in}}%
\pgfpathlineto{\pgfqpoint{1.830612in}{0.625278in}}%
\pgfpathlineto{\pgfqpoint{1.831945in}{0.633527in}}%
\pgfpathlineto{\pgfqpoint{1.831952in}{0.625278in}}%
\pgfpathlineto{\pgfqpoint{1.832789in}{0.625278in}}%
\pgfpathlineto{\pgfqpoint{1.833959in}{0.633527in}}%
\pgfpathlineto{\pgfqpoint{1.834292in}{0.625278in}}%
\pgfpathlineto{\pgfqpoint{1.834625in}{0.625278in}}%
\pgfpathlineto{\pgfqpoint{1.835810in}{0.633527in}}%
\pgfpathlineto{\pgfqpoint{1.836024in}{0.625278in}}%
\pgfpathlineto{\pgfqpoint{1.836809in}{0.625278in}}%
\pgfpathlineto{\pgfqpoint{1.838149in}{0.633527in}}%
\pgfpathlineto{\pgfqpoint{1.838320in}{0.625278in}}%
\pgfpathlineto{\pgfqpoint{1.839156in}{0.625278in}}%
\pgfpathlineto{\pgfqpoint{1.840326in}{0.633527in}}%
\pgfpathlineto{\pgfqpoint{1.840497in}{0.625278in}}%
\pgfpathlineto{\pgfqpoint{1.841178in}{0.625278in}}%
\pgfpathlineto{\pgfqpoint{1.842340in}{0.633527in}}%
\pgfpathlineto{\pgfqpoint{1.842511in}{0.625278in}}%
\pgfpathlineto{\pgfqpoint{1.843347in}{0.625278in}}%
\pgfpathlineto{\pgfqpoint{1.844524in}{0.633527in}}%
\pgfpathlineto{\pgfqpoint{1.844858in}{0.625278in}}%
\pgfpathlineto{\pgfqpoint{1.845531in}{0.625278in}}%
\pgfpathlineto{\pgfqpoint{1.846701in}{0.633527in}}%
\pgfpathlineto{\pgfqpoint{1.846864in}{0.625278in}}%
\pgfpathlineto{\pgfqpoint{1.847701in}{0.625278in}}%
\pgfpathlineto{\pgfqpoint{1.848878in}{0.633527in}}%
\pgfpathlineto{\pgfqpoint{1.848886in}{0.625278in}}%
\pgfpathlineto{\pgfqpoint{1.849893in}{0.625278in}}%
\pgfpathlineto{\pgfqpoint{1.851055in}{0.633527in}}%
\pgfpathlineto{\pgfqpoint{1.851233in}{0.625278in}}%
\pgfpathlineto{\pgfqpoint{1.851892in}{0.625278in}}%
\pgfpathlineto{\pgfqpoint{1.853232in}{0.633527in}}%
\pgfpathlineto{\pgfqpoint{1.853232in}{0.625278in}}%
\pgfpathlineto{\pgfqpoint{1.854239in}{0.625278in}}%
\pgfpathlineto{\pgfqpoint{1.855416in}{0.633527in}}%
\pgfpathlineto{\pgfqpoint{1.855416in}{0.625278in}}%
\pgfpathlineto{\pgfqpoint{1.856423in}{0.625278in}}%
\pgfpathlineto{\pgfqpoint{1.857593in}{0.633527in}}%
\pgfpathlineto{\pgfqpoint{1.857593in}{0.625278in}}%
\pgfpathlineto{\pgfqpoint{1.858600in}{0.625278in}}%
\pgfpathlineto{\pgfqpoint{1.859770in}{0.633527in}}%
\pgfpathlineto{\pgfqpoint{1.859770in}{0.625278in}}%
\pgfpathlineto{\pgfqpoint{1.860777in}{0.625278in}}%
\pgfpathlineto{\pgfqpoint{1.861954in}{0.633527in}}%
\pgfpathlineto{\pgfqpoint{1.862117in}{0.625278in}}%
\pgfpathlineto{\pgfqpoint{1.862961in}{0.625278in}}%
\pgfpathlineto{\pgfqpoint{1.864301in}{0.633527in}}%
\pgfpathlineto{\pgfqpoint{1.864634in}{0.625278in}}%
\pgfpathlineto{\pgfqpoint{1.865308in}{0.625278in}}%
\pgfpathlineto{\pgfqpoint{1.866478in}{0.633527in}}%
\pgfpathlineto{\pgfqpoint{1.866656in}{0.625278in}}%
\pgfpathlineto{\pgfqpoint{1.867485in}{0.625278in}}%
\pgfpathlineto{\pgfqpoint{1.868655in}{0.633527in}}%
\pgfpathlineto{\pgfqpoint{1.868825in}{0.625278in}}%
\pgfpathlineto{\pgfqpoint{1.869662in}{0.625278in}}%
\pgfpathlineto{\pgfqpoint{1.870839in}{0.633527in}}%
\pgfpathlineto{\pgfqpoint{1.870839in}{0.625278in}}%
\pgfpathlineto{\pgfqpoint{1.871846in}{0.625278in}}%
\pgfpathlineto{\pgfqpoint{1.873016in}{0.633527in}}%
\pgfpathlineto{\pgfqpoint{1.873349in}{0.625278in}}%
\pgfpathlineto{\pgfqpoint{1.874023in}{0.625278in}}%
\pgfpathlineto{\pgfqpoint{1.875193in}{0.633527in}}%
\pgfpathlineto{\pgfqpoint{1.875363in}{0.625278in}}%
\pgfpathlineto{\pgfqpoint{1.876200in}{0.625278in}}%
\pgfpathlineto{\pgfqpoint{1.877377in}{0.633527in}}%
\pgfpathlineto{\pgfqpoint{1.877377in}{0.625278in}}%
\pgfpathlineto{\pgfqpoint{1.878376in}{0.625278in}}%
\pgfpathlineto{\pgfqpoint{1.879554in}{0.633527in}}%
\pgfpathlineto{\pgfqpoint{1.879724in}{0.625278in}}%
\pgfpathlineto{\pgfqpoint{1.880561in}{0.625278in}}%
\pgfpathlineto{\pgfqpoint{1.881731in}{0.633527in}}%
\pgfpathlineto{\pgfqpoint{1.881901in}{0.625278in}}%
\pgfpathlineto{\pgfqpoint{1.882738in}{0.625278in}}%
\pgfpathlineto{\pgfqpoint{1.883915in}{0.633527in}}%
\pgfpathlineto{\pgfqpoint{1.884078in}{0.625278in}}%
\pgfpathlineto{\pgfqpoint{1.884751in}{0.625278in}}%
\pgfpathlineto{\pgfqpoint{1.885921in}{0.633527in}}%
\pgfpathlineto{\pgfqpoint{1.885921in}{0.625278in}}%
\pgfpathlineto{\pgfqpoint{1.886758in}{0.625278in}}%
\pgfpathlineto{\pgfqpoint{1.887935in}{0.633527in}}%
\pgfpathlineto{\pgfqpoint{1.888106in}{0.625278in}}%
\pgfpathlineto{\pgfqpoint{1.888942in}{0.625278in}}%
\pgfpathlineto{\pgfqpoint{1.890112in}{0.633527in}}%
\pgfpathlineto{\pgfqpoint{1.890282in}{0.625278in}}%
\pgfpathlineto{\pgfqpoint{1.891119in}{0.625278in}}%
\pgfpathlineto{\pgfqpoint{1.892296in}{0.633527in}}%
\pgfpathlineto{\pgfqpoint{1.892459in}{0.625278in}}%
\pgfpathlineto{\pgfqpoint{1.893296in}{0.625278in}}%
\pgfpathlineto{\pgfqpoint{1.894303in}{0.633527in}}%
\pgfpathlineto{\pgfqpoint{1.894303in}{0.625278in}}%
\pgfpathlineto{\pgfqpoint{1.895480in}{0.633527in}}%
\pgfpathlineto{\pgfqpoint{1.895480in}{0.625278in}}%
\pgfpathlineto{\pgfqpoint{1.896487in}{0.625278in}}%
\pgfpathlineto{\pgfqpoint{1.897487in}{0.633527in}}%
\pgfpathlineto{\pgfqpoint{1.897487in}{0.625278in}}%
\pgfpathlineto{\pgfqpoint{1.898664in}{0.633527in}}%
\pgfpathlineto{\pgfqpoint{1.898834in}{0.625278in}}%
\pgfpathlineto{\pgfqpoint{1.899671in}{0.625278in}}%
\pgfpathlineto{\pgfqpoint{1.900841in}{0.633527in}}%
\pgfpathlineto{\pgfqpoint{1.901011in}{0.625278in}}%
\pgfpathlineto{\pgfqpoint{1.901848in}{0.625278in}}%
\pgfpathlineto{\pgfqpoint{1.903025in}{0.633527in}}%
\pgfpathlineto{\pgfqpoint{1.903188in}{0.625278in}}%
\pgfpathlineto{\pgfqpoint{1.904025in}{0.625278in}}%
\pgfpathlineto{\pgfqpoint{1.905372in}{0.633527in}}%
\pgfpathlineto{\pgfqpoint{1.905535in}{0.625278in}}%
\pgfpathlineto{\pgfqpoint{1.906372in}{0.625278in}}%
\pgfpathlineto{\pgfqpoint{1.907712in}{0.633527in}}%
\pgfpathlineto{\pgfqpoint{1.908052in}{0.625278in}}%
\pgfpathlineto{\pgfqpoint{1.908719in}{0.625278in}}%
\pgfpathlineto{\pgfqpoint{1.909896in}{0.633527in}}%
\pgfpathlineto{\pgfqpoint{1.910229in}{0.625278in}}%
\pgfpathlineto{\pgfqpoint{1.910903in}{0.625278in}}%
\pgfpathlineto{\pgfqpoint{1.912073in}{0.633527in}}%
\pgfpathlineto{\pgfqpoint{1.912243in}{0.625278in}}%
\pgfpathlineto{\pgfqpoint{1.913080in}{0.625278in}}%
\pgfpathlineto{\pgfqpoint{1.914250in}{0.633527in}}%
\pgfpathlineto{\pgfqpoint{1.914590in}{0.625278in}}%
\pgfpathlineto{\pgfqpoint{1.915257in}{0.625278in}}%
\pgfpathlineto{\pgfqpoint{1.916264in}{0.633527in}}%
\pgfpathlineto{\pgfqpoint{1.916264in}{0.625278in}}%
\pgfpathlineto{\pgfqpoint{1.917604in}{0.633527in}}%
\pgfpathlineto{\pgfqpoint{1.917774in}{0.625278in}}%
\pgfpathlineto{\pgfqpoint{1.918611in}{0.625278in}}%
\pgfpathlineto{\pgfqpoint{1.919788in}{0.633527in}}%
\pgfpathlineto{\pgfqpoint{1.919951in}{0.625278in}}%
\pgfpathlineto{\pgfqpoint{1.920788in}{0.625278in}}%
\pgfpathlineto{\pgfqpoint{1.922135in}{0.633527in}}%
\pgfpathlineto{\pgfqpoint{1.922135in}{0.625278in}}%
\pgfpathlineto{\pgfqpoint{1.923135in}{0.625278in}}%
\pgfpathlineto{\pgfqpoint{1.924312in}{0.633527in}}%
\pgfpathlineto{\pgfqpoint{1.924475in}{0.625278in}}%
\pgfpathlineto{\pgfqpoint{1.925149in}{0.625278in}}%
\pgfpathlineto{\pgfqpoint{1.926326in}{0.633527in}}%
\pgfpathlineto{\pgfqpoint{1.926326in}{0.625278in}}%
\pgfpathlineto{\pgfqpoint{1.927326in}{0.625278in}}%
\pgfpathlineto{\pgfqpoint{1.928503in}{0.633527in}}%
\pgfpathlineto{\pgfqpoint{1.928666in}{0.625278in}}%
\pgfpathlineto{\pgfqpoint{1.929340in}{0.625278in}}%
\pgfpathlineto{\pgfqpoint{1.930517in}{0.633527in}}%
\pgfpathlineto{\pgfqpoint{1.930680in}{0.625278in}}%
\pgfpathlineto{\pgfqpoint{1.931516in}{0.625278in}}%
\pgfpathlineto{\pgfqpoint{1.932694in}{0.633527in}}%
\pgfpathlineto{\pgfqpoint{1.933027in}{0.625278in}}%
\pgfpathlineto{\pgfqpoint{1.933701in}{0.625278in}}%
\pgfpathlineto{\pgfqpoint{1.934870in}{0.633527in}}%
\pgfpathlineto{\pgfqpoint{1.935041in}{0.625278in}}%
\pgfpathlineto{\pgfqpoint{1.935707in}{0.625278in}}%
\pgfpathlineto{\pgfqpoint{1.936884in}{0.633527in}}%
\pgfpathlineto{\pgfqpoint{1.937047in}{0.625278in}}%
\pgfpathlineto{\pgfqpoint{1.937891in}{0.625278in}}%
\pgfpathlineto{\pgfqpoint{1.939061in}{0.633527in}}%
\pgfpathlineto{\pgfqpoint{1.939232in}{0.625278in}}%
\pgfpathlineto{\pgfqpoint{1.940068in}{0.625278in}}%
\pgfpathlineto{\pgfqpoint{1.941075in}{0.633527in}}%
\pgfpathlineto{\pgfqpoint{1.941075in}{0.625278in}}%
\pgfpathlineto{\pgfqpoint{1.942245in}{0.633527in}}%
\pgfpathlineto{\pgfqpoint{1.942245in}{0.625278in}}%
\pgfpathlineto{\pgfqpoint{1.943089in}{0.625278in}}%
\pgfpathlineto{\pgfqpoint{1.944429in}{0.633527in}}%
\pgfpathlineto{\pgfqpoint{1.944429in}{0.625278in}}%
\pgfpathlineto{\pgfqpoint{1.945429in}{0.625278in}}%
\pgfpathlineto{\pgfqpoint{1.946606in}{0.633527in}}%
\pgfpathlineto{\pgfqpoint{1.946776in}{0.625278in}}%
\pgfpathlineto{\pgfqpoint{1.947613in}{0.625278in}}%
\pgfpathlineto{\pgfqpoint{1.948783in}{0.633527in}}%
\pgfpathlineto{\pgfqpoint{1.948953in}{0.625278in}}%
\pgfpathlineto{\pgfqpoint{1.949790in}{0.625278in}}%
\pgfpathlineto{\pgfqpoint{1.950967in}{0.633527in}}%
\pgfpathlineto{\pgfqpoint{1.950967in}{0.625278in}}%
\pgfpathlineto{\pgfqpoint{1.951967in}{0.625278in}}%
\pgfpathlineto{\pgfqpoint{1.953314in}{0.633527in}}%
\pgfpathlineto{\pgfqpoint{1.953314in}{0.625278in}}%
\pgfpathlineto{\pgfqpoint{1.954314in}{0.625278in}}%
\pgfpathlineto{\pgfqpoint{1.955491in}{0.633527in}}%
\pgfpathlineto{\pgfqpoint{1.955661in}{0.625278in}}%
\pgfpathlineto{\pgfqpoint{1.956498in}{0.625278in}}%
\pgfpathlineto{\pgfqpoint{1.957505in}{0.633527in}}%
\pgfpathlineto{\pgfqpoint{1.957505in}{0.625278in}}%
\pgfpathlineto{\pgfqpoint{1.958675in}{0.633527in}}%
\pgfpathlineto{\pgfqpoint{1.958845in}{0.625278in}}%
\pgfpathlineto{\pgfqpoint{1.959682in}{0.625278in}}%
\pgfpathlineto{\pgfqpoint{1.960852in}{0.633527in}}%
\pgfpathlineto{\pgfqpoint{1.961022in}{0.625278in}}%
\pgfpathlineto{\pgfqpoint{1.961696in}{0.625278in}}%
\pgfpathlineto{\pgfqpoint{1.963036in}{0.633527in}}%
\pgfpathlineto{\pgfqpoint{1.963036in}{0.625278in}}%
\pgfpathlineto{\pgfqpoint{1.963873in}{0.625278in}}%
\pgfpathlineto{\pgfqpoint{1.965043in}{0.633527in}}%
\pgfpathlineto{\pgfqpoint{1.965213in}{0.625278in}}%
\pgfpathlineto{\pgfqpoint{1.966050in}{0.625278in}}%
\pgfpathlineto{\pgfqpoint{1.967227in}{0.633527in}}%
\pgfpathlineto{\pgfqpoint{1.967227in}{0.625278in}}%
\pgfpathlineto{\pgfqpoint{1.968234in}{0.625278in}}%
\pgfpathlineto{\pgfqpoint{1.969404in}{0.633527in}}%
\pgfpathlineto{\pgfqpoint{1.969404in}{0.625278in}}%
\pgfpathlineto{\pgfqpoint{1.970411in}{0.625278in}}%
\pgfpathlineto{\pgfqpoint{1.971580in}{0.633527in}}%
\pgfpathlineto{\pgfqpoint{1.971751in}{0.625278in}}%
\pgfpathlineto{\pgfqpoint{1.972610in}{0.625278in}}%
\pgfpathlineto{\pgfqpoint{1.973780in}{0.650025in}}%
\pgfpathlineto{\pgfqpoint{1.973950in}{0.641776in}}%
\pgfpathlineto{\pgfqpoint{1.974616in}{0.625278in}}%
\pgfpathlineto{\pgfqpoint{1.974786in}{0.650025in}}%
\pgfpathlineto{\pgfqpoint{1.975105in}{0.633527in}}%
\pgfpathlineto{\pgfqpoint{1.975120in}{0.650025in}}%
\pgfpathlineto{\pgfqpoint{1.976127in}{0.650025in}}%
\pgfpathlineto{\pgfqpoint{1.977282in}{0.658274in}}%
\pgfpathlineto{\pgfqpoint{1.977807in}{0.650025in}}%
\pgfpathlineto{\pgfqpoint{1.978118in}{0.683022in}}%
\pgfpathlineto{\pgfqpoint{1.978303in}{0.674773in}}%
\pgfpathlineto{\pgfqpoint{1.979125in}{0.683022in}}%
\pgfpathlineto{\pgfqpoint{1.979148in}{0.674773in}}%
\pgfpathlineto{\pgfqpoint{1.980466in}{0.658274in}}%
\pgfpathlineto{\pgfqpoint{1.981302in}{0.683022in}}%
\pgfpathlineto{\pgfqpoint{1.980821in}{0.650025in}}%
\pgfpathlineto{\pgfqpoint{1.981495in}{0.674773in}}%
\pgfpathlineto{\pgfqpoint{1.983486in}{0.732517in}}%
\pgfpathlineto{\pgfqpoint{1.984323in}{0.707769in}}%
\pgfpathlineto{\pgfqpoint{1.984508in}{0.724268in}}%
\pgfpathlineto{\pgfqpoint{1.985493in}{0.732517in}}%
\pgfpathlineto{\pgfqpoint{1.985515in}{0.724268in}}%
\pgfpathlineto{\pgfqpoint{1.986500in}{0.707769in}}%
\pgfpathlineto{\pgfqpoint{1.986522in}{0.724268in}}%
\pgfpathlineto{\pgfqpoint{1.990861in}{0.905749in}}%
\pgfpathlineto{\pgfqpoint{1.991046in}{0.897500in}}%
\pgfpathlineto{\pgfqpoint{1.991216in}{0.922247in}}%
\pgfpathlineto{\pgfqpoint{1.991883in}{0.922247in}}%
\pgfpathlineto{\pgfqpoint{1.993038in}{0.930497in}}%
\pgfpathlineto{\pgfqpoint{1.994067in}{0.922247in}}%
\pgfpathlineto{\pgfqpoint{1.995067in}{0.946995in}}%
\pgfpathlineto{\pgfqpoint{1.995237in}{0.938746in}}%
\pgfpathlineto{\pgfqpoint{1.995385in}{0.930497in}}%
\pgfpathlineto{\pgfqpoint{1.995570in}{0.971742in}}%
\pgfpathlineto{\pgfqpoint{1.996059in}{0.979992in}}%
\pgfpathlineto{\pgfqpoint{1.996244in}{0.963493in}}%
\pgfpathlineto{\pgfqpoint{1.997569in}{0.930497in}}%
\pgfpathlineto{\pgfqpoint{1.997584in}{0.946995in}}%
\pgfpathlineto{\pgfqpoint{1.999243in}{0.979992in}}%
\pgfpathlineto{\pgfqpoint{2.000264in}{0.963493in}}%
\pgfpathlineto{\pgfqpoint{2.001604in}{0.922247in}}%
\pgfpathlineto{\pgfqpoint{2.001760in}{0.938746in}}%
\pgfpathlineto{\pgfqpoint{2.002093in}{0.955244in}}%
\pgfpathlineto{\pgfqpoint{2.002449in}{0.922247in}}%
\pgfpathlineto{\pgfqpoint{2.002782in}{0.946995in}}%
\pgfpathlineto{\pgfqpoint{2.003115in}{0.971742in}}%
\pgfpathlineto{\pgfqpoint{2.003937in}{0.955244in}}%
\pgfpathlineto{\pgfqpoint{2.004292in}{0.946995in}}%
\pgfpathlineto{\pgfqpoint{2.004603in}{0.979992in}}%
\pgfpathlineto{\pgfqpoint{2.004959in}{0.963493in}}%
\pgfpathlineto{\pgfqpoint{2.005462in}{0.946995in}}%
\pgfpathlineto{\pgfqpoint{2.005632in}{0.971742in}}%
\pgfpathlineto{\pgfqpoint{2.005951in}{0.963493in}}%
\pgfpathlineto{\pgfqpoint{2.006173in}{0.955244in}}%
\pgfpathlineto{\pgfqpoint{2.006973in}{0.996490in}}%
\pgfpathlineto{\pgfqpoint{2.008483in}{0.946995in}}%
\pgfpathlineto{\pgfqpoint{2.008631in}{0.963493in}}%
\pgfpathlineto{\pgfqpoint{2.009638in}{0.979992in}}%
\pgfpathlineto{\pgfqpoint{2.009653in}{0.971742in}}%
\pgfpathlineto{\pgfqpoint{2.010156in}{0.946995in}}%
\pgfpathlineto{\pgfqpoint{2.010475in}{0.979992in}}%
\pgfpathlineto{\pgfqpoint{2.010638in}{1.004739in}}%
\pgfpathlineto{\pgfqpoint{2.011163in}{0.971742in}}%
\pgfpathlineto{\pgfqpoint{2.011496in}{0.971742in}}%
\pgfpathlineto{\pgfqpoint{2.012674in}{0.996490in}}%
\pgfpathlineto{\pgfqpoint{2.012837in}{0.988241in}}%
\pgfpathlineto{\pgfqpoint{2.013177in}{0.971742in}}%
\pgfpathlineto{\pgfqpoint{2.013488in}{1.004739in}}%
\pgfpathlineto{\pgfqpoint{2.013510in}{0.996490in}}%
\pgfpathlineto{\pgfqpoint{2.014665in}{1.004739in}}%
\pgfpathlineto{\pgfqpoint{2.015184in}{0.996490in}}%
\pgfpathlineto{\pgfqpoint{2.015502in}{1.029487in}}%
\pgfpathlineto{\pgfqpoint{2.015687in}{1.021237in}}%
\pgfpathlineto{\pgfqpoint{2.016694in}{1.045985in}}%
\pgfpathlineto{\pgfqpoint{2.016865in}{1.037736in}}%
\pgfpathlineto{\pgfqpoint{2.018538in}{0.996490in}}%
\pgfpathlineto{\pgfqpoint{2.020863in}{1.078981in}}%
\pgfpathlineto{\pgfqpoint{2.020885in}{1.070732in}}%
\pgfpathlineto{\pgfqpoint{2.021559in}{1.045985in}}%
\pgfpathlineto{\pgfqpoint{2.021722in}{1.070732in}}%
\pgfpathlineto{\pgfqpoint{2.022714in}{1.078981in}}%
\pgfpathlineto{\pgfqpoint{2.022210in}{1.054234in}}%
\pgfpathlineto{\pgfqpoint{2.022729in}{1.070732in}}%
\pgfpathlineto{\pgfqpoint{2.023047in}{1.054234in}}%
\pgfpathlineto{\pgfqpoint{2.023380in}{1.087231in}}%
\pgfpathlineto{\pgfqpoint{2.023736in}{1.120227in}}%
\pgfpathlineto{\pgfqpoint{2.024402in}{1.095480in}}%
\pgfpathlineto{\pgfqpoint{2.024891in}{1.128476in}}%
\pgfpathlineto{\pgfqpoint{2.025750in}{1.111978in}}%
\pgfpathlineto{\pgfqpoint{2.026083in}{1.095480in}}%
\pgfpathlineto{\pgfqpoint{2.026401in}{1.128476in}}%
\pgfpathlineto{\pgfqpoint{2.026564in}{1.128476in}}%
\pgfpathlineto{\pgfqpoint{2.027401in}{1.177971in}}%
\pgfpathlineto{\pgfqpoint{2.027593in}{1.169722in}}%
\pgfpathlineto{\pgfqpoint{2.028593in}{1.194470in}}%
\pgfpathlineto{\pgfqpoint{2.028763in}{1.186221in}}%
\pgfpathlineto{\pgfqpoint{2.028933in}{1.194470in}}%
\pgfpathlineto{\pgfqpoint{2.029918in}{1.177971in}}%
\pgfpathlineto{\pgfqpoint{2.030755in}{1.252214in}}%
\pgfpathlineto{\pgfqpoint{2.030940in}{1.243965in}}%
\pgfpathlineto{\pgfqpoint{2.031429in}{1.227466in}}%
\pgfpathlineto{\pgfqpoint{2.032095in}{1.252214in}}%
\pgfpathlineto{\pgfqpoint{2.032621in}{1.268712in}}%
\pgfpathlineto{\pgfqpoint{2.033124in}{1.243965in}}%
\pgfpathlineto{\pgfqpoint{2.034279in}{1.252214in}}%
\pgfpathlineto{\pgfqpoint{2.035301in}{1.194470in}}%
\pgfpathlineto{\pgfqpoint{2.035449in}{1.210968in}}%
\pgfpathlineto{\pgfqpoint{2.036286in}{1.227466in}}%
\pgfpathlineto{\pgfqpoint{2.036515in}{1.202719in}}%
\pgfpathlineto{\pgfqpoint{2.036641in}{1.219217in}}%
\pgfpathlineto{\pgfqpoint{2.037293in}{1.177971in}}%
\pgfpathlineto{\pgfqpoint{2.038152in}{1.169722in}}%
\pgfpathlineto{\pgfqpoint{2.038322in}{1.194470in}}%
\pgfpathlineto{\pgfqpoint{2.038633in}{1.202719in}}%
\pgfpathlineto{\pgfqpoint{2.038825in}{1.186221in}}%
\pgfpathlineto{\pgfqpoint{2.039492in}{1.144975in}}%
\pgfpathlineto{\pgfqpoint{2.039995in}{1.169722in}}%
\pgfpathlineto{\pgfqpoint{2.040647in}{1.153224in}}%
\pgfpathlineto{\pgfqpoint{2.040817in}{1.177971in}}%
\pgfpathlineto{\pgfqpoint{2.040832in}{1.169722in}}%
\pgfpathlineto{\pgfqpoint{2.041150in}{1.177971in}}%
\pgfpathlineto{\pgfqpoint{2.041483in}{1.153224in}}%
\pgfpathlineto{\pgfqpoint{2.041839in}{1.144975in}}%
\pgfpathlineto{\pgfqpoint{2.042157in}{1.177971in}}%
\pgfpathlineto{\pgfqpoint{2.042513in}{1.169722in}}%
\pgfpathlineto{\pgfqpoint{2.043853in}{1.144975in}}%
\pgfpathlineto{\pgfqpoint{2.045511in}{1.177971in}}%
\pgfpathlineto{\pgfqpoint{2.045526in}{1.169722in}}%
\pgfpathlineto{\pgfqpoint{2.046511in}{1.202719in}}%
\pgfpathlineto{\pgfqpoint{2.046533in}{1.194470in}}%
\pgfpathlineto{\pgfqpoint{2.047207in}{1.111978in}}%
\pgfpathlineto{\pgfqpoint{2.051227in}{0.650025in}}%
\pgfpathlineto{\pgfqpoint{2.051879in}{0.683022in}}%
\pgfpathlineto{\pgfqpoint{2.052397in}{0.666524in}}%
\pgfpathlineto{\pgfqpoint{2.053071in}{0.650025in}}%
\pgfpathlineto{\pgfqpoint{2.052568in}{0.674773in}}%
\pgfpathlineto{\pgfqpoint{2.053241in}{0.674773in}}%
\pgfpathlineto{\pgfqpoint{2.054730in}{0.707769in}}%
\pgfpathlineto{\pgfqpoint{2.054744in}{0.699520in}}%
\pgfpathlineto{\pgfqpoint{2.055418in}{0.674773in}}%
\pgfpathlineto{\pgfqpoint{2.055588in}{0.699520in}}%
\pgfpathlineto{\pgfqpoint{2.056240in}{0.707769in}}%
\pgfpathlineto{\pgfqpoint{2.056425in}{0.691271in}}%
\pgfpathlineto{\pgfqpoint{2.057092in}{0.674773in}}%
\pgfpathlineto{\pgfqpoint{2.057262in}{0.699520in}}%
\pgfpathlineto{\pgfqpoint{2.058084in}{0.707769in}}%
\pgfpathlineto{\pgfqpoint{2.058098in}{0.699520in}}%
\pgfpathlineto{\pgfqpoint{2.058750in}{0.683022in}}%
\pgfpathlineto{\pgfqpoint{2.058750in}{0.716019in}}%
\pgfpathlineto{\pgfqpoint{2.059424in}{0.732517in}}%
\pgfpathlineto{\pgfqpoint{2.059779in}{0.699520in}}%
\pgfpathlineto{\pgfqpoint{2.061097in}{0.658274in}}%
\pgfpathlineto{\pgfqpoint{2.061119in}{0.674773in}}%
\pgfpathlineto{\pgfqpoint{2.063970in}{0.798510in}}%
\pgfpathlineto{\pgfqpoint{2.064133in}{0.790261in}}%
\pgfpathlineto{\pgfqpoint{2.064622in}{0.757264in}}%
\pgfpathlineto{\pgfqpoint{2.064970in}{0.798510in}}%
\pgfpathlineto{\pgfqpoint{2.066125in}{0.806759in}}%
\pgfpathlineto{\pgfqpoint{2.067317in}{0.773763in}}%
\pgfpathlineto{\pgfqpoint{2.067635in}{0.757264in}}%
\pgfpathlineto{\pgfqpoint{2.067805in}{0.782012in}}%
\pgfpathlineto{\pgfqpoint{2.068139in}{0.782012in}}%
\pgfpathlineto{\pgfqpoint{2.068657in}{0.823258in}}%
\pgfpathlineto{\pgfqpoint{2.069160in}{0.815008in}}%
\pgfpathlineto{\pgfqpoint{2.069308in}{0.806759in}}%
\pgfpathlineto{\pgfqpoint{2.069649in}{0.839756in}}%
\pgfpathlineto{\pgfqpoint{2.070671in}{0.897500in}}%
\pgfpathlineto{\pgfqpoint{2.070989in}{0.881002in}}%
\pgfpathlineto{\pgfqpoint{2.071004in}{0.922247in}}%
\pgfpathlineto{\pgfqpoint{2.071330in}{0.913998in}}%
\pgfpathlineto{\pgfqpoint{2.072352in}{0.971742in}}%
\pgfpathlineto{\pgfqpoint{2.072670in}{0.955244in}}%
\pgfpathlineto{\pgfqpoint{2.072685in}{0.996490in}}%
\pgfpathlineto{\pgfqpoint{2.075365in}{1.169722in}}%
\pgfpathlineto{\pgfqpoint{2.076542in}{1.144975in}}%
\pgfpathlineto{\pgfqpoint{2.077882in}{1.120227in}}%
\pgfpathlineto{\pgfqpoint{2.078867in}{1.153224in}}%
\pgfpathlineto{\pgfqpoint{2.079052in}{1.136726in}}%
\pgfpathlineto{\pgfqpoint{2.080541in}{1.078981in}}%
\pgfpathlineto{\pgfqpoint{2.080733in}{1.095480in}}%
\pgfpathlineto{\pgfqpoint{2.080881in}{1.103729in}}%
\pgfpathlineto{\pgfqpoint{2.081229in}{1.062483in}}%
\pgfpathlineto{\pgfqpoint{2.084583in}{0.650025in}}%
\pgfpathlineto{\pgfqpoint{2.084924in}{0.674773in}}%
\pgfpathlineto{\pgfqpoint{2.085079in}{0.683022in}}%
\pgfpathlineto{\pgfqpoint{2.085420in}{0.650025in}}%
\pgfpathlineto{\pgfqpoint{2.085575in}{0.658274in}}%
\pgfpathlineto{\pgfqpoint{2.085923in}{0.650025in}}%
\pgfpathlineto{\pgfqpoint{2.086242in}{0.683022in}}%
\pgfpathlineto{\pgfqpoint{2.086597in}{0.674773in}}%
\pgfpathlineto{\pgfqpoint{2.087582in}{0.658274in}}%
\pgfpathlineto{\pgfqpoint{2.087604in}{0.674773in}}%
\pgfpathlineto{\pgfqpoint{2.090270in}{0.757264in}}%
\pgfpathlineto{\pgfqpoint{2.090285in}{0.749015in}}%
\pgfpathlineto{\pgfqpoint{2.090610in}{0.782012in}}%
\pgfpathlineto{\pgfqpoint{2.091291in}{0.773763in}}%
\pgfpathlineto{\pgfqpoint{2.091447in}{0.782012in}}%
\pgfpathlineto{\pgfqpoint{2.091625in}{0.765513in}}%
\pgfpathlineto{\pgfqpoint{2.093113in}{0.625278in}}%
\pgfpathlineto{\pgfqpoint{2.094290in}{0.633527in}}%
\pgfpathlineto{\pgfqpoint{2.094298in}{0.625278in}}%
\pgfpathlineto{\pgfqpoint{2.095305in}{0.625278in}}%
\pgfpathlineto{\pgfqpoint{2.096637in}{0.633527in}}%
\pgfpathlineto{\pgfqpoint{2.096808in}{0.625278in}}%
\pgfpathlineto{\pgfqpoint{2.097644in}{0.625278in}}%
\pgfpathlineto{\pgfqpoint{2.098644in}{0.633527in}}%
\pgfpathlineto{\pgfqpoint{2.098644in}{0.625278in}}%
\pgfpathlineto{\pgfqpoint{2.099821in}{0.633527in}}%
\pgfpathlineto{\pgfqpoint{2.099821in}{0.625278in}}%
\pgfpathlineto{\pgfqpoint{2.100836in}{0.625278in}}%
\pgfpathlineto{\pgfqpoint{2.102168in}{0.633527in}}%
\pgfpathlineto{\pgfqpoint{2.102176in}{0.625278in}}%
\pgfpathlineto{\pgfqpoint{2.103175in}{0.625278in}}%
\pgfpathlineto{\pgfqpoint{2.104352in}{0.633527in}}%
\pgfpathlineto{\pgfqpoint{2.104515in}{0.625278in}}%
\pgfpathlineto{\pgfqpoint{2.105359in}{0.625278in}}%
\pgfpathlineto{\pgfqpoint{2.106700in}{0.633527in}}%
\pgfpathlineto{\pgfqpoint{2.106700in}{0.625278in}}%
\pgfpathlineto{\pgfqpoint{2.107707in}{0.625278in}}%
\pgfpathlineto{\pgfqpoint{2.108869in}{0.633527in}}%
\pgfpathlineto{\pgfqpoint{2.109047in}{0.625278in}}%
\pgfpathlineto{\pgfqpoint{2.109891in}{0.625278in}}%
\pgfpathlineto{\pgfqpoint{2.110890in}{0.633527in}}%
\pgfpathlineto{\pgfqpoint{2.110890in}{0.625278in}}%
\pgfpathlineto{\pgfqpoint{2.112231in}{0.633527in}}%
\pgfpathlineto{\pgfqpoint{2.112564in}{0.625278in}}%
\pgfpathlineto{\pgfqpoint{2.113230in}{0.625278in}}%
\pgfpathlineto{\pgfqpoint{2.114415in}{0.633527in}}%
\pgfpathlineto{\pgfqpoint{2.114570in}{0.625278in}}%
\pgfpathlineto{\pgfqpoint{2.115414in}{0.625278in}}%
\pgfpathlineto{\pgfqpoint{2.116584in}{0.633527in}}%
\pgfpathlineto{\pgfqpoint{2.116755in}{0.625278in}}%
\pgfpathlineto{\pgfqpoint{2.117428in}{0.625278in}}%
\pgfpathlineto{\pgfqpoint{2.118606in}{0.633527in}}%
\pgfpathlineto{\pgfqpoint{2.118761in}{0.625278in}}%
\pgfpathlineto{\pgfqpoint{2.119605in}{0.625278in}}%
\pgfpathlineto{\pgfqpoint{2.120782in}{0.633527in}}%
\pgfpathlineto{\pgfqpoint{2.120953in}{0.625278in}}%
\pgfpathlineto{\pgfqpoint{2.121789in}{0.625278in}}%
\pgfpathlineto{\pgfqpoint{2.122959in}{0.633527in}}%
\pgfpathlineto{\pgfqpoint{2.123292in}{0.625278in}}%
\pgfpathlineto{\pgfqpoint{2.123966in}{0.625278in}}%
\pgfpathlineto{\pgfqpoint{2.125299in}{0.633527in}}%
\pgfpathlineto{\pgfqpoint{2.125469in}{0.625278in}}%
\pgfpathlineto{\pgfqpoint{2.126313in}{0.625278in}}%
\pgfpathlineto{\pgfqpoint{2.127483in}{0.633527in}}%
\pgfpathlineto{\pgfqpoint{2.127646in}{0.625278in}}%
\pgfpathlineto{\pgfqpoint{2.128483in}{0.625278in}}%
\pgfpathlineto{\pgfqpoint{2.129660in}{0.633527in}}%
\pgfpathlineto{\pgfqpoint{2.129823in}{0.625278in}}%
\pgfpathlineto{\pgfqpoint{2.130667in}{0.625278in}}%
\pgfpathlineto{\pgfqpoint{2.131674in}{0.633527in}}%
\pgfpathlineto{\pgfqpoint{2.131674in}{0.625278in}}%
\pgfpathlineto{\pgfqpoint{2.132844in}{0.633527in}}%
\pgfpathlineto{\pgfqpoint{2.133014in}{0.625278in}}%
\pgfpathlineto{\pgfqpoint{2.133851in}{0.625278in}}%
\pgfpathlineto{\pgfqpoint{2.135021in}{0.633527in}}%
\pgfpathlineto{\pgfqpoint{2.135191in}{0.625278in}}%
\pgfpathlineto{\pgfqpoint{2.135702in}{0.625278in}}%
\pgfpathlineto{\pgfqpoint{2.136872in}{0.633527in}}%
\pgfpathlineto{\pgfqpoint{2.137035in}{0.625278in}}%
\pgfpathlineto{\pgfqpoint{2.137879in}{0.625278in}}%
\pgfpathlineto{\pgfqpoint{2.139049in}{0.633527in}}%
\pgfpathlineto{\pgfqpoint{2.139211in}{0.625278in}}%
\pgfpathlineto{\pgfqpoint{2.140056in}{0.625278in}}%
\pgfpathlineto{\pgfqpoint{2.141225in}{0.633527in}}%
\pgfpathlineto{\pgfqpoint{2.141559in}{0.625278in}}%
\pgfpathlineto{\pgfqpoint{2.142232in}{0.625278in}}%
\pgfpathlineto{\pgfqpoint{2.143402in}{0.633527in}}%
\pgfpathlineto{\pgfqpoint{2.143410in}{0.625278in}}%
\pgfpathlineto{\pgfqpoint{2.144246in}{0.625278in}}%
\pgfpathlineto{\pgfqpoint{2.145586in}{0.633527in}}%
\pgfpathlineto{\pgfqpoint{2.145749in}{0.625278in}}%
\pgfpathlineto{\pgfqpoint{2.146423in}{0.625278in}}%
\pgfpathlineto{\pgfqpoint{2.147593in}{0.633527in}}%
\pgfpathlineto{\pgfqpoint{2.147763in}{0.625278in}}%
\pgfpathlineto{\pgfqpoint{2.148600in}{0.625278in}}%
\pgfpathlineto{\pgfqpoint{2.149777in}{0.633527in}}%
\pgfpathlineto{\pgfqpoint{2.149940in}{0.625278in}}%
\pgfpathlineto{\pgfqpoint{2.150777in}{0.625278in}}%
\pgfpathlineto{\pgfqpoint{2.151954in}{0.633527in}}%
\pgfpathlineto{\pgfqpoint{2.151954in}{0.625278in}}%
\pgfpathlineto{\pgfqpoint{2.152791in}{0.625278in}}%
\pgfpathlineto{\pgfqpoint{2.154131in}{0.633527in}}%
\pgfpathlineto{\pgfqpoint{2.154131in}{0.625278in}}%
\pgfpathlineto{\pgfqpoint{2.155138in}{0.625278in}}%
\pgfpathlineto{\pgfqpoint{2.156315in}{0.633527in}}%
\pgfpathlineto{\pgfqpoint{2.156648in}{0.625278in}}%
\pgfpathlineto{\pgfqpoint{2.157315in}{0.625278in}}%
\pgfpathlineto{\pgfqpoint{2.158492in}{0.633527in}}%
\pgfpathlineto{\pgfqpoint{2.158662in}{0.625278in}}%
\pgfpathlineto{\pgfqpoint{2.159499in}{0.625278in}}%
\pgfpathlineto{\pgfqpoint{2.160676in}{0.633527in}}%
\pgfpathlineto{\pgfqpoint{2.160839in}{0.625278in}}%
\pgfpathlineto{\pgfqpoint{2.161676in}{0.625278in}}%
\pgfpathlineto{\pgfqpoint{2.162683in}{0.633527in}}%
\pgfpathlineto{\pgfqpoint{2.162683in}{0.625278in}}%
\pgfpathlineto{\pgfqpoint{2.164023in}{0.633527in}}%
\pgfpathlineto{\pgfqpoint{2.164193in}{0.625278in}}%
\pgfpathlineto{\pgfqpoint{2.165030in}{0.625278in}}%
\pgfpathlineto{\pgfqpoint{2.166200in}{0.633527in}}%
\pgfpathlineto{\pgfqpoint{2.166370in}{0.625278in}}%
\pgfpathlineto{\pgfqpoint{2.167207in}{0.625278in}}%
\pgfpathlineto{\pgfqpoint{2.168384in}{0.633527in}}%
\pgfpathlineto{\pgfqpoint{2.168384in}{0.625278in}}%
\pgfpathlineto{\pgfqpoint{2.169391in}{0.625278in}}%
\pgfpathlineto{\pgfqpoint{2.170561in}{0.633527in}}%
\pgfpathlineto{\pgfqpoint{2.170731in}{0.625278in}}%
\pgfpathlineto{\pgfqpoint{2.171568in}{0.625278in}}%
\pgfpathlineto{\pgfqpoint{2.172738in}{0.633527in}}%
\pgfpathlineto{\pgfqpoint{2.172738in}{0.625278in}}%
\pgfpathlineto{\pgfqpoint{2.173745in}{0.625278in}}%
\pgfpathlineto{\pgfqpoint{2.174922in}{0.633527in}}%
\pgfpathlineto{\pgfqpoint{2.174922in}{0.625278in}}%
\pgfpathlineto{\pgfqpoint{2.175921in}{0.625278in}}%
\pgfpathlineto{\pgfqpoint{2.177099in}{0.633527in}}%
\pgfpathlineto{\pgfqpoint{2.177269in}{0.625278in}}%
\pgfpathlineto{\pgfqpoint{2.177935in}{0.625278in}}%
\pgfpathlineto{\pgfqpoint{2.179113in}{0.633527in}}%
\pgfpathlineto{\pgfqpoint{2.179120in}{0.625278in}}%
\pgfpathlineto{\pgfqpoint{2.180112in}{0.625278in}}%
\pgfpathlineto{\pgfqpoint{2.181460in}{0.633527in}}%
\pgfpathlineto{\pgfqpoint{2.181460in}{0.625278in}}%
\pgfpathlineto{\pgfqpoint{2.182459in}{0.625278in}}%
\pgfpathlineto{\pgfqpoint{2.183807in}{0.633527in}}%
\pgfpathlineto{\pgfqpoint{2.183970in}{0.625278in}}%
\pgfpathlineto{\pgfqpoint{2.184806in}{0.625278in}}%
\pgfpathlineto{\pgfqpoint{2.185984in}{0.633527in}}%
\pgfpathlineto{\pgfqpoint{2.186154in}{0.625278in}}%
\pgfpathlineto{\pgfqpoint{2.186991in}{0.625278in}}%
\pgfpathlineto{\pgfqpoint{2.188161in}{0.633527in}}%
\pgfpathlineto{\pgfqpoint{2.188331in}{0.625278in}}%
\pgfpathlineto{\pgfqpoint{2.189168in}{0.625278in}}%
\pgfpathlineto{\pgfqpoint{2.190345in}{0.633527in}}%
\pgfpathlineto{\pgfqpoint{2.190345in}{0.625278in}}%
\pgfpathlineto{\pgfqpoint{2.191344in}{0.625278in}}%
\pgfpathlineto{\pgfqpoint{2.192522in}{0.633527in}}%
\pgfpathlineto{\pgfqpoint{2.192522in}{0.625278in}}%
\pgfpathlineto{\pgfqpoint{2.193529in}{0.625278in}}%
\pgfpathlineto{\pgfqpoint{2.194869in}{0.633527in}}%
\pgfpathlineto{\pgfqpoint{2.194869in}{0.625278in}}%
\pgfpathlineto{\pgfqpoint{2.195876in}{0.625278in}}%
\pgfpathlineto{\pgfqpoint{2.197046in}{0.633527in}}%
\pgfpathlineto{\pgfqpoint{2.197216in}{0.625278in}}%
\pgfpathlineto{\pgfqpoint{2.198053in}{0.625278in}}%
\pgfpathlineto{\pgfqpoint{2.199222in}{0.633527in}}%
\pgfpathlineto{\pgfqpoint{2.199393in}{0.625278in}}%
\pgfpathlineto{\pgfqpoint{2.200229in}{0.625278in}}%
\pgfpathlineto{\pgfqpoint{2.201407in}{0.633527in}}%
\pgfpathlineto{\pgfqpoint{2.201570in}{0.625278in}}%
\pgfpathlineto{\pgfqpoint{2.202073in}{0.625278in}}%
\pgfpathlineto{\pgfqpoint{2.203250in}{0.633527in}}%
\pgfpathlineto{\pgfqpoint{2.203413in}{0.625278in}}%
\pgfpathlineto{\pgfqpoint{2.204087in}{0.625278in}}%
\pgfpathlineto{\pgfqpoint{2.205257in}{0.633527in}}%
\pgfpathlineto{\pgfqpoint{2.205427in}{0.625278in}}%
\pgfpathlineto{\pgfqpoint{2.206264in}{0.625278in}}%
\pgfpathlineto{\pgfqpoint{2.207441in}{0.633527in}}%
\pgfpathlineto{\pgfqpoint{2.207441in}{0.625278in}}%
\pgfpathlineto{\pgfqpoint{2.208448in}{0.625278in}}%
\pgfpathlineto{\pgfqpoint{2.209618in}{0.633527in}}%
\pgfpathlineto{\pgfqpoint{2.209788in}{0.625278in}}%
\pgfpathlineto{\pgfqpoint{2.210455in}{0.625278in}}%
\pgfpathlineto{\pgfqpoint{2.211632in}{0.633527in}}%
\pgfpathlineto{\pgfqpoint{2.211795in}{0.625278in}}%
\pgfpathlineto{\pgfqpoint{2.212654in}{0.625278in}}%
\pgfpathlineto{\pgfqpoint{2.214312in}{0.658274in}}%
\pgfpathlineto{\pgfqpoint{2.214838in}{0.650025in}}%
\pgfpathlineto{\pgfqpoint{2.215001in}{0.674773in}}%
\pgfpathlineto{\pgfqpoint{2.215334in}{0.674773in}}%
\pgfpathlineto{\pgfqpoint{2.215489in}{0.683022in}}%
\pgfpathlineto{\pgfqpoint{2.215823in}{0.658274in}}%
\pgfpathlineto{\pgfqpoint{2.216008in}{0.666524in}}%
\pgfpathlineto{\pgfqpoint{2.216156in}{0.658274in}}%
\pgfpathlineto{\pgfqpoint{2.216326in}{0.683022in}}%
\pgfpathlineto{\pgfqpoint{2.216682in}{0.674773in}}%
\pgfpathlineto{\pgfqpoint{2.219680in}{0.782012in}}%
\pgfpathlineto{\pgfqpoint{2.220028in}{0.773763in}}%
\pgfpathlineto{\pgfqpoint{2.220680in}{0.806759in}}%
\pgfpathlineto{\pgfqpoint{2.220702in}{0.798510in}}%
\pgfpathlineto{\pgfqpoint{2.221354in}{0.806759in}}%
\pgfpathlineto{\pgfqpoint{2.221539in}{0.790261in}}%
\pgfpathlineto{\pgfqpoint{2.221687in}{0.782012in}}%
\pgfpathlineto{\pgfqpoint{2.221857in}{0.806759in}}%
\pgfpathlineto{\pgfqpoint{2.222212in}{0.798510in}}%
\pgfpathlineto{\pgfqpoint{2.223197in}{0.831507in}}%
\pgfpathlineto{\pgfqpoint{2.223382in}{0.815008in}}%
\pgfpathlineto{\pgfqpoint{2.223553in}{0.823258in}}%
\pgfpathlineto{\pgfqpoint{2.224560in}{0.798510in}}%
\pgfpathlineto{\pgfqpoint{2.225211in}{0.831507in}}%
\pgfpathlineto{\pgfqpoint{2.225715in}{0.806759in}}%
\pgfpathlineto{\pgfqpoint{2.225900in}{0.798510in}}%
\pgfpathlineto{\pgfqpoint{2.226211in}{0.831507in}}%
\pgfpathlineto{\pgfqpoint{2.226736in}{0.815008in}}%
\pgfpathlineto{\pgfqpoint{2.228225in}{0.782012in}}%
\pgfpathlineto{\pgfqpoint{2.228225in}{0.790261in}}%
\pgfpathlineto{\pgfqpoint{2.228247in}{0.823258in}}%
\pgfpathlineto{\pgfqpoint{2.229254in}{0.798510in}}%
\pgfpathlineto{\pgfqpoint{2.229905in}{0.831507in}}%
\pgfpathlineto{\pgfqpoint{2.230402in}{0.806759in}}%
\pgfpathlineto{\pgfqpoint{2.230927in}{0.823258in}}%
\pgfpathlineto{\pgfqpoint{2.231431in}{0.798510in}}%
\pgfpathlineto{\pgfqpoint{2.232415in}{0.782012in}}%
\pgfpathlineto{\pgfqpoint{2.232438in}{0.798510in}}%
\pgfpathlineto{\pgfqpoint{2.233608in}{0.848005in}}%
\pgfpathlineto{\pgfqpoint{2.234096in}{0.831507in}}%
\pgfpathlineto{\pgfqpoint{2.234281in}{0.823258in}}%
\pgfpathlineto{\pgfqpoint{2.234948in}{0.872753in}}%
\pgfpathlineto{\pgfqpoint{2.235118in}{0.864503in}}%
\pgfpathlineto{\pgfqpoint{2.235940in}{0.831507in}}%
\pgfpathlineto{\pgfqpoint{2.236288in}{0.848005in}}%
\pgfpathlineto{\pgfqpoint{2.236606in}{0.856254in}}%
\pgfpathlineto{\pgfqpoint{2.236939in}{0.831507in}}%
\pgfpathlineto{\pgfqpoint{2.237798in}{0.823258in}}%
\pgfpathlineto{\pgfqpoint{2.237969in}{0.848005in}}%
\pgfpathlineto{\pgfqpoint{2.238302in}{0.872753in}}%
\pgfpathlineto{\pgfqpoint{2.239124in}{0.856254in}}%
\pgfpathlineto{\pgfqpoint{2.240145in}{0.848005in}}%
\pgfpathlineto{\pgfqpoint{2.241804in}{0.881002in}}%
\pgfpathlineto{\pgfqpoint{2.242307in}{0.856254in}}%
\pgfpathlineto{\pgfqpoint{2.242826in}{0.872753in}}%
\pgfpathlineto{\pgfqpoint{2.244507in}{0.946995in}}%
\pgfpathlineto{\pgfqpoint{2.244825in}{0.930497in}}%
\pgfpathlineto{\pgfqpoint{2.245321in}{0.905749in}}%
\pgfpathlineto{\pgfqpoint{2.245884in}{0.955244in}}%
\pgfpathlineto{\pgfqpoint{2.246669in}{0.930497in}}%
\pgfpathlineto{\pgfqpoint{2.246854in}{0.946995in}}%
\pgfpathlineto{\pgfqpoint{2.247187in}{0.971742in}}%
\pgfpathlineto{\pgfqpoint{2.248009in}{0.955244in}}%
\pgfpathlineto{\pgfqpoint{2.248512in}{0.930497in}}%
\pgfpathlineto{\pgfqpoint{2.249016in}{0.979992in}}%
\pgfpathlineto{\pgfqpoint{2.249030in}{0.971742in}}%
\pgfpathlineto{\pgfqpoint{2.249179in}{0.979992in}}%
\pgfpathlineto{\pgfqpoint{2.249534in}{0.946995in}}%
\pgfpathlineto{\pgfqpoint{2.249682in}{0.955244in}}%
\pgfpathlineto{\pgfqpoint{2.250186in}{0.979992in}}%
\pgfpathlineto{\pgfqpoint{2.250704in}{0.946995in}}%
\pgfpathlineto{\pgfqpoint{2.251022in}{0.930497in}}%
\pgfpathlineto{\pgfqpoint{2.251192in}{0.955244in}}%
\pgfpathlineto{\pgfqpoint{2.251526in}{0.955244in}}%
\pgfpathlineto{\pgfqpoint{2.251696in}{0.979992in}}%
\pgfpathlineto{\pgfqpoint{2.252555in}{0.971742in}}%
\pgfpathlineto{\pgfqpoint{2.253873in}{1.004739in}}%
\pgfpathlineto{\pgfqpoint{2.253895in}{0.996490in}}%
\pgfpathlineto{\pgfqpoint{2.255065in}{0.971742in}}%
\pgfpathlineto{\pgfqpoint{2.255213in}{0.979992in}}%
\pgfpathlineto{\pgfqpoint{2.255383in}{1.004739in}}%
\pgfpathlineto{\pgfqpoint{2.256220in}{0.955244in}}%
\pgfpathlineto{\pgfqpoint{2.256242in}{0.971742in}}%
\pgfpathlineto{\pgfqpoint{2.257412in}{0.946995in}}%
\pgfpathlineto{\pgfqpoint{2.257560in}{0.955244in}}%
\pgfpathlineto{\pgfqpoint{2.258064in}{1.004739in}}%
\pgfpathlineto{\pgfqpoint{2.258589in}{0.996490in}}%
\pgfpathlineto{\pgfqpoint{2.259071in}{1.029487in}}%
\pgfpathlineto{\pgfqpoint{2.259737in}{1.004739in}}%
\pgfpathlineto{\pgfqpoint{2.260766in}{0.996490in}}%
\pgfpathlineto{\pgfqpoint{2.261936in}{1.021237in}}%
\pgfpathlineto{\pgfqpoint{2.262106in}{1.012988in}}%
\pgfpathlineto{\pgfqpoint{2.262254in}{1.004739in}}%
\pgfpathlineto{\pgfqpoint{2.262425in}{1.029487in}}%
\pgfpathlineto{\pgfqpoint{2.262780in}{1.021237in}}%
\pgfpathlineto{\pgfqpoint{2.263432in}{1.004739in}}%
\pgfpathlineto{\pgfqpoint{2.263928in}{1.029487in}}%
\pgfpathlineto{\pgfqpoint{2.265127in}{0.996490in}}%
\pgfpathlineto{\pgfqpoint{2.266127in}{1.021237in}}%
\pgfpathlineto{\pgfqpoint{2.266297in}{1.012988in}}%
\pgfpathlineto{\pgfqpoint{2.266467in}{1.021237in}}%
\pgfpathlineto{\pgfqpoint{2.267467in}{0.996490in}}%
\pgfpathlineto{\pgfqpoint{2.267622in}{1.004739in}}%
\pgfpathlineto{\pgfqpoint{2.267970in}{0.971742in}}%
\pgfpathlineto{\pgfqpoint{2.268118in}{0.979992in}}%
\pgfpathlineto{\pgfqpoint{2.268474in}{0.971742in}}%
\pgfpathlineto{\pgfqpoint{2.268644in}{0.996490in}}%
\pgfpathlineto{\pgfqpoint{2.269148in}{0.988241in}}%
\pgfpathlineto{\pgfqpoint{2.270303in}{0.955244in}}%
\pgfpathlineto{\pgfqpoint{2.269318in}{0.996490in}}%
\pgfpathlineto{\pgfqpoint{2.270303in}{0.963493in}}%
\pgfpathlineto{\pgfqpoint{2.270318in}{0.996490in}}%
\pgfpathlineto{\pgfqpoint{2.271324in}{0.971742in}}%
\pgfpathlineto{\pgfqpoint{2.271813in}{1.004739in}}%
\pgfpathlineto{\pgfqpoint{2.272161in}{0.971742in}}%
\pgfpathlineto{\pgfqpoint{2.272835in}{0.946995in}}%
\pgfpathlineto{\pgfqpoint{2.273005in}{0.971742in}}%
\pgfpathlineto{\pgfqpoint{2.273487in}{1.004739in}}%
\pgfpathlineto{\pgfqpoint{2.273842in}{0.963493in}}%
\pgfpathlineto{\pgfqpoint{2.274175in}{0.946995in}}%
\pgfpathlineto{\pgfqpoint{2.274493in}{0.979992in}}%
\pgfpathlineto{\pgfqpoint{2.274545in}{0.979992in}}%
\pgfpathlineto{\pgfqpoint{2.275663in}{1.004739in}}%
\pgfpathlineto{\pgfqpoint{2.275686in}{0.996490in}}%
\pgfpathlineto{\pgfqpoint{2.276352in}{0.971742in}}%
\pgfpathlineto{\pgfqpoint{2.276522in}{0.996490in}}%
\pgfpathlineto{\pgfqpoint{2.277174in}{1.004739in}}%
\pgfpathlineto{\pgfqpoint{2.277359in}{0.988241in}}%
\pgfpathlineto{\pgfqpoint{2.277529in}{0.996490in}}%
\pgfpathlineto{\pgfqpoint{2.278536in}{0.971742in}}%
\pgfpathlineto{\pgfqpoint{2.279188in}{0.979992in}}%
\pgfpathlineto{\pgfqpoint{2.279373in}{0.963493in}}%
\pgfpathlineto{\pgfqpoint{2.280543in}{0.946995in}}%
\pgfpathlineto{\pgfqpoint{2.281535in}{0.955244in}}%
\pgfpathlineto{\pgfqpoint{2.281550in}{0.946995in}}%
\pgfpathlineto{\pgfqpoint{2.283230in}{0.897500in}}%
\pgfpathlineto{\pgfqpoint{2.283267in}{0.905749in}}%
\pgfpathlineto{\pgfqpoint{2.284067in}{0.971742in}}%
\pgfpathlineto{\pgfqpoint{2.284904in}{0.946995in}}%
\pgfpathlineto{\pgfqpoint{2.285407in}{0.922247in}}%
\pgfpathlineto{\pgfqpoint{2.285726in}{0.955244in}}%
\pgfpathlineto{\pgfqpoint{2.285740in}{0.946995in}}%
\pgfpathlineto{\pgfqpoint{2.287399in}{0.979992in}}%
\pgfpathlineto{\pgfqpoint{2.287421in}{0.971742in}}%
\pgfpathlineto{\pgfqpoint{2.288258in}{1.021237in}}%
\pgfpathlineto{\pgfqpoint{2.288591in}{1.045985in}}%
\pgfpathlineto{\pgfqpoint{2.289095in}{1.021237in}}%
\pgfpathlineto{\pgfqpoint{2.289746in}{1.004739in}}%
\pgfpathlineto{\pgfqpoint{2.289916in}{1.029487in}}%
\pgfpathlineto{\pgfqpoint{2.289931in}{1.021237in}}%
\pgfpathlineto{\pgfqpoint{2.291086in}{1.029487in}}%
\pgfpathlineto{\pgfqpoint{2.292093in}{1.004739in}}%
\pgfpathlineto{\pgfqpoint{2.292115in}{1.021237in}}%
\pgfpathlineto{\pgfqpoint{2.292597in}{1.029487in}}%
\pgfpathlineto{\pgfqpoint{2.293456in}{0.996490in}}%
\pgfpathlineto{\pgfqpoint{2.294270in}{1.004739in}}%
\pgfpathlineto{\pgfqpoint{2.294292in}{0.996490in}}%
\pgfpathlineto{\pgfqpoint{2.294796in}{0.971742in}}%
\pgfpathlineto{\pgfqpoint{2.295114in}{1.004739in}}%
\pgfpathlineto{\pgfqpoint{2.295129in}{0.996490in}}%
\pgfpathlineto{\pgfqpoint{2.295277in}{1.004739in}}%
\pgfpathlineto{\pgfqpoint{2.295632in}{0.971742in}}%
\pgfpathlineto{\pgfqpoint{2.295781in}{0.979992in}}%
\pgfpathlineto{\pgfqpoint{2.295803in}{0.971742in}}%
\pgfpathlineto{\pgfqpoint{2.296284in}{1.029487in}}%
\pgfpathlineto{\pgfqpoint{2.296639in}{1.021237in}}%
\pgfpathlineto{\pgfqpoint{2.297291in}{1.029487in}}%
\pgfpathlineto{\pgfqpoint{2.297476in}{1.012988in}}%
\pgfpathlineto{\pgfqpoint{2.297957in}{0.979992in}}%
\pgfpathlineto{\pgfqpoint{2.298313in}{1.021237in}}%
\pgfpathlineto{\pgfqpoint{2.298964in}{1.004739in}}%
\pgfpathlineto{\pgfqpoint{2.299468in}{1.029487in}}%
\pgfpathlineto{\pgfqpoint{2.299653in}{1.021237in}}%
\pgfpathlineto{\pgfqpoint{2.300475in}{1.062483in}}%
\pgfpathlineto{\pgfqpoint{2.301482in}{1.078981in}}%
\pgfpathlineto{\pgfqpoint{2.300808in}{1.054234in}}%
\pgfpathlineto{\pgfqpoint{2.301497in}{1.070732in}}%
\pgfpathlineto{\pgfqpoint{2.302652in}{1.078981in}}%
\pgfpathlineto{\pgfqpoint{2.302844in}{1.070732in}}%
\pgfpathlineto{\pgfqpoint{2.303007in}{1.095480in}}%
\pgfpathlineto{\pgfqpoint{2.303681in}{1.095480in}}%
\pgfpathlineto{\pgfqpoint{2.304851in}{1.070732in}}%
\pgfpathlineto{\pgfqpoint{2.304999in}{1.078981in}}%
\pgfpathlineto{\pgfqpoint{2.305510in}{1.128476in}}%
\pgfpathlineto{\pgfqpoint{2.306028in}{1.120227in}}%
\pgfpathlineto{\pgfqpoint{2.307686in}{1.202719in}}%
\pgfpathlineto{\pgfqpoint{2.307872in}{1.186221in}}%
\pgfpathlineto{\pgfqpoint{2.312566in}{0.625278in}}%
\pgfpathlineto{\pgfqpoint{2.313565in}{0.650025in}}%
\pgfpathlineto{\pgfqpoint{2.313736in}{0.641776in}}%
\pgfpathlineto{\pgfqpoint{2.313884in}{0.633527in}}%
\pgfpathlineto{\pgfqpoint{2.314387in}{0.666524in}}%
\pgfpathlineto{\pgfqpoint{2.314720in}{0.683022in}}%
\pgfpathlineto{\pgfqpoint{2.315076in}{0.650025in}}%
\pgfpathlineto{\pgfqpoint{2.315416in}{0.650025in}}%
\pgfpathlineto{\pgfqpoint{2.316912in}{0.707769in}}%
\pgfpathlineto{\pgfqpoint{2.317090in}{0.691271in}}%
\pgfpathlineto{\pgfqpoint{2.318430in}{0.650025in}}%
\pgfpathlineto{\pgfqpoint{2.318578in}{0.666524in}}%
\pgfpathlineto{\pgfqpoint{2.319252in}{0.683022in}}%
\pgfpathlineto{\pgfqpoint{2.319607in}{0.650025in}}%
\pgfpathlineto{\pgfqpoint{2.320592in}{0.683022in}}%
\pgfpathlineto{\pgfqpoint{2.320777in}{0.666524in}}%
\pgfpathlineto{\pgfqpoint{2.320925in}{0.658274in}}%
\pgfpathlineto{\pgfqpoint{2.321095in}{0.683022in}}%
\pgfpathlineto{\pgfqpoint{2.321451in}{0.674773in}}%
\pgfpathlineto{\pgfqpoint{2.321599in}{0.683022in}}%
\pgfpathlineto{\pgfqpoint{2.321947in}{0.650025in}}%
\pgfpathlineto{\pgfqpoint{2.322265in}{0.658274in}}%
\pgfpathlineto{\pgfqpoint{2.323124in}{0.625278in}}%
\pgfpathlineto{\pgfqpoint{2.323295in}{0.650025in}}%
\pgfpathlineto{\pgfqpoint{2.324450in}{0.683022in}}%
\pgfpathlineto{\pgfqpoint{2.324635in}{0.666524in}}%
\pgfpathlineto{\pgfqpoint{2.325138in}{0.650025in}}%
\pgfpathlineto{\pgfqpoint{2.325301in}{0.674773in}}%
\pgfpathlineto{\pgfqpoint{2.325627in}{0.666524in}}%
\pgfpathlineto{\pgfqpoint{2.325975in}{0.674773in}}%
\pgfpathlineto{\pgfqpoint{2.326641in}{0.625278in}}%
\pgfpathlineto{\pgfqpoint{2.327796in}{0.633527in}}%
\pgfpathlineto{\pgfqpoint{2.327967in}{0.625278in}}%
\pgfpathlineto{\pgfqpoint{2.328803in}{0.625278in}}%
\pgfpathlineto{\pgfqpoint{2.329981in}{0.633527in}}%
\pgfpathlineto{\pgfqpoint{2.329981in}{0.625278in}}%
\pgfpathlineto{\pgfqpoint{2.331002in}{0.625278in}}%
\pgfpathlineto{\pgfqpoint{2.332676in}{0.658274in}}%
\pgfpathlineto{\pgfqpoint{2.332676in}{0.650025in}}%
\pgfpathlineto{\pgfqpoint{2.333683in}{0.674773in}}%
\pgfpathlineto{\pgfqpoint{2.334001in}{0.658274in}}%
\pgfpathlineto{\pgfqpoint{2.334171in}{0.683022in}}%
\pgfpathlineto{\pgfqpoint{2.334519in}{0.674773in}}%
\pgfpathlineto{\pgfqpoint{2.335689in}{0.732517in}}%
\pgfpathlineto{\pgfqpoint{2.336193in}{0.707769in}}%
\pgfpathlineto{\pgfqpoint{2.336370in}{0.724268in}}%
\pgfpathlineto{\pgfqpoint{2.337207in}{0.674773in}}%
\pgfpathlineto{\pgfqpoint{2.338021in}{0.707769in}}%
\pgfpathlineto{\pgfqpoint{2.338362in}{0.683022in}}%
\pgfpathlineto{\pgfqpoint{2.338880in}{0.699520in}}%
\pgfpathlineto{\pgfqpoint{2.339384in}{0.674773in}}%
\pgfpathlineto{\pgfqpoint{2.341050in}{0.732517in}}%
\pgfpathlineto{\pgfqpoint{2.341228in}{0.716019in}}%
\pgfpathlineto{\pgfqpoint{2.341398in}{0.724268in}}%
\pgfpathlineto{\pgfqpoint{2.342383in}{0.707769in}}%
\pgfpathlineto{\pgfqpoint{2.343552in}{0.806759in}}%
\pgfpathlineto{\pgfqpoint{2.343575in}{0.798510in}}%
\pgfpathlineto{\pgfqpoint{2.344559in}{0.831507in}}%
\pgfpathlineto{\pgfqpoint{2.344752in}{0.815008in}}%
\pgfpathlineto{\pgfqpoint{2.345070in}{0.831507in}}%
\pgfpathlineto{\pgfqpoint{2.346425in}{0.773763in}}%
\pgfpathlineto{\pgfqpoint{2.347247in}{0.806759in}}%
\pgfpathlineto{\pgfqpoint{2.347765in}{0.790261in}}%
\pgfpathlineto{\pgfqpoint{2.348247in}{0.782012in}}%
\pgfpathlineto{\pgfqpoint{2.348269in}{0.823258in}}%
\pgfpathlineto{\pgfqpoint{2.350113in}{0.946995in}}%
\pgfpathlineto{\pgfqpoint{2.350772in}{0.905749in}}%
\pgfpathlineto{\pgfqpoint{2.351956in}{0.815008in}}%
\pgfpathlineto{\pgfqpoint{2.353111in}{0.782012in}}%
\pgfpathlineto{\pgfqpoint{2.353133in}{0.798510in}}%
\pgfpathlineto{\pgfqpoint{2.354785in}{0.831507in}}%
\pgfpathlineto{\pgfqpoint{2.355125in}{0.856254in}}%
\pgfpathlineto{\pgfqpoint{2.355977in}{0.798510in}}%
\pgfpathlineto{\pgfqpoint{2.356465in}{0.806759in}}%
\pgfpathlineto{\pgfqpoint{2.356650in}{0.790261in}}%
\pgfpathlineto{\pgfqpoint{2.357132in}{0.782012in}}%
\pgfpathlineto{\pgfqpoint{2.357154in}{0.823258in}}%
\pgfpathlineto{\pgfqpoint{2.358324in}{0.773763in}}%
\pgfpathlineto{\pgfqpoint{2.358494in}{0.798510in}}%
\pgfpathlineto{\pgfqpoint{2.359649in}{0.831507in}}%
\pgfpathlineto{\pgfqpoint{2.359834in}{0.815008in}}%
\pgfpathlineto{\pgfqpoint{2.362685in}{0.699520in}}%
\pgfpathlineto{\pgfqpoint{2.360153in}{0.831507in}}%
\pgfpathlineto{\pgfqpoint{2.362833in}{0.707769in}}%
\pgfpathlineto{\pgfqpoint{2.363003in}{0.732517in}}%
\pgfpathlineto{\pgfqpoint{2.363522in}{0.699520in}}%
\pgfpathlineto{\pgfqpoint{2.363855in}{0.724268in}}%
\pgfpathlineto{\pgfqpoint{2.364173in}{0.732517in}}%
\pgfpathlineto{\pgfqpoint{2.364358in}{0.716019in}}%
\pgfpathlineto{\pgfqpoint{2.364862in}{0.674773in}}%
\pgfpathlineto{\pgfqpoint{2.365535in}{0.699520in}}%
\pgfpathlineto{\pgfqpoint{2.366017in}{0.707769in}}%
\pgfpathlineto{\pgfqpoint{2.366202in}{0.691271in}}%
\pgfpathlineto{\pgfqpoint{2.366876in}{0.674773in}}%
\pgfpathlineto{\pgfqpoint{2.367046in}{0.699520in}}%
\pgfpathlineto{\pgfqpoint{2.367357in}{0.683022in}}%
\pgfpathlineto{\pgfqpoint{2.367379in}{0.699520in}}%
\pgfpathlineto{\pgfqpoint{2.367712in}{0.674773in}}%
\pgfpathlineto{\pgfqpoint{2.368386in}{0.674773in}}%
\pgfpathlineto{\pgfqpoint{2.369556in}{0.773763in}}%
\pgfpathlineto{\pgfqpoint{2.370563in}{0.740766in}}%
\pgfpathlineto{\pgfqpoint{2.370933in}{0.732517in}}%
\pgfpathlineto{\pgfqpoint{2.371400in}{0.773763in}}%
\pgfpathlineto{\pgfqpoint{2.373917in}{0.897500in}}%
\pgfpathlineto{\pgfqpoint{2.374398in}{0.881002in}}%
\pgfpathlineto{\pgfqpoint{2.375576in}{0.856254in}}%
\pgfpathlineto{\pgfqpoint{2.376597in}{0.946995in}}%
\pgfpathlineto{\pgfqpoint{2.376745in}{0.955244in}}%
\pgfpathlineto{\pgfqpoint{2.377079in}{0.930497in}}%
\pgfpathlineto{\pgfqpoint{2.377271in}{0.938746in}}%
\pgfpathlineto{\pgfqpoint{2.378108in}{0.922247in}}%
\pgfpathlineto{\pgfqpoint{2.377434in}{0.946995in}}%
\pgfpathlineto{\pgfqpoint{2.378256in}{0.938746in}}%
\pgfpathlineto{\pgfqpoint{2.379093in}{0.955244in}}%
\pgfpathlineto{\pgfqpoint{2.378589in}{0.930497in}}%
\pgfpathlineto{\pgfqpoint{2.379278in}{0.938746in}}%
\pgfpathlineto{\pgfqpoint{2.381943in}{0.831507in}}%
\pgfpathlineto{\pgfqpoint{2.381965in}{0.848005in}}%
\pgfpathlineto{\pgfqpoint{2.383639in}{0.946995in}}%
\pgfpathlineto{\pgfqpoint{2.384461in}{0.930497in}}%
\pgfpathlineto{\pgfqpoint{2.384979in}{0.922247in}}%
\pgfpathlineto{\pgfqpoint{2.385297in}{0.955244in}}%
\pgfpathlineto{\pgfqpoint{2.385482in}{0.946995in}}%
\pgfpathlineto{\pgfqpoint{2.385964in}{0.955244in}}%
\pgfpathlineto{\pgfqpoint{2.386156in}{0.938746in}}%
\pgfpathlineto{\pgfqpoint{2.386637in}{0.930497in}}%
\pgfpathlineto{\pgfqpoint{2.386660in}{0.971742in}}%
\pgfpathlineto{\pgfqpoint{2.386971in}{0.955244in}}%
\pgfpathlineto{\pgfqpoint{2.386993in}{0.971742in}}%
\pgfpathlineto{\pgfqpoint{2.388000in}{0.971742in}}%
\pgfpathlineto{\pgfqpoint{2.388985in}{0.955244in}}%
\pgfpathlineto{\pgfqpoint{2.388999in}{0.971742in}}%
\pgfpathlineto{\pgfqpoint{2.389488in}{1.004739in}}%
\pgfpathlineto{\pgfqpoint{2.390154in}{0.979992in}}%
\pgfpathlineto{\pgfqpoint{2.390347in}{0.971742in}}%
\pgfpathlineto{\pgfqpoint{2.390998in}{1.004739in}}%
\pgfpathlineto{\pgfqpoint{2.391184in}{0.996490in}}%
\pgfpathlineto{\pgfqpoint{2.393175in}{1.054234in}}%
\pgfpathlineto{\pgfqpoint{2.394197in}{1.045985in}}%
\pgfpathlineto{\pgfqpoint{2.394686in}{1.054234in}}%
\pgfpathlineto{\pgfqpoint{2.394871in}{1.037736in}}%
\pgfpathlineto{\pgfqpoint{2.396715in}{0.922247in}}%
\pgfpathlineto{\pgfqpoint{2.397699in}{0.979992in}}%
\pgfpathlineto{\pgfqpoint{2.398055in}{0.946995in}}%
\pgfpathlineto{\pgfqpoint{2.399039in}{0.930497in}}%
\pgfpathlineto{\pgfqpoint{2.399062in}{0.946995in}}%
\pgfpathlineto{\pgfqpoint{2.400046in}{0.979992in}}%
\pgfpathlineto{\pgfqpoint{2.400232in}{0.963493in}}%
\pgfpathlineto{\pgfqpoint{2.401409in}{0.946995in}}%
\pgfpathlineto{\pgfqpoint{2.401557in}{0.955244in}}%
\pgfpathlineto{\pgfqpoint{2.401890in}{0.930497in}}%
\pgfpathlineto{\pgfqpoint{2.402245in}{0.946995in}}%
\pgfpathlineto{\pgfqpoint{2.402564in}{0.930497in}}%
\pgfpathlineto{\pgfqpoint{2.402897in}{0.963493in}}%
\pgfpathlineto{\pgfqpoint{2.403904in}{0.979992in}}%
\pgfpathlineto{\pgfqpoint{2.403423in}{0.946995in}}%
\pgfpathlineto{\pgfqpoint{2.403919in}{0.971742in}}%
\pgfpathlineto{\pgfqpoint{2.405577in}{1.029487in}}%
\pgfpathlineto{\pgfqpoint{2.405762in}{1.012988in}}%
\pgfpathlineto{\pgfqpoint{2.405918in}{1.004739in}}%
\pgfpathlineto{\pgfqpoint{2.406081in}{1.029487in}}%
\pgfpathlineto{\pgfqpoint{2.406436in}{1.021237in}}%
\pgfpathlineto{\pgfqpoint{2.407258in}{1.029487in}}%
\pgfpathlineto{\pgfqpoint{2.407273in}{1.021237in}}%
\pgfpathlineto{\pgfqpoint{2.407813in}{1.004739in}}%
\pgfpathlineto{\pgfqpoint{2.407924in}{1.037736in}}%
\pgfpathlineto{\pgfqpoint{2.408931in}{1.054234in}}%
\pgfpathlineto{\pgfqpoint{2.408154in}{1.029487in}}%
\pgfpathlineto{\pgfqpoint{2.408954in}{1.045985in}}%
\pgfpathlineto{\pgfqpoint{2.409768in}{1.029487in}}%
\pgfpathlineto{\pgfqpoint{2.410109in}{1.054234in}}%
\pgfpathlineto{\pgfqpoint{2.411131in}{1.021237in}}%
\pgfpathlineto{\pgfqpoint{2.412456in}{1.004739in}}%
\pgfpathlineto{\pgfqpoint{2.412952in}{1.054234in}}%
\pgfpathlineto{\pgfqpoint{2.413478in}{1.021237in}}%
\pgfpathlineto{\pgfqpoint{2.415980in}{1.103729in}}%
\pgfpathlineto{\pgfqpoint{2.417335in}{0.938746in}}%
\pgfpathlineto{\pgfqpoint{2.419342in}{0.699520in}}%
\pgfpathlineto{\pgfqpoint{2.419682in}{0.724268in}}%
\pgfpathlineto{\pgfqpoint{2.420001in}{0.732517in}}%
\pgfpathlineto{\pgfqpoint{2.420186in}{0.716019in}}%
\pgfpathlineto{\pgfqpoint{2.420497in}{0.683022in}}%
\pgfpathlineto{\pgfqpoint{2.421171in}{0.732517in}}%
\pgfpathlineto{\pgfqpoint{2.421185in}{0.724268in}}%
\pgfpathlineto{\pgfqpoint{2.421504in}{0.732517in}}%
\pgfpathlineto{\pgfqpoint{2.421837in}{0.707769in}}%
\pgfpathlineto{\pgfqpoint{2.422363in}{0.724268in}}%
\pgfpathlineto{\pgfqpoint{2.422866in}{0.699520in}}%
\pgfpathlineto{\pgfqpoint{2.423851in}{0.683022in}}%
\pgfpathlineto{\pgfqpoint{2.423851in}{0.691271in}}%
\pgfpathlineto{\pgfqpoint{2.423873in}{0.724268in}}%
\pgfpathlineto{\pgfqpoint{2.424873in}{0.724268in}}%
\pgfpathlineto{\pgfqpoint{2.426887in}{0.823258in}}%
\pgfpathlineto{\pgfqpoint{2.427368in}{0.806759in}}%
\pgfpathlineto{\pgfqpoint{2.428042in}{0.782012in}}%
\pgfpathlineto{\pgfqpoint{2.428397in}{0.798510in}}%
\pgfpathlineto{\pgfqpoint{2.429049in}{0.856254in}}%
\pgfpathlineto{\pgfqpoint{2.430056in}{0.831507in}}%
\pgfpathlineto{\pgfqpoint{2.430241in}{0.823258in}}%
\pgfpathlineto{\pgfqpoint{2.430892in}{0.881002in}}%
\pgfpathlineto{\pgfqpoint{2.431077in}{0.872753in}}%
\pgfpathlineto{\pgfqpoint{2.432743in}{0.905749in}}%
\pgfpathlineto{\pgfqpoint{2.432921in}{0.897500in}}%
\pgfpathlineto{\pgfqpoint{2.433239in}{0.930497in}}%
\pgfpathlineto{\pgfqpoint{2.433758in}{0.922247in}}%
\pgfpathlineto{\pgfqpoint{2.434432in}{0.897500in}}%
\pgfpathlineto{\pgfqpoint{2.434920in}{0.913998in}}%
\pgfpathlineto{\pgfqpoint{2.435749in}{0.930497in}}%
\pgfpathlineto{\pgfqpoint{2.435268in}{0.897500in}}%
\pgfpathlineto{\pgfqpoint{2.435942in}{0.913998in}}%
\pgfpathlineto{\pgfqpoint{2.437282in}{0.872753in}}%
\pgfpathlineto{\pgfqpoint{2.438097in}{0.930497in}}%
\pgfpathlineto{\pgfqpoint{2.438452in}{0.889251in}}%
\pgfpathlineto{\pgfqpoint{2.439289in}{0.823258in}}%
\pgfpathlineto{\pgfqpoint{2.439792in}{0.848005in}}%
\pgfpathlineto{\pgfqpoint{2.441140in}{0.897500in}}%
\pgfpathlineto{\pgfqpoint{2.441451in}{0.881002in}}%
\pgfpathlineto{\pgfqpoint{2.441636in}{0.872753in}}%
\pgfpathlineto{\pgfqpoint{2.441969in}{0.905749in}}%
\pgfpathlineto{\pgfqpoint{2.442480in}{0.897500in}}%
\pgfpathlineto{\pgfqpoint{2.444131in}{0.930497in}}%
\pgfpathlineto{\pgfqpoint{2.444472in}{0.905749in}}%
\pgfpathlineto{\pgfqpoint{2.444827in}{0.946995in}}%
\pgfpathlineto{\pgfqpoint{2.445160in}{0.922247in}}%
\pgfpathlineto{\pgfqpoint{2.446834in}{1.021237in}}%
\pgfpathlineto{\pgfqpoint{2.447663in}{1.004739in}}%
\pgfpathlineto{\pgfqpoint{2.448322in}{0.979992in}}%
\pgfpathlineto{\pgfqpoint{2.448677in}{1.021237in}}%
\pgfpathlineto{\pgfqpoint{2.450010in}{1.054234in}}%
\pgfpathlineto{\pgfqpoint{2.450017in}{1.045985in}}%
\pgfpathlineto{\pgfqpoint{2.450669in}{1.029487in}}%
\pgfpathlineto{\pgfqpoint{2.450839in}{1.054234in}}%
\pgfpathlineto{\pgfqpoint{2.450861in}{1.045985in}}%
\pgfpathlineto{\pgfqpoint{2.451172in}{1.054234in}}%
\pgfpathlineto{\pgfqpoint{2.451365in}{1.037736in}}%
\pgfpathlineto{\pgfqpoint{2.454208in}{0.724268in}}%
\pgfpathlineto{\pgfqpoint{2.454364in}{0.740766in}}%
\pgfpathlineto{\pgfqpoint{2.454860in}{0.757264in}}%
\pgfpathlineto{\pgfqpoint{2.455200in}{0.732517in}}%
\pgfpathlineto{\pgfqpoint{2.455385in}{0.740766in}}%
\pgfpathlineto{\pgfqpoint{2.456370in}{0.707769in}}%
\pgfpathlineto{\pgfqpoint{2.456378in}{0.740766in}}%
\pgfpathlineto{\pgfqpoint{2.456703in}{0.732517in}}%
\pgfpathlineto{\pgfqpoint{2.457399in}{0.749015in}}%
\pgfpathlineto{\pgfqpoint{2.458384in}{0.707769in}}%
\pgfpathlineto{\pgfqpoint{2.458739in}{0.724268in}}%
\pgfpathlineto{\pgfqpoint{2.459228in}{0.732517in}}%
\pgfpathlineto{\pgfqpoint{2.459406in}{0.716019in}}%
\pgfpathlineto{\pgfqpoint{2.460080in}{0.699520in}}%
\pgfpathlineto{\pgfqpoint{2.460243in}{0.724268in}}%
\pgfpathlineto{\pgfqpoint{2.460909in}{0.732517in}}%
\pgfpathlineto{\pgfqpoint{2.461087in}{0.716019in}}%
\pgfpathlineto{\pgfqpoint{2.461420in}{0.699520in}}%
\pgfpathlineto{\pgfqpoint{2.461738in}{0.732517in}}%
\pgfpathlineto{\pgfqpoint{2.462079in}{0.716019in}}%
\pgfpathlineto{\pgfqpoint{2.463256in}{0.757264in}}%
\pgfpathlineto{\pgfqpoint{2.464270in}{0.749015in}}%
\pgfpathlineto{\pgfqpoint{2.464752in}{0.757264in}}%
\pgfpathlineto{\pgfqpoint{2.465085in}{0.732517in}}%
\pgfpathlineto{\pgfqpoint{2.465107in}{0.749015in}}%
\pgfpathlineto{\pgfqpoint{2.465611in}{0.724268in}}%
\pgfpathlineto{\pgfqpoint{2.465929in}{0.757264in}}%
\pgfpathlineto{\pgfqpoint{2.465944in}{0.749015in}}%
\pgfpathlineto{\pgfqpoint{2.466936in}{0.757264in}}%
\pgfpathlineto{\pgfqpoint{2.466951in}{0.749015in}}%
\pgfpathlineto{\pgfqpoint{2.469276in}{0.633527in}}%
\pgfpathlineto{\pgfqpoint{2.469624in}{0.666524in}}%
\pgfpathlineto{\pgfqpoint{2.470623in}{0.707769in}}%
\pgfpathlineto{\pgfqpoint{2.470638in}{0.699520in}}%
\pgfpathlineto{\pgfqpoint{2.472148in}{0.749015in}}%
\pgfpathlineto{\pgfqpoint{2.472319in}{0.740766in}}%
\pgfpathlineto{\pgfqpoint{2.472800in}{0.732517in}}%
\pgfpathlineto{\pgfqpoint{2.472859in}{0.757264in}}%
\pgfpathlineto{\pgfqpoint{2.472985in}{0.749015in}}%
\pgfpathlineto{\pgfqpoint{2.475836in}{0.872753in}}%
\pgfpathlineto{\pgfqpoint{2.476006in}{0.864503in}}%
\pgfpathlineto{\pgfqpoint{2.476487in}{0.856254in}}%
\pgfpathlineto{\pgfqpoint{2.476658in}{0.881002in}}%
\pgfpathlineto{\pgfqpoint{2.476843in}{0.872753in}}%
\pgfpathlineto{\pgfqpoint{2.478013in}{0.897500in}}%
\pgfpathlineto{\pgfqpoint{2.478183in}{0.889251in}}%
\pgfpathlineto{\pgfqpoint{2.478664in}{0.881002in}}%
\pgfpathlineto{\pgfqpoint{2.478834in}{0.905749in}}%
\pgfpathlineto{\pgfqpoint{2.478857in}{0.897500in}}%
\pgfpathlineto{\pgfqpoint{2.481515in}{0.979992in}}%
\pgfpathlineto{\pgfqpoint{2.482040in}{0.946995in}}%
\pgfpathlineto{\pgfqpoint{2.482544in}{0.996490in}}%
\pgfpathlineto{\pgfqpoint{2.483884in}{0.971742in}}%
\pgfpathlineto{\pgfqpoint{2.485387in}{1.021237in}}%
\pgfpathlineto{\pgfqpoint{2.485557in}{1.012988in}}%
\pgfpathlineto{\pgfqpoint{2.485706in}{1.004739in}}%
\pgfpathlineto{\pgfqpoint{2.485728in}{1.045985in}}%
\pgfpathlineto{\pgfqpoint{2.486061in}{1.045985in}}%
\pgfpathlineto{\pgfqpoint{2.487216in}{1.054234in}}%
\pgfpathlineto{\pgfqpoint{2.487734in}{1.045985in}}%
\pgfpathlineto{\pgfqpoint{2.488053in}{1.078981in}}%
\pgfpathlineto{\pgfqpoint{2.488238in}{1.070732in}}%
\pgfpathlineto{\pgfqpoint{2.489415in}{1.120227in}}%
\pgfpathlineto{\pgfqpoint{2.489896in}{1.103729in}}%
\pgfpathlineto{\pgfqpoint{2.490230in}{1.128476in}}%
\pgfpathlineto{\pgfqpoint{2.490926in}{1.095480in}}%
\pgfpathlineto{\pgfqpoint{2.491407in}{1.103729in}}%
\pgfpathlineto{\pgfqpoint{2.491592in}{1.087231in}}%
\pgfpathlineto{\pgfqpoint{2.491762in}{1.095480in}}%
\pgfpathlineto{\pgfqpoint{2.492769in}{1.070732in}}%
\pgfpathlineto{\pgfqpoint{2.492932in}{1.095480in}}%
\pgfpathlineto{\pgfqpoint{2.493606in}{1.070732in}}%
\pgfpathlineto{\pgfqpoint{2.494272in}{0.988241in}}%
\pgfpathlineto{\pgfqpoint{2.497626in}{0.650025in}}%
\pgfpathlineto{\pgfqpoint{2.498789in}{0.683022in}}%
\pgfpathlineto{\pgfqpoint{2.498804in}{0.674773in}}%
\pgfpathlineto{\pgfqpoint{2.499285in}{0.683022in}}%
\pgfpathlineto{\pgfqpoint{2.499470in}{0.666524in}}%
\pgfpathlineto{\pgfqpoint{2.500795in}{0.625278in}}%
\pgfpathlineto{\pgfqpoint{2.501965in}{0.633527in}}%
\pgfpathlineto{\pgfqpoint{2.502135in}{0.625278in}}%
\pgfpathlineto{\pgfqpoint{2.502972in}{0.625278in}}%
\pgfpathlineto{\pgfqpoint{2.504497in}{0.699520in}}%
\pgfpathlineto{\pgfqpoint{2.505149in}{0.683022in}}%
\pgfpathlineto{\pgfqpoint{2.506178in}{0.674773in}}%
\pgfpathlineto{\pgfqpoint{2.506830in}{0.683022in}}%
\pgfpathlineto{\pgfqpoint{2.507015in}{0.666524in}}%
\pgfpathlineto{\pgfqpoint{2.508503in}{0.633527in}}%
\pgfpathlineto{\pgfqpoint{2.508503in}{0.641776in}}%
\pgfpathlineto{\pgfqpoint{2.508844in}{0.633527in}}%
\pgfpathlineto{\pgfqpoint{2.509532in}{0.650025in}}%
\pgfpathlineto{\pgfqpoint{2.509843in}{0.633527in}}%
\pgfpathlineto{\pgfqpoint{2.510184in}{0.666524in}}%
\pgfpathlineto{\pgfqpoint{2.510199in}{0.674773in}}%
\pgfpathlineto{\pgfqpoint{2.510872in}{0.625278in}}%
\pgfpathlineto{\pgfqpoint{2.511206in}{0.625278in}}%
\pgfpathlineto{\pgfqpoint{2.512531in}{0.633527in}}%
\pgfpathlineto{\pgfqpoint{2.512538in}{0.625278in}}%
\pgfpathlineto{\pgfqpoint{2.513553in}{0.674773in}}%
\pgfpathlineto{\pgfqpoint{2.513760in}{0.658274in}}%
\pgfpathlineto{\pgfqpoint{2.513886in}{0.699520in}}%
\pgfpathlineto{\pgfqpoint{2.514204in}{0.691271in}}%
\pgfpathlineto{\pgfqpoint{2.514382in}{0.732517in}}%
\pgfpathlineto{\pgfqpoint{2.515226in}{0.699520in}}%
\pgfpathlineto{\pgfqpoint{2.517070in}{0.773763in}}%
\pgfpathlineto{\pgfqpoint{2.517240in}{0.765513in}}%
\pgfpathlineto{\pgfqpoint{2.517410in}{0.773763in}}%
\pgfpathlineto{\pgfqpoint{2.518395in}{0.757264in}}%
\pgfpathlineto{\pgfqpoint{2.519409in}{0.782012in}}%
\pgfpathlineto{\pgfqpoint{2.519417in}{0.773763in}}%
\pgfpathlineto{\pgfqpoint{2.520764in}{0.848005in}}%
\pgfpathlineto{\pgfqpoint{2.521261in}{0.823258in}}%
\pgfpathlineto{\pgfqpoint{2.522105in}{0.798510in}}%
\pgfpathlineto{\pgfqpoint{2.522267in}{0.823258in}}%
\pgfpathlineto{\pgfqpoint{2.523778in}{0.872753in}}%
\pgfpathlineto{\pgfqpoint{2.523948in}{0.864503in}}%
\pgfpathlineto{\pgfqpoint{2.524111in}{0.897500in}}%
\pgfpathlineto{\pgfqpoint{2.525103in}{0.856254in}}%
\pgfpathlineto{\pgfqpoint{2.525118in}{0.872753in}}%
\pgfpathlineto{\pgfqpoint{2.526125in}{0.872753in}}%
\pgfpathlineto{\pgfqpoint{2.527450in}{0.856254in}}%
\pgfpathlineto{\pgfqpoint{2.527465in}{0.872753in}}%
\pgfpathlineto{\pgfqpoint{2.528472in}{0.872753in}}%
\pgfpathlineto{\pgfqpoint{2.529124in}{0.856254in}}%
\pgfpathlineto{\pgfqpoint{2.529301in}{0.881002in}}%
\pgfpathlineto{\pgfqpoint{2.529309in}{0.872753in}}%
\pgfpathlineto{\pgfqpoint{2.531138in}{0.930497in}}%
\pgfpathlineto{\pgfqpoint{2.531153in}{0.922247in}}%
\pgfpathlineto{\pgfqpoint{2.531989in}{0.897500in}}%
\pgfpathlineto{\pgfqpoint{2.532159in}{0.922247in}}%
\pgfpathlineto{\pgfqpoint{2.536002in}{1.078981in}}%
\pgfpathlineto{\pgfqpoint{2.536506in}{1.054234in}}%
\pgfpathlineto{\pgfqpoint{2.537024in}{1.070732in}}%
\pgfpathlineto{\pgfqpoint{2.538364in}{1.045985in}}%
\pgfpathlineto{\pgfqpoint{2.539519in}{1.078981in}}%
\pgfpathlineto{\pgfqpoint{2.539704in}{1.062483in}}%
\pgfpathlineto{\pgfqpoint{2.540371in}{1.021237in}}%
\pgfpathlineto{\pgfqpoint{2.540859in}{1.054234in}}%
\pgfpathlineto{\pgfqpoint{2.541548in}{1.095480in}}%
\pgfpathlineto{\pgfqpoint{2.541881in}{1.095480in}}%
\pgfpathlineto{\pgfqpoint{2.542718in}{1.045985in}}%
\pgfpathlineto{\pgfqpoint{2.543207in}{1.078981in}}%
\pgfpathlineto{\pgfqpoint{2.543369in}{1.103729in}}%
\pgfpathlineto{\pgfqpoint{2.544228in}{1.087231in}}%
\pgfpathlineto{\pgfqpoint{2.544902in}{1.070732in}}%
\pgfpathlineto{\pgfqpoint{2.544399in}{1.095480in}}%
\pgfpathlineto{\pgfqpoint{2.545065in}{1.095480in}}%
\pgfpathlineto{\pgfqpoint{2.545406in}{1.120227in}}%
\pgfpathlineto{\pgfqpoint{2.546220in}{1.103729in}}%
\pgfpathlineto{\pgfqpoint{2.547249in}{1.095480in}}%
\pgfpathlineto{\pgfqpoint{2.548404in}{1.103729in}}%
\pgfpathlineto{\pgfqpoint{2.549426in}{1.095480in}}%
\pgfpathlineto{\pgfqpoint{2.550411in}{1.128476in}}%
\pgfpathlineto{\pgfqpoint{2.550596in}{1.111978in}}%
\pgfpathlineto{\pgfqpoint{2.553950in}{0.650025in}}%
\pgfpathlineto{\pgfqpoint{2.554772in}{0.666524in}}%
\pgfpathlineto{\pgfqpoint{2.555127in}{0.699520in}}%
\pgfpathlineto{\pgfqpoint{2.555794in}{0.650025in}}%
\pgfpathlineto{\pgfqpoint{2.557956in}{0.732517in}}%
\pgfpathlineto{\pgfqpoint{2.557978in}{0.724268in}}%
\pgfpathlineto{\pgfqpoint{2.559651in}{0.650025in}}%
\pgfpathlineto{\pgfqpoint{2.560133in}{0.666524in}}%
\pgfpathlineto{\pgfqpoint{2.560473in}{0.658274in}}%
\pgfpathlineto{\pgfqpoint{2.561162in}{0.699520in}}%
\pgfpathlineto{\pgfqpoint{2.562146in}{0.683022in}}%
\pgfpathlineto{\pgfqpoint{2.561495in}{0.724268in}}%
\pgfpathlineto{\pgfqpoint{2.562169in}{0.699520in}}%
\pgfpathlineto{\pgfqpoint{2.562650in}{0.707769in}}%
\pgfpathlineto{\pgfqpoint{2.562983in}{0.683022in}}%
\pgfpathlineto{\pgfqpoint{2.563168in}{0.674773in}}%
\pgfpathlineto{\pgfqpoint{2.563487in}{0.707769in}}%
\pgfpathlineto{\pgfqpoint{2.564012in}{0.699520in}}%
\pgfpathlineto{\pgfqpoint{2.565663in}{0.732517in}}%
\pgfpathlineto{\pgfqpoint{2.566670in}{0.707769in}}%
\pgfpathlineto{\pgfqpoint{2.566693in}{0.724268in}}%
\pgfpathlineto{\pgfqpoint{2.567011in}{0.707769in}}%
\pgfpathlineto{\pgfqpoint{2.567026in}{0.749015in}}%
\pgfpathlineto{\pgfqpoint{2.567344in}{0.740766in}}%
\pgfpathlineto{\pgfqpoint{2.567359in}{0.798510in}}%
\pgfpathlineto{\pgfqpoint{2.568366in}{0.773763in}}%
\pgfpathlineto{\pgfqpoint{2.569521in}{0.782012in}}%
\pgfpathlineto{\pgfqpoint{2.570550in}{0.773763in}}%
\pgfpathlineto{\pgfqpoint{2.571705in}{0.782012in}}%
\pgfpathlineto{\pgfqpoint{2.572727in}{0.773763in}}%
\pgfpathlineto{\pgfqpoint{2.572764in}{0.782012in}}%
\pgfpathlineto{\pgfqpoint{2.575393in}{0.905749in}}%
\pgfpathlineto{\pgfqpoint{2.575407in}{0.897500in}}%
\pgfpathlineto{\pgfqpoint{2.576244in}{0.872753in}}%
\pgfpathlineto{\pgfqpoint{2.576733in}{0.881002in}}%
\pgfpathlineto{\pgfqpoint{2.577755in}{0.946995in}}%
\pgfpathlineto{\pgfqpoint{2.579080in}{0.930497in}}%
\pgfpathlineto{\pgfqpoint{2.580087in}{0.979992in}}%
\pgfpathlineto{\pgfqpoint{2.579435in}{0.922247in}}%
\pgfpathlineto{\pgfqpoint{2.580102in}{0.971742in}}%
\pgfpathlineto{\pgfqpoint{2.580427in}{0.979992in}}%
\pgfpathlineto{\pgfqpoint{2.580753in}{0.955244in}}%
\pgfpathlineto{\pgfqpoint{2.581760in}{0.930497in}}%
\pgfpathlineto{\pgfqpoint{2.581109in}{0.971742in}}%
\pgfpathlineto{\pgfqpoint{2.581775in}{0.946995in}}%
\pgfpathlineto{\pgfqpoint{2.582271in}{0.979992in}}%
\pgfpathlineto{\pgfqpoint{2.582930in}{0.955244in}}%
\pgfpathlineto{\pgfqpoint{2.583789in}{0.946995in}}%
\pgfpathlineto{\pgfqpoint{2.583271in}{0.979992in}}%
\pgfpathlineto{\pgfqpoint{2.583959in}{0.971742in}}%
\pgfpathlineto{\pgfqpoint{2.586625in}{1.054234in}}%
\pgfpathlineto{\pgfqpoint{2.586810in}{1.045985in}}%
\pgfpathlineto{\pgfqpoint{2.587632in}{1.111978in}}%
\pgfpathlineto{\pgfqpoint{2.588809in}{1.227466in}}%
\pgfpathlineto{\pgfqpoint{2.588816in}{1.219217in}}%
\pgfpathlineto{\pgfqpoint{2.589979in}{1.227466in}}%
\pgfpathlineto{\pgfqpoint{2.591334in}{1.037736in}}%
\pgfpathlineto{\pgfqpoint{2.594347in}{0.650025in}}%
\pgfpathlineto{\pgfqpoint{2.594851in}{0.674773in}}%
\pgfpathlineto{\pgfqpoint{2.595177in}{0.683022in}}%
\pgfpathlineto{\pgfqpoint{2.595510in}{0.658274in}}%
\pgfpathlineto{\pgfqpoint{2.595695in}{0.650025in}}%
\pgfpathlineto{\pgfqpoint{2.596198in}{0.724268in}}%
\pgfpathlineto{\pgfqpoint{2.596509in}{0.716019in}}%
\pgfpathlineto{\pgfqpoint{2.597524in}{0.757264in}}%
\pgfpathlineto{\pgfqpoint{2.597539in}{0.749015in}}%
\pgfpathlineto{\pgfqpoint{2.598531in}{0.732517in}}%
\pgfpathlineto{\pgfqpoint{2.598538in}{0.749015in}}%
\pgfpathlineto{\pgfqpoint{2.598694in}{0.757264in}}%
\pgfpathlineto{\pgfqpoint{2.599042in}{0.724268in}}%
\pgfpathlineto{\pgfqpoint{2.601877in}{0.625278in}}%
\pgfpathlineto{\pgfqpoint{2.603047in}{0.633527in}}%
\pgfpathlineto{\pgfqpoint{2.603380in}{0.625278in}}%
\pgfpathlineto{\pgfqpoint{2.604054in}{0.625278in}}%
\pgfpathlineto{\pgfqpoint{2.605394in}{0.633527in}}%
\pgfpathlineto{\pgfqpoint{2.605728in}{0.625278in}}%
\pgfpathlineto{\pgfqpoint{2.606076in}{0.625278in}}%
\pgfpathlineto{\pgfqpoint{2.607238in}{0.633527in}}%
\pgfpathlineto{\pgfqpoint{2.607579in}{0.625278in}}%
\pgfpathlineto{\pgfqpoint{2.608252in}{0.625278in}}%
\pgfpathlineto{\pgfqpoint{2.609422in}{0.633527in}}%
\pgfpathlineto{\pgfqpoint{2.609593in}{0.625278in}}%
\pgfpathlineto{\pgfqpoint{2.610422in}{0.625278in}}%
\pgfpathlineto{\pgfqpoint{2.611599in}{0.633527in}}%
\pgfpathlineto{\pgfqpoint{2.611762in}{0.625278in}}%
\pgfpathlineto{\pgfqpoint{2.612613in}{0.625278in}}%
\pgfpathlineto{\pgfqpoint{2.613613in}{0.633527in}}%
\pgfpathlineto{\pgfqpoint{2.613613in}{0.625278in}}%
\pgfpathlineto{\pgfqpoint{2.614783in}{0.633527in}}%
\pgfpathlineto{\pgfqpoint{2.614961in}{0.625278in}}%
\pgfpathlineto{\pgfqpoint{2.615790in}{0.625278in}}%
\pgfpathlineto{\pgfqpoint{2.616975in}{0.633527in}}%
\pgfpathlineto{\pgfqpoint{2.617019in}{0.625278in}}%
\pgfpathlineto{\pgfqpoint{2.617981in}{0.625278in}}%
\pgfpathlineto{\pgfqpoint{2.619144in}{0.633527in}}%
\pgfpathlineto{\pgfqpoint{2.619144in}{0.625278in}}%
\pgfpathlineto{\pgfqpoint{2.620144in}{0.625278in}}%
\pgfpathlineto{\pgfqpoint{2.621150in}{0.633527in}}%
\pgfpathlineto{\pgfqpoint{2.621150in}{0.625278in}}%
\pgfpathlineto{\pgfqpoint{2.622491in}{0.633527in}}%
\pgfpathlineto{\pgfqpoint{2.622491in}{0.625278in}}%
\pgfpathlineto{\pgfqpoint{2.623498in}{0.625278in}}%
\pgfpathlineto{\pgfqpoint{2.624845in}{0.633527in}}%
\pgfpathlineto{\pgfqpoint{2.625008in}{0.625278in}}%
\pgfpathlineto{\pgfqpoint{2.625852in}{0.625278in}}%
\pgfpathlineto{\pgfqpoint{2.627022in}{0.633527in}}%
\pgfpathlineto{\pgfqpoint{2.627192in}{0.625278in}}%
\pgfpathlineto{\pgfqpoint{2.628029in}{0.625278in}}%
\pgfpathlineto{\pgfqpoint{2.629036in}{0.633527in}}%
\pgfpathlineto{\pgfqpoint{2.629036in}{0.625278in}}%
\pgfpathlineto{\pgfqpoint{2.630376in}{0.633527in}}%
\pgfpathlineto{\pgfqpoint{2.630376in}{0.625278in}}%
\pgfpathlineto{\pgfqpoint{2.631383in}{0.625278in}}%
\pgfpathlineto{\pgfqpoint{2.632560in}{0.633527in}}%
\pgfpathlineto{\pgfqpoint{2.632716in}{0.625278in}}%
\pgfpathlineto{\pgfqpoint{2.633390in}{0.625278in}}%
\pgfpathlineto{\pgfqpoint{2.634567in}{0.633527in}}%
\pgfpathlineto{\pgfqpoint{2.634730in}{0.625278in}}%
\pgfpathlineto{\pgfqpoint{2.635574in}{0.625278in}}%
\pgfpathlineto{\pgfqpoint{2.636751in}{0.633527in}}%
\pgfpathlineto{\pgfqpoint{2.636907in}{0.625278in}}%
\pgfpathlineto{\pgfqpoint{2.637751in}{0.625278in}}%
\pgfpathlineto{\pgfqpoint{2.638921in}{0.633527in}}%
\pgfpathlineto{\pgfqpoint{2.639091in}{0.625278in}}%
\pgfpathlineto{\pgfqpoint{2.639928in}{0.625278in}}%
\pgfpathlineto{\pgfqpoint{2.641097in}{0.633527in}}%
\pgfpathlineto{\pgfqpoint{2.641268in}{0.625278in}}%
\pgfpathlineto{\pgfqpoint{2.641941in}{0.625278in}}%
\pgfpathlineto{\pgfqpoint{2.643119in}{0.633527in}}%
\pgfpathlineto{\pgfqpoint{2.643282in}{0.625278in}}%
\pgfpathlineto{\pgfqpoint{2.644126in}{0.625278in}}%
\pgfpathlineto{\pgfqpoint{2.645133in}{0.633527in}}%
\pgfpathlineto{\pgfqpoint{2.645133in}{0.625278in}}%
\pgfpathlineto{\pgfqpoint{2.646295in}{0.633527in}}%
\pgfpathlineto{\pgfqpoint{2.646465in}{0.625278in}}%
\pgfpathlineto{\pgfqpoint{2.647310in}{0.625278in}}%
\pgfpathlineto{\pgfqpoint{2.648479in}{0.633527in}}%
\pgfpathlineto{\pgfqpoint{2.648642in}{0.625278in}}%
\pgfpathlineto{\pgfqpoint{2.649479in}{0.625278in}}%
\pgfpathlineto{\pgfqpoint{2.650656in}{0.633527in}}%
\pgfpathlineto{\pgfqpoint{2.650827in}{0.625278in}}%
\pgfpathlineto{\pgfqpoint{2.651493in}{0.625278in}}%
\pgfpathlineto{\pgfqpoint{2.652670in}{0.633527in}}%
\pgfpathlineto{\pgfqpoint{2.652833in}{0.625278in}}%
\pgfpathlineto{\pgfqpoint{2.653670in}{0.625278in}}%
\pgfpathlineto{\pgfqpoint{2.654847in}{0.633527in}}%
\pgfpathlineto{\pgfqpoint{2.655017in}{0.625278in}}%
\pgfpathlineto{\pgfqpoint{2.655854in}{0.625278in}}%
\pgfpathlineto{\pgfqpoint{2.657024in}{0.633527in}}%
\pgfpathlineto{\pgfqpoint{2.657194in}{0.625278in}}%
\pgfpathlineto{\pgfqpoint{2.657860in}{0.625278in}}%
\pgfpathlineto{\pgfqpoint{2.659038in}{0.633527in}}%
\pgfpathlineto{\pgfqpoint{2.659208in}{0.625278in}}%
\pgfpathlineto{\pgfqpoint{2.660045in}{0.625278in}}%
\pgfpathlineto{\pgfqpoint{2.661215in}{0.633527in}}%
\pgfpathlineto{\pgfqpoint{2.661237in}{0.625278in}}%
\pgfpathlineto{\pgfqpoint{2.662051in}{0.658274in}}%
\pgfpathlineto{\pgfqpoint{2.662244in}{0.650025in}}%
\pgfpathlineto{\pgfqpoint{2.662895in}{0.683022in}}%
\pgfpathlineto{\pgfqpoint{2.663399in}{0.658274in}}%
\pgfpathlineto{\pgfqpoint{2.663732in}{0.683022in}}%
\pgfpathlineto{\pgfqpoint{2.664421in}{0.650025in}}%
\pgfpathlineto{\pgfqpoint{2.665746in}{0.683022in}}%
\pgfpathlineto{\pgfqpoint{2.665761in}{0.674773in}}%
\pgfpathlineto{\pgfqpoint{2.666435in}{0.650025in}}%
\pgfpathlineto{\pgfqpoint{2.666597in}{0.674773in}}%
\pgfpathlineto{\pgfqpoint{2.667442in}{0.699520in}}%
\pgfpathlineto{\pgfqpoint{2.667752in}{0.683022in}}%
\pgfpathlineto{\pgfqpoint{2.668108in}{0.674773in}}%
\pgfpathlineto{\pgfqpoint{2.668426in}{0.707769in}}%
\pgfpathlineto{\pgfqpoint{2.668782in}{0.699520in}}%
\pgfpathlineto{\pgfqpoint{2.669433in}{0.707769in}}%
\pgfpathlineto{\pgfqpoint{2.669618in}{0.691271in}}%
\pgfpathlineto{\pgfqpoint{2.670122in}{0.674773in}}%
\pgfpathlineto{\pgfqpoint{2.670433in}{0.707769in}}%
\pgfpathlineto{\pgfqpoint{2.670455in}{0.699520in}}%
\pgfpathlineto{\pgfqpoint{2.670773in}{0.707769in}}%
\pgfpathlineto{\pgfqpoint{2.671107in}{0.683022in}}%
\pgfpathlineto{\pgfqpoint{2.671462in}{0.674773in}}%
\pgfpathlineto{\pgfqpoint{2.672114in}{0.707769in}}%
\pgfpathlineto{\pgfqpoint{2.672128in}{0.699520in}}%
\pgfpathlineto{\pgfqpoint{2.672632in}{0.749015in}}%
\pgfpathlineto{\pgfqpoint{2.673135in}{0.716019in}}%
\pgfpathlineto{\pgfqpoint{2.674646in}{0.650025in}}%
\pgfpathlineto{\pgfqpoint{2.675801in}{0.658274in}}%
\pgfpathlineto{\pgfqpoint{2.676993in}{0.625278in}}%
\pgfpathlineto{\pgfqpoint{2.678333in}{0.650025in}}%
\pgfpathlineto{\pgfqpoint{2.678503in}{0.641776in}}%
\pgfpathlineto{\pgfqpoint{2.678651in}{0.633527in}}%
\pgfpathlineto{\pgfqpoint{2.679007in}{0.674773in}}%
\pgfpathlineto{\pgfqpoint{2.680162in}{0.732517in}}%
\pgfpathlineto{\pgfqpoint{2.680665in}{0.707769in}}%
\pgfpathlineto{\pgfqpoint{2.680999in}{0.732517in}}%
\pgfpathlineto{\pgfqpoint{2.681687in}{0.699520in}}%
\pgfpathlineto{\pgfqpoint{2.682339in}{0.732517in}}%
\pgfpathlineto{\pgfqpoint{2.682842in}{0.707769in}}%
\pgfpathlineto{\pgfqpoint{2.683346in}{0.732517in}}%
\pgfpathlineto{\pgfqpoint{2.683864in}{0.699520in}}%
\pgfpathlineto{\pgfqpoint{2.685693in}{0.757264in}}%
\pgfpathlineto{\pgfqpoint{2.685708in}{0.749015in}}%
\pgfpathlineto{\pgfqpoint{2.687048in}{0.724268in}}%
\pgfpathlineto{\pgfqpoint{2.687870in}{0.732517in}}%
\pgfpathlineto{\pgfqpoint{2.687892in}{0.724268in}}%
\pgfpathlineto{\pgfqpoint{2.688203in}{0.707769in}}%
\pgfpathlineto{\pgfqpoint{2.688543in}{0.740766in}}%
\pgfpathlineto{\pgfqpoint{2.688706in}{0.782012in}}%
\pgfpathlineto{\pgfqpoint{2.689565in}{0.765513in}}%
\pgfpathlineto{\pgfqpoint{2.689713in}{0.757264in}}%
\pgfpathlineto{\pgfqpoint{2.690069in}{0.798510in}}%
\pgfpathlineto{\pgfqpoint{2.690387in}{0.806759in}}%
\pgfpathlineto{\pgfqpoint{2.690720in}{0.782012in}}%
\pgfpathlineto{\pgfqpoint{2.690905in}{0.790261in}}%
\pgfpathlineto{\pgfqpoint{2.691224in}{0.806759in}}%
\pgfpathlineto{\pgfqpoint{2.691912in}{0.773763in}}%
\pgfpathlineto{\pgfqpoint{2.692394in}{0.806759in}}%
\pgfpathlineto{\pgfqpoint{2.693067in}{0.782012in}}%
\pgfpathlineto{\pgfqpoint{2.694089in}{0.765513in}}%
\pgfpathlineto{\pgfqpoint{2.694578in}{0.757264in}}%
\pgfpathlineto{\pgfqpoint{2.694593in}{0.798510in}}%
\pgfpathlineto{\pgfqpoint{2.695933in}{0.773763in}}%
\pgfpathlineto{\pgfqpoint{2.696081in}{0.782012in}}%
\pgfpathlineto{\pgfqpoint{2.696436in}{0.749015in}}%
\pgfpathlineto{\pgfqpoint{2.697088in}{0.732517in}}%
\pgfpathlineto{\pgfqpoint{2.697258in}{0.757264in}}%
\pgfpathlineto{\pgfqpoint{2.697273in}{0.749015in}}%
\pgfpathlineto{\pgfqpoint{2.697925in}{0.757264in}}%
\pgfpathlineto{\pgfqpoint{2.698117in}{0.740766in}}%
\pgfpathlineto{\pgfqpoint{2.698598in}{0.732517in}}%
\pgfpathlineto{\pgfqpoint{2.698621in}{0.773763in}}%
\pgfpathlineto{\pgfqpoint{2.698954in}{0.749015in}}%
\pgfpathlineto{\pgfqpoint{2.699620in}{0.773763in}}%
\pgfpathlineto{\pgfqpoint{2.699961in}{0.749015in}}%
\pgfpathlineto{\pgfqpoint{2.700612in}{0.732517in}}%
\pgfpathlineto{\pgfqpoint{2.700612in}{0.765513in}}%
\pgfpathlineto{\pgfqpoint{2.700627in}{0.773763in}}%
\pgfpathlineto{\pgfqpoint{2.700968in}{0.749015in}}%
\pgfpathlineto{\pgfqpoint{2.701634in}{0.749015in}}%
\pgfpathlineto{\pgfqpoint{2.702308in}{0.724268in}}%
\pgfpathlineto{\pgfqpoint{2.702471in}{0.749015in}}%
\pgfpathlineto{\pgfqpoint{2.703626in}{0.757264in}}%
\pgfpathlineto{\pgfqpoint{2.703959in}{0.732517in}}%
\pgfpathlineto{\pgfqpoint{2.704655in}{0.749015in}}%
\pgfpathlineto{\pgfqpoint{2.705158in}{0.798510in}}%
\pgfpathlineto{\pgfqpoint{2.705810in}{0.757264in}}%
\pgfpathlineto{\pgfqpoint{2.706647in}{0.732517in}}%
\pgfpathlineto{\pgfqpoint{2.706832in}{0.740766in}}%
\pgfpathlineto{\pgfqpoint{2.707839in}{0.724268in}}%
\pgfpathlineto{\pgfqpoint{2.707002in}{0.749015in}}%
\pgfpathlineto{\pgfqpoint{2.707987in}{0.732517in}}%
\pgfpathlineto{\pgfqpoint{2.708002in}{0.749015in}}%
\pgfpathlineto{\pgfqpoint{2.708675in}{0.724268in}}%
\pgfpathlineto{\pgfqpoint{2.709009in}{0.724268in}}%
\pgfpathlineto{\pgfqpoint{2.710186in}{0.749015in}}%
\pgfpathlineto{\pgfqpoint{2.710349in}{0.740766in}}%
\pgfpathlineto{\pgfqpoint{2.710386in}{0.732517in}}%
\pgfpathlineto{\pgfqpoint{2.710852in}{0.798510in}}%
\pgfpathlineto{\pgfqpoint{2.711193in}{0.773763in}}%
\pgfpathlineto{\pgfqpoint{2.711341in}{0.782012in}}%
\pgfpathlineto{\pgfqpoint{2.711674in}{0.757264in}}%
\pgfpathlineto{\pgfqpoint{2.712007in}{0.757264in}}%
\pgfpathlineto{\pgfqpoint{2.713199in}{0.724268in}}%
\pgfpathlineto{\pgfqpoint{2.714021in}{0.757264in}}%
\pgfpathlineto{\pgfqpoint{2.714540in}{0.740766in}}%
\pgfpathlineto{\pgfqpoint{2.714577in}{0.732517in}}%
\pgfpathlineto{\pgfqpoint{2.714858in}{0.782012in}}%
\pgfpathlineto{\pgfqpoint{2.715384in}{0.773763in}}%
\pgfpathlineto{\pgfqpoint{2.716035in}{0.782012in}}%
\pgfpathlineto{\pgfqpoint{2.716220in}{0.765513in}}%
\pgfpathlineto{\pgfqpoint{2.717035in}{0.732517in}}%
\pgfpathlineto{\pgfqpoint{2.717390in}{0.749015in}}%
\pgfpathlineto{\pgfqpoint{2.718545in}{0.757264in}}%
\pgfpathlineto{\pgfqpoint{2.718567in}{0.749015in}}%
\pgfpathlineto{\pgfqpoint{2.719071in}{0.798510in}}%
\pgfpathlineto{\pgfqpoint{2.719574in}{0.773763in}}%
\pgfpathlineto{\pgfqpoint{2.721566in}{0.831507in}}%
\pgfpathlineto{\pgfqpoint{2.722255in}{0.798510in}}%
\pgfpathlineto{\pgfqpoint{2.722588in}{0.815008in}}%
\pgfpathlineto{\pgfqpoint{2.723913in}{0.782012in}}%
\pgfpathlineto{\pgfqpoint{2.723928in}{0.798510in}}%
\pgfpathlineto{\pgfqpoint{2.725083in}{0.806759in}}%
\pgfpathlineto{\pgfqpoint{2.725772in}{0.749015in}}%
\pgfpathlineto{\pgfqpoint{2.726275in}{0.773763in}}%
\pgfpathlineto{\pgfqpoint{2.726927in}{0.757264in}}%
\pgfpathlineto{\pgfqpoint{2.727097in}{0.782012in}}%
\pgfpathlineto{\pgfqpoint{2.727149in}{0.782012in}}%
\pgfpathlineto{\pgfqpoint{2.729104in}{0.856254in}}%
\pgfpathlineto{\pgfqpoint{2.729126in}{0.848005in}}%
\pgfpathlineto{\pgfqpoint{2.731140in}{0.922247in}}%
\pgfpathlineto{\pgfqpoint{2.731303in}{0.913998in}}%
\pgfpathlineto{\pgfqpoint{2.732480in}{0.897500in}}%
\pgfpathlineto{\pgfqpoint{2.733968in}{0.930497in}}%
\pgfpathlineto{\pgfqpoint{2.733990in}{0.922247in}}%
\pgfpathlineto{\pgfqpoint{2.735642in}{0.881002in}}%
\pgfpathlineto{\pgfqpoint{2.735642in}{0.889251in}}%
\pgfpathlineto{\pgfqpoint{2.736671in}{0.946995in}}%
\pgfpathlineto{\pgfqpoint{2.735982in}{0.881002in}}%
\pgfpathlineto{\pgfqpoint{2.736834in}{0.938746in}}%
\pgfpathlineto{\pgfqpoint{2.737004in}{0.946995in}}%
\pgfpathlineto{\pgfqpoint{2.738011in}{0.922247in}}%
\pgfpathlineto{\pgfqpoint{2.739521in}{0.971742in}}%
\pgfpathlineto{\pgfqpoint{2.739684in}{0.963493in}}%
\pgfpathlineto{\pgfqpoint{2.740025in}{0.946995in}}%
\pgfpathlineto{\pgfqpoint{2.740336in}{0.979992in}}%
\pgfpathlineto{\pgfqpoint{2.740358in}{0.971742in}}%
\pgfpathlineto{\pgfqpoint{2.741010in}{1.004739in}}%
\pgfpathlineto{\pgfqpoint{2.741513in}{0.979992in}}%
\pgfpathlineto{\pgfqpoint{2.741528in}{0.971742in}}%
\pgfpathlineto{\pgfqpoint{2.741846in}{1.004739in}}%
\pgfpathlineto{\pgfqpoint{2.742535in}{0.996490in}}%
\pgfpathlineto{\pgfqpoint{2.743186in}{0.979992in}}%
\pgfpathlineto{\pgfqpoint{2.743690in}{1.004739in}}%
\pgfpathlineto{\pgfqpoint{2.744193in}{0.979992in}}%
\pgfpathlineto{\pgfqpoint{2.744719in}{0.996490in}}%
\pgfpathlineto{\pgfqpoint{2.745371in}{1.004739in}}%
\pgfpathlineto{\pgfqpoint{2.745556in}{0.988241in}}%
\pgfpathlineto{\pgfqpoint{2.745719in}{0.971742in}}%
\pgfpathlineto{\pgfqpoint{2.746059in}{1.021237in}}%
\pgfpathlineto{\pgfqpoint{2.746370in}{1.029487in}}%
\pgfpathlineto{\pgfqpoint{2.746711in}{1.004739in}}%
\pgfpathlineto{\pgfqpoint{2.746896in}{1.012988in}}%
\pgfpathlineto{\pgfqpoint{2.747881in}{1.004739in}}%
\pgfpathlineto{\pgfqpoint{2.747214in}{1.029487in}}%
\pgfpathlineto{\pgfqpoint{2.747881in}{1.012988in}}%
\pgfpathlineto{\pgfqpoint{2.748910in}{1.070732in}}%
\pgfpathlineto{\pgfqpoint{2.748214in}{1.004739in}}%
\pgfpathlineto{\pgfqpoint{2.749073in}{1.062483in}}%
\pgfpathlineto{\pgfqpoint{2.749221in}{1.054234in}}%
\pgfpathlineto{\pgfqpoint{2.749391in}{1.078981in}}%
\pgfpathlineto{\pgfqpoint{2.749747in}{1.070732in}}%
\pgfpathlineto{\pgfqpoint{2.750398in}{1.078981in}}%
\pgfpathlineto{\pgfqpoint{2.750583in}{1.062483in}}%
\pgfpathlineto{\pgfqpoint{2.751753in}{1.045985in}}%
\pgfpathlineto{\pgfqpoint{2.752575in}{1.054234in}}%
\pgfpathlineto{\pgfqpoint{2.752597in}{1.045985in}}%
\pgfpathlineto{\pgfqpoint{2.753915in}{1.029487in}}%
\pgfpathlineto{\pgfqpoint{2.753937in}{1.045985in}}%
\pgfpathlineto{\pgfqpoint{2.754944in}{1.021237in}}%
\pgfpathlineto{\pgfqpoint{2.756262in}{1.004739in}}%
\pgfpathlineto{\pgfqpoint{2.756284in}{1.021237in}}%
\pgfpathlineto{\pgfqpoint{2.757291in}{1.021237in}}%
\pgfpathlineto{\pgfqpoint{2.759135in}{0.971742in}}%
\pgfpathlineto{\pgfqpoint{2.760623in}{1.029487in}}%
\pgfpathlineto{\pgfqpoint{2.760808in}{1.012988in}}%
\pgfpathlineto{\pgfqpoint{2.762149in}{0.971742in}}%
\pgfpathlineto{\pgfqpoint{2.761312in}{1.021237in}}%
\pgfpathlineto{\pgfqpoint{2.762297in}{0.988241in}}%
\pgfpathlineto{\pgfqpoint{2.762970in}{0.979992in}}%
\pgfpathlineto{\pgfqpoint{2.763326in}{1.021237in}}%
\pgfpathlineto{\pgfqpoint{2.764814in}{0.979992in}}%
\pgfpathlineto{\pgfqpoint{2.764829in}{0.996490in}}%
\pgfpathlineto{\pgfqpoint{2.765480in}{1.004739in}}%
\pgfpathlineto{\pgfqpoint{2.765673in}{0.988241in}}%
\pgfpathlineto{\pgfqpoint{2.765821in}{0.979992in}}%
\pgfpathlineto{\pgfqpoint{2.765984in}{1.004739in}}%
\pgfpathlineto{\pgfqpoint{2.766510in}{0.996490in}}%
\pgfpathlineto{\pgfqpoint{2.766991in}{1.029487in}}%
\pgfpathlineto{\pgfqpoint{2.767665in}{1.004739in}}%
\pgfpathlineto{\pgfqpoint{2.768501in}{0.979992in}}%
\pgfpathlineto{\pgfqpoint{2.767850in}{1.021237in}}%
\pgfpathlineto{\pgfqpoint{2.768686in}{0.988241in}}%
\pgfpathlineto{\pgfqpoint{2.769693in}{0.971742in}}%
\pgfpathlineto{\pgfqpoint{2.768857in}{0.996490in}}%
\pgfpathlineto{\pgfqpoint{2.769842in}{0.979992in}}%
\pgfpathlineto{\pgfqpoint{2.770530in}{1.021237in}}%
\pgfpathlineto{\pgfqpoint{2.770863in}{1.021237in}}%
\pgfpathlineto{\pgfqpoint{2.771182in}{1.004739in}}%
\pgfpathlineto{\pgfqpoint{2.771515in}{1.037736in}}%
\pgfpathlineto{\pgfqpoint{2.772544in}{1.095480in}}%
\pgfpathlineto{\pgfqpoint{2.772707in}{1.087231in}}%
\pgfpathlineto{\pgfqpoint{2.773884in}{1.070732in}}%
\pgfpathlineto{\pgfqpoint{2.774869in}{1.078981in}}%
\pgfpathlineto{\pgfqpoint{2.774891in}{1.070732in}}%
\pgfpathlineto{\pgfqpoint{2.776209in}{1.054234in}}%
\pgfpathlineto{\pgfqpoint{2.777238in}{1.095480in}}%
\pgfpathlineto{\pgfqpoint{2.778578in}{1.070732in}}%
\pgfpathlineto{\pgfqpoint{2.780067in}{1.128476in}}%
\pgfpathlineto{\pgfqpoint{2.780252in}{1.111978in}}%
\pgfpathlineto{\pgfqpoint{2.780422in}{1.120227in}}%
\pgfpathlineto{\pgfqpoint{2.781407in}{1.103729in}}%
\pgfpathlineto{\pgfqpoint{2.782095in}{1.169722in}}%
\pgfpathlineto{\pgfqpoint{2.782436in}{1.144975in}}%
\pgfpathlineto{\pgfqpoint{2.782917in}{1.153224in}}%
\pgfpathlineto{\pgfqpoint{2.783102in}{1.136726in}}%
\pgfpathlineto{\pgfqpoint{2.787441in}{0.633527in}}%
\pgfpathlineto{\pgfqpoint{2.787464in}{0.650025in}}%
\pgfpathlineto{\pgfqpoint{2.789129in}{0.683022in}}%
\pgfpathlineto{\pgfqpoint{2.790144in}{0.674773in}}%
\pgfpathlineto{\pgfqpoint{2.791299in}{0.683022in}}%
\pgfpathlineto{\pgfqpoint{2.792321in}{0.650025in}}%
\pgfpathlineto{\pgfqpoint{2.792972in}{0.633527in}}%
\pgfpathlineto{\pgfqpoint{2.793143in}{0.658274in}}%
\pgfpathlineto{\pgfqpoint{2.793165in}{0.650025in}}%
\pgfpathlineto{\pgfqpoint{2.793476in}{0.658274in}}%
\pgfpathlineto{\pgfqpoint{2.793831in}{0.625278in}}%
\pgfpathlineto{\pgfqpoint{2.793979in}{0.633527in}}%
\pgfpathlineto{\pgfqpoint{2.794986in}{0.625278in}}%
\pgfpathlineto{\pgfqpoint{2.796156in}{0.633527in}}%
\pgfpathlineto{\pgfqpoint{2.796326in}{0.625278in}}%
\pgfpathlineto{\pgfqpoint{2.797000in}{0.625278in}}%
\pgfpathlineto{\pgfqpoint{2.798340in}{0.633527in}}%
\pgfpathlineto{\pgfqpoint{2.798340in}{0.625278in}}%
\pgfpathlineto{\pgfqpoint{2.799029in}{0.650025in}}%
\pgfpathlineto{\pgfqpoint{2.799362in}{0.650025in}}%
\pgfpathlineto{\pgfqpoint{2.800199in}{0.674773in}}%
\pgfpathlineto{\pgfqpoint{2.800517in}{0.658274in}}%
\pgfpathlineto{\pgfqpoint{2.801524in}{0.633527in}}%
\pgfpathlineto{\pgfqpoint{2.801546in}{0.650025in}}%
\pgfpathlineto{\pgfqpoint{2.801857in}{0.633527in}}%
\pgfpathlineto{\pgfqpoint{2.801879in}{0.674773in}}%
\pgfpathlineto{\pgfqpoint{2.802198in}{0.666524in}}%
\pgfpathlineto{\pgfqpoint{2.802531in}{0.658274in}}%
\pgfpathlineto{\pgfqpoint{2.803220in}{0.674773in}}%
\pgfpathlineto{\pgfqpoint{2.804204in}{0.658274in}}%
\pgfpathlineto{\pgfqpoint{2.804227in}{0.674773in}}%
\pgfpathlineto{\pgfqpoint{2.804893in}{0.724268in}}%
\pgfpathlineto{\pgfqpoint{2.805396in}{0.691271in}}%
\pgfpathlineto{\pgfqpoint{2.805545in}{0.683022in}}%
\pgfpathlineto{\pgfqpoint{2.805900in}{0.724268in}}%
\pgfpathlineto{\pgfqpoint{2.808580in}{0.848005in}}%
\pgfpathlineto{\pgfqpoint{2.808906in}{0.831507in}}%
\pgfpathlineto{\pgfqpoint{2.809084in}{0.823258in}}%
\pgfpathlineto{\pgfqpoint{2.809410in}{0.856254in}}%
\pgfpathlineto{\pgfqpoint{2.809928in}{0.848005in}}%
\pgfpathlineto{\pgfqpoint{2.813934in}{1.004739in}}%
\pgfpathlineto{\pgfqpoint{2.813948in}{0.996490in}}%
\pgfpathlineto{\pgfqpoint{2.814607in}{1.054234in}}%
\pgfpathlineto{\pgfqpoint{2.814955in}{1.045985in}}%
\pgfpathlineto{\pgfqpoint{2.816614in}{1.004739in}}%
\pgfpathlineto{\pgfqpoint{2.815288in}{1.070732in}}%
\pgfpathlineto{\pgfqpoint{2.816621in}{1.012988in}}%
\pgfpathlineto{\pgfqpoint{2.816962in}{0.996490in}}%
\pgfpathlineto{\pgfqpoint{2.817636in}{1.021237in}}%
\pgfpathlineto{\pgfqpoint{2.818976in}{0.996490in}}%
\pgfpathlineto{\pgfqpoint{2.819798in}{1.029487in}}%
\pgfpathlineto{\pgfqpoint{2.820131in}{1.004739in}}%
\pgfpathlineto{\pgfqpoint{2.821138in}{0.979992in}}%
\pgfpathlineto{\pgfqpoint{2.821153in}{0.996490in}}%
\pgfpathlineto{\pgfqpoint{2.822656in}{1.054234in}}%
\pgfpathlineto{\pgfqpoint{2.822833in}{1.037736in}}%
\pgfpathlineto{\pgfqpoint{2.823655in}{0.979992in}}%
\pgfpathlineto{\pgfqpoint{2.824174in}{0.996490in}}%
\pgfpathlineto{\pgfqpoint{2.826343in}{1.078981in}}%
\pgfpathlineto{\pgfqpoint{2.826521in}{1.062483in}}%
\pgfpathlineto{\pgfqpoint{2.831881in}{0.823258in}}%
\pgfpathlineto{\pgfqpoint{2.832533in}{0.905749in}}%
\pgfpathlineto{\pgfqpoint{2.833725in}{0.872753in}}%
\pgfpathlineto{\pgfqpoint{2.836242in}{0.749015in}}%
\pgfpathlineto{\pgfqpoint{2.836576in}{0.773763in}}%
\pgfpathlineto{\pgfqpoint{2.837397in}{0.782012in}}%
\pgfpathlineto{\pgfqpoint{2.837412in}{0.773763in}}%
\pgfpathlineto{\pgfqpoint{2.838760in}{0.749015in}}%
\pgfpathlineto{\pgfqpoint{2.838908in}{0.757264in}}%
\pgfpathlineto{\pgfqpoint{2.839263in}{0.724268in}}%
\pgfpathlineto{\pgfqpoint{2.839574in}{0.732517in}}%
\pgfpathlineto{\pgfqpoint{2.839915in}{0.757264in}}%
\pgfpathlineto{\pgfqpoint{2.840603in}{0.724268in}}%
\pgfpathlineto{\pgfqpoint{2.841758in}{0.732517in}}%
\pgfpathlineto{\pgfqpoint{2.842780in}{0.724268in}}%
\pgfpathlineto{\pgfqpoint{2.843935in}{0.782012in}}%
\pgfpathlineto{\pgfqpoint{2.844454in}{0.765513in}}%
\pgfpathlineto{\pgfqpoint{2.846134in}{0.674773in}}%
\pgfpathlineto{\pgfqpoint{2.851332in}{0.922247in}}%
\pgfpathlineto{\pgfqpoint{2.851495in}{0.913998in}}%
\pgfpathlineto{\pgfqpoint{2.851813in}{0.930497in}}%
\pgfpathlineto{\pgfqpoint{2.853154in}{0.881002in}}%
\pgfpathlineto{\pgfqpoint{2.853990in}{0.955244in}}%
\pgfpathlineto{\pgfqpoint{2.854175in}{0.946995in}}%
\pgfpathlineto{\pgfqpoint{2.855330in}{1.029487in}}%
\pgfpathlineto{\pgfqpoint{2.855856in}{0.988241in}}%
\pgfpathlineto{\pgfqpoint{2.856004in}{0.979992in}}%
\pgfpathlineto{\pgfqpoint{2.856174in}{1.004739in}}%
\pgfpathlineto{\pgfqpoint{2.856522in}{0.996490in}}%
\pgfpathlineto{\pgfqpoint{2.857011in}{0.979992in}}%
\pgfpathlineto{\pgfqpoint{2.857685in}{1.004739in}}%
\pgfpathlineto{\pgfqpoint{2.858707in}{0.996490in}}%
\pgfpathlineto{\pgfqpoint{2.860217in}{1.045985in}}%
\pgfpathlineto{\pgfqpoint{2.860380in}{1.037736in}}%
\pgfpathlineto{\pgfqpoint{2.860869in}{1.029487in}}%
\pgfpathlineto{\pgfqpoint{2.861032in}{1.054234in}}%
\pgfpathlineto{\pgfqpoint{2.861054in}{1.045985in}}%
\pgfpathlineto{\pgfqpoint{2.863046in}{1.128476in}}%
\pgfpathlineto{\pgfqpoint{2.863231in}{1.111978in}}%
\pgfpathlineto{\pgfqpoint{2.863712in}{1.103729in}}%
\pgfpathlineto{\pgfqpoint{2.863890in}{1.128476in}}%
\pgfpathlineto{\pgfqpoint{2.863904in}{1.120227in}}%
\pgfpathlineto{\pgfqpoint{2.864053in}{1.128476in}}%
\pgfpathlineto{\pgfqpoint{2.864408in}{1.095480in}}%
\pgfpathlineto{\pgfqpoint{2.864571in}{1.095480in}}%
\pgfpathlineto{\pgfqpoint{2.866585in}{1.021237in}}%
\pgfpathlineto{\pgfqpoint{2.866733in}{1.037736in}}%
\pgfpathlineto{\pgfqpoint{2.867066in}{1.029487in}}%
\pgfpathlineto{\pgfqpoint{2.867755in}{1.045985in}}%
\pgfpathlineto{\pgfqpoint{2.869435in}{0.913998in}}%
\pgfpathlineto{\pgfqpoint{2.871427in}{0.707769in}}%
\pgfpathlineto{\pgfqpoint{2.871612in}{0.724268in}}%
\pgfpathlineto{\pgfqpoint{2.871782in}{0.749015in}}%
\pgfpathlineto{\pgfqpoint{2.872775in}{0.732517in}}%
\pgfpathlineto{\pgfqpoint{2.873789in}{0.699520in}}%
\pgfpathlineto{\pgfqpoint{2.876477in}{0.848005in}}%
\pgfpathlineto{\pgfqpoint{2.877136in}{0.831507in}}%
\pgfpathlineto{\pgfqpoint{2.878150in}{0.798510in}}%
\pgfpathlineto{\pgfqpoint{2.880142in}{0.856254in}}%
\pgfpathlineto{\pgfqpoint{2.880164in}{0.848005in}}%
\pgfpathlineto{\pgfqpoint{2.882319in}{0.683022in}}%
\pgfpathlineto{\pgfqpoint{2.883496in}{0.658274in}}%
\pgfpathlineto{\pgfqpoint{2.884518in}{0.699520in}}%
\pgfpathlineto{\pgfqpoint{2.884666in}{0.707769in}}%
\pgfpathlineto{\pgfqpoint{2.885021in}{0.674773in}}%
\pgfpathlineto{\pgfqpoint{2.885169in}{0.683022in}}%
\pgfpathlineto{\pgfqpoint{2.886176in}{0.658274in}}%
\pgfpathlineto{\pgfqpoint{2.886198in}{0.674773in}}%
\pgfpathlineto{\pgfqpoint{2.886902in}{0.707769in}}%
\pgfpathlineto{\pgfqpoint{2.887242in}{0.683022in}}%
\pgfpathlineto{\pgfqpoint{2.888375in}{0.650025in}}%
\pgfpathlineto{\pgfqpoint{2.888523in}{0.666524in}}%
\pgfpathlineto{\pgfqpoint{2.889715in}{0.724268in}}%
\pgfpathlineto{\pgfqpoint{2.890034in}{0.732517in}}%
\pgfpathlineto{\pgfqpoint{2.890219in}{0.716019in}}%
\pgfpathlineto{\pgfqpoint{2.891056in}{0.674773in}}%
\pgfpathlineto{\pgfqpoint{2.891396in}{0.699520in}}%
\pgfpathlineto{\pgfqpoint{2.892233in}{0.724268in}}%
\pgfpathlineto{\pgfqpoint{2.892551in}{0.707769in}}%
\pgfpathlineto{\pgfqpoint{2.893551in}{0.683022in}}%
\pgfpathlineto{\pgfqpoint{2.893573in}{0.699520in}}%
\pgfpathlineto{\pgfqpoint{2.894225in}{0.683022in}}%
\pgfpathlineto{\pgfqpoint{2.894232in}{0.716019in}}%
\pgfpathlineto{\pgfqpoint{2.895394in}{0.831507in}}%
\pgfpathlineto{\pgfqpoint{2.894558in}{0.707769in}}%
\pgfpathlineto{\pgfqpoint{2.895417in}{0.823258in}}%
\pgfpathlineto{\pgfqpoint{2.900089in}{1.029487in}}%
\pgfpathlineto{\pgfqpoint{2.900614in}{0.996490in}}%
\pgfpathlineto{\pgfqpoint{2.901118in}{1.021237in}}%
\pgfpathlineto{\pgfqpoint{2.902458in}{0.996490in}}%
\pgfpathlineto{\pgfqpoint{2.902962in}{0.971742in}}%
\pgfpathlineto{\pgfqpoint{2.904117in}{1.029487in}}%
\pgfpathlineto{\pgfqpoint{2.904131in}{1.021237in}}%
\pgfpathlineto{\pgfqpoint{2.905116in}{1.087231in}}%
\pgfpathlineto{\pgfqpoint{2.906301in}{1.153224in}}%
\pgfpathlineto{\pgfqpoint{2.906316in}{1.144975in}}%
\pgfpathlineto{\pgfqpoint{2.907486in}{1.169722in}}%
\pgfpathlineto{\pgfqpoint{2.907656in}{1.161473in}}%
\pgfpathlineto{\pgfqpoint{2.907804in}{1.153224in}}%
\pgfpathlineto{\pgfqpoint{2.907974in}{1.177971in}}%
\pgfpathlineto{\pgfqpoint{2.908470in}{1.177971in}}%
\pgfpathlineto{\pgfqpoint{2.909159in}{1.219217in}}%
\pgfpathlineto{\pgfqpoint{2.909499in}{1.194470in}}%
\pgfpathlineto{\pgfqpoint{2.910166in}{1.111978in}}%
\pgfpathlineto{\pgfqpoint{2.914512in}{0.625278in}}%
\pgfpathlineto{\pgfqpoint{2.915682in}{0.633527in}}%
\pgfpathlineto{\pgfqpoint{2.915682in}{0.625278in}}%
\pgfpathlineto{\pgfqpoint{2.916704in}{0.625278in}}%
\pgfpathlineto{\pgfqpoint{2.917866in}{0.633527in}}%
\pgfpathlineto{\pgfqpoint{2.918029in}{0.625278in}}%
\pgfpathlineto{\pgfqpoint{2.918866in}{0.625278in}}%
\pgfpathlineto{\pgfqpoint{2.919873in}{0.633527in}}%
\pgfpathlineto{\pgfqpoint{2.919888in}{0.625278in}}%
\pgfpathlineto{\pgfqpoint{2.920561in}{0.650025in}}%
\pgfpathlineto{\pgfqpoint{2.921213in}{0.633527in}}%
\pgfpathlineto{\pgfqpoint{2.921228in}{0.625278in}}%
\pgfpathlineto{\pgfqpoint{2.922072in}{0.699520in}}%
\pgfpathlineto{\pgfqpoint{2.922235in}{0.691271in}}%
\pgfpathlineto{\pgfqpoint{2.923412in}{0.650025in}}%
\pgfpathlineto{\pgfqpoint{2.923575in}{0.674773in}}%
\pgfpathlineto{\pgfqpoint{2.924730in}{0.683022in}}%
\pgfpathlineto{\pgfqpoint{2.924752in}{0.674773in}}%
\pgfpathlineto{\pgfqpoint{2.925404in}{0.732517in}}%
\pgfpathlineto{\pgfqpoint{2.925589in}{0.724268in}}%
\pgfpathlineto{\pgfqpoint{2.928269in}{0.897500in}}%
\pgfpathlineto{\pgfqpoint{2.928928in}{0.856254in}}%
\pgfpathlineto{\pgfqpoint{2.929780in}{0.848005in}}%
\pgfpathlineto{\pgfqpoint{2.929950in}{0.872753in}}%
\pgfpathlineto{\pgfqpoint{2.931105in}{0.881002in}}%
\pgfpathlineto{\pgfqpoint{2.932112in}{0.856254in}}%
\pgfpathlineto{\pgfqpoint{2.932127in}{0.872753in}}%
\pgfpathlineto{\pgfqpoint{2.934955in}{0.782012in}}%
\pgfpathlineto{\pgfqpoint{2.935984in}{0.922247in}}%
\pgfpathlineto{\pgfqpoint{2.936147in}{0.913998in}}%
\pgfpathlineto{\pgfqpoint{2.936303in}{0.905749in}}%
\pgfpathlineto{\pgfqpoint{2.936643in}{0.938746in}}%
\pgfpathlineto{\pgfqpoint{2.936969in}{0.930497in}}%
\pgfpathlineto{\pgfqpoint{2.937658in}{0.971742in}}%
\pgfpathlineto{\pgfqpoint{2.939842in}{1.070732in}}%
\pgfpathlineto{\pgfqpoint{2.940160in}{1.054234in}}%
\pgfpathlineto{\pgfqpoint{2.940175in}{1.045985in}}%
\pgfpathlineto{\pgfqpoint{2.940997in}{1.153224in}}%
\pgfpathlineto{\pgfqpoint{2.941012in}{1.144975in}}%
\pgfpathlineto{\pgfqpoint{2.942004in}{1.227466in}}%
\pgfpathlineto{\pgfqpoint{2.943011in}{1.202719in}}%
\pgfpathlineto{\pgfqpoint{2.944196in}{1.144975in}}%
\pgfpathlineto{\pgfqpoint{2.944869in}{1.169722in}}%
\pgfpathlineto{\pgfqpoint{2.945202in}{1.144975in}}%
\pgfpathlineto{\pgfqpoint{2.946713in}{1.070732in}}%
\pgfpathlineto{\pgfqpoint{2.947083in}{1.103729in}}%
\pgfpathlineto{\pgfqpoint{2.947216in}{1.120227in}}%
\pgfpathlineto{\pgfqpoint{2.947550in}{1.095480in}}%
\pgfpathlineto{\pgfqpoint{2.948053in}{1.111978in}}%
\pgfpathlineto{\pgfqpoint{2.951911in}{0.674773in}}%
\pgfpathlineto{\pgfqpoint{2.952059in}{0.683022in}}%
\pgfpathlineto{\pgfqpoint{2.953066in}{0.732517in}}%
\pgfpathlineto{\pgfqpoint{2.953081in}{0.724268in}}%
\pgfpathlineto{\pgfqpoint{2.953740in}{0.757264in}}%
\pgfpathlineto{\pgfqpoint{2.954088in}{0.724268in}}%
\pgfpathlineto{\pgfqpoint{2.956916in}{0.625278in}}%
\pgfpathlineto{\pgfqpoint{2.958093in}{0.633527in}}%
\pgfpathlineto{\pgfqpoint{2.958256in}{0.625278in}}%
\pgfpathlineto{\pgfqpoint{2.959100in}{0.625278in}}%
\pgfpathlineto{\pgfqpoint{2.960270in}{0.633527in}}%
\pgfpathlineto{\pgfqpoint{2.960448in}{0.625278in}}%
\pgfpathlineto{\pgfqpoint{2.961277in}{0.625278in}}%
\pgfpathlineto{\pgfqpoint{2.962454in}{0.633527in}}%
\pgfpathlineto{\pgfqpoint{2.962617in}{0.625278in}}%
\pgfpathlineto{\pgfqpoint{2.963454in}{0.625278in}}%
\pgfpathlineto{\pgfqpoint{2.964631in}{0.633527in}}%
\pgfpathlineto{\pgfqpoint{2.964794in}{0.625278in}}%
\pgfpathlineto{\pgfqpoint{2.965468in}{0.625278in}}%
\pgfpathlineto{\pgfqpoint{2.966645in}{0.633527in}}%
\pgfpathlineto{\pgfqpoint{2.966993in}{0.625278in}}%
\pgfpathlineto{\pgfqpoint{2.967319in}{0.658274in}}%
\pgfpathlineto{\pgfqpoint{2.967667in}{0.650025in}}%
\pgfpathlineto{\pgfqpoint{2.968326in}{0.683022in}}%
\pgfpathlineto{\pgfqpoint{2.968837in}{0.666524in}}%
\pgfpathlineto{\pgfqpoint{2.968985in}{0.658274in}}%
\pgfpathlineto{\pgfqpoint{2.969007in}{0.699520in}}%
\pgfpathlineto{\pgfqpoint{2.969340in}{0.699520in}}%
\pgfpathlineto{\pgfqpoint{2.969673in}{0.749015in}}%
\pgfpathlineto{\pgfqpoint{2.970680in}{0.724268in}}%
\pgfpathlineto{\pgfqpoint{2.971502in}{0.683022in}}%
\pgfpathlineto{\pgfqpoint{2.971843in}{0.716019in}}%
\pgfpathlineto{\pgfqpoint{2.972020in}{0.732517in}}%
\pgfpathlineto{\pgfqpoint{2.972361in}{0.699520in}}%
\pgfpathlineto{\pgfqpoint{2.972865in}{0.699520in}}%
\pgfpathlineto{\pgfqpoint{2.974205in}{0.674773in}}%
\pgfpathlineto{\pgfqpoint{2.975360in}{0.683022in}}%
\pgfpathlineto{\pgfqpoint{2.975545in}{0.674773in}}%
\pgfpathlineto{\pgfqpoint{2.976367in}{0.740766in}}%
\pgfpathlineto{\pgfqpoint{2.976382in}{0.773763in}}%
\pgfpathlineto{\pgfqpoint{2.977389in}{0.716019in}}%
\pgfpathlineto{\pgfqpoint{2.978395in}{0.699520in}}%
\pgfpathlineto{\pgfqpoint{2.977559in}{0.724268in}}%
\pgfpathlineto{\pgfqpoint{2.978544in}{0.707769in}}%
\pgfpathlineto{\pgfqpoint{2.978558in}{0.724268in}}%
\pgfpathlineto{\pgfqpoint{2.979565in}{0.691271in}}%
\pgfpathlineto{\pgfqpoint{2.980069in}{0.674773in}}%
\pgfpathlineto{\pgfqpoint{2.980239in}{0.699520in}}%
\pgfpathlineto{\pgfqpoint{2.980720in}{0.683022in}}%
\pgfpathlineto{\pgfqpoint{2.981394in}{0.707769in}}%
\pgfpathlineto{\pgfqpoint{2.982416in}{0.674773in}}%
\pgfpathlineto{\pgfqpoint{2.983578in}{0.707769in}}%
\pgfpathlineto{\pgfqpoint{2.983756in}{0.691271in}}%
\pgfpathlineto{\pgfqpoint{2.984933in}{0.674773in}}%
\pgfpathlineto{\pgfqpoint{2.985252in}{0.683022in}}%
\pgfpathlineto{\pgfqpoint{2.985585in}{0.658274in}}%
\pgfpathlineto{\pgfqpoint{2.985770in}{0.666524in}}%
\pgfpathlineto{\pgfqpoint{2.985940in}{0.674773in}}%
\pgfpathlineto{\pgfqpoint{2.987443in}{0.625278in}}%
\pgfpathlineto{\pgfqpoint{2.988436in}{0.658274in}}%
\pgfpathlineto{\pgfqpoint{2.988621in}{0.641776in}}%
\pgfpathlineto{\pgfqpoint{2.989287in}{0.650025in}}%
\pgfpathlineto{\pgfqpoint{2.989776in}{0.625278in}}%
\pgfpathlineto{\pgfqpoint{2.990953in}{0.633527in}}%
\pgfpathlineto{\pgfqpoint{2.991123in}{0.625278in}}%
\pgfpathlineto{\pgfqpoint{2.991967in}{0.625278in}}%
\pgfpathlineto{\pgfqpoint{2.993293in}{0.633527in}}%
\pgfpathlineto{\pgfqpoint{2.993470in}{0.625278in}}%
\pgfpathlineto{\pgfqpoint{2.994300in}{0.625278in}}%
\pgfpathlineto{\pgfqpoint{2.995470in}{0.633527in}}%
\pgfpathlineto{\pgfqpoint{2.995640in}{0.625278in}}%
\pgfpathlineto{\pgfqpoint{2.996477in}{0.625278in}}%
\pgfpathlineto{\pgfqpoint{2.997817in}{0.633527in}}%
\pgfpathlineto{\pgfqpoint{2.997817in}{0.625278in}}%
\pgfpathlineto{\pgfqpoint{2.998824in}{0.625278in}}%
\pgfpathlineto{\pgfqpoint{3.000001in}{0.633527in}}%
\pgfpathlineto{\pgfqpoint{3.000001in}{0.625278in}}%
\pgfpathlineto{\pgfqpoint{3.001008in}{0.625278in}}%
\pgfpathlineto{\pgfqpoint{3.002178in}{0.633527in}}%
\pgfpathlineto{\pgfqpoint{3.002348in}{0.625278in}}%
\pgfpathlineto{\pgfqpoint{3.003192in}{0.625278in}}%
\pgfpathlineto{\pgfqpoint{3.004192in}{0.633527in}}%
\pgfpathlineto{\pgfqpoint{3.004192in}{0.625278in}}%
\pgfpathlineto{\pgfqpoint{3.005362in}{0.633527in}}%
\pgfpathlineto{\pgfqpoint{3.005532in}{0.625278in}}%
\pgfpathlineto{\pgfqpoint{3.006369in}{0.625278in}}%
\pgfpathlineto{\pgfqpoint{3.007376in}{0.633527in}}%
\pgfpathlineto{\pgfqpoint{3.007376in}{0.625278in}}%
\pgfpathlineto{\pgfqpoint{3.008716in}{0.633527in}}%
\pgfpathlineto{\pgfqpoint{3.008886in}{0.625278in}}%
\pgfpathlineto{\pgfqpoint{3.009560in}{0.625278in}}%
\pgfpathlineto{\pgfqpoint{3.010730in}{0.633527in}}%
\pgfpathlineto{\pgfqpoint{3.010900in}{0.625278in}}%
\pgfpathlineto{\pgfqpoint{3.011737in}{0.625278in}}%
\pgfpathlineto{\pgfqpoint{3.012906in}{0.633527in}}%
\pgfpathlineto{\pgfqpoint{3.013077in}{0.625278in}}%
\pgfpathlineto{\pgfqpoint{3.013743in}{0.625278in}}%
\pgfpathlineto{\pgfqpoint{3.014920in}{0.633527in}}%
\pgfpathlineto{\pgfqpoint{3.014920in}{0.625278in}}%
\pgfpathlineto{\pgfqpoint{3.015927in}{0.625278in}}%
\pgfpathlineto{\pgfqpoint{3.017097in}{0.633527in}}%
\pgfpathlineto{\pgfqpoint{3.017097in}{0.625278in}}%
\pgfpathlineto{\pgfqpoint{3.017934in}{0.625278in}}%
\pgfpathlineto{\pgfqpoint{3.019111in}{0.633527in}}%
\pgfpathlineto{\pgfqpoint{3.019274in}{0.625278in}}%
\pgfpathlineto{\pgfqpoint{3.020118in}{0.625278in}}%
\pgfpathlineto{\pgfqpoint{3.021288in}{0.633527in}}%
\pgfpathlineto{\pgfqpoint{3.021288in}{0.625278in}}%
\pgfpathlineto{\pgfqpoint{3.022125in}{0.625278in}}%
\pgfpathlineto{\pgfqpoint{3.023302in}{0.633527in}}%
\pgfpathlineto{\pgfqpoint{3.023465in}{0.625278in}}%
\pgfpathlineto{\pgfqpoint{3.024309in}{0.625278in}}%
\pgfpathlineto{\pgfqpoint{3.025479in}{0.633527in}}%
\pgfpathlineto{\pgfqpoint{3.025649in}{0.625278in}}%
\pgfpathlineto{\pgfqpoint{3.026315in}{0.625278in}}%
\pgfpathlineto{\pgfqpoint{3.027493in}{0.633527in}}%
\pgfpathlineto{\pgfqpoint{3.027493in}{0.625278in}}%
\pgfpathlineto{\pgfqpoint{3.028500in}{0.625278in}}%
\pgfpathlineto{\pgfqpoint{3.029670in}{0.633527in}}%
\pgfpathlineto{\pgfqpoint{3.029840in}{0.625278in}}%
\pgfpathlineto{\pgfqpoint{3.030677in}{0.625278in}}%
\pgfpathlineto{\pgfqpoint{3.031846in}{0.633527in}}%
\pgfpathlineto{\pgfqpoint{3.032017in}{0.625278in}}%
\pgfpathlineto{\pgfqpoint{3.032853in}{0.625278in}}%
\pgfpathlineto{\pgfqpoint{3.033860in}{0.633527in}}%
\pgfpathlineto{\pgfqpoint{3.033860in}{0.625278in}}%
\pgfpathlineto{\pgfqpoint{3.035200in}{0.633527in}}%
\pgfpathlineto{\pgfqpoint{3.035534in}{0.625278in}}%
\pgfpathlineto{\pgfqpoint{3.036207in}{0.625278in}}%
\pgfpathlineto{\pgfqpoint{3.037377in}{0.633527in}}%
\pgfpathlineto{\pgfqpoint{3.037548in}{0.625278in}}%
\pgfpathlineto{\pgfqpoint{3.038392in}{0.625278in}}%
\pgfpathlineto{\pgfqpoint{3.039562in}{0.633527in}}%
\pgfpathlineto{\pgfqpoint{3.039724in}{0.625278in}}%
\pgfpathlineto{\pgfqpoint{3.040398in}{0.625278in}}%
\pgfpathlineto{\pgfqpoint{3.041568in}{0.633527in}}%
\pgfpathlineto{\pgfqpoint{3.041738in}{0.625278in}}%
\pgfpathlineto{\pgfqpoint{3.042412in}{0.625278in}}%
\pgfpathlineto{\pgfqpoint{3.043582in}{0.633527in}}%
\pgfpathlineto{\pgfqpoint{3.043752in}{0.625278in}}%
\pgfpathlineto{\pgfqpoint{3.044419in}{0.625278in}}%
\pgfpathlineto{\pgfqpoint{3.045759in}{0.633527in}}%
\pgfpathlineto{\pgfqpoint{3.045759in}{0.625278in}}%
\pgfpathlineto{\pgfqpoint{3.046603in}{0.625278in}}%
\pgfpathlineto{\pgfqpoint{3.047773in}{0.633527in}}%
\pgfpathlineto{\pgfqpoint{3.047943in}{0.625278in}}%
\pgfpathlineto{\pgfqpoint{3.048780in}{0.625278in}}%
\pgfpathlineto{\pgfqpoint{3.049950in}{0.633527in}}%
\pgfpathlineto{\pgfqpoint{3.049950in}{0.625278in}}%
\pgfpathlineto{\pgfqpoint{3.050957in}{0.625278in}}%
\pgfpathlineto{\pgfqpoint{3.052134in}{0.633527in}}%
\pgfpathlineto{\pgfqpoint{3.052134in}{0.625278in}}%
\pgfpathlineto{\pgfqpoint{3.053141in}{0.625278in}}%
\pgfpathlineto{\pgfqpoint{3.054311in}{0.633527in}}%
\pgfpathlineto{\pgfqpoint{3.054481in}{0.625278in}}%
\pgfpathlineto{\pgfqpoint{3.055318in}{0.625278in}}%
\pgfpathlineto{\pgfqpoint{3.056488in}{0.633527in}}%
\pgfpathlineto{\pgfqpoint{3.056658in}{0.625278in}}%
\pgfpathlineto{\pgfqpoint{3.057346in}{0.625278in}}%
\pgfpathlineto{\pgfqpoint{3.058501in}{0.633527in}}%
\pgfpathlineto{\pgfqpoint{3.058501in}{0.625278in}}%
\pgfpathlineto{\pgfqpoint{3.059508in}{0.625278in}}%
\pgfpathlineto{\pgfqpoint{3.060849in}{0.633527in}}%
\pgfpathlineto{\pgfqpoint{3.060849in}{0.625278in}}%
\pgfpathlineto{\pgfqpoint{3.061856in}{0.625278in}}%
\pgfpathlineto{\pgfqpoint{3.062863in}{0.633527in}}%
\pgfpathlineto{\pgfqpoint{3.062863in}{0.625278in}}%
\pgfpathlineto{\pgfqpoint{3.064203in}{0.633527in}}%
\pgfpathlineto{\pgfqpoint{3.064203in}{0.625278in}}%
\pgfpathlineto{\pgfqpoint{3.065210in}{0.625278in}}%
\pgfpathlineto{\pgfqpoint{3.066550in}{0.633527in}}%
\pgfpathlineto{\pgfqpoint{3.066550in}{0.625278in}}%
\pgfpathlineto{\pgfqpoint{3.067572in}{0.625278in}}%
\pgfpathlineto{\pgfqpoint{3.068727in}{0.633527in}}%
\pgfpathlineto{\pgfqpoint{3.068897in}{0.625278in}}%
\pgfpathlineto{\pgfqpoint{3.069400in}{0.658274in}}%
\pgfpathlineto{\pgfqpoint{3.069756in}{0.641776in}}%
\pgfpathlineto{\pgfqpoint{3.070755in}{0.625278in}}%
\pgfpathlineto{\pgfqpoint{3.071910in}{0.633527in}}%
\pgfpathlineto{\pgfqpoint{3.071970in}{0.625278in}}%
\pgfpathlineto{\pgfqpoint{3.072940in}{0.699520in}}%
\pgfpathlineto{\pgfqpoint{3.073591in}{0.732517in}}%
\pgfpathlineto{\pgfqpoint{3.074095in}{0.707769in}}%
\pgfpathlineto{\pgfqpoint{3.074613in}{0.699520in}}%
\pgfpathlineto{\pgfqpoint{3.075094in}{0.765513in}}%
\pgfpathlineto{\pgfqpoint{3.075435in}{0.782012in}}%
\pgfpathlineto{\pgfqpoint{3.076101in}{0.732517in}}%
\pgfpathlineto{\pgfqpoint{3.076123in}{0.749015in}}%
\pgfpathlineto{\pgfqpoint{3.077464in}{0.848005in}}%
\pgfpathlineto{\pgfqpoint{3.078115in}{0.806759in}}%
\pgfpathlineto{\pgfqpoint{3.079307in}{0.773763in}}%
\pgfpathlineto{\pgfqpoint{3.080462in}{0.782012in}}%
\pgfpathlineto{\pgfqpoint{3.081299in}{0.757264in}}%
\pgfpathlineto{\pgfqpoint{3.081484in}{0.765513in}}%
\pgfpathlineto{\pgfqpoint{3.082661in}{0.749015in}}%
\pgfpathlineto{\pgfqpoint{3.083668in}{0.798510in}}%
\pgfpathlineto{\pgfqpoint{3.084002in}{0.773763in}}%
\pgfpathlineto{\pgfqpoint{3.084320in}{0.757264in}}%
\pgfpathlineto{\pgfqpoint{3.084653in}{0.790261in}}%
\pgfpathlineto{\pgfqpoint{3.085157in}{0.806759in}}%
\pgfpathlineto{\pgfqpoint{3.085490in}{0.782012in}}%
\pgfpathlineto{\pgfqpoint{3.085675in}{0.790261in}}%
\pgfpathlineto{\pgfqpoint{3.086497in}{0.782012in}}%
\pgfpathlineto{\pgfqpoint{3.085845in}{0.798510in}}%
\pgfpathlineto{\pgfqpoint{3.086519in}{0.798510in}}%
\pgfpathlineto{\pgfqpoint{3.087171in}{0.806759in}}%
\pgfpathlineto{\pgfqpoint{3.087356in}{0.790261in}}%
\pgfpathlineto{\pgfqpoint{3.088170in}{0.782012in}}%
\pgfpathlineto{\pgfqpoint{3.087519in}{0.798510in}}%
\pgfpathlineto{\pgfqpoint{3.088170in}{0.790261in}}%
\pgfpathlineto{\pgfqpoint{3.088859in}{0.872753in}}%
\pgfpathlineto{\pgfqpoint{3.089199in}{0.872753in}}%
\pgfpathlineto{\pgfqpoint{3.090406in}{0.856254in}}%
\pgfpathlineto{\pgfqpoint{3.090873in}{0.946995in}}%
\pgfpathlineto{\pgfqpoint{3.091709in}{0.889251in}}%
\pgfpathlineto{\pgfqpoint{3.091857in}{0.881002in}}%
\pgfpathlineto{\pgfqpoint{3.092198in}{0.913998in}}%
\pgfpathlineto{\pgfqpoint{3.093035in}{1.004739in}}%
\pgfpathlineto{\pgfqpoint{3.093220in}{0.996490in}}%
\pgfpathlineto{\pgfqpoint{3.093701in}{1.004739in}}%
\pgfpathlineto{\pgfqpoint{3.093894in}{0.988241in}}%
\pgfpathlineto{\pgfqpoint{3.094042in}{0.979992in}}%
\pgfpathlineto{\pgfqpoint{3.094545in}{1.012988in}}%
\pgfpathlineto{\pgfqpoint{3.094878in}{1.004739in}}%
\pgfpathlineto{\pgfqpoint{3.095737in}{1.095480in}}%
\pgfpathlineto{\pgfqpoint{3.096907in}{1.070732in}}%
\pgfpathlineto{\pgfqpoint{3.097055in}{1.078981in}}%
\pgfpathlineto{\pgfqpoint{3.097077in}{1.095480in}}%
\pgfpathlineto{\pgfqpoint{3.097892in}{1.054234in}}%
\pgfpathlineto{\pgfqpoint{3.098084in}{1.070732in}}%
\pgfpathlineto{\pgfqpoint{3.098395in}{1.078981in}}%
\pgfpathlineto{\pgfqpoint{3.098736in}{1.054234in}}%
\pgfpathlineto{\pgfqpoint{3.098921in}{1.062483in}}%
\pgfpathlineto{\pgfqpoint{3.099091in}{1.070732in}}%
\pgfpathlineto{\pgfqpoint{3.100076in}{1.054234in}}%
\pgfpathlineto{\pgfqpoint{3.100091in}{1.070732in}}%
\pgfpathlineto{\pgfqpoint{3.101098in}{1.070732in}}%
\pgfpathlineto{\pgfqpoint{3.102423in}{1.054234in}}%
\pgfpathlineto{\pgfqpoint{3.102438in}{1.070732in}}%
\pgfpathlineto{\pgfqpoint{3.103445in}{1.070732in}}%
\pgfpathlineto{\pgfqpoint{3.104097in}{1.103729in}}%
\pgfpathlineto{\pgfqpoint{3.104600in}{1.078981in}}%
\pgfpathlineto{\pgfqpoint{3.104785in}{1.070732in}}%
\pgfpathlineto{\pgfqpoint{3.105289in}{1.120227in}}%
\pgfpathlineto{\pgfqpoint{3.105622in}{1.095480in}}%
\pgfpathlineto{\pgfqpoint{3.106777in}{1.103729in}}%
\pgfpathlineto{\pgfqpoint{3.106969in}{1.120227in}}%
\pgfpathlineto{\pgfqpoint{3.107806in}{1.095480in}}%
\pgfpathlineto{\pgfqpoint{3.108976in}{1.120227in}}%
\pgfpathlineto{\pgfqpoint{3.109146in}{1.111978in}}%
\pgfpathlineto{\pgfqpoint{3.109294in}{1.103729in}}%
\pgfpathlineto{\pgfqpoint{3.109650in}{1.144975in}}%
\pgfpathlineto{\pgfqpoint{3.110990in}{1.194470in}}%
\pgfpathlineto{\pgfqpoint{3.111308in}{1.177971in}}%
\pgfpathlineto{\pgfqpoint{3.111664in}{1.169722in}}%
\pgfpathlineto{\pgfqpoint{3.112315in}{1.227466in}}%
\pgfpathlineto{\pgfqpoint{3.112330in}{1.219217in}}%
\pgfpathlineto{\pgfqpoint{3.113337in}{1.243965in}}%
\pgfpathlineto{\pgfqpoint{3.113507in}{1.235715in}}%
\pgfpathlineto{\pgfqpoint{3.114507in}{1.219217in}}%
\pgfpathlineto{\pgfqpoint{3.114655in}{1.227466in}}%
\pgfpathlineto{\pgfqpoint{3.114825in}{1.252214in}}%
\pgfpathlineto{\pgfqpoint{3.115684in}{1.243965in}}%
\pgfpathlineto{\pgfqpoint{3.116188in}{1.219217in}}%
\pgfpathlineto{\pgfqpoint{3.116506in}{1.252214in}}%
\pgfpathlineto{\pgfqpoint{3.116521in}{1.243965in}}%
\pgfpathlineto{\pgfqpoint{3.117506in}{1.301709in}}%
\pgfpathlineto{\pgfqpoint{3.117861in}{1.268712in}}%
\pgfpathlineto{\pgfqpoint{3.118179in}{1.252214in}}%
\pgfpathlineto{\pgfqpoint{3.118350in}{1.276961in}}%
\pgfpathlineto{\pgfqpoint{3.118683in}{1.276961in}}%
\pgfpathlineto{\pgfqpoint{3.118698in}{1.293460in}}%
\pgfpathlineto{\pgfqpoint{3.119519in}{1.252214in}}%
\pgfpathlineto{\pgfqpoint{3.119705in}{1.268712in}}%
\pgfpathlineto{\pgfqpoint{3.120697in}{1.276961in}}%
\pgfpathlineto{\pgfqpoint{3.120712in}{1.268712in}}%
\pgfpathlineto{\pgfqpoint{3.122540in}{1.227466in}}%
\pgfpathlineto{\pgfqpoint{3.122555in}{1.243965in}}%
\pgfpathlineto{\pgfqpoint{3.123562in}{1.219217in}}%
\pgfpathlineto{\pgfqpoint{3.124887in}{1.202719in}}%
\pgfpathlineto{\pgfqpoint{3.125576in}{1.243965in}}%
\pgfpathlineto{\pgfqpoint{3.125909in}{1.219217in}}%
\pgfpathlineto{\pgfqpoint{3.126731in}{1.252214in}}%
\pgfpathlineto{\pgfqpoint{3.127249in}{1.235715in}}%
\pgfpathlineto{\pgfqpoint{3.127731in}{1.227466in}}%
\pgfpathlineto{\pgfqpoint{3.127753in}{1.268712in}}%
\pgfpathlineto{\pgfqpoint{3.128923in}{1.161473in}}%
\pgfpathlineto{\pgfqpoint{3.133432in}{0.633527in}}%
\pgfpathlineto{\pgfqpoint{3.133454in}{0.650025in}}%
\pgfpathlineto{\pgfqpoint{3.133765in}{0.633527in}}%
\pgfpathlineto{\pgfqpoint{3.133935in}{0.658274in}}%
\pgfpathlineto{\pgfqpoint{3.134291in}{0.650025in}}%
\pgfpathlineto{\pgfqpoint{3.136453in}{0.782012in}}%
\pgfpathlineto{\pgfqpoint{3.137290in}{0.757264in}}%
\pgfpathlineto{\pgfqpoint{3.137475in}{0.773763in}}%
\pgfpathlineto{\pgfqpoint{3.138311in}{0.749015in}}%
\pgfpathlineto{\pgfqpoint{3.139133in}{0.757264in}}%
\pgfpathlineto{\pgfqpoint{3.139148in}{0.749015in}}%
\pgfpathlineto{\pgfqpoint{3.140829in}{0.699520in}}%
\pgfpathlineto{\pgfqpoint{3.140977in}{0.707769in}}%
\pgfpathlineto{\pgfqpoint{3.141332in}{0.749015in}}%
\pgfpathlineto{\pgfqpoint{3.141999in}{0.724268in}}%
\pgfpathlineto{\pgfqpoint{3.142487in}{0.732517in}}%
\pgfpathlineto{\pgfqpoint{3.142672in}{0.716019in}}%
\pgfpathlineto{\pgfqpoint{3.143509in}{0.699520in}}%
\pgfpathlineto{\pgfqpoint{3.142843in}{0.724268in}}%
\pgfpathlineto{\pgfqpoint{3.143657in}{0.716019in}}%
\pgfpathlineto{\pgfqpoint{3.144338in}{0.757264in}}%
\pgfpathlineto{\pgfqpoint{3.144686in}{0.724268in}}%
\pgfpathlineto{\pgfqpoint{3.145841in}{0.732517in}}%
\pgfpathlineto{\pgfqpoint{3.146863in}{0.724268in}}%
\pgfpathlineto{\pgfqpoint{3.147515in}{0.757264in}}%
\pgfpathlineto{\pgfqpoint{3.148018in}{0.732517in}}%
\pgfpathlineto{\pgfqpoint{3.148855in}{0.707769in}}%
\pgfpathlineto{\pgfqpoint{3.149040in}{0.716019in}}%
\pgfpathlineto{\pgfqpoint{3.149188in}{0.707769in}}%
\pgfpathlineto{\pgfqpoint{3.149210in}{0.749015in}}%
\pgfpathlineto{\pgfqpoint{3.149543in}{0.749015in}}%
\pgfpathlineto{\pgfqpoint{3.149862in}{0.732517in}}%
\pgfpathlineto{\pgfqpoint{3.150032in}{0.757264in}}%
\pgfpathlineto{\pgfqpoint{3.150380in}{0.749015in}}%
\pgfpathlineto{\pgfqpoint{3.150528in}{0.757264in}}%
\pgfpathlineto{\pgfqpoint{3.150884in}{0.724268in}}%
\pgfpathlineto{\pgfqpoint{3.151032in}{0.732517in}}%
\pgfpathlineto{\pgfqpoint{3.151054in}{0.724268in}}%
\pgfpathlineto{\pgfqpoint{3.151557in}{0.773763in}}%
\pgfpathlineto{\pgfqpoint{3.152061in}{0.749015in}}%
\pgfpathlineto{\pgfqpoint{3.153386in}{0.831507in}}%
\pgfpathlineto{\pgfqpoint{3.153401in}{0.823258in}}%
\pgfpathlineto{\pgfqpoint{3.154571in}{0.848005in}}%
\pgfpathlineto{\pgfqpoint{3.154741in}{0.839756in}}%
\pgfpathlineto{\pgfqpoint{3.155222in}{0.831507in}}%
\pgfpathlineto{\pgfqpoint{3.155245in}{0.872753in}}%
\pgfpathlineto{\pgfqpoint{3.161279in}{1.169722in}}%
\pgfpathlineto{\pgfqpoint{3.161605in}{1.153224in}}%
\pgfpathlineto{\pgfqpoint{3.162619in}{1.120227in}}%
\pgfpathlineto{\pgfqpoint{3.163782in}{1.177971in}}%
\pgfpathlineto{\pgfqpoint{3.164285in}{1.153224in}}%
\pgfpathlineto{\pgfqpoint{3.165300in}{1.095480in}}%
\pgfpathlineto{\pgfqpoint{3.167647in}{0.938746in}}%
\pgfpathlineto{\pgfqpoint{3.169831in}{0.724268in}}%
\pgfpathlineto{\pgfqpoint{3.169994in}{0.749015in}}%
\pgfpathlineto{\pgfqpoint{3.170490in}{0.757264in}}%
\pgfpathlineto{\pgfqpoint{3.170668in}{0.740766in}}%
\pgfpathlineto{\pgfqpoint{3.170831in}{0.749015in}}%
\pgfpathlineto{\pgfqpoint{3.171837in}{0.724268in}}%
\pgfpathlineto{\pgfqpoint{3.172000in}{0.732517in}}%
\pgfpathlineto{\pgfqpoint{3.172341in}{0.699520in}}%
\pgfpathlineto{\pgfqpoint{3.172496in}{0.707769in}}%
\pgfpathlineto{\pgfqpoint{3.173518in}{0.699520in}}%
\pgfpathlineto{\pgfqpoint{3.174340in}{0.707769in}}%
\pgfpathlineto{\pgfqpoint{3.174355in}{0.699520in}}%
\pgfpathlineto{\pgfqpoint{3.176028in}{0.625278in}}%
\pgfpathlineto{\pgfqpoint{3.176199in}{0.650025in}}%
\pgfpathlineto{\pgfqpoint{3.176369in}{0.674773in}}%
\pgfpathlineto{\pgfqpoint{3.176865in}{0.625278in}}%
\pgfpathlineto{\pgfqpoint{3.177857in}{0.658274in}}%
\pgfpathlineto{\pgfqpoint{3.178042in}{0.641776in}}%
\pgfpathlineto{\pgfqpoint{3.178375in}{0.625278in}}%
\pgfpathlineto{\pgfqpoint{3.178694in}{0.658274in}}%
\pgfpathlineto{\pgfqpoint{3.178716in}{0.650025in}}%
\pgfpathlineto{\pgfqpoint{3.179197in}{0.683022in}}%
\pgfpathlineto{\pgfqpoint{3.179553in}{0.650025in}}%
\pgfpathlineto{\pgfqpoint{3.180723in}{0.625278in}}%
\pgfpathlineto{\pgfqpoint{3.180871in}{0.633527in}}%
\pgfpathlineto{\pgfqpoint{3.181559in}{0.699520in}}%
\pgfpathlineto{\pgfqpoint{3.181900in}{0.666524in}}%
\pgfpathlineto{\pgfqpoint{3.183070in}{0.625278in}}%
\pgfpathlineto{\pgfqpoint{3.183240in}{0.650025in}}%
\pgfpathlineto{\pgfqpoint{3.184410in}{0.674773in}}%
\pgfpathlineto{\pgfqpoint{3.183743in}{0.625278in}}%
\pgfpathlineto{\pgfqpoint{3.184580in}{0.666524in}}%
\pgfpathlineto{\pgfqpoint{3.185246in}{0.650025in}}%
\pgfpathlineto{\pgfqpoint{3.185417in}{0.674773in}}%
\pgfpathlineto{\pgfqpoint{3.186068in}{0.683022in}}%
\pgfpathlineto{\pgfqpoint{3.185898in}{0.658274in}}%
\pgfpathlineto{\pgfqpoint{3.186253in}{0.666524in}}%
\pgfpathlineto{\pgfqpoint{3.186572in}{0.683022in}}%
\pgfpathlineto{\pgfqpoint{3.187912in}{0.633527in}}%
\pgfpathlineto{\pgfqpoint{3.187934in}{0.650025in}}%
\pgfpathlineto{\pgfqpoint{3.188267in}{0.625278in}}%
\pgfpathlineto{\pgfqpoint{3.188919in}{0.625278in}}%
\pgfpathlineto{\pgfqpoint{3.190111in}{0.633527in}}%
\pgfpathlineto{\pgfqpoint{3.190259in}{0.625278in}}%
\pgfpathlineto{\pgfqpoint{3.191096in}{0.625278in}}%
\pgfpathlineto{\pgfqpoint{3.194280in}{0.732517in}}%
\pgfpathlineto{\pgfqpoint{3.195138in}{0.724268in}}%
\pgfpathlineto{\pgfqpoint{3.195309in}{0.749015in}}%
\pgfpathlineto{\pgfqpoint{3.198663in}{0.922247in}}%
\pgfpathlineto{\pgfqpoint{3.199322in}{0.905749in}}%
\pgfpathlineto{\pgfqpoint{3.200166in}{0.897500in}}%
\pgfpathlineto{\pgfqpoint{3.200336in}{0.922247in}}%
\pgfpathlineto{\pgfqpoint{3.202854in}{1.021237in}}%
\pgfpathlineto{\pgfqpoint{3.203017in}{1.012988in}}%
\pgfpathlineto{\pgfqpoint{3.203690in}{0.996490in}}%
\pgfpathlineto{\pgfqpoint{3.203187in}{1.021237in}}%
\pgfpathlineto{\pgfqpoint{3.203861in}{1.021237in}}%
\pgfpathlineto{\pgfqpoint{3.205201in}{1.070732in}}%
\pgfpathlineto{\pgfqpoint{3.205519in}{1.054234in}}%
\pgfpathlineto{\pgfqpoint{3.205704in}{1.045985in}}%
\pgfpathlineto{\pgfqpoint{3.206356in}{1.128476in}}%
\pgfpathlineto{\pgfqpoint{3.206371in}{1.120227in}}%
\pgfpathlineto{\pgfqpoint{3.207207in}{1.144975in}}%
\pgfpathlineto{\pgfqpoint{3.207533in}{1.128476in}}%
\pgfpathlineto{\pgfqpoint{3.208214in}{1.095480in}}%
\pgfpathlineto{\pgfqpoint{3.208547in}{1.120227in}}%
\pgfpathlineto{\pgfqpoint{3.209554in}{1.144975in}}%
\pgfpathlineto{\pgfqpoint{3.209725in}{1.136726in}}%
\pgfpathlineto{\pgfqpoint{3.210547in}{1.103729in}}%
\pgfpathlineto{\pgfqpoint{3.209895in}{1.144975in}}%
\pgfpathlineto{\pgfqpoint{3.210880in}{1.128476in}}%
\pgfpathlineto{\pgfqpoint{3.210895in}{1.144975in}}%
\pgfpathlineto{\pgfqpoint{3.211546in}{1.103729in}}%
\pgfpathlineto{\pgfqpoint{3.211902in}{1.120227in}}%
\pgfpathlineto{\pgfqpoint{3.212738in}{1.144975in}}%
\pgfpathlineto{\pgfqpoint{3.213057in}{1.128476in}}%
\pgfpathlineto{\pgfqpoint{3.214249in}{1.070732in}}%
\pgfpathlineto{\pgfqpoint{3.217603in}{0.864503in}}%
\pgfpathlineto{\pgfqpoint{3.220120in}{0.650025in}}%
\pgfpathlineto{\pgfqpoint{3.221460in}{0.749015in}}%
\pgfpathlineto{\pgfqpoint{3.222112in}{0.707769in}}%
\pgfpathlineto{\pgfqpoint{3.222127in}{0.699520in}}%
\pgfpathlineto{\pgfqpoint{3.223134in}{0.749015in}}%
\pgfpathlineto{\pgfqpoint{3.226466in}{0.625278in}}%
\pgfpathlineto{\pgfqpoint{3.228983in}{0.707769in}}%
\pgfpathlineto{\pgfqpoint{3.229338in}{0.674773in}}%
\pgfpathlineto{\pgfqpoint{3.231182in}{0.625278in}}%
\pgfpathlineto{\pgfqpoint{3.232855in}{0.699520in}}%
\pgfpathlineto{\pgfqpoint{3.233174in}{0.683022in}}%
\pgfpathlineto{\pgfqpoint{3.234181in}{0.658274in}}%
\pgfpathlineto{\pgfqpoint{3.234196in}{0.674773in}}%
\pgfpathlineto{\pgfqpoint{3.235188in}{0.658274in}}%
\pgfpathlineto{\pgfqpoint{3.235203in}{0.674773in}}%
\pgfpathlineto{\pgfqpoint{3.236883in}{0.724268in}}%
\pgfpathlineto{\pgfqpoint{3.237046in}{0.716019in}}%
\pgfpathlineto{\pgfqpoint{3.237194in}{0.707769in}}%
\pgfpathlineto{\pgfqpoint{3.237365in}{0.732517in}}%
\pgfpathlineto{\pgfqpoint{3.237720in}{0.724268in}}%
\pgfpathlineto{\pgfqpoint{3.238386in}{0.749015in}}%
\pgfpathlineto{\pgfqpoint{3.238875in}{0.732517in}}%
\pgfpathlineto{\pgfqpoint{3.238890in}{0.724268in}}%
\pgfpathlineto{\pgfqpoint{3.239393in}{0.773763in}}%
\pgfpathlineto{\pgfqpoint{3.239897in}{0.749015in}}%
\pgfpathlineto{\pgfqpoint{3.241222in}{0.831507in}}%
\pgfpathlineto{\pgfqpoint{3.241578in}{0.790261in}}%
\pgfpathlineto{\pgfqpoint{3.242229in}{0.806759in}}%
\pgfpathlineto{\pgfqpoint{3.242747in}{0.773763in}}%
\pgfpathlineto{\pgfqpoint{3.243421in}{0.798510in}}%
\pgfpathlineto{\pgfqpoint{3.243910in}{0.782012in}}%
\pgfpathlineto{\pgfqpoint{3.244924in}{0.773763in}}%
\pgfpathlineto{\pgfqpoint{3.247923in}{0.905749in}}%
\pgfpathlineto{\pgfqpoint{3.248108in}{0.889251in}}%
\pgfpathlineto{\pgfqpoint{3.248219in}{0.881002in}}%
\pgfpathlineto{\pgfqpoint{3.248426in}{0.905749in}}%
\pgfpathlineto{\pgfqpoint{3.248930in}{0.905749in}}%
\pgfpathlineto{\pgfqpoint{3.248952in}{0.922247in}}%
\pgfpathlineto{\pgfqpoint{3.249959in}{0.897500in}}%
\pgfpathlineto{\pgfqpoint{3.250611in}{0.905749in}}%
\pgfpathlineto{\pgfqpoint{3.250796in}{0.889251in}}%
\pgfpathlineto{\pgfqpoint{3.252284in}{0.856254in}}%
\pgfpathlineto{\pgfqpoint{3.250959in}{0.897500in}}%
\pgfpathlineto{\pgfqpoint{3.252284in}{0.864503in}}%
\pgfpathlineto{\pgfqpoint{3.252299in}{0.872753in}}%
\pgfpathlineto{\pgfqpoint{3.253306in}{0.823258in}}%
\pgfpathlineto{\pgfqpoint{3.257334in}{0.650025in}}%
\pgfpathlineto{\pgfqpoint{3.257482in}{0.666524in}}%
\pgfpathlineto{\pgfqpoint{3.257497in}{0.674773in}}%
\pgfpathlineto{\pgfqpoint{3.257837in}{0.650025in}}%
\pgfpathlineto{\pgfqpoint{3.258504in}{0.650025in}}%
\pgfpathlineto{\pgfqpoint{3.259155in}{0.683022in}}%
\pgfpathlineto{\pgfqpoint{3.259340in}{0.666524in}}%
\pgfpathlineto{\pgfqpoint{3.260666in}{0.625278in}}%
\pgfpathlineto{\pgfqpoint{3.260666in}{0.625278in}}%
\pgfpathlineto{\pgfqpoint{3.261673in}{0.633527in}}%
\pgfpathlineto{\pgfqpoint{3.261673in}{0.625278in}}%
\pgfpathlineto{\pgfqpoint{3.262842in}{0.633527in}}%
\pgfpathlineto{\pgfqpoint{3.263183in}{0.625278in}}%
\pgfpathlineto{\pgfqpoint{3.263679in}{0.625278in}}%
\pgfpathlineto{\pgfqpoint{3.264856in}{0.633527in}}%
\pgfpathlineto{\pgfqpoint{3.265027in}{0.625278in}}%
\pgfpathlineto{\pgfqpoint{3.265863in}{0.625278in}}%
\pgfpathlineto{\pgfqpoint{3.267033in}{0.633527in}}%
\pgfpathlineto{\pgfqpoint{3.267204in}{0.625278in}}%
\pgfpathlineto{\pgfqpoint{3.268040in}{0.625278in}}%
\pgfpathlineto{\pgfqpoint{3.269047in}{0.633527in}}%
\pgfpathlineto{\pgfqpoint{3.269062in}{0.625278in}}%
\pgfpathlineto{\pgfqpoint{3.270217in}{0.633527in}}%
\pgfpathlineto{\pgfqpoint{3.270387in}{0.625278in}}%
\pgfpathlineto{\pgfqpoint{3.271224in}{0.625278in}}%
\pgfpathlineto{\pgfqpoint{3.272231in}{0.633527in}}%
\pgfpathlineto{\pgfqpoint{3.272231in}{0.625278in}}%
\pgfpathlineto{\pgfqpoint{3.273408in}{0.633527in}}%
\pgfpathlineto{\pgfqpoint{3.273741in}{0.625278in}}%
\pgfpathlineto{\pgfqpoint{3.274408in}{0.625278in}}%
\pgfpathlineto{\pgfqpoint{3.275585in}{0.633527in}}%
\pgfpathlineto{\pgfqpoint{3.275755in}{0.625278in}}%
\pgfpathlineto{\pgfqpoint{3.276592in}{0.625278in}}%
\pgfpathlineto{\pgfqpoint{3.277762in}{0.633527in}}%
\pgfpathlineto{\pgfqpoint{3.277762in}{0.625278in}}%
\pgfpathlineto{\pgfqpoint{3.278769in}{0.625278in}}%
\pgfpathlineto{\pgfqpoint{3.279606in}{0.666524in}}%
\pgfpathlineto{\pgfqpoint{3.280109in}{0.683022in}}%
\pgfpathlineto{\pgfqpoint{3.280442in}{0.658274in}}%
\pgfpathlineto{\pgfqpoint{3.280635in}{0.666524in}}%
\pgfpathlineto{\pgfqpoint{3.280783in}{0.658274in}}%
\pgfpathlineto{\pgfqpoint{3.281286in}{0.691271in}}%
\pgfpathlineto{\pgfqpoint{3.281301in}{0.699520in}}%
\pgfpathlineto{\pgfqpoint{3.281975in}{0.650025in}}%
\pgfpathlineto{\pgfqpoint{3.282308in}{0.650025in}}%
\pgfpathlineto{\pgfqpoint{3.282456in}{0.658274in}}%
\pgfpathlineto{\pgfqpoint{3.282812in}{0.625278in}}%
\pgfpathlineto{\pgfqpoint{3.282960in}{0.633527in}}%
\pgfpathlineto{\pgfqpoint{3.282982in}{0.625278in}}%
\pgfpathlineto{\pgfqpoint{3.283293in}{0.658274in}}%
\pgfpathlineto{\pgfqpoint{3.283981in}{0.650025in}}%
\pgfpathlineto{\pgfqpoint{3.284826in}{0.699520in}}%
\pgfpathlineto{\pgfqpoint{3.285329in}{0.666524in}}%
\pgfpathlineto{\pgfqpoint{3.285477in}{0.658274in}}%
\pgfpathlineto{\pgfqpoint{3.285825in}{0.699520in}}%
\pgfpathlineto{\pgfqpoint{3.286477in}{0.683022in}}%
\pgfpathlineto{\pgfqpoint{3.287484in}{0.732517in}}%
\pgfpathlineto{\pgfqpoint{3.288676in}{0.699520in}}%
\pgfpathlineto{\pgfqpoint{3.289016in}{0.724268in}}%
\pgfpathlineto{\pgfqpoint{3.289520in}{0.699520in}}%
\pgfpathlineto{\pgfqpoint{3.290860in}{0.674773in}}%
\pgfpathlineto{\pgfqpoint{3.291171in}{0.683022in}}%
\pgfpathlineto{\pgfqpoint{3.291511in}{0.658274in}}%
\pgfpathlineto{\pgfqpoint{3.291697in}{0.666524in}}%
\pgfpathlineto{\pgfqpoint{3.291845in}{0.658274in}}%
\pgfpathlineto{\pgfqpoint{3.292015in}{0.683022in}}%
\pgfpathlineto{\pgfqpoint{3.292363in}{0.674773in}}%
\pgfpathlineto{\pgfqpoint{3.292681in}{0.683022in}}%
\pgfpathlineto{\pgfqpoint{3.293015in}{0.658274in}}%
\pgfpathlineto{\pgfqpoint{3.294207in}{0.625278in}}%
\pgfpathlineto{\pgfqpoint{3.295362in}{0.658274in}}%
\pgfpathlineto{\pgfqpoint{3.295554in}{0.641776in}}%
\pgfpathlineto{\pgfqpoint{3.296035in}{0.633527in}}%
\pgfpathlineto{\pgfqpoint{3.296206in}{0.658274in}}%
\pgfpathlineto{\pgfqpoint{3.296391in}{0.650025in}}%
\pgfpathlineto{\pgfqpoint{3.297042in}{0.683022in}}%
\pgfpathlineto{\pgfqpoint{3.297561in}{0.666524in}}%
\pgfpathlineto{\pgfqpoint{3.298049in}{0.658274in}}%
\pgfpathlineto{\pgfqpoint{3.298212in}{0.683022in}}%
\pgfpathlineto{\pgfqpoint{3.298235in}{0.674773in}}%
\pgfpathlineto{\pgfqpoint{3.298716in}{0.732517in}}%
\pgfpathlineto{\pgfqpoint{3.300078in}{0.716019in}}%
\pgfpathlineto{\pgfqpoint{3.300226in}{0.707769in}}%
\pgfpathlineto{\pgfqpoint{3.300397in}{0.732517in}}%
\pgfpathlineto{\pgfqpoint{3.300745in}{0.724268in}}%
\pgfpathlineto{\pgfqpoint{3.300900in}{0.732517in}}%
\pgfpathlineto{\pgfqpoint{3.301233in}{0.707769in}}%
\pgfpathlineto{\pgfqpoint{3.301418in}{0.716019in}}%
\pgfpathlineto{\pgfqpoint{3.302573in}{0.658274in}}%
\pgfpathlineto{\pgfqpoint{3.301589in}{0.724268in}}%
\pgfpathlineto{\pgfqpoint{3.302966in}{0.683022in}}%
\pgfpathlineto{\pgfqpoint{3.303936in}{0.773763in}}%
\pgfpathlineto{\pgfqpoint{3.304772in}{0.749015in}}%
\pgfpathlineto{\pgfqpoint{3.305424in}{0.732517in}}%
\pgfpathlineto{\pgfqpoint{3.305424in}{0.765513in}}%
\pgfpathlineto{\pgfqpoint{3.306431in}{0.782012in}}%
\pgfpathlineto{\pgfqpoint{3.305757in}{0.757264in}}%
\pgfpathlineto{\pgfqpoint{3.306446in}{0.773763in}}%
\pgfpathlineto{\pgfqpoint{3.306594in}{0.782012in}}%
\pgfpathlineto{\pgfqpoint{3.306949in}{0.749015in}}%
\pgfpathlineto{\pgfqpoint{3.307097in}{0.757264in}}%
\pgfpathlineto{\pgfqpoint{3.307934in}{0.732517in}}%
\pgfpathlineto{\pgfqpoint{3.308127in}{0.749015in}}%
\pgfpathlineto{\pgfqpoint{3.309282in}{0.782012in}}%
\pgfpathlineto{\pgfqpoint{3.309467in}{0.765513in}}%
\pgfpathlineto{\pgfqpoint{3.310637in}{0.749015in}}%
\pgfpathlineto{\pgfqpoint{3.310955in}{0.757264in}}%
\pgfpathlineto{\pgfqpoint{3.311288in}{0.732517in}}%
\pgfpathlineto{\pgfqpoint{3.311473in}{0.740766in}}%
\pgfpathlineto{\pgfqpoint{3.311621in}{0.732517in}}%
\pgfpathlineto{\pgfqpoint{3.311792in}{0.757264in}}%
\pgfpathlineto{\pgfqpoint{3.312147in}{0.749015in}}%
\pgfpathlineto{\pgfqpoint{3.312465in}{0.757264in}}%
\pgfpathlineto{\pgfqpoint{3.312799in}{0.732517in}}%
\pgfpathlineto{\pgfqpoint{3.312984in}{0.740766in}}%
\pgfpathlineto{\pgfqpoint{3.313472in}{0.732517in}}%
\pgfpathlineto{\pgfqpoint{3.313635in}{0.757264in}}%
\pgfpathlineto{\pgfqpoint{3.313657in}{0.749015in}}%
\pgfpathlineto{\pgfqpoint{3.314494in}{0.773763in}}%
\pgfpathlineto{\pgfqpoint{3.314812in}{0.757264in}}%
\pgfpathlineto{\pgfqpoint{3.315331in}{0.724268in}}%
\pgfpathlineto{\pgfqpoint{3.315834in}{0.773763in}}%
\pgfpathlineto{\pgfqpoint{3.316819in}{0.782012in}}%
\pgfpathlineto{\pgfqpoint{3.316841in}{0.773763in}}%
\pgfpathlineto{\pgfqpoint{3.317493in}{0.757264in}}%
\pgfpathlineto{\pgfqpoint{3.317663in}{0.782012in}}%
\pgfpathlineto{\pgfqpoint{3.317678in}{0.773763in}}%
\pgfpathlineto{\pgfqpoint{3.318663in}{0.831507in}}%
\pgfpathlineto{\pgfqpoint{3.319336in}{0.806759in}}%
\pgfpathlineto{\pgfqpoint{3.319522in}{0.798510in}}%
\pgfpathlineto{\pgfqpoint{3.320173in}{0.831507in}}%
\pgfpathlineto{\pgfqpoint{3.320358in}{0.823258in}}%
\pgfpathlineto{\pgfqpoint{3.320677in}{0.831507in}}%
\pgfpathlineto{\pgfqpoint{3.321010in}{0.806759in}}%
\pgfpathlineto{\pgfqpoint{3.321202in}{0.815008in}}%
\pgfpathlineto{\pgfqpoint{3.322202in}{0.773763in}}%
\pgfpathlineto{\pgfqpoint{3.321365in}{0.823258in}}%
\pgfpathlineto{\pgfqpoint{3.322372in}{0.798510in}}%
\pgfpathlineto{\pgfqpoint{3.322691in}{0.806759in}}%
\pgfpathlineto{\pgfqpoint{3.323024in}{0.782012in}}%
\pgfpathlineto{\pgfqpoint{3.323209in}{0.773763in}}%
\pgfpathlineto{\pgfqpoint{3.323860in}{0.856254in}}%
\pgfpathlineto{\pgfqpoint{3.323883in}{0.848005in}}%
\pgfpathlineto{\pgfqpoint{3.324697in}{0.856254in}}%
\pgfpathlineto{\pgfqpoint{3.324719in}{0.848005in}}%
\pgfpathlineto{\pgfqpoint{3.325874in}{0.831507in}}%
\pgfpathlineto{\pgfqpoint{3.325541in}{0.856254in}}%
\pgfpathlineto{\pgfqpoint{3.325874in}{0.839756in}}%
\pgfpathlineto{\pgfqpoint{3.326881in}{0.856254in}}%
\pgfpathlineto{\pgfqpoint{3.326208in}{0.831507in}}%
\pgfpathlineto{\pgfqpoint{3.326896in}{0.848005in}}%
\pgfpathlineto{\pgfqpoint{3.327548in}{0.856254in}}%
\pgfpathlineto{\pgfqpoint{3.327733in}{0.839756in}}%
\pgfpathlineto{\pgfqpoint{3.327888in}{0.831507in}}%
\pgfpathlineto{\pgfqpoint{3.328051in}{0.856254in}}%
\pgfpathlineto{\pgfqpoint{3.328384in}{0.856254in}}%
\pgfpathlineto{\pgfqpoint{3.328407in}{0.872753in}}%
\pgfpathlineto{\pgfqpoint{3.328740in}{0.848005in}}%
\pgfpathlineto{\pgfqpoint{3.329414in}{0.848005in}}%
\pgfpathlineto{\pgfqpoint{3.330732in}{0.831507in}}%
\pgfpathlineto{\pgfqpoint{3.330754in}{0.848005in}}%
\pgfpathlineto{\pgfqpoint{3.331087in}{0.823258in}}%
\pgfpathlineto{\pgfqpoint{3.331761in}{0.823258in}}%
\pgfpathlineto{\pgfqpoint{3.332745in}{0.806759in}}%
\pgfpathlineto{\pgfqpoint{3.332745in}{0.815008in}}%
\pgfpathlineto{\pgfqpoint{3.333775in}{0.872753in}}%
\pgfpathlineto{\pgfqpoint{3.334759in}{0.881002in}}%
\pgfpathlineto{\pgfqpoint{3.334774in}{0.872753in}}%
\pgfpathlineto{\pgfqpoint{3.335426in}{0.831507in}}%
\pgfpathlineto{\pgfqpoint{3.335781in}{0.872753in}}%
\pgfpathlineto{\pgfqpoint{3.335929in}{0.881002in}}%
\pgfpathlineto{\pgfqpoint{3.336285in}{0.848005in}}%
\pgfpathlineto{\pgfqpoint{3.336433in}{0.856254in}}%
\pgfpathlineto{\pgfqpoint{3.337610in}{0.831507in}}%
\pgfpathlineto{\pgfqpoint{3.338632in}{0.872753in}}%
\pgfpathlineto{\pgfqpoint{3.340120in}{0.930497in}}%
\pgfpathlineto{\pgfqpoint{3.340305in}{0.913998in}}%
\pgfpathlineto{\pgfqpoint{3.340461in}{0.905749in}}%
\pgfpathlineto{\pgfqpoint{3.340624in}{0.930497in}}%
\pgfpathlineto{\pgfqpoint{3.340979in}{0.922247in}}%
\pgfpathlineto{\pgfqpoint{3.341127in}{0.930497in}}%
\pgfpathlineto{\pgfqpoint{3.341460in}{0.905749in}}%
\pgfpathlineto{\pgfqpoint{3.341653in}{0.913998in}}%
\pgfpathlineto{\pgfqpoint{3.342652in}{0.872753in}}%
\pgfpathlineto{\pgfqpoint{3.341816in}{0.922247in}}%
\pgfpathlineto{\pgfqpoint{3.342823in}{0.897500in}}%
\pgfpathlineto{\pgfqpoint{3.344481in}{0.979992in}}%
\pgfpathlineto{\pgfqpoint{3.344999in}{0.963493in}}%
\pgfpathlineto{\pgfqpoint{3.345147in}{0.955244in}}%
\pgfpathlineto{\pgfqpoint{3.345651in}{0.988241in}}%
\pgfpathlineto{\pgfqpoint{3.346828in}{1.029487in}}%
\pgfpathlineto{\pgfqpoint{3.346843in}{1.021237in}}%
\pgfpathlineto{\pgfqpoint{3.348168in}{1.054234in}}%
\pgfpathlineto{\pgfqpoint{3.348191in}{1.045985in}}%
\pgfpathlineto{\pgfqpoint{3.348672in}{1.054234in}}%
\pgfpathlineto{\pgfqpoint{3.349531in}{1.021237in}}%
\pgfpathlineto{\pgfqpoint{3.350012in}{1.029487in}}%
\pgfpathlineto{\pgfqpoint{3.350345in}{1.004739in}}%
\pgfpathlineto{\pgfqpoint{3.350367in}{1.021237in}}%
\pgfpathlineto{\pgfqpoint{3.351685in}{1.004739in}}%
\pgfpathlineto{\pgfqpoint{3.351708in}{1.021237in}}%
\pgfpathlineto{\pgfqpoint{3.352715in}{0.971742in}}%
\pgfpathlineto{\pgfqpoint{3.354373in}{1.004739in}}%
\pgfpathlineto{\pgfqpoint{3.354388in}{0.996490in}}%
\pgfpathlineto{\pgfqpoint{3.355225in}{1.095480in}}%
\pgfpathlineto{\pgfqpoint{3.355543in}{1.103729in}}%
\pgfpathlineto{\pgfqpoint{3.355876in}{1.078981in}}%
\pgfpathlineto{\pgfqpoint{3.356905in}{1.070732in}}%
\pgfpathlineto{\pgfqpoint{3.358223in}{1.029487in}}%
\pgfpathlineto{\pgfqpoint{3.358246in}{1.045985in}}%
\pgfpathlineto{\pgfqpoint{3.358897in}{1.078981in}}%
\pgfpathlineto{\pgfqpoint{3.359401in}{1.054234in}}%
\pgfpathlineto{\pgfqpoint{3.359415in}{1.045985in}}%
\pgfpathlineto{\pgfqpoint{3.360067in}{1.103729in}}%
\pgfpathlineto{\pgfqpoint{3.360422in}{1.095480in}}%
\pgfpathlineto{\pgfqpoint{3.362081in}{1.128476in}}%
\pgfpathlineto{\pgfqpoint{3.363273in}{1.095480in}}%
\pgfpathlineto{\pgfqpoint{3.364783in}{1.144975in}}%
\pgfpathlineto{\pgfqpoint{3.364954in}{1.136726in}}%
\pgfpathlineto{\pgfqpoint{3.365435in}{1.128476in}}%
\pgfpathlineto{\pgfqpoint{3.365605in}{1.153224in}}%
\pgfpathlineto{\pgfqpoint{3.365768in}{1.153224in}}%
\pgfpathlineto{\pgfqpoint{3.366124in}{1.194470in}}%
\pgfpathlineto{\pgfqpoint{3.366797in}{1.169722in}}%
\pgfpathlineto{\pgfqpoint{3.370129in}{1.301709in}}%
\pgfpathlineto{\pgfqpoint{3.371151in}{1.293460in}}%
\pgfpathlineto{\pgfqpoint{3.372832in}{1.161473in}}%
\pgfpathlineto{\pgfqpoint{3.377356in}{0.625278in}}%
\pgfpathlineto{\pgfqpoint{3.378526in}{0.650025in}}%
\pgfpathlineto{\pgfqpoint{3.378696in}{0.641776in}}%
\pgfpathlineto{\pgfqpoint{3.379851in}{0.625278in}}%
\pgfpathlineto{\pgfqpoint{3.381524in}{0.658274in}}%
\pgfpathlineto{\pgfqpoint{3.381717in}{0.650025in}}%
\pgfpathlineto{\pgfqpoint{3.382213in}{0.699520in}}%
\pgfpathlineto{\pgfqpoint{3.382531in}{0.691271in}}%
\pgfpathlineto{\pgfqpoint{3.383375in}{0.782012in}}%
\pgfpathlineto{\pgfqpoint{3.383560in}{0.773763in}}%
\pgfpathlineto{\pgfqpoint{3.384716in}{0.806759in}}%
\pgfpathlineto{\pgfqpoint{3.384901in}{0.790261in}}%
\pgfpathlineto{\pgfqpoint{3.385064in}{0.798510in}}%
\pgfpathlineto{\pgfqpoint{3.386070in}{0.773763in}}%
\pgfpathlineto{\pgfqpoint{3.387226in}{0.782012in}}%
\pgfpathlineto{\pgfqpoint{3.387581in}{0.798510in}}%
\pgfpathlineto{\pgfqpoint{3.388255in}{0.773763in}}%
\pgfpathlineto{\pgfqpoint{3.389906in}{0.831507in}}%
\pgfpathlineto{\pgfqpoint{3.390098in}{0.815008in}}%
\pgfpathlineto{\pgfqpoint{3.390913in}{0.806759in}}%
\pgfpathlineto{\pgfqpoint{3.390261in}{0.823258in}}%
\pgfpathlineto{\pgfqpoint{3.390935in}{0.823258in}}%
\pgfpathlineto{\pgfqpoint{3.392142in}{0.856254in}}%
\pgfpathlineto{\pgfqpoint{3.392608in}{0.848005in}}%
\pgfpathlineto{\pgfqpoint{3.392779in}{0.872753in}}%
\pgfpathlineto{\pgfqpoint{3.394437in}{0.930497in}}%
\pgfpathlineto{\pgfqpoint{3.394622in}{0.913998in}}%
\pgfpathlineto{\pgfqpoint{3.395459in}{0.872753in}}%
\pgfpathlineto{\pgfqpoint{3.395777in}{0.905749in}}%
\pgfpathlineto{\pgfqpoint{3.396451in}{0.930497in}}%
\pgfpathlineto{\pgfqpoint{3.396799in}{0.897500in}}%
\pgfpathlineto{\pgfqpoint{3.398310in}{0.946995in}}%
\pgfpathlineto{\pgfqpoint{3.398480in}{0.938746in}}%
\pgfpathlineto{\pgfqpoint{3.399976in}{0.905749in}}%
\pgfpathlineto{\pgfqpoint{3.399976in}{0.913998in}}%
\pgfpathlineto{\pgfqpoint{3.399983in}{0.922247in}}%
\pgfpathlineto{\pgfqpoint{3.400975in}{0.856254in}}%
\pgfpathlineto{\pgfqpoint{3.400990in}{0.872753in}}%
\pgfpathlineto{\pgfqpoint{3.402330in}{0.848005in}}%
\pgfpathlineto{\pgfqpoint{3.404996in}{0.979992in}}%
\pgfpathlineto{\pgfqpoint{3.405514in}{0.963493in}}%
\pgfpathlineto{\pgfqpoint{3.406898in}{0.930497in}}%
\pgfpathlineto{\pgfqpoint{3.405684in}{0.971742in}}%
\pgfpathlineto{\pgfqpoint{3.407010in}{0.938746in}}%
\pgfpathlineto{\pgfqpoint{3.407676in}{0.955244in}}%
\pgfpathlineto{\pgfqpoint{3.408031in}{0.922247in}}%
\pgfpathlineto{\pgfqpoint{3.411052in}{1.045985in}}%
\pgfpathlineto{\pgfqpoint{3.411215in}{1.037736in}}%
\pgfpathlineto{\pgfqpoint{3.411548in}{1.021237in}}%
\pgfpathlineto{\pgfqpoint{3.411719in}{1.045985in}}%
\pgfpathlineto{\pgfqpoint{3.412052in}{1.045985in}}%
\pgfpathlineto{\pgfqpoint{3.412896in}{1.120227in}}%
\pgfpathlineto{\pgfqpoint{3.413895in}{1.087231in}}%
\pgfpathlineto{\pgfqpoint{3.414740in}{1.070732in}}%
\pgfpathlineto{\pgfqpoint{3.414895in}{1.087231in}}%
\pgfpathlineto{\pgfqpoint{3.414902in}{1.095480in}}%
\pgfpathlineto{\pgfqpoint{3.415739in}{0.988241in}}%
\pgfpathlineto{\pgfqpoint{3.418086in}{0.650025in}}%
\pgfpathlineto{\pgfqpoint{3.418930in}{0.674773in}}%
\pgfpathlineto{\pgfqpoint{3.420759in}{0.625278in}}%
\pgfpathlineto{\pgfqpoint{3.421929in}{0.633527in}}%
\pgfpathlineto{\pgfqpoint{3.421929in}{0.625278in}}%
\pgfpathlineto{\pgfqpoint{3.422929in}{0.625278in}}%
\pgfpathlineto{\pgfqpoint{3.423936in}{0.633527in}}%
\pgfpathlineto{\pgfqpoint{3.423936in}{0.625278in}}%
\pgfpathlineto{\pgfqpoint{3.425276in}{0.633527in}}%
\pgfpathlineto{\pgfqpoint{3.425616in}{0.625278in}}%
\pgfpathlineto{\pgfqpoint{3.426127in}{0.625278in}}%
\pgfpathlineto{\pgfqpoint{3.427460in}{0.633527in}}%
\pgfpathlineto{\pgfqpoint{3.427460in}{0.625278in}}%
\pgfpathlineto{\pgfqpoint{3.428474in}{0.625278in}}%
\pgfpathlineto{\pgfqpoint{3.429637in}{0.633527in}}%
\pgfpathlineto{\pgfqpoint{3.429637in}{0.625278in}}%
\pgfpathlineto{\pgfqpoint{3.430644in}{0.625278in}}%
\pgfpathlineto{\pgfqpoint{3.431821in}{0.633527in}}%
\pgfpathlineto{\pgfqpoint{3.431984in}{0.625278in}}%
\pgfpathlineto{\pgfqpoint{3.432843in}{0.625278in}}%
\pgfpathlineto{\pgfqpoint{3.434005in}{0.633527in}}%
\pgfpathlineto{\pgfqpoint{3.434686in}{0.625278in}}%
\pgfpathlineto{\pgfqpoint{3.435005in}{0.658274in}}%
\pgfpathlineto{\pgfqpoint{3.435020in}{0.650025in}}%
\pgfpathlineto{\pgfqpoint{3.435523in}{0.674773in}}%
\pgfpathlineto{\pgfqpoint{3.435856in}{0.650025in}}%
\pgfpathlineto{\pgfqpoint{3.436175in}{0.633527in}}%
\pgfpathlineto{\pgfqpoint{3.436515in}{0.666524in}}%
\pgfpathlineto{\pgfqpoint{3.437700in}{0.773763in}}%
\pgfpathlineto{\pgfqpoint{3.439040in}{0.823258in}}%
\pgfpathlineto{\pgfqpoint{3.439373in}{0.806759in}}%
\pgfpathlineto{\pgfqpoint{3.440388in}{0.773763in}}%
\pgfpathlineto{\pgfqpoint{3.445748in}{1.021237in}}%
\pgfpathlineto{\pgfqpoint{3.445919in}{1.012988in}}%
\pgfpathlineto{\pgfqpoint{3.446407in}{1.004739in}}%
\pgfpathlineto{\pgfqpoint{3.446415in}{1.037736in}}%
\pgfpathlineto{\pgfqpoint{3.446422in}{1.045985in}}%
\pgfpathlineto{\pgfqpoint{3.446755in}{1.021237in}}%
\pgfpathlineto{\pgfqpoint{3.447422in}{1.021237in}}%
\pgfpathlineto{\pgfqpoint{3.447577in}{1.029487in}}%
\pgfpathlineto{\pgfqpoint{3.447910in}{1.004739in}}%
\pgfpathlineto{\pgfqpoint{3.448095in}{1.012988in}}%
\pgfpathlineto{\pgfqpoint{3.448917in}{1.004739in}}%
\pgfpathlineto{\pgfqpoint{3.448266in}{1.021237in}}%
\pgfpathlineto{\pgfqpoint{3.448932in}{1.021237in}}%
\pgfpathlineto{\pgfqpoint{3.451427in}{1.103729in}}%
\pgfpathlineto{\pgfqpoint{3.451612in}{1.095480in}}%
\pgfpathlineto{\pgfqpoint{3.452271in}{1.153224in}}%
\pgfpathlineto{\pgfqpoint{3.452286in}{1.144975in}}%
\pgfpathlineto{\pgfqpoint{3.453278in}{1.202719in}}%
\pgfpathlineto{\pgfqpoint{3.453626in}{1.161473in}}%
\pgfpathlineto{\pgfqpoint{3.453774in}{1.153224in}}%
\pgfpathlineto{\pgfqpoint{3.454130in}{1.194470in}}%
\pgfpathlineto{\pgfqpoint{3.457484in}{1.070732in}}%
\pgfpathlineto{\pgfqpoint{3.457973in}{1.078981in}}%
\pgfpathlineto{\pgfqpoint{3.458150in}{1.062483in}}%
\pgfpathlineto{\pgfqpoint{3.459328in}{1.045985in}}%
\pgfpathlineto{\pgfqpoint{3.459483in}{1.054234in}}%
\pgfpathlineto{\pgfqpoint{3.459661in}{1.037736in}}%
\pgfpathlineto{\pgfqpoint{3.463015in}{0.699520in}}%
\pgfpathlineto{\pgfqpoint{3.463341in}{0.707769in}}%
\pgfpathlineto{\pgfqpoint{3.463518in}{0.691271in}}%
\pgfpathlineto{\pgfqpoint{3.464859in}{0.650025in}}%
\pgfpathlineto{\pgfqpoint{3.466184in}{0.625278in}}%
\pgfpathlineto{\pgfqpoint{3.466184in}{0.625278in}}%
\pgfpathlineto{\pgfqpoint{3.467354in}{0.633527in}}%
\pgfpathlineto{\pgfqpoint{3.468213in}{0.625278in}}%
\pgfpathlineto{\pgfqpoint{3.468376in}{0.650025in}}%
\pgfpathlineto{\pgfqpoint{3.468872in}{0.658274in}}%
\pgfpathlineto{\pgfqpoint{3.469086in}{0.633527in}}%
\pgfpathlineto{\pgfqpoint{3.469220in}{0.650025in}}%
\pgfpathlineto{\pgfqpoint{3.470538in}{0.625278in}}%
\pgfpathlineto{\pgfqpoint{3.471715in}{0.633527in}}%
\pgfpathlineto{\pgfqpoint{3.472218in}{0.625278in}}%
\pgfpathlineto{\pgfqpoint{3.472551in}{0.625278in}}%
\pgfpathlineto{\pgfqpoint{3.473892in}{0.633527in}}%
\pgfpathlineto{\pgfqpoint{3.474062in}{0.625278in}}%
\pgfpathlineto{\pgfqpoint{3.474899in}{0.625278in}}%
\pgfpathlineto{\pgfqpoint{3.476068in}{0.633527in}}%
\pgfpathlineto{\pgfqpoint{3.476239in}{0.625278in}}%
\pgfpathlineto{\pgfqpoint{3.477083in}{0.625278in}}%
\pgfpathlineto{\pgfqpoint{3.478423in}{0.633527in}}%
\pgfpathlineto{\pgfqpoint{3.478423in}{0.625278in}}%
\pgfpathlineto{\pgfqpoint{3.479423in}{0.625278in}}%
\pgfpathlineto{\pgfqpoint{3.480600in}{0.633527in}}%
\pgfpathlineto{\pgfqpoint{3.480763in}{0.625278in}}%
\pgfpathlineto{\pgfqpoint{3.481599in}{0.625278in}}%
\pgfpathlineto{\pgfqpoint{3.482947in}{0.633527in}}%
\pgfpathlineto{\pgfqpoint{3.482947in}{0.625278in}}%
\pgfpathlineto{\pgfqpoint{3.483947in}{0.625278in}}%
\pgfpathlineto{\pgfqpoint{3.485124in}{0.633527in}}%
\pgfpathlineto{\pgfqpoint{3.485294in}{0.625278in}}%
\pgfpathlineto{\pgfqpoint{3.486131in}{0.625278in}}%
\pgfpathlineto{\pgfqpoint{3.487301in}{0.633527in}}%
\pgfpathlineto{\pgfqpoint{3.487471in}{0.625278in}}%
\pgfpathlineto{\pgfqpoint{3.488308in}{0.625278in}}%
\pgfpathlineto{\pgfqpoint{3.489485in}{0.633527in}}%
\pgfpathlineto{\pgfqpoint{3.489648in}{0.625278in}}%
\pgfpathlineto{\pgfqpoint{3.490322in}{0.625278in}}%
\pgfpathlineto{\pgfqpoint{3.491662in}{0.633527in}}%
\pgfpathlineto{\pgfqpoint{3.491825in}{0.625278in}}%
\pgfpathlineto{\pgfqpoint{3.492676in}{0.625278in}}%
\pgfpathlineto{\pgfqpoint{3.493676in}{0.633527in}}%
\pgfpathlineto{\pgfqpoint{3.493676in}{0.625278in}}%
\pgfpathlineto{\pgfqpoint{3.494845in}{0.633527in}}%
\pgfpathlineto{\pgfqpoint{3.495016in}{0.625278in}}%
\pgfpathlineto{\pgfqpoint{3.495690in}{0.625278in}}%
\pgfpathlineto{\pgfqpoint{3.496859in}{0.633527in}}%
\pgfpathlineto{\pgfqpoint{3.497022in}{0.625278in}}%
\pgfpathlineto{\pgfqpoint{3.497866in}{0.625278in}}%
\pgfpathlineto{\pgfqpoint{3.499036in}{0.633527in}}%
\pgfpathlineto{\pgfqpoint{3.499036in}{0.625278in}}%
\pgfpathlineto{\pgfqpoint{3.499873in}{0.625278in}}%
\pgfpathlineto{\pgfqpoint{3.501213in}{0.633527in}}%
\pgfpathlineto{\pgfqpoint{3.501213in}{0.625278in}}%
\pgfpathlineto{\pgfqpoint{3.502220in}{0.625278in}}%
\pgfpathlineto{\pgfqpoint{3.503397in}{0.633527in}}%
\pgfpathlineto{\pgfqpoint{3.503397in}{0.625278in}}%
\pgfpathlineto{\pgfqpoint{3.504397in}{0.625278in}}%
\pgfpathlineto{\pgfqpoint{3.505574in}{0.633527in}}%
\pgfpathlineto{\pgfqpoint{3.505744in}{0.625278in}}%
\pgfpathlineto{\pgfqpoint{3.506581in}{0.625278in}}%
\pgfpathlineto{\pgfqpoint{3.507751in}{0.633527in}}%
\pgfpathlineto{\pgfqpoint{3.507751in}{0.625278in}}%
\pgfpathlineto{\pgfqpoint{3.508758in}{0.625278in}}%
\pgfpathlineto{\pgfqpoint{3.509935in}{0.633527in}}%
\pgfpathlineto{\pgfqpoint{3.510098in}{0.625278in}}%
\pgfpathlineto{\pgfqpoint{3.510935in}{0.625278in}}%
\pgfpathlineto{\pgfqpoint{3.512112in}{0.633527in}}%
\pgfpathlineto{\pgfqpoint{3.512112in}{0.625278in}}%
\pgfpathlineto{\pgfqpoint{3.513119in}{0.625278in}}%
\pgfpathlineto{\pgfqpoint{3.514289in}{0.633527in}}%
\pgfpathlineto{\pgfqpoint{3.514459in}{0.625278in}}%
\pgfpathlineto{\pgfqpoint{3.515296in}{0.625278in}}%
\pgfpathlineto{\pgfqpoint{3.516473in}{0.633527in}}%
\pgfpathlineto{\pgfqpoint{3.516636in}{0.625278in}}%
\pgfpathlineto{\pgfqpoint{3.517473in}{0.625278in}}%
\pgfpathlineto{\pgfqpoint{3.518480in}{0.633527in}}%
\pgfpathlineto{\pgfqpoint{3.518480in}{0.625278in}}%
\pgfpathlineto{\pgfqpoint{3.519820in}{0.633527in}}%
\pgfpathlineto{\pgfqpoint{3.519820in}{0.625278in}}%
\pgfpathlineto{\pgfqpoint{3.520827in}{0.625278in}}%
\pgfpathlineto{\pgfqpoint{3.522004in}{0.633527in}}%
\pgfpathlineto{\pgfqpoint{3.522167in}{0.625278in}}%
\pgfpathlineto{\pgfqpoint{3.523011in}{0.625278in}}%
\pgfpathlineto{\pgfqpoint{3.524011in}{0.633527in}}%
\pgfpathlineto{\pgfqpoint{3.524011in}{0.625278in}}%
\pgfpathlineto{\pgfqpoint{3.525351in}{0.633527in}}%
\pgfpathlineto{\pgfqpoint{3.525691in}{0.625278in}}%
\pgfpathlineto{\pgfqpoint{3.526358in}{0.625278in}}%
\pgfpathlineto{\pgfqpoint{3.527535in}{0.633527in}}%
\pgfpathlineto{\pgfqpoint{3.527698in}{0.625278in}}%
\pgfpathlineto{\pgfqpoint{3.528379in}{0.625278in}}%
\pgfpathlineto{\pgfqpoint{3.529542in}{0.633527in}}%
\pgfpathlineto{\pgfqpoint{3.529712in}{0.625278in}}%
\pgfpathlineto{\pgfqpoint{3.530386in}{0.625278in}}%
\pgfpathlineto{\pgfqpoint{3.531555in}{0.633527in}}%
\pgfpathlineto{\pgfqpoint{3.531555in}{0.625278in}}%
\pgfpathlineto{\pgfqpoint{3.532562in}{0.625278in}}%
\pgfpathlineto{\pgfqpoint{3.533732in}{0.633527in}}%
\pgfpathlineto{\pgfqpoint{3.533755in}{0.625278in}}%
\pgfpathlineto{\pgfqpoint{3.534739in}{0.625278in}}%
\pgfpathlineto{\pgfqpoint{3.535917in}{0.633527in}}%
\pgfpathlineto{\pgfqpoint{3.536250in}{0.625278in}}%
\pgfpathlineto{\pgfqpoint{3.536938in}{0.625278in}}%
\pgfpathlineto{\pgfqpoint{3.538093in}{0.633527in}}%
\pgfpathlineto{\pgfqpoint{3.539123in}{0.625278in}}%
\pgfpathlineto{\pgfqpoint{3.540270in}{0.633527in}}%
\pgfpathlineto{\pgfqpoint{3.540292in}{0.625278in}}%
\pgfpathlineto{\pgfqpoint{3.540944in}{0.658274in}}%
\pgfpathlineto{\pgfqpoint{3.541299in}{0.650025in}}%
\pgfpathlineto{\pgfqpoint{3.541951in}{0.633527in}}%
\pgfpathlineto{\pgfqpoint{3.542114in}{0.658274in}}%
\pgfpathlineto{\pgfqpoint{3.542136in}{0.650025in}}%
\pgfpathlineto{\pgfqpoint{3.542617in}{0.658274in}}%
\pgfpathlineto{\pgfqpoint{3.542810in}{0.641776in}}%
\pgfpathlineto{\pgfqpoint{3.543980in}{0.625278in}}%
\pgfpathlineto{\pgfqpoint{3.545135in}{0.633527in}}%
\pgfpathlineto{\pgfqpoint{3.545320in}{0.625278in}}%
\pgfpathlineto{\pgfqpoint{3.545490in}{0.650025in}}%
\pgfpathlineto{\pgfqpoint{3.546157in}{0.650025in}}%
\pgfpathlineto{\pgfqpoint{3.547164in}{0.674773in}}%
\pgfpathlineto{\pgfqpoint{3.547334in}{0.666524in}}%
\pgfpathlineto{\pgfqpoint{3.547504in}{0.674773in}}%
\pgfpathlineto{\pgfqpoint{3.549007in}{0.625278in}}%
\pgfpathlineto{\pgfqpoint{3.550162in}{0.658274in}}%
\pgfpathlineto{\pgfqpoint{3.550347in}{0.641776in}}%
\pgfpathlineto{\pgfqpoint{3.551021in}{0.650025in}}%
\pgfpathlineto{\pgfqpoint{3.551525in}{0.625278in}}%
\pgfpathlineto{\pgfqpoint{3.552509in}{0.633527in}}%
\pgfpathlineto{\pgfqpoint{3.552532in}{0.625278in}}%
\pgfpathlineto{\pgfqpoint{3.553183in}{0.666524in}}%
\pgfpathlineto{\pgfqpoint{3.553872in}{0.724268in}}%
\pgfpathlineto{\pgfqpoint{3.554353in}{0.707769in}}%
\pgfpathlineto{\pgfqpoint{3.555382in}{0.699520in}}%
\pgfpathlineto{\pgfqpoint{3.555863in}{0.683022in}}%
\pgfpathlineto{\pgfqpoint{3.556537in}{0.707769in}}%
\pgfpathlineto{\pgfqpoint{3.556870in}{0.732517in}}%
\pgfpathlineto{\pgfqpoint{3.557559in}{0.699520in}}%
\pgfpathlineto{\pgfqpoint{3.559218in}{0.732517in}}%
\pgfpathlineto{\pgfqpoint{3.559573in}{0.724268in}}%
\pgfpathlineto{\pgfqpoint{3.560076in}{0.773763in}}%
\pgfpathlineto{\pgfqpoint{3.560239in}{0.765513in}}%
\pgfpathlineto{\pgfqpoint{3.561920in}{0.724268in}}%
\pgfpathlineto{\pgfqpoint{3.562757in}{0.773763in}}%
\pgfpathlineto{\pgfqpoint{3.563571in}{0.757264in}}%
\pgfpathlineto{\pgfqpoint{3.564600in}{0.749015in}}%
\pgfpathlineto{\pgfqpoint{3.567096in}{0.856254in}}%
\pgfpathlineto{\pgfqpoint{3.567281in}{0.839756in}}%
\pgfpathlineto{\pgfqpoint{3.568954in}{0.798510in}}%
\pgfpathlineto{\pgfqpoint{3.570465in}{0.848005in}}%
\pgfpathlineto{\pgfqpoint{3.570635in}{0.839756in}}%
\pgfpathlineto{\pgfqpoint{3.570783in}{0.831507in}}%
\pgfpathlineto{\pgfqpoint{3.570953in}{0.856254in}}%
\pgfpathlineto{\pgfqpoint{3.571472in}{0.848005in}}%
\pgfpathlineto{\pgfqpoint{3.571620in}{0.856254in}}%
\pgfpathlineto{\pgfqpoint{3.571953in}{0.831507in}}%
\pgfpathlineto{\pgfqpoint{3.572145in}{0.839756in}}%
\pgfpathlineto{\pgfqpoint{3.572293in}{0.831507in}}%
\pgfpathlineto{\pgfqpoint{3.572797in}{0.864503in}}%
\pgfpathlineto{\pgfqpoint{3.573130in}{0.856254in}}%
\pgfpathlineto{\pgfqpoint{3.573819in}{0.872753in}}%
\pgfpathlineto{\pgfqpoint{3.574492in}{0.848005in}}%
\pgfpathlineto{\pgfqpoint{3.574655in}{0.872753in}}%
\pgfpathlineto{\pgfqpoint{3.575144in}{0.881002in}}%
\pgfpathlineto{\pgfqpoint{3.575329in}{0.864503in}}%
\pgfpathlineto{\pgfqpoint{3.576144in}{0.831507in}}%
\pgfpathlineto{\pgfqpoint{3.576484in}{0.856254in}}%
\pgfpathlineto{\pgfqpoint{3.576840in}{0.897500in}}%
\pgfpathlineto{\pgfqpoint{3.577506in}{0.872753in}}%
\pgfpathlineto{\pgfqpoint{3.579164in}{0.905749in}}%
\pgfpathlineto{\pgfqpoint{3.580023in}{0.897500in}}%
\pgfpathlineto{\pgfqpoint{3.580186in}{0.922247in}}%
\pgfpathlineto{\pgfqpoint{3.581845in}{0.979992in}}%
\pgfpathlineto{\pgfqpoint{3.582030in}{0.963493in}}%
\pgfpathlineto{\pgfqpoint{3.582178in}{0.955244in}}%
\pgfpathlineto{\pgfqpoint{3.582533in}{0.996490in}}%
\pgfpathlineto{\pgfqpoint{3.583022in}{1.004739in}}%
\pgfpathlineto{\pgfqpoint{3.583207in}{0.988241in}}%
\pgfpathlineto{\pgfqpoint{3.584377in}{0.971742in}}%
\pgfpathlineto{\pgfqpoint{3.585199in}{1.004739in}}%
\pgfpathlineto{\pgfqpoint{3.585384in}{0.988241in}}%
\pgfpathlineto{\pgfqpoint{3.586872in}{0.930497in}}%
\pgfpathlineto{\pgfqpoint{3.587376in}{0.979992in}}%
\pgfpathlineto{\pgfqpoint{3.587901in}{0.963493in}}%
\pgfpathlineto{\pgfqpoint{3.588064in}{0.971742in}}%
\pgfpathlineto{\pgfqpoint{3.589071in}{0.946995in}}%
\pgfpathlineto{\pgfqpoint{3.590730in}{0.979992in}}%
\pgfpathlineto{\pgfqpoint{3.591085in}{0.971742in}}%
\pgfpathlineto{\pgfqpoint{3.591404in}{1.004739in}}%
\pgfpathlineto{\pgfqpoint{3.591752in}{0.996490in}}%
\pgfpathlineto{\pgfqpoint{3.592907in}{1.004739in}}%
\pgfpathlineto{\pgfqpoint{3.593603in}{0.996490in}}%
\pgfpathlineto{\pgfqpoint{3.593766in}{1.021237in}}%
\pgfpathlineto{\pgfqpoint{3.593936in}{1.012988in}}%
\pgfpathlineto{\pgfqpoint{3.594099in}{1.021237in}}%
\pgfpathlineto{\pgfqpoint{3.595091in}{1.004739in}}%
\pgfpathlineto{\pgfqpoint{3.595254in}{1.029487in}}%
\pgfpathlineto{\pgfqpoint{3.596113in}{1.021237in}}%
\pgfpathlineto{\pgfqpoint{3.596261in}{1.029487in}}%
\pgfpathlineto{\pgfqpoint{3.596616in}{0.996490in}}%
\pgfpathlineto{\pgfqpoint{3.596764in}{1.004739in}}%
\pgfpathlineto{\pgfqpoint{3.597623in}{0.996490in}}%
\pgfpathlineto{\pgfqpoint{3.597793in}{1.021237in}}%
\pgfpathlineto{\pgfqpoint{3.598104in}{1.029487in}}%
\pgfpathlineto{\pgfqpoint{3.598460in}{0.996490in}}%
\pgfpathlineto{\pgfqpoint{3.598778in}{0.979992in}}%
\pgfpathlineto{\pgfqpoint{3.598941in}{1.004739in}}%
\pgfpathlineto{\pgfqpoint{3.599296in}{0.996490in}}%
\pgfpathlineto{\pgfqpoint{3.599785in}{0.979992in}}%
\pgfpathlineto{\pgfqpoint{3.600452in}{1.004739in}}%
\pgfpathlineto{\pgfqpoint{3.601310in}{0.996490in}}%
\pgfpathlineto{\pgfqpoint{3.601481in}{1.021237in}}%
\pgfpathlineto{\pgfqpoint{3.602295in}{1.054234in}}%
\pgfpathlineto{\pgfqpoint{3.602636in}{1.029487in}}%
\pgfpathlineto{\pgfqpoint{3.602821in}{1.021237in}}%
\pgfpathlineto{\pgfqpoint{3.603132in}{1.054234in}}%
\pgfpathlineto{\pgfqpoint{3.603658in}{1.045985in}}%
\pgfpathlineto{\pgfqpoint{3.603976in}{1.054234in}}%
\pgfpathlineto{\pgfqpoint{3.604324in}{1.021237in}}%
\pgfpathlineto{\pgfqpoint{3.605331in}{1.045985in}}%
\pgfpathlineto{\pgfqpoint{3.605501in}{1.037736in}}%
\pgfpathlineto{\pgfqpoint{3.605671in}{1.045985in}}%
\pgfpathlineto{\pgfqpoint{3.606656in}{1.029487in}}%
\pgfpathlineto{\pgfqpoint{3.607160in}{1.078981in}}%
\pgfpathlineto{\pgfqpoint{3.607678in}{1.070732in}}%
\pgfpathlineto{\pgfqpoint{3.607826in}{1.078981in}}%
\pgfpathlineto{\pgfqpoint{3.608182in}{1.045985in}}%
\pgfpathlineto{\pgfqpoint{3.609003in}{1.054234in}}%
\pgfpathlineto{\pgfqpoint{3.609522in}{1.021237in}}%
\pgfpathlineto{\pgfqpoint{3.610010in}{1.029487in}}%
\pgfpathlineto{\pgfqpoint{3.610195in}{1.012988in}}%
\pgfpathlineto{\pgfqpoint{3.611365in}{0.996490in}}%
\pgfpathlineto{\pgfqpoint{3.613357in}{1.054234in}}%
\pgfpathlineto{\pgfqpoint{3.614386in}{1.045985in}}%
\pgfpathlineto{\pgfqpoint{3.616060in}{1.095480in}}%
\pgfpathlineto{\pgfqpoint{3.616230in}{1.087231in}}%
\pgfpathlineto{\pgfqpoint{3.616400in}{1.095480in}}%
\pgfpathlineto{\pgfqpoint{3.617888in}{1.054234in}}%
\pgfpathlineto{\pgfqpoint{3.617903in}{1.070732in}}%
\pgfpathlineto{\pgfqpoint{3.618577in}{1.045985in}}%
\pgfpathlineto{\pgfqpoint{3.618910in}{1.045985in}}%
\pgfpathlineto{\pgfqpoint{3.620250in}{1.120227in}}%
\pgfpathlineto{\pgfqpoint{3.620902in}{1.103729in}}%
\pgfpathlineto{\pgfqpoint{3.622094in}{1.070732in}}%
\pgfpathlineto{\pgfqpoint{3.623086in}{1.078981in}}%
\pgfpathlineto{\pgfqpoint{3.623101in}{1.070732in}}%
\pgfpathlineto{\pgfqpoint{3.623419in}{1.054234in}}%
\pgfpathlineto{\pgfqpoint{3.623753in}{1.087231in}}%
\pgfpathlineto{\pgfqpoint{3.624782in}{1.144975in}}%
\pgfpathlineto{\pgfqpoint{3.624945in}{1.136726in}}%
\pgfpathlineto{\pgfqpoint{3.625448in}{1.120227in}}%
\pgfpathlineto{\pgfqpoint{3.625618in}{1.144975in}}%
\pgfpathlineto{\pgfqpoint{3.626100in}{1.153224in}}%
\pgfpathlineto{\pgfqpoint{3.626433in}{1.128476in}}%
\pgfpathlineto{\pgfqpoint{3.626788in}{1.120227in}}%
\pgfpathlineto{\pgfqpoint{3.627107in}{1.153224in}}%
\pgfpathlineto{\pgfqpoint{3.627462in}{1.144975in}}%
\pgfpathlineto{\pgfqpoint{3.628617in}{1.153224in}}%
\pgfpathlineto{\pgfqpoint{3.629639in}{1.144975in}}%
\pgfpathlineto{\pgfqpoint{3.630313in}{1.062483in}}%
\pgfpathlineto{\pgfqpoint{3.634000in}{0.650025in}}%
\pgfpathlineto{\pgfqpoint{3.634311in}{0.658274in}}%
\pgfpathlineto{\pgfqpoint{3.634814in}{0.633527in}}%
\pgfpathlineto{\pgfqpoint{3.634985in}{0.625278in}}%
\pgfpathlineto{\pgfqpoint{3.635673in}{0.625278in}}%
\pgfpathlineto{\pgfqpoint{3.636999in}{0.633527in}}%
\pgfpathlineto{\pgfqpoint{3.637162in}{0.625278in}}%
\pgfpathlineto{\pgfqpoint{3.638020in}{0.625278in}}%
\pgfpathlineto{\pgfqpoint{3.639175in}{0.633527in}}%
\pgfpathlineto{\pgfqpoint{3.639346in}{0.625278in}}%
\pgfpathlineto{\pgfqpoint{3.640034in}{0.650025in}}%
\pgfpathlineto{\pgfqpoint{3.640197in}{0.641776in}}%
\pgfpathlineto{\pgfqpoint{3.640368in}{0.650025in}}%
\pgfpathlineto{\pgfqpoint{3.641375in}{0.625278in}}%
\pgfpathlineto{\pgfqpoint{3.642530in}{0.633527in}}%
\pgfpathlineto{\pgfqpoint{3.643048in}{0.625278in}}%
\pgfpathlineto{\pgfqpoint{3.643366in}{0.658274in}}%
\pgfpathlineto{\pgfqpoint{3.643551in}{0.650025in}}%
\pgfpathlineto{\pgfqpoint{3.644040in}{0.633527in}}%
\pgfpathlineto{\pgfqpoint{3.644706in}{0.658274in}}%
\pgfpathlineto{\pgfqpoint{3.645884in}{0.633527in}}%
\pgfpathlineto{\pgfqpoint{3.646891in}{0.658274in}}%
\pgfpathlineto{\pgfqpoint{3.646232in}{0.625278in}}%
\pgfpathlineto{\pgfqpoint{3.646905in}{0.650025in}}%
\pgfpathlineto{\pgfqpoint{3.649586in}{0.773763in}}%
\pgfpathlineto{\pgfqpoint{3.649904in}{0.757264in}}%
\pgfpathlineto{\pgfqpoint{3.650926in}{0.749015in}}%
\pgfpathlineto{\pgfqpoint{3.653095in}{0.831507in}}%
\pgfpathlineto{\pgfqpoint{3.653110in}{0.823258in}}%
\pgfpathlineto{\pgfqpoint{3.654613in}{0.773763in}}%
\pgfpathlineto{\pgfqpoint{3.654776in}{0.790261in}}%
\pgfpathlineto{\pgfqpoint{3.655790in}{0.823258in}}%
\pgfpathlineto{\pgfqpoint{3.656116in}{0.806759in}}%
\pgfpathlineto{\pgfqpoint{3.656272in}{0.831507in}}%
\pgfpathlineto{\pgfqpoint{3.656627in}{0.823258in}}%
\pgfpathlineto{\pgfqpoint{3.658619in}{0.905749in}}%
\pgfpathlineto{\pgfqpoint{3.658804in}{0.889251in}}%
\pgfpathlineto{\pgfqpoint{3.659145in}{0.872753in}}%
\pgfpathlineto{\pgfqpoint{3.659456in}{0.905749in}}%
\pgfpathlineto{\pgfqpoint{3.659478in}{0.897500in}}%
\pgfpathlineto{\pgfqpoint{3.661136in}{0.930497in}}%
\pgfpathlineto{\pgfqpoint{3.661151in}{0.922247in}}%
\pgfpathlineto{\pgfqpoint{3.661655in}{0.971742in}}%
\pgfpathlineto{\pgfqpoint{3.662158in}{0.946995in}}%
\pgfpathlineto{\pgfqpoint{3.662832in}{1.021237in}}%
\pgfpathlineto{\pgfqpoint{3.664157in}{1.004739in}}%
\pgfpathlineto{\pgfqpoint{3.665179in}{0.971742in}}%
\pgfpathlineto{\pgfqpoint{3.665334in}{0.979992in}}%
\pgfpathlineto{\pgfqpoint{3.665512in}{0.963493in}}%
\pgfpathlineto{\pgfqpoint{3.667186in}{0.872753in}}%
\pgfpathlineto{\pgfqpoint{3.668193in}{0.897500in}}%
\pgfpathlineto{\pgfqpoint{3.668363in}{0.889251in}}%
\pgfpathlineto{\pgfqpoint{3.668851in}{0.881002in}}%
\pgfpathlineto{\pgfqpoint{3.669014in}{0.905749in}}%
\pgfpathlineto{\pgfqpoint{3.669037in}{0.897500in}}%
\pgfpathlineto{\pgfqpoint{3.671702in}{0.979992in}}%
\pgfpathlineto{\pgfqpoint{3.671717in}{0.971742in}}%
\pgfpathlineto{\pgfqpoint{3.671880in}{0.996490in}}%
\pgfpathlineto{\pgfqpoint{3.672724in}{0.988241in}}%
\pgfpathlineto{\pgfqpoint{3.672872in}{0.979992in}}%
\pgfpathlineto{\pgfqpoint{3.673213in}{1.012988in}}%
\pgfpathlineto{\pgfqpoint{3.673227in}{1.045985in}}%
\pgfpathlineto{\pgfqpoint{3.673894in}{0.996490in}}%
\pgfpathlineto{\pgfqpoint{3.674227in}{0.996490in}}%
\pgfpathlineto{\pgfqpoint{3.674545in}{1.004739in}}%
\pgfpathlineto{\pgfqpoint{3.674879in}{0.979992in}}%
\pgfpathlineto{\pgfqpoint{3.675219in}{0.955244in}}%
\pgfpathlineto{\pgfqpoint{3.675567in}{0.996490in}}%
\pgfpathlineto{\pgfqpoint{3.675908in}{0.971742in}}%
\pgfpathlineto{\pgfqpoint{3.676574in}{0.889251in}}%
\pgfpathlineto{\pgfqpoint{3.678085in}{0.823258in}}%
\pgfpathlineto{\pgfqpoint{3.678233in}{0.831507in}}%
\pgfpathlineto{\pgfqpoint{3.678566in}{0.806759in}}%
\pgfpathlineto{\pgfqpoint{3.678758in}{0.815008in}}%
\pgfpathlineto{\pgfqpoint{3.680247in}{0.782012in}}%
\pgfpathlineto{\pgfqpoint{3.678921in}{0.823258in}}%
\pgfpathlineto{\pgfqpoint{3.680247in}{0.790261in}}%
\pgfpathlineto{\pgfqpoint{3.681416in}{0.856254in}}%
\pgfpathlineto{\pgfqpoint{3.681439in}{0.848005in}}%
\pgfpathlineto{\pgfqpoint{3.682423in}{0.856254in}}%
\pgfpathlineto{\pgfqpoint{3.682446in}{0.848005in}}%
\pgfpathlineto{\pgfqpoint{3.683097in}{0.831507in}}%
\pgfpathlineto{\pgfqpoint{3.683260in}{0.856254in}}%
\pgfpathlineto{\pgfqpoint{3.683282in}{0.848005in}}%
\pgfpathlineto{\pgfqpoint{3.683601in}{0.856254in}}%
\pgfpathlineto{\pgfqpoint{3.683786in}{0.839756in}}%
\pgfpathlineto{\pgfqpoint{3.684289in}{0.798510in}}%
\pgfpathlineto{\pgfqpoint{3.684778in}{0.856254in}}%
\pgfpathlineto{\pgfqpoint{3.684793in}{0.848005in}}%
\pgfpathlineto{\pgfqpoint{3.686118in}{0.881002in}}%
\pgfpathlineto{\pgfqpoint{3.686133in}{0.872753in}}%
\pgfpathlineto{\pgfqpoint{3.687118in}{0.856254in}}%
\pgfpathlineto{\pgfqpoint{3.687140in}{0.872753in}}%
\pgfpathlineto{\pgfqpoint{3.688798in}{0.905749in}}%
\pgfpathlineto{\pgfqpoint{3.689990in}{0.815008in}}%
\pgfpathlineto{\pgfqpoint{3.691494in}{0.749015in}}%
\pgfpathlineto{\pgfqpoint{3.692649in}{0.757264in}}%
\pgfpathlineto{\pgfqpoint{3.693174in}{0.749015in}}%
\pgfpathlineto{\pgfqpoint{3.693485in}{0.782012in}}%
\pgfpathlineto{\pgfqpoint{3.693678in}{0.773763in}}%
\pgfpathlineto{\pgfqpoint{3.695166in}{0.831507in}}%
\pgfpathlineto{\pgfqpoint{3.695351in}{0.815008in}}%
\pgfpathlineto{\pgfqpoint{3.696025in}{0.798510in}}%
\pgfpathlineto{\pgfqpoint{3.696188in}{0.823258in}}%
\pgfpathlineto{\pgfqpoint{3.697343in}{0.856254in}}%
\pgfpathlineto{\pgfqpoint{3.697528in}{0.839756in}}%
\pgfpathlineto{\pgfqpoint{3.698031in}{0.798510in}}%
\pgfpathlineto{\pgfqpoint{3.698683in}{0.831507in}}%
\pgfpathlineto{\pgfqpoint{3.699520in}{0.856254in}}%
\pgfpathlineto{\pgfqpoint{3.699712in}{0.848005in}}%
\pgfpathlineto{\pgfqpoint{3.701889in}{0.946995in}}%
\pgfpathlineto{\pgfqpoint{3.702215in}{0.930497in}}%
\pgfpathlineto{\pgfqpoint{3.703229in}{0.913998in}}%
\pgfpathlineto{\pgfqpoint{3.703733in}{0.897500in}}%
\pgfpathlineto{\pgfqpoint{3.703903in}{0.922247in}}%
\pgfpathlineto{\pgfqpoint{3.704214in}{0.930497in}}%
\pgfpathlineto{\pgfqpoint{3.704569in}{0.897500in}}%
\pgfpathlineto{\pgfqpoint{3.704717in}{0.905749in}}%
\pgfpathlineto{\pgfqpoint{3.705058in}{0.930497in}}%
\pgfpathlineto{\pgfqpoint{3.705909in}{0.872753in}}%
\pgfpathlineto{\pgfqpoint{3.706561in}{0.881002in}}%
\pgfpathlineto{\pgfqpoint{3.706754in}{0.864503in}}%
\pgfpathlineto{\pgfqpoint{3.707753in}{0.848005in}}%
\pgfpathlineto{\pgfqpoint{3.707901in}{0.856254in}}%
\pgfpathlineto{\pgfqpoint{3.708908in}{0.955244in}}%
\pgfpathlineto{\pgfqpoint{3.708930in}{0.946995in}}%
\pgfpathlineto{\pgfqpoint{3.709937in}{0.971742in}}%
\pgfpathlineto{\pgfqpoint{3.710100in}{0.963493in}}%
\pgfpathlineto{\pgfqpoint{3.710441in}{0.946995in}}%
\pgfpathlineto{\pgfqpoint{3.710752in}{0.979992in}}%
\pgfpathlineto{\pgfqpoint{3.710774in}{0.971742in}}%
\pgfpathlineto{\pgfqpoint{3.711426in}{0.955244in}}%
\pgfpathlineto{\pgfqpoint{3.711929in}{0.979992in}}%
\pgfpathlineto{\pgfqpoint{3.712951in}{0.971742in}}%
\pgfpathlineto{\pgfqpoint{3.713269in}{0.979992in}}%
\pgfpathlineto{\pgfqpoint{3.713602in}{0.955244in}}%
\pgfpathlineto{\pgfqpoint{3.713958in}{0.946995in}}%
\pgfpathlineto{\pgfqpoint{3.714128in}{0.971742in}}%
\pgfpathlineto{\pgfqpoint{3.714632in}{0.963493in}}%
\pgfpathlineto{\pgfqpoint{3.716623in}{0.881002in}}%
\pgfpathlineto{\pgfqpoint{3.714965in}{0.971742in}}%
\pgfpathlineto{\pgfqpoint{3.716638in}{0.897500in}}%
\pgfpathlineto{\pgfqpoint{3.717645in}{0.971742in}}%
\pgfpathlineto{\pgfqpoint{3.718482in}{0.938746in}}%
\pgfpathlineto{\pgfqpoint{3.718630in}{0.930497in}}%
\pgfpathlineto{\pgfqpoint{3.718652in}{0.971742in}}%
\pgfpathlineto{\pgfqpoint{3.718985in}{0.971742in}}%
\pgfpathlineto{\pgfqpoint{3.722317in}{1.153224in}}%
\pgfpathlineto{\pgfqpoint{3.722510in}{1.136726in}}%
\pgfpathlineto{\pgfqpoint{3.723850in}{1.095480in}}%
\pgfpathlineto{\pgfqpoint{3.723998in}{1.103729in}}%
\pgfpathlineto{\pgfqpoint{3.724353in}{1.062483in}}%
\pgfpathlineto{\pgfqpoint{3.728018in}{0.633527in}}%
\pgfpathlineto{\pgfqpoint{3.728204in}{0.641776in}}%
\pgfpathlineto{\pgfqpoint{3.728692in}{0.633527in}}%
\pgfpathlineto{\pgfqpoint{3.728855in}{0.658274in}}%
\pgfpathlineto{\pgfqpoint{3.728877in}{0.650025in}}%
\pgfpathlineto{\pgfqpoint{3.730032in}{0.732517in}}%
\pgfpathlineto{\pgfqpoint{3.730869in}{0.707769in}}%
\pgfpathlineto{\pgfqpoint{3.730891in}{0.699520in}}%
\pgfpathlineto{\pgfqpoint{3.731202in}{0.732517in}}%
\pgfpathlineto{\pgfqpoint{3.731898in}{0.716019in}}%
\pgfpathlineto{\pgfqpoint{3.733572in}{0.674773in}}%
\pgfpathlineto{\pgfqpoint{3.734727in}{0.683022in}}%
\pgfpathlineto{\pgfqpoint{3.734912in}{0.674773in}}%
\pgfpathlineto{\pgfqpoint{3.735230in}{0.707769in}}%
\pgfpathlineto{\pgfqpoint{3.735748in}{0.699520in}}%
\pgfpathlineto{\pgfqpoint{3.736067in}{0.683022in}}%
\pgfpathlineto{\pgfqpoint{3.736400in}{0.716019in}}%
\pgfpathlineto{\pgfqpoint{3.737407in}{0.732517in}}%
\pgfpathlineto{\pgfqpoint{3.736755in}{0.699520in}}%
\pgfpathlineto{\pgfqpoint{3.737429in}{0.724268in}}%
\pgfpathlineto{\pgfqpoint{3.738244in}{0.732517in}}%
\pgfpathlineto{\pgfqpoint{3.738266in}{0.724268in}}%
\pgfpathlineto{\pgfqpoint{3.739939in}{0.650025in}}%
\pgfpathlineto{\pgfqpoint{3.740109in}{0.674773in}}%
\pgfpathlineto{\pgfqpoint{3.740443in}{0.699520in}}%
\pgfpathlineto{\pgfqpoint{3.740776in}{0.666524in}}%
\pgfpathlineto{\pgfqpoint{3.741116in}{0.650025in}}%
\pgfpathlineto{\pgfqpoint{3.741427in}{0.683022in}}%
\pgfpathlineto{\pgfqpoint{3.741450in}{0.674773in}}%
\pgfpathlineto{\pgfqpoint{3.741768in}{0.683022in}}%
\pgfpathlineto{\pgfqpoint{3.741953in}{0.666524in}}%
\pgfpathlineto{\pgfqpoint{3.743271in}{0.625278in}}%
\pgfpathlineto{\pgfqpoint{3.744448in}{0.633527in}}%
\pgfpathlineto{\pgfqpoint{3.744471in}{0.625278in}}%
\pgfpathlineto{\pgfqpoint{3.745507in}{0.625278in}}%
\pgfpathlineto{\pgfqpoint{3.745811in}{0.650025in}}%
\pgfpathlineto{\pgfqpoint{3.746477in}{0.625278in}}%
\pgfpathlineto{\pgfqpoint{3.747810in}{0.658274in}}%
\pgfpathlineto{\pgfqpoint{3.747988in}{0.641776in}}%
\pgfpathlineto{\pgfqpoint{3.748469in}{0.633527in}}%
\pgfpathlineto{\pgfqpoint{3.748491in}{0.674773in}}%
\pgfpathlineto{\pgfqpoint{3.750483in}{0.732517in}}%
\pgfpathlineto{\pgfqpoint{3.750505in}{0.724268in}}%
\pgfpathlineto{\pgfqpoint{3.752156in}{0.683022in}}%
\pgfpathlineto{\pgfqpoint{3.752163in}{0.691271in}}%
\pgfpathlineto{\pgfqpoint{3.752497in}{0.683022in}}%
\pgfpathlineto{\pgfqpoint{3.753185in}{0.699520in}}%
\pgfpathlineto{\pgfqpoint{3.754503in}{0.683022in}}%
\pgfpathlineto{\pgfqpoint{3.755518in}{0.707769in}}%
\pgfpathlineto{\pgfqpoint{3.755532in}{0.699520in}}%
\pgfpathlineto{\pgfqpoint{3.757206in}{0.773763in}}%
\pgfpathlineto{\pgfqpoint{3.757532in}{0.757264in}}%
\pgfpathlineto{\pgfqpoint{3.758546in}{0.716019in}}%
\pgfpathlineto{\pgfqpoint{3.759375in}{0.707769in}}%
\pgfpathlineto{\pgfqpoint{3.758716in}{0.724268in}}%
\pgfpathlineto{\pgfqpoint{3.759383in}{0.724268in}}%
\pgfpathlineto{\pgfqpoint{3.762389in}{0.831507in}}%
\pgfpathlineto{\pgfqpoint{3.762907in}{0.823258in}}%
\pgfpathlineto{\pgfqpoint{3.763233in}{0.856254in}}%
\pgfpathlineto{\pgfqpoint{3.763410in}{0.848005in}}%
\pgfpathlineto{\pgfqpoint{3.764573in}{0.856254in}}%
\pgfpathlineto{\pgfqpoint{3.765587in}{0.823258in}}%
\pgfpathlineto{\pgfqpoint{3.765743in}{0.831507in}}%
\pgfpathlineto{\pgfqpoint{3.766091in}{0.798510in}}%
\pgfpathlineto{\pgfqpoint{3.766246in}{0.806759in}}%
\pgfpathlineto{\pgfqpoint{3.766765in}{0.848005in}}%
\pgfpathlineto{\pgfqpoint{3.767268in}{0.798510in}}%
\pgfpathlineto{\pgfqpoint{3.771118in}{0.650025in}}%
\pgfpathlineto{\pgfqpoint{3.771459in}{0.674773in}}%
\pgfpathlineto{\pgfqpoint{3.771955in}{0.641776in}}%
\pgfpathlineto{\pgfqpoint{3.772125in}{0.650025in}}%
\pgfpathlineto{\pgfqpoint{3.773132in}{0.625278in}}%
\pgfpathlineto{\pgfqpoint{3.774791in}{0.658274in}}%
\pgfpathlineto{\pgfqpoint{3.775812in}{0.625278in}}%
\pgfpathlineto{\pgfqpoint{3.776990in}{0.650025in}}%
\pgfpathlineto{\pgfqpoint{3.777153in}{0.641776in}}%
\pgfpathlineto{\pgfqpoint{3.778145in}{0.633527in}}%
\pgfpathlineto{\pgfqpoint{3.777471in}{0.658274in}}%
\pgfpathlineto{\pgfqpoint{3.778152in}{0.641776in}}%
\pgfpathlineto{\pgfqpoint{3.778819in}{0.707769in}}%
\pgfpathlineto{\pgfqpoint{3.779167in}{0.674773in}}%
\pgfpathlineto{\pgfqpoint{3.780995in}{0.625278in}}%
\pgfpathlineto{\pgfqpoint{3.782165in}{0.633527in}}%
\pgfpathlineto{\pgfqpoint{3.782343in}{0.625278in}}%
\pgfpathlineto{\pgfqpoint{3.783172in}{0.625278in}}%
\pgfpathlineto{\pgfqpoint{3.784342in}{0.633527in}}%
\pgfpathlineto{\pgfqpoint{3.784364in}{0.625278in}}%
\pgfpathlineto{\pgfqpoint{3.785371in}{0.650025in}}%
\pgfpathlineto{\pgfqpoint{3.786038in}{0.699520in}}%
\pgfpathlineto{\pgfqpoint{3.786541in}{0.666524in}}%
\pgfpathlineto{\pgfqpoint{3.787363in}{0.658274in}}%
\pgfpathlineto{\pgfqpoint{3.786711in}{0.674773in}}%
\pgfpathlineto{\pgfqpoint{3.787378in}{0.674773in}}%
\pgfpathlineto{\pgfqpoint{3.787867in}{0.683022in}}%
\pgfpathlineto{\pgfqpoint{3.788052in}{0.666524in}}%
\pgfpathlineto{\pgfqpoint{3.788873in}{0.633527in}}%
\pgfpathlineto{\pgfqpoint{3.789221in}{0.650025in}}%
\pgfpathlineto{\pgfqpoint{3.790717in}{0.683022in}}%
\pgfpathlineto{\pgfqpoint{3.790732in}{0.674773in}}%
\pgfpathlineto{\pgfqpoint{3.792390in}{0.633527in}}%
\pgfpathlineto{\pgfqpoint{3.792390in}{0.641776in}}%
\pgfpathlineto{\pgfqpoint{3.792724in}{0.633527in}}%
\pgfpathlineto{\pgfqpoint{3.793412in}{0.650025in}}%
\pgfpathlineto{\pgfqpoint{3.794064in}{0.658274in}}%
\pgfpathlineto{\pgfqpoint{3.794256in}{0.641776in}}%
\pgfpathlineto{\pgfqpoint{3.794419in}{0.650025in}}%
\pgfpathlineto{\pgfqpoint{3.795411in}{0.633527in}}%
\pgfpathlineto{\pgfqpoint{3.796411in}{0.683022in}}%
\pgfpathlineto{\pgfqpoint{3.795759in}{0.625278in}}%
\pgfpathlineto{\pgfqpoint{3.796433in}{0.674773in}}%
\pgfpathlineto{\pgfqpoint{3.796752in}{0.683022in}}%
\pgfpathlineto{\pgfqpoint{3.797085in}{0.658274in}}%
\pgfpathlineto{\pgfqpoint{3.797270in}{0.666524in}}%
\pgfpathlineto{\pgfqpoint{3.797418in}{0.658274in}}%
\pgfpathlineto{\pgfqpoint{3.797440in}{0.699520in}}%
\pgfpathlineto{\pgfqpoint{3.797773in}{0.699520in}}%
\pgfpathlineto{\pgfqpoint{3.798092in}{0.683022in}}%
\pgfpathlineto{\pgfqpoint{3.798107in}{0.724268in}}%
\pgfpathlineto{\pgfqpoint{3.798758in}{0.757264in}}%
\pgfpathlineto{\pgfqpoint{3.799284in}{0.740766in}}%
\pgfpathlineto{\pgfqpoint{3.800098in}{0.707769in}}%
\pgfpathlineto{\pgfqpoint{3.799447in}{0.749015in}}%
\pgfpathlineto{\pgfqpoint{3.800439in}{0.732517in}}%
\pgfpathlineto{\pgfqpoint{3.800794in}{0.773763in}}%
\pgfpathlineto{\pgfqpoint{3.801461in}{0.773763in}}%
\pgfpathlineto{\pgfqpoint{3.802445in}{0.732517in}}%
\pgfpathlineto{\pgfqpoint{3.802801in}{0.749015in}}%
\pgfpathlineto{\pgfqpoint{3.803956in}{0.757264in}}%
\pgfpathlineto{\pgfqpoint{3.804141in}{0.749015in}}%
\pgfpathlineto{\pgfqpoint{3.804459in}{0.782012in}}%
\pgfpathlineto{\pgfqpoint{3.804985in}{0.765513in}}%
\pgfpathlineto{\pgfqpoint{3.805148in}{0.798510in}}%
\pgfpathlineto{\pgfqpoint{3.806155in}{0.749015in}}%
\pgfpathlineto{\pgfqpoint{3.808169in}{0.823258in}}%
\pgfpathlineto{\pgfqpoint{3.808332in}{0.815008in}}%
\pgfpathlineto{\pgfqpoint{3.808480in}{0.806759in}}%
\pgfpathlineto{\pgfqpoint{3.808650in}{0.831507in}}%
\pgfpathlineto{\pgfqpoint{3.809005in}{0.823258in}}%
\pgfpathlineto{\pgfqpoint{3.810161in}{0.856254in}}%
\pgfpathlineto{\pgfqpoint{3.810346in}{0.839756in}}%
\pgfpathlineto{\pgfqpoint{3.810849in}{0.823258in}}%
\pgfpathlineto{\pgfqpoint{3.811019in}{0.848005in}}%
\pgfpathlineto{\pgfqpoint{3.811330in}{0.839756in}}%
\pgfpathlineto{\pgfqpoint{3.811671in}{0.831507in}}%
\pgfpathlineto{\pgfqpoint{3.812360in}{0.848005in}}%
\pgfpathlineto{\pgfqpoint{3.812671in}{0.831507in}}%
\pgfpathlineto{\pgfqpoint{3.812841in}{0.856254in}}%
\pgfpathlineto{\pgfqpoint{3.813196in}{0.848005in}}%
\pgfpathlineto{\pgfqpoint{3.813344in}{0.856254in}}%
\pgfpathlineto{\pgfqpoint{3.813678in}{0.831507in}}%
\pgfpathlineto{\pgfqpoint{3.813863in}{0.839756in}}%
\pgfpathlineto{\pgfqpoint{3.814033in}{0.848005in}}%
\pgfpathlineto{\pgfqpoint{3.815543in}{0.798510in}}%
\pgfpathlineto{\pgfqpoint{3.816195in}{0.856254in}}%
\pgfpathlineto{\pgfqpoint{3.816713in}{0.815008in}}%
\pgfpathlineto{\pgfqpoint{3.816861in}{0.806759in}}%
\pgfpathlineto{\pgfqpoint{3.817365in}{0.839756in}}%
\pgfpathlineto{\pgfqpoint{3.818209in}{0.856254in}}%
\pgfpathlineto{\pgfqpoint{3.817705in}{0.831507in}}%
\pgfpathlineto{\pgfqpoint{3.818394in}{0.839756in}}%
\pgfpathlineto{\pgfqpoint{3.819712in}{0.806759in}}%
\pgfpathlineto{\pgfqpoint{3.819734in}{0.823258in}}%
\pgfpathlineto{\pgfqpoint{3.821726in}{0.881002in}}%
\pgfpathlineto{\pgfqpoint{3.822563in}{0.856254in}}%
\pgfpathlineto{\pgfqpoint{3.822748in}{0.872753in}}%
\pgfpathlineto{\pgfqpoint{3.824073in}{0.905749in}}%
\pgfpathlineto{\pgfqpoint{3.823399in}{0.856254in}}%
\pgfpathlineto{\pgfqpoint{3.824095in}{0.897500in}}%
\pgfpathlineto{\pgfqpoint{3.824747in}{0.881002in}}%
\pgfpathlineto{\pgfqpoint{3.824910in}{0.905749in}}%
\pgfpathlineto{\pgfqpoint{3.824932in}{0.897500in}}%
\pgfpathlineto{\pgfqpoint{3.825243in}{0.905749in}}%
\pgfpathlineto{\pgfqpoint{3.825598in}{0.872753in}}%
\pgfpathlineto{\pgfqpoint{3.825917in}{0.856254in}}%
\pgfpathlineto{\pgfqpoint{3.826250in}{0.889251in}}%
\pgfpathlineto{\pgfqpoint{3.826590in}{0.905749in}}%
\pgfpathlineto{\pgfqpoint{3.826938in}{0.872753in}}%
\pgfpathlineto{\pgfqpoint{3.827279in}{0.872753in}}%
\pgfpathlineto{\pgfqpoint{3.827783in}{0.922247in}}%
\pgfpathlineto{\pgfqpoint{3.828449in}{0.889251in}}%
\pgfpathlineto{\pgfqpoint{3.828938in}{0.881002in}}%
\pgfpathlineto{\pgfqpoint{3.829100in}{0.905749in}}%
\pgfpathlineto{\pgfqpoint{3.829123in}{0.897500in}}%
\pgfpathlineto{\pgfqpoint{3.830781in}{0.955244in}}%
\pgfpathlineto{\pgfqpoint{3.830966in}{0.938746in}}%
\pgfpathlineto{\pgfqpoint{3.831448in}{0.905749in}}%
\pgfpathlineto{\pgfqpoint{3.831803in}{0.946995in}}%
\pgfpathlineto{\pgfqpoint{3.833462in}{0.979992in}}%
\pgfpathlineto{\pgfqpoint{3.834298in}{0.955244in}}%
\pgfpathlineto{\pgfqpoint{3.834483in}{0.963493in}}%
\pgfpathlineto{\pgfqpoint{3.834654in}{0.971742in}}%
\pgfpathlineto{\pgfqpoint{3.835661in}{0.946995in}}%
\pgfpathlineto{\pgfqpoint{3.836816in}{0.979992in}}%
\pgfpathlineto{\pgfqpoint{3.837001in}{0.963493in}}%
\pgfpathlineto{\pgfqpoint{3.837334in}{0.946995in}}%
\pgfpathlineto{\pgfqpoint{3.837652in}{0.979992in}}%
\pgfpathlineto{\pgfqpoint{3.837815in}{0.979992in}}%
\pgfpathlineto{\pgfqpoint{3.837837in}{0.996490in}}%
\pgfpathlineto{\pgfqpoint{3.838844in}{0.946995in}}%
\pgfpathlineto{\pgfqpoint{3.839163in}{0.930497in}}%
\pgfpathlineto{\pgfqpoint{3.839326in}{0.955244in}}%
\pgfpathlineto{\pgfqpoint{3.839681in}{0.946995in}}%
\pgfpathlineto{\pgfqpoint{3.839829in}{0.955244in}}%
\pgfpathlineto{\pgfqpoint{3.840162in}{0.930497in}}%
\pgfpathlineto{\pgfqpoint{3.840355in}{0.938746in}}%
\pgfpathlineto{\pgfqpoint{3.841021in}{0.897500in}}%
\pgfpathlineto{\pgfqpoint{3.841510in}{0.930497in}}%
\pgfpathlineto{\pgfqpoint{3.841673in}{0.955244in}}%
\pgfpathlineto{\pgfqpoint{3.842361in}{0.922247in}}%
\pgfpathlineto{\pgfqpoint{3.842532in}{0.922247in}}%
\pgfpathlineto{\pgfqpoint{3.844190in}{0.955244in}}%
\pgfpathlineto{\pgfqpoint{3.845197in}{0.930497in}}%
\pgfpathlineto{\pgfqpoint{3.845212in}{0.946995in}}%
\pgfpathlineto{\pgfqpoint{3.845864in}{0.930497in}}%
\pgfpathlineto{\pgfqpoint{3.846034in}{0.955244in}}%
\pgfpathlineto{\pgfqpoint{3.846049in}{0.946995in}}%
\pgfpathlineto{\pgfqpoint{3.847204in}{0.979992in}}%
\pgfpathlineto{\pgfqpoint{3.847389in}{0.963493in}}%
\pgfpathlineto{\pgfqpoint{3.848736in}{0.922247in}}%
\pgfpathlineto{\pgfqpoint{3.848884in}{0.938746in}}%
\pgfpathlineto{\pgfqpoint{3.848899in}{0.946995in}}%
\pgfpathlineto{\pgfqpoint{3.849240in}{0.922247in}}%
\pgfpathlineto{\pgfqpoint{3.849906in}{0.922247in}}%
\pgfpathlineto{\pgfqpoint{3.851565in}{0.955244in}}%
\pgfpathlineto{\pgfqpoint{3.852735in}{0.905749in}}%
\pgfpathlineto{\pgfqpoint{3.852735in}{0.913998in}}%
\pgfpathlineto{\pgfqpoint{3.853912in}{0.955244in}}%
\pgfpathlineto{\pgfqpoint{3.853075in}{0.905749in}}%
\pgfpathlineto{\pgfqpoint{3.853927in}{0.946995in}}%
\pgfpathlineto{\pgfqpoint{3.854082in}{0.955244in}}%
\pgfpathlineto{\pgfqpoint{3.854430in}{0.922247in}}%
\pgfpathlineto{\pgfqpoint{3.854578in}{0.930497in}}%
\pgfpathlineto{\pgfqpoint{3.855607in}{0.897500in}}%
\pgfpathlineto{\pgfqpoint{3.855756in}{0.905749in}}%
\pgfpathlineto{\pgfqpoint{3.855770in}{0.922247in}}%
\pgfpathlineto{\pgfqpoint{3.856777in}{0.922247in}}%
\pgfpathlineto{\pgfqpoint{3.857429in}{0.905749in}}%
\pgfpathlineto{\pgfqpoint{3.857599in}{0.930497in}}%
\pgfpathlineto{\pgfqpoint{3.857621in}{0.922247in}}%
\pgfpathlineto{\pgfqpoint{3.858621in}{0.946995in}}%
\pgfpathlineto{\pgfqpoint{3.858791in}{0.938746in}}%
\pgfpathlineto{\pgfqpoint{3.860117in}{0.881002in}}%
\pgfpathlineto{\pgfqpoint{3.858962in}{0.946995in}}%
\pgfpathlineto{\pgfqpoint{3.860465in}{0.897500in}}%
\pgfpathlineto{\pgfqpoint{3.861620in}{0.905749in}}%
\pgfpathlineto{\pgfqpoint{3.861642in}{0.897500in}}%
\pgfpathlineto{\pgfqpoint{3.862145in}{0.946995in}}%
\pgfpathlineto{\pgfqpoint{3.862649in}{0.922247in}}%
\pgfpathlineto{\pgfqpoint{3.864137in}{0.979992in}}%
\pgfpathlineto{\pgfqpoint{3.864322in}{0.963493in}}%
\pgfpathlineto{\pgfqpoint{3.865499in}{0.946995in}}%
\pgfpathlineto{\pgfqpoint{3.867151in}{1.004739in}}%
\pgfpathlineto{\pgfqpoint{3.867343in}{0.988241in}}%
\pgfpathlineto{\pgfqpoint{3.867676in}{0.971742in}}%
\pgfpathlineto{\pgfqpoint{3.867995in}{1.012988in}}%
\pgfpathlineto{\pgfqpoint{3.868328in}{1.029487in}}%
\pgfpathlineto{\pgfqpoint{3.868683in}{0.996490in}}%
\pgfpathlineto{\pgfqpoint{3.869016in}{0.996490in}}%
\pgfpathlineto{\pgfqpoint{3.870172in}{1.004739in}}%
\pgfpathlineto{\pgfqpoint{3.871179in}{0.979992in}}%
\pgfpathlineto{\pgfqpoint{3.871193in}{0.996490in}}%
\pgfpathlineto{\pgfqpoint{3.871697in}{0.971742in}}%
\pgfpathlineto{\pgfqpoint{3.872015in}{1.004739in}}%
\pgfpathlineto{\pgfqpoint{3.872037in}{0.996490in}}%
\pgfpathlineto{\pgfqpoint{3.872533in}{0.971742in}}%
\pgfpathlineto{\pgfqpoint{3.873185in}{1.004739in}}%
\pgfpathlineto{\pgfqpoint{3.874385in}{0.971742in}}%
\pgfpathlineto{\pgfqpoint{3.874696in}{0.955244in}}%
\pgfpathlineto{\pgfqpoint{3.875036in}{0.988241in}}%
\pgfpathlineto{\pgfqpoint{3.876036in}{1.004739in}}%
\pgfpathlineto{\pgfqpoint{3.876058in}{0.996490in}}%
\pgfpathlineto{\pgfqpoint{3.876724in}{0.971742in}}%
\pgfpathlineto{\pgfqpoint{3.876895in}{0.996490in}}%
\pgfpathlineto{\pgfqpoint{3.878050in}{1.004739in}}%
\pgfpathlineto{\pgfqpoint{3.878738in}{0.996490in}}%
\pgfpathlineto{\pgfqpoint{3.879057in}{1.029487in}}%
\pgfpathlineto{\pgfqpoint{3.879071in}{1.021237in}}%
\pgfpathlineto{\pgfqpoint{3.880249in}{1.070732in}}%
\pgfpathlineto{\pgfqpoint{3.880745in}{1.054234in}}%
\pgfpathlineto{\pgfqpoint{3.881404in}{1.029487in}}%
\pgfpathlineto{\pgfqpoint{3.881759in}{1.045985in}}%
\pgfpathlineto{\pgfqpoint{3.882914in}{1.054234in}}%
\pgfpathlineto{\pgfqpoint{3.883247in}{1.078981in}}%
\pgfpathlineto{\pgfqpoint{3.883936in}{1.021237in}}%
\pgfpathlineto{\pgfqpoint{3.884610in}{1.070732in}}%
\pgfpathlineto{\pgfqpoint{3.885276in}{1.037736in}}%
\pgfpathlineto{\pgfqpoint{3.887794in}{0.839756in}}%
\pgfpathlineto{\pgfqpoint{3.889297in}{0.724268in}}%
\pgfpathlineto{\pgfqpoint{3.889467in}{0.749015in}}%
\pgfpathlineto{\pgfqpoint{3.891814in}{0.674773in}}%
\pgfpathlineto{\pgfqpoint{3.892466in}{0.683022in}}%
\pgfpathlineto{\pgfqpoint{3.892651in}{0.666524in}}%
\pgfpathlineto{\pgfqpoint{3.892799in}{0.658274in}}%
\pgfpathlineto{\pgfqpoint{3.893154in}{0.699520in}}%
\pgfpathlineto{\pgfqpoint{3.893473in}{0.707769in}}%
\pgfpathlineto{\pgfqpoint{3.893806in}{0.683022in}}%
\pgfpathlineto{\pgfqpoint{3.894139in}{0.658274in}}%
\pgfpathlineto{\pgfqpoint{3.894642in}{0.707769in}}%
\pgfpathlineto{\pgfqpoint{3.894835in}{0.699520in}}%
\pgfpathlineto{\pgfqpoint{3.896508in}{0.650025in}}%
\pgfpathlineto{\pgfqpoint{3.896656in}{0.658274in}}%
\pgfpathlineto{\pgfqpoint{3.896827in}{0.683022in}}%
\pgfpathlineto{\pgfqpoint{3.897678in}{0.674773in}}%
\pgfpathlineto{\pgfqpoint{3.898330in}{0.658274in}}%
\pgfpathlineto{\pgfqpoint{3.898500in}{0.683022in}}%
\pgfpathlineto{\pgfqpoint{3.898522in}{0.674773in}}%
\pgfpathlineto{\pgfqpoint{3.899685in}{0.683022in}}%
\pgfpathlineto{\pgfqpoint{3.900196in}{0.674773in}}%
\pgfpathlineto{\pgfqpoint{3.900699in}{0.724268in}}%
\pgfpathlineto{\pgfqpoint{3.901691in}{0.732517in}}%
\pgfpathlineto{\pgfqpoint{3.901706in}{0.724268in}}%
\pgfpathlineto{\pgfqpoint{3.902358in}{0.707769in}}%
\pgfpathlineto{\pgfqpoint{3.902520in}{0.732517in}}%
\pgfpathlineto{\pgfqpoint{3.902543in}{0.724268in}}%
\pgfpathlineto{\pgfqpoint{3.903194in}{0.806759in}}%
\pgfpathlineto{\pgfqpoint{3.904201in}{0.757264in}}%
\pgfpathlineto{\pgfqpoint{3.904720in}{0.749015in}}%
\pgfpathlineto{\pgfqpoint{3.905038in}{0.782012in}}%
\pgfpathlineto{\pgfqpoint{3.905223in}{0.765513in}}%
\pgfpathlineto{\pgfqpoint{3.906400in}{0.724268in}}%
\pgfpathlineto{\pgfqpoint{3.906563in}{0.749015in}}%
\pgfpathlineto{\pgfqpoint{3.909058in}{0.831507in}}%
\pgfpathlineto{\pgfqpoint{3.910250in}{0.798510in}}%
\pgfpathlineto{\pgfqpoint{3.911257in}{0.823258in}}%
\pgfpathlineto{\pgfqpoint{3.911428in}{0.815008in}}%
\pgfpathlineto{\pgfqpoint{3.912250in}{0.782012in}}%
\pgfpathlineto{\pgfqpoint{3.912598in}{0.798510in}}%
\pgfpathlineto{\pgfqpoint{3.915093in}{0.881002in}}%
\pgfpathlineto{\pgfqpoint{3.915937in}{0.856254in}}%
\pgfpathlineto{\pgfqpoint{3.916122in}{0.872753in}}%
\pgfpathlineto{\pgfqpoint{3.916625in}{0.897500in}}%
\pgfpathlineto{\pgfqpoint{3.917277in}{0.881002in}}%
\pgfpathlineto{\pgfqpoint{3.918469in}{0.790261in}}%
\pgfpathlineto{\pgfqpoint{3.919639in}{0.724268in}}%
\pgfpathlineto{\pgfqpoint{3.920142in}{0.749015in}}%
\pgfpathlineto{\pgfqpoint{3.921468in}{0.806759in}}%
\pgfpathlineto{\pgfqpoint{3.921986in}{0.790261in}}%
\pgfpathlineto{\pgfqpoint{3.922134in}{0.782012in}}%
\pgfpathlineto{\pgfqpoint{3.922490in}{0.823258in}}%
\pgfpathlineto{\pgfqpoint{3.923326in}{0.872753in}}%
\pgfpathlineto{\pgfqpoint{3.924156in}{0.856254in}}%
\pgfpathlineto{\pgfqpoint{3.924333in}{0.848005in}}%
\pgfpathlineto{\pgfqpoint{3.925007in}{0.897500in}}%
\pgfpathlineto{\pgfqpoint{3.925170in}{0.889251in}}%
\pgfpathlineto{\pgfqpoint{3.925325in}{0.881002in}}%
\pgfpathlineto{\pgfqpoint{3.925666in}{0.913998in}}%
\pgfpathlineto{\pgfqpoint{3.926851in}{1.021237in}}%
\pgfpathlineto{\pgfqpoint{3.927176in}{1.004739in}}%
\pgfpathlineto{\pgfqpoint{3.927510in}{1.037736in}}%
\pgfpathlineto{\pgfqpoint{3.928517in}{1.078981in}}%
\pgfpathlineto{\pgfqpoint{3.928524in}{1.070732in}}%
\pgfpathlineto{\pgfqpoint{3.929020in}{1.078981in}}%
\pgfpathlineto{\pgfqpoint{3.929864in}{1.045985in}}%
\pgfpathlineto{\pgfqpoint{3.931019in}{1.078981in}}%
\pgfpathlineto{\pgfqpoint{3.931204in}{1.062483in}}%
\pgfpathlineto{\pgfqpoint{3.932359in}{1.004739in}}%
\pgfpathlineto{\pgfqpoint{3.932885in}{1.021237in}}%
\pgfpathlineto{\pgfqpoint{3.933707in}{1.054234in}}%
\pgfpathlineto{\pgfqpoint{3.934048in}{1.029487in}}%
\pgfpathlineto{\pgfqpoint{3.935062in}{1.012988in}}%
\pgfpathlineto{\pgfqpoint{3.936720in}{0.979992in}}%
\pgfpathlineto{\pgfqpoint{3.937076in}{0.971742in}}%
\pgfpathlineto{\pgfqpoint{3.937742in}{1.021237in}}%
\pgfpathlineto{\pgfqpoint{3.937913in}{1.045985in}}%
\pgfpathlineto{\pgfqpoint{3.938749in}{1.021237in}}%
\pgfpathlineto{\pgfqpoint{3.939253in}{0.996490in}}%
\pgfpathlineto{\pgfqpoint{3.939571in}{1.029487in}}%
\pgfpathlineto{\pgfqpoint{3.939586in}{1.021237in}}%
\pgfpathlineto{\pgfqpoint{3.939741in}{1.029487in}}%
\pgfpathlineto{\pgfqpoint{3.940089in}{0.996490in}}%
\pgfpathlineto{\pgfqpoint{3.940237in}{1.004739in}}%
\pgfpathlineto{\pgfqpoint{3.941267in}{0.971742in}}%
\pgfpathlineto{\pgfqpoint{3.941415in}{0.979992in}}%
\pgfpathlineto{\pgfqpoint{3.941585in}{1.004739in}}%
\pgfpathlineto{\pgfqpoint{3.942274in}{0.971742in}}%
\pgfpathlineto{\pgfqpoint{3.942436in}{0.971742in}}%
\pgfpathlineto{\pgfqpoint{3.943429in}{1.004739in}}%
\pgfpathlineto{\pgfqpoint{3.943614in}{0.988241in}}%
\pgfpathlineto{\pgfqpoint{3.944954in}{0.864503in}}%
\pgfpathlineto{\pgfqpoint{3.945791in}{0.798510in}}%
\pgfpathlineto{\pgfqpoint{3.946294in}{0.823258in}}%
\pgfpathlineto{\pgfqpoint{3.947619in}{0.707769in}}%
\pgfpathlineto{\pgfqpoint{3.947953in}{0.683022in}}%
\pgfpathlineto{\pgfqpoint{3.948308in}{0.724268in}}%
\pgfpathlineto{\pgfqpoint{3.948641in}{0.699520in}}%
\pgfpathlineto{\pgfqpoint{3.949315in}{0.674773in}}%
\pgfpathlineto{\pgfqpoint{3.949478in}{0.699520in}}%
\pgfpathlineto{\pgfqpoint{3.951973in}{0.782012in}}%
\pgfpathlineto{\pgfqpoint{3.953165in}{0.724268in}}%
\pgfpathlineto{\pgfqpoint{3.954172in}{0.749015in}}%
\pgfpathlineto{\pgfqpoint{3.954342in}{0.740766in}}%
\pgfpathlineto{\pgfqpoint{3.954491in}{0.732517in}}%
\pgfpathlineto{\pgfqpoint{3.954994in}{0.765513in}}%
\pgfpathlineto{\pgfqpoint{3.955327in}{0.782012in}}%
\pgfpathlineto{\pgfqpoint{3.956016in}{0.724268in}}%
\pgfpathlineto{\pgfqpoint{3.956690in}{0.749015in}}%
\pgfpathlineto{\pgfqpoint{3.956852in}{0.740766in}}%
\pgfpathlineto{\pgfqpoint{3.957563in}{0.683022in}}%
\pgfpathlineto{\pgfqpoint{3.958008in}{0.716019in}}%
\pgfpathlineto{\pgfqpoint{3.958030in}{0.749015in}}%
\pgfpathlineto{\pgfqpoint{3.959014in}{0.707769in}}%
\pgfpathlineto{\pgfqpoint{3.959037in}{0.724268in}}%
\pgfpathlineto{\pgfqpoint{3.959703in}{0.773763in}}%
\pgfpathlineto{\pgfqpoint{3.960695in}{0.757264in}}%
\pgfpathlineto{\pgfqpoint{3.960710in}{0.749015in}}%
\pgfpathlineto{\pgfqpoint{3.961214in}{0.798510in}}%
\pgfpathlineto{\pgfqpoint{3.961717in}{0.790261in}}%
\pgfpathlineto{\pgfqpoint{3.962887in}{0.773763in}}%
\pgfpathlineto{\pgfqpoint{3.964545in}{0.806759in}}%
\pgfpathlineto{\pgfqpoint{3.964568in}{0.798510in}}%
\pgfpathlineto{\pgfqpoint{3.964886in}{0.831507in}}%
\pgfpathlineto{\pgfqpoint{3.965575in}{0.815008in}}%
\pgfpathlineto{\pgfqpoint{3.965723in}{0.806759in}}%
\pgfpathlineto{\pgfqpoint{3.965886in}{0.831507in}}%
\pgfpathlineto{\pgfqpoint{3.966389in}{0.831507in}}%
\pgfpathlineto{\pgfqpoint{3.966893in}{0.856254in}}%
\pgfpathlineto{\pgfqpoint{3.967418in}{0.815008in}}%
\pgfpathlineto{\pgfqpoint{3.967566in}{0.806759in}}%
\pgfpathlineto{\pgfqpoint{3.967922in}{0.848005in}}%
\pgfpathlineto{\pgfqpoint{3.969092in}{0.872753in}}%
\pgfpathlineto{\pgfqpoint{3.968410in}{0.831507in}}%
\pgfpathlineto{\pgfqpoint{3.969262in}{0.864503in}}%
\pgfpathlineto{\pgfqpoint{3.969410in}{0.856254in}}%
\pgfpathlineto{\pgfqpoint{3.969425in}{0.897500in}}%
\pgfpathlineto{\pgfqpoint{3.970076in}{0.881002in}}%
\pgfpathlineto{\pgfqpoint{3.970758in}{1.004739in}}%
\pgfpathlineto{\pgfqpoint{3.971106in}{0.988241in}}%
\pgfpathlineto{\pgfqpoint{3.972113in}{0.971742in}}%
\pgfpathlineto{\pgfqpoint{3.972261in}{0.979992in}}%
\pgfpathlineto{\pgfqpoint{3.973268in}{1.004739in}}%
\pgfpathlineto{\pgfqpoint{3.973282in}{0.996490in}}%
\pgfpathlineto{\pgfqpoint{3.974941in}{1.029487in}}%
\pgfpathlineto{\pgfqpoint{3.975629in}{1.045985in}}%
\pgfpathlineto{\pgfqpoint{3.975963in}{1.021237in}}%
\pgfpathlineto{\pgfqpoint{3.976466in}{0.996490in}}%
\pgfpathlineto{\pgfqpoint{3.976785in}{1.029487in}}%
\pgfpathlineto{\pgfqpoint{3.976807in}{1.021237in}}%
\pgfpathlineto{\pgfqpoint{3.977792in}{1.054234in}}%
\pgfpathlineto{\pgfqpoint{3.977977in}{1.037736in}}%
\pgfpathlineto{\pgfqpoint{3.978650in}{1.021237in}}%
\pgfpathlineto{\pgfqpoint{3.978813in}{1.045985in}}%
\pgfpathlineto{\pgfqpoint{3.980975in}{1.103729in}}%
\pgfpathlineto{\pgfqpoint{3.981331in}{1.095480in}}%
\pgfpathlineto{\pgfqpoint{3.981834in}{1.144975in}}%
\pgfpathlineto{\pgfqpoint{3.981997in}{1.136726in}}%
\pgfpathlineto{\pgfqpoint{3.986188in}{0.625278in}}%
\pgfpathlineto{\pgfqpoint{3.986358in}{0.650025in}}%
\pgfpathlineto{\pgfqpoint{3.986839in}{0.658274in}}%
\pgfpathlineto{\pgfqpoint{3.987032in}{0.641776in}}%
\pgfpathlineto{\pgfqpoint{3.988187in}{0.625278in}}%
\pgfpathlineto{\pgfqpoint{3.988187in}{0.625278in}}%
\pgfpathlineto{\pgfqpoint{3.989527in}{0.633527in}}%
\pgfpathlineto{\pgfqpoint{3.989690in}{0.625278in}}%
\pgfpathlineto{\pgfqpoint{3.990527in}{0.625278in}}%
\pgfpathlineto{\pgfqpoint{3.991541in}{0.633527in}}%
\pgfpathlineto{\pgfqpoint{3.991541in}{0.625278in}}%
\pgfpathlineto{\pgfqpoint{3.992874in}{0.633527in}}%
\pgfpathlineto{\pgfqpoint{3.992874in}{0.625278in}}%
\pgfpathlineto{\pgfqpoint{3.993881in}{0.625278in}}%
\pgfpathlineto{\pgfqpoint{3.994888in}{0.633527in}}%
\pgfpathlineto{\pgfqpoint{3.994888in}{0.625278in}}%
\pgfpathlineto{\pgfqpoint{3.996072in}{0.633527in}}%
\pgfpathlineto{\pgfqpoint{3.996228in}{0.625278in}}%
\pgfpathlineto{\pgfqpoint{3.997072in}{0.625278in}}%
\pgfpathlineto{\pgfqpoint{3.998079in}{0.633527in}}%
\pgfpathlineto{\pgfqpoint{3.998079in}{0.625278in}}%
\pgfpathlineto{\pgfqpoint{3.999249in}{0.633527in}}%
\pgfpathlineto{\pgfqpoint{3.999412in}{0.625278in}}%
\pgfpathlineto{\pgfqpoint{4.000093in}{0.625278in}}%
\pgfpathlineto{\pgfqpoint{4.001255in}{0.633527in}}%
\pgfpathlineto{\pgfqpoint{4.001426in}{0.625278in}}%
\pgfpathlineto{\pgfqpoint{4.002262in}{0.625278in}}%
\pgfpathlineto{\pgfqpoint{4.003447in}{0.633527in}}%
\pgfpathlineto{\pgfqpoint{4.003603in}{0.625278in}}%
\pgfpathlineto{\pgfqpoint{4.004454in}{0.625278in}}%
\pgfpathlineto{\pgfqpoint{4.005624in}{0.633527in}}%
\pgfpathlineto{\pgfqpoint{4.005794in}{0.625278in}}%
\pgfpathlineto{\pgfqpoint{4.006631in}{0.625278in}}%
\pgfpathlineto{\pgfqpoint{4.007793in}{0.633527in}}%
\pgfpathlineto{\pgfqpoint{4.007964in}{0.625278in}}%
\pgfpathlineto{\pgfqpoint{4.008645in}{0.625278in}}%
\pgfpathlineto{\pgfqpoint{4.010318in}{0.691271in}}%
\pgfpathlineto{\pgfqpoint{4.010481in}{0.707769in}}%
\pgfpathlineto{\pgfqpoint{4.011333in}{0.666524in}}%
\pgfpathlineto{\pgfqpoint{4.012991in}{0.625278in}}%
\pgfpathlineto{\pgfqpoint{4.014168in}{0.633527in}}%
\pgfpathlineto{\pgfqpoint{4.014183in}{0.625278in}}%
\pgfpathlineto{\pgfqpoint{4.015020in}{0.625278in}}%
\pgfpathlineto{\pgfqpoint{4.016175in}{0.633527in}}%
\pgfpathlineto{\pgfqpoint{4.016345in}{0.625278in}}%
\pgfpathlineto{\pgfqpoint{4.017204in}{0.625278in}}%
\pgfpathlineto{\pgfqpoint{4.018189in}{0.633527in}}%
\pgfpathlineto{\pgfqpoint{4.018211in}{0.625278in}}%
\pgfpathlineto{\pgfqpoint{4.019373in}{0.633527in}}%
\pgfpathlineto{\pgfqpoint{4.019529in}{0.625278in}}%
\pgfpathlineto{\pgfqpoint{4.020218in}{0.625278in}}%
\pgfpathlineto{\pgfqpoint{4.021387in}{0.633527in}}%
\pgfpathlineto{\pgfqpoint{4.022231in}{0.625278in}}%
\pgfpathlineto{\pgfqpoint{4.022402in}{0.650025in}}%
\pgfpathlineto{\pgfqpoint{4.022720in}{0.633527in}}%
\pgfpathlineto{\pgfqpoint{4.022735in}{0.674773in}}%
\pgfpathlineto{\pgfqpoint{4.023897in}{0.683022in}}%
\pgfpathlineto{\pgfqpoint{4.024579in}{0.674773in}}%
\pgfpathlineto{\pgfqpoint{4.024749in}{0.699520in}}%
\pgfpathlineto{\pgfqpoint{4.024912in}{0.691271in}}%
\pgfpathlineto{\pgfqpoint{4.026904in}{0.625278in}}%
\pgfpathlineto{\pgfqpoint{4.028081in}{0.633527in}}%
\pgfpathlineto{\pgfqpoint{4.028414in}{0.625278in}}%
\pgfpathlineto{\pgfqpoint{4.029088in}{0.625278in}}%
\pgfpathlineto{\pgfqpoint{4.030095in}{0.633527in}}%
\pgfpathlineto{\pgfqpoint{4.030095in}{0.625278in}}%
\pgfpathlineto{\pgfqpoint{4.031265in}{0.633527in}}%
\pgfpathlineto{\pgfqpoint{4.031435in}{0.625278in}}%
\pgfpathlineto{\pgfqpoint{4.032109in}{0.625278in}}%
\pgfpathlineto{\pgfqpoint{4.033279in}{0.633527in}}%
\pgfpathlineto{\pgfqpoint{4.033286in}{0.625278in}}%
\pgfpathlineto{\pgfqpoint{4.034286in}{0.625278in}}%
\pgfpathlineto{\pgfqpoint{4.035455in}{0.633527in}}%
\pgfpathlineto{\pgfqpoint{4.035455in}{0.625278in}}%
\pgfpathlineto{\pgfqpoint{4.036470in}{0.625278in}}%
\pgfpathlineto{\pgfqpoint{4.037632in}{0.633527in}}%
\pgfpathlineto{\pgfqpoint{4.037640in}{0.625278in}}%
\pgfpathlineto{\pgfqpoint{4.038647in}{0.625278in}}%
\pgfpathlineto{\pgfqpoint{4.039979in}{0.633527in}}%
\pgfpathlineto{\pgfqpoint{4.039979in}{0.625278in}}%
\pgfpathlineto{\pgfqpoint{4.040986in}{0.625278in}}%
\pgfpathlineto{\pgfqpoint{4.042171in}{0.633527in}}%
\pgfpathlineto{\pgfqpoint{4.042326in}{0.625278in}}%
\pgfpathlineto{\pgfqpoint{4.043185in}{0.625278in}}%
\pgfpathlineto{\pgfqpoint{4.043859in}{0.650025in}}%
\pgfpathlineto{\pgfqpoint{4.044340in}{0.633527in}}%
\pgfpathlineto{\pgfqpoint{4.045177in}{0.625278in}}%
\pgfpathlineto{\pgfqpoint{4.046517in}{0.633527in}}%
\pgfpathlineto{\pgfqpoint{4.046688in}{0.625278in}}%
\pgfpathlineto{\pgfqpoint{4.047524in}{0.625278in}}%
\pgfpathlineto{\pgfqpoint{4.048701in}{0.633527in}}%
\pgfpathlineto{\pgfqpoint{4.048864in}{0.625278in}}%
\pgfpathlineto{\pgfqpoint{4.049701in}{0.625278in}}%
\pgfpathlineto{\pgfqpoint{4.050878in}{0.658274in}}%
\pgfpathlineto{\pgfqpoint{4.051063in}{0.641776in}}%
\pgfpathlineto{\pgfqpoint{4.052070in}{0.625278in}}%
\pgfpathlineto{\pgfqpoint{4.053729in}{0.691271in}}%
\pgfpathlineto{\pgfqpoint{4.054062in}{0.707769in}}%
\pgfpathlineto{\pgfqpoint{4.054418in}{0.674773in}}%
\pgfpathlineto{\pgfqpoint{4.054751in}{0.674773in}}%
\pgfpathlineto{\pgfqpoint{4.055906in}{0.683022in}}%
\pgfpathlineto{\pgfqpoint{4.056928in}{0.625278in}}%
\pgfpathlineto{\pgfqpoint{4.058253in}{0.633527in}}%
\pgfpathlineto{\pgfqpoint{4.058586in}{0.625278in}}%
\pgfpathlineto{\pgfqpoint{4.059260in}{0.625278in}}%
\pgfpathlineto{\pgfqpoint{4.060430in}{0.633527in}}%
\pgfpathlineto{\pgfqpoint{4.060600in}{0.625278in}}%
\pgfpathlineto{\pgfqpoint{4.061437in}{0.625278in}}%
\pgfpathlineto{\pgfqpoint{4.062614in}{0.633527in}}%
\pgfpathlineto{\pgfqpoint{4.062777in}{0.625278in}}%
\pgfpathlineto{\pgfqpoint{4.063451in}{0.625278in}}%
\pgfpathlineto{\pgfqpoint{4.064791in}{0.633527in}}%
\pgfpathlineto{\pgfqpoint{4.064791in}{0.625278in}}%
\pgfpathlineto{\pgfqpoint{4.065798in}{0.625278in}}%
\pgfpathlineto{\pgfqpoint{4.066805in}{0.633527in}}%
\pgfpathlineto{\pgfqpoint{4.066805in}{0.625278in}}%
\pgfpathlineto{\pgfqpoint{4.067975in}{0.633527in}}%
\pgfpathlineto{\pgfqpoint{4.068145in}{0.625278in}}%
\pgfpathlineto{\pgfqpoint{4.068982in}{0.625278in}}%
\pgfpathlineto{\pgfqpoint{4.069989in}{0.633527in}}%
\pgfpathlineto{\pgfqpoint{4.069989in}{0.625278in}}%
\pgfpathlineto{\pgfqpoint{4.071158in}{0.633527in}}%
\pgfpathlineto{\pgfqpoint{4.071329in}{0.625278in}}%
\pgfpathlineto{\pgfqpoint{4.072188in}{0.625278in}}%
\pgfpathlineto{\pgfqpoint{4.074683in}{0.707769in}}%
\pgfpathlineto{\pgfqpoint{4.075519in}{0.683022in}}%
\pgfpathlineto{\pgfqpoint{4.075705in}{0.691271in}}%
\pgfpathlineto{\pgfqpoint{4.077548in}{0.625278in}}%
\pgfpathlineto{\pgfqpoint{4.078703in}{0.633527in}}%
\pgfpathlineto{\pgfqpoint{4.078874in}{0.625278in}}%
\pgfpathlineto{\pgfqpoint{4.079732in}{0.625278in}}%
\pgfpathlineto{\pgfqpoint{4.081050in}{0.757264in}}%
\pgfpathlineto{\pgfqpoint{4.082724in}{0.732517in}}%
\pgfpathlineto{\pgfqpoint{4.083753in}{0.724268in}}%
\pgfpathlineto{\pgfqpoint{4.084760in}{0.773763in}}%
\pgfpathlineto{\pgfqpoint{4.085093in}{0.749015in}}%
\pgfpathlineto{\pgfqpoint{4.085597in}{0.724268in}}%
\pgfpathlineto{\pgfqpoint{4.085915in}{0.757264in}}%
\pgfpathlineto{\pgfqpoint{4.085930in}{0.749015in}}%
\pgfpathlineto{\pgfqpoint{4.086752in}{0.757264in}}%
\pgfpathlineto{\pgfqpoint{4.086766in}{0.749015in}}%
\pgfpathlineto{\pgfqpoint{4.087270in}{0.724268in}}%
\pgfpathlineto{\pgfqpoint{4.087588in}{0.757264in}}%
\pgfpathlineto{\pgfqpoint{4.087611in}{0.749015in}}%
\pgfpathlineto{\pgfqpoint{4.087773in}{0.773763in}}%
\pgfpathlineto{\pgfqpoint{4.088425in}{0.732517in}}%
\pgfpathlineto{\pgfqpoint{4.089617in}{0.699520in}}%
\pgfpathlineto{\pgfqpoint{4.092956in}{0.831507in}}%
\pgfpathlineto{\pgfqpoint{4.093141in}{0.823258in}}%
\pgfpathlineto{\pgfqpoint{4.093304in}{0.848005in}}%
\pgfpathlineto{\pgfqpoint{4.093978in}{0.848005in}}%
\pgfpathlineto{\pgfqpoint{4.096303in}{0.782012in}}%
\pgfpathlineto{\pgfqpoint{4.096325in}{0.798510in}}%
\pgfpathlineto{\pgfqpoint{4.097332in}{0.798510in}}%
\pgfpathlineto{\pgfqpoint{4.097643in}{0.782012in}}%
\pgfpathlineto{\pgfqpoint{4.097665in}{0.823258in}}%
\pgfpathlineto{\pgfqpoint{4.097984in}{0.815008in}}%
\pgfpathlineto{\pgfqpoint{4.098650in}{0.782012in}}%
\pgfpathlineto{\pgfqpoint{4.099006in}{0.823258in}}%
\pgfpathlineto{\pgfqpoint{4.100331in}{0.856254in}}%
\pgfpathlineto{\pgfqpoint{4.099657in}{0.806759in}}%
\pgfpathlineto{\pgfqpoint{4.100346in}{0.848005in}}%
\pgfpathlineto{\pgfqpoint{4.100849in}{0.823258in}}%
\pgfpathlineto{\pgfqpoint{4.101168in}{0.856254in}}%
\pgfpathlineto{\pgfqpoint{4.101182in}{0.848005in}}%
\pgfpathlineto{\pgfqpoint{4.101834in}{0.938746in}}%
\pgfpathlineto{\pgfqpoint{4.102337in}{0.955244in}}%
\pgfpathlineto{\pgfqpoint{4.102693in}{0.922247in}}%
\pgfpathlineto{\pgfqpoint{4.102863in}{0.922247in}}%
\pgfpathlineto{\pgfqpoint{4.104855in}{0.979992in}}%
\pgfpathlineto{\pgfqpoint{4.105543in}{0.971742in}}%
\pgfpathlineto{\pgfqpoint{4.105862in}{1.004739in}}%
\pgfpathlineto{\pgfqpoint{4.105877in}{0.996490in}}%
\pgfpathlineto{\pgfqpoint{4.106195in}{1.004739in}}%
\pgfpathlineto{\pgfqpoint{4.106528in}{0.979992in}}%
\pgfpathlineto{\pgfqpoint{4.107372in}{0.955244in}}%
\pgfpathlineto{\pgfqpoint{4.107387in}{0.996490in}}%
\pgfpathlineto{\pgfqpoint{4.107557in}{0.988241in}}%
\pgfpathlineto{\pgfqpoint{4.107706in}{0.979992in}}%
\pgfpathlineto{\pgfqpoint{4.108209in}{1.012988in}}%
\pgfpathlineto{\pgfqpoint{4.108224in}{1.021237in}}%
\pgfpathlineto{\pgfqpoint{4.108564in}{0.996490in}}%
\pgfpathlineto{\pgfqpoint{4.109231in}{0.996490in}}%
\pgfpathlineto{\pgfqpoint{4.109734in}{0.971742in}}%
\pgfpathlineto{\pgfqpoint{4.110053in}{1.004739in}}%
\pgfpathlineto{\pgfqpoint{4.110067in}{0.996490in}}%
\pgfpathlineto{\pgfqpoint{4.111748in}{1.070732in}}%
\pgfpathlineto{\pgfqpoint{4.110719in}{0.979992in}}%
\pgfpathlineto{\pgfqpoint{4.112059in}{1.054234in}}%
\pgfpathlineto{\pgfqpoint{4.112733in}{1.029487in}}%
\pgfpathlineto{\pgfqpoint{4.112755in}{1.070732in}}%
\pgfpathlineto{\pgfqpoint{4.113088in}{1.070732in}}%
\pgfpathlineto{\pgfqpoint{4.113755in}{1.045985in}}%
\pgfpathlineto{\pgfqpoint{4.113925in}{1.070732in}}%
\pgfpathlineto{\pgfqpoint{4.115080in}{1.078981in}}%
\pgfpathlineto{\pgfqpoint{4.116272in}{1.037736in}}%
\pgfpathlineto{\pgfqpoint{4.117449in}{1.021237in}}%
\pgfpathlineto{\pgfqpoint{4.118449in}{1.045985in}}%
\pgfpathlineto{\pgfqpoint{4.118619in}{1.037736in}}%
\pgfpathlineto{\pgfqpoint{4.119441in}{1.004739in}}%
\pgfpathlineto{\pgfqpoint{4.118790in}{1.045985in}}%
\pgfpathlineto{\pgfqpoint{4.119774in}{1.029487in}}%
\pgfpathlineto{\pgfqpoint{4.120781in}{1.078981in}}%
\pgfpathlineto{\pgfqpoint{4.120796in}{1.070732in}}%
\pgfpathlineto{\pgfqpoint{4.121115in}{1.078981in}}%
\pgfpathlineto{\pgfqpoint{4.121448in}{1.054234in}}%
\pgfpathlineto{\pgfqpoint{4.122121in}{1.029487in}}%
\pgfpathlineto{\pgfqpoint{4.122455in}{1.087231in}}%
\pgfpathlineto{\pgfqpoint{4.123632in}{1.177971in}}%
\pgfpathlineto{\pgfqpoint{4.123647in}{1.169722in}}%
\pgfpathlineto{\pgfqpoint{4.124298in}{1.128476in}}%
\pgfpathlineto{\pgfqpoint{4.124654in}{1.169722in}}%
\pgfpathlineto{\pgfqpoint{4.124972in}{1.177971in}}%
\pgfpathlineto{\pgfqpoint{4.125305in}{1.153224in}}%
\pgfpathlineto{\pgfqpoint{4.125831in}{1.144975in}}%
\pgfpathlineto{\pgfqpoint{4.126312in}{1.186221in}}%
\pgfpathlineto{\pgfqpoint{4.126831in}{1.219217in}}%
\pgfpathlineto{\pgfqpoint{4.127334in}{1.194470in}}%
\pgfpathlineto{\pgfqpoint{4.127838in}{1.219217in}}%
\pgfpathlineto{\pgfqpoint{4.128489in}{1.202719in}}%
\pgfpathlineto{\pgfqpoint{4.129348in}{1.169722in}}%
\pgfpathlineto{\pgfqpoint{4.129518in}{1.194470in}}%
\pgfpathlineto{\pgfqpoint{4.129829in}{1.202719in}}%
\pgfpathlineto{\pgfqpoint{4.130185in}{1.169722in}}%
\pgfpathlineto{\pgfqpoint{4.130333in}{1.177971in}}%
\pgfpathlineto{\pgfqpoint{4.130518in}{1.169722in}}%
\pgfpathlineto{\pgfqpoint{4.130836in}{1.202719in}}%
\pgfpathlineto{\pgfqpoint{4.131362in}{1.186221in}}%
\pgfpathlineto{\pgfqpoint{4.131843in}{1.177971in}}%
\pgfpathlineto{\pgfqpoint{4.131865in}{1.219217in}}%
\pgfpathlineto{\pgfqpoint{4.132176in}{1.202719in}}%
\pgfpathlineto{\pgfqpoint{4.132199in}{1.219217in}}%
\pgfpathlineto{\pgfqpoint{4.132850in}{1.177971in}}%
\pgfpathlineto{\pgfqpoint{4.133206in}{1.219217in}}%
\pgfpathlineto{\pgfqpoint{4.133857in}{1.227466in}}%
\pgfpathlineto{\pgfqpoint{4.134042in}{1.210968in}}%
\pgfpathlineto{\pgfqpoint{4.134709in}{1.194470in}}%
\pgfpathlineto{\pgfqpoint{4.134879in}{1.219217in}}%
\pgfpathlineto{\pgfqpoint{4.135197in}{1.202719in}}%
\pgfpathlineto{\pgfqpoint{4.136034in}{1.227466in}}%
\pgfpathlineto{\pgfqpoint{4.136219in}{1.219217in}}%
\pgfpathlineto{\pgfqpoint{4.137374in}{1.227466in}}%
\pgfpathlineto{\pgfqpoint{4.138211in}{1.202719in}}%
\pgfpathlineto{\pgfqpoint{4.138403in}{1.219217in}}%
\pgfpathlineto{\pgfqpoint{4.139551in}{1.227466in}}%
\pgfpathlineto{\pgfqpoint{4.140743in}{1.194470in}}%
\pgfpathlineto{\pgfqpoint{4.141920in}{1.219217in}}%
\pgfpathlineto{\pgfqpoint{4.142091in}{1.210968in}}%
\pgfpathlineto{\pgfqpoint{4.146933in}{0.633527in}}%
\pgfpathlineto{\pgfqpoint{4.146933in}{0.641776in}}%
\pgfpathlineto{\pgfqpoint{4.147770in}{0.732517in}}%
\pgfpathlineto{\pgfqpoint{4.148125in}{0.699520in}}%
\pgfpathlineto{\pgfqpoint{4.151287in}{0.881002in}}%
\pgfpathlineto{\pgfqpoint{4.152146in}{0.864503in}}%
\pgfpathlineto{\pgfqpoint{4.152479in}{0.848005in}}%
\pgfpathlineto{\pgfqpoint{4.152797in}{0.881002in}}%
\pgfpathlineto{\pgfqpoint{4.152819in}{0.872753in}}%
\pgfpathlineto{\pgfqpoint{4.152967in}{0.881002in}}%
\pgfpathlineto{\pgfqpoint{4.153315in}{0.848005in}}%
\pgfpathlineto{\pgfqpoint{4.153967in}{0.831507in}}%
\pgfpathlineto{\pgfqpoint{4.154137in}{0.856254in}}%
\pgfpathlineto{\pgfqpoint{4.154159in}{0.848005in}}%
\pgfpathlineto{\pgfqpoint{4.155477in}{0.881002in}}%
\pgfpathlineto{\pgfqpoint{4.155500in}{0.872753in}}%
\pgfpathlineto{\pgfqpoint{4.155818in}{0.856254in}}%
\pgfpathlineto{\pgfqpoint{4.156151in}{0.889251in}}%
\pgfpathlineto{\pgfqpoint{4.156484in}{0.905749in}}%
\pgfpathlineto{\pgfqpoint{4.156818in}{0.881002in}}%
\pgfpathlineto{\pgfqpoint{4.157173in}{0.897500in}}%
\pgfpathlineto{\pgfqpoint{4.158335in}{0.905749in}}%
\pgfpathlineto{\pgfqpoint{4.159357in}{0.889251in}}%
\pgfpathlineto{\pgfqpoint{4.159520in}{0.897500in}}%
\pgfpathlineto{\pgfqpoint{4.160505in}{0.881002in}}%
\pgfpathlineto{\pgfqpoint{4.161534in}{0.922247in}}%
\pgfpathlineto{\pgfqpoint{4.162852in}{0.905749in}}%
\pgfpathlineto{\pgfqpoint{4.162874in}{0.922247in}}%
\pgfpathlineto{\pgfqpoint{4.163881in}{0.872753in}}%
\pgfpathlineto{\pgfqpoint{4.164866in}{0.905749in}}%
\pgfpathlineto{\pgfqpoint{4.165051in}{0.889251in}}%
\pgfpathlineto{\pgfqpoint{4.166206in}{0.856254in}}%
\pgfpathlineto{\pgfqpoint{4.166228in}{0.872753in}}%
\pgfpathlineto{\pgfqpoint{4.169916in}{1.095480in}}%
\pgfpathlineto{\pgfqpoint{4.170419in}{1.037736in}}%
\pgfpathlineto{\pgfqpoint{4.173425in}{0.633527in}}%
\pgfpathlineto{\pgfqpoint{4.173773in}{0.674773in}}%
\pgfpathlineto{\pgfqpoint{4.174773in}{0.699520in}}%
\pgfpathlineto{\pgfqpoint{4.174943in}{0.691271in}}%
\pgfpathlineto{\pgfqpoint{4.175765in}{0.683022in}}%
\pgfpathlineto{\pgfqpoint{4.175113in}{0.699520in}}%
\pgfpathlineto{\pgfqpoint{4.175772in}{0.691271in}}%
\pgfpathlineto{\pgfqpoint{4.176787in}{0.798510in}}%
\pgfpathlineto{\pgfqpoint{4.177794in}{0.823258in}}%
\pgfpathlineto{\pgfqpoint{4.177964in}{0.815008in}}%
\pgfpathlineto{\pgfqpoint{4.178119in}{0.806759in}}%
\pgfpathlineto{\pgfqpoint{4.178282in}{0.831507in}}%
\pgfpathlineto{\pgfqpoint{4.178630in}{0.823258in}}%
\pgfpathlineto{\pgfqpoint{4.180296in}{0.856254in}}%
\pgfpathlineto{\pgfqpoint{4.180311in}{0.848005in}}%
\pgfpathlineto{\pgfqpoint{4.180807in}{0.897500in}}%
\pgfpathlineto{\pgfqpoint{4.181311in}{0.889251in}}%
\pgfpathlineto{\pgfqpoint{4.181807in}{0.881002in}}%
\pgfpathlineto{\pgfqpoint{4.181962in}{0.905749in}}%
\pgfpathlineto{\pgfqpoint{4.181984in}{0.897500in}}%
\pgfpathlineto{\pgfqpoint{4.184154in}{0.955244in}}%
\pgfpathlineto{\pgfqpoint{4.184665in}{0.946995in}}%
\pgfpathlineto{\pgfqpoint{4.185168in}{0.996490in}}%
\pgfpathlineto{\pgfqpoint{4.185827in}{1.004739in}}%
\pgfpathlineto{\pgfqpoint{4.186005in}{0.988241in}}%
\pgfpathlineto{\pgfqpoint{4.187501in}{0.955244in}}%
\pgfpathlineto{\pgfqpoint{4.186175in}{0.996490in}}%
\pgfpathlineto{\pgfqpoint{4.187501in}{0.963493in}}%
\pgfpathlineto{\pgfqpoint{4.188189in}{0.996490in}}%
\pgfpathlineto{\pgfqpoint{4.188522in}{0.971742in}}%
\pgfpathlineto{\pgfqpoint{4.189011in}{0.979992in}}%
\pgfpathlineto{\pgfqpoint{4.189862in}{0.946995in}}%
\pgfpathlineto{\pgfqpoint{4.190344in}{0.979992in}}%
\pgfpathlineto{\pgfqpoint{4.190699in}{0.938746in}}%
\pgfpathlineto{\pgfqpoint{4.191706in}{0.798510in}}%
\pgfpathlineto{\pgfqpoint{4.192543in}{0.823258in}}%
\pgfpathlineto{\pgfqpoint{4.193705in}{0.905749in}}%
\pgfpathlineto{\pgfqpoint{4.194224in}{0.864503in}}%
\pgfpathlineto{\pgfqpoint{4.194386in}{0.872753in}}%
\pgfpathlineto{\pgfqpoint{4.199577in}{0.625278in}}%
\pgfpathlineto{\pgfqpoint{4.199910in}{0.633527in}}%
\pgfpathlineto{\pgfqpoint{4.199910in}{0.625278in}}%
\pgfpathlineto{\pgfqpoint{4.201250in}{0.633527in}}%
\pgfpathlineto{\pgfqpoint{4.201754in}{0.625278in}}%
\pgfpathlineto{\pgfqpoint{4.202250in}{0.625278in}}%
\pgfpathlineto{\pgfqpoint{4.203590in}{0.633527in}}%
\pgfpathlineto{\pgfqpoint{4.203760in}{0.625278in}}%
\pgfpathlineto{\pgfqpoint{4.204597in}{0.625278in}}%
\pgfpathlineto{\pgfqpoint{4.205767in}{0.633527in}}%
\pgfpathlineto{\pgfqpoint{4.205937in}{0.625278in}}%
\pgfpathlineto{\pgfqpoint{4.206603in}{0.625278in}}%
\pgfpathlineto{\pgfqpoint{4.207795in}{0.633527in}}%
\pgfpathlineto{\pgfqpoint{4.207951in}{0.625278in}}%
\pgfpathlineto{\pgfqpoint{4.208632in}{0.625278in}}%
\pgfpathlineto{\pgfqpoint{4.209809in}{0.633527in}}%
\pgfpathlineto{\pgfqpoint{4.209957in}{0.625278in}}%
\pgfpathlineto{\pgfqpoint{4.210809in}{0.625278in}}%
\pgfpathlineto{\pgfqpoint{4.211971in}{0.633527in}}%
\pgfpathlineto{\pgfqpoint{4.211994in}{0.625278in}}%
\pgfpathlineto{\pgfqpoint{4.212978in}{0.625278in}}%
\pgfpathlineto{\pgfqpoint{4.214156in}{0.633527in}}%
\pgfpathlineto{\pgfqpoint{4.214319in}{0.625278in}}%
\pgfpathlineto{\pgfqpoint{4.214992in}{0.625278in}}%
\pgfpathlineto{\pgfqpoint{4.216332in}{0.633527in}}%
\pgfpathlineto{\pgfqpoint{4.216332in}{0.625278in}}%
\pgfpathlineto{\pgfqpoint{4.217354in}{0.625278in}}%
\pgfpathlineto{\pgfqpoint{4.218865in}{0.658274in}}%
\pgfpathlineto{\pgfqpoint{4.219531in}{0.650025in}}%
\pgfpathlineto{\pgfqpoint{4.219679in}{0.666524in}}%
\pgfpathlineto{\pgfqpoint{4.220879in}{0.707769in}}%
\pgfpathlineto{\pgfqpoint{4.222049in}{0.699520in}}%
\pgfpathlineto{\pgfqpoint{4.223892in}{0.757264in}}%
\pgfpathlineto{\pgfqpoint{4.224225in}{0.749015in}}%
\pgfpathlineto{\pgfqpoint{4.224396in}{0.773763in}}%
\pgfpathlineto{\pgfqpoint{4.224551in}{0.782012in}}%
\pgfpathlineto{\pgfqpoint{4.224729in}{0.757264in}}%
\pgfpathlineto{\pgfqpoint{4.225232in}{0.773763in}}%
\pgfpathlineto{\pgfqpoint{4.229090in}{0.625278in}}%
\pgfpathlineto{\pgfqpoint{4.230593in}{0.658274in}}%
\pgfpathlineto{\pgfqpoint{4.230600in}{0.650025in}}%
\pgfpathlineto{\pgfqpoint{4.231770in}{0.625278in}}%
\pgfpathlineto{\pgfqpoint{4.232940in}{0.633527in}}%
\pgfpathlineto{\pgfqpoint{4.233096in}{0.625278in}}%
\pgfpathlineto{\pgfqpoint{4.233932in}{0.625278in}}%
\pgfpathlineto{\pgfqpoint{4.235102in}{0.633527in}}%
\pgfpathlineto{\pgfqpoint{4.235443in}{0.625278in}}%
\pgfpathlineto{\pgfqpoint{4.236109in}{0.625278in}}%
\pgfpathlineto{\pgfqpoint{4.237294in}{0.633527in}}%
\pgfpathlineto{\pgfqpoint{4.237449in}{0.625278in}}%
\pgfpathlineto{\pgfqpoint{4.238293in}{0.625278in}}%
\pgfpathlineto{\pgfqpoint{4.239633in}{0.633527in}}%
\pgfpathlineto{\pgfqpoint{4.239796in}{0.625278in}}%
\pgfpathlineto{\pgfqpoint{4.240130in}{0.625278in}}%
\pgfpathlineto{\pgfqpoint{4.241477in}{0.633527in}}%
\pgfpathlineto{\pgfqpoint{4.241640in}{0.625278in}}%
\pgfpathlineto{\pgfqpoint{4.242477in}{0.625278in}}%
\pgfpathlineto{\pgfqpoint{4.243484in}{0.633527in}}%
\pgfpathlineto{\pgfqpoint{4.243484in}{0.625278in}}%
\pgfpathlineto{\pgfqpoint{4.244668in}{0.633527in}}%
\pgfpathlineto{\pgfqpoint{4.244824in}{0.625278in}}%
\pgfpathlineto{\pgfqpoint{4.245675in}{0.625278in}}%
\pgfpathlineto{\pgfqpoint{4.246838in}{0.633527in}}%
\pgfpathlineto{\pgfqpoint{4.246845in}{0.625278in}}%
\pgfpathlineto{\pgfqpoint{4.247860in}{0.625278in}}%
\pgfpathlineto{\pgfqpoint{4.249022in}{0.633527in}}%
\pgfpathlineto{\pgfqpoint{4.249192in}{0.625278in}}%
\pgfpathlineto{\pgfqpoint{4.250022in}{0.625278in}}%
\pgfpathlineto{\pgfqpoint{4.251029in}{0.633527in}}%
\pgfpathlineto{\pgfqpoint{4.251029in}{0.625278in}}%
\pgfpathlineto{\pgfqpoint{4.252206in}{0.633527in}}%
\pgfpathlineto{\pgfqpoint{4.252539in}{0.625278in}}%
\pgfpathlineto{\pgfqpoint{4.253205in}{0.625278in}}%
\pgfpathlineto{\pgfqpoint{4.254212in}{0.633527in}}%
\pgfpathlineto{\pgfqpoint{4.254212in}{0.625278in}}%
\pgfpathlineto{\pgfqpoint{4.255390in}{0.633527in}}%
\pgfpathlineto{\pgfqpoint{4.255723in}{0.625278in}}%
\pgfpathlineto{\pgfqpoint{4.256397in}{0.625278in}}%
\pgfpathlineto{\pgfqpoint{4.257566in}{0.633527in}}%
\pgfpathlineto{\pgfqpoint{4.257737in}{0.625278in}}%
\pgfpathlineto{\pgfqpoint{4.258573in}{0.625278in}}%
\pgfpathlineto{\pgfqpoint{4.259743in}{0.633527in}}%
\pgfpathlineto{\pgfqpoint{4.259743in}{0.625278in}}%
\pgfpathlineto{\pgfqpoint{4.260750in}{0.625278in}}%
\pgfpathlineto{\pgfqpoint{4.261928in}{0.633527in}}%
\pgfpathlineto{\pgfqpoint{4.262431in}{0.625278in}}%
\pgfpathlineto{\pgfqpoint{4.262934in}{0.625278in}}%
\pgfpathlineto{\pgfqpoint{4.264275in}{0.633527in}}%
\pgfpathlineto{\pgfqpoint{4.264275in}{0.625278in}}%
\pgfpathlineto{\pgfqpoint{4.265274in}{0.625278in}}%
\pgfpathlineto{\pgfqpoint{4.266622in}{0.633527in}}%
\pgfpathlineto{\pgfqpoint{4.266785in}{0.625278in}}%
\pgfpathlineto{\pgfqpoint{4.267621in}{0.625278in}}%
\pgfpathlineto{\pgfqpoint{4.268799in}{0.633527in}}%
\pgfpathlineto{\pgfqpoint{4.268969in}{0.625278in}}%
\pgfpathlineto{\pgfqpoint{4.269806in}{0.625278in}}%
\pgfpathlineto{\pgfqpoint{4.270983in}{0.633527in}}%
\pgfpathlineto{\pgfqpoint{4.271146in}{0.625278in}}%
\pgfpathlineto{\pgfqpoint{4.271982in}{0.625278in}}%
\pgfpathlineto{\pgfqpoint{4.273160in}{0.633527in}}%
\pgfpathlineto{\pgfqpoint{4.273323in}{0.625278in}}%
\pgfpathlineto{\pgfqpoint{4.274159in}{0.625278in}}%
\pgfpathlineto{\pgfqpoint{4.275336in}{0.633527in}}%
\pgfpathlineto{\pgfqpoint{4.275336in}{0.625278in}}%
\pgfpathlineto{\pgfqpoint{4.276173in}{0.625278in}}%
\pgfpathlineto{\pgfqpoint{4.277513in}{0.633527in}}%
\pgfpathlineto{\pgfqpoint{4.277513in}{0.625278in}}%
\pgfpathlineto{\pgfqpoint{4.278520in}{0.625278in}}%
\pgfpathlineto{\pgfqpoint{4.279690in}{0.633527in}}%
\pgfpathlineto{\pgfqpoint{4.279860in}{0.625278in}}%
\pgfpathlineto{\pgfqpoint{4.280705in}{0.625278in}}%
\pgfpathlineto{\pgfqpoint{4.281874in}{0.633527in}}%
\pgfpathlineto{\pgfqpoint{4.281874in}{0.625278in}}%
\pgfpathlineto{\pgfqpoint{4.282881in}{0.625278in}}%
\pgfpathlineto{\pgfqpoint{4.284051in}{0.633527in}}%
\pgfpathlineto{\pgfqpoint{4.284051in}{0.625278in}}%
\pgfpathlineto{\pgfqpoint{4.285058in}{0.625278in}}%
\pgfpathlineto{\pgfqpoint{4.286287in}{0.633527in}}%
\pgfpathlineto{\pgfqpoint{4.286398in}{0.625278in}}%
\pgfpathlineto{\pgfqpoint{4.287294in}{0.625278in}}%
\pgfpathlineto{\pgfqpoint{4.288412in}{0.633527in}}%
\pgfpathlineto{\pgfqpoint{4.288412in}{0.625278in}}%
\pgfpathlineto{\pgfqpoint{4.289419in}{0.625278in}}%
\pgfpathlineto{\pgfqpoint{4.290759in}{0.633527in}}%
\pgfpathlineto{\pgfqpoint{4.290922in}{0.625278in}}%
\pgfpathlineto{\pgfqpoint{4.291766in}{0.625278in}}%
\pgfpathlineto{\pgfqpoint{4.292766in}{0.633527in}}%
\pgfpathlineto{\pgfqpoint{4.292766in}{0.625278in}}%
\pgfpathlineto{\pgfqpoint{4.293943in}{0.633527in}}%
\pgfpathlineto{\pgfqpoint{4.293943in}{0.625278in}}%
\pgfpathlineto{\pgfqpoint{4.294950in}{0.625278in}}%
\pgfpathlineto{\pgfqpoint{4.296120in}{0.633527in}}%
\pgfpathlineto{\pgfqpoint{4.296290in}{0.625278in}}%
\pgfpathlineto{\pgfqpoint{4.296957in}{0.625278in}}%
\pgfpathlineto{\pgfqpoint{4.298304in}{0.633527in}}%
\pgfpathlineto{\pgfqpoint{4.298304in}{0.625278in}}%
\pgfpathlineto{\pgfqpoint{4.299304in}{0.625278in}}%
\pgfpathlineto{\pgfqpoint{4.300481in}{0.633527in}}%
\pgfpathlineto{\pgfqpoint{4.300644in}{0.625278in}}%
\pgfpathlineto{\pgfqpoint{4.301488in}{0.625278in}}%
\pgfpathlineto{\pgfqpoint{4.302658in}{0.633527in}}%
\pgfpathlineto{\pgfqpoint{4.302828in}{0.625278in}}%
\pgfpathlineto{\pgfqpoint{4.303665in}{0.625278in}}%
\pgfpathlineto{\pgfqpoint{4.304835in}{0.633527in}}%
\pgfpathlineto{\pgfqpoint{4.305175in}{0.625278in}}%
\pgfpathlineto{\pgfqpoint{4.305842in}{0.625278in}}%
\pgfpathlineto{\pgfqpoint{4.306849in}{0.633527in}}%
\pgfpathlineto{\pgfqpoint{4.306849in}{0.625278in}}%
\pgfpathlineto{\pgfqpoint{4.308026in}{0.633527in}}%
\pgfpathlineto{\pgfqpoint{4.308189in}{0.625278in}}%
\pgfpathlineto{\pgfqpoint{4.309026in}{0.625278in}}%
\pgfpathlineto{\pgfqpoint{4.310373in}{0.633527in}}%
\pgfpathlineto{\pgfqpoint{4.310373in}{0.625278in}}%
\pgfpathlineto{\pgfqpoint{4.311210in}{0.625278in}}%
\pgfpathlineto{\pgfqpoint{4.312550in}{0.633527in}}%
\pgfpathlineto{\pgfqpoint{4.312609in}{0.625278in}}%
\pgfpathlineto{\pgfqpoint{4.313557in}{0.625278in}}%
\pgfpathlineto{\pgfqpoint{4.314564in}{0.633527in}}%
\pgfpathlineto{\pgfqpoint{4.314564in}{0.625278in}}%
\pgfpathlineto{\pgfqpoint{4.315904in}{0.633527in}}%
\pgfpathlineto{\pgfqpoint{4.315904in}{0.625278in}}%
\pgfpathlineto{\pgfqpoint{4.316911in}{0.625278in}}%
\pgfpathlineto{\pgfqpoint{4.318081in}{0.633527in}}%
\pgfpathlineto{\pgfqpoint{4.318081in}{0.625278in}}%
\pgfpathlineto{\pgfqpoint{4.319088in}{0.625278in}}%
\pgfpathlineto{\pgfqpoint{4.320258in}{0.633527in}}%
\pgfpathlineto{\pgfqpoint{4.320428in}{0.625278in}}%
\pgfpathlineto{\pgfqpoint{4.321265in}{0.625278in}}%
\pgfpathlineto{\pgfqpoint{4.322442in}{0.633527in}}%
\pgfpathlineto{\pgfqpoint{4.323449in}{0.625278in}}%
\pgfpathlineto{\pgfqpoint{4.324619in}{0.633527in}}%
\pgfpathlineto{\pgfqpoint{4.324641in}{0.625278in}}%
\pgfpathlineto{\pgfqpoint{4.325293in}{0.683022in}}%
\pgfpathlineto{\pgfqpoint{4.325478in}{0.674773in}}%
\pgfpathlineto{\pgfqpoint{4.325981in}{0.650025in}}%
\pgfpathlineto{\pgfqpoint{4.326633in}{0.683022in}}%
\pgfpathlineto{\pgfqpoint{4.327655in}{0.650025in}}%
\pgfpathlineto{\pgfqpoint{4.328832in}{0.674773in}}%
\pgfpathlineto{\pgfqpoint{4.328995in}{0.666524in}}%
\pgfpathlineto{\pgfqpoint{4.329498in}{0.650025in}}%
\pgfpathlineto{\pgfqpoint{4.329706in}{0.683022in}}%
\pgfpathlineto{\pgfqpoint{4.329831in}{0.674773in}}%
\pgfpathlineto{\pgfqpoint{4.330335in}{0.699520in}}%
\pgfpathlineto{\pgfqpoint{4.330986in}{0.683022in}}%
\pgfpathlineto{\pgfqpoint{4.331179in}{0.674773in}}%
\pgfpathlineto{\pgfqpoint{4.331831in}{0.732517in}}%
\pgfpathlineto{\pgfqpoint{4.331845in}{0.724268in}}%
\pgfpathlineto{\pgfqpoint{4.332016in}{0.749015in}}%
\pgfpathlineto{\pgfqpoint{4.333000in}{0.732517in}}%
\pgfpathlineto{\pgfqpoint{4.333519in}{0.724268in}}%
\pgfpathlineto{\pgfqpoint{4.333689in}{0.749015in}}%
\pgfpathlineto{\pgfqpoint{4.334022in}{0.749015in}}%
\pgfpathlineto{\pgfqpoint{4.335347in}{0.732517in}}%
\pgfpathlineto{\pgfqpoint{4.336184in}{0.782012in}}%
\pgfpathlineto{\pgfqpoint{4.336369in}{0.773763in}}%
\pgfpathlineto{\pgfqpoint{4.337191in}{0.782012in}}%
\pgfpathlineto{\pgfqpoint{4.337213in}{0.773763in}}%
\pgfpathlineto{\pgfqpoint{4.337865in}{0.757264in}}%
\pgfpathlineto{\pgfqpoint{4.338028in}{0.782012in}}%
\pgfpathlineto{\pgfqpoint{4.338050in}{0.773763in}}%
\pgfpathlineto{\pgfqpoint{4.339390in}{0.823258in}}%
\pgfpathlineto{\pgfqpoint{4.339709in}{0.806759in}}%
\pgfpathlineto{\pgfqpoint{4.340057in}{0.798510in}}%
\pgfpathlineto{\pgfqpoint{4.340375in}{0.831507in}}%
\pgfpathlineto{\pgfqpoint{4.340730in}{0.823258in}}%
\pgfpathlineto{\pgfqpoint{4.341049in}{0.806759in}}%
\pgfpathlineto{\pgfqpoint{4.341382in}{0.839756in}}%
\pgfpathlineto{\pgfqpoint{4.341552in}{0.881002in}}%
\pgfpathlineto{\pgfqpoint{4.342404in}{0.823258in}}%
\pgfpathlineto{\pgfqpoint{4.342722in}{0.806759in}}%
\pgfpathlineto{\pgfqpoint{4.343055in}{0.839756in}}%
\pgfpathlineto{\pgfqpoint{4.344062in}{0.881002in}}%
\pgfpathlineto{\pgfqpoint{4.343559in}{0.831507in}}%
\pgfpathlineto{\pgfqpoint{4.344084in}{0.872753in}}%
\pgfpathlineto{\pgfqpoint{4.344588in}{0.848005in}}%
\pgfpathlineto{\pgfqpoint{4.344899in}{0.881002in}}%
\pgfpathlineto{\pgfqpoint{4.344921in}{0.872753in}}%
\pgfpathlineto{\pgfqpoint{4.345069in}{0.881002in}}%
\pgfpathlineto{\pgfqpoint{4.345425in}{0.848005in}}%
\pgfpathlineto{\pgfqpoint{4.345573in}{0.856254in}}%
\pgfpathlineto{\pgfqpoint{4.345595in}{0.848005in}}%
\pgfpathlineto{\pgfqpoint{4.345906in}{0.881002in}}%
\pgfpathlineto{\pgfqpoint{4.346594in}{0.872753in}}%
\pgfpathlineto{\pgfqpoint{4.347750in}{0.881002in}}%
\pgfpathlineto{\pgfqpoint{4.348942in}{0.848005in}}%
\pgfpathlineto{\pgfqpoint{4.350282in}{0.922247in}}%
\pgfpathlineto{\pgfqpoint{4.350933in}{0.905749in}}%
\pgfpathlineto{\pgfqpoint{4.351126in}{0.897500in}}%
\pgfpathlineto{\pgfqpoint{4.351777in}{0.955244in}}%
\pgfpathlineto{\pgfqpoint{4.351792in}{0.946995in}}%
\pgfpathlineto{\pgfqpoint{4.352111in}{0.955244in}}%
\pgfpathlineto{\pgfqpoint{4.352444in}{0.930497in}}%
\pgfpathlineto{\pgfqpoint{4.353473in}{0.922247in}}%
\pgfpathlineto{\pgfqpoint{4.355124in}{0.955244in}}%
\pgfpathlineto{\pgfqpoint{4.355650in}{0.946995in}}%
\pgfpathlineto{\pgfqpoint{4.355820in}{0.971742in}}%
\pgfpathlineto{\pgfqpoint{4.356153in}{0.971742in}}%
\pgfpathlineto{\pgfqpoint{4.356361in}{0.955244in}}%
\pgfpathlineto{\pgfqpoint{4.356486in}{0.996490in}}%
\pgfpathlineto{\pgfqpoint{4.356805in}{0.988241in}}%
\pgfpathlineto{\pgfqpoint{4.357812in}{1.004739in}}%
\pgfpathlineto{\pgfqpoint{4.357827in}{0.996490in}}%
\pgfpathlineto{\pgfqpoint{4.358145in}{1.004739in}}%
\pgfpathlineto{\pgfqpoint{4.358478in}{0.979992in}}%
\pgfpathlineto{\pgfqpoint{4.359485in}{0.955244in}}%
\pgfpathlineto{\pgfqpoint{4.359507in}{0.971742in}}%
\pgfpathlineto{\pgfqpoint{4.359818in}{0.955244in}}%
\pgfpathlineto{\pgfqpoint{4.359989in}{0.979992in}}%
\pgfpathlineto{\pgfqpoint{4.360344in}{0.971742in}}%
\pgfpathlineto{\pgfqpoint{4.361514in}{0.996490in}}%
\pgfpathlineto{\pgfqpoint{4.360825in}{0.955244in}}%
\pgfpathlineto{\pgfqpoint{4.361684in}{0.988241in}}%
\pgfpathlineto{\pgfqpoint{4.363172in}{0.955244in}}%
\pgfpathlineto{\pgfqpoint{4.361855in}{0.996490in}}%
\pgfpathlineto{\pgfqpoint{4.363172in}{0.963493in}}%
\pgfpathlineto{\pgfqpoint{4.364513in}{1.054234in}}%
\pgfpathlineto{\pgfqpoint{4.365520in}{1.029487in}}%
\pgfpathlineto{\pgfqpoint{4.365542in}{1.070732in}}%
\pgfpathlineto{\pgfqpoint{4.366527in}{1.004739in}}%
\pgfpathlineto{\pgfqpoint{4.367045in}{1.045985in}}%
\pgfpathlineto{\pgfqpoint{4.368200in}{1.054234in}}%
\pgfpathlineto{\pgfqpoint{4.369229in}{1.045985in}}%
\pgfpathlineto{\pgfqpoint{4.370384in}{1.054234in}}%
\pgfpathlineto{\pgfqpoint{4.371073in}{1.021237in}}%
\pgfpathlineto{\pgfqpoint{4.371406in}{1.045985in}}%
\pgfpathlineto{\pgfqpoint{4.372228in}{1.054234in}}%
\pgfpathlineto{\pgfqpoint{4.372243in}{1.045985in}}%
\pgfpathlineto{\pgfqpoint{4.372746in}{1.021237in}}%
\pgfpathlineto{\pgfqpoint{4.373064in}{1.054234in}}%
\pgfpathlineto{\pgfqpoint{4.373087in}{1.045985in}}%
\pgfpathlineto{\pgfqpoint{4.373398in}{1.054234in}}%
\pgfpathlineto{\pgfqpoint{4.373753in}{1.021237in}}%
\pgfpathlineto{\pgfqpoint{4.375093in}{0.996490in}}%
\pgfpathlineto{\pgfqpoint{4.375241in}{1.004739in}}%
\pgfpathlineto{\pgfqpoint{4.375597in}{0.971742in}}%
\pgfpathlineto{\pgfqpoint{4.375745in}{0.979992in}}%
\pgfpathlineto{\pgfqpoint{4.376604in}{0.971742in}}%
\pgfpathlineto{\pgfqpoint{4.376774in}{0.996490in}}%
\pgfpathlineto{\pgfqpoint{4.377440in}{1.021237in}}%
\pgfpathlineto{\pgfqpoint{4.377929in}{1.004739in}}%
\pgfpathlineto{\pgfqpoint{4.378262in}{0.979992in}}%
\pgfpathlineto{\pgfqpoint{4.378951in}{0.996490in}}%
\pgfpathlineto{\pgfqpoint{4.380609in}{1.029487in}}%
\pgfpathlineto{\pgfqpoint{4.380794in}{1.045985in}}%
\pgfpathlineto{\pgfqpoint{4.381631in}{1.021237in}}%
\pgfpathlineto{\pgfqpoint{4.382283in}{1.029487in}}%
\pgfpathlineto{\pgfqpoint{4.382468in}{1.012988in}}%
\pgfpathlineto{\pgfqpoint{4.382616in}{1.004739in}}%
\pgfpathlineto{\pgfqpoint{4.382971in}{1.045985in}}%
\pgfpathlineto{\pgfqpoint{4.383645in}{1.070732in}}%
\pgfpathlineto{\pgfqpoint{4.384126in}{1.054234in}}%
\pgfpathlineto{\pgfqpoint{4.384149in}{1.045985in}}%
\pgfpathlineto{\pgfqpoint{4.384985in}{1.144975in}}%
\pgfpathlineto{\pgfqpoint{4.385304in}{1.153224in}}%
\pgfpathlineto{\pgfqpoint{4.385637in}{1.128476in}}%
\pgfpathlineto{\pgfqpoint{4.386140in}{1.153224in}}%
\pgfpathlineto{\pgfqpoint{4.386659in}{1.120227in}}%
\pgfpathlineto{\pgfqpoint{4.386807in}{1.128476in}}%
\pgfpathlineto{\pgfqpoint{4.387147in}{1.103729in}}%
\pgfpathlineto{\pgfqpoint{4.387332in}{1.111978in}}%
\pgfpathlineto{\pgfqpoint{4.387480in}{1.103729in}}%
\pgfpathlineto{\pgfqpoint{4.387651in}{1.128476in}}%
\pgfpathlineto{\pgfqpoint{4.388169in}{1.120227in}}%
\pgfpathlineto{\pgfqpoint{4.389494in}{1.153224in}}%
\pgfpathlineto{\pgfqpoint{4.388650in}{1.103729in}}%
\pgfpathlineto{\pgfqpoint{4.389509in}{1.144975in}}%
\pgfpathlineto{\pgfqpoint{4.389850in}{1.120227in}}%
\pgfpathlineto{\pgfqpoint{4.390161in}{1.153224in}}%
\pgfpathlineto{\pgfqpoint{4.390331in}{1.153224in}}%
\pgfpathlineto{\pgfqpoint{4.390501in}{1.177971in}}%
\pgfpathlineto{\pgfqpoint{4.391353in}{1.169722in}}%
\pgfpathlineto{\pgfqpoint{4.393011in}{1.202719in}}%
\pgfpathlineto{\pgfqpoint{4.393537in}{1.169722in}}%
\pgfpathlineto{\pgfqpoint{4.394041in}{1.194470in}}%
\pgfpathlineto{\pgfqpoint{4.394692in}{1.202719in}}%
\pgfpathlineto{\pgfqpoint{4.394877in}{1.186221in}}%
\pgfpathlineto{\pgfqpoint{4.395529in}{1.202719in}}%
\pgfpathlineto{\pgfqpoint{4.396047in}{1.169722in}}%
\pgfpathlineto{\pgfqpoint{4.397202in}{1.177971in}}%
\pgfpathlineto{\pgfqpoint{4.398379in}{1.153224in}}%
\pgfpathlineto{\pgfqpoint{4.398394in}{1.169722in}}%
\pgfpathlineto{\pgfqpoint{4.398727in}{1.144975in}}%
\pgfpathlineto{\pgfqpoint{4.399401in}{1.169722in}}%
\pgfpathlineto{\pgfqpoint{4.402252in}{0.988241in}}%
\pgfpathlineto{\pgfqpoint{4.405265in}{0.625278in}}%
\pgfpathlineto{\pgfqpoint{4.405606in}{0.650025in}}%
\pgfpathlineto{\pgfqpoint{4.405917in}{0.658274in}}%
\pgfpathlineto{\pgfqpoint{4.406272in}{0.625278in}}%
\pgfpathlineto{\pgfqpoint{4.407264in}{0.658274in}}%
\pgfpathlineto{\pgfqpoint{4.407598in}{0.633527in}}%
\pgfpathlineto{\pgfqpoint{4.407783in}{0.625278in}}%
\pgfpathlineto{\pgfqpoint{4.408101in}{0.658274in}}%
\pgfpathlineto{\pgfqpoint{4.408619in}{0.650025in}}%
\pgfpathlineto{\pgfqpoint{4.410456in}{0.732517in}}%
\pgfpathlineto{\pgfqpoint{4.410633in}{0.716019in}}%
\pgfpathlineto{\pgfqpoint{4.411788in}{0.658274in}}%
\pgfpathlineto{\pgfqpoint{4.412129in}{0.691271in}}%
\pgfpathlineto{\pgfqpoint{4.412647in}{0.724268in}}%
\pgfpathlineto{\pgfqpoint{4.413143in}{0.699520in}}%
\pgfpathlineto{\pgfqpoint{4.414476in}{0.732517in}}%
\pgfpathlineto{\pgfqpoint{4.414491in}{0.724268in}}%
\pgfpathlineto{\pgfqpoint{4.415476in}{0.707769in}}%
\pgfpathlineto{\pgfqpoint{4.415476in}{0.716019in}}%
\pgfpathlineto{\pgfqpoint{4.415831in}{0.773763in}}%
\pgfpathlineto{\pgfqpoint{4.416497in}{0.724268in}}%
\pgfpathlineto{\pgfqpoint{4.417660in}{0.757264in}}%
\pgfpathlineto{\pgfqpoint{4.417838in}{0.740766in}}%
\pgfpathlineto{\pgfqpoint{4.419015in}{0.724268in}}%
\pgfpathlineto{\pgfqpoint{4.420000in}{0.732517in}}%
\pgfpathlineto{\pgfqpoint{4.420022in}{0.724268in}}%
\pgfpathlineto{\pgfqpoint{4.420525in}{0.699520in}}%
\pgfpathlineto{\pgfqpoint{4.420836in}{0.732517in}}%
\pgfpathlineto{\pgfqpoint{4.420859in}{0.724268in}}%
\pgfpathlineto{\pgfqpoint{4.422680in}{0.782012in}}%
\pgfpathlineto{\pgfqpoint{4.422702in}{0.773763in}}%
\pgfpathlineto{\pgfqpoint{4.423524in}{0.732517in}}%
\pgfpathlineto{\pgfqpoint{4.423865in}{0.765513in}}%
\pgfpathlineto{\pgfqpoint{4.424716in}{0.823258in}}%
\pgfpathlineto{\pgfqpoint{4.424879in}{0.815008in}}%
\pgfpathlineto{\pgfqpoint{4.425701in}{0.806759in}}%
\pgfpathlineto{\pgfqpoint{4.425049in}{0.823258in}}%
\pgfpathlineto{\pgfqpoint{4.425716in}{0.823258in}}%
\pgfpathlineto{\pgfqpoint{4.428048in}{0.930497in}}%
\pgfpathlineto{\pgfqpoint{4.428381in}{0.905749in}}%
\pgfpathlineto{\pgfqpoint{4.428900in}{0.930497in}}%
\pgfpathlineto{\pgfqpoint{4.429558in}{0.881002in}}%
\pgfpathlineto{\pgfqpoint{4.430225in}{0.905749in}}%
\pgfpathlineto{\pgfqpoint{4.430580in}{0.872753in}}%
\pgfpathlineto{\pgfqpoint{4.430728in}{0.881002in}}%
\pgfpathlineto{\pgfqpoint{4.431084in}{0.848005in}}%
\pgfpathlineto{\pgfqpoint{4.431247in}{0.856254in}}%
\pgfpathlineto{\pgfqpoint{4.431906in}{0.806759in}}%
\pgfpathlineto{\pgfqpoint{4.432254in}{0.823258in}}%
\pgfpathlineto{\pgfqpoint{4.433912in}{0.856254in}}%
\pgfpathlineto{\pgfqpoint{4.434601in}{0.848005in}}%
\pgfpathlineto{\pgfqpoint{4.434919in}{0.881002in}}%
\pgfpathlineto{\pgfqpoint{4.434941in}{0.872753in}}%
\pgfpathlineto{\pgfqpoint{4.435104in}{0.881002in}}%
\pgfpathlineto{\pgfqpoint{4.435423in}{0.856254in}}%
\pgfpathlineto{\pgfqpoint{4.435608in}{0.864503in}}%
\pgfpathlineto{\pgfqpoint{4.435756in}{0.856254in}}%
\pgfpathlineto{\pgfqpoint{4.435778in}{0.897500in}}%
\pgfpathlineto{\pgfqpoint{4.436096in}{0.889251in}}%
\pgfpathlineto{\pgfqpoint{4.437118in}{0.971742in}}%
\pgfpathlineto{\pgfqpoint{4.438288in}{0.979992in}}%
\pgfpathlineto{\pgfqpoint{4.438621in}{1.004739in}}%
\pgfpathlineto{\pgfqpoint{4.439295in}{0.971742in}}%
\pgfpathlineto{\pgfqpoint{4.440961in}{1.029487in}}%
\pgfpathlineto{\pgfqpoint{4.441139in}{1.012988in}}%
\pgfpathlineto{\pgfqpoint{4.441812in}{0.971742in}}%
\pgfpathlineto{\pgfqpoint{4.442316in}{0.996490in}}%
\pgfpathlineto{\pgfqpoint{4.443641in}{1.029487in}}%
\pgfpathlineto{\pgfqpoint{4.443656in}{1.021237in}}%
\pgfpathlineto{\pgfqpoint{4.444315in}{1.004739in}}%
\pgfpathlineto{\pgfqpoint{4.444493in}{1.029487in}}%
\pgfpathlineto{\pgfqpoint{4.444493in}{1.021237in}}%
\pgfpathlineto{\pgfqpoint{4.445648in}{1.029487in}}%
\pgfpathlineto{\pgfqpoint{4.446173in}{1.021237in}}%
\pgfpathlineto{\pgfqpoint{4.446336in}{1.045985in}}%
\pgfpathlineto{\pgfqpoint{4.446670in}{1.045985in}}%
\pgfpathlineto{\pgfqpoint{4.447514in}{1.021237in}}%
\pgfpathlineto{\pgfqpoint{4.447995in}{1.029487in}}%
\pgfpathlineto{\pgfqpoint{4.448350in}{1.045985in}}%
\pgfpathlineto{\pgfqpoint{4.448684in}{1.021237in}}%
\pgfpathlineto{\pgfqpoint{4.449017in}{1.021237in}}%
\pgfpathlineto{\pgfqpoint{4.451364in}{0.946995in}}%
\pgfpathlineto{\pgfqpoint{4.454052in}{1.070732in}}%
\pgfpathlineto{\pgfqpoint{4.454363in}{1.054234in}}%
\pgfpathlineto{\pgfqpoint{4.456058in}{0.815008in}}%
\pgfpathlineto{\pgfqpoint{4.457739in}{0.699520in}}%
\pgfpathlineto{\pgfqpoint{4.457887in}{0.707769in}}%
\pgfpathlineto{\pgfqpoint{4.458576in}{0.749015in}}%
\pgfpathlineto{\pgfqpoint{4.458909in}{0.724268in}}%
\pgfpathlineto{\pgfqpoint{4.460923in}{0.798510in}}%
\pgfpathlineto{\pgfqpoint{4.461093in}{0.790261in}}%
\pgfpathlineto{\pgfqpoint{4.461907in}{0.757264in}}%
\pgfpathlineto{\pgfqpoint{4.462241in}{0.782012in}}%
\pgfpathlineto{\pgfqpoint{4.462751in}{0.806759in}}%
\pgfpathlineto{\pgfqpoint{4.463270in}{0.749015in}}%
\pgfpathlineto{\pgfqpoint{4.464099in}{0.658274in}}%
\pgfpathlineto{\pgfqpoint{4.464610in}{0.650025in}}%
\pgfpathlineto{\pgfqpoint{4.464780in}{0.674773in}}%
\pgfpathlineto{\pgfqpoint{4.465113in}{0.674773in}}%
\pgfpathlineto{\pgfqpoint{4.465447in}{0.699520in}}%
\pgfpathlineto{\pgfqpoint{4.465950in}{0.666524in}}%
\pgfpathlineto{\pgfqpoint{4.466120in}{0.674773in}}%
\pgfpathlineto{\pgfqpoint{4.467623in}{0.625278in}}%
\pgfpathlineto{\pgfqpoint{4.469785in}{0.683022in}}%
\pgfpathlineto{\pgfqpoint{4.470126in}{0.707769in}}%
\pgfpathlineto{\pgfqpoint{4.470815in}{0.674773in}}%
\pgfpathlineto{\pgfqpoint{4.471126in}{0.683022in}}%
\pgfpathlineto{\pgfqpoint{4.471481in}{0.650025in}}%
\pgfpathlineto{\pgfqpoint{4.471629in}{0.658274in}}%
\pgfpathlineto{\pgfqpoint{4.472318in}{0.650025in}}%
\pgfpathlineto{\pgfqpoint{4.472488in}{0.674773in}}%
\pgfpathlineto{\pgfqpoint{4.472658in}{0.666524in}}%
\pgfpathlineto{\pgfqpoint{4.473140in}{0.658274in}}%
\pgfpathlineto{\pgfqpoint{4.473162in}{0.699520in}}%
\pgfpathlineto{\pgfqpoint{4.473495in}{0.724268in}}%
\pgfpathlineto{\pgfqpoint{4.473998in}{0.691271in}}%
\pgfpathlineto{\pgfqpoint{4.474480in}{0.683022in}}%
\pgfpathlineto{\pgfqpoint{4.474502in}{0.724268in}}%
\pgfpathlineto{\pgfqpoint{4.477338in}{0.905749in}}%
\pgfpathlineto{\pgfqpoint{4.477834in}{0.856254in}}%
\pgfpathlineto{\pgfqpoint{4.478856in}{0.806759in}}%
\pgfpathlineto{\pgfqpoint{4.479026in}{0.823258in}}%
\pgfpathlineto{\pgfqpoint{4.479855in}{0.831507in}}%
\pgfpathlineto{\pgfqpoint{4.479863in}{0.823258in}}%
\pgfpathlineto{\pgfqpoint{4.480744in}{0.782012in}}%
\pgfpathlineto{\pgfqpoint{4.481077in}{0.806759in}}%
\pgfpathlineto{\pgfqpoint{4.481351in}{0.831507in}}%
\pgfpathlineto{\pgfqpoint{4.481877in}{0.790261in}}%
\pgfpathlineto{\pgfqpoint{4.482025in}{0.782012in}}%
\pgfpathlineto{\pgfqpoint{4.482380in}{0.823258in}}%
\pgfpathlineto{\pgfqpoint{4.482543in}{0.815008in}}%
\pgfpathlineto{\pgfqpoint{4.482698in}{0.806759in}}%
\pgfpathlineto{\pgfqpoint{4.483032in}{0.839756in}}%
\pgfpathlineto{\pgfqpoint{4.483046in}{0.848005in}}%
\pgfpathlineto{\pgfqpoint{4.483387in}{0.823258in}}%
\pgfpathlineto{\pgfqpoint{4.484053in}{0.823258in}}%
\pgfpathlineto{\pgfqpoint{4.485897in}{0.897500in}}%
\pgfpathlineto{\pgfqpoint{4.486067in}{0.889251in}}%
\pgfpathlineto{\pgfqpoint{4.486215in}{0.881002in}}%
\pgfpathlineto{\pgfqpoint{4.486386in}{0.905749in}}%
\pgfpathlineto{\pgfqpoint{4.486734in}{0.897500in}}%
\pgfpathlineto{\pgfqpoint{4.487911in}{0.922247in}}%
\pgfpathlineto{\pgfqpoint{4.488081in}{0.913998in}}%
\pgfpathlineto{\pgfqpoint{4.489251in}{0.897500in}}%
\pgfpathlineto{\pgfqpoint{4.490073in}{0.930497in}}%
\pgfpathlineto{\pgfqpoint{4.490406in}{0.905749in}}%
\pgfpathlineto{\pgfqpoint{4.490762in}{0.897500in}}%
\pgfpathlineto{\pgfqpoint{4.490924in}{0.922247in}}%
\pgfpathlineto{\pgfqpoint{4.491428in}{0.913998in}}%
\pgfpathlineto{\pgfqpoint{4.491917in}{0.905749in}}%
\pgfpathlineto{\pgfqpoint{4.491931in}{0.946995in}}%
\pgfpathlineto{\pgfqpoint{4.492583in}{0.930497in}}%
\pgfpathlineto{\pgfqpoint{4.492761in}{0.955244in}}%
\pgfpathlineto{\pgfqpoint{4.492768in}{0.946995in}}%
\pgfpathlineto{\pgfqpoint{4.497966in}{1.194470in}}%
\pgfpathlineto{\pgfqpoint{4.498136in}{1.186221in}}%
\pgfpathlineto{\pgfqpoint{4.499143in}{1.169722in}}%
\pgfpathlineto{\pgfqpoint{4.498306in}{1.194470in}}%
\pgfpathlineto{\pgfqpoint{4.499299in}{1.177971in}}%
\pgfpathlineto{\pgfqpoint{4.500313in}{1.219217in}}%
\pgfpathlineto{\pgfqpoint{4.501490in}{1.111978in}}%
\pgfpathlineto{\pgfqpoint{4.505326in}{0.633527in}}%
\pgfpathlineto{\pgfqpoint{4.505511in}{0.666524in}}%
\pgfpathlineto{\pgfqpoint{4.505829in}{0.683022in}}%
\pgfpathlineto{\pgfqpoint{4.507021in}{0.625278in}}%
\pgfpathlineto{\pgfqpoint{4.508176in}{0.633527in}}%
\pgfpathlineto{\pgfqpoint{4.508339in}{0.625278in}}%
\pgfpathlineto{\pgfqpoint{4.509183in}{0.625278in}}%
\pgfpathlineto{\pgfqpoint{4.510353in}{0.633527in}}%
\pgfpathlineto{\pgfqpoint{4.510523in}{0.625278in}}%
\pgfpathlineto{\pgfqpoint{4.511360in}{0.625278in}}%
\pgfpathlineto{\pgfqpoint{4.512530in}{0.633527in}}%
\pgfpathlineto{\pgfqpoint{4.512715in}{0.625278in}}%
\pgfpathlineto{\pgfqpoint{4.513537in}{0.625278in}}%
\pgfpathlineto{\pgfqpoint{4.514877in}{0.633527in}}%
\pgfpathlineto{\pgfqpoint{4.514899in}{0.625278in}}%
\pgfpathlineto{\pgfqpoint{4.515884in}{0.625278in}}%
\pgfpathlineto{\pgfqpoint{4.517061in}{0.633527in}}%
\pgfpathlineto{\pgfqpoint{4.517224in}{0.625278in}}%
\pgfpathlineto{\pgfqpoint{4.518083in}{0.625278in}}%
\pgfpathlineto{\pgfqpoint{4.520911in}{0.732517in}}%
\pgfpathlineto{\pgfqpoint{4.520934in}{0.724268in}}%
\pgfpathlineto{\pgfqpoint{4.521585in}{0.707769in}}%
\pgfpathlineto{\pgfqpoint{4.521585in}{0.740766in}}%
\pgfpathlineto{\pgfqpoint{4.522762in}{0.806759in}}%
\pgfpathlineto{\pgfqpoint{4.522777in}{0.798510in}}%
\pgfpathlineto{\pgfqpoint{4.522948in}{0.823258in}}%
\pgfpathlineto{\pgfqpoint{4.523614in}{0.790261in}}%
\pgfpathlineto{\pgfqpoint{4.524117in}{0.773763in}}%
\pgfpathlineto{\pgfqpoint{4.524288in}{0.798510in}}%
\pgfpathlineto{\pgfqpoint{4.525954in}{0.938746in}}%
\pgfpathlineto{\pgfqpoint{4.526635in}{0.971742in}}%
\pgfpathlineto{\pgfqpoint{4.526287in}{0.930497in}}%
\pgfpathlineto{\pgfqpoint{4.526968in}{0.946995in}}%
\pgfpathlineto{\pgfqpoint{4.527627in}{0.930497in}}%
\pgfpathlineto{\pgfqpoint{4.527797in}{0.955244in}}%
\pgfpathlineto{\pgfqpoint{4.527805in}{0.946995in}}%
\pgfpathlineto{\pgfqpoint{4.528301in}{0.955244in}}%
\pgfpathlineto{\pgfqpoint{4.528479in}{0.938746in}}%
\pgfpathlineto{\pgfqpoint{4.529315in}{0.946995in}}%
\pgfpathlineto{\pgfqpoint{4.529648in}{0.922247in}}%
\pgfpathlineto{\pgfqpoint{4.530648in}{0.930497in}}%
\pgfpathlineto{\pgfqpoint{4.530655in}{0.922247in}}%
\pgfpathlineto{\pgfqpoint{4.534335in}{0.782012in}}%
\pgfpathlineto{\pgfqpoint{4.531137in}{0.930497in}}%
\pgfpathlineto{\pgfqpoint{4.534335in}{0.790261in}}%
\pgfpathlineto{\pgfqpoint{4.534343in}{0.798510in}}%
\pgfpathlineto{\pgfqpoint{4.535179in}{0.699520in}}%
\pgfpathlineto{\pgfqpoint{4.535350in}{0.699520in}}%
\pgfpathlineto{\pgfqpoint{4.537008in}{0.732517in}}%
\pgfpathlineto{\pgfqpoint{4.537364in}{0.724268in}}%
\pgfpathlineto{\pgfqpoint{4.537526in}{0.749015in}}%
\pgfpathlineto{\pgfqpoint{4.538030in}{0.740766in}}%
\pgfpathlineto{\pgfqpoint{4.538178in}{0.732517in}}%
\pgfpathlineto{\pgfqpoint{4.538533in}{0.773763in}}%
\pgfpathlineto{\pgfqpoint{4.539688in}{0.806759in}}%
\pgfpathlineto{\pgfqpoint{4.539874in}{0.790261in}}%
\pgfpathlineto{\pgfqpoint{4.540214in}{0.773763in}}%
\pgfpathlineto{\pgfqpoint{4.540533in}{0.806759in}}%
\pgfpathlineto{\pgfqpoint{4.540547in}{0.798510in}}%
\pgfpathlineto{\pgfqpoint{4.541540in}{0.831507in}}%
\pgfpathlineto{\pgfqpoint{4.541717in}{0.815008in}}%
\pgfpathlineto{\pgfqpoint{4.542058in}{0.798510in}}%
\pgfpathlineto{\pgfqpoint{4.542221in}{0.823258in}}%
\pgfpathlineto{\pgfqpoint{4.542561in}{0.823258in}}%
\pgfpathlineto{\pgfqpoint{4.542717in}{0.831507in}}%
\pgfpathlineto{\pgfqpoint{4.543057in}{0.798510in}}%
\pgfpathlineto{\pgfqpoint{4.543213in}{0.806759in}}%
\pgfpathlineto{\pgfqpoint{4.544405in}{0.749015in}}%
\pgfpathlineto{\pgfqpoint{4.544560in}{0.757264in}}%
\pgfpathlineto{\pgfqpoint{4.544908in}{0.724268in}}%
\pgfpathlineto{\pgfqpoint{4.545405in}{0.699520in}}%
\pgfpathlineto{\pgfqpoint{4.545730in}{0.732517in}}%
\pgfpathlineto{\pgfqpoint{4.545745in}{0.724268in}}%
\pgfpathlineto{\pgfqpoint{4.546063in}{0.732517in}}%
\pgfpathlineto{\pgfqpoint{4.546249in}{0.716019in}}%
\pgfpathlineto{\pgfqpoint{4.547589in}{0.650025in}}%
\pgfpathlineto{\pgfqpoint{4.547737in}{0.658274in}}%
\pgfpathlineto{\pgfqpoint{4.547907in}{0.683022in}}%
\pgfpathlineto{\pgfqpoint{4.548759in}{0.666524in}}%
\pgfpathlineto{\pgfqpoint{4.550247in}{0.633527in}}%
\pgfpathlineto{\pgfqpoint{4.548929in}{0.674773in}}%
\pgfpathlineto{\pgfqpoint{4.550247in}{0.641776in}}%
\pgfpathlineto{\pgfqpoint{4.551432in}{0.683022in}}%
\pgfpathlineto{\pgfqpoint{4.550587in}{0.633527in}}%
\pgfpathlineto{\pgfqpoint{4.551439in}{0.674773in}}%
\pgfpathlineto{\pgfqpoint{4.551594in}{0.683022in}}%
\pgfpathlineto{\pgfqpoint{4.551780in}{0.666524in}}%
\pgfpathlineto{\pgfqpoint{4.553097in}{0.625278in}}%
\pgfpathlineto{\pgfqpoint{4.554275in}{0.633527in}}%
\pgfpathlineto{\pgfqpoint{4.554949in}{0.625278in}}%
\pgfpathlineto{\pgfqpoint{4.555111in}{0.625278in}}%
\pgfpathlineto{\pgfqpoint{4.556281in}{0.633527in}}%
\pgfpathlineto{\pgfqpoint{4.556452in}{0.625278in}}%
\pgfpathlineto{\pgfqpoint{4.557125in}{0.625278in}}%
\pgfpathlineto{\pgfqpoint{4.558295in}{0.633527in}}%
\pgfpathlineto{\pgfqpoint{4.558466in}{0.625278in}}%
\pgfpathlineto{\pgfqpoint{4.559302in}{0.625278in}}%
\pgfpathlineto{\pgfqpoint{4.560472in}{0.633527in}}%
\pgfpathlineto{\pgfqpoint{4.560642in}{0.625278in}}%
\pgfpathlineto{\pgfqpoint{4.561479in}{0.625278in}}%
\pgfpathlineto{\pgfqpoint{4.562486in}{0.633527in}}%
\pgfpathlineto{\pgfqpoint{4.562486in}{0.625278in}}%
\pgfpathlineto{\pgfqpoint{4.563663in}{0.633527in}}%
\pgfpathlineto{\pgfqpoint{4.563826in}{0.625278in}}%
\pgfpathlineto{\pgfqpoint{4.564663in}{0.625278in}}%
\pgfpathlineto{\pgfqpoint{4.565840in}{0.633527in}}%
\pgfpathlineto{\pgfqpoint{4.566010in}{0.625278in}}%
\pgfpathlineto{\pgfqpoint{4.566847in}{0.625278in}}%
\pgfpathlineto{\pgfqpoint{4.568017in}{0.633527in}}%
\pgfpathlineto{\pgfqpoint{4.568187in}{0.625278in}}%
\pgfpathlineto{\pgfqpoint{4.569024in}{0.625278in}}%
\pgfpathlineto{\pgfqpoint{4.570201in}{0.633527in}}%
\pgfpathlineto{\pgfqpoint{4.570364in}{0.625278in}}%
\pgfpathlineto{\pgfqpoint{4.571201in}{0.625278in}}%
\pgfpathlineto{\pgfqpoint{4.572378in}{0.633527in}}%
\pgfpathlineto{\pgfqpoint{4.572378in}{0.625278in}}%
\pgfpathlineto{\pgfqpoint{4.573400in}{0.625278in}}%
\pgfpathlineto{\pgfqpoint{4.574555in}{0.633527in}}%
\pgfpathlineto{\pgfqpoint{4.574725in}{0.625278in}}%
\pgfpathlineto{\pgfqpoint{4.575562in}{0.625278in}}%
\pgfpathlineto{\pgfqpoint{4.576902in}{0.633527in}}%
\pgfpathlineto{\pgfqpoint{4.576902in}{0.625278in}}%
\pgfpathlineto{\pgfqpoint{4.577909in}{0.625278in}}%
\pgfpathlineto{\pgfqpoint{4.579086in}{0.633527in}}%
\pgfpathlineto{\pgfqpoint{4.579249in}{0.625278in}}%
\pgfpathlineto{\pgfqpoint{4.580086in}{0.625278in}}%
\pgfpathlineto{\pgfqpoint{4.580930in}{0.633527in}}%
\pgfpathlineto{\pgfqpoint{4.580930in}{0.625278in}}%
\pgfpathlineto{\pgfqpoint{4.582100in}{0.633527in}}%
\pgfpathlineto{\pgfqpoint{4.582270in}{0.625278in}}%
\pgfpathlineto{\pgfqpoint{4.582936in}{0.625278in}}%
\pgfpathlineto{\pgfqpoint{4.584114in}{0.633527in}}%
\pgfpathlineto{\pgfqpoint{4.585121in}{0.625278in}}%
\pgfpathlineto{\pgfqpoint{4.586290in}{0.633527in}}%
\pgfpathlineto{\pgfqpoint{4.587312in}{0.625278in}}%
\pgfpathlineto{\pgfqpoint{4.587801in}{0.658274in}}%
\pgfpathlineto{\pgfqpoint{4.588660in}{0.641776in}}%
\pgfpathlineto{\pgfqpoint{4.588808in}{0.633527in}}%
\pgfpathlineto{\pgfqpoint{4.589156in}{0.674773in}}%
\pgfpathlineto{\pgfqpoint{4.591659in}{0.757264in}}%
\pgfpathlineto{\pgfqpoint{4.592325in}{0.707769in}}%
\pgfpathlineto{\pgfqpoint{4.592680in}{0.724268in}}%
\pgfpathlineto{\pgfqpoint{4.593347in}{0.773763in}}%
\pgfpathlineto{\pgfqpoint{4.594020in}{0.740766in}}%
\pgfpathlineto{\pgfqpoint{4.594502in}{0.732517in}}%
\pgfpathlineto{\pgfqpoint{4.594672in}{0.757264in}}%
\pgfpathlineto{\pgfqpoint{4.594694in}{0.749015in}}%
\pgfpathlineto{\pgfqpoint{4.596686in}{0.806759in}}%
\pgfpathlineto{\pgfqpoint{4.597708in}{0.773763in}}%
\pgfpathlineto{\pgfqpoint{4.599033in}{0.806759in}}%
\pgfpathlineto{\pgfqpoint{4.599048in}{0.798510in}}%
\pgfpathlineto{\pgfqpoint{4.599699in}{0.782012in}}%
\pgfpathlineto{\pgfqpoint{4.599870in}{0.806759in}}%
\pgfpathlineto{\pgfqpoint{4.599885in}{0.798510in}}%
\pgfpathlineto{\pgfqpoint{4.600706in}{0.831507in}}%
\pgfpathlineto{\pgfqpoint{4.601040in}{0.806759in}}%
\pgfpathlineto{\pgfqpoint{4.602232in}{0.765513in}}%
\pgfpathlineto{\pgfqpoint{4.602720in}{0.757264in}}%
\pgfpathlineto{\pgfqpoint{4.602735in}{0.798510in}}%
\pgfpathlineto{\pgfqpoint{4.603054in}{0.782012in}}%
\pgfpathlineto{\pgfqpoint{4.603224in}{0.806759in}}%
\pgfpathlineto{\pgfqpoint{4.603579in}{0.798510in}}%
\pgfpathlineto{\pgfqpoint{4.604416in}{0.823258in}}%
\pgfpathlineto{\pgfqpoint{4.604727in}{0.806759in}}%
\pgfpathlineto{\pgfqpoint{4.605756in}{0.798510in}}%
\pgfpathlineto{\pgfqpoint{4.606741in}{0.806759in}}%
\pgfpathlineto{\pgfqpoint{4.606763in}{0.798510in}}%
\pgfpathlineto{\pgfqpoint{4.608422in}{0.757264in}}%
\pgfpathlineto{\pgfqpoint{4.608422in}{0.765513in}}%
\pgfpathlineto{\pgfqpoint{4.609110in}{0.773763in}}%
\pgfpathlineto{\pgfqpoint{4.609443in}{0.749015in}}%
\pgfpathlineto{\pgfqpoint{4.610265in}{0.782012in}}%
\pgfpathlineto{\pgfqpoint{4.610598in}{0.757264in}}%
\pgfpathlineto{\pgfqpoint{4.610954in}{0.749015in}}%
\pgfpathlineto{\pgfqpoint{4.611154in}{0.782012in}}%
\pgfpathlineto{\pgfqpoint{4.611620in}{0.765513in}}%
\pgfpathlineto{\pgfqpoint{4.612946in}{0.732517in}}%
\pgfpathlineto{\pgfqpoint{4.611791in}{0.773763in}}%
\pgfpathlineto{\pgfqpoint{4.612960in}{0.749015in}}%
\pgfpathlineto{\pgfqpoint{4.614138in}{0.773763in}}%
\pgfpathlineto{\pgfqpoint{4.614301in}{0.765513in}}%
\pgfpathlineto{\pgfqpoint{4.614789in}{0.757264in}}%
\pgfpathlineto{\pgfqpoint{4.614952in}{0.782012in}}%
\pgfpathlineto{\pgfqpoint{4.614974in}{0.773763in}}%
\pgfpathlineto{\pgfqpoint{4.615626in}{0.831507in}}%
\pgfpathlineto{\pgfqpoint{4.616818in}{0.815008in}}%
\pgfpathlineto{\pgfqpoint{4.616966in}{0.806759in}}%
\pgfpathlineto{\pgfqpoint{4.616988in}{0.848005in}}%
\pgfpathlineto{\pgfqpoint{4.617321in}{0.848005in}}%
\pgfpathlineto{\pgfqpoint{4.618980in}{0.806759in}}%
\pgfpathlineto{\pgfqpoint{4.618980in}{0.815008in}}%
\pgfpathlineto{\pgfqpoint{4.620172in}{0.897500in}}%
\pgfpathlineto{\pgfqpoint{4.620676in}{0.922247in}}%
\pgfpathlineto{\pgfqpoint{4.621327in}{0.905749in}}%
\pgfpathlineto{\pgfqpoint{4.622349in}{0.872753in}}%
\pgfpathlineto{\pgfqpoint{4.622497in}{0.881002in}}%
\pgfpathlineto{\pgfqpoint{4.622838in}{0.856254in}}%
\pgfpathlineto{\pgfqpoint{4.623023in}{0.864503in}}%
\pgfpathlineto{\pgfqpoint{4.623171in}{0.856254in}}%
\pgfpathlineto{\pgfqpoint{4.623674in}{0.889251in}}%
\pgfpathlineto{\pgfqpoint{4.624681in}{0.905749in}}%
\pgfpathlineto{\pgfqpoint{4.624696in}{0.897500in}}%
\pgfpathlineto{\pgfqpoint{4.624844in}{0.905749in}}%
\pgfpathlineto{\pgfqpoint{4.625185in}{0.881002in}}%
\pgfpathlineto{\pgfqpoint{4.625370in}{0.889251in}}%
\pgfpathlineto{\pgfqpoint{4.626540in}{0.872753in}}%
\pgfpathlineto{\pgfqpoint{4.627524in}{0.881002in}}%
\pgfpathlineto{\pgfqpoint{4.627547in}{0.872753in}}%
\pgfpathlineto{\pgfqpoint{4.628872in}{0.856254in}}%
\pgfpathlineto{\pgfqpoint{4.628887in}{0.897500in}}%
\pgfpathlineto{\pgfqpoint{4.629894in}{0.872753in}}%
\pgfpathlineto{\pgfqpoint{4.631049in}{0.881002in}}%
\pgfpathlineto{\pgfqpoint{4.631886in}{0.856254in}}%
\pgfpathlineto{\pgfqpoint{4.631234in}{0.897500in}}%
\pgfpathlineto{\pgfqpoint{4.632071in}{0.864503in}}%
\pgfpathlineto{\pgfqpoint{4.632915in}{0.848005in}}%
\pgfpathlineto{\pgfqpoint{4.632241in}{0.872753in}}%
\pgfpathlineto{\pgfqpoint{4.633063in}{0.864503in}}%
\pgfpathlineto{\pgfqpoint{4.633581in}{0.922247in}}%
\pgfpathlineto{\pgfqpoint{4.634240in}{0.905749in}}%
\pgfpathlineto{\pgfqpoint{4.634921in}{0.897500in}}%
\pgfpathlineto{\pgfqpoint{4.635254in}{0.946995in}}%
\pgfpathlineto{\pgfqpoint{4.635595in}{0.971742in}}%
\pgfpathlineto{\pgfqpoint{4.636098in}{0.946995in}}%
\pgfpathlineto{\pgfqpoint{4.637416in}{0.930497in}}%
\pgfpathlineto{\pgfqpoint{4.637587in}{0.955244in}}%
\pgfpathlineto{\pgfqpoint{4.638446in}{0.922247in}}%
\pgfpathlineto{\pgfqpoint{4.639764in}{0.905749in}}%
\pgfpathlineto{\pgfqpoint{4.640622in}{0.971742in}}%
\pgfpathlineto{\pgfqpoint{4.640793in}{0.963493in}}%
\pgfpathlineto{\pgfqpoint{4.641629in}{0.922247in}}%
\pgfpathlineto{\pgfqpoint{4.641778in}{0.963493in}}%
\pgfpathlineto{\pgfqpoint{4.641792in}{0.971742in}}%
\pgfpathlineto{\pgfqpoint{4.642466in}{0.922247in}}%
\pgfpathlineto{\pgfqpoint{4.642799in}{0.922247in}}%
\pgfpathlineto{\pgfqpoint{4.643977in}{0.996490in}}%
\pgfpathlineto{\pgfqpoint{4.644458in}{0.955244in}}%
\pgfpathlineto{\pgfqpoint{4.644643in}{0.946995in}}%
\pgfpathlineto{\pgfqpoint{4.644961in}{0.979992in}}%
\pgfpathlineto{\pgfqpoint{4.645487in}{0.971742in}}%
\pgfpathlineto{\pgfqpoint{4.646487in}{0.996490in}}%
\pgfpathlineto{\pgfqpoint{4.646657in}{0.988241in}}%
\pgfpathlineto{\pgfqpoint{4.647827in}{0.971742in}}%
\pgfpathlineto{\pgfqpoint{4.648315in}{0.979992in}}%
\pgfpathlineto{\pgfqpoint{4.648501in}{0.963493in}}%
\pgfpathlineto{\pgfqpoint{4.648671in}{0.971742in}}%
\pgfpathlineto{\pgfqpoint{4.649656in}{0.955244in}}%
\pgfpathlineto{\pgfqpoint{4.650677in}{1.021237in}}%
\pgfpathlineto{\pgfqpoint{4.652003in}{1.054234in}}%
\pgfpathlineto{\pgfqpoint{4.652018in}{1.045985in}}%
\pgfpathlineto{\pgfqpoint{4.653676in}{1.004739in}}%
\pgfpathlineto{\pgfqpoint{4.653676in}{1.012988in}}%
\pgfpathlineto{\pgfqpoint{4.654017in}{1.004739in}}%
\pgfpathlineto{\pgfqpoint{4.654705in}{1.045985in}}%
\pgfpathlineto{\pgfqpoint{4.656045in}{0.996490in}}%
\pgfpathlineto{\pgfqpoint{4.656527in}{1.004739in}}%
\pgfpathlineto{\pgfqpoint{4.657363in}{1.054234in}}%
\pgfpathlineto{\pgfqpoint{4.657556in}{1.037736in}}%
\pgfpathlineto{\pgfqpoint{4.657719in}{1.045985in}}%
\pgfpathlineto{\pgfqpoint{4.658711in}{1.029487in}}%
\pgfpathlineto{\pgfqpoint{4.659044in}{1.078981in}}%
\pgfpathlineto{\pgfqpoint{4.659733in}{1.045985in}}%
\pgfpathlineto{\pgfqpoint{4.660888in}{1.054234in}}%
\pgfpathlineto{\pgfqpoint{4.661724in}{1.029487in}}%
\pgfpathlineto{\pgfqpoint{4.661910in}{1.045985in}}%
\pgfpathlineto{\pgfqpoint{4.663901in}{1.103729in}}%
\pgfpathlineto{\pgfqpoint{4.665093in}{1.070732in}}%
\pgfpathlineto{\pgfqpoint{4.665767in}{1.045985in}}%
\pgfpathlineto{\pgfqpoint{4.665937in}{1.070732in}}%
\pgfpathlineto{\pgfqpoint{4.666271in}{1.095480in}}%
\pgfpathlineto{\pgfqpoint{4.667092in}{1.078981in}}%
\pgfpathlineto{\pgfqpoint{4.668114in}{1.062483in}}%
\pgfpathlineto{\pgfqpoint{4.669284in}{1.045985in}}%
\pgfpathlineto{\pgfqpoint{4.670106in}{1.054234in}}%
\pgfpathlineto{\pgfqpoint{4.670128in}{1.045985in}}%
\pgfpathlineto{\pgfqpoint{4.671446in}{1.029487in}}%
\pgfpathlineto{\pgfqpoint{4.671802in}{1.021237in}}%
\pgfpathlineto{\pgfqpoint{4.672475in}{1.070732in}}%
\pgfpathlineto{\pgfqpoint{4.673127in}{1.078981in}}%
\pgfpathlineto{\pgfqpoint{4.673312in}{1.062483in}}%
\pgfpathlineto{\pgfqpoint{4.674800in}{1.029487in}}%
\pgfpathlineto{\pgfqpoint{4.673793in}{1.078981in}}%
\pgfpathlineto{\pgfqpoint{4.674800in}{1.037736in}}%
\pgfpathlineto{\pgfqpoint{4.675133in}{1.029487in}}%
\pgfpathlineto{\pgfqpoint{4.675822in}{1.045985in}}%
\pgfpathlineto{\pgfqpoint{4.676140in}{1.029487in}}%
\pgfpathlineto{\pgfqpoint{4.676163in}{1.070732in}}%
\pgfpathlineto{\pgfqpoint{4.676666in}{1.095480in}}%
\pgfpathlineto{\pgfqpoint{4.677162in}{1.070732in}}%
\pgfpathlineto{\pgfqpoint{4.677814in}{1.054234in}}%
\pgfpathlineto{\pgfqpoint{4.677984in}{1.078981in}}%
\pgfpathlineto{\pgfqpoint{4.678006in}{1.070732in}}%
\pgfpathlineto{\pgfqpoint{4.678843in}{1.120227in}}%
\pgfpathlineto{\pgfqpoint{4.679176in}{1.087231in}}%
\pgfpathlineto{\pgfqpoint{4.682738in}{0.633527in}}%
\pgfpathlineto{\pgfqpoint{4.683204in}{0.650025in}}%
\pgfpathlineto{\pgfqpoint{4.684204in}{0.625278in}}%
\pgfpathlineto{\pgfqpoint{4.684352in}{0.641776in}}%
\pgfpathlineto{\pgfqpoint{4.684877in}{0.674773in}}%
\pgfpathlineto{\pgfqpoint{4.685381in}{0.641776in}}%
\pgfpathlineto{\pgfqpoint{4.685714in}{0.625278in}}%
\pgfpathlineto{\pgfqpoint{4.686032in}{0.658274in}}%
\pgfpathlineto{\pgfqpoint{4.686092in}{0.658274in}}%
\pgfpathlineto{\pgfqpoint{4.687202in}{0.707769in}}%
\pgfpathlineto{\pgfqpoint{4.687395in}{0.691271in}}%
\pgfpathlineto{\pgfqpoint{4.687543in}{0.683022in}}%
\pgfpathlineto{\pgfqpoint{4.688046in}{0.716019in}}%
\pgfpathlineto{\pgfqpoint{4.688061in}{0.724268in}}%
\pgfpathlineto{\pgfqpoint{4.688735in}{0.674773in}}%
\pgfpathlineto{\pgfqpoint{4.689068in}{0.674773in}}%
\pgfpathlineto{\pgfqpoint{4.690238in}{0.699520in}}%
\pgfpathlineto{\pgfqpoint{4.689720in}{0.658274in}}%
\pgfpathlineto{\pgfqpoint{4.690408in}{0.691271in}}%
\pgfpathlineto{\pgfqpoint{4.691586in}{0.674773in}}%
\pgfpathlineto{\pgfqpoint{4.692755in}{0.699520in}}%
\pgfpathlineto{\pgfqpoint{4.692926in}{0.691271in}}%
\pgfpathlineto{\pgfqpoint{4.693407in}{0.683022in}}%
\pgfpathlineto{\pgfqpoint{4.693585in}{0.707769in}}%
\pgfpathlineto{\pgfqpoint{4.693762in}{0.699520in}}%
\pgfpathlineto{\pgfqpoint{4.694577in}{0.732517in}}%
\pgfpathlineto{\pgfqpoint{4.694769in}{0.716019in}}%
\pgfpathlineto{\pgfqpoint{4.695436in}{0.674773in}}%
\pgfpathlineto{\pgfqpoint{4.695939in}{0.699520in}}%
\pgfpathlineto{\pgfqpoint{4.697598in}{0.732517in}}%
\pgfpathlineto{\pgfqpoint{4.698116in}{0.724268in}}%
\pgfpathlineto{\pgfqpoint{4.698457in}{0.773763in}}%
\pgfpathlineto{\pgfqpoint{4.699626in}{0.798510in}}%
\pgfpathlineto{\pgfqpoint{4.699797in}{0.790261in}}%
\pgfpathlineto{\pgfqpoint{4.700463in}{0.773763in}}%
\pgfpathlineto{\pgfqpoint{4.699967in}{0.798510in}}%
\pgfpathlineto{\pgfqpoint{4.700952in}{0.782012in}}%
\pgfpathlineto{\pgfqpoint{4.701959in}{0.831507in}}%
\pgfpathlineto{\pgfqpoint{4.701974in}{0.823258in}}%
\pgfpathlineto{\pgfqpoint{4.702470in}{0.856254in}}%
\pgfpathlineto{\pgfqpoint{4.703314in}{0.839756in}}%
\pgfpathlineto{\pgfqpoint{4.703802in}{0.831507in}}%
\pgfpathlineto{\pgfqpoint{4.703965in}{0.856254in}}%
\pgfpathlineto{\pgfqpoint{4.703988in}{0.848005in}}%
\pgfpathlineto{\pgfqpoint{4.706483in}{0.979992in}}%
\pgfpathlineto{\pgfqpoint{4.707149in}{0.955244in}}%
\pgfpathlineto{\pgfqpoint{4.707342in}{0.946995in}}%
\pgfpathlineto{\pgfqpoint{4.708008in}{0.996490in}}%
\pgfpathlineto{\pgfqpoint{4.709163in}{1.004739in}}%
\pgfpathlineto{\pgfqpoint{4.709348in}{0.996490in}}%
\pgfpathlineto{\pgfqpoint{4.709852in}{1.045985in}}%
\pgfpathlineto{\pgfqpoint{4.710185in}{1.037736in}}%
\pgfpathlineto{\pgfqpoint{4.710503in}{1.029487in}}%
\pgfpathlineto{\pgfqpoint{4.711192in}{1.045985in}}%
\pgfpathlineto{\pgfqpoint{4.711510in}{1.029487in}}%
\pgfpathlineto{\pgfqpoint{4.711681in}{1.054234in}}%
\pgfpathlineto{\pgfqpoint{4.712014in}{1.054234in}}%
\pgfpathlineto{\pgfqpoint{4.712702in}{1.070732in}}%
\pgfpathlineto{\pgfqpoint{4.713035in}{1.045985in}}%
\pgfpathlineto{\pgfqpoint{4.715383in}{0.971742in}}%
\pgfpathlineto{\pgfqpoint{4.715553in}{0.996490in}}%
\pgfpathlineto{\pgfqpoint{4.716227in}{0.963493in}}%
\pgfpathlineto{\pgfqpoint{4.717397in}{0.946995in}}%
\pgfpathlineto{\pgfqpoint{4.717885in}{0.955244in}}%
\pgfpathlineto{\pgfqpoint{4.718070in}{0.938746in}}%
\pgfpathlineto{\pgfqpoint{4.718381in}{0.905749in}}%
\pgfpathlineto{\pgfqpoint{4.719240in}{0.922247in}}%
\pgfpathlineto{\pgfqpoint{4.719566in}{0.930497in}}%
\pgfpathlineto{\pgfqpoint{4.719899in}{0.905749in}}%
\pgfpathlineto{\pgfqpoint{4.720921in}{0.897500in}}%
\pgfpathlineto{\pgfqpoint{4.721076in}{0.905749in}}%
\pgfpathlineto{\pgfqpoint{4.721417in}{0.872753in}}%
\pgfpathlineto{\pgfqpoint{4.721573in}{0.881002in}}%
\pgfpathlineto{\pgfqpoint{4.722765in}{0.798510in}}%
\pgfpathlineto{\pgfqpoint{4.723098in}{0.823258in}}%
\pgfpathlineto{\pgfqpoint{4.723912in}{0.806759in}}%
\pgfpathlineto{\pgfqpoint{4.725112in}{0.765513in}}%
\pgfpathlineto{\pgfqpoint{4.726785in}{0.724268in}}%
\pgfpathlineto{\pgfqpoint{4.726933in}{0.732517in}}%
\pgfpathlineto{\pgfqpoint{4.727289in}{0.691271in}}%
\pgfpathlineto{\pgfqpoint{4.728125in}{0.625278in}}%
\pgfpathlineto{\pgfqpoint{4.728629in}{0.650025in}}%
\pgfpathlineto{\pgfqpoint{4.729280in}{0.658274in}}%
\pgfpathlineto{\pgfqpoint{4.729110in}{0.633527in}}%
\pgfpathlineto{\pgfqpoint{4.729465in}{0.641776in}}%
\pgfpathlineto{\pgfqpoint{4.729954in}{0.633527in}}%
\pgfpathlineto{\pgfqpoint{4.730117in}{0.658274in}}%
\pgfpathlineto{\pgfqpoint{4.730139in}{0.650025in}}%
\pgfpathlineto{\pgfqpoint{4.731124in}{0.658274in}}%
\pgfpathlineto{\pgfqpoint{4.731146in}{0.650025in}}%
\pgfpathlineto{\pgfqpoint{4.731642in}{0.625278in}}%
\pgfpathlineto{\pgfqpoint{4.731961in}{0.658274in}}%
\pgfpathlineto{\pgfqpoint{4.731983in}{0.650025in}}%
\pgfpathlineto{\pgfqpoint{4.733138in}{0.658274in}}%
\pgfpathlineto{\pgfqpoint{4.733493in}{0.650025in}}%
\pgfpathlineto{\pgfqpoint{4.734145in}{0.707769in}}%
\pgfpathlineto{\pgfqpoint{4.734160in}{0.699520in}}%
\pgfpathlineto{\pgfqpoint{4.735485in}{0.757264in}}%
\pgfpathlineto{\pgfqpoint{4.736003in}{0.740766in}}%
\pgfpathlineto{\pgfqpoint{4.736151in}{0.732517in}}%
\pgfpathlineto{\pgfqpoint{4.736507in}{0.773763in}}%
\pgfpathlineto{\pgfqpoint{4.736655in}{0.782012in}}%
\pgfpathlineto{\pgfqpoint{4.737010in}{0.749015in}}%
\pgfpathlineto{\pgfqpoint{4.737166in}{0.757264in}}%
\pgfpathlineto{\pgfqpoint{4.737669in}{0.782012in}}%
\pgfpathlineto{\pgfqpoint{4.738180in}{0.749015in}}%
\pgfpathlineto{\pgfqpoint{4.739335in}{0.757264in}}%
\pgfpathlineto{\pgfqpoint{4.739528in}{0.749015in}}%
\pgfpathlineto{\pgfqpoint{4.740016in}{0.806759in}}%
\pgfpathlineto{\pgfqpoint{4.740194in}{0.798510in}}%
\pgfpathlineto{\pgfqpoint{4.741682in}{0.831507in}}%
\pgfpathlineto{\pgfqpoint{4.741705in}{0.823258in}}%
\pgfpathlineto{\pgfqpoint{4.743045in}{0.798510in}}%
\pgfpathlineto{\pgfqpoint{4.744215in}{0.823258in}}%
\pgfpathlineto{\pgfqpoint{4.744385in}{0.815008in}}%
\pgfpathlineto{\pgfqpoint{4.744866in}{0.806759in}}%
\pgfpathlineto{\pgfqpoint{4.745044in}{0.831507in}}%
\pgfpathlineto{\pgfqpoint{4.745059in}{0.823258in}}%
\pgfpathlineto{\pgfqpoint{4.747406in}{0.922247in}}%
\pgfpathlineto{\pgfqpoint{4.747569in}{0.913998in}}%
\pgfpathlineto{\pgfqpoint{4.747894in}{0.930497in}}%
\pgfpathlineto{\pgfqpoint{4.750086in}{0.823258in}}%
\pgfpathlineto{\pgfqpoint{4.752078in}{0.881002in}}%
\pgfpathlineto{\pgfqpoint{4.753100in}{0.872753in}}%
\pgfpathlineto{\pgfqpoint{4.753418in}{0.856254in}}%
\pgfpathlineto{\pgfqpoint{4.753751in}{0.889251in}}%
\pgfpathlineto{\pgfqpoint{4.753921in}{0.905749in}}%
\pgfpathlineto{\pgfqpoint{4.754780in}{0.864503in}}%
\pgfpathlineto{\pgfqpoint{4.755114in}{0.848005in}}%
\pgfpathlineto{\pgfqpoint{4.755432in}{0.881002in}}%
\pgfpathlineto{\pgfqpoint{4.755595in}{0.881002in}}%
\pgfpathlineto{\pgfqpoint{4.755617in}{0.897500in}}%
\pgfpathlineto{\pgfqpoint{4.756291in}{0.848005in}}%
\pgfpathlineto{\pgfqpoint{4.756624in}{0.872753in}}%
\pgfpathlineto{\pgfqpoint{4.758623in}{0.930497in}}%
\pgfpathlineto{\pgfqpoint{4.758956in}{0.955244in}}%
\pgfpathlineto{\pgfqpoint{4.759637in}{0.922247in}}%
\pgfpathlineto{\pgfqpoint{4.759786in}{0.930497in}}%
\pgfpathlineto{\pgfqpoint{4.760141in}{0.897500in}}%
\pgfpathlineto{\pgfqpoint{4.760289in}{0.905749in}}%
\pgfpathlineto{\pgfqpoint{4.760482in}{0.897500in}}%
\pgfpathlineto{\pgfqpoint{4.761296in}{0.938746in}}%
\pgfpathlineto{\pgfqpoint{4.762473in}{0.979992in}}%
\pgfpathlineto{\pgfqpoint{4.762488in}{0.971742in}}%
\pgfpathlineto{\pgfqpoint{4.763643in}{0.979992in}}%
\pgfpathlineto{\pgfqpoint{4.764835in}{0.946995in}}%
\pgfpathlineto{\pgfqpoint{4.766494in}{0.979992in}}%
\pgfpathlineto{\pgfqpoint{4.767686in}{0.897500in}}%
\pgfpathlineto{\pgfqpoint{4.767834in}{0.905749in}}%
\pgfpathlineto{\pgfqpoint{4.768167in}{0.881002in}}%
\pgfpathlineto{\pgfqpoint{4.768360in}{0.889251in}}%
\pgfpathlineto{\pgfqpoint{4.768863in}{0.872753in}}%
\pgfpathlineto{\pgfqpoint{4.769174in}{0.905749in}}%
\pgfpathlineto{\pgfqpoint{4.769196in}{0.897500in}}%
\pgfpathlineto{\pgfqpoint{4.770181in}{0.905749in}}%
\pgfpathlineto{\pgfqpoint{4.770203in}{0.897500in}}%
\pgfpathlineto{\pgfqpoint{4.771188in}{0.881002in}}%
\pgfpathlineto{\pgfqpoint{4.771210in}{0.897500in}}%
\pgfpathlineto{\pgfqpoint{4.771692in}{0.905749in}}%
\pgfpathlineto{\pgfqpoint{4.772025in}{0.881002in}}%
\pgfpathlineto{\pgfqpoint{4.773217in}{0.848005in}}%
\pgfpathlineto{\pgfqpoint{4.774727in}{0.897500in}}%
\pgfpathlineto{\pgfqpoint{4.774898in}{0.889251in}}%
\pgfpathlineto{\pgfqpoint{4.776067in}{0.872753in}}%
\pgfpathlineto{\pgfqpoint{4.777245in}{0.897500in}}%
\pgfpathlineto{\pgfqpoint{4.777408in}{0.889251in}}%
\pgfpathlineto{\pgfqpoint{4.777741in}{0.872753in}}%
\pgfpathlineto{\pgfqpoint{4.778059in}{0.905749in}}%
\pgfpathlineto{\pgfqpoint{4.778081in}{0.897500in}}%
\pgfpathlineto{\pgfqpoint{4.779399in}{0.930497in}}%
\pgfpathlineto{\pgfqpoint{4.778585in}{0.872753in}}%
\pgfpathlineto{\pgfqpoint{4.779421in}{0.922247in}}%
\pgfpathlineto{\pgfqpoint{4.780258in}{0.897500in}}%
\pgfpathlineto{\pgfqpoint{4.780739in}{0.905749in}}%
\pgfpathlineto{\pgfqpoint{4.781584in}{0.930497in}}%
\pgfpathlineto{\pgfqpoint{4.781769in}{0.913998in}}%
\pgfpathlineto{\pgfqpoint{4.783427in}{0.881002in}}%
\pgfpathlineto{\pgfqpoint{4.784116in}{0.922247in}}%
\pgfpathlineto{\pgfqpoint{4.784449in}{0.897500in}}%
\pgfpathlineto{\pgfqpoint{4.786278in}{0.979992in}}%
\pgfpathlineto{\pgfqpoint{4.786463in}{0.963493in}}%
\pgfpathlineto{\pgfqpoint{4.786626in}{0.971742in}}%
\pgfpathlineto{\pgfqpoint{4.787618in}{0.955244in}}%
\pgfpathlineto{\pgfqpoint{4.788121in}{1.004739in}}%
\pgfpathlineto{\pgfqpoint{4.788640in}{0.996490in}}%
\pgfpathlineto{\pgfqpoint{4.790469in}{1.054234in}}%
\pgfpathlineto{\pgfqpoint{4.790483in}{1.045985in}}%
\pgfpathlineto{\pgfqpoint{4.791320in}{0.996490in}}%
\pgfpathlineto{\pgfqpoint{4.792149in}{1.012988in}}%
\pgfpathlineto{\pgfqpoint{4.793312in}{1.078981in}}%
\pgfpathlineto{\pgfqpoint{4.793334in}{1.070732in}}%
\pgfpathlineto{\pgfqpoint{4.794008in}{0.988241in}}%
\pgfpathlineto{\pgfqpoint{4.796851in}{0.699520in}}%
\pgfpathlineto{\pgfqpoint{4.797169in}{0.707769in}}%
\pgfpathlineto{\pgfqpoint{4.797503in}{0.683022in}}%
\pgfpathlineto{\pgfqpoint{4.797695in}{0.674773in}}%
\pgfpathlineto{\pgfqpoint{4.798006in}{0.707769in}}%
\pgfpathlineto{\pgfqpoint{4.798532in}{0.699520in}}%
\pgfpathlineto{\pgfqpoint{4.799346in}{0.683022in}}%
\pgfpathlineto{\pgfqpoint{4.799687in}{0.707769in}}%
\pgfpathlineto{\pgfqpoint{4.799872in}{0.724268in}}%
\pgfpathlineto{\pgfqpoint{4.800709in}{0.674773in}}%
\pgfpathlineto{\pgfqpoint{4.801864in}{0.707769in}}%
\pgfpathlineto{\pgfqpoint{4.802049in}{0.691271in}}%
\pgfpathlineto{\pgfqpoint{4.802537in}{0.683022in}}%
\pgfpathlineto{\pgfqpoint{4.802552in}{0.724268in}}%
\pgfpathlineto{\pgfqpoint{4.802871in}{0.707769in}}%
\pgfpathlineto{\pgfqpoint{4.803211in}{0.740766in}}%
\pgfpathlineto{\pgfqpoint{4.803707in}{0.757264in}}%
\pgfpathlineto{\pgfqpoint{4.804233in}{0.716019in}}%
\pgfpathlineto{\pgfqpoint{4.804381in}{0.707769in}}%
\pgfpathlineto{\pgfqpoint{4.804736in}{0.749015in}}%
\pgfpathlineto{\pgfqpoint{4.805047in}{0.757264in}}%
\pgfpathlineto{\pgfqpoint{4.805388in}{0.732517in}}%
\pgfpathlineto{\pgfqpoint{4.805573in}{0.740766in}}%
\pgfpathlineto{\pgfqpoint{4.805721in}{0.732517in}}%
\pgfpathlineto{\pgfqpoint{4.806077in}{0.773763in}}%
\pgfpathlineto{\pgfqpoint{4.806580in}{0.823258in}}%
\pgfpathlineto{\pgfqpoint{4.807735in}{0.806759in}}%
\pgfpathlineto{\pgfqpoint{4.808735in}{0.782012in}}%
\pgfpathlineto{\pgfqpoint{4.808757in}{0.798510in}}%
\pgfpathlineto{\pgfqpoint{4.809742in}{0.782012in}}%
\pgfpathlineto{\pgfqpoint{4.809742in}{0.790261in}}%
\pgfpathlineto{\pgfqpoint{4.810771in}{0.872753in}}%
\pgfpathlineto{\pgfqpoint{4.812111in}{0.848005in}}%
\pgfpathlineto{\pgfqpoint{4.812429in}{0.856254in}}%
\pgfpathlineto{\pgfqpoint{4.812777in}{0.823258in}}%
\pgfpathlineto{\pgfqpoint{4.812925in}{0.831507in}}%
\pgfpathlineto{\pgfqpoint{4.813955in}{0.823258in}}%
\pgfpathlineto{\pgfqpoint{4.815628in}{0.872753in}}%
\pgfpathlineto{\pgfqpoint{4.815798in}{0.864503in}}%
\pgfpathlineto{\pgfqpoint{4.815946in}{0.856254in}}%
\pgfpathlineto{\pgfqpoint{4.815961in}{0.897500in}}%
\pgfpathlineto{\pgfqpoint{4.816465in}{0.889251in}}%
\pgfpathlineto{\pgfqpoint{4.817620in}{0.856254in}}%
\pgfpathlineto{\pgfqpoint{4.816635in}{0.897500in}}%
\pgfpathlineto{\pgfqpoint{4.817642in}{0.872753in}}%
\pgfpathlineto{\pgfqpoint{4.818797in}{0.881002in}}%
\pgfpathlineto{\pgfqpoint{4.819486in}{0.872753in}}%
\pgfpathlineto{\pgfqpoint{4.819648in}{0.897500in}}%
\pgfpathlineto{\pgfqpoint{4.819819in}{0.889251in}}%
\pgfpathlineto{\pgfqpoint{4.819967in}{0.881002in}}%
\pgfpathlineto{\pgfqpoint{4.820145in}{0.905749in}}%
\pgfpathlineto{\pgfqpoint{4.820493in}{0.897500in}}%
\pgfpathlineto{\pgfqpoint{4.821500in}{0.922247in}}%
\pgfpathlineto{\pgfqpoint{4.821662in}{0.913998in}}%
\pgfpathlineto{\pgfqpoint{4.822151in}{0.905749in}}%
\pgfpathlineto{\pgfqpoint{4.822329in}{0.930497in}}%
\pgfpathlineto{\pgfqpoint{4.822336in}{0.922247in}}%
\pgfpathlineto{\pgfqpoint{4.823506in}{0.930497in}}%
\pgfpathlineto{\pgfqpoint{4.824328in}{0.881002in}}%
\pgfpathlineto{\pgfqpoint{4.824513in}{0.889251in}}%
\pgfpathlineto{\pgfqpoint{4.825668in}{0.856254in}}%
\pgfpathlineto{\pgfqpoint{4.824683in}{0.897500in}}%
\pgfpathlineto{\pgfqpoint{4.825690in}{0.872753in}}%
\pgfpathlineto{\pgfqpoint{4.826527in}{0.897500in}}%
\pgfpathlineto{\pgfqpoint{4.826690in}{0.889251in}}%
\pgfpathlineto{\pgfqpoint{4.828704in}{0.749015in}}%
\pgfpathlineto{\pgfqpoint{4.829526in}{0.757264in}}%
\pgfpathlineto{\pgfqpoint{4.829540in}{0.749015in}}%
\pgfpathlineto{\pgfqpoint{4.830044in}{0.724268in}}%
\pgfpathlineto{\pgfqpoint{4.830362in}{0.757264in}}%
\pgfpathlineto{\pgfqpoint{4.830377in}{0.749015in}}%
\pgfpathlineto{\pgfqpoint{4.831554in}{0.773763in}}%
\pgfpathlineto{\pgfqpoint{4.831725in}{0.765513in}}%
\pgfpathlineto{\pgfqpoint{4.832206in}{0.757264in}}%
\pgfpathlineto{\pgfqpoint{4.832221in}{0.798510in}}%
\pgfpathlineto{\pgfqpoint{4.832561in}{0.773763in}}%
\pgfpathlineto{\pgfqpoint{4.833050in}{0.782012in}}%
\pgfpathlineto{\pgfqpoint{4.833228in}{0.765513in}}%
\pgfpathlineto{\pgfqpoint{4.834909in}{0.724268in}}%
\pgfpathlineto{\pgfqpoint{4.836752in}{0.798510in}}%
\pgfpathlineto{\pgfqpoint{4.837071in}{0.782012in}}%
\pgfpathlineto{\pgfqpoint{4.837907in}{0.757264in}}%
\pgfpathlineto{\pgfqpoint{4.838092in}{0.765513in}}%
\pgfpathlineto{\pgfqpoint{4.838263in}{0.773763in}}%
\pgfpathlineto{\pgfqpoint{4.839262in}{0.749015in}}%
\pgfpathlineto{\pgfqpoint{4.839751in}{0.757264in}}%
\pgfpathlineto{\pgfqpoint{4.839936in}{0.740766in}}%
\pgfpathlineto{\pgfqpoint{4.842787in}{0.625278in}}%
\pgfpathlineto{\pgfqpoint{4.844778in}{0.683022in}}%
\pgfpathlineto{\pgfqpoint{4.845112in}{0.658274in}}%
\pgfpathlineto{\pgfqpoint{4.845467in}{0.650025in}}%
\pgfpathlineto{\pgfqpoint{4.845785in}{0.683022in}}%
\pgfpathlineto{\pgfqpoint{4.846141in}{0.674773in}}%
\pgfpathlineto{\pgfqpoint{4.847481in}{0.749015in}}%
\pgfpathlineto{\pgfqpoint{4.847984in}{0.716019in}}%
\pgfpathlineto{\pgfqpoint{4.848799in}{0.707769in}}%
\pgfpathlineto{\pgfqpoint{4.848147in}{0.724268in}}%
\pgfpathlineto{\pgfqpoint{4.848799in}{0.716019in}}%
\pgfpathlineto{\pgfqpoint{4.849828in}{0.773763in}}%
\pgfpathlineto{\pgfqpoint{4.850139in}{0.757264in}}%
\pgfpathlineto{\pgfqpoint{4.850161in}{0.798510in}}%
\pgfpathlineto{\pgfqpoint{4.850487in}{0.790261in}}%
\pgfpathlineto{\pgfqpoint{4.851487in}{0.806759in}}%
\pgfpathlineto{\pgfqpoint{4.850813in}{0.782012in}}%
\pgfpathlineto{\pgfqpoint{4.851501in}{0.798510in}}%
\pgfpathlineto{\pgfqpoint{4.852331in}{0.856254in}}%
\pgfpathlineto{\pgfqpoint{4.853160in}{0.831507in}}%
\pgfpathlineto{\pgfqpoint{4.854182in}{0.798510in}}%
\pgfpathlineto{\pgfqpoint{4.854678in}{0.806759in}}%
\pgfpathlineto{\pgfqpoint{4.856025in}{0.749015in}}%
\pgfpathlineto{\pgfqpoint{4.856344in}{0.782012in}}%
\pgfpathlineto{\pgfqpoint{4.857180in}{0.757264in}}%
\pgfpathlineto{\pgfqpoint{4.857854in}{0.732517in}}%
\pgfpathlineto{\pgfqpoint{4.858210in}{0.749015in}}%
\pgfpathlineto{\pgfqpoint{4.859194in}{0.683022in}}%
\pgfpathlineto{\pgfqpoint{4.859550in}{0.724268in}}%
\pgfpathlineto{\pgfqpoint{4.860216in}{0.749015in}}%
\pgfpathlineto{\pgfqpoint{4.860705in}{0.732517in}}%
\pgfpathlineto{\pgfqpoint{4.861556in}{0.699520in}}%
\pgfpathlineto{\pgfqpoint{4.861727in}{0.724268in}}%
\pgfpathlineto{\pgfqpoint{4.862733in}{0.699520in}}%
\pgfpathlineto{\pgfqpoint{4.862889in}{0.716019in}}%
\pgfpathlineto{\pgfqpoint{4.863563in}{0.757264in}}%
\pgfpathlineto{\pgfqpoint{4.863903in}{0.749015in}}%
\pgfpathlineto{\pgfqpoint{4.866754in}{0.650025in}}%
\pgfpathlineto{\pgfqpoint{4.867250in}{0.683022in}}%
\pgfpathlineto{\pgfqpoint{4.867761in}{0.650025in}}%
\pgfpathlineto{\pgfqpoint{4.868435in}{0.625278in}}%
\pgfpathlineto{\pgfqpoint{4.868598in}{0.650025in}}%
\pgfpathlineto{\pgfqpoint{4.870278in}{0.724268in}}%
\pgfpathlineto{\pgfqpoint{4.870589in}{0.707769in}}%
\pgfpathlineto{\pgfqpoint{4.871619in}{0.699520in}}%
\pgfpathlineto{\pgfqpoint{4.873610in}{0.757264in}}%
\pgfpathlineto{\pgfqpoint{4.874802in}{0.716019in}}%
\pgfpathlineto{\pgfqpoint{4.875284in}{0.707769in}}%
\pgfpathlineto{\pgfqpoint{4.875306in}{0.749015in}}%
\pgfpathlineto{\pgfqpoint{4.875957in}{0.732517in}}%
\pgfpathlineto{\pgfqpoint{4.876128in}{0.757264in}}%
\pgfpathlineto{\pgfqpoint{4.876142in}{0.749015in}}%
\pgfpathlineto{\pgfqpoint{4.876461in}{0.757264in}}%
\pgfpathlineto{\pgfqpoint{4.876794in}{0.732517in}}%
\pgfpathlineto{\pgfqpoint{4.877149in}{0.724268in}}%
\pgfpathlineto{\pgfqpoint{4.877468in}{0.757264in}}%
\pgfpathlineto{\pgfqpoint{4.877823in}{0.749015in}}%
\pgfpathlineto{\pgfqpoint{4.879326in}{0.798510in}}%
\pgfpathlineto{\pgfqpoint{4.879497in}{0.790261in}}%
\pgfpathlineto{\pgfqpoint{4.883332in}{0.625278in}}%
\pgfpathlineto{\pgfqpoint{4.884339in}{0.633527in}}%
\pgfpathlineto{\pgfqpoint{4.884361in}{0.625278in}}%
\pgfpathlineto{\pgfqpoint{4.885509in}{0.633527in}}%
\pgfpathlineto{\pgfqpoint{4.885679in}{0.625278in}}%
\pgfpathlineto{\pgfqpoint{4.886538in}{0.625278in}}%
\pgfpathlineto{\pgfqpoint{4.888552in}{0.699520in}}%
\pgfpathlineto{\pgfqpoint{4.888715in}{0.691271in}}%
\pgfpathlineto{\pgfqpoint{4.889203in}{0.683022in}}%
\pgfpathlineto{\pgfqpoint{4.889218in}{0.724268in}}%
\pgfpathlineto{\pgfqpoint{4.890396in}{0.699520in}}%
\pgfpathlineto{\pgfqpoint{4.890544in}{0.707769in}}%
\pgfpathlineto{\pgfqpoint{4.890707in}{0.732517in}}%
\pgfpathlineto{\pgfqpoint{4.891565in}{0.691271in}}%
\pgfpathlineto{\pgfqpoint{4.891714in}{0.683022in}}%
\pgfpathlineto{\pgfqpoint{4.892069in}{0.724268in}}%
\pgfpathlineto{\pgfqpoint{4.893224in}{0.732517in}}%
\pgfpathlineto{\pgfqpoint{4.894231in}{0.707769in}}%
\pgfpathlineto{\pgfqpoint{4.894246in}{0.724268in}}%
\pgfpathlineto{\pgfqpoint{4.894564in}{0.707769in}}%
\pgfpathlineto{\pgfqpoint{4.894734in}{0.732517in}}%
\pgfpathlineto{\pgfqpoint{4.895082in}{0.724268in}}%
\pgfpathlineto{\pgfqpoint{4.896741in}{0.757264in}}%
\pgfpathlineto{\pgfqpoint{4.897933in}{0.724268in}}%
\pgfpathlineto{\pgfqpoint{4.899088in}{0.732517in}}%
\pgfpathlineto{\pgfqpoint{4.899925in}{0.707769in}}%
\pgfpathlineto{\pgfqpoint{4.900117in}{0.724268in}}%
\pgfpathlineto{\pgfqpoint{4.900450in}{0.749015in}}%
\pgfpathlineto{\pgfqpoint{4.900954in}{0.716019in}}%
\pgfpathlineto{\pgfqpoint{4.901961in}{0.674773in}}%
\pgfpathlineto{\pgfqpoint{4.902124in}{0.699520in}}%
\pgfpathlineto{\pgfqpoint{4.903131in}{0.724268in}}%
\pgfpathlineto{\pgfqpoint{4.903301in}{0.716019in}}%
\pgfpathlineto{\pgfqpoint{4.904789in}{0.683022in}}%
\pgfpathlineto{\pgfqpoint{4.903464in}{0.724268in}}%
\pgfpathlineto{\pgfqpoint{4.904789in}{0.691271in}}%
\pgfpathlineto{\pgfqpoint{4.905145in}{0.724268in}}%
\pgfpathlineto{\pgfqpoint{4.905123in}{0.683022in}}%
\pgfpathlineto{\pgfqpoint{4.905811in}{0.699520in}}%
\pgfpathlineto{\pgfqpoint{4.906152in}{0.724268in}}%
\pgfpathlineto{\pgfqpoint{4.906966in}{0.707769in}}%
\pgfpathlineto{\pgfqpoint{4.907995in}{0.699520in}}%
\pgfpathlineto{\pgfqpoint{4.909654in}{0.732517in}}%
\pgfpathlineto{\pgfqpoint{4.910342in}{0.724268in}}%
\pgfpathlineto{\pgfqpoint{4.910653in}{0.757264in}}%
\pgfpathlineto{\pgfqpoint{4.910676in}{0.749015in}}%
\pgfpathlineto{\pgfqpoint{4.911846in}{0.798510in}}%
\pgfpathlineto{\pgfqpoint{4.912186in}{0.773763in}}%
\pgfpathlineto{\pgfqpoint{4.913859in}{0.724268in}}%
\pgfpathlineto{\pgfqpoint{4.914008in}{0.732517in}}%
\pgfpathlineto{\pgfqpoint{4.914844in}{0.782012in}}%
\pgfpathlineto{\pgfqpoint{4.915037in}{0.773763in}}%
\pgfpathlineto{\pgfqpoint{4.915688in}{0.806759in}}%
\pgfpathlineto{\pgfqpoint{4.916192in}{0.782012in}}%
\pgfpathlineto{\pgfqpoint{4.916207in}{0.773763in}}%
\pgfpathlineto{\pgfqpoint{4.916525in}{0.806759in}}%
\pgfpathlineto{\pgfqpoint{4.917214in}{0.798510in}}%
\pgfpathlineto{\pgfqpoint{4.918539in}{0.856254in}}%
\pgfpathlineto{\pgfqpoint{4.918872in}{0.831507in}}%
\pgfpathlineto{\pgfqpoint{4.919228in}{0.823258in}}%
\pgfpathlineto{\pgfqpoint{4.919390in}{0.848005in}}%
\pgfpathlineto{\pgfqpoint{4.919894in}{0.839756in}}%
\pgfpathlineto{\pgfqpoint{4.921382in}{0.806759in}}%
\pgfpathlineto{\pgfqpoint{4.921382in}{0.815008in}}%
\pgfpathlineto{\pgfqpoint{4.922389in}{0.831507in}}%
\pgfpathlineto{\pgfqpoint{4.921738in}{0.798510in}}%
\pgfpathlineto{\pgfqpoint{4.922411in}{0.823258in}}%
\pgfpathlineto{\pgfqpoint{4.923226in}{0.831507in}}%
\pgfpathlineto{\pgfqpoint{4.923248in}{0.823258in}}%
\pgfpathlineto{\pgfqpoint{4.924588in}{0.798510in}}%
\pgfpathlineto{\pgfqpoint{4.925913in}{0.856254in}}%
\pgfpathlineto{\pgfqpoint{4.926247in}{0.831507in}}%
\pgfpathlineto{\pgfqpoint{4.927254in}{0.806759in}}%
\pgfpathlineto{\pgfqpoint{4.927268in}{0.823258in}}%
\pgfpathlineto{\pgfqpoint{4.929112in}{0.897500in}}%
\pgfpathlineto{\pgfqpoint{4.929282in}{0.889251in}}%
\pgfpathlineto{\pgfqpoint{4.930104in}{0.856254in}}%
\pgfpathlineto{\pgfqpoint{4.930119in}{0.897500in}}%
\pgfpathlineto{\pgfqpoint{4.930437in}{0.881002in}}%
\pgfpathlineto{\pgfqpoint{4.930460in}{0.897500in}}%
\pgfpathlineto{\pgfqpoint{4.930793in}{0.872753in}}%
\pgfpathlineto{\pgfqpoint{4.931459in}{0.872753in}}%
\pgfpathlineto{\pgfqpoint{4.933451in}{0.930497in}}%
\pgfpathlineto{\pgfqpoint{4.933792in}{0.905749in}}%
\pgfpathlineto{\pgfqpoint{4.934295in}{0.955244in}}%
\pgfpathlineto{\pgfqpoint{4.934480in}{0.938746in}}%
\pgfpathlineto{\pgfqpoint{4.934813in}{0.922247in}}%
\pgfpathlineto{\pgfqpoint{4.935132in}{0.955244in}}%
\pgfpathlineto{\pgfqpoint{4.935147in}{0.946995in}}%
\pgfpathlineto{\pgfqpoint{4.935798in}{0.930497in}}%
\pgfpathlineto{\pgfqpoint{4.936302in}{0.955244in}}%
\pgfpathlineto{\pgfqpoint{4.936827in}{0.946995in}}%
\pgfpathlineto{\pgfqpoint{4.937146in}{0.979992in}}%
\pgfpathlineto{\pgfqpoint{4.937331in}{0.971742in}}%
\pgfpathlineto{\pgfqpoint{4.938989in}{1.004739in}}%
\pgfpathlineto{\pgfqpoint{4.939174in}{0.996490in}}%
\pgfpathlineto{\pgfqpoint{4.939493in}{1.029487in}}%
\pgfpathlineto{\pgfqpoint{4.940011in}{1.021237in}}%
\pgfpathlineto{\pgfqpoint{4.941166in}{1.054234in}}%
\pgfpathlineto{\pgfqpoint{4.941351in}{1.037736in}}%
\pgfpathlineto{\pgfqpoint{4.942506in}{1.004739in}}%
\pgfpathlineto{\pgfqpoint{4.942529in}{1.021237in}}%
\pgfpathlineto{\pgfqpoint{4.943684in}{1.054234in}}%
\pgfpathlineto{\pgfqpoint{4.943869in}{1.037736in}}%
\pgfpathlineto{\pgfqpoint{4.944032in}{1.045985in}}%
\pgfpathlineto{\pgfqpoint{4.945024in}{1.029487in}}%
\pgfpathlineto{\pgfqpoint{4.945527in}{1.054234in}}%
\pgfpathlineto{\pgfqpoint{4.946045in}{1.037736in}}%
\pgfpathlineto{\pgfqpoint{4.946719in}{1.045985in}}%
\pgfpathlineto{\pgfqpoint{4.947704in}{1.004739in}}%
\pgfpathlineto{\pgfqpoint{4.948541in}{1.029487in}}%
\pgfpathlineto{\pgfqpoint{4.948059in}{0.996490in}}%
\pgfpathlineto{\pgfqpoint{4.948726in}{1.021237in}}%
\pgfpathlineto{\pgfqpoint{4.950066in}{1.070732in}}%
\pgfpathlineto{\pgfqpoint{4.950384in}{1.054234in}}%
\pgfpathlineto{\pgfqpoint{4.950910in}{1.021237in}}%
\pgfpathlineto{\pgfqpoint{4.951414in}{1.070732in}}%
\pgfpathlineto{\pgfqpoint{4.953065in}{1.029487in}}%
\pgfpathlineto{\pgfqpoint{4.953065in}{1.037736in}}%
\pgfpathlineto{\pgfqpoint{4.953568in}{1.103729in}}%
\pgfpathlineto{\pgfqpoint{4.954094in}{1.070732in}}%
\pgfpathlineto{\pgfqpoint{4.954575in}{1.078981in}}%
\pgfpathlineto{\pgfqpoint{4.955434in}{1.045985in}}%
\pgfpathlineto{\pgfqpoint{4.957596in}{1.128476in}}%
\pgfpathlineto{\pgfqpoint{4.957611in}{1.120227in}}%
\pgfpathlineto{\pgfqpoint{4.958448in}{1.095480in}}%
\pgfpathlineto{\pgfqpoint{4.958618in}{1.120227in}}%
\pgfpathlineto{\pgfqpoint{4.959269in}{1.128476in}}%
\pgfpathlineto{\pgfqpoint{4.959454in}{1.111978in}}%
\pgfpathlineto{\pgfqpoint{4.959943in}{1.103729in}}%
\pgfpathlineto{\pgfqpoint{4.960106in}{1.128476in}}%
\pgfpathlineto{\pgfqpoint{4.960128in}{1.120227in}}%
\pgfpathlineto{\pgfqpoint{4.961787in}{1.153224in}}%
\pgfpathlineto{\pgfqpoint{4.962809in}{1.144975in}}%
\pgfpathlineto{\pgfqpoint{4.963460in}{1.153224in}}%
\pgfpathlineto{\pgfqpoint{4.963645in}{1.136726in}}%
\pgfpathlineto{\pgfqpoint{4.964823in}{1.120227in}}%
\pgfpathlineto{\pgfqpoint{4.965978in}{1.153224in}}%
\pgfpathlineto{\pgfqpoint{4.966163in}{1.136726in}}%
\pgfpathlineto{\pgfqpoint{4.966326in}{1.144975in}}%
\pgfpathlineto{\pgfqpoint{4.967836in}{0.988241in}}%
\pgfpathlineto{\pgfqpoint{4.970353in}{0.650025in}}%
\pgfpathlineto{\pgfqpoint{4.971020in}{0.674773in}}%
\pgfpathlineto{\pgfqpoint{4.971842in}{0.707769in}}%
\pgfpathlineto{\pgfqpoint{4.972367in}{0.691271in}}%
\pgfpathlineto{\pgfqpoint{4.973034in}{0.674773in}}%
\pgfpathlineto{\pgfqpoint{4.972530in}{0.699520in}}%
\pgfpathlineto{\pgfqpoint{4.973204in}{0.699520in}}%
\pgfpathlineto{\pgfqpoint{4.973537in}{0.724268in}}%
\pgfpathlineto{\pgfqpoint{4.974041in}{0.691271in}}%
\pgfpathlineto{\pgfqpoint{4.975048in}{0.674773in}}%
\pgfpathlineto{\pgfqpoint{4.975196in}{0.683022in}}%
\pgfpathlineto{\pgfqpoint{4.975551in}{0.724268in}}%
\pgfpathlineto{\pgfqpoint{4.976218in}{0.724268in}}%
\pgfpathlineto{\pgfqpoint{4.977373in}{0.757264in}}%
\pgfpathlineto{\pgfqpoint{4.977558in}{0.740766in}}%
\pgfpathlineto{\pgfqpoint{4.977706in}{0.732517in}}%
\pgfpathlineto{\pgfqpoint{4.978061in}{0.773763in}}%
\pgfpathlineto{\pgfqpoint{4.979216in}{0.782012in}}%
\pgfpathlineto{\pgfqpoint{4.979401in}{0.773763in}}%
\pgfpathlineto{\pgfqpoint{4.979897in}{0.831507in}}%
\pgfpathlineto{\pgfqpoint{4.980075in}{0.823258in}}%
\pgfpathlineto{\pgfqpoint{4.981230in}{0.831507in}}%
\pgfpathlineto{\pgfqpoint{4.982252in}{0.823258in}}%
\pgfpathlineto{\pgfqpoint{4.983074in}{0.881002in}}%
\pgfpathlineto{\pgfqpoint{4.983592in}{0.839756in}}%
\pgfpathlineto{\pgfqpoint{4.983933in}{0.823258in}}%
\pgfpathlineto{\pgfqpoint{4.984244in}{0.864503in}}%
\pgfpathlineto{\pgfqpoint{4.984769in}{0.897500in}}%
\pgfpathlineto{\pgfqpoint{4.985273in}{0.872753in}}%
\pgfpathlineto{\pgfqpoint{4.986087in}{0.905749in}}%
\pgfpathlineto{\pgfqpoint{4.986280in}{0.889251in}}%
\pgfpathlineto{\pgfqpoint{4.987094in}{0.831507in}}%
\pgfpathlineto{\pgfqpoint{4.987620in}{0.848005in}}%
\pgfpathlineto{\pgfqpoint{4.988605in}{0.881002in}}%
\pgfpathlineto{\pgfqpoint{4.988790in}{0.864503in}}%
\pgfpathlineto{\pgfqpoint{4.989279in}{0.831507in}}%
\pgfpathlineto{\pgfqpoint{4.989627in}{0.872753in}}%
\pgfpathlineto{\pgfqpoint{4.989782in}{0.881002in}}%
\pgfpathlineto{\pgfqpoint{4.990130in}{0.848005in}}%
\pgfpathlineto{\pgfqpoint{4.990278in}{0.856254in}}%
\pgfpathlineto{\pgfqpoint{4.990804in}{0.848005in}}%
\pgfpathlineto{\pgfqpoint{4.991293in}{0.889251in}}%
\pgfpathlineto{\pgfqpoint{4.992292in}{0.905749in}}%
\pgfpathlineto{\pgfqpoint{4.991626in}{0.881002in}}%
\pgfpathlineto{\pgfqpoint{4.992314in}{0.897500in}}%
\pgfpathlineto{\pgfqpoint{4.997993in}{1.177971in}}%
\pgfpathlineto{\pgfqpoint{4.998178in}{1.161473in}}%
\pgfpathlineto{\pgfqpoint{4.998349in}{1.169722in}}%
\pgfpathlineto{\pgfqpoint{4.999852in}{1.012988in}}%
\pgfpathlineto{\pgfqpoint{5.003206in}{0.625278in}}%
\pgfpathlineto{\pgfqpoint{5.003376in}{0.650025in}}%
\pgfpathlineto{\pgfqpoint{5.004539in}{0.633527in}}%
\pgfpathlineto{\pgfqpoint{5.003865in}{0.658274in}}%
\pgfpathlineto{\pgfqpoint{5.004539in}{0.641776in}}%
\pgfpathlineto{\pgfqpoint{5.005546in}{0.658274in}}%
\pgfpathlineto{\pgfqpoint{5.005553in}{0.650025in}}%
\pgfpathlineto{\pgfqpoint{5.006893in}{0.724268in}}%
\pgfpathlineto{\pgfqpoint{5.007722in}{0.707769in}}%
\pgfpathlineto{\pgfqpoint{5.007737in}{0.699520in}}%
\pgfpathlineto{\pgfqpoint{5.008396in}{0.757264in}}%
\pgfpathlineto{\pgfqpoint{5.008737in}{0.749015in}}%
\pgfpathlineto{\pgfqpoint{5.009396in}{0.757264in}}%
\pgfpathlineto{\pgfqpoint{5.009581in}{0.740766in}}%
\pgfpathlineto{\pgfqpoint{5.010573in}{0.683022in}}%
\pgfpathlineto{\pgfqpoint{5.010921in}{0.724268in}}%
\pgfpathlineto{\pgfqpoint{5.011232in}{0.732517in}}%
\pgfpathlineto{\pgfqpoint{5.011573in}{0.707769in}}%
\pgfpathlineto{\pgfqpoint{5.012580in}{0.683022in}}%
\pgfpathlineto{\pgfqpoint{5.012594in}{0.699520in}}%
\pgfpathlineto{\pgfqpoint{5.013601in}{0.724268in}}%
\pgfpathlineto{\pgfqpoint{5.013772in}{0.716019in}}%
\pgfpathlineto{\pgfqpoint{5.016267in}{0.625278in}}%
\pgfpathlineto{\pgfqpoint{5.016267in}{0.625278in}}%
\pgfpathlineto{\pgfqpoint{5.017607in}{0.633527in}}%
\pgfpathlineto{\pgfqpoint{5.017607in}{0.625278in}}%
\pgfpathlineto{\pgfqpoint{5.018614in}{0.625278in}}%
\pgfpathlineto{\pgfqpoint{5.019962in}{0.633527in}}%
\pgfpathlineto{\pgfqpoint{5.020117in}{0.625278in}}%
\pgfpathlineto{\pgfqpoint{5.020961in}{0.625278in}}%
\pgfpathlineto{\pgfqpoint{5.022131in}{0.633527in}}%
\pgfpathlineto{\pgfqpoint{5.022131in}{0.625278in}}%
\pgfpathlineto{\pgfqpoint{5.023153in}{0.625278in}}%
\pgfpathlineto{\pgfqpoint{5.024308in}{0.633527in}}%
\pgfpathlineto{\pgfqpoint{5.024663in}{0.650025in}}%
\pgfpathlineto{\pgfqpoint{5.025337in}{0.625278in}}%
\pgfpathlineto{\pgfqpoint{5.026670in}{0.658274in}}%
\pgfpathlineto{\pgfqpoint{5.026847in}{0.641776in}}%
\pgfpathlineto{\pgfqpoint{5.028017in}{0.625278in}}%
\pgfpathlineto{\pgfqpoint{5.029343in}{0.658274in}}%
\pgfpathlineto{\pgfqpoint{5.029358in}{0.650025in}}%
\pgfpathlineto{\pgfqpoint{5.030350in}{0.633527in}}%
\pgfpathlineto{\pgfqpoint{5.030364in}{0.650025in}}%
\pgfpathlineto{\pgfqpoint{5.031201in}{0.699520in}}%
\pgfpathlineto{\pgfqpoint{5.031705in}{0.666524in}}%
\pgfpathlineto{\pgfqpoint{5.032186in}{0.658274in}}%
\pgfpathlineto{\pgfqpoint{5.032356in}{0.683022in}}%
\pgfpathlineto{\pgfqpoint{5.032378in}{0.674773in}}%
\pgfpathlineto{\pgfqpoint{5.033874in}{0.732517in}}%
\pgfpathlineto{\pgfqpoint{5.034052in}{0.716019in}}%
\pgfpathlineto{\pgfqpoint{5.036569in}{0.625278in}}%
\pgfpathlineto{\pgfqpoint{5.038228in}{0.691271in}}%
\pgfpathlineto{\pgfqpoint{5.039398in}{0.782012in}}%
\pgfpathlineto{\pgfqpoint{5.039420in}{0.773763in}}%
\pgfpathlineto{\pgfqpoint{5.043766in}{1.062483in}}%
\pgfpathlineto{\pgfqpoint{5.044107in}{1.070732in}}%
\pgfpathlineto{\pgfqpoint{5.044780in}{1.021237in}}%
\pgfpathlineto{\pgfqpoint{5.045284in}{0.996490in}}%
\pgfpathlineto{\pgfqpoint{5.045602in}{1.029487in}}%
\pgfpathlineto{\pgfqpoint{5.045935in}{1.078981in}}%
\pgfpathlineto{\pgfqpoint{5.046624in}{1.070732in}}%
\pgfpathlineto{\pgfqpoint{5.046957in}{1.095480in}}%
\pgfpathlineto{\pgfqpoint{5.047461in}{1.062483in}}%
\pgfpathlineto{\pgfqpoint{5.047609in}{1.054234in}}%
\pgfpathlineto{\pgfqpoint{5.047964in}{1.095480in}}%
\pgfpathlineto{\pgfqpoint{5.048283in}{1.103729in}}%
\pgfpathlineto{\pgfqpoint{5.048616in}{1.078981in}}%
\pgfpathlineto{\pgfqpoint{5.048801in}{1.087231in}}%
\pgfpathlineto{\pgfqpoint{5.049623in}{1.078981in}}%
\pgfpathlineto{\pgfqpoint{5.048971in}{1.095480in}}%
\pgfpathlineto{\pgfqpoint{5.049623in}{1.087231in}}%
\pgfpathlineto{\pgfqpoint{5.050311in}{1.144975in}}%
\pgfpathlineto{\pgfqpoint{5.050645in}{1.120227in}}%
\pgfpathlineto{\pgfqpoint{5.050985in}{1.144975in}}%
\pgfpathlineto{\pgfqpoint{5.051800in}{1.128476in}}%
\pgfpathlineto{\pgfqpoint{5.052992in}{1.062483in}}%
\pgfpathlineto{\pgfqpoint{5.057353in}{0.872753in}}%
\pgfpathlineto{\pgfqpoint{5.058338in}{0.930497in}}%
\pgfpathlineto{\pgfqpoint{5.059019in}{0.905749in}}%
\pgfpathlineto{\pgfqpoint{5.059530in}{0.897500in}}%
\pgfpathlineto{\pgfqpoint{5.059848in}{0.930497in}}%
\pgfpathlineto{\pgfqpoint{5.060033in}{0.922247in}}%
\pgfpathlineto{\pgfqpoint{5.062025in}{0.979992in}}%
\pgfpathlineto{\pgfqpoint{5.062047in}{0.971742in}}%
\pgfpathlineto{\pgfqpoint{5.063217in}{0.897500in}}%
\pgfpathlineto{\pgfqpoint{5.063891in}{0.922247in}}%
\pgfpathlineto{\pgfqpoint{5.065379in}{1.062483in}}%
\pgfpathlineto{\pgfqpoint{5.065712in}{1.078981in}}%
\pgfpathlineto{\pgfqpoint{5.066067in}{1.045985in}}%
\pgfpathlineto{\pgfqpoint{5.066408in}{1.045985in}}%
\pgfpathlineto{\pgfqpoint{5.067245in}{1.070732in}}%
\pgfpathlineto{\pgfqpoint{5.067556in}{1.054234in}}%
\pgfpathlineto{\pgfqpoint{5.068252in}{1.021237in}}%
\pgfpathlineto{\pgfqpoint{5.068733in}{1.029487in}}%
\pgfpathlineto{\pgfqpoint{5.068755in}{1.070732in}}%
\pgfpathlineto{\pgfqpoint{5.069755in}{1.045985in}}%
\pgfpathlineto{\pgfqpoint{5.070258in}{1.070732in}}%
\pgfpathlineto{\pgfqpoint{5.070599in}{1.045985in}}%
\pgfpathlineto{\pgfqpoint{5.071939in}{1.021237in}}%
\pgfpathlineto{\pgfqpoint{5.072420in}{1.004739in}}%
\pgfpathlineto{\pgfqpoint{5.073931in}{1.078981in}}%
\pgfpathlineto{\pgfqpoint{5.074767in}{1.054234in}}%
\pgfpathlineto{\pgfqpoint{5.074953in}{1.062483in}}%
\pgfpathlineto{\pgfqpoint{5.076130in}{1.045985in}}%
\pgfpathlineto{\pgfqpoint{5.077448in}{1.078981in}}%
\pgfpathlineto{\pgfqpoint{5.077470in}{1.070732in}}%
\pgfpathlineto{\pgfqpoint{5.077951in}{1.078981in}}%
\pgfpathlineto{\pgfqpoint{5.079476in}{0.913998in}}%
\pgfpathlineto{\pgfqpoint{5.082327in}{0.650025in}}%
\pgfpathlineto{\pgfqpoint{5.083504in}{0.699520in}}%
\pgfpathlineto{\pgfqpoint{5.083986in}{0.683022in}}%
\pgfpathlineto{\pgfqpoint{5.084008in}{0.674773in}}%
\pgfpathlineto{\pgfqpoint{5.084511in}{0.749015in}}%
\pgfpathlineto{\pgfqpoint{5.084845in}{0.749015in}}%
\pgfpathlineto{\pgfqpoint{5.087169in}{0.831507in}}%
\pgfpathlineto{\pgfqpoint{5.087192in}{0.823258in}}%
\pgfpathlineto{\pgfqpoint{5.088014in}{0.806759in}}%
\pgfpathlineto{\pgfqpoint{5.087673in}{0.831507in}}%
\pgfpathlineto{\pgfqpoint{5.088028in}{0.823258in}}%
\pgfpathlineto{\pgfqpoint{5.088510in}{0.831507in}}%
\pgfpathlineto{\pgfqpoint{5.088702in}{0.815008in}}%
\pgfpathlineto{\pgfqpoint{5.089539in}{0.773763in}}%
\pgfpathlineto{\pgfqpoint{5.089857in}{0.806759in}}%
\pgfpathlineto{\pgfqpoint{5.090205in}{0.848005in}}%
\pgfpathlineto{\pgfqpoint{5.090879in}{0.823258in}}%
\pgfpathlineto{\pgfqpoint{5.091701in}{0.831507in}}%
\pgfpathlineto{\pgfqpoint{5.091716in}{0.823258in}}%
\pgfpathlineto{\pgfqpoint{5.093878in}{0.757264in}}%
\pgfpathlineto{\pgfqpoint{5.093878in}{0.765513in}}%
\pgfpathlineto{\pgfqpoint{5.094899in}{0.872753in}}%
\pgfpathlineto{\pgfqpoint{5.096240in}{0.848005in}}%
\pgfpathlineto{\pgfqpoint{5.096913in}{0.872753in}}%
\pgfpathlineto{\pgfqpoint{5.097247in}{0.848005in}}%
\pgfpathlineto{\pgfqpoint{5.098572in}{0.831507in}}%
\pgfpathlineto{\pgfqpoint{5.099075in}{0.856254in}}%
\pgfpathlineto{\pgfqpoint{5.099594in}{0.848005in}}%
\pgfpathlineto{\pgfqpoint{5.101082in}{0.806759in}}%
\pgfpathlineto{\pgfqpoint{5.101104in}{0.823258in}}%
\pgfpathlineto{\pgfqpoint{5.101585in}{0.831507in}}%
\pgfpathlineto{\pgfqpoint{5.101778in}{0.815008in}}%
\pgfpathlineto{\pgfqpoint{5.102430in}{0.831507in}}%
\pgfpathlineto{\pgfqpoint{5.102948in}{0.798510in}}%
\pgfpathlineto{\pgfqpoint{5.103784in}{0.823258in}}%
\pgfpathlineto{\pgfqpoint{5.104103in}{0.806759in}}%
\pgfpathlineto{\pgfqpoint{5.104458in}{0.823258in}}%
\pgfpathlineto{\pgfqpoint{5.105295in}{0.773763in}}%
\pgfpathlineto{\pgfqpoint{5.107309in}{0.848005in}}%
\pgfpathlineto{\pgfqpoint{5.107472in}{0.839756in}}%
\pgfpathlineto{\pgfqpoint{5.108627in}{0.782012in}}%
\pgfpathlineto{\pgfqpoint{5.108975in}{0.815008in}}%
\pgfpathlineto{\pgfqpoint{5.109315in}{0.823258in}}%
\pgfpathlineto{\pgfqpoint{5.109656in}{0.798510in}}%
\pgfpathlineto{\pgfqpoint{5.109989in}{0.798510in}}%
\pgfpathlineto{\pgfqpoint{5.112825in}{0.905749in}}%
\pgfpathlineto{\pgfqpoint{5.112840in}{0.897500in}}%
\pgfpathlineto{\pgfqpoint{5.113166in}{0.881002in}}%
\pgfpathlineto{\pgfqpoint{5.113321in}{0.905749in}}%
\pgfpathlineto{\pgfqpoint{5.113676in}{0.897500in}}%
\pgfpathlineto{\pgfqpoint{5.114854in}{0.922247in}}%
\pgfpathlineto{\pgfqpoint{5.115017in}{0.913998in}}%
\pgfpathlineto{\pgfqpoint{5.116860in}{0.823258in}}%
\pgfpathlineto{\pgfqpoint{5.115690in}{0.922247in}}%
\pgfpathlineto{\pgfqpoint{5.117193in}{0.848005in}}%
\pgfpathlineto{\pgfqpoint{5.117852in}{0.856254in}}%
\pgfpathlineto{\pgfqpoint{5.118038in}{0.839756in}}%
\pgfpathlineto{\pgfqpoint{5.118741in}{0.831507in}}%
\pgfpathlineto{\pgfqpoint{5.118200in}{0.848005in}}%
\pgfpathlineto{\pgfqpoint{5.118852in}{0.839756in}}%
\pgfpathlineto{\pgfqpoint{5.119541in}{0.872753in}}%
\pgfpathlineto{\pgfqpoint{5.119881in}{0.848005in}}%
\pgfpathlineto{\pgfqpoint{5.121207in}{0.881002in}}%
\pgfpathlineto{\pgfqpoint{5.121221in}{0.872753in}}%
\pgfpathlineto{\pgfqpoint{5.124383in}{0.757264in}}%
\pgfpathlineto{\pgfqpoint{5.124383in}{0.765513in}}%
\pgfpathlineto{\pgfqpoint{5.124738in}{0.798510in}}%
\pgfpathlineto{\pgfqpoint{5.125412in}{0.740766in}}%
\pgfpathlineto{\pgfqpoint{5.125575in}{0.749015in}}%
\pgfpathlineto{\pgfqpoint{5.127085in}{0.699520in}}%
\pgfpathlineto{\pgfqpoint{5.127426in}{0.724268in}}%
\pgfpathlineto{\pgfqpoint{5.127922in}{0.699520in}}%
\pgfpathlineto{\pgfqpoint{5.129270in}{0.625278in}}%
\pgfpathlineto{\pgfqpoint{5.129766in}{0.650025in}}%
\pgfpathlineto{\pgfqpoint{5.131091in}{0.683022in}}%
\pgfpathlineto{\pgfqpoint{5.131113in}{0.674773in}}%
\pgfpathlineto{\pgfqpoint{5.131780in}{0.650025in}}%
\pgfpathlineto{\pgfqpoint{5.131950in}{0.674773in}}%
\pgfpathlineto{\pgfqpoint{5.132787in}{0.699520in}}%
\pgfpathlineto{\pgfqpoint{5.133112in}{0.683022in}}%
\pgfpathlineto{\pgfqpoint{5.134119in}{0.658274in}}%
\pgfpathlineto{\pgfqpoint{5.134127in}{0.674773in}}%
\pgfpathlineto{\pgfqpoint{5.134949in}{0.683022in}}%
\pgfpathlineto{\pgfqpoint{5.134964in}{0.674773in}}%
\pgfpathlineto{\pgfqpoint{5.136644in}{0.625278in}}%
\pgfpathlineto{\pgfqpoint{5.136792in}{0.633527in}}%
\pgfpathlineto{\pgfqpoint{5.137814in}{0.674773in}}%
\pgfpathlineto{\pgfqpoint{5.138969in}{0.683022in}}%
\pgfpathlineto{\pgfqpoint{5.140161in}{0.625278in}}%
\pgfpathlineto{\pgfqpoint{5.141324in}{0.633527in}}%
\pgfpathlineto{\pgfqpoint{5.141494in}{0.625278in}}%
\pgfpathlineto{\pgfqpoint{5.142331in}{0.625278in}}%
\pgfpathlineto{\pgfqpoint{5.143493in}{0.633527in}}%
\pgfpathlineto{\pgfqpoint{5.143834in}{0.625278in}}%
\pgfpathlineto{\pgfqpoint{5.144500in}{0.625278in}}%
\pgfpathlineto{\pgfqpoint{5.145507in}{0.633527in}}%
\pgfpathlineto{\pgfqpoint{5.145507in}{0.625278in}}%
\pgfpathlineto{\pgfqpoint{5.146684in}{0.633527in}}%
\pgfpathlineto{\pgfqpoint{5.146855in}{0.625278in}}%
\pgfpathlineto{\pgfqpoint{5.147691in}{0.625278in}}%
\pgfpathlineto{\pgfqpoint{5.148861in}{0.633527in}}%
\pgfpathlineto{\pgfqpoint{5.148869in}{0.625278in}}%
\pgfpathlineto{\pgfqpoint{5.149868in}{0.625278in}}%
\pgfpathlineto{\pgfqpoint{5.150875in}{0.633527in}}%
\pgfpathlineto{\pgfqpoint{5.150875in}{0.625278in}}%
\pgfpathlineto{\pgfqpoint{5.152045in}{0.633527in}}%
\pgfpathlineto{\pgfqpoint{5.152215in}{0.625278in}}%
\pgfpathlineto{\pgfqpoint{5.153059in}{0.625278in}}%
\pgfpathlineto{\pgfqpoint{5.154222in}{0.633527in}}%
\pgfpathlineto{\pgfqpoint{5.154229in}{0.625278in}}%
\pgfpathlineto{\pgfqpoint{5.155066in}{0.625278in}}%
\pgfpathlineto{\pgfqpoint{5.156236in}{0.633527in}}%
\pgfpathlineto{\pgfqpoint{5.156236in}{0.625278in}}%
\pgfpathlineto{\pgfqpoint{5.157243in}{0.625278in}}%
\pgfpathlineto{\pgfqpoint{5.158413in}{0.633527in}}%
\pgfpathlineto{\pgfqpoint{5.158583in}{0.625278in}}%
\pgfpathlineto{\pgfqpoint{5.159257in}{0.625278in}}%
\pgfpathlineto{\pgfqpoint{5.160597in}{0.633527in}}%
\pgfpathlineto{\pgfqpoint{5.160760in}{0.625278in}}%
\pgfpathlineto{\pgfqpoint{5.161434in}{0.625278in}}%
\pgfpathlineto{\pgfqpoint{5.162774in}{0.633527in}}%
\pgfpathlineto{\pgfqpoint{5.162774in}{0.625278in}}%
\pgfpathlineto{\pgfqpoint{5.163781in}{0.625278in}}%
\pgfpathlineto{\pgfqpoint{5.164951in}{0.633527in}}%
\pgfpathlineto{\pgfqpoint{5.165121in}{0.625278in}}%
\pgfpathlineto{\pgfqpoint{5.165957in}{0.625278in}}%
\pgfpathlineto{\pgfqpoint{5.167142in}{0.633527in}}%
\pgfpathlineto{\pgfqpoint{5.167298in}{0.625278in}}%
\pgfpathlineto{\pgfqpoint{5.168134in}{0.625278in}}%
\pgfpathlineto{\pgfqpoint{5.169141in}{0.633527in}}%
\pgfpathlineto{\pgfqpoint{5.169141in}{0.625278in}}%
\pgfpathlineto{\pgfqpoint{5.170319in}{0.633527in}}%
\pgfpathlineto{\pgfqpoint{5.170481in}{0.625278in}}%
\pgfpathlineto{\pgfqpoint{5.171155in}{0.625278in}}%
\pgfpathlineto{\pgfqpoint{5.172325in}{0.633527in}}%
\pgfpathlineto{\pgfqpoint{5.172495in}{0.625278in}}%
\pgfpathlineto{\pgfqpoint{5.172666in}{0.625278in}}%
\pgfusepath{stroke}%
\end{pgfscope}%
\begin{pgfscope}%
\pgfpathrectangle{\pgfqpoint{0.586112in}{0.502778in}}{\pgfqpoint{4.805000in}{2.695000in}}%
\pgfusepath{clip}%
\pgfsetbuttcap%
\pgfsetroundjoin%
\pgfsetlinewidth{1.505625pt}%
\definecolor{currentstroke}{rgb}{0.400000,0.400000,0.400000}%
\pgfsetstrokecolor{currentstroke}%
\pgfsetdash{{5.550000pt}{2.400000pt}}{0.000000pt}%
\pgfpathmoveto{\pgfqpoint{0.804521in}{0.633527in}}%
\pgfpathlineto{\pgfqpoint{0.804521in}{0.625278in}}%
\pgfpathlineto{\pgfqpoint{0.804521in}{0.625278in}}%
\pgfpathlineto{\pgfqpoint{0.806399in}{0.633527in}}%
\pgfpathlineto{\pgfqpoint{0.806416in}{0.625278in}}%
\pgfpathlineto{\pgfqpoint{0.807508in}{0.625278in}}%
\pgfpathlineto{\pgfqpoint{0.808565in}{0.633527in}}%
\pgfpathlineto{\pgfqpoint{0.808565in}{0.625278in}}%
\pgfpathlineto{\pgfqpoint{0.809751in}{0.633527in}}%
\pgfpathlineto{\pgfqpoint{0.809769in}{0.625278in}}%
\pgfpathlineto{\pgfqpoint{0.810738in}{0.625278in}}%
\pgfpathlineto{\pgfqpoint{0.810906in}{0.633527in}}%
\pgfpathlineto{\pgfqpoint{0.810906in}{0.625278in}}%
\pgfpathlineto{\pgfqpoint{0.812415in}{0.633527in}}%
\pgfpathlineto{\pgfqpoint{0.812583in}{0.625278in}}%
\pgfpathlineto{\pgfqpoint{0.813421in}{0.625278in}}%
\pgfpathlineto{\pgfqpoint{0.813923in}{0.691271in}}%
\pgfpathlineto{\pgfqpoint{0.814762in}{0.881002in}}%
\pgfpathlineto{\pgfqpoint{0.814948in}{0.872753in}}%
\pgfpathlineto{\pgfqpoint{0.815767in}{0.782012in}}%
\pgfpathlineto{\pgfqpoint{0.815954in}{0.773763in}}%
\pgfpathlineto{\pgfqpoint{0.816605in}{0.889251in}}%
\pgfpathlineto{\pgfqpoint{0.817630in}{1.318207in}}%
\pgfpathlineto{\pgfqpoint{0.818282in}{1.375951in}}%
\pgfpathlineto{\pgfqpoint{0.818636in}{1.334705in}}%
\pgfpathlineto{\pgfqpoint{0.819306in}{1.268712in}}%
\pgfpathlineto{\pgfqpoint{0.819791in}{1.326456in}}%
\pgfpathlineto{\pgfqpoint{0.820815in}{1.763662in}}%
\pgfpathlineto{\pgfqpoint{0.821299in}{1.846153in}}%
\pgfpathlineto{\pgfqpoint{0.821467in}{1.870901in}}%
\pgfpathlineto{\pgfqpoint{0.822324in}{1.780160in}}%
\pgfpathlineto{\pgfqpoint{0.822659in}{1.763662in}}%
\pgfpathlineto{\pgfqpoint{0.822975in}{1.804907in}}%
\pgfpathlineto{\pgfqpoint{0.824987in}{2.365850in}}%
\pgfpathlineto{\pgfqpoint{0.825173in}{2.349352in}}%
\pgfpathlineto{\pgfqpoint{0.826012in}{2.258611in}}%
\pgfpathlineto{\pgfqpoint{0.826496in}{2.316355in}}%
\pgfpathlineto{\pgfqpoint{0.827521in}{2.728813in}}%
\pgfpathlineto{\pgfqpoint{0.828172in}{2.786557in}}%
\pgfpathlineto{\pgfqpoint{0.828675in}{2.737062in}}%
\pgfpathlineto{\pgfqpoint{0.829197in}{2.704066in}}%
\pgfpathlineto{\pgfqpoint{0.829680in}{2.770059in}}%
\pgfpathlineto{\pgfqpoint{0.831209in}{3.001035in}}%
\pgfpathlineto{\pgfqpoint{0.833052in}{3.075278in}}%
\pgfpathlineto{\pgfqpoint{0.833220in}{3.067029in}}%
\pgfpathlineto{\pgfqpoint{0.833388in}{3.075278in}}%
\pgfpathlineto{\pgfqpoint{0.834375in}{3.058779in}}%
\pgfpathlineto{\pgfqpoint{0.834393in}{3.075278in}}%
\pgfpathlineto{\pgfqpoint{0.835399in}{3.025783in}}%
\pgfpathlineto{\pgfqpoint{0.835716in}{3.034032in}}%
\pgfpathlineto{\pgfqpoint{0.836051in}{3.009285in}}%
\pgfpathlineto{\pgfqpoint{0.836238in}{3.017534in}}%
\pgfpathlineto{\pgfqpoint{0.837056in}{3.009285in}}%
\pgfpathlineto{\pgfqpoint{0.836405in}{3.025783in}}%
\pgfpathlineto{\pgfqpoint{0.837076in}{3.025783in}}%
\pgfpathlineto{\pgfqpoint{0.838230in}{3.058779in}}%
\pgfpathlineto{\pgfqpoint{0.838417in}{3.042281in}}%
\pgfpathlineto{\pgfqpoint{0.855515in}{0.650025in}}%
\pgfpathlineto{\pgfqpoint{0.856501in}{0.658274in}}%
\pgfpathlineto{\pgfqpoint{0.856521in}{0.650025in}}%
\pgfpathlineto{\pgfqpoint{0.856837in}{0.633527in}}%
\pgfpathlineto{\pgfqpoint{0.857005in}{0.658274in}}%
\pgfpathlineto{\pgfqpoint{0.857359in}{0.650025in}}%
\pgfpathlineto{\pgfqpoint{0.857843in}{0.658274in}}%
\pgfpathlineto{\pgfqpoint{0.858030in}{0.641776in}}%
\pgfpathlineto{\pgfqpoint{0.858197in}{0.650025in}}%
\pgfpathlineto{\pgfqpoint{0.859184in}{0.633527in}}%
\pgfpathlineto{\pgfqpoint{0.859202in}{0.650025in}}%
\pgfpathlineto{\pgfqpoint{0.860209in}{0.650025in}}%
\pgfpathlineto{\pgfqpoint{0.861530in}{0.625278in}}%
\pgfpathlineto{\pgfqpoint{0.861530in}{0.625278in}}%
\pgfpathlineto{\pgfqpoint{0.862704in}{0.633527in}}%
\pgfpathlineto{\pgfqpoint{0.862704in}{0.625278in}}%
\pgfpathlineto{\pgfqpoint{0.863710in}{0.625278in}}%
\pgfpathlineto{\pgfqpoint{0.864883in}{0.633527in}}%
\pgfpathlineto{\pgfqpoint{0.864902in}{0.625278in}}%
\pgfpathlineto{\pgfqpoint{0.865386in}{0.658274in}}%
\pgfpathlineto{\pgfqpoint{0.865908in}{0.641776in}}%
\pgfpathlineto{\pgfqpoint{0.866243in}{0.625278in}}%
\pgfpathlineto{\pgfqpoint{0.866559in}{0.658274in}}%
\pgfpathlineto{\pgfqpoint{0.866578in}{0.650025in}}%
\pgfpathlineto{\pgfqpoint{0.867230in}{0.658274in}}%
\pgfpathlineto{\pgfqpoint{0.867063in}{0.633527in}}%
\pgfpathlineto{\pgfqpoint{0.867417in}{0.641776in}}%
\pgfpathlineto{\pgfqpoint{0.868422in}{0.625278in}}%
\pgfpathlineto{\pgfqpoint{0.869088in}{0.666524in}}%
\pgfpathlineto{\pgfqpoint{0.869764in}{0.699520in}}%
\pgfpathlineto{\pgfqpoint{0.869416in}{0.658274in}}%
\pgfpathlineto{\pgfqpoint{0.870098in}{0.674773in}}%
\pgfpathlineto{\pgfqpoint{0.871272in}{0.650025in}}%
\pgfpathlineto{\pgfqpoint{0.871274in}{0.658274in}}%
\pgfpathlineto{\pgfqpoint{0.871775in}{0.699520in}}%
\pgfpathlineto{\pgfqpoint{0.872278in}{0.666524in}}%
\pgfpathlineto{\pgfqpoint{0.873116in}{0.650025in}}%
\pgfpathlineto{\pgfqpoint{0.872446in}{0.674773in}}%
\pgfpathlineto{\pgfqpoint{0.873279in}{0.666524in}}%
\pgfpathlineto{\pgfqpoint{0.873782in}{0.683022in}}%
\pgfpathlineto{\pgfqpoint{0.874122in}{0.650025in}}%
\pgfpathlineto{\pgfqpoint{0.874292in}{0.658274in}}%
\pgfpathlineto{\pgfqpoint{0.875295in}{0.650025in}}%
\pgfpathlineto{\pgfqpoint{0.874625in}{0.674773in}}%
\pgfpathlineto{\pgfqpoint{0.875298in}{0.658274in}}%
\pgfpathlineto{\pgfqpoint{0.875800in}{0.683022in}}%
\pgfpathlineto{\pgfqpoint{0.876303in}{0.683022in}}%
\pgfpathlineto{\pgfqpoint{0.877307in}{0.674773in}}%
\pgfpathlineto{\pgfqpoint{0.877309in}{0.683022in}}%
\pgfpathlineto{\pgfqpoint{0.877812in}{0.707769in}}%
\pgfpathlineto{\pgfqpoint{0.878147in}{0.683022in}}%
\pgfpathlineto{\pgfqpoint{0.878818in}{0.658274in}}%
\pgfpathlineto{\pgfqpoint{0.878818in}{0.691271in}}%
\pgfpathlineto{\pgfqpoint{0.879153in}{0.683022in}}%
\pgfpathlineto{\pgfqpoint{0.880159in}{0.790261in}}%
\pgfpathlineto{\pgfqpoint{0.881165in}{0.938746in}}%
\pgfpathlineto{\pgfqpoint{0.881500in}{0.930497in}}%
\pgfpathlineto{\pgfqpoint{0.884350in}{1.153224in}}%
\pgfpathlineto{\pgfqpoint{0.886864in}{1.252214in}}%
\pgfpathlineto{\pgfqpoint{0.887030in}{1.243965in}}%
\pgfpathlineto{\pgfqpoint{0.888205in}{1.227466in}}%
\pgfpathlineto{\pgfqpoint{0.887528in}{1.252214in}}%
\pgfpathlineto{\pgfqpoint{0.888205in}{1.235715in}}%
\pgfpathlineto{\pgfqpoint{0.889041in}{1.268712in}}%
\pgfpathlineto{\pgfqpoint{0.889379in}{1.252214in}}%
\pgfpathlineto{\pgfqpoint{0.890552in}{1.227466in}}%
\pgfpathlineto{\pgfqpoint{0.889711in}{1.268712in}}%
\pgfpathlineto{\pgfqpoint{0.890552in}{1.235715in}}%
\pgfpathlineto{\pgfqpoint{0.891388in}{1.243965in}}%
\pgfpathlineto{\pgfqpoint{0.890887in}{1.227466in}}%
\pgfpathlineto{\pgfqpoint{0.891556in}{1.235715in}}%
\pgfpathlineto{\pgfqpoint{0.891723in}{1.243965in}}%
\pgfpathlineto{\pgfqpoint{0.892563in}{1.227466in}}%
\pgfpathlineto{\pgfqpoint{0.892729in}{1.243965in}}%
\pgfpathlineto{\pgfqpoint{0.893232in}{1.210968in}}%
\pgfpathlineto{\pgfqpoint{0.894238in}{1.194470in}}%
\pgfpathlineto{\pgfqpoint{0.893402in}{1.227466in}}%
\pgfpathlineto{\pgfqpoint{0.894240in}{1.202719in}}%
\pgfpathlineto{\pgfqpoint{0.896081in}{1.062483in}}%
\pgfpathlineto{\pgfqpoint{0.899101in}{0.732517in}}%
\pgfpathlineto{\pgfqpoint{0.899266in}{0.749015in}}%
\pgfpathlineto{\pgfqpoint{0.900105in}{0.666524in}}%
\pgfpathlineto{\pgfqpoint{0.901106in}{0.625278in}}%
\pgfpathlineto{\pgfqpoint{0.902279in}{0.633527in}}%
\pgfpathlineto{\pgfqpoint{0.902621in}{0.625278in}}%
\pgfpathlineto{\pgfqpoint{0.902957in}{0.625278in}}%
\pgfpathlineto{\pgfqpoint{0.904466in}{0.633527in}}%
\pgfpathlineto{\pgfqpoint{0.905471in}{0.625278in}}%
\pgfpathlineto{\pgfqpoint{0.905974in}{0.633527in}}%
\pgfpathlineto{\pgfqpoint{0.905974in}{0.625278in}}%
\pgfpathlineto{\pgfqpoint{0.908314in}{0.633527in}}%
\pgfpathlineto{\pgfqpoint{0.908649in}{0.625278in}}%
\pgfpathlineto{\pgfqpoint{0.908984in}{0.625278in}}%
\pgfpathlineto{\pgfqpoint{0.910325in}{0.633527in}}%
\pgfpathlineto{\pgfqpoint{0.910661in}{0.625278in}}%
\pgfpathlineto{\pgfqpoint{0.910996in}{0.625278in}}%
\pgfpathlineto{\pgfqpoint{0.912504in}{0.633527in}}%
\pgfpathlineto{\pgfqpoint{0.912504in}{0.625278in}}%
\pgfpathlineto{\pgfqpoint{0.917708in}{0.633527in}}%
\pgfpathlineto{\pgfqpoint{0.917876in}{0.625278in}}%
\pgfpathlineto{\pgfqpoint{0.919552in}{0.633527in}}%
\pgfpathlineto{\pgfqpoint{0.919887in}{0.625278in}}%
\pgfpathlineto{\pgfqpoint{0.920558in}{0.625278in}}%
\pgfpathlineto{\pgfqpoint{0.921899in}{0.633527in}}%
\pgfpathlineto{\pgfqpoint{0.922234in}{0.625278in}}%
\pgfpathlineto{\pgfqpoint{0.922905in}{0.625278in}}%
\pgfpathlineto{\pgfqpoint{0.924078in}{0.633527in}}%
\pgfpathlineto{\pgfqpoint{0.924246in}{0.625278in}}%
\pgfpathlineto{\pgfqpoint{0.924916in}{0.625278in}}%
\pgfpathlineto{\pgfqpoint{0.926092in}{0.633527in}}%
\pgfpathlineto{\pgfqpoint{0.926258in}{0.625278in}}%
\pgfpathlineto{\pgfqpoint{0.926930in}{0.625278in}}%
\pgfpathlineto{\pgfqpoint{0.928269in}{0.633527in}}%
\pgfpathlineto{\pgfqpoint{0.928604in}{0.625278in}}%
\pgfpathlineto{\pgfqpoint{0.929275in}{0.625278in}}%
\pgfpathlineto{\pgfqpoint{0.930616in}{0.633527in}}%
\pgfpathlineto{\pgfqpoint{0.930951in}{0.625278in}}%
\pgfpathlineto{\pgfqpoint{0.931621in}{0.625278in}}%
\pgfpathlineto{\pgfqpoint{0.932797in}{0.633527in}}%
\pgfpathlineto{\pgfqpoint{0.932963in}{0.625278in}}%
\pgfpathlineto{\pgfqpoint{0.933810in}{0.625278in}}%
\pgfpathlineto{\pgfqpoint{0.934974in}{0.633527in}}%
\pgfpathlineto{\pgfqpoint{0.935146in}{0.625278in}}%
\pgfpathlineto{\pgfqpoint{0.935980in}{0.625278in}}%
\pgfpathlineto{\pgfqpoint{0.937166in}{0.633527in}}%
\pgfpathlineto{\pgfqpoint{0.937321in}{0.625278in}}%
\pgfpathlineto{\pgfqpoint{0.938169in}{0.625278in}}%
\pgfpathlineto{\pgfqpoint{0.939333in}{0.633527in}}%
\pgfpathlineto{\pgfqpoint{0.939514in}{0.625278in}}%
\pgfpathlineto{\pgfqpoint{0.940338in}{0.625278in}}%
\pgfpathlineto{\pgfqpoint{0.941528in}{0.633527in}}%
\pgfpathlineto{\pgfqpoint{0.941679in}{0.625278in}}%
\pgfpathlineto{\pgfqpoint{0.942366in}{0.625278in}}%
\pgfpathlineto{\pgfqpoint{0.943691in}{0.633527in}}%
\pgfpathlineto{\pgfqpoint{0.943708in}{0.625278in}}%
\pgfpathlineto{\pgfqpoint{0.944713in}{0.625278in}}%
\pgfpathlineto{\pgfqpoint{0.945873in}{0.633527in}}%
\pgfpathlineto{\pgfqpoint{0.946038in}{0.625278in}}%
\pgfpathlineto{\pgfqpoint{0.946878in}{0.625278in}}%
\pgfpathlineto{\pgfqpoint{0.948050in}{0.633527in}}%
\pgfpathlineto{\pgfqpoint{0.948059in}{0.625278in}}%
\pgfpathlineto{\pgfqpoint{0.949072in}{0.625278in}}%
\pgfpathlineto{\pgfqpoint{0.950231in}{0.633527in}}%
\pgfpathlineto{\pgfqpoint{0.950245in}{0.625278in}}%
\pgfpathlineto{\pgfqpoint{0.951251in}{0.625278in}}%
\pgfpathlineto{\pgfqpoint{0.952424in}{0.633527in}}%
\pgfpathlineto{\pgfqpoint{0.952578in}{0.625278in}}%
\pgfpathlineto{\pgfqpoint{0.953262in}{0.625278in}}%
\pgfpathlineto{\pgfqpoint{0.954414in}{0.633527in}}%
\pgfpathlineto{\pgfqpoint{0.954604in}{0.625278in}}%
\pgfpathlineto{\pgfqpoint{0.955442in}{0.625278in}}%
\pgfpathlineto{\pgfqpoint{0.956594in}{0.633527in}}%
\pgfpathlineto{\pgfqpoint{0.956762in}{0.625278in}}%
\pgfpathlineto{\pgfqpoint{0.957432in}{0.625278in}}%
\pgfpathlineto{\pgfqpoint{0.958773in}{0.633527in}}%
\pgfpathlineto{\pgfqpoint{0.958773in}{0.625278in}}%
\pgfpathlineto{\pgfqpoint{0.959800in}{0.625278in}}%
\pgfpathlineto{\pgfqpoint{0.960974in}{0.633527in}}%
\pgfpathlineto{\pgfqpoint{0.961120in}{0.625278in}}%
\pgfpathlineto{\pgfqpoint{0.962079in}{0.625278in}}%
\pgfpathlineto{\pgfqpoint{0.963153in}{0.633527in}}%
\pgfpathlineto{\pgfqpoint{0.963153in}{0.625278in}}%
\pgfpathlineto{\pgfqpoint{0.964472in}{0.633527in}}%
\pgfpathlineto{\pgfqpoint{0.964640in}{0.625278in}}%
\pgfpathlineto{\pgfqpoint{0.965478in}{0.625278in}}%
\pgfpathlineto{\pgfqpoint{0.966673in}{0.633527in}}%
\pgfpathlineto{\pgfqpoint{0.966766in}{0.625278in}}%
\pgfpathlineto{\pgfqpoint{0.967679in}{0.625278in}}%
\pgfpathlineto{\pgfqpoint{0.968831in}{0.633527in}}%
\pgfpathlineto{\pgfqpoint{0.968998in}{0.625278in}}%
\pgfpathlineto{\pgfqpoint{0.969837in}{0.625278in}}%
\pgfpathlineto{\pgfqpoint{0.971010in}{0.633527in}}%
\pgfpathlineto{\pgfqpoint{0.971177in}{0.625278in}}%
\pgfpathlineto{\pgfqpoint{0.972037in}{0.625278in}}%
\pgfpathlineto{\pgfqpoint{0.973189in}{0.633527in}}%
\pgfpathlineto{\pgfqpoint{0.973357in}{0.625278in}}%
\pgfpathlineto{\pgfqpoint{0.974195in}{0.625278in}}%
\pgfpathlineto{\pgfqpoint{0.975368in}{0.633527in}}%
\pgfpathlineto{\pgfqpoint{0.975536in}{0.625278in}}%
\pgfpathlineto{\pgfqpoint{0.976374in}{0.625278in}}%
\pgfpathlineto{\pgfqpoint{0.977548in}{0.633527in}}%
\pgfpathlineto{\pgfqpoint{0.977715in}{0.625278in}}%
\pgfpathlineto{\pgfqpoint{0.978554in}{0.625278in}}%
\pgfpathlineto{\pgfqpoint{0.979749in}{0.633527in}}%
\pgfpathlineto{\pgfqpoint{0.979894in}{0.625278in}}%
\pgfpathlineto{\pgfqpoint{0.980565in}{0.625278in}}%
\pgfpathlineto{\pgfqpoint{0.981739in}{0.633527in}}%
\pgfpathlineto{\pgfqpoint{0.981928in}{0.625278in}}%
\pgfpathlineto{\pgfqpoint{0.982744in}{0.625278in}}%
\pgfpathlineto{\pgfqpoint{0.983918in}{0.633527in}}%
\pgfpathlineto{\pgfqpoint{0.984085in}{0.625278in}}%
\pgfpathlineto{\pgfqpoint{0.984923in}{0.625278in}}%
\pgfpathlineto{\pgfqpoint{0.986097in}{0.633527in}}%
\pgfpathlineto{\pgfqpoint{0.986286in}{0.625278in}}%
\pgfpathlineto{\pgfqpoint{0.987117in}{0.625278in}}%
\pgfpathlineto{\pgfqpoint{0.988276in}{0.633527in}}%
\pgfpathlineto{\pgfqpoint{0.988444in}{0.625278in}}%
\pgfpathlineto{\pgfqpoint{0.989282in}{0.625278in}}%
\pgfpathlineto{\pgfqpoint{0.990455in}{0.633527in}}%
\pgfpathlineto{\pgfqpoint{0.990623in}{0.625278in}}%
\pgfpathlineto{\pgfqpoint{0.991475in}{0.625278in}}%
\pgfpathlineto{\pgfqpoint{0.992635in}{0.633527in}}%
\pgfpathlineto{\pgfqpoint{0.992656in}{0.625278in}}%
\pgfpathlineto{\pgfqpoint{0.993662in}{0.625278in}}%
\pgfpathlineto{\pgfqpoint{0.994814in}{0.633527in}}%
\pgfpathlineto{\pgfqpoint{0.994835in}{0.625278in}}%
\pgfpathlineto{\pgfqpoint{0.995841in}{0.625278in}}%
\pgfpathlineto{\pgfqpoint{0.997000in}{0.633527in}}%
\pgfpathlineto{\pgfqpoint{0.997014in}{0.625278in}}%
\pgfpathlineto{\pgfqpoint{0.997998in}{0.625278in}}%
\pgfpathlineto{\pgfqpoint{0.999179in}{0.633527in}}%
\pgfpathlineto{\pgfqpoint{0.999340in}{0.625278in}}%
\pgfpathlineto{\pgfqpoint{1.000185in}{0.625278in}}%
\pgfpathlineto{\pgfqpoint{1.001358in}{0.633527in}}%
\pgfpathlineto{\pgfqpoint{1.001373in}{0.625278in}}%
\pgfpathlineto{\pgfqpoint{1.002379in}{0.625278in}}%
\pgfpathlineto{\pgfqpoint{1.003531in}{0.633527in}}%
\pgfpathlineto{\pgfqpoint{1.003712in}{0.625278in}}%
\pgfpathlineto{\pgfqpoint{1.004536in}{0.625278in}}%
\pgfpathlineto{\pgfqpoint{1.005710in}{0.633527in}}%
\pgfpathlineto{\pgfqpoint{1.005724in}{0.625278in}}%
\pgfpathlineto{\pgfqpoint{1.006715in}{0.625278in}}%
\pgfpathlineto{\pgfqpoint{1.007889in}{0.633527in}}%
\pgfpathlineto{\pgfqpoint{1.008056in}{0.625278in}}%
\pgfpathlineto{\pgfqpoint{1.008727in}{0.625278in}}%
\pgfpathlineto{\pgfqpoint{1.010068in}{0.633527in}}%
\pgfpathlineto{\pgfqpoint{1.010087in}{0.625278in}}%
\pgfpathlineto{\pgfqpoint{1.011093in}{0.674773in}}%
\pgfpathlineto{\pgfqpoint{1.014278in}{0.823258in}}%
\pgfpathlineto{\pgfqpoint{1.014445in}{0.815008in}}%
\pgfpathlineto{\pgfqpoint{1.014930in}{0.806759in}}%
\pgfpathlineto{\pgfqpoint{1.014948in}{0.848005in}}%
\pgfpathlineto{\pgfqpoint{1.015265in}{0.831507in}}%
\pgfpathlineto{\pgfqpoint{1.015284in}{0.872753in}}%
\pgfpathlineto{\pgfqpoint{1.015600in}{0.864503in}}%
\pgfpathlineto{\pgfqpoint{1.016624in}{0.996490in}}%
\pgfpathlineto{\pgfqpoint{1.020294in}{1.260463in}}%
\pgfpathlineto{\pgfqpoint{1.021474in}{1.326456in}}%
\pgfpathlineto{\pgfqpoint{1.020964in}{1.252214in}}%
\pgfpathlineto{\pgfqpoint{1.021486in}{1.318207in}}%
\pgfpathlineto{\pgfqpoint{1.025162in}{1.582180in}}%
\pgfpathlineto{\pgfqpoint{1.026328in}{1.648173in}}%
\pgfpathlineto{\pgfqpoint{1.025379in}{1.573931in}}%
\pgfpathlineto{\pgfqpoint{1.026348in}{1.639924in}}%
\pgfpathlineto{\pgfqpoint{1.028005in}{1.697668in}}%
\pgfpathlineto{\pgfqpoint{1.028191in}{1.681170in}}%
\pgfpathlineto{\pgfqpoint{1.028340in}{1.672921in}}%
\pgfpathlineto{\pgfqpoint{1.028694in}{1.714167in}}%
\pgfpathlineto{\pgfqpoint{1.033034in}{1.895648in}}%
\pgfpathlineto{\pgfqpoint{1.033053in}{1.887399in}}%
\pgfpathlineto{\pgfqpoint{1.033872in}{2.002887in}}%
\pgfpathlineto{\pgfqpoint{1.034710in}{2.068880in}}%
\pgfpathlineto{\pgfqpoint{1.035064in}{2.060631in}}%
\pgfpathlineto{\pgfqpoint{1.035381in}{2.044133in}}%
\pgfpathlineto{\pgfqpoint{1.035715in}{2.077130in}}%
\pgfpathlineto{\pgfqpoint{1.036070in}{2.159621in}}%
\pgfpathlineto{\pgfqpoint{1.036740in}{2.134874in}}%
\pgfpathlineto{\pgfqpoint{1.039236in}{2.217365in}}%
\pgfpathlineto{\pgfqpoint{1.040428in}{2.159621in}}%
\pgfpathlineto{\pgfqpoint{1.041099in}{2.184369in}}%
\pgfpathlineto{\pgfqpoint{1.041434in}{2.159621in}}%
\pgfpathlineto{\pgfqpoint{1.041751in}{2.143123in}}%
\pgfpathlineto{\pgfqpoint{1.041769in}{2.184369in}}%
\pgfpathlineto{\pgfqpoint{1.042421in}{2.192618in}}%
\pgfpathlineto{\pgfqpoint{1.042607in}{2.176120in}}%
\pgfpathlineto{\pgfqpoint{1.045103in}{2.068880in}}%
\pgfpathlineto{\pgfqpoint{1.045122in}{2.085379in}}%
\pgfpathlineto{\pgfqpoint{1.045439in}{2.093628in}}%
\pgfpathlineto{\pgfqpoint{1.045625in}{2.077130in}}%
\pgfpathlineto{\pgfqpoint{1.046966in}{2.035884in}}%
\pgfpathlineto{\pgfqpoint{1.057340in}{1.474941in}}%
\pgfpathlineto{\pgfqpoint{1.058532in}{1.392449in}}%
\pgfpathlineto{\pgfqpoint{1.060860in}{1.227466in}}%
\pgfpathlineto{\pgfqpoint{1.062053in}{1.186221in}}%
\pgfpathlineto{\pgfqpoint{1.064232in}{1.021237in}}%
\pgfpathlineto{\pgfqpoint{1.064381in}{1.029487in}}%
\pgfpathlineto{\pgfqpoint{1.064735in}{0.996490in}}%
\pgfpathlineto{\pgfqpoint{1.072613in}{0.650025in}}%
\pgfpathlineto{\pgfqpoint{1.073440in}{0.683022in}}%
\pgfpathlineto{\pgfqpoint{1.073775in}{0.658274in}}%
\pgfpathlineto{\pgfqpoint{1.074445in}{0.633527in}}%
\pgfpathlineto{\pgfqpoint{1.074781in}{0.683022in}}%
\pgfpathlineto{\pgfqpoint{1.074793in}{0.674773in}}%
\pgfpathlineto{\pgfqpoint{1.075116in}{0.707769in}}%
\pgfpathlineto{\pgfqpoint{1.076133in}{0.691271in}}%
\pgfpathlineto{\pgfqpoint{1.076301in}{0.674773in}}%
\pgfpathlineto{\pgfqpoint{1.076632in}{0.707769in}}%
\pgfpathlineto{\pgfqpoint{1.076639in}{0.707769in}}%
\pgfpathlineto{\pgfqpoint{1.077295in}{0.757264in}}%
\pgfpathlineto{\pgfqpoint{1.077810in}{0.716019in}}%
\pgfpathlineto{\pgfqpoint{1.078818in}{0.683022in}}%
\pgfpathlineto{\pgfqpoint{1.078145in}{0.724268in}}%
\pgfpathlineto{\pgfqpoint{1.078818in}{0.691271in}}%
\pgfpathlineto{\pgfqpoint{1.078983in}{0.699520in}}%
\pgfpathlineto{\pgfqpoint{1.079319in}{0.666524in}}%
\pgfpathlineto{\pgfqpoint{1.080327in}{0.625278in}}%
\pgfpathlineto{\pgfqpoint{1.081493in}{0.633527in}}%
\pgfpathlineto{\pgfqpoint{1.081654in}{0.625278in}}%
\pgfpathlineto{\pgfqpoint{1.082317in}{0.658274in}}%
\pgfpathlineto{\pgfqpoint{1.082504in}{0.650025in}}%
\pgfpathlineto{\pgfqpoint{1.083009in}{0.625278in}}%
\pgfpathlineto{\pgfqpoint{1.083337in}{0.658274in}}%
\pgfpathlineto{\pgfqpoint{1.083342in}{0.650025in}}%
\pgfpathlineto{\pgfqpoint{1.083833in}{0.658274in}}%
\pgfpathlineto{\pgfqpoint{1.084012in}{0.641776in}}%
\pgfpathlineto{\pgfqpoint{1.084348in}{0.625278in}}%
\pgfpathlineto{\pgfqpoint{1.084518in}{0.658274in}}%
\pgfpathlineto{\pgfqpoint{1.084683in}{0.650025in}}%
\pgfpathlineto{\pgfqpoint{1.085857in}{0.749015in}}%
\pgfpathlineto{\pgfqpoint{1.086527in}{0.699520in}}%
\pgfpathlineto{\pgfqpoint{1.087868in}{0.650025in}}%
\pgfpathlineto{\pgfqpoint{1.088036in}{0.674773in}}%
\pgfpathlineto{\pgfqpoint{1.088538in}{0.699520in}}%
\pgfpathlineto{\pgfqpoint{1.089043in}{0.683022in}}%
\pgfpathlineto{\pgfqpoint{1.090217in}{0.633527in}}%
\pgfpathlineto{\pgfqpoint{1.090552in}{0.658274in}}%
\pgfpathlineto{\pgfqpoint{1.091391in}{0.707769in}}%
\pgfpathlineto{\pgfqpoint{1.091558in}{0.707769in}}%
\pgfpathlineto{\pgfqpoint{1.092550in}{0.732517in}}%
\pgfpathlineto{\pgfqpoint{1.092729in}{0.716019in}}%
\pgfpathlineto{\pgfqpoint{1.093065in}{0.749015in}}%
\pgfpathlineto{\pgfqpoint{1.093738in}{0.683022in}}%
\pgfpathlineto{\pgfqpoint{1.094072in}{0.658274in}}%
\pgfpathlineto{\pgfqpoint{1.094238in}{0.699520in}}%
\pgfpathlineto{\pgfqpoint{1.094568in}{0.691271in}}%
\pgfpathlineto{\pgfqpoint{1.095239in}{0.707769in}}%
\pgfpathlineto{\pgfqpoint{1.095579in}{0.674773in}}%
\pgfpathlineto{\pgfqpoint{1.095581in}{0.691271in}}%
\pgfpathlineto{\pgfqpoint{1.097423in}{0.773763in}}%
\pgfpathlineto{\pgfqpoint{1.097760in}{0.757264in}}%
\pgfpathlineto{\pgfqpoint{1.101113in}{0.633527in}}%
\pgfpathlineto{\pgfqpoint{1.098429in}{0.773763in}}%
\pgfpathlineto{\pgfqpoint{1.101113in}{0.641776in}}%
\pgfpathlineto{\pgfqpoint{1.101278in}{0.674773in}}%
\pgfpathlineto{\pgfqpoint{1.101782in}{0.625278in}}%
\pgfpathlineto{\pgfqpoint{1.102119in}{0.625278in}}%
\pgfpathlineto{\pgfqpoint{1.103122in}{0.650025in}}%
\pgfpathlineto{\pgfqpoint{1.103290in}{0.641776in}}%
\pgfpathlineto{\pgfqpoint{1.104128in}{0.650025in}}%
\pgfpathlineto{\pgfqpoint{1.104631in}{0.625278in}}%
\pgfpathlineto{\pgfqpoint{1.105472in}{0.658274in}}%
\pgfpathlineto{\pgfqpoint{1.105637in}{0.650025in}}%
\pgfpathlineto{\pgfqpoint{1.105975in}{0.658274in}}%
\pgfpathlineto{\pgfqpoint{1.106978in}{0.625278in}}%
\pgfpathlineto{\pgfqpoint{1.108154in}{0.633527in}}%
\pgfpathlineto{\pgfqpoint{1.108321in}{0.625278in}}%
\pgfpathlineto{\pgfqpoint{1.108992in}{0.625278in}}%
\pgfpathlineto{\pgfqpoint{1.110326in}{0.633527in}}%
\pgfpathlineto{\pgfqpoint{1.110333in}{0.625278in}}%
\pgfpathlineto{\pgfqpoint{1.111164in}{0.625278in}}%
\pgfpathlineto{\pgfqpoint{1.112344in}{0.633527in}}%
\pgfpathlineto{\pgfqpoint{1.112680in}{0.625278in}}%
\pgfpathlineto{\pgfqpoint{1.113350in}{0.625278in}}%
\pgfpathlineto{\pgfqpoint{1.114684in}{0.633527in}}%
\pgfpathlineto{\pgfqpoint{1.115194in}{0.625278in}}%
\pgfpathlineto{\pgfqpoint{1.115530in}{0.625278in}}%
\pgfpathlineto{\pgfqpoint{1.116696in}{0.633527in}}%
\pgfpathlineto{\pgfqpoint{1.116863in}{0.625278in}}%
\pgfpathlineto{\pgfqpoint{1.117709in}{0.625278in}}%
\pgfpathlineto{\pgfqpoint{1.118708in}{0.633527in}}%
\pgfpathlineto{\pgfqpoint{1.118708in}{0.625278in}}%
\pgfpathlineto{\pgfqpoint{1.119880in}{0.633527in}}%
\pgfpathlineto{\pgfqpoint{1.120041in}{0.625278in}}%
\pgfpathlineto{\pgfqpoint{1.120893in}{0.625278in}}%
\pgfpathlineto{\pgfqpoint{1.122220in}{0.633527in}}%
\pgfpathlineto{\pgfqpoint{1.122563in}{0.625278in}}%
\pgfpathlineto{\pgfqpoint{1.123058in}{0.625278in}}%
\pgfpathlineto{\pgfqpoint{1.124239in}{0.633527in}}%
\pgfpathlineto{\pgfqpoint{1.124407in}{0.625278in}}%
\pgfpathlineto{\pgfqpoint{1.125245in}{0.625278in}}%
\pgfpathlineto{\pgfqpoint{1.126411in}{0.633527in}}%
\pgfpathlineto{\pgfqpoint{1.126746in}{0.625278in}}%
\pgfpathlineto{\pgfqpoint{1.127424in}{0.625278in}}%
\pgfpathlineto{\pgfqpoint{1.128758in}{0.633527in}}%
\pgfpathlineto{\pgfqpoint{1.128919in}{0.625278in}}%
\pgfpathlineto{\pgfqpoint{1.129771in}{0.625278in}}%
\pgfpathlineto{\pgfqpoint{1.130930in}{0.633527in}}%
\pgfpathlineto{\pgfqpoint{1.131265in}{0.625278in}}%
\pgfpathlineto{\pgfqpoint{1.131943in}{0.625278in}}%
\pgfpathlineto{\pgfqpoint{1.133131in}{0.633527in}}%
\pgfpathlineto{\pgfqpoint{1.133459in}{0.625278in}}%
\pgfpathlineto{\pgfqpoint{1.134122in}{0.625278in}}%
\pgfpathlineto{\pgfqpoint{1.135296in}{0.633527in}}%
\pgfpathlineto{\pgfqpoint{1.135471in}{0.625278in}}%
\pgfpathlineto{\pgfqpoint{1.136309in}{0.625278in}}%
\pgfpathlineto{\pgfqpoint{1.137475in}{0.633527in}}%
\pgfpathlineto{\pgfqpoint{1.137636in}{0.625278in}}%
\pgfpathlineto{\pgfqpoint{1.138306in}{0.625278in}}%
\pgfpathlineto{\pgfqpoint{1.139493in}{0.633527in}}%
\pgfpathlineto{\pgfqpoint{1.139647in}{0.625278in}}%
\pgfpathlineto{\pgfqpoint{1.140485in}{0.625278in}}%
\pgfpathlineto{\pgfqpoint{1.141658in}{0.633527in}}%
\pgfpathlineto{\pgfqpoint{1.141994in}{0.625278in}}%
\pgfpathlineto{\pgfqpoint{1.142671in}{0.625278in}}%
\pgfpathlineto{\pgfqpoint{1.143670in}{0.633527in}}%
\pgfpathlineto{\pgfqpoint{1.143670in}{0.625278in}}%
\pgfpathlineto{\pgfqpoint{1.144844in}{0.633527in}}%
\pgfpathlineto{\pgfqpoint{1.144865in}{0.625278in}}%
\pgfpathlineto{\pgfqpoint{1.145857in}{0.625278in}}%
\pgfpathlineto{\pgfqpoint{1.147023in}{0.633527in}}%
\pgfpathlineto{\pgfqpoint{1.147023in}{0.625278in}}%
\pgfpathlineto{\pgfqpoint{1.148036in}{0.625278in}}%
\pgfpathlineto{\pgfqpoint{1.149202in}{0.633527in}}%
\pgfpathlineto{\pgfqpoint{1.149537in}{0.625278in}}%
\pgfpathlineto{\pgfqpoint{1.150208in}{0.625278in}}%
\pgfpathlineto{\pgfqpoint{1.151381in}{0.633527in}}%
\pgfpathlineto{\pgfqpoint{1.151549in}{0.625278in}}%
\pgfpathlineto{\pgfqpoint{1.152226in}{0.625278in}}%
\pgfpathlineto{\pgfqpoint{1.153392in}{0.633527in}}%
\pgfpathlineto{\pgfqpoint{1.153400in}{0.625278in}}%
\pgfpathlineto{\pgfqpoint{1.154399in}{0.625278in}}%
\pgfpathlineto{\pgfqpoint{1.155418in}{0.633527in}}%
\pgfpathlineto{\pgfqpoint{1.155418in}{0.625278in}}%
\pgfpathlineto{\pgfqpoint{1.156578in}{0.633527in}}%
\pgfpathlineto{\pgfqpoint{1.156913in}{0.625278in}}%
\pgfpathlineto{\pgfqpoint{1.157430in}{0.625278in}}%
\pgfpathlineto{\pgfqpoint{1.158589in}{0.633527in}}%
\pgfpathlineto{\pgfqpoint{1.158764in}{0.625278in}}%
\pgfpathlineto{\pgfqpoint{1.159602in}{0.625278in}}%
\pgfpathlineto{\pgfqpoint{1.160769in}{0.633527in}}%
\pgfpathlineto{\pgfqpoint{1.160943in}{0.625278in}}%
\pgfpathlineto{\pgfqpoint{1.161774in}{0.625278in}}%
\pgfpathlineto{\pgfqpoint{1.162948in}{0.633527in}}%
\pgfpathlineto{\pgfqpoint{1.163115in}{0.625278in}}%
\pgfpathlineto{\pgfqpoint{1.163961in}{0.625278in}}%
\pgfpathlineto{\pgfqpoint{1.165127in}{0.633527in}}%
\pgfpathlineto{\pgfqpoint{1.165295in}{0.625278in}}%
\pgfpathlineto{\pgfqpoint{1.166133in}{0.625278in}}%
\pgfpathlineto{\pgfqpoint{1.167306in}{0.633527in}}%
\pgfpathlineto{\pgfqpoint{1.167321in}{0.625278in}}%
\pgfpathlineto{\pgfqpoint{1.168312in}{0.625278in}}%
\pgfpathlineto{\pgfqpoint{1.169456in}{0.633527in}}%
\pgfpathlineto{\pgfqpoint{1.169492in}{0.625278in}}%
\pgfpathlineto{\pgfqpoint{1.170498in}{0.625278in}}%
\pgfpathlineto{\pgfqpoint{1.171665in}{0.633527in}}%
\pgfpathlineto{\pgfqpoint{1.171839in}{0.625278in}}%
\pgfpathlineto{\pgfqpoint{1.172670in}{0.625278in}}%
\pgfpathlineto{\pgfqpoint{1.173844in}{0.633527in}}%
\pgfpathlineto{\pgfqpoint{1.174011in}{0.625278in}}%
\pgfpathlineto{\pgfqpoint{1.174857in}{0.625278in}}%
\pgfpathlineto{\pgfqpoint{1.175993in}{0.633527in}}%
\pgfpathlineto{\pgfqpoint{1.176023in}{0.625278in}}%
\pgfpathlineto{\pgfqpoint{1.177036in}{0.625278in}}%
\pgfpathlineto{\pgfqpoint{1.178202in}{0.633527in}}%
\pgfpathlineto{\pgfqpoint{1.178370in}{0.625278in}}%
\pgfpathlineto{\pgfqpoint{1.178705in}{0.625278in}}%
\pgfpathlineto{\pgfqpoint{1.180046in}{0.633527in}}%
\pgfpathlineto{\pgfqpoint{1.180382in}{0.625278in}}%
\pgfpathlineto{\pgfqpoint{1.180898in}{0.625278in}}%
\pgfpathlineto{\pgfqpoint{1.182058in}{0.633527in}}%
\pgfpathlineto{\pgfqpoint{1.182393in}{0.625278in}}%
\pgfpathlineto{\pgfqpoint{1.183063in}{0.625278in}}%
\pgfpathlineto{\pgfqpoint{1.184237in}{0.633527in}}%
\pgfpathlineto{\pgfqpoint{1.184368in}{0.625278in}}%
\pgfpathlineto{\pgfqpoint{1.185242in}{0.625278in}}%
\pgfpathlineto{\pgfqpoint{1.186249in}{0.633527in}}%
\pgfpathlineto{\pgfqpoint{1.186249in}{0.625278in}}%
\pgfpathlineto{\pgfqpoint{1.187421in}{0.633527in}}%
\pgfpathlineto{\pgfqpoint{1.187757in}{0.625278in}}%
\pgfpathlineto{\pgfqpoint{1.188428in}{0.625278in}}%
\pgfpathlineto{\pgfqpoint{1.189601in}{0.633527in}}%
\pgfpathlineto{\pgfqpoint{1.189769in}{0.625278in}}%
\pgfpathlineto{\pgfqpoint{1.190607in}{0.625278in}}%
\pgfpathlineto{\pgfqpoint{1.191787in}{0.633527in}}%
\pgfpathlineto{\pgfqpoint{1.191948in}{0.625278in}}%
\pgfpathlineto{\pgfqpoint{1.192618in}{0.625278in}}%
\pgfpathlineto{\pgfqpoint{1.193792in}{0.633527in}}%
\pgfpathlineto{\pgfqpoint{1.193959in}{0.625278in}}%
\pgfpathlineto{\pgfqpoint{1.194797in}{0.625278in}}%
\pgfpathlineto{\pgfqpoint{1.195971in}{0.633527in}}%
\pgfpathlineto{\pgfqpoint{1.195971in}{0.625278in}}%
\pgfpathlineto{\pgfqpoint{1.196977in}{0.625278in}}%
\pgfpathlineto{\pgfqpoint{1.198150in}{0.633527in}}%
\pgfpathlineto{\pgfqpoint{1.198150in}{0.625278in}}%
\pgfpathlineto{\pgfqpoint{1.199156in}{0.625278in}}%
\pgfpathlineto{\pgfqpoint{1.200329in}{0.633527in}}%
\pgfpathlineto{\pgfqpoint{1.200497in}{0.625278in}}%
\pgfpathlineto{\pgfqpoint{1.201335in}{0.625278in}}%
\pgfpathlineto{\pgfqpoint{1.202508in}{0.633527in}}%
\pgfpathlineto{\pgfqpoint{1.202676in}{0.625278in}}%
\pgfpathlineto{\pgfqpoint{1.203514in}{0.625278in}}%
\pgfpathlineto{\pgfqpoint{1.204688in}{0.633527in}}%
\pgfpathlineto{\pgfqpoint{1.204855in}{0.625278in}}%
\pgfpathlineto{\pgfqpoint{1.205694in}{0.625278in}}%
\pgfpathlineto{\pgfqpoint{1.206867in}{0.633527in}}%
\pgfpathlineto{\pgfqpoint{1.207034in}{0.625278in}}%
\pgfpathlineto{\pgfqpoint{1.207873in}{0.625278in}}%
\pgfpathlineto{\pgfqpoint{1.209046in}{0.633527in}}%
\pgfpathlineto{\pgfqpoint{1.209213in}{0.625278in}}%
\pgfpathlineto{\pgfqpoint{1.210052in}{0.625278in}}%
\pgfpathlineto{\pgfqpoint{1.211225in}{0.633527in}}%
\pgfpathlineto{\pgfqpoint{1.211561in}{0.625278in}}%
\pgfpathlineto{\pgfqpoint{1.212231in}{0.625278in}}%
\pgfpathlineto{\pgfqpoint{1.213404in}{0.633527in}}%
\pgfpathlineto{\pgfqpoint{1.213572in}{0.625278in}}%
\pgfpathlineto{\pgfqpoint{1.214410in}{0.625278in}}%
\pgfpathlineto{\pgfqpoint{1.215416in}{0.633527in}}%
\pgfpathlineto{\pgfqpoint{1.215416in}{0.625278in}}%
\pgfpathlineto{\pgfqpoint{1.216757in}{0.633527in}}%
\pgfpathlineto{\pgfqpoint{1.217092in}{0.625278in}}%
\pgfpathlineto{\pgfqpoint{1.217595in}{0.625278in}}%
\pgfpathlineto{\pgfqpoint{1.218769in}{0.633527in}}%
\pgfpathlineto{\pgfqpoint{1.218937in}{0.625278in}}%
\pgfpathlineto{\pgfqpoint{1.219607in}{0.625278in}}%
\pgfpathlineto{\pgfqpoint{1.220948in}{0.633527in}}%
\pgfpathlineto{\pgfqpoint{1.220948in}{0.625278in}}%
\pgfpathlineto{\pgfqpoint{1.221954in}{0.625278in}}%
\pgfpathlineto{\pgfqpoint{1.223127in}{0.633527in}}%
\pgfpathlineto{\pgfqpoint{1.223295in}{0.625278in}}%
\pgfpathlineto{\pgfqpoint{1.223966in}{0.625278in}}%
\pgfpathlineto{\pgfqpoint{1.225138in}{0.633527in}}%
\pgfpathlineto{\pgfqpoint{1.225306in}{0.625278in}}%
\pgfpathlineto{\pgfqpoint{1.225977in}{0.625278in}}%
\pgfpathlineto{\pgfqpoint{1.227318in}{0.633527in}}%
\pgfpathlineto{\pgfqpoint{1.227318in}{0.625278in}}%
\pgfpathlineto{\pgfqpoint{1.228324in}{0.625278in}}%
\pgfpathlineto{\pgfqpoint{1.229497in}{0.633527in}}%
\pgfpathlineto{\pgfqpoint{1.229497in}{0.625278in}}%
\pgfpathlineto{\pgfqpoint{1.230503in}{0.625278in}}%
\pgfpathlineto{\pgfqpoint{1.231844in}{0.633527in}}%
\pgfpathlineto{\pgfqpoint{1.232012in}{0.625278in}}%
\pgfpathlineto{\pgfqpoint{1.232850in}{0.625278in}}%
\pgfpathlineto{\pgfqpoint{1.234023in}{0.633527in}}%
\pgfpathlineto{\pgfqpoint{1.234191in}{0.625278in}}%
\pgfpathlineto{\pgfqpoint{1.234862in}{0.625278in}}%
\pgfpathlineto{\pgfqpoint{1.236034in}{0.633527in}}%
\pgfpathlineto{\pgfqpoint{1.236202in}{0.625278in}}%
\pgfpathlineto{\pgfqpoint{1.237041in}{0.625278in}}%
\pgfpathlineto{\pgfqpoint{1.238214in}{0.633527in}}%
\pgfpathlineto{\pgfqpoint{1.238549in}{0.625278in}}%
\pgfpathlineto{\pgfqpoint{1.239220in}{0.625278in}}%
\pgfpathlineto{\pgfqpoint{1.240393in}{0.633527in}}%
\pgfpathlineto{\pgfqpoint{1.240561in}{0.625278in}}%
\pgfpathlineto{\pgfqpoint{1.241399in}{0.625278in}}%
\pgfpathlineto{\pgfqpoint{1.242405in}{0.633527in}}%
\pgfpathlineto{\pgfqpoint{1.242405in}{0.625278in}}%
\pgfpathlineto{\pgfqpoint{1.243578in}{0.633527in}}%
\pgfpathlineto{\pgfqpoint{1.243578in}{0.625278in}}%
\pgfpathlineto{\pgfqpoint{1.244584in}{0.625278in}}%
\pgfpathlineto{\pgfqpoint{1.245758in}{0.633527in}}%
\pgfpathlineto{\pgfqpoint{1.245925in}{0.625278in}}%
\pgfpathlineto{\pgfqpoint{1.246763in}{0.625278in}}%
\pgfpathlineto{\pgfqpoint{1.247937in}{0.633527in}}%
\pgfpathlineto{\pgfqpoint{1.248104in}{0.625278in}}%
\pgfpathlineto{\pgfqpoint{1.248942in}{0.625278in}}%
\pgfpathlineto{\pgfqpoint{1.249948in}{0.633527in}}%
\pgfpathlineto{\pgfqpoint{1.249948in}{0.625278in}}%
\pgfpathlineto{\pgfqpoint{1.251121in}{0.633527in}}%
\pgfpathlineto{\pgfqpoint{1.251289in}{0.625278in}}%
\pgfpathlineto{\pgfqpoint{1.252127in}{0.625278in}}%
\pgfpathlineto{\pgfqpoint{1.253301in}{0.633527in}}%
\pgfpathlineto{\pgfqpoint{1.253301in}{0.625278in}}%
\pgfpathlineto{\pgfqpoint{1.254306in}{0.625278in}}%
\pgfpathlineto{\pgfqpoint{1.255312in}{0.633527in}}%
\pgfpathlineto{\pgfqpoint{1.255312in}{0.625278in}}%
\pgfpathlineto{\pgfqpoint{1.256654in}{0.633527in}}%
\pgfpathlineto{\pgfqpoint{1.256821in}{0.625278in}}%
\pgfpathlineto{\pgfqpoint{1.257678in}{0.625278in}}%
\pgfpathlineto{\pgfqpoint{1.259019in}{0.674773in}}%
\pgfpathlineto{\pgfqpoint{1.259335in}{0.658274in}}%
\pgfpathlineto{\pgfqpoint{1.259522in}{0.650025in}}%
\pgfpathlineto{\pgfqpoint{1.260174in}{0.707769in}}%
\pgfpathlineto{\pgfqpoint{1.260193in}{0.699520in}}%
\pgfpathlineto{\pgfqpoint{1.260676in}{0.707769in}}%
\pgfpathlineto{\pgfqpoint{1.260863in}{0.691271in}}%
\pgfpathlineto{\pgfqpoint{1.261012in}{0.683022in}}%
\pgfpathlineto{\pgfqpoint{1.261366in}{0.724268in}}%
\pgfpathlineto{\pgfqpoint{1.261683in}{0.732517in}}%
\pgfpathlineto{\pgfqpoint{1.262017in}{0.707769in}}%
\pgfpathlineto{\pgfqpoint{1.263042in}{0.674773in}}%
\pgfpathlineto{\pgfqpoint{1.263526in}{0.707769in}}%
\pgfpathlineto{\pgfqpoint{1.264197in}{0.683022in}}%
\pgfpathlineto{\pgfqpoint{1.264886in}{0.674773in}}%
\pgfpathlineto{\pgfqpoint{1.265203in}{0.707769in}}%
\pgfpathlineto{\pgfqpoint{1.265222in}{0.699520in}}%
\pgfpathlineto{\pgfqpoint{1.265389in}{0.724268in}}%
\pgfpathlineto{\pgfqpoint{1.266376in}{0.707769in}}%
\pgfpathlineto{\pgfqpoint{1.266395in}{0.699520in}}%
\pgfpathlineto{\pgfqpoint{1.266711in}{0.732517in}}%
\pgfpathlineto{\pgfqpoint{1.267401in}{0.724268in}}%
\pgfpathlineto{\pgfqpoint{1.268220in}{0.732517in}}%
\pgfpathlineto{\pgfqpoint{1.268239in}{0.724268in}}%
\pgfpathlineto{\pgfqpoint{1.268742in}{0.699520in}}%
\pgfpathlineto{\pgfqpoint{1.269058in}{0.732517in}}%
\pgfpathlineto{\pgfqpoint{1.269077in}{0.724268in}}%
\pgfpathlineto{\pgfqpoint{1.270734in}{0.757264in}}%
\pgfpathlineto{\pgfqpoint{1.271405in}{0.732517in}}%
\pgfpathlineto{\pgfqpoint{1.271759in}{0.749015in}}%
\pgfpathlineto{\pgfqpoint{1.272579in}{0.782012in}}%
\pgfpathlineto{\pgfqpoint{1.272913in}{0.757264in}}%
\pgfpathlineto{\pgfqpoint{1.273417in}{0.732517in}}%
\pgfpathlineto{\pgfqpoint{1.273938in}{0.749015in}}%
\pgfpathlineto{\pgfqpoint{1.275596in}{0.782012in}}%
\pgfpathlineto{\pgfqpoint{1.276452in}{0.749015in}}%
\pgfpathlineto{\pgfqpoint{1.276621in}{0.773763in}}%
\pgfpathlineto{\pgfqpoint{1.276937in}{0.782012in}}%
\pgfpathlineto{\pgfqpoint{1.277272in}{0.757264in}}%
\pgfpathlineto{\pgfqpoint{1.277942in}{0.732517in}}%
\pgfpathlineto{\pgfqpoint{1.277961in}{0.773763in}}%
\pgfpathlineto{\pgfqpoint{1.278297in}{0.773763in}}%
\pgfpathlineto{\pgfqpoint{1.279638in}{0.724268in}}%
\pgfpathlineto{\pgfqpoint{1.279805in}{0.749015in}}%
\pgfpathlineto{\pgfqpoint{1.280289in}{0.782012in}}%
\pgfpathlineto{\pgfqpoint{1.280643in}{0.749015in}}%
\pgfpathlineto{\pgfqpoint{1.281147in}{0.724268in}}%
\pgfpathlineto{\pgfqpoint{1.281463in}{0.757264in}}%
\pgfpathlineto{\pgfqpoint{1.281481in}{0.749015in}}%
\pgfpathlineto{\pgfqpoint{1.281649in}{0.773763in}}%
\pgfpathlineto{\pgfqpoint{1.282636in}{0.757264in}}%
\pgfpathlineto{\pgfqpoint{1.283661in}{0.724268in}}%
\pgfpathlineto{\pgfqpoint{1.284815in}{0.732517in}}%
\pgfpathlineto{\pgfqpoint{1.285840in}{0.699520in}}%
\pgfpathlineto{\pgfqpoint{1.286343in}{0.724268in}}%
\pgfpathlineto{\pgfqpoint{1.286995in}{0.707769in}}%
\pgfpathlineto{\pgfqpoint{1.288187in}{0.666524in}}%
\pgfpathlineto{\pgfqpoint{1.288671in}{0.658274in}}%
\pgfpathlineto{\pgfqpoint{1.288690in}{0.699520in}}%
\pgfpathlineto{\pgfqpoint{1.289696in}{0.674773in}}%
\pgfpathlineto{\pgfqpoint{1.289844in}{0.691271in}}%
\pgfpathlineto{\pgfqpoint{1.289863in}{0.699520in}}%
\pgfpathlineto{\pgfqpoint{1.290198in}{0.674773in}}%
\pgfpathlineto{\pgfqpoint{1.290869in}{0.674773in}}%
\pgfpathlineto{\pgfqpoint{1.292694in}{0.732517in}}%
\pgfpathlineto{\pgfqpoint{1.292713in}{0.724268in}}%
\pgfpathlineto{\pgfqpoint{1.294054in}{0.699520in}}%
\pgfpathlineto{\pgfqpoint{1.294892in}{0.773763in}}%
\pgfpathlineto{\pgfqpoint{1.296047in}{0.757264in}}%
\pgfpathlineto{\pgfqpoint{1.296885in}{0.707769in}}%
\pgfpathlineto{\pgfqpoint{1.297072in}{0.724268in}}%
\pgfpathlineto{\pgfqpoint{1.299064in}{0.782012in}}%
\pgfpathlineto{\pgfqpoint{1.299586in}{0.773763in}}%
\pgfpathlineto{\pgfqpoint{1.299902in}{0.806759in}}%
\pgfpathlineto{\pgfqpoint{1.300089in}{0.790261in}}%
\pgfpathlineto{\pgfqpoint{1.300592in}{0.773763in}}%
\pgfpathlineto{\pgfqpoint{1.300759in}{0.798510in}}%
\pgfpathlineto{\pgfqpoint{1.301262in}{0.823258in}}%
\pgfpathlineto{\pgfqpoint{1.301765in}{0.798510in}}%
\pgfpathlineto{\pgfqpoint{1.302752in}{0.782012in}}%
\pgfpathlineto{\pgfqpoint{1.302771in}{0.798510in}}%
\pgfpathlineto{\pgfqpoint{1.303925in}{0.806759in}}%
\pgfpathlineto{\pgfqpoint{1.304447in}{0.798510in}}%
\pgfpathlineto{\pgfqpoint{1.304782in}{0.848005in}}%
\pgfpathlineto{\pgfqpoint{1.305937in}{0.856254in}}%
\pgfpathlineto{\pgfqpoint{1.305956in}{0.848005in}}%
\pgfpathlineto{\pgfqpoint{1.306272in}{0.881002in}}%
\pgfpathlineto{\pgfqpoint{1.306961in}{0.872753in}}%
\pgfpathlineto{\pgfqpoint{1.307129in}{0.897500in}}%
\pgfpathlineto{\pgfqpoint{1.308116in}{0.881002in}}%
\pgfpathlineto{\pgfqpoint{1.309140in}{0.872753in}}%
\pgfpathlineto{\pgfqpoint{1.310147in}{0.897500in}}%
\pgfpathlineto{\pgfqpoint{1.310314in}{0.889251in}}%
\pgfpathlineto{\pgfqpoint{1.310482in}{0.897500in}}%
\pgfpathlineto{\pgfqpoint{1.311468in}{0.881002in}}%
\pgfpathlineto{\pgfqpoint{1.311488in}{0.897500in}}%
\pgfpathlineto{\pgfqpoint{1.312158in}{0.848005in}}%
\pgfpathlineto{\pgfqpoint{1.312493in}{0.848005in}}%
\pgfpathlineto{\pgfqpoint{1.313499in}{0.872753in}}%
\pgfpathlineto{\pgfqpoint{1.313667in}{0.864503in}}%
\pgfpathlineto{\pgfqpoint{1.314151in}{0.856254in}}%
\pgfpathlineto{\pgfqpoint{1.314318in}{0.881002in}}%
\pgfpathlineto{\pgfqpoint{1.314337in}{0.872753in}}%
\pgfpathlineto{\pgfqpoint{1.315995in}{0.905749in}}%
\pgfpathlineto{\pgfqpoint{1.317019in}{0.872753in}}%
\pgfpathlineto{\pgfqpoint{1.318677in}{0.930497in}}%
\pgfpathlineto{\pgfqpoint{1.318864in}{0.913998in}}%
\pgfpathlineto{\pgfqpoint{1.319198in}{0.897500in}}%
\pgfpathlineto{\pgfqpoint{1.319515in}{0.938746in}}%
\pgfpathlineto{\pgfqpoint{1.320018in}{0.955244in}}%
\pgfpathlineto{\pgfqpoint{1.320521in}{0.905749in}}%
\pgfpathlineto{\pgfqpoint{1.321043in}{0.922247in}}%
\pgfpathlineto{\pgfqpoint{1.321545in}{0.897500in}}%
\pgfpathlineto{\pgfqpoint{1.322719in}{0.922247in}}%
\pgfpathlineto{\pgfqpoint{1.322886in}{0.913998in}}%
\pgfpathlineto{\pgfqpoint{1.323035in}{0.905749in}}%
\pgfpathlineto{\pgfqpoint{1.323203in}{0.930497in}}%
\pgfpathlineto{\pgfqpoint{1.323557in}{0.922247in}}%
\pgfpathlineto{\pgfqpoint{1.324711in}{0.930497in}}%
\pgfpathlineto{\pgfqpoint{1.325401in}{0.897500in}}%
\pgfpathlineto{\pgfqpoint{1.325736in}{0.922247in}}%
\pgfpathlineto{\pgfqpoint{1.326723in}{0.930497in}}%
\pgfpathlineto{\pgfqpoint{1.326742in}{0.922247in}}%
\pgfpathlineto{\pgfqpoint{1.327059in}{0.905749in}}%
\pgfpathlineto{\pgfqpoint{1.327226in}{0.930497in}}%
\pgfpathlineto{\pgfqpoint{1.327580in}{0.922247in}}%
\pgfpathlineto{\pgfqpoint{1.328902in}{0.955244in}}%
\pgfpathlineto{\pgfqpoint{1.328922in}{0.946995in}}%
\pgfpathlineto{\pgfqpoint{1.329238in}{0.930497in}}%
\pgfpathlineto{\pgfqpoint{1.329405in}{0.955244in}}%
\pgfpathlineto{\pgfqpoint{1.329760in}{0.946995in}}%
\pgfpathlineto{\pgfqpoint{1.330598in}{0.971742in}}%
\pgfpathlineto{\pgfqpoint{1.330914in}{0.955244in}}%
\pgfpathlineto{\pgfqpoint{1.331939in}{0.946995in}}%
\pgfpathlineto{\pgfqpoint{1.331976in}{0.955244in}}%
\pgfpathlineto{\pgfqpoint{1.333260in}{0.979992in}}%
\pgfpathlineto{\pgfqpoint{1.333280in}{0.971742in}}%
\pgfpathlineto{\pgfqpoint{1.334937in}{1.004739in}}%
\pgfpathlineto{\pgfqpoint{1.335459in}{0.996490in}}%
\pgfpathlineto{\pgfqpoint{1.335943in}{1.037736in}}%
\pgfpathlineto{\pgfqpoint{1.336297in}{1.045985in}}%
\pgfpathlineto{\pgfqpoint{1.336632in}{1.021237in}}%
\pgfpathlineto{\pgfqpoint{1.336968in}{1.021237in}}%
\pgfpathlineto{\pgfqpoint{1.338289in}{1.054234in}}%
\pgfpathlineto{\pgfqpoint{1.338309in}{1.045985in}}%
\pgfpathlineto{\pgfqpoint{1.338811in}{1.021237in}}%
\pgfpathlineto{\pgfqpoint{1.339127in}{1.054234in}}%
\pgfpathlineto{\pgfqpoint{1.339147in}{1.045985in}}%
\pgfpathlineto{\pgfqpoint{1.340152in}{1.095480in}}%
\pgfpathlineto{\pgfqpoint{1.340804in}{1.078981in}}%
\pgfpathlineto{\pgfqpoint{1.341326in}{1.070732in}}%
\pgfpathlineto{\pgfqpoint{1.341661in}{1.120227in}}%
\pgfpathlineto{\pgfqpoint{1.341977in}{1.128476in}}%
\pgfpathlineto{\pgfqpoint{1.342313in}{1.103729in}}%
\pgfpathlineto{\pgfqpoint{1.342648in}{1.078981in}}%
\pgfpathlineto{\pgfqpoint{1.343170in}{1.120227in}}%
\pgfpathlineto{\pgfqpoint{1.343337in}{1.111978in}}%
\pgfpathlineto{\pgfqpoint{1.343505in}{1.120227in}}%
\pgfpathlineto{\pgfqpoint{1.344511in}{1.095480in}}%
\pgfpathlineto{\pgfqpoint{1.346168in}{1.128476in}}%
\pgfpathlineto{\pgfqpoint{1.346690in}{1.120227in}}%
\pgfpathlineto{\pgfqpoint{1.347006in}{1.153224in}}%
\pgfpathlineto{\pgfqpoint{1.347193in}{1.144975in}}%
\pgfpathlineto{\pgfqpoint{1.348180in}{1.153224in}}%
\pgfpathlineto{\pgfqpoint{1.348199in}{1.144975in}}%
\pgfpathlineto{\pgfqpoint{1.348702in}{1.120227in}}%
\pgfpathlineto{\pgfqpoint{1.349018in}{1.153224in}}%
\pgfpathlineto{\pgfqpoint{1.349037in}{1.144975in}}%
\pgfpathlineto{\pgfqpoint{1.350192in}{1.153224in}}%
\pgfpathlineto{\pgfqpoint{1.350862in}{1.128476in}}%
\pgfpathlineto{\pgfqpoint{1.351216in}{1.144975in}}%
\pgfpathlineto{\pgfqpoint{1.352203in}{1.153224in}}%
\pgfpathlineto{\pgfqpoint{1.352222in}{1.144975in}}%
\pgfpathlineto{\pgfqpoint{1.354066in}{1.095480in}}%
\pgfpathlineto{\pgfqpoint{1.355556in}{1.128476in}}%
\pgfpathlineto{\pgfqpoint{1.355574in}{1.120227in}}%
\pgfpathlineto{\pgfqpoint{1.356077in}{1.095480in}}%
\pgfpathlineto{\pgfqpoint{1.356394in}{1.128476in}}%
\pgfpathlineto{\pgfqpoint{1.356413in}{1.120227in}}%
\pgfpathlineto{\pgfqpoint{1.356729in}{1.128476in}}%
\pgfpathlineto{\pgfqpoint{1.357064in}{1.103729in}}%
\pgfpathlineto{\pgfqpoint{1.357400in}{1.078981in}}%
\pgfpathlineto{\pgfqpoint{1.358070in}{1.128476in}}%
\pgfpathlineto{\pgfqpoint{1.358089in}{1.120227in}}%
\pgfpathlineto{\pgfqpoint{1.358405in}{1.128476in}}%
\pgfpathlineto{\pgfqpoint{1.358740in}{1.103729in}}%
\pgfpathlineto{\pgfqpoint{1.359095in}{1.095480in}}%
\pgfpathlineto{\pgfqpoint{1.359747in}{1.136726in}}%
\pgfpathlineto{\pgfqpoint{1.359765in}{1.144975in}}%
\pgfpathlineto{\pgfqpoint{1.360101in}{1.120227in}}%
\pgfpathlineto{\pgfqpoint{1.360771in}{1.120227in}}%
\pgfpathlineto{\pgfqpoint{1.362261in}{1.177971in}}%
\pgfpathlineto{\pgfqpoint{1.362448in}{1.161473in}}%
\pgfpathlineto{\pgfqpoint{1.363937in}{1.103729in}}%
\pgfpathlineto{\pgfqpoint{1.364440in}{1.153224in}}%
\pgfpathlineto{\pgfqpoint{1.364962in}{1.136726in}}%
\pgfpathlineto{\pgfqpoint{1.365110in}{1.128476in}}%
\pgfpathlineto{\pgfqpoint{1.365278in}{1.153224in}}%
\pgfpathlineto{\pgfqpoint{1.365632in}{1.144975in}}%
\pgfpathlineto{\pgfqpoint{1.365781in}{1.153224in}}%
\pgfpathlineto{\pgfqpoint{1.366135in}{1.120227in}}%
\pgfpathlineto{\pgfqpoint{1.366284in}{1.128476in}}%
\pgfpathlineto{\pgfqpoint{1.366638in}{1.120227in}}%
\pgfpathlineto{\pgfqpoint{1.366955in}{1.153224in}}%
\pgfpathlineto{\pgfqpoint{1.367309in}{1.144975in}}%
\pgfpathlineto{\pgfqpoint{1.368649in}{1.120227in}}%
\pgfpathlineto{\pgfqpoint{1.369301in}{1.128476in}}%
\pgfpathlineto{\pgfqpoint{1.369488in}{1.111978in}}%
\pgfpathlineto{\pgfqpoint{1.370158in}{1.095480in}}%
\pgfpathlineto{\pgfqpoint{1.370326in}{1.120227in}}%
\pgfpathlineto{\pgfqpoint{1.370810in}{1.128476in}}%
\pgfpathlineto{\pgfqpoint{1.370997in}{1.111978in}}%
\pgfpathlineto{\pgfqpoint{1.371481in}{1.078981in}}%
\pgfpathlineto{\pgfqpoint{1.372170in}{1.095480in}}%
\pgfpathlineto{\pgfqpoint{1.373157in}{1.128476in}}%
\pgfpathlineto{\pgfqpoint{1.373344in}{1.111978in}}%
\pgfpathlineto{\pgfqpoint{1.374182in}{1.095480in}}%
\pgfpathlineto{\pgfqpoint{1.373511in}{1.120227in}}%
\pgfpathlineto{\pgfqpoint{1.374330in}{1.111978in}}%
\pgfpathlineto{\pgfqpoint{1.374665in}{1.128476in}}%
\pgfpathlineto{\pgfqpoint{1.375020in}{1.095480in}}%
\pgfpathlineto{\pgfqpoint{1.375355in}{1.095480in}}%
\pgfpathlineto{\pgfqpoint{1.377031in}{1.144975in}}%
\pgfpathlineto{\pgfqpoint{1.377199in}{1.136726in}}%
\pgfpathlineto{\pgfqpoint{1.378018in}{1.128476in}}%
\pgfpathlineto{\pgfqpoint{1.377366in}{1.144975in}}%
\pgfpathlineto{\pgfqpoint{1.378037in}{1.144975in}}%
\pgfpathlineto{\pgfqpoint{1.378689in}{1.177971in}}%
\pgfpathlineto{\pgfqpoint{1.379192in}{1.153224in}}%
\pgfpathlineto{\pgfqpoint{1.379862in}{1.128476in}}%
\pgfpathlineto{\pgfqpoint{1.380216in}{1.144975in}}%
\pgfpathlineto{\pgfqpoint{1.380365in}{1.153224in}}%
\pgfpathlineto{\pgfqpoint{1.380701in}{1.128476in}}%
\pgfpathlineto{\pgfqpoint{1.380887in}{1.136726in}}%
\pgfpathlineto{\pgfqpoint{1.381035in}{1.128476in}}%
\pgfpathlineto{\pgfqpoint{1.381203in}{1.153224in}}%
\pgfpathlineto{\pgfqpoint{1.381557in}{1.144975in}}%
\pgfpathlineto{\pgfqpoint{1.381706in}{1.153224in}}%
\pgfpathlineto{\pgfqpoint{1.382041in}{1.128476in}}%
\pgfpathlineto{\pgfqpoint{1.382228in}{1.136726in}}%
\pgfpathlineto{\pgfqpoint{1.382377in}{1.128476in}}%
\pgfpathlineto{\pgfqpoint{1.382544in}{1.153224in}}%
\pgfpathlineto{\pgfqpoint{1.382898in}{1.144975in}}%
\pgfpathlineto{\pgfqpoint{1.385729in}{1.301709in}}%
\pgfpathlineto{\pgfqpoint{1.386083in}{1.268712in}}%
\pgfpathlineto{\pgfqpoint{1.387070in}{1.227466in}}%
\pgfpathlineto{\pgfqpoint{1.387406in}{1.252214in}}%
\pgfpathlineto{\pgfqpoint{1.388244in}{1.301709in}}%
\pgfpathlineto{\pgfqpoint{1.388431in}{1.285210in}}%
\pgfpathlineto{\pgfqpoint{1.389269in}{1.268712in}}%
\pgfpathlineto{\pgfqpoint{1.388598in}{1.293460in}}%
\pgfpathlineto{\pgfqpoint{1.389417in}{1.285210in}}%
\pgfpathlineto{\pgfqpoint{1.390598in}{1.351204in}}%
\pgfpathlineto{\pgfqpoint{1.390610in}{1.342955in}}%
\pgfpathlineto{\pgfqpoint{1.390758in}{1.351204in}}%
\pgfpathlineto{\pgfqpoint{1.391112in}{1.318207in}}%
\pgfpathlineto{\pgfqpoint{1.391268in}{1.326456in}}%
\pgfpathlineto{\pgfqpoint{1.391950in}{1.318207in}}%
\pgfpathlineto{\pgfqpoint{1.392274in}{1.351204in}}%
\pgfpathlineto{\pgfqpoint{1.392286in}{1.342955in}}%
\pgfpathlineto{\pgfqpoint{1.393440in}{1.375951in}}%
\pgfpathlineto{\pgfqpoint{1.393627in}{1.359453in}}%
\pgfpathlineto{\pgfqpoint{1.393776in}{1.351204in}}%
\pgfpathlineto{\pgfqpoint{1.393943in}{1.375951in}}%
\pgfpathlineto{\pgfqpoint{1.394298in}{1.367702in}}%
\pgfpathlineto{\pgfqpoint{1.395117in}{1.375951in}}%
\pgfpathlineto{\pgfqpoint{1.395136in}{1.367702in}}%
\pgfpathlineto{\pgfqpoint{1.396625in}{1.301709in}}%
\pgfpathlineto{\pgfqpoint{1.396979in}{1.318207in}}%
\pgfpathlineto{\pgfqpoint{1.397966in}{1.375951in}}%
\pgfpathlineto{\pgfqpoint{1.398320in}{1.342955in}}%
\pgfpathlineto{\pgfqpoint{1.400667in}{1.268712in}}%
\pgfpathlineto{\pgfqpoint{1.401486in}{1.301709in}}%
\pgfpathlineto{\pgfqpoint{1.401822in}{1.276961in}}%
\pgfpathlineto{\pgfqpoint{1.402344in}{1.243965in}}%
\pgfpathlineto{\pgfqpoint{1.402846in}{1.268712in}}%
\pgfpathlineto{\pgfqpoint{1.402995in}{1.276961in}}%
\pgfpathlineto{\pgfqpoint{1.403349in}{1.243965in}}%
\pgfpathlineto{\pgfqpoint{1.403665in}{1.227466in}}%
\pgfpathlineto{\pgfqpoint{1.404001in}{1.260463in}}%
\pgfpathlineto{\pgfqpoint{1.404336in}{1.276961in}}%
\pgfpathlineto{\pgfqpoint{1.404672in}{1.252214in}}%
\pgfpathlineto{\pgfqpoint{1.405026in}{1.268712in}}%
\pgfpathlineto{\pgfqpoint{1.406702in}{1.318207in}}%
\pgfpathlineto{\pgfqpoint{1.406870in}{1.309958in}}%
\pgfpathlineto{\pgfqpoint{1.408546in}{1.268712in}}%
\pgfpathlineto{\pgfqpoint{1.408862in}{1.276961in}}%
\pgfpathlineto{\pgfqpoint{1.409204in}{1.252214in}}%
\pgfpathlineto{\pgfqpoint{1.409719in}{1.243965in}}%
\pgfpathlineto{\pgfqpoint{1.409887in}{1.268712in}}%
\pgfpathlineto{\pgfqpoint{1.410223in}{1.268712in}}%
\pgfpathlineto{\pgfqpoint{1.411899in}{1.219217in}}%
\pgfpathlineto{\pgfqpoint{1.411049in}{1.276961in}}%
\pgfpathlineto{\pgfqpoint{1.412047in}{1.227466in}}%
\pgfpathlineto{\pgfqpoint{1.412215in}{1.252214in}}%
\pgfpathlineto{\pgfqpoint{1.412904in}{1.219217in}}%
\pgfpathlineto{\pgfqpoint{1.413072in}{1.219217in}}%
\pgfpathlineto{\pgfqpoint{1.413389in}{1.252214in}}%
\pgfpathlineto{\pgfqpoint{1.413910in}{1.210968in}}%
\pgfpathlineto{\pgfqpoint{1.415735in}{1.153224in}}%
\pgfpathlineto{\pgfqpoint{1.415754in}{1.169722in}}%
\pgfpathlineto{\pgfqpoint{1.416070in}{1.153224in}}%
\pgfpathlineto{\pgfqpoint{1.416090in}{1.194470in}}%
\pgfpathlineto{\pgfqpoint{1.416406in}{1.186221in}}%
\pgfpathlineto{\pgfqpoint{1.416424in}{1.194470in}}%
\pgfpathlineto{\pgfqpoint{1.416760in}{1.169722in}}%
\pgfpathlineto{\pgfqpoint{1.417431in}{1.169722in}}%
\pgfpathlineto{\pgfqpoint{1.417579in}{1.177971in}}%
\pgfpathlineto{\pgfqpoint{1.417933in}{1.144975in}}%
\pgfpathlineto{\pgfqpoint{1.418089in}{1.153224in}}%
\pgfpathlineto{\pgfqpoint{1.418269in}{1.144975in}}%
\pgfpathlineto{\pgfqpoint{1.418771in}{1.194470in}}%
\pgfpathlineto{\pgfqpoint{1.419107in}{1.194470in}}%
\pgfpathlineto{\pgfqpoint{1.419590in}{1.227466in}}%
\pgfpathlineto{\pgfqpoint{1.420448in}{1.210968in}}%
\pgfpathlineto{\pgfqpoint{1.421957in}{1.144975in}}%
\pgfpathlineto{\pgfqpoint{1.422273in}{1.177971in}}%
\pgfpathlineto{\pgfqpoint{1.422627in}{1.136726in}}%
\pgfpathlineto{\pgfqpoint{1.423298in}{1.120227in}}%
\pgfpathlineto{\pgfqpoint{1.423465in}{1.144975in}}%
\pgfpathlineto{\pgfqpoint{1.423789in}{1.128476in}}%
\pgfpathlineto{\pgfqpoint{1.423800in}{1.144975in}}%
\pgfpathlineto{\pgfqpoint{1.424459in}{1.103729in}}%
\pgfpathlineto{\pgfqpoint{1.424806in}{1.144975in}}%
\pgfpathlineto{\pgfqpoint{1.425465in}{1.128476in}}%
\pgfpathlineto{\pgfqpoint{1.425625in}{1.153224in}}%
\pgfpathlineto{\pgfqpoint{1.425644in}{1.144975in}}%
\pgfpathlineto{\pgfqpoint{1.426631in}{1.202719in}}%
\pgfpathlineto{\pgfqpoint{1.427488in}{1.186221in}}%
\pgfpathlineto{\pgfqpoint{1.430654in}{1.054234in}}%
\pgfpathlineto{\pgfqpoint{1.430673in}{1.070732in}}%
\pgfpathlineto{\pgfqpoint{1.431550in}{1.054234in}}%
\pgfpathlineto{\pgfqpoint{1.431660in}{1.062483in}}%
\pgfpathlineto{\pgfqpoint{1.431884in}{1.103729in}}%
\pgfpathlineto{\pgfqpoint{1.432685in}{1.070732in}}%
\pgfpathlineto{\pgfqpoint{1.433187in}{1.045985in}}%
\pgfpathlineto{\pgfqpoint{1.433504in}{1.078981in}}%
\pgfpathlineto{\pgfqpoint{1.433523in}{1.070732in}}%
\pgfpathlineto{\pgfqpoint{1.434182in}{1.078981in}}%
\pgfpathlineto{\pgfqpoint{1.434361in}{1.062483in}}%
\pgfpathlineto{\pgfqpoint{1.436373in}{0.996490in}}%
\pgfpathlineto{\pgfqpoint{1.434529in}{1.070732in}}%
\pgfpathlineto{\pgfqpoint{1.436528in}{1.004739in}}%
\pgfpathlineto{\pgfqpoint{1.436689in}{1.029487in}}%
\pgfpathlineto{\pgfqpoint{1.437546in}{0.988241in}}%
\pgfpathlineto{\pgfqpoint{1.440061in}{0.897500in}}%
\pgfpathlineto{\pgfqpoint{1.440377in}{0.905749in}}%
\pgfpathlineto{\pgfqpoint{1.440712in}{0.881002in}}%
\pgfpathlineto{\pgfqpoint{1.440899in}{0.889251in}}%
\pgfpathlineto{\pgfqpoint{1.441893in}{0.881002in}}%
\pgfpathlineto{\pgfqpoint{1.441229in}{0.905749in}}%
\pgfpathlineto{\pgfqpoint{1.441900in}{0.889251in}}%
\pgfpathlineto{\pgfqpoint{1.441904in}{0.897500in}}%
\pgfpathlineto{\pgfqpoint{1.442563in}{0.856254in}}%
\pgfpathlineto{\pgfqpoint{1.442911in}{0.872753in}}%
\pgfpathlineto{\pgfqpoint{1.444400in}{0.930497in}}%
\pgfpathlineto{\pgfqpoint{1.444587in}{0.913998in}}%
\pgfpathlineto{\pgfqpoint{1.448107in}{0.773763in}}%
\pgfpathlineto{\pgfqpoint{1.448262in}{0.782012in}}%
\pgfpathlineto{\pgfqpoint{1.448444in}{0.757264in}}%
\pgfpathlineto{\pgfqpoint{1.448780in}{0.757264in}}%
\pgfpathlineto{\pgfqpoint{1.450274in}{0.732517in}}%
\pgfpathlineto{\pgfqpoint{1.450777in}{0.782012in}}%
\pgfpathlineto{\pgfqpoint{1.451292in}{0.773763in}}%
\pgfpathlineto{\pgfqpoint{1.453124in}{0.683022in}}%
\pgfpathlineto{\pgfqpoint{1.453639in}{0.699520in}}%
\pgfpathlineto{\pgfqpoint{1.456644in}{0.806759in}}%
\pgfpathlineto{\pgfqpoint{1.456824in}{0.823258in}}%
\pgfpathlineto{\pgfqpoint{1.457829in}{0.749015in}}%
\pgfpathlineto{\pgfqpoint{1.458656in}{0.782012in}}%
\pgfpathlineto{\pgfqpoint{1.458320in}{0.732517in}}%
\pgfpathlineto{\pgfqpoint{1.459170in}{0.765513in}}%
\pgfpathlineto{\pgfqpoint{1.462022in}{0.658274in}}%
\pgfpathlineto{\pgfqpoint{1.462169in}{0.666524in}}%
\pgfpathlineto{\pgfqpoint{1.462858in}{0.674773in}}%
\pgfpathlineto{\pgfqpoint{1.462693in}{0.633527in}}%
\pgfpathlineto{\pgfqpoint{1.463194in}{0.650025in}}%
\pgfpathlineto{\pgfqpoint{1.464535in}{0.625278in}}%
\pgfpathlineto{\pgfqpoint{1.465864in}{0.658274in}}%
\pgfpathlineto{\pgfqpoint{1.465875in}{0.650025in}}%
\pgfpathlineto{\pgfqpoint{1.466717in}{0.633527in}}%
\pgfpathlineto{\pgfqpoint{1.466870in}{0.641776in}}%
\pgfpathlineto{\pgfqpoint{1.467887in}{0.724268in}}%
\pgfpathlineto{\pgfqpoint{1.468558in}{0.798510in}}%
\pgfpathlineto{\pgfqpoint{1.469228in}{0.749015in}}%
\pgfpathlineto{\pgfqpoint{1.472419in}{0.625278in}}%
\pgfpathlineto{\pgfqpoint{1.473582in}{0.633527in}}%
\pgfpathlineto{\pgfqpoint{1.474093in}{0.625278in}}%
\pgfpathlineto{\pgfqpoint{1.474256in}{0.625278in}}%
\pgfpathlineto{\pgfqpoint{1.475588in}{0.633527in}}%
\pgfpathlineto{\pgfqpoint{1.475596in}{0.625278in}}%
\pgfpathlineto{\pgfqpoint{1.476595in}{0.625278in}}%
\pgfpathlineto{\pgfqpoint{1.477772in}{0.633527in}}%
\pgfpathlineto{\pgfqpoint{1.477935in}{0.625278in}}%
\pgfpathlineto{\pgfqpoint{1.478772in}{0.625278in}}%
\pgfpathlineto{\pgfqpoint{1.479942in}{0.633527in}}%
\pgfpathlineto{\pgfqpoint{1.479949in}{0.625278in}}%
\pgfpathlineto{\pgfqpoint{1.480949in}{0.625278in}}%
\pgfpathlineto{\pgfqpoint{1.481963in}{0.633527in}}%
\pgfpathlineto{\pgfqpoint{1.481963in}{0.625278in}}%
\pgfpathlineto{\pgfqpoint{1.483289in}{0.633527in}}%
\pgfpathlineto{\pgfqpoint{1.483303in}{0.625278in}}%
\pgfpathlineto{\pgfqpoint{1.484303in}{0.625278in}}%
\pgfpathlineto{\pgfqpoint{1.485480in}{0.633527in}}%
\pgfpathlineto{\pgfqpoint{1.485813in}{0.625278in}}%
\pgfpathlineto{\pgfqpoint{1.486147in}{0.625278in}}%
\pgfpathlineto{\pgfqpoint{1.487317in}{0.633527in}}%
\pgfpathlineto{\pgfqpoint{1.487650in}{0.625278in}}%
\pgfpathlineto{\pgfqpoint{1.488331in}{0.625278in}}%
\pgfpathlineto{\pgfqpoint{1.489508in}{0.633527in}}%
\pgfpathlineto{\pgfqpoint{1.489671in}{0.625278in}}%
\pgfpathlineto{\pgfqpoint{1.490515in}{0.625278in}}%
\pgfpathlineto{\pgfqpoint{1.491515in}{0.633527in}}%
\pgfpathlineto{\pgfqpoint{1.491515in}{0.625278in}}%
\pgfpathlineto{\pgfqpoint{1.492855in}{0.633527in}}%
\pgfpathlineto{\pgfqpoint{1.493514in}{0.625278in}}%
\pgfpathlineto{\pgfqpoint{1.494854in}{0.633527in}}%
\pgfpathlineto{\pgfqpoint{1.495024in}{0.625278in}}%
\pgfpathlineto{\pgfqpoint{1.495861in}{0.625278in}}%
\pgfpathlineto{\pgfqpoint{1.497038in}{0.633527in}}%
\pgfpathlineto{\pgfqpoint{1.497201in}{0.625278in}}%
\pgfpathlineto{\pgfqpoint{1.498045in}{0.625278in}}%
\pgfpathlineto{\pgfqpoint{1.499215in}{0.633527in}}%
\pgfpathlineto{\pgfqpoint{1.499385in}{0.625278in}}%
\pgfpathlineto{\pgfqpoint{1.500222in}{0.625278in}}%
\pgfpathlineto{\pgfqpoint{1.501562in}{0.633527in}}%
\pgfpathlineto{\pgfqpoint{1.501562in}{0.625278in}}%
\pgfpathlineto{\pgfqpoint{1.502399in}{0.625278in}}%
\pgfpathlineto{\pgfqpoint{1.503576in}{0.633527in}}%
\pgfpathlineto{\pgfqpoint{1.503739in}{0.625278in}}%
\pgfpathlineto{\pgfqpoint{1.504583in}{0.625278in}}%
\pgfpathlineto{\pgfqpoint{1.505583in}{0.633527in}}%
\pgfpathlineto{\pgfqpoint{1.505583in}{0.625278in}}%
\pgfpathlineto{\pgfqpoint{1.506760in}{0.633527in}}%
\pgfpathlineto{\pgfqpoint{1.506923in}{0.625278in}}%
\pgfpathlineto{\pgfqpoint{1.507767in}{0.625278in}}%
\pgfpathlineto{\pgfqpoint{1.508937in}{0.633527in}}%
\pgfpathlineto{\pgfqpoint{1.509107in}{0.625278in}}%
\pgfpathlineto{\pgfqpoint{1.509944in}{0.625278in}}%
\pgfpathlineto{\pgfqpoint{1.511114in}{0.633527in}}%
\pgfpathlineto{\pgfqpoint{1.511284in}{0.625278in}}%
\pgfpathlineto{\pgfqpoint{1.512121in}{0.625278in}}%
\pgfpathlineto{\pgfqpoint{1.513128in}{0.633527in}}%
\pgfpathlineto{\pgfqpoint{1.513128in}{0.625278in}}%
\pgfpathlineto{\pgfqpoint{1.514468in}{0.633527in}}%
\pgfpathlineto{\pgfqpoint{1.514468in}{0.625278in}}%
\pgfpathlineto{\pgfqpoint{1.515475in}{0.625278in}}%
\pgfpathlineto{\pgfqpoint{1.516652in}{0.633527in}}%
\pgfpathlineto{\pgfqpoint{1.516985in}{0.625278in}}%
\pgfpathlineto{\pgfqpoint{1.517651in}{0.625278in}}%
\pgfpathlineto{\pgfqpoint{1.518829in}{0.633527in}}%
\pgfpathlineto{\pgfqpoint{1.518999in}{0.625278in}}%
\pgfpathlineto{\pgfqpoint{1.519836in}{0.625278in}}%
\pgfpathlineto{\pgfqpoint{1.521006in}{0.633527in}}%
\pgfpathlineto{\pgfqpoint{1.521176in}{0.625278in}}%
\pgfpathlineto{\pgfqpoint{1.522013in}{0.625278in}}%
\pgfpathlineto{\pgfqpoint{1.523353in}{0.633527in}}%
\pgfpathlineto{\pgfqpoint{1.523523in}{0.625278in}}%
\pgfpathlineto{\pgfqpoint{1.524360in}{0.625278in}}%
\pgfpathlineto{\pgfqpoint{1.525537in}{0.633527in}}%
\pgfpathlineto{\pgfqpoint{1.525700in}{0.625278in}}%
\pgfpathlineto{\pgfqpoint{1.526537in}{0.625278in}}%
\pgfpathlineto{\pgfqpoint{1.527714in}{0.633527in}}%
\pgfpathlineto{\pgfqpoint{1.527714in}{0.625278in}}%
\pgfpathlineto{\pgfqpoint{1.528550in}{0.625278in}}%
\pgfpathlineto{\pgfqpoint{1.529891in}{0.633527in}}%
\pgfpathlineto{\pgfqpoint{1.530061in}{0.625278in}}%
\pgfpathlineto{\pgfqpoint{1.530727in}{0.625278in}}%
\pgfpathlineto{\pgfqpoint{1.531905in}{0.633527in}}%
\pgfpathlineto{\pgfqpoint{1.532067in}{0.625278in}}%
\pgfpathlineto{\pgfqpoint{1.532912in}{0.625278in}}%
\pgfpathlineto{\pgfqpoint{1.534081in}{0.633527in}}%
\pgfpathlineto{\pgfqpoint{1.534252in}{0.625278in}}%
\pgfpathlineto{\pgfqpoint{1.535088in}{0.625278in}}%
\pgfpathlineto{\pgfqpoint{1.536258in}{0.633527in}}%
\pgfpathlineto{\pgfqpoint{1.536429in}{0.625278in}}%
\pgfpathlineto{\pgfqpoint{1.537265in}{0.625278in}}%
\pgfpathlineto{\pgfqpoint{1.538442in}{0.633527in}}%
\pgfpathlineto{\pgfqpoint{1.538605in}{0.625278in}}%
\pgfpathlineto{\pgfqpoint{1.539449in}{0.625278in}}%
\pgfpathlineto{\pgfqpoint{1.540619in}{0.633527in}}%
\pgfpathlineto{\pgfqpoint{1.540790in}{0.625278in}}%
\pgfpathlineto{\pgfqpoint{1.541626in}{0.625278in}}%
\pgfpathlineto{\pgfqpoint{1.542796in}{0.633527in}}%
\pgfpathlineto{\pgfqpoint{1.542966in}{0.625278in}}%
\pgfpathlineto{\pgfqpoint{1.543803in}{0.625278in}}%
\pgfpathlineto{\pgfqpoint{1.544980in}{0.633527in}}%
\pgfpathlineto{\pgfqpoint{1.545143in}{0.625278in}}%
\pgfpathlineto{\pgfqpoint{1.545987in}{0.625278in}}%
\pgfpathlineto{\pgfqpoint{1.547157in}{0.633527in}}%
\pgfpathlineto{\pgfqpoint{1.547490in}{0.625278in}}%
\pgfpathlineto{\pgfqpoint{1.548164in}{0.625278in}}%
\pgfpathlineto{\pgfqpoint{1.549334in}{0.633527in}}%
\pgfpathlineto{\pgfqpoint{1.549504in}{0.625278in}}%
\pgfpathlineto{\pgfqpoint{1.550178in}{0.625278in}}%
\pgfpathlineto{\pgfqpoint{1.551518in}{0.633527in}}%
\pgfpathlineto{\pgfqpoint{1.551851in}{0.625278in}}%
\pgfpathlineto{\pgfqpoint{1.552525in}{0.625278in}}%
\pgfpathlineto{\pgfqpoint{1.553525in}{0.633527in}}%
\pgfpathlineto{\pgfqpoint{1.553525in}{0.625278in}}%
\pgfpathlineto{\pgfqpoint{1.554702in}{0.633527in}}%
\pgfpathlineto{\pgfqpoint{1.554872in}{0.625278in}}%
\pgfpathlineto{\pgfqpoint{1.555709in}{0.625278in}}%
\pgfpathlineto{\pgfqpoint{1.556879in}{0.633527in}}%
\pgfpathlineto{\pgfqpoint{1.557049in}{0.625278in}}%
\pgfpathlineto{\pgfqpoint{1.557886in}{0.625278in}}%
\pgfpathlineto{\pgfqpoint{1.559063in}{0.633527in}}%
\pgfpathlineto{\pgfqpoint{1.559063in}{0.625278in}}%
\pgfpathlineto{\pgfqpoint{1.560063in}{0.625278in}}%
\pgfpathlineto{\pgfqpoint{1.561240in}{0.633527in}}%
\pgfpathlineto{\pgfqpoint{1.561403in}{0.625278in}}%
\pgfpathlineto{\pgfqpoint{1.562077in}{0.625278in}}%
\pgfpathlineto{\pgfqpoint{1.563254in}{0.633527in}}%
\pgfpathlineto{\pgfqpoint{1.563417in}{0.625278in}}%
\pgfpathlineto{\pgfqpoint{1.564253in}{0.625278in}}%
\pgfpathlineto{\pgfqpoint{1.565260in}{0.633527in}}%
\pgfpathlineto{\pgfqpoint{1.565260in}{0.625278in}}%
\pgfpathlineto{\pgfqpoint{1.566438in}{0.633527in}}%
\pgfpathlineto{\pgfqpoint{1.566438in}{0.625278in}}%
\pgfpathlineto{\pgfqpoint{1.567445in}{0.625278in}}%
\pgfpathlineto{\pgfqpoint{1.568615in}{0.633527in}}%
\pgfpathlineto{\pgfqpoint{1.568785in}{0.625278in}}%
\pgfpathlineto{\pgfqpoint{1.569622in}{0.625278in}}%
\pgfpathlineto{\pgfqpoint{1.570791in}{0.633527in}}%
\pgfpathlineto{\pgfqpoint{1.570962in}{0.625278in}}%
\pgfpathlineto{\pgfqpoint{1.571635in}{0.625278in}}%
\pgfpathlineto{\pgfqpoint{1.572805in}{0.633527in}}%
\pgfpathlineto{\pgfqpoint{1.572976in}{0.625278in}}%
\pgfpathlineto{\pgfqpoint{1.573812in}{0.625278in}}%
\pgfpathlineto{\pgfqpoint{1.574819in}{0.633527in}}%
\pgfpathlineto{\pgfqpoint{1.574819in}{0.625278in}}%
\pgfpathlineto{\pgfqpoint{1.575989in}{0.633527in}}%
\pgfpathlineto{\pgfqpoint{1.576159in}{0.625278in}}%
\pgfpathlineto{\pgfqpoint{1.576996in}{0.625278in}}%
\pgfpathlineto{\pgfqpoint{1.578166in}{0.633527in}}%
\pgfpathlineto{\pgfqpoint{1.578336in}{0.625278in}}%
\pgfpathlineto{\pgfqpoint{1.579173in}{0.625278in}}%
\pgfpathlineto{\pgfqpoint{1.580350in}{0.633527in}}%
\pgfpathlineto{\pgfqpoint{1.580513in}{0.625278in}}%
\pgfpathlineto{\pgfqpoint{1.581357in}{0.625278in}}%
\pgfpathlineto{\pgfqpoint{1.582527in}{0.633527in}}%
\pgfpathlineto{\pgfqpoint{1.582697in}{0.625278in}}%
\pgfpathlineto{\pgfqpoint{1.583364in}{0.625278in}}%
\pgfpathlineto{\pgfqpoint{1.584541in}{0.633527in}}%
\pgfpathlineto{\pgfqpoint{1.584704in}{0.625278in}}%
\pgfpathlineto{\pgfqpoint{1.585548in}{0.625278in}}%
\pgfpathlineto{\pgfqpoint{1.586718in}{0.633527in}}%
\pgfpathlineto{\pgfqpoint{1.586888in}{0.625278in}}%
\pgfpathlineto{\pgfqpoint{1.587725in}{0.625278in}}%
\pgfpathlineto{\pgfqpoint{1.588895in}{0.633527in}}%
\pgfpathlineto{\pgfqpoint{1.589065in}{0.625278in}}%
\pgfpathlineto{\pgfqpoint{1.589902in}{0.625278in}}%
\pgfpathlineto{\pgfqpoint{1.591079in}{0.633527in}}%
\pgfpathlineto{\pgfqpoint{1.591242in}{0.625278in}}%
\pgfpathlineto{\pgfqpoint{1.591916in}{0.625278in}}%
\pgfpathlineto{\pgfqpoint{1.593085in}{0.633527in}}%
\pgfpathlineto{\pgfqpoint{1.593256in}{0.625278in}}%
\pgfpathlineto{\pgfqpoint{1.594092in}{0.625278in}}%
\pgfpathlineto{\pgfqpoint{1.595270in}{0.633527in}}%
\pgfpathlineto{\pgfqpoint{1.595433in}{0.625278in}}%
\pgfpathlineto{\pgfqpoint{1.596277in}{0.625278in}}%
\pgfpathlineto{\pgfqpoint{1.597447in}{0.633527in}}%
\pgfpathlineto{\pgfqpoint{1.597617in}{0.625278in}}%
\pgfpathlineto{\pgfqpoint{1.598453in}{0.625278in}}%
\pgfpathlineto{\pgfqpoint{1.599623in}{0.633527in}}%
\pgfpathlineto{\pgfqpoint{1.599623in}{0.625278in}}%
\pgfpathlineto{\pgfqpoint{1.600630in}{0.625278in}}%
\pgfpathlineto{\pgfqpoint{1.601637in}{0.633527in}}%
\pgfpathlineto{\pgfqpoint{1.601637in}{0.625278in}}%
\pgfpathlineto{\pgfqpoint{1.602815in}{0.633527in}}%
\pgfpathlineto{\pgfqpoint{1.602977in}{0.625278in}}%
\pgfpathlineto{\pgfqpoint{1.603814in}{0.625278in}}%
\pgfpathlineto{\pgfqpoint{1.604991in}{0.633527in}}%
\pgfpathlineto{\pgfqpoint{1.605162in}{0.625278in}}%
\pgfpathlineto{\pgfqpoint{1.605998in}{0.625278in}}%
\pgfpathlineto{\pgfqpoint{1.607168in}{0.633527in}}%
\pgfpathlineto{\pgfqpoint{1.607168in}{0.625278in}}%
\pgfpathlineto{\pgfqpoint{1.608175in}{0.625278in}}%
\pgfpathlineto{\pgfqpoint{1.609352in}{0.633527in}}%
\pgfpathlineto{\pgfqpoint{1.610374in}{0.625278in}}%
\pgfpathlineto{\pgfqpoint{1.611529in}{0.633527in}}%
\pgfpathlineto{\pgfqpoint{1.612048in}{0.625278in}}%
\pgfpathlineto{\pgfqpoint{1.612218in}{0.650025in}}%
\pgfpathlineto{\pgfqpoint{1.612551in}{0.650025in}}%
\pgfpathlineto{\pgfqpoint{1.613055in}{0.625278in}}%
\pgfpathlineto{\pgfqpoint{1.613373in}{0.658274in}}%
\pgfpathlineto{\pgfqpoint{1.613388in}{0.650025in}}%
\pgfpathlineto{\pgfqpoint{1.615046in}{0.683022in}}%
\pgfpathlineto{\pgfqpoint{1.616238in}{0.650025in}}%
\pgfpathlineto{\pgfqpoint{1.617897in}{0.683022in}}%
\pgfpathlineto{\pgfqpoint{1.618585in}{0.650025in}}%
\pgfpathlineto{\pgfqpoint{1.618926in}{0.674773in}}%
\pgfpathlineto{\pgfqpoint{1.619237in}{0.683022in}}%
\pgfpathlineto{\pgfqpoint{1.619592in}{0.650025in}}%
\pgfpathlineto{\pgfqpoint{1.619741in}{0.658274in}}%
\pgfpathlineto{\pgfqpoint{1.620770in}{0.650025in}}%
\pgfpathlineto{\pgfqpoint{1.621606in}{0.625278in}}%
\pgfpathlineto{\pgfqpoint{1.621769in}{0.650025in}}%
\pgfpathlineto{\pgfqpoint{1.623095in}{0.683022in}}%
\pgfpathlineto{\pgfqpoint{1.623117in}{0.674773in}}%
\pgfpathlineto{\pgfqpoint{1.624287in}{0.650025in}}%
\pgfpathlineto{\pgfqpoint{1.624435in}{0.658274in}}%
\pgfpathlineto{\pgfqpoint{1.624605in}{0.683022in}}%
\pgfpathlineto{\pgfqpoint{1.625442in}{0.633527in}}%
\pgfpathlineto{\pgfqpoint{1.625464in}{0.650025in}}%
\pgfpathlineto{\pgfqpoint{1.626116in}{0.633527in}}%
\pgfpathlineto{\pgfqpoint{1.626278in}{0.658274in}}%
\pgfpathlineto{\pgfqpoint{1.626301in}{0.650025in}}%
\pgfpathlineto{\pgfqpoint{1.627789in}{0.707769in}}%
\pgfpathlineto{\pgfqpoint{1.627974in}{0.691271in}}%
\pgfpathlineto{\pgfqpoint{1.628122in}{0.683022in}}%
\pgfpathlineto{\pgfqpoint{1.628292in}{0.707769in}}%
\pgfpathlineto{\pgfqpoint{1.628648in}{0.699520in}}%
\pgfpathlineto{\pgfqpoint{1.628796in}{0.707769in}}%
\pgfpathlineto{\pgfqpoint{1.629129in}{0.683022in}}%
\pgfpathlineto{\pgfqpoint{1.629462in}{0.683022in}}%
\pgfpathlineto{\pgfqpoint{1.629818in}{0.674773in}}%
\pgfpathlineto{\pgfqpoint{1.630136in}{0.707769in}}%
\pgfpathlineto{\pgfqpoint{1.630491in}{0.699520in}}%
\pgfpathlineto{\pgfqpoint{1.631306in}{0.732517in}}%
\pgfpathlineto{\pgfqpoint{1.631646in}{0.707769in}}%
\pgfpathlineto{\pgfqpoint{1.632483in}{0.683022in}}%
\pgfpathlineto{\pgfqpoint{1.632668in}{0.691271in}}%
\pgfpathlineto{\pgfqpoint{1.632816in}{0.683022in}}%
\pgfpathlineto{\pgfqpoint{1.633150in}{0.716019in}}%
\pgfpathlineto{\pgfqpoint{1.633320in}{0.757264in}}%
\pgfpathlineto{\pgfqpoint{1.634179in}{0.749015in}}%
\pgfpathlineto{\pgfqpoint{1.635163in}{0.757264in}}%
\pgfpathlineto{\pgfqpoint{1.635186in}{0.749015in}}%
\pgfpathlineto{\pgfqpoint{1.635689in}{0.724268in}}%
\pgfpathlineto{\pgfqpoint{1.636000in}{0.757264in}}%
\pgfpathlineto{\pgfqpoint{1.636022in}{0.749015in}}%
\pgfpathlineto{\pgfqpoint{1.636341in}{0.757264in}}%
\pgfpathlineto{\pgfqpoint{1.636674in}{0.732517in}}%
\pgfpathlineto{\pgfqpoint{1.636859in}{0.740766in}}%
\pgfpathlineto{\pgfqpoint{1.637007in}{0.732517in}}%
\pgfpathlineto{\pgfqpoint{1.637029in}{0.773763in}}%
\pgfpathlineto{\pgfqpoint{1.637363in}{0.773763in}}%
\pgfpathlineto{\pgfqpoint{1.638014in}{0.732517in}}%
\pgfpathlineto{\pgfqpoint{1.638347in}{0.765513in}}%
\pgfpathlineto{\pgfqpoint{1.638518in}{0.806759in}}%
\pgfpathlineto{\pgfqpoint{1.639354in}{0.757264in}}%
\pgfpathlineto{\pgfqpoint{1.639376in}{0.773763in}}%
\pgfpathlineto{\pgfqpoint{1.639576in}{0.757264in}}%
\pgfpathlineto{\pgfqpoint{1.639710in}{0.798510in}}%
\pgfpathlineto{\pgfqpoint{1.640028in}{0.790261in}}%
\pgfpathlineto{\pgfqpoint{1.640376in}{0.798510in}}%
\pgfpathlineto{\pgfqpoint{1.641028in}{0.757264in}}%
\pgfpathlineto{\pgfqpoint{1.641050in}{0.773763in}}%
\pgfpathlineto{\pgfqpoint{1.642538in}{0.806759in}}%
\pgfpathlineto{\pgfqpoint{1.642560in}{0.798510in}}%
\pgfpathlineto{\pgfqpoint{1.643900in}{0.773763in}}%
\pgfpathlineto{\pgfqpoint{1.645218in}{0.831507in}}%
\pgfpathlineto{\pgfqpoint{1.645744in}{0.815008in}}%
\pgfpathlineto{\pgfqpoint{1.647255in}{0.749015in}}%
\pgfpathlineto{\pgfqpoint{1.645914in}{0.823258in}}%
\pgfpathlineto{\pgfqpoint{1.647417in}{0.773763in}}%
\pgfpathlineto{\pgfqpoint{1.647906in}{0.782012in}}%
\pgfpathlineto{\pgfqpoint{1.648091in}{0.765513in}}%
\pgfpathlineto{\pgfqpoint{1.648239in}{0.757264in}}%
\pgfpathlineto{\pgfqpoint{1.648595in}{0.798510in}}%
\pgfpathlineto{\pgfqpoint{1.649913in}{0.831507in}}%
\pgfpathlineto{\pgfqpoint{1.649246in}{0.782012in}}%
\pgfpathlineto{\pgfqpoint{1.649935in}{0.823258in}}%
\pgfpathlineto{\pgfqpoint{1.651275in}{0.798510in}}%
\pgfpathlineto{\pgfqpoint{1.652430in}{0.806759in}}%
\pgfpathlineto{\pgfqpoint{1.653622in}{0.773763in}}%
\pgfpathlineto{\pgfqpoint{1.653941in}{0.757264in}}%
\pgfpathlineto{\pgfqpoint{1.654274in}{0.790261in}}%
\pgfpathlineto{\pgfqpoint{1.655110in}{0.831507in}}%
\pgfpathlineto{\pgfqpoint{1.655295in}{0.815008in}}%
\pgfpathlineto{\pgfqpoint{1.656954in}{0.782012in}}%
\pgfpathlineto{\pgfqpoint{1.657124in}{0.806759in}}%
\pgfpathlineto{\pgfqpoint{1.657983in}{0.798510in}}%
\pgfpathlineto{\pgfqpoint{1.658635in}{0.831507in}}%
\pgfpathlineto{\pgfqpoint{1.658990in}{0.798510in}}%
\pgfpathlineto{\pgfqpoint{1.660330in}{0.773763in}}%
\pgfpathlineto{\pgfqpoint{1.662322in}{0.831507in}}%
\pgfpathlineto{\pgfqpoint{1.663011in}{0.798510in}}%
\pgfpathlineto{\pgfqpoint{1.663344in}{0.815008in}}%
\pgfpathlineto{\pgfqpoint{1.664832in}{0.782012in}}%
\pgfpathlineto{\pgfqpoint{1.663514in}{0.823258in}}%
\pgfpathlineto{\pgfqpoint{1.664832in}{0.790261in}}%
\pgfpathlineto{\pgfqpoint{1.665676in}{0.856254in}}%
\pgfpathlineto{\pgfqpoint{1.666024in}{0.823258in}}%
\pgfpathlineto{\pgfqpoint{1.666676in}{0.856254in}}%
\pgfpathlineto{\pgfqpoint{1.667031in}{0.823258in}}%
\pgfpathlineto{\pgfqpoint{1.667683in}{0.806759in}}%
\pgfpathlineto{\pgfqpoint{1.667683in}{0.839756in}}%
\pgfpathlineto{\pgfqpoint{1.668690in}{0.856254in}}%
\pgfpathlineto{\pgfqpoint{1.668023in}{0.831507in}}%
\pgfpathlineto{\pgfqpoint{1.668712in}{0.848005in}}%
\pgfpathlineto{\pgfqpoint{1.670052in}{0.897500in}}%
\pgfpathlineto{\pgfqpoint{1.670363in}{0.881002in}}%
\pgfpathlineto{\pgfqpoint{1.671562in}{0.848005in}}%
\pgfpathlineto{\pgfqpoint{1.673717in}{0.930497in}}%
\pgfpathlineto{\pgfqpoint{1.673739in}{0.922247in}}%
\pgfpathlineto{\pgfqpoint{1.674058in}{0.905749in}}%
\pgfpathlineto{\pgfqpoint{1.674391in}{0.938746in}}%
\pgfpathlineto{\pgfqpoint{1.674724in}{0.955244in}}%
\pgfpathlineto{\pgfqpoint{1.675079in}{0.922247in}}%
\pgfpathlineto{\pgfqpoint{1.675413in}{0.922247in}}%
\pgfpathlineto{\pgfqpoint{1.675561in}{0.930497in}}%
\pgfpathlineto{\pgfqpoint{1.675916in}{0.897500in}}%
\pgfpathlineto{\pgfqpoint{1.676064in}{0.905749in}}%
\pgfpathlineto{\pgfqpoint{1.677242in}{0.881002in}}%
\pgfpathlineto{\pgfqpoint{1.677404in}{0.905749in}}%
\pgfpathlineto{\pgfqpoint{1.678263in}{0.897500in}}%
\pgfpathlineto{\pgfqpoint{1.678915in}{0.930497in}}%
\pgfpathlineto{\pgfqpoint{1.679270in}{0.897500in}}%
\pgfpathlineto{\pgfqpoint{1.680596in}{0.881002in}}%
\pgfpathlineto{\pgfqpoint{1.680610in}{0.897500in}}%
\pgfpathlineto{\pgfqpoint{1.681617in}{0.897500in}}%
\pgfpathlineto{\pgfqpoint{1.682935in}{0.955244in}}%
\pgfpathlineto{\pgfqpoint{1.683276in}{0.930497in}}%
\pgfpathlineto{\pgfqpoint{1.683461in}{0.922247in}}%
\pgfpathlineto{\pgfqpoint{1.683779in}{0.955244in}}%
\pgfpathlineto{\pgfqpoint{1.684298in}{0.946995in}}%
\pgfpathlineto{\pgfqpoint{1.684786in}{0.955244in}}%
\pgfpathlineto{\pgfqpoint{1.684971in}{0.938746in}}%
\pgfpathlineto{\pgfqpoint{1.685120in}{0.930497in}}%
\pgfpathlineto{\pgfqpoint{1.685475in}{0.971742in}}%
\pgfpathlineto{\pgfqpoint{1.687970in}{1.054234in}}%
\pgfpathlineto{\pgfqpoint{1.688977in}{1.029487in}}%
\pgfpathlineto{\pgfqpoint{1.688992in}{1.045985in}}%
\pgfpathlineto{\pgfqpoint{1.689310in}{1.029487in}}%
\pgfpathlineto{\pgfqpoint{1.689644in}{1.062483in}}%
\pgfpathlineto{\pgfqpoint{1.690651in}{1.078981in}}%
\pgfpathlineto{\pgfqpoint{1.690480in}{1.054234in}}%
\pgfpathlineto{\pgfqpoint{1.690665in}{1.070732in}}%
\pgfpathlineto{\pgfqpoint{1.691006in}{1.095480in}}%
\pgfpathlineto{\pgfqpoint{1.691509in}{1.062483in}}%
\pgfpathlineto{\pgfqpoint{1.692013in}{1.045985in}}%
\pgfpathlineto{\pgfqpoint{1.692324in}{1.078981in}}%
\pgfpathlineto{\pgfqpoint{1.692346in}{1.070732in}}%
\pgfpathlineto{\pgfqpoint{1.693353in}{1.095480in}}%
\pgfpathlineto{\pgfqpoint{1.693516in}{1.087231in}}%
\pgfpathlineto{\pgfqpoint{1.694167in}{1.054234in}}%
\pgfpathlineto{\pgfqpoint{1.694508in}{1.087231in}}%
\pgfpathlineto{\pgfqpoint{1.695530in}{1.120227in}}%
\pgfpathlineto{\pgfqpoint{1.696181in}{1.103729in}}%
\pgfpathlineto{\pgfqpoint{1.696181in}{1.136726in}}%
\pgfpathlineto{\pgfqpoint{1.696204in}{1.144975in}}%
\pgfpathlineto{\pgfqpoint{1.697018in}{1.103729in}}%
\pgfpathlineto{\pgfqpoint{1.697203in}{1.111978in}}%
\pgfpathlineto{\pgfqpoint{1.697374in}{1.120227in}}%
\pgfpathlineto{\pgfqpoint{1.698380in}{1.095480in}}%
\pgfpathlineto{\pgfqpoint{1.700039in}{1.128476in}}%
\pgfpathlineto{\pgfqpoint{1.700705in}{1.103729in}}%
\pgfpathlineto{\pgfqpoint{1.701061in}{1.120227in}}%
\pgfpathlineto{\pgfqpoint{1.702216in}{1.128476in}}%
\pgfpathlineto{\pgfqpoint{1.703393in}{1.078981in}}%
\pgfpathlineto{\pgfqpoint{1.703408in}{1.095480in}}%
\pgfpathlineto{\pgfqpoint{1.704733in}{1.128476in}}%
\pgfpathlineto{\pgfqpoint{1.704748in}{1.120227in}}%
\pgfpathlineto{\pgfqpoint{1.706088in}{1.095480in}}%
\pgfpathlineto{\pgfqpoint{1.706762in}{1.120227in}}%
\pgfpathlineto{\pgfqpoint{1.707243in}{1.103729in}}%
\pgfpathlineto{\pgfqpoint{1.707266in}{1.095480in}}%
\pgfpathlineto{\pgfqpoint{1.707769in}{1.144975in}}%
\pgfpathlineto{\pgfqpoint{1.708272in}{1.111978in}}%
\pgfpathlineto{\pgfqpoint{1.709109in}{1.095480in}}%
\pgfpathlineto{\pgfqpoint{1.708435in}{1.120227in}}%
\pgfpathlineto{\pgfqpoint{1.709257in}{1.111978in}}%
\pgfpathlineto{\pgfqpoint{1.710264in}{1.128476in}}%
\pgfpathlineto{\pgfqpoint{1.710279in}{1.120227in}}%
\pgfpathlineto{\pgfqpoint{1.711604in}{1.078981in}}%
\pgfpathlineto{\pgfqpoint{1.711619in}{1.095480in}}%
\pgfpathlineto{\pgfqpoint{1.713618in}{1.153224in}}%
\pgfpathlineto{\pgfqpoint{1.713633in}{1.144975in}}%
\pgfpathlineto{\pgfqpoint{1.714285in}{1.103729in}}%
\pgfpathlineto{\pgfqpoint{1.714640in}{1.144975in}}%
\pgfpathlineto{\pgfqpoint{1.714958in}{1.153224in}}%
\pgfpathlineto{\pgfqpoint{1.715292in}{1.128476in}}%
\pgfpathlineto{\pgfqpoint{1.715477in}{1.136726in}}%
\pgfpathlineto{\pgfqpoint{1.716313in}{1.120227in}}%
\pgfpathlineto{\pgfqpoint{1.715647in}{1.144975in}}%
\pgfpathlineto{\pgfqpoint{1.716462in}{1.136726in}}%
\pgfpathlineto{\pgfqpoint{1.717661in}{1.194470in}}%
\pgfpathlineto{\pgfqpoint{1.718646in}{1.153224in}}%
\pgfpathlineto{\pgfqpoint{1.718979in}{1.177971in}}%
\pgfpathlineto{\pgfqpoint{1.719149in}{1.202719in}}%
\pgfpathlineto{\pgfqpoint{1.720001in}{1.194470in}}%
\pgfpathlineto{\pgfqpoint{1.721659in}{1.153224in}}%
\pgfpathlineto{\pgfqpoint{1.721659in}{1.161473in}}%
\pgfpathlineto{\pgfqpoint{1.722837in}{1.202719in}}%
\pgfpathlineto{\pgfqpoint{1.722000in}{1.153224in}}%
\pgfpathlineto{\pgfqpoint{1.722851in}{1.194470in}}%
\pgfpathlineto{\pgfqpoint{1.723170in}{1.202719in}}%
\pgfpathlineto{\pgfqpoint{1.723503in}{1.177971in}}%
\pgfpathlineto{\pgfqpoint{1.724532in}{1.169722in}}%
\pgfpathlineto{\pgfqpoint{1.726191in}{1.202719in}}%
\pgfpathlineto{\pgfqpoint{1.726879in}{1.169722in}}%
\pgfpathlineto{\pgfqpoint{1.727212in}{1.194470in}}%
\pgfpathlineto{\pgfqpoint{1.727361in}{1.202719in}}%
\pgfpathlineto{\pgfqpoint{1.727716in}{1.169722in}}%
\pgfpathlineto{\pgfqpoint{1.727864in}{1.177971in}}%
\pgfpathlineto{\pgfqpoint{1.728886in}{1.169722in}}%
\pgfpathlineto{\pgfqpoint{1.730063in}{1.219217in}}%
\pgfpathlineto{\pgfqpoint{1.730396in}{1.194470in}}%
\pgfpathlineto{\pgfqpoint{1.731070in}{1.169722in}}%
\pgfpathlineto{\pgfqpoint{1.731559in}{1.186221in}}%
\pgfpathlineto{\pgfqpoint{1.732055in}{1.202719in}}%
\pgfpathlineto{\pgfqpoint{1.732410in}{1.169722in}}%
\pgfpathlineto{\pgfqpoint{1.732573in}{1.169722in}}%
\pgfpathlineto{\pgfqpoint{1.733062in}{1.177971in}}%
\pgfpathlineto{\pgfqpoint{1.733247in}{1.161473in}}%
\pgfpathlineto{\pgfqpoint{1.734402in}{1.128476in}}%
\pgfpathlineto{\pgfqpoint{1.734424in}{1.144975in}}%
\pgfpathlineto{\pgfqpoint{1.735239in}{1.177971in}}%
\pgfpathlineto{\pgfqpoint{1.735572in}{1.153224in}}%
\pgfpathlineto{\pgfqpoint{1.736075in}{1.177971in}}%
\pgfpathlineto{\pgfqpoint{1.736601in}{1.144975in}}%
\pgfpathlineto{\pgfqpoint{1.737253in}{1.128476in}}%
\pgfpathlineto{\pgfqpoint{1.737423in}{1.153224in}}%
\pgfpathlineto{\pgfqpoint{1.737438in}{1.144975in}}%
\pgfpathlineto{\pgfqpoint{1.738615in}{1.169722in}}%
\pgfpathlineto{\pgfqpoint{1.738778in}{1.161473in}}%
\pgfpathlineto{\pgfqpoint{1.738926in}{1.153224in}}%
\pgfpathlineto{\pgfqpoint{1.739096in}{1.177971in}}%
\pgfpathlineto{\pgfqpoint{1.739452in}{1.169722in}}%
\pgfpathlineto{\pgfqpoint{1.739600in}{1.177971in}}%
\pgfpathlineto{\pgfqpoint{1.739955in}{1.144975in}}%
\pgfpathlineto{\pgfqpoint{1.740459in}{1.120227in}}%
\pgfpathlineto{\pgfqpoint{1.740770in}{1.153224in}}%
\pgfpathlineto{\pgfqpoint{1.740792in}{1.144975in}}%
\pgfpathlineto{\pgfqpoint{1.741110in}{1.153224in}}%
\pgfpathlineto{\pgfqpoint{1.741458in}{1.120227in}}%
\pgfpathlineto{\pgfqpoint{1.741606in}{1.128476in}}%
\pgfpathlineto{\pgfqpoint{1.741628in}{1.120227in}}%
\pgfpathlineto{\pgfqpoint{1.742465in}{1.169722in}}%
\pgfpathlineto{\pgfqpoint{1.742954in}{1.177971in}}%
\pgfpathlineto{\pgfqpoint{1.743139in}{1.161473in}}%
\pgfpathlineto{\pgfqpoint{1.744797in}{1.128476in}}%
\pgfpathlineto{\pgfqpoint{1.745819in}{1.144975in}}%
\pgfpathlineto{\pgfqpoint{1.746471in}{1.128476in}}%
\pgfpathlineto{\pgfqpoint{1.746478in}{1.161473in}}%
\pgfpathlineto{\pgfqpoint{1.747159in}{1.194470in}}%
\pgfpathlineto{\pgfqpoint{1.746804in}{1.153224in}}%
\pgfpathlineto{\pgfqpoint{1.747493in}{1.169722in}}%
\pgfpathlineto{\pgfqpoint{1.748485in}{1.153224in}}%
\pgfpathlineto{\pgfqpoint{1.748485in}{1.161473in}}%
\pgfpathlineto{\pgfqpoint{1.749506in}{1.194470in}}%
\pgfpathlineto{\pgfqpoint{1.750328in}{1.202719in}}%
\pgfpathlineto{\pgfqpoint{1.750343in}{1.194470in}}%
\pgfpathlineto{\pgfqpoint{1.750884in}{1.177971in}}%
\pgfpathlineto{\pgfqpoint{1.750995in}{1.210968in}}%
\pgfpathlineto{\pgfqpoint{1.751017in}{1.219217in}}%
\pgfpathlineto{\pgfqpoint{1.751668in}{1.177971in}}%
\pgfpathlineto{\pgfqpoint{1.752024in}{1.194470in}}%
\pgfpathlineto{\pgfqpoint{1.752335in}{1.177971in}}%
\pgfpathlineto{\pgfqpoint{1.752357in}{1.219217in}}%
\pgfpathlineto{\pgfqpoint{1.752675in}{1.210968in}}%
\pgfpathlineto{\pgfqpoint{1.753853in}{1.252214in}}%
\pgfpathlineto{\pgfqpoint{1.755023in}{1.202719in}}%
\pgfpathlineto{\pgfqpoint{1.755023in}{1.210968in}}%
\pgfpathlineto{\pgfqpoint{1.755037in}{1.219217in}}%
\pgfpathlineto{\pgfqpoint{1.755689in}{1.177971in}}%
\pgfpathlineto{\pgfqpoint{1.756044in}{1.194470in}}%
\pgfpathlineto{\pgfqpoint{1.757207in}{1.202719in}}%
\pgfpathlineto{\pgfqpoint{1.757555in}{1.169722in}}%
\pgfpathlineto{\pgfqpoint{1.758206in}{1.252214in}}%
\pgfpathlineto{\pgfqpoint{1.758221in}{1.243965in}}%
\pgfpathlineto{\pgfqpoint{1.759384in}{1.252214in}}%
\pgfpathlineto{\pgfqpoint{1.760383in}{1.227466in}}%
\pgfpathlineto{\pgfqpoint{1.760405in}{1.243965in}}%
\pgfpathlineto{\pgfqpoint{1.762079in}{1.293460in}}%
\pgfpathlineto{\pgfqpoint{1.762249in}{1.285210in}}%
\pgfpathlineto{\pgfqpoint{1.763064in}{1.276961in}}%
\pgfpathlineto{\pgfqpoint{1.762412in}{1.293460in}}%
\pgfpathlineto{\pgfqpoint{1.763071in}{1.285210in}}%
\pgfpathlineto{\pgfqpoint{1.763086in}{1.318207in}}%
\pgfpathlineto{\pgfqpoint{1.764093in}{1.293460in}}%
\pgfpathlineto{\pgfqpoint{1.764915in}{1.351204in}}%
\pgfpathlineto{\pgfqpoint{1.765751in}{1.326456in}}%
\pgfpathlineto{\pgfqpoint{1.766270in}{1.293460in}}%
\pgfpathlineto{\pgfqpoint{1.766773in}{1.318207in}}%
\pgfpathlineto{\pgfqpoint{1.766929in}{1.326456in}}%
\pgfpathlineto{\pgfqpoint{1.767277in}{1.293460in}}%
\pgfpathlineto{\pgfqpoint{1.770942in}{1.153224in}}%
\pgfpathlineto{\pgfqpoint{1.770949in}{1.161473in}}%
\pgfpathlineto{\pgfqpoint{1.771638in}{1.194470in}}%
\pgfpathlineto{\pgfqpoint{1.771282in}{1.153224in}}%
\pgfpathlineto{\pgfqpoint{1.771971in}{1.169722in}}%
\pgfpathlineto{\pgfqpoint{1.772289in}{1.153224in}}%
\pgfpathlineto{\pgfqpoint{1.772630in}{1.186221in}}%
\pgfpathlineto{\pgfqpoint{1.772963in}{1.202719in}}%
\pgfpathlineto{\pgfqpoint{1.773289in}{1.177971in}}%
\pgfpathlineto{\pgfqpoint{1.773644in}{1.194470in}}%
\pgfpathlineto{\pgfqpoint{1.775303in}{1.153224in}}%
\pgfpathlineto{\pgfqpoint{1.775310in}{1.161473in}}%
\pgfpathlineto{\pgfqpoint{1.776487in}{1.202719in}}%
\pgfpathlineto{\pgfqpoint{1.776665in}{1.219217in}}%
\pgfpathlineto{\pgfqpoint{1.777502in}{1.169722in}}%
\pgfpathlineto{\pgfqpoint{1.777665in}{1.177971in}}%
\pgfpathlineto{\pgfqpoint{1.778005in}{1.144975in}}%
\pgfpathlineto{\pgfqpoint{1.778153in}{1.153224in}}%
\pgfpathlineto{\pgfqpoint{1.779345in}{1.120227in}}%
\pgfpathlineto{\pgfqpoint{1.780182in}{1.144975in}}%
\pgfpathlineto{\pgfqpoint{1.780500in}{1.128476in}}%
\pgfpathlineto{\pgfqpoint{1.780686in}{1.120227in}}%
\pgfpathlineto{\pgfqpoint{1.781337in}{1.153224in}}%
\pgfpathlineto{\pgfqpoint{1.781522in}{1.144975in}}%
\pgfpathlineto{\pgfqpoint{1.781670in}{1.153224in}}%
\pgfpathlineto{\pgfqpoint{1.782026in}{1.120227in}}%
\pgfpathlineto{\pgfqpoint{1.785713in}{0.946995in}}%
\pgfpathlineto{\pgfqpoint{1.786054in}{0.971742in}}%
\pgfpathlineto{\pgfqpoint{1.786557in}{0.946995in}}%
\pgfpathlineto{\pgfqpoint{1.787209in}{0.979992in}}%
\pgfpathlineto{\pgfqpoint{1.788401in}{0.922247in}}%
\pgfpathlineto{\pgfqpoint{1.790222in}{0.856254in}}%
\pgfpathlineto{\pgfqpoint{1.789052in}{0.930497in}}%
\pgfpathlineto{\pgfqpoint{1.790244in}{0.872753in}}%
\pgfpathlineto{\pgfqpoint{1.791584in}{0.922247in}}%
\pgfpathlineto{\pgfqpoint{1.791903in}{0.905749in}}%
\pgfpathlineto{\pgfqpoint{1.792421in}{0.897500in}}%
\pgfpathlineto{\pgfqpoint{1.792754in}{0.930497in}}%
\pgfpathlineto{\pgfqpoint{1.792925in}{0.922247in}}%
\pgfpathlineto{\pgfqpoint{1.794413in}{0.955244in}}%
\pgfpathlineto{\pgfqpoint{1.794435in}{0.946995in}}%
\pgfpathlineto{\pgfqpoint{1.795272in}{0.897500in}}%
\pgfpathlineto{\pgfqpoint{1.795753in}{0.930497in}}%
\pgfpathlineto{\pgfqpoint{1.796086in}{0.979992in}}%
\pgfpathlineto{\pgfqpoint{1.796782in}{0.971742in}}%
\pgfpathlineto{\pgfqpoint{1.797286in}{0.946995in}}%
\pgfpathlineto{\pgfqpoint{1.798100in}{0.955244in}}%
\pgfpathlineto{\pgfqpoint{1.798122in}{0.971742in}}%
\pgfpathlineto{\pgfqpoint{1.798959in}{0.889251in}}%
\pgfpathlineto{\pgfqpoint{1.801136in}{0.823258in}}%
\pgfpathlineto{\pgfqpoint{1.801462in}{0.831507in}}%
\pgfpathlineto{\pgfqpoint{1.801787in}{0.806759in}}%
\pgfpathlineto{\pgfqpoint{1.801973in}{0.815008in}}%
\pgfpathlineto{\pgfqpoint{1.802143in}{0.848005in}}%
\pgfpathlineto{\pgfqpoint{1.803128in}{0.806759in}}%
\pgfpathlineto{\pgfqpoint{1.803987in}{0.848005in}}%
\pgfpathlineto{\pgfqpoint{1.804157in}{0.839756in}}%
\pgfpathlineto{\pgfqpoint{1.805149in}{0.831507in}}%
\pgfpathlineto{\pgfqpoint{1.804475in}{0.856254in}}%
\pgfpathlineto{\pgfqpoint{1.805149in}{0.839756in}}%
\pgfpathlineto{\pgfqpoint{1.806163in}{0.905749in}}%
\pgfpathlineto{\pgfqpoint{1.806334in}{0.889251in}}%
\pgfpathlineto{\pgfqpoint{1.806837in}{0.848005in}}%
\pgfpathlineto{\pgfqpoint{1.807333in}{0.889251in}}%
\pgfpathlineto{\pgfqpoint{1.808496in}{0.955244in}}%
\pgfpathlineto{\pgfqpoint{1.808681in}{0.946995in}}%
\pgfpathlineto{\pgfqpoint{1.809014in}{0.979992in}}%
\pgfpathlineto{\pgfqpoint{1.809517in}{0.971742in}}%
\pgfpathlineto{\pgfqpoint{1.811865in}{0.897500in}}%
\pgfpathlineto{\pgfqpoint{1.813523in}{0.930497in}}%
\pgfpathlineto{\pgfqpoint{1.814049in}{0.881002in}}%
\pgfpathlineto{\pgfqpoint{1.814715in}{0.889251in}}%
\pgfpathlineto{\pgfqpoint{1.816559in}{0.823258in}}%
\pgfpathlineto{\pgfqpoint{1.816707in}{0.831507in}}%
\pgfpathlineto{\pgfqpoint{1.817396in}{0.872753in}}%
\pgfpathlineto{\pgfqpoint{1.817736in}{0.848005in}}%
\pgfpathlineto{\pgfqpoint{1.818573in}{0.806759in}}%
\pgfpathlineto{\pgfqpoint{1.819076in}{0.823258in}}%
\pgfpathlineto{\pgfqpoint{1.819743in}{0.831507in}}%
\pgfpathlineto{\pgfqpoint{1.819913in}{0.815008in}}%
\pgfpathlineto{\pgfqpoint{1.820061in}{0.806759in}}%
\pgfpathlineto{\pgfqpoint{1.820416in}{0.848005in}}%
\pgfpathlineto{\pgfqpoint{1.820579in}{0.839756in}}%
\pgfpathlineto{\pgfqpoint{1.821253in}{0.806759in}}%
\pgfpathlineto{\pgfqpoint{1.821594in}{0.831507in}}%
\pgfpathlineto{\pgfqpoint{1.821757in}{0.848005in}}%
\pgfpathlineto{\pgfqpoint{1.822260in}{0.806759in}}%
\pgfpathlineto{\pgfqpoint{1.822586in}{0.831507in}}%
\pgfpathlineto{\pgfqpoint{1.822764in}{0.823258in}}%
\pgfpathlineto{\pgfqpoint{1.823267in}{0.872753in}}%
\pgfpathlineto{\pgfqpoint{1.823600in}{0.848005in}}%
\pgfpathlineto{\pgfqpoint{1.829124in}{0.625278in}}%
\pgfpathlineto{\pgfqpoint{1.830116in}{0.633527in}}%
\pgfpathlineto{\pgfqpoint{1.830116in}{0.625278in}}%
\pgfpathlineto{\pgfqpoint{1.831123in}{0.633527in}}%
\pgfpathlineto{\pgfqpoint{1.831123in}{0.625278in}}%
\pgfpathlineto{\pgfqpoint{1.832300in}{0.633527in}}%
\pgfpathlineto{\pgfqpoint{1.832308in}{0.625278in}}%
\pgfpathlineto{\pgfqpoint{1.833307in}{0.625278in}}%
\pgfpathlineto{\pgfqpoint{1.834484in}{0.633527in}}%
\pgfpathlineto{\pgfqpoint{1.834647in}{0.625278in}}%
\pgfpathlineto{\pgfqpoint{1.835314in}{0.625278in}}%
\pgfpathlineto{\pgfqpoint{1.836654in}{0.633527in}}%
\pgfpathlineto{\pgfqpoint{1.836994in}{0.625278in}}%
\pgfpathlineto{\pgfqpoint{1.837676in}{0.625278in}}%
\pgfpathlineto{\pgfqpoint{1.839023in}{0.633527in}}%
\pgfpathlineto{\pgfqpoint{1.839060in}{0.625278in}}%
\pgfpathlineto{\pgfqpoint{1.840067in}{0.625278in}}%
\pgfpathlineto{\pgfqpoint{1.841348in}{0.633527in}}%
\pgfpathlineto{\pgfqpoint{1.841348in}{0.625278in}}%
\pgfpathlineto{\pgfqpoint{1.842355in}{0.625278in}}%
\pgfpathlineto{\pgfqpoint{1.843532in}{0.633527in}}%
\pgfpathlineto{\pgfqpoint{1.843695in}{0.625278in}}%
\pgfpathlineto{\pgfqpoint{1.844539in}{0.625278in}}%
\pgfpathlineto{\pgfqpoint{1.845709in}{0.633527in}}%
\pgfpathlineto{\pgfqpoint{1.845887in}{0.625278in}}%
\pgfpathlineto{\pgfqpoint{1.846716in}{0.625278in}}%
\pgfpathlineto{\pgfqpoint{1.847886in}{0.633527in}}%
\pgfpathlineto{\pgfqpoint{1.848056in}{0.625278in}}%
\pgfpathlineto{\pgfqpoint{1.848900in}{0.625278in}}%
\pgfpathlineto{\pgfqpoint{1.850070in}{0.633527in}}%
\pgfpathlineto{\pgfqpoint{1.850233in}{0.625278in}}%
\pgfpathlineto{\pgfqpoint{1.851070in}{0.625278in}}%
\pgfpathlineto{\pgfqpoint{1.852417in}{0.633527in}}%
\pgfpathlineto{\pgfqpoint{1.852580in}{0.625278in}}%
\pgfpathlineto{\pgfqpoint{1.853417in}{0.625278in}}%
\pgfpathlineto{\pgfqpoint{1.854757in}{0.633527in}}%
\pgfpathlineto{\pgfqpoint{1.854757in}{0.625278in}}%
\pgfpathlineto{\pgfqpoint{1.855764in}{0.625278in}}%
\pgfpathlineto{\pgfqpoint{1.856771in}{0.633527in}}%
\pgfpathlineto{\pgfqpoint{1.856771in}{0.625278in}}%
\pgfpathlineto{\pgfqpoint{1.857956in}{0.633527in}}%
\pgfpathlineto{\pgfqpoint{1.858111in}{0.625278in}}%
\pgfpathlineto{\pgfqpoint{1.858955in}{0.625278in}}%
\pgfpathlineto{\pgfqpoint{1.860295in}{0.633527in}}%
\pgfpathlineto{\pgfqpoint{1.860303in}{0.625278in}}%
\pgfpathlineto{\pgfqpoint{1.861132in}{0.625278in}}%
\pgfpathlineto{\pgfqpoint{1.862309in}{0.633527in}}%
\pgfpathlineto{\pgfqpoint{1.862309in}{0.625278in}}%
\pgfpathlineto{\pgfqpoint{1.863316in}{0.625278in}}%
\pgfpathlineto{\pgfqpoint{1.864649in}{0.633527in}}%
\pgfpathlineto{\pgfqpoint{1.864819in}{0.625278in}}%
\pgfpathlineto{\pgfqpoint{1.865656in}{0.625278in}}%
\pgfpathlineto{\pgfqpoint{1.866833in}{0.633527in}}%
\pgfpathlineto{\pgfqpoint{1.866833in}{0.625278in}}%
\pgfpathlineto{\pgfqpoint{1.867833in}{0.625278in}}%
\pgfpathlineto{\pgfqpoint{1.869010in}{0.633527in}}%
\pgfpathlineto{\pgfqpoint{1.869010in}{0.625278in}}%
\pgfpathlineto{\pgfqpoint{1.870017in}{0.625278in}}%
\pgfpathlineto{\pgfqpoint{1.871357in}{0.633527in}}%
\pgfpathlineto{\pgfqpoint{1.871520in}{0.625278in}}%
\pgfpathlineto{\pgfqpoint{1.872364in}{0.625278in}}%
\pgfpathlineto{\pgfqpoint{1.873371in}{0.633527in}}%
\pgfpathlineto{\pgfqpoint{1.873371in}{0.625278in}}%
\pgfpathlineto{\pgfqpoint{1.874711in}{0.633527in}}%
\pgfpathlineto{\pgfqpoint{1.874711in}{0.625278in}}%
\pgfpathlineto{\pgfqpoint{1.875711in}{0.625278in}}%
\pgfpathlineto{\pgfqpoint{1.876888in}{0.633527in}}%
\pgfpathlineto{\pgfqpoint{1.877059in}{0.625278in}}%
\pgfpathlineto{\pgfqpoint{1.877895in}{0.625278in}}%
\pgfpathlineto{\pgfqpoint{1.879065in}{0.633527in}}%
\pgfpathlineto{\pgfqpoint{1.879235in}{0.625278in}}%
\pgfpathlineto{\pgfqpoint{1.880072in}{0.625278in}}%
\pgfpathlineto{\pgfqpoint{1.881412in}{0.633527in}}%
\pgfpathlineto{\pgfqpoint{1.881582in}{0.625278in}}%
\pgfpathlineto{\pgfqpoint{1.882419in}{0.625278in}}%
\pgfpathlineto{\pgfqpoint{1.883426in}{0.633527in}}%
\pgfpathlineto{\pgfqpoint{1.883426in}{0.625278in}}%
\pgfpathlineto{\pgfqpoint{1.884766in}{0.633527in}}%
\pgfpathlineto{\pgfqpoint{1.884937in}{0.625278in}}%
\pgfpathlineto{\pgfqpoint{1.885773in}{0.625278in}}%
\pgfpathlineto{\pgfqpoint{1.886943in}{0.633527in}}%
\pgfpathlineto{\pgfqpoint{1.887113in}{0.625278in}}%
\pgfpathlineto{\pgfqpoint{1.887950in}{0.625278in}}%
\pgfpathlineto{\pgfqpoint{1.889290in}{0.633527in}}%
\pgfpathlineto{\pgfqpoint{1.889461in}{0.625278in}}%
\pgfpathlineto{\pgfqpoint{1.890297in}{0.625278in}}%
\pgfpathlineto{\pgfqpoint{1.891304in}{0.633527in}}%
\pgfpathlineto{\pgfqpoint{1.891304in}{0.625278in}}%
\pgfpathlineto{\pgfqpoint{1.892474in}{0.633527in}}%
\pgfpathlineto{\pgfqpoint{1.892644in}{0.625278in}}%
\pgfpathlineto{\pgfqpoint{1.893481in}{0.625278in}}%
\pgfpathlineto{\pgfqpoint{1.894658in}{0.633527in}}%
\pgfpathlineto{\pgfqpoint{1.894821in}{0.625278in}}%
\pgfpathlineto{\pgfqpoint{1.895665in}{0.625278in}}%
\pgfpathlineto{\pgfqpoint{1.896835in}{0.633527in}}%
\pgfpathlineto{\pgfqpoint{1.897005in}{0.625278in}}%
\pgfpathlineto{\pgfqpoint{1.897901in}{0.625278in}}%
\pgfpathlineto{\pgfqpoint{1.899012in}{0.633527in}}%
\pgfpathlineto{\pgfqpoint{1.899012in}{0.625278in}}%
\pgfpathlineto{\pgfqpoint{1.900189in}{0.633527in}}%
\pgfpathlineto{\pgfqpoint{1.900360in}{0.625278in}}%
\pgfpathlineto{\pgfqpoint{1.901211in}{0.625278in}}%
\pgfpathlineto{\pgfqpoint{1.902203in}{0.633527in}}%
\pgfpathlineto{\pgfqpoint{1.902218in}{0.625278in}}%
\pgfpathlineto{\pgfqpoint{1.903558in}{0.650025in}}%
\pgfpathlineto{\pgfqpoint{1.903728in}{0.641776in}}%
\pgfpathlineto{\pgfqpoint{1.903877in}{0.633527in}}%
\pgfpathlineto{\pgfqpoint{1.904232in}{0.674773in}}%
\pgfpathlineto{\pgfqpoint{1.905890in}{0.707769in}}%
\pgfpathlineto{\pgfqpoint{1.906076in}{0.724268in}}%
\pgfpathlineto{\pgfqpoint{1.906912in}{0.691271in}}%
\pgfpathlineto{\pgfqpoint{1.908089in}{0.674773in}}%
\pgfpathlineto{\pgfqpoint{1.908400in}{0.683022in}}%
\pgfpathlineto{\pgfqpoint{1.908741in}{0.658274in}}%
\pgfpathlineto{\pgfqpoint{1.908926in}{0.666524in}}%
\pgfpathlineto{\pgfqpoint{1.909089in}{0.674773in}}%
\pgfpathlineto{\pgfqpoint{1.910096in}{0.650025in}}%
\pgfpathlineto{\pgfqpoint{1.911251in}{0.658274in}}%
\pgfpathlineto{\pgfqpoint{1.911606in}{0.650025in}}%
\pgfpathlineto{\pgfqpoint{1.912110in}{0.699520in}}%
\pgfpathlineto{\pgfqpoint{1.913095in}{0.707769in}}%
\pgfpathlineto{\pgfqpoint{1.913117in}{0.699520in}}%
\pgfpathlineto{\pgfqpoint{1.913324in}{0.683022in}}%
\pgfpathlineto{\pgfqpoint{1.913450in}{0.724268in}}%
\pgfpathlineto{\pgfqpoint{1.916279in}{0.881002in}}%
\pgfpathlineto{\pgfqpoint{1.916782in}{0.856254in}}%
\pgfpathlineto{\pgfqpoint{1.916804in}{0.848005in}}%
\pgfpathlineto{\pgfqpoint{1.917308in}{0.897500in}}%
\pgfpathlineto{\pgfqpoint{1.917811in}{0.864503in}}%
\pgfpathlineto{\pgfqpoint{1.917974in}{0.872753in}}%
\pgfpathlineto{\pgfqpoint{1.918966in}{0.856254in}}%
\pgfpathlineto{\pgfqpoint{1.918981in}{0.897500in}}%
\pgfpathlineto{\pgfqpoint{1.919988in}{0.872753in}}%
\pgfpathlineto{\pgfqpoint{1.920136in}{0.881002in}}%
\pgfpathlineto{\pgfqpoint{1.920492in}{0.848005in}}%
\pgfpathlineto{\pgfqpoint{1.920640in}{0.856254in}}%
\pgfpathlineto{\pgfqpoint{1.921647in}{0.831507in}}%
\pgfpathlineto{\pgfqpoint{1.921661in}{0.848005in}}%
\pgfpathlineto{\pgfqpoint{1.922987in}{0.881002in}}%
\pgfpathlineto{\pgfqpoint{1.923009in}{0.872753in}}%
\pgfpathlineto{\pgfqpoint{1.923320in}{0.856254in}}%
\pgfpathlineto{\pgfqpoint{1.923661in}{0.889251in}}%
\pgfpathlineto{\pgfqpoint{1.924660in}{0.905749in}}%
\pgfpathlineto{\pgfqpoint{1.924682in}{0.897500in}}%
\pgfpathlineto{\pgfqpoint{1.925334in}{0.856254in}}%
\pgfpathlineto{\pgfqpoint{1.925689in}{0.897500in}}%
\pgfpathlineto{\pgfqpoint{1.926341in}{0.955244in}}%
\pgfpathlineto{\pgfqpoint{1.927533in}{0.938746in}}%
\pgfpathlineto{\pgfqpoint{1.927696in}{0.946995in}}%
\pgfpathlineto{\pgfqpoint{1.928688in}{0.930497in}}%
\pgfpathlineto{\pgfqpoint{1.928851in}{0.955244in}}%
\pgfpathlineto{\pgfqpoint{1.929710in}{0.946995in}}%
\pgfpathlineto{\pgfqpoint{1.930865in}{0.955244in}}%
\pgfpathlineto{\pgfqpoint{1.931701in}{0.930497in}}%
\pgfpathlineto{\pgfqpoint{1.931887in}{0.946995in}}%
\pgfpathlineto{\pgfqpoint{1.933042in}{0.979992in}}%
\pgfpathlineto{\pgfqpoint{1.933234in}{0.963493in}}%
\pgfpathlineto{\pgfqpoint{1.933382in}{0.955244in}}%
\pgfpathlineto{\pgfqpoint{1.933730in}{0.996490in}}%
\pgfpathlineto{\pgfqpoint{1.935056in}{1.029487in}}%
\pgfpathlineto{\pgfqpoint{1.935078in}{1.021237in}}%
\pgfpathlineto{\pgfqpoint{1.935581in}{0.996490in}}%
\pgfpathlineto{\pgfqpoint{1.935892in}{1.029487in}}%
\pgfpathlineto{\pgfqpoint{1.935914in}{1.021237in}}%
\pgfpathlineto{\pgfqpoint{1.937573in}{1.103729in}}%
\pgfpathlineto{\pgfqpoint{1.937906in}{1.078981in}}%
\pgfpathlineto{\pgfqpoint{1.938410in}{1.103729in}}%
\pgfpathlineto{\pgfqpoint{1.938928in}{1.070732in}}%
\pgfpathlineto{\pgfqpoint{1.940083in}{1.078981in}}%
\pgfpathlineto{\pgfqpoint{1.941112in}{1.070732in}}%
\pgfpathlineto{\pgfqpoint{1.942452in}{1.045985in}}%
\pgfpathlineto{\pgfqpoint{1.943770in}{1.128476in}}%
\pgfpathlineto{\pgfqpoint{1.944296in}{1.087231in}}%
\pgfpathlineto{\pgfqpoint{1.944444in}{1.078981in}}%
\pgfpathlineto{\pgfqpoint{1.944799in}{1.120227in}}%
\pgfpathlineto{\pgfqpoint{1.945614in}{1.153224in}}%
\pgfpathlineto{\pgfqpoint{1.945955in}{1.128476in}}%
\pgfpathlineto{\pgfqpoint{1.946806in}{1.120227in}}%
\pgfpathlineto{\pgfqpoint{1.946976in}{1.144975in}}%
\pgfpathlineto{\pgfqpoint{1.949323in}{1.243965in}}%
\pgfpathlineto{\pgfqpoint{1.949494in}{1.235715in}}%
\pgfpathlineto{\pgfqpoint{1.949657in}{1.243965in}}%
\pgfpathlineto{\pgfqpoint{1.950664in}{1.219217in}}%
\pgfpathlineto{\pgfqpoint{1.951819in}{1.227466in}}%
\pgfpathlineto{\pgfqpoint{1.953011in}{1.194470in}}%
\pgfpathlineto{\pgfqpoint{1.953514in}{1.169722in}}%
\pgfpathlineto{\pgfqpoint{1.953833in}{1.202719in}}%
\pgfpathlineto{\pgfqpoint{1.953847in}{1.194470in}}%
\pgfpathlineto{\pgfqpoint{1.955839in}{1.252214in}}%
\pgfpathlineto{\pgfqpoint{1.955861in}{1.243965in}}%
\pgfpathlineto{\pgfqpoint{1.956365in}{1.293460in}}%
\pgfpathlineto{\pgfqpoint{1.956698in}{1.293460in}}%
\pgfpathlineto{\pgfqpoint{1.958023in}{1.227466in}}%
\pgfpathlineto{\pgfqpoint{1.958357in}{1.260463in}}%
\pgfpathlineto{\pgfqpoint{1.959549in}{1.318207in}}%
\pgfpathlineto{\pgfqpoint{1.961207in}{1.351204in}}%
\pgfpathlineto{\pgfqpoint{1.962214in}{1.326456in}}%
\pgfpathlineto{\pgfqpoint{1.962229in}{1.342955in}}%
\pgfpathlineto{\pgfqpoint{1.963888in}{1.375951in}}%
\pgfpathlineto{\pgfqpoint{1.964783in}{1.351204in}}%
\pgfpathlineto{\pgfqpoint{1.964917in}{1.392449in}}%
\pgfpathlineto{\pgfqpoint{1.965583in}{1.417197in}}%
\pgfpathlineto{\pgfqpoint{1.965916in}{1.392449in}}%
\pgfpathlineto{\pgfqpoint{1.966760in}{1.367702in}}%
\pgfpathlineto{\pgfqpoint{1.966923in}{1.392449in}}%
\pgfpathlineto{\pgfqpoint{1.967597in}{1.441944in}}%
\pgfpathlineto{\pgfqpoint{1.968434in}{1.417197in}}%
\pgfpathlineto{\pgfqpoint{1.969107in}{1.392449in}}%
\pgfpathlineto{\pgfqpoint{1.969270in}{1.417197in}}%
\pgfpathlineto{\pgfqpoint{1.970781in}{1.466692in}}%
\pgfpathlineto{\pgfqpoint{1.970951in}{1.458443in}}%
\pgfpathlineto{\pgfqpoint{1.971447in}{1.441944in}}%
\pgfpathlineto{\pgfqpoint{1.971766in}{1.474941in}}%
\pgfpathlineto{\pgfqpoint{1.971788in}{1.466692in}}%
\pgfpathlineto{\pgfqpoint{1.972602in}{1.474941in}}%
\pgfpathlineto{\pgfqpoint{1.972624in}{1.466692in}}%
\pgfpathlineto{\pgfqpoint{1.973965in}{1.441944in}}%
\pgfpathlineto{\pgfqpoint{1.975120in}{1.450194in}}%
\pgfpathlineto{\pgfqpoint{1.976312in}{1.417197in}}%
\pgfpathlineto{\pgfqpoint{1.976460in}{1.425446in}}%
\pgfpathlineto{\pgfqpoint{1.976793in}{1.400699in}}%
\pgfpathlineto{\pgfqpoint{1.976986in}{1.408948in}}%
\pgfpathlineto{\pgfqpoint{1.977134in}{1.400699in}}%
\pgfpathlineto{\pgfqpoint{1.977148in}{1.441944in}}%
\pgfpathlineto{\pgfqpoint{1.977822in}{1.417197in}}%
\pgfpathlineto{\pgfqpoint{1.978489in}{1.466692in}}%
\pgfpathlineto{\pgfqpoint{1.978992in}{1.433695in}}%
\pgfpathlineto{\pgfqpoint{1.979481in}{1.425446in}}%
\pgfpathlineto{\pgfqpoint{1.979496in}{1.466692in}}%
\pgfpathlineto{\pgfqpoint{1.982161in}{1.375951in}}%
\pgfpathlineto{\pgfqpoint{1.979992in}{1.474941in}}%
\pgfpathlineto{\pgfqpoint{1.982161in}{1.384200in}}%
\pgfpathlineto{\pgfqpoint{1.982516in}{1.367702in}}%
\pgfpathlineto{\pgfqpoint{1.983183in}{1.392449in}}%
\pgfpathlineto{\pgfqpoint{1.986537in}{1.565682in}}%
\pgfpathlineto{\pgfqpoint{1.987018in}{1.549184in}}%
\pgfpathlineto{\pgfqpoint{1.988047in}{1.540934in}}%
\pgfpathlineto{\pgfqpoint{1.988084in}{1.549184in}}%
\pgfpathlineto{\pgfqpoint{1.988714in}{1.598678in}}%
\pgfpathlineto{\pgfqpoint{1.989388in}{1.565682in}}%
\pgfpathlineto{\pgfqpoint{1.995422in}{1.318207in}}%
\pgfpathlineto{\pgfqpoint{1.996407in}{1.326456in}}%
\pgfpathlineto{\pgfqpoint{1.995903in}{1.301709in}}%
\pgfpathlineto{\pgfqpoint{1.996429in}{1.318207in}}%
\pgfpathlineto{\pgfqpoint{1.999442in}{1.219217in}}%
\pgfpathlineto{\pgfqpoint{1.999939in}{1.227466in}}%
\pgfpathlineto{\pgfqpoint{2.000287in}{1.194470in}}%
\pgfpathlineto{\pgfqpoint{2.000449in}{1.194470in}}%
\pgfpathlineto{\pgfqpoint{2.000605in}{1.202719in}}%
\pgfpathlineto{\pgfqpoint{2.000953in}{1.169722in}}%
\pgfpathlineto{\pgfqpoint{2.002797in}{1.120227in}}%
\pgfpathlineto{\pgfqpoint{2.003137in}{1.128476in}}%
\pgfpathlineto{\pgfqpoint{2.003470in}{1.111978in}}%
\pgfpathlineto{\pgfqpoint{2.005314in}{0.996490in}}%
\pgfpathlineto{\pgfqpoint{2.011015in}{0.699520in}}%
\pgfpathlineto{\pgfqpoint{2.011511in}{0.724268in}}%
\pgfpathlineto{\pgfqpoint{2.012511in}{0.732517in}}%
\pgfpathlineto{\pgfqpoint{2.012518in}{0.724268in}}%
\pgfpathlineto{\pgfqpoint{2.015206in}{0.633527in}}%
\pgfpathlineto{\pgfqpoint{2.015369in}{0.650025in}}%
\pgfpathlineto{\pgfqpoint{2.015702in}{0.625278in}}%
\pgfpathlineto{\pgfqpoint{2.016191in}{0.633527in}}%
\pgfpathlineto{\pgfqpoint{2.017213in}{0.625278in}}%
\pgfpathlineto{\pgfqpoint{2.017220in}{0.633527in}}%
\pgfpathlineto{\pgfqpoint{2.018390in}{0.625278in}}%
\pgfpathlineto{\pgfqpoint{2.019552in}{0.633527in}}%
\pgfpathlineto{\pgfqpoint{2.020567in}{0.625278in}}%
\pgfpathlineto{\pgfqpoint{2.021892in}{0.633527in}}%
\pgfpathlineto{\pgfqpoint{2.022077in}{0.625278in}}%
\pgfpathlineto{\pgfqpoint{2.022892in}{0.625278in}}%
\pgfpathlineto{\pgfqpoint{2.024239in}{0.633527in}}%
\pgfpathlineto{\pgfqpoint{2.024254in}{0.625278in}}%
\pgfpathlineto{\pgfqpoint{2.025246in}{0.625278in}}%
\pgfpathlineto{\pgfqpoint{2.026423in}{0.633527in}}%
\pgfpathlineto{\pgfqpoint{2.026438in}{0.625278in}}%
\pgfpathlineto{\pgfqpoint{2.027445in}{0.625278in}}%
\pgfpathlineto{\pgfqpoint{2.033472in}{0.922247in}}%
\pgfpathlineto{\pgfqpoint{2.033642in}{0.913998in}}%
\pgfpathlineto{\pgfqpoint{2.034146in}{0.905749in}}%
\pgfpathlineto{\pgfqpoint{2.033813in}{0.930497in}}%
\pgfpathlineto{\pgfqpoint{2.034649in}{0.922247in}}%
\pgfpathlineto{\pgfqpoint{2.038166in}{0.798510in}}%
\pgfpathlineto{\pgfqpoint{2.039166in}{0.806759in}}%
\pgfpathlineto{\pgfqpoint{2.039173in}{0.798510in}}%
\pgfpathlineto{\pgfqpoint{2.041017in}{0.749015in}}%
\pgfpathlineto{\pgfqpoint{2.041180in}{0.757264in}}%
\pgfpathlineto{\pgfqpoint{2.041506in}{0.732517in}}%
\pgfpathlineto{\pgfqpoint{2.041691in}{0.740766in}}%
\pgfpathlineto{\pgfqpoint{2.043364in}{0.699520in}}%
\pgfpathlineto{\pgfqpoint{2.043520in}{0.707769in}}%
\pgfpathlineto{\pgfqpoint{2.043705in}{0.683022in}}%
\pgfpathlineto{\pgfqpoint{2.044038in}{0.683022in}}%
\pgfpathlineto{\pgfqpoint{2.044541in}{0.699520in}}%
\pgfpathlineto{\pgfqpoint{2.046044in}{0.625278in}}%
\pgfpathlineto{\pgfqpoint{2.047044in}{0.633527in}}%
\pgfpathlineto{\pgfqpoint{2.047044in}{0.625278in}}%
\pgfpathlineto{\pgfqpoint{2.048399in}{0.633527in}}%
\pgfpathlineto{\pgfqpoint{2.048540in}{0.625278in}}%
\pgfpathlineto{\pgfqpoint{2.049399in}{0.625278in}}%
\pgfpathlineto{\pgfqpoint{2.050568in}{0.633527in}}%
\pgfpathlineto{\pgfqpoint{2.050731in}{0.625278in}}%
\pgfpathlineto{\pgfqpoint{2.051568in}{0.625278in}}%
\pgfpathlineto{\pgfqpoint{2.052575in}{0.633527in}}%
\pgfpathlineto{\pgfqpoint{2.052575in}{0.625278in}}%
\pgfpathlineto{\pgfqpoint{2.053745in}{0.633527in}}%
\pgfpathlineto{\pgfqpoint{2.053915in}{0.625278in}}%
\pgfpathlineto{\pgfqpoint{2.054574in}{0.625278in}}%
\pgfpathlineto{\pgfqpoint{2.055914in}{0.633527in}}%
\pgfpathlineto{\pgfqpoint{2.056588in}{0.625278in}}%
\pgfpathlineto{\pgfqpoint{2.056921in}{0.625278in}}%
\pgfpathlineto{\pgfqpoint{2.058106in}{0.633527in}}%
\pgfpathlineto{\pgfqpoint{2.058261in}{0.625278in}}%
\pgfpathlineto{\pgfqpoint{2.059113in}{0.625278in}}%
\pgfpathlineto{\pgfqpoint{2.060105in}{0.633527in}}%
\pgfpathlineto{\pgfqpoint{2.060105in}{0.625278in}}%
\pgfpathlineto{\pgfqpoint{2.061282in}{0.633527in}}%
\pgfpathlineto{\pgfqpoint{2.061282in}{0.625278in}}%
\pgfpathlineto{\pgfqpoint{2.062289in}{0.625278in}}%
\pgfpathlineto{\pgfqpoint{2.063459in}{0.633527in}}%
\pgfpathlineto{\pgfqpoint{2.063629in}{0.625278in}}%
\pgfpathlineto{\pgfqpoint{2.064466in}{0.625278in}}%
\pgfpathlineto{\pgfqpoint{2.065643in}{0.633527in}}%
\pgfpathlineto{\pgfqpoint{2.065806in}{0.625278in}}%
\pgfpathlineto{\pgfqpoint{2.066650in}{0.625278in}}%
\pgfpathlineto{\pgfqpoint{2.067820in}{0.633527in}}%
\pgfpathlineto{\pgfqpoint{2.067998in}{0.625278in}}%
\pgfpathlineto{\pgfqpoint{2.068827in}{0.625278in}}%
\pgfpathlineto{\pgfqpoint{2.069997in}{0.633527in}}%
\pgfpathlineto{\pgfqpoint{2.070167in}{0.625278in}}%
\pgfpathlineto{\pgfqpoint{2.070834in}{0.625278in}}%
\pgfpathlineto{\pgfqpoint{2.072018in}{0.633527in}}%
\pgfpathlineto{\pgfqpoint{2.072352in}{0.625278in}}%
\pgfpathlineto{\pgfqpoint{2.073025in}{0.625278in}}%
\pgfpathlineto{\pgfqpoint{2.074188in}{0.633527in}}%
\pgfpathlineto{\pgfqpoint{2.074366in}{0.625278in}}%
\pgfpathlineto{\pgfqpoint{2.075195in}{0.625278in}}%
\pgfpathlineto{\pgfqpoint{2.076372in}{0.633527in}}%
\pgfpathlineto{\pgfqpoint{2.076705in}{0.625278in}}%
\pgfpathlineto{\pgfqpoint{2.077379in}{0.625278in}}%
\pgfpathlineto{\pgfqpoint{2.078379in}{0.633527in}}%
\pgfpathlineto{\pgfqpoint{2.078379in}{0.625278in}}%
\pgfpathlineto{\pgfqpoint{2.079563in}{0.633527in}}%
\pgfpathlineto{\pgfqpoint{2.079719in}{0.625278in}}%
\pgfpathlineto{\pgfqpoint{2.080563in}{0.625278in}}%
\pgfpathlineto{\pgfqpoint{2.081733in}{0.633527in}}%
\pgfpathlineto{\pgfqpoint{2.081733in}{0.625278in}}%
\pgfpathlineto{\pgfqpoint{2.082740in}{0.625278in}}%
\pgfpathlineto{\pgfqpoint{2.083910in}{0.633527in}}%
\pgfpathlineto{\pgfqpoint{2.084250in}{0.625278in}}%
\pgfpathlineto{\pgfqpoint{2.084916in}{0.625278in}}%
\pgfpathlineto{\pgfqpoint{2.086094in}{0.633527in}}%
\pgfpathlineto{\pgfqpoint{2.086094in}{0.625278in}}%
\pgfpathlineto{\pgfqpoint{2.087101in}{0.625278in}}%
\pgfpathlineto{\pgfqpoint{2.088271in}{0.633527in}}%
\pgfpathlineto{\pgfqpoint{2.088604in}{0.625278in}}%
\pgfpathlineto{\pgfqpoint{2.089107in}{0.625278in}}%
\pgfpathlineto{\pgfqpoint{2.090285in}{0.633527in}}%
\pgfpathlineto{\pgfqpoint{2.090447in}{0.625278in}}%
\pgfpathlineto{\pgfqpoint{2.091291in}{0.625278in}}%
\pgfpathlineto{\pgfqpoint{2.092461in}{0.633527in}}%
\pgfpathlineto{\pgfqpoint{2.092632in}{0.625278in}}%
\pgfpathlineto{\pgfqpoint{2.093468in}{0.625278in}}%
\pgfpathlineto{\pgfqpoint{2.094808in}{0.633527in}}%
\pgfpathlineto{\pgfqpoint{2.094808in}{0.625278in}}%
\pgfpathlineto{\pgfqpoint{2.095645in}{0.625278in}}%
\pgfpathlineto{\pgfqpoint{2.096985in}{0.633527in}}%
\pgfpathlineto{\pgfqpoint{2.096985in}{0.625278in}}%
\pgfpathlineto{\pgfqpoint{2.097992in}{0.625278in}}%
\pgfpathlineto{\pgfqpoint{2.098999in}{0.633527in}}%
\pgfpathlineto{\pgfqpoint{2.098999in}{0.625278in}}%
\pgfpathlineto{\pgfqpoint{2.100169in}{0.633527in}}%
\pgfpathlineto{\pgfqpoint{2.100339in}{0.625278in}}%
\pgfpathlineto{\pgfqpoint{2.101176in}{0.625278in}}%
\pgfpathlineto{\pgfqpoint{2.102361in}{0.633527in}}%
\pgfpathlineto{\pgfqpoint{2.102516in}{0.625278in}}%
\pgfpathlineto{\pgfqpoint{2.103368in}{0.625278in}}%
\pgfpathlineto{\pgfqpoint{2.104530in}{0.633527in}}%
\pgfpathlineto{\pgfqpoint{2.104700in}{0.625278in}}%
\pgfpathlineto{\pgfqpoint{2.105537in}{0.625278in}}%
\pgfpathlineto{\pgfqpoint{2.106707in}{0.633527in}}%
\pgfpathlineto{\pgfqpoint{2.106877in}{0.625278in}}%
\pgfpathlineto{\pgfqpoint{2.107714in}{0.625278in}}%
\pgfpathlineto{\pgfqpoint{2.108891in}{0.633527in}}%
\pgfpathlineto{\pgfqpoint{2.109054in}{0.625278in}}%
\pgfpathlineto{\pgfqpoint{2.109728in}{0.625278in}}%
\pgfpathlineto{\pgfqpoint{2.111068in}{0.633527in}}%
\pgfpathlineto{\pgfqpoint{2.111068in}{0.625278in}}%
\pgfpathlineto{\pgfqpoint{2.112075in}{0.625278in}}%
\pgfpathlineto{\pgfqpoint{2.113245in}{0.633527in}}%
\pgfpathlineto{\pgfqpoint{2.113415in}{0.625278in}}%
\pgfpathlineto{\pgfqpoint{2.114252in}{0.625278in}}%
\pgfpathlineto{\pgfqpoint{2.115429in}{0.633527in}}%
\pgfpathlineto{\pgfqpoint{2.115429in}{0.625278in}}%
\pgfpathlineto{\pgfqpoint{2.116436in}{0.625278in}}%
\pgfpathlineto{\pgfqpoint{2.117606in}{0.633527in}}%
\pgfpathlineto{\pgfqpoint{2.117776in}{0.625278in}}%
\pgfpathlineto{\pgfqpoint{2.118613in}{0.625278in}}%
\pgfpathlineto{\pgfqpoint{2.119620in}{0.633527in}}%
\pgfpathlineto{\pgfqpoint{2.119620in}{0.625278in}}%
\pgfpathlineto{\pgfqpoint{2.120790in}{0.633527in}}%
\pgfpathlineto{\pgfqpoint{2.120960in}{0.625278in}}%
\pgfpathlineto{\pgfqpoint{2.121797in}{0.625278in}}%
\pgfpathlineto{\pgfqpoint{2.122967in}{0.633527in}}%
\pgfpathlineto{\pgfqpoint{2.123137in}{0.625278in}}%
\pgfpathlineto{\pgfqpoint{2.123974in}{0.625278in}}%
\pgfpathlineto{\pgfqpoint{2.124981in}{0.633527in}}%
\pgfpathlineto{\pgfqpoint{2.124981in}{0.625278in}}%
\pgfpathlineto{\pgfqpoint{2.126158in}{0.633527in}}%
\pgfpathlineto{\pgfqpoint{2.126321in}{0.625278in}}%
\pgfpathlineto{\pgfqpoint{2.127157in}{0.625278in}}%
\pgfpathlineto{\pgfqpoint{2.128335in}{0.633527in}}%
\pgfpathlineto{\pgfqpoint{2.128335in}{0.625278in}}%
\pgfpathlineto{\pgfqpoint{2.129342in}{0.625278in}}%
\pgfpathlineto{\pgfqpoint{2.130512in}{0.633527in}}%
\pgfpathlineto{\pgfqpoint{2.130682in}{0.625278in}}%
\pgfpathlineto{\pgfqpoint{2.131348in}{0.625278in}}%
\pgfpathlineto{\pgfqpoint{2.132525in}{0.633527in}}%
\pgfpathlineto{\pgfqpoint{2.132696in}{0.625278in}}%
\pgfpathlineto{\pgfqpoint{2.133532in}{0.625278in}}%
\pgfpathlineto{\pgfqpoint{2.134539in}{0.633527in}}%
\pgfpathlineto{\pgfqpoint{2.134539in}{0.625278in}}%
\pgfpathlineto{\pgfqpoint{2.135709in}{0.633527in}}%
\pgfpathlineto{\pgfqpoint{2.135709in}{0.625278in}}%
\pgfpathlineto{\pgfqpoint{2.136716in}{0.625278in}}%
\pgfpathlineto{\pgfqpoint{2.137886in}{0.633527in}}%
\pgfpathlineto{\pgfqpoint{2.138056in}{0.625278in}}%
\pgfpathlineto{\pgfqpoint{2.138730in}{0.625278in}}%
\pgfpathlineto{\pgfqpoint{2.139900in}{0.633527in}}%
\pgfpathlineto{\pgfqpoint{2.140070in}{0.625278in}}%
\pgfpathlineto{\pgfqpoint{2.140907in}{0.625278in}}%
\pgfpathlineto{\pgfqpoint{2.142077in}{0.633527in}}%
\pgfpathlineto{\pgfqpoint{2.142077in}{0.625278in}}%
\pgfpathlineto{\pgfqpoint{2.143084in}{0.625278in}}%
\pgfpathlineto{\pgfqpoint{2.144261in}{0.633527in}}%
\pgfpathlineto{\pgfqpoint{2.144424in}{0.625278in}}%
\pgfpathlineto{\pgfqpoint{2.145268in}{0.625278in}}%
\pgfpathlineto{\pgfqpoint{2.146438in}{0.633527in}}%
\pgfpathlineto{\pgfqpoint{2.146771in}{0.625278in}}%
\pgfpathlineto{\pgfqpoint{2.147445in}{0.625278in}}%
\pgfpathlineto{\pgfqpoint{2.148452in}{0.633527in}}%
\pgfpathlineto{\pgfqpoint{2.148452in}{0.625278in}}%
\pgfpathlineto{\pgfqpoint{2.149792in}{0.633527in}}%
\pgfpathlineto{\pgfqpoint{2.149792in}{0.625278in}}%
\pgfpathlineto{\pgfqpoint{2.150799in}{0.625278in}}%
\pgfpathlineto{\pgfqpoint{2.151969in}{0.633527in}}%
\pgfpathlineto{\pgfqpoint{2.152139in}{0.625278in}}%
\pgfpathlineto{\pgfqpoint{2.152976in}{0.625278in}}%
\pgfpathlineto{\pgfqpoint{2.154153in}{0.633527in}}%
\pgfpathlineto{\pgfqpoint{2.154153in}{0.625278in}}%
\pgfpathlineto{\pgfqpoint{2.155153in}{0.625278in}}%
\pgfpathlineto{\pgfqpoint{2.156493in}{0.633527in}}%
\pgfpathlineto{\pgfqpoint{2.156493in}{0.625278in}}%
\pgfpathlineto{\pgfqpoint{2.157337in}{0.625278in}}%
\pgfpathlineto{\pgfqpoint{2.158677in}{0.633527in}}%
\pgfpathlineto{\pgfqpoint{2.158840in}{0.625278in}}%
\pgfpathlineto{\pgfqpoint{2.159514in}{0.625278in}}%
\pgfpathlineto{\pgfqpoint{2.160684in}{0.633527in}}%
\pgfpathlineto{\pgfqpoint{2.160854in}{0.625278in}}%
\pgfpathlineto{\pgfqpoint{2.161528in}{0.625278in}}%
\pgfpathlineto{\pgfqpoint{2.162698in}{0.633527in}}%
\pgfpathlineto{\pgfqpoint{2.163031in}{0.625278in}}%
\pgfpathlineto{\pgfqpoint{2.163705in}{0.625278in}}%
\pgfpathlineto{\pgfqpoint{2.164874in}{0.633527in}}%
\pgfpathlineto{\pgfqpoint{2.165045in}{0.625278in}}%
\pgfpathlineto{\pgfqpoint{2.165881in}{0.625278in}}%
\pgfpathlineto{\pgfqpoint{2.167059in}{0.633527in}}%
\pgfpathlineto{\pgfqpoint{2.167222in}{0.625278in}}%
\pgfpathlineto{\pgfqpoint{2.168066in}{0.625278in}}%
\pgfpathlineto{\pgfqpoint{2.169235in}{0.633527in}}%
\pgfpathlineto{\pgfqpoint{2.169235in}{0.625278in}}%
\pgfpathlineto{\pgfqpoint{2.170242in}{0.625278in}}%
\pgfpathlineto{\pgfqpoint{2.171412in}{0.633527in}}%
\pgfpathlineto{\pgfqpoint{2.171412in}{0.625278in}}%
\pgfpathlineto{\pgfqpoint{2.172419in}{0.625278in}}%
\pgfpathlineto{\pgfqpoint{2.173597in}{0.633527in}}%
\pgfpathlineto{\pgfqpoint{2.173759in}{0.625278in}}%
\pgfpathlineto{\pgfqpoint{2.174433in}{0.625278in}}%
\pgfpathlineto{\pgfqpoint{2.175603in}{0.633527in}}%
\pgfpathlineto{\pgfqpoint{2.175773in}{0.625278in}}%
\pgfpathlineto{\pgfqpoint{2.176447in}{0.625278in}}%
\pgfpathlineto{\pgfqpoint{2.177617in}{0.633527in}}%
\pgfpathlineto{\pgfqpoint{2.177787in}{0.625278in}}%
\pgfpathlineto{\pgfqpoint{2.178624in}{0.625278in}}%
\pgfpathlineto{\pgfqpoint{2.179794in}{0.633527in}}%
\pgfpathlineto{\pgfqpoint{2.179794in}{0.625278in}}%
\pgfpathlineto{\pgfqpoint{2.180801in}{0.625278in}}%
\pgfpathlineto{\pgfqpoint{2.181978in}{0.633527in}}%
\pgfpathlineto{\pgfqpoint{2.182141in}{0.625278in}}%
\pgfpathlineto{\pgfqpoint{2.182985in}{0.625278in}}%
\pgfpathlineto{\pgfqpoint{2.184155in}{0.633527in}}%
\pgfpathlineto{\pgfqpoint{2.184325in}{0.625278in}}%
\pgfpathlineto{\pgfqpoint{2.185162in}{0.625278in}}%
\pgfpathlineto{\pgfqpoint{2.186169in}{0.633527in}}%
\pgfpathlineto{\pgfqpoint{2.186169in}{0.625278in}}%
\pgfpathlineto{\pgfqpoint{2.187339in}{0.633527in}}%
\pgfpathlineto{\pgfqpoint{2.187509in}{0.625278in}}%
\pgfpathlineto{\pgfqpoint{2.188346in}{0.625278in}}%
\pgfpathlineto{\pgfqpoint{2.189523in}{0.633527in}}%
\pgfpathlineto{\pgfqpoint{2.189523in}{0.625278in}}%
\pgfpathlineto{\pgfqpoint{2.190545in}{0.625278in}}%
\pgfpathlineto{\pgfqpoint{2.192055in}{0.674773in}}%
\pgfpathlineto{\pgfqpoint{2.192218in}{0.666524in}}%
\pgfpathlineto{\pgfqpoint{2.192366in}{0.658274in}}%
\pgfpathlineto{\pgfqpoint{2.192388in}{0.699520in}}%
\pgfpathlineto{\pgfqpoint{2.192722in}{0.699520in}}%
\pgfpathlineto{\pgfqpoint{2.194062in}{0.674773in}}%
\pgfpathlineto{\pgfqpoint{2.195217in}{0.683022in}}%
\pgfpathlineto{\pgfqpoint{2.195905in}{0.674773in}}%
\pgfpathlineto{\pgfqpoint{2.196224in}{0.707769in}}%
\pgfpathlineto{\pgfqpoint{2.196246in}{0.699520in}}%
\pgfpathlineto{\pgfqpoint{2.197401in}{0.732517in}}%
\pgfpathlineto{\pgfqpoint{2.196749in}{0.674773in}}%
\pgfpathlineto{\pgfqpoint{2.197586in}{0.716019in}}%
\pgfpathlineto{\pgfqpoint{2.199074in}{0.658274in}}%
\pgfpathlineto{\pgfqpoint{2.199259in}{0.674773in}}%
\pgfpathlineto{\pgfqpoint{2.199933in}{0.724268in}}%
\pgfpathlineto{\pgfqpoint{2.200600in}{0.691271in}}%
\pgfpathlineto{\pgfqpoint{2.202443in}{0.625278in}}%
\pgfpathlineto{\pgfqpoint{2.202591in}{0.641776in}}%
\pgfpathlineto{\pgfqpoint{2.203598in}{0.658274in}}%
\pgfpathlineto{\pgfqpoint{2.203621in}{0.650025in}}%
\pgfpathlineto{\pgfqpoint{2.203769in}{0.658274in}}%
\pgfpathlineto{\pgfqpoint{2.204124in}{0.625278in}}%
\pgfpathlineto{\pgfqpoint{2.204287in}{0.625278in}}%
\pgfpathlineto{\pgfqpoint{2.205279in}{0.633527in}}%
\pgfpathlineto{\pgfqpoint{2.205294in}{0.625278in}}%
\pgfpathlineto{\pgfqpoint{2.206619in}{0.658274in}}%
\pgfpathlineto{\pgfqpoint{2.206804in}{0.641776in}}%
\pgfpathlineto{\pgfqpoint{2.206952in}{0.633527in}}%
\pgfpathlineto{\pgfqpoint{2.207123in}{0.658274in}}%
\pgfpathlineto{\pgfqpoint{2.207478in}{0.650025in}}%
\pgfpathlineto{\pgfqpoint{2.207959in}{0.633527in}}%
\pgfpathlineto{\pgfqpoint{2.208633in}{0.658274in}}%
\pgfpathlineto{\pgfqpoint{2.209655in}{0.625278in}}%
\pgfpathlineto{\pgfqpoint{2.211313in}{0.683022in}}%
\pgfpathlineto{\pgfqpoint{2.211499in}{0.666524in}}%
\pgfpathlineto{\pgfqpoint{2.212483in}{0.633527in}}%
\pgfpathlineto{\pgfqpoint{2.212668in}{0.650025in}}%
\pgfpathlineto{\pgfqpoint{2.213824in}{0.658274in}}%
\pgfpathlineto{\pgfqpoint{2.214853in}{0.650025in}}%
\pgfpathlineto{\pgfqpoint{2.216008in}{0.732517in}}%
\pgfpathlineto{\pgfqpoint{2.216526in}{0.691271in}}%
\pgfpathlineto{\pgfqpoint{2.216674in}{0.683022in}}%
\pgfpathlineto{\pgfqpoint{2.216844in}{0.707769in}}%
\pgfpathlineto{\pgfqpoint{2.217200in}{0.699520in}}%
\pgfpathlineto{\pgfqpoint{2.217348in}{0.707769in}}%
\pgfpathlineto{\pgfqpoint{2.217703in}{0.674773in}}%
\pgfpathlineto{\pgfqpoint{2.218014in}{0.683022in}}%
\pgfpathlineto{\pgfqpoint{2.218873in}{0.674773in}}%
\pgfpathlineto{\pgfqpoint{2.218207in}{0.699520in}}%
\pgfpathlineto{\pgfqpoint{2.219043in}{0.699520in}}%
\pgfpathlineto{\pgfqpoint{2.220532in}{0.732517in}}%
\pgfpathlineto{\pgfqpoint{2.220554in}{0.724268in}}%
\pgfpathlineto{\pgfqpoint{2.221872in}{0.707769in}}%
\pgfpathlineto{\pgfqpoint{2.222227in}{0.749015in}}%
\pgfpathlineto{\pgfqpoint{2.222894in}{0.749015in}}%
\pgfpathlineto{\pgfqpoint{2.223382in}{0.757264in}}%
\pgfpathlineto{\pgfqpoint{2.223567in}{0.740766in}}%
\pgfpathlineto{\pgfqpoint{2.225056in}{0.707769in}}%
\pgfpathlineto{\pgfqpoint{2.225056in}{0.716019in}}%
\pgfpathlineto{\pgfqpoint{2.226248in}{0.773763in}}%
\pgfpathlineto{\pgfqpoint{2.227588in}{0.749015in}}%
\pgfpathlineto{\pgfqpoint{2.228077in}{0.782012in}}%
\pgfpathlineto{\pgfqpoint{2.228432in}{0.749015in}}%
\pgfpathlineto{\pgfqpoint{2.229098in}{0.724268in}}%
\pgfpathlineto{\pgfqpoint{2.229269in}{0.749015in}}%
\pgfpathlineto{\pgfqpoint{2.230779in}{0.798510in}}%
\pgfpathlineto{\pgfqpoint{2.230942in}{0.790261in}}%
\pgfpathlineto{\pgfqpoint{2.231927in}{0.757264in}}%
\pgfpathlineto{\pgfqpoint{2.232119in}{0.773763in}}%
\pgfpathlineto{\pgfqpoint{2.233126in}{0.798510in}}%
\pgfpathlineto{\pgfqpoint{2.233289in}{0.790261in}}%
\pgfpathlineto{\pgfqpoint{2.233459in}{0.798510in}}%
\pgfpathlineto{\pgfqpoint{2.234466in}{0.773763in}}%
\pgfpathlineto{\pgfqpoint{2.235281in}{0.782012in}}%
\pgfpathlineto{\pgfqpoint{2.235303in}{0.773763in}}%
\pgfpathlineto{\pgfqpoint{2.235621in}{0.757264in}}%
\pgfpathlineto{\pgfqpoint{2.235955in}{0.790261in}}%
\pgfpathlineto{\pgfqpoint{2.236288in}{0.806759in}}%
\pgfpathlineto{\pgfqpoint{2.236643in}{0.773763in}}%
\pgfpathlineto{\pgfqpoint{2.236976in}{0.798510in}}%
\pgfpathlineto{\pgfqpoint{2.238131in}{0.831507in}}%
\pgfpathlineto{\pgfqpoint{2.237628in}{0.782012in}}%
\pgfpathlineto{\pgfqpoint{2.238317in}{0.815008in}}%
\pgfpathlineto{\pgfqpoint{2.239309in}{0.806759in}}%
\pgfpathlineto{\pgfqpoint{2.238968in}{0.831507in}}%
\pgfpathlineto{\pgfqpoint{2.239309in}{0.815008in}}%
\pgfpathlineto{\pgfqpoint{2.240308in}{0.831507in}}%
\pgfpathlineto{\pgfqpoint{2.240331in}{0.823258in}}%
\pgfpathlineto{\pgfqpoint{2.242322in}{0.881002in}}%
\pgfpathlineto{\pgfqpoint{2.242848in}{0.872753in}}%
\pgfpathlineto{\pgfqpoint{2.243159in}{0.905749in}}%
\pgfpathlineto{\pgfqpoint{2.243351in}{0.897500in}}%
\pgfpathlineto{\pgfqpoint{2.243500in}{0.905749in}}%
\pgfpathlineto{\pgfqpoint{2.243848in}{0.872753in}}%
\pgfpathlineto{\pgfqpoint{2.244166in}{0.881002in}}%
\pgfpathlineto{\pgfqpoint{2.244351in}{0.872753in}}%
\pgfpathlineto{\pgfqpoint{2.245173in}{0.913998in}}%
\pgfpathlineto{\pgfqpoint{2.246195in}{0.971742in}}%
\pgfpathlineto{\pgfqpoint{2.246365in}{0.963493in}}%
\pgfpathlineto{\pgfqpoint{2.247372in}{0.946995in}}%
\pgfpathlineto{\pgfqpoint{2.246535in}{0.971742in}}%
\pgfpathlineto{\pgfqpoint{2.247520in}{0.955244in}}%
\pgfpathlineto{\pgfqpoint{2.247690in}{0.979992in}}%
\pgfpathlineto{\pgfqpoint{2.248542in}{0.971742in}}%
\pgfpathlineto{\pgfqpoint{2.249030in}{0.979992in}}%
\pgfpathlineto{\pgfqpoint{2.249216in}{0.963493in}}%
\pgfpathlineto{\pgfqpoint{2.249549in}{0.946995in}}%
\pgfpathlineto{\pgfqpoint{2.249867in}{0.979992in}}%
\pgfpathlineto{\pgfqpoint{2.249889in}{0.971742in}}%
\pgfpathlineto{\pgfqpoint{2.250541in}{0.955244in}}%
\pgfpathlineto{\pgfqpoint{2.251037in}{0.979992in}}%
\pgfpathlineto{\pgfqpoint{2.252214in}{0.955244in}}%
\pgfpathlineto{\pgfqpoint{2.252385in}{0.979992in}}%
\pgfpathlineto{\pgfqpoint{2.253236in}{0.971742in}}%
\pgfpathlineto{\pgfqpoint{2.254391in}{0.979992in}}%
\pgfpathlineto{\pgfqpoint{2.255420in}{0.971742in}}%
\pgfpathlineto{\pgfqpoint{2.256072in}{0.979992in}}%
\pgfpathlineto{\pgfqpoint{2.256257in}{0.963493in}}%
\pgfpathlineto{\pgfqpoint{2.256405in}{0.955244in}}%
\pgfpathlineto{\pgfqpoint{2.256575in}{0.979992in}}%
\pgfpathlineto{\pgfqpoint{2.257094in}{0.971742in}}%
\pgfpathlineto{\pgfqpoint{2.257915in}{1.004739in}}%
\pgfpathlineto{\pgfqpoint{2.258249in}{0.979992in}}%
\pgfpathlineto{\pgfqpoint{2.258434in}{0.971742in}}%
\pgfpathlineto{\pgfqpoint{2.258752in}{1.004739in}}%
\pgfpathlineto{\pgfqpoint{2.259270in}{0.996490in}}%
\pgfpathlineto{\pgfqpoint{2.259589in}{0.979992in}}%
\pgfpathlineto{\pgfqpoint{2.259611in}{1.021237in}}%
\pgfpathlineto{\pgfqpoint{2.259922in}{1.012988in}}%
\pgfpathlineto{\pgfqpoint{2.260929in}{1.029487in}}%
\pgfpathlineto{\pgfqpoint{2.260951in}{1.021237in}}%
\pgfpathlineto{\pgfqpoint{2.261099in}{1.029487in}}%
\pgfpathlineto{\pgfqpoint{2.261432in}{1.004739in}}%
\pgfpathlineto{\pgfqpoint{2.261788in}{1.021237in}}%
\pgfpathlineto{\pgfqpoint{2.263128in}{0.996490in}}%
\pgfpathlineto{\pgfqpoint{2.263780in}{1.029487in}}%
\pgfpathlineto{\pgfqpoint{2.264283in}{1.004739in}}%
\pgfpathlineto{\pgfqpoint{2.264305in}{0.996490in}}%
\pgfpathlineto{\pgfqpoint{2.264787in}{1.054234in}}%
\pgfpathlineto{\pgfqpoint{2.265305in}{1.037736in}}%
\pgfpathlineto{\pgfqpoint{2.266312in}{0.996490in}}%
\pgfpathlineto{\pgfqpoint{2.265475in}{1.045985in}}%
\pgfpathlineto{\pgfqpoint{2.266482in}{1.021237in}}%
\pgfpathlineto{\pgfqpoint{2.267489in}{1.045985in}}%
\pgfpathlineto{\pgfqpoint{2.267652in}{1.037736in}}%
\pgfpathlineto{\pgfqpoint{2.267800in}{1.029487in}}%
\pgfpathlineto{\pgfqpoint{2.267822in}{1.070732in}}%
\pgfpathlineto{\pgfqpoint{2.268156in}{1.070732in}}%
\pgfpathlineto{\pgfqpoint{2.269496in}{1.045985in}}%
\pgfpathlineto{\pgfqpoint{2.270488in}{1.078981in}}%
\pgfpathlineto{\pgfqpoint{2.270673in}{1.062483in}}%
\pgfpathlineto{\pgfqpoint{2.271006in}{1.045985in}}%
\pgfpathlineto{\pgfqpoint{2.271324in}{1.078981in}}%
\pgfpathlineto{\pgfqpoint{2.271339in}{1.070732in}}%
\pgfpathlineto{\pgfqpoint{2.271510in}{1.095480in}}%
\pgfpathlineto{\pgfqpoint{2.272013in}{1.062483in}}%
\pgfpathlineto{\pgfqpoint{2.272346in}{1.045985in}}%
\pgfpathlineto{\pgfqpoint{2.272665in}{1.078981in}}%
\pgfpathlineto{\pgfqpoint{2.272687in}{1.070732in}}%
\pgfpathlineto{\pgfqpoint{2.273524in}{1.095480in}}%
\pgfpathlineto{\pgfqpoint{2.273835in}{1.078981in}}%
\pgfpathlineto{\pgfqpoint{2.274027in}{1.070732in}}%
\pgfpathlineto{\pgfqpoint{2.274338in}{1.103729in}}%
\pgfpathlineto{\pgfqpoint{2.274864in}{1.095480in}}%
\pgfpathlineto{\pgfqpoint{2.276522in}{1.128476in}}%
\pgfpathlineto{\pgfqpoint{2.277544in}{1.111978in}}%
\pgfpathlineto{\pgfqpoint{2.277692in}{1.103729in}}%
\pgfpathlineto{\pgfqpoint{2.278048in}{1.144975in}}%
\pgfpathlineto{\pgfqpoint{2.278366in}{1.153224in}}%
\pgfpathlineto{\pgfqpoint{2.278699in}{1.128476in}}%
\pgfpathlineto{\pgfqpoint{2.279891in}{1.095480in}}%
\pgfpathlineto{\pgfqpoint{2.280728in}{1.070732in}}%
\pgfpathlineto{\pgfqpoint{2.281217in}{1.078981in}}%
\pgfpathlineto{\pgfqpoint{2.281720in}{1.103729in}}%
\pgfpathlineto{\pgfqpoint{2.282238in}{1.095480in}}%
\pgfpathlineto{\pgfqpoint{2.282890in}{1.103729in}}%
\pgfpathlineto{\pgfqpoint{2.283075in}{1.087231in}}%
\pgfpathlineto{\pgfqpoint{2.284756in}{1.045985in}}%
\pgfpathlineto{\pgfqpoint{2.285237in}{1.029487in}}%
\pgfpathlineto{\pgfqpoint{2.285911in}{1.054234in}}%
\pgfpathlineto{\pgfqpoint{2.286259in}{1.021237in}}%
\pgfpathlineto{\pgfqpoint{2.286933in}{1.045985in}}%
\pgfpathlineto{\pgfqpoint{2.287917in}{1.054234in}}%
\pgfpathlineto{\pgfqpoint{2.287940in}{1.045985in}}%
\pgfpathlineto{\pgfqpoint{2.288443in}{1.021237in}}%
\pgfpathlineto{\pgfqpoint{2.288754in}{1.054234in}}%
\pgfpathlineto{\pgfqpoint{2.288776in}{1.045985in}}%
\pgfpathlineto{\pgfqpoint{2.289761in}{1.078981in}}%
\pgfpathlineto{\pgfqpoint{2.289946in}{1.062483in}}%
\pgfpathlineto{\pgfqpoint{2.290938in}{1.054234in}}%
\pgfpathlineto{\pgfqpoint{2.290264in}{1.078981in}}%
\pgfpathlineto{\pgfqpoint{2.290938in}{1.062483in}}%
\pgfpathlineto{\pgfqpoint{2.291627in}{1.045985in}}%
\pgfpathlineto{\pgfqpoint{2.291960in}{1.070732in}}%
\pgfpathlineto{\pgfqpoint{2.292449in}{1.078981in}}%
\pgfpathlineto{\pgfqpoint{2.292634in}{1.062483in}}%
\pgfpathlineto{\pgfqpoint{2.293641in}{1.045985in}}%
\pgfpathlineto{\pgfqpoint{2.292797in}{1.070732in}}%
\pgfpathlineto{\pgfqpoint{2.293789in}{1.054234in}}%
\pgfpathlineto{\pgfqpoint{2.293952in}{1.078981in}}%
\pgfpathlineto{\pgfqpoint{2.294811in}{1.070732in}}%
\pgfpathlineto{\pgfqpoint{2.295129in}{1.078981in}}%
\pgfpathlineto{\pgfqpoint{2.295462in}{1.054234in}}%
\pgfpathlineto{\pgfqpoint{2.295966in}{1.029487in}}%
\pgfpathlineto{\pgfqpoint{2.296484in}{1.045985in}}%
\pgfpathlineto{\pgfqpoint{2.296973in}{1.054234in}}%
\pgfpathlineto{\pgfqpoint{2.297158in}{1.037736in}}%
\pgfpathlineto{\pgfqpoint{2.297491in}{1.021237in}}%
\pgfpathlineto{\pgfqpoint{2.297809in}{1.054234in}}%
\pgfpathlineto{\pgfqpoint{2.297869in}{1.054234in}}%
\pgfpathlineto{\pgfqpoint{2.298483in}{1.103729in}}%
\pgfpathlineto{\pgfqpoint{2.299335in}{1.087231in}}%
\pgfpathlineto{\pgfqpoint{2.300156in}{1.078981in}}%
\pgfpathlineto{\pgfqpoint{2.299505in}{1.095480in}}%
\pgfpathlineto{\pgfqpoint{2.300156in}{1.087231in}}%
\pgfpathlineto{\pgfqpoint{2.300179in}{1.120227in}}%
\pgfpathlineto{\pgfqpoint{2.301178in}{1.095480in}}%
\pgfpathlineto{\pgfqpoint{2.301349in}{1.120227in}}%
\pgfpathlineto{\pgfqpoint{2.302341in}{1.103729in}}%
\pgfpathlineto{\pgfqpoint{2.303022in}{1.120227in}}%
\pgfpathlineto{\pgfqpoint{2.303362in}{1.095480in}}%
\pgfpathlineto{\pgfqpoint{2.304369in}{1.169722in}}%
\pgfpathlineto{\pgfqpoint{2.305206in}{1.144975in}}%
\pgfpathlineto{\pgfqpoint{2.305517in}{1.128476in}}%
\pgfpathlineto{\pgfqpoint{2.305858in}{1.161473in}}%
\pgfpathlineto{\pgfqpoint{2.305872in}{1.194470in}}%
\pgfpathlineto{\pgfqpoint{2.306879in}{1.194470in}}%
\pgfpathlineto{\pgfqpoint{2.308220in}{1.169722in}}%
\pgfpathlineto{\pgfqpoint{2.308723in}{1.194470in}}%
\pgfpathlineto{\pgfqpoint{2.309375in}{1.177971in}}%
\pgfpathlineto{\pgfqpoint{2.309730in}{1.169722in}}%
\pgfpathlineto{\pgfqpoint{2.310234in}{1.219217in}}%
\pgfpathlineto{\pgfqpoint{2.312077in}{1.169722in}}%
\pgfpathlineto{\pgfqpoint{2.312247in}{1.194470in}}%
\pgfpathlineto{\pgfqpoint{2.312914in}{1.161473in}}%
\pgfpathlineto{\pgfqpoint{2.313403in}{1.153224in}}%
\pgfpathlineto{\pgfqpoint{2.313565in}{1.177971in}}%
\pgfpathlineto{\pgfqpoint{2.313588in}{1.169722in}}%
\pgfpathlineto{\pgfqpoint{2.314069in}{1.202719in}}%
\pgfpathlineto{\pgfqpoint{2.314743in}{1.177971in}}%
\pgfpathlineto{\pgfqpoint{2.314758in}{1.169722in}}%
\pgfpathlineto{\pgfqpoint{2.315594in}{1.243965in}}%
\pgfpathlineto{\pgfqpoint{2.315742in}{1.252214in}}%
\pgfpathlineto{\pgfqpoint{2.316098in}{1.219217in}}%
\pgfpathlineto{\pgfqpoint{2.316438in}{1.235715in}}%
\pgfpathlineto{\pgfqpoint{2.316920in}{1.227466in}}%
\pgfpathlineto{\pgfqpoint{2.317090in}{1.252214in}}%
\pgfpathlineto{\pgfqpoint{2.317275in}{1.243965in}}%
\pgfpathlineto{\pgfqpoint{2.318430in}{1.276961in}}%
\pgfpathlineto{\pgfqpoint{2.318615in}{1.260463in}}%
\pgfpathlineto{\pgfqpoint{2.319267in}{1.227466in}}%
\pgfpathlineto{\pgfqpoint{2.319600in}{1.260463in}}%
\pgfpathlineto{\pgfqpoint{2.319622in}{1.268712in}}%
\pgfpathlineto{\pgfqpoint{2.320288in}{1.219217in}}%
\pgfpathlineto{\pgfqpoint{2.320629in}{1.243965in}}%
\pgfpathlineto{\pgfqpoint{2.322280in}{1.276961in}}%
\pgfpathlineto{\pgfqpoint{2.322473in}{1.268712in}}%
\pgfpathlineto{\pgfqpoint{2.322784in}{1.301709in}}%
\pgfpathlineto{\pgfqpoint{2.323309in}{1.293460in}}%
\pgfpathlineto{\pgfqpoint{2.323791in}{1.301709in}}%
\pgfpathlineto{\pgfqpoint{2.324124in}{1.276961in}}%
\pgfpathlineto{\pgfqpoint{2.324650in}{1.268712in}}%
\pgfpathlineto{\pgfqpoint{2.324975in}{1.301709in}}%
\pgfpathlineto{\pgfqpoint{2.325153in}{1.293460in}}%
\pgfpathlineto{\pgfqpoint{2.326160in}{1.219217in}}%
\pgfpathlineto{\pgfqpoint{2.326997in}{1.243965in}}%
\pgfpathlineto{\pgfqpoint{2.327367in}{1.227466in}}%
\pgfpathlineto{\pgfqpoint{2.328152in}{1.252214in}}%
\pgfpathlineto{\pgfqpoint{2.329159in}{1.227466in}}%
\pgfpathlineto{\pgfqpoint{2.328507in}{1.268712in}}%
\pgfpathlineto{\pgfqpoint{2.329173in}{1.243965in}}%
\pgfpathlineto{\pgfqpoint{2.330832in}{1.202719in}}%
\pgfpathlineto{\pgfqpoint{2.330832in}{1.210968in}}%
\pgfpathlineto{\pgfqpoint{2.332009in}{1.276961in}}%
\pgfpathlineto{\pgfqpoint{2.331165in}{1.202719in}}%
\pgfpathlineto{\pgfqpoint{2.332024in}{1.268712in}}%
\pgfpathlineto{\pgfqpoint{2.333201in}{1.219217in}}%
\pgfpathlineto{\pgfqpoint{2.333364in}{1.243965in}}%
\pgfpathlineto{\pgfqpoint{2.333683in}{1.252214in}}%
\pgfpathlineto{\pgfqpoint{2.334023in}{1.227466in}}%
\pgfpathlineto{\pgfqpoint{2.334542in}{1.219217in}}%
\pgfpathlineto{\pgfqpoint{2.334852in}{1.252214in}}%
\pgfpathlineto{\pgfqpoint{2.335045in}{1.243965in}}%
\pgfpathlineto{\pgfqpoint{2.337052in}{1.318207in}}%
\pgfpathlineto{\pgfqpoint{2.337222in}{1.309958in}}%
\pgfpathlineto{\pgfqpoint{2.337725in}{1.268712in}}%
\pgfpathlineto{\pgfqpoint{2.338392in}{1.293460in}}%
\pgfpathlineto{\pgfqpoint{2.339236in}{1.342955in}}%
\pgfpathlineto{\pgfqpoint{2.340065in}{1.326456in}}%
\pgfpathlineto{\pgfqpoint{2.340739in}{1.293460in}}%
\pgfpathlineto{\pgfqpoint{2.341079in}{1.342955in}}%
\pgfpathlineto{\pgfqpoint{2.342064in}{1.276961in}}%
\pgfpathlineto{\pgfqpoint{2.342923in}{1.293460in}}%
\pgfpathlineto{\pgfqpoint{2.345078in}{1.375951in}}%
\pgfpathlineto{\pgfqpoint{2.345100in}{1.367702in}}%
\pgfpathlineto{\pgfqpoint{2.346440in}{1.342955in}}%
\pgfpathlineto{\pgfqpoint{2.346944in}{1.367702in}}%
\pgfpathlineto{\pgfqpoint{2.347277in}{1.342955in}}%
\pgfpathlineto{\pgfqpoint{2.347595in}{1.326456in}}%
\pgfpathlineto{\pgfqpoint{2.347617in}{1.367702in}}%
\pgfpathlineto{\pgfqpoint{2.349291in}{1.417197in}}%
\pgfpathlineto{\pgfqpoint{2.349461in}{1.408948in}}%
\pgfpathlineto{\pgfqpoint{2.350468in}{1.392449in}}%
\pgfpathlineto{\pgfqpoint{2.350616in}{1.400699in}}%
\pgfpathlineto{\pgfqpoint{2.351638in}{1.417197in}}%
\pgfpathlineto{\pgfqpoint{2.352289in}{1.375951in}}%
\pgfpathlineto{\pgfqpoint{2.352978in}{1.392449in}}%
\pgfpathlineto{\pgfqpoint{2.353630in}{1.400699in}}%
\pgfpathlineto{\pgfqpoint{2.353815in}{1.384200in}}%
\pgfpathlineto{\pgfqpoint{2.354303in}{1.375951in}}%
\pgfpathlineto{\pgfqpoint{2.354481in}{1.400699in}}%
\pgfpathlineto{\pgfqpoint{2.354488in}{1.392449in}}%
\pgfpathlineto{\pgfqpoint{2.355814in}{1.425446in}}%
\pgfpathlineto{\pgfqpoint{2.354970in}{1.375951in}}%
\pgfpathlineto{\pgfqpoint{2.355829in}{1.417197in}}%
\pgfpathlineto{\pgfqpoint{2.357843in}{1.342955in}}%
\pgfpathlineto{\pgfqpoint{2.357998in}{1.359453in}}%
\pgfpathlineto{\pgfqpoint{2.358005in}{1.367702in}}%
\pgfpathlineto{\pgfqpoint{2.359012in}{1.309958in}}%
\pgfpathlineto{\pgfqpoint{2.360693in}{1.219217in}}%
\pgfpathlineto{\pgfqpoint{2.361019in}{1.227466in}}%
\pgfpathlineto{\pgfqpoint{2.361345in}{1.202719in}}%
\pgfpathlineto{\pgfqpoint{2.362366in}{1.194470in}}%
\pgfpathlineto{\pgfqpoint{2.366372in}{1.029487in}}%
\pgfpathlineto{\pgfqpoint{2.366387in}{1.045985in}}%
\pgfpathlineto{\pgfqpoint{2.366542in}{1.054234in}}%
\pgfpathlineto{\pgfqpoint{2.366728in}{1.037736in}}%
\pgfpathlineto{\pgfqpoint{2.368045in}{0.979992in}}%
\pgfpathlineto{\pgfqpoint{2.368053in}{0.988241in}}%
\pgfpathlineto{\pgfqpoint{2.368556in}{1.029487in}}%
\pgfpathlineto{\pgfqpoint{2.369075in}{0.996490in}}%
\pgfpathlineto{\pgfqpoint{2.373251in}{0.732517in}}%
\pgfpathlineto{\pgfqpoint{2.374265in}{0.724268in}}%
\pgfpathlineto{\pgfqpoint{2.376449in}{0.650025in}}%
\pgfpathlineto{\pgfqpoint{2.376605in}{0.658274in}}%
\pgfpathlineto{\pgfqpoint{2.376953in}{0.650025in}}%
\pgfpathlineto{\pgfqpoint{2.377619in}{0.699520in}}%
\pgfpathlineto{\pgfqpoint{2.378604in}{0.658274in}}%
\pgfpathlineto{\pgfqpoint{2.378833in}{0.683022in}}%
\pgfpathlineto{\pgfqpoint{2.380640in}{0.782012in}}%
\pgfpathlineto{\pgfqpoint{2.380973in}{0.765513in}}%
\pgfpathlineto{\pgfqpoint{2.381817in}{0.732517in}}%
\pgfpathlineto{\pgfqpoint{2.381144in}{0.773763in}}%
\pgfpathlineto{\pgfqpoint{2.382143in}{0.749015in}}%
\pgfpathlineto{\pgfqpoint{2.382306in}{0.757264in}}%
\pgfpathlineto{\pgfqpoint{2.382484in}{0.732517in}}%
\pgfpathlineto{\pgfqpoint{2.382817in}{0.740766in}}%
\pgfpathlineto{\pgfqpoint{2.383491in}{0.707769in}}%
\pgfpathlineto{\pgfqpoint{2.383654in}{0.749015in}}%
\pgfpathlineto{\pgfqpoint{2.383824in}{0.732517in}}%
\pgfpathlineto{\pgfqpoint{2.383994in}{0.773763in}}%
\pgfpathlineto{\pgfqpoint{2.384661in}{0.724268in}}%
\pgfpathlineto{\pgfqpoint{2.384816in}{0.732517in}}%
\pgfpathlineto{\pgfqpoint{2.385667in}{0.724268in}}%
\pgfpathlineto{\pgfqpoint{2.384994in}{0.749015in}}%
\pgfpathlineto{\pgfqpoint{2.385838in}{0.749015in}}%
\pgfpathlineto{\pgfqpoint{2.387008in}{0.724268in}}%
\pgfpathlineto{\pgfqpoint{2.386319in}{0.757264in}}%
\pgfpathlineto{\pgfqpoint{2.387008in}{0.732517in}}%
\pgfpathlineto{\pgfqpoint{2.387511in}{0.757264in}}%
\pgfpathlineto{\pgfqpoint{2.387852in}{0.732517in}}%
\pgfpathlineto{\pgfqpoint{2.388518in}{0.724268in}}%
\pgfpathlineto{\pgfqpoint{2.388674in}{0.740766in}}%
\pgfpathlineto{\pgfqpoint{2.388681in}{0.773763in}}%
\pgfpathlineto{\pgfqpoint{2.389355in}{0.699520in}}%
\pgfpathlineto{\pgfqpoint{2.389688in}{0.699520in}}%
\pgfpathlineto{\pgfqpoint{2.390191in}{0.707769in}}%
\pgfpathlineto{\pgfqpoint{2.391702in}{0.625278in}}%
\pgfpathlineto{\pgfqpoint{2.392879in}{0.633527in}}%
\pgfpathlineto{\pgfqpoint{2.393198in}{0.625278in}}%
\pgfpathlineto{\pgfqpoint{2.393886in}{0.625278in}}%
\pgfpathlineto{\pgfqpoint{2.395049in}{0.633527in}}%
\pgfpathlineto{\pgfqpoint{2.395374in}{0.625278in}}%
\pgfpathlineto{\pgfqpoint{2.396063in}{0.625278in}}%
\pgfpathlineto{\pgfqpoint{2.397225in}{0.633527in}}%
\pgfpathlineto{\pgfqpoint{2.397559in}{0.625278in}}%
\pgfpathlineto{\pgfqpoint{2.398225in}{0.625278in}}%
\pgfpathlineto{\pgfqpoint{2.399232in}{0.633527in}}%
\pgfpathlineto{\pgfqpoint{2.399232in}{0.625278in}}%
\pgfpathlineto{\pgfqpoint{2.400424in}{0.633527in}}%
\pgfpathlineto{\pgfqpoint{2.400572in}{0.625278in}}%
\pgfpathlineto{\pgfqpoint{2.401927in}{0.633527in}}%
\pgfpathlineto{\pgfqpoint{2.402268in}{0.625278in}}%
\pgfpathlineto{\pgfqpoint{2.402771in}{0.625278in}}%
\pgfpathlineto{\pgfqpoint{2.403941in}{0.633527in}}%
\pgfpathlineto{\pgfqpoint{2.404097in}{0.625278in}}%
\pgfpathlineto{\pgfqpoint{2.404948in}{0.625278in}}%
\pgfpathlineto{\pgfqpoint{2.406110in}{0.633527in}}%
\pgfpathlineto{\pgfqpoint{2.406273in}{0.625278in}}%
\pgfpathlineto{\pgfqpoint{2.407125in}{0.625278in}}%
\pgfpathlineto{\pgfqpoint{2.408280in}{0.633527in}}%
\pgfpathlineto{\pgfqpoint{2.408465in}{0.625278in}}%
\pgfpathlineto{\pgfqpoint{2.409139in}{0.625278in}}%
\pgfpathlineto{\pgfqpoint{2.410301in}{0.633527in}}%
\pgfpathlineto{\pgfqpoint{2.410479in}{0.625278in}}%
\pgfpathlineto{\pgfqpoint{2.411316in}{0.625278in}}%
\pgfpathlineto{\pgfqpoint{2.412478in}{0.633527in}}%
\pgfpathlineto{\pgfqpoint{2.412648in}{0.625278in}}%
\pgfpathlineto{\pgfqpoint{2.413492in}{0.625278in}}%
\pgfpathlineto{\pgfqpoint{2.414655in}{0.633527in}}%
\pgfpathlineto{\pgfqpoint{2.414670in}{0.625278in}}%
\pgfpathlineto{\pgfqpoint{2.415654in}{0.625278in}}%
\pgfpathlineto{\pgfqpoint{2.416832in}{0.633527in}}%
\pgfpathlineto{\pgfqpoint{2.416995in}{0.625278in}}%
\pgfpathlineto{\pgfqpoint{2.417839in}{0.625278in}}%
\pgfpathlineto{\pgfqpoint{2.419179in}{0.633527in}}%
\pgfpathlineto{\pgfqpoint{2.419186in}{0.625278in}}%
\pgfpathlineto{\pgfqpoint{2.420186in}{0.625278in}}%
\pgfpathlineto{\pgfqpoint{2.421208in}{0.633527in}}%
\pgfpathlineto{\pgfqpoint{2.421208in}{0.625278in}}%
\pgfpathlineto{\pgfqpoint{2.422363in}{0.633527in}}%
\pgfpathlineto{\pgfqpoint{2.422548in}{0.625278in}}%
\pgfpathlineto{\pgfqpoint{2.423222in}{0.625278in}}%
\pgfpathlineto{\pgfqpoint{2.424214in}{0.633527in}}%
\pgfpathlineto{\pgfqpoint{2.424214in}{0.625278in}}%
\pgfpathlineto{\pgfqpoint{2.425384in}{0.633527in}}%
\pgfpathlineto{\pgfqpoint{2.425546in}{0.625278in}}%
\pgfpathlineto{\pgfqpoint{2.426405in}{0.625278in}}%
\pgfpathlineto{\pgfqpoint{2.427553in}{0.633527in}}%
\pgfpathlineto{\pgfqpoint{2.427553in}{0.625278in}}%
\pgfpathlineto{\pgfqpoint{2.428560in}{0.625278in}}%
\pgfpathlineto{\pgfqpoint{2.429567in}{0.633527in}}%
\pgfpathlineto{\pgfqpoint{2.429567in}{0.625278in}}%
\pgfpathlineto{\pgfqpoint{2.430915in}{0.633527in}}%
\pgfpathlineto{\pgfqpoint{2.430929in}{0.625278in}}%
\pgfpathlineto{\pgfqpoint{2.431744in}{0.625278in}}%
\pgfpathlineto{\pgfqpoint{2.433084in}{0.633527in}}%
\pgfpathlineto{\pgfqpoint{2.433254in}{0.625278in}}%
\pgfpathlineto{\pgfqpoint{2.434113in}{0.625278in}}%
\pgfpathlineto{\pgfqpoint{2.435098in}{0.633527in}}%
\pgfpathlineto{\pgfqpoint{2.435098in}{0.625278in}}%
\pgfpathlineto{\pgfqpoint{2.436438in}{0.633527in}}%
\pgfpathlineto{\pgfqpoint{2.436438in}{0.625278in}}%
\pgfpathlineto{\pgfqpoint{2.437297in}{0.625278in}}%
\pgfpathlineto{\pgfqpoint{2.438452in}{0.633527in}}%
\pgfpathlineto{\pgfqpoint{2.438622in}{0.625278in}}%
\pgfpathlineto{\pgfqpoint{2.439459in}{0.625278in}}%
\pgfpathlineto{\pgfqpoint{2.440473in}{0.633527in}}%
\pgfpathlineto{\pgfqpoint{2.440473in}{0.625278in}}%
\pgfpathlineto{\pgfqpoint{2.441636in}{0.633527in}}%
\pgfpathlineto{\pgfqpoint{2.441658in}{0.625278in}}%
\pgfpathlineto{\pgfqpoint{2.442643in}{0.625278in}}%
\pgfpathlineto{\pgfqpoint{2.443820in}{0.633527in}}%
\pgfpathlineto{\pgfqpoint{2.444005in}{0.625278in}}%
\pgfpathlineto{\pgfqpoint{2.444842in}{0.625278in}}%
\pgfpathlineto{\pgfqpoint{2.446160in}{0.633527in}}%
\pgfpathlineto{\pgfqpoint{2.446500in}{0.625278in}}%
\pgfpathlineto{\pgfqpoint{2.447167in}{0.625278in}}%
\pgfpathlineto{\pgfqpoint{2.448344in}{0.633527in}}%
\pgfpathlineto{\pgfqpoint{2.448507in}{0.625278in}}%
\pgfpathlineto{\pgfqpoint{2.449181in}{0.625278in}}%
\pgfpathlineto{\pgfqpoint{2.450351in}{0.633527in}}%
\pgfpathlineto{\pgfqpoint{2.450521in}{0.625278in}}%
\pgfpathlineto{\pgfqpoint{2.451357in}{0.625278in}}%
\pgfpathlineto{\pgfqpoint{2.452364in}{0.633527in}}%
\pgfpathlineto{\pgfqpoint{2.452364in}{0.625278in}}%
\pgfpathlineto{\pgfqpoint{2.453542in}{0.633527in}}%
\pgfpathlineto{\pgfqpoint{2.453875in}{0.625278in}}%
\pgfpathlineto{\pgfqpoint{2.454208in}{0.625278in}}%
\pgfpathlineto{\pgfqpoint{2.455385in}{0.633527in}}%
\pgfpathlineto{\pgfqpoint{2.455719in}{0.625278in}}%
\pgfpathlineto{\pgfqpoint{2.456385in}{0.625278in}}%
\pgfpathlineto{\pgfqpoint{2.457562in}{0.633527in}}%
\pgfpathlineto{\pgfqpoint{2.457740in}{0.625278in}}%
\pgfpathlineto{\pgfqpoint{2.458569in}{0.625278in}}%
\pgfpathlineto{\pgfqpoint{2.459746in}{0.633527in}}%
\pgfpathlineto{\pgfqpoint{2.459909in}{0.625278in}}%
\pgfpathlineto{\pgfqpoint{2.460753in}{0.625278in}}%
\pgfpathlineto{\pgfqpoint{2.461923in}{0.633527in}}%
\pgfpathlineto{\pgfqpoint{2.462086in}{0.625278in}}%
\pgfpathlineto{\pgfqpoint{2.462923in}{0.625278in}}%
\pgfpathlineto{\pgfqpoint{2.464100in}{0.633527in}}%
\pgfpathlineto{\pgfqpoint{2.464270in}{0.625278in}}%
\pgfpathlineto{\pgfqpoint{2.464937in}{0.625278in}}%
\pgfpathlineto{\pgfqpoint{2.466114in}{0.633527in}}%
\pgfpathlineto{\pgfqpoint{2.466277in}{0.625278in}}%
\pgfpathlineto{\pgfqpoint{2.467069in}{0.625278in}}%
\pgfpathlineto{\pgfqpoint{2.468291in}{0.633527in}}%
\pgfpathlineto{\pgfqpoint{2.468461in}{0.625278in}}%
\pgfpathlineto{\pgfqpoint{2.469298in}{0.625278in}}%
\pgfpathlineto{\pgfqpoint{2.470468in}{0.633527in}}%
\pgfpathlineto{\pgfqpoint{2.470638in}{0.625278in}}%
\pgfpathlineto{\pgfqpoint{2.471304in}{0.625278in}}%
\pgfpathlineto{\pgfqpoint{2.472652in}{0.633527in}}%
\pgfpathlineto{\pgfqpoint{2.472985in}{0.625278in}}%
\pgfpathlineto{\pgfqpoint{2.473652in}{0.625278in}}%
\pgfpathlineto{\pgfqpoint{2.474658in}{0.633527in}}%
\pgfpathlineto{\pgfqpoint{2.474658in}{0.625278in}}%
\pgfpathlineto{\pgfqpoint{2.475836in}{0.633527in}}%
\pgfpathlineto{\pgfqpoint{2.475999in}{0.625278in}}%
\pgfpathlineto{\pgfqpoint{2.476843in}{0.625278in}}%
\pgfpathlineto{\pgfqpoint{2.478013in}{0.633527in}}%
\pgfpathlineto{\pgfqpoint{2.478183in}{0.625278in}}%
\pgfpathlineto{\pgfqpoint{2.478849in}{0.625278in}}%
\pgfpathlineto{\pgfqpoint{2.480189in}{0.633527in}}%
\pgfpathlineto{\pgfqpoint{2.480360in}{0.625278in}}%
\pgfpathlineto{\pgfqpoint{2.481034in}{0.625278in}}%
\pgfpathlineto{\pgfqpoint{2.482203in}{0.633527in}}%
\pgfpathlineto{\pgfqpoint{2.482374in}{0.625278in}}%
\pgfpathlineto{\pgfqpoint{2.483210in}{0.625278in}}%
\pgfpathlineto{\pgfqpoint{2.484380in}{0.633527in}}%
\pgfpathlineto{\pgfqpoint{2.484551in}{0.625278in}}%
\pgfpathlineto{\pgfqpoint{2.485387in}{0.625278in}}%
\pgfpathlineto{\pgfqpoint{2.486727in}{0.633527in}}%
\pgfpathlineto{\pgfqpoint{2.486727in}{0.625278in}}%
\pgfpathlineto{\pgfqpoint{2.487734in}{0.625278in}}%
\pgfpathlineto{\pgfqpoint{2.488912in}{0.633527in}}%
\pgfpathlineto{\pgfqpoint{2.489074in}{0.625278in}}%
\pgfpathlineto{\pgfqpoint{2.489748in}{0.625278in}}%
\pgfpathlineto{\pgfqpoint{2.490918in}{0.633527in}}%
\pgfpathlineto{\pgfqpoint{2.491088in}{0.625278in}}%
\pgfpathlineto{\pgfqpoint{2.491925in}{0.625278in}}%
\pgfpathlineto{\pgfqpoint{2.493102in}{0.633527in}}%
\pgfpathlineto{\pgfqpoint{2.493265in}{0.625278in}}%
\pgfpathlineto{\pgfqpoint{2.494102in}{0.625278in}}%
\pgfpathlineto{\pgfqpoint{2.495279in}{0.633527in}}%
\pgfpathlineto{\pgfqpoint{2.495449in}{0.625278in}}%
\pgfpathlineto{\pgfqpoint{2.496286in}{0.625278in}}%
\pgfpathlineto{\pgfqpoint{2.497626in}{0.633527in}}%
\pgfpathlineto{\pgfqpoint{2.497959in}{0.625278in}}%
\pgfpathlineto{\pgfqpoint{2.498633in}{0.625278in}}%
\pgfpathlineto{\pgfqpoint{2.499973in}{0.633527in}}%
\pgfpathlineto{\pgfqpoint{2.499973in}{0.625278in}}%
\pgfpathlineto{\pgfqpoint{2.500980in}{0.625278in}}%
\pgfpathlineto{\pgfqpoint{2.502150in}{0.633527in}}%
\pgfpathlineto{\pgfqpoint{2.502321in}{0.625278in}}%
\pgfpathlineto{\pgfqpoint{2.503157in}{0.625278in}}%
\pgfpathlineto{\pgfqpoint{2.504327in}{0.633527in}}%
\pgfpathlineto{\pgfqpoint{2.504497in}{0.625278in}}%
\pgfpathlineto{\pgfqpoint{2.505334in}{0.625278in}}%
\pgfpathlineto{\pgfqpoint{2.506341in}{0.633527in}}%
\pgfpathlineto{\pgfqpoint{2.506341in}{0.625278in}}%
\pgfpathlineto{\pgfqpoint{2.507681in}{0.633527in}}%
\pgfpathlineto{\pgfqpoint{2.507852in}{0.625278in}}%
\pgfpathlineto{\pgfqpoint{2.508688in}{0.625278in}}%
\pgfpathlineto{\pgfqpoint{2.509865in}{0.633527in}}%
\pgfpathlineto{\pgfqpoint{2.510199in}{0.625278in}}%
\pgfpathlineto{\pgfqpoint{2.510865in}{0.625278in}}%
\pgfpathlineto{\pgfqpoint{2.512042in}{0.633527in}}%
\pgfpathlineto{\pgfqpoint{2.512213in}{0.625278in}}%
\pgfpathlineto{\pgfqpoint{2.513049in}{0.625278in}}%
\pgfpathlineto{\pgfqpoint{2.514219in}{0.633527in}}%
\pgfpathlineto{\pgfqpoint{2.514219in}{0.625278in}}%
\pgfpathlineto{\pgfqpoint{2.515226in}{0.625278in}}%
\pgfpathlineto{\pgfqpoint{2.516403in}{0.633527in}}%
\pgfpathlineto{\pgfqpoint{2.516403in}{0.625278in}}%
\pgfpathlineto{\pgfqpoint{2.517403in}{0.625278in}}%
\pgfpathlineto{\pgfqpoint{2.518750in}{0.633527in}}%
\pgfpathlineto{\pgfqpoint{2.518750in}{0.625278in}}%
\pgfpathlineto{\pgfqpoint{2.519750in}{0.625278in}}%
\pgfpathlineto{\pgfqpoint{2.520927in}{0.633527in}}%
\pgfpathlineto{\pgfqpoint{2.521090in}{0.625278in}}%
\pgfpathlineto{\pgfqpoint{2.521934in}{0.625278in}}%
\pgfpathlineto{\pgfqpoint{2.523104in}{0.633527in}}%
\pgfpathlineto{\pgfqpoint{2.523274in}{0.625278in}}%
\pgfpathlineto{\pgfqpoint{2.524111in}{0.625278in}}%
\pgfpathlineto{\pgfqpoint{2.525281in}{0.633527in}}%
\pgfpathlineto{\pgfqpoint{2.525451in}{0.625278in}}%
\pgfpathlineto{\pgfqpoint{2.526288in}{0.625278in}}%
\pgfpathlineto{\pgfqpoint{2.527465in}{0.633527in}}%
\pgfpathlineto{\pgfqpoint{2.527628in}{0.625278in}}%
\pgfpathlineto{\pgfqpoint{2.528472in}{0.625278in}}%
\pgfpathlineto{\pgfqpoint{2.529642in}{0.633527in}}%
\pgfpathlineto{\pgfqpoint{2.529642in}{0.625278in}}%
\pgfpathlineto{\pgfqpoint{2.530649in}{0.625278in}}%
\pgfpathlineto{\pgfqpoint{2.531989in}{0.633527in}}%
\pgfpathlineto{\pgfqpoint{2.531989in}{0.625278in}}%
\pgfpathlineto{\pgfqpoint{2.532996in}{0.625278in}}%
\pgfpathlineto{\pgfqpoint{2.534166in}{0.633527in}}%
\pgfpathlineto{\pgfqpoint{2.534507in}{0.625278in}}%
\pgfpathlineto{\pgfqpoint{2.535173in}{0.625278in}}%
\pgfpathlineto{\pgfqpoint{2.536350in}{0.633527in}}%
\pgfpathlineto{\pgfqpoint{2.536513in}{0.625278in}}%
\pgfpathlineto{\pgfqpoint{2.537357in}{0.625278in}}%
\pgfpathlineto{\pgfqpoint{2.538527in}{0.633527in}}%
\pgfpathlineto{\pgfqpoint{2.538697in}{0.625278in}}%
\pgfpathlineto{\pgfqpoint{2.539534in}{0.625278in}}%
\pgfpathlineto{\pgfqpoint{2.540704in}{0.633527in}}%
\pgfpathlineto{\pgfqpoint{2.540874in}{0.625278in}}%
\pgfpathlineto{\pgfqpoint{2.541711in}{0.625278in}}%
\pgfpathlineto{\pgfqpoint{2.542888in}{0.633527in}}%
\pgfpathlineto{\pgfqpoint{2.543051in}{0.625278in}}%
\pgfpathlineto{\pgfqpoint{2.543895in}{0.625278in}}%
\pgfpathlineto{\pgfqpoint{2.545065in}{0.633527in}}%
\pgfpathlineto{\pgfqpoint{2.545398in}{0.625278in}}%
\pgfpathlineto{\pgfqpoint{2.546072in}{0.625278in}}%
\pgfpathlineto{\pgfqpoint{2.547412in}{0.633527in}}%
\pgfpathlineto{\pgfqpoint{2.547412in}{0.625278in}}%
\pgfpathlineto{\pgfqpoint{2.548419in}{0.625278in}}%
\pgfpathlineto{\pgfqpoint{2.549759in}{0.633527in}}%
\pgfpathlineto{\pgfqpoint{2.549759in}{0.625278in}}%
\pgfpathlineto{\pgfqpoint{2.550766in}{0.625278in}}%
\pgfpathlineto{\pgfqpoint{2.551936in}{0.633527in}}%
\pgfpathlineto{\pgfqpoint{2.552106in}{0.625278in}}%
\pgfpathlineto{\pgfqpoint{2.552943in}{0.625278in}}%
\pgfpathlineto{\pgfqpoint{2.554120in}{0.633527in}}%
\pgfpathlineto{\pgfqpoint{2.554283in}{0.625278in}}%
\pgfpathlineto{\pgfqpoint{2.555120in}{0.625278in}}%
\pgfpathlineto{\pgfqpoint{2.556297in}{0.633527in}}%
\pgfpathlineto{\pgfqpoint{2.556630in}{0.625278in}}%
\pgfpathlineto{\pgfqpoint{2.557304in}{0.625278in}}%
\pgfpathlineto{\pgfqpoint{2.558644in}{0.633527in}}%
\pgfpathlineto{\pgfqpoint{2.558807in}{0.625278in}}%
\pgfpathlineto{\pgfqpoint{2.559651in}{0.625278in}}%
\pgfpathlineto{\pgfqpoint{2.560991in}{0.633527in}}%
\pgfpathlineto{\pgfqpoint{2.560991in}{0.625278in}}%
\pgfpathlineto{\pgfqpoint{2.561998in}{0.625278in}}%
\pgfpathlineto{\pgfqpoint{2.562998in}{0.633527in}}%
\pgfpathlineto{\pgfqpoint{2.562998in}{0.625278in}}%
\pgfpathlineto{\pgfqpoint{2.564346in}{0.633527in}}%
\pgfpathlineto{\pgfqpoint{2.564360in}{0.625278in}}%
\pgfpathlineto{\pgfqpoint{2.565345in}{0.625278in}}%
\pgfpathlineto{\pgfqpoint{2.566352in}{0.633527in}}%
\pgfpathlineto{\pgfqpoint{2.566352in}{0.625278in}}%
\pgfpathlineto{\pgfqpoint{2.567692in}{0.633527in}}%
\pgfpathlineto{\pgfqpoint{2.567863in}{0.625278in}}%
\pgfpathlineto{\pgfqpoint{2.568699in}{0.625278in}}%
\pgfpathlineto{\pgfqpoint{2.569876in}{0.633527in}}%
\pgfpathlineto{\pgfqpoint{2.570039in}{0.625278in}}%
\pgfpathlineto{\pgfqpoint{2.570883in}{0.625278in}}%
\pgfpathlineto{\pgfqpoint{2.572053in}{0.633527in}}%
\pgfpathlineto{\pgfqpoint{2.572224in}{0.625278in}}%
\pgfpathlineto{\pgfqpoint{2.573075in}{0.650025in}}%
\pgfpathlineto{\pgfqpoint{2.573393in}{0.658274in}}%
\pgfpathlineto{\pgfqpoint{2.573727in}{0.633527in}}%
\pgfpathlineto{\pgfqpoint{2.573919in}{0.625278in}}%
\pgfpathlineto{\pgfqpoint{2.574230in}{0.658274in}}%
\pgfpathlineto{\pgfqpoint{2.574586in}{0.641776in}}%
\pgfpathlineto{\pgfqpoint{2.575926in}{0.625278in}}%
\pgfpathlineto{\pgfqpoint{2.577081in}{0.658274in}}%
\pgfpathlineto{\pgfqpoint{2.577266in}{0.641776in}}%
\pgfpathlineto{\pgfqpoint{2.577421in}{0.633527in}}%
\pgfpathlineto{\pgfqpoint{2.577755in}{0.666524in}}%
\pgfpathlineto{\pgfqpoint{2.578088in}{0.683022in}}%
\pgfpathlineto{\pgfqpoint{2.578421in}{0.658274in}}%
\pgfpathlineto{\pgfqpoint{2.578776in}{0.674773in}}%
\pgfpathlineto{\pgfqpoint{2.579761in}{0.658274in}}%
\pgfpathlineto{\pgfqpoint{2.579783in}{0.674773in}}%
\pgfpathlineto{\pgfqpoint{2.580938in}{0.707769in}}%
\pgfpathlineto{\pgfqpoint{2.580961in}{0.699520in}}%
\pgfpathlineto{\pgfqpoint{2.582967in}{0.773763in}}%
\pgfpathlineto{\pgfqpoint{2.583137in}{0.765513in}}%
\pgfpathlineto{\pgfqpoint{2.584144in}{0.749015in}}%
\pgfpathlineto{\pgfqpoint{2.584292in}{0.757264in}}%
\pgfpathlineto{\pgfqpoint{2.584455in}{0.782012in}}%
\pgfpathlineto{\pgfqpoint{2.585314in}{0.773763in}}%
\pgfpathlineto{\pgfqpoint{2.586640in}{0.806759in}}%
\pgfpathlineto{\pgfqpoint{2.585966in}{0.757264in}}%
\pgfpathlineto{\pgfqpoint{2.586654in}{0.798510in}}%
\pgfpathlineto{\pgfqpoint{2.587995in}{0.773763in}}%
\pgfpathlineto{\pgfqpoint{2.589994in}{0.831507in}}%
\pgfpathlineto{\pgfqpoint{2.590342in}{0.823258in}}%
\pgfpathlineto{\pgfqpoint{2.590660in}{0.856254in}}%
\pgfpathlineto{\pgfqpoint{2.591015in}{0.848005in}}%
\pgfpathlineto{\pgfqpoint{2.591519in}{0.872753in}}%
\pgfpathlineto{\pgfqpoint{2.591852in}{0.848005in}}%
\pgfpathlineto{\pgfqpoint{2.593533in}{0.798510in}}%
\pgfpathlineto{\pgfqpoint{2.593681in}{0.806759in}}%
\pgfpathlineto{\pgfqpoint{2.594680in}{0.856254in}}%
\pgfpathlineto{\pgfqpoint{2.594703in}{0.848005in}}%
\pgfpathlineto{\pgfqpoint{2.595525in}{0.856254in}}%
\pgfpathlineto{\pgfqpoint{2.595539in}{0.848005in}}%
\pgfpathlineto{\pgfqpoint{2.596043in}{0.823258in}}%
\pgfpathlineto{\pgfqpoint{2.596361in}{0.856254in}}%
\pgfpathlineto{\pgfqpoint{2.596376in}{0.848005in}}%
\pgfpathlineto{\pgfqpoint{2.596694in}{0.856254in}}%
\pgfpathlineto{\pgfqpoint{2.597028in}{0.831507in}}%
\pgfpathlineto{\pgfqpoint{2.597220in}{0.839756in}}%
\pgfpathlineto{\pgfqpoint{2.598205in}{0.831507in}}%
\pgfpathlineto{\pgfqpoint{2.597531in}{0.856254in}}%
\pgfpathlineto{\pgfqpoint{2.598205in}{0.839756in}}%
\pgfpathlineto{\pgfqpoint{2.598538in}{0.831507in}}%
\pgfpathlineto{\pgfqpoint{2.599227in}{0.848005in}}%
\pgfpathlineto{\pgfqpoint{2.600404in}{0.823258in}}%
\pgfpathlineto{\pgfqpoint{2.600552in}{0.831507in}}%
\pgfpathlineto{\pgfqpoint{2.601574in}{0.897500in}}%
\pgfpathlineto{\pgfqpoint{2.602410in}{0.946995in}}%
\pgfpathlineto{\pgfqpoint{2.603232in}{0.930497in}}%
\pgfpathlineto{\pgfqpoint{2.604254in}{0.922247in}}%
\pgfpathlineto{\pgfqpoint{2.605913in}{0.955244in}}%
\pgfpathlineto{\pgfqpoint{2.606601in}{0.922247in}}%
\pgfpathlineto{\pgfqpoint{2.606942in}{0.946995in}}%
\pgfpathlineto{\pgfqpoint{2.608260in}{0.979992in}}%
\pgfpathlineto{\pgfqpoint{2.608282in}{0.971742in}}%
\pgfpathlineto{\pgfqpoint{2.610126in}{0.922247in}}%
\pgfpathlineto{\pgfqpoint{2.611110in}{0.955244in}}%
\pgfpathlineto{\pgfqpoint{2.611296in}{0.938746in}}%
\pgfpathlineto{\pgfqpoint{2.611636in}{0.922247in}}%
\pgfpathlineto{\pgfqpoint{2.611947in}{0.955244in}}%
\pgfpathlineto{\pgfqpoint{2.611969in}{0.946995in}}%
\pgfpathlineto{\pgfqpoint{2.612288in}{0.955244in}}%
\pgfpathlineto{\pgfqpoint{2.612636in}{0.922247in}}%
\pgfpathlineto{\pgfqpoint{2.613458in}{0.955244in}}%
\pgfpathlineto{\pgfqpoint{2.613983in}{0.938746in}}%
\pgfpathlineto{\pgfqpoint{2.614465in}{0.930497in}}%
\pgfpathlineto{\pgfqpoint{2.614635in}{0.955244in}}%
\pgfpathlineto{\pgfqpoint{2.614798in}{0.955244in}}%
\pgfpathlineto{\pgfqpoint{2.614968in}{0.979992in}}%
\pgfpathlineto{\pgfqpoint{2.615486in}{0.946995in}}%
\pgfpathlineto{\pgfqpoint{2.615827in}{0.971742in}}%
\pgfpathlineto{\pgfqpoint{2.616138in}{0.979992in}}%
\pgfpathlineto{\pgfqpoint{2.616330in}{0.963493in}}%
\pgfpathlineto{\pgfqpoint{2.617671in}{0.922247in}}%
\pgfpathlineto{\pgfqpoint{2.618337in}{0.971742in}}%
\pgfpathlineto{\pgfqpoint{2.618840in}{0.938746in}}%
\pgfpathlineto{\pgfqpoint{2.620018in}{0.922247in}}%
\pgfpathlineto{\pgfqpoint{2.621173in}{0.930497in}}%
\pgfpathlineto{\pgfqpoint{2.621691in}{0.946995in}}%
\pgfpathlineto{\pgfqpoint{2.622194in}{0.922247in}}%
\pgfpathlineto{\pgfqpoint{2.622676in}{0.930497in}}%
\pgfpathlineto{\pgfqpoint{2.622868in}{0.913998in}}%
\pgfpathlineto{\pgfqpoint{2.623179in}{0.930497in}}%
\pgfpathlineto{\pgfqpoint{2.624038in}{0.897500in}}%
\pgfpathlineto{\pgfqpoint{2.625193in}{0.905749in}}%
\pgfpathlineto{\pgfqpoint{2.625711in}{0.897500in}}%
\pgfpathlineto{\pgfqpoint{2.626215in}{0.946995in}}%
\pgfpathlineto{\pgfqpoint{2.627370in}{0.955244in}}%
\pgfpathlineto{\pgfqpoint{2.627555in}{0.946995in}}%
\pgfpathlineto{\pgfqpoint{2.627725in}{0.971742in}}%
\pgfpathlineto{\pgfqpoint{2.628399in}{0.971742in}}%
\pgfpathlineto{\pgfqpoint{2.628710in}{0.955244in}}%
\pgfpathlineto{\pgfqpoint{2.628880in}{0.979992in}}%
\pgfpathlineto{\pgfqpoint{2.629214in}{0.979992in}}%
\pgfpathlineto{\pgfqpoint{2.629384in}{1.004739in}}%
\pgfpathlineto{\pgfqpoint{2.630243in}{0.971742in}}%
\pgfpathlineto{\pgfqpoint{2.631901in}{1.004739in}}%
\pgfpathlineto{\pgfqpoint{2.632235in}{1.029487in}}%
\pgfpathlineto{\pgfqpoint{2.632923in}{0.996490in}}%
\pgfpathlineto{\pgfqpoint{2.633575in}{1.029487in}}%
\pgfpathlineto{\pgfqpoint{2.634093in}{1.012988in}}%
\pgfpathlineto{\pgfqpoint{2.634241in}{1.004739in}}%
\pgfpathlineto{\pgfqpoint{2.634263in}{1.045985in}}%
\pgfpathlineto{\pgfqpoint{2.634597in}{1.045985in}}%
\pgfpathlineto{\pgfqpoint{2.635774in}{1.021237in}}%
\pgfpathlineto{\pgfqpoint{2.635922in}{1.029487in}}%
\pgfpathlineto{\pgfqpoint{2.636277in}{1.070732in}}%
\pgfpathlineto{\pgfqpoint{2.636944in}{1.021237in}}%
\pgfpathlineto{\pgfqpoint{2.638269in}{1.054234in}}%
\pgfpathlineto{\pgfqpoint{2.638284in}{1.045985in}}%
\pgfpathlineto{\pgfqpoint{2.639276in}{1.004739in}}%
\pgfpathlineto{\pgfqpoint{2.639609in}{1.029487in}}%
\pgfpathlineto{\pgfqpoint{2.639631in}{1.045985in}}%
\pgfpathlineto{\pgfqpoint{2.640283in}{1.004739in}}%
\pgfpathlineto{\pgfqpoint{2.640631in}{1.021237in}}%
\pgfpathlineto{\pgfqpoint{2.641956in}{1.004739in}}%
\pgfpathlineto{\pgfqpoint{2.642460in}{1.029487in}}%
\pgfpathlineto{\pgfqpoint{2.642978in}{1.021237in}}%
\pgfpathlineto{\pgfqpoint{2.644155in}{1.070732in}}%
\pgfpathlineto{\pgfqpoint{2.644637in}{1.054234in}}%
\pgfpathlineto{\pgfqpoint{2.645814in}{1.029487in}}%
\pgfpathlineto{\pgfqpoint{2.645829in}{1.045985in}}%
\pgfpathlineto{\pgfqpoint{2.646836in}{1.021237in}}%
\pgfpathlineto{\pgfqpoint{2.647487in}{1.054234in}}%
\pgfpathlineto{\pgfqpoint{2.647672in}{1.037736in}}%
\pgfpathlineto{\pgfqpoint{2.649012in}{0.996490in}}%
\pgfpathlineto{\pgfqpoint{2.650168in}{1.004739in}}%
\pgfpathlineto{\pgfqpoint{2.650353in}{1.021237in}}%
\pgfpathlineto{\pgfqpoint{2.651360in}{0.971742in}}%
\pgfpathlineto{\pgfqpoint{2.652537in}{0.996490in}}%
\pgfpathlineto{\pgfqpoint{2.652700in}{0.988241in}}%
\pgfpathlineto{\pgfqpoint{2.653188in}{0.979992in}}%
\pgfpathlineto{\pgfqpoint{2.653351in}{1.004739in}}%
\pgfpathlineto{\pgfqpoint{2.653374in}{0.996490in}}%
\pgfpathlineto{\pgfqpoint{2.653692in}{1.004739in}}%
\pgfpathlineto{\pgfqpoint{2.654025in}{0.979992in}}%
\pgfpathlineto{\pgfqpoint{2.654529in}{1.004739in}}%
\pgfpathlineto{\pgfqpoint{2.655047in}{0.971742in}}%
\pgfpathlineto{\pgfqpoint{2.656557in}{1.021237in}}%
\pgfpathlineto{\pgfqpoint{2.656728in}{1.012988in}}%
\pgfpathlineto{\pgfqpoint{2.656876in}{1.004739in}}%
\pgfpathlineto{\pgfqpoint{2.657061in}{1.045985in}}%
\pgfpathlineto{\pgfqpoint{2.658719in}{1.078981in}}%
\pgfpathlineto{\pgfqpoint{2.659556in}{1.054234in}}%
\pgfpathlineto{\pgfqpoint{2.659741in}{1.070732in}}%
\pgfpathlineto{\pgfqpoint{2.660230in}{1.103729in}}%
\pgfpathlineto{\pgfqpoint{2.660896in}{1.078981in}}%
\pgfpathlineto{\pgfqpoint{2.661252in}{1.095480in}}%
\pgfpathlineto{\pgfqpoint{2.661925in}{1.070732in}}%
\pgfpathlineto{\pgfqpoint{2.662429in}{1.095480in}}%
\pgfpathlineto{\pgfqpoint{2.663080in}{1.078981in}}%
\pgfpathlineto{\pgfqpoint{2.664102in}{1.070732in}}%
\pgfpathlineto{\pgfqpoint{2.666116in}{1.144975in}}%
\pgfpathlineto{\pgfqpoint{2.666279in}{1.136726in}}%
\pgfpathlineto{\pgfqpoint{2.666427in}{1.128476in}}%
\pgfpathlineto{\pgfqpoint{2.666783in}{1.169722in}}%
\pgfpathlineto{\pgfqpoint{2.667101in}{1.177971in}}%
\pgfpathlineto{\pgfqpoint{2.667434in}{1.153224in}}%
\pgfpathlineto{\pgfqpoint{2.667619in}{1.161473in}}%
\pgfpathlineto{\pgfqpoint{2.667790in}{1.169722in}}%
\pgfpathlineto{\pgfqpoint{2.668774in}{1.153224in}}%
\pgfpathlineto{\pgfqpoint{2.669803in}{1.219217in}}%
\pgfpathlineto{\pgfqpoint{2.671121in}{1.202719in}}%
\pgfpathlineto{\pgfqpoint{2.671477in}{1.243965in}}%
\pgfpathlineto{\pgfqpoint{2.672151in}{1.243965in}}%
\pgfpathlineto{\pgfqpoint{2.673320in}{1.219217in}}%
\pgfpathlineto{\pgfqpoint{2.673469in}{1.227466in}}%
\pgfpathlineto{\pgfqpoint{2.674476in}{1.276961in}}%
\pgfpathlineto{\pgfqpoint{2.674498in}{1.268712in}}%
\pgfpathlineto{\pgfqpoint{2.675482in}{1.301709in}}%
\pgfpathlineto{\pgfqpoint{2.674979in}{1.252214in}}%
\pgfpathlineto{\pgfqpoint{2.675668in}{1.285210in}}%
\pgfpathlineto{\pgfqpoint{2.676001in}{1.268712in}}%
\pgfpathlineto{\pgfqpoint{2.676319in}{1.301709in}}%
\pgfpathlineto{\pgfqpoint{2.676341in}{1.293460in}}%
\pgfpathlineto{\pgfqpoint{2.676489in}{1.301709in}}%
\pgfpathlineto{\pgfqpoint{2.676845in}{1.268712in}}%
\pgfpathlineto{\pgfqpoint{2.677156in}{1.276961in}}%
\pgfpathlineto{\pgfqpoint{2.677844in}{1.268712in}}%
\pgfpathlineto{\pgfqpoint{2.678163in}{1.301709in}}%
\pgfpathlineto{\pgfqpoint{2.678185in}{1.293460in}}%
\pgfpathlineto{\pgfqpoint{2.678666in}{1.301709in}}%
\pgfpathlineto{\pgfqpoint{2.678999in}{1.276961in}}%
\pgfpathlineto{\pgfqpoint{2.680029in}{1.194470in}}%
\pgfpathlineto{\pgfqpoint{2.680192in}{1.219217in}}%
\pgfpathlineto{\pgfqpoint{2.680510in}{1.227466in}}%
\pgfpathlineto{\pgfqpoint{2.680843in}{1.202719in}}%
\pgfpathlineto{\pgfqpoint{2.681369in}{1.219217in}}%
\pgfpathlineto{\pgfqpoint{2.681872in}{1.194470in}}%
\pgfpathlineto{\pgfqpoint{2.682687in}{1.227466in}}%
\pgfpathlineto{\pgfqpoint{2.683027in}{1.202719in}}%
\pgfpathlineto{\pgfqpoint{2.683383in}{1.194470in}}%
\pgfpathlineto{\pgfqpoint{2.683694in}{1.227466in}}%
\pgfpathlineto{\pgfqpoint{2.684049in}{1.219217in}}%
\pgfpathlineto{\pgfqpoint{2.685034in}{1.153224in}}%
\pgfpathlineto{\pgfqpoint{2.685893in}{1.169722in}}%
\pgfpathlineto{\pgfqpoint{2.688410in}{1.268712in}}%
\pgfpathlineto{\pgfqpoint{2.688573in}{1.260463in}}%
\pgfpathlineto{\pgfqpoint{2.689062in}{1.252214in}}%
\pgfpathlineto{\pgfqpoint{2.689232in}{1.276961in}}%
\pgfpathlineto{\pgfqpoint{2.689247in}{1.268712in}}%
\pgfpathlineto{\pgfqpoint{2.691246in}{1.326456in}}%
\pgfpathlineto{\pgfqpoint{2.692431in}{1.285210in}}%
\pgfpathlineto{\pgfqpoint{2.692764in}{1.268712in}}%
\pgfpathlineto{\pgfqpoint{2.693082in}{1.301709in}}%
\pgfpathlineto{\pgfqpoint{2.693104in}{1.293460in}}%
\pgfpathlineto{\pgfqpoint{2.693415in}{1.301709in}}%
\pgfpathlineto{\pgfqpoint{2.693756in}{1.276961in}}%
\pgfpathlineto{\pgfqpoint{2.694926in}{1.202719in}}%
\pgfpathlineto{\pgfqpoint{2.694948in}{1.219217in}}%
\pgfpathlineto{\pgfqpoint{2.696288in}{1.194470in}}%
\pgfpathlineto{\pgfqpoint{2.696947in}{1.227466in}}%
\pgfpathlineto{\pgfqpoint{2.697125in}{1.210968in}}%
\pgfpathlineto{\pgfqpoint{2.697458in}{1.169722in}}%
\pgfpathlineto{\pgfqpoint{2.697962in}{1.219217in}}%
\pgfpathlineto{\pgfqpoint{2.698302in}{1.194470in}}%
\pgfpathlineto{\pgfqpoint{2.700146in}{1.268712in}}%
\pgfpathlineto{\pgfqpoint{2.700309in}{1.260463in}}%
\pgfpathlineto{\pgfqpoint{2.701486in}{1.219217in}}%
\pgfpathlineto{\pgfqpoint{2.701634in}{1.227466in}}%
\pgfpathlineto{\pgfqpoint{2.701989in}{1.268712in}}%
\pgfpathlineto{\pgfqpoint{2.702656in}{1.268712in}}%
\pgfpathlineto{\pgfqpoint{2.703159in}{1.243965in}}%
\pgfpathlineto{\pgfqpoint{2.703478in}{1.276961in}}%
\pgfpathlineto{\pgfqpoint{2.703493in}{1.268712in}}%
\pgfpathlineto{\pgfqpoint{2.704492in}{1.276961in}}%
\pgfpathlineto{\pgfqpoint{2.704500in}{1.268712in}}%
\pgfpathlineto{\pgfqpoint{2.705151in}{1.227466in}}%
\pgfpathlineto{\pgfqpoint{2.705825in}{1.252214in}}%
\pgfpathlineto{\pgfqpoint{2.706847in}{1.293460in}}%
\pgfpathlineto{\pgfqpoint{2.707172in}{1.301709in}}%
\pgfpathlineto{\pgfqpoint{2.707350in}{1.285210in}}%
\pgfpathlineto{\pgfqpoint{2.709023in}{1.194470in}}%
\pgfpathlineto{\pgfqpoint{2.713051in}{1.021237in}}%
\pgfpathlineto{\pgfqpoint{2.713199in}{1.037736in}}%
\pgfpathlineto{\pgfqpoint{2.714221in}{1.120227in}}%
\pgfpathlineto{\pgfqpoint{2.714392in}{1.111978in}}%
\pgfpathlineto{\pgfqpoint{2.715547in}{1.078981in}}%
\pgfpathlineto{\pgfqpoint{2.714562in}{1.120227in}}%
\pgfpathlineto{\pgfqpoint{2.715561in}{1.095480in}}%
\pgfpathlineto{\pgfqpoint{2.716739in}{1.144975in}}%
\pgfpathlineto{\pgfqpoint{2.717072in}{1.120227in}}%
\pgfpathlineto{\pgfqpoint{2.718730in}{1.078981in}}%
\pgfpathlineto{\pgfqpoint{2.718738in}{1.087231in}}%
\pgfpathlineto{\pgfqpoint{2.719737in}{1.153224in}}%
\pgfpathlineto{\pgfqpoint{2.719064in}{1.078981in}}%
\pgfpathlineto{\pgfqpoint{2.719922in}{1.136726in}}%
\pgfpathlineto{\pgfqpoint{2.721433in}{1.070732in}}%
\pgfpathlineto{\pgfqpoint{2.722255in}{1.078981in}}%
\pgfpathlineto{\pgfqpoint{2.722270in}{1.070732in}}%
\pgfpathlineto{\pgfqpoint{2.722588in}{1.054234in}}%
\pgfpathlineto{\pgfqpoint{2.722929in}{1.087231in}}%
\pgfpathlineto{\pgfqpoint{2.724098in}{1.128476in}}%
\pgfpathlineto{\pgfqpoint{2.724935in}{1.078981in}}%
\pgfpathlineto{\pgfqpoint{2.725120in}{1.095480in}}%
\pgfpathlineto{\pgfqpoint{2.726949in}{1.153224in}}%
\pgfpathlineto{\pgfqpoint{2.726964in}{1.144975in}}%
\pgfpathlineto{\pgfqpoint{2.728474in}{1.045985in}}%
\pgfpathlineto{\pgfqpoint{2.728807in}{1.095480in}}%
\pgfpathlineto{\pgfqpoint{2.729296in}{1.103729in}}%
\pgfpathlineto{\pgfqpoint{2.729481in}{1.087231in}}%
\pgfpathlineto{\pgfqpoint{2.730970in}{1.054234in}}%
\pgfpathlineto{\pgfqpoint{2.730977in}{1.062483in}}%
\pgfpathlineto{\pgfqpoint{2.730984in}{1.095480in}}%
\pgfpathlineto{\pgfqpoint{2.731991in}{1.062483in}}%
\pgfpathlineto{\pgfqpoint{2.733480in}{1.029487in}}%
\pgfpathlineto{\pgfqpoint{2.732828in}{1.070732in}}%
\pgfpathlineto{\pgfqpoint{2.733487in}{1.037736in}}%
\pgfpathlineto{\pgfqpoint{2.734168in}{1.070732in}}%
\pgfpathlineto{\pgfqpoint{2.733709in}{1.029487in}}%
\pgfpathlineto{\pgfqpoint{2.734509in}{1.045985in}}%
\pgfpathlineto{\pgfqpoint{2.736182in}{1.095480in}}%
\pgfpathlineto{\pgfqpoint{2.736352in}{1.087231in}}%
\pgfpathlineto{\pgfqpoint{2.738529in}{1.021237in}}%
\pgfpathlineto{\pgfqpoint{2.739692in}{1.029487in}}%
\pgfpathlineto{\pgfqpoint{2.740188in}{1.054234in}}%
\pgfpathlineto{\pgfqpoint{2.740706in}{1.012988in}}%
\pgfpathlineto{\pgfqpoint{2.741883in}{0.996490in}}%
\pgfpathlineto{\pgfqpoint{2.742365in}{1.004739in}}%
\pgfpathlineto{\pgfqpoint{2.742550in}{0.988241in}}%
\pgfpathlineto{\pgfqpoint{2.743897in}{0.946995in}}%
\pgfpathlineto{\pgfqpoint{2.744386in}{0.955244in}}%
\pgfpathlineto{\pgfqpoint{2.744564in}{0.938746in}}%
\pgfpathlineto{\pgfqpoint{2.745904in}{0.897500in}}%
\pgfpathlineto{\pgfqpoint{2.746563in}{0.930497in}}%
\pgfpathlineto{\pgfqpoint{2.747081in}{0.913998in}}%
\pgfpathlineto{\pgfqpoint{2.747229in}{0.905749in}}%
\pgfpathlineto{\pgfqpoint{2.747407in}{0.930497in}}%
\pgfpathlineto{\pgfqpoint{2.747918in}{0.922247in}}%
\pgfpathlineto{\pgfqpoint{2.748066in}{0.930497in}}%
\pgfpathlineto{\pgfqpoint{2.748421in}{0.897500in}}%
\pgfpathlineto{\pgfqpoint{2.748569in}{0.905749in}}%
\pgfpathlineto{\pgfqpoint{2.749591in}{0.897500in}}%
\pgfpathlineto{\pgfqpoint{2.750768in}{0.823258in}}%
\pgfpathlineto{\pgfqpoint{2.751435in}{0.848005in}}%
\pgfpathlineto{\pgfqpoint{2.751590in}{0.856254in}}%
\pgfpathlineto{\pgfqpoint{2.751938in}{0.823258in}}%
\pgfpathlineto{\pgfqpoint{2.755463in}{0.674773in}}%
\pgfpathlineto{\pgfqpoint{2.755611in}{0.691271in}}%
\pgfpathlineto{\pgfqpoint{2.756795in}{0.757264in}}%
\pgfpathlineto{\pgfqpoint{2.755944in}{0.683022in}}%
\pgfpathlineto{\pgfqpoint{2.756803in}{0.749015in}}%
\pgfpathlineto{\pgfqpoint{2.757973in}{0.773763in}}%
\pgfpathlineto{\pgfqpoint{2.758143in}{0.765513in}}%
\pgfpathlineto{\pgfqpoint{2.758987in}{0.757264in}}%
\pgfpathlineto{\pgfqpoint{2.758313in}{0.798510in}}%
\pgfpathlineto{\pgfqpoint{2.759150in}{0.773763in}}%
\pgfpathlineto{\pgfqpoint{2.762993in}{0.625278in}}%
\pgfpathlineto{\pgfqpoint{2.764325in}{0.633527in}}%
\pgfpathlineto{\pgfqpoint{2.764325in}{0.625278in}}%
\pgfpathlineto{\pgfqpoint{2.764681in}{0.674773in}}%
\pgfpathlineto{\pgfqpoint{2.765355in}{0.674773in}}%
\pgfpathlineto{\pgfqpoint{2.765666in}{0.658274in}}%
\pgfpathlineto{\pgfqpoint{2.766014in}{0.691271in}}%
\pgfpathlineto{\pgfqpoint{2.767028in}{0.724268in}}%
\pgfpathlineto{\pgfqpoint{2.768346in}{0.683022in}}%
\pgfpathlineto{\pgfqpoint{2.767517in}{0.732517in}}%
\pgfpathlineto{\pgfqpoint{2.768368in}{0.699520in}}%
\pgfpathlineto{\pgfqpoint{2.770041in}{0.749015in}}%
\pgfpathlineto{\pgfqpoint{2.770212in}{0.740766in}}%
\pgfpathlineto{\pgfqpoint{2.771722in}{0.674773in}}%
\pgfpathlineto{\pgfqpoint{2.772211in}{0.683022in}}%
\pgfpathlineto{\pgfqpoint{2.772389in}{0.666524in}}%
\pgfpathlineto{\pgfqpoint{2.772714in}{0.683022in}}%
\pgfpathlineto{\pgfqpoint{2.774069in}{0.625278in}}%
\pgfpathlineto{\pgfqpoint{2.775061in}{0.633527in}}%
\pgfpathlineto{\pgfqpoint{2.775076in}{0.625278in}}%
\pgfpathlineto{\pgfqpoint{2.776231in}{0.633527in}}%
\pgfpathlineto{\pgfqpoint{2.776920in}{0.625278in}}%
\pgfpathlineto{\pgfqpoint{2.776416in}{0.650025in}}%
\pgfpathlineto{\pgfqpoint{2.777253in}{0.641776in}}%
\pgfpathlineto{\pgfqpoint{2.778416in}{0.625278in}}%
\pgfpathlineto{\pgfqpoint{2.779593in}{0.633527in}}%
\pgfpathlineto{\pgfqpoint{2.779763in}{0.625278in}}%
\pgfpathlineto{\pgfqpoint{2.780600in}{0.625278in}}%
\pgfpathlineto{\pgfqpoint{2.781762in}{0.633527in}}%
\pgfpathlineto{\pgfqpoint{2.781925in}{0.625278in}}%
\pgfpathlineto{\pgfqpoint{2.782769in}{0.625278in}}%
\pgfpathlineto{\pgfqpoint{2.784109in}{0.633527in}}%
\pgfpathlineto{\pgfqpoint{2.784280in}{0.625278in}}%
\pgfpathlineto{\pgfqpoint{2.785109in}{0.625278in}}%
\pgfpathlineto{\pgfqpoint{2.786294in}{0.633527in}}%
\pgfpathlineto{\pgfqpoint{2.786457in}{0.625278in}}%
\pgfpathlineto{\pgfqpoint{2.787301in}{0.625278in}}%
\pgfpathlineto{\pgfqpoint{2.788463in}{0.633527in}}%
\pgfpathlineto{\pgfqpoint{2.788641in}{0.625278in}}%
\pgfpathlineto{\pgfqpoint{2.789470in}{0.625278in}}%
\pgfpathlineto{\pgfqpoint{2.790647in}{0.633527in}}%
\pgfpathlineto{\pgfqpoint{2.790818in}{0.625278in}}%
\pgfpathlineto{\pgfqpoint{2.791647in}{0.625278in}}%
\pgfpathlineto{\pgfqpoint{2.792832in}{0.633527in}}%
\pgfpathlineto{\pgfqpoint{2.792994in}{0.625278in}}%
\pgfpathlineto{\pgfqpoint{2.793839in}{0.625278in}}%
\pgfpathlineto{\pgfqpoint{2.795008in}{0.633527in}}%
\pgfpathlineto{\pgfqpoint{2.795171in}{0.625278in}}%
\pgfpathlineto{\pgfqpoint{2.796015in}{0.625278in}}%
\pgfpathlineto{\pgfqpoint{2.797348in}{0.633527in}}%
\pgfpathlineto{\pgfqpoint{2.797859in}{0.625278in}}%
\pgfpathlineto{\pgfqpoint{2.798192in}{0.625278in}}%
\pgfpathlineto{\pgfqpoint{2.799369in}{0.633527in}}%
\pgfpathlineto{\pgfqpoint{2.799532in}{0.625278in}}%
\pgfpathlineto{\pgfqpoint{2.800369in}{0.625278in}}%
\pgfpathlineto{\pgfqpoint{2.801539in}{0.633527in}}%
\pgfpathlineto{\pgfqpoint{2.801709in}{0.625278in}}%
\pgfpathlineto{\pgfqpoint{2.802546in}{0.625278in}}%
\pgfpathlineto{\pgfqpoint{2.803886in}{0.633527in}}%
\pgfpathlineto{\pgfqpoint{2.803893in}{0.625278in}}%
\pgfpathlineto{\pgfqpoint{2.804900in}{0.625278in}}%
\pgfpathlineto{\pgfqpoint{2.805900in}{0.633527in}}%
\pgfpathlineto{\pgfqpoint{2.805900in}{0.625278in}}%
\pgfpathlineto{\pgfqpoint{2.807070in}{0.633527in}}%
\pgfpathlineto{\pgfqpoint{2.807070in}{0.625278in}}%
\pgfpathlineto{\pgfqpoint{2.808084in}{0.625278in}}%
\pgfpathlineto{\pgfqpoint{2.808913in}{0.633527in}}%
\pgfpathlineto{\pgfqpoint{2.808913in}{0.625278in}}%
\pgfpathlineto{\pgfqpoint{2.810424in}{0.633527in}}%
\pgfpathlineto{\pgfqpoint{2.810594in}{0.625278in}}%
\pgfpathlineto{\pgfqpoint{2.811431in}{0.625278in}}%
\pgfpathlineto{\pgfqpoint{2.812601in}{0.633527in}}%
\pgfpathlineto{\pgfqpoint{2.812601in}{0.625278in}}%
\pgfpathlineto{\pgfqpoint{2.813608in}{0.625278in}}%
\pgfpathlineto{\pgfqpoint{2.814615in}{0.633527in}}%
\pgfpathlineto{\pgfqpoint{2.814615in}{0.625278in}}%
\pgfpathlineto{\pgfqpoint{2.815792in}{0.633527in}}%
\pgfpathlineto{\pgfqpoint{2.815955in}{0.625278in}}%
\pgfpathlineto{\pgfqpoint{2.816792in}{0.625278in}}%
\pgfpathlineto{\pgfqpoint{2.817969in}{0.633527in}}%
\pgfpathlineto{\pgfqpoint{2.817969in}{0.625278in}}%
\pgfpathlineto{\pgfqpoint{2.818976in}{0.625278in}}%
\pgfpathlineto{\pgfqpoint{2.820146in}{0.633527in}}%
\pgfpathlineto{\pgfqpoint{2.820316in}{0.625278in}}%
\pgfpathlineto{\pgfqpoint{2.821153in}{0.625278in}}%
\pgfpathlineto{\pgfqpoint{2.822330in}{0.633527in}}%
\pgfpathlineto{\pgfqpoint{2.822493in}{0.625278in}}%
\pgfpathlineto{\pgfqpoint{2.823329in}{0.625278in}}%
\pgfpathlineto{\pgfqpoint{2.824507in}{0.633527in}}%
\pgfpathlineto{\pgfqpoint{2.824840in}{0.625278in}}%
\pgfpathlineto{\pgfqpoint{2.825514in}{0.625278in}}%
\pgfpathlineto{\pgfqpoint{2.826521in}{0.633527in}}%
\pgfpathlineto{\pgfqpoint{2.826521in}{0.625278in}}%
\pgfpathlineto{\pgfqpoint{2.827691in}{0.633527in}}%
\pgfpathlineto{\pgfqpoint{2.827861in}{0.625278in}}%
\pgfpathlineto{\pgfqpoint{2.828697in}{0.625278in}}%
\pgfpathlineto{\pgfqpoint{2.830038in}{0.633527in}}%
\pgfpathlineto{\pgfqpoint{2.830038in}{0.625278in}}%
\pgfpathlineto{\pgfqpoint{2.831045in}{0.625278in}}%
\pgfpathlineto{\pgfqpoint{2.832214in}{0.633527in}}%
\pgfpathlineto{\pgfqpoint{2.832385in}{0.625278in}}%
\pgfpathlineto{\pgfqpoint{2.833059in}{0.625278in}}%
\pgfpathlineto{\pgfqpoint{2.834228in}{0.633527in}}%
\pgfpathlineto{\pgfqpoint{2.834399in}{0.625278in}}%
\pgfpathlineto{\pgfqpoint{2.835065in}{0.625278in}}%
\pgfpathlineto{\pgfqpoint{2.836405in}{0.633527in}}%
\pgfpathlineto{\pgfqpoint{2.836405in}{0.625278in}}%
\pgfpathlineto{\pgfqpoint{2.837412in}{0.625278in}}%
\pgfpathlineto{\pgfqpoint{2.838589in}{0.633527in}}%
\pgfpathlineto{\pgfqpoint{2.838589in}{0.625278in}}%
\pgfpathlineto{\pgfqpoint{2.839648in}{0.625278in}}%
\pgfpathlineto{\pgfqpoint{2.840766in}{0.633527in}}%
\pgfpathlineto{\pgfqpoint{2.840937in}{0.625278in}}%
\pgfpathlineto{\pgfqpoint{2.841603in}{0.625278in}}%
\pgfpathlineto{\pgfqpoint{2.842780in}{0.633527in}}%
\pgfpathlineto{\pgfqpoint{2.843802in}{0.625278in}}%
\pgfpathlineto{\pgfqpoint{2.845290in}{0.658274in}}%
\pgfpathlineto{\pgfqpoint{2.845312in}{0.650025in}}%
\pgfpathlineto{\pgfqpoint{2.846630in}{0.625278in}}%
\pgfpathlineto{\pgfqpoint{2.847808in}{0.633527in}}%
\pgfpathlineto{\pgfqpoint{2.847971in}{0.625278in}}%
\pgfpathlineto{\pgfqpoint{2.848667in}{0.625278in}}%
\pgfpathlineto{\pgfqpoint{2.850155in}{0.683022in}}%
\pgfpathlineto{\pgfqpoint{2.850510in}{0.650025in}}%
\pgfpathlineto{\pgfqpoint{2.851347in}{0.625278in}}%
\pgfpathlineto{\pgfqpoint{2.851510in}{0.650025in}}%
\pgfpathlineto{\pgfqpoint{2.852665in}{0.658274in}}%
\pgfpathlineto{\pgfqpoint{2.853672in}{0.633527in}}%
\pgfpathlineto{\pgfqpoint{2.853694in}{0.650025in}}%
\pgfpathlineto{\pgfqpoint{2.854012in}{0.633527in}}%
\pgfpathlineto{\pgfqpoint{2.854027in}{0.674773in}}%
\pgfpathlineto{\pgfqpoint{2.854346in}{0.666524in}}%
\pgfpathlineto{\pgfqpoint{2.855182in}{0.683022in}}%
\pgfpathlineto{\pgfqpoint{2.854679in}{0.658274in}}%
\pgfpathlineto{\pgfqpoint{2.855367in}{0.666524in}}%
\pgfpathlineto{\pgfqpoint{2.855515in}{0.658274in}}%
\pgfpathlineto{\pgfqpoint{2.855686in}{0.683022in}}%
\pgfpathlineto{\pgfqpoint{2.856041in}{0.674773in}}%
\pgfpathlineto{\pgfqpoint{2.857700in}{0.707769in}}%
\pgfpathlineto{\pgfqpoint{2.858388in}{0.674773in}}%
\pgfpathlineto{\pgfqpoint{2.858721in}{0.699520in}}%
\pgfpathlineto{\pgfqpoint{2.861217in}{0.782012in}}%
\pgfpathlineto{\pgfqpoint{2.861239in}{0.773763in}}%
\pgfpathlineto{\pgfqpoint{2.861720in}{0.831507in}}%
\pgfpathlineto{\pgfqpoint{2.862238in}{0.815008in}}%
\pgfpathlineto{\pgfqpoint{2.863749in}{0.773763in}}%
\pgfpathlineto{\pgfqpoint{2.863897in}{0.782012in}}%
\pgfpathlineto{\pgfqpoint{2.864401in}{0.831507in}}%
\pgfpathlineto{\pgfqpoint{2.864926in}{0.798510in}}%
\pgfpathlineto{\pgfqpoint{2.865911in}{0.782012in}}%
\pgfpathlineto{\pgfqpoint{2.865933in}{0.798510in}}%
\pgfpathlineto{\pgfqpoint{2.866585in}{0.831507in}}%
\pgfpathlineto{\pgfqpoint{2.867081in}{0.806759in}}%
\pgfpathlineto{\pgfqpoint{2.868110in}{0.798510in}}%
\pgfpathlineto{\pgfqpoint{2.869954in}{0.872753in}}%
\pgfpathlineto{\pgfqpoint{2.870124in}{0.864503in}}%
\pgfpathlineto{\pgfqpoint{2.870620in}{0.848005in}}%
\pgfpathlineto{\pgfqpoint{2.870938in}{0.881002in}}%
\pgfpathlineto{\pgfqpoint{2.870961in}{0.872753in}}%
\pgfpathlineto{\pgfqpoint{2.871294in}{0.897500in}}%
\pgfpathlineto{\pgfqpoint{2.872116in}{0.881002in}}%
\pgfpathlineto{\pgfqpoint{2.872301in}{0.897500in}}%
\pgfpathlineto{\pgfqpoint{2.873137in}{0.864503in}}%
\pgfpathlineto{\pgfqpoint{2.873619in}{0.856254in}}%
\pgfpathlineto{\pgfqpoint{2.873789in}{0.881002in}}%
\pgfpathlineto{\pgfqpoint{2.873811in}{0.872753in}}%
\pgfpathlineto{\pgfqpoint{2.874463in}{0.881002in}}%
\pgfpathlineto{\pgfqpoint{2.874648in}{0.864503in}}%
\pgfpathlineto{\pgfqpoint{2.874796in}{0.856254in}}%
\pgfpathlineto{\pgfqpoint{2.874966in}{0.881002in}}%
\pgfpathlineto{\pgfqpoint{2.875485in}{0.872753in}}%
\pgfpathlineto{\pgfqpoint{2.876654in}{0.922247in}}%
\pgfpathlineto{\pgfqpoint{2.877143in}{0.905749in}}%
\pgfpathlineto{\pgfqpoint{2.877980in}{0.881002in}}%
\pgfpathlineto{\pgfqpoint{2.878165in}{0.889251in}}%
\pgfpathlineto{\pgfqpoint{2.878313in}{0.881002in}}%
\pgfpathlineto{\pgfqpoint{2.878483in}{0.930497in}}%
\pgfpathlineto{\pgfqpoint{2.878506in}{0.922247in}}%
\pgfpathlineto{\pgfqpoint{2.879157in}{0.905749in}}%
\pgfpathlineto{\pgfqpoint{2.882341in}{1.078981in}}%
\pgfpathlineto{\pgfqpoint{2.882526in}{1.095480in}}%
\pgfpathlineto{\pgfqpoint{2.883363in}{1.045985in}}%
\pgfpathlineto{\pgfqpoint{2.884518in}{1.078981in}}%
\pgfpathlineto{\pgfqpoint{2.884703in}{1.062483in}}%
\pgfpathlineto{\pgfqpoint{2.886213in}{1.021237in}}%
\pgfpathlineto{\pgfqpoint{2.886361in}{1.029487in}}%
\pgfpathlineto{\pgfqpoint{2.886717in}{1.095480in}}%
\pgfpathlineto{\pgfqpoint{2.887383in}{1.045985in}}%
\pgfpathlineto{\pgfqpoint{2.888035in}{1.029487in}}%
\pgfpathlineto{\pgfqpoint{2.888035in}{1.062483in}}%
\pgfpathlineto{\pgfqpoint{2.888390in}{1.070732in}}%
\pgfpathlineto{\pgfqpoint{2.889042in}{1.029487in}}%
\pgfpathlineto{\pgfqpoint{2.889064in}{1.045985in}}%
\pgfpathlineto{\pgfqpoint{2.889382in}{1.029487in}}%
\pgfpathlineto{\pgfqpoint{2.889397in}{1.070732in}}%
\pgfpathlineto{\pgfqpoint{2.890049in}{1.103729in}}%
\pgfpathlineto{\pgfqpoint{2.890404in}{1.070732in}}%
\pgfpathlineto{\pgfqpoint{2.891389in}{1.054234in}}%
\pgfpathlineto{\pgfqpoint{2.891411in}{1.070732in}}%
\pgfpathlineto{\pgfqpoint{2.893736in}{1.153224in}}%
\pgfpathlineto{\pgfqpoint{2.893758in}{1.144975in}}%
\pgfpathlineto{\pgfqpoint{2.894743in}{1.103729in}}%
\pgfpathlineto{\pgfqpoint{2.895098in}{1.120227in}}%
\pgfpathlineto{\pgfqpoint{2.896253in}{1.128476in}}%
\pgfpathlineto{\pgfqpoint{2.897275in}{1.120227in}}%
\pgfpathlineto{\pgfqpoint{2.897779in}{1.169722in}}%
\pgfpathlineto{\pgfqpoint{2.898615in}{1.136726in}}%
\pgfpathlineto{\pgfqpoint{2.899289in}{1.120227in}}%
\pgfpathlineto{\pgfqpoint{2.898786in}{1.144975in}}%
\pgfpathlineto{\pgfqpoint{2.899459in}{1.144975in}}%
\pgfpathlineto{\pgfqpoint{2.901451in}{1.202719in}}%
\pgfpathlineto{\pgfqpoint{2.901784in}{1.227466in}}%
\pgfpathlineto{\pgfqpoint{2.902473in}{1.194470in}}%
\pgfpathlineto{\pgfqpoint{2.903961in}{1.252214in}}%
\pgfpathlineto{\pgfqpoint{2.904146in}{1.235715in}}%
\pgfpathlineto{\pgfqpoint{2.904820in}{1.243965in}}%
\pgfpathlineto{\pgfqpoint{2.905827in}{1.194470in}}%
\pgfpathlineto{\pgfqpoint{2.907486in}{1.227466in}}%
\pgfpathlineto{\pgfqpoint{2.908322in}{1.202719in}}%
\pgfpathlineto{\pgfqpoint{2.908507in}{1.210968in}}%
\pgfpathlineto{\pgfqpoint{2.910018in}{1.144975in}}%
\pgfpathlineto{\pgfqpoint{2.910181in}{1.169722in}}%
\pgfpathlineto{\pgfqpoint{2.910521in}{1.194470in}}%
\pgfpathlineto{\pgfqpoint{2.911025in}{1.161473in}}%
\pgfpathlineto{\pgfqpoint{2.911528in}{1.144975in}}%
\pgfpathlineto{\pgfqpoint{2.911839in}{1.177971in}}%
\pgfpathlineto{\pgfqpoint{2.911861in}{1.169722in}}%
\pgfpathlineto{\pgfqpoint{2.912513in}{1.202719in}}%
\pgfpathlineto{\pgfqpoint{2.912868in}{1.169722in}}%
\pgfpathlineto{\pgfqpoint{2.914038in}{1.144975in}}%
\pgfpathlineto{\pgfqpoint{2.914186in}{1.153224in}}%
\pgfpathlineto{\pgfqpoint{2.915023in}{1.202719in}}%
\pgfpathlineto{\pgfqpoint{2.915216in}{1.194470in}}%
\pgfpathlineto{\pgfqpoint{2.915526in}{1.202719in}}%
\pgfpathlineto{\pgfqpoint{2.915867in}{1.177971in}}%
\pgfpathlineto{\pgfqpoint{2.916052in}{1.186221in}}%
\pgfpathlineto{\pgfqpoint{2.917540in}{1.128476in}}%
\pgfpathlineto{\pgfqpoint{2.916222in}{1.194470in}}%
\pgfpathlineto{\pgfqpoint{2.917726in}{1.144975in}}%
\pgfpathlineto{\pgfqpoint{2.918881in}{1.202719in}}%
\pgfpathlineto{\pgfqpoint{2.919384in}{1.177971in}}%
\pgfpathlineto{\pgfqpoint{2.919569in}{1.169722in}}%
\pgfpathlineto{\pgfqpoint{2.919888in}{1.202719in}}%
\pgfpathlineto{\pgfqpoint{2.920413in}{1.186221in}}%
\pgfpathlineto{\pgfqpoint{2.920576in}{1.194470in}}%
\pgfpathlineto{\pgfqpoint{2.921583in}{1.169722in}}%
\pgfpathlineto{\pgfqpoint{2.924582in}{1.326456in}}%
\pgfpathlineto{\pgfqpoint{2.925100in}{1.309958in}}%
\pgfpathlineto{\pgfqpoint{2.925441in}{1.293460in}}%
\pgfpathlineto{\pgfqpoint{2.925752in}{1.326456in}}%
\pgfpathlineto{\pgfqpoint{2.925774in}{1.318207in}}%
\pgfpathlineto{\pgfqpoint{2.927432in}{1.400699in}}%
\pgfpathlineto{\pgfqpoint{2.927766in}{1.375951in}}%
\pgfpathlineto{\pgfqpoint{2.928795in}{1.367702in}}%
\pgfpathlineto{\pgfqpoint{2.929446in}{1.400699in}}%
\pgfpathlineto{\pgfqpoint{2.929631in}{1.384200in}}%
\pgfpathlineto{\pgfqpoint{2.929942in}{1.351204in}}%
\pgfpathlineto{\pgfqpoint{2.930446in}{1.400699in}}%
\pgfpathlineto{\pgfqpoint{2.930638in}{1.392449in}}%
\pgfpathlineto{\pgfqpoint{2.931120in}{1.400699in}}%
\pgfpathlineto{\pgfqpoint{2.931453in}{1.375951in}}%
\pgfpathlineto{\pgfqpoint{2.931475in}{1.392449in}}%
\pgfpathlineto{\pgfqpoint{2.931786in}{1.375951in}}%
\pgfpathlineto{\pgfqpoint{2.932127in}{1.408948in}}%
\pgfpathlineto{\pgfqpoint{2.932460in}{1.425446in}}%
\pgfpathlineto{\pgfqpoint{2.932815in}{1.392449in}}%
\pgfpathlineto{\pgfqpoint{2.933148in}{1.417197in}}%
\pgfpathlineto{\pgfqpoint{2.934326in}{1.367702in}}%
\pgfpathlineto{\pgfqpoint{2.934489in}{1.392449in}}%
\pgfpathlineto{\pgfqpoint{2.935814in}{1.425446in}}%
\pgfpathlineto{\pgfqpoint{2.935829in}{1.417197in}}%
\pgfpathlineto{\pgfqpoint{2.937487in}{1.375951in}}%
\pgfpathlineto{\pgfqpoint{2.937487in}{1.384200in}}%
\pgfpathlineto{\pgfqpoint{2.937991in}{1.425446in}}%
\pgfpathlineto{\pgfqpoint{2.938517in}{1.417197in}}%
\pgfpathlineto{\pgfqpoint{2.940168in}{1.351204in}}%
\pgfpathlineto{\pgfqpoint{2.940360in}{1.367702in}}%
\pgfpathlineto{\pgfqpoint{2.941863in}{1.417197in}}%
\pgfpathlineto{\pgfqpoint{2.942033in}{1.408948in}}%
\pgfpathlineto{\pgfqpoint{2.942515in}{1.375951in}}%
\pgfpathlineto{\pgfqpoint{2.942870in}{1.417197in}}%
\pgfpathlineto{\pgfqpoint{2.943189in}{1.400699in}}%
\pgfpathlineto{\pgfqpoint{2.943211in}{1.417197in}}%
\pgfpathlineto{\pgfqpoint{2.943862in}{1.375951in}}%
\pgfpathlineto{\pgfqpoint{2.944210in}{1.392449in}}%
\pgfpathlineto{\pgfqpoint{2.946876in}{1.276961in}}%
\pgfpathlineto{\pgfqpoint{2.946876in}{1.309958in}}%
\pgfpathlineto{\pgfqpoint{2.947231in}{1.318207in}}%
\pgfpathlineto{\pgfqpoint{2.947883in}{1.276961in}}%
\pgfpathlineto{\pgfqpoint{2.948068in}{1.268712in}}%
\pgfpathlineto{\pgfqpoint{2.948905in}{1.318207in}}%
\pgfpathlineto{\pgfqpoint{2.950230in}{1.375951in}}%
\pgfpathlineto{\pgfqpoint{2.950571in}{1.351204in}}%
\pgfpathlineto{\pgfqpoint{2.951089in}{1.318207in}}%
\pgfpathlineto{\pgfqpoint{2.951592in}{1.342955in}}%
\pgfpathlineto{\pgfqpoint{2.951740in}{1.351204in}}%
\pgfpathlineto{\pgfqpoint{2.952088in}{1.318207in}}%
\pgfpathlineto{\pgfqpoint{2.952407in}{1.326456in}}%
\pgfpathlineto{\pgfqpoint{2.952762in}{1.318207in}}%
\pgfpathlineto{\pgfqpoint{2.953088in}{1.351204in}}%
\pgfpathlineto{\pgfqpoint{2.953436in}{1.334705in}}%
\pgfpathlineto{\pgfqpoint{2.953939in}{1.318207in}}%
\pgfpathlineto{\pgfqpoint{2.954102in}{1.342955in}}%
\pgfpathlineto{\pgfqpoint{2.954421in}{1.351204in}}%
\pgfpathlineto{\pgfqpoint{2.954754in}{1.326456in}}%
\pgfpathlineto{\pgfqpoint{2.955946in}{1.285210in}}%
\pgfpathlineto{\pgfqpoint{2.956094in}{1.276961in}}%
\pgfpathlineto{\pgfqpoint{2.956449in}{1.318207in}}%
\pgfpathlineto{\pgfqpoint{2.957271in}{1.326456in}}%
\pgfpathlineto{\pgfqpoint{2.957286in}{1.318207in}}%
\pgfpathlineto{\pgfqpoint{2.958967in}{1.268712in}}%
\pgfpathlineto{\pgfqpoint{2.959122in}{1.276961in}}%
\pgfpathlineto{\pgfqpoint{2.959952in}{1.326456in}}%
\pgfpathlineto{\pgfqpoint{2.960137in}{1.309958in}}%
\pgfpathlineto{\pgfqpoint{2.962654in}{1.219217in}}%
\pgfpathlineto{\pgfqpoint{2.963143in}{1.227466in}}%
\pgfpathlineto{\pgfqpoint{2.963321in}{1.210968in}}%
\pgfpathlineto{\pgfqpoint{2.966993in}{1.054234in}}%
\pgfpathlineto{\pgfqpoint{2.966993in}{1.062483in}}%
\pgfpathlineto{\pgfqpoint{2.967348in}{1.095480in}}%
\pgfpathlineto{\pgfqpoint{2.968015in}{1.045985in}}%
\pgfpathlineto{\pgfqpoint{2.968333in}{1.054234in}}%
\pgfpathlineto{\pgfqpoint{2.968666in}{1.029487in}}%
\pgfpathlineto{\pgfqpoint{2.968851in}{1.037736in}}%
\pgfpathlineto{\pgfqpoint{2.969022in}{1.045985in}}%
\pgfpathlineto{\pgfqpoint{2.970532in}{0.996490in}}%
\pgfpathlineto{\pgfqpoint{2.971369in}{1.021237in}}%
\pgfpathlineto{\pgfqpoint{2.971687in}{1.004739in}}%
\pgfpathlineto{\pgfqpoint{2.972694in}{0.979992in}}%
\pgfpathlineto{\pgfqpoint{2.972043in}{1.021237in}}%
\pgfpathlineto{\pgfqpoint{2.972709in}{0.996490in}}%
\pgfpathlineto{\pgfqpoint{2.974701in}{1.078981in}}%
\pgfpathlineto{\pgfqpoint{2.974893in}{1.062483in}}%
\pgfpathlineto{\pgfqpoint{2.975375in}{1.054234in}}%
\pgfpathlineto{\pgfqpoint{2.975545in}{1.078981in}}%
\pgfpathlineto{\pgfqpoint{2.975708in}{1.078981in}}%
\pgfpathlineto{\pgfqpoint{2.976737in}{1.144975in}}%
\pgfpathlineto{\pgfqpoint{2.977233in}{1.120227in}}%
\pgfpathlineto{\pgfqpoint{2.977737in}{1.136726in}}%
\pgfpathlineto{\pgfqpoint{2.978388in}{1.202719in}}%
\pgfpathlineto{\pgfqpoint{2.978743in}{1.169722in}}%
\pgfpathlineto{\pgfqpoint{2.978906in}{1.177971in}}%
\pgfpathlineto{\pgfqpoint{2.979247in}{1.144975in}}%
\pgfpathlineto{\pgfqpoint{2.979402in}{1.153224in}}%
\pgfpathlineto{\pgfqpoint{2.979750in}{1.120227in}}%
\pgfpathlineto{\pgfqpoint{2.980424in}{1.136726in}}%
\pgfpathlineto{\pgfqpoint{2.986288in}{0.872753in}}%
\pgfpathlineto{\pgfqpoint{2.986288in}{0.881002in}}%
\pgfpathlineto{\pgfqpoint{2.986459in}{0.905749in}}%
\pgfpathlineto{\pgfqpoint{2.986962in}{0.872753in}}%
\pgfpathlineto{\pgfqpoint{2.990479in}{0.732517in}}%
\pgfpathlineto{\pgfqpoint{2.990479in}{0.740766in}}%
\pgfpathlineto{\pgfqpoint{2.990649in}{0.773763in}}%
\pgfpathlineto{\pgfqpoint{2.991153in}{0.732517in}}%
\pgfpathlineto{\pgfqpoint{2.991486in}{0.732517in}}%
\pgfpathlineto{\pgfqpoint{2.992160in}{0.757264in}}%
\pgfpathlineto{\pgfqpoint{2.992493in}{0.740766in}}%
\pgfpathlineto{\pgfqpoint{2.993996in}{0.674773in}}%
\pgfpathlineto{\pgfqpoint{2.994670in}{0.699520in}}%
\pgfpathlineto{\pgfqpoint{2.995003in}{0.674773in}}%
\pgfpathlineto{\pgfqpoint{2.997017in}{0.625278in}}%
\pgfpathlineto{\pgfqpoint{2.998187in}{0.658274in}}%
\pgfpathlineto{\pgfqpoint{2.998357in}{0.641776in}}%
\pgfpathlineto{\pgfqpoint{2.999534in}{0.625278in}}%
\pgfpathlineto{\pgfqpoint{3.000023in}{0.633527in}}%
\pgfpathlineto{\pgfqpoint{3.000038in}{0.625278in}}%
\pgfpathlineto{\pgfqpoint{3.001378in}{0.633527in}}%
\pgfpathlineto{\pgfqpoint{3.001526in}{0.625278in}}%
\pgfpathlineto{\pgfqpoint{3.002385in}{0.625278in}}%
\pgfpathlineto{\pgfqpoint{3.003533in}{0.633527in}}%
\pgfpathlineto{\pgfqpoint{3.003555in}{0.625278in}}%
\pgfpathlineto{\pgfqpoint{3.004562in}{0.625278in}}%
\pgfpathlineto{\pgfqpoint{3.005902in}{0.633527in}}%
\pgfpathlineto{\pgfqpoint{3.005902in}{0.625278in}}%
\pgfpathlineto{\pgfqpoint{3.006909in}{0.625278in}}%
\pgfpathlineto{\pgfqpoint{3.008086in}{0.633527in}}%
\pgfpathlineto{\pgfqpoint{3.008420in}{0.625278in}}%
\pgfpathlineto{\pgfqpoint{3.009086in}{0.625278in}}%
\pgfpathlineto{\pgfqpoint{3.010241in}{0.633527in}}%
\pgfpathlineto{\pgfqpoint{3.010433in}{0.625278in}}%
\pgfpathlineto{\pgfqpoint{3.011270in}{0.625278in}}%
\pgfpathlineto{\pgfqpoint{3.012425in}{0.633527in}}%
\pgfpathlineto{\pgfqpoint{3.012610in}{0.625278in}}%
\pgfpathlineto{\pgfqpoint{3.013447in}{0.625278in}}%
\pgfpathlineto{\pgfqpoint{3.014624in}{0.633527in}}%
\pgfpathlineto{\pgfqpoint{3.014787in}{0.625278in}}%
\pgfpathlineto{\pgfqpoint{3.015461in}{0.625278in}}%
\pgfpathlineto{\pgfqpoint{3.016631in}{0.633527in}}%
\pgfpathlineto{\pgfqpoint{3.016964in}{0.625278in}}%
\pgfpathlineto{\pgfqpoint{3.017616in}{0.625278in}}%
\pgfpathlineto{\pgfqpoint{3.018474in}{0.633527in}}%
\pgfpathlineto{\pgfqpoint{3.018474in}{0.625278in}}%
\pgfpathlineto{\pgfqpoint{3.019652in}{0.633527in}}%
\pgfpathlineto{\pgfqpoint{3.019792in}{0.625278in}}%
\pgfpathlineto{\pgfqpoint{3.020488in}{0.625278in}}%
\pgfpathlineto{\pgfqpoint{3.021806in}{0.633527in}}%
\pgfpathlineto{\pgfqpoint{3.021821in}{0.625278in}}%
\pgfpathlineto{\pgfqpoint{3.022835in}{0.625278in}}%
\pgfpathlineto{\pgfqpoint{3.023983in}{0.633527in}}%
\pgfpathlineto{\pgfqpoint{3.024176in}{0.625278in}}%
\pgfpathlineto{\pgfqpoint{3.024849in}{0.625278in}}%
\pgfpathlineto{\pgfqpoint{3.026175in}{0.633527in}}%
\pgfpathlineto{\pgfqpoint{3.026190in}{0.625278in}}%
\pgfpathlineto{\pgfqpoint{3.027197in}{0.625278in}}%
\pgfpathlineto{\pgfqpoint{3.028344in}{0.633527in}}%
\pgfpathlineto{\pgfqpoint{3.028677in}{0.625278in}}%
\pgfpathlineto{\pgfqpoint{3.029373in}{0.625278in}}%
\pgfpathlineto{\pgfqpoint{3.030543in}{0.633527in}}%
\pgfpathlineto{\pgfqpoint{3.030691in}{0.625278in}}%
\pgfpathlineto{\pgfqpoint{3.031550in}{0.625278in}}%
\pgfpathlineto{\pgfqpoint{3.032727in}{0.633527in}}%
\pgfpathlineto{\pgfqpoint{3.033046in}{0.625278in}}%
\pgfpathlineto{\pgfqpoint{3.033727in}{0.625278in}}%
\pgfpathlineto{\pgfqpoint{3.034904in}{0.633527in}}%
\pgfpathlineto{\pgfqpoint{3.035215in}{0.625278in}}%
\pgfpathlineto{\pgfqpoint{3.035719in}{0.625278in}}%
\pgfpathlineto{\pgfqpoint{3.036896in}{0.633527in}}%
\pgfpathlineto{\pgfqpoint{3.036911in}{0.625278in}}%
\pgfpathlineto{\pgfqpoint{3.037903in}{0.625278in}}%
\pgfpathlineto{\pgfqpoint{3.039073in}{0.633527in}}%
\pgfpathlineto{\pgfqpoint{3.039265in}{0.625278in}}%
\pgfpathlineto{\pgfqpoint{3.039932in}{0.625278in}}%
\pgfpathlineto{\pgfqpoint{3.041102in}{0.633527in}}%
\pgfpathlineto{\pgfqpoint{3.041272in}{0.625278in}}%
\pgfpathlineto{\pgfqpoint{3.042205in}{0.625278in}}%
\pgfpathlineto{\pgfqpoint{3.043286in}{0.633527in}}%
\pgfpathlineto{\pgfqpoint{3.043286in}{0.625278in}}%
\pgfpathlineto{\pgfqpoint{3.044448in}{0.633527in}}%
\pgfpathlineto{\pgfqpoint{3.044626in}{0.625278in}}%
\pgfpathlineto{\pgfqpoint{3.045300in}{0.625278in}}%
\pgfpathlineto{\pgfqpoint{3.046470in}{0.633527in}}%
\pgfpathlineto{\pgfqpoint{3.046618in}{0.625278in}}%
\pgfpathlineto{\pgfqpoint{3.047477in}{0.625278in}}%
\pgfpathlineto{\pgfqpoint{3.048632in}{0.633527in}}%
\pgfpathlineto{\pgfqpoint{3.048646in}{0.625278in}}%
\pgfpathlineto{\pgfqpoint{3.049468in}{0.625278in}}%
\pgfpathlineto{\pgfqpoint{3.050809in}{0.633527in}}%
\pgfpathlineto{\pgfqpoint{3.050823in}{0.625278in}}%
\pgfpathlineto{\pgfqpoint{3.051815in}{0.625278in}}%
\pgfpathlineto{\pgfqpoint{3.052985in}{0.633527in}}%
\pgfpathlineto{\pgfqpoint{3.053008in}{0.625278in}}%
\pgfpathlineto{\pgfqpoint{3.053992in}{0.625278in}}%
\pgfpathlineto{\pgfqpoint{3.055170in}{0.633527in}}%
\pgfpathlineto{\pgfqpoint{3.055184in}{0.625278in}}%
\pgfpathlineto{\pgfqpoint{3.056191in}{0.625278in}}%
\pgfpathlineto{\pgfqpoint{3.057509in}{0.633527in}}%
\pgfpathlineto{\pgfqpoint{3.057532in}{0.625278in}}%
\pgfpathlineto{\pgfqpoint{3.058524in}{0.625278in}}%
\pgfpathlineto{\pgfqpoint{3.059664in}{0.633527in}}%
\pgfpathlineto{\pgfqpoint{3.059694in}{0.625278in}}%
\pgfpathlineto{\pgfqpoint{3.060552in}{0.625278in}}%
\pgfpathlineto{\pgfqpoint{3.061722in}{0.633527in}}%
\pgfpathlineto{\pgfqpoint{3.061870in}{0.625278in}}%
\pgfpathlineto{\pgfqpoint{3.062544in}{0.625278in}}%
\pgfpathlineto{\pgfqpoint{3.063714in}{0.633527in}}%
\pgfpathlineto{\pgfqpoint{3.063714in}{0.625278in}}%
\pgfpathlineto{\pgfqpoint{3.064551in}{0.625278in}}%
\pgfpathlineto{\pgfqpoint{3.065728in}{0.633527in}}%
\pgfpathlineto{\pgfqpoint{3.065891in}{0.625278in}}%
\pgfpathlineto{\pgfqpoint{3.066735in}{0.625278in}}%
\pgfpathlineto{\pgfqpoint{3.067905in}{0.633527in}}%
\pgfpathlineto{\pgfqpoint{3.067927in}{0.625278in}}%
\pgfpathlineto{\pgfqpoint{3.068912in}{0.625278in}}%
\pgfpathlineto{\pgfqpoint{3.070252in}{0.633527in}}%
\pgfpathlineto{\pgfqpoint{3.070585in}{0.625278in}}%
\pgfpathlineto{\pgfqpoint{3.071089in}{0.625278in}}%
\pgfpathlineto{\pgfqpoint{3.072429in}{0.633527in}}%
\pgfpathlineto{\pgfqpoint{3.072769in}{0.625278in}}%
\pgfpathlineto{\pgfqpoint{3.073273in}{0.625278in}}%
\pgfpathlineto{\pgfqpoint{3.074450in}{0.633527in}}%
\pgfpathlineto{\pgfqpoint{3.074946in}{0.625278in}}%
\pgfpathlineto{\pgfqpoint{3.075450in}{0.625278in}}%
\pgfpathlineto{\pgfqpoint{3.076620in}{0.633527in}}%
\pgfpathlineto{\pgfqpoint{3.076960in}{0.625278in}}%
\pgfpathlineto{\pgfqpoint{3.077627in}{0.625278in}}%
\pgfpathlineto{\pgfqpoint{3.078804in}{0.633527in}}%
\pgfpathlineto{\pgfqpoint{3.078967in}{0.625278in}}%
\pgfpathlineto{\pgfqpoint{3.079811in}{0.625278in}}%
\pgfpathlineto{\pgfqpoint{3.080981in}{0.633527in}}%
\pgfpathlineto{\pgfqpoint{3.081484in}{0.625278in}}%
\pgfpathlineto{\pgfqpoint{3.081988in}{0.625278in}}%
\pgfpathlineto{\pgfqpoint{3.082995in}{0.633527in}}%
\pgfpathlineto{\pgfqpoint{3.082995in}{0.625278in}}%
\pgfpathlineto{\pgfqpoint{3.084335in}{0.633527in}}%
\pgfpathlineto{\pgfqpoint{3.084335in}{0.625278in}}%
\pgfpathlineto{\pgfqpoint{3.085231in}{0.625278in}}%
\pgfpathlineto{\pgfqpoint{3.086349in}{0.633527in}}%
\pgfpathlineto{\pgfqpoint{3.086349in}{0.625278in}}%
\pgfpathlineto{\pgfqpoint{3.087185in}{0.625278in}}%
\pgfpathlineto{\pgfqpoint{3.088355in}{0.633527in}}%
\pgfpathlineto{\pgfqpoint{3.088525in}{0.625278in}}%
\pgfpathlineto{\pgfqpoint{3.089192in}{0.625278in}}%
\pgfpathlineto{\pgfqpoint{3.090539in}{0.633527in}}%
\pgfpathlineto{\pgfqpoint{3.090873in}{0.625278in}}%
\pgfpathlineto{\pgfqpoint{3.091539in}{0.625278in}}%
\pgfpathlineto{\pgfqpoint{3.092887in}{0.633527in}}%
\pgfpathlineto{\pgfqpoint{3.092887in}{0.625278in}}%
\pgfpathlineto{\pgfqpoint{3.093886in}{0.625278in}}%
\pgfpathlineto{\pgfqpoint{3.095063in}{0.633527in}}%
\pgfpathlineto{\pgfqpoint{3.095226in}{0.625278in}}%
\pgfpathlineto{\pgfqpoint{3.096070in}{0.625278in}}%
\pgfpathlineto{\pgfqpoint{3.097240in}{0.633527in}}%
\pgfpathlineto{\pgfqpoint{3.097240in}{0.625278in}}%
\pgfpathlineto{\pgfqpoint{3.098247in}{0.625278in}}%
\pgfpathlineto{\pgfqpoint{3.099417in}{0.633527in}}%
\pgfpathlineto{\pgfqpoint{3.099587in}{0.625278in}}%
\pgfpathlineto{\pgfqpoint{3.100424in}{0.625278in}}%
\pgfpathlineto{\pgfqpoint{3.101601in}{0.633527in}}%
\pgfpathlineto{\pgfqpoint{3.101934in}{0.625278in}}%
\pgfpathlineto{\pgfqpoint{3.102608in}{0.625278in}}%
\pgfpathlineto{\pgfqpoint{3.103778in}{0.633527in}}%
\pgfpathlineto{\pgfqpoint{3.103948in}{0.625278in}}%
\pgfpathlineto{\pgfqpoint{3.104785in}{0.625278in}}%
\pgfpathlineto{\pgfqpoint{3.105955in}{0.633527in}}%
\pgfpathlineto{\pgfqpoint{3.106125in}{0.625278in}}%
\pgfpathlineto{\pgfqpoint{3.106962in}{0.625278in}}%
\pgfpathlineto{\pgfqpoint{3.108139in}{0.633527in}}%
\pgfpathlineto{\pgfqpoint{3.108139in}{0.625278in}}%
\pgfpathlineto{\pgfqpoint{3.109146in}{0.625278in}}%
\pgfpathlineto{\pgfqpoint{3.110316in}{0.633527in}}%
\pgfpathlineto{\pgfqpoint{3.110486in}{0.625278in}}%
\pgfpathlineto{\pgfqpoint{3.111323in}{0.625278in}}%
\pgfpathlineto{\pgfqpoint{3.112493in}{0.633527in}}%
\pgfpathlineto{\pgfqpoint{3.112493in}{0.625278in}}%
\pgfpathlineto{\pgfqpoint{3.113500in}{0.625278in}}%
\pgfpathlineto{\pgfqpoint{3.114677in}{0.633527in}}%
\pgfpathlineto{\pgfqpoint{3.114677in}{0.625278in}}%
\pgfpathlineto{\pgfqpoint{3.115684in}{0.625278in}}%
\pgfpathlineto{\pgfqpoint{3.116854in}{0.633527in}}%
\pgfpathlineto{\pgfqpoint{3.116854in}{0.625278in}}%
\pgfpathlineto{\pgfqpoint{3.117861in}{0.625278in}}%
\pgfpathlineto{\pgfqpoint{3.118868in}{0.633527in}}%
\pgfpathlineto{\pgfqpoint{3.118868in}{0.625278in}}%
\pgfpathlineto{\pgfqpoint{3.120208in}{0.633527in}}%
\pgfpathlineto{\pgfqpoint{3.120208in}{0.625278in}}%
\pgfpathlineto{\pgfqpoint{3.121215in}{0.625278in}}%
\pgfpathlineto{\pgfqpoint{3.122385in}{0.633527in}}%
\pgfpathlineto{\pgfqpoint{3.122555in}{0.625278in}}%
\pgfpathlineto{\pgfqpoint{3.123392in}{0.625278in}}%
\pgfpathlineto{\pgfqpoint{3.124562in}{0.633527in}}%
\pgfpathlineto{\pgfqpoint{3.124732in}{0.625278in}}%
\pgfpathlineto{\pgfqpoint{3.125406in}{0.625278in}}%
\pgfpathlineto{\pgfqpoint{3.126576in}{0.633527in}}%
\pgfpathlineto{\pgfqpoint{3.126746in}{0.625278in}}%
\pgfpathlineto{\pgfqpoint{3.127583in}{0.625278in}}%
\pgfpathlineto{\pgfqpoint{3.128752in}{0.633527in}}%
\pgfpathlineto{\pgfqpoint{3.128752in}{0.625278in}}%
\pgfpathlineto{\pgfqpoint{3.129759in}{0.625278in}}%
\pgfpathlineto{\pgfqpoint{3.130937in}{0.633527in}}%
\pgfpathlineto{\pgfqpoint{3.130937in}{0.625278in}}%
\pgfpathlineto{\pgfqpoint{3.131944in}{0.625278in}}%
\pgfpathlineto{\pgfqpoint{3.133114in}{0.633527in}}%
\pgfpathlineto{\pgfqpoint{3.133114in}{0.625278in}}%
\pgfpathlineto{\pgfqpoint{3.134121in}{0.625278in}}%
\pgfpathlineto{\pgfqpoint{3.135290in}{0.633527in}}%
\pgfpathlineto{\pgfqpoint{3.135461in}{0.625278in}}%
\pgfpathlineto{\pgfqpoint{3.136297in}{0.625278in}}%
\pgfpathlineto{\pgfqpoint{3.137475in}{0.633527in}}%
\pgfpathlineto{\pgfqpoint{3.137638in}{0.625278in}}%
\pgfpathlineto{\pgfqpoint{3.138482in}{0.625278in}}%
\pgfpathlineto{\pgfqpoint{3.139481in}{0.633527in}}%
\pgfpathlineto{\pgfqpoint{3.139481in}{0.625278in}}%
\pgfpathlineto{\pgfqpoint{3.140658in}{0.633527in}}%
\pgfpathlineto{\pgfqpoint{3.140829in}{0.625278in}}%
\pgfpathlineto{\pgfqpoint{3.141665in}{0.625278in}}%
\pgfpathlineto{\pgfqpoint{3.142835in}{0.633527in}}%
\pgfpathlineto{\pgfqpoint{3.143006in}{0.625278in}}%
\pgfpathlineto{\pgfqpoint{3.143672in}{0.625278in}}%
\pgfpathlineto{\pgfqpoint{3.144849in}{0.633527in}}%
\pgfpathlineto{\pgfqpoint{3.145020in}{0.625278in}}%
\pgfpathlineto{\pgfqpoint{3.145856in}{0.625278in}}%
\pgfpathlineto{\pgfqpoint{3.147026in}{0.633527in}}%
\pgfpathlineto{\pgfqpoint{3.147026in}{0.625278in}}%
\pgfpathlineto{\pgfqpoint{3.148033in}{0.625278in}}%
\pgfpathlineto{\pgfqpoint{3.149210in}{0.633527in}}%
\pgfpathlineto{\pgfqpoint{3.149373in}{0.625278in}}%
\pgfpathlineto{\pgfqpoint{3.150210in}{0.625278in}}%
\pgfpathlineto{\pgfqpoint{3.151387in}{0.633527in}}%
\pgfpathlineto{\pgfqpoint{3.151557in}{0.625278in}}%
\pgfpathlineto{\pgfqpoint{3.152409in}{0.699520in}}%
\pgfpathlineto{\pgfqpoint{3.153586in}{0.749015in}}%
\pgfpathlineto{\pgfqpoint{3.154067in}{0.732517in}}%
\pgfpathlineto{\pgfqpoint{3.155097in}{0.724268in}}%
\pgfpathlineto{\pgfqpoint{3.156244in}{0.732517in}}%
\pgfpathlineto{\pgfqpoint{3.156770in}{0.724268in}}%
\pgfpathlineto{\pgfqpoint{3.157088in}{0.757264in}}%
\pgfpathlineto{\pgfqpoint{3.157273in}{0.749015in}}%
\pgfpathlineto{\pgfqpoint{3.157755in}{0.732517in}}%
\pgfpathlineto{\pgfqpoint{3.158428in}{0.757264in}}%
\pgfpathlineto{\pgfqpoint{3.159450in}{0.749015in}}%
\pgfpathlineto{\pgfqpoint{3.160961in}{0.823258in}}%
\pgfpathlineto{\pgfqpoint{3.161442in}{0.806759in}}%
\pgfpathlineto{\pgfqpoint{3.161464in}{0.798510in}}%
\pgfpathlineto{\pgfqpoint{3.161968in}{0.848005in}}%
\pgfpathlineto{\pgfqpoint{3.162301in}{0.848005in}}%
\pgfpathlineto{\pgfqpoint{3.163478in}{0.798510in}}%
\pgfpathlineto{\pgfqpoint{3.163641in}{0.823258in}}%
\pgfpathlineto{\pgfqpoint{3.163974in}{0.848005in}}%
\pgfpathlineto{\pgfqpoint{3.164478in}{0.815008in}}%
\pgfpathlineto{\pgfqpoint{3.165152in}{0.798510in}}%
\pgfpathlineto{\pgfqpoint{3.165322in}{0.823258in}}%
\pgfpathlineto{\pgfqpoint{3.165633in}{0.806759in}}%
\pgfpathlineto{\pgfqpoint{3.166662in}{0.848005in}}%
\pgfpathlineto{\pgfqpoint{3.166973in}{0.856254in}}%
\pgfpathlineto{\pgfqpoint{3.167314in}{0.831507in}}%
\pgfpathlineto{\pgfqpoint{3.167499in}{0.839756in}}%
\pgfpathlineto{\pgfqpoint{3.167647in}{0.831507in}}%
\pgfpathlineto{\pgfqpoint{3.167669in}{0.872753in}}%
\pgfpathlineto{\pgfqpoint{3.168165in}{0.864503in}}%
\pgfpathlineto{\pgfqpoint{3.169513in}{0.823258in}}%
\pgfpathlineto{\pgfqpoint{3.169661in}{0.839756in}}%
\pgfpathlineto{\pgfqpoint{3.170164in}{0.881002in}}%
\pgfpathlineto{\pgfqpoint{3.170719in}{0.831507in}}%
\pgfpathlineto{\pgfqpoint{3.171186in}{0.872753in}}%
\pgfpathlineto{\pgfqpoint{3.171689in}{0.839756in}}%
\pgfpathlineto{\pgfqpoint{3.172023in}{0.823258in}}%
\pgfpathlineto{\pgfqpoint{3.172341in}{0.856254in}}%
\pgfpathlineto{\pgfqpoint{3.172356in}{0.848005in}}%
\pgfpathlineto{\pgfqpoint{3.173348in}{0.856254in}}%
\pgfpathlineto{\pgfqpoint{3.173363in}{0.848005in}}%
\pgfpathlineto{\pgfqpoint{3.174703in}{0.823258in}}%
\pgfpathlineto{\pgfqpoint{3.175858in}{0.831507in}}%
\pgfpathlineto{\pgfqpoint{3.176050in}{0.823258in}}%
\pgfpathlineto{\pgfqpoint{3.176361in}{0.856254in}}%
\pgfpathlineto{\pgfqpoint{3.176887in}{0.839756in}}%
\pgfpathlineto{\pgfqpoint{3.177894in}{0.823258in}}%
\pgfpathlineto{\pgfqpoint{3.178042in}{0.831507in}}%
\pgfpathlineto{\pgfqpoint{3.178205in}{0.856254in}}%
\pgfpathlineto{\pgfqpoint{3.178731in}{0.823258in}}%
\pgfpathlineto{\pgfqpoint{3.179064in}{0.848005in}}%
\pgfpathlineto{\pgfqpoint{3.180071in}{0.872753in}}%
\pgfpathlineto{\pgfqpoint{3.180241in}{0.864503in}}%
\pgfpathlineto{\pgfqpoint{3.180723in}{0.856254in}}%
\pgfpathlineto{\pgfqpoint{3.180893in}{0.881002in}}%
\pgfpathlineto{\pgfqpoint{3.181056in}{0.881002in}}%
\pgfpathlineto{\pgfqpoint{3.181226in}{0.905749in}}%
\pgfpathlineto{\pgfqpoint{3.181915in}{0.872753in}}%
\pgfpathlineto{\pgfqpoint{3.182085in}{0.872753in}}%
\pgfpathlineto{\pgfqpoint{3.183403in}{0.905749in}}%
\pgfpathlineto{\pgfqpoint{3.183425in}{0.897500in}}%
\pgfpathlineto{\pgfqpoint{3.184077in}{0.881002in}}%
\pgfpathlineto{\pgfqpoint{3.184077in}{0.913998in}}%
\pgfpathlineto{\pgfqpoint{3.184432in}{0.922247in}}%
\pgfpathlineto{\pgfqpoint{3.185084in}{0.881002in}}%
\pgfpathlineto{\pgfqpoint{3.185098in}{0.897500in}}%
\pgfpathlineto{\pgfqpoint{3.186757in}{0.856254in}}%
\pgfpathlineto{\pgfqpoint{3.186757in}{0.864503in}}%
\pgfpathlineto{\pgfqpoint{3.187446in}{0.922247in}}%
\pgfpathlineto{\pgfqpoint{3.187927in}{0.905749in}}%
\pgfpathlineto{\pgfqpoint{3.188282in}{0.897500in}}%
\pgfpathlineto{\pgfqpoint{3.188453in}{0.922247in}}%
\pgfpathlineto{\pgfqpoint{3.188956in}{0.913998in}}%
\pgfpathlineto{\pgfqpoint{3.189770in}{0.881002in}}%
\pgfpathlineto{\pgfqpoint{3.189793in}{0.922247in}}%
\pgfpathlineto{\pgfqpoint{3.190963in}{0.946995in}}%
\pgfpathlineto{\pgfqpoint{3.191133in}{0.938746in}}%
\pgfpathlineto{\pgfqpoint{3.191510in}{0.905749in}}%
\pgfpathlineto{\pgfqpoint{3.191636in}{0.946995in}}%
\pgfpathlineto{\pgfqpoint{3.192310in}{0.922247in}}%
\pgfpathlineto{\pgfqpoint{3.192621in}{0.930497in}}%
\pgfpathlineto{\pgfqpoint{3.192976in}{0.897500in}}%
\pgfpathlineto{\pgfqpoint{3.193125in}{0.905749in}}%
\pgfpathlineto{\pgfqpoint{3.193147in}{0.897500in}}%
\pgfpathlineto{\pgfqpoint{3.193813in}{0.971742in}}%
\pgfpathlineto{\pgfqpoint{3.193983in}{0.963493in}}%
\pgfpathlineto{\pgfqpoint{3.195153in}{0.922247in}}%
\pgfpathlineto{\pgfqpoint{3.195324in}{0.946995in}}%
\pgfpathlineto{\pgfqpoint{3.196331in}{0.971742in}}%
\pgfpathlineto{\pgfqpoint{3.196501in}{0.963493in}}%
\pgfpathlineto{\pgfqpoint{3.196982in}{0.955244in}}%
\pgfpathlineto{\pgfqpoint{3.197152in}{0.979992in}}%
\pgfpathlineto{\pgfqpoint{3.197338in}{0.971742in}}%
\pgfpathlineto{\pgfqpoint{3.198996in}{1.004739in}}%
\pgfpathlineto{\pgfqpoint{3.199181in}{0.996490in}}%
\pgfpathlineto{\pgfqpoint{3.199833in}{1.029487in}}%
\pgfpathlineto{\pgfqpoint{3.200018in}{1.021237in}}%
\pgfpathlineto{\pgfqpoint{3.201173in}{1.029487in}}%
\pgfpathlineto{\pgfqpoint{3.201862in}{1.021237in}}%
\pgfpathlineto{\pgfqpoint{3.202180in}{1.054234in}}%
\pgfpathlineto{\pgfqpoint{3.202195in}{1.045985in}}%
\pgfpathlineto{\pgfqpoint{3.203535in}{1.095480in}}%
\pgfpathlineto{\pgfqpoint{3.203853in}{1.078981in}}%
\pgfpathlineto{\pgfqpoint{3.204527in}{1.054234in}}%
\pgfpathlineto{\pgfqpoint{3.204919in}{1.078981in}}%
\pgfpathlineto{\pgfqpoint{3.206200in}{1.103729in}}%
\pgfpathlineto{\pgfqpoint{3.206223in}{1.095480in}}%
\pgfpathlineto{\pgfqpoint{3.207207in}{1.103729in}}%
\pgfpathlineto{\pgfqpoint{3.207230in}{1.095480in}}%
\pgfpathlineto{\pgfqpoint{3.208547in}{1.078981in}}%
\pgfpathlineto{\pgfqpoint{3.208570in}{1.120227in}}%
\pgfpathlineto{\pgfqpoint{3.209577in}{1.120227in}}%
\pgfpathlineto{\pgfqpoint{3.210228in}{1.128476in}}%
\pgfpathlineto{\pgfqpoint{3.210413in}{1.111978in}}%
\pgfpathlineto{\pgfqpoint{3.211080in}{1.095480in}}%
\pgfpathlineto{\pgfqpoint{3.211250in}{1.120227in}}%
\pgfpathlineto{\pgfqpoint{3.211568in}{1.103729in}}%
\pgfpathlineto{\pgfqpoint{3.211916in}{1.169722in}}%
\pgfpathlineto{\pgfqpoint{3.212590in}{1.169722in}}%
\pgfpathlineto{\pgfqpoint{3.213767in}{1.144975in}}%
\pgfpathlineto{\pgfqpoint{3.213916in}{1.153224in}}%
\pgfpathlineto{\pgfqpoint{3.214078in}{1.177971in}}%
\pgfpathlineto{\pgfqpoint{3.214915in}{1.128476in}}%
\pgfpathlineto{\pgfqpoint{3.214937in}{1.144975in}}%
\pgfpathlineto{\pgfqpoint{3.215944in}{1.120227in}}%
\pgfpathlineto{\pgfqpoint{3.216092in}{1.136726in}}%
\pgfpathlineto{\pgfqpoint{3.216426in}{1.153224in}}%
\pgfpathlineto{\pgfqpoint{3.216781in}{1.120227in}}%
\pgfpathlineto{\pgfqpoint{3.217114in}{1.144975in}}%
\pgfpathlineto{\pgfqpoint{3.218269in}{1.153224in}}%
\pgfpathlineto{\pgfqpoint{3.218291in}{1.144975in}}%
\pgfpathlineto{\pgfqpoint{3.219128in}{1.194470in}}%
\pgfpathlineto{\pgfqpoint{3.220283in}{1.202719in}}%
\pgfpathlineto{\pgfqpoint{3.221290in}{1.153224in}}%
\pgfpathlineto{\pgfqpoint{3.221305in}{1.169722in}}%
\pgfpathlineto{\pgfqpoint{3.221623in}{1.153224in}}%
\pgfpathlineto{\pgfqpoint{3.221794in}{1.177971in}}%
\pgfpathlineto{\pgfqpoint{3.222149in}{1.169722in}}%
\pgfpathlineto{\pgfqpoint{3.222801in}{1.177971in}}%
\pgfpathlineto{\pgfqpoint{3.222986in}{1.161473in}}%
\pgfpathlineto{\pgfqpoint{3.223652in}{1.144975in}}%
\pgfpathlineto{\pgfqpoint{3.223822in}{1.169722in}}%
\pgfpathlineto{\pgfqpoint{3.224141in}{1.153224in}}%
\pgfpathlineto{\pgfqpoint{3.224304in}{1.177971in}}%
\pgfpathlineto{\pgfqpoint{3.225163in}{1.144975in}}%
\pgfpathlineto{\pgfqpoint{3.225984in}{1.177971in}}%
\pgfpathlineto{\pgfqpoint{3.226503in}{1.161473in}}%
\pgfpathlineto{\pgfqpoint{3.227658in}{1.128476in}}%
\pgfpathlineto{\pgfqpoint{3.226673in}{1.169722in}}%
\pgfpathlineto{\pgfqpoint{3.227680in}{1.144975in}}%
\pgfpathlineto{\pgfqpoint{3.227828in}{1.153224in}}%
\pgfpathlineto{\pgfqpoint{3.228183in}{1.120227in}}%
\pgfpathlineto{\pgfqpoint{3.229524in}{1.095480in}}%
\pgfpathlineto{\pgfqpoint{3.230864in}{1.169722in}}%
\pgfpathlineto{\pgfqpoint{3.231515in}{1.153224in}}%
\pgfpathlineto{\pgfqpoint{3.232034in}{1.144975in}}%
\pgfpathlineto{\pgfqpoint{3.232352in}{1.177971in}}%
\pgfpathlineto{\pgfqpoint{3.232537in}{1.169722in}}%
\pgfpathlineto{\pgfqpoint{3.232685in}{1.177971in}}%
\pgfpathlineto{\pgfqpoint{3.233041in}{1.144975in}}%
\pgfpathlineto{\pgfqpoint{3.233692in}{1.128476in}}%
\pgfpathlineto{\pgfqpoint{3.233692in}{1.161473in}}%
\pgfpathlineto{\pgfqpoint{3.234699in}{1.177971in}}%
\pgfpathlineto{\pgfqpoint{3.234025in}{1.153224in}}%
\pgfpathlineto{\pgfqpoint{3.234721in}{1.169722in}}%
\pgfpathlineto{\pgfqpoint{3.235032in}{1.177971in}}%
\pgfpathlineto{\pgfqpoint{3.235373in}{1.153224in}}%
\pgfpathlineto{\pgfqpoint{3.235558in}{1.161473in}}%
\pgfpathlineto{\pgfqpoint{3.236372in}{1.153224in}}%
\pgfpathlineto{\pgfqpoint{3.235721in}{1.169722in}}%
\pgfpathlineto{\pgfqpoint{3.236395in}{1.169722in}}%
\pgfpathlineto{\pgfqpoint{3.238224in}{1.252214in}}%
\pgfpathlineto{\pgfqpoint{3.236876in}{1.153224in}}%
\pgfpathlineto{\pgfqpoint{3.238409in}{1.235715in}}%
\pgfpathlineto{\pgfqpoint{3.238890in}{1.202719in}}%
\pgfpathlineto{\pgfqpoint{3.239578in}{1.219217in}}%
\pgfpathlineto{\pgfqpoint{3.240756in}{1.243965in}}%
\pgfpathlineto{\pgfqpoint{3.240919in}{1.235715in}}%
\pgfpathlineto{\pgfqpoint{3.242599in}{1.169722in}}%
\pgfpathlineto{\pgfqpoint{3.241089in}{1.243965in}}%
\pgfpathlineto{\pgfqpoint{3.242762in}{1.194470in}}%
\pgfpathlineto{\pgfqpoint{3.245931in}{1.326456in}}%
\pgfpathlineto{\pgfqpoint{3.245946in}{1.318207in}}%
\pgfpathlineto{\pgfqpoint{3.246598in}{1.301709in}}%
\pgfpathlineto{\pgfqpoint{3.246768in}{1.326456in}}%
\pgfpathlineto{\pgfqpoint{3.246790in}{1.318207in}}%
\pgfpathlineto{\pgfqpoint{3.246938in}{1.326456in}}%
\pgfpathlineto{\pgfqpoint{3.247294in}{1.293460in}}%
\pgfpathlineto{\pgfqpoint{3.247605in}{1.301709in}}%
\pgfpathlineto{\pgfqpoint{3.248612in}{1.276961in}}%
\pgfpathlineto{\pgfqpoint{3.247790in}{1.318207in}}%
\pgfpathlineto{\pgfqpoint{3.248634in}{1.293460in}}%
\pgfpathlineto{\pgfqpoint{3.249115in}{1.301709in}}%
\pgfpathlineto{\pgfqpoint{3.249448in}{1.276961in}}%
\pgfpathlineto{\pgfqpoint{3.249470in}{1.293460in}}%
\pgfpathlineto{\pgfqpoint{3.250455in}{1.276961in}}%
\pgfpathlineto{\pgfqpoint{3.250477in}{1.293460in}}%
\pgfpathlineto{\pgfqpoint{3.252136in}{1.326456in}}%
\pgfpathlineto{\pgfqpoint{3.253158in}{1.293460in}}%
\pgfpathlineto{\pgfqpoint{3.254483in}{1.351204in}}%
\pgfpathlineto{\pgfqpoint{3.254824in}{1.326456in}}%
\pgfpathlineto{\pgfqpoint{3.255157in}{1.301709in}}%
\pgfpathlineto{\pgfqpoint{3.255838in}{1.318207in}}%
\pgfpathlineto{\pgfqpoint{3.256497in}{1.326456in}}%
\pgfpathlineto{\pgfqpoint{3.256675in}{1.309958in}}%
\pgfpathlineto{\pgfqpoint{3.256712in}{1.301709in}}%
\pgfpathlineto{\pgfqpoint{3.256845in}{1.342955in}}%
\pgfpathlineto{\pgfqpoint{3.257519in}{1.342955in}}%
\pgfpathlineto{\pgfqpoint{3.257830in}{1.326456in}}%
\pgfpathlineto{\pgfqpoint{3.257852in}{1.367702in}}%
\pgfpathlineto{\pgfqpoint{3.258170in}{1.359453in}}%
\pgfpathlineto{\pgfqpoint{3.259192in}{1.392449in}}%
\pgfpathlineto{\pgfqpoint{3.261710in}{1.491439in}}%
\pgfpathlineto{\pgfqpoint{3.261873in}{1.483190in}}%
\pgfpathlineto{\pgfqpoint{3.262361in}{1.450194in}}%
\pgfpathlineto{\pgfqpoint{3.262709in}{1.491439in}}%
\pgfpathlineto{\pgfqpoint{3.263035in}{1.499689in}}%
\pgfpathlineto{\pgfqpoint{3.263213in}{1.483190in}}%
\pgfpathlineto{\pgfqpoint{3.264553in}{1.441944in}}%
\pgfpathlineto{\pgfqpoint{3.265397in}{1.491439in}}%
\pgfpathlineto{\pgfqpoint{3.265730in}{1.458443in}}%
\pgfpathlineto{\pgfqpoint{3.271431in}{1.169722in}}%
\pgfpathlineto{\pgfqpoint{3.271594in}{1.194470in}}%
\pgfpathlineto{\pgfqpoint{3.271742in}{1.202719in}}%
\pgfpathlineto{\pgfqpoint{3.272083in}{1.177971in}}%
\pgfpathlineto{\pgfqpoint{3.272268in}{1.186221in}}%
\pgfpathlineto{\pgfqpoint{3.273941in}{1.144975in}}%
\pgfpathlineto{\pgfqpoint{3.274097in}{1.153224in}}%
\pgfpathlineto{\pgfqpoint{3.274445in}{1.120227in}}%
\pgfpathlineto{\pgfqpoint{3.274593in}{1.128476in}}%
\pgfpathlineto{\pgfqpoint{3.275622in}{1.095480in}}%
\pgfpathlineto{\pgfqpoint{3.275778in}{1.103729in}}%
\pgfpathlineto{\pgfqpoint{3.276126in}{1.070732in}}%
\pgfpathlineto{\pgfqpoint{3.276281in}{1.078981in}}%
\pgfpathlineto{\pgfqpoint{3.276629in}{1.095480in}}%
\pgfpathlineto{\pgfqpoint{3.277295in}{1.037736in}}%
\pgfpathlineto{\pgfqpoint{3.277969in}{1.021237in}}%
\pgfpathlineto{\pgfqpoint{3.278132in}{1.045985in}}%
\pgfpathlineto{\pgfqpoint{3.278288in}{1.054234in}}%
\pgfpathlineto{\pgfqpoint{3.278636in}{1.021237in}}%
\pgfpathlineto{\pgfqpoint{3.278791in}{1.029487in}}%
\pgfpathlineto{\pgfqpoint{3.279976in}{0.988241in}}%
\pgfpathlineto{\pgfqpoint{3.280316in}{0.971742in}}%
\pgfpathlineto{\pgfqpoint{3.280479in}{1.021237in}}%
\pgfpathlineto{\pgfqpoint{3.281316in}{0.971742in}}%
\pgfpathlineto{\pgfqpoint{3.281819in}{0.996490in}}%
\pgfpathlineto{\pgfqpoint{3.281967in}{1.004739in}}%
\pgfpathlineto{\pgfqpoint{3.282323in}{0.971742in}}%
\pgfpathlineto{\pgfqpoint{3.282486in}{0.979992in}}%
\pgfpathlineto{\pgfqpoint{3.283500in}{0.922247in}}%
\pgfpathlineto{\pgfqpoint{3.283656in}{0.930497in}}%
\pgfpathlineto{\pgfqpoint{3.283989in}{0.905749in}}%
\pgfpathlineto{\pgfqpoint{3.284167in}{0.913998in}}%
\pgfpathlineto{\pgfqpoint{3.286158in}{0.856254in}}%
\pgfpathlineto{\pgfqpoint{3.284507in}{0.922247in}}%
\pgfpathlineto{\pgfqpoint{3.286166in}{0.864503in}}%
\pgfpathlineto{\pgfqpoint{3.287187in}{0.922247in}}%
\pgfpathlineto{\pgfqpoint{3.287350in}{0.913998in}}%
\pgfpathlineto{\pgfqpoint{3.287691in}{0.872753in}}%
\pgfpathlineto{\pgfqpoint{3.288528in}{0.897500in}}%
\pgfpathlineto{\pgfqpoint{3.289016in}{0.905749in}}%
\pgfpathlineto{\pgfqpoint{3.289201in}{0.889251in}}%
\pgfpathlineto{\pgfqpoint{3.289697in}{0.872753in}}%
\pgfpathlineto{\pgfqpoint{3.289868in}{0.897500in}}%
\pgfpathlineto{\pgfqpoint{3.290186in}{0.889251in}}%
\pgfpathlineto{\pgfqpoint{3.291208in}{0.946995in}}%
\pgfpathlineto{\pgfqpoint{3.293392in}{0.872753in}}%
\pgfpathlineto{\pgfqpoint{3.293548in}{0.881002in}}%
\pgfpathlineto{\pgfqpoint{3.293555in}{0.897500in}}%
\pgfpathlineto{\pgfqpoint{3.293888in}{0.872753in}}%
\pgfpathlineto{\pgfqpoint{3.294562in}{0.872753in}}%
\pgfpathlineto{\pgfqpoint{3.295399in}{0.848005in}}%
\pgfpathlineto{\pgfqpoint{3.295569in}{0.872753in}}%
\pgfpathlineto{\pgfqpoint{3.297746in}{0.971742in}}%
\pgfpathlineto{\pgfqpoint{3.298079in}{0.946995in}}%
\pgfpathlineto{\pgfqpoint{3.300093in}{0.831507in}}%
\pgfpathlineto{\pgfqpoint{3.301100in}{0.848005in}}%
\pgfpathlineto{\pgfqpoint{3.301433in}{0.872753in}}%
\pgfpathlineto{\pgfqpoint{3.301937in}{0.839756in}}%
\pgfpathlineto{\pgfqpoint{3.302277in}{0.831507in}}%
\pgfpathlineto{\pgfqpoint{3.302596in}{0.856254in}}%
\pgfpathlineto{\pgfqpoint{3.302944in}{0.848005in}}%
\pgfpathlineto{\pgfqpoint{3.304121in}{0.806759in}}%
\pgfpathlineto{\pgfqpoint{3.303432in}{0.856254in}}%
\pgfpathlineto{\pgfqpoint{3.304276in}{0.815008in}}%
\pgfpathlineto{\pgfqpoint{3.304284in}{0.848005in}}%
\pgfpathlineto{\pgfqpoint{3.305291in}{0.798510in}}%
\pgfpathlineto{\pgfqpoint{3.306283in}{0.881002in}}%
\pgfpathlineto{\pgfqpoint{3.306801in}{0.823258in}}%
\pgfpathlineto{\pgfqpoint{3.308460in}{0.782012in}}%
\pgfpathlineto{\pgfqpoint{3.308467in}{0.790261in}}%
\pgfpathlineto{\pgfqpoint{3.309481in}{0.848005in}}%
\pgfpathlineto{\pgfqpoint{3.311814in}{0.782012in}}%
\pgfpathlineto{\pgfqpoint{3.311977in}{0.806759in}}%
\pgfpathlineto{\pgfqpoint{3.312665in}{0.773763in}}%
\pgfpathlineto{\pgfqpoint{3.312836in}{0.773763in}}%
\pgfpathlineto{\pgfqpoint{3.313169in}{0.798510in}}%
\pgfpathlineto{\pgfqpoint{3.313672in}{0.765513in}}%
\pgfpathlineto{\pgfqpoint{3.316856in}{0.633527in}}%
\pgfpathlineto{\pgfqpoint{3.317012in}{0.641776in}}%
\pgfpathlineto{\pgfqpoint{3.317360in}{0.674773in}}%
\pgfpathlineto{\pgfqpoint{3.318033in}{0.633527in}}%
\pgfpathlineto{\pgfqpoint{3.318196in}{0.650025in}}%
\pgfpathlineto{\pgfqpoint{3.318537in}{0.625278in}}%
\pgfpathlineto{\pgfqpoint{3.319025in}{0.633527in}}%
\pgfpathlineto{\pgfqpoint{3.320047in}{0.625278in}}%
\pgfpathlineto{\pgfqpoint{3.320877in}{0.650025in}}%
\pgfpathlineto{\pgfqpoint{3.321210in}{0.633527in}}%
\pgfpathlineto{\pgfqpoint{3.322224in}{0.625278in}}%
\pgfpathlineto{\pgfqpoint{3.323216in}{0.658274in}}%
\pgfpathlineto{\pgfqpoint{3.323394in}{0.633527in}}%
\pgfpathlineto{\pgfqpoint{3.324216in}{0.625278in}}%
\pgfpathlineto{\pgfqpoint{3.325578in}{0.633527in}}%
\pgfpathlineto{\pgfqpoint{3.325911in}{0.625278in}}%
\pgfpathlineto{\pgfqpoint{3.326585in}{0.625278in}}%
\pgfpathlineto{\pgfqpoint{3.327903in}{0.633527in}}%
\pgfpathlineto{\pgfqpoint{3.327911in}{0.625278in}}%
\pgfpathlineto{\pgfqpoint{3.328747in}{0.625278in}}%
\pgfpathlineto{\pgfqpoint{3.329917in}{0.633527in}}%
\pgfpathlineto{\pgfqpoint{3.329924in}{0.625278in}}%
\pgfpathlineto{\pgfqpoint{3.330931in}{0.625278in}}%
\pgfpathlineto{\pgfqpoint{3.332094in}{0.633527in}}%
\pgfpathlineto{\pgfqpoint{3.332101in}{0.625278in}}%
\pgfpathlineto{\pgfqpoint{3.332931in}{0.625278in}}%
\pgfpathlineto{\pgfqpoint{3.334278in}{0.633527in}}%
\pgfpathlineto{\pgfqpoint{3.334293in}{0.625278in}}%
\pgfpathlineto{\pgfqpoint{3.335285in}{0.625278in}}%
\pgfpathlineto{\pgfqpoint{3.336625in}{0.633527in}}%
\pgfpathlineto{\pgfqpoint{3.336625in}{0.625278in}}%
\pgfpathlineto{\pgfqpoint{3.337625in}{0.625278in}}%
\pgfpathlineto{\pgfqpoint{3.338802in}{0.633527in}}%
\pgfpathlineto{\pgfqpoint{3.338980in}{0.625278in}}%
\pgfpathlineto{\pgfqpoint{3.339816in}{0.625278in}}%
\pgfpathlineto{\pgfqpoint{3.340994in}{0.633527in}}%
\pgfpathlineto{\pgfqpoint{3.341142in}{0.625278in}}%
\pgfpathlineto{\pgfqpoint{3.341823in}{0.625278in}}%
\pgfpathlineto{\pgfqpoint{3.342993in}{0.633527in}}%
\pgfpathlineto{\pgfqpoint{3.343000in}{0.625278in}}%
\pgfpathlineto{\pgfqpoint{3.344000in}{0.625278in}}%
\pgfpathlineto{\pgfqpoint{3.345170in}{0.633527in}}%
\pgfpathlineto{\pgfqpoint{3.345347in}{0.625278in}}%
\pgfpathlineto{\pgfqpoint{3.346177in}{0.625278in}}%
\pgfpathlineto{\pgfqpoint{3.347354in}{0.633527in}}%
\pgfpathlineto{\pgfqpoint{3.347517in}{0.625278in}}%
\pgfpathlineto{\pgfqpoint{3.348354in}{0.625278in}}%
\pgfpathlineto{\pgfqpoint{3.349523in}{0.633527in}}%
\pgfpathlineto{\pgfqpoint{3.349649in}{0.625278in}}%
\pgfpathlineto{\pgfqpoint{3.350538in}{0.625278in}}%
\pgfpathlineto{\pgfqpoint{3.351708in}{0.633527in}}%
\pgfpathlineto{\pgfqpoint{3.351708in}{0.625278in}}%
\pgfpathlineto{\pgfqpoint{3.352715in}{0.625278in}}%
\pgfpathlineto{\pgfqpoint{3.354055in}{0.633527in}}%
\pgfpathlineto{\pgfqpoint{3.354388in}{0.625278in}}%
\pgfpathlineto{\pgfqpoint{3.355054in}{0.625278in}}%
\pgfpathlineto{\pgfqpoint{3.356232in}{0.633527in}}%
\pgfpathlineto{\pgfqpoint{3.356409in}{0.625278in}}%
\pgfpathlineto{\pgfqpoint{3.356905in}{0.625278in}}%
\pgfpathlineto{\pgfqpoint{3.358246in}{0.633527in}}%
\pgfpathlineto{\pgfqpoint{3.358579in}{0.625278in}}%
\pgfpathlineto{\pgfqpoint{3.359252in}{0.625278in}}%
\pgfpathlineto{\pgfqpoint{3.360593in}{0.633527in}}%
\pgfpathlineto{\pgfqpoint{3.360763in}{0.625278in}}%
\pgfpathlineto{\pgfqpoint{3.361592in}{0.625278in}}%
\pgfpathlineto{\pgfqpoint{3.362769in}{0.633527in}}%
\pgfpathlineto{\pgfqpoint{3.363110in}{0.625278in}}%
\pgfpathlineto{\pgfqpoint{3.363776in}{0.625278in}}%
\pgfpathlineto{\pgfqpoint{3.364946in}{0.633527in}}%
\pgfpathlineto{\pgfqpoint{3.365117in}{0.625278in}}%
\pgfpathlineto{\pgfqpoint{3.365961in}{0.625278in}}%
\pgfpathlineto{\pgfqpoint{3.367131in}{0.633527in}}%
\pgfpathlineto{\pgfqpoint{3.367131in}{0.625278in}}%
\pgfpathlineto{\pgfqpoint{3.367967in}{0.625278in}}%
\pgfpathlineto{\pgfqpoint{3.369137in}{0.633527in}}%
\pgfpathlineto{\pgfqpoint{3.369307in}{0.625278in}}%
\pgfpathlineto{\pgfqpoint{3.370144in}{0.625278in}}%
\pgfpathlineto{\pgfqpoint{3.371321in}{0.633527in}}%
\pgfpathlineto{\pgfqpoint{3.371321in}{0.625278in}}%
\pgfpathlineto{\pgfqpoint{3.372328in}{0.625278in}}%
\pgfpathlineto{\pgfqpoint{3.373498in}{0.633527in}}%
\pgfpathlineto{\pgfqpoint{3.373668in}{0.625278in}}%
\pgfpathlineto{\pgfqpoint{3.374513in}{0.625278in}}%
\pgfpathlineto{\pgfqpoint{3.375675in}{0.633527in}}%
\pgfpathlineto{\pgfqpoint{3.375682in}{0.625278in}}%
\pgfpathlineto{\pgfqpoint{3.376512in}{0.625278in}}%
\pgfpathlineto{\pgfqpoint{3.377859in}{0.633527in}}%
\pgfpathlineto{\pgfqpoint{3.378022in}{0.625278in}}%
\pgfpathlineto{\pgfqpoint{3.378859in}{0.625278in}}%
\pgfpathlineto{\pgfqpoint{3.379866in}{0.633527in}}%
\pgfpathlineto{\pgfqpoint{3.379866in}{0.625278in}}%
\pgfpathlineto{\pgfqpoint{3.381206in}{0.633527in}}%
\pgfpathlineto{\pgfqpoint{3.381206in}{0.625278in}}%
\pgfpathlineto{\pgfqpoint{3.382213in}{0.625278in}}%
\pgfpathlineto{\pgfqpoint{3.383390in}{0.633527in}}%
\pgfpathlineto{\pgfqpoint{3.383723in}{0.625278in}}%
\pgfpathlineto{\pgfqpoint{3.384390in}{0.625278in}}%
\pgfpathlineto{\pgfqpoint{3.385567in}{0.633527in}}%
\pgfpathlineto{\pgfqpoint{3.385737in}{0.625278in}}%
\pgfpathlineto{\pgfqpoint{3.386574in}{0.625278in}}%
\pgfpathlineto{\pgfqpoint{3.387744in}{0.633527in}}%
\pgfpathlineto{\pgfqpoint{3.388247in}{0.625278in}}%
\pgfpathlineto{\pgfqpoint{3.388580in}{0.625278in}}%
\pgfpathlineto{\pgfqpoint{3.389758in}{0.633527in}}%
\pgfpathlineto{\pgfqpoint{3.389928in}{0.625278in}}%
\pgfpathlineto{\pgfqpoint{3.390765in}{0.625278in}}%
\pgfpathlineto{\pgfqpoint{3.391772in}{0.633527in}}%
\pgfpathlineto{\pgfqpoint{3.391772in}{0.625278in}}%
\pgfpathlineto{\pgfqpoint{3.392942in}{0.633527in}}%
\pgfpathlineto{\pgfqpoint{3.393112in}{0.625278in}}%
\pgfpathlineto{\pgfqpoint{3.393949in}{0.625278in}}%
\pgfpathlineto{\pgfqpoint{3.395118in}{0.633527in}}%
\pgfpathlineto{\pgfqpoint{3.395118in}{0.625278in}}%
\pgfpathlineto{\pgfqpoint{3.395962in}{0.625278in}}%
\pgfpathlineto{\pgfqpoint{3.397303in}{0.633527in}}%
\pgfpathlineto{\pgfqpoint{3.397303in}{0.625278in}}%
\pgfpathlineto{\pgfqpoint{3.398310in}{0.625278in}}%
\pgfpathlineto{\pgfqpoint{3.399650in}{0.633527in}}%
\pgfpathlineto{\pgfqpoint{3.399650in}{0.625278in}}%
\pgfpathlineto{\pgfqpoint{3.400486in}{0.625278in}}%
\pgfpathlineto{\pgfqpoint{3.401656in}{0.633527in}}%
\pgfpathlineto{\pgfqpoint{3.401827in}{0.625278in}}%
\pgfpathlineto{\pgfqpoint{3.402663in}{0.625278in}}%
\pgfpathlineto{\pgfqpoint{3.404003in}{0.633527in}}%
\pgfpathlineto{\pgfqpoint{3.404344in}{0.625278in}}%
\pgfpathlineto{\pgfqpoint{3.405010in}{0.625278in}}%
\pgfpathlineto{\pgfqpoint{3.406188in}{0.633527in}}%
\pgfpathlineto{\pgfqpoint{3.406351in}{0.625278in}}%
\pgfpathlineto{\pgfqpoint{3.407195in}{0.625278in}}%
\pgfpathlineto{\pgfqpoint{3.408365in}{0.633527in}}%
\pgfpathlineto{\pgfqpoint{3.408698in}{0.625278in}}%
\pgfpathlineto{\pgfqpoint{3.409371in}{0.625278in}}%
\pgfpathlineto{\pgfqpoint{3.410378in}{0.633527in}}%
\pgfpathlineto{\pgfqpoint{3.410378in}{0.625278in}}%
\pgfpathlineto{\pgfqpoint{3.411548in}{0.633527in}}%
\pgfpathlineto{\pgfqpoint{3.411719in}{0.625278in}}%
\pgfpathlineto{\pgfqpoint{3.412555in}{0.625278in}}%
\pgfpathlineto{\pgfqpoint{3.413562in}{0.633527in}}%
\pgfpathlineto{\pgfqpoint{3.413562in}{0.625278in}}%
\pgfpathlineto{\pgfqpoint{3.414732in}{0.633527in}}%
\pgfpathlineto{\pgfqpoint{3.414902in}{0.625278in}}%
\pgfpathlineto{\pgfqpoint{3.415798in}{0.625278in}}%
\pgfpathlineto{\pgfqpoint{3.416916in}{0.633527in}}%
\pgfpathlineto{\pgfqpoint{3.416916in}{0.625278in}}%
\pgfpathlineto{\pgfqpoint{3.417916in}{0.625278in}}%
\pgfpathlineto{\pgfqpoint{3.419263in}{0.633527in}}%
\pgfpathlineto{\pgfqpoint{3.419263in}{0.625278in}}%
\pgfpathlineto{\pgfqpoint{3.420263in}{0.625278in}}%
\pgfpathlineto{\pgfqpoint{3.421440in}{0.633527in}}%
\pgfpathlineto{\pgfqpoint{3.421611in}{0.625278in}}%
\pgfpathlineto{\pgfqpoint{3.422447in}{0.625278in}}%
\pgfpathlineto{\pgfqpoint{3.423787in}{0.633527in}}%
\pgfpathlineto{\pgfqpoint{3.423787in}{0.625278in}}%
\pgfpathlineto{\pgfqpoint{3.424794in}{0.625278in}}%
\pgfpathlineto{\pgfqpoint{3.425964in}{0.633527in}}%
\pgfpathlineto{\pgfqpoint{3.426297in}{0.625278in}}%
\pgfpathlineto{\pgfqpoint{3.426971in}{0.625278in}}%
\pgfpathlineto{\pgfqpoint{3.428149in}{0.633527in}}%
\pgfpathlineto{\pgfqpoint{3.428311in}{0.625278in}}%
\pgfpathlineto{\pgfqpoint{3.428985in}{0.625278in}}%
\pgfpathlineto{\pgfqpoint{3.430155in}{0.633527in}}%
\pgfpathlineto{\pgfqpoint{3.430325in}{0.625278in}}%
\pgfpathlineto{\pgfqpoint{3.431162in}{0.625278in}}%
\pgfpathlineto{\pgfqpoint{3.432339in}{0.633527in}}%
\pgfpathlineto{\pgfqpoint{3.432672in}{0.625278in}}%
\pgfpathlineto{\pgfqpoint{3.433339in}{0.625278in}}%
\pgfpathlineto{\pgfqpoint{3.434346in}{0.633527in}}%
\pgfpathlineto{\pgfqpoint{3.434346in}{0.625278in}}%
\pgfpathlineto{\pgfqpoint{3.435523in}{0.633527in}}%
\pgfpathlineto{\pgfqpoint{3.435686in}{0.625278in}}%
\pgfpathlineto{\pgfqpoint{3.436530in}{0.625278in}}%
\pgfpathlineto{\pgfqpoint{3.437700in}{0.633527in}}%
\pgfpathlineto{\pgfqpoint{3.437870in}{0.625278in}}%
\pgfpathlineto{\pgfqpoint{3.438707in}{0.625278in}}%
\pgfpathlineto{\pgfqpoint{3.439877in}{0.633527in}}%
\pgfpathlineto{\pgfqpoint{3.440047in}{0.625278in}}%
\pgfpathlineto{\pgfqpoint{3.440884in}{0.625278in}}%
\pgfpathlineto{\pgfqpoint{3.442061in}{0.633527in}}%
\pgfpathlineto{\pgfqpoint{3.442224in}{0.625278in}}%
\pgfpathlineto{\pgfqpoint{3.443061in}{0.625278in}}%
\pgfpathlineto{\pgfqpoint{3.444238in}{0.633527in}}%
\pgfpathlineto{\pgfqpoint{3.444408in}{0.625278in}}%
\pgfpathlineto{\pgfqpoint{3.445245in}{0.625278in}}%
\pgfpathlineto{\pgfqpoint{3.446415in}{0.633527in}}%
\pgfpathlineto{\pgfqpoint{3.446585in}{0.625278in}}%
\pgfpathlineto{\pgfqpoint{3.447422in}{0.625278in}}%
\pgfpathlineto{\pgfqpoint{3.448599in}{0.633527in}}%
\pgfpathlineto{\pgfqpoint{3.448762in}{0.625278in}}%
\pgfpathlineto{\pgfqpoint{3.449598in}{0.625278in}}%
\pgfpathlineto{\pgfqpoint{3.450776in}{0.633527in}}%
\pgfpathlineto{\pgfqpoint{3.450776in}{0.625278in}}%
\pgfpathlineto{\pgfqpoint{3.451783in}{0.625278in}}%
\pgfpathlineto{\pgfqpoint{3.452953in}{0.633527in}}%
\pgfpathlineto{\pgfqpoint{3.453123in}{0.625278in}}%
\pgfpathlineto{\pgfqpoint{3.453960in}{0.625278in}}%
\pgfpathlineto{\pgfqpoint{3.455137in}{0.633527in}}%
\pgfpathlineto{\pgfqpoint{3.455137in}{0.625278in}}%
\pgfpathlineto{\pgfqpoint{3.456136in}{0.625278in}}%
\pgfpathlineto{\pgfqpoint{3.457314in}{0.633527in}}%
\pgfpathlineto{\pgfqpoint{3.457647in}{0.625278in}}%
\pgfpathlineto{\pgfqpoint{3.458321in}{0.625278in}}%
\pgfpathlineto{\pgfqpoint{3.459490in}{0.633527in}}%
\pgfpathlineto{\pgfqpoint{3.459661in}{0.625278in}}%
\pgfpathlineto{\pgfqpoint{3.460497in}{0.625278in}}%
\pgfpathlineto{\pgfqpoint{3.461675in}{0.633527in}}%
\pgfpathlineto{\pgfqpoint{3.461838in}{0.625278in}}%
\pgfpathlineto{\pgfqpoint{3.462674in}{0.625278in}}%
\pgfpathlineto{\pgfqpoint{3.463681in}{0.633527in}}%
\pgfpathlineto{\pgfqpoint{3.463681in}{0.625278in}}%
\pgfpathlineto{\pgfqpoint{3.465021in}{0.633527in}}%
\pgfpathlineto{\pgfqpoint{3.465021in}{0.625278in}}%
\pgfpathlineto{\pgfqpoint{3.466028in}{0.625278in}}%
\pgfpathlineto{\pgfqpoint{3.467206in}{0.633527in}}%
\pgfpathlineto{\pgfqpoint{3.467369in}{0.625278in}}%
\pgfpathlineto{\pgfqpoint{3.468205in}{0.625278in}}%
\pgfpathlineto{\pgfqpoint{3.469382in}{0.633527in}}%
\pgfpathlineto{\pgfqpoint{3.470389in}{0.625278in}}%
\pgfpathlineto{\pgfqpoint{3.470389in}{0.625278in}}%
\pgfpathlineto{\pgfqpoint{3.471559in}{0.633527in}}%
\pgfpathlineto{\pgfqpoint{3.471730in}{0.625278in}}%
\pgfpathlineto{\pgfqpoint{3.472566in}{0.625278in}}%
\pgfpathlineto{\pgfqpoint{3.473744in}{0.633527in}}%
\pgfpathlineto{\pgfqpoint{3.474802in}{0.625278in}}%
\pgfpathlineto{\pgfqpoint{3.476091in}{0.658274in}}%
\pgfpathlineto{\pgfqpoint{3.476105in}{0.650025in}}%
\pgfpathlineto{\pgfqpoint{3.477446in}{0.625278in}}%
\pgfpathlineto{\pgfqpoint{3.478660in}{0.633527in}}%
\pgfpathlineto{\pgfqpoint{3.478771in}{0.625278in}}%
\pgfpathlineto{\pgfqpoint{3.479608in}{0.625278in}}%
\pgfpathlineto{\pgfqpoint{3.480778in}{0.633527in}}%
\pgfpathlineto{\pgfqpoint{3.481784in}{0.625278in}}%
\pgfpathlineto{\pgfqpoint{3.483125in}{0.658274in}}%
\pgfpathlineto{\pgfqpoint{3.483147in}{0.650025in}}%
\pgfpathlineto{\pgfqpoint{3.484487in}{0.625278in}}%
\pgfpathlineto{\pgfqpoint{3.485812in}{0.633527in}}%
\pgfpathlineto{\pgfqpoint{3.485975in}{0.625278in}}%
\pgfpathlineto{\pgfqpoint{3.486479in}{0.658274in}}%
\pgfpathlineto{\pgfqpoint{3.486834in}{0.641776in}}%
\pgfpathlineto{\pgfqpoint{3.487671in}{0.625278in}}%
\pgfpathlineto{\pgfqpoint{3.487819in}{0.641776in}}%
\pgfpathlineto{\pgfqpoint{3.488322in}{0.683022in}}%
\pgfpathlineto{\pgfqpoint{3.488848in}{0.650025in}}%
\pgfpathlineto{\pgfqpoint{3.489833in}{0.658274in}}%
\pgfpathlineto{\pgfqpoint{3.489855in}{0.650025in}}%
\pgfpathlineto{\pgfqpoint{3.490166in}{0.633527in}}%
\pgfpathlineto{\pgfqpoint{3.490507in}{0.666524in}}%
\pgfpathlineto{\pgfqpoint{3.490670in}{0.683022in}}%
\pgfpathlineto{\pgfqpoint{3.491343in}{0.633527in}}%
\pgfpathlineto{\pgfqpoint{3.491528in}{0.641776in}}%
\pgfpathlineto{\pgfqpoint{3.491677in}{0.633527in}}%
\pgfpathlineto{\pgfqpoint{3.491699in}{0.674773in}}%
\pgfpathlineto{\pgfqpoint{3.492032in}{0.674773in}}%
\pgfpathlineto{\pgfqpoint{3.493017in}{0.683022in}}%
\pgfpathlineto{\pgfqpoint{3.492513in}{0.658274in}}%
\pgfpathlineto{\pgfqpoint{3.493039in}{0.674773in}}%
\pgfpathlineto{\pgfqpoint{3.493246in}{0.658274in}}%
\pgfpathlineto{\pgfqpoint{3.493690in}{0.691271in}}%
\pgfpathlineto{\pgfqpoint{3.494697in}{0.707769in}}%
\pgfpathlineto{\pgfqpoint{3.494712in}{0.699520in}}%
\pgfpathlineto{\pgfqpoint{3.495364in}{0.707769in}}%
\pgfpathlineto{\pgfqpoint{3.495549in}{0.691271in}}%
\pgfpathlineto{\pgfqpoint{3.496704in}{0.658274in}}%
\pgfpathlineto{\pgfqpoint{3.496726in}{0.674773in}}%
\pgfpathlineto{\pgfqpoint{3.498214in}{0.732517in}}%
\pgfpathlineto{\pgfqpoint{3.498400in}{0.716019in}}%
\pgfpathlineto{\pgfqpoint{3.499221in}{0.707769in}}%
\pgfpathlineto{\pgfqpoint{3.498570in}{0.724268in}}%
\pgfpathlineto{\pgfqpoint{3.499236in}{0.724268in}}%
\pgfpathlineto{\pgfqpoint{3.500058in}{0.757264in}}%
\pgfpathlineto{\pgfqpoint{3.500391in}{0.732517in}}%
\pgfpathlineto{\pgfqpoint{3.500732in}{0.707769in}}%
\pgfpathlineto{\pgfqpoint{3.501235in}{0.757264in}}%
\pgfpathlineto{\pgfqpoint{3.501420in}{0.749015in}}%
\pgfpathlineto{\pgfqpoint{3.502072in}{0.782012in}}%
\pgfpathlineto{\pgfqpoint{3.502427in}{0.749015in}}%
\pgfpathlineto{\pgfqpoint{3.503427in}{0.724268in}}%
\pgfpathlineto{\pgfqpoint{3.503582in}{0.740766in}}%
\pgfpathlineto{\pgfqpoint{3.503916in}{0.757264in}}%
\pgfpathlineto{\pgfqpoint{3.504271in}{0.724268in}}%
\pgfpathlineto{\pgfqpoint{3.504604in}{0.724268in}}%
\pgfpathlineto{\pgfqpoint{3.505271in}{0.773763in}}%
\pgfpathlineto{\pgfqpoint{3.505774in}{0.740766in}}%
\pgfpathlineto{\pgfqpoint{3.506618in}{0.699520in}}%
\pgfpathlineto{\pgfqpoint{3.506951in}{0.724268in}}%
\pgfpathlineto{\pgfqpoint{3.508610in}{0.757264in}}%
\pgfpathlineto{\pgfqpoint{3.509447in}{0.732517in}}%
\pgfpathlineto{\pgfqpoint{3.509632in}{0.749015in}}%
\pgfpathlineto{\pgfqpoint{3.509950in}{0.782012in}}%
\pgfpathlineto{\pgfqpoint{3.510468in}{0.740766in}}%
\pgfpathlineto{\pgfqpoint{3.510957in}{0.732517in}}%
\pgfpathlineto{\pgfqpoint{3.511120in}{0.757264in}}%
\pgfpathlineto{\pgfqpoint{3.511142in}{0.749015in}}%
\pgfpathlineto{\pgfqpoint{3.511461in}{0.757264in}}%
\pgfpathlineto{\pgfqpoint{3.511809in}{0.724268in}}%
\pgfpathlineto{\pgfqpoint{3.512460in}{0.757264in}}%
\pgfpathlineto{\pgfqpoint{3.512653in}{0.740766in}}%
\pgfpathlineto{\pgfqpoint{3.512853in}{0.707769in}}%
\pgfpathlineto{\pgfqpoint{3.513467in}{0.757264in}}%
\pgfpathlineto{\pgfqpoint{3.513637in}{0.757264in}}%
\pgfpathlineto{\pgfqpoint{3.513652in}{0.773763in}}%
\pgfpathlineto{\pgfqpoint{3.514474in}{0.732517in}}%
\pgfpathlineto{\pgfqpoint{3.514659in}{0.749015in}}%
\pgfpathlineto{\pgfqpoint{3.514978in}{0.757264in}}%
\pgfpathlineto{\pgfqpoint{3.515311in}{0.732517in}}%
\pgfpathlineto{\pgfqpoint{3.516503in}{0.691271in}}%
\pgfpathlineto{\pgfqpoint{3.516651in}{0.683022in}}%
\pgfpathlineto{\pgfqpoint{3.516843in}{0.724268in}}%
\pgfpathlineto{\pgfqpoint{3.518332in}{0.806759in}}%
\pgfpathlineto{\pgfqpoint{3.518850in}{0.790261in}}%
\pgfpathlineto{\pgfqpoint{3.518998in}{0.782012in}}%
\pgfpathlineto{\pgfqpoint{3.519020in}{0.823258in}}%
\pgfpathlineto{\pgfqpoint{3.519339in}{0.815008in}}%
\pgfpathlineto{\pgfqpoint{3.519353in}{0.848005in}}%
\pgfpathlineto{\pgfqpoint{3.520360in}{0.823258in}}%
\pgfpathlineto{\pgfqpoint{3.521012in}{0.806759in}}%
\pgfpathlineto{\pgfqpoint{3.521012in}{0.839756in}}%
\pgfpathlineto{\pgfqpoint{3.521345in}{0.831507in}}%
\pgfpathlineto{\pgfqpoint{3.522034in}{0.848005in}}%
\pgfpathlineto{\pgfqpoint{3.523026in}{0.881002in}}%
\pgfpathlineto{\pgfqpoint{3.523211in}{0.864503in}}%
\pgfpathlineto{\pgfqpoint{3.523692in}{0.856254in}}%
\pgfpathlineto{\pgfqpoint{3.523863in}{0.881002in}}%
\pgfpathlineto{\pgfqpoint{3.523885in}{0.872753in}}%
\pgfpathlineto{\pgfqpoint{3.525388in}{0.922247in}}%
\pgfpathlineto{\pgfqpoint{3.525558in}{0.913998in}}%
\pgfpathlineto{\pgfqpoint{3.526380in}{0.905749in}}%
\pgfpathlineto{\pgfqpoint{3.525728in}{0.922247in}}%
\pgfpathlineto{\pgfqpoint{3.526395in}{0.922247in}}%
\pgfpathlineto{\pgfqpoint{3.527231in}{0.946995in}}%
\pgfpathlineto{\pgfqpoint{3.527550in}{0.930497in}}%
\pgfpathlineto{\pgfqpoint{3.528572in}{0.922247in}}%
\pgfpathlineto{\pgfqpoint{3.529060in}{0.955244in}}%
\pgfpathlineto{\pgfqpoint{3.529919in}{0.938746in}}%
\pgfpathlineto{\pgfqpoint{3.530415in}{0.922247in}}%
\pgfpathlineto{\pgfqpoint{3.530586in}{0.946995in}}%
\pgfpathlineto{\pgfqpoint{3.530904in}{0.938746in}}%
\pgfpathlineto{\pgfqpoint{3.531911in}{0.955244in}}%
\pgfpathlineto{\pgfqpoint{3.531259in}{0.922247in}}%
\pgfpathlineto{\pgfqpoint{3.531926in}{0.946995in}}%
\pgfpathlineto{\pgfqpoint{3.532096in}{0.971742in}}%
\pgfpathlineto{\pgfqpoint{3.532599in}{0.938746in}}%
\pgfpathlineto{\pgfqpoint{3.533606in}{0.897500in}}%
\pgfpathlineto{\pgfqpoint{3.534088in}{0.913998in}}%
\pgfpathlineto{\pgfqpoint{3.534443in}{0.922247in}}%
\pgfpathlineto{\pgfqpoint{3.534776in}{0.897500in}}%
\pgfpathlineto{\pgfqpoint{3.535110in}{0.897500in}}%
\pgfpathlineto{\pgfqpoint{3.536265in}{0.905749in}}%
\pgfpathlineto{\pgfqpoint{3.536457in}{0.897500in}}%
\pgfpathlineto{\pgfqpoint{3.536953in}{0.946995in}}%
\pgfpathlineto{\pgfqpoint{3.537294in}{0.922247in}}%
\pgfpathlineto{\pgfqpoint{3.537945in}{0.930497in}}%
\pgfpathlineto{\pgfqpoint{3.538130in}{0.913998in}}%
\pgfpathlineto{\pgfqpoint{3.539285in}{0.881002in}}%
\pgfpathlineto{\pgfqpoint{3.539300in}{0.897500in}}%
\pgfpathlineto{\pgfqpoint{3.540959in}{0.930497in}}%
\pgfpathlineto{\pgfqpoint{3.540981in}{0.922247in}}%
\pgfpathlineto{\pgfqpoint{3.541633in}{0.955244in}}%
\pgfpathlineto{\pgfqpoint{3.541988in}{0.946995in}}%
\pgfpathlineto{\pgfqpoint{3.543158in}{0.922247in}}%
\pgfpathlineto{\pgfqpoint{3.543306in}{0.930497in}}%
\pgfpathlineto{\pgfqpoint{3.543639in}{0.979992in}}%
\pgfpathlineto{\pgfqpoint{3.544335in}{0.963493in}}%
\pgfpathlineto{\pgfqpoint{3.545172in}{0.946995in}}%
\pgfpathlineto{\pgfqpoint{3.544498in}{0.971742in}}%
\pgfpathlineto{\pgfqpoint{3.545320in}{0.963493in}}%
\pgfpathlineto{\pgfqpoint{3.545823in}{0.979992in}}%
\pgfpathlineto{\pgfqpoint{3.546179in}{0.946995in}}%
\pgfpathlineto{\pgfqpoint{3.546342in}{0.946995in}}%
\pgfpathlineto{\pgfqpoint{3.548171in}{1.004739in}}%
\pgfpathlineto{\pgfqpoint{3.548185in}{0.996490in}}%
\pgfpathlineto{\pgfqpoint{3.548837in}{0.979992in}}%
\pgfpathlineto{\pgfqpoint{3.548837in}{1.012988in}}%
\pgfpathlineto{\pgfqpoint{3.549192in}{1.021237in}}%
\pgfpathlineto{\pgfqpoint{3.549525in}{0.996490in}}%
\pgfpathlineto{\pgfqpoint{3.549866in}{0.996490in}}%
\pgfpathlineto{\pgfqpoint{3.550347in}{1.029487in}}%
\pgfpathlineto{\pgfqpoint{3.551206in}{1.012988in}}%
\pgfpathlineto{\pgfqpoint{3.551710in}{0.996490in}}%
\pgfpathlineto{\pgfqpoint{3.551858in}{1.037736in}}%
\pgfpathlineto{\pgfqpoint{3.552361in}{1.054234in}}%
\pgfpathlineto{\pgfqpoint{3.552880in}{0.996490in}}%
\pgfpathlineto{\pgfqpoint{3.553553in}{1.070732in}}%
\pgfpathlineto{\pgfqpoint{3.554057in}{1.037736in}}%
\pgfpathlineto{\pgfqpoint{3.554538in}{1.004739in}}%
\pgfpathlineto{\pgfqpoint{3.554894in}{1.045985in}}%
\pgfpathlineto{\pgfqpoint{3.555212in}{1.054234in}}%
\pgfpathlineto{\pgfqpoint{3.555545in}{1.029487in}}%
\pgfpathlineto{\pgfqpoint{3.555730in}{1.037736in}}%
\pgfpathlineto{\pgfqpoint{3.556552in}{1.029487in}}%
\pgfpathlineto{\pgfqpoint{3.555900in}{1.045985in}}%
\pgfpathlineto{\pgfqpoint{3.556567in}{1.045985in}}%
\pgfpathlineto{\pgfqpoint{3.557722in}{1.054234in}}%
\pgfpathlineto{\pgfqpoint{3.558914in}{1.021237in}}%
\pgfpathlineto{\pgfqpoint{3.561913in}{1.128476in}}%
\pgfpathlineto{\pgfqpoint{3.563105in}{1.095480in}}%
\pgfpathlineto{\pgfqpoint{3.563253in}{1.103729in}}%
\pgfpathlineto{\pgfqpoint{3.563593in}{1.078981in}}%
\pgfpathlineto{\pgfqpoint{3.563779in}{1.087231in}}%
\pgfpathlineto{\pgfqpoint{3.564593in}{1.078981in}}%
\pgfpathlineto{\pgfqpoint{3.563941in}{1.095480in}}%
\pgfpathlineto{\pgfqpoint{3.564593in}{1.087231in}}%
\pgfpathlineto{\pgfqpoint{3.564615in}{1.120227in}}%
\pgfpathlineto{\pgfqpoint{3.565622in}{1.120227in}}%
\pgfpathlineto{\pgfqpoint{3.566607in}{1.078981in}}%
\pgfpathlineto{\pgfqpoint{3.566962in}{1.095480in}}%
\pgfpathlineto{\pgfqpoint{3.567281in}{1.103729in}}%
\pgfpathlineto{\pgfqpoint{3.567614in}{1.078981in}}%
\pgfpathlineto{\pgfqpoint{3.568132in}{1.070732in}}%
\pgfpathlineto{\pgfqpoint{3.568621in}{1.111978in}}%
\pgfpathlineto{\pgfqpoint{3.568784in}{1.128476in}}%
\pgfpathlineto{\pgfqpoint{3.569458in}{1.078981in}}%
\pgfpathlineto{\pgfqpoint{3.569643in}{1.095480in}}%
\pgfpathlineto{\pgfqpoint{3.571131in}{1.128476in}}%
\pgfpathlineto{\pgfqpoint{3.571153in}{1.120227in}}%
\pgfpathlineto{\pgfqpoint{3.571657in}{1.095480in}}%
\pgfpathlineto{\pgfqpoint{3.571975in}{1.128476in}}%
\pgfpathlineto{\pgfqpoint{3.571990in}{1.120227in}}%
\pgfpathlineto{\pgfqpoint{3.572138in}{1.128476in}}%
\pgfpathlineto{\pgfqpoint{3.572493in}{1.095480in}}%
\pgfpathlineto{\pgfqpoint{3.572701in}{1.078981in}}%
\pgfpathlineto{\pgfqpoint{3.572826in}{1.120227in}}%
\pgfpathlineto{\pgfqpoint{3.573145in}{1.111978in}}%
\pgfpathlineto{\pgfqpoint{3.573167in}{1.144975in}}%
\pgfpathlineto{\pgfqpoint{3.574004in}{1.095480in}}%
\pgfpathlineto{\pgfqpoint{3.574174in}{1.095480in}}%
\pgfpathlineto{\pgfqpoint{3.574507in}{1.120227in}}%
\pgfpathlineto{\pgfqpoint{3.575011in}{1.087231in}}%
\pgfpathlineto{\pgfqpoint{3.576329in}{1.054234in}}%
\pgfpathlineto{\pgfqpoint{3.575677in}{1.095480in}}%
\pgfpathlineto{\pgfqpoint{3.576351in}{1.070732in}}%
\pgfpathlineto{\pgfqpoint{3.576832in}{1.103729in}}%
\pgfpathlineto{\pgfqpoint{3.577188in}{1.070732in}}%
\pgfpathlineto{\pgfqpoint{3.578513in}{1.054234in}}%
\pgfpathlineto{\pgfqpoint{3.579201in}{1.120227in}}%
\pgfpathlineto{\pgfqpoint{3.579535in}{1.095480in}}%
\pgfpathlineto{\pgfqpoint{3.580690in}{1.103729in}}%
\pgfpathlineto{\pgfqpoint{3.581712in}{1.095480in}}%
\pgfpathlineto{\pgfqpoint{3.583370in}{1.128476in}}%
\pgfpathlineto{\pgfqpoint{3.583392in}{1.120227in}}%
\pgfpathlineto{\pgfqpoint{3.584044in}{1.153224in}}%
\pgfpathlineto{\pgfqpoint{3.584399in}{1.144975in}}%
\pgfpathlineto{\pgfqpoint{3.585051in}{1.153224in}}%
\pgfpathlineto{\pgfqpoint{3.585236in}{1.136726in}}%
\pgfpathlineto{\pgfqpoint{3.585717in}{1.128476in}}%
\pgfpathlineto{\pgfqpoint{3.585739in}{1.169722in}}%
\pgfpathlineto{\pgfqpoint{3.586073in}{1.144975in}}%
\pgfpathlineto{\pgfqpoint{3.586894in}{1.153224in}}%
\pgfpathlineto{\pgfqpoint{3.586909in}{1.144975in}}%
\pgfpathlineto{\pgfqpoint{3.587228in}{1.128476in}}%
\pgfpathlineto{\pgfqpoint{3.587398in}{1.153224in}}%
\pgfpathlineto{\pgfqpoint{3.587746in}{1.144975in}}%
\pgfpathlineto{\pgfqpoint{3.587894in}{1.153224in}}%
\pgfpathlineto{\pgfqpoint{3.588249in}{1.120227in}}%
\pgfpathlineto{\pgfqpoint{3.588398in}{1.128476in}}%
\pgfpathlineto{\pgfqpoint{3.588738in}{1.103729in}}%
\pgfpathlineto{\pgfqpoint{3.589427in}{1.144975in}}%
\pgfpathlineto{\pgfqpoint{3.590597in}{1.120227in}}%
\pgfpathlineto{\pgfqpoint{3.589797in}{1.153224in}}%
\pgfpathlineto{\pgfqpoint{3.590745in}{1.128476in}}%
\pgfpathlineto{\pgfqpoint{3.591774in}{1.219217in}}%
\pgfpathlineto{\pgfqpoint{3.592440in}{1.194470in}}%
\pgfpathlineto{\pgfqpoint{3.592610in}{1.219217in}}%
\pgfpathlineto{\pgfqpoint{3.593766in}{1.227466in}}%
\pgfpathlineto{\pgfqpoint{3.594099in}{1.202719in}}%
\pgfpathlineto{\pgfqpoint{3.594773in}{1.276961in}}%
\pgfpathlineto{\pgfqpoint{3.594787in}{1.268712in}}%
\pgfpathlineto{\pgfqpoint{3.594935in}{1.276961in}}%
\pgfpathlineto{\pgfqpoint{3.595291in}{1.243965in}}%
\pgfpathlineto{\pgfqpoint{3.595942in}{1.227466in}}%
\pgfpathlineto{\pgfqpoint{3.596113in}{1.252214in}}%
\pgfpathlineto{\pgfqpoint{3.596127in}{1.243965in}}%
\pgfpathlineto{\pgfqpoint{3.596779in}{1.252214in}}%
\pgfpathlineto{\pgfqpoint{3.596972in}{1.235715in}}%
\pgfpathlineto{\pgfqpoint{3.597009in}{1.227466in}}%
\pgfpathlineto{\pgfqpoint{3.597134in}{1.268712in}}%
\pgfpathlineto{\pgfqpoint{3.597971in}{1.243965in}}%
\pgfpathlineto{\pgfqpoint{3.599126in}{1.252214in}}%
\pgfpathlineto{\pgfqpoint{3.599467in}{1.227466in}}%
\pgfpathlineto{\pgfqpoint{3.600155in}{1.243965in}}%
\pgfpathlineto{\pgfqpoint{3.600822in}{1.219217in}}%
\pgfpathlineto{\pgfqpoint{3.600992in}{1.243965in}}%
\pgfpathlineto{\pgfqpoint{3.601310in}{1.252214in}}%
\pgfpathlineto{\pgfqpoint{3.601658in}{1.219217in}}%
\pgfpathlineto{\pgfqpoint{3.602310in}{1.202719in}}%
\pgfpathlineto{\pgfqpoint{3.602369in}{1.227466in}}%
\pgfpathlineto{\pgfqpoint{3.602502in}{1.219217in}}%
\pgfpathlineto{\pgfqpoint{3.603672in}{1.243965in}}%
\pgfpathlineto{\pgfqpoint{3.603843in}{1.235715in}}%
\pgfpathlineto{\pgfqpoint{3.604324in}{1.227466in}}%
\pgfpathlineto{\pgfqpoint{3.604494in}{1.252214in}}%
\pgfpathlineto{\pgfqpoint{3.604679in}{1.243965in}}%
\pgfpathlineto{\pgfqpoint{3.606508in}{1.301709in}}%
\pgfpathlineto{\pgfqpoint{3.606523in}{1.293460in}}%
\pgfpathlineto{\pgfqpoint{3.606841in}{1.276961in}}%
\pgfpathlineto{\pgfqpoint{3.607004in}{1.301709in}}%
\pgfpathlineto{\pgfqpoint{3.607360in}{1.293460in}}%
\pgfpathlineto{\pgfqpoint{3.608367in}{1.318207in}}%
\pgfpathlineto{\pgfqpoint{3.608537in}{1.309958in}}%
\pgfpathlineto{\pgfqpoint{3.609544in}{1.293460in}}%
\pgfpathlineto{\pgfqpoint{3.608700in}{1.318207in}}%
\pgfpathlineto{\pgfqpoint{3.609692in}{1.301709in}}%
\pgfpathlineto{\pgfqpoint{3.609707in}{1.318207in}}%
\pgfpathlineto{\pgfqpoint{3.610381in}{1.293460in}}%
\pgfpathlineto{\pgfqpoint{3.610714in}{1.318207in}}%
\pgfpathlineto{\pgfqpoint{3.611550in}{1.293460in}}%
\pgfpathlineto{\pgfqpoint{3.612039in}{1.301709in}}%
\pgfpathlineto{\pgfqpoint{3.612202in}{1.326456in}}%
\pgfpathlineto{\pgfqpoint{3.613061in}{1.243965in}}%
\pgfpathlineto{\pgfqpoint{3.614216in}{1.252214in}}%
\pgfpathlineto{\pgfqpoint{3.614905in}{1.219217in}}%
\pgfpathlineto{\pgfqpoint{3.615238in}{1.268712in}}%
\pgfpathlineto{\pgfqpoint{3.616082in}{1.293460in}}%
\pgfpathlineto{\pgfqpoint{3.616400in}{1.276961in}}%
\pgfpathlineto{\pgfqpoint{3.617229in}{1.252214in}}%
\pgfpathlineto{\pgfqpoint{3.617422in}{1.268712in}}%
\pgfpathlineto{\pgfqpoint{3.618074in}{1.326456in}}%
\pgfpathlineto{\pgfqpoint{3.618592in}{1.285210in}}%
\pgfpathlineto{\pgfqpoint{3.620080in}{1.252214in}}%
\pgfpathlineto{\pgfqpoint{3.620080in}{1.260463in}}%
\pgfpathlineto{\pgfqpoint{3.620102in}{1.268712in}}%
\pgfpathlineto{\pgfqpoint{3.620924in}{1.202719in}}%
\pgfpathlineto{\pgfqpoint{3.621087in}{1.227466in}}%
\pgfpathlineto{\pgfqpoint{3.621109in}{1.219217in}}%
\pgfpathlineto{\pgfqpoint{3.621924in}{1.276961in}}%
\pgfpathlineto{\pgfqpoint{3.621946in}{1.268712in}}%
\pgfpathlineto{\pgfqpoint{3.623279in}{1.301709in}}%
\pgfpathlineto{\pgfqpoint{3.623286in}{1.293460in}}%
\pgfpathlineto{\pgfqpoint{3.624293in}{1.268712in}}%
\pgfpathlineto{\pgfqpoint{3.624456in}{1.285210in}}%
\pgfpathlineto{\pgfqpoint{3.624611in}{1.301709in}}%
\pgfpathlineto{\pgfqpoint{3.624959in}{1.268712in}}%
\pgfpathlineto{\pgfqpoint{3.625463in}{1.268712in}}%
\pgfpathlineto{\pgfqpoint{3.628314in}{1.169722in}}%
\pgfpathlineto{\pgfqpoint{3.629150in}{1.243965in}}%
\pgfpathlineto{\pgfqpoint{3.629994in}{1.210968in}}%
\pgfpathlineto{\pgfqpoint{3.630142in}{1.202719in}}%
\pgfpathlineto{\pgfqpoint{3.630157in}{1.243965in}}%
\pgfpathlineto{\pgfqpoint{3.630809in}{1.227466in}}%
\pgfpathlineto{\pgfqpoint{3.630831in}{1.268712in}}%
\pgfpathlineto{\pgfqpoint{3.631838in}{1.268712in}}%
\pgfpathlineto{\pgfqpoint{3.634163in}{1.202719in}}%
\pgfpathlineto{\pgfqpoint{3.634518in}{1.268712in}}%
\pgfpathlineto{\pgfqpoint{3.635185in}{1.243965in}}%
\pgfpathlineto{\pgfqpoint{3.635503in}{1.227466in}}%
\pgfpathlineto{\pgfqpoint{3.635673in}{1.252214in}}%
\pgfpathlineto{\pgfqpoint{3.636029in}{1.243965in}}%
\pgfpathlineto{\pgfqpoint{3.636184in}{1.252214in}}%
\pgfpathlineto{\pgfqpoint{3.636510in}{1.227466in}}%
\pgfpathlineto{\pgfqpoint{3.636695in}{1.235715in}}%
\pgfpathlineto{\pgfqpoint{3.636843in}{1.227466in}}%
\pgfpathlineto{\pgfqpoint{3.637021in}{1.252214in}}%
\pgfpathlineto{\pgfqpoint{3.637369in}{1.243965in}}%
\pgfpathlineto{\pgfqpoint{3.637702in}{1.268712in}}%
\pgfpathlineto{\pgfqpoint{3.638206in}{1.235715in}}%
\pgfpathlineto{\pgfqpoint{3.639879in}{1.144975in}}%
\pgfpathlineto{\pgfqpoint{3.641034in}{1.128476in}}%
\pgfpathlineto{\pgfqpoint{3.640375in}{1.153224in}}%
\pgfpathlineto{\pgfqpoint{3.641041in}{1.136726in}}%
\pgfpathlineto{\pgfqpoint{3.642048in}{1.153224in}}%
\pgfpathlineto{\pgfqpoint{3.642063in}{1.144975in}}%
\pgfpathlineto{\pgfqpoint{3.643736in}{1.070732in}}%
\pgfpathlineto{\pgfqpoint{3.643907in}{1.095480in}}%
\pgfpathlineto{\pgfqpoint{3.644729in}{1.153224in}}%
\pgfpathlineto{\pgfqpoint{3.645580in}{1.136726in}}%
\pgfpathlineto{\pgfqpoint{3.648098in}{1.045985in}}%
\pgfpathlineto{\pgfqpoint{3.648749in}{1.078981in}}%
\pgfpathlineto{\pgfqpoint{3.649260in}{1.054234in}}%
\pgfpathlineto{\pgfqpoint{3.650274in}{1.045985in}}%
\pgfpathlineto{\pgfqpoint{3.651111in}{1.021237in}}%
\pgfpathlineto{\pgfqpoint{3.651281in}{1.045985in}}%
\pgfpathlineto{\pgfqpoint{3.651615in}{1.070732in}}%
\pgfpathlineto{\pgfqpoint{3.651948in}{1.037736in}}%
\pgfpathlineto{\pgfqpoint{3.652103in}{1.029487in}}%
\pgfpathlineto{\pgfqpoint{3.652444in}{1.062483in}}%
\pgfpathlineto{\pgfqpoint{3.652940in}{1.078981in}}%
\pgfpathlineto{\pgfqpoint{3.653295in}{1.045985in}}%
\pgfpathlineto{\pgfqpoint{3.653458in}{1.045985in}}%
\pgfpathlineto{\pgfqpoint{3.654110in}{1.029487in}}%
\pgfpathlineto{\pgfqpoint{3.654287in}{1.054234in}}%
\pgfpathlineto{\pgfqpoint{3.654295in}{1.045985in}}%
\pgfpathlineto{\pgfqpoint{3.655642in}{1.095480in}}%
\pgfpathlineto{\pgfqpoint{3.655961in}{1.078981in}}%
\pgfpathlineto{\pgfqpoint{3.656983in}{1.045985in}}%
\pgfpathlineto{\pgfqpoint{3.657819in}{1.070732in}}%
\pgfpathlineto{\pgfqpoint{3.658138in}{1.054234in}}%
\pgfpathlineto{\pgfqpoint{3.658804in}{1.029487in}}%
\pgfpathlineto{\pgfqpoint{3.658982in}{1.078981in}}%
\pgfpathlineto{\pgfqpoint{3.659159in}{1.054234in}}%
\pgfpathlineto{\pgfqpoint{3.659330in}{1.070732in}}%
\pgfpathlineto{\pgfqpoint{3.659833in}{1.037736in}}%
\pgfpathlineto{\pgfqpoint{3.661507in}{0.996490in}}%
\pgfpathlineto{\pgfqpoint{3.661825in}{1.004739in}}%
\pgfpathlineto{\pgfqpoint{3.662165in}{0.979992in}}%
\pgfpathlineto{\pgfqpoint{3.662513in}{0.996490in}}%
\pgfpathlineto{\pgfqpoint{3.663350in}{0.930497in}}%
\pgfpathlineto{\pgfqpoint{3.664009in}{0.955244in}}%
\pgfpathlineto{\pgfqpoint{3.664342in}{0.930497in}}%
\pgfpathlineto{\pgfqpoint{3.664690in}{0.905749in}}%
\pgfpathlineto{\pgfqpoint{3.664861in}{0.946995in}}%
\pgfpathlineto{\pgfqpoint{3.665364in}{0.938746in}}%
\pgfpathlineto{\pgfqpoint{3.666030in}{0.897500in}}%
\pgfpathlineto{\pgfqpoint{3.666534in}{0.922247in}}%
\pgfpathlineto{\pgfqpoint{3.667037in}{0.946995in}}%
\pgfpathlineto{\pgfqpoint{3.666875in}{0.905749in}}%
\pgfpathlineto{\pgfqpoint{3.667371in}{0.913998in}}%
\pgfpathlineto{\pgfqpoint{3.667689in}{0.930497in}}%
\pgfpathlineto{\pgfqpoint{3.668385in}{0.856254in}}%
\pgfpathlineto{\pgfqpoint{3.668881in}{0.872753in}}%
\pgfpathlineto{\pgfqpoint{3.669214in}{0.848005in}}%
\pgfpathlineto{\pgfqpoint{3.672065in}{0.749015in}}%
\pgfpathlineto{\pgfqpoint{3.672568in}{0.757264in}}%
\pgfpathlineto{\pgfqpoint{3.672739in}{0.740766in}}%
\pgfpathlineto{\pgfqpoint{3.674079in}{0.699520in}}%
\pgfpathlineto{\pgfqpoint{3.674249in}{0.724268in}}%
\pgfpathlineto{\pgfqpoint{3.675086in}{0.724268in}}%
\pgfpathlineto{\pgfqpoint{3.677922in}{0.625278in}}%
\pgfpathlineto{\pgfqpoint{3.678929in}{0.633527in}}%
\pgfpathlineto{\pgfqpoint{3.678929in}{0.625278in}}%
\pgfpathlineto{\pgfqpoint{3.680261in}{0.633527in}}%
\pgfpathlineto{\pgfqpoint{3.680284in}{0.625278in}}%
\pgfpathlineto{\pgfqpoint{3.681120in}{0.650025in}}%
\pgfpathlineto{\pgfqpoint{3.681291in}{0.633527in}}%
\pgfpathlineto{\pgfqpoint{3.682127in}{0.650025in}}%
\pgfpathlineto{\pgfqpoint{3.682290in}{0.641776in}}%
\pgfpathlineto{\pgfqpoint{3.682964in}{0.625278in}}%
\pgfpathlineto{\pgfqpoint{3.683134in}{0.650025in}}%
\pgfpathlineto{\pgfqpoint{3.683297in}{0.633527in}}%
\pgfpathlineto{\pgfqpoint{3.684134in}{0.674773in}}%
\pgfpathlineto{\pgfqpoint{3.684474in}{0.650025in}}%
\pgfpathlineto{\pgfqpoint{3.685814in}{0.625278in}}%
\pgfpathlineto{\pgfqpoint{3.686488in}{0.658274in}}%
\pgfpathlineto{\pgfqpoint{3.686821in}{0.633527in}}%
\pgfpathlineto{\pgfqpoint{3.687658in}{0.625278in}}%
\pgfpathlineto{\pgfqpoint{3.688821in}{0.633527in}}%
\pgfpathlineto{\pgfqpoint{3.689154in}{0.625278in}}%
\pgfpathlineto{\pgfqpoint{3.689665in}{0.625278in}}%
\pgfpathlineto{\pgfqpoint{3.690997in}{0.633527in}}%
\pgfpathlineto{\pgfqpoint{3.691012in}{0.625278in}}%
\pgfpathlineto{\pgfqpoint{3.692019in}{0.625278in}}%
\pgfpathlineto{\pgfqpoint{3.693174in}{0.633527in}}%
\pgfpathlineto{\pgfqpoint{3.693352in}{0.625278in}}%
\pgfpathlineto{\pgfqpoint{3.694018in}{0.625278in}}%
\pgfpathlineto{\pgfqpoint{3.695196in}{0.633527in}}%
\pgfpathlineto{\pgfqpoint{3.695373in}{0.625278in}}%
\pgfpathlineto{\pgfqpoint{3.696203in}{0.625278in}}%
\pgfpathlineto{\pgfqpoint{3.697372in}{0.633527in}}%
\pgfpathlineto{\pgfqpoint{3.697535in}{0.625278in}}%
\pgfpathlineto{\pgfqpoint{3.698209in}{0.625278in}}%
\pgfpathlineto{\pgfqpoint{3.699386in}{0.633527in}}%
\pgfpathlineto{\pgfqpoint{3.699549in}{0.625278in}}%
\pgfpathlineto{\pgfqpoint{3.700216in}{0.625278in}}%
\pgfpathlineto{\pgfqpoint{3.701393in}{0.633527in}}%
\pgfpathlineto{\pgfqpoint{3.701400in}{0.625278in}}%
\pgfpathlineto{\pgfqpoint{3.702400in}{0.625278in}}%
\pgfpathlineto{\pgfqpoint{3.703577in}{0.633527in}}%
\pgfpathlineto{\pgfqpoint{3.703740in}{0.625278in}}%
\pgfpathlineto{\pgfqpoint{3.704584in}{0.625278in}}%
\pgfpathlineto{\pgfqpoint{3.705747in}{0.633527in}}%
\pgfpathlineto{\pgfqpoint{3.706080in}{0.625278in}}%
\pgfpathlineto{\pgfqpoint{3.706746in}{0.625278in}}%
\pgfpathlineto{\pgfqpoint{3.707427in}{0.633527in}}%
\pgfpathlineto{\pgfqpoint{3.707427in}{0.625278in}}%
\pgfpathlineto{\pgfqpoint{3.708590in}{0.633527in}}%
\pgfpathlineto{\pgfqpoint{3.709264in}{0.625278in}}%
\pgfpathlineto{\pgfqpoint{3.709612in}{0.625278in}}%
\pgfpathlineto{\pgfqpoint{3.710781in}{0.633527in}}%
\pgfpathlineto{\pgfqpoint{3.710937in}{0.625278in}}%
\pgfpathlineto{\pgfqpoint{3.711618in}{0.625278in}}%
\pgfpathlineto{\pgfqpoint{3.712788in}{0.633527in}}%
\pgfpathlineto{\pgfqpoint{3.712958in}{0.625278in}}%
\pgfpathlineto{\pgfqpoint{3.713795in}{0.625278in}}%
\pgfpathlineto{\pgfqpoint{3.714802in}{0.633527in}}%
\pgfpathlineto{\pgfqpoint{3.714802in}{0.625278in}}%
\pgfpathlineto{\pgfqpoint{3.716135in}{0.633527in}}%
\pgfpathlineto{\pgfqpoint{3.716305in}{0.625278in}}%
\pgfpathlineto{\pgfqpoint{3.717149in}{0.625278in}}%
\pgfpathlineto{\pgfqpoint{3.718312in}{0.633527in}}%
\pgfpathlineto{\pgfqpoint{3.718319in}{0.625278in}}%
\pgfpathlineto{\pgfqpoint{3.719163in}{0.625278in}}%
\pgfpathlineto{\pgfqpoint{3.720325in}{0.633527in}}%
\pgfpathlineto{\pgfqpoint{3.720503in}{0.625278in}}%
\pgfpathlineto{\pgfqpoint{3.721332in}{0.625278in}}%
\pgfpathlineto{\pgfqpoint{3.722502in}{0.633527in}}%
\pgfpathlineto{\pgfqpoint{3.722510in}{0.625278in}}%
\pgfpathlineto{\pgfqpoint{3.723509in}{0.625278in}}%
\pgfpathlineto{\pgfqpoint{3.724687in}{0.633527in}}%
\pgfpathlineto{\pgfqpoint{3.724849in}{0.625278in}}%
\pgfpathlineto{\pgfqpoint{3.725693in}{0.625278in}}%
\pgfpathlineto{\pgfqpoint{3.727034in}{0.633527in}}%
\pgfpathlineto{\pgfqpoint{3.727034in}{0.625278in}}%
\pgfpathlineto{\pgfqpoint{3.728041in}{0.625278in}}%
\pgfpathlineto{\pgfqpoint{3.729210in}{0.633527in}}%
\pgfpathlineto{\pgfqpoint{3.729388in}{0.625278in}}%
\pgfpathlineto{\pgfqpoint{3.730217in}{0.625278in}}%
\pgfpathlineto{\pgfqpoint{3.731387in}{0.633527in}}%
\pgfpathlineto{\pgfqpoint{3.731558in}{0.625278in}}%
\pgfpathlineto{\pgfqpoint{3.732394in}{0.625278in}}%
\pgfpathlineto{\pgfqpoint{3.733572in}{0.633527in}}%
\pgfpathlineto{\pgfqpoint{3.733572in}{0.625278in}}%
\pgfpathlineto{\pgfqpoint{3.734571in}{0.625278in}}%
\pgfpathlineto{\pgfqpoint{3.735748in}{0.633527in}}%
\pgfpathlineto{\pgfqpoint{3.735919in}{0.625278in}}%
\pgfpathlineto{\pgfqpoint{3.736755in}{0.625278in}}%
\pgfpathlineto{\pgfqpoint{3.737925in}{0.633527in}}%
\pgfpathlineto{\pgfqpoint{3.737933in}{0.625278in}}%
\pgfpathlineto{\pgfqpoint{3.738932in}{0.625278in}}%
\pgfpathlineto{\pgfqpoint{3.740272in}{0.633527in}}%
\pgfpathlineto{\pgfqpoint{3.740606in}{0.625278in}}%
\pgfpathlineto{\pgfqpoint{3.741279in}{0.625278in}}%
\pgfpathlineto{\pgfqpoint{3.742457in}{0.633527in}}%
\pgfpathlineto{\pgfqpoint{3.742457in}{0.625278in}}%
\pgfpathlineto{\pgfqpoint{3.743456in}{0.625278in}}%
\pgfpathlineto{\pgfqpoint{3.744796in}{0.633527in}}%
\pgfpathlineto{\pgfqpoint{3.745137in}{0.625278in}}%
\pgfpathlineto{\pgfqpoint{3.745803in}{0.625278in}}%
\pgfpathlineto{\pgfqpoint{3.746981in}{0.633527in}}%
\pgfpathlineto{\pgfqpoint{3.747143in}{0.625278in}}%
\pgfpathlineto{\pgfqpoint{3.747988in}{0.625278in}}%
\pgfpathlineto{\pgfqpoint{3.749157in}{0.633527in}}%
\pgfpathlineto{\pgfqpoint{3.749157in}{0.625278in}}%
\pgfpathlineto{\pgfqpoint{3.750164in}{0.625278in}}%
\pgfpathlineto{\pgfqpoint{3.751334in}{0.633527in}}%
\pgfpathlineto{\pgfqpoint{3.751675in}{0.625278in}}%
\pgfpathlineto{\pgfqpoint{3.752341in}{0.625278in}}%
\pgfpathlineto{\pgfqpoint{3.753681in}{0.633527in}}%
\pgfpathlineto{\pgfqpoint{3.753852in}{0.625278in}}%
\pgfpathlineto{\pgfqpoint{3.754688in}{0.625278in}}%
\pgfpathlineto{\pgfqpoint{3.755866in}{0.633527in}}%
\pgfpathlineto{\pgfqpoint{3.756028in}{0.625278in}}%
\pgfpathlineto{\pgfqpoint{3.756873in}{0.625278in}}%
\pgfpathlineto{\pgfqpoint{3.758042in}{0.633527in}}%
\pgfpathlineto{\pgfqpoint{3.758213in}{0.625278in}}%
\pgfpathlineto{\pgfqpoint{3.759049in}{0.625278in}}%
\pgfpathlineto{\pgfqpoint{3.760219in}{0.633527in}}%
\pgfpathlineto{\pgfqpoint{3.760219in}{0.625278in}}%
\pgfpathlineto{\pgfqpoint{3.761226in}{0.625278in}}%
\pgfpathlineto{\pgfqpoint{3.762566in}{0.633527in}}%
\pgfpathlineto{\pgfqpoint{3.762737in}{0.625278in}}%
\pgfpathlineto{\pgfqpoint{3.763573in}{0.625278in}}%
\pgfpathlineto{\pgfqpoint{3.764580in}{0.633527in}}%
\pgfpathlineto{\pgfqpoint{3.764580in}{0.625278in}}%
\pgfpathlineto{\pgfqpoint{3.765750in}{0.633527in}}%
\pgfpathlineto{\pgfqpoint{3.765920in}{0.625278in}}%
\pgfpathlineto{\pgfqpoint{3.766757in}{0.625278in}}%
\pgfpathlineto{\pgfqpoint{3.767934in}{0.633527in}}%
\pgfpathlineto{\pgfqpoint{3.767934in}{0.625278in}}%
\pgfpathlineto{\pgfqpoint{3.768941in}{0.625278in}}%
\pgfpathlineto{\pgfqpoint{3.770111in}{0.633527in}}%
\pgfpathlineto{\pgfqpoint{3.770111in}{0.625278in}}%
\pgfpathlineto{\pgfqpoint{3.771118in}{0.625278in}}%
\pgfpathlineto{\pgfqpoint{3.772458in}{0.633527in}}%
\pgfpathlineto{\pgfqpoint{3.772458in}{0.625278in}}%
\pgfpathlineto{\pgfqpoint{3.773465in}{0.625278in}}%
\pgfpathlineto{\pgfqpoint{3.774806in}{0.633527in}}%
\pgfpathlineto{\pgfqpoint{3.774806in}{0.625278in}}%
\pgfpathlineto{\pgfqpoint{3.775812in}{0.625278in}}%
\pgfpathlineto{\pgfqpoint{3.777042in}{0.633527in}}%
\pgfpathlineto{\pgfqpoint{3.777153in}{0.625278in}}%
\pgfpathlineto{\pgfqpoint{3.777989in}{0.625278in}}%
\pgfpathlineto{\pgfqpoint{3.779167in}{0.633527in}}%
\pgfpathlineto{\pgfqpoint{3.779500in}{0.625278in}}%
\pgfpathlineto{\pgfqpoint{3.780174in}{0.625278in}}%
\pgfpathlineto{\pgfqpoint{3.781343in}{0.633527in}}%
\pgfpathlineto{\pgfqpoint{3.781403in}{0.625278in}}%
\pgfpathlineto{\pgfqpoint{3.782350in}{0.625278in}}%
\pgfpathlineto{\pgfqpoint{3.783520in}{0.633527in}}%
\pgfpathlineto{\pgfqpoint{3.783691in}{0.625278in}}%
\pgfpathlineto{\pgfqpoint{3.784364in}{0.625278in}}%
\pgfpathlineto{\pgfqpoint{3.785704in}{0.633527in}}%
\pgfpathlineto{\pgfqpoint{3.785704in}{0.625278in}}%
\pgfpathlineto{\pgfqpoint{3.786704in}{0.625278in}}%
\pgfpathlineto{\pgfqpoint{3.787881in}{0.633527in}}%
\pgfpathlineto{\pgfqpoint{3.788052in}{0.625278in}}%
\pgfpathlineto{\pgfqpoint{3.788888in}{0.625278in}}%
\pgfpathlineto{\pgfqpoint{3.790058in}{0.633527in}}%
\pgfpathlineto{\pgfqpoint{3.790228in}{0.625278in}}%
\pgfpathlineto{\pgfqpoint{3.791065in}{0.625278in}}%
\pgfpathlineto{\pgfqpoint{3.792072in}{0.633527in}}%
\pgfpathlineto{\pgfqpoint{3.792072in}{0.625278in}}%
\pgfpathlineto{\pgfqpoint{3.793242in}{0.633527in}}%
\pgfpathlineto{\pgfqpoint{3.793242in}{0.625278in}}%
\pgfpathlineto{\pgfqpoint{3.794249in}{0.625278in}}%
\pgfpathlineto{\pgfqpoint{3.795426in}{0.633527in}}%
\pgfpathlineto{\pgfqpoint{3.795589in}{0.625278in}}%
\pgfpathlineto{\pgfqpoint{3.796433in}{0.625278in}}%
\pgfpathlineto{\pgfqpoint{3.797773in}{0.633527in}}%
\pgfpathlineto{\pgfqpoint{3.797773in}{0.625278in}}%
\pgfpathlineto{\pgfqpoint{3.798780in}{0.625278in}}%
\pgfpathlineto{\pgfqpoint{3.799950in}{0.633527in}}%
\pgfpathlineto{\pgfqpoint{3.800120in}{0.625278in}}%
\pgfpathlineto{\pgfqpoint{3.800957in}{0.625278in}}%
\pgfpathlineto{\pgfqpoint{3.802127in}{0.633527in}}%
\pgfpathlineto{\pgfqpoint{3.802297in}{0.625278in}}%
\pgfpathlineto{\pgfqpoint{3.803134in}{0.625278in}}%
\pgfpathlineto{\pgfqpoint{3.804311in}{0.633527in}}%
\pgfpathlineto{\pgfqpoint{3.804311in}{0.625278in}}%
\pgfpathlineto{\pgfqpoint{3.805318in}{0.625278in}}%
\pgfpathlineto{\pgfqpoint{3.806488in}{0.633527in}}%
\pgfpathlineto{\pgfqpoint{3.806488in}{0.625278in}}%
\pgfpathlineto{\pgfqpoint{3.807495in}{0.625278in}}%
\pgfpathlineto{\pgfqpoint{3.808665in}{0.633527in}}%
\pgfpathlineto{\pgfqpoint{3.808835in}{0.625278in}}%
\pgfpathlineto{\pgfqpoint{3.809509in}{0.625278in}}%
\pgfpathlineto{\pgfqpoint{3.810679in}{0.633527in}}%
\pgfpathlineto{\pgfqpoint{3.810849in}{0.625278in}}%
\pgfpathlineto{\pgfqpoint{3.811686in}{0.625278in}}%
\pgfpathlineto{\pgfqpoint{3.813026in}{0.633527in}}%
\pgfpathlineto{\pgfqpoint{3.813026in}{0.625278in}}%
\pgfpathlineto{\pgfqpoint{3.814033in}{0.625278in}}%
\pgfpathlineto{\pgfqpoint{3.815203in}{0.633527in}}%
\pgfpathlineto{\pgfqpoint{3.815373in}{0.625278in}}%
\pgfpathlineto{\pgfqpoint{3.816039in}{0.625278in}}%
\pgfpathlineto{\pgfqpoint{3.817217in}{0.633527in}}%
\pgfpathlineto{\pgfqpoint{3.817387in}{0.625278in}}%
\pgfpathlineto{\pgfqpoint{3.818053in}{0.625278in}}%
\pgfpathlineto{\pgfqpoint{3.819394in}{0.633527in}}%
\pgfpathlineto{\pgfqpoint{3.819394in}{0.625278in}}%
\pgfpathlineto{\pgfqpoint{3.820401in}{0.625278in}}%
\pgfpathlineto{\pgfqpoint{3.821578in}{0.633527in}}%
\pgfpathlineto{\pgfqpoint{3.821578in}{0.625278in}}%
\pgfpathlineto{\pgfqpoint{3.822577in}{0.625278in}}%
\pgfpathlineto{\pgfqpoint{3.823584in}{0.633527in}}%
\pgfpathlineto{\pgfqpoint{3.823584in}{0.625278in}}%
\pgfpathlineto{\pgfqpoint{3.824925in}{0.633527in}}%
\pgfpathlineto{\pgfqpoint{3.824925in}{0.625278in}}%
\pgfpathlineto{\pgfqpoint{3.825769in}{0.625278in}}%
\pgfpathlineto{\pgfqpoint{3.826938in}{0.633527in}}%
\pgfpathlineto{\pgfqpoint{3.827109in}{0.625278in}}%
\pgfpathlineto{\pgfqpoint{3.827775in}{0.625278in}}%
\pgfpathlineto{\pgfqpoint{3.829115in}{0.633527in}}%
\pgfpathlineto{\pgfqpoint{3.829456in}{0.625278in}}%
\pgfpathlineto{\pgfqpoint{3.830122in}{0.625278in}}%
\pgfpathlineto{\pgfqpoint{3.831300in}{0.633527in}}%
\pgfpathlineto{\pgfqpoint{3.831633in}{0.625278in}}%
\pgfpathlineto{\pgfqpoint{3.832306in}{0.625278in}}%
\pgfpathlineto{\pgfqpoint{3.833647in}{0.633527in}}%
\pgfpathlineto{\pgfqpoint{3.833810in}{0.625278in}}%
\pgfpathlineto{\pgfqpoint{3.834654in}{0.625278in}}%
\pgfpathlineto{\pgfqpoint{3.835653in}{0.633527in}}%
\pgfpathlineto{\pgfqpoint{3.835653in}{0.625278in}}%
\pgfpathlineto{\pgfqpoint{3.836993in}{0.633527in}}%
\pgfpathlineto{\pgfqpoint{3.836993in}{0.625278in}}%
\pgfpathlineto{\pgfqpoint{3.838000in}{0.625278in}}%
\pgfpathlineto{\pgfqpoint{3.839178in}{0.633527in}}%
\pgfpathlineto{\pgfqpoint{3.839340in}{0.625278in}}%
\pgfpathlineto{\pgfqpoint{3.840185in}{0.625278in}}%
\pgfpathlineto{\pgfqpoint{3.841354in}{0.633527in}}%
\pgfpathlineto{\pgfqpoint{3.841688in}{0.625278in}}%
\pgfpathlineto{\pgfqpoint{3.842361in}{0.625278in}}%
\pgfpathlineto{\pgfqpoint{3.843531in}{0.633527in}}%
\pgfpathlineto{\pgfqpoint{3.843702in}{0.625278in}}%
\pgfpathlineto{\pgfqpoint{3.844538in}{0.625278in}}%
\pgfpathlineto{\pgfqpoint{3.845715in}{0.633527in}}%
\pgfpathlineto{\pgfqpoint{3.845715in}{0.625278in}}%
\pgfpathlineto{\pgfqpoint{3.846722in}{0.625278in}}%
\pgfpathlineto{\pgfqpoint{3.848063in}{0.633527in}}%
\pgfpathlineto{\pgfqpoint{3.848063in}{0.625278in}}%
\pgfpathlineto{\pgfqpoint{3.849084in}{0.625278in}}%
\pgfpathlineto{\pgfqpoint{3.850410in}{0.633527in}}%
\pgfpathlineto{\pgfqpoint{3.850573in}{0.625278in}}%
\pgfpathlineto{\pgfqpoint{3.851417in}{0.625278in}}%
\pgfpathlineto{\pgfqpoint{3.852587in}{0.633527in}}%
\pgfpathlineto{\pgfqpoint{3.852587in}{0.625278in}}%
\pgfpathlineto{\pgfqpoint{3.853594in}{0.625278in}}%
\pgfpathlineto{\pgfqpoint{3.854763in}{0.633527in}}%
\pgfpathlineto{\pgfqpoint{3.854786in}{0.625278in}}%
\pgfpathlineto{\pgfqpoint{3.855793in}{0.625278in}}%
\pgfpathlineto{\pgfqpoint{3.857451in}{0.658274in}}%
\pgfpathlineto{\pgfqpoint{3.858118in}{0.633527in}}%
\pgfpathlineto{\pgfqpoint{3.858473in}{0.650025in}}%
\pgfpathlineto{\pgfqpoint{3.859643in}{0.674773in}}%
\pgfpathlineto{\pgfqpoint{3.859813in}{0.666524in}}%
\pgfpathlineto{\pgfqpoint{3.860487in}{0.650025in}}%
\pgfpathlineto{\pgfqpoint{3.860650in}{0.674773in}}%
\pgfpathlineto{\pgfqpoint{3.860968in}{0.658274in}}%
\pgfpathlineto{\pgfqpoint{3.861694in}{0.683022in}}%
\pgfpathlineto{\pgfqpoint{3.861990in}{0.674773in}}%
\pgfpathlineto{\pgfqpoint{3.863337in}{0.724268in}}%
\pgfpathlineto{\pgfqpoint{3.863648in}{0.707769in}}%
\pgfpathlineto{\pgfqpoint{3.864322in}{0.683022in}}%
\pgfpathlineto{\pgfqpoint{3.864337in}{0.724268in}}%
\pgfpathlineto{\pgfqpoint{3.864678in}{0.699520in}}%
\pgfpathlineto{\pgfqpoint{3.865492in}{0.707769in}}%
\pgfpathlineto{\pgfqpoint{3.865514in}{0.699520in}}%
\pgfpathlineto{\pgfqpoint{3.866521in}{0.674773in}}%
\pgfpathlineto{\pgfqpoint{3.866669in}{0.691271in}}%
\pgfpathlineto{\pgfqpoint{3.867173in}{0.732517in}}%
\pgfpathlineto{\pgfqpoint{3.867691in}{0.699520in}}%
\pgfpathlineto{\pgfqpoint{3.868846in}{0.707769in}}%
\pgfpathlineto{\pgfqpoint{3.869179in}{0.732517in}}%
\pgfpathlineto{\pgfqpoint{3.869868in}{0.699520in}}%
\pgfpathlineto{\pgfqpoint{3.870520in}{0.732517in}}%
\pgfpathlineto{\pgfqpoint{3.871045in}{0.716019in}}%
\pgfpathlineto{\pgfqpoint{3.871527in}{0.707769in}}%
\pgfpathlineto{\pgfqpoint{3.871697in}{0.732517in}}%
\pgfpathlineto{\pgfqpoint{3.871719in}{0.724268in}}%
\pgfpathlineto{\pgfqpoint{3.872867in}{0.757264in}}%
\pgfpathlineto{\pgfqpoint{3.873059in}{0.740766in}}%
\pgfpathlineto{\pgfqpoint{3.873874in}{0.732517in}}%
\pgfpathlineto{\pgfqpoint{3.873222in}{0.749015in}}%
\pgfpathlineto{\pgfqpoint{3.873896in}{0.749015in}}%
\pgfpathlineto{\pgfqpoint{3.874881in}{0.806759in}}%
\pgfpathlineto{\pgfqpoint{3.875066in}{0.790261in}}%
\pgfpathlineto{\pgfqpoint{3.875406in}{0.773763in}}%
\pgfpathlineto{\pgfqpoint{3.875717in}{0.806759in}}%
\pgfpathlineto{\pgfqpoint{3.875740in}{0.798510in}}%
\pgfpathlineto{\pgfqpoint{3.875888in}{0.806759in}}%
\pgfpathlineto{\pgfqpoint{3.876221in}{0.782012in}}%
\pgfpathlineto{\pgfqpoint{3.876406in}{0.790261in}}%
\pgfpathlineto{\pgfqpoint{3.876561in}{0.782012in}}%
\pgfpathlineto{\pgfqpoint{3.876724in}{0.806759in}}%
\pgfpathlineto{\pgfqpoint{3.877080in}{0.798510in}}%
\pgfpathlineto{\pgfqpoint{3.877250in}{0.823258in}}%
\pgfpathlineto{\pgfqpoint{3.878235in}{0.806759in}}%
\pgfpathlineto{\pgfqpoint{3.878590in}{0.798510in}}%
\pgfpathlineto{\pgfqpoint{3.878901in}{0.831507in}}%
\pgfpathlineto{\pgfqpoint{3.879256in}{0.823258in}}%
\pgfpathlineto{\pgfqpoint{3.879575in}{0.806759in}}%
\pgfpathlineto{\pgfqpoint{3.879745in}{0.831507in}}%
\pgfpathlineto{\pgfqpoint{3.880101in}{0.823258in}}%
\pgfpathlineto{\pgfqpoint{3.882092in}{0.881002in}}%
\pgfpathlineto{\pgfqpoint{3.883114in}{0.872753in}}%
\pgfpathlineto{\pgfqpoint{3.883595in}{0.881002in}}%
\pgfpathlineto{\pgfqpoint{3.883788in}{0.864503in}}%
\pgfpathlineto{\pgfqpoint{3.883936in}{0.856254in}}%
\pgfpathlineto{\pgfqpoint{3.884439in}{0.889251in}}%
\pgfpathlineto{\pgfqpoint{3.884787in}{0.872753in}}%
\pgfpathlineto{\pgfqpoint{3.885461in}{0.897500in}}%
\pgfpathlineto{\pgfqpoint{3.886298in}{0.872753in}}%
\pgfpathlineto{\pgfqpoint{3.886787in}{0.881002in}}%
\pgfpathlineto{\pgfqpoint{3.886801in}{0.922247in}}%
\pgfpathlineto{\pgfqpoint{3.887808in}{0.922247in}}%
\pgfpathlineto{\pgfqpoint{3.889467in}{0.881002in}}%
\pgfpathlineto{\pgfqpoint{3.889467in}{0.889251in}}%
\pgfpathlineto{\pgfqpoint{3.890304in}{0.930497in}}%
\pgfpathlineto{\pgfqpoint{3.889800in}{0.881002in}}%
\pgfpathlineto{\pgfqpoint{3.890489in}{0.913998in}}%
\pgfpathlineto{\pgfqpoint{3.890659in}{0.922247in}}%
\pgfpathlineto{\pgfqpoint{3.891666in}{0.897500in}}%
\pgfpathlineto{\pgfqpoint{3.893154in}{0.930497in}}%
\pgfpathlineto{\pgfqpoint{3.893169in}{0.922247in}}%
\pgfpathlineto{\pgfqpoint{3.894346in}{0.897500in}}%
\pgfpathlineto{\pgfqpoint{3.894494in}{0.905749in}}%
\pgfpathlineto{\pgfqpoint{3.894665in}{0.930497in}}%
\pgfpathlineto{\pgfqpoint{3.895353in}{0.897500in}}%
\pgfpathlineto{\pgfqpoint{3.895516in}{0.897500in}}%
\pgfpathlineto{\pgfqpoint{3.896864in}{0.946995in}}%
\pgfpathlineto{\pgfqpoint{3.897175in}{0.930497in}}%
\pgfpathlineto{\pgfqpoint{3.897530in}{0.922247in}}%
\pgfpathlineto{\pgfqpoint{3.897848in}{0.955244in}}%
\pgfpathlineto{\pgfqpoint{3.898204in}{0.946995in}}%
\pgfpathlineto{\pgfqpoint{3.898515in}{0.930497in}}%
\pgfpathlineto{\pgfqpoint{3.898855in}{0.963493in}}%
\pgfpathlineto{\pgfqpoint{3.899189in}{0.979992in}}%
\pgfpathlineto{\pgfqpoint{3.899855in}{0.930497in}}%
\pgfpathlineto{\pgfqpoint{3.899877in}{0.946995in}}%
\pgfpathlineto{\pgfqpoint{3.900084in}{0.930497in}}%
\pgfpathlineto{\pgfqpoint{3.900210in}{0.971742in}}%
\pgfpathlineto{\pgfqpoint{3.900699in}{0.979992in}}%
\pgfpathlineto{\pgfqpoint{3.900884in}{0.963493in}}%
\pgfpathlineto{\pgfqpoint{3.902039in}{0.930497in}}%
\pgfpathlineto{\pgfqpoint{3.902054in}{0.946995in}}%
\pgfpathlineto{\pgfqpoint{3.903713in}{0.979992in}}%
\pgfpathlineto{\pgfqpoint{3.904905in}{0.946995in}}%
\pgfpathlineto{\pgfqpoint{3.906082in}{0.971742in}}%
\pgfpathlineto{\pgfqpoint{3.906245in}{0.963493in}}%
\pgfpathlineto{\pgfqpoint{3.907740in}{0.930497in}}%
\pgfpathlineto{\pgfqpoint{3.906415in}{0.971742in}}%
\pgfpathlineto{\pgfqpoint{3.907740in}{0.938746in}}%
\pgfpathlineto{\pgfqpoint{3.908088in}{0.971742in}}%
\pgfpathlineto{\pgfqpoint{3.908074in}{0.930497in}}%
\pgfpathlineto{\pgfqpoint{3.908762in}{0.946995in}}%
\pgfpathlineto{\pgfqpoint{3.909917in}{0.955244in}}%
\pgfpathlineto{\pgfqpoint{3.910754in}{0.930497in}}%
\pgfpathlineto{\pgfqpoint{3.910939in}{0.938746in}}%
\pgfpathlineto{\pgfqpoint{3.911257in}{0.955244in}}%
\pgfpathlineto{\pgfqpoint{3.912116in}{0.922247in}}%
\pgfpathlineto{\pgfqpoint{3.913790in}{0.971742in}}%
\pgfpathlineto{\pgfqpoint{3.913960in}{0.963493in}}%
\pgfpathlineto{\pgfqpoint{3.914108in}{0.955244in}}%
\pgfpathlineto{\pgfqpoint{3.914463in}{0.996490in}}%
\pgfpathlineto{\pgfqpoint{3.915618in}{1.004739in}}%
\pgfpathlineto{\pgfqpoint{3.915804in}{0.996490in}}%
\pgfpathlineto{\pgfqpoint{3.916470in}{1.045985in}}%
\pgfpathlineto{\pgfqpoint{3.916640in}{1.037736in}}%
\pgfpathlineto{\pgfqpoint{3.916811in}{1.045985in}}%
\pgfpathlineto{\pgfqpoint{3.917818in}{1.021237in}}%
\pgfpathlineto{\pgfqpoint{3.918987in}{1.045985in}}%
\pgfpathlineto{\pgfqpoint{3.919158in}{1.037736in}}%
\pgfpathlineto{\pgfqpoint{3.919809in}{1.004739in}}%
\pgfpathlineto{\pgfqpoint{3.920313in}{1.029487in}}%
\pgfpathlineto{\pgfqpoint{3.920328in}{1.045985in}}%
\pgfpathlineto{\pgfqpoint{3.921001in}{0.996490in}}%
\pgfpathlineto{\pgfqpoint{3.921335in}{1.021237in}}%
\pgfpathlineto{\pgfqpoint{3.921653in}{1.029487in}}%
\pgfpathlineto{\pgfqpoint{3.921986in}{1.004739in}}%
\pgfpathlineto{\pgfqpoint{3.923178in}{0.971742in}}%
\pgfpathlineto{\pgfqpoint{3.924689in}{1.045985in}}%
\pgfpathlineto{\pgfqpoint{3.925170in}{1.029487in}}%
\pgfpathlineto{\pgfqpoint{3.926362in}{0.996490in}}%
\pgfpathlineto{\pgfqpoint{3.927539in}{1.021237in}}%
\pgfpathlineto{\pgfqpoint{3.927702in}{1.012988in}}%
\pgfpathlineto{\pgfqpoint{3.927872in}{0.996490in}}%
\pgfpathlineto{\pgfqpoint{3.928191in}{1.029487in}}%
\pgfpathlineto{\pgfqpoint{3.928376in}{1.021237in}}%
\pgfpathlineto{\pgfqpoint{3.928694in}{1.029487in}}%
\pgfpathlineto{\pgfqpoint{3.929042in}{0.996490in}}%
\pgfpathlineto{\pgfqpoint{3.929694in}{0.979992in}}%
\pgfpathlineto{\pgfqpoint{3.929694in}{1.012988in}}%
\pgfpathlineto{\pgfqpoint{3.930034in}{1.004739in}}%
\pgfpathlineto{\pgfqpoint{3.930723in}{1.021237in}}%
\pgfpathlineto{\pgfqpoint{3.931878in}{1.029487in}}%
\pgfpathlineto{\pgfqpoint{3.932715in}{0.979992in}}%
\pgfpathlineto{\pgfqpoint{3.932900in}{1.012988in}}%
\pgfpathlineto{\pgfqpoint{3.933048in}{1.004739in}}%
\pgfpathlineto{\pgfqpoint{3.933070in}{1.045985in}}%
\pgfpathlineto{\pgfqpoint{3.933574in}{1.021237in}}%
\pgfpathlineto{\pgfqpoint{3.934055in}{1.029487in}}%
\pgfpathlineto{\pgfqpoint{3.934388in}{1.004739in}}%
\pgfpathlineto{\pgfqpoint{3.934744in}{0.996490in}}%
\pgfpathlineto{\pgfqpoint{3.935232in}{1.054234in}}%
\pgfpathlineto{\pgfqpoint{3.935284in}{1.054234in}}%
\pgfpathlineto{\pgfqpoint{3.936254in}{1.095480in}}%
\pgfpathlineto{\pgfqpoint{3.936587in}{1.070732in}}%
\pgfpathlineto{\pgfqpoint{3.937239in}{1.078981in}}%
\pgfpathlineto{\pgfqpoint{3.937424in}{1.062483in}}%
\pgfpathlineto{\pgfqpoint{3.937764in}{1.045985in}}%
\pgfpathlineto{\pgfqpoint{3.938075in}{1.078981in}}%
\pgfpathlineto{\pgfqpoint{3.938098in}{1.070732in}}%
\pgfpathlineto{\pgfqpoint{3.938246in}{1.078981in}}%
\pgfpathlineto{\pgfqpoint{3.938601in}{1.045985in}}%
\pgfpathlineto{\pgfqpoint{3.938749in}{1.054234in}}%
\pgfpathlineto{\pgfqpoint{3.938934in}{1.045985in}}%
\pgfpathlineto{\pgfqpoint{3.939253in}{1.078981in}}%
\pgfpathlineto{\pgfqpoint{3.939771in}{1.070732in}}%
\pgfpathlineto{\pgfqpoint{3.940926in}{1.078981in}}%
\pgfpathlineto{\pgfqpoint{3.940948in}{1.070732in}}%
\pgfpathlineto{\pgfqpoint{3.941489in}{1.103729in}}%
\pgfpathlineto{\pgfqpoint{3.941955in}{1.095480in}}%
\pgfpathlineto{\pgfqpoint{3.942622in}{1.144975in}}%
\pgfpathlineto{\pgfqpoint{3.943614in}{1.128476in}}%
\pgfpathlineto{\pgfqpoint{3.943962in}{1.120227in}}%
\pgfpathlineto{\pgfqpoint{3.944132in}{1.144975in}}%
\pgfpathlineto{\pgfqpoint{3.944636in}{1.136726in}}%
\pgfpathlineto{\pgfqpoint{3.944806in}{1.144975in}}%
\pgfpathlineto{\pgfqpoint{3.945791in}{1.128476in}}%
\pgfpathlineto{\pgfqpoint{3.946309in}{1.169722in}}%
\pgfpathlineto{\pgfqpoint{3.946812in}{1.144975in}}%
\pgfpathlineto{\pgfqpoint{3.947301in}{1.153224in}}%
\pgfpathlineto{\pgfqpoint{3.947486in}{1.136726in}}%
\pgfpathlineto{\pgfqpoint{3.948301in}{1.128476in}}%
\pgfpathlineto{\pgfqpoint{3.947649in}{1.144975in}}%
\pgfpathlineto{\pgfqpoint{3.948323in}{1.144975in}}%
\pgfpathlineto{\pgfqpoint{3.948974in}{1.177971in}}%
\pgfpathlineto{\pgfqpoint{3.949493in}{1.161473in}}%
\pgfpathlineto{\pgfqpoint{3.949663in}{1.169722in}}%
\pgfpathlineto{\pgfqpoint{3.951173in}{1.120227in}}%
\pgfpathlineto{\pgfqpoint{3.952514in}{1.194470in}}%
\pgfpathlineto{\pgfqpoint{3.953165in}{1.177971in}}%
\pgfpathlineto{\pgfqpoint{3.954357in}{1.144975in}}%
\pgfpathlineto{\pgfqpoint{3.955349in}{1.177971in}}%
\pgfpathlineto{\pgfqpoint{3.954839in}{1.128476in}}%
\pgfpathlineto{\pgfqpoint{3.955535in}{1.161473in}}%
\pgfpathlineto{\pgfqpoint{3.956349in}{1.153224in}}%
\pgfpathlineto{\pgfqpoint{3.955697in}{1.169722in}}%
\pgfpathlineto{\pgfqpoint{3.956371in}{1.169722in}}%
\pgfpathlineto{\pgfqpoint{3.958718in}{1.268712in}}%
\pgfpathlineto{\pgfqpoint{3.958881in}{1.260463in}}%
\pgfpathlineto{\pgfqpoint{3.959555in}{1.243965in}}%
\pgfpathlineto{\pgfqpoint{3.959725in}{1.268712in}}%
\pgfpathlineto{\pgfqpoint{3.960036in}{1.252214in}}%
\pgfpathlineto{\pgfqpoint{3.960058in}{1.268712in}}%
\pgfpathlineto{\pgfqpoint{3.961065in}{1.268712in}}%
\pgfpathlineto{\pgfqpoint{3.961717in}{1.301709in}}%
\pgfpathlineto{\pgfqpoint{3.962220in}{1.276961in}}%
\pgfpathlineto{\pgfqpoint{3.963242in}{1.268712in}}%
\pgfpathlineto{\pgfqpoint{3.963724in}{1.276961in}}%
\pgfpathlineto{\pgfqpoint{3.963916in}{1.260463in}}%
\pgfpathlineto{\pgfqpoint{3.964901in}{1.252214in}}%
\pgfpathlineto{\pgfqpoint{3.964568in}{1.276961in}}%
\pgfpathlineto{\pgfqpoint{3.964901in}{1.260463in}}%
\pgfpathlineto{\pgfqpoint{3.965234in}{1.252214in}}%
\pgfpathlineto{\pgfqpoint{3.966071in}{1.301709in}}%
\pgfpathlineto{\pgfqpoint{3.967100in}{1.293460in}}%
\pgfpathlineto{\pgfqpoint{3.967137in}{1.301709in}}%
\pgfpathlineto{\pgfqpoint{3.968418in}{1.326456in}}%
\pgfpathlineto{\pgfqpoint{3.968440in}{1.318207in}}%
\pgfpathlineto{\pgfqpoint{3.970269in}{1.375951in}}%
\pgfpathlineto{\pgfqpoint{3.970284in}{1.367702in}}%
\pgfpathlineto{\pgfqpoint{3.971602in}{1.351204in}}%
\pgfpathlineto{\pgfqpoint{3.971624in}{1.392449in}}%
\pgfpathlineto{\pgfqpoint{3.972631in}{1.367702in}}%
\pgfpathlineto{\pgfqpoint{3.973949in}{1.450194in}}%
\pgfpathlineto{\pgfqpoint{3.974474in}{1.408948in}}%
\pgfpathlineto{\pgfqpoint{3.974637in}{1.417197in}}%
\pgfpathlineto{\pgfqpoint{3.975629in}{1.400699in}}%
\pgfpathlineto{\pgfqpoint{3.975644in}{1.417197in}}%
\pgfpathlineto{\pgfqpoint{3.976636in}{1.375951in}}%
\pgfpathlineto{\pgfqpoint{3.976651in}{1.392449in}}%
\pgfpathlineto{\pgfqpoint{3.977991in}{1.367702in}}%
\pgfpathlineto{\pgfqpoint{3.978984in}{1.425446in}}%
\pgfpathlineto{\pgfqpoint{3.979332in}{1.392449in}}%
\pgfpathlineto{\pgfqpoint{3.980005in}{1.367702in}}%
\pgfpathlineto{\pgfqpoint{3.980176in}{1.392449in}}%
\pgfpathlineto{\pgfqpoint{3.980324in}{1.400699in}}%
\pgfpathlineto{\pgfqpoint{3.980679in}{1.359453in}}%
\pgfpathlineto{\pgfqpoint{3.982019in}{1.318207in}}%
\pgfpathlineto{\pgfqpoint{3.983019in}{1.293460in}}%
\pgfpathlineto{\pgfqpoint{3.983174in}{1.309958in}}%
\pgfpathlineto{\pgfqpoint{3.983189in}{1.318207in}}%
\pgfpathlineto{\pgfqpoint{3.984026in}{1.268712in}}%
\pgfpathlineto{\pgfqpoint{3.984196in}{1.268712in}}%
\pgfpathlineto{\pgfqpoint{3.984515in}{1.252214in}}%
\pgfpathlineto{\pgfqpoint{3.984848in}{1.285210in}}%
\pgfpathlineto{\pgfqpoint{3.985181in}{1.301709in}}%
\pgfpathlineto{\pgfqpoint{3.985855in}{1.252214in}}%
\pgfpathlineto{\pgfqpoint{3.985869in}{1.268712in}}%
\pgfpathlineto{\pgfqpoint{3.986366in}{1.276961in}}%
\pgfpathlineto{\pgfqpoint{3.986543in}{1.260463in}}%
\pgfpathlineto{\pgfqpoint{3.987550in}{1.243965in}}%
\pgfpathlineto{\pgfqpoint{3.987210in}{1.268712in}}%
\pgfpathlineto{\pgfqpoint{3.987706in}{1.252214in}}%
\pgfpathlineto{\pgfqpoint{3.988054in}{1.243965in}}%
\pgfpathlineto{\pgfqpoint{3.988720in}{1.293460in}}%
\pgfpathlineto{\pgfqpoint{3.989897in}{1.243965in}}%
\pgfpathlineto{\pgfqpoint{3.990060in}{1.268712in}}%
\pgfpathlineto{\pgfqpoint{3.990549in}{1.276961in}}%
\pgfpathlineto{\pgfqpoint{3.990734in}{1.260463in}}%
\pgfpathlineto{\pgfqpoint{3.992555in}{1.202719in}}%
\pgfpathlineto{\pgfqpoint{3.992578in}{1.219217in}}%
\pgfpathlineto{\pgfqpoint{3.993237in}{1.252214in}}%
\pgfpathlineto{\pgfqpoint{3.993585in}{1.219217in}}%
\pgfpathlineto{\pgfqpoint{3.996917in}{1.054234in}}%
\pgfpathlineto{\pgfqpoint{3.994073in}{1.227466in}}%
\pgfpathlineto{\pgfqpoint{3.997257in}{1.087231in}}%
\pgfpathlineto{\pgfqpoint{3.997598in}{1.103729in}}%
\pgfpathlineto{\pgfqpoint{3.998257in}{1.054234in}}%
\pgfpathlineto{\pgfqpoint{3.998279in}{1.070732in}}%
\pgfpathlineto{\pgfqpoint{3.998938in}{1.078981in}}%
\pgfpathlineto{\pgfqpoint{3.999116in}{1.062483in}}%
\pgfpathlineto{\pgfqpoint{4.000456in}{1.021237in}}%
\pgfpathlineto{\pgfqpoint{4.000774in}{1.029487in}}%
\pgfpathlineto{\pgfqpoint{4.001107in}{1.004739in}}%
\pgfpathlineto{\pgfqpoint{4.002129in}{0.946995in}}%
\pgfpathlineto{\pgfqpoint{4.002299in}{0.971742in}}%
\pgfpathlineto{\pgfqpoint{4.002447in}{0.979992in}}%
\pgfpathlineto{\pgfqpoint{4.002803in}{0.946995in}}%
\pgfpathlineto{\pgfqpoint{4.003121in}{0.930497in}}%
\pgfpathlineto{\pgfqpoint{4.003454in}{0.963493in}}%
\pgfpathlineto{\pgfqpoint{4.004476in}{1.021237in}}%
\pgfpathlineto{\pgfqpoint{4.004647in}{1.012988in}}%
\pgfpathlineto{\pgfqpoint{4.005320in}{0.996490in}}%
\pgfpathlineto{\pgfqpoint{4.004817in}{1.021237in}}%
\pgfpathlineto{\pgfqpoint{4.005468in}{1.012988in}}%
\pgfpathlineto{\pgfqpoint{4.005824in}{1.070732in}}%
\pgfpathlineto{\pgfqpoint{4.006490in}{1.045985in}}%
\pgfpathlineto{\pgfqpoint{4.007319in}{1.054234in}}%
\pgfpathlineto{\pgfqpoint{4.007327in}{1.045985in}}%
\pgfpathlineto{\pgfqpoint{4.008001in}{1.021237in}}%
\pgfpathlineto{\pgfqpoint{4.008164in}{1.045985in}}%
\pgfpathlineto{\pgfqpoint{4.008326in}{1.054234in}}%
\pgfpathlineto{\pgfqpoint{4.008667in}{1.021237in}}%
\pgfpathlineto{\pgfqpoint{4.008815in}{1.029487in}}%
\pgfpathlineto{\pgfqpoint{4.009674in}{1.021237in}}%
\pgfpathlineto{\pgfqpoint{4.009844in}{1.045985in}}%
\pgfpathlineto{\pgfqpoint{4.010177in}{1.054234in}}%
\pgfpathlineto{\pgfqpoint{4.010496in}{1.029487in}}%
\pgfpathlineto{\pgfqpoint{4.010851in}{1.021237in}}%
\pgfpathlineto{\pgfqpoint{4.011347in}{1.078981in}}%
\pgfpathlineto{\pgfqpoint{4.011518in}{1.070732in}}%
\pgfpathlineto{\pgfqpoint{4.012687in}{1.103729in}}%
\pgfpathlineto{\pgfqpoint{4.012858in}{1.087231in}}%
\pgfpathlineto{\pgfqpoint{4.015020in}{1.029487in}}%
\pgfpathlineto{\pgfqpoint{4.015042in}{1.045985in}}%
\pgfpathlineto{\pgfqpoint{4.015708in}{0.996490in}}%
\pgfpathlineto{\pgfqpoint{4.016049in}{0.996490in}}%
\pgfpathlineto{\pgfqpoint{4.016715in}{1.045985in}}%
\pgfpathlineto{\pgfqpoint{4.017389in}{1.021237in}}%
\pgfpathlineto{\pgfqpoint{4.018226in}{0.996490in}}%
\pgfpathlineto{\pgfqpoint{4.018396in}{1.021237in}}%
\pgfpathlineto{\pgfqpoint{4.019899in}{1.103729in}}%
\pgfpathlineto{\pgfqpoint{4.020240in}{1.070732in}}%
\pgfpathlineto{\pgfqpoint{4.020728in}{1.078981in}}%
\pgfpathlineto{\pgfqpoint{4.022083in}{1.021237in}}%
\pgfpathlineto{\pgfqpoint{4.022231in}{1.029487in}}%
\pgfpathlineto{\pgfqpoint{4.022587in}{0.996490in}}%
\pgfpathlineto{\pgfqpoint{4.024764in}{0.897500in}}%
\pgfpathlineto{\pgfqpoint{4.025097in}{0.922247in}}%
\pgfpathlineto{\pgfqpoint{4.028110in}{1.029487in}}%
\pgfpathlineto{\pgfqpoint{4.029117in}{0.988241in}}%
\pgfpathlineto{\pgfqpoint{4.029451in}{0.955244in}}%
\pgfpathlineto{\pgfqpoint{4.029939in}{1.004739in}}%
\pgfpathlineto{\pgfqpoint{4.030295in}{0.971742in}}%
\pgfpathlineto{\pgfqpoint{4.031805in}{0.881002in}}%
\pgfpathlineto{\pgfqpoint{4.032309in}{0.922247in}}%
\pgfpathlineto{\pgfqpoint{4.033649in}{0.971742in}}%
\pgfpathlineto{\pgfqpoint{4.033975in}{0.955244in}}%
\pgfpathlineto{\pgfqpoint{4.034989in}{0.922247in}}%
\pgfpathlineto{\pgfqpoint{4.035648in}{0.905749in}}%
\pgfpathlineto{\pgfqpoint{4.035818in}{0.930497in}}%
\pgfpathlineto{\pgfqpoint{4.035826in}{0.922247in}}%
\pgfpathlineto{\pgfqpoint{4.036322in}{0.930497in}}%
\pgfpathlineto{\pgfqpoint{4.036499in}{0.913998in}}%
\pgfpathlineto{\pgfqpoint{4.040690in}{0.749015in}}%
\pgfpathlineto{\pgfqpoint{4.041009in}{0.757264in}}%
\pgfpathlineto{\pgfqpoint{4.041357in}{0.732517in}}%
\pgfpathlineto{\pgfqpoint{4.041357in}{0.724268in}}%
\pgfpathlineto{\pgfqpoint{4.041860in}{0.749015in}}%
\pgfpathlineto{\pgfqpoint{4.042364in}{0.740766in}}%
\pgfpathlineto{\pgfqpoint{4.044207in}{0.683022in}}%
\pgfpathlineto{\pgfqpoint{4.042534in}{0.749015in}}%
\pgfpathlineto{\pgfqpoint{4.044355in}{0.691271in}}%
\pgfpathlineto{\pgfqpoint{4.045362in}{0.707769in}}%
\pgfpathlineto{\pgfqpoint{4.044711in}{0.674773in}}%
\pgfpathlineto{\pgfqpoint{4.045384in}{0.699520in}}%
\pgfpathlineto{\pgfqpoint{4.045880in}{0.674773in}}%
\pgfpathlineto{\pgfqpoint{4.046702in}{0.683022in}}%
\pgfpathlineto{\pgfqpoint{4.047376in}{0.732517in}}%
\pgfpathlineto{\pgfqpoint{4.047732in}{0.724268in}}%
\pgfpathlineto{\pgfqpoint{4.048043in}{0.732517in}}%
\pgfpathlineto{\pgfqpoint{4.048383in}{0.707769in}}%
\pgfpathlineto{\pgfqpoint{4.048568in}{0.716019in}}%
\pgfpathlineto{\pgfqpoint{4.048568in}{0.707769in}}%
\pgfpathlineto{\pgfqpoint{4.048887in}{0.757264in}}%
\pgfpathlineto{\pgfqpoint{4.049575in}{0.749015in}}%
\pgfpathlineto{\pgfqpoint{4.051752in}{0.683022in}}%
\pgfpathlineto{\pgfqpoint{4.054099in}{0.757264in}}%
\pgfpathlineto{\pgfqpoint{4.056950in}{0.625278in}}%
\pgfpathlineto{\pgfqpoint{4.057113in}{0.650025in}}%
\pgfpathlineto{\pgfqpoint{4.057276in}{0.658274in}}%
\pgfpathlineto{\pgfqpoint{4.057616in}{0.625278in}}%
\pgfpathlineto{\pgfqpoint{4.057772in}{0.633527in}}%
\pgfpathlineto{\pgfqpoint{4.058793in}{0.625278in}}%
\pgfpathlineto{\pgfqpoint{4.058793in}{0.633527in}}%
\pgfpathlineto{\pgfqpoint{4.058956in}{0.625278in}}%
\pgfpathlineto{\pgfqpoint{4.059127in}{0.658274in}}%
\pgfpathlineto{\pgfqpoint{4.059297in}{0.650025in}}%
\pgfpathlineto{\pgfqpoint{4.062473in}{0.806759in}}%
\pgfpathlineto{\pgfqpoint{4.062644in}{0.790261in}}%
\pgfpathlineto{\pgfqpoint{4.063991in}{0.749015in}}%
\pgfpathlineto{\pgfqpoint{4.064650in}{0.732517in}}%
\pgfpathlineto{\pgfqpoint{4.064820in}{0.757264in}}%
\pgfpathlineto{\pgfqpoint{4.064828in}{0.749015in}}%
\pgfpathlineto{\pgfqpoint{4.065650in}{0.757264in}}%
\pgfpathlineto{\pgfqpoint{4.065665in}{0.749015in}}%
\pgfpathlineto{\pgfqpoint{4.068848in}{0.625278in}}%
\pgfpathlineto{\pgfqpoint{4.070018in}{0.633527in}}%
\pgfpathlineto{\pgfqpoint{4.070359in}{0.625278in}}%
\pgfpathlineto{\pgfqpoint{4.071033in}{0.625278in}}%
\pgfpathlineto{\pgfqpoint{4.072195in}{0.633527in}}%
\pgfpathlineto{\pgfqpoint{4.072239in}{0.625278in}}%
\pgfpathlineto{\pgfqpoint{4.073187in}{0.625278in}}%
\pgfpathlineto{\pgfqpoint{4.074364in}{0.633527in}}%
\pgfpathlineto{\pgfqpoint{4.074542in}{0.625278in}}%
\pgfpathlineto{\pgfqpoint{4.075371in}{0.625278in}}%
\pgfpathlineto{\pgfqpoint{4.076541in}{0.633527in}}%
\pgfpathlineto{\pgfqpoint{4.076874in}{0.625278in}}%
\pgfpathlineto{\pgfqpoint{4.077556in}{0.625278in}}%
\pgfpathlineto{\pgfqpoint{4.078903in}{0.633527in}}%
\pgfpathlineto{\pgfqpoint{4.079066in}{0.625278in}}%
\pgfpathlineto{\pgfqpoint{4.079732in}{0.625278in}}%
\pgfpathlineto{\pgfqpoint{4.081065in}{0.633527in}}%
\pgfpathlineto{\pgfqpoint{4.081065in}{0.625278in}}%
\pgfpathlineto{\pgfqpoint{4.082072in}{0.625278in}}%
\pgfpathlineto{\pgfqpoint{4.083257in}{0.633527in}}%
\pgfpathlineto{\pgfqpoint{4.083412in}{0.625278in}}%
\pgfpathlineto{\pgfqpoint{4.084086in}{0.625278in}}%
\pgfpathlineto{\pgfqpoint{4.085426in}{0.633527in}}%
\pgfpathlineto{\pgfqpoint{4.085759in}{0.625278in}}%
\pgfpathlineto{\pgfqpoint{4.086107in}{0.625278in}}%
\pgfpathlineto{\pgfqpoint{4.087270in}{0.633527in}}%
\pgfpathlineto{\pgfqpoint{4.087270in}{0.625278in}}%
\pgfpathlineto{\pgfqpoint{4.088277in}{0.625278in}}%
\pgfpathlineto{\pgfqpoint{4.089454in}{0.633527in}}%
\pgfpathlineto{\pgfqpoint{4.089617in}{0.625278in}}%
\pgfpathlineto{\pgfqpoint{4.090454in}{0.625278in}}%
\pgfpathlineto{\pgfqpoint{4.091631in}{0.633527in}}%
\pgfpathlineto{\pgfqpoint{4.091801in}{0.625278in}}%
\pgfpathlineto{\pgfqpoint{4.092631in}{0.625278in}}%
\pgfpathlineto{\pgfqpoint{4.093978in}{0.633527in}}%
\pgfpathlineto{\pgfqpoint{4.093986in}{0.625278in}}%
\pgfpathlineto{\pgfqpoint{4.094985in}{0.625278in}}%
\pgfpathlineto{\pgfqpoint{4.096155in}{0.633527in}}%
\pgfpathlineto{\pgfqpoint{4.096155in}{0.625278in}}%
\pgfpathlineto{\pgfqpoint{4.097162in}{0.625278in}}%
\pgfpathlineto{\pgfqpoint{4.098502in}{0.633527in}}%
\pgfpathlineto{\pgfqpoint{4.098672in}{0.625278in}}%
\pgfpathlineto{\pgfqpoint{4.099509in}{0.625278in}}%
\pgfpathlineto{\pgfqpoint{4.100516in}{0.633527in}}%
\pgfpathlineto{\pgfqpoint{4.100516in}{0.625278in}}%
\pgfpathlineto{\pgfqpoint{4.101856in}{0.633527in}}%
\pgfpathlineto{\pgfqpoint{4.102189in}{0.625278in}}%
\pgfpathlineto{\pgfqpoint{4.102693in}{0.625278in}}%
\pgfpathlineto{\pgfqpoint{4.103863in}{0.633527in}}%
\pgfpathlineto{\pgfqpoint{4.104203in}{0.625278in}}%
\pgfpathlineto{\pgfqpoint{4.104714in}{0.625278in}}%
\pgfpathlineto{\pgfqpoint{4.105877in}{0.633527in}}%
\pgfpathlineto{\pgfqpoint{4.106054in}{0.625278in}}%
\pgfpathlineto{\pgfqpoint{4.106884in}{0.625278in}}%
\pgfpathlineto{\pgfqpoint{4.108054in}{0.633527in}}%
\pgfpathlineto{\pgfqpoint{4.108224in}{0.625278in}}%
\pgfpathlineto{\pgfqpoint{4.109068in}{0.625278in}}%
\pgfpathlineto{\pgfqpoint{4.109897in}{0.633527in}}%
\pgfpathlineto{\pgfqpoint{4.109897in}{0.625278in}}%
\pgfpathlineto{\pgfqpoint{4.111074in}{0.633527in}}%
\pgfpathlineto{\pgfqpoint{4.111578in}{0.625278in}}%
\pgfpathlineto{\pgfqpoint{4.112081in}{0.625278in}}%
\pgfpathlineto{\pgfqpoint{4.113251in}{0.633527in}}%
\pgfpathlineto{\pgfqpoint{4.113251in}{0.625278in}}%
\pgfpathlineto{\pgfqpoint{4.114258in}{0.625278in}}%
\pgfpathlineto{\pgfqpoint{4.115436in}{0.633527in}}%
\pgfpathlineto{\pgfqpoint{4.115598in}{0.625278in}}%
\pgfpathlineto{\pgfqpoint{4.116280in}{0.625278in}}%
\pgfpathlineto{\pgfqpoint{4.117442in}{0.633527in}}%
\pgfpathlineto{\pgfqpoint{4.117612in}{0.625278in}}%
\pgfpathlineto{\pgfqpoint{4.118456in}{0.625278in}}%
\pgfpathlineto{\pgfqpoint{4.119626in}{0.633527in}}%
\pgfpathlineto{\pgfqpoint{4.119789in}{0.625278in}}%
\pgfpathlineto{\pgfqpoint{4.120633in}{0.625278in}}%
\pgfpathlineto{\pgfqpoint{4.121803in}{0.633527in}}%
\pgfpathlineto{\pgfqpoint{4.121966in}{0.625278in}}%
\pgfpathlineto{\pgfqpoint{4.122640in}{0.625278in}}%
\pgfpathlineto{\pgfqpoint{4.123817in}{0.633527in}}%
\pgfpathlineto{\pgfqpoint{4.124150in}{0.625278in}}%
\pgfpathlineto{\pgfqpoint{4.124817in}{0.625278in}}%
\pgfpathlineto{\pgfqpoint{4.125994in}{0.633527in}}%
\pgfpathlineto{\pgfqpoint{4.126157in}{0.625278in}}%
\pgfpathlineto{\pgfqpoint{4.127001in}{0.625278in}}%
\pgfpathlineto{\pgfqpoint{4.128171in}{0.633527in}}%
\pgfpathlineto{\pgfqpoint{4.128504in}{0.625278in}}%
\pgfpathlineto{\pgfqpoint{4.129178in}{0.625278in}}%
\pgfpathlineto{\pgfqpoint{4.130518in}{0.633527in}}%
\pgfpathlineto{\pgfqpoint{4.130688in}{0.625278in}}%
\pgfpathlineto{\pgfqpoint{4.131525in}{0.625278in}}%
\pgfpathlineto{\pgfqpoint{4.132702in}{0.633527in}}%
\pgfpathlineto{\pgfqpoint{4.132865in}{0.625278in}}%
\pgfpathlineto{\pgfqpoint{4.133709in}{0.625278in}}%
\pgfpathlineto{\pgfqpoint{4.134879in}{0.633527in}}%
\pgfpathlineto{\pgfqpoint{4.134879in}{0.625278in}}%
\pgfpathlineto{\pgfqpoint{4.135886in}{0.625278in}}%
\pgfpathlineto{\pgfqpoint{4.137056in}{0.633527in}}%
\pgfpathlineto{\pgfqpoint{4.137056in}{0.625278in}}%
\pgfpathlineto{\pgfqpoint{4.138063in}{0.625278in}}%
\pgfpathlineto{\pgfqpoint{4.139233in}{0.633527in}}%
\pgfpathlineto{\pgfqpoint{4.139233in}{0.625278in}}%
\pgfpathlineto{\pgfqpoint{4.140077in}{0.625278in}}%
\pgfpathlineto{\pgfqpoint{4.141247in}{0.633527in}}%
\pgfpathlineto{\pgfqpoint{4.141417in}{0.625278in}}%
\pgfpathlineto{\pgfqpoint{4.142253in}{0.625278in}}%
\pgfpathlineto{\pgfqpoint{4.143594in}{0.633527in}}%
\pgfpathlineto{\pgfqpoint{4.143594in}{0.625278in}}%
\pgfpathlineto{\pgfqpoint{4.144601in}{0.625278in}}%
\pgfpathlineto{\pgfqpoint{4.145941in}{0.633527in}}%
\pgfpathlineto{\pgfqpoint{4.146111in}{0.625278in}}%
\pgfpathlineto{\pgfqpoint{4.146948in}{0.625278in}}%
\pgfpathlineto{\pgfqpoint{4.147955in}{0.633527in}}%
\pgfpathlineto{\pgfqpoint{4.147955in}{0.625278in}}%
\pgfpathlineto{\pgfqpoint{4.149295in}{0.633527in}}%
\pgfpathlineto{\pgfqpoint{4.149458in}{0.625278in}}%
\pgfpathlineto{\pgfqpoint{4.150132in}{0.625278in}}%
\pgfpathlineto{\pgfqpoint{4.151301in}{0.633527in}}%
\pgfpathlineto{\pgfqpoint{4.151642in}{0.625278in}}%
\pgfpathlineto{\pgfqpoint{4.152308in}{0.625278in}}%
\pgfpathlineto{\pgfqpoint{4.153486in}{0.633527in}}%
\pgfpathlineto{\pgfqpoint{4.153649in}{0.625278in}}%
\pgfpathlineto{\pgfqpoint{4.154493in}{0.625278in}}%
\pgfpathlineto{\pgfqpoint{4.155662in}{0.633527in}}%
\pgfpathlineto{\pgfqpoint{4.155833in}{0.625278in}}%
\pgfpathlineto{\pgfqpoint{4.156669in}{0.625278in}}%
\pgfpathlineto{\pgfqpoint{4.157839in}{0.633527in}}%
\pgfpathlineto{\pgfqpoint{4.158010in}{0.625278in}}%
\pgfpathlineto{\pgfqpoint{4.158846in}{0.625278in}}%
\pgfpathlineto{\pgfqpoint{4.159853in}{0.633527in}}%
\pgfpathlineto{\pgfqpoint{4.159853in}{0.625278in}}%
\pgfpathlineto{\pgfqpoint{4.161193in}{0.633527in}}%
\pgfpathlineto{\pgfqpoint{4.161534in}{0.625278in}}%
\pgfpathlineto{\pgfqpoint{4.162200in}{0.625278in}}%
\pgfpathlineto{\pgfqpoint{4.163207in}{0.633527in}}%
\pgfpathlineto{\pgfqpoint{4.163207in}{0.625278in}}%
\pgfpathlineto{\pgfqpoint{4.164548in}{0.633527in}}%
\pgfpathlineto{\pgfqpoint{4.164548in}{0.625278in}}%
\pgfpathlineto{\pgfqpoint{4.165554in}{0.625278in}}%
\pgfpathlineto{\pgfqpoint{4.166724in}{0.633527in}}%
\pgfpathlineto{\pgfqpoint{4.166724in}{0.625278in}}%
\pgfpathlineto{\pgfqpoint{4.167731in}{0.625278in}}%
\pgfpathlineto{\pgfqpoint{4.168909in}{0.633527in}}%
\pgfpathlineto{\pgfqpoint{4.169071in}{0.625278in}}%
\pgfpathlineto{\pgfqpoint{4.169916in}{0.625278in}}%
\pgfpathlineto{\pgfqpoint{4.171085in}{0.633527in}}%
\pgfpathlineto{\pgfqpoint{4.171256in}{0.625278in}}%
\pgfpathlineto{\pgfqpoint{4.172092in}{0.625278in}}%
\pgfpathlineto{\pgfqpoint{4.173262in}{0.633527in}}%
\pgfpathlineto{\pgfqpoint{4.173433in}{0.625278in}}%
\pgfpathlineto{\pgfqpoint{4.174269in}{0.625278in}}%
\pgfpathlineto{\pgfqpoint{4.175447in}{0.633527in}}%
\pgfpathlineto{\pgfqpoint{4.175447in}{0.625278in}}%
\pgfpathlineto{\pgfqpoint{4.176446in}{0.625278in}}%
\pgfpathlineto{\pgfqpoint{4.177794in}{0.633527in}}%
\pgfpathlineto{\pgfqpoint{4.177794in}{0.625278in}}%
\pgfpathlineto{\pgfqpoint{4.178630in}{0.625278in}}%
\pgfpathlineto{\pgfqpoint{4.179800in}{0.633527in}}%
\pgfpathlineto{\pgfqpoint{4.179970in}{0.625278in}}%
\pgfpathlineto{\pgfqpoint{4.180807in}{0.625278in}}%
\pgfpathlineto{\pgfqpoint{4.181984in}{0.633527in}}%
\pgfpathlineto{\pgfqpoint{4.182147in}{0.625278in}}%
\pgfpathlineto{\pgfqpoint{4.182984in}{0.625278in}}%
\pgfpathlineto{\pgfqpoint{4.184161in}{0.633527in}}%
\pgfpathlineto{\pgfqpoint{4.184332in}{0.625278in}}%
\pgfpathlineto{\pgfqpoint{4.185168in}{0.625278in}}%
\pgfpathlineto{\pgfqpoint{4.186338in}{0.633527in}}%
\pgfpathlineto{\pgfqpoint{4.186508in}{0.625278in}}%
\pgfpathlineto{\pgfqpoint{4.187345in}{0.625278in}}%
\pgfpathlineto{\pgfqpoint{4.188522in}{0.633527in}}%
\pgfpathlineto{\pgfqpoint{4.188685in}{0.625278in}}%
\pgfpathlineto{\pgfqpoint{4.189522in}{0.625278in}}%
\pgfpathlineto{\pgfqpoint{4.190699in}{0.633527in}}%
\pgfpathlineto{\pgfqpoint{4.190869in}{0.625278in}}%
\pgfpathlineto{\pgfqpoint{4.191706in}{0.625278in}}%
\pgfpathlineto{\pgfqpoint{4.192876in}{0.633527in}}%
\pgfpathlineto{\pgfqpoint{4.192876in}{0.625278in}}%
\pgfpathlineto{\pgfqpoint{4.193883in}{0.625278in}}%
\pgfpathlineto{\pgfqpoint{4.195060in}{0.633527in}}%
\pgfpathlineto{\pgfqpoint{4.195223in}{0.625278in}}%
\pgfpathlineto{\pgfqpoint{4.196060in}{0.625278in}}%
\pgfpathlineto{\pgfqpoint{4.197237in}{0.633527in}}%
\pgfpathlineto{\pgfqpoint{4.197400in}{0.625278in}}%
\pgfpathlineto{\pgfqpoint{4.198244in}{0.625278in}}%
\pgfpathlineto{\pgfqpoint{4.199414in}{0.633527in}}%
\pgfpathlineto{\pgfqpoint{4.199584in}{0.625278in}}%
\pgfpathlineto{\pgfqpoint{4.200421in}{0.625278in}}%
\pgfpathlineto{\pgfqpoint{4.201591in}{0.633527in}}%
\pgfpathlineto{\pgfqpoint{4.201761in}{0.625278in}}%
\pgfpathlineto{\pgfqpoint{4.202598in}{0.625278in}}%
\pgfpathlineto{\pgfqpoint{4.203775in}{0.633527in}}%
\pgfpathlineto{\pgfqpoint{4.203775in}{0.625278in}}%
\pgfpathlineto{\pgfqpoint{4.204612in}{0.625278in}}%
\pgfpathlineto{\pgfqpoint{4.205952in}{0.633527in}}%
\pgfpathlineto{\pgfqpoint{4.205952in}{0.625278in}}%
\pgfpathlineto{\pgfqpoint{4.206959in}{0.625278in}}%
\pgfpathlineto{\pgfqpoint{4.208129in}{0.633527in}}%
\pgfpathlineto{\pgfqpoint{4.208299in}{0.625278in}}%
\pgfpathlineto{\pgfqpoint{4.209136in}{0.625278in}}%
\pgfpathlineto{\pgfqpoint{4.210313in}{0.633527in}}%
\pgfpathlineto{\pgfqpoint{4.210476in}{0.625278in}}%
\pgfpathlineto{\pgfqpoint{4.211150in}{0.625278in}}%
\pgfpathlineto{\pgfqpoint{4.212490in}{0.633527in}}%
\pgfpathlineto{\pgfqpoint{4.212823in}{0.625278in}}%
\pgfpathlineto{\pgfqpoint{4.213497in}{0.625278in}}%
\pgfpathlineto{\pgfqpoint{4.214667in}{0.633527in}}%
\pgfpathlineto{\pgfqpoint{4.214837in}{0.625278in}}%
\pgfpathlineto{\pgfqpoint{4.215511in}{0.625278in}}%
\pgfpathlineto{\pgfqpoint{4.216680in}{0.633527in}}%
\pgfpathlineto{\pgfqpoint{4.216851in}{0.625278in}}%
\pgfpathlineto{\pgfqpoint{4.217687in}{0.625278in}}%
\pgfpathlineto{\pgfqpoint{4.218857in}{0.633527in}}%
\pgfpathlineto{\pgfqpoint{4.219028in}{0.625278in}}%
\pgfpathlineto{\pgfqpoint{4.219864in}{0.625278in}}%
\pgfpathlineto{\pgfqpoint{4.221042in}{0.633527in}}%
\pgfpathlineto{\pgfqpoint{4.221204in}{0.625278in}}%
\pgfpathlineto{\pgfqpoint{4.222049in}{0.625278in}}%
\pgfpathlineto{\pgfqpoint{4.223218in}{0.633527in}}%
\pgfpathlineto{\pgfqpoint{4.223218in}{0.625278in}}%
\pgfpathlineto{\pgfqpoint{4.224225in}{0.625278in}}%
\pgfpathlineto{\pgfqpoint{4.225395in}{0.633527in}}%
\pgfpathlineto{\pgfqpoint{4.225417in}{0.625278in}}%
\pgfpathlineto{\pgfqpoint{4.226402in}{0.625278in}}%
\pgfpathlineto{\pgfqpoint{4.227579in}{0.633527in}}%
\pgfpathlineto{\pgfqpoint{4.227742in}{0.625278in}}%
\pgfpathlineto{\pgfqpoint{4.228601in}{0.625278in}}%
\pgfpathlineto{\pgfqpoint{4.228934in}{0.650025in}}%
\pgfpathlineto{\pgfqpoint{4.229756in}{0.633527in}}%
\pgfpathlineto{\pgfqpoint{4.229778in}{0.625278in}}%
\pgfpathlineto{\pgfqpoint{4.230763in}{0.625278in}}%
\pgfpathlineto{\pgfqpoint{4.231933in}{0.633527in}}%
\pgfpathlineto{\pgfqpoint{4.232792in}{0.625278in}}%
\pgfpathlineto{\pgfqpoint{4.233969in}{0.650025in}}%
\pgfpathlineto{\pgfqpoint{4.234132in}{0.641776in}}%
\pgfpathlineto{\pgfqpoint{4.234302in}{0.650025in}}%
\pgfpathlineto{\pgfqpoint{4.235287in}{0.633527in}}%
\pgfpathlineto{\pgfqpoint{4.235643in}{0.674773in}}%
\pgfpathlineto{\pgfqpoint{4.236316in}{0.674773in}}%
\pgfpathlineto{\pgfqpoint{4.237657in}{0.650025in}}%
\pgfpathlineto{\pgfqpoint{4.238323in}{0.724268in}}%
\pgfpathlineto{\pgfqpoint{4.239308in}{0.683022in}}%
\pgfpathlineto{\pgfqpoint{4.240337in}{0.650025in}}%
\pgfpathlineto{\pgfqpoint{4.240485in}{0.658274in}}%
\pgfpathlineto{\pgfqpoint{4.241492in}{0.732517in}}%
\pgfpathlineto{\pgfqpoint{4.241507in}{0.724268in}}%
\pgfpathlineto{\pgfqpoint{4.243002in}{0.806759in}}%
\pgfpathlineto{\pgfqpoint{4.243521in}{0.790261in}}%
\pgfpathlineto{\pgfqpoint{4.243669in}{0.782012in}}%
\pgfpathlineto{\pgfqpoint{4.243839in}{0.806759in}}%
\pgfpathlineto{\pgfqpoint{4.244194in}{0.798510in}}%
\pgfpathlineto{\pgfqpoint{4.244846in}{0.806759in}}%
\pgfpathlineto{\pgfqpoint{4.245031in}{0.790261in}}%
\pgfpathlineto{\pgfqpoint{4.245364in}{0.773763in}}%
\pgfpathlineto{\pgfqpoint{4.245683in}{0.806759in}}%
\pgfpathlineto{\pgfqpoint{4.245846in}{0.806759in}}%
\pgfpathlineto{\pgfqpoint{4.246016in}{0.831507in}}%
\pgfpathlineto{\pgfqpoint{4.246875in}{0.823258in}}%
\pgfpathlineto{\pgfqpoint{4.247526in}{0.856254in}}%
\pgfpathlineto{\pgfqpoint{4.248045in}{0.839756in}}%
\pgfpathlineto{\pgfqpoint{4.248193in}{0.831507in}}%
\pgfpathlineto{\pgfqpoint{4.248696in}{0.864503in}}%
\pgfpathlineto{\pgfqpoint{4.248718in}{0.872753in}}%
\pgfpathlineto{\pgfqpoint{4.249052in}{0.848005in}}%
\pgfpathlineto{\pgfqpoint{4.249725in}{0.848005in}}%
\pgfpathlineto{\pgfqpoint{4.250540in}{0.881002in}}%
\pgfpathlineto{\pgfqpoint{4.250880in}{0.856254in}}%
\pgfpathlineto{\pgfqpoint{4.251880in}{0.831507in}}%
\pgfpathlineto{\pgfqpoint{4.251902in}{0.848005in}}%
\pgfpathlineto{\pgfqpoint{4.253079in}{0.823258in}}%
\pgfpathlineto{\pgfqpoint{4.253228in}{0.831507in}}%
\pgfpathlineto{\pgfqpoint{4.253916in}{0.872753in}}%
\pgfpathlineto{\pgfqpoint{4.254249in}{0.872753in}}%
\pgfpathlineto{\pgfqpoint{4.255234in}{0.856254in}}%
\pgfpathlineto{\pgfqpoint{4.255256in}{0.872753in}}%
\pgfpathlineto{\pgfqpoint{4.256093in}{0.897500in}}%
\pgfpathlineto{\pgfqpoint{4.256411in}{0.881002in}}%
\pgfpathlineto{\pgfqpoint{4.257433in}{0.872753in}}%
\pgfpathlineto{\pgfqpoint{4.257752in}{0.881002in}}%
\pgfpathlineto{\pgfqpoint{4.258085in}{0.856254in}}%
\pgfpathlineto{\pgfqpoint{4.258610in}{0.872753in}}%
\pgfpathlineto{\pgfqpoint{4.259114in}{0.848005in}}%
\pgfpathlineto{\pgfqpoint{4.260617in}{0.798510in}}%
\pgfpathlineto{\pgfqpoint{4.260765in}{0.815008in}}%
\pgfpathlineto{\pgfqpoint{4.261106in}{0.831507in}}%
\pgfpathlineto{\pgfqpoint{4.261461in}{0.798510in}}%
\pgfpathlineto{\pgfqpoint{4.261794in}{0.798510in}}%
\pgfpathlineto{\pgfqpoint{4.262779in}{0.831507in}}%
\pgfpathlineto{\pgfqpoint{4.262964in}{0.815008in}}%
\pgfpathlineto{\pgfqpoint{4.264475in}{0.773763in}}%
\pgfpathlineto{\pgfqpoint{4.264623in}{0.782012in}}%
\pgfpathlineto{\pgfqpoint{4.264793in}{0.806759in}}%
\pgfpathlineto{\pgfqpoint{4.265652in}{0.798510in}}%
\pgfpathlineto{\pgfqpoint{4.266148in}{0.773763in}}%
\pgfpathlineto{\pgfqpoint{4.266466in}{0.806759in}}%
\pgfpathlineto{\pgfqpoint{4.266488in}{0.798510in}}%
\pgfpathlineto{\pgfqpoint{4.268984in}{0.881002in}}%
\pgfpathlineto{\pgfqpoint{4.268999in}{0.872753in}}%
\pgfpathlineto{\pgfqpoint{4.269991in}{0.930497in}}%
\pgfpathlineto{\pgfqpoint{4.270005in}{0.922247in}}%
\pgfpathlineto{\pgfqpoint{4.270827in}{0.955244in}}%
\pgfpathlineto{\pgfqpoint{4.271161in}{0.930497in}}%
\pgfpathlineto{\pgfqpoint{4.272182in}{0.922247in}}%
\pgfpathlineto{\pgfqpoint{4.273678in}{0.979992in}}%
\pgfpathlineto{\pgfqpoint{4.273863in}{0.963493in}}%
\pgfpathlineto{\pgfqpoint{4.275033in}{0.946995in}}%
\pgfpathlineto{\pgfqpoint{4.276543in}{0.996490in}}%
\pgfpathlineto{\pgfqpoint{4.276714in}{0.988241in}}%
\pgfpathlineto{\pgfqpoint{4.276862in}{0.979992in}}%
\pgfpathlineto{\pgfqpoint{4.277365in}{1.012988in}}%
\pgfpathlineto{\pgfqpoint{4.278387in}{1.070732in}}%
\pgfpathlineto{\pgfqpoint{4.278557in}{1.062483in}}%
\pgfpathlineto{\pgfqpoint{4.279372in}{1.054234in}}%
\pgfpathlineto{\pgfqpoint{4.278720in}{1.070732in}}%
\pgfpathlineto{\pgfqpoint{4.279372in}{1.062483in}}%
\pgfpathlineto{\pgfqpoint{4.279394in}{1.095480in}}%
\pgfpathlineto{\pgfqpoint{4.280401in}{1.070732in}}%
\pgfpathlineto{\pgfqpoint{4.280549in}{1.078981in}}%
\pgfpathlineto{\pgfqpoint{4.280904in}{1.045985in}}%
\pgfpathlineto{\pgfqpoint{4.281067in}{1.045985in}}%
\pgfpathlineto{\pgfqpoint{4.281386in}{1.029487in}}%
\pgfpathlineto{\pgfqpoint{4.281556in}{1.054234in}}%
\pgfpathlineto{\pgfqpoint{4.281911in}{1.045985in}}%
\pgfpathlineto{\pgfqpoint{4.282726in}{1.054234in}}%
\pgfpathlineto{\pgfqpoint{4.282748in}{1.045985in}}%
\pgfpathlineto{\pgfqpoint{4.283066in}{1.029487in}}%
\pgfpathlineto{\pgfqpoint{4.283229in}{1.054234in}}%
\pgfpathlineto{\pgfqpoint{4.283563in}{1.054234in}}%
\pgfpathlineto{\pgfqpoint{4.283585in}{1.070732in}}%
\pgfpathlineto{\pgfqpoint{4.283918in}{1.045985in}}%
\pgfpathlineto{\pgfqpoint{4.284592in}{1.045985in}}%
\pgfpathlineto{\pgfqpoint{4.285073in}{1.029487in}}%
\pgfpathlineto{\pgfqpoint{4.285747in}{1.054234in}}%
\pgfpathlineto{\pgfqpoint{4.286769in}{1.045985in}}%
\pgfpathlineto{\pgfqpoint{4.287420in}{1.078981in}}%
\pgfpathlineto{\pgfqpoint{4.287946in}{1.062483in}}%
\pgfpathlineto{\pgfqpoint{4.288760in}{1.054234in}}%
\pgfpathlineto{\pgfqpoint{4.288109in}{1.070732in}}%
\pgfpathlineto{\pgfqpoint{4.288783in}{1.070732in}}%
\pgfpathlineto{\pgfqpoint{4.289938in}{1.078981in}}%
\pgfpathlineto{\pgfqpoint{4.291130in}{1.021237in}}%
\pgfpathlineto{\pgfqpoint{4.291611in}{1.029487in}}%
\pgfpathlineto{\pgfqpoint{4.291944in}{1.004739in}}%
\pgfpathlineto{\pgfqpoint{4.292300in}{0.996490in}}%
\pgfpathlineto{\pgfqpoint{4.292618in}{1.029487in}}%
\pgfpathlineto{\pgfqpoint{4.292973in}{1.021237in}}%
\pgfpathlineto{\pgfqpoint{4.294291in}{1.054234in}}%
\pgfpathlineto{\pgfqpoint{4.294313in}{1.045985in}}%
\pgfpathlineto{\pgfqpoint{4.294817in}{1.021237in}}%
\pgfpathlineto{\pgfqpoint{4.295135in}{1.054234in}}%
\pgfpathlineto{\pgfqpoint{4.295150in}{1.045985in}}%
\pgfpathlineto{\pgfqpoint{4.296305in}{1.054234in}}%
\pgfpathlineto{\pgfqpoint{4.297327in}{1.045985in}}%
\pgfpathlineto{\pgfqpoint{4.298482in}{1.103729in}}%
\pgfpathlineto{\pgfqpoint{4.298985in}{1.078981in}}%
\pgfpathlineto{\pgfqpoint{4.300178in}{1.045985in}}%
\pgfpathlineto{\pgfqpoint{4.300496in}{1.054234in}}%
\pgfpathlineto{\pgfqpoint{4.300829in}{1.029487in}}%
\pgfpathlineto{\pgfqpoint{4.301333in}{1.054234in}}%
\pgfpathlineto{\pgfqpoint{4.301858in}{1.021237in}}%
\pgfpathlineto{\pgfqpoint{4.303028in}{1.045985in}}%
\pgfpathlineto{\pgfqpoint{4.303198in}{1.037736in}}%
\pgfpathlineto{\pgfqpoint{4.303347in}{1.029487in}}%
\pgfpathlineto{\pgfqpoint{4.303517in}{1.054234in}}%
\pgfpathlineto{\pgfqpoint{4.303865in}{1.045985in}}%
\pgfpathlineto{\pgfqpoint{4.304687in}{1.054234in}}%
\pgfpathlineto{\pgfqpoint{4.304709in}{1.045985in}}%
\pgfpathlineto{\pgfqpoint{4.305546in}{1.021237in}}%
\pgfpathlineto{\pgfqpoint{4.305916in}{1.029487in}}%
\pgfpathlineto{\pgfqpoint{4.307034in}{1.054234in}}%
\pgfpathlineto{\pgfqpoint{4.307056in}{1.045985in}}%
\pgfpathlineto{\pgfqpoint{4.307560in}{1.021237in}}%
\pgfpathlineto{\pgfqpoint{4.307871in}{1.054234in}}%
\pgfpathlineto{\pgfqpoint{4.307893in}{1.045985in}}%
\pgfpathlineto{\pgfqpoint{4.308878in}{1.054234in}}%
\pgfpathlineto{\pgfqpoint{4.308900in}{1.045985in}}%
\pgfpathlineto{\pgfqpoint{4.310240in}{1.021237in}}%
\pgfpathlineto{\pgfqpoint{4.312061in}{1.078981in}}%
\pgfpathlineto{\pgfqpoint{4.312084in}{1.070732in}}%
\pgfpathlineto{\pgfqpoint{4.313068in}{1.054234in}}%
\pgfpathlineto{\pgfqpoint{4.313090in}{1.070732in}}%
\pgfpathlineto{\pgfqpoint{4.314742in}{1.103729in}}%
\pgfpathlineto{\pgfqpoint{4.315941in}{1.070732in}}%
\pgfpathlineto{\pgfqpoint{4.316756in}{1.078981in}}%
\pgfpathlineto{\pgfqpoint{4.316778in}{1.070732in}}%
\pgfpathlineto{\pgfqpoint{4.317089in}{1.054234in}}%
\pgfpathlineto{\pgfqpoint{4.317259in}{1.078981in}}%
\pgfpathlineto{\pgfqpoint{4.317614in}{1.070732in}}%
\pgfpathlineto{\pgfqpoint{4.318770in}{1.078981in}}%
\pgfpathlineto{\pgfqpoint{4.319628in}{1.070732in}}%
\pgfpathlineto{\pgfqpoint{4.319791in}{1.095480in}}%
\pgfpathlineto{\pgfqpoint{4.320783in}{1.128476in}}%
\pgfpathlineto{\pgfqpoint{4.320969in}{1.111978in}}%
\pgfpathlineto{\pgfqpoint{4.321131in}{1.120227in}}%
\pgfpathlineto{\pgfqpoint{4.322138in}{1.095480in}}%
\pgfpathlineto{\pgfqpoint{4.323464in}{1.153224in}}%
\pgfpathlineto{\pgfqpoint{4.323649in}{1.136726in}}%
\pgfpathlineto{\pgfqpoint{4.323982in}{1.095480in}}%
\pgfpathlineto{\pgfqpoint{4.324819in}{1.120227in}}%
\pgfpathlineto{\pgfqpoint{4.326499in}{1.169722in}}%
\pgfpathlineto{\pgfqpoint{4.326662in}{1.161473in}}%
\pgfpathlineto{\pgfqpoint{4.327840in}{1.144975in}}%
\pgfpathlineto{\pgfqpoint{4.328321in}{1.177971in}}%
\pgfpathlineto{\pgfqpoint{4.328995in}{1.153224in}}%
\pgfpathlineto{\pgfqpoint{4.329180in}{1.169722in}}%
\pgfpathlineto{\pgfqpoint{4.330016in}{1.144975in}}%
\pgfpathlineto{\pgfqpoint{4.330838in}{1.153224in}}%
\pgfpathlineto{\pgfqpoint{4.330853in}{1.144975in}}%
\pgfpathlineto{\pgfqpoint{4.331172in}{1.128476in}}%
\pgfpathlineto{\pgfqpoint{4.331505in}{1.161473in}}%
\pgfpathlineto{\pgfqpoint{4.332534in}{1.219217in}}%
\pgfpathlineto{\pgfqpoint{4.332704in}{1.210968in}}%
\pgfpathlineto{\pgfqpoint{4.332867in}{1.219217in}}%
\pgfpathlineto{\pgfqpoint{4.333874in}{1.194470in}}%
\pgfpathlineto{\pgfqpoint{4.334526in}{1.202719in}}%
\pgfpathlineto{\pgfqpoint{4.334711in}{1.186221in}}%
\pgfpathlineto{\pgfqpoint{4.335199in}{1.177971in}}%
\pgfpathlineto{\pgfqpoint{4.335362in}{1.202719in}}%
\pgfpathlineto{\pgfqpoint{4.335385in}{1.194470in}}%
\pgfpathlineto{\pgfqpoint{4.336540in}{1.252214in}}%
\pgfpathlineto{\pgfqpoint{4.337043in}{1.227466in}}%
\pgfpathlineto{\pgfqpoint{4.337895in}{1.219217in}}%
\pgfpathlineto{\pgfqpoint{4.338065in}{1.243965in}}%
\pgfpathlineto{\pgfqpoint{4.338568in}{1.219217in}}%
\pgfpathlineto{\pgfqpoint{4.338887in}{1.252214in}}%
\pgfpathlineto{\pgfqpoint{4.338902in}{1.243965in}}%
\pgfpathlineto{\pgfqpoint{4.339398in}{1.227466in}}%
\pgfpathlineto{\pgfqpoint{4.340560in}{1.276961in}}%
\pgfpathlineto{\pgfqpoint{4.340582in}{1.268712in}}%
\pgfpathlineto{\pgfqpoint{4.341234in}{1.326456in}}%
\pgfpathlineto{\pgfqpoint{4.341419in}{1.318207in}}%
\pgfpathlineto{\pgfqpoint{4.341737in}{1.326456in}}%
\pgfpathlineto{\pgfqpoint{4.342071in}{1.301709in}}%
\pgfpathlineto{\pgfqpoint{4.342256in}{1.309958in}}%
\pgfpathlineto{\pgfqpoint{4.342426in}{1.318207in}}%
\pgfpathlineto{\pgfqpoint{4.343425in}{1.293460in}}%
\pgfpathlineto{\pgfqpoint{4.344270in}{1.318207in}}%
\pgfpathlineto{\pgfqpoint{4.344581in}{1.301709in}}%
\pgfpathlineto{\pgfqpoint{4.344921in}{1.276961in}}%
\pgfpathlineto{\pgfqpoint{4.345439in}{1.318207in}}%
\pgfpathlineto{\pgfqpoint{4.345610in}{1.309958in}}%
\pgfpathlineto{\pgfqpoint{4.347120in}{1.243965in}}%
\pgfpathlineto{\pgfqpoint{4.347601in}{1.252214in}}%
\pgfpathlineto{\pgfqpoint{4.347935in}{1.227466in}}%
\pgfpathlineto{\pgfqpoint{4.348964in}{1.219217in}}%
\pgfpathlineto{\pgfqpoint{4.349615in}{1.227466in}}%
\pgfpathlineto{\pgfqpoint{4.349800in}{1.210968in}}%
\pgfpathlineto{\pgfqpoint{4.350134in}{1.194470in}}%
\pgfpathlineto{\pgfqpoint{4.350452in}{1.227466in}}%
\pgfpathlineto{\pgfqpoint{4.350956in}{1.202719in}}%
\pgfpathlineto{\pgfqpoint{4.351474in}{1.243965in}}%
\pgfpathlineto{\pgfqpoint{4.351977in}{1.235715in}}%
\pgfpathlineto{\pgfqpoint{4.352459in}{1.227466in}}%
\pgfpathlineto{\pgfqpoint{4.352629in}{1.252214in}}%
\pgfpathlineto{\pgfqpoint{4.352651in}{1.243965in}}%
\pgfpathlineto{\pgfqpoint{4.353806in}{1.252214in}}%
\pgfpathlineto{\pgfqpoint{4.354658in}{1.219217in}}%
\pgfpathlineto{\pgfqpoint{4.354828in}{1.243965in}}%
\pgfpathlineto{\pgfqpoint{4.356338in}{1.293460in}}%
\pgfpathlineto{\pgfqpoint{4.356501in}{1.285210in}}%
\pgfpathlineto{\pgfqpoint{4.356649in}{1.276961in}}%
\pgfpathlineto{\pgfqpoint{4.356820in}{1.301709in}}%
\pgfpathlineto{\pgfqpoint{4.357175in}{1.293460in}}%
\pgfpathlineto{\pgfqpoint{4.359167in}{1.351204in}}%
\pgfpathlineto{\pgfqpoint{4.359352in}{1.342955in}}%
\pgfpathlineto{\pgfqpoint{4.359678in}{1.400699in}}%
\pgfpathlineto{\pgfqpoint{4.360189in}{1.384200in}}%
\pgfpathlineto{\pgfqpoint{4.360359in}{1.392449in}}%
\pgfpathlineto{\pgfqpoint{4.361344in}{1.375951in}}%
\pgfpathlineto{\pgfqpoint{4.361366in}{1.392449in}}%
\pgfpathlineto{\pgfqpoint{4.362373in}{1.367702in}}%
\pgfpathlineto{\pgfqpoint{4.362861in}{1.375951in}}%
\pgfpathlineto{\pgfqpoint{4.364216in}{1.318207in}}%
\pgfpathlineto{\pgfqpoint{4.364379in}{1.342955in}}%
\pgfpathlineto{\pgfqpoint{4.365372in}{1.326456in}}%
\pgfpathlineto{\pgfqpoint{4.366208in}{1.301709in}}%
\pgfpathlineto{\pgfqpoint{4.366393in}{1.309958in}}%
\pgfpathlineto{\pgfqpoint{4.366882in}{1.301709in}}%
\pgfpathlineto{\pgfqpoint{4.367045in}{1.326456in}}%
\pgfpathlineto{\pgfqpoint{4.367067in}{1.318207in}}%
\pgfpathlineto{\pgfqpoint{4.368230in}{1.351204in}}%
\pgfpathlineto{\pgfqpoint{4.368407in}{1.334705in}}%
\pgfpathlineto{\pgfqpoint{4.370081in}{1.293460in}}%
\pgfpathlineto{\pgfqpoint{4.370229in}{1.301709in}}%
\pgfpathlineto{\pgfqpoint{4.370569in}{1.276961in}}%
\pgfpathlineto{\pgfqpoint{4.370754in}{1.285210in}}%
\pgfpathlineto{\pgfqpoint{4.372265in}{1.243965in}}%
\pgfpathlineto{\pgfqpoint{4.371073in}{1.301709in}}%
\pgfpathlineto{\pgfqpoint{4.372413in}{1.252214in}}%
\pgfpathlineto{\pgfqpoint{4.373087in}{1.276961in}}%
\pgfpathlineto{\pgfqpoint{4.373435in}{1.243965in}}%
\pgfpathlineto{\pgfqpoint{4.376115in}{1.367702in}}%
\pgfpathlineto{\pgfqpoint{4.376433in}{1.351204in}}%
\pgfpathlineto{\pgfqpoint{4.376767in}{1.375951in}}%
\pgfpathlineto{\pgfqpoint{4.377625in}{1.309958in}}%
\pgfpathlineto{\pgfqpoint{4.377796in}{1.318207in}}%
\pgfpathlineto{\pgfqpoint{4.381987in}{1.120227in}}%
\pgfpathlineto{\pgfqpoint{4.382490in}{1.144975in}}%
\pgfpathlineto{\pgfqpoint{4.382468in}{1.103729in}}%
\pgfpathlineto{\pgfqpoint{4.382823in}{1.120227in}}%
\pgfpathlineto{\pgfqpoint{4.384334in}{1.070732in}}%
\pgfpathlineto{\pgfqpoint{4.384489in}{1.087231in}}%
\pgfpathlineto{\pgfqpoint{4.384837in}{1.120227in}}%
\pgfpathlineto{\pgfqpoint{4.385504in}{1.095480in}}%
\pgfpathlineto{\pgfqpoint{4.386496in}{1.103729in}}%
\pgfpathlineto{\pgfqpoint{4.386510in}{1.095480in}}%
\pgfpathlineto{\pgfqpoint{4.389176in}{1.004739in}}%
\pgfpathlineto{\pgfqpoint{4.389183in}{1.012988in}}%
\pgfpathlineto{\pgfqpoint{4.389687in}{1.029487in}}%
\pgfpathlineto{\pgfqpoint{4.390027in}{0.996490in}}%
\pgfpathlineto{\pgfqpoint{4.390198in}{0.996490in}}%
\pgfpathlineto{\pgfqpoint{4.391693in}{1.029487in}}%
\pgfpathlineto{\pgfqpoint{4.391708in}{1.021237in}}%
\pgfpathlineto{\pgfqpoint{4.392027in}{1.004739in}}%
\pgfpathlineto{\pgfqpoint{4.392197in}{1.029487in}}%
\pgfpathlineto{\pgfqpoint{4.392545in}{1.021237in}}%
\pgfpathlineto{\pgfqpoint{4.392700in}{1.029487in}}%
\pgfpathlineto{\pgfqpoint{4.393048in}{0.996490in}}%
\pgfpathlineto{\pgfqpoint{4.393204in}{1.004739in}}%
\pgfpathlineto{\pgfqpoint{4.393537in}{1.029487in}}%
\pgfpathlineto{\pgfqpoint{4.394389in}{0.963493in}}%
\pgfpathlineto{\pgfqpoint{4.395729in}{0.922247in}}%
\pgfpathlineto{\pgfqpoint{4.394714in}{0.979992in}}%
\pgfpathlineto{\pgfqpoint{4.395884in}{0.938746in}}%
\pgfpathlineto{\pgfqpoint{4.395899in}{0.946995in}}%
\pgfpathlineto{\pgfqpoint{4.396906in}{0.897500in}}%
\pgfpathlineto{\pgfqpoint{4.397891in}{0.881002in}}%
\pgfpathlineto{\pgfqpoint{4.397906in}{0.897500in}}%
\pgfpathlineto{\pgfqpoint{4.399061in}{0.905749in}}%
\pgfpathlineto{\pgfqpoint{4.400090in}{0.897500in}}%
\pgfpathlineto{\pgfqpoint{4.400593in}{0.946995in}}%
\pgfpathlineto{\pgfqpoint{4.401245in}{0.905749in}}%
\pgfpathlineto{\pgfqpoint{4.402437in}{0.815008in}}%
\pgfpathlineto{\pgfqpoint{4.402918in}{0.806759in}}%
\pgfpathlineto{\pgfqpoint{4.402940in}{0.831507in}}%
\pgfpathlineto{\pgfqpoint{4.403274in}{0.823258in}}%
\pgfpathlineto{\pgfqpoint{4.403762in}{0.881002in}}%
\pgfpathlineto{\pgfqpoint{4.404281in}{0.872753in}}%
\pgfpathlineto{\pgfqpoint{4.405102in}{0.930497in}}%
\pgfpathlineto{\pgfqpoint{4.405621in}{0.889251in}}%
\pgfpathlineto{\pgfqpoint{4.405961in}{0.856254in}}%
\pgfpathlineto{\pgfqpoint{4.406124in}{0.897500in}}%
\pgfpathlineto{\pgfqpoint{4.406628in}{0.881002in}}%
\pgfpathlineto{\pgfqpoint{4.407464in}{0.971742in}}%
\pgfpathlineto{\pgfqpoint{4.408805in}{0.938746in}}%
\pgfpathlineto{\pgfqpoint{4.409308in}{0.897500in}}%
\pgfpathlineto{\pgfqpoint{4.409811in}{0.897500in}}%
\pgfpathlineto{\pgfqpoint{4.409960in}{0.905749in}}%
\pgfpathlineto{\pgfqpoint{4.410315in}{0.872753in}}%
\pgfpathlineto{\pgfqpoint{4.411825in}{0.848005in}}%
\pgfpathlineto{\pgfqpoint{4.412995in}{0.872753in}}%
\pgfpathlineto{\pgfqpoint{4.413166in}{0.864503in}}%
\pgfpathlineto{\pgfqpoint{4.413839in}{0.806759in}}%
\pgfpathlineto{\pgfqpoint{4.413328in}{0.872753in}}%
\pgfpathlineto{\pgfqpoint{4.414173in}{0.831507in}}%
\pgfpathlineto{\pgfqpoint{4.415009in}{0.922247in}}%
\pgfpathlineto{\pgfqpoint{4.416349in}{0.905749in}}%
\pgfpathlineto{\pgfqpoint{4.417690in}{0.881002in}}%
\pgfpathlineto{\pgfqpoint{4.419030in}{0.922247in}}%
\pgfpathlineto{\pgfqpoint{4.419200in}{0.913998in}}%
\pgfpathlineto{\pgfqpoint{4.419866in}{0.897500in}}%
\pgfpathlineto{\pgfqpoint{4.419363in}{0.922247in}}%
\pgfpathlineto{\pgfqpoint{4.420037in}{0.922247in}}%
\pgfpathlineto{\pgfqpoint{4.421199in}{0.930497in}}%
\pgfpathlineto{\pgfqpoint{4.422214in}{0.897500in}}%
\pgfpathlineto{\pgfqpoint{4.422221in}{0.905749in}}%
\pgfpathlineto{\pgfqpoint{4.423894in}{0.823258in}}%
\pgfpathlineto{\pgfqpoint{4.424553in}{0.839756in}}%
\pgfpathlineto{\pgfqpoint{4.425234in}{0.848005in}}%
\pgfpathlineto{\pgfqpoint{4.425568in}{0.823258in}}%
\pgfpathlineto{\pgfqpoint{4.426745in}{0.716019in}}%
\pgfpathlineto{\pgfqpoint{4.427248in}{0.658274in}}%
\pgfpathlineto{\pgfqpoint{4.427915in}{0.683022in}}%
\pgfpathlineto{\pgfqpoint{4.428589in}{0.724268in}}%
\pgfpathlineto{\pgfqpoint{4.429092in}{0.691271in}}%
\pgfpathlineto{\pgfqpoint{4.429255in}{0.674773in}}%
\pgfpathlineto{\pgfqpoint{4.430092in}{0.732517in}}%
\pgfpathlineto{\pgfqpoint{4.430092in}{0.724268in}}%
\pgfpathlineto{\pgfqpoint{4.430765in}{0.699520in}}%
\pgfpathlineto{\pgfqpoint{4.430936in}{0.724268in}}%
\pgfpathlineto{\pgfqpoint{4.432098in}{0.732517in}}%
\pgfpathlineto{\pgfqpoint{4.433112in}{0.683022in}}%
\pgfpathlineto{\pgfqpoint{4.433616in}{0.724268in}}%
\pgfpathlineto{\pgfqpoint{4.434119in}{0.691271in}}%
\pgfpathlineto{\pgfqpoint{4.434282in}{0.699520in}}%
\pgfpathlineto{\pgfqpoint{4.436126in}{0.625278in}}%
\pgfpathlineto{\pgfqpoint{4.437288in}{0.633527in}}%
\pgfpathlineto{\pgfqpoint{4.437303in}{0.625278in}}%
\pgfpathlineto{\pgfqpoint{4.438310in}{0.625278in}}%
\pgfpathlineto{\pgfqpoint{4.439643in}{0.633527in}}%
\pgfpathlineto{\pgfqpoint{4.439798in}{0.625278in}}%
\pgfpathlineto{\pgfqpoint{4.440465in}{0.625278in}}%
\pgfpathlineto{\pgfqpoint{4.441657in}{0.633527in}}%
\pgfpathlineto{\pgfqpoint{4.441812in}{0.625278in}}%
\pgfpathlineto{\pgfqpoint{4.442486in}{0.625278in}}%
\pgfpathlineto{\pgfqpoint{4.443663in}{0.633527in}}%
\pgfpathlineto{\pgfqpoint{4.443678in}{0.625278in}}%
\pgfpathlineto{\pgfqpoint{4.444656in}{0.625278in}}%
\pgfpathlineto{\pgfqpoint{4.446011in}{0.633527in}}%
\pgfpathlineto{\pgfqpoint{4.446344in}{0.625278in}}%
\pgfpathlineto{\pgfqpoint{4.447003in}{0.625278in}}%
\pgfpathlineto{\pgfqpoint{4.448180in}{0.633527in}}%
\pgfpathlineto{\pgfqpoint{4.448195in}{0.625278in}}%
\pgfpathlineto{\pgfqpoint{4.449194in}{0.625278in}}%
\pgfpathlineto{\pgfqpoint{4.450535in}{0.633527in}}%
\pgfpathlineto{\pgfqpoint{4.450705in}{0.625278in}}%
\pgfpathlineto{\pgfqpoint{4.451542in}{0.625278in}}%
\pgfpathlineto{\pgfqpoint{4.452711in}{0.633527in}}%
\pgfpathlineto{\pgfqpoint{4.452896in}{0.625278in}}%
\pgfpathlineto{\pgfqpoint{4.453733in}{0.625278in}}%
\pgfpathlineto{\pgfqpoint{4.454903in}{0.633527in}}%
\pgfpathlineto{\pgfqpoint{4.455244in}{0.625278in}}%
\pgfpathlineto{\pgfqpoint{4.455725in}{0.625278in}}%
\pgfpathlineto{\pgfqpoint{4.456895in}{0.633527in}}%
\pgfpathlineto{\pgfqpoint{4.457087in}{0.625278in}}%
\pgfpathlineto{\pgfqpoint{4.457739in}{0.625278in}}%
\pgfpathlineto{\pgfqpoint{4.458909in}{0.633527in}}%
\pgfpathlineto{\pgfqpoint{4.459079in}{0.625278in}}%
\pgfpathlineto{\pgfqpoint{4.459923in}{0.625278in}}%
\pgfpathlineto{\pgfqpoint{4.461256in}{0.633527in}}%
\pgfpathlineto{\pgfqpoint{4.461426in}{0.625278in}}%
\pgfpathlineto{\pgfqpoint{4.462263in}{0.625278in}}%
\pgfpathlineto{\pgfqpoint{4.463433in}{0.633527in}}%
\pgfpathlineto{\pgfqpoint{4.463603in}{0.625278in}}%
\pgfpathlineto{\pgfqpoint{4.464447in}{0.625278in}}%
\pgfpathlineto{\pgfqpoint{4.465610in}{0.633527in}}%
\pgfpathlineto{\pgfqpoint{4.465780in}{0.625278in}}%
\pgfpathlineto{\pgfqpoint{4.466639in}{0.625278in}}%
\pgfpathlineto{\pgfqpoint{4.467794in}{0.633527in}}%
\pgfpathlineto{\pgfqpoint{4.467957in}{0.625278in}}%
\pgfpathlineto{\pgfqpoint{4.468630in}{0.625278in}}%
\pgfpathlineto{\pgfqpoint{4.469800in}{0.633527in}}%
\pgfpathlineto{\pgfqpoint{4.470141in}{0.625278in}}%
\pgfpathlineto{\pgfqpoint{4.470807in}{0.625278in}}%
\pgfpathlineto{\pgfqpoint{4.471985in}{0.633527in}}%
\pgfpathlineto{\pgfqpoint{4.472318in}{0.625278in}}%
\pgfpathlineto{\pgfqpoint{4.472999in}{0.625278in}}%
\pgfpathlineto{\pgfqpoint{4.474161in}{0.633527in}}%
\pgfpathlineto{\pgfqpoint{4.474339in}{0.625278in}}%
\pgfpathlineto{\pgfqpoint{4.475176in}{0.625278in}}%
\pgfpathlineto{\pgfqpoint{4.476338in}{0.633527in}}%
\pgfpathlineto{\pgfqpoint{4.476346in}{0.625278in}}%
\pgfpathlineto{\pgfqpoint{4.477353in}{0.625278in}}%
\pgfpathlineto{\pgfqpoint{4.478522in}{0.633527in}}%
\pgfpathlineto{\pgfqpoint{4.478685in}{0.625278in}}%
\pgfpathlineto{\pgfqpoint{4.479529in}{0.625278in}}%
\pgfpathlineto{\pgfqpoint{4.480707in}{0.633527in}}%
\pgfpathlineto{\pgfqpoint{4.480870in}{0.625278in}}%
\pgfpathlineto{\pgfqpoint{4.481706in}{0.625278in}}%
\pgfpathlineto{\pgfqpoint{4.482876in}{0.633527in}}%
\pgfpathlineto{\pgfqpoint{4.483054in}{0.625278in}}%
\pgfpathlineto{\pgfqpoint{4.483883in}{0.625278in}}%
\pgfpathlineto{\pgfqpoint{4.484897in}{0.633527in}}%
\pgfpathlineto{\pgfqpoint{4.484897in}{0.625278in}}%
\pgfpathlineto{\pgfqpoint{4.486067in}{0.633527in}}%
\pgfpathlineto{\pgfqpoint{4.486400in}{0.625278in}}%
\pgfpathlineto{\pgfqpoint{4.487570in}{0.633527in}}%
\pgfpathlineto{\pgfqpoint{4.488074in}{0.625278in}}%
\pgfpathlineto{\pgfqpoint{4.488414in}{0.625278in}}%
\pgfpathlineto{\pgfqpoint{4.489584in}{0.633527in}}%
\pgfpathlineto{\pgfqpoint{4.489762in}{0.625278in}}%
\pgfpathlineto{\pgfqpoint{4.490591in}{0.625278in}}%
\pgfpathlineto{\pgfqpoint{4.491761in}{0.633527in}}%
\pgfpathlineto{\pgfqpoint{4.491769in}{0.625278in}}%
\pgfpathlineto{\pgfqpoint{4.492768in}{0.625278in}}%
\pgfpathlineto{\pgfqpoint{4.493945in}{0.633527in}}%
\pgfpathlineto{\pgfqpoint{4.494108in}{0.625278in}}%
\pgfpathlineto{\pgfqpoint{4.494945in}{0.625278in}}%
\pgfpathlineto{\pgfqpoint{4.496122in}{0.633527in}}%
\pgfpathlineto{\pgfqpoint{4.496130in}{0.625278in}}%
\pgfpathlineto{\pgfqpoint{4.497129in}{0.625278in}}%
\pgfpathlineto{\pgfqpoint{4.498469in}{0.633527in}}%
\pgfpathlineto{\pgfqpoint{4.498640in}{0.625278in}}%
\pgfpathlineto{\pgfqpoint{4.499476in}{0.625278in}}%
\pgfpathlineto{\pgfqpoint{4.500646in}{0.633527in}}%
\pgfpathlineto{\pgfqpoint{4.500987in}{0.625278in}}%
\pgfpathlineto{\pgfqpoint{4.501653in}{0.625278in}}%
\pgfpathlineto{\pgfqpoint{4.502660in}{0.633527in}}%
\pgfpathlineto{\pgfqpoint{4.502660in}{0.625278in}}%
\pgfpathlineto{\pgfqpoint{4.504000in}{0.633527in}}%
\pgfpathlineto{\pgfqpoint{4.504000in}{0.625278in}}%
\pgfpathlineto{\pgfqpoint{4.505007in}{0.625278in}}%
\pgfpathlineto{\pgfqpoint{4.506177in}{0.633527in}}%
\pgfpathlineto{\pgfqpoint{4.506518in}{0.625278in}}%
\pgfpathlineto{\pgfqpoint{4.507184in}{0.625278in}}%
\pgfpathlineto{\pgfqpoint{4.508361in}{0.633527in}}%
\pgfpathlineto{\pgfqpoint{4.508524in}{0.625278in}}%
\pgfpathlineto{\pgfqpoint{4.509368in}{0.625278in}}%
\pgfpathlineto{\pgfqpoint{4.510546in}{0.633527in}}%
\pgfpathlineto{\pgfqpoint{4.510708in}{0.625278in}}%
\pgfpathlineto{\pgfqpoint{4.511545in}{0.625278in}}%
\pgfpathlineto{\pgfqpoint{4.512715in}{0.633527in}}%
\pgfpathlineto{\pgfqpoint{4.512885in}{0.625278in}}%
\pgfpathlineto{\pgfqpoint{4.513722in}{0.625278in}}%
\pgfpathlineto{\pgfqpoint{4.514729in}{0.633527in}}%
\pgfpathlineto{\pgfqpoint{4.514729in}{0.625278in}}%
\pgfpathlineto{\pgfqpoint{4.516025in}{0.633527in}}%
\pgfpathlineto{\pgfqpoint{4.516069in}{0.625278in}}%
\pgfpathlineto{\pgfqpoint{4.517076in}{0.625278in}}%
\pgfpathlineto{\pgfqpoint{4.518246in}{0.633527in}}%
\pgfpathlineto{\pgfqpoint{4.518416in}{0.625278in}}%
\pgfpathlineto{\pgfqpoint{4.519253in}{0.625278in}}%
\pgfpathlineto{\pgfqpoint{4.520430in}{0.633527in}}%
\pgfpathlineto{\pgfqpoint{4.520593in}{0.625278in}}%
\pgfpathlineto{\pgfqpoint{4.521267in}{0.625278in}}%
\pgfpathlineto{\pgfqpoint{4.522607in}{0.633527in}}%
\pgfpathlineto{\pgfqpoint{4.522607in}{0.625278in}}%
\pgfpathlineto{\pgfqpoint{4.523614in}{0.625278in}}%
\pgfpathlineto{\pgfqpoint{4.524954in}{0.633527in}}%
\pgfpathlineto{\pgfqpoint{4.525124in}{0.625278in}}%
\pgfpathlineto{\pgfqpoint{4.525791in}{0.625278in}}%
\pgfpathlineto{\pgfqpoint{4.526968in}{0.633527in}}%
\pgfpathlineto{\pgfqpoint{4.527301in}{0.625278in}}%
\pgfpathlineto{\pgfqpoint{4.527975in}{0.625278in}}%
\pgfpathlineto{\pgfqpoint{4.529145in}{0.633527in}}%
\pgfpathlineto{\pgfqpoint{4.529315in}{0.625278in}}%
\pgfpathlineto{\pgfqpoint{4.530211in}{0.625278in}}%
\pgfpathlineto{\pgfqpoint{4.531322in}{0.633527in}}%
\pgfpathlineto{\pgfqpoint{4.531322in}{0.625278in}}%
\pgfpathlineto{\pgfqpoint{4.532499in}{0.633527in}}%
\pgfpathlineto{\pgfqpoint{4.532662in}{0.625278in}}%
\pgfpathlineto{\pgfqpoint{4.533506in}{0.625278in}}%
\pgfpathlineto{\pgfqpoint{4.534676in}{0.633527in}}%
\pgfpathlineto{\pgfqpoint{4.534846in}{0.625278in}}%
\pgfpathlineto{\pgfqpoint{4.535683in}{0.625278in}}%
\pgfpathlineto{\pgfqpoint{4.536853in}{0.633527in}}%
\pgfpathlineto{\pgfqpoint{4.536853in}{0.625278in}}%
\pgfpathlineto{\pgfqpoint{4.537860in}{0.625278in}}%
\pgfpathlineto{\pgfqpoint{4.539037in}{0.633527in}}%
\pgfpathlineto{\pgfqpoint{4.539200in}{0.625278in}}%
\pgfpathlineto{\pgfqpoint{4.539874in}{0.625278in}}%
\pgfpathlineto{\pgfqpoint{4.541043in}{0.633527in}}%
\pgfpathlineto{\pgfqpoint{4.541043in}{0.625278in}}%
\pgfpathlineto{\pgfqpoint{4.542050in}{0.625278in}}%
\pgfpathlineto{\pgfqpoint{4.543228in}{0.633527in}}%
\pgfpathlineto{\pgfqpoint{4.543391in}{0.625278in}}%
\pgfpathlineto{\pgfqpoint{4.544235in}{0.625278in}}%
\pgfpathlineto{\pgfqpoint{4.545405in}{0.633527in}}%
\pgfpathlineto{\pgfqpoint{4.545575in}{0.625278in}}%
\pgfpathlineto{\pgfqpoint{4.546411in}{0.625278in}}%
\pgfpathlineto{\pgfqpoint{4.547581in}{0.633527in}}%
\pgfpathlineto{\pgfqpoint{4.547752in}{0.625278in}}%
\pgfpathlineto{\pgfqpoint{4.548588in}{0.625278in}}%
\pgfpathlineto{\pgfqpoint{4.549766in}{0.633527in}}%
\pgfpathlineto{\pgfqpoint{4.549928in}{0.625278in}}%
\pgfpathlineto{\pgfqpoint{4.550773in}{0.625278in}}%
\pgfpathlineto{\pgfqpoint{4.551942in}{0.633527in}}%
\pgfpathlineto{\pgfqpoint{4.552113in}{0.625278in}}%
\pgfpathlineto{\pgfqpoint{4.552949in}{0.625278in}}%
\pgfpathlineto{\pgfqpoint{4.554119in}{0.633527in}}%
\pgfpathlineto{\pgfqpoint{4.554178in}{0.625278in}}%
\pgfpathlineto{\pgfqpoint{4.555126in}{0.625278in}}%
\pgfpathlineto{\pgfqpoint{4.556303in}{0.633527in}}%
\pgfpathlineto{\pgfqpoint{4.556303in}{0.625278in}}%
\pgfpathlineto{\pgfqpoint{4.557310in}{0.625278in}}%
\pgfpathlineto{\pgfqpoint{4.558480in}{0.633527in}}%
\pgfpathlineto{\pgfqpoint{4.558651in}{0.625278in}}%
\pgfpathlineto{\pgfqpoint{4.559487in}{0.625278in}}%
\pgfpathlineto{\pgfqpoint{4.560657in}{0.633527in}}%
\pgfpathlineto{\pgfqpoint{4.560827in}{0.625278in}}%
\pgfpathlineto{\pgfqpoint{4.561664in}{0.625278in}}%
\pgfpathlineto{\pgfqpoint{4.562841in}{0.633527in}}%
\pgfpathlineto{\pgfqpoint{4.563175in}{0.625278in}}%
\pgfpathlineto{\pgfqpoint{4.563848in}{0.625278in}}%
\pgfpathlineto{\pgfqpoint{4.565189in}{0.633527in}}%
\pgfpathlineto{\pgfqpoint{4.565351in}{0.625278in}}%
\pgfpathlineto{\pgfqpoint{4.566025in}{0.625278in}}%
\pgfpathlineto{\pgfqpoint{4.567195in}{0.633527in}}%
\pgfpathlineto{\pgfqpoint{4.567365in}{0.625278in}}%
\pgfpathlineto{\pgfqpoint{4.568202in}{0.625278in}}%
\pgfpathlineto{\pgfqpoint{4.569379in}{0.633527in}}%
\pgfpathlineto{\pgfqpoint{4.569542in}{0.625278in}}%
\pgfpathlineto{\pgfqpoint{4.570216in}{0.625278in}}%
\pgfpathlineto{\pgfqpoint{4.571556in}{0.633527in}}%
\pgfpathlineto{\pgfqpoint{4.571556in}{0.625278in}}%
\pgfpathlineto{\pgfqpoint{4.572563in}{0.625278in}}%
\pgfpathlineto{\pgfqpoint{4.573733in}{0.633527in}}%
\pgfpathlineto{\pgfqpoint{4.573733in}{0.625278in}}%
\pgfpathlineto{\pgfqpoint{4.574740in}{0.625278in}}%
\pgfpathlineto{\pgfqpoint{4.576080in}{0.633527in}}%
\pgfpathlineto{\pgfqpoint{4.576080in}{0.625278in}}%
\pgfpathlineto{\pgfqpoint{4.577087in}{0.625278in}}%
\pgfpathlineto{\pgfqpoint{4.578264in}{0.633527in}}%
\pgfpathlineto{\pgfqpoint{4.578427in}{0.625278in}}%
\pgfpathlineto{\pgfqpoint{4.579101in}{0.625278in}}%
\pgfpathlineto{\pgfqpoint{4.580271in}{0.633527in}}%
\pgfpathlineto{\pgfqpoint{4.580441in}{0.625278in}}%
\pgfpathlineto{\pgfqpoint{4.581278in}{0.625278in}}%
\pgfpathlineto{\pgfqpoint{4.582455in}{0.633527in}}%
\pgfpathlineto{\pgfqpoint{4.582455in}{0.625278in}}%
\pgfpathlineto{\pgfqpoint{4.583292in}{0.625278in}}%
\pgfpathlineto{\pgfqpoint{4.584462in}{0.633527in}}%
\pgfpathlineto{\pgfqpoint{4.584632in}{0.625278in}}%
\pgfpathlineto{\pgfqpoint{4.585469in}{0.625278in}}%
\pgfpathlineto{\pgfqpoint{4.586809in}{0.633527in}}%
\pgfpathlineto{\pgfqpoint{4.586979in}{0.625278in}}%
\pgfpathlineto{\pgfqpoint{4.587816in}{0.625278in}}%
\pgfpathlineto{\pgfqpoint{4.588823in}{0.633527in}}%
\pgfpathlineto{\pgfqpoint{4.588823in}{0.625278in}}%
\pgfpathlineto{\pgfqpoint{4.589993in}{0.633527in}}%
\pgfpathlineto{\pgfqpoint{4.589993in}{0.625278in}}%
\pgfpathlineto{\pgfqpoint{4.591000in}{0.625278in}}%
\pgfpathlineto{\pgfqpoint{4.592177in}{0.633527in}}%
\pgfpathlineto{\pgfqpoint{4.592340in}{0.625278in}}%
\pgfpathlineto{\pgfqpoint{4.593184in}{0.625278in}}%
\pgfpathlineto{\pgfqpoint{4.594354in}{0.633527in}}%
\pgfpathlineto{\pgfqpoint{4.594524in}{0.625278in}}%
\pgfpathlineto{\pgfqpoint{4.595361in}{0.625278in}}%
\pgfpathlineto{\pgfqpoint{4.596368in}{0.633527in}}%
\pgfpathlineto{\pgfqpoint{4.596368in}{0.625278in}}%
\pgfpathlineto{\pgfqpoint{4.597708in}{0.633527in}}%
\pgfpathlineto{\pgfqpoint{4.597708in}{0.625278in}}%
\pgfpathlineto{\pgfqpoint{4.598715in}{0.625278in}}%
\pgfpathlineto{\pgfqpoint{4.599885in}{0.633527in}}%
\pgfpathlineto{\pgfqpoint{4.600055in}{0.625278in}}%
\pgfpathlineto{\pgfqpoint{4.600892in}{0.625278in}}%
\pgfpathlineto{\pgfqpoint{4.602061in}{0.633527in}}%
\pgfpathlineto{\pgfqpoint{4.602061in}{0.625278in}}%
\pgfpathlineto{\pgfqpoint{4.603068in}{0.625278in}}%
\pgfpathlineto{\pgfqpoint{4.604246in}{0.633527in}}%
\pgfpathlineto{\pgfqpoint{4.604409in}{0.625278in}}%
\pgfpathlineto{\pgfqpoint{4.605267in}{0.650025in}}%
\pgfpathlineto{\pgfqpoint{4.605586in}{0.633527in}}%
\pgfpathlineto{\pgfqpoint{4.605756in}{0.658274in}}%
\pgfpathlineto{\pgfqpoint{4.606089in}{0.658274in}}%
\pgfpathlineto{\pgfqpoint{4.606104in}{0.674773in}}%
\pgfpathlineto{\pgfqpoint{4.606445in}{0.650025in}}%
\pgfpathlineto{\pgfqpoint{4.607111in}{0.650025in}}%
\pgfpathlineto{\pgfqpoint{4.607259in}{0.658274in}}%
\pgfpathlineto{\pgfqpoint{4.607600in}{0.633527in}}%
\pgfpathlineto{\pgfqpoint{4.607785in}{0.641776in}}%
\pgfpathlineto{\pgfqpoint{4.607933in}{0.633527in}}%
\pgfpathlineto{\pgfqpoint{4.608436in}{0.666524in}}%
\pgfpathlineto{\pgfqpoint{4.608451in}{0.674773in}}%
\pgfpathlineto{\pgfqpoint{4.608792in}{0.650025in}}%
\pgfpathlineto{\pgfqpoint{4.609458in}{0.650025in}}%
\pgfpathlineto{\pgfqpoint{4.610798in}{0.625278in}}%
\pgfpathlineto{\pgfqpoint{4.611953in}{0.633527in}}%
\pgfpathlineto{\pgfqpoint{4.612983in}{0.625278in}}%
\pgfpathlineto{\pgfqpoint{4.614301in}{0.683022in}}%
\pgfpathlineto{\pgfqpoint{4.614634in}{0.658274in}}%
\pgfpathlineto{\pgfqpoint{4.615811in}{0.633527in}}%
\pgfpathlineto{\pgfqpoint{4.616648in}{0.658274in}}%
\pgfpathlineto{\pgfqpoint{4.616833in}{0.650025in}}%
\pgfpathlineto{\pgfqpoint{4.619165in}{0.732517in}}%
\pgfpathlineto{\pgfqpoint{4.619180in}{0.724268in}}%
\pgfpathlineto{\pgfqpoint{4.619831in}{0.707769in}}%
\pgfpathlineto{\pgfqpoint{4.620002in}{0.732517in}}%
\pgfpathlineto{\pgfqpoint{4.620017in}{0.724268in}}%
\pgfpathlineto{\pgfqpoint{4.620357in}{0.749015in}}%
\pgfpathlineto{\pgfqpoint{4.620861in}{0.724268in}}%
\pgfpathlineto{\pgfqpoint{4.621527in}{0.699520in}}%
\pgfpathlineto{\pgfqpoint{4.621697in}{0.724268in}}%
\pgfpathlineto{\pgfqpoint{4.622179in}{0.757264in}}%
\pgfpathlineto{\pgfqpoint{4.622852in}{0.732517in}}%
\pgfpathlineto{\pgfqpoint{4.622867in}{0.724268in}}%
\pgfpathlineto{\pgfqpoint{4.623356in}{0.782012in}}%
\pgfpathlineto{\pgfqpoint{4.623874in}{0.749015in}}%
\pgfpathlineto{\pgfqpoint{4.624859in}{0.757264in}}%
\pgfpathlineto{\pgfqpoint{4.624881in}{0.749015in}}%
\pgfpathlineto{\pgfqpoint{4.625555in}{0.724268in}}%
\pgfpathlineto{\pgfqpoint{4.625718in}{0.749015in}}%
\pgfpathlineto{\pgfqpoint{4.626369in}{0.757264in}}%
\pgfpathlineto{\pgfqpoint{4.626554in}{0.740766in}}%
\pgfpathlineto{\pgfqpoint{4.627376in}{0.732517in}}%
\pgfpathlineto{\pgfqpoint{4.626725in}{0.749015in}}%
\pgfpathlineto{\pgfqpoint{4.627399in}{0.749015in}}%
\pgfpathlineto{\pgfqpoint{4.628568in}{0.773763in}}%
\pgfpathlineto{\pgfqpoint{4.628739in}{0.765513in}}%
\pgfpathlineto{\pgfqpoint{4.628902in}{0.773763in}}%
\pgfpathlineto{\pgfqpoint{4.629909in}{0.749015in}}%
\pgfpathlineto{\pgfqpoint{4.631064in}{0.757264in}}%
\pgfpathlineto{\pgfqpoint{4.632093in}{0.749015in}}%
\pgfpathlineto{\pgfqpoint{4.633240in}{0.782012in}}%
\pgfpathlineto{\pgfqpoint{4.633433in}{0.765513in}}%
\pgfpathlineto{\pgfqpoint{4.633596in}{0.749015in}}%
\pgfpathlineto{\pgfqpoint{4.633936in}{0.798510in}}%
\pgfpathlineto{\pgfqpoint{4.634418in}{0.831507in}}%
\pgfpathlineto{\pgfqpoint{4.634936in}{0.790261in}}%
\pgfpathlineto{\pgfqpoint{4.635758in}{0.782012in}}%
\pgfpathlineto{\pgfqpoint{4.635106in}{0.798510in}}%
\pgfpathlineto{\pgfqpoint{4.635780in}{0.798510in}}%
\pgfpathlineto{\pgfqpoint{4.636261in}{0.831507in}}%
\pgfpathlineto{\pgfqpoint{4.636935in}{0.806759in}}%
\pgfpathlineto{\pgfqpoint{4.637957in}{0.798510in}}%
\pgfpathlineto{\pgfqpoint{4.637994in}{0.806759in}}%
\pgfpathlineto{\pgfqpoint{4.638631in}{0.798510in}}%
\pgfpathlineto{\pgfqpoint{4.638794in}{0.823258in}}%
\pgfpathlineto{\pgfqpoint{4.641289in}{0.905749in}}%
\pgfpathlineto{\pgfqpoint{4.642296in}{0.881002in}}%
\pgfpathlineto{\pgfqpoint{4.642318in}{0.897500in}}%
\pgfpathlineto{\pgfqpoint{4.643658in}{0.872753in}}%
\pgfpathlineto{\pgfqpoint{4.644325in}{0.897500in}}%
\pgfpathlineto{\pgfqpoint{4.644813in}{0.881002in}}%
\pgfpathlineto{\pgfqpoint{4.645317in}{0.905749in}}%
\pgfpathlineto{\pgfqpoint{4.645835in}{0.864503in}}%
\pgfpathlineto{\pgfqpoint{4.646168in}{0.848005in}}%
\pgfpathlineto{\pgfqpoint{4.646487in}{0.881002in}}%
\pgfpathlineto{\pgfqpoint{4.646509in}{0.872753in}}%
\pgfpathlineto{\pgfqpoint{4.648160in}{0.905749in}}%
\pgfpathlineto{\pgfqpoint{4.649004in}{0.881002in}}%
\pgfpathlineto{\pgfqpoint{4.649189in}{0.897500in}}%
\pgfpathlineto{\pgfqpoint{4.650359in}{0.922247in}}%
\pgfpathlineto{\pgfqpoint{4.650529in}{0.913998in}}%
\pgfpathlineto{\pgfqpoint{4.650700in}{0.922247in}}%
\pgfpathlineto{\pgfqpoint{4.651684in}{0.905749in}}%
\pgfpathlineto{\pgfqpoint{4.652040in}{0.946995in}}%
\pgfpathlineto{\pgfqpoint{4.652706in}{0.946995in}}%
\pgfpathlineto{\pgfqpoint{4.653713in}{0.996490in}}%
\pgfpathlineto{\pgfqpoint{4.654046in}{0.971742in}}%
\pgfpathlineto{\pgfqpoint{4.654365in}{0.955244in}}%
\pgfpathlineto{\pgfqpoint{4.654698in}{0.988241in}}%
\pgfpathlineto{\pgfqpoint{4.655201in}{1.004739in}}%
\pgfpathlineto{\pgfqpoint{4.655557in}{0.971742in}}%
\pgfpathlineto{\pgfqpoint{4.655727in}{0.971742in}}%
\pgfpathlineto{\pgfqpoint{4.656541in}{0.955244in}}%
\pgfpathlineto{\pgfqpoint{4.657386in}{1.004739in}}%
\pgfpathlineto{\pgfqpoint{4.657734in}{0.996490in}}%
\pgfpathlineto{\pgfqpoint{4.658052in}{1.029487in}}%
\pgfpathlineto{\pgfqpoint{4.658407in}{1.021237in}}%
\pgfpathlineto{\pgfqpoint{4.659392in}{1.054234in}}%
\pgfpathlineto{\pgfqpoint{4.658889in}{1.004739in}}%
\pgfpathlineto{\pgfqpoint{4.659585in}{1.037736in}}%
\pgfpathlineto{\pgfqpoint{4.660066in}{1.029487in}}%
\pgfpathlineto{\pgfqpoint{4.660236in}{1.054234in}}%
\pgfpathlineto{\pgfqpoint{4.660421in}{1.045985in}}%
\pgfpathlineto{\pgfqpoint{4.660732in}{1.054234in}}%
\pgfpathlineto{\pgfqpoint{4.661073in}{1.029487in}}%
\pgfpathlineto{\pgfqpoint{4.661258in}{1.037736in}}%
\pgfpathlineto{\pgfqpoint{4.661295in}{1.029487in}}%
\pgfpathlineto{\pgfqpoint{4.661910in}{1.078981in}}%
\pgfpathlineto{\pgfqpoint{4.662095in}{1.070732in}}%
\pgfpathlineto{\pgfqpoint{4.662598in}{1.095480in}}%
\pgfpathlineto{\pgfqpoint{4.663250in}{1.078981in}}%
\pgfpathlineto{\pgfqpoint{4.664109in}{1.045985in}}%
\pgfpathlineto{\pgfqpoint{4.664271in}{1.070732in}}%
\pgfpathlineto{\pgfqpoint{4.664590in}{1.078981in}}%
\pgfpathlineto{\pgfqpoint{4.664923in}{1.054234in}}%
\pgfpathlineto{\pgfqpoint{4.665116in}{1.062483in}}%
\pgfpathlineto{\pgfqpoint{4.666271in}{1.029487in}}%
\pgfpathlineto{\pgfqpoint{4.665427in}{1.078981in}}%
\pgfpathlineto{\pgfqpoint{4.666285in}{1.045985in}}%
\pgfpathlineto{\pgfqpoint{4.666937in}{1.078981in}}%
\pgfpathlineto{\pgfqpoint{4.667440in}{1.054234in}}%
\pgfpathlineto{\pgfqpoint{4.668462in}{1.045985in}}%
\pgfpathlineto{\pgfqpoint{4.669810in}{1.021237in}}%
\pgfpathlineto{\pgfqpoint{4.670306in}{1.070732in}}%
\pgfpathlineto{\pgfqpoint{4.671461in}{1.054234in}}%
\pgfpathlineto{\pgfqpoint{4.672490in}{1.045985in}}%
\pgfpathlineto{\pgfqpoint{4.672808in}{1.054234in}}%
\pgfpathlineto{\pgfqpoint{4.673142in}{1.029487in}}%
\pgfpathlineto{\pgfqpoint{4.673327in}{1.037736in}}%
\pgfpathlineto{\pgfqpoint{4.674001in}{1.021237in}}%
\pgfpathlineto{\pgfqpoint{4.674163in}{1.045985in}}%
\pgfpathlineto{\pgfqpoint{4.674482in}{1.029487in}}%
\pgfpathlineto{\pgfqpoint{4.674497in}{1.045985in}}%
\pgfpathlineto{\pgfqpoint{4.675170in}{1.021237in}}%
\pgfpathlineto{\pgfqpoint{4.675504in}{1.021237in}}%
\pgfpathlineto{\pgfqpoint{4.676659in}{1.029487in}}%
\pgfpathlineto{\pgfqpoint{4.677836in}{1.004739in}}%
\pgfpathlineto{\pgfqpoint{4.677851in}{1.021237in}}%
\pgfpathlineto{\pgfqpoint{4.678502in}{0.979992in}}%
\pgfpathlineto{\pgfqpoint{4.678858in}{1.021237in}}%
\pgfpathlineto{\pgfqpoint{4.680516in}{1.054234in}}%
\pgfpathlineto{\pgfqpoint{4.681538in}{1.021237in}}%
\pgfpathlineto{\pgfqpoint{4.681686in}{1.029487in}}%
\pgfpathlineto{\pgfqpoint{4.681708in}{1.045985in}}%
\pgfpathlineto{\pgfqpoint{4.682382in}{1.021237in}}%
\pgfpathlineto{\pgfqpoint{4.682715in}{1.045985in}}%
\pgfpathlineto{\pgfqpoint{4.683885in}{1.070732in}}%
\pgfpathlineto{\pgfqpoint{4.684055in}{1.062483in}}%
\pgfpathlineto{\pgfqpoint{4.684389in}{1.045985in}}%
\pgfpathlineto{\pgfqpoint{4.684707in}{1.078981in}}%
\pgfpathlineto{\pgfqpoint{4.684729in}{1.070732in}}%
\pgfpathlineto{\pgfqpoint{4.685040in}{1.078981in}}%
\pgfpathlineto{\pgfqpoint{4.685225in}{1.062483in}}%
\pgfpathlineto{\pgfqpoint{4.685899in}{1.021237in}}%
\pgfpathlineto{\pgfqpoint{4.686380in}{1.054234in}}%
\pgfpathlineto{\pgfqpoint{4.686551in}{1.078981in}}%
\pgfpathlineto{\pgfqpoint{4.687387in}{1.029487in}}%
\pgfpathlineto{\pgfqpoint{4.687410in}{1.045985in}}%
\pgfpathlineto{\pgfqpoint{4.689401in}{1.103729in}}%
\pgfpathlineto{\pgfqpoint{4.690068in}{1.078981in}}%
\pgfpathlineto{\pgfqpoint{4.690423in}{1.095480in}}%
\pgfpathlineto{\pgfqpoint{4.691578in}{1.103729in}}%
\pgfpathlineto{\pgfqpoint{4.691934in}{1.120227in}}%
\pgfpathlineto{\pgfqpoint{4.692770in}{1.070732in}}%
\pgfpathlineto{\pgfqpoint{4.694096in}{1.054234in}}%
\pgfpathlineto{\pgfqpoint{4.694110in}{1.070732in}}%
\pgfpathlineto{\pgfqpoint{4.695117in}{1.070732in}}%
\pgfpathlineto{\pgfqpoint{4.696102in}{1.054234in}}%
\pgfpathlineto{\pgfqpoint{4.696124in}{1.070732in}}%
\pgfpathlineto{\pgfqpoint{4.696798in}{1.144975in}}%
\pgfpathlineto{\pgfqpoint{4.697968in}{1.111978in}}%
\pgfpathlineto{\pgfqpoint{4.699145in}{1.095480in}}%
\pgfpathlineto{\pgfqpoint{4.699293in}{1.103729in}}%
\pgfpathlineto{\pgfqpoint{4.699626in}{1.078981in}}%
\pgfpathlineto{\pgfqpoint{4.699982in}{1.095480in}}%
\pgfpathlineto{\pgfqpoint{4.700293in}{1.078981in}}%
\pgfpathlineto{\pgfqpoint{4.700633in}{1.111978in}}%
\pgfpathlineto{\pgfqpoint{4.700967in}{1.128476in}}%
\pgfpathlineto{\pgfqpoint{4.701322in}{1.095480in}}%
\pgfpathlineto{\pgfqpoint{4.701655in}{1.095480in}}%
\pgfpathlineto{\pgfqpoint{4.702981in}{1.128476in}}%
\pgfpathlineto{\pgfqpoint{4.702995in}{1.120227in}}%
\pgfpathlineto{\pgfqpoint{4.703314in}{1.103729in}}%
\pgfpathlineto{\pgfqpoint{4.703647in}{1.136726in}}%
\pgfpathlineto{\pgfqpoint{4.704839in}{1.194470in}}%
\pgfpathlineto{\pgfqpoint{4.706498in}{1.227466in}}%
\pgfpathlineto{\pgfqpoint{4.706890in}{1.202719in}}%
\pgfpathlineto{\pgfqpoint{4.707505in}{1.252214in}}%
\pgfpathlineto{\pgfqpoint{4.707527in}{1.243965in}}%
\pgfpathlineto{\pgfqpoint{4.707675in}{1.252214in}}%
\pgfpathlineto{\pgfqpoint{4.708023in}{1.210968in}}%
\pgfpathlineto{\pgfqpoint{4.708512in}{1.202719in}}%
\pgfpathlineto{\pgfqpoint{4.708512in}{1.235715in}}%
\pgfpathlineto{\pgfqpoint{4.708867in}{1.243965in}}%
\pgfpathlineto{\pgfqpoint{4.709533in}{1.194470in}}%
\pgfpathlineto{\pgfqpoint{4.709852in}{1.177971in}}%
\pgfpathlineto{\pgfqpoint{4.710185in}{1.210968in}}%
\pgfpathlineto{\pgfqpoint{4.711214in}{1.268712in}}%
\pgfpathlineto{\pgfqpoint{4.712029in}{1.301709in}}%
\pgfpathlineto{\pgfqpoint{4.712214in}{1.285210in}}%
\pgfpathlineto{\pgfqpoint{4.712887in}{1.219217in}}%
\pgfpathlineto{\pgfqpoint{4.713561in}{1.243965in}}%
\pgfpathlineto{\pgfqpoint{4.714213in}{1.252214in}}%
\pgfpathlineto{\pgfqpoint{4.714398in}{1.235715in}}%
\pgfpathlineto{\pgfqpoint{4.714546in}{1.227466in}}%
\pgfpathlineto{\pgfqpoint{4.714901in}{1.268712in}}%
\pgfpathlineto{\pgfqpoint{4.715049in}{1.276961in}}%
\pgfpathlineto{\pgfqpoint{4.715405in}{1.243965in}}%
\pgfpathlineto{\pgfqpoint{4.715716in}{1.252214in}}%
\pgfpathlineto{\pgfqpoint{4.716745in}{1.243965in}}%
\pgfpathlineto{\pgfqpoint{4.717900in}{1.252214in}}%
\pgfpathlineto{\pgfqpoint{4.719092in}{1.194470in}}%
\pgfpathlineto{\pgfqpoint{4.720751in}{1.227466in}}%
\pgfpathlineto{\pgfqpoint{4.721099in}{1.243965in}}%
\pgfpathlineto{\pgfqpoint{4.721772in}{1.210968in}}%
\pgfpathlineto{\pgfqpoint{4.722594in}{1.202719in}}%
\pgfpathlineto{\pgfqpoint{4.721943in}{1.219217in}}%
\pgfpathlineto{\pgfqpoint{4.722609in}{1.219217in}}%
\pgfpathlineto{\pgfqpoint{4.724268in}{1.252214in}}%
\pgfpathlineto{\pgfqpoint{4.725460in}{1.219217in}}%
\pgfpathlineto{\pgfqpoint{4.725963in}{1.194470in}}%
\pgfpathlineto{\pgfqpoint{4.726282in}{1.227466in}}%
\pgfpathlineto{\pgfqpoint{4.726296in}{1.219217in}}%
\pgfpathlineto{\pgfqpoint{4.727807in}{1.268712in}}%
\pgfpathlineto{\pgfqpoint{4.726948in}{1.202719in}}%
\pgfpathlineto{\pgfqpoint{4.727977in}{1.260463in}}%
\pgfpathlineto{\pgfqpoint{4.728792in}{1.252214in}}%
\pgfpathlineto{\pgfqpoint{4.728140in}{1.293460in}}%
\pgfpathlineto{\pgfqpoint{4.728977in}{1.260463in}}%
\pgfpathlineto{\pgfqpoint{4.729132in}{1.252214in}}%
\pgfpathlineto{\pgfqpoint{4.729147in}{1.293460in}}%
\pgfpathlineto{\pgfqpoint{4.729465in}{1.285210in}}%
\pgfpathlineto{\pgfqpoint{4.729799in}{1.301709in}}%
\pgfpathlineto{\pgfqpoint{4.730154in}{1.268712in}}%
\pgfpathlineto{\pgfqpoint{4.730487in}{1.268712in}}%
\pgfpathlineto{\pgfqpoint{4.730806in}{1.252214in}}%
\pgfpathlineto{\pgfqpoint{4.730976in}{1.276961in}}%
\pgfpathlineto{\pgfqpoint{4.731324in}{1.268712in}}%
\pgfpathlineto{\pgfqpoint{4.731642in}{1.276961in}}%
\pgfpathlineto{\pgfqpoint{4.731975in}{1.252214in}}%
\pgfpathlineto{\pgfqpoint{4.732982in}{1.227466in}}%
\pgfpathlineto{\pgfqpoint{4.733005in}{1.243965in}}%
\pgfpathlineto{\pgfqpoint{4.733841in}{1.169722in}}%
\pgfpathlineto{\pgfqpoint{4.734848in}{1.194470in}}%
\pgfpathlineto{\pgfqpoint{4.736003in}{1.202719in}}%
\pgfpathlineto{\pgfqpoint{4.736862in}{1.194470in}}%
\pgfpathlineto{\pgfqpoint{4.736507in}{1.227466in}}%
\pgfpathlineto{\pgfqpoint{4.737025in}{1.219217in}}%
\pgfpathlineto{\pgfqpoint{4.738180in}{1.227466in}}%
\pgfpathlineto{\pgfqpoint{4.738365in}{1.219217in}}%
\pgfpathlineto{\pgfqpoint{4.738684in}{1.252214in}}%
\pgfpathlineto{\pgfqpoint{4.739209in}{1.243965in}}%
\pgfpathlineto{\pgfqpoint{4.739409in}{1.227466in}}%
\pgfpathlineto{\pgfqpoint{4.739543in}{1.268712in}}%
\pgfpathlineto{\pgfqpoint{4.739861in}{1.260463in}}%
\pgfpathlineto{\pgfqpoint{4.739876in}{1.293460in}}%
\pgfpathlineto{\pgfqpoint{4.740883in}{1.293460in}}%
\pgfpathlineto{\pgfqpoint{4.742541in}{1.252214in}}%
\pgfpathlineto{\pgfqpoint{4.742541in}{1.260463in}}%
\pgfpathlineto{\pgfqpoint{4.742897in}{1.318207in}}%
\pgfpathlineto{\pgfqpoint{4.743563in}{1.293460in}}%
\pgfpathlineto{\pgfqpoint{4.745407in}{1.243965in}}%
\pgfpathlineto{\pgfqpoint{4.746058in}{1.252214in}}%
\pgfpathlineto{\pgfqpoint{4.746243in}{1.235715in}}%
\pgfpathlineto{\pgfqpoint{4.746562in}{1.252214in}}%
\pgfpathlineto{\pgfqpoint{4.747421in}{1.219217in}}%
\pgfpathlineto{\pgfqpoint{4.747569in}{1.227466in}}%
\pgfpathlineto{\pgfqpoint{4.747924in}{1.194470in}}%
\pgfpathlineto{\pgfqpoint{4.748242in}{1.177971in}}%
\pgfpathlineto{\pgfqpoint{4.748583in}{1.210968in}}%
\pgfpathlineto{\pgfqpoint{4.749242in}{1.202719in}}%
\pgfpathlineto{\pgfqpoint{4.749597in}{1.243965in}}%
\pgfpathlineto{\pgfqpoint{4.749923in}{1.252214in}}%
\pgfpathlineto{\pgfqpoint{4.750256in}{1.227466in}}%
\pgfpathlineto{\pgfqpoint{4.750434in}{1.235715in}}%
\pgfpathlineto{\pgfqpoint{4.750582in}{1.227466in}}%
\pgfpathlineto{\pgfqpoint{4.750604in}{1.268712in}}%
\pgfpathlineto{\pgfqpoint{4.750923in}{1.260463in}}%
\pgfpathlineto{\pgfqpoint{4.751945in}{1.318207in}}%
\pgfpathlineto{\pgfqpoint{4.752433in}{1.326456in}}%
\pgfpathlineto{\pgfqpoint{4.752618in}{1.309958in}}%
\pgfpathlineto{\pgfqpoint{4.752766in}{1.301709in}}%
\pgfpathlineto{\pgfqpoint{4.753122in}{1.342955in}}%
\pgfpathlineto{\pgfqpoint{4.754277in}{1.375951in}}%
\pgfpathlineto{\pgfqpoint{4.754462in}{1.359453in}}%
\pgfpathlineto{\pgfqpoint{4.755469in}{1.318207in}}%
\pgfpathlineto{\pgfqpoint{4.755632in}{1.342955in}}%
\pgfpathlineto{\pgfqpoint{4.756972in}{1.392449in}}%
\pgfpathlineto{\pgfqpoint{4.757290in}{1.375951in}}%
\pgfpathlineto{\pgfqpoint{4.758149in}{1.342955in}}%
\pgfpathlineto{\pgfqpoint{4.758312in}{1.367702in}}%
\pgfpathlineto{\pgfqpoint{4.761814in}{1.202719in}}%
\pgfpathlineto{\pgfqpoint{4.762318in}{1.227466in}}%
\pgfpathlineto{\pgfqpoint{4.762340in}{1.243965in}}%
\pgfpathlineto{\pgfqpoint{4.763177in}{1.169722in}}%
\pgfpathlineto{\pgfqpoint{4.763347in}{1.169722in}}%
\pgfpathlineto{\pgfqpoint{4.764332in}{1.202719in}}%
\pgfpathlineto{\pgfqpoint{4.764517in}{1.186221in}}%
\pgfpathlineto{\pgfqpoint{4.766198in}{1.095480in}}%
\pgfpathlineto{\pgfqpoint{4.766849in}{1.103729in}}%
\pgfpathlineto{\pgfqpoint{4.767034in}{1.087231in}}%
\pgfpathlineto{\pgfqpoint{4.768204in}{1.045985in}}%
\pgfpathlineto{\pgfqpoint{4.768693in}{1.054234in}}%
\pgfpathlineto{\pgfqpoint{4.769359in}{1.029487in}}%
\pgfpathlineto{\pgfqpoint{4.769715in}{1.070732in}}%
\pgfpathlineto{\pgfqpoint{4.771706in}{0.955244in}}%
\pgfpathlineto{\pgfqpoint{4.772047in}{0.988241in}}%
\pgfpathlineto{\pgfqpoint{4.772062in}{1.021237in}}%
\pgfpathlineto{\pgfqpoint{4.772736in}{0.971742in}}%
\pgfpathlineto{\pgfqpoint{4.773069in}{0.971742in}}%
\pgfpathlineto{\pgfqpoint{4.773557in}{0.979992in}}%
\pgfpathlineto{\pgfqpoint{4.774409in}{0.946995in}}%
\pgfpathlineto{\pgfqpoint{4.774727in}{0.979992in}}%
\pgfpathlineto{\pgfqpoint{4.775075in}{0.938746in}}%
\pgfpathlineto{\pgfqpoint{4.776252in}{0.897500in}}%
\pgfpathlineto{\pgfqpoint{4.776415in}{0.905749in}}%
\pgfpathlineto{\pgfqpoint{4.776756in}{0.872753in}}%
\pgfpathlineto{\pgfqpoint{4.776911in}{0.881002in}}%
\pgfpathlineto{\pgfqpoint{4.778096in}{0.823258in}}%
\pgfpathlineto{\pgfqpoint{4.778763in}{0.831507in}}%
\pgfpathlineto{\pgfqpoint{4.778933in}{0.815008in}}%
\pgfpathlineto{\pgfqpoint{4.779940in}{0.773763in}}%
\pgfpathlineto{\pgfqpoint{4.780614in}{0.749015in}}%
\pgfpathlineto{\pgfqpoint{4.780776in}{0.773763in}}%
\pgfpathlineto{\pgfqpoint{4.781435in}{0.782012in}}%
\pgfpathlineto{\pgfqpoint{4.781613in}{0.765513in}}%
\pgfpathlineto{\pgfqpoint{4.782287in}{0.749015in}}%
\pgfpathlineto{\pgfqpoint{4.782457in}{0.773763in}}%
\pgfpathlineto{\pgfqpoint{4.782783in}{0.782012in}}%
\pgfpathlineto{\pgfqpoint{4.783124in}{0.749015in}}%
\pgfpathlineto{\pgfqpoint{4.783627in}{0.724268in}}%
\pgfpathlineto{\pgfqpoint{4.783953in}{0.757264in}}%
\pgfpathlineto{\pgfqpoint{4.783960in}{0.749015in}}%
\pgfpathlineto{\pgfqpoint{4.786307in}{0.848005in}}%
\pgfpathlineto{\pgfqpoint{4.786478in}{0.839756in}}%
\pgfpathlineto{\pgfqpoint{4.786811in}{0.823258in}}%
\pgfpathlineto{\pgfqpoint{4.787144in}{0.856254in}}%
\pgfpathlineto{\pgfqpoint{4.787151in}{0.848005in}}%
\pgfpathlineto{\pgfqpoint{4.787810in}{0.881002in}}%
\pgfpathlineto{\pgfqpoint{4.788151in}{0.848005in}}%
\pgfpathlineto{\pgfqpoint{4.789661in}{0.823258in}}%
\pgfpathlineto{\pgfqpoint{4.790817in}{0.856254in}}%
\pgfpathlineto{\pgfqpoint{4.791002in}{0.839756in}}%
\pgfpathlineto{\pgfqpoint{4.792349in}{0.806759in}}%
\pgfpathlineto{\pgfqpoint{4.792490in}{0.815008in}}%
\pgfpathlineto{\pgfqpoint{4.792512in}{0.823258in}}%
\pgfpathlineto{\pgfqpoint{4.793186in}{0.773763in}}%
\pgfpathlineto{\pgfqpoint{4.793497in}{0.782012in}}%
\pgfpathlineto{\pgfqpoint{4.794185in}{0.773763in}}%
\pgfpathlineto{\pgfqpoint{4.794504in}{0.806759in}}%
\pgfpathlineto{\pgfqpoint{4.794526in}{0.806759in}}%
\pgfpathlineto{\pgfqpoint{4.795696in}{0.897500in}}%
\pgfpathlineto{\pgfqpoint{4.796370in}{0.848005in}}%
\pgfpathlineto{\pgfqpoint{4.796540in}{0.831507in}}%
\pgfpathlineto{\pgfqpoint{4.796703in}{0.872753in}}%
\pgfpathlineto{\pgfqpoint{4.797021in}{0.864503in}}%
\pgfpathlineto{\pgfqpoint{4.798043in}{0.897500in}}%
\pgfpathlineto{\pgfqpoint{4.798709in}{0.905749in}}%
\pgfpathlineto{\pgfqpoint{4.798880in}{0.889251in}}%
\pgfpathlineto{\pgfqpoint{4.800220in}{0.872753in}}%
\pgfpathlineto{\pgfqpoint{4.800546in}{0.881002in}}%
\pgfpathlineto{\pgfqpoint{4.800731in}{0.856254in}}%
\pgfpathlineto{\pgfqpoint{4.801064in}{0.864503in}}%
\pgfpathlineto{\pgfqpoint{4.801064in}{0.856254in}}%
\pgfpathlineto{\pgfqpoint{4.802064in}{0.905749in}}%
\pgfpathlineto{\pgfqpoint{4.802071in}{0.897500in}}%
\pgfpathlineto{\pgfqpoint{4.804085in}{0.831507in}}%
\pgfpathlineto{\pgfqpoint{4.804233in}{0.839756in}}%
\pgfpathlineto{\pgfqpoint{4.804736in}{0.856254in}}%
\pgfpathlineto{\pgfqpoint{4.805255in}{0.790261in}}%
\pgfpathlineto{\pgfqpoint{4.807602in}{0.707769in}}%
\pgfpathlineto{\pgfqpoint{4.807765in}{0.716019in}}%
\pgfpathlineto{\pgfqpoint{4.808431in}{0.732517in}}%
\pgfpathlineto{\pgfqpoint{4.808105in}{0.699520in}}%
\pgfpathlineto{\pgfqpoint{4.808772in}{0.732517in}}%
\pgfpathlineto{\pgfqpoint{4.811793in}{0.625278in}}%
\pgfpathlineto{\pgfqpoint{4.812955in}{0.633527in}}%
\pgfpathlineto{\pgfqpoint{4.813296in}{0.625278in}}%
\pgfpathlineto{\pgfqpoint{4.813962in}{0.625278in}}%
\pgfpathlineto{\pgfqpoint{4.814466in}{0.633527in}}%
\pgfpathlineto{\pgfqpoint{4.814466in}{0.625278in}}%
\pgfpathlineto{\pgfqpoint{4.815643in}{0.633527in}}%
\pgfpathlineto{\pgfqpoint{4.815806in}{0.625278in}}%
\pgfpathlineto{\pgfqpoint{4.816642in}{0.625278in}}%
\pgfpathlineto{\pgfqpoint{4.817657in}{0.633527in}}%
\pgfpathlineto{\pgfqpoint{4.817657in}{0.625278in}}%
\pgfpathlineto{\pgfqpoint{4.818990in}{0.633527in}}%
\pgfpathlineto{\pgfqpoint{4.818997in}{0.625278in}}%
\pgfpathlineto{\pgfqpoint{4.819826in}{0.625278in}}%
\pgfpathlineto{\pgfqpoint{4.821159in}{0.633527in}}%
\pgfpathlineto{\pgfqpoint{4.821344in}{0.625278in}}%
\pgfpathlineto{\pgfqpoint{4.822166in}{0.625278in}}%
\pgfpathlineto{\pgfqpoint{4.823173in}{0.633527in}}%
\pgfpathlineto{\pgfqpoint{4.823173in}{0.625278in}}%
\pgfpathlineto{\pgfqpoint{4.824343in}{0.633527in}}%
\pgfpathlineto{\pgfqpoint{4.824528in}{0.625278in}}%
\pgfpathlineto{\pgfqpoint{4.825365in}{0.625278in}}%
\pgfpathlineto{\pgfqpoint{4.826527in}{0.633527in}}%
\pgfpathlineto{\pgfqpoint{4.826697in}{0.625278in}}%
\pgfpathlineto{\pgfqpoint{4.827534in}{0.625278in}}%
\pgfpathlineto{\pgfqpoint{4.828711in}{0.633527in}}%
\pgfpathlineto{\pgfqpoint{4.828874in}{0.625278in}}%
\pgfpathlineto{\pgfqpoint{4.829718in}{0.625278in}}%
\pgfpathlineto{\pgfqpoint{4.830888in}{0.633527in}}%
\pgfpathlineto{\pgfqpoint{4.830895in}{0.625278in}}%
\pgfpathlineto{\pgfqpoint{4.831895in}{0.625278in}}%
\pgfpathlineto{\pgfqpoint{4.833065in}{0.633527in}}%
\pgfpathlineto{\pgfqpoint{4.833235in}{0.625278in}}%
\pgfpathlineto{\pgfqpoint{4.834072in}{0.625278in}}%
\pgfpathlineto{\pgfqpoint{4.835079in}{0.633527in}}%
\pgfpathlineto{\pgfqpoint{4.835079in}{0.625278in}}%
\pgfpathlineto{\pgfqpoint{4.836249in}{0.633527in}}%
\pgfpathlineto{\pgfqpoint{4.836412in}{0.625278in}}%
\pgfpathlineto{\pgfqpoint{4.837248in}{0.625278in}}%
\pgfpathlineto{\pgfqpoint{4.838588in}{0.633527in}}%
\pgfpathlineto{\pgfqpoint{4.838588in}{0.625278in}}%
\pgfpathlineto{\pgfqpoint{4.839455in}{0.625278in}}%
\pgfpathlineto{\pgfqpoint{4.840602in}{0.633527in}}%
\pgfpathlineto{\pgfqpoint{4.840610in}{0.625278in}}%
\pgfpathlineto{\pgfqpoint{4.841461in}{0.625278in}}%
\pgfpathlineto{\pgfqpoint{4.842787in}{0.633527in}}%
\pgfpathlineto{\pgfqpoint{4.842957in}{0.625278in}}%
\pgfpathlineto{\pgfqpoint{4.843794in}{0.625278in}}%
\pgfpathlineto{\pgfqpoint{4.844971in}{0.633527in}}%
\pgfpathlineto{\pgfqpoint{4.844971in}{0.625278in}}%
\pgfpathlineto{\pgfqpoint{4.847140in}{0.633527in}}%
\pgfpathlineto{\pgfqpoint{4.847311in}{0.625278in}}%
\pgfpathlineto{\pgfqpoint{4.848147in}{0.625278in}}%
\pgfpathlineto{\pgfqpoint{4.849317in}{0.633527in}}%
\pgfpathlineto{\pgfqpoint{4.849658in}{0.625278in}}%
\pgfpathlineto{\pgfqpoint{4.850324in}{0.625278in}}%
\pgfpathlineto{\pgfqpoint{4.851501in}{0.633527in}}%
\pgfpathlineto{\pgfqpoint{4.851664in}{0.625278in}}%
\pgfpathlineto{\pgfqpoint{4.852508in}{0.625278in}}%
\pgfpathlineto{\pgfqpoint{4.853678in}{0.633527in}}%
\pgfpathlineto{\pgfqpoint{4.853848in}{0.625278in}}%
\pgfpathlineto{\pgfqpoint{4.854685in}{0.625278in}}%
\pgfpathlineto{\pgfqpoint{4.855855in}{0.633527in}}%
\pgfpathlineto{\pgfqpoint{4.856033in}{0.625278in}}%
\pgfpathlineto{\pgfqpoint{4.856862in}{0.625278in}}%
\pgfpathlineto{\pgfqpoint{4.858039in}{0.633527in}}%
\pgfpathlineto{\pgfqpoint{4.858217in}{0.625278in}}%
\pgfpathlineto{\pgfqpoint{4.858876in}{0.625278in}}%
\pgfpathlineto{\pgfqpoint{4.860053in}{0.633527in}}%
\pgfpathlineto{\pgfqpoint{4.860216in}{0.625278in}}%
\pgfpathlineto{\pgfqpoint{4.861053in}{0.625278in}}%
\pgfpathlineto{\pgfqpoint{4.862230in}{0.633527in}}%
\pgfpathlineto{\pgfqpoint{4.862393in}{0.625278in}}%
\pgfpathlineto{\pgfqpoint{4.863237in}{0.625278in}}%
\pgfpathlineto{\pgfqpoint{4.864407in}{0.633527in}}%
\pgfpathlineto{\pgfqpoint{4.864577in}{0.625278in}}%
\pgfpathlineto{\pgfqpoint{4.865414in}{0.625278in}}%
\pgfpathlineto{\pgfqpoint{4.866584in}{0.633527in}}%
\pgfpathlineto{\pgfqpoint{4.866754in}{0.625278in}}%
\pgfpathlineto{\pgfqpoint{4.867591in}{0.625278in}}%
\pgfpathlineto{\pgfqpoint{4.868783in}{0.633527in}}%
\pgfpathlineto{\pgfqpoint{4.868931in}{0.625278in}}%
\pgfpathlineto{\pgfqpoint{4.869775in}{0.625278in}}%
\pgfpathlineto{\pgfqpoint{4.870960in}{0.633527in}}%
\pgfpathlineto{\pgfqpoint{4.871122in}{0.625278in}}%
\pgfpathlineto{\pgfqpoint{4.871959in}{0.625278in}}%
\pgfpathlineto{\pgfqpoint{4.872959in}{0.633527in}}%
\pgfpathlineto{\pgfqpoint{4.872959in}{0.625278in}}%
\pgfpathlineto{\pgfqpoint{4.874129in}{0.633527in}}%
\pgfpathlineto{\pgfqpoint{4.874306in}{0.625278in}}%
\pgfpathlineto{\pgfqpoint{4.875136in}{0.625278in}}%
\pgfpathlineto{\pgfqpoint{4.876305in}{0.633527in}}%
\pgfpathlineto{\pgfqpoint{4.876476in}{0.625278in}}%
\pgfpathlineto{\pgfqpoint{4.877312in}{0.625278in}}%
\pgfpathlineto{\pgfqpoint{4.878319in}{0.633527in}}%
\pgfpathlineto{\pgfqpoint{4.878319in}{0.625278in}}%
\pgfpathlineto{\pgfqpoint{4.879497in}{0.633527in}}%
\pgfpathlineto{\pgfqpoint{4.879497in}{0.625278in}}%
\pgfpathlineto{\pgfqpoint{4.880333in}{0.625278in}}%
\pgfpathlineto{\pgfqpoint{4.881673in}{0.633527in}}%
\pgfpathlineto{\pgfqpoint{4.881844in}{0.625278in}}%
\pgfpathlineto{\pgfqpoint{4.882510in}{0.625278in}}%
\pgfpathlineto{\pgfqpoint{4.883850in}{0.633527in}}%
\pgfpathlineto{\pgfqpoint{4.884191in}{0.625278in}}%
\pgfpathlineto{\pgfqpoint{4.884857in}{0.625278in}}%
\pgfpathlineto{\pgfqpoint{4.886197in}{0.633527in}}%
\pgfpathlineto{\pgfqpoint{4.886538in}{0.625278in}}%
\pgfpathlineto{\pgfqpoint{4.887204in}{0.625278in}}%
\pgfpathlineto{\pgfqpoint{4.888382in}{0.633527in}}%
\pgfpathlineto{\pgfqpoint{4.888715in}{0.625278in}}%
\pgfpathlineto{\pgfqpoint{4.889381in}{0.625278in}}%
\pgfpathlineto{\pgfqpoint{4.890558in}{0.633527in}}%
\pgfpathlineto{\pgfqpoint{4.890892in}{0.625278in}}%
\pgfpathlineto{\pgfqpoint{4.891565in}{0.625278in}}%
\pgfpathlineto{\pgfqpoint{4.892735in}{0.633527in}}%
\pgfpathlineto{\pgfqpoint{4.892735in}{0.625278in}}%
\pgfpathlineto{\pgfqpoint{4.893742in}{0.625278in}}%
\pgfpathlineto{\pgfqpoint{4.894920in}{0.633527in}}%
\pgfpathlineto{\pgfqpoint{4.895082in}{0.625278in}}%
\pgfpathlineto{\pgfqpoint{4.895919in}{0.625278in}}%
\pgfpathlineto{\pgfqpoint{4.897096in}{0.633527in}}%
\pgfpathlineto{\pgfqpoint{4.897259in}{0.625278in}}%
\pgfpathlineto{\pgfqpoint{4.898103in}{0.625278in}}%
\pgfpathlineto{\pgfqpoint{4.899273in}{0.633527in}}%
\pgfpathlineto{\pgfqpoint{4.899443in}{0.625278in}}%
\pgfpathlineto{\pgfqpoint{4.900280in}{0.625278in}}%
\pgfpathlineto{\pgfqpoint{4.901450in}{0.633527in}}%
\pgfpathlineto{\pgfqpoint{4.901620in}{0.625278in}}%
\pgfpathlineto{\pgfqpoint{4.902457in}{0.625278in}}%
\pgfpathlineto{\pgfqpoint{4.903634in}{0.633527in}}%
\pgfpathlineto{\pgfqpoint{4.903797in}{0.625278in}}%
\pgfpathlineto{\pgfqpoint{4.904641in}{0.625278in}}%
\pgfpathlineto{\pgfqpoint{4.905811in}{0.633527in}}%
\pgfpathlineto{\pgfqpoint{4.906144in}{0.625278in}}%
\pgfpathlineto{\pgfqpoint{4.906648in}{0.625278in}}%
\pgfpathlineto{\pgfqpoint{4.907988in}{0.633527in}}%
\pgfpathlineto{\pgfqpoint{4.907988in}{0.625278in}}%
\pgfpathlineto{\pgfqpoint{4.908995in}{0.625278in}}%
\pgfpathlineto{\pgfqpoint{4.910002in}{0.633527in}}%
\pgfpathlineto{\pgfqpoint{4.910002in}{0.625278in}}%
\pgfpathlineto{\pgfqpoint{4.911179in}{0.633527in}}%
\pgfpathlineto{\pgfqpoint{4.911342in}{0.625278in}}%
\pgfpathlineto{\pgfqpoint{4.912179in}{0.625278in}}%
\pgfpathlineto{\pgfqpoint{4.913356in}{0.633527in}}%
\pgfpathlineto{\pgfqpoint{4.913526in}{0.625278in}}%
\pgfpathlineto{\pgfqpoint{4.914363in}{0.625278in}}%
\pgfpathlineto{\pgfqpoint{4.915533in}{0.633527in}}%
\pgfpathlineto{\pgfqpoint{4.915703in}{0.625278in}}%
\pgfpathlineto{\pgfqpoint{4.916540in}{0.625278in}}%
\pgfpathlineto{\pgfqpoint{4.917717in}{0.633527in}}%
\pgfpathlineto{\pgfqpoint{4.917880in}{0.625278in}}%
\pgfpathlineto{\pgfqpoint{4.918717in}{0.625278in}}%
\pgfpathlineto{\pgfqpoint{4.919894in}{0.633527in}}%
\pgfpathlineto{\pgfqpoint{4.920064in}{0.625278in}}%
\pgfpathlineto{\pgfqpoint{4.920901in}{0.625278in}}%
\pgfpathlineto{\pgfqpoint{4.921908in}{0.633527in}}%
\pgfpathlineto{\pgfqpoint{4.921908in}{0.625278in}}%
\pgfpathlineto{\pgfqpoint{4.923078in}{0.633527in}}%
\pgfpathlineto{\pgfqpoint{4.923248in}{0.625278in}}%
\pgfpathlineto{\pgfqpoint{4.923914in}{0.625278in}}%
\pgfpathlineto{\pgfqpoint{4.925255in}{0.633527in}}%
\pgfpathlineto{\pgfqpoint{4.925255in}{0.625278in}}%
\pgfpathlineto{\pgfqpoint{4.926261in}{0.625278in}}%
\pgfpathlineto{\pgfqpoint{4.927602in}{0.633527in}}%
\pgfpathlineto{\pgfqpoint{4.927602in}{0.625278in}}%
\pgfpathlineto{\pgfqpoint{4.928609in}{0.625278in}}%
\pgfpathlineto{\pgfqpoint{4.929786in}{0.633527in}}%
\pgfpathlineto{\pgfqpoint{4.929949in}{0.625278in}}%
\pgfpathlineto{\pgfqpoint{4.930785in}{0.625278in}}%
\pgfpathlineto{\pgfqpoint{4.931963in}{0.633527in}}%
\pgfpathlineto{\pgfqpoint{4.932133in}{0.625278in}}%
\pgfpathlineto{\pgfqpoint{4.932970in}{0.625278in}}%
\pgfpathlineto{\pgfqpoint{4.934140in}{0.633527in}}%
\pgfpathlineto{\pgfqpoint{4.934310in}{0.625278in}}%
\pgfpathlineto{\pgfqpoint{4.935147in}{0.625278in}}%
\pgfpathlineto{\pgfqpoint{4.936324in}{0.633527in}}%
\pgfpathlineto{\pgfqpoint{4.936487in}{0.625278in}}%
\pgfpathlineto{\pgfqpoint{4.937323in}{0.625278in}}%
\pgfpathlineto{\pgfqpoint{4.938501in}{0.633527in}}%
\pgfpathlineto{\pgfqpoint{4.938671in}{0.625278in}}%
\pgfpathlineto{\pgfqpoint{4.939337in}{0.625278in}}%
\pgfpathlineto{\pgfqpoint{4.940515in}{0.633527in}}%
\pgfpathlineto{\pgfqpoint{4.940515in}{0.625278in}}%
\pgfpathlineto{\pgfqpoint{4.941514in}{0.625278in}}%
\pgfpathlineto{\pgfqpoint{4.942521in}{0.633527in}}%
\pgfpathlineto{\pgfqpoint{4.942521in}{0.625278in}}%
\pgfpathlineto{\pgfqpoint{4.943698in}{0.633527in}}%
\pgfpathlineto{\pgfqpoint{4.943861in}{0.625278in}}%
\pgfpathlineto{\pgfqpoint{4.944705in}{0.625278in}}%
\pgfpathlineto{\pgfqpoint{4.945875in}{0.633527in}}%
\pgfpathlineto{\pgfqpoint{4.946045in}{0.625278in}}%
\pgfpathlineto{\pgfqpoint{4.946882in}{0.625278in}}%
\pgfpathlineto{\pgfqpoint{4.947889in}{0.633527in}}%
\pgfpathlineto{\pgfqpoint{4.947889in}{0.625278in}}%
\pgfpathlineto{\pgfqpoint{4.949059in}{0.633527in}}%
\pgfpathlineto{\pgfqpoint{4.949400in}{0.625278in}}%
\pgfpathlineto{\pgfqpoint{4.950066in}{0.625278in}}%
\pgfpathlineto{\pgfqpoint{4.951243in}{0.633527in}}%
\pgfpathlineto{\pgfqpoint{4.951406in}{0.625278in}}%
\pgfpathlineto{\pgfqpoint{4.952243in}{0.625278in}}%
\pgfpathlineto{\pgfqpoint{4.953420in}{0.633527in}}%
\pgfpathlineto{\pgfqpoint{4.953590in}{0.625278in}}%
\pgfpathlineto{\pgfqpoint{4.954427in}{0.625278in}}%
\pgfpathlineto{\pgfqpoint{4.955597in}{0.633527in}}%
\pgfpathlineto{\pgfqpoint{4.955767in}{0.625278in}}%
\pgfpathlineto{\pgfqpoint{4.956604in}{0.625278in}}%
\pgfpathlineto{\pgfqpoint{4.957781in}{0.633527in}}%
\pgfpathlineto{\pgfqpoint{4.957781in}{0.625278in}}%
\pgfpathlineto{\pgfqpoint{4.958781in}{0.625278in}}%
\pgfpathlineto{\pgfqpoint{4.959958in}{0.633527in}}%
\pgfpathlineto{\pgfqpoint{4.960121in}{0.625278in}}%
\pgfpathlineto{\pgfqpoint{4.960965in}{0.625278in}}%
\pgfpathlineto{\pgfqpoint{4.962135in}{0.633527in}}%
\pgfpathlineto{\pgfqpoint{4.962305in}{0.625278in}}%
\pgfpathlineto{\pgfqpoint{4.963142in}{0.625278in}}%
\pgfpathlineto{\pgfqpoint{4.964312in}{0.633527in}}%
\pgfpathlineto{\pgfqpoint{4.964482in}{0.625278in}}%
\pgfpathlineto{\pgfqpoint{4.965319in}{0.625278in}}%
\pgfpathlineto{\pgfqpoint{4.966496in}{0.633527in}}%
\pgfpathlineto{\pgfqpoint{4.966659in}{0.625278in}}%
\pgfpathlineto{\pgfqpoint{4.967503in}{0.625278in}}%
\pgfpathlineto{\pgfqpoint{4.968502in}{0.633527in}}%
\pgfpathlineto{\pgfqpoint{4.968502in}{0.625278in}}%
\pgfpathlineto{\pgfqpoint{4.969680in}{0.633527in}}%
\pgfpathlineto{\pgfqpoint{4.969702in}{0.625278in}}%
\pgfpathlineto{\pgfqpoint{4.970687in}{0.658274in}}%
\pgfpathlineto{\pgfqpoint{4.970701in}{0.650025in}}%
\pgfpathlineto{\pgfqpoint{4.972027in}{0.683022in}}%
\pgfpathlineto{\pgfqpoint{4.971353in}{0.633527in}}%
\pgfpathlineto{\pgfqpoint{4.972042in}{0.674773in}}%
\pgfpathlineto{\pgfqpoint{4.972545in}{0.650025in}}%
\pgfpathlineto{\pgfqpoint{4.972863in}{0.683022in}}%
\pgfpathlineto{\pgfqpoint{4.972886in}{0.674773in}}%
\pgfpathlineto{\pgfqpoint{4.974056in}{0.699520in}}%
\pgfpathlineto{\pgfqpoint{4.973537in}{0.658274in}}%
\pgfpathlineto{\pgfqpoint{4.974226in}{0.691271in}}%
\pgfpathlineto{\pgfqpoint{4.975566in}{0.650025in}}%
\pgfpathlineto{\pgfqpoint{4.974389in}{0.699520in}}%
\pgfpathlineto{\pgfqpoint{4.975714in}{0.666524in}}%
\pgfpathlineto{\pgfqpoint{4.976721in}{0.683022in}}%
\pgfpathlineto{\pgfqpoint{4.976736in}{0.674773in}}%
\pgfpathlineto{\pgfqpoint{4.978061in}{0.658274in}}%
\pgfpathlineto{\pgfqpoint{4.978083in}{0.674773in}}%
\pgfpathlineto{\pgfqpoint{4.978417in}{0.650025in}}%
\pgfpathlineto{\pgfqpoint{4.979083in}{0.650025in}}%
\pgfpathlineto{\pgfqpoint{4.980912in}{0.707769in}}%
\pgfpathlineto{\pgfqpoint{4.980927in}{0.699520in}}%
\pgfpathlineto{\pgfqpoint{4.982082in}{0.683022in}}%
\pgfpathlineto{\pgfqpoint{4.981304in}{0.707769in}}%
\pgfpathlineto{\pgfqpoint{4.982082in}{0.691271in}}%
\pgfpathlineto{\pgfqpoint{4.983089in}{0.707769in}}%
\pgfpathlineto{\pgfqpoint{4.983111in}{0.699520in}}%
\pgfpathlineto{\pgfqpoint{4.983777in}{0.724268in}}%
\pgfpathlineto{\pgfqpoint{4.984118in}{0.699520in}}%
\pgfpathlineto{\pgfqpoint{4.984429in}{0.683022in}}%
\pgfpathlineto{\pgfqpoint{4.984599in}{0.707769in}}%
\pgfpathlineto{\pgfqpoint{4.984955in}{0.699520in}}%
\pgfpathlineto{\pgfqpoint{4.985621in}{0.724268in}}%
\pgfpathlineto{\pgfqpoint{4.986110in}{0.707769in}}%
\pgfpathlineto{\pgfqpoint{4.986628in}{0.724268in}}%
\pgfpathlineto{\pgfqpoint{4.987131in}{0.699520in}}%
\pgfpathlineto{\pgfqpoint{4.988116in}{0.683022in}}%
\pgfpathlineto{\pgfqpoint{4.988138in}{0.699520in}}%
\pgfpathlineto{\pgfqpoint{4.988642in}{0.724268in}}%
\pgfpathlineto{\pgfqpoint{4.988975in}{0.699520in}}%
\pgfpathlineto{\pgfqpoint{4.990634in}{0.658274in}}%
\pgfpathlineto{\pgfqpoint{4.990634in}{0.666524in}}%
\pgfpathlineto{\pgfqpoint{4.990656in}{0.674773in}}%
\pgfpathlineto{\pgfqpoint{4.990989in}{0.650025in}}%
\pgfpathlineto{\pgfqpoint{4.991655in}{0.650025in}}%
\pgfpathlineto{\pgfqpoint{4.992662in}{0.674773in}}%
\pgfpathlineto{\pgfqpoint{4.992144in}{0.633527in}}%
\pgfpathlineto{\pgfqpoint{4.992833in}{0.666524in}}%
\pgfpathlineto{\pgfqpoint{4.993499in}{0.650025in}}%
\pgfpathlineto{\pgfqpoint{4.993669in}{0.674773in}}%
\pgfpathlineto{\pgfqpoint{4.994676in}{0.699520in}}%
\pgfpathlineto{\pgfqpoint{4.994847in}{0.691271in}}%
\pgfpathlineto{\pgfqpoint{4.995498in}{0.707769in}}%
\pgfpathlineto{\pgfqpoint{4.996016in}{0.674773in}}%
\pgfpathlineto{\pgfqpoint{4.996498in}{0.683022in}}%
\pgfpathlineto{\pgfqpoint{4.996690in}{0.666524in}}%
\pgfpathlineto{\pgfqpoint{4.997171in}{0.658274in}}%
\pgfpathlineto{\pgfqpoint{4.997186in}{0.699520in}}%
\pgfpathlineto{\pgfqpoint{4.997505in}{0.683022in}}%
\pgfpathlineto{\pgfqpoint{4.997527in}{0.699520in}}%
\pgfpathlineto{\pgfqpoint{4.998178in}{0.658274in}}%
\pgfpathlineto{\pgfqpoint{4.998534in}{0.699520in}}%
\pgfpathlineto{\pgfqpoint{5.000022in}{0.658274in}}%
\pgfpathlineto{\pgfqpoint{5.000022in}{0.666524in}}%
\pgfpathlineto{\pgfqpoint{5.000859in}{0.707769in}}%
\pgfpathlineto{\pgfqpoint{5.001044in}{0.699520in}}%
\pgfpathlineto{\pgfqpoint{5.001192in}{0.707769in}}%
\pgfpathlineto{\pgfqpoint{5.001547in}{0.674773in}}%
\pgfpathlineto{\pgfqpoint{5.001695in}{0.683022in}}%
\pgfpathlineto{\pgfqpoint{5.001718in}{0.674773in}}%
\pgfpathlineto{\pgfqpoint{5.002029in}{0.707769in}}%
\pgfpathlineto{\pgfqpoint{5.002725in}{0.699520in}}%
\pgfpathlineto{\pgfqpoint{5.003724in}{0.724268in}}%
\pgfpathlineto{\pgfqpoint{5.003894in}{0.716019in}}%
\pgfpathlineto{\pgfqpoint{5.004731in}{0.699520in}}%
\pgfpathlineto{\pgfqpoint{5.004065in}{0.724268in}}%
\pgfpathlineto{\pgfqpoint{5.004879in}{0.716019in}}%
\pgfpathlineto{\pgfqpoint{5.006056in}{0.757264in}}%
\pgfpathlineto{\pgfqpoint{5.005553in}{0.707769in}}%
\pgfpathlineto{\pgfqpoint{5.006071in}{0.749015in}}%
\pgfpathlineto{\pgfqpoint{5.006390in}{0.757264in}}%
\pgfpathlineto{\pgfqpoint{5.006723in}{0.732517in}}%
\pgfpathlineto{\pgfqpoint{5.006915in}{0.740766in}}%
\pgfpathlineto{\pgfqpoint{5.007619in}{0.732517in}}%
\pgfpathlineto{\pgfqpoint{5.007078in}{0.749015in}}%
\pgfpathlineto{\pgfqpoint{5.007730in}{0.740766in}}%
\pgfpathlineto{\pgfqpoint{5.008233in}{0.806759in}}%
\pgfpathlineto{\pgfqpoint{5.008759in}{0.773763in}}%
\pgfpathlineto{\pgfqpoint{5.009263in}{0.749015in}}%
\pgfpathlineto{\pgfqpoint{5.009573in}{0.782012in}}%
\pgfpathlineto{\pgfqpoint{5.009596in}{0.773763in}}%
\pgfpathlineto{\pgfqpoint{5.009914in}{0.806759in}}%
\pgfpathlineto{\pgfqpoint{5.010751in}{0.782012in}}%
\pgfpathlineto{\pgfqpoint{5.011758in}{0.757264in}}%
\pgfpathlineto{\pgfqpoint{5.011773in}{0.773763in}}%
\pgfpathlineto{\pgfqpoint{5.013764in}{0.856254in}}%
\pgfpathlineto{\pgfqpoint{5.013949in}{0.839756in}}%
\pgfpathlineto{\pgfqpoint{5.013994in}{0.831507in}}%
\pgfpathlineto{\pgfqpoint{5.014771in}{0.881002in}}%
\pgfpathlineto{\pgfqpoint{5.014793in}{0.872753in}}%
\pgfpathlineto{\pgfqpoint{5.014942in}{0.881002in}}%
\pgfpathlineto{\pgfqpoint{5.015297in}{0.848005in}}%
\pgfpathlineto{\pgfqpoint{5.016615in}{0.806759in}}%
\pgfpathlineto{\pgfqpoint{5.016637in}{0.823258in}}%
\pgfpathlineto{\pgfqpoint{5.017622in}{0.856254in}}%
\pgfpathlineto{\pgfqpoint{5.017807in}{0.839756in}}%
\pgfpathlineto{\pgfqpoint{5.017955in}{0.831507in}}%
\pgfpathlineto{\pgfqpoint{5.018310in}{0.872753in}}%
\pgfpathlineto{\pgfqpoint{5.018962in}{0.856254in}}%
\pgfpathlineto{\pgfqpoint{5.019969in}{0.905749in}}%
\pgfpathlineto{\pgfqpoint{5.020154in}{0.922247in}}%
\pgfpathlineto{\pgfqpoint{5.020991in}{0.897500in}}%
\pgfpathlineto{\pgfqpoint{5.022649in}{0.955244in}}%
\pgfpathlineto{\pgfqpoint{5.022834in}{0.938746in}}%
\pgfpathlineto{\pgfqpoint{5.023338in}{0.922247in}}%
\pgfpathlineto{\pgfqpoint{5.023508in}{0.946995in}}%
\pgfpathlineto{\pgfqpoint{5.024663in}{0.955244in}}%
\pgfpathlineto{\pgfqpoint{5.025670in}{0.930497in}}%
\pgfpathlineto{\pgfqpoint{5.025685in}{0.946995in}}%
\pgfpathlineto{\pgfqpoint{5.027010in}{0.979992in}}%
\pgfpathlineto{\pgfqpoint{5.027025in}{0.971742in}}%
\pgfpathlineto{\pgfqpoint{5.029372in}{0.897500in}}%
\pgfpathlineto{\pgfqpoint{5.030194in}{0.930497in}}%
\pgfpathlineto{\pgfqpoint{5.030712in}{0.913998in}}%
\pgfpathlineto{\pgfqpoint{5.030868in}{0.905749in}}%
\pgfpathlineto{\pgfqpoint{5.031031in}{0.930497in}}%
\pgfpathlineto{\pgfqpoint{5.031386in}{0.922247in}}%
\pgfpathlineto{\pgfqpoint{5.031868in}{0.905749in}}%
\pgfpathlineto{\pgfqpoint{5.032541in}{0.930497in}}%
\pgfpathlineto{\pgfqpoint{5.033060in}{0.922247in}}%
\pgfpathlineto{\pgfqpoint{5.033548in}{0.963493in}}%
\pgfpathlineto{\pgfqpoint{5.034718in}{1.004739in}}%
\pgfpathlineto{\pgfqpoint{5.035244in}{0.971742in}}%
\pgfpathlineto{\pgfqpoint{5.035747in}{0.996490in}}%
\pgfpathlineto{\pgfqpoint{5.035895in}{1.004739in}}%
\pgfpathlineto{\pgfqpoint{5.036251in}{0.971742in}}%
\pgfpathlineto{\pgfqpoint{5.036562in}{0.979992in}}%
\pgfpathlineto{\pgfqpoint{5.037087in}{0.971742in}}%
\pgfpathlineto{\pgfqpoint{5.037569in}{1.012988in}}%
\pgfpathlineto{\pgfqpoint{5.037591in}{1.021237in}}%
\pgfpathlineto{\pgfqpoint{5.037924in}{0.996490in}}%
\pgfpathlineto{\pgfqpoint{5.038598in}{0.996490in}}%
\pgfpathlineto{\pgfqpoint{5.039264in}{0.971742in}}%
\pgfpathlineto{\pgfqpoint{5.039435in}{0.996490in}}%
\pgfpathlineto{\pgfqpoint{5.039746in}{1.004739in}}%
\pgfpathlineto{\pgfqpoint{5.040101in}{0.971742in}}%
\pgfpathlineto{\pgfqpoint{5.040249in}{0.979992in}}%
\pgfpathlineto{\pgfqpoint{5.040271in}{0.971742in}}%
\pgfpathlineto{\pgfqpoint{5.040590in}{1.004739in}}%
\pgfpathlineto{\pgfqpoint{5.041278in}{0.996490in}}%
\pgfpathlineto{\pgfqpoint{5.042618in}{0.971742in}}%
\pgfpathlineto{\pgfqpoint{5.043122in}{1.021237in}}%
\pgfpathlineto{\pgfqpoint{5.043959in}{0.996490in}}%
\pgfpathlineto{\pgfqpoint{5.044610in}{0.955244in}}%
\pgfpathlineto{\pgfqpoint{5.044966in}{0.996490in}}%
\pgfpathlineto{\pgfqpoint{5.045973in}{1.021237in}}%
\pgfpathlineto{\pgfqpoint{5.046135in}{1.012988in}}%
\pgfpathlineto{\pgfqpoint{5.046809in}{0.996490in}}%
\pgfpathlineto{\pgfqpoint{5.046979in}{1.021237in}}%
\pgfpathlineto{\pgfqpoint{5.047290in}{1.004739in}}%
\pgfpathlineto{\pgfqpoint{5.048127in}{1.029487in}}%
\pgfpathlineto{\pgfqpoint{5.047646in}{0.996490in}}%
\pgfpathlineto{\pgfqpoint{5.048320in}{1.021237in}}%
\pgfpathlineto{\pgfqpoint{5.048823in}{1.045985in}}%
\pgfpathlineto{\pgfqpoint{5.049475in}{1.029487in}}%
\pgfpathlineto{\pgfqpoint{5.049490in}{1.021237in}}%
\pgfpathlineto{\pgfqpoint{5.049808in}{1.054234in}}%
\pgfpathlineto{\pgfqpoint{5.050496in}{1.045985in}}%
\pgfpathlineto{\pgfqpoint{5.050815in}{1.029487in}}%
\pgfpathlineto{\pgfqpoint{5.050978in}{1.054234in}}%
\pgfpathlineto{\pgfqpoint{5.051333in}{1.045985in}}%
\pgfpathlineto{\pgfqpoint{5.052659in}{1.078981in}}%
\pgfpathlineto{\pgfqpoint{5.052673in}{1.070732in}}%
\pgfpathlineto{\pgfqpoint{5.053177in}{1.045985in}}%
\pgfpathlineto{\pgfqpoint{5.053495in}{1.078981in}}%
\pgfpathlineto{\pgfqpoint{5.053517in}{1.070732in}}%
\pgfpathlineto{\pgfqpoint{5.053665in}{1.078981in}}%
\pgfpathlineto{\pgfqpoint{5.053999in}{1.054234in}}%
\pgfpathlineto{\pgfqpoint{5.054354in}{1.070732in}}%
\pgfpathlineto{\pgfqpoint{5.054835in}{1.078981in}}%
\pgfpathlineto{\pgfqpoint{5.055694in}{1.045985in}}%
\pgfpathlineto{\pgfqpoint{5.056346in}{1.078981in}}%
\pgfpathlineto{\pgfqpoint{5.056864in}{1.062483in}}%
\pgfpathlineto{\pgfqpoint{5.057686in}{1.054234in}}%
\pgfpathlineto{\pgfqpoint{5.057034in}{1.070732in}}%
\pgfpathlineto{\pgfqpoint{5.057708in}{1.070732in}}%
\pgfpathlineto{\pgfqpoint{5.059026in}{1.128476in}}%
\pgfpathlineto{\pgfqpoint{5.059552in}{1.111978in}}%
\pgfpathlineto{\pgfqpoint{5.059715in}{1.120227in}}%
\pgfpathlineto{\pgfqpoint{5.060722in}{1.095480in}}%
\pgfpathlineto{\pgfqpoint{5.061729in}{1.120227in}}%
\pgfpathlineto{\pgfqpoint{5.061899in}{1.111978in}}%
\pgfpathlineto{\pgfqpoint{5.062232in}{1.095480in}}%
\pgfpathlineto{\pgfqpoint{5.062551in}{1.128476in}}%
\pgfpathlineto{\pgfqpoint{5.062713in}{1.128476in}}%
\pgfpathlineto{\pgfqpoint{5.062736in}{1.144975in}}%
\pgfpathlineto{\pgfqpoint{5.063402in}{1.120227in}}%
\pgfpathlineto{\pgfqpoint{5.063743in}{1.120227in}}%
\pgfpathlineto{\pgfqpoint{5.064912in}{1.144975in}}%
\pgfpathlineto{\pgfqpoint{5.065083in}{1.136726in}}%
\pgfpathlineto{\pgfqpoint{5.065231in}{1.128476in}}%
\pgfpathlineto{\pgfqpoint{5.065394in}{1.153224in}}%
\pgfpathlineto{\pgfqpoint{5.065919in}{1.144975in}}%
\pgfpathlineto{\pgfqpoint{5.067074in}{1.153224in}}%
\pgfpathlineto{\pgfqpoint{5.068096in}{1.144975in}}%
\pgfpathlineto{\pgfqpoint{5.069140in}{1.153224in}}%
\pgfpathlineto{\pgfqpoint{5.070614in}{1.120227in}}%
\pgfpathlineto{\pgfqpoint{5.071095in}{1.128476in}}%
\pgfpathlineto{\pgfqpoint{5.071428in}{1.103729in}}%
\pgfpathlineto{\pgfqpoint{5.071769in}{1.128476in}}%
\pgfpathlineto{\pgfqpoint{5.072457in}{1.095480in}}%
\pgfpathlineto{\pgfqpoint{5.073109in}{1.078981in}}%
\pgfpathlineto{\pgfqpoint{5.073272in}{1.103729in}}%
\pgfpathlineto{\pgfqpoint{5.073294in}{1.095480in}}%
\pgfpathlineto{\pgfqpoint{5.074953in}{1.128476in}}%
\pgfpathlineto{\pgfqpoint{5.075974in}{1.095480in}}%
\pgfpathlineto{\pgfqpoint{5.077633in}{1.128476in}}%
\pgfpathlineto{\pgfqpoint{5.077655in}{1.120227in}}%
\pgfpathlineto{\pgfqpoint{5.078307in}{1.177971in}}%
\pgfpathlineto{\pgfqpoint{5.078662in}{1.136726in}}%
\pgfpathlineto{\pgfqpoint{5.078825in}{1.144975in}}%
\pgfpathlineto{\pgfqpoint{5.079832in}{1.120227in}}%
\pgfpathlineto{\pgfqpoint{5.079980in}{1.128476in}}%
\pgfpathlineto{\pgfqpoint{5.080335in}{1.087231in}}%
\pgfpathlineto{\pgfqpoint{5.081505in}{1.070732in}}%
\pgfpathlineto{\pgfqpoint{5.082512in}{1.095480in}}%
\pgfpathlineto{\pgfqpoint{5.082683in}{1.087231in}}%
\pgfpathlineto{\pgfqpoint{5.082853in}{1.095480in}}%
\pgfpathlineto{\pgfqpoint{5.083727in}{1.078981in}}%
\pgfpathlineto{\pgfqpoint{5.084859in}{1.169722in}}%
\pgfpathlineto{\pgfqpoint{5.085533in}{1.136726in}}%
\pgfpathlineto{\pgfqpoint{5.087021in}{1.103729in}}%
\pgfpathlineto{\pgfqpoint{5.085696in}{1.144975in}}%
\pgfpathlineto{\pgfqpoint{5.087021in}{1.111978in}}%
\pgfpathlineto{\pgfqpoint{5.088213in}{1.194470in}}%
\pgfpathlineto{\pgfqpoint{5.088532in}{1.177971in}}%
\pgfpathlineto{\pgfqpoint{5.088695in}{1.202719in}}%
\pgfpathlineto{\pgfqpoint{5.089050in}{1.194470in}}%
\pgfpathlineto{\pgfqpoint{5.090205in}{1.202719in}}%
\pgfpathlineto{\pgfqpoint{5.091234in}{1.194470in}}%
\pgfpathlineto{\pgfqpoint{5.090731in}{1.219217in}}%
\pgfpathlineto{\pgfqpoint{5.091271in}{1.202719in}}%
\pgfpathlineto{\pgfqpoint{5.091397in}{1.194470in}}%
\pgfpathlineto{\pgfqpoint{5.091716in}{1.227466in}}%
\pgfpathlineto{\pgfqpoint{5.092382in}{1.202719in}}%
\pgfpathlineto{\pgfqpoint{5.093396in}{1.227466in}}%
\pgfpathlineto{\pgfqpoint{5.093411in}{1.219217in}}%
\pgfpathlineto{\pgfqpoint{5.095077in}{1.252214in}}%
\pgfpathlineto{\pgfqpoint{5.095425in}{1.243965in}}%
\pgfpathlineto{\pgfqpoint{5.096077in}{1.301709in}}%
\pgfpathlineto{\pgfqpoint{5.096092in}{1.293460in}}%
\pgfpathlineto{\pgfqpoint{5.097750in}{1.326456in}}%
\pgfpathlineto{\pgfqpoint{5.098439in}{1.293460in}}%
\pgfpathlineto{\pgfqpoint{5.098772in}{1.318207in}}%
\pgfpathlineto{\pgfqpoint{5.100453in}{1.417197in}}%
\pgfpathlineto{\pgfqpoint{5.101267in}{1.400699in}}%
\pgfpathlineto{\pgfqpoint{5.102296in}{1.392449in}}%
\pgfpathlineto{\pgfqpoint{5.102800in}{1.367702in}}%
\pgfpathlineto{\pgfqpoint{5.103118in}{1.400699in}}%
\pgfpathlineto{\pgfqpoint{5.103133in}{1.392449in}}%
\pgfpathlineto{\pgfqpoint{5.104288in}{1.400699in}}%
\pgfpathlineto{\pgfqpoint{5.105310in}{1.367702in}}%
\pgfpathlineto{\pgfqpoint{5.105813in}{1.342955in}}%
\pgfpathlineto{\pgfqpoint{5.106132in}{1.375951in}}%
\pgfpathlineto{\pgfqpoint{5.106146in}{1.367702in}}%
\pgfpathlineto{\pgfqpoint{5.106317in}{1.417197in}}%
\pgfpathlineto{\pgfqpoint{5.107301in}{1.375951in}}%
\pgfpathlineto{\pgfqpoint{5.107494in}{1.367702in}}%
\pgfpathlineto{\pgfqpoint{5.107997in}{1.417197in}}%
\pgfpathlineto{\pgfqpoint{5.108331in}{1.384200in}}%
\pgfpathlineto{\pgfqpoint{5.109501in}{1.367702in}}%
\pgfpathlineto{\pgfqpoint{5.111996in}{1.450194in}}%
\pgfpathlineto{\pgfqpoint{5.113025in}{1.441944in}}%
\pgfpathlineto{\pgfqpoint{5.113358in}{1.417197in}}%
\pgfpathlineto{\pgfqpoint{5.113684in}{1.450194in}}%
\pgfpathlineto{\pgfqpoint{5.113862in}{1.441944in}}%
\pgfpathlineto{\pgfqpoint{5.114513in}{1.474941in}}%
\pgfpathlineto{\pgfqpoint{5.115031in}{1.458443in}}%
\pgfpathlineto{\pgfqpoint{5.117379in}{1.268712in}}%
\pgfpathlineto{\pgfqpoint{5.117882in}{1.293460in}}%
\pgfpathlineto{\pgfqpoint{5.118223in}{1.260463in}}%
\pgfpathlineto{\pgfqpoint{5.120733in}{1.169722in}}%
\pgfpathlineto{\pgfqpoint{5.121895in}{1.177971in}}%
\pgfpathlineto{\pgfqpoint{5.122391in}{1.153224in}}%
\pgfpathlineto{\pgfqpoint{5.122910in}{1.169722in}}%
\pgfpathlineto{\pgfqpoint{5.123228in}{1.177971in}}%
\pgfpathlineto{\pgfqpoint{5.123413in}{1.161473in}}%
\pgfpathlineto{\pgfqpoint{5.124087in}{1.120227in}}%
\pgfpathlineto{\pgfqpoint{5.124590in}{1.144975in}}%
\pgfpathlineto{\pgfqpoint{5.125412in}{1.103729in}}%
\pgfpathlineto{\pgfqpoint{5.125930in}{1.120227in}}%
\pgfpathlineto{\pgfqpoint{5.126078in}{1.128476in}}%
\pgfpathlineto{\pgfqpoint{5.126434in}{1.095480in}}%
\pgfpathlineto{\pgfqpoint{5.128278in}{1.045985in}}%
\pgfpathlineto{\pgfqpoint{5.129262in}{1.078981in}}%
\pgfpathlineto{\pgfqpoint{5.129447in}{1.062483in}}%
\pgfpathlineto{\pgfqpoint{5.129951in}{1.021237in}}%
\pgfpathlineto{\pgfqpoint{5.130625in}{1.045985in}}%
\pgfpathlineto{\pgfqpoint{5.130943in}{1.029487in}}%
\pgfpathlineto{\pgfqpoint{5.131284in}{1.062483in}}%
\pgfpathlineto{\pgfqpoint{5.131617in}{1.078981in}}%
\pgfpathlineto{\pgfqpoint{5.131965in}{1.045985in}}%
\pgfpathlineto{\pgfqpoint{5.132298in}{1.045985in}}%
\pgfpathlineto{\pgfqpoint{5.133127in}{1.054234in}}%
\pgfpathlineto{\pgfqpoint{5.133142in}{1.045985in}}%
\pgfpathlineto{\pgfqpoint{5.134793in}{0.979992in}}%
\pgfpathlineto{\pgfqpoint{5.134986in}{0.996490in}}%
\pgfpathlineto{\pgfqpoint{5.135141in}{1.004739in}}%
\pgfpathlineto{\pgfqpoint{5.135482in}{0.971742in}}%
\pgfpathlineto{\pgfqpoint{5.135985in}{0.946995in}}%
\pgfpathlineto{\pgfqpoint{5.136311in}{0.979992in}}%
\pgfpathlineto{\pgfqpoint{5.136326in}{0.971742in}}%
\pgfpathlineto{\pgfqpoint{5.136489in}{0.996490in}}%
\pgfpathlineto{\pgfqpoint{5.136992in}{0.946995in}}%
\pgfpathlineto{\pgfqpoint{5.137140in}{0.955244in}}%
\pgfpathlineto{\pgfqpoint{5.137481in}{0.979992in}}%
\pgfpathlineto{\pgfqpoint{5.138170in}{0.946995in}}%
\pgfpathlineto{\pgfqpoint{5.140332in}{0.881002in}}%
\pgfpathlineto{\pgfqpoint{5.140332in}{0.889251in}}%
\pgfpathlineto{\pgfqpoint{5.140672in}{0.881002in}}%
\pgfpathlineto{\pgfqpoint{5.141524in}{0.946995in}}%
\pgfpathlineto{\pgfqpoint{5.141679in}{0.955244in}}%
\pgfpathlineto{\pgfqpoint{5.142020in}{0.922247in}}%
\pgfpathlineto{\pgfqpoint{5.142175in}{0.930497in}}%
\pgfpathlineto{\pgfqpoint{5.142864in}{0.946995in}}%
\pgfpathlineto{\pgfqpoint{5.143197in}{0.922247in}}%
\pgfpathlineto{\pgfqpoint{5.143863in}{0.872753in}}%
\pgfpathlineto{\pgfqpoint{5.144537in}{0.897500in}}%
\pgfpathlineto{\pgfqpoint{5.145018in}{0.905749in}}%
\pgfpathlineto{\pgfqpoint{5.145211in}{0.889251in}}%
\pgfpathlineto{\pgfqpoint{5.145359in}{0.881002in}}%
\pgfpathlineto{\pgfqpoint{5.145714in}{0.922247in}}%
\pgfpathlineto{\pgfqpoint{5.146699in}{0.955244in}}%
\pgfpathlineto{\pgfqpoint{5.146884in}{0.938746in}}%
\pgfpathlineto{\pgfqpoint{5.147032in}{0.930497in}}%
\pgfpathlineto{\pgfqpoint{5.147210in}{0.955244in}}%
\pgfpathlineto{\pgfqpoint{5.147558in}{0.946995in}}%
\pgfpathlineto{\pgfqpoint{5.148395in}{0.971742in}}%
\pgfpathlineto{\pgfqpoint{5.148713in}{0.955244in}}%
\pgfpathlineto{\pgfqpoint{5.149054in}{0.930497in}}%
\pgfpathlineto{\pgfqpoint{5.149557in}{0.979992in}}%
\pgfpathlineto{\pgfqpoint{5.149735in}{0.963493in}}%
\pgfpathlineto{\pgfqpoint{5.150238in}{0.946995in}}%
\pgfpathlineto{\pgfqpoint{5.150401in}{0.971742in}}%
\pgfpathlineto{\pgfqpoint{5.150897in}{0.979992in}}%
\pgfpathlineto{\pgfqpoint{5.151075in}{0.963493in}}%
\pgfpathlineto{\pgfqpoint{5.154592in}{0.823258in}}%
\pgfpathlineto{\pgfqpoint{5.155592in}{0.856254in}}%
\pgfpathlineto{\pgfqpoint{5.155769in}{0.839756in}}%
\pgfpathlineto{\pgfqpoint{5.156939in}{0.823258in}}%
\pgfpathlineto{\pgfqpoint{5.158613in}{0.881002in}}%
\pgfpathlineto{\pgfqpoint{5.158783in}{0.864503in}}%
\pgfpathlineto{\pgfqpoint{5.159627in}{0.831507in}}%
\pgfpathlineto{\pgfqpoint{5.159790in}{0.872753in}}%
\pgfpathlineto{\pgfqpoint{5.163647in}{0.683022in}}%
\pgfpathlineto{\pgfqpoint{5.163818in}{0.724268in}}%
\pgfpathlineto{\pgfqpoint{5.163973in}{0.732517in}}%
\pgfpathlineto{\pgfqpoint{5.164321in}{0.699520in}}%
\pgfpathlineto{\pgfqpoint{5.164817in}{0.674773in}}%
\pgfpathlineto{\pgfqpoint{5.165136in}{0.707769in}}%
\pgfpathlineto{\pgfqpoint{5.165158in}{0.699520in}}%
\pgfpathlineto{\pgfqpoint{5.165661in}{0.724268in}}%
\pgfpathlineto{\pgfqpoint{5.165995in}{0.699520in}}%
\pgfpathlineto{\pgfqpoint{5.166320in}{0.683022in}}%
\pgfpathlineto{\pgfqpoint{5.166328in}{0.724268in}}%
\pgfpathlineto{\pgfqpoint{5.166661in}{0.716019in}}%
\pgfpathlineto{\pgfqpoint{5.167823in}{0.782012in}}%
\pgfpathlineto{\pgfqpoint{5.168342in}{0.749015in}}%
\pgfpathlineto{\pgfqpoint{5.168845in}{0.782012in}}%
\pgfpathlineto{\pgfqpoint{5.169008in}{0.798510in}}%
\pgfpathlineto{\pgfqpoint{5.169349in}{0.765513in}}%
\pgfpathlineto{\pgfqpoint{5.171192in}{0.650025in}}%
\pgfpathlineto{\pgfqpoint{5.171674in}{0.658274in}}%
\pgfpathlineto{\pgfqpoint{5.171866in}{0.633527in}}%
\pgfpathlineto{\pgfqpoint{5.172029in}{0.625278in}}%
\pgfpathlineto{\pgfqpoint{5.172199in}{0.650025in}}%
\pgfpathlineto{\pgfqpoint{5.172532in}{0.650025in}}%
\pgfpathlineto{\pgfqpoint{5.172532in}{0.658274in}}%
\pgfpathlineto{\pgfqpoint{5.172703in}{0.658274in}}%
\pgfusepath{stroke}%
\end{pgfscope}%
\begin{pgfscope}%
\pgfpathrectangle{\pgfqpoint{0.586112in}{0.502778in}}{\pgfqpoint{4.805000in}{2.695000in}}%
\pgfusepath{clip}%
\pgfsetbuttcap%
\pgfsetroundjoin%
\pgfsetlinewidth{1.505625pt}%
\definecolor{currentstroke}{rgb}{0.650000,0.650000,0.650000}%
\pgfsetstrokecolor{currentstroke}%
\pgfsetdash{{1.500000pt}{2.475000pt}}{0.000000pt}%
\pgfpathmoveto{\pgfqpoint{0.804521in}{0.633527in}}%
\pgfpathlineto{\pgfqpoint{0.804521in}{0.625278in}}%
\pgfpathlineto{\pgfqpoint{0.804521in}{0.625278in}}%
\pgfpathlineto{\pgfqpoint{0.806399in}{0.633527in}}%
\pgfpathlineto{\pgfqpoint{0.806416in}{0.625278in}}%
\pgfpathlineto{\pgfqpoint{0.807508in}{0.625278in}}%
\pgfpathlineto{\pgfqpoint{0.808565in}{0.633527in}}%
\pgfpathlineto{\pgfqpoint{0.808565in}{0.625278in}}%
\pgfpathlineto{\pgfqpoint{0.809751in}{0.633527in}}%
\pgfpathlineto{\pgfqpoint{0.809769in}{0.625278in}}%
\pgfpathlineto{\pgfqpoint{0.810738in}{0.625278in}}%
\pgfpathlineto{\pgfqpoint{0.810906in}{0.633527in}}%
\pgfpathlineto{\pgfqpoint{0.810906in}{0.625278in}}%
\pgfpathlineto{\pgfqpoint{0.812415in}{0.633527in}}%
\pgfpathlineto{\pgfqpoint{0.812583in}{0.625278in}}%
\pgfpathlineto{\pgfqpoint{0.813421in}{0.625278in}}%
\pgfpathlineto{\pgfqpoint{0.813923in}{0.691271in}}%
\pgfpathlineto{\pgfqpoint{0.814762in}{0.881002in}}%
\pgfpathlineto{\pgfqpoint{0.814948in}{0.872753in}}%
\pgfpathlineto{\pgfqpoint{0.815767in}{0.782012in}}%
\pgfpathlineto{\pgfqpoint{0.815954in}{0.773763in}}%
\pgfpathlineto{\pgfqpoint{0.816605in}{0.889251in}}%
\pgfpathlineto{\pgfqpoint{0.817630in}{1.318207in}}%
\pgfpathlineto{\pgfqpoint{0.818282in}{1.375951in}}%
\pgfpathlineto{\pgfqpoint{0.818636in}{1.334705in}}%
\pgfpathlineto{\pgfqpoint{0.819306in}{1.268712in}}%
\pgfpathlineto{\pgfqpoint{0.819791in}{1.326456in}}%
\pgfpathlineto{\pgfqpoint{0.820815in}{1.763662in}}%
\pgfpathlineto{\pgfqpoint{0.821299in}{1.846153in}}%
\pgfpathlineto{\pgfqpoint{0.821467in}{1.870901in}}%
\pgfpathlineto{\pgfqpoint{0.822324in}{1.780160in}}%
\pgfpathlineto{\pgfqpoint{0.822659in}{1.763662in}}%
\pgfpathlineto{\pgfqpoint{0.822975in}{1.804907in}}%
\pgfpathlineto{\pgfqpoint{0.824987in}{2.365850in}}%
\pgfpathlineto{\pgfqpoint{0.825173in}{2.349352in}}%
\pgfpathlineto{\pgfqpoint{0.826012in}{2.258611in}}%
\pgfpathlineto{\pgfqpoint{0.826496in}{2.316355in}}%
\pgfpathlineto{\pgfqpoint{0.827521in}{2.654571in}}%
\pgfpathlineto{\pgfqpoint{0.828172in}{2.687567in}}%
\pgfpathlineto{\pgfqpoint{0.828359in}{2.671069in}}%
\pgfpathlineto{\pgfqpoint{0.829680in}{2.588577in}}%
\pgfpathlineto{\pgfqpoint{0.829867in}{2.605076in}}%
\pgfpathlineto{\pgfqpoint{0.832381in}{2.704066in}}%
\pgfpathlineto{\pgfqpoint{0.832549in}{2.695816in}}%
\pgfpathlineto{\pgfqpoint{0.832717in}{2.704066in}}%
\pgfpathlineto{\pgfqpoint{0.833890in}{2.572079in}}%
\pgfpathlineto{\pgfqpoint{0.847785in}{0.633527in}}%
\pgfpathlineto{\pgfqpoint{0.847804in}{0.650025in}}%
\pgfpathlineto{\pgfqpoint{0.848958in}{0.658274in}}%
\pgfpathlineto{\pgfqpoint{0.849983in}{0.650025in}}%
\pgfpathlineto{\pgfqpoint{0.851138in}{0.658274in}}%
\pgfpathlineto{\pgfqpoint{0.852143in}{0.633527in}}%
\pgfpathlineto{\pgfqpoint{0.852162in}{0.650025in}}%
\pgfpathlineto{\pgfqpoint{0.853484in}{0.633527in}}%
\pgfpathlineto{\pgfqpoint{0.853503in}{0.650025in}}%
\pgfpathlineto{\pgfqpoint{0.854509in}{0.650025in}}%
\pgfpathlineto{\pgfqpoint{0.855663in}{0.658274in}}%
\pgfpathlineto{\pgfqpoint{0.856688in}{0.641776in}}%
\pgfpathlineto{\pgfqpoint{0.856837in}{0.633527in}}%
\pgfpathlineto{\pgfqpoint{0.857005in}{0.658274in}}%
\pgfpathlineto{\pgfqpoint{0.857359in}{0.650025in}}%
\pgfpathlineto{\pgfqpoint{0.857843in}{0.658274in}}%
\pgfpathlineto{\pgfqpoint{0.858030in}{0.641776in}}%
\pgfpathlineto{\pgfqpoint{0.859202in}{0.625278in}}%
\pgfpathlineto{\pgfqpoint{0.859854in}{0.658274in}}%
\pgfpathlineto{\pgfqpoint{0.860364in}{0.633527in}}%
\pgfpathlineto{\pgfqpoint{0.861047in}{0.625278in}}%
\pgfpathlineto{\pgfqpoint{0.861363in}{0.658274in}}%
\pgfpathlineto{\pgfqpoint{0.861382in}{0.650025in}}%
\pgfpathlineto{\pgfqpoint{0.862220in}{0.625278in}}%
\pgfpathlineto{\pgfqpoint{0.862388in}{0.650025in}}%
\pgfpathlineto{\pgfqpoint{0.863058in}{0.674773in}}%
\pgfpathlineto{\pgfqpoint{0.863542in}{0.658274in}}%
\pgfpathlineto{\pgfqpoint{0.863729in}{0.674773in}}%
\pgfpathlineto{\pgfqpoint{0.864569in}{0.633527in}}%
\pgfpathlineto{\pgfqpoint{0.865908in}{0.674773in}}%
\pgfpathlineto{\pgfqpoint{0.866076in}{0.666524in}}%
\pgfpathlineto{\pgfqpoint{0.866914in}{0.650025in}}%
\pgfpathlineto{\pgfqpoint{0.866243in}{0.674773in}}%
\pgfpathlineto{\pgfqpoint{0.867077in}{0.666524in}}%
\pgfpathlineto{\pgfqpoint{0.867919in}{0.674773in}}%
\pgfpathlineto{\pgfqpoint{0.867417in}{0.650025in}}%
\pgfpathlineto{\pgfqpoint{0.868090in}{0.658274in}}%
\pgfpathlineto{\pgfqpoint{0.869431in}{0.683022in}}%
\pgfpathlineto{\pgfqpoint{0.870098in}{0.650025in}}%
\pgfpathlineto{\pgfqpoint{0.870602in}{0.674773in}}%
\pgfpathlineto{\pgfqpoint{0.871105in}{0.650025in}}%
\pgfpathlineto{\pgfqpoint{0.871610in}{0.683022in}}%
\pgfpathlineto{\pgfqpoint{0.872110in}{0.674773in}}%
\pgfpathlineto{\pgfqpoint{0.872278in}{0.699520in}}%
\pgfpathlineto{\pgfqpoint{0.873286in}{0.806759in}}%
\pgfpathlineto{\pgfqpoint{0.873451in}{0.798510in}}%
\pgfpathlineto{\pgfqpoint{0.873454in}{0.839756in}}%
\pgfpathlineto{\pgfqpoint{0.874962in}{0.881002in}}%
\pgfpathlineto{\pgfqpoint{0.875127in}{0.872753in}}%
\pgfpathlineto{\pgfqpoint{0.875798in}{0.848005in}}%
\pgfpathlineto{\pgfqpoint{0.875968in}{0.881002in}}%
\pgfpathlineto{\pgfqpoint{0.876134in}{0.872753in}}%
\pgfpathlineto{\pgfqpoint{0.877141in}{0.905749in}}%
\pgfpathlineto{\pgfqpoint{0.877474in}{0.897500in}}%
\pgfpathlineto{\pgfqpoint{0.877642in}{0.922247in}}%
\pgfpathlineto{\pgfqpoint{0.877812in}{0.913998in}}%
\pgfpathlineto{\pgfqpoint{0.878313in}{0.897500in}}%
\pgfpathlineto{\pgfqpoint{0.878986in}{0.930497in}}%
\pgfpathlineto{\pgfqpoint{0.879822in}{0.897500in}}%
\pgfpathlineto{\pgfqpoint{0.880159in}{0.913998in}}%
\pgfpathlineto{\pgfqpoint{0.880327in}{0.930497in}}%
\pgfpathlineto{\pgfqpoint{0.880827in}{0.897500in}}%
\pgfpathlineto{\pgfqpoint{0.880829in}{0.905749in}}%
\pgfpathlineto{\pgfqpoint{0.882003in}{0.856254in}}%
\pgfpathlineto{\pgfqpoint{0.882503in}{0.872753in}}%
\pgfpathlineto{\pgfqpoint{0.883509in}{0.922247in}}%
\pgfpathlineto{\pgfqpoint{0.884012in}{0.897500in}}%
\pgfpathlineto{\pgfqpoint{0.884182in}{0.930497in}}%
\pgfpathlineto{\pgfqpoint{0.884350in}{0.930497in}}%
\pgfpathlineto{\pgfqpoint{0.885355in}{0.979992in}}%
\pgfpathlineto{\pgfqpoint{0.885689in}{0.971742in}}%
\pgfpathlineto{\pgfqpoint{0.885856in}{0.996490in}}%
\pgfpathlineto{\pgfqpoint{0.886026in}{0.988241in}}%
\pgfpathlineto{\pgfqpoint{0.886861in}{0.996490in}}%
\pgfpathlineto{\pgfqpoint{0.886361in}{0.979992in}}%
\pgfpathlineto{\pgfqpoint{0.887030in}{0.988241in}}%
\pgfpathlineto{\pgfqpoint{0.887535in}{1.004739in}}%
\pgfpathlineto{\pgfqpoint{0.888037in}{0.979992in}}%
\pgfpathlineto{\pgfqpoint{0.890885in}{0.790261in}}%
\pgfpathlineto{\pgfqpoint{0.891890in}{0.724268in}}%
\pgfpathlineto{\pgfqpoint{0.891893in}{0.732517in}}%
\pgfpathlineto{\pgfqpoint{0.893904in}{0.625278in}}%
\pgfpathlineto{\pgfqpoint{0.895078in}{0.633527in}}%
\pgfpathlineto{\pgfqpoint{0.895413in}{0.625278in}}%
\pgfpathlineto{\pgfqpoint{0.896754in}{0.633527in}}%
\pgfpathlineto{\pgfqpoint{0.896922in}{0.625278in}}%
\pgfpathlineto{\pgfqpoint{0.897760in}{0.625278in}}%
\pgfpathlineto{\pgfqpoint{0.898766in}{0.633527in}}%
\pgfpathlineto{\pgfqpoint{0.898766in}{0.625278in}}%
\pgfpathlineto{\pgfqpoint{0.899101in}{0.633527in}}%
\pgfpathlineto{\pgfqpoint{0.899101in}{0.625278in}}%
\pgfpathlineto{\pgfqpoint{0.900275in}{0.633527in}}%
\pgfpathlineto{\pgfqpoint{0.900610in}{0.625278in}}%
\pgfpathlineto{\pgfqpoint{0.901951in}{0.633527in}}%
\pgfpathlineto{\pgfqpoint{0.902119in}{0.625278in}}%
\pgfpathlineto{\pgfqpoint{0.903459in}{0.633527in}}%
\pgfpathlineto{\pgfqpoint{0.903795in}{0.625278in}}%
\pgfpathlineto{\pgfqpoint{0.904968in}{0.633527in}}%
\pgfpathlineto{\pgfqpoint{0.905304in}{0.625278in}}%
\pgfpathlineto{\pgfqpoint{0.907483in}{0.633527in}}%
\pgfpathlineto{\pgfqpoint{0.907986in}{0.625278in}}%
\pgfpathlineto{\pgfqpoint{0.909159in}{0.633527in}}%
\pgfpathlineto{\pgfqpoint{0.909327in}{0.625278in}}%
\pgfpathlineto{\pgfqpoint{0.910165in}{0.625278in}}%
\pgfpathlineto{\pgfqpoint{0.911338in}{0.633527in}}%
\pgfpathlineto{\pgfqpoint{0.911506in}{0.625278in}}%
\pgfpathlineto{\pgfqpoint{0.912176in}{0.625278in}}%
\pgfpathlineto{\pgfqpoint{0.913517in}{0.633527in}}%
\pgfpathlineto{\pgfqpoint{0.913853in}{0.625278in}}%
\pgfpathlineto{\pgfqpoint{0.914524in}{0.625278in}}%
\pgfpathlineto{\pgfqpoint{0.915696in}{0.633527in}}%
\pgfpathlineto{\pgfqpoint{0.915864in}{0.625278in}}%
\pgfpathlineto{\pgfqpoint{0.916703in}{0.625278in}}%
\pgfpathlineto{\pgfqpoint{0.917876in}{0.633527in}}%
\pgfpathlineto{\pgfqpoint{0.918211in}{0.625278in}}%
\pgfpathlineto{\pgfqpoint{0.918882in}{0.625278in}}%
\pgfpathlineto{\pgfqpoint{0.920055in}{0.633527in}}%
\pgfpathlineto{\pgfqpoint{0.920391in}{0.625278in}}%
\pgfpathlineto{\pgfqpoint{0.921054in}{0.625278in}}%
\pgfpathlineto{\pgfqpoint{0.922402in}{0.633527in}}%
\pgfpathlineto{\pgfqpoint{0.922562in}{0.625278in}}%
\pgfpathlineto{\pgfqpoint{0.923408in}{0.625278in}}%
\pgfpathlineto{\pgfqpoint{0.924574in}{0.633527in}}%
\pgfpathlineto{\pgfqpoint{0.924749in}{0.625278in}}%
\pgfpathlineto{\pgfqpoint{0.925412in}{0.625278in}}%
\pgfpathlineto{\pgfqpoint{0.926586in}{0.633527in}}%
\pgfpathlineto{\pgfqpoint{0.926921in}{0.625278in}}%
\pgfpathlineto{\pgfqpoint{0.927599in}{0.625278in}}%
\pgfpathlineto{\pgfqpoint{0.928765in}{0.633527in}}%
\pgfpathlineto{\pgfqpoint{0.928932in}{0.625278in}}%
\pgfpathlineto{\pgfqpoint{0.929778in}{0.625278in}}%
\pgfpathlineto{\pgfqpoint{0.930783in}{0.633527in}}%
\pgfpathlineto{\pgfqpoint{0.930783in}{0.625278in}}%
\pgfpathlineto{\pgfqpoint{0.932125in}{0.633527in}}%
\pgfpathlineto{\pgfqpoint{0.932125in}{0.625278in}}%
\pgfpathlineto{\pgfqpoint{0.933130in}{0.625278in}}%
\pgfpathlineto{\pgfqpoint{0.934304in}{0.633527in}}%
\pgfpathlineto{\pgfqpoint{0.934304in}{0.625278in}}%
\pgfpathlineto{\pgfqpoint{0.935309in}{0.625278in}}%
\pgfpathlineto{\pgfqpoint{0.936644in}{0.633527in}}%
\pgfpathlineto{\pgfqpoint{0.936650in}{0.625278in}}%
\pgfpathlineto{\pgfqpoint{0.937584in}{0.625278in}}%
\pgfpathlineto{\pgfqpoint{0.938830in}{0.633527in}}%
\pgfpathlineto{\pgfqpoint{0.939165in}{0.625278in}}%
\pgfpathlineto{\pgfqpoint{0.939828in}{0.625278in}}%
\pgfpathlineto{\pgfqpoint{0.941002in}{0.633527in}}%
\pgfpathlineto{\pgfqpoint{0.941009in}{0.625278in}}%
\pgfpathlineto{\pgfqpoint{0.942015in}{0.625278in}}%
\pgfpathlineto{\pgfqpoint{0.943181in}{0.633527in}}%
\pgfpathlineto{\pgfqpoint{0.943349in}{0.625278in}}%
\pgfpathlineto{\pgfqpoint{0.944194in}{0.625278in}}%
\pgfpathlineto{\pgfqpoint{0.945367in}{0.633527in}}%
\pgfpathlineto{\pgfqpoint{0.945535in}{0.625278in}}%
\pgfpathlineto{\pgfqpoint{0.946373in}{0.625278in}}%
\pgfpathlineto{\pgfqpoint{0.947372in}{0.633527in}}%
\pgfpathlineto{\pgfqpoint{0.947372in}{0.625278in}}%
\pgfpathlineto{\pgfqpoint{0.948713in}{0.633527in}}%
\pgfpathlineto{\pgfqpoint{0.948720in}{0.625278in}}%
\pgfpathlineto{\pgfqpoint{0.949726in}{0.625278in}}%
\pgfpathlineto{\pgfqpoint{0.950899in}{0.633527in}}%
\pgfpathlineto{\pgfqpoint{0.950899in}{0.625278in}}%
\pgfpathlineto{\pgfqpoint{0.951898in}{0.625278in}}%
\pgfpathlineto{\pgfqpoint{0.953071in}{0.633527in}}%
\pgfpathlineto{\pgfqpoint{0.953407in}{0.625278in}}%
\pgfpathlineto{\pgfqpoint{0.954084in}{0.625278in}}%
\pgfpathlineto{\pgfqpoint{0.955258in}{0.633527in}}%
\pgfpathlineto{\pgfqpoint{0.955404in}{0.625278in}}%
\pgfpathlineto{\pgfqpoint{0.956263in}{0.625278in}}%
\pgfpathlineto{\pgfqpoint{0.957415in}{0.633527in}}%
\pgfpathlineto{\pgfqpoint{0.957437in}{0.625278in}}%
\pgfpathlineto{\pgfqpoint{0.958442in}{0.625278in}}%
\pgfpathlineto{\pgfqpoint{0.959784in}{0.633527in}}%
\pgfpathlineto{\pgfqpoint{0.959900in}{0.625278in}}%
\pgfpathlineto{\pgfqpoint{0.960775in}{0.625278in}}%
\pgfpathlineto{\pgfqpoint{0.961941in}{0.633527in}}%
\pgfpathlineto{\pgfqpoint{0.962116in}{0.625278in}}%
\pgfpathlineto{\pgfqpoint{0.962947in}{0.625278in}}%
\pgfpathlineto{\pgfqpoint{0.964121in}{0.633527in}}%
\pgfpathlineto{\pgfqpoint{0.964303in}{0.625278in}}%
\pgfpathlineto{\pgfqpoint{0.965127in}{0.625278in}}%
\pgfpathlineto{\pgfqpoint{0.966321in}{0.633527in}}%
\pgfpathlineto{\pgfqpoint{0.966468in}{0.625278in}}%
\pgfpathlineto{\pgfqpoint{0.967327in}{0.625278in}}%
\pgfpathlineto{\pgfqpoint{0.968500in}{0.633527in}}%
\pgfpathlineto{\pgfqpoint{0.968647in}{0.625278in}}%
\pgfpathlineto{\pgfqpoint{0.969485in}{0.625278in}}%
\pgfpathlineto{\pgfqpoint{0.970658in}{0.633527in}}%
\pgfpathlineto{\pgfqpoint{0.970658in}{0.625278in}}%
\pgfpathlineto{\pgfqpoint{0.971664in}{0.625278in}}%
\pgfpathlineto{\pgfqpoint{0.972837in}{0.633527in}}%
\pgfpathlineto{\pgfqpoint{0.973020in}{0.625278in}}%
\pgfpathlineto{\pgfqpoint{0.973843in}{0.625278in}}%
\pgfpathlineto{\pgfqpoint{0.975017in}{0.633527in}}%
\pgfpathlineto{\pgfqpoint{0.975205in}{0.625278in}}%
\pgfpathlineto{\pgfqpoint{0.976044in}{0.625278in}}%
\pgfpathlineto{\pgfqpoint{0.977196in}{0.633527in}}%
\pgfpathlineto{\pgfqpoint{0.977385in}{0.625278in}}%
\pgfpathlineto{\pgfqpoint{0.978223in}{0.625278in}}%
\pgfpathlineto{\pgfqpoint{0.979543in}{0.633527in}}%
\pgfpathlineto{\pgfqpoint{0.979543in}{0.625278in}}%
\pgfpathlineto{\pgfqpoint{0.980570in}{0.625278in}}%
\pgfpathlineto{\pgfqpoint{0.981722in}{0.633527in}}%
\pgfpathlineto{\pgfqpoint{0.981890in}{0.625278in}}%
\pgfpathlineto{\pgfqpoint{0.982742in}{0.625278in}}%
\pgfpathlineto{\pgfqpoint{0.983901in}{0.633527in}}%
\pgfpathlineto{\pgfqpoint{0.983901in}{0.625278in}}%
\pgfpathlineto{\pgfqpoint{0.984907in}{0.625278in}}%
\pgfpathlineto{\pgfqpoint{0.986095in}{0.633527in}}%
\pgfpathlineto{\pgfqpoint{0.986101in}{0.625278in}}%
\pgfpathlineto{\pgfqpoint{0.987100in}{0.625278in}}%
\pgfpathlineto{\pgfqpoint{0.988260in}{0.633527in}}%
\pgfpathlineto{\pgfqpoint{0.988427in}{0.625278in}}%
\pgfpathlineto{\pgfqpoint{0.989119in}{0.625278in}}%
\pgfpathlineto{\pgfqpoint{0.990271in}{0.633527in}}%
\pgfpathlineto{\pgfqpoint{0.990439in}{0.625278in}}%
\pgfpathlineto{\pgfqpoint{0.991298in}{0.625278in}}%
\pgfpathlineto{\pgfqpoint{0.992450in}{0.633527in}}%
\pgfpathlineto{\pgfqpoint{0.992639in}{0.625278in}}%
\pgfpathlineto{\pgfqpoint{0.993310in}{0.625278in}}%
\pgfpathlineto{\pgfqpoint{0.994469in}{0.633527in}}%
\pgfpathlineto{\pgfqpoint{0.994797in}{0.625278in}}%
\pgfpathlineto{\pgfqpoint{0.995475in}{0.625278in}}%
\pgfpathlineto{\pgfqpoint{0.996816in}{0.633527in}}%
\pgfpathlineto{\pgfqpoint{0.997151in}{0.625278in}}%
\pgfpathlineto{\pgfqpoint{0.997821in}{0.625278in}}%
\pgfpathlineto{\pgfqpoint{0.998988in}{0.633527in}}%
\pgfpathlineto{\pgfqpoint{0.999002in}{0.625278in}}%
\pgfpathlineto{\pgfqpoint{0.999678in}{0.674773in}}%
\pgfpathlineto{\pgfqpoint{1.000012in}{0.674773in}}%
\pgfpathlineto{\pgfqpoint{1.000664in}{0.658274in}}%
\pgfpathlineto{\pgfqpoint{1.000832in}{0.683022in}}%
\pgfpathlineto{\pgfqpoint{1.000851in}{0.674773in}}%
\pgfpathlineto{\pgfqpoint{1.001335in}{0.740766in}}%
\pgfpathlineto{\pgfqpoint{1.002360in}{0.823258in}}%
\pgfpathlineto{\pgfqpoint{1.002530in}{0.806759in}}%
\pgfpathlineto{\pgfqpoint{1.002695in}{0.848005in}}%
\pgfpathlineto{\pgfqpoint{1.003011in}{0.839756in}}%
\pgfpathlineto{\pgfqpoint{1.004036in}{0.971742in}}%
\pgfpathlineto{\pgfqpoint{1.011393in}{1.433695in}}%
\pgfpathlineto{\pgfqpoint{1.011895in}{1.474941in}}%
\pgfpathlineto{\pgfqpoint{1.012417in}{1.466692in}}%
\pgfpathlineto{\pgfqpoint{1.012566in}{1.474941in}}%
\pgfpathlineto{\pgfqpoint{1.012902in}{1.450194in}}%
\pgfpathlineto{\pgfqpoint{1.013088in}{1.458443in}}%
\pgfpathlineto{\pgfqpoint{1.013236in}{1.450194in}}%
\pgfpathlineto{\pgfqpoint{1.013404in}{1.474941in}}%
\pgfpathlineto{\pgfqpoint{1.013758in}{1.466692in}}%
\pgfpathlineto{\pgfqpoint{1.018769in}{1.747163in}}%
\pgfpathlineto{\pgfqpoint{1.019439in}{1.722416in}}%
\pgfpathlineto{\pgfqpoint{1.020128in}{1.689419in}}%
\pgfpathlineto{\pgfqpoint{1.020464in}{1.714167in}}%
\pgfpathlineto{\pgfqpoint{1.022959in}{1.821406in}}%
\pgfpathlineto{\pgfqpoint{1.023146in}{1.804907in}}%
\pgfpathlineto{\pgfqpoint{1.023481in}{1.763662in}}%
\pgfpathlineto{\pgfqpoint{1.024319in}{1.788409in}}%
\pgfpathlineto{\pgfqpoint{1.024468in}{1.796658in}}%
\pgfpathlineto{\pgfqpoint{1.024803in}{1.771911in}}%
\pgfpathlineto{\pgfqpoint{1.024990in}{1.780160in}}%
\pgfpathlineto{\pgfqpoint{1.025325in}{1.763662in}}%
\pgfpathlineto{\pgfqpoint{1.025530in}{1.796658in}}%
\pgfpathlineto{\pgfqpoint{1.025661in}{1.788409in}}%
\pgfpathlineto{\pgfqpoint{1.026144in}{1.821406in}}%
\pgfpathlineto{\pgfqpoint{1.026499in}{1.780160in}}%
\pgfpathlineto{\pgfqpoint{1.028491in}{1.722416in}}%
\pgfpathlineto{\pgfqpoint{1.028491in}{1.730665in}}%
\pgfpathlineto{\pgfqpoint{1.028845in}{1.738914in}}%
\pgfpathlineto{\pgfqpoint{1.029497in}{1.697668in}}%
\pgfpathlineto{\pgfqpoint{1.030689in}{1.631675in}}%
\pgfpathlineto{\pgfqpoint{1.034042in}{1.491439in}}%
\pgfpathlineto{\pgfqpoint{1.034198in}{1.499689in}}%
\pgfpathlineto{\pgfqpoint{1.034545in}{1.466692in}}%
\pgfpathlineto{\pgfqpoint{1.038884in}{1.202719in}}%
\pgfpathlineto{\pgfqpoint{1.040076in}{1.111978in}}%
\pgfpathlineto{\pgfqpoint{1.044081in}{0.831507in}}%
\pgfpathlineto{\pgfqpoint{1.045105in}{0.749015in}}%
\pgfpathlineto{\pgfqpoint{1.045273in}{0.773763in}}%
\pgfpathlineto{\pgfqpoint{1.048291in}{0.650025in}}%
\pgfpathlineto{\pgfqpoint{1.048454in}{0.666524in}}%
\pgfpathlineto{\pgfqpoint{1.048793in}{0.674773in}}%
\pgfpathlineto{\pgfqpoint{1.049129in}{0.650025in}}%
\pgfpathlineto{\pgfqpoint{1.049464in}{0.650025in}}%
\pgfpathlineto{\pgfqpoint{1.049799in}{0.674773in}}%
\pgfpathlineto{\pgfqpoint{1.050302in}{0.641776in}}%
\pgfpathlineto{\pgfqpoint{1.051140in}{0.625278in}}%
\pgfpathlineto{\pgfqpoint{1.051296in}{0.641776in}}%
\pgfpathlineto{\pgfqpoint{1.051638in}{0.683022in}}%
\pgfpathlineto{\pgfqpoint{1.052313in}{0.625278in}}%
\pgfpathlineto{\pgfqpoint{1.053642in}{0.683022in}}%
\pgfpathlineto{\pgfqpoint{1.054158in}{0.666524in}}%
\pgfpathlineto{\pgfqpoint{1.054996in}{0.650025in}}%
\pgfpathlineto{\pgfqpoint{1.054325in}{0.674773in}}%
\pgfpathlineto{\pgfqpoint{1.055151in}{0.666524in}}%
\pgfpathlineto{\pgfqpoint{1.055319in}{0.683022in}}%
\pgfpathlineto{\pgfqpoint{1.055666in}{0.650025in}}%
\pgfpathlineto{\pgfqpoint{1.056169in}{0.650025in}}%
\pgfpathlineto{\pgfqpoint{1.057680in}{0.683022in}}%
\pgfpathlineto{\pgfqpoint{1.058013in}{0.674773in}}%
\pgfpathlineto{\pgfqpoint{1.058180in}{0.699520in}}%
\pgfpathlineto{\pgfqpoint{1.058846in}{0.732517in}}%
\pgfpathlineto{\pgfqpoint{1.059189in}{0.707769in}}%
\pgfpathlineto{\pgfqpoint{1.060865in}{0.658274in}}%
\pgfpathlineto{\pgfqpoint{1.060865in}{0.666524in}}%
\pgfpathlineto{\pgfqpoint{1.061030in}{0.699520in}}%
\pgfpathlineto{\pgfqpoint{1.061871in}{0.691271in}}%
\pgfpathlineto{\pgfqpoint{1.063366in}{0.782012in}}%
\pgfpathlineto{\pgfqpoint{1.063377in}{0.773763in}}%
\pgfpathlineto{\pgfqpoint{1.064886in}{0.749015in}}%
\pgfpathlineto{\pgfqpoint{1.065056in}{0.757264in}}%
\pgfpathlineto{\pgfqpoint{1.065391in}{0.732517in}}%
\pgfpathlineto{\pgfqpoint{1.066564in}{0.707769in}}%
\pgfpathlineto{\pgfqpoint{1.065724in}{0.749015in}}%
\pgfpathlineto{\pgfqpoint{1.066718in}{0.716019in}}%
\pgfpathlineto{\pgfqpoint{1.067570in}{0.757264in}}%
\pgfpathlineto{\pgfqpoint{1.067736in}{0.749015in}}%
\pgfpathlineto{\pgfqpoint{1.067906in}{0.732517in}}%
\pgfpathlineto{\pgfqpoint{1.068071in}{0.773763in}}%
\pgfpathlineto{\pgfqpoint{1.068241in}{0.765513in}}%
\pgfpathlineto{\pgfqpoint{1.069747in}{0.823258in}}%
\pgfpathlineto{\pgfqpoint{1.068576in}{0.757264in}}%
\pgfpathlineto{\pgfqpoint{1.069917in}{0.806759in}}%
\pgfpathlineto{\pgfqpoint{1.074443in}{0.625278in}}%
\pgfpathlineto{\pgfqpoint{1.075617in}{0.633527in}}%
\pgfpathlineto{\pgfqpoint{1.075952in}{0.625278in}}%
\pgfpathlineto{\pgfqpoint{1.076622in}{0.625278in}}%
\pgfpathlineto{\pgfqpoint{1.077796in}{0.633527in}}%
\pgfpathlineto{\pgfqpoint{1.077796in}{0.625278in}}%
\pgfpathlineto{\pgfqpoint{1.078634in}{0.625278in}}%
\pgfpathlineto{\pgfqpoint{1.079472in}{0.633527in}}%
\pgfpathlineto{\pgfqpoint{1.079472in}{0.625278in}}%
\pgfpathlineto{\pgfqpoint{1.080646in}{0.633527in}}%
\pgfpathlineto{\pgfqpoint{1.080813in}{0.625278in}}%
\pgfpathlineto{\pgfqpoint{1.081651in}{0.625278in}}%
\pgfpathlineto{\pgfqpoint{1.082825in}{0.633527in}}%
\pgfpathlineto{\pgfqpoint{1.082985in}{0.625278in}}%
\pgfpathlineto{\pgfqpoint{1.083831in}{0.625278in}}%
\pgfpathlineto{\pgfqpoint{1.085172in}{0.633527in}}%
\pgfpathlineto{\pgfqpoint{1.085674in}{0.625278in}}%
\pgfpathlineto{\pgfqpoint{1.086010in}{0.625278in}}%
\pgfpathlineto{\pgfqpoint{1.087183in}{0.633527in}}%
\pgfpathlineto{\pgfqpoint{1.087351in}{0.625278in}}%
\pgfpathlineto{\pgfqpoint{1.088189in}{0.625278in}}%
\pgfpathlineto{\pgfqpoint{1.088860in}{0.633527in}}%
\pgfpathlineto{\pgfqpoint{1.088860in}{0.625278in}}%
\pgfpathlineto{\pgfqpoint{1.090201in}{0.633527in}}%
\pgfpathlineto{\pgfqpoint{1.090201in}{0.625278in}}%
\pgfpathlineto{\pgfqpoint{1.091206in}{0.625278in}}%
\pgfpathlineto{\pgfqpoint{1.092380in}{0.633527in}}%
\pgfpathlineto{\pgfqpoint{1.092540in}{0.625278in}}%
\pgfpathlineto{\pgfqpoint{1.093385in}{0.625278in}}%
\pgfpathlineto{\pgfqpoint{1.094552in}{0.633527in}}%
\pgfpathlineto{\pgfqpoint{1.094887in}{0.625278in}}%
\pgfpathlineto{\pgfqpoint{1.095558in}{0.625278in}}%
\pgfpathlineto{\pgfqpoint{1.096570in}{0.633527in}}%
\pgfpathlineto{\pgfqpoint{1.096570in}{0.625278in}}%
\pgfpathlineto{\pgfqpoint{1.097744in}{0.633527in}}%
\pgfpathlineto{\pgfqpoint{1.097911in}{0.625278in}}%
\pgfpathlineto{\pgfqpoint{1.098750in}{0.625278in}}%
\pgfpathlineto{\pgfqpoint{1.100090in}{0.633527in}}%
\pgfpathlineto{\pgfqpoint{1.100090in}{0.625278in}}%
\pgfpathlineto{\pgfqpoint{1.100929in}{0.625278in}}%
\pgfpathlineto{\pgfqpoint{1.102270in}{0.633527in}}%
\pgfpathlineto{\pgfqpoint{1.102438in}{0.625278in}}%
\pgfpathlineto{\pgfqpoint{1.103276in}{0.625278in}}%
\pgfpathlineto{\pgfqpoint{1.104267in}{0.633527in}}%
\pgfpathlineto{\pgfqpoint{1.104267in}{0.625278in}}%
\pgfpathlineto{\pgfqpoint{1.105455in}{0.633527in}}%
\pgfpathlineto{\pgfqpoint{1.105608in}{0.625278in}}%
\pgfpathlineto{\pgfqpoint{1.106286in}{0.625278in}}%
\pgfpathlineto{\pgfqpoint{1.107620in}{0.633527in}}%
\pgfpathlineto{\pgfqpoint{1.107795in}{0.625278in}}%
\pgfpathlineto{\pgfqpoint{1.108633in}{0.625278in}}%
\pgfpathlineto{\pgfqpoint{1.109959in}{0.633527in}}%
\pgfpathlineto{\pgfqpoint{1.110295in}{0.625278in}}%
\pgfpathlineto{\pgfqpoint{1.110972in}{0.625278in}}%
\pgfpathlineto{\pgfqpoint{1.112146in}{0.633527in}}%
\pgfpathlineto{\pgfqpoint{1.112321in}{0.625278in}}%
\pgfpathlineto{\pgfqpoint{1.113151in}{0.625278in}}%
\pgfpathlineto{\pgfqpoint{1.114339in}{0.633527in}}%
\pgfpathlineto{\pgfqpoint{1.114653in}{0.625278in}}%
\pgfpathlineto{\pgfqpoint{1.115324in}{0.625278in}}%
\pgfpathlineto{\pgfqpoint{1.116504in}{0.633527in}}%
\pgfpathlineto{\pgfqpoint{1.116665in}{0.625278in}}%
\pgfpathlineto{\pgfqpoint{1.117503in}{0.625278in}}%
\pgfpathlineto{\pgfqpoint{1.118676in}{0.633527in}}%
\pgfpathlineto{\pgfqpoint{1.118844in}{0.625278in}}%
\pgfpathlineto{\pgfqpoint{1.119682in}{0.625278in}}%
\pgfpathlineto{\pgfqpoint{1.121024in}{0.633527in}}%
\pgfpathlineto{\pgfqpoint{1.121024in}{0.625278in}}%
\pgfpathlineto{\pgfqpoint{1.122029in}{0.625278in}}%
\pgfpathlineto{\pgfqpoint{1.123035in}{0.633527in}}%
\pgfpathlineto{\pgfqpoint{1.123035in}{0.625278in}}%
\pgfpathlineto{\pgfqpoint{1.124208in}{0.633527in}}%
\pgfpathlineto{\pgfqpoint{1.124376in}{0.625278in}}%
\pgfpathlineto{\pgfqpoint{1.125221in}{0.625278in}}%
\pgfpathlineto{\pgfqpoint{1.126395in}{0.633527in}}%
\pgfpathlineto{\pgfqpoint{1.126730in}{0.625278in}}%
\pgfpathlineto{\pgfqpoint{1.127225in}{0.625278in}}%
\pgfpathlineto{\pgfqpoint{1.128399in}{0.633527in}}%
\pgfpathlineto{\pgfqpoint{1.128574in}{0.625278in}}%
\pgfpathlineto{\pgfqpoint{1.129244in}{0.625278in}}%
\pgfpathlineto{\pgfqpoint{1.130411in}{0.633527in}}%
\pgfpathlineto{\pgfqpoint{1.130593in}{0.625278in}}%
\pgfpathlineto{\pgfqpoint{1.131416in}{0.625278in}}%
\pgfpathlineto{\pgfqpoint{1.132590in}{0.633527in}}%
\pgfpathlineto{\pgfqpoint{1.132758in}{0.625278in}}%
\pgfpathlineto{\pgfqpoint{1.133596in}{0.625278in}}%
\pgfpathlineto{\pgfqpoint{1.134601in}{0.633527in}}%
\pgfpathlineto{\pgfqpoint{1.134601in}{0.625278in}}%
\pgfpathlineto{\pgfqpoint{1.135775in}{0.633527in}}%
\pgfpathlineto{\pgfqpoint{1.135942in}{0.625278in}}%
\pgfpathlineto{\pgfqpoint{1.136788in}{0.625278in}}%
\pgfpathlineto{\pgfqpoint{1.137961in}{0.633527in}}%
\pgfpathlineto{\pgfqpoint{1.138129in}{0.625278in}}%
\pgfpathlineto{\pgfqpoint{1.138959in}{0.625278in}}%
\pgfpathlineto{\pgfqpoint{1.140133in}{0.633527in}}%
\pgfpathlineto{\pgfqpoint{1.140468in}{0.625278in}}%
\pgfpathlineto{\pgfqpoint{1.141139in}{0.625278in}}%
\pgfpathlineto{\pgfqpoint{1.142312in}{0.633527in}}%
\pgfpathlineto{\pgfqpoint{1.142320in}{0.625278in}}%
\pgfpathlineto{\pgfqpoint{1.143318in}{0.625278in}}%
\pgfpathlineto{\pgfqpoint{1.144492in}{0.633527in}}%
\pgfpathlineto{\pgfqpoint{1.144659in}{0.625278in}}%
\pgfpathlineto{\pgfqpoint{1.145497in}{0.625278in}}%
\pgfpathlineto{\pgfqpoint{1.146671in}{0.633527in}}%
\pgfpathlineto{\pgfqpoint{1.146692in}{0.625278in}}%
\pgfpathlineto{\pgfqpoint{1.147676in}{0.625278in}}%
\pgfpathlineto{\pgfqpoint{1.148871in}{0.633527in}}%
\pgfpathlineto{\pgfqpoint{1.149207in}{0.625278in}}%
\pgfpathlineto{\pgfqpoint{1.149528in}{0.625278in}}%
\pgfpathlineto{\pgfqpoint{1.150862in}{0.633527in}}%
\pgfpathlineto{\pgfqpoint{1.151029in}{0.625278in}}%
\pgfpathlineto{\pgfqpoint{1.151867in}{0.625278in}}%
\pgfpathlineto{\pgfqpoint{1.153011in}{0.633527in}}%
\pgfpathlineto{\pgfqpoint{1.153041in}{0.625278in}}%
\pgfpathlineto{\pgfqpoint{1.154046in}{0.625278in}}%
\pgfpathlineto{\pgfqpoint{1.155220in}{0.633527in}}%
\pgfpathlineto{\pgfqpoint{1.155388in}{0.625278in}}%
\pgfpathlineto{\pgfqpoint{1.156226in}{0.625278in}}%
\pgfpathlineto{\pgfqpoint{1.157399in}{0.633527in}}%
\pgfpathlineto{\pgfqpoint{1.157567in}{0.625278in}}%
\pgfpathlineto{\pgfqpoint{1.158405in}{0.625278in}}%
\pgfpathlineto{\pgfqpoint{1.159579in}{0.633527in}}%
\pgfpathlineto{\pgfqpoint{1.159913in}{0.625278in}}%
\pgfpathlineto{\pgfqpoint{1.160584in}{0.625278in}}%
\pgfpathlineto{\pgfqpoint{1.161590in}{0.633527in}}%
\pgfpathlineto{\pgfqpoint{1.161590in}{0.625278in}}%
\pgfpathlineto{\pgfqpoint{1.162763in}{0.633527in}}%
\pgfpathlineto{\pgfqpoint{1.162931in}{0.625278in}}%
\pgfpathlineto{\pgfqpoint{1.163601in}{0.625278in}}%
\pgfpathlineto{\pgfqpoint{1.164775in}{0.633527in}}%
\pgfpathlineto{\pgfqpoint{1.164942in}{0.625278in}}%
\pgfpathlineto{\pgfqpoint{1.165613in}{0.625278in}}%
\pgfpathlineto{\pgfqpoint{1.166787in}{0.633527in}}%
\pgfpathlineto{\pgfqpoint{1.166954in}{0.625278in}}%
\pgfpathlineto{\pgfqpoint{1.167792in}{0.625278in}}%
\pgfpathlineto{\pgfqpoint{1.168966in}{0.633527in}}%
\pgfpathlineto{\pgfqpoint{1.169133in}{0.625278in}}%
\pgfpathlineto{\pgfqpoint{1.169971in}{0.625278in}}%
\pgfpathlineto{\pgfqpoint{1.171313in}{0.633527in}}%
\pgfpathlineto{\pgfqpoint{1.171480in}{0.625278in}}%
\pgfpathlineto{\pgfqpoint{1.172318in}{0.625278in}}%
\pgfpathlineto{\pgfqpoint{1.173492in}{0.633527in}}%
\pgfpathlineto{\pgfqpoint{1.173659in}{0.625278in}}%
\pgfpathlineto{\pgfqpoint{1.174497in}{0.625278in}}%
\pgfpathlineto{\pgfqpoint{1.175671in}{0.633527in}}%
\pgfpathlineto{\pgfqpoint{1.175838in}{0.625278in}}%
\pgfpathlineto{\pgfqpoint{1.176676in}{0.625278in}}%
\pgfpathlineto{\pgfqpoint{1.177850in}{0.633527in}}%
\pgfpathlineto{\pgfqpoint{1.178018in}{0.625278in}}%
\pgfpathlineto{\pgfqpoint{1.178856in}{0.625278in}}%
\pgfpathlineto{\pgfqpoint{1.180029in}{0.633527in}}%
\pgfpathlineto{\pgfqpoint{1.180029in}{0.625278in}}%
\pgfpathlineto{\pgfqpoint{1.180867in}{0.625278in}}%
\pgfpathlineto{\pgfqpoint{1.182041in}{0.633527in}}%
\pgfpathlineto{\pgfqpoint{1.182209in}{0.625278in}}%
\pgfpathlineto{\pgfqpoint{1.183047in}{0.625278in}}%
\pgfpathlineto{\pgfqpoint{1.184220in}{0.633527in}}%
\pgfpathlineto{\pgfqpoint{1.184388in}{0.625278in}}%
\pgfpathlineto{\pgfqpoint{1.185226in}{0.625278in}}%
\pgfpathlineto{\pgfqpoint{1.186400in}{0.633527in}}%
\pgfpathlineto{\pgfqpoint{1.186456in}{0.625278in}}%
\pgfpathlineto{\pgfqpoint{1.187405in}{0.625278in}}%
\pgfpathlineto{\pgfqpoint{1.188636in}{0.633527in}}%
\pgfpathlineto{\pgfqpoint{1.188746in}{0.625278in}}%
\pgfpathlineto{\pgfqpoint{1.189417in}{0.625278in}}%
\pgfpathlineto{\pgfqpoint{1.190758in}{0.633527in}}%
\pgfpathlineto{\pgfqpoint{1.190758in}{0.625278in}}%
\pgfpathlineto{\pgfqpoint{1.191763in}{0.625278in}}%
\pgfpathlineto{\pgfqpoint{1.192937in}{0.633527in}}%
\pgfpathlineto{\pgfqpoint{1.192937in}{0.625278in}}%
\pgfpathlineto{\pgfqpoint{1.193943in}{0.625278in}}%
\pgfpathlineto{\pgfqpoint{1.195116in}{0.633527in}}%
\pgfpathlineto{\pgfqpoint{1.195116in}{0.625278in}}%
\pgfpathlineto{\pgfqpoint{1.196122in}{0.625278in}}%
\pgfpathlineto{\pgfqpoint{1.197128in}{0.633527in}}%
\pgfpathlineto{\pgfqpoint{1.197128in}{0.625278in}}%
\pgfpathlineto{\pgfqpoint{1.198301in}{0.633527in}}%
\pgfpathlineto{\pgfqpoint{1.198468in}{0.625278in}}%
\pgfpathlineto{\pgfqpoint{1.199307in}{0.625278in}}%
\pgfpathlineto{\pgfqpoint{1.200480in}{0.633527in}}%
\pgfpathlineto{\pgfqpoint{1.200480in}{0.625278in}}%
\pgfpathlineto{\pgfqpoint{1.201486in}{0.625278in}}%
\pgfpathlineto{\pgfqpoint{1.202659in}{0.633527in}}%
\pgfpathlineto{\pgfqpoint{1.202659in}{0.625278in}}%
\pgfpathlineto{\pgfqpoint{1.203665in}{0.625278in}}%
\pgfpathlineto{\pgfqpoint{1.204839in}{0.633527in}}%
\pgfpathlineto{\pgfqpoint{1.205006in}{0.625278in}}%
\pgfpathlineto{\pgfqpoint{1.205845in}{0.625278in}}%
\pgfpathlineto{\pgfqpoint{1.207018in}{0.633527in}}%
\pgfpathlineto{\pgfqpoint{1.207018in}{0.625278in}}%
\pgfpathlineto{\pgfqpoint{1.208024in}{0.625278in}}%
\pgfpathlineto{\pgfqpoint{1.209197in}{0.633527in}}%
\pgfpathlineto{\pgfqpoint{1.209197in}{0.625278in}}%
\pgfpathlineto{\pgfqpoint{1.210203in}{0.625278in}}%
\pgfpathlineto{\pgfqpoint{1.211544in}{0.633527in}}%
\pgfpathlineto{\pgfqpoint{1.211544in}{0.625278in}}%
\pgfpathlineto{\pgfqpoint{1.212550in}{0.625278in}}%
\pgfpathlineto{\pgfqpoint{1.213723in}{0.633527in}}%
\pgfpathlineto{\pgfqpoint{1.213891in}{0.625278in}}%
\pgfpathlineto{\pgfqpoint{1.214729in}{0.625278in}}%
\pgfpathlineto{\pgfqpoint{1.215902in}{0.633527in}}%
\pgfpathlineto{\pgfqpoint{1.216070in}{0.625278in}}%
\pgfpathlineto{\pgfqpoint{1.216741in}{0.625278in}}%
\pgfpathlineto{\pgfqpoint{1.217914in}{0.633527in}}%
\pgfpathlineto{\pgfqpoint{1.218081in}{0.625278in}}%
\pgfpathlineto{\pgfqpoint{1.218920in}{0.625278in}}%
\pgfpathlineto{\pgfqpoint{1.220093in}{0.633527in}}%
\pgfpathlineto{\pgfqpoint{1.220261in}{0.625278in}}%
\pgfpathlineto{\pgfqpoint{1.221099in}{0.625278in}}%
\pgfpathlineto{\pgfqpoint{1.222272in}{0.633527in}}%
\pgfpathlineto{\pgfqpoint{1.222440in}{0.625278in}}%
\pgfpathlineto{\pgfqpoint{1.223278in}{0.625278in}}%
\pgfpathlineto{\pgfqpoint{1.224451in}{0.633527in}}%
\pgfpathlineto{\pgfqpoint{1.224619in}{0.625278in}}%
\pgfpathlineto{\pgfqpoint{1.225289in}{0.625278in}}%
\pgfpathlineto{\pgfqpoint{1.226463in}{0.633527in}}%
\pgfpathlineto{\pgfqpoint{1.226631in}{0.625278in}}%
\pgfpathlineto{\pgfqpoint{1.227301in}{0.625278in}}%
\pgfpathlineto{\pgfqpoint{1.228475in}{0.633527in}}%
\pgfpathlineto{\pgfqpoint{1.228642in}{0.625278in}}%
\pgfpathlineto{\pgfqpoint{1.229313in}{0.625278in}}%
\pgfpathlineto{\pgfqpoint{1.230654in}{0.633527in}}%
\pgfpathlineto{\pgfqpoint{1.230654in}{0.625278in}}%
\pgfpathlineto{\pgfqpoint{1.231679in}{0.625278in}}%
\pgfpathlineto{\pgfqpoint{1.233188in}{0.724268in}}%
\pgfpathlineto{\pgfqpoint{1.234174in}{0.707769in}}%
\pgfpathlineto{\pgfqpoint{1.234361in}{0.699520in}}%
\pgfpathlineto{\pgfqpoint{1.235013in}{0.782012in}}%
\pgfpathlineto{\pgfqpoint{1.235031in}{0.773763in}}%
\pgfpathlineto{\pgfqpoint{1.235180in}{0.782012in}}%
\pgfpathlineto{\pgfqpoint{1.235534in}{0.749015in}}%
\pgfpathlineto{\pgfqpoint{1.235683in}{0.757264in}}%
\pgfpathlineto{\pgfqpoint{1.236707in}{0.749015in}}%
\pgfpathlineto{\pgfqpoint{1.237714in}{0.773763in}}%
\pgfpathlineto{\pgfqpoint{1.237881in}{0.765513in}}%
\pgfpathlineto{\pgfqpoint{1.238030in}{0.757264in}}%
\pgfpathlineto{\pgfqpoint{1.238048in}{0.798510in}}%
\pgfpathlineto{\pgfqpoint{1.238384in}{0.798510in}}%
\pgfpathlineto{\pgfqpoint{1.239035in}{0.782012in}}%
\pgfpathlineto{\pgfqpoint{1.239203in}{0.806759in}}%
\pgfpathlineto{\pgfqpoint{1.239222in}{0.798510in}}%
\pgfpathlineto{\pgfqpoint{1.240880in}{0.831507in}}%
\pgfpathlineto{\pgfqpoint{1.241066in}{0.823258in}}%
\pgfpathlineto{\pgfqpoint{1.241382in}{0.856254in}}%
\pgfpathlineto{\pgfqpoint{1.241904in}{0.848005in}}%
\pgfpathlineto{\pgfqpoint{1.243059in}{0.856254in}}%
\pgfpathlineto{\pgfqpoint{1.243897in}{0.831507in}}%
\pgfpathlineto{\pgfqpoint{1.244084in}{0.839756in}}%
\pgfpathlineto{\pgfqpoint{1.244902in}{0.831507in}}%
\pgfpathlineto{\pgfqpoint{1.244251in}{0.848005in}}%
\pgfpathlineto{\pgfqpoint{1.244902in}{0.839756in}}%
\pgfpathlineto{\pgfqpoint{1.244922in}{0.872753in}}%
\pgfpathlineto{\pgfqpoint{1.245592in}{0.823258in}}%
\pgfpathlineto{\pgfqpoint{1.245927in}{0.823258in}}%
\pgfpathlineto{\pgfqpoint{1.246243in}{0.831507in}}%
\pgfpathlineto{\pgfqpoint{1.246579in}{0.806759in}}%
\pgfpathlineto{\pgfqpoint{1.246765in}{0.815008in}}%
\pgfpathlineto{\pgfqpoint{1.246914in}{0.806759in}}%
\pgfpathlineto{\pgfqpoint{1.247268in}{0.848005in}}%
\pgfpathlineto{\pgfqpoint{1.247603in}{0.872753in}}%
\pgfpathlineto{\pgfqpoint{1.248423in}{0.856254in}}%
\pgfpathlineto{\pgfqpoint{1.249429in}{0.831507in}}%
\pgfpathlineto{\pgfqpoint{1.249448in}{0.848005in}}%
\pgfpathlineto{\pgfqpoint{1.250770in}{0.831507in}}%
\pgfpathlineto{\pgfqpoint{1.251440in}{0.905749in}}%
\pgfpathlineto{\pgfqpoint{1.251794in}{0.889251in}}%
\pgfpathlineto{\pgfqpoint{1.253135in}{0.848005in}}%
\pgfpathlineto{\pgfqpoint{1.253284in}{0.864503in}}%
\pgfpathlineto{\pgfqpoint{1.254458in}{0.905749in}}%
\pgfpathlineto{\pgfqpoint{1.255631in}{0.856254in}}%
\pgfpathlineto{\pgfqpoint{1.254644in}{0.922247in}}%
\pgfpathlineto{\pgfqpoint{1.255631in}{0.864503in}}%
\pgfpathlineto{\pgfqpoint{1.256134in}{0.881002in}}%
\pgfpathlineto{\pgfqpoint{1.256488in}{0.848005in}}%
\pgfpathlineto{\pgfqpoint{1.256656in}{0.848005in}}%
\pgfpathlineto{\pgfqpoint{1.256972in}{0.856254in}}%
\pgfpathlineto{\pgfqpoint{1.257307in}{0.831507in}}%
\pgfpathlineto{\pgfqpoint{1.257661in}{0.823258in}}%
\pgfpathlineto{\pgfqpoint{1.258164in}{0.872753in}}%
\pgfpathlineto{\pgfqpoint{1.258313in}{0.881002in}}%
\pgfpathlineto{\pgfqpoint{1.258648in}{0.856254in}}%
\pgfpathlineto{\pgfqpoint{1.258835in}{0.864503in}}%
\pgfpathlineto{\pgfqpoint{1.258984in}{0.856254in}}%
\pgfpathlineto{\pgfqpoint{1.259002in}{0.897500in}}%
\pgfpathlineto{\pgfqpoint{1.259338in}{0.897500in}}%
\pgfpathlineto{\pgfqpoint{1.259840in}{0.872753in}}%
\pgfpathlineto{\pgfqpoint{1.260046in}{0.905749in}}%
\pgfpathlineto{\pgfqpoint{1.260847in}{0.946995in}}%
\pgfpathlineto{\pgfqpoint{1.261349in}{0.922247in}}%
\pgfpathlineto{\pgfqpoint{1.261666in}{0.930497in}}%
\pgfpathlineto{\pgfqpoint{1.262001in}{0.905749in}}%
\pgfpathlineto{\pgfqpoint{1.262188in}{0.897500in}}%
\pgfpathlineto{\pgfqpoint{1.262504in}{0.930497in}}%
\pgfpathlineto{\pgfqpoint{1.263026in}{0.922247in}}%
\pgfpathlineto{\pgfqpoint{1.263510in}{0.930497in}}%
\pgfpathlineto{\pgfqpoint{1.263696in}{0.913998in}}%
\pgfpathlineto{\pgfqpoint{1.263845in}{0.905749in}}%
\pgfpathlineto{\pgfqpoint{1.264013in}{0.930497in}}%
\pgfpathlineto{\pgfqpoint{1.264367in}{0.922247in}}%
\pgfpathlineto{\pgfqpoint{1.264515in}{0.930497in}}%
\pgfpathlineto{\pgfqpoint{1.264851in}{0.905749in}}%
\pgfpathlineto{\pgfqpoint{1.265037in}{0.913998in}}%
\pgfpathlineto{\pgfqpoint{1.265521in}{0.905749in}}%
\pgfpathlineto{\pgfqpoint{1.265689in}{0.930497in}}%
\pgfpathlineto{\pgfqpoint{1.265707in}{0.922247in}}%
\pgfpathlineto{\pgfqpoint{1.267365in}{0.955244in}}%
\pgfpathlineto{\pgfqpoint{1.267887in}{0.946995in}}%
\pgfpathlineto{\pgfqpoint{1.268203in}{0.979992in}}%
\pgfpathlineto{\pgfqpoint{1.268390in}{0.971742in}}%
\pgfpathlineto{\pgfqpoint{1.269209in}{0.979992in}}%
\pgfpathlineto{\pgfqpoint{1.269228in}{0.971742in}}%
\pgfpathlineto{\pgfqpoint{1.269564in}{0.946995in}}%
\pgfpathlineto{\pgfqpoint{1.270382in}{0.963493in}}%
\pgfpathlineto{\pgfqpoint{1.270402in}{0.971742in}}%
\pgfpathlineto{\pgfqpoint{1.271053in}{0.930497in}}%
\pgfpathlineto{\pgfqpoint{1.271407in}{0.946995in}}%
\pgfpathlineto{\pgfqpoint{1.271910in}{0.922247in}}%
\pgfpathlineto{\pgfqpoint{1.272226in}{0.955244in}}%
\pgfpathlineto{\pgfqpoint{1.272394in}{0.979992in}}%
\pgfpathlineto{\pgfqpoint{1.272916in}{0.946995in}}%
\pgfpathlineto{\pgfqpoint{1.273251in}{0.946995in}}%
\pgfpathlineto{\pgfqpoint{1.274089in}{0.971742in}}%
\pgfpathlineto{\pgfqpoint{1.274406in}{0.955244in}}%
\pgfpathlineto{\pgfqpoint{1.275431in}{0.946995in}}%
\pgfpathlineto{\pgfqpoint{1.277088in}{1.004739in}}%
\pgfpathlineto{\pgfqpoint{1.277274in}{0.988241in}}%
\pgfpathlineto{\pgfqpoint{1.278261in}{0.930497in}}%
\pgfpathlineto{\pgfqpoint{1.278615in}{0.971742in}}%
\pgfpathlineto{\pgfqpoint{1.278764in}{0.979992in}}%
\pgfpathlineto{\pgfqpoint{1.279118in}{0.946995in}}%
\pgfpathlineto{\pgfqpoint{1.279770in}{0.930497in}}%
\pgfpathlineto{\pgfqpoint{1.279938in}{0.955244in}}%
\pgfpathlineto{\pgfqpoint{1.279956in}{0.946995in}}%
\pgfpathlineto{\pgfqpoint{1.280943in}{0.955244in}}%
\pgfpathlineto{\pgfqpoint{1.280665in}{0.930497in}}%
\pgfpathlineto{\pgfqpoint{1.280962in}{0.946995in}}%
\pgfpathlineto{\pgfqpoint{1.282303in}{0.922247in}}%
\pgfpathlineto{\pgfqpoint{1.284631in}{1.004739in}}%
\pgfpathlineto{\pgfqpoint{1.284650in}{0.996490in}}%
\pgfpathlineto{\pgfqpoint{1.285991in}{0.971742in}}%
\pgfpathlineto{\pgfqpoint{1.287146in}{0.979992in}}%
\pgfpathlineto{\pgfqpoint{1.287165in}{0.971742in}}%
\pgfpathlineto{\pgfqpoint{1.287481in}{1.004739in}}%
\pgfpathlineto{\pgfqpoint{1.288170in}{0.996490in}}%
\pgfpathlineto{\pgfqpoint{1.289325in}{1.029487in}}%
\pgfpathlineto{\pgfqpoint{1.289511in}{1.012988in}}%
\pgfpathlineto{\pgfqpoint{1.289660in}{1.004739in}}%
\pgfpathlineto{\pgfqpoint{1.290014in}{1.045985in}}%
\pgfpathlineto{\pgfqpoint{1.291001in}{1.078981in}}%
\pgfpathlineto{\pgfqpoint{1.291188in}{1.062483in}}%
\pgfpathlineto{\pgfqpoint{1.292677in}{1.029487in}}%
\pgfpathlineto{\pgfqpoint{1.292677in}{1.037736in}}%
\pgfpathlineto{\pgfqpoint{1.293740in}{1.103729in}}%
\pgfpathlineto{\pgfqpoint{1.293013in}{1.029487in}}%
\pgfpathlineto{\pgfqpoint{1.293870in}{1.095480in}}%
\pgfpathlineto{\pgfqpoint{1.294354in}{1.078981in}}%
\pgfpathlineto{\pgfqpoint{1.295024in}{1.103729in}}%
\pgfpathlineto{\pgfqpoint{1.295378in}{1.095480in}}%
\pgfpathlineto{\pgfqpoint{1.295881in}{1.144975in}}%
\pgfpathlineto{\pgfqpoint{1.296049in}{1.136726in}}%
\pgfpathlineto{\pgfqpoint{1.296701in}{1.103729in}}%
\pgfpathlineto{\pgfqpoint{1.297203in}{1.128476in}}%
\pgfpathlineto{\pgfqpoint{1.297223in}{1.144975in}}%
\pgfpathlineto{\pgfqpoint{1.297893in}{1.095480in}}%
\pgfpathlineto{\pgfqpoint{1.298228in}{1.095480in}}%
\pgfpathlineto{\pgfqpoint{1.298544in}{1.103729in}}%
\pgfpathlineto{\pgfqpoint{1.298880in}{1.078981in}}%
\pgfpathlineto{\pgfqpoint{1.299551in}{1.054234in}}%
\pgfpathlineto{\pgfqpoint{1.299904in}{1.095480in}}%
\pgfpathlineto{\pgfqpoint{1.300240in}{1.120227in}}%
\pgfpathlineto{\pgfqpoint{1.301059in}{1.103729in}}%
\pgfpathlineto{\pgfqpoint{1.302084in}{1.070732in}}%
\pgfpathlineto{\pgfqpoint{1.302735in}{1.078981in}}%
\pgfpathlineto{\pgfqpoint{1.302922in}{1.062483in}}%
\pgfpathlineto{\pgfqpoint{1.303424in}{1.045985in}}%
\pgfpathlineto{\pgfqpoint{1.303593in}{1.070732in}}%
\pgfpathlineto{\pgfqpoint{1.304747in}{1.078981in}}%
\pgfpathlineto{\pgfqpoint{1.305772in}{1.045985in}}%
\pgfpathlineto{\pgfqpoint{1.306256in}{1.054234in}}%
\pgfpathlineto{\pgfqpoint{1.306442in}{1.037736in}}%
\pgfpathlineto{\pgfqpoint{1.306610in}{1.045985in}}%
\pgfpathlineto{\pgfqpoint{1.307597in}{1.029487in}}%
\pgfpathlineto{\pgfqpoint{1.308602in}{1.054234in}}%
\pgfpathlineto{\pgfqpoint{1.307951in}{1.021237in}}%
\pgfpathlineto{\pgfqpoint{1.308621in}{1.045985in}}%
\pgfpathlineto{\pgfqpoint{1.308938in}{1.054234in}}%
\pgfpathlineto{\pgfqpoint{1.309273in}{1.029487in}}%
\pgfpathlineto{\pgfqpoint{1.309460in}{1.037736in}}%
\pgfpathlineto{\pgfqpoint{1.309608in}{1.029487in}}%
\pgfpathlineto{\pgfqpoint{1.309627in}{1.070732in}}%
\pgfpathlineto{\pgfqpoint{1.309962in}{1.070732in}}%
\pgfpathlineto{\pgfqpoint{1.311303in}{1.045985in}}%
\pgfpathlineto{\pgfqpoint{1.313799in}{1.128476in}}%
\pgfpathlineto{\pgfqpoint{1.314637in}{1.103729in}}%
\pgfpathlineto{\pgfqpoint{1.314824in}{1.111978in}}%
\pgfpathlineto{\pgfqpoint{1.314991in}{1.095480in}}%
\pgfpathlineto{\pgfqpoint{1.315327in}{1.144975in}}%
\pgfpathlineto{\pgfqpoint{1.315643in}{1.153224in}}%
\pgfpathlineto{\pgfqpoint{1.315978in}{1.128476in}}%
\pgfpathlineto{\pgfqpoint{1.317003in}{1.095480in}}%
\pgfpathlineto{\pgfqpoint{1.318008in}{1.120227in}}%
\pgfpathlineto{\pgfqpoint{1.318177in}{1.111978in}}%
\pgfpathlineto{\pgfqpoint{1.319331in}{1.078981in}}%
\pgfpathlineto{\pgfqpoint{1.318344in}{1.120227in}}%
\pgfpathlineto{\pgfqpoint{1.319349in}{1.095480in}}%
\pgfpathlineto{\pgfqpoint{1.320356in}{1.144975in}}%
\pgfpathlineto{\pgfqpoint{1.320691in}{1.120227in}}%
\pgfpathlineto{\pgfqpoint{1.321194in}{1.095480in}}%
\pgfpathlineto{\pgfqpoint{1.321343in}{1.136726in}}%
\pgfpathlineto{\pgfqpoint{1.322515in}{1.177971in}}%
\pgfpathlineto{\pgfqpoint{1.323708in}{1.111978in}}%
\pgfpathlineto{\pgfqpoint{1.324211in}{1.095480in}}%
\pgfpathlineto{\pgfqpoint{1.324360in}{1.136726in}}%
\pgfpathlineto{\pgfqpoint{1.324882in}{1.169722in}}%
\pgfpathlineto{\pgfqpoint{1.325385in}{1.136726in}}%
\pgfpathlineto{\pgfqpoint{1.325868in}{1.128476in}}%
\pgfpathlineto{\pgfqpoint{1.326036in}{1.153224in}}%
\pgfpathlineto{\pgfqpoint{1.326055in}{1.144975in}}%
\pgfpathlineto{\pgfqpoint{1.326874in}{1.153224in}}%
\pgfpathlineto{\pgfqpoint{1.326893in}{1.144975in}}%
\pgfpathlineto{\pgfqpoint{1.328234in}{1.120227in}}%
\pgfpathlineto{\pgfqpoint{1.328718in}{1.128476in}}%
\pgfpathlineto{\pgfqpoint{1.328904in}{1.111978in}}%
\pgfpathlineto{\pgfqpoint{1.330562in}{1.054234in}}%
\pgfpathlineto{\pgfqpoint{1.330581in}{1.070732in}}%
\pgfpathlineto{\pgfqpoint{1.331232in}{1.078981in}}%
\pgfpathlineto{\pgfqpoint{1.331419in}{1.062483in}}%
\pgfpathlineto{\pgfqpoint{1.331587in}{1.070732in}}%
\pgfpathlineto{\pgfqpoint{1.333095in}{1.021237in}}%
\pgfpathlineto{\pgfqpoint{1.333915in}{1.029487in}}%
\pgfpathlineto{\pgfqpoint{1.333933in}{1.021237in}}%
\pgfpathlineto{\pgfqpoint{1.334604in}{0.996490in}}%
\pgfpathlineto{\pgfqpoint{1.335088in}{1.012988in}}%
\pgfpathlineto{\pgfqpoint{1.335423in}{1.029487in}}%
\pgfpathlineto{\pgfqpoint{1.335945in}{0.971742in}}%
\pgfpathlineto{\pgfqpoint{1.336094in}{0.979992in}}%
\pgfpathlineto{\pgfqpoint{1.336951in}{0.971742in}}%
\pgfpathlineto{\pgfqpoint{1.337119in}{0.996490in}}%
\pgfpathlineto{\pgfqpoint{1.337435in}{0.979992in}}%
\pgfpathlineto{\pgfqpoint{1.337602in}{1.004739in}}%
\pgfpathlineto{\pgfqpoint{1.337957in}{0.996490in}}%
\pgfpathlineto{\pgfqpoint{1.338795in}{1.021237in}}%
\pgfpathlineto{\pgfqpoint{1.339119in}{1.004739in}}%
\pgfpathlineto{\pgfqpoint{1.339298in}{0.996490in}}%
\pgfpathlineto{\pgfqpoint{1.339614in}{1.029487in}}%
\pgfpathlineto{\pgfqpoint{1.340136in}{1.021237in}}%
\pgfpathlineto{\pgfqpoint{1.340459in}{1.029487in}}%
\pgfpathlineto{\pgfqpoint{1.340787in}{1.004739in}}%
\pgfpathlineto{\pgfqpoint{1.341626in}{0.979992in}}%
\pgfpathlineto{\pgfqpoint{1.341645in}{1.021237in}}%
\pgfpathlineto{\pgfqpoint{1.341812in}{1.012988in}}%
\pgfpathlineto{\pgfqpoint{1.342967in}{0.979992in}}%
\pgfpathlineto{\pgfqpoint{1.342967in}{0.988241in}}%
\pgfpathlineto{\pgfqpoint{1.343637in}{1.029487in}}%
\pgfpathlineto{\pgfqpoint{1.343991in}{0.988241in}}%
\pgfpathlineto{\pgfqpoint{1.346209in}{0.905749in}}%
\pgfpathlineto{\pgfqpoint{1.346319in}{0.913998in}}%
\pgfpathlineto{\pgfqpoint{1.346338in}{0.946995in}}%
\pgfpathlineto{\pgfqpoint{1.346990in}{0.905749in}}%
\pgfpathlineto{\pgfqpoint{1.347344in}{0.922247in}}%
\pgfpathlineto{\pgfqpoint{1.348517in}{0.897500in}}%
\pgfpathlineto{\pgfqpoint{1.348666in}{0.905749in}}%
\pgfpathlineto{\pgfqpoint{1.348685in}{0.922247in}}%
\pgfpathlineto{\pgfqpoint{1.349336in}{0.881002in}}%
\pgfpathlineto{\pgfqpoint{1.349691in}{0.897500in}}%
\pgfpathlineto{\pgfqpoint{1.350845in}{0.930497in}}%
\pgfpathlineto{\pgfqpoint{1.350194in}{0.872753in}}%
\pgfpathlineto{\pgfqpoint{1.351032in}{0.913998in}}%
\pgfpathlineto{\pgfqpoint{1.352541in}{0.872753in}}%
\pgfpathlineto{\pgfqpoint{1.352689in}{0.881002in}}%
\pgfpathlineto{\pgfqpoint{1.352864in}{0.905749in}}%
\pgfpathlineto{\pgfqpoint{1.353714in}{0.897500in}}%
\pgfpathlineto{\pgfqpoint{1.355211in}{0.930497in}}%
\pgfpathlineto{\pgfqpoint{1.355223in}{0.922247in}}%
\pgfpathlineto{\pgfqpoint{1.355725in}{0.897500in}}%
\pgfpathlineto{\pgfqpoint{1.356042in}{0.930497in}}%
\pgfpathlineto{\pgfqpoint{1.356061in}{0.922247in}}%
\pgfpathlineto{\pgfqpoint{1.357223in}{0.955244in}}%
\pgfpathlineto{\pgfqpoint{1.356552in}{0.905749in}}%
\pgfpathlineto{\pgfqpoint{1.357402in}{0.938746in}}%
\pgfpathlineto{\pgfqpoint{1.357558in}{0.930497in}}%
\pgfpathlineto{\pgfqpoint{1.357718in}{0.955244in}}%
\pgfpathlineto{\pgfqpoint{1.358073in}{0.946995in}}%
\pgfpathlineto{\pgfqpoint{1.358228in}{0.955244in}}%
\pgfpathlineto{\pgfqpoint{1.358575in}{0.922247in}}%
\pgfpathlineto{\pgfqpoint{1.358724in}{0.930497in}}%
\pgfpathlineto{\pgfqpoint{1.359749in}{0.897500in}}%
\pgfpathlineto{\pgfqpoint{1.359916in}{0.922247in}}%
\pgfpathlineto{\pgfqpoint{1.360754in}{0.897500in}}%
\pgfpathlineto{\pgfqpoint{1.362412in}{0.856254in}}%
\pgfpathlineto{\pgfqpoint{1.362419in}{0.864503in}}%
\pgfpathlineto{\pgfqpoint{1.362766in}{0.872753in}}%
\pgfpathlineto{\pgfqpoint{1.363271in}{0.831507in}}%
\pgfpathlineto{\pgfqpoint{1.363437in}{0.848005in}}%
\pgfpathlineto{\pgfqpoint{1.364442in}{0.872753in}}%
\pgfpathlineto{\pgfqpoint{1.364610in}{0.864503in}}%
\pgfpathlineto{\pgfqpoint{1.364778in}{0.872753in}}%
\pgfpathlineto{\pgfqpoint{1.365951in}{0.848005in}}%
\pgfpathlineto{\pgfqpoint{1.366456in}{0.831507in}}%
\pgfpathlineto{\pgfqpoint{1.367292in}{0.872753in}}%
\pgfpathlineto{\pgfqpoint{1.367462in}{0.856254in}}%
\pgfpathlineto{\pgfqpoint{1.367628in}{0.897500in}}%
\pgfpathlineto{\pgfqpoint{1.368298in}{0.872753in}}%
\pgfpathlineto{\pgfqpoint{1.368300in}{0.881002in}}%
\pgfpathlineto{\pgfqpoint{1.369138in}{0.831507in}}%
\pgfpathlineto{\pgfqpoint{1.369304in}{0.848005in}}%
\pgfpathlineto{\pgfqpoint{1.369809in}{0.806759in}}%
\pgfpathlineto{\pgfqpoint{1.370812in}{0.823258in}}%
\pgfpathlineto{\pgfqpoint{1.370961in}{0.831507in}}%
\pgfpathlineto{\pgfqpoint{1.371296in}{0.806759in}}%
\pgfpathlineto{\pgfqpoint{1.371483in}{0.815008in}}%
\pgfpathlineto{\pgfqpoint{1.374165in}{0.724268in}}%
\pgfpathlineto{\pgfqpoint{1.374328in}{0.732517in}}%
\pgfpathlineto{\pgfqpoint{1.374668in}{0.699520in}}%
\pgfpathlineto{\pgfqpoint{1.375003in}{0.716019in}}%
\pgfpathlineto{\pgfqpoint{1.377674in}{0.633527in}}%
\pgfpathlineto{\pgfqpoint{1.377841in}{0.658274in}}%
\pgfpathlineto{\pgfqpoint{1.378188in}{0.625278in}}%
\pgfpathlineto{\pgfqpoint{1.378691in}{0.641776in}}%
\pgfpathlineto{\pgfqpoint{1.379859in}{0.625278in}}%
\pgfpathlineto{\pgfqpoint{1.381033in}{0.633527in}}%
\pgfpathlineto{\pgfqpoint{1.381876in}{0.625278in}}%
\pgfpathlineto{\pgfqpoint{1.382044in}{0.650025in}}%
\pgfpathlineto{\pgfqpoint{1.383380in}{0.625278in}}%
\pgfpathlineto{\pgfqpoint{1.384554in}{0.633527in}}%
\pgfpathlineto{\pgfqpoint{1.384714in}{0.625278in}}%
\pgfpathlineto{\pgfqpoint{1.385552in}{0.625278in}}%
\pgfpathlineto{\pgfqpoint{1.386405in}{0.633527in}}%
\pgfpathlineto{\pgfqpoint{1.386405in}{0.625278in}}%
\pgfpathlineto{\pgfqpoint{1.387578in}{0.633527in}}%
\pgfpathlineto{\pgfqpoint{1.387738in}{0.625278in}}%
\pgfpathlineto{\pgfqpoint{1.388584in}{0.625278in}}%
\pgfpathlineto{\pgfqpoint{1.389757in}{0.633527in}}%
\pgfpathlineto{\pgfqpoint{1.389925in}{0.625278in}}%
\pgfpathlineto{\pgfqpoint{1.390763in}{0.625278in}}%
\pgfpathlineto{\pgfqpoint{1.391768in}{0.633527in}}%
\pgfpathlineto{\pgfqpoint{1.391768in}{0.625278in}}%
\pgfpathlineto{\pgfqpoint{1.392935in}{0.633527in}}%
\pgfpathlineto{\pgfqpoint{1.393095in}{0.625278in}}%
\pgfpathlineto{\pgfqpoint{1.393933in}{0.625278in}}%
\pgfpathlineto{\pgfqpoint{1.395121in}{0.633527in}}%
\pgfpathlineto{\pgfqpoint{1.395121in}{0.625278in}}%
\pgfpathlineto{\pgfqpoint{1.396120in}{0.658274in}}%
\pgfpathlineto{\pgfqpoint{1.396127in}{0.658274in}}%
\pgfpathlineto{\pgfqpoint{1.398136in}{0.724268in}}%
\pgfpathlineto{\pgfqpoint{1.398304in}{0.716019in}}%
\pgfpathlineto{\pgfqpoint{1.399640in}{0.625278in}}%
\pgfpathlineto{\pgfqpoint{1.400813in}{0.633527in}}%
\pgfpathlineto{\pgfqpoint{1.400988in}{0.625278in}}%
\pgfpathlineto{\pgfqpoint{1.401820in}{0.625278in}}%
\pgfpathlineto{\pgfqpoint{1.402978in}{0.633527in}}%
\pgfpathlineto{\pgfqpoint{1.402993in}{0.625278in}}%
\pgfpathlineto{\pgfqpoint{1.403999in}{0.625278in}}%
\pgfpathlineto{\pgfqpoint{1.404837in}{0.633527in}}%
\pgfpathlineto{\pgfqpoint{1.404837in}{0.625278in}}%
\pgfpathlineto{\pgfqpoint{1.406010in}{0.633527in}}%
\pgfpathlineto{\pgfqpoint{1.406178in}{0.625278in}}%
\pgfpathlineto{\pgfqpoint{1.406848in}{0.625278in}}%
\pgfpathlineto{\pgfqpoint{1.408022in}{0.633527in}}%
\pgfpathlineto{\pgfqpoint{1.408685in}{0.625278in}}%
\pgfpathlineto{\pgfqpoint{1.408853in}{0.625278in}}%
\pgfpathlineto{\pgfqpoint{1.410187in}{0.633527in}}%
\pgfpathlineto{\pgfqpoint{1.410522in}{0.625278in}}%
\pgfpathlineto{\pgfqpoint{1.411192in}{0.625278in}}%
\pgfpathlineto{\pgfqpoint{1.412031in}{0.633527in}}%
\pgfpathlineto{\pgfqpoint{1.412031in}{0.625278in}}%
\pgfpathlineto{\pgfqpoint{1.413204in}{0.633527in}}%
\pgfpathlineto{\pgfqpoint{1.413379in}{0.625278in}}%
\pgfpathlineto{\pgfqpoint{1.414217in}{0.625278in}}%
\pgfpathlineto{\pgfqpoint{1.415551in}{0.633527in}}%
\pgfpathlineto{\pgfqpoint{1.415719in}{0.625278in}}%
\pgfpathlineto{\pgfqpoint{1.416389in}{0.625278in}}%
\pgfpathlineto{\pgfqpoint{1.417562in}{0.633527in}}%
\pgfpathlineto{\pgfqpoint{1.417562in}{0.625278in}}%
\pgfpathlineto{\pgfqpoint{1.418575in}{0.625278in}}%
\pgfpathlineto{\pgfqpoint{1.419756in}{0.633527in}}%
\pgfpathlineto{\pgfqpoint{1.419916in}{0.625278in}}%
\pgfpathlineto{\pgfqpoint{1.420748in}{0.625278in}}%
\pgfpathlineto{\pgfqpoint{1.421921in}{0.633527in}}%
\pgfpathlineto{\pgfqpoint{1.422256in}{0.625278in}}%
\pgfpathlineto{\pgfqpoint{1.422927in}{0.625278in}}%
\pgfpathlineto{\pgfqpoint{1.424100in}{0.633527in}}%
\pgfpathlineto{\pgfqpoint{1.424268in}{0.625278in}}%
\pgfpathlineto{\pgfqpoint{1.425106in}{0.625278in}}%
\pgfpathlineto{\pgfqpoint{1.426279in}{0.633527in}}%
\pgfpathlineto{\pgfqpoint{1.426279in}{0.625278in}}%
\pgfpathlineto{\pgfqpoint{1.427285in}{0.625278in}}%
\pgfpathlineto{\pgfqpoint{1.428458in}{0.633527in}}%
\pgfpathlineto{\pgfqpoint{1.428794in}{0.625278in}}%
\pgfpathlineto{\pgfqpoint{1.429296in}{0.625278in}}%
\pgfpathlineto{\pgfqpoint{1.430637in}{0.633527in}}%
\pgfpathlineto{\pgfqpoint{1.430637in}{0.625278in}}%
\pgfpathlineto{\pgfqpoint{1.431644in}{0.625278in}}%
\pgfpathlineto{\pgfqpoint{1.432817in}{0.633527in}}%
\pgfpathlineto{\pgfqpoint{1.432985in}{0.625278in}}%
\pgfpathlineto{\pgfqpoint{1.433655in}{0.625278in}}%
\pgfpathlineto{\pgfqpoint{1.434828in}{0.633527in}}%
\pgfpathlineto{\pgfqpoint{1.434828in}{0.625278in}}%
\pgfpathlineto{\pgfqpoint{1.435834in}{0.625278in}}%
\pgfpathlineto{\pgfqpoint{1.437175in}{0.633527in}}%
\pgfpathlineto{\pgfqpoint{1.437175in}{0.625278in}}%
\pgfpathlineto{\pgfqpoint{1.438181in}{0.625278in}}%
\pgfpathlineto{\pgfqpoint{1.439354in}{0.633527in}}%
\pgfpathlineto{\pgfqpoint{1.439522in}{0.625278in}}%
\pgfpathlineto{\pgfqpoint{1.440361in}{0.625278in}}%
\pgfpathlineto{\pgfqpoint{1.441533in}{0.633527in}}%
\pgfpathlineto{\pgfqpoint{1.441701in}{0.625278in}}%
\pgfpathlineto{\pgfqpoint{1.442540in}{0.625278in}}%
\pgfpathlineto{\pgfqpoint{1.443713in}{0.633527in}}%
\pgfpathlineto{\pgfqpoint{1.443881in}{0.625278in}}%
\pgfpathlineto{\pgfqpoint{1.444551in}{0.625278in}}%
\pgfpathlineto{\pgfqpoint{1.445892in}{0.633527in}}%
\pgfpathlineto{\pgfqpoint{1.445892in}{0.625278in}}%
\pgfpathlineto{\pgfqpoint{1.446898in}{0.625278in}}%
\pgfpathlineto{\pgfqpoint{1.448071in}{0.633527in}}%
\pgfpathlineto{\pgfqpoint{1.448239in}{0.625278in}}%
\pgfpathlineto{\pgfqpoint{1.449077in}{0.625278in}}%
\pgfpathlineto{\pgfqpoint{1.450083in}{0.633527in}}%
\pgfpathlineto{\pgfqpoint{1.450083in}{0.625278in}}%
\pgfpathlineto{\pgfqpoint{1.451424in}{0.633527in}}%
\pgfpathlineto{\pgfqpoint{1.451424in}{0.625278in}}%
\pgfpathlineto{\pgfqpoint{1.452429in}{0.625278in}}%
\pgfpathlineto{\pgfqpoint{1.453771in}{0.633527in}}%
\pgfpathlineto{\pgfqpoint{1.453771in}{0.625278in}}%
\pgfpathlineto{\pgfqpoint{1.454777in}{0.625278in}}%
\pgfpathlineto{\pgfqpoint{1.455950in}{0.633527in}}%
\pgfpathlineto{\pgfqpoint{1.456117in}{0.625278in}}%
\pgfpathlineto{\pgfqpoint{1.456956in}{0.625278in}}%
\pgfpathlineto{\pgfqpoint{1.458129in}{0.633527in}}%
\pgfpathlineto{\pgfqpoint{1.458186in}{0.625278in}}%
\pgfpathlineto{\pgfqpoint{1.459135in}{0.625278in}}%
\pgfpathlineto{\pgfqpoint{1.460476in}{0.633527in}}%
\pgfpathlineto{\pgfqpoint{1.460476in}{0.625278in}}%
\pgfpathlineto{\pgfqpoint{1.461482in}{0.625278in}}%
\pgfpathlineto{\pgfqpoint{1.462655in}{0.633527in}}%
\pgfpathlineto{\pgfqpoint{1.462823in}{0.625278in}}%
\pgfpathlineto{\pgfqpoint{1.463494in}{0.625278in}}%
\pgfpathlineto{\pgfqpoint{1.464667in}{0.633527in}}%
\pgfpathlineto{\pgfqpoint{1.464834in}{0.625278in}}%
\pgfpathlineto{\pgfqpoint{1.465673in}{0.625278in}}%
\pgfpathlineto{\pgfqpoint{1.466846in}{0.633527in}}%
\pgfpathlineto{\pgfqpoint{1.467013in}{0.625278in}}%
\pgfpathlineto{\pgfqpoint{1.467684in}{0.625278in}}%
\pgfpathlineto{\pgfqpoint{1.468858in}{0.633527in}}%
\pgfpathlineto{\pgfqpoint{1.469025in}{0.625278in}}%
\pgfpathlineto{\pgfqpoint{1.469696in}{0.625278in}}%
\pgfpathlineto{\pgfqpoint{1.471035in}{0.633527in}}%
\pgfpathlineto{\pgfqpoint{1.471035in}{0.625278in}}%
\pgfpathlineto{\pgfqpoint{1.472042in}{0.625278in}}%
\pgfpathlineto{\pgfqpoint{1.473219in}{0.633527in}}%
\pgfpathlineto{\pgfqpoint{1.473219in}{0.625278in}}%
\pgfpathlineto{\pgfqpoint{1.474056in}{0.625278in}}%
\pgfpathlineto{\pgfqpoint{1.475225in}{0.633527in}}%
\pgfpathlineto{\pgfqpoint{1.475225in}{0.625278in}}%
\pgfpathlineto{\pgfqpoint{1.476232in}{0.625278in}}%
\pgfpathlineto{\pgfqpoint{1.477573in}{0.633527in}}%
\pgfpathlineto{\pgfqpoint{1.477573in}{0.625278in}}%
\pgfpathlineto{\pgfqpoint{1.478580in}{0.625278in}}%
\pgfpathlineto{\pgfqpoint{1.479757in}{0.633527in}}%
\pgfpathlineto{\pgfqpoint{1.480090in}{0.625278in}}%
\pgfpathlineto{\pgfqpoint{1.480756in}{0.625278in}}%
\pgfpathlineto{\pgfqpoint{1.481934in}{0.633527in}}%
\pgfpathlineto{\pgfqpoint{1.482104in}{0.625278in}}%
\pgfpathlineto{\pgfqpoint{1.482770in}{0.625278in}}%
\pgfpathlineto{\pgfqpoint{1.483948in}{0.633527in}}%
\pgfpathlineto{\pgfqpoint{1.484110in}{0.625278in}}%
\pgfpathlineto{\pgfqpoint{1.484784in}{0.625278in}}%
\pgfpathlineto{\pgfqpoint{1.485954in}{0.633527in}}%
\pgfpathlineto{\pgfqpoint{1.486124in}{0.625278in}}%
\pgfpathlineto{\pgfqpoint{1.486791in}{0.625278in}}%
\pgfpathlineto{\pgfqpoint{1.488138in}{0.633527in}}%
\pgfpathlineto{\pgfqpoint{1.488138in}{0.625278in}}%
\pgfpathlineto{\pgfqpoint{1.489138in}{0.625278in}}%
\pgfpathlineto{\pgfqpoint{1.490315in}{0.633527in}}%
\pgfpathlineto{\pgfqpoint{1.490315in}{0.625278in}}%
\pgfpathlineto{\pgfqpoint{1.491152in}{0.625278in}}%
\pgfpathlineto{\pgfqpoint{1.492329in}{0.633527in}}%
\pgfpathlineto{\pgfqpoint{1.492492in}{0.625278in}}%
\pgfpathlineto{\pgfqpoint{1.493329in}{0.625278in}}%
\pgfpathlineto{\pgfqpoint{1.494506in}{0.633527in}}%
\pgfpathlineto{\pgfqpoint{1.494676in}{0.625278in}}%
\pgfpathlineto{\pgfqpoint{1.495528in}{0.625278in}}%
\pgfpathlineto{\pgfqpoint{1.496705in}{0.650025in}}%
\pgfpathlineto{\pgfqpoint{1.496868in}{0.641776in}}%
\pgfpathlineto{\pgfqpoint{1.497016in}{0.633527in}}%
\pgfpathlineto{\pgfqpoint{1.497371in}{0.674773in}}%
\pgfpathlineto{\pgfqpoint{1.498193in}{0.683022in}}%
\pgfpathlineto{\pgfqpoint{1.498215in}{0.674773in}}%
\pgfpathlineto{\pgfqpoint{1.498867in}{0.658274in}}%
\pgfpathlineto{\pgfqpoint{1.498867in}{0.691271in}}%
\pgfpathlineto{\pgfqpoint{1.498882in}{0.699520in}}%
\pgfpathlineto{\pgfqpoint{1.499704in}{0.658274in}}%
\pgfpathlineto{\pgfqpoint{1.499889in}{0.666524in}}%
\pgfpathlineto{\pgfqpoint{1.500392in}{0.650025in}}%
\pgfpathlineto{\pgfqpoint{1.500555in}{0.674773in}}%
\pgfpathlineto{\pgfqpoint{1.500896in}{0.699520in}}%
\pgfpathlineto{\pgfqpoint{1.501399in}{0.666524in}}%
\pgfpathlineto{\pgfqpoint{1.502569in}{0.650025in}}%
\pgfpathlineto{\pgfqpoint{1.502717in}{0.658274in}}%
\pgfpathlineto{\pgfqpoint{1.503073in}{0.625278in}}%
\pgfpathlineto{\pgfqpoint{1.504057in}{0.658274in}}%
\pgfpathlineto{\pgfqpoint{1.504250in}{0.641776in}}%
\pgfpathlineto{\pgfqpoint{1.505590in}{0.625278in}}%
\pgfpathlineto{\pgfqpoint{1.506908in}{0.633527in}}%
\pgfpathlineto{\pgfqpoint{1.506908in}{0.625278in}}%
\pgfpathlineto{\pgfqpoint{1.507915in}{0.625278in}}%
\pgfpathlineto{\pgfqpoint{1.509092in}{0.633527in}}%
\pgfpathlineto{\pgfqpoint{1.509255in}{0.625278in}}%
\pgfpathlineto{\pgfqpoint{1.509929in}{0.625278in}}%
\pgfpathlineto{\pgfqpoint{1.511099in}{0.633527in}}%
\pgfpathlineto{\pgfqpoint{1.511439in}{0.658274in}}%
\pgfpathlineto{\pgfqpoint{1.512128in}{0.625278in}}%
\pgfpathlineto{\pgfqpoint{1.513779in}{0.658274in}}%
\pgfpathlineto{\pgfqpoint{1.513972in}{0.650025in}}%
\pgfpathlineto{\pgfqpoint{1.514638in}{0.699520in}}%
\pgfpathlineto{\pgfqpoint{1.515793in}{0.707769in}}%
\pgfpathlineto{\pgfqpoint{1.516985in}{0.650025in}}%
\pgfpathlineto{\pgfqpoint{1.518977in}{0.707769in}}%
\pgfpathlineto{\pgfqpoint{1.520006in}{0.699520in}}%
\pgfpathlineto{\pgfqpoint{1.520658in}{0.707769in}}%
\pgfpathlineto{\pgfqpoint{1.520843in}{0.691271in}}%
\pgfpathlineto{\pgfqpoint{1.520991in}{0.683022in}}%
\pgfpathlineto{\pgfqpoint{1.521161in}{0.707769in}}%
\pgfpathlineto{\pgfqpoint{1.521509in}{0.699520in}}%
\pgfpathlineto{\pgfqpoint{1.523020in}{0.749015in}}%
\pgfpathlineto{\pgfqpoint{1.523190in}{0.740766in}}%
\pgfpathlineto{\pgfqpoint{1.524345in}{0.683022in}}%
\pgfpathlineto{\pgfqpoint{1.524678in}{0.716019in}}%
\pgfpathlineto{\pgfqpoint{1.524700in}{0.724268in}}%
\pgfpathlineto{\pgfqpoint{1.525367in}{0.674773in}}%
\pgfpathlineto{\pgfqpoint{1.525685in}{0.683022in}}%
\pgfpathlineto{\pgfqpoint{1.525870in}{0.674773in}}%
\pgfpathlineto{\pgfqpoint{1.526522in}{0.732517in}}%
\pgfpathlineto{\pgfqpoint{1.526707in}{0.724268in}}%
\pgfpathlineto{\pgfqpoint{1.526855in}{0.732517in}}%
\pgfpathlineto{\pgfqpoint{1.527210in}{0.699520in}}%
\pgfpathlineto{\pgfqpoint{1.527358in}{0.707769in}}%
\pgfpathlineto{\pgfqpoint{1.527714in}{0.699520in}}%
\pgfpathlineto{\pgfqpoint{1.527884in}{0.724268in}}%
\pgfpathlineto{\pgfqpoint{1.528388in}{0.716019in}}%
\pgfpathlineto{\pgfqpoint{1.530046in}{0.683022in}}%
\pgfpathlineto{\pgfqpoint{1.530209in}{0.707769in}}%
\pgfpathlineto{\pgfqpoint{1.531068in}{0.666524in}}%
\pgfpathlineto{\pgfqpoint{1.531216in}{0.658274in}}%
\pgfpathlineto{\pgfqpoint{1.531571in}{0.699520in}}%
\pgfpathlineto{\pgfqpoint{1.531719in}{0.707769in}}%
\pgfpathlineto{\pgfqpoint{1.532075in}{0.674773in}}%
\pgfpathlineto{\pgfqpoint{1.533415in}{0.650025in}}%
\pgfpathlineto{\pgfqpoint{1.535910in}{0.732517in}}%
\pgfpathlineto{\pgfqpoint{1.536599in}{0.699520in}}%
\pgfpathlineto{\pgfqpoint{1.536932in}{0.724268in}}%
\pgfpathlineto{\pgfqpoint{1.537102in}{0.749015in}}%
\pgfpathlineto{\pgfqpoint{1.538087in}{0.732517in}}%
\pgfpathlineto{\pgfqpoint{1.538428in}{0.707769in}}%
\pgfpathlineto{\pgfqpoint{1.539116in}{0.724268in}}%
\pgfpathlineto{\pgfqpoint{1.539620in}{0.749015in}}%
\pgfpathlineto{\pgfqpoint{1.540123in}{0.724268in}}%
\pgfpathlineto{\pgfqpoint{1.540434in}{0.707769in}}%
\pgfpathlineto{\pgfqpoint{1.540604in}{0.732517in}}%
\pgfpathlineto{\pgfqpoint{1.540960in}{0.724268in}}%
\pgfpathlineto{\pgfqpoint{1.542781in}{0.806759in}}%
\pgfpathlineto{\pgfqpoint{1.542966in}{0.790261in}}%
\pgfpathlineto{\pgfqpoint{1.543973in}{0.749015in}}%
\pgfpathlineto{\pgfqpoint{1.544462in}{0.765513in}}%
\pgfpathlineto{\pgfqpoint{1.544477in}{0.773763in}}%
\pgfpathlineto{\pgfqpoint{1.544810in}{0.749015in}}%
\pgfpathlineto{\pgfqpoint{1.545484in}{0.749015in}}%
\pgfpathlineto{\pgfqpoint{1.546994in}{0.798510in}}%
\pgfpathlineto{\pgfqpoint{1.547157in}{0.790261in}}%
\pgfpathlineto{\pgfqpoint{1.547328in}{0.798510in}}%
\pgfpathlineto{\pgfqpoint{1.548312in}{0.782012in}}%
\pgfpathlineto{\pgfqpoint{1.548334in}{0.798510in}}%
\pgfpathlineto{\pgfqpoint{1.549341in}{0.798510in}}%
\pgfpathlineto{\pgfqpoint{1.550156in}{0.806759in}}%
\pgfpathlineto{\pgfqpoint{1.550682in}{0.773763in}}%
\pgfpathlineto{\pgfqpoint{1.551185in}{0.823258in}}%
\pgfpathlineto{\pgfqpoint{1.551837in}{0.782012in}}%
\pgfpathlineto{\pgfqpoint{1.552022in}{0.773763in}}%
\pgfpathlineto{\pgfqpoint{1.552340in}{0.806759in}}%
\pgfpathlineto{\pgfqpoint{1.552858in}{0.798510in}}%
\pgfpathlineto{\pgfqpoint{1.554347in}{0.831507in}}%
\pgfpathlineto{\pgfqpoint{1.554369in}{0.823258in}}%
\pgfpathlineto{\pgfqpoint{1.555709in}{0.798510in}}%
\pgfpathlineto{\pgfqpoint{1.556886in}{0.823258in}}%
\pgfpathlineto{\pgfqpoint{1.557049in}{0.815008in}}%
\pgfpathlineto{\pgfqpoint{1.557723in}{0.798510in}}%
\pgfpathlineto{\pgfqpoint{1.557886in}{0.823258in}}%
\pgfpathlineto{\pgfqpoint{1.558204in}{0.806759in}}%
\pgfpathlineto{\pgfqpoint{1.558226in}{0.823258in}}%
\pgfpathlineto{\pgfqpoint{1.558560in}{0.798510in}}%
\pgfpathlineto{\pgfqpoint{1.559226in}{0.798510in}}%
\pgfpathlineto{\pgfqpoint{1.560574in}{0.848005in}}%
\pgfpathlineto{\pgfqpoint{1.560885in}{0.831507in}}%
\pgfpathlineto{\pgfqpoint{1.562077in}{0.798510in}}%
\pgfpathlineto{\pgfqpoint{1.562728in}{0.831507in}}%
\pgfpathlineto{\pgfqpoint{1.563232in}{0.806759in}}%
\pgfpathlineto{\pgfqpoint{1.564239in}{0.782012in}}%
\pgfpathlineto{\pgfqpoint{1.564298in}{0.806759in}}%
\pgfpathlineto{\pgfqpoint{1.565579in}{0.782012in}}%
\pgfpathlineto{\pgfqpoint{1.564594in}{0.823258in}}%
\pgfpathlineto{\pgfqpoint{1.565579in}{0.790261in}}%
\pgfpathlineto{\pgfqpoint{1.565601in}{0.823258in}}%
\pgfpathlineto{\pgfqpoint{1.566608in}{0.798510in}}%
\pgfpathlineto{\pgfqpoint{1.567089in}{0.831507in}}%
\pgfpathlineto{\pgfqpoint{1.567763in}{0.806759in}}%
\pgfpathlineto{\pgfqpoint{1.568452in}{0.798510in}}%
\pgfpathlineto{\pgfqpoint{1.568615in}{0.823258in}}%
\pgfpathlineto{\pgfqpoint{1.568785in}{0.815008in}}%
\pgfpathlineto{\pgfqpoint{1.569792in}{0.798510in}}%
\pgfpathlineto{\pgfqpoint{1.569118in}{0.823258in}}%
\pgfpathlineto{\pgfqpoint{1.569940in}{0.806759in}}%
\pgfpathlineto{\pgfqpoint{1.570110in}{0.831507in}}%
\pgfpathlineto{\pgfqpoint{1.570962in}{0.790261in}}%
\pgfpathlineto{\pgfqpoint{1.571110in}{0.782012in}}%
\pgfpathlineto{\pgfqpoint{1.571465in}{0.823258in}}%
\pgfpathlineto{\pgfqpoint{1.572472in}{0.848005in}}%
\pgfpathlineto{\pgfqpoint{1.572642in}{0.839756in}}%
\pgfpathlineto{\pgfqpoint{1.572791in}{0.831507in}}%
\pgfpathlineto{\pgfqpoint{1.572953in}{0.856254in}}%
\pgfpathlineto{\pgfqpoint{1.573479in}{0.848005in}}%
\pgfpathlineto{\pgfqpoint{1.575471in}{0.905749in}}%
\pgfpathlineto{\pgfqpoint{1.576663in}{0.872753in}}%
\pgfpathlineto{\pgfqpoint{1.578677in}{0.946995in}}%
\pgfpathlineto{\pgfqpoint{1.578840in}{0.938746in}}%
\pgfpathlineto{\pgfqpoint{1.580017in}{0.922247in}}%
\pgfpathlineto{\pgfqpoint{1.581335in}{0.955244in}}%
\pgfpathlineto{\pgfqpoint{1.581357in}{0.946995in}}%
\pgfpathlineto{\pgfqpoint{1.581861in}{0.922247in}}%
\pgfpathlineto{\pgfqpoint{1.582683in}{0.930497in}}%
\pgfpathlineto{\pgfqpoint{1.583371in}{0.996490in}}%
\pgfpathlineto{\pgfqpoint{1.583704in}{0.996490in}}%
\pgfpathlineto{\pgfqpoint{1.584541in}{0.971742in}}%
\pgfpathlineto{\pgfqpoint{1.585022in}{0.979992in}}%
\pgfpathlineto{\pgfqpoint{1.585193in}{1.004739in}}%
\pgfpathlineto{\pgfqpoint{1.586051in}{0.988241in}}%
\pgfpathlineto{\pgfqpoint{1.586200in}{0.979992in}}%
\pgfpathlineto{\pgfqpoint{1.586370in}{1.004739in}}%
\pgfpathlineto{\pgfqpoint{1.586718in}{0.996490in}}%
\pgfpathlineto{\pgfqpoint{1.587873in}{1.029487in}}%
\pgfpathlineto{\pgfqpoint{1.588065in}{1.012988in}}%
\pgfpathlineto{\pgfqpoint{1.588228in}{1.021237in}}%
\pgfpathlineto{\pgfqpoint{1.589213in}{1.004739in}}%
\pgfpathlineto{\pgfqpoint{1.589235in}{1.021237in}}%
\pgfpathlineto{\pgfqpoint{1.590072in}{0.971742in}}%
\pgfpathlineto{\pgfqpoint{1.590242in}{0.971742in}}%
\pgfpathlineto{\pgfqpoint{1.591227in}{1.029487in}}%
\pgfpathlineto{\pgfqpoint{1.591582in}{0.996490in}}%
\pgfpathlineto{\pgfqpoint{1.591901in}{0.979992in}}%
\pgfpathlineto{\pgfqpoint{1.592234in}{1.012988in}}%
\pgfpathlineto{\pgfqpoint{1.592589in}{1.045985in}}%
\pgfpathlineto{\pgfqpoint{1.593256in}{1.021237in}}%
\pgfpathlineto{\pgfqpoint{1.595270in}{1.120227in}}%
\pgfpathlineto{\pgfqpoint{1.595751in}{1.103729in}}%
\pgfpathlineto{\pgfqpoint{1.595943in}{1.095480in}}%
\pgfpathlineto{\pgfqpoint{1.596254in}{1.128476in}}%
\pgfpathlineto{\pgfqpoint{1.596780in}{1.111978in}}%
\pgfpathlineto{\pgfqpoint{1.597787in}{1.070732in}}%
\pgfpathlineto{\pgfqpoint{1.597950in}{1.095480in}}%
\pgfpathlineto{\pgfqpoint{1.598268in}{1.103729in}}%
\pgfpathlineto{\pgfqpoint{1.598602in}{1.078981in}}%
\pgfpathlineto{\pgfqpoint{1.599794in}{1.045985in}}%
\pgfpathlineto{\pgfqpoint{1.600786in}{1.078981in}}%
\pgfpathlineto{\pgfqpoint{1.600971in}{1.062483in}}%
\pgfpathlineto{\pgfqpoint{1.602141in}{1.021237in}}%
\pgfpathlineto{\pgfqpoint{1.602311in}{1.045985in}}%
\pgfpathlineto{\pgfqpoint{1.602644in}{1.070732in}}%
\pgfpathlineto{\pgfqpoint{1.603148in}{1.037736in}}%
\pgfpathlineto{\pgfqpoint{1.603970in}{1.029487in}}%
\pgfpathlineto{\pgfqpoint{1.603318in}{1.045985in}}%
\pgfpathlineto{\pgfqpoint{1.603970in}{1.037736in}}%
\pgfpathlineto{\pgfqpoint{1.604473in}{1.103729in}}%
\pgfpathlineto{\pgfqpoint{1.604991in}{1.062483in}}%
\pgfpathlineto{\pgfqpoint{1.605976in}{1.054234in}}%
\pgfpathlineto{\pgfqpoint{1.605310in}{1.078981in}}%
\pgfpathlineto{\pgfqpoint{1.605976in}{1.062483in}}%
\pgfpathlineto{\pgfqpoint{1.606983in}{1.078981in}}%
\pgfpathlineto{\pgfqpoint{1.606317in}{1.054234in}}%
\pgfpathlineto{\pgfqpoint{1.607005in}{1.070732in}}%
\pgfpathlineto{\pgfqpoint{1.607487in}{1.078981in}}%
\pgfpathlineto{\pgfqpoint{1.607672in}{1.062483in}}%
\pgfpathlineto{\pgfqpoint{1.608160in}{1.054234in}}%
\pgfpathlineto{\pgfqpoint{1.608323in}{1.078981in}}%
\pgfpathlineto{\pgfqpoint{1.608345in}{1.070732in}}%
\pgfpathlineto{\pgfqpoint{1.609671in}{1.103729in}}%
\pgfpathlineto{\pgfqpoint{1.609686in}{1.095480in}}%
\pgfpathlineto{\pgfqpoint{1.611514in}{1.029487in}}%
\pgfpathlineto{\pgfqpoint{1.611529in}{1.045985in}}%
\pgfpathlineto{\pgfqpoint{1.612366in}{1.070732in}}%
\pgfpathlineto{\pgfqpoint{1.612684in}{1.054234in}}%
\pgfpathlineto{\pgfqpoint{1.613691in}{0.979992in}}%
\pgfpathlineto{\pgfqpoint{1.613876in}{0.996490in}}%
\pgfpathlineto{\pgfqpoint{1.614358in}{1.029487in}}%
\pgfpathlineto{\pgfqpoint{1.615217in}{1.012988in}}%
\pgfpathlineto{\pgfqpoint{1.615557in}{0.996490in}}%
\pgfpathlineto{\pgfqpoint{1.615868in}{1.029487in}}%
\pgfpathlineto{\pgfqpoint{1.615890in}{1.021237in}}%
\pgfpathlineto{\pgfqpoint{1.617045in}{1.029487in}}%
\pgfpathlineto{\pgfqpoint{1.618067in}{1.021237in}}%
\pgfpathlineto{\pgfqpoint{1.619052in}{1.078981in}}%
\pgfpathlineto{\pgfqpoint{1.619911in}{1.062483in}}%
\pgfpathlineto{\pgfqpoint{1.620059in}{1.054234in}}%
\pgfpathlineto{\pgfqpoint{1.620414in}{1.095480in}}%
\pgfpathlineto{\pgfqpoint{1.621569in}{1.103729in}}%
\pgfpathlineto{\pgfqpoint{1.621903in}{1.078981in}}%
\pgfpathlineto{\pgfqpoint{1.622428in}{1.120227in}}%
\pgfpathlineto{\pgfqpoint{1.622591in}{1.111978in}}%
\pgfpathlineto{\pgfqpoint{1.624272in}{1.045985in}}%
\pgfpathlineto{\pgfqpoint{1.624435in}{1.070732in}}%
\pgfpathlineto{\pgfqpoint{1.625257in}{1.103729in}}%
\pgfpathlineto{\pgfqpoint{1.625590in}{1.078981in}}%
\pgfpathlineto{\pgfqpoint{1.626767in}{1.054234in}}%
\pgfpathlineto{\pgfqpoint{1.627789in}{1.120227in}}%
\pgfpathlineto{\pgfqpoint{1.629114in}{1.103729in}}%
\pgfpathlineto{\pgfqpoint{1.629277in}{1.128476in}}%
\pgfpathlineto{\pgfqpoint{1.629973in}{1.095480in}}%
\pgfpathlineto{\pgfqpoint{1.630136in}{1.095480in}}%
\pgfpathlineto{\pgfqpoint{1.630625in}{1.128476in}}%
\pgfpathlineto{\pgfqpoint{1.631291in}{1.103729in}}%
\pgfpathlineto{\pgfqpoint{1.631817in}{1.070732in}}%
\pgfpathlineto{\pgfqpoint{1.632320in}{1.095480in}}%
\pgfpathlineto{\pgfqpoint{1.633809in}{1.177971in}}%
\pgfpathlineto{\pgfqpoint{1.634327in}{1.161473in}}%
\pgfpathlineto{\pgfqpoint{1.635312in}{1.128476in}}%
\pgfpathlineto{\pgfqpoint{1.635504in}{1.144975in}}%
\pgfpathlineto{\pgfqpoint{1.636659in}{1.153224in}}%
\pgfpathlineto{\pgfqpoint{1.637851in}{1.095480in}}%
\pgfpathlineto{\pgfqpoint{1.638518in}{1.120227in}}%
\pgfpathlineto{\pgfqpoint{1.639006in}{1.103729in}}%
\pgfpathlineto{\pgfqpoint{1.640028in}{1.095480in}}%
\pgfpathlineto{\pgfqpoint{1.640176in}{1.103729in}}%
\pgfpathlineto{\pgfqpoint{1.640532in}{1.070732in}}%
\pgfpathlineto{\pgfqpoint{1.640865in}{1.087231in}}%
\pgfpathlineto{\pgfqpoint{1.642190in}{1.054234in}}%
\pgfpathlineto{\pgfqpoint{1.641538in}{1.095480in}}%
\pgfpathlineto{\pgfqpoint{1.642190in}{1.062483in}}%
\pgfpathlineto{\pgfqpoint{1.642205in}{1.095480in}}%
\pgfpathlineto{\pgfqpoint{1.643212in}{1.062483in}}%
\pgfpathlineto{\pgfqpoint{1.644367in}{1.029487in}}%
\pgfpathlineto{\pgfqpoint{1.644389in}{1.045985in}}%
\pgfpathlineto{\pgfqpoint{1.644893in}{1.070732in}}%
\pgfpathlineto{\pgfqpoint{1.645226in}{1.045985in}}%
\pgfpathlineto{\pgfqpoint{1.645559in}{1.021237in}}%
\pgfpathlineto{\pgfqpoint{1.646381in}{1.037736in}}%
\pgfpathlineto{\pgfqpoint{1.646936in}{1.004739in}}%
\pgfpathlineto{\pgfqpoint{1.647403in}{1.045985in}}%
\pgfpathlineto{\pgfqpoint{1.647943in}{0.979992in}}%
\pgfpathlineto{\pgfqpoint{1.648898in}{1.012988in}}%
\pgfpathlineto{\pgfqpoint{1.648913in}{1.045985in}}%
\pgfpathlineto{\pgfqpoint{1.649750in}{0.996490in}}%
\pgfpathlineto{\pgfqpoint{1.649920in}{0.996490in}}%
\pgfpathlineto{\pgfqpoint{1.650927in}{1.021237in}}%
\pgfpathlineto{\pgfqpoint{1.651090in}{1.012988in}}%
\pgfpathlineto{\pgfqpoint{1.653755in}{0.905749in}}%
\pgfpathlineto{\pgfqpoint{1.651260in}{1.021237in}}%
\pgfpathlineto{\pgfqpoint{1.653770in}{0.922247in}}%
\pgfpathlineto{\pgfqpoint{1.654259in}{0.930497in}}%
\pgfpathlineto{\pgfqpoint{1.654444in}{0.913998in}}%
\pgfpathlineto{\pgfqpoint{1.655784in}{0.872753in}}%
\pgfpathlineto{\pgfqpoint{1.656265in}{0.930497in}}%
\pgfpathlineto{\pgfqpoint{1.656791in}{0.872753in}}%
\pgfpathlineto{\pgfqpoint{1.657443in}{0.856254in}}%
\pgfpathlineto{\pgfqpoint{1.657613in}{0.881002in}}%
\pgfpathlineto{\pgfqpoint{1.657628in}{0.872753in}}%
\pgfpathlineto{\pgfqpoint{1.658783in}{0.905749in}}%
\pgfpathlineto{\pgfqpoint{1.658968in}{0.889251in}}%
\pgfpathlineto{\pgfqpoint{1.660649in}{0.848005in}}%
\pgfpathlineto{\pgfqpoint{1.661130in}{0.856254in}}%
\pgfpathlineto{\pgfqpoint{1.661463in}{0.831507in}}%
\pgfpathlineto{\pgfqpoint{1.661485in}{0.848005in}}%
\pgfpathlineto{\pgfqpoint{1.662826in}{0.798510in}}%
\pgfpathlineto{\pgfqpoint{1.662996in}{0.823258in}}%
\pgfpathlineto{\pgfqpoint{1.663151in}{0.831507in}}%
\pgfpathlineto{\pgfqpoint{1.663499in}{0.798510in}}%
\pgfpathlineto{\pgfqpoint{1.663647in}{0.806759in}}%
\pgfpathlineto{\pgfqpoint{1.664336in}{0.757264in}}%
\pgfpathlineto{\pgfqpoint{1.664669in}{0.782012in}}%
\pgfpathlineto{\pgfqpoint{1.665010in}{0.806759in}}%
\pgfpathlineto{\pgfqpoint{1.665506in}{0.773763in}}%
\pgfpathlineto{\pgfqpoint{1.666009in}{0.757264in}}%
\pgfpathlineto{\pgfqpoint{1.666172in}{0.790261in}}%
\pgfpathlineto{\pgfqpoint{1.667016in}{0.831507in}}%
\pgfpathlineto{\pgfqpoint{1.666350in}{0.782012in}}%
\pgfpathlineto{\pgfqpoint{1.667187in}{0.823258in}}%
\pgfpathlineto{\pgfqpoint{1.668697in}{0.782012in}}%
\pgfpathlineto{\pgfqpoint{1.668838in}{0.790261in}}%
\pgfpathlineto{\pgfqpoint{1.669704in}{0.782012in}}%
\pgfpathlineto{\pgfqpoint{1.669867in}{0.823258in}}%
\pgfpathlineto{\pgfqpoint{1.670874in}{0.749015in}}%
\pgfpathlineto{\pgfqpoint{1.671377in}{0.798510in}}%
\pgfpathlineto{\pgfqpoint{1.671866in}{0.806759in}}%
\pgfpathlineto{\pgfqpoint{1.672044in}{0.790261in}}%
\pgfpathlineto{\pgfqpoint{1.673725in}{0.749015in}}%
\pgfpathlineto{\pgfqpoint{1.673887in}{0.773763in}}%
\pgfpathlineto{\pgfqpoint{1.674391in}{0.740766in}}%
\pgfpathlineto{\pgfqpoint{1.674561in}{0.749015in}}%
\pgfpathlineto{\pgfqpoint{1.675553in}{0.732517in}}%
\pgfpathlineto{\pgfqpoint{1.676568in}{0.773763in}}%
\pgfpathlineto{\pgfqpoint{1.676908in}{0.823258in}}%
\pgfpathlineto{\pgfqpoint{1.677915in}{0.798510in}}%
\pgfpathlineto{\pgfqpoint{1.678234in}{0.782012in}}%
\pgfpathlineto{\pgfqpoint{1.678582in}{0.815008in}}%
\pgfpathlineto{\pgfqpoint{1.679070in}{0.856254in}}%
\pgfpathlineto{\pgfqpoint{1.679589in}{0.848005in}}%
\pgfpathlineto{\pgfqpoint{1.681417in}{0.806759in}}%
\pgfpathlineto{\pgfqpoint{1.682439in}{0.856254in}}%
\pgfpathlineto{\pgfqpoint{1.682610in}{0.848005in}}%
\pgfpathlineto{\pgfqpoint{1.682772in}{0.889251in}}%
\pgfpathlineto{\pgfqpoint{1.683609in}{0.946995in}}%
\pgfpathlineto{\pgfqpoint{1.683779in}{0.930497in}}%
\pgfpathlineto{\pgfqpoint{1.684283in}{0.996490in}}%
\pgfpathlineto{\pgfqpoint{1.685290in}{0.963493in}}%
\pgfpathlineto{\pgfqpoint{1.688644in}{0.831507in}}%
\pgfpathlineto{\pgfqpoint{1.685453in}{0.971742in}}%
\pgfpathlineto{\pgfqpoint{1.688807in}{0.839756in}}%
\pgfpathlineto{\pgfqpoint{1.688807in}{0.848005in}}%
\pgfpathlineto{\pgfqpoint{1.689644in}{0.798510in}}%
\pgfpathlineto{\pgfqpoint{1.689814in}{0.798510in}}%
\pgfpathlineto{\pgfqpoint{1.693842in}{0.633527in}}%
\pgfpathlineto{\pgfqpoint{1.693997in}{0.641776in}}%
\pgfpathlineto{\pgfqpoint{1.694005in}{0.650025in}}%
\pgfpathlineto{\pgfqpoint{1.694338in}{0.625278in}}%
\pgfpathlineto{\pgfqpoint{1.695004in}{0.625278in}}%
\pgfpathlineto{\pgfqpoint{1.695833in}{0.633527in}}%
\pgfpathlineto{\pgfqpoint{1.695833in}{0.625278in}}%
\pgfpathlineto{\pgfqpoint{1.697011in}{0.633527in}}%
\pgfpathlineto{\pgfqpoint{1.697529in}{0.625278in}}%
\pgfpathlineto{\pgfqpoint{1.697684in}{0.625278in}}%
\pgfpathlineto{\pgfqpoint{1.698869in}{0.633527in}}%
\pgfpathlineto{\pgfqpoint{1.699039in}{0.625278in}}%
\pgfpathlineto{\pgfqpoint{1.699861in}{0.625278in}}%
\pgfpathlineto{\pgfqpoint{1.700705in}{0.633527in}}%
\pgfpathlineto{\pgfqpoint{1.700705in}{0.625278in}}%
\pgfpathlineto{\pgfqpoint{1.701875in}{0.633527in}}%
\pgfpathlineto{\pgfqpoint{1.702038in}{0.625278in}}%
\pgfpathlineto{\pgfqpoint{1.702712in}{0.625278in}}%
\pgfpathlineto{\pgfqpoint{1.703889in}{0.633527in}}%
\pgfpathlineto{\pgfqpoint{1.704052in}{0.625278in}}%
\pgfpathlineto{\pgfqpoint{1.704726in}{0.625278in}}%
\pgfpathlineto{\pgfqpoint{1.705911in}{0.633527in}}%
\pgfpathlineto{\pgfqpoint{1.706244in}{0.625278in}}%
\pgfpathlineto{\pgfqpoint{1.706903in}{0.625278in}}%
\pgfpathlineto{\pgfqpoint{1.708065in}{0.633527in}}%
\pgfpathlineto{\pgfqpoint{1.708406in}{0.625278in}}%
\pgfpathlineto{\pgfqpoint{1.709080in}{0.625278in}}%
\pgfpathlineto{\pgfqpoint{1.710257in}{0.633527in}}%
\pgfpathlineto{\pgfqpoint{1.710420in}{0.625278in}}%
\pgfpathlineto{\pgfqpoint{1.711093in}{0.625278in}}%
\pgfpathlineto{\pgfqpoint{1.712271in}{0.633527in}}%
\pgfpathlineto{\pgfqpoint{1.712426in}{0.625278in}}%
\pgfpathlineto{\pgfqpoint{1.713285in}{0.625278in}}%
\pgfpathlineto{\pgfqpoint{1.714440in}{0.633527in}}%
\pgfpathlineto{\pgfqpoint{1.714455in}{0.625278in}}%
\pgfpathlineto{\pgfqpoint{1.715462in}{0.625278in}}%
\pgfpathlineto{\pgfqpoint{1.716624in}{0.633527in}}%
\pgfpathlineto{\pgfqpoint{1.716787in}{0.625278in}}%
\pgfpathlineto{\pgfqpoint{1.717631in}{0.625278in}}%
\pgfpathlineto{\pgfqpoint{1.718794in}{0.633527in}}%
\pgfpathlineto{\pgfqpoint{1.719127in}{0.625278in}}%
\pgfpathlineto{\pgfqpoint{1.719638in}{0.625278in}}%
\pgfpathlineto{\pgfqpoint{1.720808in}{0.633527in}}%
\pgfpathlineto{\pgfqpoint{1.720978in}{0.625278in}}%
\pgfpathlineto{\pgfqpoint{1.721652in}{0.625278in}}%
\pgfpathlineto{\pgfqpoint{1.722829in}{0.633527in}}%
\pgfpathlineto{\pgfqpoint{1.722985in}{0.625278in}}%
\pgfpathlineto{\pgfqpoint{1.723829in}{0.625278in}}%
\pgfpathlineto{\pgfqpoint{1.724836in}{0.633527in}}%
\pgfpathlineto{\pgfqpoint{1.724836in}{0.625278in}}%
\pgfpathlineto{\pgfqpoint{1.726020in}{0.633527in}}%
\pgfpathlineto{\pgfqpoint{1.726354in}{0.625278in}}%
\pgfpathlineto{\pgfqpoint{1.726842in}{0.625278in}}%
\pgfpathlineto{\pgfqpoint{1.728012in}{0.633527in}}%
\pgfpathlineto{\pgfqpoint{1.728182in}{0.625278in}}%
\pgfpathlineto{\pgfqpoint{1.729019in}{0.625278in}}%
\pgfpathlineto{\pgfqpoint{1.730359in}{0.633527in}}%
\pgfpathlineto{\pgfqpoint{1.730367in}{0.625278in}}%
\pgfpathlineto{\pgfqpoint{1.731366in}{0.625278in}}%
\pgfpathlineto{\pgfqpoint{1.732543in}{0.633527in}}%
\pgfpathlineto{\pgfqpoint{1.732543in}{0.625278in}}%
\pgfpathlineto{\pgfqpoint{1.733550in}{0.625278in}}%
\pgfpathlineto{\pgfqpoint{1.734720in}{0.633527in}}%
\pgfpathlineto{\pgfqpoint{1.734891in}{0.625278in}}%
\pgfpathlineto{\pgfqpoint{1.735557in}{0.625278in}}%
\pgfpathlineto{\pgfqpoint{1.736897in}{0.633527in}}%
\pgfpathlineto{\pgfqpoint{1.736905in}{0.625278in}}%
\pgfpathlineto{\pgfqpoint{1.737904in}{0.625278in}}%
\pgfpathlineto{\pgfqpoint{1.739089in}{0.633527in}}%
\pgfpathlineto{\pgfqpoint{1.739244in}{0.625278in}}%
\pgfpathlineto{\pgfqpoint{1.739918in}{0.625278in}}%
\pgfpathlineto{\pgfqpoint{1.741088in}{0.633527in}}%
\pgfpathlineto{\pgfqpoint{1.741428in}{0.625278in}}%
\pgfpathlineto{\pgfqpoint{1.742102in}{0.625278in}}%
\pgfpathlineto{\pgfqpoint{1.743280in}{0.633527in}}%
\pgfpathlineto{\pgfqpoint{1.743435in}{0.625278in}}%
\pgfpathlineto{\pgfqpoint{1.744279in}{0.625278in}}%
\pgfpathlineto{\pgfqpoint{1.745449in}{0.633527in}}%
\pgfpathlineto{\pgfqpoint{1.745619in}{0.625278in}}%
\pgfpathlineto{\pgfqpoint{1.746286in}{0.625278in}}%
\pgfpathlineto{\pgfqpoint{1.747626in}{0.633527in}}%
\pgfpathlineto{\pgfqpoint{1.747626in}{0.625278in}}%
\pgfpathlineto{\pgfqpoint{1.748633in}{0.625278in}}%
\pgfpathlineto{\pgfqpoint{1.749973in}{0.633527in}}%
\pgfpathlineto{\pgfqpoint{1.750314in}{0.625278in}}%
\pgfpathlineto{\pgfqpoint{1.750980in}{0.625278in}}%
\pgfpathlineto{\pgfqpoint{1.752165in}{0.633527in}}%
\pgfpathlineto{\pgfqpoint{1.752320in}{0.625278in}}%
\pgfpathlineto{\pgfqpoint{1.753157in}{0.625278in}}%
\pgfpathlineto{\pgfqpoint{1.754334in}{0.633527in}}%
\pgfpathlineto{\pgfqpoint{1.754667in}{0.625278in}}%
\pgfpathlineto{\pgfqpoint{1.755178in}{0.625278in}}%
\pgfpathlineto{\pgfqpoint{1.756348in}{0.633527in}}%
\pgfpathlineto{\pgfqpoint{1.756681in}{0.625278in}}%
\pgfpathlineto{\pgfqpoint{1.757347in}{0.625278in}}%
\pgfpathlineto{\pgfqpoint{1.758695in}{0.633527in}}%
\pgfpathlineto{\pgfqpoint{1.758858in}{0.625278in}}%
\pgfpathlineto{\pgfqpoint{1.759695in}{0.625278in}}%
\pgfpathlineto{\pgfqpoint{1.760872in}{0.633527in}}%
\pgfpathlineto{\pgfqpoint{1.760872in}{0.625278in}}%
\pgfpathlineto{\pgfqpoint{1.761879in}{0.625278in}}%
\pgfpathlineto{\pgfqpoint{1.763049in}{0.633527in}}%
\pgfpathlineto{\pgfqpoint{1.763219in}{0.625278in}}%
\pgfpathlineto{\pgfqpoint{1.763885in}{0.625278in}}%
\pgfpathlineto{\pgfqpoint{1.765063in}{0.633527in}}%
\pgfpathlineto{\pgfqpoint{1.765396in}{0.625278in}}%
\pgfpathlineto{\pgfqpoint{1.766070in}{0.625278in}}%
\pgfpathlineto{\pgfqpoint{1.767410in}{0.633527in}}%
\pgfpathlineto{\pgfqpoint{1.767417in}{0.625278in}}%
\pgfpathlineto{\pgfqpoint{1.768417in}{0.625278in}}%
\pgfpathlineto{\pgfqpoint{1.769587in}{0.633527in}}%
\pgfpathlineto{\pgfqpoint{1.769587in}{0.625278in}}%
\pgfpathlineto{\pgfqpoint{1.770423in}{0.625278in}}%
\pgfpathlineto{\pgfqpoint{1.771601in}{0.633527in}}%
\pgfpathlineto{\pgfqpoint{1.771934in}{0.625278in}}%
\pgfpathlineto{\pgfqpoint{1.772608in}{0.625278in}}%
\pgfpathlineto{\pgfqpoint{1.773777in}{0.633527in}}%
\pgfpathlineto{\pgfqpoint{1.773777in}{0.625278in}}%
\pgfpathlineto{\pgfqpoint{1.774614in}{0.625278in}}%
\pgfpathlineto{\pgfqpoint{1.775791in}{0.633527in}}%
\pgfpathlineto{\pgfqpoint{1.775954in}{0.625278in}}%
\pgfpathlineto{\pgfqpoint{1.776798in}{0.625278in}}%
\pgfpathlineto{\pgfqpoint{1.777968in}{0.633527in}}%
\pgfpathlineto{\pgfqpoint{1.778138in}{0.625278in}}%
\pgfpathlineto{\pgfqpoint{1.778975in}{0.625278in}}%
\pgfpathlineto{\pgfqpoint{1.780315in}{0.633527in}}%
\pgfpathlineto{\pgfqpoint{1.780315in}{0.625278in}}%
\pgfpathlineto{\pgfqpoint{1.781322in}{0.625278in}}%
\pgfpathlineto{\pgfqpoint{1.782492in}{0.633527in}}%
\pgfpathlineto{\pgfqpoint{1.782662in}{0.625278in}}%
\pgfpathlineto{\pgfqpoint{1.783336in}{0.625278in}}%
\pgfpathlineto{\pgfqpoint{1.784506in}{0.633527in}}%
\pgfpathlineto{\pgfqpoint{1.784676in}{0.625278in}}%
\pgfpathlineto{\pgfqpoint{1.785343in}{0.625278in}}%
\pgfpathlineto{\pgfqpoint{1.786520in}{0.633527in}}%
\pgfpathlineto{\pgfqpoint{1.786683in}{0.625278in}}%
\pgfpathlineto{\pgfqpoint{1.787527in}{0.625278in}}%
\pgfpathlineto{\pgfqpoint{1.788697in}{0.633527in}}%
\pgfpathlineto{\pgfqpoint{1.788697in}{0.625278in}}%
\pgfpathlineto{\pgfqpoint{1.789534in}{0.625278in}}%
\pgfpathlineto{\pgfqpoint{1.790874in}{0.633527in}}%
\pgfpathlineto{\pgfqpoint{1.791044in}{0.625278in}}%
\pgfpathlineto{\pgfqpoint{1.791881in}{0.625278in}}%
\pgfpathlineto{\pgfqpoint{1.793058in}{0.633527in}}%
\pgfpathlineto{\pgfqpoint{1.793058in}{0.625278in}}%
\pgfpathlineto{\pgfqpoint{1.794065in}{0.625278in}}%
\pgfpathlineto{\pgfqpoint{1.795235in}{0.633527in}}%
\pgfpathlineto{\pgfqpoint{1.795235in}{0.625278in}}%
\pgfpathlineto{\pgfqpoint{1.796242in}{0.625278in}}%
\pgfpathlineto{\pgfqpoint{1.797412in}{0.633527in}}%
\pgfpathlineto{\pgfqpoint{1.797582in}{0.625278in}}%
\pgfpathlineto{\pgfqpoint{1.798419in}{0.625278in}}%
\pgfpathlineto{\pgfqpoint{1.799596in}{0.633527in}}%
\pgfpathlineto{\pgfqpoint{1.799929in}{0.625278in}}%
\pgfpathlineto{\pgfqpoint{1.800603in}{0.625278in}}%
\pgfpathlineto{\pgfqpoint{1.801773in}{0.633527in}}%
\pgfpathlineto{\pgfqpoint{1.801943in}{0.625278in}}%
\pgfpathlineto{\pgfqpoint{1.802609in}{0.625278in}}%
\pgfpathlineto{\pgfqpoint{1.803787in}{0.633527in}}%
\pgfpathlineto{\pgfqpoint{1.803949in}{0.625278in}}%
\pgfpathlineto{\pgfqpoint{1.804794in}{0.625278in}}%
\pgfpathlineto{\pgfqpoint{1.805963in}{0.633527in}}%
\pgfpathlineto{\pgfqpoint{1.805963in}{0.625278in}}%
\pgfpathlineto{\pgfqpoint{1.806970in}{0.625278in}}%
\pgfpathlineto{\pgfqpoint{1.808140in}{0.633527in}}%
\pgfpathlineto{\pgfqpoint{1.808140in}{0.625278in}}%
\pgfpathlineto{\pgfqpoint{1.809147in}{0.625278in}}%
\pgfpathlineto{\pgfqpoint{1.810154in}{0.633527in}}%
\pgfpathlineto{\pgfqpoint{1.810154in}{0.625278in}}%
\pgfpathlineto{\pgfqpoint{1.811494in}{0.633527in}}%
\pgfpathlineto{\pgfqpoint{1.811494in}{0.625278in}}%
\pgfpathlineto{\pgfqpoint{1.812501in}{0.625278in}}%
\pgfpathlineto{\pgfqpoint{1.813671in}{0.633527in}}%
\pgfpathlineto{\pgfqpoint{1.813842in}{0.625278in}}%
\pgfpathlineto{\pgfqpoint{1.814678in}{0.625278in}}%
\pgfpathlineto{\pgfqpoint{1.815855in}{0.633527in}}%
\pgfpathlineto{\pgfqpoint{1.816018in}{0.625278in}}%
\pgfpathlineto{\pgfqpoint{1.816862in}{0.625278in}}%
\pgfpathlineto{\pgfqpoint{1.818032in}{0.633527in}}%
\pgfpathlineto{\pgfqpoint{1.818032in}{0.625278in}}%
\pgfpathlineto{\pgfqpoint{1.819039in}{0.625278in}}%
\pgfpathlineto{\pgfqpoint{1.820209in}{0.633527in}}%
\pgfpathlineto{\pgfqpoint{1.820379in}{0.625278in}}%
\pgfpathlineto{\pgfqpoint{1.821216in}{0.625278in}}%
\pgfpathlineto{\pgfqpoint{1.822393in}{0.633527in}}%
\pgfpathlineto{\pgfqpoint{1.822393in}{0.625278in}}%
\pgfpathlineto{\pgfqpoint{1.823400in}{0.625278in}}%
\pgfpathlineto{\pgfqpoint{1.824570in}{0.633527in}}%
\pgfpathlineto{\pgfqpoint{1.824740in}{0.625278in}}%
\pgfpathlineto{\pgfqpoint{1.825577in}{0.625278in}}%
\pgfpathlineto{\pgfqpoint{1.826747in}{0.633527in}}%
\pgfpathlineto{\pgfqpoint{1.826917in}{0.625278in}}%
\pgfpathlineto{\pgfqpoint{1.827754in}{0.625278in}}%
\pgfpathlineto{\pgfqpoint{1.828931in}{0.633527in}}%
\pgfpathlineto{\pgfqpoint{1.828931in}{0.625278in}}%
\pgfpathlineto{\pgfqpoint{1.829938in}{0.625278in}}%
\pgfpathlineto{\pgfqpoint{1.831108in}{0.633527in}}%
\pgfpathlineto{\pgfqpoint{1.831108in}{0.625278in}}%
\pgfpathlineto{\pgfqpoint{1.832115in}{0.625278in}}%
\pgfpathlineto{\pgfqpoint{1.833455in}{0.633527in}}%
\pgfpathlineto{\pgfqpoint{1.833626in}{0.625278in}}%
\pgfpathlineto{\pgfqpoint{1.834462in}{0.625278in}}%
\pgfpathlineto{\pgfqpoint{1.835469in}{0.633527in}}%
\pgfpathlineto{\pgfqpoint{1.835469in}{0.625278in}}%
\pgfpathlineto{\pgfqpoint{1.836639in}{0.633527in}}%
\pgfpathlineto{\pgfqpoint{1.836639in}{0.625278in}}%
\pgfpathlineto{\pgfqpoint{1.837646in}{0.625278in}}%
\pgfpathlineto{\pgfqpoint{1.838816in}{0.633527in}}%
\pgfpathlineto{\pgfqpoint{1.838986in}{0.625278in}}%
\pgfpathlineto{\pgfqpoint{1.839823in}{0.625278in}}%
\pgfpathlineto{\pgfqpoint{1.841000in}{0.633527in}}%
\pgfpathlineto{\pgfqpoint{1.841333in}{0.625278in}}%
\pgfpathlineto{\pgfqpoint{1.842007in}{0.625278in}}%
\pgfpathlineto{\pgfqpoint{1.843177in}{0.633527in}}%
\pgfpathlineto{\pgfqpoint{1.843177in}{0.625278in}}%
\pgfpathlineto{\pgfqpoint{1.844184in}{0.625278in}}%
\pgfpathlineto{\pgfqpoint{1.845354in}{0.633527in}}%
\pgfpathlineto{\pgfqpoint{1.845354in}{0.625278in}}%
\pgfpathlineto{\pgfqpoint{1.846361in}{0.625278in}}%
\pgfpathlineto{\pgfqpoint{1.847538in}{0.633527in}}%
\pgfpathlineto{\pgfqpoint{1.847701in}{0.625278in}}%
\pgfpathlineto{\pgfqpoint{1.848375in}{0.625278in}}%
\pgfpathlineto{\pgfqpoint{1.849715in}{0.633527in}}%
\pgfpathlineto{\pgfqpoint{1.849715in}{0.625278in}}%
\pgfpathlineto{\pgfqpoint{1.850722in}{0.625278in}}%
\pgfpathlineto{\pgfqpoint{1.851892in}{0.633527in}}%
\pgfpathlineto{\pgfqpoint{1.852062in}{0.625278in}}%
\pgfpathlineto{\pgfqpoint{1.852899in}{0.625278in}}%
\pgfpathlineto{\pgfqpoint{1.854076in}{0.633527in}}%
\pgfpathlineto{\pgfqpoint{1.854076in}{0.625278in}}%
\pgfpathlineto{\pgfqpoint{1.855083in}{0.625278in}}%
\pgfpathlineto{\pgfqpoint{1.856253in}{0.633527in}}%
\pgfpathlineto{\pgfqpoint{1.856586in}{0.625278in}}%
\pgfpathlineto{\pgfqpoint{1.857089in}{0.625278in}}%
\pgfpathlineto{\pgfqpoint{1.858430in}{0.633527in}}%
\pgfpathlineto{\pgfqpoint{1.858430in}{0.625278in}}%
\pgfpathlineto{\pgfqpoint{1.859437in}{0.625278in}}%
\pgfpathlineto{\pgfqpoint{1.860614in}{0.633527in}}%
\pgfpathlineto{\pgfqpoint{1.860777in}{0.625278in}}%
\pgfpathlineto{\pgfqpoint{1.861613in}{0.625278in}}%
\pgfpathlineto{\pgfqpoint{1.862791in}{0.633527in}}%
\pgfpathlineto{\pgfqpoint{1.862961in}{0.625278in}}%
\pgfpathlineto{\pgfqpoint{1.863798in}{0.625278in}}%
\pgfpathlineto{\pgfqpoint{1.864967in}{0.633527in}}%
\pgfpathlineto{\pgfqpoint{1.864967in}{0.625278in}}%
\pgfpathlineto{\pgfqpoint{1.865974in}{0.625278in}}%
\pgfpathlineto{\pgfqpoint{1.867152in}{0.633527in}}%
\pgfpathlineto{\pgfqpoint{1.867315in}{0.625278in}}%
\pgfpathlineto{\pgfqpoint{1.868151in}{0.625278in}}%
\pgfpathlineto{\pgfqpoint{1.869329in}{0.633527in}}%
\pgfpathlineto{\pgfqpoint{1.869329in}{0.625278in}}%
\pgfpathlineto{\pgfqpoint{1.870336in}{0.625278in}}%
\pgfpathlineto{\pgfqpoint{1.871505in}{0.633527in}}%
\pgfpathlineto{\pgfqpoint{1.871735in}{0.625278in}}%
\pgfpathlineto{\pgfqpoint{1.872535in}{0.625278in}}%
\pgfpathlineto{\pgfqpoint{1.873408in}{0.658274in}}%
\pgfpathlineto{\pgfqpoint{1.873704in}{0.641776in}}%
\pgfpathlineto{\pgfqpoint{1.874038in}{0.625278in}}%
\pgfpathlineto{\pgfqpoint{1.874356in}{0.658274in}}%
\pgfpathlineto{\pgfqpoint{1.874378in}{0.650025in}}%
\pgfpathlineto{\pgfqpoint{1.875030in}{0.683022in}}%
\pgfpathlineto{\pgfqpoint{1.875533in}{0.658274in}}%
\pgfpathlineto{\pgfqpoint{1.875718in}{0.650025in}}%
\pgfpathlineto{\pgfqpoint{1.875881in}{0.674773in}}%
\pgfpathlineto{\pgfqpoint{1.876555in}{0.674773in}}%
\pgfpathlineto{\pgfqpoint{1.877895in}{0.650025in}}%
\pgfpathlineto{\pgfqpoint{1.879050in}{0.658274in}}%
\pgfpathlineto{\pgfqpoint{1.880242in}{0.625278in}}%
\pgfpathlineto{\pgfqpoint{1.881568in}{0.683022in}}%
\pgfpathlineto{\pgfqpoint{1.881901in}{0.658274in}}%
\pgfpathlineto{\pgfqpoint{1.882589in}{0.625278in}}%
\pgfpathlineto{\pgfqpoint{1.882923in}{0.650025in}}%
\pgfpathlineto{\pgfqpoint{1.883767in}{0.674773in}}%
\pgfpathlineto{\pgfqpoint{1.884078in}{0.658274in}}%
\pgfpathlineto{\pgfqpoint{1.884937in}{0.625278in}}%
\pgfpathlineto{\pgfqpoint{1.885107in}{0.650025in}}%
\pgfpathlineto{\pgfqpoint{1.886262in}{0.658274in}}%
\pgfpathlineto{\pgfqpoint{1.886928in}{0.633527in}}%
\pgfpathlineto{\pgfqpoint{1.887284in}{0.650025in}}%
\pgfpathlineto{\pgfqpoint{1.887602in}{0.658274in}}%
\pgfpathlineto{\pgfqpoint{1.887935in}{0.633527in}}%
\pgfpathlineto{\pgfqpoint{1.888794in}{0.625278in}}%
\pgfpathlineto{\pgfqpoint{1.888957in}{0.650025in}}%
\pgfpathlineto{\pgfqpoint{1.889949in}{0.633527in}}%
\pgfpathlineto{\pgfqpoint{1.889964in}{0.650025in}}%
\pgfpathlineto{\pgfqpoint{1.891452in}{0.683022in}}%
\pgfpathlineto{\pgfqpoint{1.891474in}{0.674773in}}%
\pgfpathlineto{\pgfqpoint{1.891682in}{0.658274in}}%
\pgfpathlineto{\pgfqpoint{1.892126in}{0.691271in}}%
\pgfpathlineto{\pgfqpoint{1.892459in}{0.707769in}}%
\pgfpathlineto{\pgfqpoint{1.892800in}{0.683022in}}%
\pgfpathlineto{\pgfqpoint{1.893148in}{0.699520in}}%
\pgfpathlineto{\pgfqpoint{1.893651in}{0.674773in}}%
\pgfpathlineto{\pgfqpoint{1.893970in}{0.707769in}}%
\pgfpathlineto{\pgfqpoint{1.894029in}{0.707769in}}%
\pgfpathlineto{\pgfqpoint{1.896317in}{0.806759in}}%
\pgfpathlineto{\pgfqpoint{1.896339in}{0.798510in}}%
\pgfpathlineto{\pgfqpoint{1.896487in}{0.806759in}}%
\pgfpathlineto{\pgfqpoint{1.896835in}{0.773763in}}%
\pgfpathlineto{\pgfqpoint{1.897154in}{0.782012in}}%
\pgfpathlineto{\pgfqpoint{1.898183in}{0.773763in}}%
\pgfpathlineto{\pgfqpoint{1.899353in}{0.798510in}}%
\pgfpathlineto{\pgfqpoint{1.899523in}{0.790261in}}%
\pgfpathlineto{\pgfqpoint{1.899686in}{0.798510in}}%
\pgfpathlineto{\pgfqpoint{1.900678in}{0.782012in}}%
\pgfpathlineto{\pgfqpoint{1.901196in}{0.823258in}}%
\pgfpathlineto{\pgfqpoint{1.901700in}{0.815008in}}%
\pgfpathlineto{\pgfqpoint{1.902033in}{0.798510in}}%
\pgfpathlineto{\pgfqpoint{1.902351in}{0.831507in}}%
\pgfpathlineto{\pgfqpoint{1.902373in}{0.823258in}}%
\pgfpathlineto{\pgfqpoint{1.903543in}{0.872753in}}%
\pgfpathlineto{\pgfqpoint{1.904025in}{0.856254in}}%
\pgfpathlineto{\pgfqpoint{1.905032in}{0.831507in}}%
\pgfpathlineto{\pgfqpoint{1.905054in}{0.848005in}}%
\pgfpathlineto{\pgfqpoint{1.905535in}{0.856254in}}%
\pgfpathlineto{\pgfqpoint{1.905868in}{0.831507in}}%
\pgfpathlineto{\pgfqpoint{1.906224in}{0.823258in}}%
\pgfpathlineto{\pgfqpoint{1.906542in}{0.856254in}}%
\pgfpathlineto{\pgfqpoint{1.906897in}{0.848005in}}%
\pgfpathlineto{\pgfqpoint{1.908238in}{0.823258in}}%
\pgfpathlineto{\pgfqpoint{1.909563in}{0.856254in}}%
\pgfpathlineto{\pgfqpoint{1.909578in}{0.848005in}}%
\pgfpathlineto{\pgfqpoint{1.910229in}{0.806759in}}%
\pgfpathlineto{\pgfqpoint{1.910918in}{0.823258in}}%
\pgfpathlineto{\pgfqpoint{1.911407in}{0.831507in}}%
\pgfpathlineto{\pgfqpoint{1.911592in}{0.815008in}}%
\pgfpathlineto{\pgfqpoint{1.913413in}{0.757264in}}%
\pgfpathlineto{\pgfqpoint{1.913435in}{0.773763in}}%
\pgfpathlineto{\pgfqpoint{1.914087in}{0.806759in}}%
\pgfpathlineto{\pgfqpoint{1.914605in}{0.790261in}}%
\pgfpathlineto{\pgfqpoint{1.915427in}{0.757264in}}%
\pgfpathlineto{\pgfqpoint{1.914775in}{0.798510in}}%
\pgfpathlineto{\pgfqpoint{1.915760in}{0.782012in}}%
\pgfpathlineto{\pgfqpoint{1.915931in}{0.806759in}}%
\pgfpathlineto{\pgfqpoint{1.916286in}{0.773763in}}%
\pgfpathlineto{\pgfqpoint{1.916789in}{0.773763in}}%
\pgfpathlineto{\pgfqpoint{1.917774in}{0.732517in}}%
\pgfpathlineto{\pgfqpoint{1.918107in}{0.757264in}}%
\pgfpathlineto{\pgfqpoint{1.918130in}{0.773763in}}%
\pgfpathlineto{\pgfqpoint{1.919137in}{0.773763in}}%
\pgfpathlineto{\pgfqpoint{1.921128in}{0.831507in}}%
\pgfpathlineto{\pgfqpoint{1.921484in}{0.823258in}}%
\pgfpathlineto{\pgfqpoint{1.921980in}{0.872753in}}%
\pgfpathlineto{\pgfqpoint{1.922150in}{0.864503in}}%
\pgfpathlineto{\pgfqpoint{1.922320in}{0.872753in}}%
\pgfpathlineto{\pgfqpoint{1.923327in}{0.848005in}}%
\pgfpathlineto{\pgfqpoint{1.924645in}{0.881002in}}%
\pgfpathlineto{\pgfqpoint{1.924667in}{0.872753in}}%
\pgfpathlineto{\pgfqpoint{1.925171in}{0.848005in}}%
\pgfpathlineto{\pgfqpoint{1.925482in}{0.881002in}}%
\pgfpathlineto{\pgfqpoint{1.925504in}{0.872753in}}%
\pgfpathlineto{\pgfqpoint{1.925652in}{0.881002in}}%
\pgfpathlineto{\pgfqpoint{1.926008in}{0.848005in}}%
\pgfpathlineto{\pgfqpoint{1.926171in}{0.848005in}}%
\pgfpathlineto{\pgfqpoint{1.926844in}{0.872753in}}%
\pgfpathlineto{\pgfqpoint{1.927326in}{0.856254in}}%
\pgfpathlineto{\pgfqpoint{1.928503in}{0.831507in}}%
\pgfpathlineto{\pgfqpoint{1.928518in}{0.872753in}}%
\pgfpathlineto{\pgfqpoint{1.929525in}{0.872753in}}%
\pgfpathlineto{\pgfqpoint{1.930176in}{0.856254in}}%
\pgfpathlineto{\pgfqpoint{1.930176in}{0.889251in}}%
\pgfpathlineto{\pgfqpoint{1.930702in}{0.946995in}}%
\pgfpathlineto{\pgfqpoint{1.931368in}{0.922247in}}%
\pgfpathlineto{\pgfqpoint{1.932694in}{1.004739in}}%
\pgfpathlineto{\pgfqpoint{1.932708in}{0.996490in}}%
\pgfpathlineto{\pgfqpoint{1.933864in}{1.004739in}}%
\pgfpathlineto{\pgfqpoint{1.934056in}{0.996490in}}%
\pgfpathlineto{\pgfqpoint{1.934219in}{1.021237in}}%
\pgfpathlineto{\pgfqpoint{1.934893in}{1.021237in}}%
\pgfpathlineto{\pgfqpoint{1.936063in}{0.971742in}}%
\pgfpathlineto{\pgfqpoint{1.936233in}{0.996490in}}%
\pgfpathlineto{\pgfqpoint{1.937906in}{1.045985in}}%
\pgfpathlineto{\pgfqpoint{1.938076in}{1.037736in}}%
\pgfpathlineto{\pgfqpoint{1.938247in}{1.045985in}}%
\pgfpathlineto{\pgfqpoint{1.939232in}{1.029487in}}%
\pgfpathlineto{\pgfqpoint{1.939246in}{1.045985in}}%
\pgfpathlineto{\pgfqpoint{1.940253in}{1.045985in}}%
\pgfpathlineto{\pgfqpoint{1.941579in}{1.029487in}}%
\pgfpathlineto{\pgfqpoint{1.941934in}{1.070732in}}%
\pgfpathlineto{\pgfqpoint{1.942600in}{1.070732in}}%
\pgfpathlineto{\pgfqpoint{1.943607in}{1.045985in}}%
\pgfpathlineto{\pgfqpoint{1.943756in}{1.062483in}}%
\pgfpathlineto{\pgfqpoint{1.944762in}{1.078981in}}%
\pgfpathlineto{\pgfqpoint{1.944777in}{1.070732in}}%
\pgfpathlineto{\pgfqpoint{1.945932in}{1.078981in}}%
\pgfpathlineto{\pgfqpoint{1.946436in}{1.103729in}}%
\pgfpathlineto{\pgfqpoint{1.946962in}{1.070732in}}%
\pgfpathlineto{\pgfqpoint{1.947776in}{1.103729in}}%
\pgfpathlineto{\pgfqpoint{1.948117in}{1.078981in}}%
\pgfpathlineto{\pgfqpoint{1.948635in}{1.070732in}}%
\pgfpathlineto{\pgfqpoint{1.948805in}{1.095480in}}%
\pgfpathlineto{\pgfqpoint{1.949138in}{1.095480in}}%
\pgfpathlineto{\pgfqpoint{1.950293in}{1.103729in}}%
\pgfpathlineto{\pgfqpoint{1.951485in}{1.070732in}}%
\pgfpathlineto{\pgfqpoint{1.951967in}{1.103729in}}%
\pgfpathlineto{\pgfqpoint{1.952641in}{1.078981in}}%
\pgfpathlineto{\pgfqpoint{1.953329in}{1.070732in}}%
\pgfpathlineto{\pgfqpoint{1.953499in}{1.095480in}}%
\pgfpathlineto{\pgfqpoint{1.953662in}{1.087231in}}%
\pgfpathlineto{\pgfqpoint{1.954166in}{1.070732in}}%
\pgfpathlineto{\pgfqpoint{1.954336in}{1.095480in}}%
\pgfpathlineto{\pgfqpoint{1.956009in}{1.144975in}}%
\pgfpathlineto{\pgfqpoint{1.956180in}{1.136726in}}%
\pgfpathlineto{\pgfqpoint{1.958001in}{1.078981in}}%
\pgfpathlineto{\pgfqpoint{1.958023in}{1.095480in}}%
\pgfpathlineto{\pgfqpoint{1.958342in}{1.078981in}}%
\pgfpathlineto{\pgfqpoint{1.958357in}{1.120227in}}%
\pgfpathlineto{\pgfqpoint{1.958675in}{1.111978in}}%
\pgfpathlineto{\pgfqpoint{1.959682in}{1.153224in}}%
\pgfpathlineto{\pgfqpoint{1.959697in}{1.144975in}}%
\pgfpathlineto{\pgfqpoint{1.961044in}{1.120227in}}%
\pgfpathlineto{\pgfqpoint{1.963036in}{1.177971in}}%
\pgfpathlineto{\pgfqpoint{1.964058in}{1.144975in}}%
\pgfpathlineto{\pgfqpoint{1.964539in}{1.177971in}}%
\pgfpathlineto{\pgfqpoint{1.964894in}{1.144975in}}%
\pgfpathlineto{\pgfqpoint{1.965901in}{1.120227in}}%
\pgfpathlineto{\pgfqpoint{1.966050in}{1.136726in}}%
\pgfpathlineto{\pgfqpoint{1.967057in}{1.153224in}}%
\pgfpathlineto{\pgfqpoint{1.967079in}{1.144975in}}%
\pgfpathlineto{\pgfqpoint{1.968234in}{1.153224in}}%
\pgfpathlineto{\pgfqpoint{1.968922in}{1.144975in}}%
\pgfpathlineto{\pgfqpoint{1.969233in}{1.177971in}}%
\pgfpathlineto{\pgfqpoint{1.969256in}{1.169722in}}%
\pgfpathlineto{\pgfqpoint{1.969574in}{1.177971in}}%
\pgfpathlineto{\pgfqpoint{1.969907in}{1.153224in}}%
\pgfpathlineto{\pgfqpoint{1.970092in}{1.161473in}}%
\pgfpathlineto{\pgfqpoint{1.970240in}{1.153224in}}%
\pgfpathlineto{\pgfqpoint{1.970411in}{1.177971in}}%
\pgfpathlineto{\pgfqpoint{1.970766in}{1.169722in}}%
\pgfpathlineto{\pgfqpoint{1.970914in}{1.177971in}}%
\pgfpathlineto{\pgfqpoint{1.971269in}{1.144975in}}%
\pgfpathlineto{\pgfqpoint{1.971580in}{1.153224in}}%
\pgfpathlineto{\pgfqpoint{1.971773in}{1.144975in}}%
\pgfpathlineto{\pgfqpoint{1.971936in}{1.169722in}}%
\pgfpathlineto{\pgfqpoint{1.972610in}{1.169722in}}%
\pgfpathlineto{\pgfqpoint{1.973780in}{1.144975in}}%
\pgfpathlineto{\pgfqpoint{1.973928in}{1.153224in}}%
\pgfpathlineto{\pgfqpoint{1.974431in}{1.202719in}}%
\pgfpathlineto{\pgfqpoint{1.974957in}{1.194470in}}%
\pgfpathlineto{\pgfqpoint{1.975793in}{1.144975in}}%
\pgfpathlineto{\pgfqpoint{1.976615in}{1.161473in}}%
\pgfpathlineto{\pgfqpoint{1.977282in}{1.202719in}}%
\pgfpathlineto{\pgfqpoint{1.977637in}{1.169722in}}%
\pgfpathlineto{\pgfqpoint{1.977785in}{1.177971in}}%
\pgfpathlineto{\pgfqpoint{1.978141in}{1.144975in}}%
\pgfpathlineto{\pgfqpoint{1.978303in}{1.144975in}}%
\pgfpathlineto{\pgfqpoint{1.979799in}{1.177971in}}%
\pgfpathlineto{\pgfqpoint{1.979814in}{1.169722in}}%
\pgfpathlineto{\pgfqpoint{1.980991in}{1.120227in}}%
\pgfpathlineto{\pgfqpoint{1.981472in}{1.136726in}}%
\pgfpathlineto{\pgfqpoint{1.982650in}{1.177971in}}%
\pgfpathlineto{\pgfqpoint{1.982665in}{1.169722in}}%
\pgfpathlineto{\pgfqpoint{1.983205in}{1.128476in}}%
\pgfpathlineto{\pgfqpoint{1.983990in}{1.153224in}}%
\pgfpathlineto{\pgfqpoint{1.984005in}{1.169722in}}%
\pgfpathlineto{\pgfqpoint{1.985012in}{1.144975in}}%
\pgfpathlineto{\pgfqpoint{1.986167in}{1.153224in}}%
\pgfpathlineto{\pgfqpoint{1.987003in}{1.078981in}}%
\pgfpathlineto{\pgfqpoint{1.987359in}{1.095480in}}%
\pgfpathlineto{\pgfqpoint{1.988536in}{1.120227in}}%
\pgfpathlineto{\pgfqpoint{1.988699in}{1.111978in}}%
\pgfpathlineto{\pgfqpoint{1.989854in}{1.054234in}}%
\pgfpathlineto{\pgfqpoint{1.990380in}{1.070732in}}%
\pgfpathlineto{\pgfqpoint{1.991372in}{1.103729in}}%
\pgfpathlineto{\pgfqpoint{1.991550in}{1.087231in}}%
\pgfpathlineto{\pgfqpoint{1.991698in}{1.078981in}}%
\pgfpathlineto{\pgfqpoint{1.992053in}{1.120227in}}%
\pgfpathlineto{\pgfqpoint{1.992371in}{1.128476in}}%
\pgfpathlineto{\pgfqpoint{1.992557in}{1.111978in}}%
\pgfpathlineto{\pgfqpoint{1.994067in}{1.045985in}}%
\pgfpathlineto{\pgfqpoint{1.995718in}{1.004739in}}%
\pgfpathlineto{\pgfqpoint{1.995718in}{1.012988in}}%
\pgfpathlineto{\pgfqpoint{1.996074in}{1.021237in}}%
\pgfpathlineto{\pgfqpoint{1.996725in}{0.979992in}}%
\pgfpathlineto{\pgfqpoint{1.996747in}{0.996490in}}%
\pgfpathlineto{\pgfqpoint{1.997081in}{1.021237in}}%
\pgfpathlineto{\pgfqpoint{1.997902in}{1.004739in}}%
\pgfpathlineto{\pgfqpoint{1.999094in}{0.946995in}}%
\pgfpathlineto{\pgfqpoint{1.999243in}{0.955244in}}%
\pgfpathlineto{\pgfqpoint{1.999598in}{0.922247in}}%
\pgfpathlineto{\pgfqpoint{2.001945in}{0.848005in}}%
\pgfpathlineto{\pgfqpoint{2.002093in}{0.856254in}}%
\pgfpathlineto{\pgfqpoint{2.002449in}{0.823258in}}%
\pgfpathlineto{\pgfqpoint{2.002760in}{0.831507in}}%
\pgfpathlineto{\pgfqpoint{2.003952in}{0.790261in}}%
\pgfpathlineto{\pgfqpoint{2.004122in}{0.773763in}}%
\pgfpathlineto{\pgfqpoint{2.004292in}{0.798510in}}%
\pgfpathlineto{\pgfqpoint{2.004611in}{0.790261in}}%
\pgfpathlineto{\pgfqpoint{2.004625in}{0.823258in}}%
\pgfpathlineto{\pgfqpoint{2.005129in}{0.773763in}}%
\pgfpathlineto{\pgfqpoint{2.005632in}{0.823258in}}%
\pgfpathlineto{\pgfqpoint{2.006121in}{0.856254in}}%
\pgfpathlineto{\pgfqpoint{2.006973in}{0.839756in}}%
\pgfpathlineto{\pgfqpoint{2.008313in}{0.798510in}}%
\pgfpathlineto{\pgfqpoint{2.007143in}{0.848005in}}%
\pgfpathlineto{\pgfqpoint{2.008468in}{0.815008in}}%
\pgfpathlineto{\pgfqpoint{2.008979in}{0.856254in}}%
\pgfpathlineto{\pgfqpoint{2.009490in}{0.790261in}}%
\pgfpathlineto{\pgfqpoint{2.009645in}{0.782012in}}%
\pgfpathlineto{\pgfqpoint{2.009986in}{0.815008in}}%
\pgfpathlineto{\pgfqpoint{2.010660in}{0.848005in}}%
\pgfpathlineto{\pgfqpoint{2.010993in}{0.848005in}}%
\pgfpathlineto{\pgfqpoint{2.011496in}{0.823258in}}%
\pgfpathlineto{\pgfqpoint{2.011815in}{0.856254in}}%
\pgfpathlineto{\pgfqpoint{2.011830in}{0.848005in}}%
\pgfpathlineto{\pgfqpoint{2.012481in}{0.856254in}}%
\pgfpathlineto{\pgfqpoint{2.012674in}{0.839756in}}%
\pgfpathlineto{\pgfqpoint{2.014014in}{0.773763in}}%
\pgfpathlineto{\pgfqpoint{2.014851in}{0.798510in}}%
\pgfpathlineto{\pgfqpoint{2.014503in}{0.757264in}}%
\pgfpathlineto{\pgfqpoint{2.015169in}{0.782012in}}%
\pgfpathlineto{\pgfqpoint{2.015354in}{0.773763in}}%
\pgfpathlineto{\pgfqpoint{2.015680in}{0.806759in}}%
\pgfpathlineto{\pgfqpoint{2.016191in}{0.798510in}}%
\pgfpathlineto{\pgfqpoint{2.016346in}{0.806759in}}%
\pgfpathlineto{\pgfqpoint{2.016694in}{0.773763in}}%
\pgfpathlineto{\pgfqpoint{2.016850in}{0.782012in}}%
\pgfpathlineto{\pgfqpoint{2.016865in}{0.773763in}}%
\pgfpathlineto{\pgfqpoint{2.017531in}{0.823258in}}%
\pgfpathlineto{\pgfqpoint{2.017871in}{0.798510in}}%
\pgfpathlineto{\pgfqpoint{2.018190in}{0.782012in}}%
\pgfpathlineto{\pgfqpoint{2.018523in}{0.815008in}}%
\pgfpathlineto{\pgfqpoint{2.019582in}{0.905749in}}%
\pgfpathlineto{\pgfqpoint{2.019715in}{0.897500in}}%
\pgfpathlineto{\pgfqpoint{2.019878in}{0.946995in}}%
\pgfpathlineto{\pgfqpoint{2.020382in}{0.971742in}}%
\pgfpathlineto{\pgfqpoint{2.020715in}{0.938746in}}%
\pgfpathlineto{\pgfqpoint{2.021892in}{0.881002in}}%
\pgfpathlineto{\pgfqpoint{2.021055in}{0.946995in}}%
\pgfpathlineto{\pgfqpoint{2.022062in}{0.897500in}}%
\pgfpathlineto{\pgfqpoint{2.022395in}{0.872753in}}%
\pgfpathlineto{\pgfqpoint{2.022566in}{0.905749in}}%
\pgfpathlineto{\pgfqpoint{2.022884in}{0.905749in}}%
\pgfpathlineto{\pgfqpoint{2.023736in}{0.946995in}}%
\pgfpathlineto{\pgfqpoint{2.023906in}{0.938746in}}%
\pgfpathlineto{\pgfqpoint{2.025402in}{0.905749in}}%
\pgfpathlineto{\pgfqpoint{2.025402in}{0.913998in}}%
\pgfpathlineto{\pgfqpoint{2.025409in}{0.946995in}}%
\pgfpathlineto{\pgfqpoint{2.026238in}{0.905749in}}%
\pgfpathlineto{\pgfqpoint{2.026416in}{0.922247in}}%
\pgfpathlineto{\pgfqpoint{2.027401in}{0.905749in}}%
\pgfpathlineto{\pgfqpoint{2.027423in}{0.922247in}}%
\pgfpathlineto{\pgfqpoint{2.027423in}{0.930497in}}%
\pgfpathlineto{\pgfqpoint{2.028430in}{0.848005in}}%
\pgfpathlineto{\pgfqpoint{2.029940in}{0.806759in}}%
\pgfpathlineto{\pgfqpoint{2.030103in}{0.815008in}}%
\pgfpathlineto{\pgfqpoint{2.030103in}{0.823258in}}%
\pgfpathlineto{\pgfqpoint{2.030444in}{0.782012in}}%
\pgfpathlineto{\pgfqpoint{2.031110in}{0.798510in}}%
\pgfpathlineto{\pgfqpoint{2.035449in}{0.625278in}}%
\pgfpathlineto{\pgfqpoint{2.036626in}{0.633527in}}%
\pgfpathlineto{\pgfqpoint{2.036634in}{0.625278in}}%
\pgfpathlineto{\pgfqpoint{2.037648in}{0.625278in}}%
\pgfpathlineto{\pgfqpoint{2.038492in}{0.633527in}}%
\pgfpathlineto{\pgfqpoint{2.038492in}{0.625278in}}%
\pgfpathlineto{\pgfqpoint{2.039647in}{0.633527in}}%
\pgfpathlineto{\pgfqpoint{2.039832in}{0.625278in}}%
\pgfpathlineto{\pgfqpoint{2.040654in}{0.625278in}}%
\pgfpathlineto{\pgfqpoint{2.041839in}{0.633527in}}%
\pgfpathlineto{\pgfqpoint{2.042165in}{0.625278in}}%
\pgfpathlineto{\pgfqpoint{2.042513in}{0.625278in}}%
\pgfpathlineto{\pgfqpoint{2.043668in}{0.633527in}}%
\pgfpathlineto{\pgfqpoint{2.044001in}{0.625278in}}%
\pgfpathlineto{\pgfqpoint{2.044689in}{0.625278in}}%
\pgfpathlineto{\pgfqpoint{2.045852in}{0.633527in}}%
\pgfpathlineto{\pgfqpoint{2.045859in}{0.625278in}}%
\pgfpathlineto{\pgfqpoint{2.046689in}{0.625278in}}%
\pgfpathlineto{\pgfqpoint{2.047866in}{0.633527in}}%
\pgfpathlineto{\pgfqpoint{2.048044in}{0.625278in}}%
\pgfpathlineto{\pgfqpoint{2.048540in}{0.625278in}}%
\pgfpathlineto{\pgfqpoint{2.049710in}{0.633527in}}%
\pgfpathlineto{\pgfqpoint{2.049880in}{0.625278in}}%
\pgfpathlineto{\pgfqpoint{2.050724in}{0.625278in}}%
\pgfpathlineto{\pgfqpoint{2.051723in}{0.633527in}}%
\pgfpathlineto{\pgfqpoint{2.051723in}{0.625278in}}%
\pgfpathlineto{\pgfqpoint{2.052901in}{0.633527in}}%
\pgfpathlineto{\pgfqpoint{2.052901in}{0.625278in}}%
\pgfpathlineto{\pgfqpoint{2.053900in}{0.625278in}}%
\pgfpathlineto{\pgfqpoint{2.054915in}{0.633527in}}%
\pgfpathlineto{\pgfqpoint{2.054915in}{0.625278in}}%
\pgfpathlineto{\pgfqpoint{2.056240in}{0.633527in}}%
\pgfpathlineto{\pgfqpoint{2.056247in}{0.625278in}}%
\pgfpathlineto{\pgfqpoint{2.057262in}{0.625278in}}%
\pgfpathlineto{\pgfqpoint{2.058595in}{0.633527in}}%
\pgfpathlineto{\pgfqpoint{2.058750in}{0.625278in}}%
\pgfpathlineto{\pgfqpoint{2.059587in}{0.625278in}}%
\pgfpathlineto{\pgfqpoint{2.060601in}{0.633527in}}%
\pgfpathlineto{\pgfqpoint{2.060601in}{0.625278in}}%
\pgfpathlineto{\pgfqpoint{2.061771in}{0.633527in}}%
\pgfpathlineto{\pgfqpoint{2.061941in}{0.625278in}}%
\pgfpathlineto{\pgfqpoint{2.062615in}{0.625278in}}%
\pgfpathlineto{\pgfqpoint{2.063955in}{0.633527in}}%
\pgfpathlineto{\pgfqpoint{2.064466in}{0.625278in}}%
\pgfpathlineto{\pgfqpoint{2.064955in}{0.625278in}}%
\pgfpathlineto{\pgfqpoint{2.066147in}{0.633527in}}%
\pgfpathlineto{\pgfqpoint{2.066295in}{0.625278in}}%
\pgfpathlineto{\pgfqpoint{2.067139in}{0.625278in}}%
\pgfpathlineto{\pgfqpoint{2.068316in}{0.633527in}}%
\pgfpathlineto{\pgfqpoint{2.068324in}{0.625278in}}%
\pgfpathlineto{\pgfqpoint{2.069323in}{0.625278in}}%
\pgfpathlineto{\pgfqpoint{2.070486in}{0.633527in}}%
\pgfpathlineto{\pgfqpoint{2.070501in}{0.625278in}}%
\pgfpathlineto{\pgfqpoint{2.071500in}{0.625278in}}%
\pgfpathlineto{\pgfqpoint{2.072670in}{0.633527in}}%
\pgfpathlineto{\pgfqpoint{2.072833in}{0.625278in}}%
\pgfpathlineto{\pgfqpoint{2.073670in}{0.625278in}}%
\pgfpathlineto{\pgfqpoint{2.074847in}{0.633527in}}%
\pgfpathlineto{\pgfqpoint{2.075010in}{0.625278in}}%
\pgfpathlineto{\pgfqpoint{2.075683in}{0.625278in}}%
\pgfpathlineto{\pgfqpoint{2.076861in}{0.633527in}}%
\pgfpathlineto{\pgfqpoint{2.077024in}{0.625278in}}%
\pgfpathlineto{\pgfqpoint{2.077690in}{0.625278in}}%
\pgfpathlineto{\pgfqpoint{2.078697in}{0.633527in}}%
\pgfpathlineto{\pgfqpoint{2.078697in}{0.625278in}}%
\pgfpathlineto{\pgfqpoint{2.080052in}{0.633527in}}%
\pgfpathlineto{\pgfqpoint{2.080378in}{0.625278in}}%
\pgfpathlineto{\pgfqpoint{2.080726in}{0.625278in}}%
\pgfpathlineto{\pgfqpoint{2.082066in}{0.633527in}}%
\pgfpathlineto{\pgfqpoint{2.082221in}{0.625278in}}%
\pgfpathlineto{\pgfqpoint{2.083065in}{0.625278in}}%
\pgfpathlineto{\pgfqpoint{2.084235in}{0.633527in}}%
\pgfpathlineto{\pgfqpoint{2.084398in}{0.625278in}}%
\pgfpathlineto{\pgfqpoint{2.085079in}{0.625278in}}%
\pgfpathlineto{\pgfqpoint{2.086242in}{0.633527in}}%
\pgfpathlineto{\pgfqpoint{2.086249in}{0.625278in}}%
\pgfpathlineto{\pgfqpoint{2.087256in}{0.625278in}}%
\pgfpathlineto{\pgfqpoint{2.088419in}{0.633527in}}%
\pgfpathlineto{\pgfqpoint{2.088426in}{0.625278in}}%
\pgfpathlineto{\pgfqpoint{2.089433in}{0.625278in}}%
\pgfpathlineto{\pgfqpoint{2.090603in}{0.633527in}}%
\pgfpathlineto{\pgfqpoint{2.090610in}{0.625278in}}%
\pgfpathlineto{\pgfqpoint{2.091610in}{0.625278in}}%
\pgfpathlineto{\pgfqpoint{2.092950in}{0.633527in}}%
\pgfpathlineto{\pgfqpoint{2.093283in}{0.625278in}}%
\pgfpathlineto{\pgfqpoint{2.093957in}{0.625278in}}%
\pgfpathlineto{\pgfqpoint{2.095127in}{0.633527in}}%
\pgfpathlineto{\pgfqpoint{2.095312in}{0.625278in}}%
\pgfpathlineto{\pgfqpoint{2.096149in}{0.625278in}}%
\pgfpathlineto{\pgfqpoint{2.097311in}{0.633527in}}%
\pgfpathlineto{\pgfqpoint{2.097474in}{0.625278in}}%
\pgfpathlineto{\pgfqpoint{2.098148in}{0.625278in}}%
\pgfpathlineto{\pgfqpoint{2.099318in}{0.633527in}}%
\pgfpathlineto{\pgfqpoint{2.099821in}{0.625278in}}%
\pgfpathlineto{\pgfqpoint{2.100154in}{0.625278in}}%
\pgfpathlineto{\pgfqpoint{2.101332in}{0.633527in}}%
\pgfpathlineto{\pgfqpoint{2.101494in}{0.625278in}}%
\pgfpathlineto{\pgfqpoint{2.102339in}{0.625278in}}%
\pgfpathlineto{\pgfqpoint{2.103508in}{0.633527in}}%
\pgfpathlineto{\pgfqpoint{2.103679in}{0.625278in}}%
\pgfpathlineto{\pgfqpoint{2.104515in}{0.625278in}}%
\pgfpathlineto{\pgfqpoint{2.105685in}{0.633527in}}%
\pgfpathlineto{\pgfqpoint{2.106026in}{0.625278in}}%
\pgfpathlineto{\pgfqpoint{2.106692in}{0.625278in}}%
\pgfpathlineto{\pgfqpoint{2.107869in}{0.633527in}}%
\pgfpathlineto{\pgfqpoint{2.108032in}{0.625278in}}%
\pgfpathlineto{\pgfqpoint{2.108706in}{0.625278in}}%
\pgfpathlineto{\pgfqpoint{2.109846in}{0.633527in}}%
\pgfpathlineto{\pgfqpoint{2.110046in}{0.625278in}}%
\pgfpathlineto{\pgfqpoint{2.110720in}{0.625278in}}%
\pgfpathlineto{\pgfqpoint{2.112060in}{0.633527in}}%
\pgfpathlineto{\pgfqpoint{2.112060in}{0.625278in}}%
\pgfpathlineto{\pgfqpoint{2.113060in}{0.625278in}}%
\pgfpathlineto{\pgfqpoint{2.114237in}{0.633527in}}%
\pgfpathlineto{\pgfqpoint{2.114570in}{0.625278in}}%
\pgfpathlineto{\pgfqpoint{2.115244in}{0.625278in}}%
\pgfpathlineto{\pgfqpoint{2.116414in}{0.633527in}}%
\pgfpathlineto{\pgfqpoint{2.116755in}{0.625278in}}%
\pgfpathlineto{\pgfqpoint{2.117421in}{0.625278in}}%
\pgfpathlineto{\pgfqpoint{2.118761in}{0.633527in}}%
\pgfpathlineto{\pgfqpoint{2.118931in}{0.625278in}}%
\pgfpathlineto{\pgfqpoint{2.119768in}{0.625278in}}%
\pgfpathlineto{\pgfqpoint{2.120945in}{0.633527in}}%
\pgfpathlineto{\pgfqpoint{2.121108in}{0.625278in}}%
\pgfpathlineto{\pgfqpoint{2.121782in}{0.625278in}}%
\pgfpathlineto{\pgfqpoint{2.122952in}{0.633527in}}%
\pgfpathlineto{\pgfqpoint{2.123292in}{0.625278in}}%
\pgfpathlineto{\pgfqpoint{2.123959in}{0.625278in}}%
\pgfpathlineto{\pgfqpoint{2.125136in}{0.633527in}}%
\pgfpathlineto{\pgfqpoint{2.125469in}{0.625278in}}%
\pgfpathlineto{\pgfqpoint{2.126136in}{0.625278in}}%
\pgfpathlineto{\pgfqpoint{2.127313in}{0.633527in}}%
\pgfpathlineto{\pgfqpoint{2.127483in}{0.625278in}}%
\pgfpathlineto{\pgfqpoint{2.128320in}{0.625278in}}%
\pgfpathlineto{\pgfqpoint{2.129490in}{0.633527in}}%
\pgfpathlineto{\pgfqpoint{2.129660in}{0.625278in}}%
\pgfpathlineto{\pgfqpoint{2.130497in}{0.625278in}}%
\pgfpathlineto{\pgfqpoint{2.131674in}{0.633527in}}%
\pgfpathlineto{\pgfqpoint{2.131837in}{0.625278in}}%
\pgfpathlineto{\pgfqpoint{2.132674in}{0.625278in}}%
\pgfpathlineto{\pgfqpoint{2.133851in}{0.633527in}}%
\pgfpathlineto{\pgfqpoint{2.134014in}{0.625278in}}%
\pgfpathlineto{\pgfqpoint{2.134858in}{0.625278in}}%
\pgfpathlineto{\pgfqpoint{2.136028in}{0.633527in}}%
\pgfpathlineto{\pgfqpoint{2.136198in}{0.625278in}}%
\pgfpathlineto{\pgfqpoint{2.137094in}{0.625278in}}%
\pgfpathlineto{\pgfqpoint{2.138375in}{0.633527in}}%
\pgfpathlineto{\pgfqpoint{2.138375in}{0.625278in}}%
\pgfpathlineto{\pgfqpoint{2.139382in}{0.625278in}}%
\pgfpathlineto{\pgfqpoint{2.140389in}{0.633527in}}%
\pgfpathlineto{\pgfqpoint{2.140389in}{0.625278in}}%
\pgfpathlineto{\pgfqpoint{2.141559in}{0.633527in}}%
\pgfpathlineto{\pgfqpoint{2.141559in}{0.625278in}}%
\pgfpathlineto{\pgfqpoint{2.142566in}{0.625278in}}%
\pgfpathlineto{\pgfqpoint{2.143743in}{0.633527in}}%
\pgfpathlineto{\pgfqpoint{2.143906in}{0.625278in}}%
\pgfpathlineto{\pgfqpoint{2.144742in}{0.625278in}}%
\pgfpathlineto{\pgfqpoint{2.145920in}{0.633527in}}%
\pgfpathlineto{\pgfqpoint{2.146253in}{0.625278in}}%
\pgfpathlineto{\pgfqpoint{2.146927in}{0.625278in}}%
\pgfpathlineto{\pgfqpoint{2.148096in}{0.633527in}}%
\pgfpathlineto{\pgfqpoint{2.148267in}{0.625278in}}%
\pgfpathlineto{\pgfqpoint{2.149103in}{0.625278in}}%
\pgfpathlineto{\pgfqpoint{2.150281in}{0.633527in}}%
\pgfpathlineto{\pgfqpoint{2.150614in}{0.625278in}}%
\pgfpathlineto{\pgfqpoint{2.151280in}{0.625278in}}%
\pgfpathlineto{\pgfqpoint{2.152458in}{0.633527in}}%
\pgfpathlineto{\pgfqpoint{2.152628in}{0.625278in}}%
\pgfpathlineto{\pgfqpoint{2.153465in}{0.625278in}}%
\pgfpathlineto{\pgfqpoint{2.154634in}{0.633527in}}%
\pgfpathlineto{\pgfqpoint{2.154805in}{0.625278in}}%
\pgfpathlineto{\pgfqpoint{2.155641in}{0.625278in}}%
\pgfpathlineto{\pgfqpoint{2.156819in}{0.633527in}}%
\pgfpathlineto{\pgfqpoint{2.156982in}{0.625278in}}%
\pgfpathlineto{\pgfqpoint{2.157818in}{0.625278in}}%
\pgfpathlineto{\pgfqpoint{2.158825in}{0.633527in}}%
\pgfpathlineto{\pgfqpoint{2.158825in}{0.625278in}}%
\pgfpathlineto{\pgfqpoint{2.160002in}{0.633527in}}%
\pgfpathlineto{\pgfqpoint{2.160165in}{0.625278in}}%
\pgfpathlineto{\pgfqpoint{2.161009in}{0.625278in}}%
\pgfpathlineto{\pgfqpoint{2.162179in}{0.633527in}}%
\pgfpathlineto{\pgfqpoint{2.162179in}{0.625278in}}%
\pgfpathlineto{\pgfqpoint{2.163186in}{0.625278in}}%
\pgfpathlineto{\pgfqpoint{2.164526in}{0.633527in}}%
\pgfpathlineto{\pgfqpoint{2.164526in}{0.625278in}}%
\pgfpathlineto{\pgfqpoint{2.165533in}{0.625278in}}%
\pgfpathlineto{\pgfqpoint{2.166703in}{0.633527in}}%
\pgfpathlineto{\pgfqpoint{2.166703in}{0.625278in}}%
\pgfpathlineto{\pgfqpoint{2.167710in}{0.625278in}}%
\pgfpathlineto{\pgfqpoint{2.168887in}{0.633527in}}%
\pgfpathlineto{\pgfqpoint{2.168887in}{0.625278in}}%
\pgfpathlineto{\pgfqpoint{2.169887in}{0.625278in}}%
\pgfpathlineto{\pgfqpoint{2.171064in}{0.633527in}}%
\pgfpathlineto{\pgfqpoint{2.171064in}{0.625278in}}%
\pgfpathlineto{\pgfqpoint{2.172071in}{0.625278in}}%
\pgfpathlineto{\pgfqpoint{2.173078in}{0.633527in}}%
\pgfpathlineto{\pgfqpoint{2.173078in}{0.625278in}}%
\pgfpathlineto{\pgfqpoint{2.174248in}{0.633527in}}%
\pgfpathlineto{\pgfqpoint{2.174418in}{0.625278in}}%
\pgfpathlineto{\pgfqpoint{2.175255in}{0.625278in}}%
\pgfpathlineto{\pgfqpoint{2.176595in}{0.633527in}}%
\pgfpathlineto{\pgfqpoint{2.176595in}{0.625278in}}%
\pgfpathlineto{\pgfqpoint{2.177602in}{0.625278in}}%
\pgfpathlineto{\pgfqpoint{2.178772in}{0.633527in}}%
\pgfpathlineto{\pgfqpoint{2.178942in}{0.625278in}}%
\pgfpathlineto{\pgfqpoint{2.179446in}{0.625278in}}%
\pgfpathlineto{\pgfqpoint{2.180616in}{0.633527in}}%
\pgfpathlineto{\pgfqpoint{2.180786in}{0.625278in}}%
\pgfpathlineto{\pgfqpoint{2.181623in}{0.625278in}}%
\pgfpathlineto{\pgfqpoint{2.182963in}{0.633527in}}%
\pgfpathlineto{\pgfqpoint{2.182963in}{0.625278in}}%
\pgfpathlineto{\pgfqpoint{2.183970in}{0.625278in}}%
\pgfpathlineto{\pgfqpoint{2.185147in}{0.633527in}}%
\pgfpathlineto{\pgfqpoint{2.185310in}{0.625278in}}%
\pgfpathlineto{\pgfqpoint{2.186154in}{0.625278in}}%
\pgfpathlineto{\pgfqpoint{2.187324in}{0.633527in}}%
\pgfpathlineto{\pgfqpoint{2.187324in}{0.625278in}}%
\pgfpathlineto{\pgfqpoint{2.188331in}{0.625278in}}%
\pgfpathlineto{\pgfqpoint{2.189501in}{0.633527in}}%
\pgfpathlineto{\pgfqpoint{2.189671in}{0.625278in}}%
\pgfpathlineto{\pgfqpoint{2.190345in}{0.625278in}}%
\pgfpathlineto{\pgfqpoint{2.191515in}{0.633527in}}%
\pgfpathlineto{\pgfqpoint{2.191685in}{0.625278in}}%
\pgfpathlineto{\pgfqpoint{2.192351in}{0.625278in}}%
\pgfpathlineto{\pgfqpoint{2.193529in}{0.633527in}}%
\pgfpathlineto{\pgfqpoint{2.193529in}{0.625278in}}%
\pgfpathlineto{\pgfqpoint{2.194536in}{0.625278in}}%
\pgfpathlineto{\pgfqpoint{2.195705in}{0.633527in}}%
\pgfpathlineto{\pgfqpoint{2.195705in}{0.625278in}}%
\pgfpathlineto{\pgfqpoint{2.196712in}{0.625278in}}%
\pgfpathlineto{\pgfqpoint{2.197882in}{0.633527in}}%
\pgfpathlineto{\pgfqpoint{2.198053in}{0.625278in}}%
\pgfpathlineto{\pgfqpoint{2.198889in}{0.625278in}}%
\pgfpathlineto{\pgfqpoint{2.200067in}{0.633527in}}%
\pgfpathlineto{\pgfqpoint{2.200229in}{0.625278in}}%
\pgfpathlineto{\pgfqpoint{2.200903in}{0.625278in}}%
\pgfpathlineto{\pgfqpoint{2.202243in}{0.633527in}}%
\pgfpathlineto{\pgfqpoint{2.202414in}{0.625278in}}%
\pgfpathlineto{\pgfqpoint{2.203250in}{0.625278in}}%
\pgfpathlineto{\pgfqpoint{2.204420in}{0.633527in}}%
\pgfpathlineto{\pgfqpoint{2.204420in}{0.625278in}}%
\pgfpathlineto{\pgfqpoint{2.205427in}{0.625278in}}%
\pgfpathlineto{\pgfqpoint{2.206604in}{0.633527in}}%
\pgfpathlineto{\pgfqpoint{2.206767in}{0.625278in}}%
\pgfpathlineto{\pgfqpoint{2.207604in}{0.625278in}}%
\pgfpathlineto{\pgfqpoint{2.208781in}{0.633527in}}%
\pgfpathlineto{\pgfqpoint{2.208952in}{0.625278in}}%
\pgfpathlineto{\pgfqpoint{2.209788in}{0.625278in}}%
\pgfpathlineto{\pgfqpoint{2.210958in}{0.633527in}}%
\pgfpathlineto{\pgfqpoint{2.210958in}{0.625278in}}%
\pgfpathlineto{\pgfqpoint{2.211965in}{0.625278in}}%
\pgfpathlineto{\pgfqpoint{2.212972in}{0.633527in}}%
\pgfpathlineto{\pgfqpoint{2.212972in}{0.625278in}}%
\pgfpathlineto{\pgfqpoint{2.214312in}{0.633527in}}%
\pgfpathlineto{\pgfqpoint{2.214482in}{0.625278in}}%
\pgfpathlineto{\pgfqpoint{2.215319in}{0.625278in}}%
\pgfpathlineto{\pgfqpoint{2.216489in}{0.633527in}}%
\pgfpathlineto{\pgfqpoint{2.216659in}{0.625278in}}%
\pgfpathlineto{\pgfqpoint{2.217496in}{0.625278in}}%
\pgfpathlineto{\pgfqpoint{2.218673in}{0.633527in}}%
\pgfpathlineto{\pgfqpoint{2.218673in}{0.625278in}}%
\pgfpathlineto{\pgfqpoint{2.219680in}{0.625278in}}%
\pgfpathlineto{\pgfqpoint{2.220850in}{0.633527in}}%
\pgfpathlineto{\pgfqpoint{2.221020in}{0.625278in}}%
\pgfpathlineto{\pgfqpoint{2.221687in}{0.625278in}}%
\pgfpathlineto{\pgfqpoint{2.222864in}{0.633527in}}%
\pgfpathlineto{\pgfqpoint{2.223027in}{0.625278in}}%
\pgfpathlineto{\pgfqpoint{2.223871in}{0.625278in}}%
\pgfpathlineto{\pgfqpoint{2.225041in}{0.633527in}}%
\pgfpathlineto{\pgfqpoint{2.225211in}{0.625278in}}%
\pgfpathlineto{\pgfqpoint{2.226063in}{0.625278in}}%
\pgfpathlineto{\pgfqpoint{2.227218in}{0.633527in}}%
\pgfpathlineto{\pgfqpoint{2.227240in}{0.625278in}}%
\pgfpathlineto{\pgfqpoint{2.227743in}{0.674773in}}%
\pgfpathlineto{\pgfqpoint{2.228077in}{0.674773in}}%
\pgfpathlineto{\pgfqpoint{2.229232in}{0.683022in}}%
\pgfpathlineto{\pgfqpoint{2.229735in}{0.658274in}}%
\pgfpathlineto{\pgfqpoint{2.230253in}{0.674773in}}%
\pgfpathlineto{\pgfqpoint{2.231601in}{0.625278in}}%
\pgfpathlineto{\pgfqpoint{2.231764in}{0.650025in}}%
\pgfpathlineto{\pgfqpoint{2.232919in}{0.707769in}}%
\pgfpathlineto{\pgfqpoint{2.233274in}{0.674773in}}%
\pgfpathlineto{\pgfqpoint{2.234614in}{0.650025in}}%
\pgfpathlineto{\pgfqpoint{2.236777in}{0.732517in}}%
\pgfpathlineto{\pgfqpoint{2.236791in}{0.724268in}}%
\pgfpathlineto{\pgfqpoint{2.237295in}{0.699520in}}%
\pgfpathlineto{\pgfqpoint{2.237613in}{0.732517in}}%
\pgfpathlineto{\pgfqpoint{2.237635in}{0.724268in}}%
\pgfpathlineto{\pgfqpoint{2.238783in}{0.732517in}}%
\pgfpathlineto{\pgfqpoint{2.239812in}{0.724268in}}%
\pgfpathlineto{\pgfqpoint{2.240967in}{0.732517in}}%
\pgfpathlineto{\pgfqpoint{2.241989in}{0.724268in}}%
\pgfpathlineto{\pgfqpoint{2.242663in}{0.749015in}}%
\pgfpathlineto{\pgfqpoint{2.243144in}{0.732517in}}%
\pgfpathlineto{\pgfqpoint{2.243329in}{0.724268in}}%
\pgfpathlineto{\pgfqpoint{2.243648in}{0.757264in}}%
\pgfpathlineto{\pgfqpoint{2.244173in}{0.740766in}}%
\pgfpathlineto{\pgfqpoint{2.244877in}{0.732517in}}%
\pgfpathlineto{\pgfqpoint{2.244336in}{0.749015in}}%
\pgfpathlineto{\pgfqpoint{2.244988in}{0.740766in}}%
\pgfpathlineto{\pgfqpoint{2.245158in}{0.782012in}}%
\pgfpathlineto{\pgfqpoint{2.246017in}{0.749015in}}%
\pgfpathlineto{\pgfqpoint{2.247668in}{0.782012in}}%
\pgfpathlineto{\pgfqpoint{2.248697in}{0.773763in}}%
\pgfpathlineto{\pgfqpoint{2.249030in}{0.798510in}}%
\pgfpathlineto{\pgfqpoint{2.249852in}{0.782012in}}%
\pgfpathlineto{\pgfqpoint{2.250037in}{0.773763in}}%
\pgfpathlineto{\pgfqpoint{2.250356in}{0.806759in}}%
\pgfpathlineto{\pgfqpoint{2.250874in}{0.790261in}}%
\pgfpathlineto{\pgfqpoint{2.251378in}{0.773763in}}%
\pgfpathlineto{\pgfqpoint{2.251548in}{0.798510in}}%
\pgfpathlineto{\pgfqpoint{2.252029in}{0.806759in}}%
\pgfpathlineto{\pgfqpoint{2.252362in}{0.782012in}}%
\pgfpathlineto{\pgfqpoint{2.253369in}{0.757264in}}%
\pgfpathlineto{\pgfqpoint{2.253392in}{0.773763in}}%
\pgfpathlineto{\pgfqpoint{2.254547in}{0.782012in}}%
\pgfpathlineto{\pgfqpoint{2.255568in}{0.749015in}}%
\pgfpathlineto{\pgfqpoint{2.256050in}{0.757264in}}%
\pgfpathlineto{\pgfqpoint{2.256242in}{0.740766in}}%
\pgfpathlineto{\pgfqpoint{2.257057in}{0.732517in}}%
\pgfpathlineto{\pgfqpoint{2.256405in}{0.749015in}}%
\pgfpathlineto{\pgfqpoint{2.257079in}{0.749015in}}%
\pgfpathlineto{\pgfqpoint{2.259071in}{0.806759in}}%
\pgfpathlineto{\pgfqpoint{2.260092in}{0.798510in}}%
\pgfpathlineto{\pgfqpoint{2.262084in}{0.856254in}}%
\pgfpathlineto{\pgfqpoint{2.262780in}{0.823258in}}%
\pgfpathlineto{\pgfqpoint{2.263113in}{0.848005in}}%
\pgfpathlineto{\pgfqpoint{2.263432in}{0.856254in}}%
\pgfpathlineto{\pgfqpoint{2.263765in}{0.831507in}}%
\pgfpathlineto{\pgfqpoint{2.263950in}{0.839756in}}%
\pgfpathlineto{\pgfqpoint{2.264431in}{0.831507in}}%
\pgfpathlineto{\pgfqpoint{2.264601in}{0.856254in}}%
\pgfpathlineto{\pgfqpoint{2.264624in}{0.848005in}}%
\pgfpathlineto{\pgfqpoint{2.264935in}{0.856254in}}%
\pgfpathlineto{\pgfqpoint{2.265290in}{0.823258in}}%
\pgfpathlineto{\pgfqpoint{2.265438in}{0.831507in}}%
\pgfpathlineto{\pgfqpoint{2.266297in}{0.823258in}}%
\pgfpathlineto{\pgfqpoint{2.266467in}{0.848005in}}%
\pgfpathlineto{\pgfqpoint{2.267008in}{0.831507in}}%
\pgfpathlineto{\pgfqpoint{2.267282in}{0.856254in}}%
\pgfpathlineto{\pgfqpoint{2.267304in}{0.848005in}}%
\pgfpathlineto{\pgfqpoint{2.268977in}{0.897500in}}%
\pgfpathlineto{\pgfqpoint{2.269148in}{0.889251in}}%
\pgfpathlineto{\pgfqpoint{2.269984in}{0.872753in}}%
\pgfpathlineto{\pgfqpoint{2.270132in}{0.889251in}}%
\pgfpathlineto{\pgfqpoint{2.270969in}{0.905749in}}%
\pgfpathlineto{\pgfqpoint{2.270466in}{0.881002in}}%
\pgfpathlineto{\pgfqpoint{2.271162in}{0.889251in}}%
\pgfpathlineto{\pgfqpoint{2.271643in}{0.881002in}}%
\pgfpathlineto{\pgfqpoint{2.271813in}{0.905749in}}%
\pgfpathlineto{\pgfqpoint{2.271998in}{0.897500in}}%
\pgfpathlineto{\pgfqpoint{2.273657in}{0.930497in}}%
\pgfpathlineto{\pgfqpoint{2.274679in}{0.922247in}}%
\pgfpathlineto{\pgfqpoint{2.276004in}{0.955244in}}%
\pgfpathlineto{\pgfqpoint{2.276019in}{0.946995in}}%
\pgfpathlineto{\pgfqpoint{2.276670in}{0.905749in}}%
\pgfpathlineto{\pgfqpoint{2.277026in}{0.946995in}}%
\pgfpathlineto{\pgfqpoint{2.277677in}{0.955244in}}%
\pgfpathlineto{\pgfqpoint{2.277862in}{0.938746in}}%
\pgfpathlineto{\pgfqpoint{2.279543in}{0.897500in}}%
\pgfpathlineto{\pgfqpoint{2.280713in}{0.922247in}}%
\pgfpathlineto{\pgfqpoint{2.280883in}{0.913998in}}%
\pgfpathlineto{\pgfqpoint{2.281365in}{0.905749in}}%
\pgfpathlineto{\pgfqpoint{2.281535in}{0.930497in}}%
\pgfpathlineto{\pgfqpoint{2.281720in}{0.922247in}}%
\pgfpathlineto{\pgfqpoint{2.283038in}{0.955244in}}%
\pgfpathlineto{\pgfqpoint{2.283060in}{0.946995in}}%
\pgfpathlineto{\pgfqpoint{2.283379in}{0.930497in}}%
\pgfpathlineto{\pgfqpoint{2.283712in}{0.963493in}}%
\pgfpathlineto{\pgfqpoint{2.284045in}{0.979992in}}%
\pgfpathlineto{\pgfqpoint{2.284719in}{0.930497in}}%
\pgfpathlineto{\pgfqpoint{2.284778in}{0.955244in}}%
\pgfpathlineto{\pgfqpoint{2.284904in}{0.946995in}}%
\pgfpathlineto{\pgfqpoint{2.285385in}{0.988241in}}%
\pgfpathlineto{\pgfqpoint{2.285726in}{1.004739in}}%
\pgfpathlineto{\pgfqpoint{2.286081in}{0.971742in}}%
\pgfpathlineto{\pgfqpoint{2.286414in}{0.971742in}}%
\pgfpathlineto{\pgfqpoint{2.287399in}{1.004739in}}%
\pgfpathlineto{\pgfqpoint{2.287584in}{0.988241in}}%
\pgfpathlineto{\pgfqpoint{2.288591in}{0.971742in}}%
\pgfpathlineto{\pgfqpoint{2.287754in}{0.996490in}}%
\pgfpathlineto{\pgfqpoint{2.288739in}{0.979992in}}%
\pgfpathlineto{\pgfqpoint{2.288761in}{0.996490in}}%
\pgfpathlineto{\pgfqpoint{2.289428in}{0.971742in}}%
\pgfpathlineto{\pgfqpoint{2.289768in}{0.971742in}}%
\pgfpathlineto{\pgfqpoint{2.290753in}{1.004739in}}%
\pgfpathlineto{\pgfqpoint{2.290938in}{0.988241in}}%
\pgfpathlineto{\pgfqpoint{2.292597in}{0.955244in}}%
\pgfpathlineto{\pgfqpoint{2.292930in}{1.004739in}}%
\pgfpathlineto{\pgfqpoint{2.293619in}{0.996490in}}%
\pgfpathlineto{\pgfqpoint{2.294611in}{1.054234in}}%
\pgfpathlineto{\pgfqpoint{2.295277in}{1.029487in}}%
\pgfpathlineto{\pgfqpoint{2.295966in}{0.996490in}}%
\pgfpathlineto{\pgfqpoint{2.296306in}{1.021237in}}%
\pgfpathlineto{\pgfqpoint{2.297624in}{1.054234in}}%
\pgfpathlineto{\pgfqpoint{2.296802in}{0.996490in}}%
\pgfpathlineto{\pgfqpoint{2.297646in}{1.045985in}}%
\pgfpathlineto{\pgfqpoint{2.298631in}{1.004739in}}%
\pgfpathlineto{\pgfqpoint{2.298987in}{1.021237in}}%
\pgfpathlineto{\pgfqpoint{2.300497in}{1.095480in}}%
\pgfpathlineto{\pgfqpoint{2.300830in}{1.070732in}}%
\pgfpathlineto{\pgfqpoint{2.302148in}{1.054234in}}%
\pgfpathlineto{\pgfqpoint{2.302992in}{1.078981in}}%
\pgfpathlineto{\pgfqpoint{2.303177in}{1.070732in}}%
\pgfpathlineto{\pgfqpoint{2.304666in}{1.128476in}}%
\pgfpathlineto{\pgfqpoint{2.304851in}{1.111978in}}%
\pgfpathlineto{\pgfqpoint{2.305339in}{1.103729in}}%
\pgfpathlineto{\pgfqpoint{2.305502in}{1.136726in}}%
\pgfpathlineto{\pgfqpoint{2.305835in}{1.177971in}}%
\pgfpathlineto{\pgfqpoint{2.306509in}{1.128476in}}%
\pgfpathlineto{\pgfqpoint{2.306531in}{1.144975in}}%
\pgfpathlineto{\pgfqpoint{2.307183in}{1.128476in}}%
\pgfpathlineto{\pgfqpoint{2.307346in}{1.153224in}}%
\pgfpathlineto{\pgfqpoint{2.307368in}{1.144975in}}%
\pgfpathlineto{\pgfqpoint{2.308686in}{1.177971in}}%
\pgfpathlineto{\pgfqpoint{2.308708in}{1.169722in}}%
\pgfpathlineto{\pgfqpoint{2.309545in}{1.120227in}}%
\pgfpathlineto{\pgfqpoint{2.310367in}{1.136726in}}%
\pgfpathlineto{\pgfqpoint{2.310722in}{1.194470in}}%
\pgfpathlineto{\pgfqpoint{2.311389in}{1.194470in}}%
\pgfpathlineto{\pgfqpoint{2.313232in}{1.144975in}}%
\pgfpathlineto{\pgfqpoint{2.314069in}{1.169722in}}%
\pgfpathlineto{\pgfqpoint{2.314387in}{1.153224in}}%
\pgfpathlineto{\pgfqpoint{2.315061in}{1.128476in}}%
\pgfpathlineto{\pgfqpoint{2.315416in}{1.144975in}}%
\pgfpathlineto{\pgfqpoint{2.316734in}{1.202719in}}%
\pgfpathlineto{\pgfqpoint{2.317068in}{1.177971in}}%
\pgfpathlineto{\pgfqpoint{2.317912in}{1.153224in}}%
\pgfpathlineto{\pgfqpoint{2.317926in}{1.194470in}}%
\pgfpathlineto{\pgfqpoint{2.318097in}{1.186221in}}%
\pgfpathlineto{\pgfqpoint{2.318260in}{1.194470in}}%
\pgfpathlineto{\pgfqpoint{2.319252in}{1.177971in}}%
\pgfpathlineto{\pgfqpoint{2.319415in}{1.202719in}}%
\pgfpathlineto{\pgfqpoint{2.320274in}{1.194470in}}%
\pgfpathlineto{\pgfqpoint{2.321429in}{1.227466in}}%
\pgfpathlineto{\pgfqpoint{2.321614in}{1.210968in}}%
\pgfpathlineto{\pgfqpoint{2.322102in}{1.202719in}}%
\pgfpathlineto{\pgfqpoint{2.322265in}{1.227466in}}%
\pgfpathlineto{\pgfqpoint{2.322288in}{1.219217in}}%
\pgfpathlineto{\pgfqpoint{2.322599in}{1.227466in}}%
\pgfpathlineto{\pgfqpoint{2.322954in}{1.194470in}}%
\pgfpathlineto{\pgfqpoint{2.323102in}{1.202719in}}%
\pgfpathlineto{\pgfqpoint{2.324279in}{1.177971in}}%
\pgfpathlineto{\pgfqpoint{2.324294in}{1.194470in}}%
\pgfpathlineto{\pgfqpoint{2.325301in}{1.194470in}}%
\pgfpathlineto{\pgfqpoint{2.325642in}{1.169722in}}%
\pgfpathlineto{\pgfqpoint{2.326456in}{1.186221in}}%
\pgfpathlineto{\pgfqpoint{2.327463in}{1.202719in}}%
\pgfpathlineto{\pgfqpoint{2.326812in}{1.169722in}}%
\pgfpathlineto{\pgfqpoint{2.327485in}{1.194470in}}%
\pgfpathlineto{\pgfqpoint{2.329832in}{1.120227in}}%
\pgfpathlineto{\pgfqpoint{2.331002in}{1.169722in}}%
\pgfpathlineto{\pgfqpoint{2.330484in}{1.103729in}}%
\pgfpathlineto{\pgfqpoint{2.331484in}{1.153224in}}%
\pgfpathlineto{\pgfqpoint{2.332328in}{1.103729in}}%
\pgfpathlineto{\pgfqpoint{2.332676in}{1.120227in}}%
\pgfpathlineto{\pgfqpoint{2.334838in}{1.177971in}}%
\pgfpathlineto{\pgfqpoint{2.335511in}{1.128476in}}%
\pgfpathlineto{\pgfqpoint{2.336030in}{1.144975in}}%
\pgfpathlineto{\pgfqpoint{2.336518in}{1.153224in}}%
\pgfpathlineto{\pgfqpoint{2.336704in}{1.136726in}}%
\pgfpathlineto{\pgfqpoint{2.338377in}{1.070732in}}%
\pgfpathlineto{\pgfqpoint{2.338547in}{1.095480in}}%
\pgfpathlineto{\pgfqpoint{2.342212in}{0.930497in}}%
\pgfpathlineto{\pgfqpoint{2.342405in}{0.946995in}}%
\pgfpathlineto{\pgfqpoint{2.342886in}{0.955244in}}%
\pgfpathlineto{\pgfqpoint{2.343219in}{0.930497in}}%
\pgfpathlineto{\pgfqpoint{2.344078in}{0.897500in}}%
\pgfpathlineto{\pgfqpoint{2.344248in}{0.922247in}}%
\pgfpathlineto{\pgfqpoint{2.344730in}{0.979992in}}%
\pgfpathlineto{\pgfqpoint{2.345589in}{0.938746in}}%
\pgfpathlineto{\pgfqpoint{2.347595in}{0.872753in}}%
\pgfpathlineto{\pgfqpoint{2.347743in}{0.881002in}}%
\pgfpathlineto{\pgfqpoint{2.348099in}{0.848005in}}%
\pgfpathlineto{\pgfqpoint{2.348247in}{0.856254in}}%
\pgfpathlineto{\pgfqpoint{2.348750in}{0.881002in}}%
\pgfpathlineto{\pgfqpoint{2.349276in}{0.848005in}}%
\pgfpathlineto{\pgfqpoint{2.349439in}{0.872753in}}%
\pgfpathlineto{\pgfqpoint{2.350090in}{0.831507in}}%
\pgfpathlineto{\pgfqpoint{2.350113in}{0.848005in}}%
\pgfpathlineto{\pgfqpoint{2.352126in}{0.782012in}}%
\pgfpathlineto{\pgfqpoint{2.352275in}{0.790261in}}%
\pgfpathlineto{\pgfqpoint{2.352800in}{0.757264in}}%
\pgfpathlineto{\pgfqpoint{2.353296in}{0.798510in}}%
\pgfpathlineto{\pgfqpoint{2.354474in}{0.773763in}}%
\pgfpathlineto{\pgfqpoint{2.353800in}{0.823258in}}%
\pgfpathlineto{\pgfqpoint{2.354474in}{0.782012in}}%
\pgfpathlineto{\pgfqpoint{2.355133in}{0.831507in}}%
\pgfpathlineto{\pgfqpoint{2.355814in}{0.798510in}}%
\pgfpathlineto{\pgfqpoint{2.356125in}{0.782012in}}%
\pgfpathlineto{\pgfqpoint{2.356147in}{0.823258in}}%
\pgfpathlineto{\pgfqpoint{2.356473in}{0.815008in}}%
\pgfpathlineto{\pgfqpoint{2.356821in}{0.872753in}}%
\pgfpathlineto{\pgfqpoint{2.357487in}{0.848005in}}%
\pgfpathlineto{\pgfqpoint{2.357820in}{0.856254in}}%
\pgfpathlineto{\pgfqpoint{2.358153in}{0.831507in}}%
\pgfpathlineto{\pgfqpoint{2.358161in}{0.823258in}}%
\pgfpathlineto{\pgfqpoint{2.358812in}{0.856254in}}%
\pgfpathlineto{\pgfqpoint{2.359168in}{0.848005in}}%
\pgfpathlineto{\pgfqpoint{2.361663in}{0.930497in}}%
\pgfpathlineto{\pgfqpoint{2.362507in}{0.905749in}}%
\pgfpathlineto{\pgfqpoint{2.362685in}{0.922247in}}%
\pgfpathlineto{\pgfqpoint{2.364195in}{1.021237in}}%
\pgfpathlineto{\pgfqpoint{2.365032in}{0.988241in}}%
\pgfpathlineto{\pgfqpoint{2.365032in}{0.979992in}}%
\pgfpathlineto{\pgfqpoint{2.366039in}{1.045985in}}%
\pgfpathlineto{\pgfqpoint{2.366187in}{1.054234in}}%
\pgfpathlineto{\pgfqpoint{2.366542in}{1.012988in}}%
\pgfpathlineto{\pgfqpoint{2.366876in}{1.004739in}}%
\pgfpathlineto{\pgfqpoint{2.367527in}{1.054234in}}%
\pgfpathlineto{\pgfqpoint{2.367549in}{1.045985in}}%
\pgfpathlineto{\pgfqpoint{2.369393in}{0.971742in}}%
\pgfpathlineto{\pgfqpoint{2.369556in}{0.996490in}}%
\pgfpathlineto{\pgfqpoint{2.369704in}{1.004739in}}%
\pgfpathlineto{\pgfqpoint{2.370059in}{0.971742in}}%
\pgfpathlineto{\pgfqpoint{2.370208in}{0.979992in}}%
\pgfpathlineto{\pgfqpoint{2.370563in}{0.971742in}}%
\pgfpathlineto{\pgfqpoint{2.371214in}{1.004739in}}%
\pgfpathlineto{\pgfqpoint{2.371237in}{0.996490in}}%
\pgfpathlineto{\pgfqpoint{2.372577in}{1.070732in}}%
\pgfpathlineto{\pgfqpoint{2.373080in}{1.037736in}}%
\pgfpathlineto{\pgfqpoint{2.373584in}{1.004739in}}%
\pgfpathlineto{\pgfqpoint{2.373243in}{1.045985in}}%
\pgfpathlineto{\pgfqpoint{2.374087in}{1.021237in}}%
\pgfpathlineto{\pgfqpoint{2.376101in}{0.881002in}}%
\pgfpathlineto{\pgfqpoint{2.377264in}{0.913998in}}%
\pgfpathlineto{\pgfqpoint{2.377271in}{0.922247in}}%
\pgfpathlineto{\pgfqpoint{2.378278in}{0.831507in}}%
\pgfpathlineto{\pgfqpoint{2.378441in}{0.848005in}}%
\pgfpathlineto{\pgfqpoint{2.378944in}{0.815008in}}%
\pgfpathlineto{\pgfqpoint{2.382136in}{0.683022in}}%
\pgfpathlineto{\pgfqpoint{2.382284in}{0.691271in}}%
\pgfpathlineto{\pgfqpoint{2.382624in}{0.707769in}}%
\pgfpathlineto{\pgfqpoint{2.383135in}{0.641776in}}%
\pgfpathlineto{\pgfqpoint{2.384313in}{0.625278in}}%
\pgfpathlineto{\pgfqpoint{2.385475in}{0.633527in}}%
\pgfpathlineto{\pgfqpoint{2.385645in}{0.625278in}}%
\pgfpathlineto{\pgfqpoint{2.386482in}{0.625278in}}%
\pgfpathlineto{\pgfqpoint{2.387652in}{0.633527in}}%
\pgfpathlineto{\pgfqpoint{2.388155in}{0.625278in}}%
\pgfpathlineto{\pgfqpoint{2.388659in}{0.625278in}}%
\pgfpathlineto{\pgfqpoint{2.389843in}{0.633527in}}%
\pgfpathlineto{\pgfqpoint{2.389999in}{0.625278in}}%
\pgfpathlineto{\pgfqpoint{2.390850in}{0.625278in}}%
\pgfpathlineto{\pgfqpoint{2.392020in}{0.633527in}}%
\pgfpathlineto{\pgfqpoint{2.392176in}{0.625278in}}%
\pgfpathlineto{\pgfqpoint{2.393012in}{0.625278in}}%
\pgfpathlineto{\pgfqpoint{2.394205in}{0.633527in}}%
\pgfpathlineto{\pgfqpoint{2.394360in}{0.625278in}}%
\pgfpathlineto{\pgfqpoint{2.395204in}{0.625278in}}%
\pgfpathlineto{\pgfqpoint{2.396211in}{0.633527in}}%
\pgfpathlineto{\pgfqpoint{2.396211in}{0.625278in}}%
\pgfpathlineto{\pgfqpoint{2.397381in}{0.633527in}}%
\pgfpathlineto{\pgfqpoint{2.397551in}{0.625278in}}%
\pgfpathlineto{\pgfqpoint{2.398395in}{0.625278in}}%
\pgfpathlineto{\pgfqpoint{2.399550in}{0.633527in}}%
\pgfpathlineto{\pgfqpoint{2.399728in}{0.625278in}}%
\pgfpathlineto{\pgfqpoint{2.400557in}{0.625278in}}%
\pgfpathlineto{\pgfqpoint{2.401742in}{0.633527in}}%
\pgfpathlineto{\pgfqpoint{2.401905in}{0.625278in}}%
\pgfpathlineto{\pgfqpoint{2.402749in}{0.625278in}}%
\pgfpathlineto{\pgfqpoint{2.403756in}{0.633527in}}%
\pgfpathlineto{\pgfqpoint{2.403756in}{0.625278in}}%
\pgfpathlineto{\pgfqpoint{2.404911in}{0.633527in}}%
\pgfpathlineto{\pgfqpoint{2.405081in}{0.625278in}}%
\pgfpathlineto{\pgfqpoint{2.405925in}{0.625278in}}%
\pgfpathlineto{\pgfqpoint{2.407088in}{0.633527in}}%
\pgfpathlineto{\pgfqpoint{2.407258in}{0.625278in}}%
\pgfpathlineto{\pgfqpoint{2.408117in}{0.625278in}}%
\pgfpathlineto{\pgfqpoint{2.409279in}{0.633527in}}%
\pgfpathlineto{\pgfqpoint{2.409457in}{0.625278in}}%
\pgfpathlineto{\pgfqpoint{2.410286in}{0.625278in}}%
\pgfpathlineto{\pgfqpoint{2.411464in}{0.633527in}}%
\pgfpathlineto{\pgfqpoint{2.411789in}{0.625278in}}%
\pgfpathlineto{\pgfqpoint{2.412286in}{0.625278in}}%
\pgfpathlineto{\pgfqpoint{2.413463in}{0.633527in}}%
\pgfpathlineto{\pgfqpoint{2.413641in}{0.625278in}}%
\pgfpathlineto{\pgfqpoint{2.414470in}{0.625278in}}%
\pgfpathlineto{\pgfqpoint{2.415640in}{0.633527in}}%
\pgfpathlineto{\pgfqpoint{2.415973in}{0.625278in}}%
\pgfpathlineto{\pgfqpoint{2.416647in}{0.625278in}}%
\pgfpathlineto{\pgfqpoint{2.417839in}{0.633527in}}%
\pgfpathlineto{\pgfqpoint{2.417987in}{0.625278in}}%
\pgfpathlineto{\pgfqpoint{2.418823in}{0.625278in}}%
\pgfpathlineto{\pgfqpoint{2.419993in}{0.633527in}}%
\pgfpathlineto{\pgfqpoint{2.419993in}{0.625278in}}%
\pgfpathlineto{\pgfqpoint{2.421000in}{0.625278in}}%
\pgfpathlineto{\pgfqpoint{2.422340in}{0.633527in}}%
\pgfpathlineto{\pgfqpoint{2.422511in}{0.625278in}}%
\pgfpathlineto{\pgfqpoint{2.423347in}{0.625278in}}%
\pgfpathlineto{\pgfqpoint{2.424354in}{0.633527in}}%
\pgfpathlineto{\pgfqpoint{2.424354in}{0.625278in}}%
\pgfpathlineto{\pgfqpoint{2.425206in}{0.633527in}}%
\pgfpathlineto{\pgfqpoint{2.425206in}{0.625278in}}%
\pgfpathlineto{\pgfqpoint{2.426368in}{0.633527in}}%
\pgfpathlineto{\pgfqpoint{2.426546in}{0.625278in}}%
\pgfpathlineto{\pgfqpoint{2.427205in}{0.625278in}}%
\pgfpathlineto{\pgfqpoint{2.428382in}{0.633527in}}%
\pgfpathlineto{\pgfqpoint{2.428545in}{0.625278in}}%
\pgfpathlineto{\pgfqpoint{2.429382in}{0.625278in}}%
\pgfpathlineto{\pgfqpoint{2.430722in}{0.633527in}}%
\pgfpathlineto{\pgfqpoint{2.430729in}{0.625278in}}%
\pgfpathlineto{\pgfqpoint{2.431566in}{0.625278in}}%
\pgfpathlineto{\pgfqpoint{2.432736in}{0.633527in}}%
\pgfpathlineto{\pgfqpoint{2.432906in}{0.625278in}}%
\pgfpathlineto{\pgfqpoint{2.433743in}{0.625278in}}%
\pgfpathlineto{\pgfqpoint{2.434913in}{0.633527in}}%
\pgfpathlineto{\pgfqpoint{2.435083in}{0.625278in}}%
\pgfpathlineto{\pgfqpoint{2.435920in}{0.625278in}}%
\pgfpathlineto{\pgfqpoint{2.437097in}{0.633527in}}%
\pgfpathlineto{\pgfqpoint{2.437097in}{0.625278in}}%
\pgfpathlineto{\pgfqpoint{2.438097in}{0.625278in}}%
\pgfpathlineto{\pgfqpoint{2.439111in}{0.633527in}}%
\pgfpathlineto{\pgfqpoint{2.439111in}{0.625278in}}%
\pgfpathlineto{\pgfqpoint{2.440281in}{0.633527in}}%
\pgfpathlineto{\pgfqpoint{2.440444in}{0.625278in}}%
\pgfpathlineto{\pgfqpoint{2.441125in}{0.625278in}}%
\pgfpathlineto{\pgfqpoint{2.442287in}{0.633527in}}%
\pgfpathlineto{\pgfqpoint{2.442458in}{0.625278in}}%
\pgfpathlineto{\pgfqpoint{2.443302in}{0.625278in}}%
\pgfpathlineto{\pgfqpoint{2.444634in}{0.633527in}}%
\pgfpathlineto{\pgfqpoint{2.444642in}{0.625278in}}%
\pgfpathlineto{\pgfqpoint{2.445641in}{0.625278in}}%
\pgfpathlineto{\pgfqpoint{2.446819in}{0.633527in}}%
\pgfpathlineto{\pgfqpoint{2.446982in}{0.625278in}}%
\pgfpathlineto{\pgfqpoint{2.447826in}{0.625278in}}%
\pgfpathlineto{\pgfqpoint{2.448996in}{0.633527in}}%
\pgfpathlineto{\pgfqpoint{2.449166in}{0.625278in}}%
\pgfpathlineto{\pgfqpoint{2.450003in}{0.625278in}}%
\pgfpathlineto{\pgfqpoint{2.451172in}{0.633527in}}%
\pgfpathlineto{\pgfqpoint{2.451343in}{0.625278in}}%
\pgfpathlineto{\pgfqpoint{2.452179in}{0.625278in}}%
\pgfpathlineto{\pgfqpoint{2.453357in}{0.633527in}}%
\pgfpathlineto{\pgfqpoint{2.453690in}{0.625278in}}%
\pgfpathlineto{\pgfqpoint{2.454415in}{0.625278in}}%
\pgfpathlineto{\pgfqpoint{2.455363in}{0.633527in}}%
\pgfpathlineto{\pgfqpoint{2.455363in}{0.625278in}}%
\pgfpathlineto{\pgfqpoint{2.456703in}{0.633527in}}%
\pgfpathlineto{\pgfqpoint{2.456711in}{0.625278in}}%
\pgfpathlineto{\pgfqpoint{2.457710in}{0.625278in}}%
\pgfpathlineto{\pgfqpoint{2.458888in}{0.633527in}}%
\pgfpathlineto{\pgfqpoint{2.459050in}{0.625278in}}%
\pgfpathlineto{\pgfqpoint{2.459895in}{0.625278in}}%
\pgfpathlineto{\pgfqpoint{2.461064in}{0.633527in}}%
\pgfpathlineto{\pgfqpoint{2.461235in}{0.625278in}}%
\pgfpathlineto{\pgfqpoint{2.461901in}{0.625278in}}%
\pgfpathlineto{\pgfqpoint{2.463078in}{0.633527in}}%
\pgfpathlineto{\pgfqpoint{2.463241in}{0.625278in}}%
\pgfpathlineto{\pgfqpoint{2.464085in}{0.625278in}}%
\pgfpathlineto{\pgfqpoint{2.465255in}{0.633527in}}%
\pgfpathlineto{\pgfqpoint{2.465425in}{0.625278in}}%
\pgfpathlineto{\pgfqpoint{2.466262in}{0.625278in}}%
\pgfpathlineto{\pgfqpoint{2.467432in}{0.633527in}}%
\pgfpathlineto{\pgfqpoint{2.467602in}{0.625278in}}%
\pgfpathlineto{\pgfqpoint{2.468439in}{0.625278in}}%
\pgfpathlineto{\pgfqpoint{2.469624in}{0.633527in}}%
\pgfpathlineto{\pgfqpoint{2.469779in}{0.625278in}}%
\pgfpathlineto{\pgfqpoint{2.470623in}{0.625278in}}%
\pgfpathlineto{\pgfqpoint{2.471793in}{0.633527in}}%
\pgfpathlineto{\pgfqpoint{2.471963in}{0.625278in}}%
\pgfpathlineto{\pgfqpoint{2.472630in}{0.625278in}}%
\pgfpathlineto{\pgfqpoint{2.473807in}{0.633527in}}%
\pgfpathlineto{\pgfqpoint{2.473970in}{0.625278in}}%
\pgfpathlineto{\pgfqpoint{2.474814in}{0.625278in}}%
\pgfpathlineto{\pgfqpoint{2.475984in}{0.633527in}}%
\pgfpathlineto{\pgfqpoint{2.476154in}{0.625278in}}%
\pgfpathlineto{\pgfqpoint{2.476991in}{0.625278in}}%
\pgfpathlineto{\pgfqpoint{2.477998in}{0.633527in}}%
\pgfpathlineto{\pgfqpoint{2.477998in}{0.625278in}}%
\pgfpathlineto{\pgfqpoint{2.479168in}{0.633527in}}%
\pgfpathlineto{\pgfqpoint{2.479338in}{0.625278in}}%
\pgfpathlineto{\pgfqpoint{2.480004in}{0.625278in}}%
\pgfpathlineto{\pgfqpoint{2.481182in}{0.633527in}}%
\pgfpathlineto{\pgfqpoint{2.481352in}{0.625278in}}%
\pgfpathlineto{\pgfqpoint{2.482189in}{0.625278in}}%
\pgfpathlineto{\pgfqpoint{2.483358in}{0.633527in}}%
\pgfpathlineto{\pgfqpoint{2.483529in}{0.625278in}}%
\pgfpathlineto{\pgfqpoint{2.484365in}{0.625278in}}%
\pgfpathlineto{\pgfqpoint{2.485543in}{0.633527in}}%
\pgfpathlineto{\pgfqpoint{2.485594in}{0.625278in}}%
\pgfpathlineto{\pgfqpoint{2.486542in}{0.625278in}}%
\pgfpathlineto{\pgfqpoint{2.487719in}{0.633527in}}%
\pgfpathlineto{\pgfqpoint{2.487890in}{0.625278in}}%
\pgfpathlineto{\pgfqpoint{2.488726in}{0.625278in}}%
\pgfpathlineto{\pgfqpoint{2.489896in}{0.633527in}}%
\pgfpathlineto{\pgfqpoint{2.489896in}{0.625278in}}%
\pgfpathlineto{\pgfqpoint{2.490903in}{0.625278in}}%
\pgfpathlineto{\pgfqpoint{2.491910in}{0.633527in}}%
\pgfpathlineto{\pgfqpoint{2.491910in}{0.625278in}}%
\pgfpathlineto{\pgfqpoint{2.493250in}{0.633527in}}%
\pgfpathlineto{\pgfqpoint{2.493421in}{0.625278in}}%
\pgfpathlineto{\pgfqpoint{2.494257in}{0.625278in}}%
\pgfpathlineto{\pgfqpoint{2.495427in}{0.633527in}}%
\pgfpathlineto{\pgfqpoint{2.495427in}{0.625278in}}%
\pgfpathlineto{\pgfqpoint{2.496434in}{0.625278in}}%
\pgfpathlineto{\pgfqpoint{2.497612in}{0.633527in}}%
\pgfpathlineto{\pgfqpoint{2.497774in}{0.625278in}}%
\pgfpathlineto{\pgfqpoint{2.498611in}{0.625278in}}%
\pgfpathlineto{\pgfqpoint{2.499788in}{0.633527in}}%
\pgfpathlineto{\pgfqpoint{2.500122in}{0.625278in}}%
\pgfpathlineto{\pgfqpoint{2.500795in}{0.625278in}}%
\pgfpathlineto{\pgfqpoint{2.501965in}{0.633527in}}%
\pgfpathlineto{\pgfqpoint{2.501965in}{0.625278in}}%
\pgfpathlineto{\pgfqpoint{2.502972in}{0.625278in}}%
\pgfpathlineto{\pgfqpoint{2.504149in}{0.633527in}}%
\pgfpathlineto{\pgfqpoint{2.504312in}{0.625278in}}%
\pgfpathlineto{\pgfqpoint{2.504986in}{0.625278in}}%
\pgfpathlineto{\pgfqpoint{2.506156in}{0.633527in}}%
\pgfpathlineto{\pgfqpoint{2.506326in}{0.625278in}}%
\pgfpathlineto{\pgfqpoint{2.507163in}{0.625278in}}%
\pgfpathlineto{\pgfqpoint{2.508340in}{0.633527in}}%
\pgfpathlineto{\pgfqpoint{2.508340in}{0.625278in}}%
\pgfpathlineto{\pgfqpoint{2.509340in}{0.625278in}}%
\pgfpathlineto{\pgfqpoint{2.510517in}{0.633527in}}%
\pgfpathlineto{\pgfqpoint{2.510739in}{0.625278in}}%
\pgfpathlineto{\pgfqpoint{2.511524in}{0.625278in}}%
\pgfpathlineto{\pgfqpoint{2.512694in}{0.633527in}}%
\pgfpathlineto{\pgfqpoint{2.512864in}{0.625278in}}%
\pgfpathlineto{\pgfqpoint{2.513701in}{0.625278in}}%
\pgfpathlineto{\pgfqpoint{2.514878in}{0.633527in}}%
\pgfpathlineto{\pgfqpoint{2.515211in}{0.625278in}}%
\pgfpathlineto{\pgfqpoint{2.515878in}{0.625278in}}%
\pgfpathlineto{\pgfqpoint{2.516885in}{0.633527in}}%
\pgfpathlineto{\pgfqpoint{2.516885in}{0.625278in}}%
\pgfpathlineto{\pgfqpoint{2.518062in}{0.633527in}}%
\pgfpathlineto{\pgfqpoint{2.518225in}{0.625278in}}%
\pgfpathlineto{\pgfqpoint{2.519069in}{0.625278in}}%
\pgfpathlineto{\pgfqpoint{2.520239in}{0.633527in}}%
\pgfpathlineto{\pgfqpoint{2.520239in}{0.625278in}}%
\pgfpathlineto{\pgfqpoint{2.521246in}{0.625278in}}%
\pgfpathlineto{\pgfqpoint{2.522416in}{0.633527in}}%
\pgfpathlineto{\pgfqpoint{2.522586in}{0.625278in}}%
\pgfpathlineto{\pgfqpoint{2.523423in}{0.625278in}}%
\pgfpathlineto{\pgfqpoint{2.524600in}{0.633527in}}%
\pgfpathlineto{\pgfqpoint{2.524763in}{0.625278in}}%
\pgfpathlineto{\pgfqpoint{2.525607in}{0.625278in}}%
\pgfpathlineto{\pgfqpoint{2.526777in}{0.633527in}}%
\pgfpathlineto{\pgfqpoint{2.527110in}{0.625278in}}%
\pgfpathlineto{\pgfqpoint{2.527784in}{0.625278in}}%
\pgfpathlineto{\pgfqpoint{2.528953in}{0.633527in}}%
\pgfpathlineto{\pgfqpoint{2.529124in}{0.625278in}}%
\pgfpathlineto{\pgfqpoint{2.529960in}{0.625278in}}%
\pgfpathlineto{\pgfqpoint{2.531138in}{0.633527in}}%
\pgfpathlineto{\pgfqpoint{2.531301in}{0.625278in}}%
\pgfpathlineto{\pgfqpoint{2.531974in}{0.625278in}}%
\pgfpathlineto{\pgfqpoint{2.533144in}{0.633527in}}%
\pgfpathlineto{\pgfqpoint{2.533315in}{0.625278in}}%
\pgfpathlineto{\pgfqpoint{2.534151in}{0.625278in}}%
\pgfpathlineto{\pgfqpoint{2.535328in}{0.633527in}}%
\pgfpathlineto{\pgfqpoint{2.535328in}{0.625278in}}%
\pgfpathlineto{\pgfqpoint{2.536328in}{0.625278in}}%
\pgfpathlineto{\pgfqpoint{2.537505in}{0.633527in}}%
\pgfpathlineto{\pgfqpoint{2.537676in}{0.625278in}}%
\pgfpathlineto{\pgfqpoint{2.538512in}{0.625278in}}%
\pgfpathlineto{\pgfqpoint{2.539682in}{0.633527in}}%
\pgfpathlineto{\pgfqpoint{2.539852in}{0.625278in}}%
\pgfpathlineto{\pgfqpoint{2.540689in}{0.625278in}}%
\pgfpathlineto{\pgfqpoint{2.541866in}{0.633527in}}%
\pgfpathlineto{\pgfqpoint{2.542029in}{0.625278in}}%
\pgfpathlineto{\pgfqpoint{2.542866in}{0.625278in}}%
\pgfpathlineto{\pgfqpoint{2.543873in}{0.633527in}}%
\pgfpathlineto{\pgfqpoint{2.543873in}{0.625278in}}%
\pgfpathlineto{\pgfqpoint{2.545050in}{0.633527in}}%
\pgfpathlineto{\pgfqpoint{2.545213in}{0.625278in}}%
\pgfpathlineto{\pgfqpoint{2.546057in}{0.625278in}}%
\pgfpathlineto{\pgfqpoint{2.547227in}{0.633527in}}%
\pgfpathlineto{\pgfqpoint{2.547397in}{0.625278in}}%
\pgfpathlineto{\pgfqpoint{2.548234in}{0.625278in}}%
\pgfpathlineto{\pgfqpoint{2.549241in}{0.633527in}}%
\pgfpathlineto{\pgfqpoint{2.549241in}{0.625278in}}%
\pgfpathlineto{\pgfqpoint{2.550411in}{0.633527in}}%
\pgfpathlineto{\pgfqpoint{2.550581in}{0.625278in}}%
\pgfpathlineto{\pgfqpoint{2.551418in}{0.625278in}}%
\pgfpathlineto{\pgfqpoint{2.552595in}{0.633527in}}%
\pgfpathlineto{\pgfqpoint{2.552595in}{0.625278in}}%
\pgfpathlineto{\pgfqpoint{2.553595in}{0.625278in}}%
\pgfpathlineto{\pgfqpoint{2.554772in}{0.633527in}}%
\pgfpathlineto{\pgfqpoint{2.554942in}{0.625278in}}%
\pgfpathlineto{\pgfqpoint{2.555779in}{0.625278in}}%
\pgfpathlineto{\pgfqpoint{2.556949in}{0.633527in}}%
\pgfpathlineto{\pgfqpoint{2.557119in}{0.625278in}}%
\pgfpathlineto{\pgfqpoint{2.557956in}{0.625278in}}%
\pgfpathlineto{\pgfqpoint{2.559133in}{0.633527in}}%
\pgfpathlineto{\pgfqpoint{2.559296in}{0.625278in}}%
\pgfpathlineto{\pgfqpoint{2.560133in}{0.625278in}}%
\pgfpathlineto{\pgfqpoint{2.561310in}{0.633527in}}%
\pgfpathlineto{\pgfqpoint{2.561473in}{0.625278in}}%
\pgfpathlineto{\pgfqpoint{2.562317in}{0.625278in}}%
\pgfpathlineto{\pgfqpoint{2.563487in}{0.633527in}}%
\pgfpathlineto{\pgfqpoint{2.563657in}{0.625278in}}%
\pgfpathlineto{\pgfqpoint{2.564323in}{0.625278in}}%
\pgfpathlineto{\pgfqpoint{2.565501in}{0.633527in}}%
\pgfpathlineto{\pgfqpoint{2.565663in}{0.625278in}}%
\pgfpathlineto{\pgfqpoint{2.566508in}{0.625278in}}%
\pgfpathlineto{\pgfqpoint{2.567677in}{0.633527in}}%
\pgfpathlineto{\pgfqpoint{2.567677in}{0.625278in}}%
\pgfpathlineto{\pgfqpoint{2.568684in}{0.625278in}}%
\pgfpathlineto{\pgfqpoint{2.569854in}{0.633527in}}%
\pgfpathlineto{\pgfqpoint{2.570025in}{0.625278in}}%
\pgfpathlineto{\pgfqpoint{2.570861in}{0.625278in}}%
\pgfpathlineto{\pgfqpoint{2.572201in}{0.633527in}}%
\pgfpathlineto{\pgfqpoint{2.572201in}{0.625278in}}%
\pgfpathlineto{\pgfqpoint{2.573208in}{0.625278in}}%
\pgfpathlineto{\pgfqpoint{2.574386in}{0.633527in}}%
\pgfpathlineto{\pgfqpoint{2.574437in}{0.625278in}}%
\pgfpathlineto{\pgfqpoint{2.575393in}{0.625278in}}%
\pgfpathlineto{\pgfqpoint{2.576562in}{0.633527in}}%
\pgfpathlineto{\pgfqpoint{2.576733in}{0.625278in}}%
\pgfpathlineto{\pgfqpoint{2.577569in}{0.625278in}}%
\pgfpathlineto{\pgfqpoint{2.578910in}{0.633527in}}%
\pgfpathlineto{\pgfqpoint{2.579080in}{0.625278in}}%
\pgfpathlineto{\pgfqpoint{2.579917in}{0.625278in}}%
\pgfpathlineto{\pgfqpoint{2.581086in}{0.633527in}}%
\pgfpathlineto{\pgfqpoint{2.581257in}{0.625278in}}%
\pgfpathlineto{\pgfqpoint{2.582093in}{0.625278in}}%
\pgfpathlineto{\pgfqpoint{2.583271in}{0.633527in}}%
\pgfpathlineto{\pgfqpoint{2.583434in}{0.625278in}}%
\pgfpathlineto{\pgfqpoint{2.584278in}{0.625278in}}%
\pgfpathlineto{\pgfqpoint{2.585447in}{0.633527in}}%
\pgfpathlineto{\pgfqpoint{2.585618in}{0.625278in}}%
\pgfpathlineto{\pgfqpoint{2.586454in}{0.625278in}}%
\pgfpathlineto{\pgfqpoint{2.587624in}{0.633527in}}%
\pgfpathlineto{\pgfqpoint{2.587795in}{0.625278in}}%
\pgfpathlineto{\pgfqpoint{2.588468in}{0.625278in}}%
\pgfpathlineto{\pgfqpoint{2.589638in}{0.633527in}}%
\pgfpathlineto{\pgfqpoint{2.589660in}{0.625278in}}%
\pgfpathlineto{\pgfqpoint{2.590660in}{0.625278in}}%
\pgfpathlineto{\pgfqpoint{2.591985in}{0.658274in}}%
\pgfpathlineto{\pgfqpoint{2.592170in}{0.641776in}}%
\pgfpathlineto{\pgfqpoint{2.593326in}{0.625278in}}%
\pgfpathlineto{\pgfqpoint{2.593326in}{0.633527in}}%
\pgfpathlineto{\pgfqpoint{2.593829in}{0.658274in}}%
\pgfpathlineto{\pgfqpoint{2.594347in}{0.650025in}}%
\pgfpathlineto{\pgfqpoint{2.594666in}{0.633527in}}%
\pgfpathlineto{\pgfqpoint{2.594836in}{0.658274in}}%
\pgfpathlineto{\pgfqpoint{2.595169in}{0.658274in}}%
\pgfpathlineto{\pgfqpoint{2.595191in}{0.674773in}}%
\pgfpathlineto{\pgfqpoint{2.595858in}{0.650025in}}%
\pgfpathlineto{\pgfqpoint{2.596198in}{0.674773in}}%
\pgfpathlineto{\pgfqpoint{2.596532in}{0.699520in}}%
\pgfpathlineto{\pgfqpoint{2.597035in}{0.666524in}}%
\pgfpathlineto{\pgfqpoint{2.597539in}{0.650025in}}%
\pgfpathlineto{\pgfqpoint{2.597849in}{0.683022in}}%
\pgfpathlineto{\pgfqpoint{2.597872in}{0.674773in}}%
\pgfpathlineto{\pgfqpoint{2.598190in}{0.683022in}}%
\pgfpathlineto{\pgfqpoint{2.598538in}{0.650025in}}%
\pgfpathlineto{\pgfqpoint{2.599693in}{0.658274in}}%
\pgfpathlineto{\pgfqpoint{2.600389in}{0.674773in}}%
\pgfpathlineto{\pgfqpoint{2.600722in}{0.650025in}}%
\pgfpathlineto{\pgfqpoint{2.602211in}{0.683022in}}%
\pgfpathlineto{\pgfqpoint{2.602233in}{0.674773in}}%
\pgfpathlineto{\pgfqpoint{2.603047in}{0.683022in}}%
\pgfpathlineto{\pgfqpoint{2.603573in}{0.650025in}}%
\pgfpathlineto{\pgfqpoint{2.605565in}{0.707769in}}%
\pgfpathlineto{\pgfqpoint{2.606231in}{0.683022in}}%
\pgfpathlineto{\pgfqpoint{2.606586in}{0.699520in}}%
\pgfpathlineto{\pgfqpoint{2.606735in}{0.707769in}}%
\pgfpathlineto{\pgfqpoint{2.607090in}{0.674773in}}%
\pgfpathlineto{\pgfqpoint{2.607238in}{0.683022in}}%
\pgfpathlineto{\pgfqpoint{2.607260in}{0.674773in}}%
\pgfpathlineto{\pgfqpoint{2.607571in}{0.707769in}}%
\pgfpathlineto{\pgfqpoint{2.608267in}{0.691271in}}%
\pgfpathlineto{\pgfqpoint{2.608578in}{0.707769in}}%
\pgfpathlineto{\pgfqpoint{2.609437in}{0.674773in}}%
\pgfpathlineto{\pgfqpoint{2.611784in}{0.773763in}}%
\pgfpathlineto{\pgfqpoint{2.611954in}{0.765513in}}%
\pgfpathlineto{\pgfqpoint{2.612436in}{0.757264in}}%
\pgfpathlineto{\pgfqpoint{2.612606in}{0.782012in}}%
\pgfpathlineto{\pgfqpoint{2.612769in}{0.782012in}}%
\pgfpathlineto{\pgfqpoint{2.612939in}{0.806759in}}%
\pgfpathlineto{\pgfqpoint{2.613295in}{0.773763in}}%
\pgfpathlineto{\pgfqpoint{2.613798in}{0.798510in}}%
\pgfpathlineto{\pgfqpoint{2.614109in}{0.782012in}}%
\pgfpathlineto{\pgfqpoint{2.614279in}{0.806759in}}%
\pgfpathlineto{\pgfqpoint{2.614635in}{0.798510in}}%
\pgfpathlineto{\pgfqpoint{2.616960in}{0.881002in}}%
\pgfpathlineto{\pgfqpoint{2.616982in}{0.872753in}}%
\pgfpathlineto{\pgfqpoint{2.618322in}{0.848005in}}%
\pgfpathlineto{\pgfqpoint{2.619477in}{0.856254in}}%
\pgfpathlineto{\pgfqpoint{2.619995in}{0.872753in}}%
\pgfpathlineto{\pgfqpoint{2.620499in}{0.848005in}}%
\pgfpathlineto{\pgfqpoint{2.621676in}{0.872753in}}%
\pgfpathlineto{\pgfqpoint{2.621839in}{0.864503in}}%
\pgfpathlineto{\pgfqpoint{2.622661in}{0.856254in}}%
\pgfpathlineto{\pgfqpoint{2.622009in}{0.872753in}}%
\pgfpathlineto{\pgfqpoint{2.622661in}{0.864503in}}%
\pgfpathlineto{\pgfqpoint{2.623498in}{0.905749in}}%
\pgfpathlineto{\pgfqpoint{2.623683in}{0.897500in}}%
\pgfpathlineto{\pgfqpoint{2.624334in}{0.905749in}}%
\pgfpathlineto{\pgfqpoint{2.624527in}{0.889251in}}%
\pgfpathlineto{\pgfqpoint{2.624860in}{0.872753in}}%
\pgfpathlineto{\pgfqpoint{2.625178in}{0.905749in}}%
\pgfpathlineto{\pgfqpoint{2.625341in}{0.905749in}}%
\pgfpathlineto{\pgfqpoint{2.625512in}{0.930497in}}%
\pgfpathlineto{\pgfqpoint{2.626200in}{0.897500in}}%
\pgfpathlineto{\pgfqpoint{2.626370in}{0.897500in}}%
\pgfpathlineto{\pgfqpoint{2.627022in}{0.930497in}}%
\pgfpathlineto{\pgfqpoint{2.627377in}{0.897500in}}%
\pgfpathlineto{\pgfqpoint{2.628718in}{0.872753in}}%
\pgfpathlineto{\pgfqpoint{2.629369in}{0.881002in}}%
\pgfpathlineto{\pgfqpoint{2.629554in}{0.864503in}}%
\pgfpathlineto{\pgfqpoint{2.629887in}{0.848005in}}%
\pgfpathlineto{\pgfqpoint{2.630206in}{0.881002in}}%
\pgfpathlineto{\pgfqpoint{2.630221in}{0.872753in}}%
\pgfpathlineto{\pgfqpoint{2.630376in}{0.881002in}}%
\pgfpathlineto{\pgfqpoint{2.630724in}{0.848005in}}%
\pgfpathlineto{\pgfqpoint{2.631731in}{0.823258in}}%
\pgfpathlineto{\pgfqpoint{2.631879in}{0.839756in}}%
\pgfpathlineto{\pgfqpoint{2.632220in}{0.856254in}}%
\pgfpathlineto{\pgfqpoint{2.632553in}{0.831507in}}%
\pgfpathlineto{\pgfqpoint{2.632908in}{0.848005in}}%
\pgfpathlineto{\pgfqpoint{2.633893in}{0.806759in}}%
\pgfpathlineto{\pgfqpoint{2.634249in}{0.823258in}}%
\pgfpathlineto{\pgfqpoint{2.636240in}{0.881002in}}%
\pgfpathlineto{\pgfqpoint{2.636759in}{0.872753in}}%
\pgfpathlineto{\pgfqpoint{2.636929in}{0.897500in}}%
\pgfpathlineto{\pgfqpoint{2.637262in}{0.897500in}}%
\pgfpathlineto{\pgfqpoint{2.637580in}{0.881002in}}%
\pgfpathlineto{\pgfqpoint{2.637914in}{0.913998in}}%
\pgfpathlineto{\pgfqpoint{2.638921in}{0.930497in}}%
\pgfpathlineto{\pgfqpoint{2.638439in}{0.897500in}}%
\pgfpathlineto{\pgfqpoint{2.638943in}{0.922247in}}%
\pgfpathlineto{\pgfqpoint{2.640113in}{0.946995in}}%
\pgfpathlineto{\pgfqpoint{2.639594in}{0.905749in}}%
\pgfpathlineto{\pgfqpoint{2.640283in}{0.938746in}}%
\pgfpathlineto{\pgfqpoint{2.640935in}{0.905749in}}%
\pgfpathlineto{\pgfqpoint{2.641453in}{0.922247in}}%
\pgfpathlineto{\pgfqpoint{2.642275in}{0.955244in}}%
\pgfpathlineto{\pgfqpoint{2.642608in}{0.930497in}}%
\pgfpathlineto{\pgfqpoint{2.642963in}{0.922247in}}%
\pgfpathlineto{\pgfqpoint{2.643615in}{0.955244in}}%
\pgfpathlineto{\pgfqpoint{2.643637in}{0.946995in}}%
\pgfpathlineto{\pgfqpoint{2.643948in}{0.955244in}}%
\pgfpathlineto{\pgfqpoint{2.644303in}{0.922247in}}%
\pgfpathlineto{\pgfqpoint{2.644451in}{0.930497in}}%
\pgfpathlineto{\pgfqpoint{2.645310in}{0.922247in}}%
\pgfpathlineto{\pgfqpoint{2.645481in}{0.946995in}}%
\pgfpathlineto{\pgfqpoint{2.647302in}{1.004739in}}%
\pgfpathlineto{\pgfqpoint{2.647324in}{0.996490in}}%
\pgfpathlineto{\pgfqpoint{2.647828in}{0.971742in}}%
\pgfpathlineto{\pgfqpoint{2.648139in}{1.004739in}}%
\pgfpathlineto{\pgfqpoint{2.648161in}{0.996490in}}%
\pgfpathlineto{\pgfqpoint{2.648479in}{1.004739in}}%
\pgfpathlineto{\pgfqpoint{2.648827in}{0.971742in}}%
\pgfpathlineto{\pgfqpoint{2.650338in}{1.021237in}}%
\pgfpathlineto{\pgfqpoint{2.650508in}{1.012988in}}%
\pgfpathlineto{\pgfqpoint{2.651330in}{1.004739in}}%
\pgfpathlineto{\pgfqpoint{2.650678in}{1.021237in}}%
\pgfpathlineto{\pgfqpoint{2.651330in}{1.012988in}}%
\pgfpathlineto{\pgfqpoint{2.651345in}{1.045985in}}%
\pgfpathlineto{\pgfqpoint{2.652352in}{1.045985in}}%
\pgfpathlineto{\pgfqpoint{2.653840in}{0.979992in}}%
\pgfpathlineto{\pgfqpoint{2.654195in}{0.996490in}}%
\pgfpathlineto{\pgfqpoint{2.655017in}{1.029487in}}%
\pgfpathlineto{\pgfqpoint{2.655202in}{1.012988in}}%
\pgfpathlineto{\pgfqpoint{2.656372in}{0.971742in}}%
\pgfpathlineto{\pgfqpoint{2.656520in}{0.979992in}}%
\pgfpathlineto{\pgfqpoint{2.657046in}{1.045985in}}%
\pgfpathlineto{\pgfqpoint{2.657550in}{1.021237in}}%
\pgfpathlineto{\pgfqpoint{2.658364in}{1.054234in}}%
\pgfpathlineto{\pgfqpoint{2.658890in}{1.037736in}}%
\pgfpathlineto{\pgfqpoint{2.659371in}{1.029487in}}%
\pgfpathlineto{\pgfqpoint{2.659541in}{1.054234in}}%
\pgfpathlineto{\pgfqpoint{2.659556in}{1.045985in}}%
\pgfpathlineto{\pgfqpoint{2.659897in}{1.070732in}}%
\pgfpathlineto{\pgfqpoint{2.660400in}{1.045985in}}%
\pgfpathlineto{\pgfqpoint{2.660711in}{1.029487in}}%
\pgfpathlineto{\pgfqpoint{2.660881in}{1.054234in}}%
\pgfpathlineto{\pgfqpoint{2.661237in}{1.045985in}}%
\pgfpathlineto{\pgfqpoint{2.661718in}{1.029487in}}%
\pgfpathlineto{\pgfqpoint{2.662392in}{1.054234in}}%
\pgfpathlineto{\pgfqpoint{2.662910in}{1.045985in}}%
\pgfpathlineto{\pgfqpoint{2.663414in}{1.095480in}}%
\pgfpathlineto{\pgfqpoint{2.663732in}{1.078981in}}%
\pgfpathlineto{\pgfqpoint{2.663902in}{1.103729in}}%
\pgfpathlineto{\pgfqpoint{2.664250in}{1.095480in}}%
\pgfpathlineto{\pgfqpoint{2.664902in}{1.103729in}}%
\pgfpathlineto{\pgfqpoint{2.665094in}{1.087231in}}%
\pgfpathlineto{\pgfqpoint{2.666583in}{1.054234in}}%
\pgfpathlineto{\pgfqpoint{2.666583in}{1.062483in}}%
\pgfpathlineto{\pgfqpoint{2.667590in}{1.078981in}}%
\pgfpathlineto{\pgfqpoint{2.666938in}{1.045985in}}%
\pgfpathlineto{\pgfqpoint{2.667604in}{1.070732in}}%
\pgfpathlineto{\pgfqpoint{2.669263in}{1.103729in}}%
\pgfpathlineto{\pgfqpoint{2.670270in}{1.078981in}}%
\pgfpathlineto{\pgfqpoint{2.670285in}{1.095480in}}%
\pgfpathlineto{\pgfqpoint{2.670959in}{1.070732in}}%
\pgfpathlineto{\pgfqpoint{2.671129in}{1.095480in}}%
\pgfpathlineto{\pgfqpoint{2.672284in}{1.103729in}}%
\pgfpathlineto{\pgfqpoint{2.673306in}{1.070732in}}%
\pgfpathlineto{\pgfqpoint{2.675127in}{1.128476in}}%
\pgfpathlineto{\pgfqpoint{2.675149in}{1.120227in}}%
\pgfpathlineto{\pgfqpoint{2.675801in}{1.103729in}}%
\pgfpathlineto{\pgfqpoint{2.675801in}{1.136726in}}%
\pgfpathlineto{\pgfqpoint{2.675823in}{1.144975in}}%
\pgfpathlineto{\pgfqpoint{2.676638in}{1.103729in}}%
\pgfpathlineto{\pgfqpoint{2.676823in}{1.120227in}}%
\pgfpathlineto{\pgfqpoint{2.677978in}{1.128476in}}%
\pgfpathlineto{\pgfqpoint{2.678985in}{1.103729in}}%
\pgfpathlineto{\pgfqpoint{2.679007in}{1.120227in}}%
\pgfpathlineto{\pgfqpoint{2.680347in}{1.095480in}}%
\pgfpathlineto{\pgfqpoint{2.682006in}{1.128476in}}%
\pgfpathlineto{\pgfqpoint{2.683027in}{1.120227in}}%
\pgfpathlineto{\pgfqpoint{2.684034in}{1.144975in}}%
\pgfpathlineto{\pgfqpoint{2.684205in}{1.136726in}}%
\pgfpathlineto{\pgfqpoint{2.684538in}{1.120227in}}%
\pgfpathlineto{\pgfqpoint{2.684856in}{1.153224in}}%
\pgfpathlineto{\pgfqpoint{2.685041in}{1.144975in}}%
\pgfpathlineto{\pgfqpoint{2.685523in}{1.128476in}}%
\pgfpathlineto{\pgfqpoint{2.686700in}{1.177971in}}%
\pgfpathlineto{\pgfqpoint{2.687388in}{1.194470in}}%
\pgfpathlineto{\pgfqpoint{2.687722in}{1.169722in}}%
\pgfpathlineto{\pgfqpoint{2.689380in}{1.128476in}}%
\pgfpathlineto{\pgfqpoint{2.689380in}{1.136726in}}%
\pgfpathlineto{\pgfqpoint{2.690069in}{1.194470in}}%
\pgfpathlineto{\pgfqpoint{2.689713in}{1.128476in}}%
\pgfpathlineto{\pgfqpoint{2.690550in}{1.177971in}}%
\pgfpathlineto{\pgfqpoint{2.690905in}{1.194470in}}%
\pgfpathlineto{\pgfqpoint{2.691579in}{1.169722in}}%
\pgfpathlineto{\pgfqpoint{2.691742in}{1.194470in}}%
\pgfpathlineto{\pgfqpoint{2.692394in}{1.153224in}}%
\pgfpathlineto{\pgfqpoint{2.693423in}{1.120227in}}%
\pgfpathlineto{\pgfqpoint{2.693571in}{1.128476in}}%
\pgfpathlineto{\pgfqpoint{2.693756in}{1.169722in}}%
\pgfpathlineto{\pgfqpoint{2.694260in}{1.120227in}}%
\pgfpathlineto{\pgfqpoint{2.694593in}{1.120227in}}%
\pgfpathlineto{\pgfqpoint{2.695600in}{1.144975in}}%
\pgfpathlineto{\pgfqpoint{2.695770in}{1.136726in}}%
\pgfpathlineto{\pgfqpoint{2.695933in}{1.169722in}}%
\pgfpathlineto{\pgfqpoint{2.696925in}{1.128476in}}%
\pgfpathlineto{\pgfqpoint{2.697429in}{1.153224in}}%
\pgfpathlineto{\pgfqpoint{2.697947in}{1.144975in}}%
\pgfpathlineto{\pgfqpoint{2.698598in}{1.128476in}}%
\pgfpathlineto{\pgfqpoint{2.698606in}{1.161473in}}%
\pgfpathlineto{\pgfqpoint{2.698621in}{1.169722in}}%
\pgfpathlineto{\pgfqpoint{2.699287in}{1.120227in}}%
\pgfpathlineto{\pgfqpoint{2.699620in}{1.120227in}}%
\pgfpathlineto{\pgfqpoint{2.700449in}{1.153224in}}%
\pgfpathlineto{\pgfqpoint{2.700783in}{1.128476in}}%
\pgfpathlineto{\pgfqpoint{2.701131in}{1.120227in}}%
\pgfpathlineto{\pgfqpoint{2.701790in}{1.186221in}}%
\pgfpathlineto{\pgfqpoint{2.702138in}{1.194470in}}%
\pgfpathlineto{\pgfqpoint{2.702811in}{1.144975in}}%
\pgfpathlineto{\pgfqpoint{2.703626in}{1.153224in}}%
\pgfpathlineto{\pgfqpoint{2.704655in}{1.095480in}}%
\pgfpathlineto{\pgfqpoint{2.705810in}{1.103729in}}%
\pgfpathlineto{\pgfqpoint{2.706662in}{1.070732in}}%
\pgfpathlineto{\pgfqpoint{2.706832in}{1.120227in}}%
\pgfpathlineto{\pgfqpoint{2.709682in}{1.021237in}}%
\pgfpathlineto{\pgfqpoint{2.710667in}{1.029487in}}%
\pgfpathlineto{\pgfqpoint{2.710334in}{1.004739in}}%
\pgfpathlineto{\pgfqpoint{2.710689in}{1.021237in}}%
\pgfpathlineto{\pgfqpoint{2.713014in}{0.955244in}}%
\pgfpathlineto{\pgfqpoint{2.713858in}{0.979992in}}%
\pgfpathlineto{\pgfqpoint{2.714036in}{0.971742in}}%
\pgfpathlineto{\pgfqpoint{2.714540in}{0.996490in}}%
\pgfpathlineto{\pgfqpoint{2.715191in}{0.979992in}}%
\pgfpathlineto{\pgfqpoint{2.716383in}{0.938746in}}%
\pgfpathlineto{\pgfqpoint{2.717731in}{0.897500in}}%
\pgfpathlineto{\pgfqpoint{2.719574in}{0.848005in}}%
\pgfpathlineto{\pgfqpoint{2.720063in}{0.856254in}}%
\pgfpathlineto{\pgfqpoint{2.720241in}{0.839756in}}%
\pgfpathlineto{\pgfqpoint{2.721751in}{0.773763in}}%
\pgfpathlineto{\pgfqpoint{2.720411in}{0.848005in}}%
\pgfpathlineto{\pgfqpoint{2.721922in}{0.798510in}}%
\pgfpathlineto{\pgfqpoint{2.722240in}{0.806759in}}%
\pgfpathlineto{\pgfqpoint{2.722573in}{0.782012in}}%
\pgfpathlineto{\pgfqpoint{2.723410in}{0.757264in}}%
\pgfpathlineto{\pgfqpoint{2.723595in}{0.765513in}}%
\pgfpathlineto{\pgfqpoint{2.725942in}{0.674773in}}%
\pgfpathlineto{\pgfqpoint{2.726098in}{0.691271in}}%
\pgfpathlineto{\pgfqpoint{2.727282in}{0.773763in}}%
\pgfpathlineto{\pgfqpoint{2.727956in}{0.749015in}}%
\pgfpathlineto{\pgfqpoint{2.728119in}{0.773763in}}%
\pgfpathlineto{\pgfqpoint{2.729126in}{0.798510in}}%
\pgfpathlineto{\pgfqpoint{2.729296in}{0.790261in}}%
\pgfpathlineto{\pgfqpoint{2.730799in}{0.749015in}}%
\pgfpathlineto{\pgfqpoint{2.730844in}{0.757264in}}%
\pgfpathlineto{\pgfqpoint{2.731125in}{0.806759in}}%
\pgfpathlineto{\pgfqpoint{2.732147in}{0.773763in}}%
\pgfpathlineto{\pgfqpoint{2.733154in}{0.757264in}}%
\pgfpathlineto{\pgfqpoint{2.732480in}{0.798510in}}%
\pgfpathlineto{\pgfqpoint{2.733294in}{0.765513in}}%
\pgfpathlineto{\pgfqpoint{2.733650in}{0.773763in}}%
\pgfpathlineto{\pgfqpoint{2.734301in}{0.732517in}}%
\pgfpathlineto{\pgfqpoint{2.734324in}{0.749015in}}%
\pgfpathlineto{\pgfqpoint{2.735479in}{0.757264in}}%
\pgfpathlineto{\pgfqpoint{2.736500in}{0.749015in}}%
\pgfpathlineto{\pgfqpoint{2.735834in}{0.773763in}}%
\pgfpathlineto{\pgfqpoint{2.736508in}{0.757264in}}%
\pgfpathlineto{\pgfqpoint{2.737167in}{0.782012in}}%
\pgfpathlineto{\pgfqpoint{2.737500in}{0.757264in}}%
\pgfpathlineto{\pgfqpoint{2.737678in}{0.732517in}}%
\pgfpathlineto{\pgfqpoint{2.738337in}{0.806759in}}%
\pgfpathlineto{\pgfqpoint{2.738514in}{0.798510in}}%
\pgfpathlineto{\pgfqpoint{2.740528in}{0.699520in}}%
\pgfpathlineto{\pgfqpoint{2.740847in}{0.732517in}}%
\pgfpathlineto{\pgfqpoint{2.741195in}{0.773763in}}%
\pgfpathlineto{\pgfqpoint{2.741868in}{0.749015in}}%
\pgfpathlineto{\pgfqpoint{2.743038in}{0.699520in}}%
\pgfpathlineto{\pgfqpoint{2.743527in}{0.716019in}}%
\pgfpathlineto{\pgfqpoint{2.744216in}{0.757264in}}%
\pgfpathlineto{\pgfqpoint{2.744549in}{0.740766in}}%
\pgfpathlineto{\pgfqpoint{2.745889in}{0.683022in}}%
\pgfpathlineto{\pgfqpoint{2.746044in}{0.716019in}}%
\pgfpathlineto{\pgfqpoint{2.746896in}{0.749015in}}%
\pgfpathlineto{\pgfqpoint{2.746230in}{0.707769in}}%
\pgfpathlineto{\pgfqpoint{2.747066in}{0.732517in}}%
\pgfpathlineto{\pgfqpoint{2.747562in}{0.724268in}}%
\pgfpathlineto{\pgfqpoint{2.747733in}{0.749015in}}%
\pgfpathlineto{\pgfqpoint{2.747888in}{0.757264in}}%
\pgfpathlineto{\pgfqpoint{2.748066in}{0.740766in}}%
\pgfpathlineto{\pgfqpoint{2.749406in}{0.699520in}}%
\pgfpathlineto{\pgfqpoint{2.751760in}{0.625278in}}%
\pgfpathlineto{\pgfqpoint{2.752767in}{0.633527in}}%
\pgfpathlineto{\pgfqpoint{2.752767in}{0.625278in}}%
\pgfpathlineto{\pgfqpoint{2.754100in}{0.633527in}}%
\pgfpathlineto{\pgfqpoint{2.754108in}{0.625278in}}%
\pgfpathlineto{\pgfqpoint{2.755115in}{0.625278in}}%
\pgfpathlineto{\pgfqpoint{2.756277in}{0.633527in}}%
\pgfpathlineto{\pgfqpoint{2.756447in}{0.625278in}}%
\pgfpathlineto{\pgfqpoint{2.757114in}{0.625278in}}%
\pgfpathlineto{\pgfqpoint{2.758284in}{0.633527in}}%
\pgfpathlineto{\pgfqpoint{2.758454in}{0.625278in}}%
\pgfpathlineto{\pgfqpoint{2.759335in}{0.625278in}}%
\pgfpathlineto{\pgfqpoint{2.760475in}{0.633527in}}%
\pgfpathlineto{\pgfqpoint{2.760638in}{0.625278in}}%
\pgfpathlineto{\pgfqpoint{2.761460in}{0.625278in}}%
\pgfpathlineto{\pgfqpoint{2.762489in}{0.633527in}}%
\pgfpathlineto{\pgfqpoint{2.762489in}{0.625278in}}%
\pgfpathlineto{\pgfqpoint{2.763637in}{0.633527in}}%
\pgfpathlineto{\pgfqpoint{2.763829in}{0.625278in}}%
\pgfpathlineto{\pgfqpoint{2.764481in}{0.625278in}}%
\pgfpathlineto{\pgfqpoint{2.765666in}{0.633527in}}%
\pgfpathlineto{\pgfqpoint{2.765828in}{0.625278in}}%
\pgfpathlineto{\pgfqpoint{2.766673in}{0.625278in}}%
\pgfpathlineto{\pgfqpoint{2.767835in}{0.633527in}}%
\pgfpathlineto{\pgfqpoint{2.768005in}{0.625278in}}%
\pgfpathlineto{\pgfqpoint{2.768679in}{0.625278in}}%
\pgfpathlineto{\pgfqpoint{2.769856in}{0.633527in}}%
\pgfpathlineto{\pgfqpoint{2.770019in}{0.625278in}}%
\pgfpathlineto{\pgfqpoint{2.770856in}{0.625278in}}%
\pgfpathlineto{\pgfqpoint{2.772026in}{0.633527in}}%
\pgfpathlineto{\pgfqpoint{2.772196in}{0.625278in}}%
\pgfpathlineto{\pgfqpoint{2.773033in}{0.625278in}}%
\pgfpathlineto{\pgfqpoint{2.774040in}{0.633527in}}%
\pgfpathlineto{\pgfqpoint{2.774040in}{0.625278in}}%
\pgfpathlineto{\pgfqpoint{2.775380in}{0.633527in}}%
\pgfpathlineto{\pgfqpoint{2.776217in}{0.625278in}}%
\pgfpathlineto{\pgfqpoint{2.777386in}{0.633527in}}%
\pgfpathlineto{\pgfqpoint{2.777549in}{0.625278in}}%
\pgfpathlineto{\pgfqpoint{2.778230in}{0.625278in}}%
\pgfpathlineto{\pgfqpoint{2.779415in}{0.633527in}}%
\pgfpathlineto{\pgfqpoint{2.779563in}{0.625278in}}%
\pgfpathlineto{\pgfqpoint{2.780237in}{0.625278in}}%
\pgfpathlineto{\pgfqpoint{2.781577in}{0.633527in}}%
\pgfpathlineto{\pgfqpoint{2.781592in}{0.625278in}}%
\pgfpathlineto{\pgfqpoint{2.782599in}{0.625278in}}%
\pgfpathlineto{\pgfqpoint{2.783924in}{0.633527in}}%
\pgfpathlineto{\pgfqpoint{2.783932in}{0.625278in}}%
\pgfpathlineto{\pgfqpoint{2.784931in}{0.625278in}}%
\pgfpathlineto{\pgfqpoint{2.786101in}{0.633527in}}%
\pgfpathlineto{\pgfqpoint{2.786775in}{0.625278in}}%
\pgfpathlineto{\pgfqpoint{2.787108in}{0.625278in}}%
\pgfpathlineto{\pgfqpoint{2.788278in}{0.633527in}}%
\pgfpathlineto{\pgfqpoint{2.788448in}{0.625278in}}%
\pgfpathlineto{\pgfqpoint{2.789285in}{0.625278in}}%
\pgfpathlineto{\pgfqpoint{2.790462in}{0.633527in}}%
\pgfpathlineto{\pgfqpoint{2.790462in}{0.625278in}}%
\pgfpathlineto{\pgfqpoint{2.791469in}{0.625278in}}%
\pgfpathlineto{\pgfqpoint{2.792639in}{0.633527in}}%
\pgfpathlineto{\pgfqpoint{2.792980in}{0.625278in}}%
\pgfpathlineto{\pgfqpoint{2.793653in}{0.625278in}}%
\pgfpathlineto{\pgfqpoint{2.794986in}{0.633527in}}%
\pgfpathlineto{\pgfqpoint{2.795319in}{0.625278in}}%
\pgfpathlineto{\pgfqpoint{2.795993in}{0.625278in}}%
\pgfpathlineto{\pgfqpoint{2.797163in}{0.633527in}}%
\pgfpathlineto{\pgfqpoint{2.797333in}{0.625278in}}%
\pgfpathlineto{\pgfqpoint{2.798007in}{0.625278in}}%
\pgfpathlineto{\pgfqpoint{2.799347in}{0.633527in}}%
\pgfpathlineto{\pgfqpoint{2.799347in}{0.625278in}}%
\pgfpathlineto{\pgfqpoint{2.800347in}{0.625278in}}%
\pgfpathlineto{\pgfqpoint{2.801524in}{0.633527in}}%
\pgfpathlineto{\pgfqpoint{2.801694in}{0.625278in}}%
\pgfpathlineto{\pgfqpoint{2.802531in}{0.625278in}}%
\pgfpathlineto{\pgfqpoint{2.803871in}{0.633527in}}%
\pgfpathlineto{\pgfqpoint{2.803871in}{0.625278in}}%
\pgfpathlineto{\pgfqpoint{2.804878in}{0.625278in}}%
\pgfpathlineto{\pgfqpoint{2.805885in}{0.633527in}}%
\pgfpathlineto{\pgfqpoint{2.805885in}{0.625278in}}%
\pgfpathlineto{\pgfqpoint{2.807225in}{0.633527in}}%
\pgfpathlineto{\pgfqpoint{2.807388in}{0.625278in}}%
\pgfpathlineto{\pgfqpoint{2.808232in}{0.625278in}}%
\pgfpathlineto{\pgfqpoint{2.809402in}{0.633527in}}%
\pgfpathlineto{\pgfqpoint{2.809735in}{0.625278in}}%
\pgfpathlineto{\pgfqpoint{2.810409in}{0.625278in}}%
\pgfpathlineto{\pgfqpoint{2.811579in}{0.633527in}}%
\pgfpathlineto{\pgfqpoint{2.811579in}{0.625278in}}%
\pgfpathlineto{\pgfqpoint{2.812423in}{0.625278in}}%
\pgfpathlineto{\pgfqpoint{2.813593in}{0.633527in}}%
\pgfpathlineto{\pgfqpoint{2.813763in}{0.625278in}}%
\pgfpathlineto{\pgfqpoint{2.814600in}{0.625278in}}%
\pgfpathlineto{\pgfqpoint{2.815770in}{0.633527in}}%
\pgfpathlineto{\pgfqpoint{2.815940in}{0.625278in}}%
\pgfpathlineto{\pgfqpoint{2.816777in}{0.625278in}}%
\pgfpathlineto{\pgfqpoint{2.817954in}{0.633527in}}%
\pgfpathlineto{\pgfqpoint{2.818117in}{0.625278in}}%
\pgfpathlineto{\pgfqpoint{2.818791in}{0.625278in}}%
\pgfpathlineto{\pgfqpoint{2.819961in}{0.633527in}}%
\pgfpathlineto{\pgfqpoint{2.820301in}{0.625278in}}%
\pgfpathlineto{\pgfqpoint{2.820805in}{0.625278in}}%
\pgfpathlineto{\pgfqpoint{2.821974in}{0.633527in}}%
\pgfpathlineto{\pgfqpoint{2.821974in}{0.625278in}}%
\pgfpathlineto{\pgfqpoint{2.822981in}{0.625278in}}%
\pgfpathlineto{\pgfqpoint{2.824322in}{0.633527in}}%
\pgfpathlineto{\pgfqpoint{2.824655in}{0.625278in}}%
\pgfpathlineto{\pgfqpoint{2.825329in}{0.625278in}}%
\pgfpathlineto{\pgfqpoint{2.826498in}{0.633527in}}%
\pgfpathlineto{\pgfqpoint{2.826669in}{0.625278in}}%
\pgfpathlineto{\pgfqpoint{2.827505in}{0.625278in}}%
\pgfpathlineto{\pgfqpoint{2.828683in}{0.633527in}}%
\pgfpathlineto{\pgfqpoint{2.828846in}{0.625278in}}%
\pgfpathlineto{\pgfqpoint{2.829682in}{0.625278in}}%
\pgfpathlineto{\pgfqpoint{2.830860in}{0.633527in}}%
\pgfpathlineto{\pgfqpoint{2.831030in}{0.625278in}}%
\pgfpathlineto{\pgfqpoint{2.831866in}{0.625278in}}%
\pgfpathlineto{\pgfqpoint{2.833036in}{0.633527in}}%
\pgfpathlineto{\pgfqpoint{2.833207in}{0.625278in}}%
\pgfpathlineto{\pgfqpoint{2.834043in}{0.625278in}}%
\pgfpathlineto{\pgfqpoint{2.835050in}{0.633527in}}%
\pgfpathlineto{\pgfqpoint{2.835050in}{0.625278in}}%
\pgfpathlineto{\pgfqpoint{2.836390in}{0.633527in}}%
\pgfpathlineto{\pgfqpoint{2.836390in}{0.625278in}}%
\pgfpathlineto{\pgfqpoint{2.837397in}{0.625278in}}%
\pgfpathlineto{\pgfqpoint{2.838738in}{0.633527in}}%
\pgfpathlineto{\pgfqpoint{2.838908in}{0.625278in}}%
\pgfpathlineto{\pgfqpoint{2.839745in}{0.625278in}}%
\pgfpathlineto{\pgfqpoint{2.840752in}{0.633527in}}%
\pgfpathlineto{\pgfqpoint{2.840752in}{0.625278in}}%
\pgfpathlineto{\pgfqpoint{2.841921in}{0.633527in}}%
\pgfpathlineto{\pgfqpoint{2.842092in}{0.625278in}}%
\pgfpathlineto{\pgfqpoint{2.842928in}{0.625278in}}%
\pgfpathlineto{\pgfqpoint{2.844106in}{0.633527in}}%
\pgfpathlineto{\pgfqpoint{2.844439in}{0.625278in}}%
\pgfpathlineto{\pgfqpoint{2.845105in}{0.625278in}}%
\pgfpathlineto{\pgfqpoint{2.846445in}{0.633527in}}%
\pgfpathlineto{\pgfqpoint{2.846445in}{0.625278in}}%
\pgfpathlineto{\pgfqpoint{2.847452in}{0.625278in}}%
\pgfpathlineto{\pgfqpoint{2.848630in}{0.633527in}}%
\pgfpathlineto{\pgfqpoint{2.848963in}{0.625278in}}%
\pgfpathlineto{\pgfqpoint{2.849637in}{0.625278in}}%
\pgfpathlineto{\pgfqpoint{2.850806in}{0.633527in}}%
\pgfpathlineto{\pgfqpoint{2.850977in}{0.625278in}}%
\pgfpathlineto{\pgfqpoint{2.851813in}{0.625278in}}%
\pgfpathlineto{\pgfqpoint{2.852820in}{0.633527in}}%
\pgfpathlineto{\pgfqpoint{2.852820in}{0.625278in}}%
\pgfpathlineto{\pgfqpoint{2.854161in}{0.633527in}}%
\pgfpathlineto{\pgfqpoint{2.854331in}{0.625278in}}%
\pgfpathlineto{\pgfqpoint{2.855167in}{0.625278in}}%
\pgfpathlineto{\pgfqpoint{2.856337in}{0.633527in}}%
\pgfpathlineto{\pgfqpoint{2.856337in}{0.625278in}}%
\pgfpathlineto{\pgfqpoint{2.857344in}{0.625278in}}%
\pgfpathlineto{\pgfqpoint{2.858522in}{0.633527in}}%
\pgfpathlineto{\pgfqpoint{2.858684in}{0.625278in}}%
\pgfpathlineto{\pgfqpoint{2.859521in}{0.625278in}}%
\pgfpathlineto{\pgfqpoint{2.860698in}{0.633527in}}%
\pgfpathlineto{\pgfqpoint{2.861032in}{0.625278in}}%
\pgfpathlineto{\pgfqpoint{2.861705in}{0.625278in}}%
\pgfpathlineto{\pgfqpoint{2.862875in}{0.633527in}}%
\pgfpathlineto{\pgfqpoint{2.863046in}{0.625278in}}%
\pgfpathlineto{\pgfqpoint{2.863882in}{0.625278in}}%
\pgfpathlineto{\pgfqpoint{2.865059in}{0.633527in}}%
\pgfpathlineto{\pgfqpoint{2.865222in}{0.625278in}}%
\pgfpathlineto{\pgfqpoint{2.866059in}{0.625278in}}%
\pgfpathlineto{\pgfqpoint{2.867066in}{0.633527in}}%
\pgfpathlineto{\pgfqpoint{2.867066in}{0.625278in}}%
\pgfpathlineto{\pgfqpoint{2.868243in}{0.633527in}}%
\pgfpathlineto{\pgfqpoint{2.868243in}{0.625278in}}%
\pgfpathlineto{\pgfqpoint{2.869250in}{0.625278in}}%
\pgfpathlineto{\pgfqpoint{2.870250in}{0.633527in}}%
\pgfpathlineto{\pgfqpoint{2.870250in}{0.625278in}}%
\pgfpathlineto{\pgfqpoint{2.871590in}{0.633527in}}%
\pgfpathlineto{\pgfqpoint{2.871760in}{0.625278in}}%
\pgfpathlineto{\pgfqpoint{2.872597in}{0.625278in}}%
\pgfpathlineto{\pgfqpoint{2.873604in}{0.633527in}}%
\pgfpathlineto{\pgfqpoint{2.873604in}{0.625278in}}%
\pgfpathlineto{\pgfqpoint{2.874944in}{0.633527in}}%
\pgfpathlineto{\pgfqpoint{2.874944in}{0.625278in}}%
\pgfpathlineto{\pgfqpoint{2.875951in}{0.625278in}}%
\pgfpathlineto{\pgfqpoint{2.876958in}{0.633527in}}%
\pgfpathlineto{\pgfqpoint{2.876958in}{0.625278in}}%
\pgfpathlineto{\pgfqpoint{2.878298in}{0.633527in}}%
\pgfpathlineto{\pgfqpoint{2.878298in}{0.625278in}}%
\pgfpathlineto{\pgfqpoint{2.879305in}{0.625278in}}%
\pgfpathlineto{\pgfqpoint{2.880475in}{0.633527in}}%
\pgfpathlineto{\pgfqpoint{2.880475in}{0.625278in}}%
\pgfpathlineto{\pgfqpoint{2.881482in}{0.625278in}}%
\pgfpathlineto{\pgfqpoint{2.882659in}{0.633527in}}%
\pgfpathlineto{\pgfqpoint{2.882822in}{0.625278in}}%
\pgfpathlineto{\pgfqpoint{2.883496in}{0.625278in}}%
\pgfpathlineto{\pgfqpoint{2.884666in}{0.633527in}}%
\pgfpathlineto{\pgfqpoint{2.884666in}{0.625278in}}%
\pgfpathlineto{\pgfqpoint{2.885673in}{0.625278in}}%
\pgfpathlineto{\pgfqpoint{2.886680in}{0.633527in}}%
\pgfpathlineto{\pgfqpoint{2.886680in}{0.625278in}}%
\pgfpathlineto{\pgfqpoint{2.888020in}{0.633527in}}%
\pgfpathlineto{\pgfqpoint{2.888020in}{0.625278in}}%
\pgfpathlineto{\pgfqpoint{2.889027in}{0.625278in}}%
\pgfpathlineto{\pgfqpoint{2.890204in}{0.633527in}}%
\pgfpathlineto{\pgfqpoint{2.890204in}{0.625278in}}%
\pgfpathlineto{\pgfqpoint{2.891204in}{0.625278in}}%
\pgfpathlineto{\pgfqpoint{2.892381in}{0.633527in}}%
\pgfpathlineto{\pgfqpoint{2.892544in}{0.625278in}}%
\pgfpathlineto{\pgfqpoint{2.893388in}{0.625278in}}%
\pgfpathlineto{\pgfqpoint{2.894395in}{0.633527in}}%
\pgfpathlineto{\pgfqpoint{2.894395in}{0.625278in}}%
\pgfpathlineto{\pgfqpoint{2.895735in}{0.633527in}}%
\pgfpathlineto{\pgfqpoint{2.895735in}{0.625278in}}%
\pgfpathlineto{\pgfqpoint{2.896735in}{0.625278in}}%
\pgfpathlineto{\pgfqpoint{2.897912in}{0.633527in}}%
\pgfpathlineto{\pgfqpoint{2.898082in}{0.625278in}}%
\pgfpathlineto{\pgfqpoint{2.898919in}{0.625278in}}%
\pgfpathlineto{\pgfqpoint{2.900089in}{0.633527in}}%
\pgfpathlineto{\pgfqpoint{2.900259in}{0.625278in}}%
\pgfpathlineto{\pgfqpoint{2.900925in}{0.625278in}}%
\pgfpathlineto{\pgfqpoint{2.902273in}{0.633527in}}%
\pgfpathlineto{\pgfqpoint{2.902273in}{0.625278in}}%
\pgfpathlineto{\pgfqpoint{2.903273in}{0.625278in}}%
\pgfpathlineto{\pgfqpoint{2.904620in}{0.633527in}}%
\pgfpathlineto{\pgfqpoint{2.904620in}{0.625278in}}%
\pgfpathlineto{\pgfqpoint{2.905620in}{0.625278in}}%
\pgfpathlineto{\pgfqpoint{2.906797in}{0.633527in}}%
\pgfpathlineto{\pgfqpoint{2.906797in}{0.625278in}}%
\pgfpathlineto{\pgfqpoint{2.907804in}{0.625278in}}%
\pgfpathlineto{\pgfqpoint{2.908974in}{0.633527in}}%
\pgfpathlineto{\pgfqpoint{2.909144in}{0.625278in}}%
\pgfpathlineto{\pgfqpoint{2.909981in}{0.625278in}}%
\pgfpathlineto{\pgfqpoint{2.911158in}{0.633527in}}%
\pgfpathlineto{\pgfqpoint{2.911321in}{0.625278in}}%
\pgfpathlineto{\pgfqpoint{2.912158in}{0.625278in}}%
\pgfpathlineto{\pgfqpoint{2.913335in}{0.633527in}}%
\pgfpathlineto{\pgfqpoint{2.913498in}{0.625278in}}%
\pgfpathlineto{\pgfqpoint{2.914342in}{0.625278in}}%
\pgfpathlineto{\pgfqpoint{2.915512in}{0.633527in}}%
\pgfpathlineto{\pgfqpoint{2.915682in}{0.625278in}}%
\pgfpathlineto{\pgfqpoint{2.916541in}{0.625278in}}%
\pgfpathlineto{\pgfqpoint{2.917689in}{0.633527in}}%
\pgfpathlineto{\pgfqpoint{2.917859in}{0.625278in}}%
\pgfpathlineto{\pgfqpoint{2.918547in}{0.625278in}}%
\pgfpathlineto{\pgfqpoint{2.919873in}{0.633527in}}%
\pgfpathlineto{\pgfqpoint{2.920036in}{0.625278in}}%
\pgfpathlineto{\pgfqpoint{2.920895in}{0.625278in}}%
\pgfpathlineto{\pgfqpoint{2.922050in}{0.633527in}}%
\pgfpathlineto{\pgfqpoint{2.922220in}{0.625278in}}%
\pgfpathlineto{\pgfqpoint{2.923057in}{0.625278in}}%
\pgfpathlineto{\pgfqpoint{2.924226in}{0.633527in}}%
\pgfpathlineto{\pgfqpoint{2.924397in}{0.625278in}}%
\pgfpathlineto{\pgfqpoint{2.925233in}{0.625278in}}%
\pgfpathlineto{\pgfqpoint{2.926411in}{0.633527in}}%
\pgfpathlineto{\pgfqpoint{2.926425in}{0.625278in}}%
\pgfpathlineto{\pgfqpoint{2.927099in}{0.674773in}}%
\pgfpathlineto{\pgfqpoint{2.927432in}{0.674773in}}%
\pgfpathlineto{\pgfqpoint{2.927751in}{0.658274in}}%
\pgfpathlineto{\pgfqpoint{2.927921in}{0.683022in}}%
\pgfpathlineto{\pgfqpoint{2.928269in}{0.674773in}}%
\pgfpathlineto{\pgfqpoint{2.928921in}{0.683022in}}%
\pgfpathlineto{\pgfqpoint{2.929113in}{0.666524in}}%
\pgfpathlineto{\pgfqpoint{2.930113in}{0.650025in}}%
\pgfpathlineto{\pgfqpoint{2.930157in}{0.658274in}}%
\pgfpathlineto{\pgfqpoint{2.931438in}{0.707769in}}%
\pgfpathlineto{\pgfqpoint{2.931623in}{0.691271in}}%
\pgfpathlineto{\pgfqpoint{2.931793in}{0.699520in}}%
\pgfpathlineto{\pgfqpoint{2.932778in}{0.683022in}}%
\pgfpathlineto{\pgfqpoint{2.933467in}{0.724268in}}%
\pgfpathlineto{\pgfqpoint{2.933800in}{0.724268in}}%
\pgfpathlineto{\pgfqpoint{2.934452in}{0.707769in}}%
\pgfpathlineto{\pgfqpoint{2.934452in}{0.740766in}}%
\pgfpathlineto{\pgfqpoint{2.935459in}{0.757264in}}%
\pgfpathlineto{\pgfqpoint{2.934807in}{0.724268in}}%
\pgfpathlineto{\pgfqpoint{2.935481in}{0.749015in}}%
\pgfpathlineto{\pgfqpoint{2.935799in}{0.757264in}}%
\pgfpathlineto{\pgfqpoint{2.936147in}{0.724268in}}%
\pgfpathlineto{\pgfqpoint{2.937302in}{0.757264in}}%
\pgfpathlineto{\pgfqpoint{2.936636in}{0.707769in}}%
\pgfpathlineto{\pgfqpoint{2.937495in}{0.740766in}}%
\pgfpathlineto{\pgfqpoint{2.938983in}{0.707769in}}%
\pgfpathlineto{\pgfqpoint{2.937658in}{0.749015in}}%
\pgfpathlineto{\pgfqpoint{2.938983in}{0.716019in}}%
\pgfpathlineto{\pgfqpoint{2.939316in}{0.707769in}}%
\pgfpathlineto{\pgfqpoint{2.940153in}{0.757264in}}%
\pgfpathlineto{\pgfqpoint{2.941012in}{0.724268in}}%
\pgfpathlineto{\pgfqpoint{2.941182in}{0.749015in}}%
\pgfpathlineto{\pgfqpoint{2.941493in}{0.757264in}}%
\pgfpathlineto{\pgfqpoint{2.941834in}{0.732517in}}%
\pgfpathlineto{\pgfqpoint{2.942019in}{0.740766in}}%
\pgfpathlineto{\pgfqpoint{2.942182in}{0.749015in}}%
\pgfpathlineto{\pgfqpoint{2.943174in}{0.732517in}}%
\pgfpathlineto{\pgfqpoint{2.944196in}{0.798510in}}%
\pgfpathlineto{\pgfqpoint{2.945876in}{0.724268in}}%
\pgfpathlineto{\pgfqpoint{2.946039in}{0.749015in}}%
\pgfpathlineto{\pgfqpoint{2.946358in}{0.757264in}}%
\pgfpathlineto{\pgfqpoint{2.946691in}{0.732517in}}%
\pgfpathlineto{\pgfqpoint{2.947216in}{0.749015in}}%
\pgfpathlineto{\pgfqpoint{2.947720in}{0.724268in}}%
\pgfpathlineto{\pgfqpoint{2.948371in}{0.732517in}}%
\pgfpathlineto{\pgfqpoint{2.948557in}{0.716019in}}%
\pgfpathlineto{\pgfqpoint{2.949060in}{0.749015in}}%
\pgfpathlineto{\pgfqpoint{2.949726in}{0.699520in}}%
\pgfpathlineto{\pgfqpoint{2.950733in}{0.773763in}}%
\pgfpathlineto{\pgfqpoint{2.951570in}{0.740766in}}%
\pgfpathlineto{\pgfqpoint{2.951740in}{0.749015in}}%
\pgfpathlineto{\pgfqpoint{2.952725in}{0.732517in}}%
\pgfpathlineto{\pgfqpoint{2.952747in}{0.749015in}}%
\pgfpathlineto{\pgfqpoint{2.953414in}{0.724268in}}%
\pgfpathlineto{\pgfqpoint{2.953754in}{0.724268in}}%
\pgfpathlineto{\pgfqpoint{2.955746in}{0.806759in}}%
\pgfpathlineto{\pgfqpoint{2.955931in}{0.790261in}}%
\pgfpathlineto{\pgfqpoint{2.956264in}{0.773763in}}%
\pgfpathlineto{\pgfqpoint{2.956583in}{0.815008in}}%
\pgfpathlineto{\pgfqpoint{2.956938in}{0.848005in}}%
\pgfpathlineto{\pgfqpoint{2.957590in}{0.806759in}}%
\pgfpathlineto{\pgfqpoint{2.957605in}{0.823258in}}%
\pgfpathlineto{\pgfqpoint{2.958945in}{0.773763in}}%
\pgfpathlineto{\pgfqpoint{2.959115in}{0.798510in}}%
\pgfpathlineto{\pgfqpoint{2.961107in}{0.856254in}}%
\pgfpathlineto{\pgfqpoint{2.961780in}{0.831507in}}%
\pgfpathlineto{\pgfqpoint{2.962136in}{0.872753in}}%
\pgfpathlineto{\pgfqpoint{2.963454in}{0.856254in}}%
\pgfpathlineto{\pgfqpoint{2.964483in}{0.897500in}}%
\pgfpathlineto{\pgfqpoint{2.964631in}{0.905749in}}%
\pgfpathlineto{\pgfqpoint{2.964987in}{0.872753in}}%
\pgfpathlineto{\pgfqpoint{2.965297in}{0.881002in}}%
\pgfpathlineto{\pgfqpoint{2.965986in}{0.848005in}}%
\pgfpathlineto{\pgfqpoint{2.966327in}{0.872753in}}%
\pgfpathlineto{\pgfqpoint{2.967497in}{0.897500in}}%
\pgfpathlineto{\pgfqpoint{2.966808in}{0.856254in}}%
\pgfpathlineto{\pgfqpoint{2.967667in}{0.889251in}}%
\pgfpathlineto{\pgfqpoint{2.967830in}{0.897500in}}%
\pgfpathlineto{\pgfqpoint{2.968822in}{0.881002in}}%
\pgfpathlineto{\pgfqpoint{2.968837in}{0.897500in}}%
\pgfpathlineto{\pgfqpoint{2.969177in}{0.872753in}}%
\pgfpathlineto{\pgfqpoint{2.969844in}{0.872753in}}%
\pgfpathlineto{\pgfqpoint{2.970347in}{0.922247in}}%
\pgfpathlineto{\pgfqpoint{2.971021in}{0.889251in}}%
\pgfpathlineto{\pgfqpoint{2.972191in}{0.872753in}}%
\pgfpathlineto{\pgfqpoint{2.973176in}{0.881002in}}%
\pgfpathlineto{\pgfqpoint{2.973198in}{0.872753in}}%
\pgfpathlineto{\pgfqpoint{2.974368in}{0.848005in}}%
\pgfpathlineto{\pgfqpoint{2.974516in}{0.856254in}}%
\pgfpathlineto{\pgfqpoint{2.975523in}{0.930497in}}%
\pgfpathlineto{\pgfqpoint{2.975545in}{0.922247in}}%
\pgfpathlineto{\pgfqpoint{2.975693in}{0.930497in}}%
\pgfpathlineto{\pgfqpoint{2.976048in}{0.897500in}}%
\pgfpathlineto{\pgfqpoint{2.977033in}{0.881002in}}%
\pgfpathlineto{\pgfqpoint{2.977055in}{0.897500in}}%
\pgfpathlineto{\pgfqpoint{2.977559in}{0.922247in}}%
\pgfpathlineto{\pgfqpoint{2.977892in}{0.897500in}}%
\pgfpathlineto{\pgfqpoint{2.979232in}{0.872753in}}%
\pgfpathlineto{\pgfqpoint{2.979269in}{0.881002in}}%
\pgfpathlineto{\pgfqpoint{2.981054in}{0.955244in}}%
\pgfpathlineto{\pgfqpoint{2.981076in}{0.946995in}}%
\pgfpathlineto{\pgfqpoint{2.981579in}{0.922247in}}%
\pgfpathlineto{\pgfqpoint{2.981898in}{0.955244in}}%
\pgfpathlineto{\pgfqpoint{2.981912in}{0.946995in}}%
\pgfpathlineto{\pgfqpoint{2.982083in}{0.971742in}}%
\pgfpathlineto{\pgfqpoint{2.982586in}{0.938746in}}%
\pgfpathlineto{\pgfqpoint{2.983756in}{0.922247in}}%
\pgfpathlineto{\pgfqpoint{2.984911in}{0.930497in}}%
\pgfpathlineto{\pgfqpoint{2.985096in}{0.922247in}}%
\pgfpathlineto{\pgfqpoint{2.985437in}{0.971742in}}%
\pgfpathlineto{\pgfqpoint{2.985770in}{0.971742in}}%
\pgfpathlineto{\pgfqpoint{2.986251in}{0.979992in}}%
\pgfpathlineto{\pgfqpoint{2.986436in}{0.963493in}}%
\pgfpathlineto{\pgfqpoint{2.986777in}{0.946995in}}%
\pgfpathlineto{\pgfqpoint{2.987088in}{0.979992in}}%
\pgfpathlineto{\pgfqpoint{2.987110in}{0.971742in}}%
\pgfpathlineto{\pgfqpoint{2.987258in}{0.979992in}}%
\pgfpathlineto{\pgfqpoint{2.987614in}{0.946995in}}%
\pgfpathlineto{\pgfqpoint{2.987932in}{0.930497in}}%
\pgfpathlineto{\pgfqpoint{2.988265in}{0.963493in}}%
\pgfpathlineto{\pgfqpoint{2.989124in}{0.996490in}}%
\pgfpathlineto{\pgfqpoint{2.988991in}{0.955244in}}%
\pgfpathlineto{\pgfqpoint{2.989287in}{0.988241in}}%
\pgfpathlineto{\pgfqpoint{2.990109in}{0.955244in}}%
\pgfpathlineto{\pgfqpoint{2.990131in}{0.996490in}}%
\pgfpathlineto{\pgfqpoint{2.990442in}{0.979992in}}%
\pgfpathlineto{\pgfqpoint{2.990464in}{0.996490in}}%
\pgfpathlineto{\pgfqpoint{2.991471in}{0.996490in}}%
\pgfpathlineto{\pgfqpoint{2.991975in}{1.045985in}}%
\pgfpathlineto{\pgfqpoint{2.993122in}{1.029487in}}%
\pgfpathlineto{\pgfqpoint{2.994152in}{1.021237in}}%
\pgfpathlineto{\pgfqpoint{2.995810in}{1.078981in}}%
\pgfpathlineto{\pgfqpoint{2.995995in}{1.062483in}}%
\pgfpathlineto{\pgfqpoint{2.996477in}{1.054234in}}%
\pgfpathlineto{\pgfqpoint{2.996499in}{1.095480in}}%
\pgfpathlineto{\pgfqpoint{2.996662in}{1.087231in}}%
\pgfpathlineto{\pgfqpoint{2.996832in}{1.095480in}}%
\pgfpathlineto{\pgfqpoint{2.998342in}{1.045985in}}%
\pgfpathlineto{\pgfqpoint{2.998994in}{1.078981in}}%
\pgfpathlineto{\pgfqpoint{2.999497in}{1.054234in}}%
\pgfpathlineto{\pgfqpoint{2.999512in}{1.045985in}}%
\pgfpathlineto{\pgfqpoint{2.999831in}{1.078981in}}%
\pgfpathlineto{\pgfqpoint{3.000519in}{1.062483in}}%
\pgfpathlineto{\pgfqpoint{3.000667in}{1.054234in}}%
\pgfpathlineto{\pgfqpoint{3.001008in}{1.087231in}}%
\pgfpathlineto{\pgfqpoint{3.001171in}{1.128476in}}%
\pgfpathlineto{\pgfqpoint{3.001859in}{1.070732in}}%
\pgfpathlineto{\pgfqpoint{3.002030in}{1.070732in}}%
\pgfpathlineto{\pgfqpoint{3.002348in}{1.078981in}}%
\pgfpathlineto{\pgfqpoint{3.002681in}{1.054234in}}%
\pgfpathlineto{\pgfqpoint{3.002866in}{1.062483in}}%
\pgfpathlineto{\pgfqpoint{3.003037in}{1.095480in}}%
\pgfpathlineto{\pgfqpoint{3.004044in}{1.045985in}}%
\pgfpathlineto{\pgfqpoint{3.005199in}{1.054234in}}%
\pgfpathlineto{\pgfqpoint{3.006220in}{1.045985in}}%
\pgfpathlineto{\pgfqpoint{3.007227in}{1.070732in}}%
\pgfpathlineto{\pgfqpoint{3.007390in}{1.062483in}}%
\pgfpathlineto{\pgfqpoint{3.007731in}{1.045985in}}%
\pgfpathlineto{\pgfqpoint{3.008042in}{1.078981in}}%
\pgfpathlineto{\pgfqpoint{3.008212in}{1.103729in}}%
\pgfpathlineto{\pgfqpoint{3.009071in}{1.045985in}}%
\pgfpathlineto{\pgfqpoint{3.010744in}{1.144975in}}%
\pgfpathlineto{\pgfqpoint{3.011248in}{1.111978in}}%
\pgfpathlineto{\pgfqpoint{3.011922in}{1.070732in}}%
\pgfpathlineto{\pgfqpoint{3.012736in}{1.087231in}}%
\pgfpathlineto{\pgfqpoint{3.013913in}{1.128476in}}%
\pgfpathlineto{\pgfqpoint{3.013077in}{1.078981in}}%
\pgfpathlineto{\pgfqpoint{3.013928in}{1.120227in}}%
\pgfpathlineto{\pgfqpoint{3.014099in}{1.144975in}}%
\pgfpathlineto{\pgfqpoint{3.015083in}{1.128476in}}%
\pgfpathlineto{\pgfqpoint{3.016112in}{1.120227in}}%
\pgfpathlineto{\pgfqpoint{3.017430in}{1.153224in}}%
\pgfpathlineto{\pgfqpoint{3.017453in}{1.144975in}}%
\pgfpathlineto{\pgfqpoint{3.017771in}{1.128476in}}%
\pgfpathlineto{\pgfqpoint{3.018104in}{1.161473in}}%
\pgfpathlineto{\pgfqpoint{3.018608in}{1.153224in}}%
\pgfpathlineto{\pgfqpoint{3.019274in}{1.202719in}}%
\pgfpathlineto{\pgfqpoint{3.020466in}{1.169722in}}%
\pgfpathlineto{\pgfqpoint{3.021806in}{1.219217in}}%
\pgfpathlineto{\pgfqpoint{3.022125in}{1.202719in}}%
\pgfpathlineto{\pgfqpoint{3.022798in}{1.177971in}}%
\pgfpathlineto{\pgfqpoint{3.023154in}{1.194470in}}%
\pgfpathlineto{\pgfqpoint{3.025146in}{1.103729in}}%
\pgfpathlineto{\pgfqpoint{3.025479in}{1.128476in}}%
\pgfpathlineto{\pgfqpoint{3.025664in}{1.169722in}}%
\pgfpathlineto{\pgfqpoint{3.026501in}{1.161473in}}%
\pgfpathlineto{\pgfqpoint{3.026671in}{1.169722in}}%
\pgfpathlineto{\pgfqpoint{3.027656in}{1.153224in}}%
\pgfpathlineto{\pgfqpoint{3.027678in}{1.169722in}}%
\pgfpathlineto{\pgfqpoint{3.028663in}{1.128476in}}%
\pgfpathlineto{\pgfqpoint{3.028685in}{1.144975in}}%
\pgfpathlineto{\pgfqpoint{3.028996in}{1.128476in}}%
\pgfpathlineto{\pgfqpoint{3.029018in}{1.169722in}}%
\pgfpathlineto{\pgfqpoint{3.029336in}{1.161473in}}%
\pgfpathlineto{\pgfqpoint{3.030003in}{1.128476in}}%
\pgfpathlineto{\pgfqpoint{3.030358in}{1.169722in}}%
\pgfpathlineto{\pgfqpoint{3.031180in}{1.177971in}}%
\pgfpathlineto{\pgfqpoint{3.031195in}{1.169722in}}%
\pgfpathlineto{\pgfqpoint{3.031698in}{1.144975in}}%
\pgfpathlineto{\pgfqpoint{3.032017in}{1.177971in}}%
\pgfpathlineto{\pgfqpoint{3.032039in}{1.169722in}}%
\pgfpathlineto{\pgfqpoint{3.032535in}{1.194470in}}%
\pgfpathlineto{\pgfqpoint{3.032876in}{1.169722in}}%
\pgfpathlineto{\pgfqpoint{3.035223in}{1.095480in}}%
\pgfpathlineto{\pgfqpoint{3.035874in}{1.103729in}}%
\pgfpathlineto{\pgfqpoint{3.036059in}{1.087231in}}%
\pgfpathlineto{\pgfqpoint{3.036207in}{1.078981in}}%
\pgfpathlineto{\pgfqpoint{3.036563in}{1.120227in}}%
\pgfpathlineto{\pgfqpoint{3.036726in}{1.111978in}}%
\pgfpathlineto{\pgfqpoint{3.037044in}{1.128476in}}%
\pgfpathlineto{\pgfqpoint{3.037903in}{1.095480in}}%
\pgfpathlineto{\pgfqpoint{3.038725in}{1.103729in}}%
\pgfpathlineto{\pgfqpoint{3.038740in}{1.095480in}}%
\pgfpathlineto{\pgfqpoint{3.040420in}{1.045985in}}%
\pgfpathlineto{\pgfqpoint{3.040569in}{1.054234in}}%
\pgfpathlineto{\pgfqpoint{3.040902in}{1.078981in}}%
\pgfpathlineto{\pgfqpoint{3.041590in}{1.070732in}}%
\pgfpathlineto{\pgfqpoint{3.043249in}{1.029487in}}%
\pgfpathlineto{\pgfqpoint{3.041923in}{1.095480in}}%
\pgfpathlineto{\pgfqpoint{3.043249in}{1.037736in}}%
\pgfpathlineto{\pgfqpoint{3.043767in}{1.070732in}}%
\pgfpathlineto{\pgfqpoint{3.044271in}{1.012988in}}%
\pgfpathlineto{\pgfqpoint{3.045611in}{0.971742in}}%
\pgfpathlineto{\pgfqpoint{3.044922in}{1.029487in}}%
\pgfpathlineto{\pgfqpoint{3.045759in}{0.988241in}}%
\pgfpathlineto{\pgfqpoint{3.045929in}{1.004739in}}%
\pgfpathlineto{\pgfqpoint{3.046262in}{0.979992in}}%
\pgfpathlineto{\pgfqpoint{3.046788in}{0.988241in}}%
\pgfpathlineto{\pgfqpoint{3.047795in}{0.922247in}}%
\pgfpathlineto{\pgfqpoint{3.046951in}{0.996490in}}%
\pgfpathlineto{\pgfqpoint{3.048284in}{0.963493in}}%
\pgfpathlineto{\pgfqpoint{3.048950in}{1.004739in}}%
\pgfpathlineto{\pgfqpoint{3.049298in}{0.971742in}}%
\pgfpathlineto{\pgfqpoint{3.050290in}{0.979992in}}%
\pgfpathlineto{\pgfqpoint{3.050305in}{0.971742in}}%
\pgfpathlineto{\pgfqpoint{3.051986in}{0.922247in}}%
\pgfpathlineto{\pgfqpoint{3.052134in}{0.930497in}}%
\pgfpathlineto{\pgfqpoint{3.052993in}{0.996490in}}%
\pgfpathlineto{\pgfqpoint{3.053156in}{0.988241in}}%
\pgfpathlineto{\pgfqpoint{3.053644in}{0.955244in}}%
\pgfpathlineto{\pgfqpoint{3.054140in}{0.988241in}}%
\pgfpathlineto{\pgfqpoint{3.054163in}{0.996490in}}%
\pgfpathlineto{\pgfqpoint{3.054496in}{0.971742in}}%
\pgfpathlineto{\pgfqpoint{3.055170in}{0.971742in}}%
\pgfpathlineto{\pgfqpoint{3.055673in}{0.996490in}}%
\pgfpathlineto{\pgfqpoint{3.056006in}{0.971742in}}%
\pgfpathlineto{\pgfqpoint{3.056325in}{0.955244in}}%
\pgfpathlineto{\pgfqpoint{3.056339in}{0.996490in}}%
\pgfpathlineto{\pgfqpoint{3.056658in}{1.004739in}}%
\pgfpathlineto{\pgfqpoint{3.056991in}{0.979992in}}%
\pgfpathlineto{\pgfqpoint{3.057184in}{0.988241in}}%
\pgfpathlineto{\pgfqpoint{3.058672in}{0.955244in}}%
\pgfpathlineto{\pgfqpoint{3.058672in}{0.963493in}}%
\pgfpathlineto{\pgfqpoint{3.059027in}{0.971742in}}%
\pgfpathlineto{\pgfqpoint{3.059523in}{0.922247in}}%
\pgfpathlineto{\pgfqpoint{3.059694in}{0.922247in}}%
\pgfpathlineto{\pgfqpoint{3.060012in}{0.930497in}}%
\pgfpathlineto{\pgfqpoint{3.060345in}{0.905749in}}%
\pgfpathlineto{\pgfqpoint{3.061537in}{0.839756in}}%
\pgfpathlineto{\pgfqpoint{3.062692in}{0.806759in}}%
\pgfpathlineto{\pgfqpoint{3.061707in}{0.848005in}}%
\pgfpathlineto{\pgfqpoint{3.062692in}{0.815008in}}%
\pgfpathlineto{\pgfqpoint{3.063048in}{0.848005in}}%
\pgfpathlineto{\pgfqpoint{3.063714in}{0.823258in}}%
\pgfpathlineto{\pgfqpoint{3.063877in}{0.831507in}}%
\pgfpathlineto{\pgfqpoint{3.064218in}{0.798510in}}%
\pgfpathlineto{\pgfqpoint{3.064366in}{0.806759in}}%
\pgfpathlineto{\pgfqpoint{3.065565in}{0.740766in}}%
\pgfpathlineto{\pgfqpoint{3.067053in}{0.707769in}}%
\pgfpathlineto{\pgfqpoint{3.065898in}{0.749015in}}%
\pgfpathlineto{\pgfqpoint{3.067061in}{0.716019in}}%
\pgfpathlineto{\pgfqpoint{3.067742in}{0.749015in}}%
\pgfpathlineto{\pgfqpoint{3.067387in}{0.707769in}}%
\pgfpathlineto{\pgfqpoint{3.068075in}{0.724268in}}%
\pgfpathlineto{\pgfqpoint{3.068579in}{0.707769in}}%
\pgfpathlineto{\pgfqpoint{3.068904in}{0.732517in}}%
\pgfpathlineto{\pgfqpoint{3.068919in}{0.732517in}}%
\pgfpathlineto{\pgfqpoint{3.070082in}{0.806759in}}%
\pgfpathlineto{\pgfqpoint{3.070422in}{0.765513in}}%
\pgfpathlineto{\pgfqpoint{3.073103in}{0.650025in}}%
\pgfpathlineto{\pgfqpoint{3.073265in}{0.666524in}}%
\pgfpathlineto{\pgfqpoint{3.074110in}{0.699520in}}%
\pgfpathlineto{\pgfqpoint{3.073947in}{0.658274in}}%
\pgfpathlineto{\pgfqpoint{3.074280in}{0.691271in}}%
\pgfpathlineto{\pgfqpoint{3.074450in}{0.658274in}}%
\pgfpathlineto{\pgfqpoint{3.075116in}{0.699520in}}%
\pgfpathlineto{\pgfqpoint{3.075450in}{0.674773in}}%
\pgfpathlineto{\pgfqpoint{3.075598in}{0.683022in}}%
\pgfpathlineto{\pgfqpoint{3.075953in}{0.650025in}}%
\pgfpathlineto{\pgfqpoint{3.076286in}{0.666524in}}%
\pgfpathlineto{\pgfqpoint{3.076960in}{0.650025in}}%
\pgfpathlineto{\pgfqpoint{3.077130in}{0.674773in}}%
\pgfpathlineto{\pgfqpoint{3.077449in}{0.658274in}}%
\pgfpathlineto{\pgfqpoint{3.077464in}{0.674773in}}%
\pgfpathlineto{\pgfqpoint{3.078137in}{0.633527in}}%
\pgfpathlineto{\pgfqpoint{3.078471in}{0.658274in}}%
\pgfpathlineto{\pgfqpoint{3.080818in}{0.749015in}}%
\pgfpathlineto{\pgfqpoint{3.080981in}{0.740766in}}%
\pgfpathlineto{\pgfqpoint{3.081136in}{0.732517in}}%
\pgfpathlineto{\pgfqpoint{3.081306in}{0.757264in}}%
\pgfpathlineto{\pgfqpoint{3.081654in}{0.757264in}}%
\pgfpathlineto{\pgfqpoint{3.084994in}{0.881002in}}%
\pgfpathlineto{\pgfqpoint{3.085009in}{0.872753in}}%
\pgfpathlineto{\pgfqpoint{3.087526in}{0.782012in}}%
\pgfpathlineto{\pgfqpoint{3.087681in}{0.790261in}}%
\pgfpathlineto{\pgfqpoint{3.088022in}{0.773763in}}%
\pgfpathlineto{\pgfqpoint{3.088696in}{0.798510in}}%
\pgfpathlineto{\pgfqpoint{3.089518in}{0.806759in}}%
\pgfpathlineto{\pgfqpoint{3.090036in}{0.773763in}}%
\pgfpathlineto{\pgfqpoint{3.090191in}{0.782012in}}%
\pgfpathlineto{\pgfqpoint{3.090539in}{0.749015in}}%
\pgfpathlineto{\pgfqpoint{3.090539in}{0.757264in}}%
\pgfpathlineto{\pgfqpoint{3.092724in}{0.633527in}}%
\pgfpathlineto{\pgfqpoint{3.093220in}{0.674773in}}%
\pgfpathlineto{\pgfqpoint{3.093560in}{0.658274in}}%
\pgfpathlineto{\pgfqpoint{3.095397in}{0.732517in}}%
\pgfpathlineto{\pgfqpoint{3.095397in}{0.724268in}}%
\pgfpathlineto{\pgfqpoint{3.095567in}{0.749015in}}%
\pgfpathlineto{\pgfqpoint{3.096411in}{0.732517in}}%
\pgfpathlineto{\pgfqpoint{3.097248in}{0.707769in}}%
\pgfpathlineto{\pgfqpoint{3.097396in}{0.716019in}}%
\pgfpathlineto{\pgfqpoint{3.097411in}{0.749015in}}%
\pgfpathlineto{\pgfqpoint{3.098417in}{0.699520in}}%
\pgfpathlineto{\pgfqpoint{3.098573in}{0.707769in}}%
\pgfpathlineto{\pgfqpoint{3.098758in}{0.683022in}}%
\pgfpathlineto{\pgfqpoint{3.099254in}{0.699520in}}%
\pgfpathlineto{\pgfqpoint{3.101594in}{0.625278in}}%
\pgfpathlineto{\pgfqpoint{3.102764in}{0.633527in}}%
\pgfpathlineto{\pgfqpoint{3.102771in}{0.625278in}}%
\pgfpathlineto{\pgfqpoint{3.103778in}{0.625278in}}%
\pgfpathlineto{\pgfqpoint{3.104785in}{0.633527in}}%
\pgfpathlineto{\pgfqpoint{3.104785in}{0.625278in}}%
\pgfpathlineto{\pgfqpoint{3.106118in}{0.633527in}}%
\pgfpathlineto{\pgfqpoint{3.106118in}{0.625278in}}%
\pgfpathlineto{\pgfqpoint{3.106955in}{0.625278in}}%
\pgfpathlineto{\pgfqpoint{3.108295in}{0.633527in}}%
\pgfpathlineto{\pgfqpoint{3.108295in}{0.625278in}}%
\pgfpathlineto{\pgfqpoint{3.109131in}{0.625278in}}%
\pgfpathlineto{\pgfqpoint{3.110309in}{0.633527in}}%
\pgfpathlineto{\pgfqpoint{3.110309in}{0.625278in}}%
\pgfpathlineto{\pgfqpoint{3.111316in}{0.625278in}}%
\pgfpathlineto{\pgfqpoint{3.112478in}{0.633527in}}%
\pgfpathlineto{\pgfqpoint{3.112819in}{0.625278in}}%
\pgfpathlineto{\pgfqpoint{3.113485in}{0.625278in}}%
\pgfpathlineto{\pgfqpoint{3.114655in}{0.633527in}}%
\pgfpathlineto{\pgfqpoint{3.114655in}{0.625278in}}%
\pgfpathlineto{\pgfqpoint{3.115499in}{0.625278in}}%
\pgfpathlineto{\pgfqpoint{3.116669in}{0.633527in}}%
\pgfpathlineto{\pgfqpoint{3.116839in}{0.625278in}}%
\pgfpathlineto{\pgfqpoint{3.117506in}{0.625278in}}%
\pgfpathlineto{\pgfqpoint{3.118683in}{0.633527in}}%
\pgfpathlineto{\pgfqpoint{3.118683in}{0.625278in}}%
\pgfpathlineto{\pgfqpoint{3.119690in}{0.625278in}}%
\pgfpathlineto{\pgfqpoint{3.120534in}{0.633527in}}%
\pgfpathlineto{\pgfqpoint{3.120534in}{0.625278in}}%
\pgfpathlineto{\pgfqpoint{3.121867in}{0.633527in}}%
\pgfpathlineto{\pgfqpoint{3.121874in}{0.625278in}}%
\pgfpathlineto{\pgfqpoint{3.122711in}{0.625278in}}%
\pgfpathlineto{\pgfqpoint{3.123718in}{0.633527in}}%
\pgfpathlineto{\pgfqpoint{3.123718in}{0.625278in}}%
\pgfpathlineto{\pgfqpoint{3.124887in}{0.633527in}}%
\pgfpathlineto{\pgfqpoint{3.125050in}{0.625278in}}%
\pgfpathlineto{\pgfqpoint{3.125887in}{0.625278in}}%
\pgfpathlineto{\pgfqpoint{3.127227in}{0.633527in}}%
\pgfpathlineto{\pgfqpoint{3.127227in}{0.625278in}}%
\pgfpathlineto{\pgfqpoint{3.128242in}{0.625278in}}%
\pgfpathlineto{\pgfqpoint{3.129411in}{0.633527in}}%
\pgfpathlineto{\pgfqpoint{3.129574in}{0.625278in}}%
\pgfpathlineto{\pgfqpoint{3.130418in}{0.625278in}}%
\pgfpathlineto{\pgfqpoint{3.131588in}{0.633527in}}%
\pgfpathlineto{\pgfqpoint{3.131759in}{0.625278in}}%
\pgfpathlineto{\pgfqpoint{3.132595in}{0.625278in}}%
\pgfpathlineto{\pgfqpoint{3.133765in}{0.633527in}}%
\pgfpathlineto{\pgfqpoint{3.133765in}{0.625278in}}%
\pgfpathlineto{\pgfqpoint{3.134772in}{0.625278in}}%
\pgfpathlineto{\pgfqpoint{3.135949in}{0.633527in}}%
\pgfpathlineto{\pgfqpoint{3.136112in}{0.625278in}}%
\pgfpathlineto{\pgfqpoint{3.136956in}{0.625278in}}%
\pgfpathlineto{\pgfqpoint{3.138126in}{0.633527in}}%
\pgfpathlineto{\pgfqpoint{3.138134in}{0.625278in}}%
\pgfpathlineto{\pgfqpoint{3.139141in}{0.625278in}}%
\pgfpathlineto{\pgfqpoint{3.140303in}{0.633527in}}%
\pgfpathlineto{\pgfqpoint{3.140481in}{0.625278in}}%
\pgfpathlineto{\pgfqpoint{3.141310in}{0.625278in}}%
\pgfpathlineto{\pgfqpoint{3.142487in}{0.633527in}}%
\pgfpathlineto{\pgfqpoint{3.142650in}{0.625278in}}%
\pgfpathlineto{\pgfqpoint{3.143494in}{0.625278in}}%
\pgfpathlineto{\pgfqpoint{3.144494in}{0.633527in}}%
\pgfpathlineto{\pgfqpoint{3.144494in}{0.625278in}}%
\pgfpathlineto{\pgfqpoint{3.145841in}{0.633527in}}%
\pgfpathlineto{\pgfqpoint{3.146175in}{0.625278in}}%
\pgfpathlineto{\pgfqpoint{3.146841in}{0.625278in}}%
\pgfpathlineto{\pgfqpoint{3.148018in}{0.633527in}}%
\pgfpathlineto{\pgfqpoint{3.148181in}{0.625278in}}%
\pgfpathlineto{\pgfqpoint{3.149025in}{0.625278in}}%
\pgfpathlineto{\pgfqpoint{3.150195in}{0.633527in}}%
\pgfpathlineto{\pgfqpoint{3.150365in}{0.625278in}}%
\pgfpathlineto{\pgfqpoint{3.150869in}{0.625278in}}%
\pgfpathlineto{\pgfqpoint{3.152039in}{0.633527in}}%
\pgfpathlineto{\pgfqpoint{3.152209in}{0.625278in}}%
\pgfpathlineto{\pgfqpoint{3.153046in}{0.625278in}}%
\pgfpathlineto{\pgfqpoint{3.154386in}{0.633527in}}%
\pgfpathlineto{\pgfqpoint{3.154386in}{0.625278in}}%
\pgfpathlineto{\pgfqpoint{3.155400in}{0.625278in}}%
\pgfpathlineto{\pgfqpoint{3.156563in}{0.633527in}}%
\pgfpathlineto{\pgfqpoint{3.156570in}{0.625278in}}%
\pgfpathlineto{\pgfqpoint{3.157407in}{0.625278in}}%
\pgfpathlineto{\pgfqpoint{3.158577in}{0.633527in}}%
\pgfpathlineto{\pgfqpoint{3.158577in}{0.625278in}}%
\pgfpathlineto{\pgfqpoint{3.159584in}{0.625278in}}%
\pgfpathlineto{\pgfqpoint{3.160753in}{0.633527in}}%
\pgfpathlineto{\pgfqpoint{3.160924in}{0.625278in}}%
\pgfpathlineto{\pgfqpoint{3.161760in}{0.625278in}}%
\pgfpathlineto{\pgfqpoint{3.162938in}{0.633527in}}%
\pgfpathlineto{\pgfqpoint{3.163271in}{0.625278in}}%
\pgfpathlineto{\pgfqpoint{3.163945in}{0.625278in}}%
\pgfpathlineto{\pgfqpoint{3.165114in}{0.633527in}}%
\pgfpathlineto{\pgfqpoint{3.165285in}{0.625278in}}%
\pgfpathlineto{\pgfqpoint{3.166121in}{0.625278in}}%
\pgfpathlineto{\pgfqpoint{3.167291in}{0.633527in}}%
\pgfpathlineto{\pgfqpoint{3.167462in}{0.625278in}}%
\pgfpathlineto{\pgfqpoint{3.168298in}{0.625278in}}%
\pgfpathlineto{\pgfqpoint{3.169305in}{0.633527in}}%
\pgfpathlineto{\pgfqpoint{3.169305in}{0.625278in}}%
\pgfpathlineto{\pgfqpoint{3.170483in}{0.633527in}}%
\pgfpathlineto{\pgfqpoint{3.170816in}{0.625278in}}%
\pgfpathlineto{\pgfqpoint{3.171482in}{0.625278in}}%
\pgfpathlineto{\pgfqpoint{3.172659in}{0.633527in}}%
\pgfpathlineto{\pgfqpoint{3.172830in}{0.625278in}}%
\pgfpathlineto{\pgfqpoint{3.173666in}{0.625278in}}%
\pgfpathlineto{\pgfqpoint{3.174673in}{0.633527in}}%
\pgfpathlineto{\pgfqpoint{3.174673in}{0.625278in}}%
\pgfpathlineto{\pgfqpoint{3.176013in}{0.633527in}}%
\pgfpathlineto{\pgfqpoint{3.176013in}{0.625278in}}%
\pgfpathlineto{\pgfqpoint{3.177020in}{0.625278in}}%
\pgfpathlineto{\pgfqpoint{3.178020in}{0.633527in}}%
\pgfpathlineto{\pgfqpoint{3.178020in}{0.625278in}}%
\pgfpathlineto{\pgfqpoint{3.179197in}{0.633527in}}%
\pgfpathlineto{\pgfqpoint{3.179368in}{0.625278in}}%
\pgfpathlineto{\pgfqpoint{3.180204in}{0.625278in}}%
\pgfpathlineto{\pgfqpoint{3.181374in}{0.633527in}}%
\pgfpathlineto{\pgfqpoint{3.181374in}{0.625278in}}%
\pgfpathlineto{\pgfqpoint{3.182211in}{0.625278in}}%
\pgfpathlineto{\pgfqpoint{3.183388in}{0.633527in}}%
\pgfpathlineto{\pgfqpoint{3.183558in}{0.625278in}}%
\pgfpathlineto{\pgfqpoint{3.184395in}{0.625278in}}%
\pgfpathlineto{\pgfqpoint{3.185565in}{0.633527in}}%
\pgfpathlineto{\pgfqpoint{3.185735in}{0.625278in}}%
\pgfpathlineto{\pgfqpoint{3.186572in}{0.625278in}}%
\pgfpathlineto{\pgfqpoint{3.187749in}{0.633527in}}%
\pgfpathlineto{\pgfqpoint{3.187749in}{0.625278in}}%
\pgfpathlineto{\pgfqpoint{3.188749in}{0.625278in}}%
\pgfpathlineto{\pgfqpoint{3.190089in}{0.633527in}}%
\pgfpathlineto{\pgfqpoint{3.190089in}{0.625278in}}%
\pgfpathlineto{\pgfqpoint{3.191096in}{0.625278in}}%
\pgfpathlineto{\pgfqpoint{3.192103in}{0.633527in}}%
\pgfpathlineto{\pgfqpoint{3.192103in}{0.625278in}}%
\pgfpathlineto{\pgfqpoint{3.193280in}{0.633527in}}%
\pgfpathlineto{\pgfqpoint{3.193443in}{0.625278in}}%
\pgfpathlineto{\pgfqpoint{3.194280in}{0.625278in}}%
\pgfpathlineto{\pgfqpoint{3.195457in}{0.633527in}}%
\pgfpathlineto{\pgfqpoint{3.195627in}{0.625278in}}%
\pgfpathlineto{\pgfqpoint{3.196464in}{0.625278in}}%
\pgfpathlineto{\pgfqpoint{3.197634in}{0.633527in}}%
\pgfpathlineto{\pgfqpoint{3.197804in}{0.625278in}}%
\pgfpathlineto{\pgfqpoint{3.198641in}{0.625278in}}%
\pgfpathlineto{\pgfqpoint{3.199818in}{0.633527in}}%
\pgfpathlineto{\pgfqpoint{3.200151in}{0.625278in}}%
\pgfpathlineto{\pgfqpoint{3.200818in}{0.625278in}}%
\pgfpathlineto{\pgfqpoint{3.201995in}{0.633527in}}%
\pgfpathlineto{\pgfqpoint{3.202165in}{0.625278in}}%
\pgfpathlineto{\pgfqpoint{3.202831in}{0.625278in}}%
\pgfpathlineto{\pgfqpoint{3.204009in}{0.633527in}}%
\pgfpathlineto{\pgfqpoint{3.204172in}{0.625278in}}%
\pgfpathlineto{\pgfqpoint{3.205008in}{0.625278in}}%
\pgfpathlineto{\pgfqpoint{3.206186in}{0.633527in}}%
\pgfpathlineto{\pgfqpoint{3.206186in}{0.625278in}}%
\pgfpathlineto{\pgfqpoint{3.207022in}{0.625278in}}%
\pgfpathlineto{\pgfqpoint{3.208199in}{0.633527in}}%
\pgfpathlineto{\pgfqpoint{3.208362in}{0.625278in}}%
\pgfpathlineto{\pgfqpoint{3.209199in}{0.625278in}}%
\pgfpathlineto{\pgfqpoint{3.210376in}{0.633527in}}%
\pgfpathlineto{\pgfqpoint{3.210547in}{0.625278in}}%
\pgfpathlineto{\pgfqpoint{3.211383in}{0.625278in}}%
\pgfpathlineto{\pgfqpoint{3.212390in}{0.633527in}}%
\pgfpathlineto{\pgfqpoint{3.212390in}{0.625278in}}%
\pgfpathlineto{\pgfqpoint{3.213619in}{0.633527in}}%
\pgfpathlineto{\pgfqpoint{3.213730in}{0.625278in}}%
\pgfpathlineto{\pgfqpoint{3.214397in}{0.625278in}}%
\pgfpathlineto{\pgfqpoint{3.215574in}{0.633527in}}%
\pgfpathlineto{\pgfqpoint{3.215737in}{0.625278in}}%
\pgfpathlineto{\pgfqpoint{3.216581in}{0.625278in}}%
\pgfpathlineto{\pgfqpoint{3.217751in}{0.633527in}}%
\pgfpathlineto{\pgfqpoint{3.217921in}{0.625278in}}%
\pgfpathlineto{\pgfqpoint{3.218758in}{0.625278in}}%
\pgfpathlineto{\pgfqpoint{3.219928in}{0.633527in}}%
\pgfpathlineto{\pgfqpoint{3.220098in}{0.625278in}}%
\pgfpathlineto{\pgfqpoint{3.220935in}{0.625278in}}%
\pgfpathlineto{\pgfqpoint{3.222112in}{0.633527in}}%
\pgfpathlineto{\pgfqpoint{3.222275in}{0.625278in}}%
\pgfpathlineto{\pgfqpoint{3.223119in}{0.625278in}}%
\pgfpathlineto{\pgfqpoint{3.224289in}{0.633527in}}%
\pgfpathlineto{\pgfqpoint{3.224459in}{0.625278in}}%
\pgfpathlineto{\pgfqpoint{3.225125in}{0.625278in}}%
\pgfpathlineto{\pgfqpoint{3.226303in}{0.633527in}}%
\pgfpathlineto{\pgfqpoint{3.226303in}{0.625278in}}%
\pgfpathlineto{\pgfqpoint{3.227310in}{0.625278in}}%
\pgfpathlineto{\pgfqpoint{3.228480in}{0.633527in}}%
\pgfpathlineto{\pgfqpoint{3.228650in}{0.625278in}}%
\pgfpathlineto{\pgfqpoint{3.229487in}{0.625278in}}%
\pgfpathlineto{\pgfqpoint{3.230656in}{0.633527in}}%
\pgfpathlineto{\pgfqpoint{3.230827in}{0.625278in}}%
\pgfpathlineto{\pgfqpoint{3.231663in}{0.625278in}}%
\pgfpathlineto{\pgfqpoint{3.232841in}{0.633527in}}%
\pgfpathlineto{\pgfqpoint{3.233004in}{0.625278in}}%
\pgfpathlineto{\pgfqpoint{3.233848in}{0.625278in}}%
\pgfpathlineto{\pgfqpoint{3.235017in}{0.633527in}}%
\pgfpathlineto{\pgfqpoint{3.235188in}{0.625278in}}%
\pgfpathlineto{\pgfqpoint{3.236024in}{0.625278in}}%
\pgfpathlineto{\pgfqpoint{3.237031in}{0.633527in}}%
\pgfpathlineto{\pgfqpoint{3.237031in}{0.625278in}}%
\pgfpathlineto{\pgfqpoint{3.238201in}{0.633527in}}%
\pgfpathlineto{\pgfqpoint{3.238534in}{0.625278in}}%
\pgfpathlineto{\pgfqpoint{3.239038in}{0.625278in}}%
\pgfpathlineto{\pgfqpoint{3.240215in}{0.633527in}}%
\pgfpathlineto{\pgfqpoint{3.240215in}{0.625278in}}%
\pgfpathlineto{\pgfqpoint{3.241222in}{0.625278in}}%
\pgfpathlineto{\pgfqpoint{3.242392in}{0.633527in}}%
\pgfpathlineto{\pgfqpoint{3.242562in}{0.625278in}}%
\pgfpathlineto{\pgfqpoint{3.243399in}{0.625278in}}%
\pgfpathlineto{\pgfqpoint{3.244569in}{0.633527in}}%
\pgfpathlineto{\pgfqpoint{3.244739in}{0.625278in}}%
\pgfpathlineto{\pgfqpoint{3.245576in}{0.625278in}}%
\pgfpathlineto{\pgfqpoint{3.246753in}{0.633527in}}%
\pgfpathlineto{\pgfqpoint{3.246753in}{0.625278in}}%
\pgfpathlineto{\pgfqpoint{3.247760in}{0.625278in}}%
\pgfpathlineto{\pgfqpoint{3.248760in}{0.633527in}}%
\pgfpathlineto{\pgfqpoint{3.248760in}{0.625278in}}%
\pgfpathlineto{\pgfqpoint{3.249937in}{0.633527in}}%
\pgfpathlineto{\pgfqpoint{3.250107in}{0.625278in}}%
\pgfpathlineto{\pgfqpoint{3.250959in}{0.625278in}}%
\pgfpathlineto{\pgfqpoint{3.254128in}{0.782012in}}%
\pgfpathlineto{\pgfqpoint{3.254313in}{0.765513in}}%
\pgfpathlineto{\pgfqpoint{3.254801in}{0.732517in}}%
\pgfpathlineto{\pgfqpoint{3.255149in}{0.773763in}}%
\pgfpathlineto{\pgfqpoint{3.255801in}{0.757264in}}%
\pgfpathlineto{\pgfqpoint{3.256305in}{0.782012in}}%
\pgfpathlineto{\pgfqpoint{3.257497in}{0.749015in}}%
\pgfpathlineto{\pgfqpoint{3.258674in}{0.798510in}}%
\pgfpathlineto{\pgfqpoint{3.258148in}{0.732517in}}%
\pgfpathlineto{\pgfqpoint{3.259007in}{0.773763in}}%
\pgfpathlineto{\pgfqpoint{3.259681in}{0.749015in}}%
\pgfpathlineto{\pgfqpoint{3.259844in}{0.773763in}}%
\pgfpathlineto{\pgfqpoint{3.260495in}{0.782012in}}%
\pgfpathlineto{\pgfqpoint{3.260680in}{0.765513in}}%
\pgfpathlineto{\pgfqpoint{3.262176in}{0.732517in}}%
\pgfpathlineto{\pgfqpoint{3.262176in}{0.740766in}}%
\pgfpathlineto{\pgfqpoint{3.262531in}{0.724268in}}%
\pgfpathlineto{\pgfqpoint{3.263198in}{0.749015in}}%
\pgfpathlineto{\pgfqpoint{3.264856in}{0.782012in}}%
\pgfpathlineto{\pgfqpoint{3.264871in}{0.773763in}}%
\pgfpathlineto{\pgfqpoint{3.265042in}{0.798510in}}%
\pgfpathlineto{\pgfqpoint{3.265878in}{0.790261in}}%
\pgfpathlineto{\pgfqpoint{3.266026in}{0.782012in}}%
\pgfpathlineto{\pgfqpoint{3.266197in}{0.806759in}}%
\pgfpathlineto{\pgfqpoint{3.266552in}{0.798510in}}%
\pgfpathlineto{\pgfqpoint{3.267707in}{0.806759in}}%
\pgfpathlineto{\pgfqpoint{3.268544in}{0.782012in}}%
\pgfpathlineto{\pgfqpoint{3.268729in}{0.790261in}}%
\pgfpathlineto{\pgfqpoint{3.268899in}{0.798510in}}%
\pgfpathlineto{\pgfqpoint{3.269906in}{0.773763in}}%
\pgfpathlineto{\pgfqpoint{3.270410in}{0.749015in}}%
\pgfpathlineto{\pgfqpoint{3.271061in}{0.782012in}}%
\pgfpathlineto{\pgfqpoint{3.271579in}{0.749015in}}%
\pgfpathlineto{\pgfqpoint{3.272083in}{0.773763in}}%
\pgfpathlineto{\pgfqpoint{3.272586in}{0.749015in}}%
\pgfpathlineto{\pgfqpoint{3.272905in}{0.782012in}}%
\pgfpathlineto{\pgfqpoint{3.272920in}{0.773763in}}%
\pgfpathlineto{\pgfqpoint{3.273068in}{0.782012in}}%
\pgfpathlineto{\pgfqpoint{3.273423in}{0.749015in}}%
\pgfpathlineto{\pgfqpoint{3.273571in}{0.757264in}}%
\pgfpathlineto{\pgfqpoint{3.273593in}{0.749015in}}%
\pgfpathlineto{\pgfqpoint{3.274075in}{0.806759in}}%
\pgfpathlineto{\pgfqpoint{3.274430in}{0.798510in}}%
\pgfpathlineto{\pgfqpoint{3.274748in}{0.782012in}}%
\pgfpathlineto{\pgfqpoint{3.275082in}{0.815008in}}%
\pgfpathlineto{\pgfqpoint{3.275585in}{0.831507in}}%
\pgfpathlineto{\pgfqpoint{3.275940in}{0.798510in}}%
\pgfpathlineto{\pgfqpoint{3.276103in}{0.798510in}}%
\pgfpathlineto{\pgfqpoint{3.277258in}{0.806759in}}%
\pgfpathlineto{\pgfqpoint{3.277947in}{0.798510in}}%
\pgfpathlineto{\pgfqpoint{3.278265in}{0.831507in}}%
\pgfpathlineto{\pgfqpoint{3.278288in}{0.823258in}}%
\pgfpathlineto{\pgfqpoint{3.278939in}{0.856254in}}%
\pgfpathlineto{\pgfqpoint{3.279443in}{0.831507in}}%
\pgfpathlineto{\pgfqpoint{3.279791in}{0.823258in}}%
\pgfpathlineto{\pgfqpoint{3.280109in}{0.856254in}}%
\pgfpathlineto{\pgfqpoint{3.280464in}{0.848005in}}%
\pgfpathlineto{\pgfqpoint{3.280613in}{0.856254in}}%
\pgfpathlineto{\pgfqpoint{3.280946in}{0.831507in}}%
\pgfpathlineto{\pgfqpoint{3.281138in}{0.839756in}}%
\pgfpathlineto{\pgfqpoint{3.282812in}{0.798510in}}%
\pgfpathlineto{\pgfqpoint{3.283796in}{0.856254in}}%
\pgfpathlineto{\pgfqpoint{3.284152in}{0.815008in}}%
\pgfpathlineto{\pgfqpoint{3.284485in}{0.798510in}}%
\pgfpathlineto{\pgfqpoint{3.284803in}{0.831507in}}%
\pgfpathlineto{\pgfqpoint{3.284826in}{0.823258in}}%
\pgfpathlineto{\pgfqpoint{3.285640in}{0.831507in}}%
\pgfpathlineto{\pgfqpoint{3.285662in}{0.823258in}}%
\pgfpathlineto{\pgfqpoint{3.286314in}{0.782012in}}%
\pgfpathlineto{\pgfqpoint{3.286669in}{0.823258in}}%
\pgfpathlineto{\pgfqpoint{3.288157in}{0.856254in}}%
\pgfpathlineto{\pgfqpoint{3.288172in}{0.848005in}}%
\pgfpathlineto{\pgfqpoint{3.289498in}{0.831507in}}%
\pgfpathlineto{\pgfqpoint{3.289520in}{0.848005in}}%
\pgfpathlineto{\pgfqpoint{3.289853in}{0.823258in}}%
\pgfpathlineto{\pgfqpoint{3.290519in}{0.823258in}}%
\pgfpathlineto{\pgfqpoint{3.291023in}{0.798510in}}%
\pgfpathlineto{\pgfqpoint{3.291341in}{0.831507in}}%
\pgfpathlineto{\pgfqpoint{3.291363in}{0.823258in}}%
\pgfpathlineto{\pgfqpoint{3.292178in}{0.856254in}}%
\pgfpathlineto{\pgfqpoint{3.292518in}{0.831507in}}%
\pgfpathlineto{\pgfqpoint{3.293037in}{0.823258in}}%
\pgfpathlineto{\pgfqpoint{3.293518in}{0.889251in}}%
\pgfpathlineto{\pgfqpoint{3.294207in}{0.946995in}}%
\pgfpathlineto{\pgfqpoint{3.293859in}{0.881002in}}%
\pgfpathlineto{\pgfqpoint{3.294547in}{0.922247in}}%
\pgfpathlineto{\pgfqpoint{3.294858in}{0.905749in}}%
\pgfpathlineto{\pgfqpoint{3.295199in}{0.938746in}}%
\pgfpathlineto{\pgfqpoint{3.295532in}{0.955244in}}%
\pgfpathlineto{\pgfqpoint{3.295887in}{0.922247in}}%
\pgfpathlineto{\pgfqpoint{3.296221in}{0.922247in}}%
\pgfpathlineto{\pgfqpoint{3.296391in}{0.946995in}}%
\pgfpathlineto{\pgfqpoint{3.297376in}{0.930497in}}%
\pgfpathlineto{\pgfqpoint{3.298397in}{0.922247in}}%
\pgfpathlineto{\pgfqpoint{3.299893in}{0.955244in}}%
\pgfpathlineto{\pgfqpoint{3.299908in}{0.946995in}}%
\pgfpathlineto{\pgfqpoint{3.300411in}{0.922247in}}%
\pgfpathlineto{\pgfqpoint{3.300730in}{0.955244in}}%
\pgfpathlineto{\pgfqpoint{3.300745in}{0.946995in}}%
\pgfpathlineto{\pgfqpoint{3.301737in}{0.955244in}}%
\pgfpathlineto{\pgfqpoint{3.301751in}{0.946995in}}%
\pgfpathlineto{\pgfqpoint{3.302758in}{0.922247in}}%
\pgfpathlineto{\pgfqpoint{3.302907in}{0.938746in}}%
\pgfpathlineto{\pgfqpoint{3.304247in}{1.029487in}}%
\pgfpathlineto{\pgfqpoint{3.304935in}{0.996490in}}%
\pgfpathlineto{\pgfqpoint{3.305276in}{1.021237in}}%
\pgfpathlineto{\pgfqpoint{3.305587in}{1.029487in}}%
\pgfpathlineto{\pgfqpoint{3.305942in}{0.996490in}}%
\pgfpathlineto{\pgfqpoint{3.306090in}{1.004739in}}%
\pgfpathlineto{\pgfqpoint{3.306764in}{0.955244in}}%
\pgfpathlineto{\pgfqpoint{3.307120in}{0.996490in}}%
\pgfpathlineto{\pgfqpoint{3.308778in}{1.054234in}}%
\pgfpathlineto{\pgfqpoint{3.308963in}{1.037736in}}%
\pgfpathlineto{\pgfqpoint{3.309126in}{1.045985in}}%
\pgfpathlineto{\pgfqpoint{3.310118in}{1.029487in}}%
\pgfpathlineto{\pgfqpoint{3.310133in}{1.045985in}}%
\pgfpathlineto{\pgfqpoint{3.310785in}{1.004739in}}%
\pgfpathlineto{\pgfqpoint{3.311140in}{1.021237in}}%
\pgfpathlineto{\pgfqpoint{3.311288in}{1.029487in}}%
\pgfpathlineto{\pgfqpoint{3.311644in}{0.996490in}}%
\pgfpathlineto{\pgfqpoint{3.311792in}{1.004739in}}%
\pgfpathlineto{\pgfqpoint{3.311814in}{0.996490in}}%
\pgfpathlineto{\pgfqpoint{3.312295in}{1.054234in}}%
\pgfpathlineto{\pgfqpoint{3.312650in}{1.045985in}}%
\pgfpathlineto{\pgfqpoint{3.312984in}{1.070732in}}%
\pgfpathlineto{\pgfqpoint{3.313806in}{1.054234in}}%
\pgfpathlineto{\pgfqpoint{3.314494in}{1.021237in}}%
\pgfpathlineto{\pgfqpoint{3.314827in}{1.045985in}}%
\pgfpathlineto{\pgfqpoint{3.315146in}{1.054234in}}%
\pgfpathlineto{\pgfqpoint{3.315479in}{1.029487in}}%
\pgfpathlineto{\pgfqpoint{3.316671in}{0.996490in}}%
\pgfpathlineto{\pgfqpoint{3.317012in}{1.021237in}}%
\pgfpathlineto{\pgfqpoint{3.317826in}{1.004739in}}%
\pgfpathlineto{\pgfqpoint{3.319018in}{0.971742in}}%
\pgfpathlineto{\pgfqpoint{3.320195in}{0.996490in}}%
\pgfpathlineto{\pgfqpoint{3.320358in}{0.988241in}}%
\pgfpathlineto{\pgfqpoint{3.320506in}{0.979992in}}%
\pgfpathlineto{\pgfqpoint{3.321010in}{1.012988in}}%
\pgfpathlineto{\pgfqpoint{3.321350in}{1.029487in}}%
\pgfpathlineto{\pgfqpoint{3.321698in}{0.996490in}}%
\pgfpathlineto{\pgfqpoint{3.322039in}{1.021237in}}%
\pgfpathlineto{\pgfqpoint{3.323194in}{1.029487in}}%
\pgfpathlineto{\pgfqpoint{3.324216in}{0.996490in}}%
\pgfpathlineto{\pgfqpoint{3.325371in}{1.054234in}}%
\pgfpathlineto{\pgfqpoint{3.325726in}{1.021237in}}%
\pgfpathlineto{\pgfqpoint{3.326045in}{1.004739in}}%
\pgfpathlineto{\pgfqpoint{3.326208in}{1.029487in}}%
\pgfpathlineto{\pgfqpoint{3.326541in}{1.029487in}}%
\pgfpathlineto{\pgfqpoint{3.326711in}{1.054234in}}%
\pgfpathlineto{\pgfqpoint{3.327570in}{1.045985in}}%
\pgfpathlineto{\pgfqpoint{3.328577in}{1.021237in}}%
\pgfpathlineto{\pgfqpoint{3.328725in}{1.037736in}}%
\pgfpathlineto{\pgfqpoint{3.329058in}{1.078981in}}%
\pgfpathlineto{\pgfqpoint{3.329732in}{1.029487in}}%
\pgfpathlineto{\pgfqpoint{3.329747in}{1.045985in}}%
\pgfpathlineto{\pgfqpoint{3.331405in}{1.004739in}}%
\pgfpathlineto{\pgfqpoint{3.331405in}{1.012988in}}%
\pgfpathlineto{\pgfqpoint{3.332412in}{1.029487in}}%
\pgfpathlineto{\pgfqpoint{3.331738in}{1.004739in}}%
\pgfpathlineto{\pgfqpoint{3.332427in}{1.021237in}}%
\pgfpathlineto{\pgfqpoint{3.333434in}{1.045985in}}%
\pgfpathlineto{\pgfqpoint{3.333604in}{1.037736in}}%
\pgfpathlineto{\pgfqpoint{3.333752in}{1.029487in}}%
\pgfpathlineto{\pgfqpoint{3.334086in}{1.062483in}}%
\pgfpathlineto{\pgfqpoint{3.334256in}{1.103729in}}%
\pgfpathlineto{\pgfqpoint{3.335115in}{1.070732in}}%
\pgfpathlineto{\pgfqpoint{3.337943in}{1.177971in}}%
\pgfpathlineto{\pgfqpoint{3.337965in}{1.169722in}}%
\pgfpathlineto{\pgfqpoint{3.339306in}{1.144975in}}%
\pgfpathlineto{\pgfqpoint{3.340957in}{1.177971in}}%
\pgfpathlineto{\pgfqpoint{3.342156in}{1.144975in}}%
\pgfpathlineto{\pgfqpoint{3.343326in}{1.169722in}}%
\pgfpathlineto{\pgfqpoint{3.343496in}{1.161473in}}%
\pgfpathlineto{\pgfqpoint{3.344666in}{1.144975in}}%
\pgfpathlineto{\pgfqpoint{3.345673in}{1.169722in}}%
\pgfpathlineto{\pgfqpoint{3.345843in}{1.161473in}}%
\pgfpathlineto{\pgfqpoint{3.347332in}{1.103729in}}%
\pgfpathlineto{\pgfqpoint{3.347517in}{1.120227in}}%
\pgfpathlineto{\pgfqpoint{3.349005in}{1.177971in}}%
\pgfpathlineto{\pgfqpoint{3.349190in}{1.161473in}}%
\pgfpathlineto{\pgfqpoint{3.350516in}{1.128476in}}%
\pgfpathlineto{\pgfqpoint{3.350538in}{1.144975in}}%
\pgfpathlineto{\pgfqpoint{3.351189in}{1.153224in}}%
\pgfpathlineto{\pgfqpoint{3.351374in}{1.136726in}}%
\pgfpathlineto{\pgfqpoint{3.353381in}{1.070732in}}%
\pgfpathlineto{\pgfqpoint{3.354033in}{1.078981in}}%
\pgfpathlineto{\pgfqpoint{3.354225in}{1.062483in}}%
\pgfpathlineto{\pgfqpoint{3.355395in}{1.045985in}}%
\pgfpathlineto{\pgfqpoint{3.356069in}{1.095480in}}%
\pgfpathlineto{\pgfqpoint{3.356402in}{1.062483in}}%
\pgfpathlineto{\pgfqpoint{3.357387in}{1.004739in}}%
\pgfpathlineto{\pgfqpoint{3.357742in}{1.021237in}}%
\pgfpathlineto{\pgfqpoint{3.359571in}{1.078981in}}%
\pgfpathlineto{\pgfqpoint{3.359586in}{1.070732in}}%
\pgfpathlineto{\pgfqpoint{3.360763in}{1.045985in}}%
\pgfpathlineto{\pgfqpoint{3.360911in}{1.054234in}}%
\pgfpathlineto{\pgfqpoint{3.361415in}{1.078981in}}%
\pgfpathlineto{\pgfqpoint{3.361762in}{1.045985in}}%
\pgfpathlineto{\pgfqpoint{3.361933in}{1.070732in}}%
\pgfpathlineto{\pgfqpoint{3.362584in}{1.054234in}}%
\pgfpathlineto{\pgfqpoint{3.362755in}{1.078981in}}%
\pgfpathlineto{\pgfqpoint{3.362769in}{1.070732in}}%
\pgfpathlineto{\pgfqpoint{3.363273in}{1.120227in}}%
\pgfpathlineto{\pgfqpoint{3.364280in}{1.095480in}}%
\pgfpathlineto{\pgfqpoint{3.365953in}{1.045985in}}%
\pgfpathlineto{\pgfqpoint{3.366101in}{1.054234in}}%
\pgfpathlineto{\pgfqpoint{3.366124in}{1.070732in}}%
\pgfpathlineto{\pgfqpoint{3.367131in}{1.037736in}}%
\pgfpathlineto{\pgfqpoint{3.367967in}{1.021237in}}%
\pgfpathlineto{\pgfqpoint{3.368115in}{1.037736in}}%
\pgfpathlineto{\pgfqpoint{3.368138in}{1.045985in}}%
\pgfpathlineto{\pgfqpoint{3.368804in}{0.996490in}}%
\pgfpathlineto{\pgfqpoint{3.369144in}{1.021237in}}%
\pgfpathlineto{\pgfqpoint{3.369293in}{1.029487in}}%
\pgfpathlineto{\pgfqpoint{3.369641in}{0.996490in}}%
\pgfpathlineto{\pgfqpoint{3.369959in}{1.004739in}}%
\pgfpathlineto{\pgfqpoint{3.370648in}{0.971742in}}%
\pgfpathlineto{\pgfqpoint{3.370988in}{0.996490in}}%
\pgfpathlineto{\pgfqpoint{3.372143in}{1.029487in}}%
\pgfpathlineto{\pgfqpoint{3.371640in}{0.979992in}}%
\pgfpathlineto{\pgfqpoint{3.372328in}{1.012988in}}%
\pgfpathlineto{\pgfqpoint{3.373335in}{0.971742in}}%
\pgfpathlineto{\pgfqpoint{3.373498in}{0.996490in}}%
\pgfpathlineto{\pgfqpoint{3.373654in}{1.004739in}}%
\pgfpathlineto{\pgfqpoint{3.374002in}{0.971742in}}%
\pgfpathlineto{\pgfqpoint{3.374150in}{0.979992in}}%
\pgfpathlineto{\pgfqpoint{3.374505in}{0.971742in}}%
\pgfpathlineto{\pgfqpoint{3.374675in}{0.996490in}}%
\pgfpathlineto{\pgfqpoint{3.375179in}{0.988241in}}%
\pgfpathlineto{\pgfqpoint{3.375342in}{0.996490in}}%
\pgfpathlineto{\pgfqpoint{3.376852in}{0.946995in}}%
\pgfpathlineto{\pgfqpoint{3.377348in}{0.979992in}}%
\pgfpathlineto{\pgfqpoint{3.377689in}{0.946995in}}%
\pgfpathlineto{\pgfqpoint{3.379370in}{0.872753in}}%
\pgfpathlineto{\pgfqpoint{3.379533in}{0.897500in}}%
\pgfpathlineto{\pgfqpoint{3.379688in}{0.905749in}}%
\pgfpathlineto{\pgfqpoint{3.380036in}{0.872753in}}%
\pgfpathlineto{\pgfqpoint{3.381361in}{0.856254in}}%
\pgfpathlineto{\pgfqpoint{3.382050in}{0.897500in}}%
\pgfpathlineto{\pgfqpoint{3.382383in}{0.872753in}}%
\pgfpathlineto{\pgfqpoint{3.383375in}{0.905749in}}%
\pgfpathlineto{\pgfqpoint{3.383560in}{0.889251in}}%
\pgfpathlineto{\pgfqpoint{3.385404in}{0.765513in}}%
\pgfpathlineto{\pgfqpoint{3.386574in}{0.749015in}}%
\pgfpathlineto{\pgfqpoint{3.386744in}{0.773763in}}%
\pgfpathlineto{\pgfqpoint{3.387248in}{0.740766in}}%
\pgfpathlineto{\pgfqpoint{3.387411in}{0.749015in}}%
\pgfpathlineto{\pgfqpoint{3.388418in}{0.724268in}}%
\pgfpathlineto{\pgfqpoint{3.389595in}{0.823258in}}%
\pgfpathlineto{\pgfqpoint{3.390765in}{0.806759in}}%
\pgfpathlineto{\pgfqpoint{3.391942in}{0.773763in}}%
\pgfpathlineto{\pgfqpoint{3.392105in}{0.790261in}}%
\pgfpathlineto{\pgfqpoint{3.392601in}{0.806759in}}%
\pgfpathlineto{\pgfqpoint{3.392275in}{0.782012in}}%
\pgfpathlineto{\pgfqpoint{3.393112in}{0.782012in}}%
\pgfpathlineto{\pgfqpoint{3.395296in}{0.707769in}}%
\pgfpathlineto{\pgfqpoint{3.395444in}{0.716019in}}%
\pgfpathlineto{\pgfqpoint{3.396296in}{0.749015in}}%
\pgfpathlineto{\pgfqpoint{3.396466in}{0.740766in}}%
\pgfpathlineto{\pgfqpoint{3.398310in}{0.625278in}}%
\pgfpathlineto{\pgfqpoint{3.399309in}{0.633527in}}%
\pgfpathlineto{\pgfqpoint{3.399317in}{0.625278in}}%
\pgfpathlineto{\pgfqpoint{3.400486in}{0.633527in}}%
\pgfpathlineto{\pgfqpoint{3.400657in}{0.650025in}}%
\pgfpathlineto{\pgfqpoint{3.401323in}{0.625278in}}%
\pgfpathlineto{\pgfqpoint{3.402478in}{0.633527in}}%
\pgfpathlineto{\pgfqpoint{3.402656in}{0.625278in}}%
\pgfpathlineto{\pgfqpoint{3.403507in}{0.625278in}}%
\pgfpathlineto{\pgfqpoint{3.404670in}{0.633527in}}%
\pgfpathlineto{\pgfqpoint{3.404848in}{0.625278in}}%
\pgfpathlineto{\pgfqpoint{3.405669in}{0.625278in}}%
\pgfpathlineto{\pgfqpoint{3.406847in}{0.633527in}}%
\pgfpathlineto{\pgfqpoint{3.407180in}{0.625278in}}%
\pgfpathlineto{\pgfqpoint{3.407854in}{0.625278in}}%
\pgfpathlineto{\pgfqpoint{3.408705in}{0.633527in}}%
\pgfpathlineto{\pgfqpoint{3.408705in}{0.625278in}}%
\pgfpathlineto{\pgfqpoint{3.409875in}{0.633527in}}%
\pgfpathlineto{\pgfqpoint{3.410038in}{0.625278in}}%
\pgfpathlineto{\pgfqpoint{3.410882in}{0.625278in}}%
\pgfpathlineto{\pgfqpoint{3.412059in}{0.633527in}}%
\pgfpathlineto{\pgfqpoint{3.412392in}{0.625278in}}%
\pgfpathlineto{\pgfqpoint{3.413066in}{0.625278in}}%
\pgfpathlineto{\pgfqpoint{3.414221in}{0.633527in}}%
\pgfpathlineto{\pgfqpoint{3.414236in}{0.625278in}}%
\pgfpathlineto{\pgfqpoint{3.415243in}{0.625278in}}%
\pgfpathlineto{\pgfqpoint{3.416405in}{0.633527in}}%
\pgfpathlineto{\pgfqpoint{3.416568in}{0.625278in}}%
\pgfpathlineto{\pgfqpoint{3.417412in}{0.625278in}}%
\pgfpathlineto{\pgfqpoint{3.418412in}{0.633527in}}%
\pgfpathlineto{\pgfqpoint{3.418412in}{0.625278in}}%
\pgfpathlineto{\pgfqpoint{3.419604in}{0.633527in}}%
\pgfpathlineto{\pgfqpoint{3.419922in}{0.625278in}}%
\pgfpathlineto{\pgfqpoint{3.420589in}{0.625278in}}%
\pgfpathlineto{\pgfqpoint{3.421781in}{0.633527in}}%
\pgfpathlineto{\pgfqpoint{3.421929in}{0.625278in}}%
\pgfpathlineto{\pgfqpoint{3.422432in}{0.625278in}}%
\pgfpathlineto{\pgfqpoint{3.423787in}{0.633527in}}%
\pgfpathlineto{\pgfqpoint{3.424121in}{0.625278in}}%
\pgfpathlineto{\pgfqpoint{3.424617in}{0.625278in}}%
\pgfpathlineto{\pgfqpoint{3.425972in}{0.633527in}}%
\pgfpathlineto{\pgfqpoint{3.426127in}{0.625278in}}%
\pgfpathlineto{\pgfqpoint{3.426979in}{0.625278in}}%
\pgfpathlineto{\pgfqpoint{3.428134in}{0.633527in}}%
\pgfpathlineto{\pgfqpoint{3.428134in}{0.625278in}}%
\pgfpathlineto{\pgfqpoint{3.429141in}{0.625278in}}%
\pgfpathlineto{\pgfqpoint{3.430325in}{0.633527in}}%
\pgfpathlineto{\pgfqpoint{3.431310in}{0.625278in}}%
\pgfpathlineto{\pgfqpoint{3.432495in}{0.633527in}}%
\pgfpathlineto{\pgfqpoint{3.432665in}{0.625278in}}%
\pgfpathlineto{\pgfqpoint{3.433494in}{0.625278in}}%
\pgfpathlineto{\pgfqpoint{3.434501in}{0.633527in}}%
\pgfpathlineto{\pgfqpoint{3.434501in}{0.625278in}}%
\pgfpathlineto{\pgfqpoint{3.435671in}{0.633527in}}%
\pgfpathlineto{\pgfqpoint{3.435671in}{0.625278in}}%
\pgfpathlineto{\pgfqpoint{3.436508in}{0.625278in}}%
\pgfpathlineto{\pgfqpoint{3.437685in}{0.633527in}}%
\pgfpathlineto{\pgfqpoint{3.437685in}{0.625278in}}%
\pgfpathlineto{\pgfqpoint{3.438692in}{0.625278in}}%
\pgfpathlineto{\pgfqpoint{3.439862in}{0.633527in}}%
\pgfpathlineto{\pgfqpoint{3.439862in}{0.625278in}}%
\pgfpathlineto{\pgfqpoint{3.440699in}{0.625278in}}%
\pgfpathlineto{\pgfqpoint{3.441876in}{0.633527in}}%
\pgfpathlineto{\pgfqpoint{3.442039in}{0.625278in}}%
\pgfpathlineto{\pgfqpoint{3.442883in}{0.625278in}}%
\pgfpathlineto{\pgfqpoint{3.444053in}{0.633527in}}%
\pgfpathlineto{\pgfqpoint{3.444053in}{0.625278in}}%
\pgfpathlineto{\pgfqpoint{3.445060in}{0.625278in}}%
\pgfpathlineto{\pgfqpoint{3.446230in}{0.633527in}}%
\pgfpathlineto{\pgfqpoint{3.446230in}{0.625278in}}%
\pgfpathlineto{\pgfqpoint{3.447081in}{0.625278in}}%
\pgfpathlineto{\pgfqpoint{3.448243in}{0.633527in}}%
\pgfpathlineto{\pgfqpoint{3.448577in}{0.625278in}}%
\pgfpathlineto{\pgfqpoint{3.449250in}{0.625278in}}%
\pgfpathlineto{\pgfqpoint{3.450420in}{0.633527in}}%
\pgfpathlineto{\pgfqpoint{3.450591in}{0.625278in}}%
\pgfpathlineto{\pgfqpoint{3.451427in}{0.625278in}}%
\pgfpathlineto{\pgfqpoint{3.452605in}{0.633527in}}%
\pgfpathlineto{\pgfqpoint{3.452767in}{0.625278in}}%
\pgfpathlineto{\pgfqpoint{3.453612in}{0.625278in}}%
\pgfpathlineto{\pgfqpoint{3.454611in}{0.633527in}}%
\pgfpathlineto{\pgfqpoint{3.454611in}{0.625278in}}%
\pgfpathlineto{\pgfqpoint{3.455959in}{0.633527in}}%
\pgfpathlineto{\pgfqpoint{3.456292in}{0.625278in}}%
\pgfpathlineto{\pgfqpoint{3.456958in}{0.625278in}}%
\pgfpathlineto{\pgfqpoint{3.458136in}{0.633527in}}%
\pgfpathlineto{\pgfqpoint{3.458298in}{0.625278in}}%
\pgfpathlineto{\pgfqpoint{3.459142in}{0.625278in}}%
\pgfpathlineto{\pgfqpoint{3.460312in}{0.633527in}}%
\pgfpathlineto{\pgfqpoint{3.460312in}{0.625278in}}%
\pgfpathlineto{\pgfqpoint{3.461149in}{0.625278in}}%
\pgfpathlineto{\pgfqpoint{3.462489in}{0.633527in}}%
\pgfpathlineto{\pgfqpoint{3.462489in}{0.625278in}}%
\pgfpathlineto{\pgfqpoint{3.463496in}{0.625278in}}%
\pgfpathlineto{\pgfqpoint{3.464673in}{0.633527in}}%
\pgfpathlineto{\pgfqpoint{3.464836in}{0.625278in}}%
\pgfpathlineto{\pgfqpoint{3.465510in}{0.625278in}}%
\pgfpathlineto{\pgfqpoint{3.466850in}{0.633527in}}%
\pgfpathlineto{\pgfqpoint{3.467021in}{0.625278in}}%
\pgfpathlineto{\pgfqpoint{3.467857in}{0.625278in}}%
\pgfpathlineto{\pgfqpoint{3.469027in}{0.633527in}}%
\pgfpathlineto{\pgfqpoint{3.469197in}{0.625278in}}%
\pgfpathlineto{\pgfqpoint{3.470034in}{0.625278in}}%
\pgfpathlineto{\pgfqpoint{3.471211in}{0.633527in}}%
\pgfpathlineto{\pgfqpoint{3.471544in}{0.625278in}}%
\pgfpathlineto{\pgfqpoint{3.472218in}{0.625278in}}%
\pgfpathlineto{\pgfqpoint{3.473558in}{0.633527in}}%
\pgfpathlineto{\pgfqpoint{3.473558in}{0.625278in}}%
\pgfpathlineto{\pgfqpoint{3.474565in}{0.625278in}}%
\pgfpathlineto{\pgfqpoint{3.475735in}{0.633527in}}%
\pgfpathlineto{\pgfqpoint{3.475735in}{0.625278in}}%
\pgfpathlineto{\pgfqpoint{3.476742in}{0.625278in}}%
\pgfpathlineto{\pgfqpoint{3.477912in}{0.633527in}}%
\pgfpathlineto{\pgfqpoint{3.478082in}{0.625278in}}%
\pgfpathlineto{\pgfqpoint{3.478919in}{0.625278in}}%
\pgfpathlineto{\pgfqpoint{3.479926in}{0.633527in}}%
\pgfpathlineto{\pgfqpoint{3.479926in}{0.625278in}}%
\pgfpathlineto{\pgfqpoint{3.481103in}{0.633527in}}%
\pgfpathlineto{\pgfqpoint{3.481266in}{0.625278in}}%
\pgfpathlineto{\pgfqpoint{3.481940in}{0.625278in}}%
\pgfpathlineto{\pgfqpoint{3.483110in}{0.633527in}}%
\pgfpathlineto{\pgfqpoint{3.483280in}{0.625278in}}%
\pgfpathlineto{\pgfqpoint{3.484117in}{0.625278in}}%
\pgfpathlineto{\pgfqpoint{3.485294in}{0.633527in}}%
\pgfpathlineto{\pgfqpoint{3.485457in}{0.625278in}}%
\pgfpathlineto{\pgfqpoint{3.486294in}{0.625278in}}%
\pgfpathlineto{\pgfqpoint{3.487634in}{0.633527in}}%
\pgfpathlineto{\pgfqpoint{3.487974in}{0.625278in}}%
\pgfpathlineto{\pgfqpoint{3.488641in}{0.625278in}}%
\pgfpathlineto{\pgfqpoint{3.489648in}{0.633527in}}%
\pgfpathlineto{\pgfqpoint{3.489648in}{0.625278in}}%
\pgfpathlineto{\pgfqpoint{3.490825in}{0.633527in}}%
\pgfpathlineto{\pgfqpoint{3.490825in}{0.625278in}}%
\pgfpathlineto{\pgfqpoint{3.491825in}{0.625278in}}%
\pgfpathlineto{\pgfqpoint{3.493002in}{0.633527in}}%
\pgfpathlineto{\pgfqpoint{3.493002in}{0.625278in}}%
\pgfpathlineto{\pgfqpoint{3.494009in}{0.625278in}}%
\pgfpathlineto{\pgfqpoint{3.495016in}{0.633527in}}%
\pgfpathlineto{\pgfqpoint{3.495016in}{0.625278in}}%
\pgfpathlineto{\pgfqpoint{3.496356in}{0.633527in}}%
\pgfpathlineto{\pgfqpoint{3.496356in}{0.625278in}}%
\pgfpathlineto{\pgfqpoint{3.497363in}{0.625278in}}%
\pgfpathlineto{\pgfqpoint{3.498533in}{0.633527in}}%
\pgfpathlineto{\pgfqpoint{3.498866in}{0.625278in}}%
\pgfpathlineto{\pgfqpoint{3.499540in}{0.625278in}}%
\pgfpathlineto{\pgfqpoint{3.500547in}{0.633527in}}%
\pgfpathlineto{\pgfqpoint{3.500547in}{0.625278in}}%
\pgfpathlineto{\pgfqpoint{3.501887in}{0.633527in}}%
\pgfpathlineto{\pgfqpoint{3.501887in}{0.625278in}}%
\pgfpathlineto{\pgfqpoint{3.502894in}{0.625278in}}%
\pgfpathlineto{\pgfqpoint{3.504064in}{0.633527in}}%
\pgfpathlineto{\pgfqpoint{3.504234in}{0.625278in}}%
\pgfpathlineto{\pgfqpoint{3.505071in}{0.625278in}}%
\pgfpathlineto{\pgfqpoint{3.506248in}{0.633527in}}%
\pgfpathlineto{\pgfqpoint{3.506248in}{0.625278in}}%
\pgfpathlineto{\pgfqpoint{3.507085in}{0.625278in}}%
\pgfpathlineto{\pgfqpoint{3.508254in}{0.633527in}}%
\pgfpathlineto{\pgfqpoint{3.508425in}{0.625278in}}%
\pgfpathlineto{\pgfqpoint{3.509261in}{0.625278in}}%
\pgfpathlineto{\pgfqpoint{3.510602in}{0.633527in}}%
\pgfpathlineto{\pgfqpoint{3.510772in}{0.625278in}}%
\pgfpathlineto{\pgfqpoint{3.511609in}{0.625278in}}%
\pgfpathlineto{\pgfqpoint{3.512778in}{0.633527in}}%
\pgfpathlineto{\pgfqpoint{3.512949in}{0.625278in}}%
\pgfpathlineto{\pgfqpoint{3.513785in}{0.625278in}}%
\pgfpathlineto{\pgfqpoint{3.514963in}{0.633527in}}%
\pgfpathlineto{\pgfqpoint{3.515126in}{0.625278in}}%
\pgfpathlineto{\pgfqpoint{3.515970in}{0.625278in}}%
\pgfpathlineto{\pgfqpoint{3.517140in}{0.633527in}}%
\pgfpathlineto{\pgfqpoint{3.517473in}{0.625278in}}%
\pgfpathlineto{\pgfqpoint{3.518147in}{0.625278in}}%
\pgfpathlineto{\pgfqpoint{3.519487in}{0.633527in}}%
\pgfpathlineto{\pgfqpoint{3.519487in}{0.625278in}}%
\pgfpathlineto{\pgfqpoint{3.520323in}{0.625278in}}%
\pgfpathlineto{\pgfqpoint{3.521501in}{0.633527in}}%
\pgfpathlineto{\pgfqpoint{3.521663in}{0.625278in}}%
\pgfpathlineto{\pgfqpoint{3.522508in}{0.625278in}}%
\pgfpathlineto{\pgfqpoint{3.523677in}{0.633527in}}%
\pgfpathlineto{\pgfqpoint{3.523848in}{0.625278in}}%
\pgfpathlineto{\pgfqpoint{3.524684in}{0.625278in}}%
\pgfpathlineto{\pgfqpoint{3.525854in}{0.633527in}}%
\pgfpathlineto{\pgfqpoint{3.526025in}{0.625278in}}%
\pgfpathlineto{\pgfqpoint{3.526861in}{0.625278in}}%
\pgfpathlineto{\pgfqpoint{3.528201in}{0.633527in}}%
\pgfpathlineto{\pgfqpoint{3.528201in}{0.625278in}}%
\pgfpathlineto{\pgfqpoint{3.529208in}{0.625278in}}%
\pgfpathlineto{\pgfqpoint{3.530386in}{0.633527in}}%
\pgfpathlineto{\pgfqpoint{3.530549in}{0.625278in}}%
\pgfpathlineto{\pgfqpoint{3.531393in}{0.625278in}}%
\pgfpathlineto{\pgfqpoint{3.532562in}{0.633527in}}%
\pgfpathlineto{\pgfqpoint{3.532733in}{0.625278in}}%
\pgfpathlineto{\pgfqpoint{3.533569in}{0.625278in}}%
\pgfpathlineto{\pgfqpoint{3.534739in}{0.633527in}}%
\pgfpathlineto{\pgfqpoint{3.534739in}{0.625278in}}%
\pgfpathlineto{\pgfqpoint{3.535746in}{0.625278in}}%
\pgfpathlineto{\pgfqpoint{3.536753in}{0.633527in}}%
\pgfpathlineto{\pgfqpoint{3.536753in}{0.625278in}}%
\pgfpathlineto{\pgfqpoint{3.538093in}{0.633527in}}%
\pgfpathlineto{\pgfqpoint{3.538093in}{0.625278in}}%
\pgfpathlineto{\pgfqpoint{3.539100in}{0.625278in}}%
\pgfpathlineto{\pgfqpoint{3.540270in}{0.633527in}}%
\pgfpathlineto{\pgfqpoint{3.540441in}{0.625278in}}%
\pgfpathlineto{\pgfqpoint{3.541277in}{0.625278in}}%
\pgfpathlineto{\pgfqpoint{3.542454in}{0.633527in}}%
\pgfpathlineto{\pgfqpoint{3.542617in}{0.625278in}}%
\pgfpathlineto{\pgfqpoint{3.543461in}{0.625278in}}%
\pgfpathlineto{\pgfqpoint{3.544631in}{0.633527in}}%
\pgfpathlineto{\pgfqpoint{3.544802in}{0.625278in}}%
\pgfpathlineto{\pgfqpoint{3.545638in}{0.625278in}}%
\pgfpathlineto{\pgfqpoint{3.546808in}{0.633527in}}%
\pgfpathlineto{\pgfqpoint{3.546978in}{0.625278in}}%
\pgfpathlineto{\pgfqpoint{3.547815in}{0.625278in}}%
\pgfpathlineto{\pgfqpoint{3.548992in}{0.633527in}}%
\pgfpathlineto{\pgfqpoint{3.549155in}{0.625278in}}%
\pgfpathlineto{\pgfqpoint{3.549999in}{0.625278in}}%
\pgfpathlineto{\pgfqpoint{3.551169in}{0.633527in}}%
\pgfpathlineto{\pgfqpoint{3.551502in}{0.625278in}}%
\pgfpathlineto{\pgfqpoint{3.552176in}{0.625278in}}%
\pgfpathlineto{\pgfqpoint{3.553516in}{0.633527in}}%
\pgfpathlineto{\pgfqpoint{3.553516in}{0.625278in}}%
\pgfpathlineto{\pgfqpoint{3.554523in}{0.625278in}}%
\pgfpathlineto{\pgfqpoint{3.555693in}{0.633527in}}%
\pgfpathlineto{\pgfqpoint{3.555863in}{0.625278in}}%
\pgfpathlineto{\pgfqpoint{3.556700in}{0.625278in}}%
\pgfpathlineto{\pgfqpoint{3.557877in}{0.633527in}}%
\pgfpathlineto{\pgfqpoint{3.558040in}{0.625278in}}%
\pgfpathlineto{\pgfqpoint{3.558877in}{0.625278in}}%
\pgfpathlineto{\pgfqpoint{3.559884in}{0.633527in}}%
\pgfpathlineto{\pgfqpoint{3.559884in}{0.625278in}}%
\pgfpathlineto{\pgfqpoint{3.560950in}{0.633527in}}%
\pgfpathlineto{\pgfqpoint{3.560950in}{0.625278in}}%
\pgfpathlineto{\pgfqpoint{3.562068in}{0.633527in}}%
\pgfpathlineto{\pgfqpoint{3.562401in}{0.625278in}}%
\pgfpathlineto{\pgfqpoint{3.563068in}{0.625278in}}%
\pgfpathlineto{\pgfqpoint{3.564415in}{0.633527in}}%
\pgfpathlineto{\pgfqpoint{3.564415in}{0.625278in}}%
\pgfpathlineto{\pgfqpoint{3.565415in}{0.625278in}}%
\pgfpathlineto{\pgfqpoint{3.566592in}{0.633527in}}%
\pgfpathlineto{\pgfqpoint{3.566607in}{0.625278in}}%
\pgfpathlineto{\pgfqpoint{3.567614in}{0.625278in}}%
\pgfpathlineto{\pgfqpoint{3.568606in}{0.658274in}}%
\pgfpathlineto{\pgfqpoint{3.568791in}{0.641776in}}%
\pgfpathlineto{\pgfqpoint{3.570131in}{0.625278in}}%
\pgfpathlineto{\pgfqpoint{3.571116in}{0.707769in}}%
\pgfpathlineto{\pgfqpoint{3.572308in}{0.691271in}}%
\pgfpathlineto{\pgfqpoint{3.573130in}{0.683022in}}%
\pgfpathlineto{\pgfqpoint{3.572478in}{0.699520in}}%
\pgfpathlineto{\pgfqpoint{3.573130in}{0.691271in}}%
\pgfpathlineto{\pgfqpoint{3.573967in}{0.732517in}}%
\pgfpathlineto{\pgfqpoint{3.574152in}{0.724268in}}%
\pgfpathlineto{\pgfqpoint{3.575307in}{0.732517in}}%
\pgfpathlineto{\pgfqpoint{3.576144in}{0.707769in}}%
\pgfpathlineto{\pgfqpoint{3.576336in}{0.724268in}}%
\pgfpathlineto{\pgfqpoint{3.576647in}{0.732517in}}%
\pgfpathlineto{\pgfqpoint{3.577002in}{0.699520in}}%
\pgfpathlineto{\pgfqpoint{3.577151in}{0.707769in}}%
\pgfpathlineto{\pgfqpoint{3.577491in}{0.732517in}}%
\pgfpathlineto{\pgfqpoint{3.578180in}{0.699520in}}%
\pgfpathlineto{\pgfqpoint{3.579668in}{0.757264in}}%
\pgfpathlineto{\pgfqpoint{3.579853in}{0.740766in}}%
\pgfpathlineto{\pgfqpoint{3.580357in}{0.724268in}}%
\pgfpathlineto{\pgfqpoint{3.580527in}{0.749015in}}%
\pgfpathlineto{\pgfqpoint{3.580838in}{0.757264in}}%
\pgfpathlineto{\pgfqpoint{3.581178in}{0.732517in}}%
\pgfpathlineto{\pgfqpoint{3.581364in}{0.740766in}}%
\pgfpathlineto{\pgfqpoint{3.581526in}{0.749015in}}%
\pgfpathlineto{\pgfqpoint{3.582533in}{0.724268in}}%
\pgfpathlineto{\pgfqpoint{3.583688in}{0.732517in}}%
\pgfpathlineto{\pgfqpoint{3.584718in}{0.724268in}}%
\pgfpathlineto{\pgfqpoint{3.585029in}{0.732517in}}%
\pgfpathlineto{\pgfqpoint{3.585384in}{0.699520in}}%
\pgfpathlineto{\pgfqpoint{3.585532in}{0.707769in}}%
\pgfpathlineto{\pgfqpoint{3.585873in}{0.732517in}}%
\pgfpathlineto{\pgfqpoint{3.586561in}{0.699520in}}%
\pgfpathlineto{\pgfqpoint{3.587901in}{0.749015in}}%
\pgfpathlineto{\pgfqpoint{3.588212in}{0.732517in}}%
\pgfpathlineto{\pgfqpoint{3.588886in}{0.707769in}}%
\pgfpathlineto{\pgfqpoint{3.589242in}{0.749015in}}%
\pgfpathlineto{\pgfqpoint{3.590249in}{0.773763in}}%
\pgfpathlineto{\pgfqpoint{3.590411in}{0.765513in}}%
\pgfpathlineto{\pgfqpoint{3.592403in}{0.707769in}}%
\pgfpathlineto{\pgfqpoint{3.592403in}{0.716019in}}%
\pgfpathlineto{\pgfqpoint{3.593580in}{0.782012in}}%
\pgfpathlineto{\pgfqpoint{3.592744in}{0.707769in}}%
\pgfpathlineto{\pgfqpoint{3.593603in}{0.773763in}}%
\pgfpathlineto{\pgfqpoint{3.595424in}{0.831507in}}%
\pgfpathlineto{\pgfqpoint{3.595446in}{0.823258in}}%
\pgfpathlineto{\pgfqpoint{3.595942in}{0.798510in}}%
\pgfpathlineto{\pgfqpoint{3.596261in}{0.831507in}}%
\pgfpathlineto{\pgfqpoint{3.596283in}{0.823258in}}%
\pgfpathlineto{\pgfqpoint{3.597120in}{0.848005in}}%
\pgfpathlineto{\pgfqpoint{3.596764in}{0.806759in}}%
\pgfpathlineto{\pgfqpoint{3.597438in}{0.831507in}}%
\pgfpathlineto{\pgfqpoint{3.598445in}{0.782012in}}%
\pgfpathlineto{\pgfqpoint{3.598460in}{0.798510in}}%
\pgfpathlineto{\pgfqpoint{3.599134in}{0.823258in}}%
\pgfpathlineto{\pgfqpoint{3.599615in}{0.806759in}}%
\pgfpathlineto{\pgfqpoint{3.599948in}{0.782012in}}%
\pgfpathlineto{\pgfqpoint{3.600474in}{0.823258in}}%
\pgfpathlineto{\pgfqpoint{3.600637in}{0.815008in}}%
\pgfpathlineto{\pgfqpoint{3.602132in}{0.782012in}}%
\pgfpathlineto{\pgfqpoint{3.602132in}{0.790261in}}%
\pgfpathlineto{\pgfqpoint{3.602465in}{0.782012in}}%
\pgfpathlineto{\pgfqpoint{3.603154in}{0.798510in}}%
\pgfpathlineto{\pgfqpoint{3.603472in}{0.782012in}}%
\pgfpathlineto{\pgfqpoint{3.603806in}{0.815008in}}%
\pgfpathlineto{\pgfqpoint{3.603976in}{0.856254in}}%
\pgfpathlineto{\pgfqpoint{3.604827in}{0.848005in}}%
\pgfpathlineto{\pgfqpoint{3.605501in}{0.872753in}}%
\pgfpathlineto{\pgfqpoint{3.605982in}{0.856254in}}%
\pgfpathlineto{\pgfqpoint{3.607175in}{0.823258in}}%
\pgfpathlineto{\pgfqpoint{3.608019in}{0.848005in}}%
\pgfpathlineto{\pgfqpoint{3.608330in}{0.831507in}}%
\pgfpathlineto{\pgfqpoint{3.608670in}{0.806759in}}%
\pgfpathlineto{\pgfqpoint{3.609166in}{0.856254in}}%
\pgfpathlineto{\pgfqpoint{3.609359in}{0.839756in}}%
\pgfpathlineto{\pgfqpoint{3.610025in}{0.848005in}}%
\pgfpathlineto{\pgfqpoint{3.610529in}{0.823258in}}%
\pgfpathlineto{\pgfqpoint{3.610862in}{0.848005in}}%
\pgfpathlineto{\pgfqpoint{3.611365in}{0.815008in}}%
\pgfpathlineto{\pgfqpoint{3.611536in}{0.823258in}}%
\pgfpathlineto{\pgfqpoint{3.612520in}{0.806759in}}%
\pgfpathlineto{\pgfqpoint{3.613194in}{0.856254in}}%
\pgfpathlineto{\pgfqpoint{3.613550in}{0.839756in}}%
\pgfpathlineto{\pgfqpoint{3.614201in}{0.856254in}}%
\pgfpathlineto{\pgfqpoint{3.614719in}{0.823258in}}%
\pgfpathlineto{\pgfqpoint{3.615874in}{0.831507in}}%
\pgfpathlineto{\pgfqpoint{3.616563in}{0.848005in}}%
\pgfpathlineto{\pgfqpoint{3.616896in}{0.823258in}}%
\pgfpathlineto{\pgfqpoint{3.617718in}{0.831507in}}%
\pgfpathlineto{\pgfqpoint{3.618244in}{0.798510in}}%
\pgfpathlineto{\pgfqpoint{3.620739in}{0.881002in}}%
\pgfpathlineto{\pgfqpoint{3.621739in}{0.831507in}}%
\pgfpathlineto{\pgfqpoint{3.621761in}{0.848005in}}%
\pgfpathlineto{\pgfqpoint{3.622079in}{0.831507in}}%
\pgfpathlineto{\pgfqpoint{3.622094in}{0.872753in}}%
\pgfpathlineto{\pgfqpoint{3.622412in}{0.864503in}}%
\pgfpathlineto{\pgfqpoint{3.622746in}{0.856254in}}%
\pgfpathlineto{\pgfqpoint{3.623434in}{0.872753in}}%
\pgfpathlineto{\pgfqpoint{3.623753in}{0.856254in}}%
\pgfpathlineto{\pgfqpoint{3.623923in}{0.881002in}}%
\pgfpathlineto{\pgfqpoint{3.624256in}{0.881002in}}%
\pgfpathlineto{\pgfqpoint{3.624278in}{0.897500in}}%
\pgfpathlineto{\pgfqpoint{3.624611in}{0.872753in}}%
\pgfpathlineto{\pgfqpoint{3.625278in}{0.872753in}}%
\pgfpathlineto{\pgfqpoint{3.626100in}{0.905749in}}%
\pgfpathlineto{\pgfqpoint{3.626433in}{0.881002in}}%
\pgfpathlineto{\pgfqpoint{3.627462in}{0.872753in}}%
\pgfpathlineto{\pgfqpoint{3.628617in}{0.881002in}}%
\pgfpathlineto{\pgfqpoint{3.628972in}{0.872753in}}%
\pgfpathlineto{\pgfqpoint{3.629454in}{0.955244in}}%
\pgfpathlineto{\pgfqpoint{3.629469in}{0.946995in}}%
\pgfpathlineto{\pgfqpoint{3.630646in}{0.971742in}}%
\pgfpathlineto{\pgfqpoint{3.630816in}{0.963493in}}%
\pgfpathlineto{\pgfqpoint{3.631297in}{0.930497in}}%
\pgfpathlineto{\pgfqpoint{3.631653in}{0.971742in}}%
\pgfpathlineto{\pgfqpoint{3.631801in}{0.979992in}}%
\pgfpathlineto{\pgfqpoint{3.632156in}{0.946995in}}%
\pgfpathlineto{\pgfqpoint{3.632467in}{0.955244in}}%
\pgfpathlineto{\pgfqpoint{3.632993in}{0.946995in}}%
\pgfpathlineto{\pgfqpoint{3.633474in}{0.988241in}}%
\pgfpathlineto{\pgfqpoint{3.634481in}{1.004739in}}%
\pgfpathlineto{\pgfqpoint{3.633830in}{0.971742in}}%
\pgfpathlineto{\pgfqpoint{3.634503in}{0.996490in}}%
\pgfpathlineto{\pgfqpoint{3.635007in}{0.971742in}}%
\pgfpathlineto{\pgfqpoint{3.635318in}{1.004739in}}%
\pgfpathlineto{\pgfqpoint{3.635658in}{1.054234in}}%
\pgfpathlineto{\pgfqpoint{3.636347in}{1.021237in}}%
\pgfpathlineto{\pgfqpoint{3.636999in}{1.054234in}}%
\pgfpathlineto{\pgfqpoint{3.637354in}{1.021237in}}%
\pgfpathlineto{\pgfqpoint{3.638339in}{1.004739in}}%
\pgfpathlineto{\pgfqpoint{3.638354in}{1.021237in}}%
\pgfpathlineto{\pgfqpoint{3.640034in}{1.095480in}}%
\pgfpathlineto{\pgfqpoint{3.640353in}{1.078981in}}%
\pgfpathlineto{\pgfqpoint{3.640849in}{1.054234in}}%
\pgfpathlineto{\pgfqpoint{3.641375in}{1.070732in}}%
\pgfpathlineto{\pgfqpoint{3.642041in}{1.095480in}}%
\pgfpathlineto{\pgfqpoint{3.642381in}{1.070732in}}%
\pgfpathlineto{\pgfqpoint{3.643033in}{1.054234in}}%
\pgfpathlineto{\pgfqpoint{3.643196in}{1.078981in}}%
\pgfpathlineto{\pgfqpoint{3.643218in}{1.070732in}}%
\pgfpathlineto{\pgfqpoint{3.644373in}{1.078981in}}%
\pgfpathlineto{\pgfqpoint{3.644892in}{1.070732in}}%
\pgfpathlineto{\pgfqpoint{3.645395in}{1.120227in}}%
\pgfpathlineto{\pgfqpoint{3.645884in}{1.103729in}}%
\pgfpathlineto{\pgfqpoint{3.646550in}{1.128476in}}%
\pgfpathlineto{\pgfqpoint{3.646883in}{1.103729in}}%
\pgfpathlineto{\pgfqpoint{3.647579in}{1.120227in}}%
\pgfpathlineto{\pgfqpoint{3.649230in}{1.153224in}}%
\pgfpathlineto{\pgfqpoint{3.649571in}{1.177971in}}%
\pgfpathlineto{\pgfqpoint{3.650422in}{1.120227in}}%
\pgfpathlineto{\pgfqpoint{3.651074in}{1.103729in}}%
\pgfpathlineto{\pgfqpoint{3.651244in}{1.128476in}}%
\pgfpathlineto{\pgfqpoint{3.651267in}{1.120227in}}%
\pgfpathlineto{\pgfqpoint{3.651918in}{1.177971in}}%
\pgfpathlineto{\pgfqpoint{3.652436in}{1.136726in}}%
\pgfpathlineto{\pgfqpoint{3.652607in}{1.144975in}}%
\pgfpathlineto{\pgfqpoint{3.653591in}{1.128476in}}%
\pgfpathlineto{\pgfqpoint{3.654613in}{1.194470in}}%
\pgfpathlineto{\pgfqpoint{3.655620in}{1.169722in}}%
\pgfpathlineto{\pgfqpoint{3.655768in}{1.186221in}}%
\pgfpathlineto{\pgfqpoint{3.655790in}{1.194470in}}%
\pgfpathlineto{\pgfqpoint{3.656627in}{1.144975in}}%
\pgfpathlineto{\pgfqpoint{3.656775in}{1.153224in}}%
\pgfpathlineto{\pgfqpoint{3.657301in}{1.169722in}}%
\pgfpathlineto{\pgfqpoint{3.657804in}{1.144975in}}%
\pgfpathlineto{\pgfqpoint{3.658619in}{1.128476in}}%
\pgfpathlineto{\pgfqpoint{3.658959in}{1.153224in}}%
\pgfpathlineto{\pgfqpoint{3.659981in}{1.144975in}}%
\pgfpathlineto{\pgfqpoint{3.661136in}{1.153224in}}%
\pgfpathlineto{\pgfqpoint{3.661321in}{1.144975in}}%
\pgfpathlineto{\pgfqpoint{3.661825in}{1.194470in}}%
\pgfpathlineto{\pgfqpoint{3.662158in}{1.169722in}}%
\pgfpathlineto{\pgfqpoint{3.663313in}{1.177971in}}%
\pgfpathlineto{\pgfqpoint{3.663817in}{1.153224in}}%
\pgfpathlineto{\pgfqpoint{3.664342in}{1.169722in}}%
\pgfpathlineto{\pgfqpoint{3.664490in}{1.177971in}}%
\pgfpathlineto{\pgfqpoint{3.664824in}{1.153224in}}%
\pgfpathlineto{\pgfqpoint{3.665179in}{1.169722in}}%
\pgfpathlineto{\pgfqpoint{3.666519in}{1.144975in}}%
\pgfpathlineto{\pgfqpoint{3.667504in}{1.153224in}}%
\pgfpathlineto{\pgfqpoint{3.667526in}{1.144975in}}%
\pgfpathlineto{\pgfqpoint{3.669873in}{1.070732in}}%
\pgfpathlineto{\pgfqpoint{3.671865in}{1.128476in}}%
\pgfpathlineto{\pgfqpoint{3.672872in}{1.103729in}}%
\pgfpathlineto{\pgfqpoint{3.672887in}{1.120227in}}%
\pgfpathlineto{\pgfqpoint{3.673879in}{1.078981in}}%
\pgfpathlineto{\pgfqpoint{3.674212in}{1.103729in}}%
\pgfpathlineto{\pgfqpoint{3.674375in}{1.128476in}}%
\pgfpathlineto{\pgfqpoint{3.675234in}{1.087231in}}%
\pgfpathlineto{\pgfqpoint{3.675737in}{1.070732in}}%
\pgfpathlineto{\pgfqpoint{3.675908in}{1.095480in}}%
\pgfpathlineto{\pgfqpoint{3.676219in}{1.103729in}}%
\pgfpathlineto{\pgfqpoint{3.676574in}{1.070732in}}%
\pgfpathlineto{\pgfqpoint{3.676722in}{1.078981in}}%
\pgfpathlineto{\pgfqpoint{3.677751in}{1.070732in}}%
\pgfpathlineto{\pgfqpoint{3.677914in}{1.095480in}}%
\pgfpathlineto{\pgfqpoint{3.678566in}{1.054234in}}%
\pgfpathlineto{\pgfqpoint{3.678906in}{1.078981in}}%
\pgfpathlineto{\pgfqpoint{3.679758in}{1.070732in}}%
\pgfpathlineto{\pgfqpoint{3.679965in}{1.103729in}}%
\pgfpathlineto{\pgfqpoint{3.681772in}{1.045985in}}%
\pgfpathlineto{\pgfqpoint{3.683097in}{1.078981in}}%
\pgfpathlineto{\pgfqpoint{3.683112in}{1.070732in}}%
\pgfpathlineto{\pgfqpoint{3.683949in}{1.021237in}}%
\pgfpathlineto{\pgfqpoint{3.684437in}{1.054234in}}%
\pgfpathlineto{\pgfqpoint{3.684452in}{1.070732in}}%
\pgfpathlineto{\pgfqpoint{3.685126in}{1.045985in}}%
\pgfpathlineto{\pgfqpoint{3.685459in}{1.045985in}}%
\pgfpathlineto{\pgfqpoint{3.685948in}{1.054234in}}%
\pgfpathlineto{\pgfqpoint{3.686133in}{1.037736in}}%
\pgfpathlineto{\pgfqpoint{3.686466in}{1.021237in}}%
\pgfpathlineto{\pgfqpoint{3.686784in}{1.054234in}}%
\pgfpathlineto{\pgfqpoint{3.686970in}{1.045985in}}%
\pgfpathlineto{\pgfqpoint{3.687118in}{1.054234in}}%
\pgfpathlineto{\pgfqpoint{3.687473in}{1.012988in}}%
\pgfpathlineto{\pgfqpoint{3.688813in}{0.971742in}}%
\pgfpathlineto{\pgfqpoint{3.689465in}{0.979992in}}%
\pgfpathlineto{\pgfqpoint{3.689650in}{0.963493in}}%
\pgfpathlineto{\pgfqpoint{3.690153in}{0.946995in}}%
\pgfpathlineto{\pgfqpoint{3.690324in}{0.971742in}}%
\pgfpathlineto{\pgfqpoint{3.690657in}{0.996490in}}%
\pgfpathlineto{\pgfqpoint{3.691160in}{0.963493in}}%
\pgfpathlineto{\pgfqpoint{3.691331in}{0.971742in}}%
\pgfpathlineto{\pgfqpoint{3.692315in}{0.955244in}}%
\pgfpathlineto{\pgfqpoint{3.692330in}{0.996490in}}%
\pgfpathlineto{\pgfqpoint{3.693337in}{0.971742in}}%
\pgfpathlineto{\pgfqpoint{3.695669in}{0.905749in}}%
\pgfpathlineto{\pgfqpoint{3.696506in}{0.979992in}}%
\pgfpathlineto{\pgfqpoint{3.696691in}{0.963493in}}%
\pgfpathlineto{\pgfqpoint{3.697024in}{0.946995in}}%
\pgfpathlineto{\pgfqpoint{3.697350in}{0.979992in}}%
\pgfpathlineto{\pgfqpoint{3.697365in}{0.971742in}}%
\pgfpathlineto{\pgfqpoint{3.697528in}{0.996490in}}%
\pgfpathlineto{\pgfqpoint{3.698031in}{0.963493in}}%
\pgfpathlineto{\pgfqpoint{3.698520in}{0.955244in}}%
\pgfpathlineto{\pgfqpoint{3.698527in}{0.988241in}}%
\pgfpathlineto{\pgfqpoint{3.698535in}{0.996490in}}%
\pgfpathlineto{\pgfqpoint{3.699186in}{0.955244in}}%
\pgfpathlineto{\pgfqpoint{3.699542in}{0.971742in}}%
\pgfpathlineto{\pgfqpoint{3.700193in}{0.930497in}}%
\pgfpathlineto{\pgfqpoint{3.700875in}{0.955244in}}%
\pgfpathlineto{\pgfqpoint{3.700882in}{0.971742in}}%
\pgfpathlineto{\pgfqpoint{3.701889in}{0.946995in}}%
\pgfpathlineto{\pgfqpoint{3.703059in}{0.922247in}}%
\pgfpathlineto{\pgfqpoint{3.703214in}{0.930497in}}%
\pgfpathlineto{\pgfqpoint{3.704236in}{0.996490in}}%
\pgfpathlineto{\pgfqpoint{3.704392in}{1.004739in}}%
\pgfpathlineto{\pgfqpoint{3.704740in}{0.971742in}}%
\pgfpathlineto{\pgfqpoint{3.705909in}{0.922247in}}%
\pgfpathlineto{\pgfqpoint{3.706080in}{0.946995in}}%
\pgfpathlineto{\pgfqpoint{3.706583in}{0.971742in}}%
\pgfpathlineto{\pgfqpoint{3.707242in}{0.955244in}}%
\pgfpathlineto{\pgfqpoint{3.707923in}{0.946995in}}%
\pgfpathlineto{\pgfqpoint{3.708094in}{0.971742in}}%
\pgfpathlineto{\pgfqpoint{3.708257in}{0.963493in}}%
\pgfpathlineto{\pgfqpoint{3.708405in}{0.955244in}}%
\pgfpathlineto{\pgfqpoint{3.708760in}{0.996490in}}%
\pgfpathlineto{\pgfqpoint{3.709264in}{0.971742in}}%
\pgfpathlineto{\pgfqpoint{3.709589in}{1.004739in}}%
\pgfpathlineto{\pgfqpoint{3.709597in}{0.996490in}}%
\pgfpathlineto{\pgfqpoint{3.710426in}{1.004739in}}%
\pgfpathlineto{\pgfqpoint{3.710093in}{0.979992in}}%
\pgfpathlineto{\pgfqpoint{3.710441in}{0.996490in}}%
\pgfpathlineto{\pgfqpoint{3.711781in}{0.971742in}}%
\pgfpathlineto{\pgfqpoint{3.712788in}{1.021237in}}%
\pgfpathlineto{\pgfqpoint{3.713121in}{0.996490in}}%
\pgfpathlineto{\pgfqpoint{3.714121in}{0.979992in}}%
\pgfpathlineto{\pgfqpoint{3.714128in}{0.996490in}}%
\pgfpathlineto{\pgfqpoint{3.714276in}{1.004739in}}%
\pgfpathlineto{\pgfqpoint{3.714632in}{0.971742in}}%
\pgfpathlineto{\pgfqpoint{3.714795in}{0.971742in}}%
\pgfpathlineto{\pgfqpoint{3.717127in}{0.881002in}}%
\pgfpathlineto{\pgfqpoint{3.717142in}{0.897500in}}%
\pgfpathlineto{\pgfqpoint{3.717638in}{0.905749in}}%
\pgfpathlineto{\pgfqpoint{3.717815in}{0.889251in}}%
\pgfpathlineto{\pgfqpoint{3.723346in}{0.650025in}}%
\pgfpathlineto{\pgfqpoint{3.723346in}{0.658274in}}%
\pgfpathlineto{\pgfqpoint{3.723517in}{0.674773in}}%
\pgfpathlineto{\pgfqpoint{3.724013in}{0.641776in}}%
\pgfpathlineto{\pgfqpoint{3.725027in}{0.625278in}}%
\pgfpathlineto{\pgfqpoint{3.726197in}{0.633527in}}%
\pgfpathlineto{\pgfqpoint{3.726530in}{0.625278in}}%
\pgfpathlineto{\pgfqpoint{3.727204in}{0.625278in}}%
\pgfpathlineto{\pgfqpoint{3.728366in}{0.633527in}}%
\pgfpathlineto{\pgfqpoint{3.728544in}{0.625278in}}%
\pgfpathlineto{\pgfqpoint{3.729048in}{0.625278in}}%
\pgfpathlineto{\pgfqpoint{3.730388in}{0.633527in}}%
\pgfpathlineto{\pgfqpoint{3.730706in}{0.625278in}}%
\pgfpathlineto{\pgfqpoint{3.731395in}{0.625278in}}%
\pgfpathlineto{\pgfqpoint{3.732402in}{0.633527in}}%
\pgfpathlineto{\pgfqpoint{3.732402in}{0.625278in}}%
\pgfpathlineto{\pgfqpoint{3.733557in}{0.633527in}}%
\pgfpathlineto{\pgfqpoint{3.733742in}{0.625278in}}%
\pgfpathlineto{\pgfqpoint{3.734571in}{0.625278in}}%
\pgfpathlineto{\pgfqpoint{3.735748in}{0.633527in}}%
\pgfpathlineto{\pgfqpoint{3.735911in}{0.625278in}}%
\pgfpathlineto{\pgfqpoint{3.736755in}{0.625278in}}%
\pgfpathlineto{\pgfqpoint{3.737762in}{0.633527in}}%
\pgfpathlineto{\pgfqpoint{3.737762in}{0.625278in}}%
\pgfpathlineto{\pgfqpoint{3.738932in}{0.633527in}}%
\pgfpathlineto{\pgfqpoint{3.739102in}{0.625278in}}%
\pgfpathlineto{\pgfqpoint{3.739776in}{0.625278in}}%
\pgfpathlineto{\pgfqpoint{3.741116in}{0.633527in}}%
\pgfpathlineto{\pgfqpoint{3.741116in}{0.625278in}}%
\pgfpathlineto{\pgfqpoint{3.742109in}{0.625278in}}%
\pgfpathlineto{\pgfqpoint{3.743464in}{0.633527in}}%
\pgfpathlineto{\pgfqpoint{3.743464in}{0.625278in}}%
\pgfpathlineto{\pgfqpoint{3.744471in}{0.625278in}}%
\pgfpathlineto{\pgfqpoint{3.745626in}{0.633527in}}%
\pgfpathlineto{\pgfqpoint{3.745811in}{0.625278in}}%
\pgfpathlineto{\pgfqpoint{3.746640in}{0.625278in}}%
\pgfpathlineto{\pgfqpoint{3.747640in}{0.633527in}}%
\pgfpathlineto{\pgfqpoint{3.747640in}{0.625278in}}%
\pgfpathlineto{\pgfqpoint{3.748980in}{0.633527in}}%
\pgfpathlineto{\pgfqpoint{3.748980in}{0.625278in}}%
\pgfpathlineto{\pgfqpoint{3.749987in}{0.625278in}}%
\pgfpathlineto{\pgfqpoint{3.751342in}{0.633527in}}%
\pgfpathlineto{\pgfqpoint{3.752015in}{0.625278in}}%
\pgfpathlineto{\pgfqpoint{3.752349in}{0.625278in}}%
\pgfpathlineto{\pgfqpoint{3.753348in}{0.633527in}}%
\pgfpathlineto{\pgfqpoint{3.753348in}{0.625278in}}%
\pgfpathlineto{\pgfqpoint{3.754696in}{0.633527in}}%
\pgfpathlineto{\pgfqpoint{3.754851in}{0.625278in}}%
\pgfpathlineto{\pgfqpoint{3.755688in}{0.625278in}}%
\pgfpathlineto{\pgfqpoint{3.757043in}{0.633527in}}%
\pgfpathlineto{\pgfqpoint{3.757376in}{0.625278in}}%
\pgfpathlineto{\pgfqpoint{3.758050in}{0.625278in}}%
\pgfpathlineto{\pgfqpoint{3.759197in}{0.633527in}}%
\pgfpathlineto{\pgfqpoint{3.759205in}{0.625278in}}%
\pgfpathlineto{\pgfqpoint{3.760212in}{0.625278in}}%
\pgfpathlineto{\pgfqpoint{3.761219in}{0.633527in}}%
\pgfpathlineto{\pgfqpoint{3.761219in}{0.625278in}}%
\pgfpathlineto{\pgfqpoint{3.762381in}{0.633527in}}%
\pgfpathlineto{\pgfqpoint{3.762559in}{0.625278in}}%
\pgfpathlineto{\pgfqpoint{3.763396in}{0.625278in}}%
\pgfpathlineto{\pgfqpoint{3.764566in}{0.633527in}}%
\pgfpathlineto{\pgfqpoint{3.764728in}{0.625278in}}%
\pgfpathlineto{\pgfqpoint{3.765580in}{0.625278in}}%
\pgfpathlineto{\pgfqpoint{3.766742in}{0.633527in}}%
\pgfpathlineto{\pgfqpoint{3.766750in}{0.625278in}}%
\pgfpathlineto{\pgfqpoint{3.767757in}{0.625278in}}%
\pgfpathlineto{\pgfqpoint{3.768927in}{0.633527in}}%
\pgfpathlineto{\pgfqpoint{3.769089in}{0.625278in}}%
\pgfpathlineto{\pgfqpoint{3.769934in}{0.625278in}}%
\pgfpathlineto{\pgfqpoint{3.770941in}{0.633527in}}%
\pgfpathlineto{\pgfqpoint{3.770941in}{0.625278in}}%
\pgfpathlineto{\pgfqpoint{3.772118in}{0.633527in}}%
\pgfpathlineto{\pgfqpoint{3.772273in}{0.625278in}}%
\pgfpathlineto{\pgfqpoint{3.773117in}{0.625278in}}%
\pgfpathlineto{\pgfqpoint{3.774124in}{0.633527in}}%
\pgfpathlineto{\pgfqpoint{3.774124in}{0.625278in}}%
\pgfpathlineto{\pgfqpoint{3.775302in}{0.633527in}}%
\pgfpathlineto{\pgfqpoint{3.775464in}{0.625278in}}%
\pgfpathlineto{\pgfqpoint{3.776301in}{0.625278in}}%
\pgfpathlineto{\pgfqpoint{3.777471in}{0.633527in}}%
\pgfpathlineto{\pgfqpoint{3.777478in}{0.625278in}}%
\pgfpathlineto{\pgfqpoint{3.778478in}{0.625278in}}%
\pgfpathlineto{\pgfqpoint{3.779655in}{0.633527in}}%
\pgfpathlineto{\pgfqpoint{3.779655in}{0.625278in}}%
\pgfpathlineto{\pgfqpoint{3.780662in}{0.625278in}}%
\pgfpathlineto{\pgfqpoint{3.781832in}{0.633527in}}%
\pgfpathlineto{\pgfqpoint{3.781840in}{0.625278in}}%
\pgfpathlineto{\pgfqpoint{3.782839in}{0.625278in}}%
\pgfpathlineto{\pgfqpoint{3.784009in}{0.633527in}}%
\pgfpathlineto{\pgfqpoint{3.784187in}{0.625278in}}%
\pgfpathlineto{\pgfqpoint{3.785016in}{0.625278in}}%
\pgfpathlineto{\pgfqpoint{3.786186in}{0.633527in}}%
\pgfpathlineto{\pgfqpoint{3.786363in}{0.625278in}}%
\pgfpathlineto{\pgfqpoint{3.787200in}{0.625278in}}%
\pgfpathlineto{\pgfqpoint{3.788318in}{0.633527in}}%
\pgfpathlineto{\pgfqpoint{3.788370in}{0.625278in}}%
\pgfpathlineto{\pgfqpoint{3.789207in}{0.625278in}}%
\pgfpathlineto{\pgfqpoint{3.790377in}{0.633527in}}%
\pgfpathlineto{\pgfqpoint{3.790547in}{0.625278in}}%
\pgfpathlineto{\pgfqpoint{3.791384in}{0.625278in}}%
\pgfpathlineto{\pgfqpoint{3.792390in}{0.633527in}}%
\pgfpathlineto{\pgfqpoint{3.792390in}{0.625278in}}%
\pgfpathlineto{\pgfqpoint{3.793568in}{0.633527in}}%
\pgfpathlineto{\pgfqpoint{3.793731in}{0.625278in}}%
\pgfpathlineto{\pgfqpoint{3.794567in}{0.625278in}}%
\pgfpathlineto{\pgfqpoint{3.795745in}{0.633527in}}%
\pgfpathlineto{\pgfqpoint{3.795915in}{0.625278in}}%
\pgfpathlineto{\pgfqpoint{3.796759in}{0.625278in}}%
\pgfpathlineto{\pgfqpoint{3.797921in}{0.633527in}}%
\pgfpathlineto{\pgfqpoint{3.798099in}{0.625278in}}%
\pgfpathlineto{\pgfqpoint{3.798758in}{0.625278in}}%
\pgfpathlineto{\pgfqpoint{3.799935in}{0.633527in}}%
\pgfpathlineto{\pgfqpoint{3.800098in}{0.625278in}}%
\pgfpathlineto{\pgfqpoint{3.800772in}{0.625278in}}%
\pgfpathlineto{\pgfqpoint{3.802112in}{0.633527in}}%
\pgfpathlineto{\pgfqpoint{3.802445in}{0.625278in}}%
\pgfpathlineto{\pgfqpoint{3.802949in}{0.625278in}}%
\pgfpathlineto{\pgfqpoint{3.804289in}{0.633527in}}%
\pgfpathlineto{\pgfqpoint{3.804459in}{0.625278in}}%
\pgfpathlineto{\pgfqpoint{3.805133in}{0.625278in}}%
\pgfpathlineto{\pgfqpoint{3.806310in}{0.633527in}}%
\pgfpathlineto{\pgfqpoint{3.806473in}{0.625278in}}%
\pgfpathlineto{\pgfqpoint{3.807140in}{0.625278in}}%
\pgfpathlineto{\pgfqpoint{3.808317in}{0.633527in}}%
\pgfpathlineto{\pgfqpoint{3.808480in}{0.625278in}}%
\pgfpathlineto{\pgfqpoint{3.809324in}{0.625278in}}%
\pgfpathlineto{\pgfqpoint{3.810664in}{0.633527in}}%
\pgfpathlineto{\pgfqpoint{3.810664in}{0.625278in}}%
\pgfpathlineto{\pgfqpoint{3.811671in}{0.625278in}}%
\pgfpathlineto{\pgfqpoint{3.812841in}{0.633527in}}%
\pgfpathlineto{\pgfqpoint{3.813011in}{0.625278in}}%
\pgfpathlineto{\pgfqpoint{3.813848in}{0.625278in}}%
\pgfpathlineto{\pgfqpoint{3.815018in}{0.633527in}}%
\pgfpathlineto{\pgfqpoint{3.815018in}{0.625278in}}%
\pgfpathlineto{\pgfqpoint{3.816025in}{0.625278in}}%
\pgfpathlineto{\pgfqpoint{3.817202in}{0.633527in}}%
\pgfpathlineto{\pgfqpoint{3.817365in}{0.625278in}}%
\pgfpathlineto{\pgfqpoint{3.818209in}{0.625278in}}%
\pgfpathlineto{\pgfqpoint{3.819549in}{0.633527in}}%
\pgfpathlineto{\pgfqpoint{3.819549in}{0.625278in}}%
\pgfpathlineto{\pgfqpoint{3.820556in}{0.625278in}}%
\pgfpathlineto{\pgfqpoint{3.821726in}{0.633527in}}%
\pgfpathlineto{\pgfqpoint{3.821896in}{0.625278in}}%
\pgfpathlineto{\pgfqpoint{3.822733in}{0.625278in}}%
\pgfpathlineto{\pgfqpoint{3.823903in}{0.633527in}}%
\pgfpathlineto{\pgfqpoint{3.824073in}{0.625278in}}%
\pgfpathlineto{\pgfqpoint{3.824910in}{0.625278in}}%
\pgfpathlineto{\pgfqpoint{3.826087in}{0.633527in}}%
\pgfpathlineto{\pgfqpoint{3.826087in}{0.625278in}}%
\pgfpathlineto{\pgfqpoint{3.827087in}{0.625278in}}%
\pgfpathlineto{\pgfqpoint{3.828264in}{0.633527in}}%
\pgfpathlineto{\pgfqpoint{3.828434in}{0.625278in}}%
\pgfpathlineto{\pgfqpoint{3.829271in}{0.625278in}}%
\pgfpathlineto{\pgfqpoint{3.830278in}{0.633527in}}%
\pgfpathlineto{\pgfqpoint{3.830278in}{0.625278in}}%
\pgfpathlineto{\pgfqpoint{3.831618in}{0.633527in}}%
\pgfpathlineto{\pgfqpoint{3.831781in}{0.625278in}}%
\pgfpathlineto{\pgfqpoint{3.832625in}{0.625278in}}%
\pgfpathlineto{\pgfqpoint{3.833795in}{0.633527in}}%
\pgfpathlineto{\pgfqpoint{3.833965in}{0.625278in}}%
\pgfpathlineto{\pgfqpoint{3.834631in}{0.625278in}}%
\pgfpathlineto{\pgfqpoint{3.835809in}{0.633527in}}%
\pgfpathlineto{\pgfqpoint{3.835972in}{0.625278in}}%
\pgfpathlineto{\pgfqpoint{3.836816in}{0.625278in}}%
\pgfpathlineto{\pgfqpoint{3.837986in}{0.633527in}}%
\pgfpathlineto{\pgfqpoint{3.838156in}{0.625278in}}%
\pgfpathlineto{\pgfqpoint{3.838992in}{0.625278in}}%
\pgfpathlineto{\pgfqpoint{3.840162in}{0.633527in}}%
\pgfpathlineto{\pgfqpoint{3.840162in}{0.625278in}}%
\pgfpathlineto{\pgfqpoint{3.841169in}{0.625278in}}%
\pgfpathlineto{\pgfqpoint{3.842347in}{0.633527in}}%
\pgfpathlineto{\pgfqpoint{3.842509in}{0.625278in}}%
\pgfpathlineto{\pgfqpoint{3.843354in}{0.625278in}}%
\pgfpathlineto{\pgfqpoint{3.844523in}{0.633527in}}%
\pgfpathlineto{\pgfqpoint{3.844857in}{0.625278in}}%
\pgfpathlineto{\pgfqpoint{3.845530in}{0.625278in}}%
\pgfpathlineto{\pgfqpoint{3.846537in}{0.633527in}}%
\pgfpathlineto{\pgfqpoint{3.846537in}{0.625278in}}%
\pgfpathlineto{\pgfqpoint{3.847707in}{0.633527in}}%
\pgfpathlineto{\pgfqpoint{3.847878in}{0.625278in}}%
\pgfpathlineto{\pgfqpoint{3.848544in}{0.625278in}}%
\pgfpathlineto{\pgfqpoint{3.849721in}{0.633527in}}%
\pgfpathlineto{\pgfqpoint{3.849891in}{0.625278in}}%
\pgfpathlineto{\pgfqpoint{3.850558in}{0.625278in}}%
\pgfpathlineto{\pgfqpoint{3.851898in}{0.633527in}}%
\pgfpathlineto{\pgfqpoint{3.851898in}{0.625278in}}%
\pgfpathlineto{\pgfqpoint{3.852905in}{0.625278in}}%
\pgfpathlineto{\pgfqpoint{3.854082in}{0.633527in}}%
\pgfpathlineto{\pgfqpoint{3.854415in}{0.625278in}}%
\pgfpathlineto{\pgfqpoint{3.855082in}{0.625278in}}%
\pgfpathlineto{\pgfqpoint{3.856259in}{0.633527in}}%
\pgfpathlineto{\pgfqpoint{3.856422in}{0.625278in}}%
\pgfpathlineto{\pgfqpoint{3.857266in}{0.625278in}}%
\pgfpathlineto{\pgfqpoint{3.858606in}{0.633527in}}%
\pgfpathlineto{\pgfqpoint{3.858606in}{0.625278in}}%
\pgfpathlineto{\pgfqpoint{3.859613in}{0.625278in}}%
\pgfpathlineto{\pgfqpoint{3.860783in}{0.633527in}}%
\pgfpathlineto{\pgfqpoint{3.860953in}{0.625278in}}%
\pgfpathlineto{\pgfqpoint{3.861620in}{0.625278in}}%
\pgfpathlineto{\pgfqpoint{3.862797in}{0.633527in}}%
\pgfpathlineto{\pgfqpoint{3.862960in}{0.625278in}}%
\pgfpathlineto{\pgfqpoint{3.863804in}{0.625278in}}%
\pgfpathlineto{\pgfqpoint{3.864974in}{0.633527in}}%
\pgfpathlineto{\pgfqpoint{3.864974in}{0.625278in}}%
\pgfpathlineto{\pgfqpoint{3.865981in}{0.625278in}}%
\pgfpathlineto{\pgfqpoint{3.867321in}{0.633527in}}%
\pgfpathlineto{\pgfqpoint{3.867491in}{0.625278in}}%
\pgfpathlineto{\pgfqpoint{3.868158in}{0.625278in}}%
\pgfpathlineto{\pgfqpoint{3.869335in}{0.633527in}}%
\pgfpathlineto{\pgfqpoint{3.869335in}{0.625278in}}%
\pgfpathlineto{\pgfqpoint{3.870342in}{0.625278in}}%
\pgfpathlineto{\pgfqpoint{3.871512in}{0.633527in}}%
\pgfpathlineto{\pgfqpoint{3.871682in}{0.625278in}}%
\pgfpathlineto{\pgfqpoint{3.872519in}{0.625278in}}%
\pgfpathlineto{\pgfqpoint{3.873689in}{0.633527in}}%
\pgfpathlineto{\pgfqpoint{3.873859in}{0.625278in}}%
\pgfpathlineto{\pgfqpoint{3.874696in}{0.625278in}}%
\pgfpathlineto{\pgfqpoint{3.875873in}{0.633527in}}%
\pgfpathlineto{\pgfqpoint{3.876036in}{0.625278in}}%
\pgfpathlineto{\pgfqpoint{3.876880in}{0.625278in}}%
\pgfpathlineto{\pgfqpoint{3.878050in}{0.633527in}}%
\pgfpathlineto{\pgfqpoint{3.878220in}{0.625278in}}%
\pgfpathlineto{\pgfqpoint{3.878886in}{0.625278in}}%
\pgfpathlineto{\pgfqpoint{3.880064in}{0.633527in}}%
\pgfpathlineto{\pgfqpoint{3.880064in}{0.625278in}}%
\pgfpathlineto{\pgfqpoint{3.880900in}{0.625278in}}%
\pgfpathlineto{\pgfqpoint{3.882070in}{0.633527in}}%
\pgfpathlineto{\pgfqpoint{3.882240in}{0.625278in}}%
\pgfpathlineto{\pgfqpoint{3.883077in}{0.625278in}}%
\pgfpathlineto{\pgfqpoint{3.884254in}{0.633527in}}%
\pgfpathlineto{\pgfqpoint{3.884417in}{0.625278in}}%
\pgfpathlineto{\pgfqpoint{3.885261in}{0.625278in}}%
\pgfpathlineto{\pgfqpoint{3.886601in}{0.633527in}}%
\pgfpathlineto{\pgfqpoint{3.886601in}{0.625278in}}%
\pgfpathlineto{\pgfqpoint{3.887608in}{0.625278in}}%
\pgfpathlineto{\pgfqpoint{3.888778in}{0.633527in}}%
\pgfpathlineto{\pgfqpoint{3.888949in}{0.625278in}}%
\pgfpathlineto{\pgfqpoint{3.889785in}{0.625278in}}%
\pgfpathlineto{\pgfqpoint{3.890955in}{0.633527in}}%
\pgfpathlineto{\pgfqpoint{3.891296in}{0.625278in}}%
\pgfpathlineto{\pgfqpoint{3.891962in}{0.625278in}}%
\pgfpathlineto{\pgfqpoint{3.892969in}{0.633527in}}%
\pgfpathlineto{\pgfqpoint{3.892969in}{0.625278in}}%
\pgfpathlineto{\pgfqpoint{3.894139in}{0.633527in}}%
\pgfpathlineto{\pgfqpoint{3.894309in}{0.625278in}}%
\pgfpathlineto{\pgfqpoint{3.895146in}{0.625278in}}%
\pgfpathlineto{\pgfqpoint{3.896323in}{0.633527in}}%
\pgfpathlineto{\pgfqpoint{3.896486in}{0.625278in}}%
\pgfpathlineto{\pgfqpoint{3.897330in}{0.625278in}}%
\pgfpathlineto{\pgfqpoint{3.898330in}{0.633527in}}%
\pgfpathlineto{\pgfqpoint{3.898330in}{0.625278in}}%
\pgfpathlineto{\pgfqpoint{3.899677in}{0.633527in}}%
\pgfpathlineto{\pgfqpoint{3.899677in}{0.625278in}}%
\pgfpathlineto{\pgfqpoint{3.900677in}{0.625278in}}%
\pgfpathlineto{\pgfqpoint{3.901854in}{0.633527in}}%
\pgfpathlineto{\pgfqpoint{3.902024in}{0.625278in}}%
\pgfpathlineto{\pgfqpoint{3.902861in}{0.625278in}}%
\pgfpathlineto{\pgfqpoint{3.904031in}{0.633527in}}%
\pgfpathlineto{\pgfqpoint{3.904031in}{0.625278in}}%
\pgfpathlineto{\pgfqpoint{3.904705in}{0.625278in}}%
\pgfpathlineto{\pgfqpoint{3.905875in}{0.633527in}}%
\pgfpathlineto{\pgfqpoint{3.906045in}{0.625278in}}%
\pgfpathlineto{\pgfqpoint{3.906882in}{0.625278in}}%
\pgfpathlineto{\pgfqpoint{3.908059in}{0.633527in}}%
\pgfpathlineto{\pgfqpoint{3.908392in}{0.625278in}}%
\pgfpathlineto{\pgfqpoint{3.909058in}{0.625278in}}%
\pgfpathlineto{\pgfqpoint{3.910236in}{0.633527in}}%
\pgfpathlineto{\pgfqpoint{3.910406in}{0.625278in}}%
\pgfpathlineto{\pgfqpoint{3.911243in}{0.625278in}}%
\pgfpathlineto{\pgfqpoint{3.912412in}{0.633527in}}%
\pgfpathlineto{\pgfqpoint{3.912753in}{0.625278in}}%
\pgfpathlineto{\pgfqpoint{3.913419in}{0.625278in}}%
\pgfpathlineto{\pgfqpoint{3.914597in}{0.633527in}}%
\pgfpathlineto{\pgfqpoint{3.914760in}{0.625278in}}%
\pgfpathlineto{\pgfqpoint{3.915596in}{0.625278in}}%
\pgfpathlineto{\pgfqpoint{3.916774in}{0.633527in}}%
\pgfpathlineto{\pgfqpoint{3.916944in}{0.625278in}}%
\pgfpathlineto{\pgfqpoint{3.917781in}{0.625278in}}%
\pgfpathlineto{\pgfqpoint{3.918950in}{0.633527in}}%
\pgfpathlineto{\pgfqpoint{3.918950in}{0.625278in}}%
\pgfpathlineto{\pgfqpoint{3.919957in}{0.625278in}}%
\pgfpathlineto{\pgfqpoint{3.921135in}{0.633527in}}%
\pgfpathlineto{\pgfqpoint{3.921135in}{0.625278in}}%
\pgfpathlineto{\pgfqpoint{3.922134in}{0.625278in}}%
\pgfpathlineto{\pgfqpoint{3.923311in}{0.633527in}}%
\pgfpathlineto{\pgfqpoint{3.923474in}{0.625278in}}%
\pgfpathlineto{\pgfqpoint{3.924318in}{0.625278in}}%
\pgfpathlineto{\pgfqpoint{3.925488in}{0.633527in}}%
\pgfpathlineto{\pgfqpoint{3.925821in}{0.625278in}}%
\pgfpathlineto{\pgfqpoint{3.926495in}{0.625278in}}%
\pgfpathlineto{\pgfqpoint{3.927665in}{0.633527in}}%
\pgfpathlineto{\pgfqpoint{3.927835in}{0.625278in}}%
\pgfpathlineto{\pgfqpoint{3.928672in}{0.625278in}}%
\pgfpathlineto{\pgfqpoint{3.929849in}{0.633527in}}%
\pgfpathlineto{\pgfqpoint{3.930012in}{0.625278in}}%
\pgfpathlineto{\pgfqpoint{3.930856in}{0.625278in}}%
\pgfpathlineto{\pgfqpoint{3.932026in}{0.633527in}}%
\pgfpathlineto{\pgfqpoint{3.932026in}{0.625278in}}%
\pgfpathlineto{\pgfqpoint{3.933033in}{0.625278in}}%
\pgfpathlineto{\pgfqpoint{3.934373in}{0.633527in}}%
\pgfpathlineto{\pgfqpoint{3.934373in}{0.625278in}}%
\pgfpathlineto{\pgfqpoint{3.935380in}{0.625278in}}%
\pgfpathlineto{\pgfqpoint{3.936387in}{0.633527in}}%
\pgfpathlineto{\pgfqpoint{3.936387in}{0.625278in}}%
\pgfpathlineto{\pgfqpoint{3.937557in}{0.633527in}}%
\pgfpathlineto{\pgfqpoint{3.937727in}{0.625278in}}%
\pgfpathlineto{\pgfqpoint{3.938564in}{0.625278in}}%
\pgfpathlineto{\pgfqpoint{3.939904in}{0.633527in}}%
\pgfpathlineto{\pgfqpoint{3.939904in}{0.625278in}}%
\pgfpathlineto{\pgfqpoint{3.940911in}{0.625278in}}%
\pgfpathlineto{\pgfqpoint{3.942088in}{0.633527in}}%
\pgfpathlineto{\pgfqpoint{3.942140in}{0.625278in}}%
\pgfpathlineto{\pgfqpoint{3.943110in}{0.625278in}}%
\pgfpathlineto{\pgfqpoint{3.944095in}{0.633527in}}%
\pgfpathlineto{\pgfqpoint{3.944117in}{0.625278in}}%
\pgfpathlineto{\pgfqpoint{3.945272in}{0.633527in}}%
\pgfpathlineto{\pgfqpoint{3.945457in}{0.650025in}}%
\pgfpathlineto{\pgfqpoint{3.946279in}{0.625278in}}%
\pgfpathlineto{\pgfqpoint{3.947131in}{0.650025in}}%
\pgfpathlineto{\pgfqpoint{3.947449in}{0.633527in}}%
\pgfpathlineto{\pgfqpoint{3.948471in}{0.625278in}}%
\pgfpathlineto{\pgfqpoint{3.950129in}{0.658274in}}%
\pgfpathlineto{\pgfqpoint{3.951159in}{0.625278in}}%
\pgfpathlineto{\pgfqpoint{3.953313in}{0.658274in}}%
\pgfpathlineto{\pgfqpoint{3.953506in}{0.650025in}}%
\pgfpathlineto{\pgfqpoint{3.954157in}{0.707769in}}%
\pgfpathlineto{\pgfqpoint{3.954172in}{0.699520in}}%
\pgfpathlineto{\pgfqpoint{3.954491in}{0.707769in}}%
\pgfpathlineto{\pgfqpoint{3.954824in}{0.683022in}}%
\pgfpathlineto{\pgfqpoint{3.955009in}{0.691271in}}%
\pgfpathlineto{\pgfqpoint{3.955157in}{0.683022in}}%
\pgfpathlineto{\pgfqpoint{3.955512in}{0.724268in}}%
\pgfpathlineto{\pgfqpoint{3.955831in}{0.732517in}}%
\pgfpathlineto{\pgfqpoint{3.956164in}{0.707769in}}%
\pgfpathlineto{\pgfqpoint{3.956690in}{0.724268in}}%
\pgfpathlineto{\pgfqpoint{3.957341in}{0.683022in}}%
\pgfpathlineto{\pgfqpoint{3.957356in}{0.699520in}}%
\pgfpathlineto{\pgfqpoint{3.958363in}{0.699520in}}%
\pgfpathlineto{\pgfqpoint{3.959370in}{0.749015in}}%
\pgfpathlineto{\pgfqpoint{3.960021in}{0.732517in}}%
\pgfpathlineto{\pgfqpoint{3.961028in}{0.707769in}}%
\pgfpathlineto{\pgfqpoint{3.961043in}{0.724268in}}%
\pgfpathlineto{\pgfqpoint{3.961695in}{0.732517in}}%
\pgfpathlineto{\pgfqpoint{3.961887in}{0.716019in}}%
\pgfpathlineto{\pgfqpoint{3.962391in}{0.699520in}}%
\pgfpathlineto{\pgfqpoint{3.962702in}{0.732517in}}%
\pgfpathlineto{\pgfqpoint{3.962724in}{0.724268in}}%
\pgfpathlineto{\pgfqpoint{3.964397in}{0.773763in}}%
\pgfpathlineto{\pgfqpoint{3.964568in}{0.765513in}}%
\pgfpathlineto{\pgfqpoint{3.965404in}{0.749015in}}%
\pgfpathlineto{\pgfqpoint{3.965552in}{0.765513in}}%
\pgfpathlineto{\pgfqpoint{3.966559in}{0.782012in}}%
\pgfpathlineto{\pgfqpoint{3.966582in}{0.773763in}}%
\pgfpathlineto{\pgfqpoint{3.967900in}{0.806759in}}%
\pgfpathlineto{\pgfqpoint{3.967922in}{0.798510in}}%
\pgfpathlineto{\pgfqpoint{3.968573in}{0.782012in}}%
\pgfpathlineto{\pgfqpoint{3.968736in}{0.806759in}}%
\pgfpathlineto{\pgfqpoint{3.968758in}{0.798510in}}%
\pgfpathlineto{\pgfqpoint{3.970417in}{0.831507in}}%
\pgfpathlineto{\pgfqpoint{3.971106in}{0.798510in}}%
\pgfpathlineto{\pgfqpoint{3.971439in}{0.823258in}}%
\pgfpathlineto{\pgfqpoint{3.971609in}{0.848005in}}%
\pgfpathlineto{\pgfqpoint{3.972113in}{0.815008in}}%
\pgfpathlineto{\pgfqpoint{3.972594in}{0.782012in}}%
\pgfpathlineto{\pgfqpoint{3.972949in}{0.823258in}}%
\pgfpathlineto{\pgfqpoint{3.974119in}{0.848005in}}%
\pgfpathlineto{\pgfqpoint{3.973453in}{0.798510in}}%
\pgfpathlineto{\pgfqpoint{3.974289in}{0.839756in}}%
\pgfpathlineto{\pgfqpoint{3.975111in}{0.806759in}}%
\pgfpathlineto{\pgfqpoint{3.975444in}{0.831507in}}%
\pgfpathlineto{\pgfqpoint{3.975615in}{0.856254in}}%
\pgfpathlineto{\pgfqpoint{3.976133in}{0.823258in}}%
\pgfpathlineto{\pgfqpoint{3.976466in}{0.823258in}}%
\pgfpathlineto{\pgfqpoint{3.976636in}{0.848005in}}%
\pgfpathlineto{\pgfqpoint{3.977140in}{0.815008in}}%
\pgfpathlineto{\pgfqpoint{3.977954in}{0.806759in}}%
\pgfpathlineto{\pgfqpoint{3.977303in}{0.823258in}}%
\pgfpathlineto{\pgfqpoint{3.977977in}{0.823258in}}%
\pgfpathlineto{\pgfqpoint{3.979635in}{0.856254in}}%
\pgfpathlineto{\pgfqpoint{3.980472in}{0.831507in}}%
\pgfpathlineto{\pgfqpoint{3.980657in}{0.839756in}}%
\pgfpathlineto{\pgfqpoint{3.981479in}{0.831507in}}%
\pgfpathlineto{\pgfqpoint{3.980827in}{0.848005in}}%
\pgfpathlineto{\pgfqpoint{3.981479in}{0.839756in}}%
\pgfpathlineto{\pgfqpoint{3.981494in}{0.872753in}}%
\pgfpathlineto{\pgfqpoint{3.982501in}{0.848005in}}%
\pgfpathlineto{\pgfqpoint{3.983152in}{0.881002in}}%
\pgfpathlineto{\pgfqpoint{3.983508in}{0.848005in}}%
\pgfpathlineto{\pgfqpoint{3.984011in}{0.823258in}}%
\pgfpathlineto{\pgfqpoint{3.984329in}{0.856254in}}%
\pgfpathlineto{\pgfqpoint{3.984344in}{0.848005in}}%
\pgfpathlineto{\pgfqpoint{3.984515in}{0.872753in}}%
\pgfpathlineto{\pgfqpoint{3.985499in}{0.856254in}}%
\pgfpathlineto{\pgfqpoint{3.986528in}{0.848005in}}%
\pgfpathlineto{\pgfqpoint{3.988187in}{0.905749in}}%
\pgfpathlineto{\pgfqpoint{3.988372in}{0.889251in}}%
\pgfpathlineto{\pgfqpoint{3.989187in}{0.881002in}}%
\pgfpathlineto{\pgfqpoint{3.988535in}{0.897500in}}%
\pgfpathlineto{\pgfqpoint{3.989209in}{0.897500in}}%
\pgfpathlineto{\pgfqpoint{3.990364in}{0.905749in}}%
\pgfpathlineto{\pgfqpoint{3.991386in}{0.889251in}}%
\pgfpathlineto{\pgfqpoint{3.992207in}{0.881002in}}%
\pgfpathlineto{\pgfqpoint{3.991556in}{0.897500in}}%
\pgfpathlineto{\pgfqpoint{3.992222in}{0.897500in}}%
\pgfpathlineto{\pgfqpoint{3.993044in}{0.905749in}}%
\pgfpathlineto{\pgfqpoint{3.993066in}{0.897500in}}%
\pgfpathlineto{\pgfqpoint{3.993377in}{0.881002in}}%
\pgfpathlineto{\pgfqpoint{3.993718in}{0.913998in}}%
\pgfpathlineto{\pgfqpoint{3.994384in}{0.930497in}}%
\pgfpathlineto{\pgfqpoint{3.994718in}{0.905749in}}%
\pgfpathlineto{\pgfqpoint{3.994740in}{0.922247in}}%
\pgfpathlineto{\pgfqpoint{3.996065in}{0.881002in}}%
\pgfpathlineto{\pgfqpoint{3.996080in}{0.897500in}}%
\pgfpathlineto{\pgfqpoint{3.997257in}{0.946995in}}%
\pgfpathlineto{\pgfqpoint{3.997738in}{0.930497in}}%
\pgfpathlineto{\pgfqpoint{3.998760in}{0.922247in}}%
\pgfpathlineto{\pgfqpoint{4.000759in}{0.979992in}}%
\pgfpathlineto{\pgfqpoint{4.001759in}{0.955244in}}%
\pgfpathlineto{\pgfqpoint{4.001781in}{0.971742in}}%
\pgfpathlineto{\pgfqpoint{4.002099in}{0.955244in}}%
\pgfpathlineto{\pgfqpoint{4.002433in}{0.988241in}}%
\pgfpathlineto{\pgfqpoint{4.002766in}{1.004739in}}%
\pgfpathlineto{\pgfqpoint{4.003121in}{0.971742in}}%
\pgfpathlineto{\pgfqpoint{4.003491in}{0.979992in}}%
\pgfpathlineto{\pgfqpoint{4.004276in}{1.029487in}}%
\pgfpathlineto{\pgfqpoint{4.004795in}{1.012988in}}%
\pgfpathlineto{\pgfqpoint{4.004950in}{1.004739in}}%
\pgfpathlineto{\pgfqpoint{4.005283in}{1.037736in}}%
\pgfpathlineto{\pgfqpoint{4.005616in}{1.054234in}}%
\pgfpathlineto{\pgfqpoint{4.005972in}{1.021237in}}%
\pgfpathlineto{\pgfqpoint{4.006305in}{1.021237in}}%
\pgfpathlineto{\pgfqpoint{4.007793in}{1.103729in}}%
\pgfpathlineto{\pgfqpoint{4.008297in}{1.078981in}}%
\pgfpathlineto{\pgfqpoint{4.009489in}{1.045985in}}%
\pgfpathlineto{\pgfqpoint{4.010140in}{1.029487in}}%
\pgfpathlineto{\pgfqpoint{4.010311in}{1.054234in}}%
\pgfpathlineto{\pgfqpoint{4.010333in}{1.045985in}}%
\pgfpathlineto{\pgfqpoint{4.011503in}{1.070732in}}%
\pgfpathlineto{\pgfqpoint{4.011673in}{1.062483in}}%
\pgfpathlineto{\pgfqpoint{4.012488in}{1.054234in}}%
\pgfpathlineto{\pgfqpoint{4.011836in}{1.070732in}}%
\pgfpathlineto{\pgfqpoint{4.012488in}{1.062483in}}%
\pgfpathlineto{\pgfqpoint{4.013517in}{1.120227in}}%
\pgfpathlineto{\pgfqpoint{4.014857in}{1.095480in}}%
\pgfpathlineto{\pgfqpoint{4.015338in}{1.128476in}}%
\pgfpathlineto{\pgfqpoint{4.016012in}{1.103729in}}%
\pgfpathlineto{\pgfqpoint{4.016197in}{1.120227in}}%
\pgfpathlineto{\pgfqpoint{4.017034in}{1.087231in}}%
\pgfpathlineto{\pgfqpoint{4.017522in}{1.054234in}}%
\pgfpathlineto{\pgfqpoint{4.018019in}{1.087231in}}%
\pgfpathlineto{\pgfqpoint{4.019025in}{1.103729in}}%
\pgfpathlineto{\pgfqpoint{4.019048in}{1.095480in}}%
\pgfpathlineto{\pgfqpoint{4.020203in}{1.103729in}}%
\pgfpathlineto{\pgfqpoint{4.020869in}{1.078981in}}%
\pgfpathlineto{\pgfqpoint{4.021225in}{1.095480in}}%
\pgfpathlineto{\pgfqpoint{4.021543in}{1.103729in}}%
\pgfpathlineto{\pgfqpoint{4.021876in}{1.078981in}}%
\pgfpathlineto{\pgfqpoint{4.022061in}{1.087231in}}%
\pgfpathlineto{\pgfqpoint{4.022209in}{1.078981in}}%
\pgfpathlineto{\pgfqpoint{4.022380in}{1.103729in}}%
\pgfpathlineto{\pgfqpoint{4.022735in}{1.095480in}}%
\pgfpathlineto{\pgfqpoint{4.022883in}{1.103729in}}%
\pgfpathlineto{\pgfqpoint{4.023216in}{1.078981in}}%
\pgfpathlineto{\pgfqpoint{4.023401in}{1.087231in}}%
\pgfpathlineto{\pgfqpoint{4.023557in}{1.078981in}}%
\pgfpathlineto{\pgfqpoint{4.023572in}{1.120227in}}%
\pgfpathlineto{\pgfqpoint{4.023890in}{1.111978in}}%
\pgfpathlineto{\pgfqpoint{4.024223in}{1.103729in}}%
\pgfpathlineto{\pgfqpoint{4.024912in}{1.144975in}}%
\pgfpathlineto{\pgfqpoint{4.025119in}{1.128476in}}%
\pgfpathlineto{\pgfqpoint{4.025563in}{1.161473in}}%
\pgfpathlineto{\pgfqpoint{4.026067in}{1.177971in}}%
\pgfpathlineto{\pgfqpoint{4.026400in}{1.153224in}}%
\pgfpathlineto{\pgfqpoint{4.026593in}{1.161473in}}%
\pgfpathlineto{\pgfqpoint{4.027762in}{1.120227in}}%
\pgfpathlineto{\pgfqpoint{4.026755in}{1.169722in}}%
\pgfpathlineto{\pgfqpoint{4.027933in}{1.144975in}}%
\pgfpathlineto{\pgfqpoint{4.029273in}{1.194470in}}%
\pgfpathlineto{\pgfqpoint{4.029591in}{1.177971in}}%
\pgfpathlineto{\pgfqpoint{4.029776in}{1.169722in}}%
\pgfpathlineto{\pgfqpoint{4.030095in}{1.202719in}}%
\pgfpathlineto{\pgfqpoint{4.030613in}{1.194470in}}%
\pgfpathlineto{\pgfqpoint{4.031768in}{1.202719in}}%
\pgfpathlineto{\pgfqpoint{4.032790in}{1.169722in}}%
\pgfpathlineto{\pgfqpoint{4.034782in}{1.227466in}}%
\pgfpathlineto{\pgfqpoint{4.035137in}{1.243965in}}%
\pgfpathlineto{\pgfqpoint{4.035811in}{1.219217in}}%
\pgfpathlineto{\pgfqpoint{4.036966in}{1.227466in}}%
\pgfpathlineto{\pgfqpoint{4.037632in}{1.202719in}}%
\pgfpathlineto{\pgfqpoint{4.037988in}{1.268712in}}%
\pgfpathlineto{\pgfqpoint{4.038306in}{1.276961in}}%
\pgfpathlineto{\pgfqpoint{4.038491in}{1.260463in}}%
\pgfpathlineto{\pgfqpoint{4.039816in}{1.202719in}}%
\pgfpathlineto{\pgfqpoint{4.039831in}{1.219217in}}%
\pgfpathlineto{\pgfqpoint{4.040320in}{1.252214in}}%
\pgfpathlineto{\pgfqpoint{4.041171in}{1.235715in}}%
\pgfpathlineto{\pgfqpoint{4.041512in}{1.219217in}}%
\pgfpathlineto{\pgfqpoint{4.041675in}{1.243965in}}%
\pgfpathlineto{\pgfqpoint{4.042326in}{1.227466in}}%
\pgfpathlineto{\pgfqpoint{4.043163in}{1.276961in}}%
\pgfpathlineto{\pgfqpoint{4.043356in}{1.268712in}}%
\pgfpathlineto{\pgfqpoint{4.044696in}{1.243965in}}%
\pgfpathlineto{\pgfqpoint{4.045510in}{1.252214in}}%
\pgfpathlineto{\pgfqpoint{4.045532in}{1.243965in}}%
\pgfpathlineto{\pgfqpoint{4.046184in}{1.227466in}}%
\pgfpathlineto{\pgfqpoint{4.046354in}{1.252214in}}%
\pgfpathlineto{\pgfqpoint{4.046369in}{1.243965in}}%
\pgfpathlineto{\pgfqpoint{4.047080in}{1.276961in}}%
\pgfpathlineto{\pgfqpoint{4.047524in}{1.252214in}}%
\pgfpathlineto{\pgfqpoint{4.048028in}{1.276961in}}%
\pgfpathlineto{\pgfqpoint{4.048546in}{1.243965in}}%
\pgfpathlineto{\pgfqpoint{4.049035in}{1.252214in}}%
\pgfpathlineto{\pgfqpoint{4.049220in}{1.235715in}}%
\pgfpathlineto{\pgfqpoint{4.050893in}{1.194470in}}%
\pgfpathlineto{\pgfqpoint{4.051212in}{1.202719in}}%
\pgfpathlineto{\pgfqpoint{4.051545in}{1.177971in}}%
\pgfpathlineto{\pgfqpoint{4.051737in}{1.186221in}}%
\pgfpathlineto{\pgfqpoint{4.053240in}{1.144975in}}%
\pgfpathlineto{\pgfqpoint{4.053388in}{1.153224in}}%
\pgfpathlineto{\pgfqpoint{4.053411in}{1.169722in}}%
\pgfpathlineto{\pgfqpoint{4.054418in}{1.120227in}}%
\pgfpathlineto{\pgfqpoint{4.055069in}{1.128476in}}%
\pgfpathlineto{\pgfqpoint{4.055254in}{1.111978in}}%
\pgfpathlineto{\pgfqpoint{4.056742in}{1.078981in}}%
\pgfpathlineto{\pgfqpoint{4.055425in}{1.120227in}}%
\pgfpathlineto{\pgfqpoint{4.056742in}{1.087231in}}%
\pgfpathlineto{\pgfqpoint{4.056913in}{1.103729in}}%
\pgfpathlineto{\pgfqpoint{4.057268in}{1.070732in}}%
\pgfpathlineto{\pgfqpoint{4.057772in}{1.070732in}}%
\pgfpathlineto{\pgfqpoint{4.061459in}{0.922247in}}%
\pgfpathlineto{\pgfqpoint{4.061607in}{0.930497in}}%
\pgfpathlineto{\pgfqpoint{4.061770in}{0.955244in}}%
\pgfpathlineto{\pgfqpoint{4.062296in}{0.897500in}}%
\pgfpathlineto{\pgfqpoint{4.062629in}{0.897500in}}%
\pgfpathlineto{\pgfqpoint{4.062947in}{0.905749in}}%
\pgfpathlineto{\pgfqpoint{4.063280in}{0.881002in}}%
\pgfpathlineto{\pgfqpoint{4.063465in}{0.889251in}}%
\pgfpathlineto{\pgfqpoint{4.064961in}{0.831507in}}%
\pgfpathlineto{\pgfqpoint{4.063636in}{0.897500in}}%
\pgfpathlineto{\pgfqpoint{4.064969in}{0.864503in}}%
\pgfpathlineto{\pgfqpoint{4.065650in}{0.872753in}}%
\pgfpathlineto{\pgfqpoint{4.065983in}{0.848005in}}%
\pgfpathlineto{\pgfqpoint{4.066990in}{0.872753in}}%
\pgfpathlineto{\pgfqpoint{4.067160in}{0.864503in}}%
\pgfpathlineto{\pgfqpoint{4.068308in}{0.806759in}}%
\pgfpathlineto{\pgfqpoint{4.068648in}{0.839756in}}%
\pgfpathlineto{\pgfqpoint{4.069167in}{0.897500in}}%
\pgfpathlineto{\pgfqpoint{4.069670in}{0.864503in}}%
\pgfpathlineto{\pgfqpoint{4.070344in}{0.848005in}}%
\pgfpathlineto{\pgfqpoint{4.070507in}{0.872753in}}%
\pgfpathlineto{\pgfqpoint{4.070833in}{0.856254in}}%
\pgfpathlineto{\pgfqpoint{4.071662in}{0.905749in}}%
\pgfpathlineto{\pgfqpoint{4.071847in}{0.889251in}}%
\pgfpathlineto{\pgfqpoint{4.073357in}{0.848005in}}%
\pgfpathlineto{\pgfqpoint{4.073506in}{0.856254in}}%
\pgfpathlineto{\pgfqpoint{4.074024in}{0.881002in}}%
\pgfpathlineto{\pgfqpoint{4.074520in}{0.831507in}}%
\pgfpathlineto{\pgfqpoint{4.074535in}{0.848005in}}%
\pgfpathlineto{\pgfqpoint{4.075860in}{0.806759in}}%
\pgfpathlineto{\pgfqpoint{4.075912in}{0.831507in}}%
\pgfpathlineto{\pgfqpoint{4.076038in}{0.848005in}}%
\pgfpathlineto{\pgfqpoint{4.076378in}{0.815008in}}%
\pgfpathlineto{\pgfqpoint{4.077719in}{0.707769in}}%
\pgfpathlineto{\pgfqpoint{4.078052in}{0.716019in}}%
\pgfpathlineto{\pgfqpoint{4.079547in}{0.683022in}}%
\pgfpathlineto{\pgfqpoint{4.078711in}{0.732517in}}%
\pgfpathlineto{\pgfqpoint{4.079547in}{0.691271in}}%
\pgfpathlineto{\pgfqpoint{4.079895in}{0.674773in}}%
\pgfpathlineto{\pgfqpoint{4.080569in}{0.699520in}}%
\pgfpathlineto{\pgfqpoint{4.081739in}{0.749015in}}%
\pgfpathlineto{\pgfqpoint{4.081065in}{0.683022in}}%
\pgfpathlineto{\pgfqpoint{4.082072in}{0.724268in}}%
\pgfpathlineto{\pgfqpoint{4.082405in}{0.707769in}}%
\pgfpathlineto{\pgfqpoint{4.082739in}{0.740766in}}%
\pgfpathlineto{\pgfqpoint{4.083420in}{0.823258in}}%
\pgfpathlineto{\pgfqpoint{4.083753in}{0.798510in}}%
\pgfpathlineto{\pgfqpoint{4.084575in}{0.806759in}}%
\pgfpathlineto{\pgfqpoint{4.085093in}{0.773763in}}%
\pgfpathlineto{\pgfqpoint{4.085589in}{0.782012in}}%
\pgfpathlineto{\pgfqpoint{4.085767in}{0.765513in}}%
\pgfpathlineto{\pgfqpoint{4.087440in}{0.724268in}}%
\pgfpathlineto{\pgfqpoint{4.087603in}{0.732517in}}%
\pgfpathlineto{\pgfqpoint{4.087929in}{0.707769in}}%
\pgfpathlineto{\pgfqpoint{4.088114in}{0.716019in}}%
\pgfpathlineto{\pgfqpoint{4.088277in}{0.724268in}}%
\pgfpathlineto{\pgfqpoint{4.089284in}{0.699520in}}%
\pgfpathlineto{\pgfqpoint{4.090957in}{0.732517in}}%
\pgfpathlineto{\pgfqpoint{4.091298in}{0.707769in}}%
\pgfpathlineto{\pgfqpoint{4.091801in}{0.773763in}}%
\pgfpathlineto{\pgfqpoint{4.091972in}{0.757264in}}%
\pgfpathlineto{\pgfqpoint{4.093127in}{0.782012in}}%
\pgfpathlineto{\pgfqpoint{4.093141in}{0.773763in}}%
\pgfpathlineto{\pgfqpoint{4.094815in}{0.699520in}}%
\pgfpathlineto{\pgfqpoint{4.095311in}{0.716019in}}%
\pgfpathlineto{\pgfqpoint{4.096162in}{0.732517in}}%
\pgfpathlineto{\pgfqpoint{4.095651in}{0.699520in}}%
\pgfpathlineto{\pgfqpoint{4.096325in}{0.724268in}}%
\pgfpathlineto{\pgfqpoint{4.098843in}{0.633527in}}%
\pgfpathlineto{\pgfqpoint{4.098843in}{0.641776in}}%
\pgfpathlineto{\pgfqpoint{4.099006in}{0.650025in}}%
\pgfpathlineto{\pgfqpoint{4.099339in}{0.625278in}}%
\pgfpathlineto{\pgfqpoint{4.099679in}{0.633527in}}%
\pgfpathlineto{\pgfqpoint{4.099842in}{0.625278in}}%
\pgfpathlineto{\pgfqpoint{4.100686in}{0.625278in}}%
\pgfpathlineto{\pgfqpoint{4.101679in}{0.633527in}}%
\pgfpathlineto{\pgfqpoint{4.101679in}{0.625278in}}%
\pgfpathlineto{\pgfqpoint{4.102863in}{0.633527in}}%
\pgfpathlineto{\pgfqpoint{4.103033in}{0.625278in}}%
\pgfpathlineto{\pgfqpoint{4.103700in}{0.625278in}}%
\pgfpathlineto{\pgfqpoint{4.104877in}{0.633527in}}%
\pgfpathlineto{\pgfqpoint{4.105040in}{0.625278in}}%
\pgfpathlineto{\pgfqpoint{4.105884in}{0.625278in}}%
\pgfpathlineto{\pgfqpoint{4.107054in}{0.633527in}}%
\pgfpathlineto{\pgfqpoint{4.107054in}{0.625278in}}%
\pgfpathlineto{\pgfqpoint{4.108061in}{0.625278in}}%
\pgfpathlineto{\pgfqpoint{4.109386in}{0.633527in}}%
\pgfpathlineto{\pgfqpoint{4.109557in}{0.625278in}}%
\pgfpathlineto{\pgfqpoint{4.110223in}{0.625278in}}%
\pgfpathlineto{\pgfqpoint{4.111400in}{0.633527in}}%
\pgfpathlineto{\pgfqpoint{4.111748in}{0.625278in}}%
\pgfpathlineto{\pgfqpoint{4.112252in}{0.625278in}}%
\pgfpathlineto{\pgfqpoint{4.113422in}{0.633527in}}%
\pgfpathlineto{\pgfqpoint{4.113592in}{0.625278in}}%
\pgfpathlineto{\pgfqpoint{4.114421in}{0.625278in}}%
\pgfpathlineto{\pgfqpoint{4.115606in}{0.633527in}}%
\pgfpathlineto{\pgfqpoint{4.115761in}{0.625278in}}%
\pgfpathlineto{\pgfqpoint{4.116598in}{0.625278in}}%
\pgfpathlineto{\pgfqpoint{4.117612in}{0.633527in}}%
\pgfpathlineto{\pgfqpoint{4.117612in}{0.625278in}}%
\pgfpathlineto{\pgfqpoint{4.118767in}{0.633527in}}%
\pgfpathlineto{\pgfqpoint{4.118960in}{0.625278in}}%
\pgfpathlineto{\pgfqpoint{4.119797in}{0.625278in}}%
\pgfpathlineto{\pgfqpoint{4.120959in}{0.633527in}}%
\pgfpathlineto{\pgfqpoint{4.121300in}{0.650025in}}%
\pgfpathlineto{\pgfqpoint{4.121951in}{0.625278in}}%
\pgfpathlineto{\pgfqpoint{4.123136in}{0.633527in}}%
\pgfpathlineto{\pgfqpoint{4.123136in}{0.625278in}}%
\pgfpathlineto{\pgfqpoint{4.124143in}{0.625278in}}%
\pgfpathlineto{\pgfqpoint{4.125483in}{0.633527in}}%
\pgfpathlineto{\pgfqpoint{4.125483in}{0.625278in}}%
\pgfpathlineto{\pgfqpoint{4.126320in}{0.625278in}}%
\pgfpathlineto{\pgfqpoint{4.127497in}{0.633527in}}%
\pgfpathlineto{\pgfqpoint{4.127660in}{0.625278in}}%
\pgfpathlineto{\pgfqpoint{4.128334in}{0.625278in}}%
\pgfpathlineto{\pgfqpoint{4.129503in}{0.633527in}}%
\pgfpathlineto{\pgfqpoint{4.129837in}{0.625278in}}%
\pgfpathlineto{\pgfqpoint{4.130333in}{0.625278in}}%
\pgfpathlineto{\pgfqpoint{4.131680in}{0.633527in}}%
\pgfpathlineto{\pgfqpoint{4.132021in}{0.625278in}}%
\pgfpathlineto{\pgfqpoint{4.132687in}{0.625278in}}%
\pgfpathlineto{\pgfqpoint{4.133857in}{0.633527in}}%
\pgfpathlineto{\pgfqpoint{4.134020in}{0.625278in}}%
\pgfpathlineto{\pgfqpoint{4.134864in}{0.625278in}}%
\pgfpathlineto{\pgfqpoint{4.136041in}{0.633527in}}%
\pgfpathlineto{\pgfqpoint{4.136204in}{0.625278in}}%
\pgfpathlineto{\pgfqpoint{4.136545in}{0.625278in}}%
\pgfpathlineto{\pgfqpoint{4.137878in}{0.633527in}}%
\pgfpathlineto{\pgfqpoint{4.138211in}{0.625278in}}%
\pgfpathlineto{\pgfqpoint{4.138885in}{0.625278in}}%
\pgfpathlineto{\pgfqpoint{4.140054in}{0.633527in}}%
\pgfpathlineto{\pgfqpoint{4.140225in}{0.625278in}}%
\pgfpathlineto{\pgfqpoint{4.140899in}{0.625278in}}%
\pgfpathlineto{\pgfqpoint{4.142068in}{0.633527in}}%
\pgfpathlineto{\pgfqpoint{4.142402in}{0.625278in}}%
\pgfpathlineto{\pgfqpoint{4.143075in}{0.625278in}}%
\pgfpathlineto{\pgfqpoint{4.144245in}{0.633527in}}%
\pgfpathlineto{\pgfqpoint{4.144416in}{0.625278in}}%
\pgfpathlineto{\pgfqpoint{4.145252in}{0.625278in}}%
\pgfpathlineto{\pgfqpoint{4.146429in}{0.633527in}}%
\pgfpathlineto{\pgfqpoint{4.146592in}{0.625278in}}%
\pgfpathlineto{\pgfqpoint{4.147436in}{0.625278in}}%
\pgfpathlineto{\pgfqpoint{4.148606in}{0.633527in}}%
\pgfpathlineto{\pgfqpoint{4.148777in}{0.625278in}}%
\pgfpathlineto{\pgfqpoint{4.149613in}{0.625278in}}%
\pgfpathlineto{\pgfqpoint{4.150783in}{0.633527in}}%
\pgfpathlineto{\pgfqpoint{4.150953in}{0.625278in}}%
\pgfpathlineto{\pgfqpoint{4.151790in}{0.625278in}}%
\pgfpathlineto{\pgfqpoint{4.152967in}{0.633527in}}%
\pgfpathlineto{\pgfqpoint{4.153130in}{0.625278in}}%
\pgfpathlineto{\pgfqpoint{4.153967in}{0.625278in}}%
\pgfpathlineto{\pgfqpoint{4.155144in}{0.633527in}}%
\pgfpathlineto{\pgfqpoint{4.155144in}{0.625278in}}%
\pgfpathlineto{\pgfqpoint{4.156151in}{0.625278in}}%
\pgfpathlineto{\pgfqpoint{4.157321in}{0.633527in}}%
\pgfpathlineto{\pgfqpoint{4.157662in}{0.625278in}}%
\pgfpathlineto{\pgfqpoint{4.158328in}{0.625278in}}%
\pgfpathlineto{\pgfqpoint{4.159505in}{0.633527in}}%
\pgfpathlineto{\pgfqpoint{4.159505in}{0.625278in}}%
\pgfpathlineto{\pgfqpoint{4.160505in}{0.625278in}}%
\pgfpathlineto{\pgfqpoint{4.161682in}{0.633527in}}%
\pgfpathlineto{\pgfqpoint{4.162015in}{0.625278in}}%
\pgfpathlineto{\pgfqpoint{4.162696in}{0.625278in}}%
\pgfpathlineto{\pgfqpoint{4.163859in}{0.633527in}}%
\pgfpathlineto{\pgfqpoint{4.164037in}{0.625278in}}%
\pgfpathlineto{\pgfqpoint{4.164866in}{0.625278in}}%
\pgfpathlineto{\pgfqpoint{4.166043in}{0.633527in}}%
\pgfpathlineto{\pgfqpoint{4.166206in}{0.625278in}}%
\pgfpathlineto{\pgfqpoint{4.167043in}{0.625278in}}%
\pgfpathlineto{\pgfqpoint{4.168227in}{0.633527in}}%
\pgfpathlineto{\pgfqpoint{4.168390in}{0.625278in}}%
\pgfpathlineto{\pgfqpoint{4.169057in}{0.625278in}}%
\pgfpathlineto{\pgfqpoint{4.170234in}{0.633527in}}%
\pgfpathlineto{\pgfqpoint{4.170567in}{0.625278in}}%
\pgfpathlineto{\pgfqpoint{4.171234in}{0.625278in}}%
\pgfpathlineto{\pgfqpoint{4.172411in}{0.633527in}}%
\pgfpathlineto{\pgfqpoint{4.172581in}{0.625278in}}%
\pgfpathlineto{\pgfqpoint{4.173418in}{0.625278in}}%
\pgfpathlineto{\pgfqpoint{4.174588in}{0.633527in}}%
\pgfpathlineto{\pgfqpoint{4.174588in}{0.625278in}}%
\pgfpathlineto{\pgfqpoint{4.175424in}{0.625278in}}%
\pgfpathlineto{\pgfqpoint{4.176602in}{0.633527in}}%
\pgfpathlineto{\pgfqpoint{4.176772in}{0.625278in}}%
\pgfpathlineto{\pgfqpoint{4.177609in}{0.625278in}}%
\pgfpathlineto{\pgfqpoint{4.178778in}{0.633527in}}%
\pgfpathlineto{\pgfqpoint{4.178949in}{0.625278in}}%
\pgfpathlineto{\pgfqpoint{4.179785in}{0.625278in}}%
\pgfpathlineto{\pgfqpoint{4.180963in}{0.633527in}}%
\pgfpathlineto{\pgfqpoint{4.181126in}{0.625278in}}%
\pgfpathlineto{\pgfqpoint{4.181962in}{0.625278in}}%
\pgfpathlineto{\pgfqpoint{4.183139in}{0.633527in}}%
\pgfpathlineto{\pgfqpoint{4.183473in}{0.625278in}}%
\pgfpathlineto{\pgfqpoint{4.184146in}{0.625278in}}%
\pgfpathlineto{\pgfqpoint{4.185316in}{0.633527in}}%
\pgfpathlineto{\pgfqpoint{4.185316in}{0.625278in}}%
\pgfpathlineto{\pgfqpoint{4.186323in}{0.625278in}}%
\pgfpathlineto{\pgfqpoint{4.187493in}{0.633527in}}%
\pgfpathlineto{\pgfqpoint{4.187663in}{0.625278in}}%
\pgfpathlineto{\pgfqpoint{4.188500in}{0.625278in}}%
\pgfpathlineto{\pgfqpoint{4.189677in}{0.633527in}}%
\pgfpathlineto{\pgfqpoint{4.189840in}{0.625278in}}%
\pgfpathlineto{\pgfqpoint{4.190684in}{0.625278in}}%
\pgfpathlineto{\pgfqpoint{4.191854in}{0.633527in}}%
\pgfpathlineto{\pgfqpoint{4.192024in}{0.625278in}}%
\pgfpathlineto{\pgfqpoint{4.192861in}{0.625278in}}%
\pgfpathlineto{\pgfqpoint{4.194031in}{0.633527in}}%
\pgfpathlineto{\pgfqpoint{4.194201in}{0.625278in}}%
\pgfpathlineto{\pgfqpoint{4.195038in}{0.625278in}}%
\pgfpathlineto{\pgfqpoint{4.196045in}{0.633527in}}%
\pgfpathlineto{\pgfqpoint{4.196045in}{0.625278in}}%
\pgfpathlineto{\pgfqpoint{4.197222in}{0.633527in}}%
\pgfpathlineto{\pgfqpoint{4.197385in}{0.625278in}}%
\pgfpathlineto{\pgfqpoint{4.198222in}{0.625278in}}%
\pgfpathlineto{\pgfqpoint{4.199569in}{0.633527in}}%
\pgfpathlineto{\pgfqpoint{4.199732in}{0.625278in}}%
\pgfpathlineto{\pgfqpoint{4.200569in}{0.625278in}}%
\pgfpathlineto{\pgfqpoint{4.201746in}{0.633527in}}%
\pgfpathlineto{\pgfqpoint{4.201917in}{0.625278in}}%
\pgfpathlineto{\pgfqpoint{4.202753in}{0.625278in}}%
\pgfpathlineto{\pgfqpoint{4.203923in}{0.633527in}}%
\pgfpathlineto{\pgfqpoint{4.203923in}{0.625278in}}%
\pgfpathlineto{\pgfqpoint{4.204930in}{0.625278in}}%
\pgfpathlineto{\pgfqpoint{4.206107in}{0.633527in}}%
\pgfpathlineto{\pgfqpoint{4.206107in}{0.625278in}}%
\pgfpathlineto{\pgfqpoint{4.207107in}{0.625278in}}%
\pgfpathlineto{\pgfqpoint{4.208284in}{0.633527in}}%
\pgfpathlineto{\pgfqpoint{4.208447in}{0.625278in}}%
\pgfpathlineto{\pgfqpoint{4.209291in}{0.625278in}}%
\pgfpathlineto{\pgfqpoint{4.210461in}{0.633527in}}%
\pgfpathlineto{\pgfqpoint{4.210461in}{0.625278in}}%
\pgfpathlineto{\pgfqpoint{4.211468in}{0.625278in}}%
\pgfpathlineto{\pgfqpoint{4.212638in}{0.633527in}}%
\pgfpathlineto{\pgfqpoint{4.212808in}{0.625278in}}%
\pgfpathlineto{\pgfqpoint{4.213645in}{0.625278in}}%
\pgfpathlineto{\pgfqpoint{4.214822in}{0.633527in}}%
\pgfpathlineto{\pgfqpoint{4.215155in}{0.625278in}}%
\pgfpathlineto{\pgfqpoint{4.215881in}{0.625278in}}%
\pgfpathlineto{\pgfqpoint{4.216999in}{0.633527in}}%
\pgfpathlineto{\pgfqpoint{4.217169in}{0.625278in}}%
\pgfpathlineto{\pgfqpoint{4.218006in}{0.625278in}}%
\pgfpathlineto{\pgfqpoint{4.219176in}{0.633527in}}%
\pgfpathlineto{\pgfqpoint{4.219346in}{0.625278in}}%
\pgfpathlineto{\pgfqpoint{4.220183in}{0.625278in}}%
\pgfpathlineto{\pgfqpoint{4.221360in}{0.633527in}}%
\pgfpathlineto{\pgfqpoint{4.221360in}{0.625278in}}%
\pgfpathlineto{\pgfqpoint{4.222367in}{0.625278in}}%
\pgfpathlineto{\pgfqpoint{4.223537in}{0.633527in}}%
\pgfpathlineto{\pgfqpoint{4.223707in}{0.625278in}}%
\pgfpathlineto{\pgfqpoint{4.224544in}{0.625278in}}%
\pgfpathlineto{\pgfqpoint{4.225714in}{0.633527in}}%
\pgfpathlineto{\pgfqpoint{4.225714in}{0.625278in}}%
\pgfpathlineto{\pgfqpoint{4.226721in}{0.625278in}}%
\pgfpathlineto{\pgfqpoint{4.227898in}{0.633527in}}%
\pgfpathlineto{\pgfqpoint{4.228061in}{0.625278in}}%
\pgfpathlineto{\pgfqpoint{4.228905in}{0.625278in}}%
\pgfpathlineto{\pgfqpoint{4.230075in}{0.633527in}}%
\pgfpathlineto{\pgfqpoint{4.230075in}{0.625278in}}%
\pgfpathlineto{\pgfqpoint{4.231082in}{0.625278in}}%
\pgfpathlineto{\pgfqpoint{4.232251in}{0.633527in}}%
\pgfpathlineto{\pgfqpoint{4.232251in}{0.625278in}}%
\pgfpathlineto{\pgfqpoint{4.233258in}{0.625278in}}%
\pgfpathlineto{\pgfqpoint{4.234436in}{0.633527in}}%
\pgfpathlineto{\pgfqpoint{4.234769in}{0.625278in}}%
\pgfpathlineto{\pgfqpoint{4.235443in}{0.625278in}}%
\pgfpathlineto{\pgfqpoint{4.236613in}{0.633527in}}%
\pgfpathlineto{\pgfqpoint{4.236783in}{0.625278in}}%
\pgfpathlineto{\pgfqpoint{4.237449in}{0.625278in}}%
\pgfpathlineto{\pgfqpoint{4.238626in}{0.633527in}}%
\pgfpathlineto{\pgfqpoint{4.238789in}{0.625278in}}%
\pgfpathlineto{\pgfqpoint{4.239633in}{0.625278in}}%
\pgfpathlineto{\pgfqpoint{4.240633in}{0.633527in}}%
\pgfpathlineto{\pgfqpoint{4.240633in}{0.625278in}}%
\pgfpathlineto{\pgfqpoint{4.241810in}{0.633527in}}%
\pgfpathlineto{\pgfqpoint{4.242143in}{0.625278in}}%
\pgfpathlineto{\pgfqpoint{4.242817in}{0.625278in}}%
\pgfpathlineto{\pgfqpoint{4.243987in}{0.633527in}}%
\pgfpathlineto{\pgfqpoint{4.243987in}{0.625278in}}%
\pgfpathlineto{\pgfqpoint{4.244994in}{0.625278in}}%
\pgfpathlineto{\pgfqpoint{4.246001in}{0.633527in}}%
\pgfpathlineto{\pgfqpoint{4.246001in}{0.625278in}}%
\pgfpathlineto{\pgfqpoint{4.247171in}{0.633527in}}%
\pgfpathlineto{\pgfqpoint{4.247341in}{0.625278in}}%
\pgfpathlineto{\pgfqpoint{4.248015in}{0.625278in}}%
\pgfpathlineto{\pgfqpoint{4.249185in}{0.633527in}}%
\pgfpathlineto{\pgfqpoint{4.249355in}{0.625278in}}%
\pgfpathlineto{\pgfqpoint{4.250192in}{0.625278in}}%
\pgfpathlineto{\pgfqpoint{4.251362in}{0.633527in}}%
\pgfpathlineto{\pgfqpoint{4.251362in}{0.625278in}}%
\pgfpathlineto{\pgfqpoint{4.252369in}{0.625278in}}%
\pgfpathlineto{\pgfqpoint{4.253709in}{0.633527in}}%
\pgfpathlineto{\pgfqpoint{4.253709in}{0.625278in}}%
\pgfpathlineto{\pgfqpoint{4.254716in}{0.625278in}}%
\pgfpathlineto{\pgfqpoint{4.255893in}{0.633527in}}%
\pgfpathlineto{\pgfqpoint{4.256056in}{0.625278in}}%
\pgfpathlineto{\pgfqpoint{4.256893in}{0.625278in}}%
\pgfpathlineto{\pgfqpoint{4.257900in}{0.633527in}}%
\pgfpathlineto{\pgfqpoint{4.257900in}{0.625278in}}%
\pgfpathlineto{\pgfqpoint{4.259077in}{0.633527in}}%
\pgfpathlineto{\pgfqpoint{4.259077in}{0.625278in}}%
\pgfpathlineto{\pgfqpoint{4.260084in}{0.625278in}}%
\pgfpathlineto{\pgfqpoint{4.261254in}{0.633527in}}%
\pgfpathlineto{\pgfqpoint{4.261424in}{0.625278in}}%
\pgfpathlineto{\pgfqpoint{4.262261in}{0.625278in}}%
\pgfpathlineto{\pgfqpoint{4.263431in}{0.633527in}}%
\pgfpathlineto{\pgfqpoint{4.263431in}{0.625278in}}%
\pgfpathlineto{\pgfqpoint{4.264438in}{0.625278in}}%
\pgfpathlineto{\pgfqpoint{4.265444in}{0.633527in}}%
\pgfpathlineto{\pgfqpoint{4.265444in}{0.625278in}}%
\pgfpathlineto{\pgfqpoint{4.266622in}{0.633527in}}%
\pgfpathlineto{\pgfqpoint{4.266785in}{0.625278in}}%
\pgfpathlineto{\pgfqpoint{4.267458in}{0.625278in}}%
\pgfpathlineto{\pgfqpoint{4.268628in}{0.633527in}}%
\pgfpathlineto{\pgfqpoint{4.268628in}{0.625278in}}%
\pgfpathlineto{\pgfqpoint{4.269465in}{0.625278in}}%
\pgfpathlineto{\pgfqpoint{4.270642in}{0.633527in}}%
\pgfpathlineto{\pgfqpoint{4.270813in}{0.625278in}}%
\pgfpathlineto{\pgfqpoint{4.271649in}{0.625278in}}%
\pgfpathlineto{\pgfqpoint{4.272819in}{0.633527in}}%
\pgfpathlineto{\pgfqpoint{4.272841in}{0.625278in}}%
\pgfpathlineto{\pgfqpoint{4.273656in}{0.658274in}}%
\pgfpathlineto{\pgfqpoint{4.273848in}{0.650025in}}%
\pgfpathlineto{\pgfqpoint{4.275188in}{0.625278in}}%
\pgfpathlineto{\pgfqpoint{4.276343in}{0.633527in}}%
\pgfpathlineto{\pgfqpoint{4.276529in}{0.625278in}}%
\pgfpathlineto{\pgfqpoint{4.276699in}{0.650025in}}%
\pgfpathlineto{\pgfqpoint{4.277365in}{0.650025in}}%
\pgfpathlineto{\pgfqpoint{4.278691in}{0.633527in}}%
\pgfpathlineto{\pgfqpoint{4.279690in}{0.707769in}}%
\pgfpathlineto{\pgfqpoint{4.279712in}{0.699520in}}%
\pgfpathlineto{\pgfqpoint{4.280194in}{0.732517in}}%
\pgfpathlineto{\pgfqpoint{4.281053in}{0.716019in}}%
\pgfpathlineto{\pgfqpoint{4.281726in}{0.699520in}}%
\pgfpathlineto{\pgfqpoint{4.281223in}{0.724268in}}%
\pgfpathlineto{\pgfqpoint{4.281889in}{0.724268in}}%
\pgfpathlineto{\pgfqpoint{4.282230in}{0.749015in}}%
\pgfpathlineto{\pgfqpoint{4.282733in}{0.716019in}}%
\pgfpathlineto{\pgfqpoint{4.283066in}{0.699520in}}%
\pgfpathlineto{\pgfqpoint{4.283385in}{0.732517in}}%
\pgfpathlineto{\pgfqpoint{4.283881in}{0.707769in}}%
\pgfpathlineto{\pgfqpoint{4.283903in}{0.724268in}}%
\pgfpathlineto{\pgfqpoint{4.284888in}{0.683022in}}%
\pgfpathlineto{\pgfqpoint{4.284910in}{0.699520in}}%
\pgfpathlineto{\pgfqpoint{4.285562in}{0.683022in}}%
\pgfpathlineto{\pgfqpoint{4.285562in}{0.716019in}}%
\pgfpathlineto{\pgfqpoint{4.285895in}{0.707769in}}%
\pgfpathlineto{\pgfqpoint{4.286583in}{0.724268in}}%
\pgfpathlineto{\pgfqpoint{4.287072in}{0.732517in}}%
\pgfpathlineto{\pgfqpoint{4.287257in}{0.716019in}}%
\pgfpathlineto{\pgfqpoint{4.287590in}{0.699520in}}%
\pgfpathlineto{\pgfqpoint{4.287909in}{0.732517in}}%
\pgfpathlineto{\pgfqpoint{4.288072in}{0.732517in}}%
\pgfpathlineto{\pgfqpoint{4.288094in}{0.749015in}}%
\pgfpathlineto{\pgfqpoint{4.289101in}{0.749015in}}%
\pgfpathlineto{\pgfqpoint{4.289604in}{0.724268in}}%
\pgfpathlineto{\pgfqpoint{4.289923in}{0.757264in}}%
\pgfpathlineto{\pgfqpoint{4.289938in}{0.749015in}}%
\pgfpathlineto{\pgfqpoint{4.290086in}{0.757264in}}%
\pgfpathlineto{\pgfqpoint{4.290441in}{0.724268in}}%
\pgfpathlineto{\pgfqpoint{4.290589in}{0.732517in}}%
\pgfpathlineto{\pgfqpoint{4.290945in}{0.724268in}}%
\pgfpathlineto{\pgfqpoint{4.291115in}{0.749015in}}%
\pgfpathlineto{\pgfqpoint{4.291611in}{0.740766in}}%
\pgfpathlineto{\pgfqpoint{4.291952in}{0.724268in}}%
\pgfpathlineto{\pgfqpoint{4.292262in}{0.757264in}}%
\pgfpathlineto{\pgfqpoint{4.292603in}{0.806759in}}%
\pgfpathlineto{\pgfqpoint{4.293292in}{0.790261in}}%
\pgfpathlineto{\pgfqpoint{4.293440in}{0.782012in}}%
\pgfpathlineto{\pgfqpoint{4.293462in}{0.823258in}}%
\pgfpathlineto{\pgfqpoint{4.293958in}{0.815008in}}%
\pgfpathlineto{\pgfqpoint{4.294447in}{0.782012in}}%
\pgfpathlineto{\pgfqpoint{4.294802in}{0.823258in}}%
\pgfpathlineto{\pgfqpoint{4.295846in}{0.831507in}}%
\pgfpathlineto{\pgfqpoint{4.295972in}{0.823258in}}%
\pgfpathlineto{\pgfqpoint{4.296290in}{0.856254in}}%
\pgfpathlineto{\pgfqpoint{4.296475in}{0.848005in}}%
\pgfpathlineto{\pgfqpoint{4.296794in}{0.856254in}}%
\pgfpathlineto{\pgfqpoint{4.297127in}{0.831507in}}%
\pgfpathlineto{\pgfqpoint{4.297482in}{0.823258in}}%
\pgfpathlineto{\pgfqpoint{4.297801in}{0.856254in}}%
\pgfpathlineto{\pgfqpoint{4.298149in}{0.839756in}}%
\pgfpathlineto{\pgfqpoint{4.298489in}{0.823258in}}%
\pgfpathlineto{\pgfqpoint{4.298800in}{0.856254in}}%
\pgfpathlineto{\pgfqpoint{4.298823in}{0.848005in}}%
\pgfpathlineto{\pgfqpoint{4.300333in}{0.897500in}}%
\pgfpathlineto{\pgfqpoint{4.300496in}{0.889251in}}%
\pgfpathlineto{\pgfqpoint{4.300666in}{0.897500in}}%
\pgfpathlineto{\pgfqpoint{4.301651in}{0.881002in}}%
\pgfpathlineto{\pgfqpoint{4.302495in}{0.905749in}}%
\pgfpathlineto{\pgfqpoint{4.302680in}{0.897500in}}%
\pgfpathlineto{\pgfqpoint{4.303665in}{0.905749in}}%
\pgfpathlineto{\pgfqpoint{4.303050in}{0.881002in}}%
\pgfpathlineto{\pgfqpoint{4.303687in}{0.897500in}}%
\pgfpathlineto{\pgfqpoint{4.305005in}{0.881002in}}%
\pgfpathlineto{\pgfqpoint{4.305842in}{0.930497in}}%
\pgfpathlineto{\pgfqpoint{4.306034in}{0.922247in}}%
\pgfpathlineto{\pgfqpoint{4.306849in}{0.930497in}}%
\pgfpathlineto{\pgfqpoint{4.306871in}{0.922247in}}%
\pgfpathlineto{\pgfqpoint{4.307182in}{0.905749in}}%
\pgfpathlineto{\pgfqpoint{4.307352in}{0.930497in}}%
\pgfpathlineto{\pgfqpoint{4.307708in}{0.922247in}}%
\pgfpathlineto{\pgfqpoint{4.308863in}{0.930497in}}%
\pgfpathlineto{\pgfqpoint{4.310055in}{0.872753in}}%
\pgfpathlineto{\pgfqpoint{4.311210in}{0.881002in}}%
\pgfpathlineto{\pgfqpoint{4.312232in}{0.872753in}}%
\pgfpathlineto{\pgfqpoint{4.312720in}{0.856254in}}%
\pgfpathlineto{\pgfqpoint{4.313387in}{0.881002in}}%
\pgfpathlineto{\pgfqpoint{4.314060in}{0.856254in}}%
\pgfpathlineto{\pgfqpoint{4.314416in}{0.872753in}}%
\pgfpathlineto{\pgfqpoint{4.314727in}{0.881002in}}%
\pgfpathlineto{\pgfqpoint{4.315082in}{0.848005in}}%
\pgfpathlineto{\pgfqpoint{4.315230in}{0.856254in}}%
\pgfpathlineto{\pgfqpoint{4.315415in}{0.848005in}}%
\pgfpathlineto{\pgfqpoint{4.315904in}{0.905749in}}%
\pgfpathlineto{\pgfqpoint{4.316089in}{0.897500in}}%
\pgfpathlineto{\pgfqpoint{4.316570in}{0.905749in}}%
\pgfpathlineto{\pgfqpoint{4.316756in}{0.889251in}}%
\pgfpathlineto{\pgfqpoint{4.317259in}{0.848005in}}%
\pgfpathlineto{\pgfqpoint{4.318081in}{0.856254in}}%
\pgfpathlineto{\pgfqpoint{4.318584in}{0.881002in}}%
\pgfpathlineto{\pgfqpoint{4.319103in}{0.872753in}}%
\pgfpathlineto{\pgfqpoint{4.320761in}{0.905749in}}%
\pgfpathlineto{\pgfqpoint{4.321790in}{0.897500in}}%
\pgfpathlineto{\pgfqpoint{4.321939in}{0.905749in}}%
\pgfpathlineto{\pgfqpoint{4.322272in}{0.881002in}}%
\pgfpathlineto{\pgfqpoint{4.322627in}{0.897500in}}%
\pgfpathlineto{\pgfqpoint{4.323131in}{0.872753in}}%
\pgfpathlineto{\pgfqpoint{4.323449in}{0.905749in}}%
\pgfpathlineto{\pgfqpoint{4.323612in}{0.930497in}}%
\pgfpathlineto{\pgfqpoint{4.324471in}{0.922247in}}%
\pgfpathlineto{\pgfqpoint{4.324952in}{0.930497in}}%
\pgfpathlineto{\pgfqpoint{4.325137in}{0.913998in}}%
\pgfpathlineto{\pgfqpoint{4.325626in}{0.905749in}}%
\pgfpathlineto{\pgfqpoint{4.325626in}{0.938746in}}%
\pgfpathlineto{\pgfqpoint{4.325959in}{0.930497in}}%
\pgfpathlineto{\pgfqpoint{4.326796in}{0.979992in}}%
\pgfpathlineto{\pgfqpoint{4.327321in}{0.946995in}}%
\pgfpathlineto{\pgfqpoint{4.327825in}{0.971742in}}%
\pgfpathlineto{\pgfqpoint{4.328832in}{0.946995in}}%
\pgfpathlineto{\pgfqpoint{4.328980in}{0.963493in}}%
\pgfpathlineto{\pgfqpoint{4.329646in}{1.004739in}}%
\pgfpathlineto{\pgfqpoint{4.330002in}{0.971742in}}%
\pgfpathlineto{\pgfqpoint{4.330653in}{0.979992in}}%
\pgfpathlineto{\pgfqpoint{4.330838in}{0.963493in}}%
\pgfpathlineto{\pgfqpoint{4.330986in}{0.955244in}}%
\pgfpathlineto{\pgfqpoint{4.331490in}{0.988241in}}%
\pgfpathlineto{\pgfqpoint{4.332497in}{1.004739in}}%
\pgfpathlineto{\pgfqpoint{4.332519in}{0.996490in}}%
\pgfpathlineto{\pgfqpoint{4.333674in}{1.004739in}}%
\pgfpathlineto{\pgfqpoint{4.334363in}{0.996490in}}%
\pgfpathlineto{\pgfqpoint{4.334526in}{1.021237in}}%
\pgfpathlineto{\pgfqpoint{4.334696in}{1.012988in}}%
\pgfpathlineto{\pgfqpoint{4.335370in}{0.996490in}}%
\pgfpathlineto{\pgfqpoint{4.334866in}{1.021237in}}%
\pgfpathlineto{\pgfqpoint{4.335533in}{1.021237in}}%
\pgfpathlineto{\pgfqpoint{4.335851in}{1.029487in}}%
\pgfpathlineto{\pgfqpoint{4.336184in}{1.004739in}}%
\pgfpathlineto{\pgfqpoint{4.336369in}{1.012988in}}%
\pgfpathlineto{\pgfqpoint{4.336517in}{1.004739in}}%
\pgfpathlineto{\pgfqpoint{4.336688in}{1.054234in}}%
\pgfpathlineto{\pgfqpoint{4.336710in}{1.045985in}}%
\pgfpathlineto{\pgfqpoint{4.338028in}{1.078981in}}%
\pgfpathlineto{\pgfqpoint{4.338050in}{1.070732in}}%
\pgfpathlineto{\pgfqpoint{4.339205in}{1.054234in}}%
\pgfpathlineto{\pgfqpoint{4.338531in}{1.078981in}}%
\pgfpathlineto{\pgfqpoint{4.339205in}{1.062483in}}%
\pgfpathlineto{\pgfqpoint{4.339538in}{1.054234in}}%
\pgfpathlineto{\pgfqpoint{4.340375in}{1.103729in}}%
\pgfpathlineto{\pgfqpoint{4.341567in}{1.045985in}}%
\pgfpathlineto{\pgfqpoint{4.342389in}{1.103729in}}%
\pgfpathlineto{\pgfqpoint{4.342907in}{1.070732in}}%
\pgfpathlineto{\pgfqpoint{4.343899in}{1.029487in}}%
\pgfpathlineto{\pgfqpoint{4.343914in}{1.070732in}}%
\pgfpathlineto{\pgfqpoint{4.344566in}{1.103729in}}%
\pgfpathlineto{\pgfqpoint{4.345069in}{1.078981in}}%
\pgfpathlineto{\pgfqpoint{4.345254in}{1.070732in}}%
\pgfpathlineto{\pgfqpoint{4.345573in}{1.103729in}}%
\pgfpathlineto{\pgfqpoint{4.346091in}{1.095480in}}%
\pgfpathlineto{\pgfqpoint{4.347439in}{1.070732in}}%
\pgfpathlineto{\pgfqpoint{4.347942in}{1.045985in}}%
\pgfpathlineto{\pgfqpoint{4.348594in}{1.078981in}}%
\pgfpathlineto{\pgfqpoint{4.349430in}{1.029487in}}%
\pgfpathlineto{\pgfqpoint{4.349615in}{1.045985in}}%
\pgfpathlineto{\pgfqpoint{4.350452in}{1.070732in}}%
\pgfpathlineto{\pgfqpoint{4.350770in}{1.054234in}}%
\pgfpathlineto{\pgfqpoint{4.351792in}{1.037736in}}%
\pgfpathlineto{\pgfqpoint{4.352281in}{1.029487in}}%
\pgfpathlineto{\pgfqpoint{4.352444in}{1.054234in}}%
\pgfpathlineto{\pgfqpoint{4.352466in}{1.045985in}}%
\pgfpathlineto{\pgfqpoint{4.352947in}{1.078981in}}%
\pgfpathlineto{\pgfqpoint{4.353621in}{1.054234in}}%
\pgfpathlineto{\pgfqpoint{4.354643in}{1.045985in}}%
\pgfpathlineto{\pgfqpoint{4.356635in}{1.103729in}}%
\pgfpathlineto{\pgfqpoint{4.357530in}{1.078981in}}%
\pgfpathlineto{\pgfqpoint{4.357664in}{1.095480in}}%
\pgfpathlineto{\pgfqpoint{4.358649in}{1.078981in}}%
\pgfpathlineto{\pgfqpoint{4.358663in}{1.095480in}}%
\pgfpathlineto{\pgfqpoint{4.359818in}{1.103729in}}%
\pgfpathlineto{\pgfqpoint{4.360514in}{1.120227in}}%
\pgfpathlineto{\pgfqpoint{4.361010in}{1.070732in}}%
\pgfpathlineto{\pgfqpoint{4.361662in}{1.054234in}}%
\pgfpathlineto{\pgfqpoint{4.361832in}{1.078981in}}%
\pgfpathlineto{\pgfqpoint{4.361855in}{1.070732in}}%
\pgfpathlineto{\pgfqpoint{4.363676in}{1.153224in}}%
\pgfpathlineto{\pgfqpoint{4.363861in}{1.136726in}}%
\pgfpathlineto{\pgfqpoint{4.364202in}{1.120227in}}%
\pgfpathlineto{\pgfqpoint{4.364513in}{1.153224in}}%
\pgfpathlineto{\pgfqpoint{4.364535in}{1.144975in}}%
\pgfpathlineto{\pgfqpoint{4.364683in}{1.153224in}}%
\pgfpathlineto{\pgfqpoint{4.365016in}{1.128476in}}%
\pgfpathlineto{\pgfqpoint{4.365201in}{1.136726in}}%
\pgfpathlineto{\pgfqpoint{4.365875in}{1.120227in}}%
\pgfpathlineto{\pgfqpoint{4.365372in}{1.144975in}}%
\pgfpathlineto{\pgfqpoint{4.366045in}{1.144975in}}%
\pgfpathlineto{\pgfqpoint{4.368037in}{1.202719in}}%
\pgfpathlineto{\pgfqpoint{4.369044in}{1.177971in}}%
\pgfpathlineto{\pgfqpoint{4.369059in}{1.194470in}}%
\pgfpathlineto{\pgfqpoint{4.371406in}{1.120227in}}%
\pgfpathlineto{\pgfqpoint{4.371909in}{1.169722in}}%
\pgfpathlineto{\pgfqpoint{4.372746in}{1.136726in}}%
\pgfpathlineto{\pgfqpoint{4.373235in}{1.128476in}}%
\pgfpathlineto{\pgfqpoint{4.373398in}{1.153224in}}%
\pgfpathlineto{\pgfqpoint{4.373420in}{1.144975in}}%
\pgfpathlineto{\pgfqpoint{4.375241in}{1.202719in}}%
\pgfpathlineto{\pgfqpoint{4.375264in}{1.194470in}}%
\pgfpathlineto{\pgfqpoint{4.375930in}{1.169722in}}%
\pgfpathlineto{\pgfqpoint{4.376100in}{1.194470in}}%
\pgfpathlineto{\pgfqpoint{4.376922in}{1.202719in}}%
\pgfpathlineto{\pgfqpoint{4.376937in}{1.194470in}}%
\pgfpathlineto{\pgfqpoint{4.377774in}{1.169722in}}%
\pgfpathlineto{\pgfqpoint{4.377944in}{1.194470in}}%
\pgfpathlineto{\pgfqpoint{4.378277in}{1.219217in}}%
\pgfpathlineto{\pgfqpoint{4.378618in}{1.186221in}}%
\pgfpathlineto{\pgfqpoint{4.379284in}{1.194470in}}%
\pgfpathlineto{\pgfqpoint{4.380624in}{1.120227in}}%
\pgfpathlineto{\pgfqpoint{4.381468in}{1.169722in}}%
\pgfpathlineto{\pgfqpoint{4.382135in}{1.144975in}}%
\pgfpathlineto{\pgfqpoint{4.382808in}{1.095480in}}%
\pgfpathlineto{\pgfqpoint{4.383793in}{1.111978in}}%
\pgfpathlineto{\pgfqpoint{4.384149in}{1.144975in}}%
\pgfpathlineto{\pgfqpoint{4.384815in}{1.120227in}}%
\pgfpathlineto{\pgfqpoint{4.385970in}{1.128476in}}%
\pgfpathlineto{\pgfqpoint{4.387162in}{1.070732in}}%
\pgfpathlineto{\pgfqpoint{4.387651in}{1.078981in}}%
\pgfpathlineto{\pgfqpoint{4.387836in}{1.062483in}}%
\pgfpathlineto{\pgfqpoint{4.388317in}{1.054234in}}%
\pgfpathlineto{\pgfqpoint{4.388339in}{1.095480in}}%
\pgfpathlineto{\pgfqpoint{4.388502in}{1.087231in}}%
\pgfpathlineto{\pgfqpoint{4.390183in}{1.045985in}}%
\pgfpathlineto{\pgfqpoint{4.391671in}{1.078981in}}%
\pgfpathlineto{\pgfqpoint{4.391693in}{1.070732in}}%
\pgfpathlineto{\pgfqpoint{4.398061in}{0.773763in}}%
\pgfpathlineto{\pgfqpoint{4.398209in}{0.790261in}}%
\pgfpathlineto{\pgfqpoint{4.398542in}{0.806759in}}%
\pgfpathlineto{\pgfqpoint{4.399231in}{0.773763in}}%
\pgfpathlineto{\pgfqpoint{4.399549in}{0.782012in}}%
\pgfpathlineto{\pgfqpoint{4.399882in}{0.757264in}}%
\pgfpathlineto{\pgfqpoint{4.401075in}{0.724268in}}%
\pgfpathlineto{\pgfqpoint{4.402074in}{0.732517in}}%
\pgfpathlineto{\pgfqpoint{4.401726in}{0.707769in}}%
\pgfpathlineto{\pgfqpoint{4.402082in}{0.724268in}}%
\pgfpathlineto{\pgfqpoint{4.403762in}{0.650025in}}%
\pgfpathlineto{\pgfqpoint{4.404258in}{0.666524in}}%
\pgfpathlineto{\pgfqpoint{4.404925in}{0.707769in}}%
\pgfpathlineto{\pgfqpoint{4.405265in}{0.674773in}}%
\pgfpathlineto{\pgfqpoint{4.405584in}{0.658274in}}%
\pgfpathlineto{\pgfqpoint{4.405924in}{0.691271in}}%
\pgfpathlineto{\pgfqpoint{4.405939in}{0.724268in}}%
\pgfpathlineto{\pgfqpoint{4.406946in}{0.691271in}}%
\pgfpathlineto{\pgfqpoint{4.407109in}{0.699520in}}%
\pgfpathlineto{\pgfqpoint{4.408109in}{0.683022in}}%
\pgfpathlineto{\pgfqpoint{4.408790in}{0.749015in}}%
\pgfpathlineto{\pgfqpoint{4.409123in}{0.716019in}}%
\pgfpathlineto{\pgfqpoint{4.409271in}{0.707769in}}%
\pgfpathlineto{\pgfqpoint{4.409626in}{0.749015in}}%
\pgfpathlineto{\pgfqpoint{4.410463in}{0.773763in}}%
\pgfpathlineto{\pgfqpoint{4.410633in}{0.765513in}}%
\pgfpathlineto{\pgfqpoint{4.411974in}{0.683022in}}%
\pgfpathlineto{\pgfqpoint{4.412144in}{0.699520in}}%
\pgfpathlineto{\pgfqpoint{4.412980in}{0.683022in}}%
\pgfpathlineto{\pgfqpoint{4.412477in}{0.724268in}}%
\pgfpathlineto{\pgfqpoint{4.413143in}{0.699520in}}%
\pgfpathlineto{\pgfqpoint{4.414321in}{0.724268in}}%
\pgfpathlineto{\pgfqpoint{4.414491in}{0.716019in}}%
\pgfpathlineto{\pgfqpoint{4.414491in}{0.707769in}}%
\pgfpathlineto{\pgfqpoint{4.414654in}{0.749015in}}%
\pgfpathlineto{\pgfqpoint{4.415491in}{0.740766in}}%
\pgfpathlineto{\pgfqpoint{4.415994in}{0.732517in}}%
\pgfpathlineto{\pgfqpoint{4.416031in}{0.757264in}}%
\pgfpathlineto{\pgfqpoint{4.416164in}{0.749015in}}%
\pgfpathlineto{\pgfqpoint{4.416505in}{0.782012in}}%
\pgfpathlineto{\pgfqpoint{4.417171in}{0.757264in}}%
\pgfpathlineto{\pgfqpoint{4.417675in}{0.749015in}}%
\pgfpathlineto{\pgfqpoint{4.417838in}{0.773763in}}%
\pgfpathlineto{\pgfqpoint{4.418008in}{0.765513in}}%
\pgfpathlineto{\pgfqpoint{4.419518in}{0.848005in}}%
\pgfpathlineto{\pgfqpoint{4.421695in}{0.749015in}}%
\pgfpathlineto{\pgfqpoint{4.421866in}{0.773763in}}%
\pgfpathlineto{\pgfqpoint{4.422702in}{0.798510in}}%
\pgfpathlineto{\pgfqpoint{4.423028in}{0.782012in}}%
\pgfpathlineto{\pgfqpoint{4.424042in}{0.724268in}}%
\pgfpathlineto{\pgfqpoint{4.426715in}{0.625278in}}%
\pgfpathlineto{\pgfqpoint{4.427885in}{0.633527in}}%
\pgfpathlineto{\pgfqpoint{4.428055in}{0.625278in}}%
\pgfpathlineto{\pgfqpoint{4.428900in}{0.625278in}}%
\pgfpathlineto{\pgfqpoint{4.430232in}{0.633527in}}%
\pgfpathlineto{\pgfqpoint{4.430247in}{0.625278in}}%
\pgfpathlineto{\pgfqpoint{4.431254in}{0.625278in}}%
\pgfpathlineto{\pgfqpoint{4.432594in}{0.633527in}}%
\pgfpathlineto{\pgfqpoint{4.432594in}{0.625278in}}%
\pgfpathlineto{\pgfqpoint{4.433261in}{0.625278in}}%
\pgfpathlineto{\pgfqpoint{4.434438in}{0.633527in}}%
\pgfpathlineto{\pgfqpoint{4.434608in}{0.625278in}}%
\pgfpathlineto{\pgfqpoint{4.435275in}{0.625278in}}%
\pgfpathlineto{\pgfqpoint{4.436452in}{0.633527in}}%
\pgfpathlineto{\pgfqpoint{4.436607in}{0.625278in}}%
\pgfpathlineto{\pgfqpoint{4.437451in}{0.625278in}}%
\pgfpathlineto{\pgfqpoint{4.438621in}{0.633527in}}%
\pgfpathlineto{\pgfqpoint{4.438784in}{0.625278in}}%
\pgfpathlineto{\pgfqpoint{4.439621in}{0.625278in}}%
\pgfpathlineto{\pgfqpoint{4.440791in}{0.633527in}}%
\pgfpathlineto{\pgfqpoint{4.441139in}{0.625278in}}%
\pgfpathlineto{\pgfqpoint{4.441812in}{0.625278in}}%
\pgfpathlineto{\pgfqpoint{4.442819in}{0.633527in}}%
\pgfpathlineto{\pgfqpoint{4.442819in}{0.625278in}}%
\pgfpathlineto{\pgfqpoint{4.443989in}{0.633527in}}%
\pgfpathlineto{\pgfqpoint{4.444145in}{0.625278in}}%
\pgfpathlineto{\pgfqpoint{4.444989in}{0.625278in}}%
\pgfpathlineto{\pgfqpoint{4.446159in}{0.633527in}}%
\pgfpathlineto{\pgfqpoint{4.446159in}{0.625278in}}%
\pgfpathlineto{\pgfqpoint{4.447166in}{0.625278in}}%
\pgfpathlineto{\pgfqpoint{4.448498in}{0.633527in}}%
\pgfpathlineto{\pgfqpoint{4.448669in}{0.625278in}}%
\pgfpathlineto{\pgfqpoint{4.449513in}{0.625278in}}%
\pgfpathlineto{\pgfqpoint{4.450683in}{0.633527in}}%
\pgfpathlineto{\pgfqpoint{4.451023in}{0.625278in}}%
\pgfpathlineto{\pgfqpoint{4.451527in}{0.625278in}}%
\pgfpathlineto{\pgfqpoint{4.452697in}{0.633527in}}%
\pgfpathlineto{\pgfqpoint{4.452867in}{0.625278in}}%
\pgfpathlineto{\pgfqpoint{4.453704in}{0.625278in}}%
\pgfpathlineto{\pgfqpoint{4.454873in}{0.633527in}}%
\pgfpathlineto{\pgfqpoint{4.455044in}{0.625278in}}%
\pgfpathlineto{\pgfqpoint{4.455880in}{0.625278in}}%
\pgfpathlineto{\pgfqpoint{4.456998in}{0.633527in}}%
\pgfpathlineto{\pgfqpoint{4.457050in}{0.625278in}}%
\pgfpathlineto{\pgfqpoint{4.457717in}{0.625278in}}%
\pgfpathlineto{\pgfqpoint{4.458894in}{0.633527in}}%
\pgfpathlineto{\pgfqpoint{4.459072in}{0.625278in}}%
\pgfpathlineto{\pgfqpoint{4.459901in}{0.625278in}}%
\pgfpathlineto{\pgfqpoint{4.461071in}{0.633527in}}%
\pgfpathlineto{\pgfqpoint{4.461241in}{0.625278in}}%
\pgfpathlineto{\pgfqpoint{4.462078in}{0.625278in}}%
\pgfpathlineto{\pgfqpoint{4.463248in}{0.633527in}}%
\pgfpathlineto{\pgfqpoint{4.463248in}{0.625278in}}%
\pgfpathlineto{\pgfqpoint{4.464255in}{0.625278in}}%
\pgfpathlineto{\pgfqpoint{4.465262in}{0.633527in}}%
\pgfpathlineto{\pgfqpoint{4.465262in}{0.625278in}}%
\pgfpathlineto{\pgfqpoint{4.466276in}{0.633527in}}%
\pgfpathlineto{\pgfqpoint{4.466276in}{0.625278in}}%
\pgfpathlineto{\pgfqpoint{4.467438in}{0.633527in}}%
\pgfpathlineto{\pgfqpoint{4.467446in}{0.625278in}}%
\pgfpathlineto{\pgfqpoint{4.468445in}{0.625278in}}%
\pgfpathlineto{\pgfqpoint{4.469623in}{0.633527in}}%
\pgfpathlineto{\pgfqpoint{4.469630in}{0.625278in}}%
\pgfpathlineto{\pgfqpoint{4.470622in}{0.625278in}}%
\pgfpathlineto{\pgfqpoint{4.471799in}{0.633527in}}%
\pgfpathlineto{\pgfqpoint{4.471807in}{0.625278in}}%
\pgfpathlineto{\pgfqpoint{4.472806in}{0.625278in}}%
\pgfpathlineto{\pgfqpoint{4.473821in}{0.633527in}}%
\pgfpathlineto{\pgfqpoint{4.473821in}{0.625278in}}%
\pgfpathlineto{\pgfqpoint{4.474983in}{0.633527in}}%
\pgfpathlineto{\pgfqpoint{4.475154in}{0.625278in}}%
\pgfpathlineto{\pgfqpoint{4.475990in}{0.625278in}}%
\pgfpathlineto{\pgfqpoint{4.477160in}{0.633527in}}%
\pgfpathlineto{\pgfqpoint{4.477501in}{0.625278in}}%
\pgfpathlineto{\pgfqpoint{4.478167in}{0.625278in}}%
\pgfpathlineto{\pgfqpoint{4.479344in}{0.633527in}}%
\pgfpathlineto{\pgfqpoint{4.479352in}{0.625278in}}%
\pgfpathlineto{\pgfqpoint{4.480359in}{0.625278in}}%
\pgfpathlineto{\pgfqpoint{4.481358in}{0.633527in}}%
\pgfpathlineto{\pgfqpoint{4.481358in}{0.625278in}}%
\pgfpathlineto{\pgfqpoint{4.482528in}{0.633527in}}%
\pgfpathlineto{\pgfqpoint{4.482698in}{0.625278in}}%
\pgfpathlineto{\pgfqpoint{4.483535in}{0.625278in}}%
\pgfpathlineto{\pgfqpoint{4.484705in}{0.633527in}}%
\pgfpathlineto{\pgfqpoint{4.484705in}{0.625278in}}%
\pgfpathlineto{\pgfqpoint{4.485712in}{0.625278in}}%
\pgfpathlineto{\pgfqpoint{4.486719in}{0.633527in}}%
\pgfpathlineto{\pgfqpoint{4.486719in}{0.625278in}}%
\pgfpathlineto{\pgfqpoint{4.488059in}{0.633527in}}%
\pgfpathlineto{\pgfqpoint{4.488059in}{0.625278in}}%
\pgfpathlineto{\pgfqpoint{4.489066in}{0.625278in}}%
\pgfpathlineto{\pgfqpoint{4.490236in}{0.633527in}}%
\pgfpathlineto{\pgfqpoint{4.490406in}{0.625278in}}%
\pgfpathlineto{\pgfqpoint{4.491080in}{0.625278in}}%
\pgfpathlineto{\pgfqpoint{4.492250in}{0.633527in}}%
\pgfpathlineto{\pgfqpoint{4.492583in}{0.625278in}}%
\pgfpathlineto{\pgfqpoint{4.493257in}{0.625278in}}%
\pgfpathlineto{\pgfqpoint{4.494427in}{0.633527in}}%
\pgfpathlineto{\pgfqpoint{4.494434in}{0.625278in}}%
\pgfpathlineto{\pgfqpoint{4.495434in}{0.625278in}}%
\pgfpathlineto{\pgfqpoint{4.496611in}{0.633527in}}%
\pgfpathlineto{\pgfqpoint{4.496774in}{0.625278in}}%
\pgfpathlineto{\pgfqpoint{4.497448in}{0.625278in}}%
\pgfpathlineto{\pgfqpoint{4.498617in}{0.633527in}}%
\pgfpathlineto{\pgfqpoint{4.498788in}{0.625278in}}%
\pgfpathlineto{\pgfqpoint{4.499624in}{0.625278in}}%
\pgfpathlineto{\pgfqpoint{4.500631in}{0.633527in}}%
\pgfpathlineto{\pgfqpoint{4.500631in}{0.625278in}}%
\pgfpathlineto{\pgfqpoint{4.501801in}{0.633527in}}%
\pgfpathlineto{\pgfqpoint{4.501801in}{0.625278in}}%
\pgfpathlineto{\pgfqpoint{4.502808in}{0.625278in}}%
\pgfpathlineto{\pgfqpoint{4.503985in}{0.633527in}}%
\pgfpathlineto{\pgfqpoint{4.503993in}{0.625278in}}%
\pgfpathlineto{\pgfqpoint{4.504992in}{0.625278in}}%
\pgfpathlineto{\pgfqpoint{4.506162in}{0.633527in}}%
\pgfpathlineto{\pgfqpoint{4.506340in}{0.625278in}}%
\pgfpathlineto{\pgfqpoint{4.507169in}{0.625278in}}%
\pgfpathlineto{\pgfqpoint{4.508339in}{0.633527in}}%
\pgfpathlineto{\pgfqpoint{4.508339in}{0.625278in}}%
\pgfpathlineto{\pgfqpoint{4.509346in}{0.625278in}}%
\pgfpathlineto{\pgfqpoint{4.510523in}{0.633527in}}%
\pgfpathlineto{\pgfqpoint{4.510686in}{0.625278in}}%
\pgfpathlineto{\pgfqpoint{4.511360in}{0.625278in}}%
\pgfpathlineto{\pgfqpoint{4.512530in}{0.633527in}}%
\pgfpathlineto{\pgfqpoint{4.512700in}{0.625278in}}%
\pgfpathlineto{\pgfqpoint{4.513537in}{0.625278in}}%
\pgfpathlineto{\pgfqpoint{4.514544in}{0.633527in}}%
\pgfpathlineto{\pgfqpoint{4.514544in}{0.625278in}}%
\pgfpathlineto{\pgfqpoint{4.515721in}{0.633527in}}%
\pgfpathlineto{\pgfqpoint{4.515884in}{0.625278in}}%
\pgfpathlineto{\pgfqpoint{4.516721in}{0.625278in}}%
\pgfpathlineto{\pgfqpoint{4.517898in}{0.633527in}}%
\pgfpathlineto{\pgfqpoint{4.518068in}{0.625278in}}%
\pgfpathlineto{\pgfqpoint{4.518735in}{0.625278in}}%
\pgfpathlineto{\pgfqpoint{4.520075in}{0.633527in}}%
\pgfpathlineto{\pgfqpoint{4.520415in}{0.625278in}}%
\pgfpathlineto{\pgfqpoint{4.521082in}{0.625278in}}%
\pgfpathlineto{\pgfqpoint{4.522089in}{0.633527in}}%
\pgfpathlineto{\pgfqpoint{4.522089in}{0.625278in}}%
\pgfpathlineto{\pgfqpoint{4.523429in}{0.633527in}}%
\pgfpathlineto{\pgfqpoint{4.523599in}{0.625278in}}%
\pgfpathlineto{\pgfqpoint{4.524436in}{0.625278in}}%
\pgfpathlineto{\pgfqpoint{4.525443in}{0.633527in}}%
\pgfpathlineto{\pgfqpoint{4.525443in}{0.625278in}}%
\pgfpathlineto{\pgfqpoint{4.526613in}{0.633527in}}%
\pgfpathlineto{\pgfqpoint{4.526783in}{0.625278in}}%
\pgfpathlineto{\pgfqpoint{4.527620in}{0.625278in}}%
\pgfpathlineto{\pgfqpoint{4.528797in}{0.633527in}}%
\pgfpathlineto{\pgfqpoint{4.528960in}{0.625278in}}%
\pgfpathlineto{\pgfqpoint{4.529796in}{0.625278in}}%
\pgfpathlineto{\pgfqpoint{4.530803in}{0.633527in}}%
\pgfpathlineto{\pgfqpoint{4.530803in}{0.625278in}}%
\pgfpathlineto{\pgfqpoint{4.531981in}{0.633527in}}%
\pgfpathlineto{\pgfqpoint{4.532144in}{0.625278in}}%
\pgfpathlineto{\pgfqpoint{4.532988in}{0.625278in}}%
\pgfpathlineto{\pgfqpoint{4.534158in}{0.633527in}}%
\pgfpathlineto{\pgfqpoint{4.534158in}{0.625278in}}%
\pgfpathlineto{\pgfqpoint{4.535165in}{0.625278in}}%
\pgfpathlineto{\pgfqpoint{4.536334in}{0.633527in}}%
\pgfpathlineto{\pgfqpoint{4.536505in}{0.625278in}}%
\pgfpathlineto{\pgfqpoint{4.537341in}{0.625278in}}%
\pgfpathlineto{\pgfqpoint{4.538519in}{0.633527in}}%
\pgfpathlineto{\pgfqpoint{4.538682in}{0.625278in}}%
\pgfpathlineto{\pgfqpoint{4.539518in}{0.625278in}}%
\pgfpathlineto{\pgfqpoint{4.540695in}{0.633527in}}%
\pgfpathlineto{\pgfqpoint{4.540866in}{0.625278in}}%
\pgfpathlineto{\pgfqpoint{4.541702in}{0.625278in}}%
\pgfpathlineto{\pgfqpoint{4.542872in}{0.633527in}}%
\pgfpathlineto{\pgfqpoint{4.542872in}{0.625278in}}%
\pgfpathlineto{\pgfqpoint{4.543879in}{0.625278in}}%
\pgfpathlineto{\pgfqpoint{4.545057in}{0.633527in}}%
\pgfpathlineto{\pgfqpoint{4.545219in}{0.625278in}}%
\pgfpathlineto{\pgfqpoint{4.546056in}{0.625278in}}%
\pgfpathlineto{\pgfqpoint{4.547233in}{0.633527in}}%
\pgfpathlineto{\pgfqpoint{4.547233in}{0.625278in}}%
\pgfpathlineto{\pgfqpoint{4.548292in}{0.625278in}}%
\pgfpathlineto{\pgfqpoint{4.549410in}{0.633527in}}%
\pgfpathlineto{\pgfqpoint{4.549580in}{0.625278in}}%
\pgfpathlineto{\pgfqpoint{4.550417in}{0.625278in}}%
\pgfpathlineto{\pgfqpoint{4.551594in}{0.633527in}}%
\pgfpathlineto{\pgfqpoint{4.551594in}{0.625278in}}%
\pgfpathlineto{\pgfqpoint{4.552594in}{0.625278in}}%
\pgfpathlineto{\pgfqpoint{4.553771in}{0.633527in}}%
\pgfpathlineto{\pgfqpoint{4.553942in}{0.625278in}}%
\pgfpathlineto{\pgfqpoint{4.554778in}{0.625278in}}%
\pgfpathlineto{\pgfqpoint{4.555948in}{0.633527in}}%
\pgfpathlineto{\pgfqpoint{4.556118in}{0.625278in}}%
\pgfpathlineto{\pgfqpoint{4.556955in}{0.625278in}}%
\pgfpathlineto{\pgfqpoint{4.558132in}{0.633527in}}%
\pgfpathlineto{\pgfqpoint{4.558295in}{0.625278in}}%
\pgfpathlineto{\pgfqpoint{4.559132in}{0.625278in}}%
\pgfpathlineto{\pgfqpoint{4.560472in}{0.633527in}}%
\pgfpathlineto{\pgfqpoint{4.560472in}{0.625278in}}%
\pgfpathlineto{\pgfqpoint{4.561479in}{0.625278in}}%
\pgfpathlineto{\pgfqpoint{4.562656in}{0.633527in}}%
\pgfpathlineto{\pgfqpoint{4.562819in}{0.625278in}}%
\pgfpathlineto{\pgfqpoint{4.563663in}{0.625278in}}%
\pgfpathlineto{\pgfqpoint{4.564833in}{0.633527in}}%
\pgfpathlineto{\pgfqpoint{4.564833in}{0.625278in}}%
\pgfpathlineto{\pgfqpoint{4.565840in}{0.625278in}}%
\pgfpathlineto{\pgfqpoint{4.567010in}{0.633527in}}%
\pgfpathlineto{\pgfqpoint{4.567180in}{0.625278in}}%
\pgfpathlineto{\pgfqpoint{4.568017in}{0.625278in}}%
\pgfpathlineto{\pgfqpoint{4.569194in}{0.633527in}}%
\pgfpathlineto{\pgfqpoint{4.569357in}{0.625278in}}%
\pgfpathlineto{\pgfqpoint{4.570201in}{0.625278in}}%
\pgfpathlineto{\pgfqpoint{4.571371in}{0.633527in}}%
\pgfpathlineto{\pgfqpoint{4.571371in}{0.625278in}}%
\pgfpathlineto{\pgfqpoint{4.572208in}{0.625278in}}%
\pgfpathlineto{\pgfqpoint{4.573385in}{0.633527in}}%
\pgfpathlineto{\pgfqpoint{4.573548in}{0.625278in}}%
\pgfpathlineto{\pgfqpoint{4.574222in}{0.625278in}}%
\pgfpathlineto{\pgfqpoint{4.575562in}{0.633527in}}%
\pgfpathlineto{\pgfqpoint{4.575562in}{0.625278in}}%
\pgfpathlineto{\pgfqpoint{4.576398in}{0.625278in}}%
\pgfpathlineto{\pgfqpoint{4.577739in}{0.633527in}}%
\pgfpathlineto{\pgfqpoint{4.577739in}{0.625278in}}%
\pgfpathlineto{\pgfqpoint{4.578583in}{0.625278in}}%
\pgfpathlineto{\pgfqpoint{4.579753in}{0.633527in}}%
\pgfpathlineto{\pgfqpoint{4.579923in}{0.625278in}}%
\pgfpathlineto{\pgfqpoint{4.580760in}{0.625278in}}%
\pgfpathlineto{\pgfqpoint{4.581929in}{0.633527in}}%
\pgfpathlineto{\pgfqpoint{4.582100in}{0.625278in}}%
\pgfpathlineto{\pgfqpoint{4.582936in}{0.625278in}}%
\pgfpathlineto{\pgfqpoint{4.584114in}{0.633527in}}%
\pgfpathlineto{\pgfqpoint{4.584277in}{0.625278in}}%
\pgfpathlineto{\pgfqpoint{4.584950in}{0.625278in}}%
\pgfpathlineto{\pgfqpoint{4.586120in}{0.633527in}}%
\pgfpathlineto{\pgfqpoint{4.586290in}{0.625278in}}%
\pgfpathlineto{\pgfqpoint{4.587127in}{0.625278in}}%
\pgfpathlineto{\pgfqpoint{4.588304in}{0.633527in}}%
\pgfpathlineto{\pgfqpoint{4.588304in}{0.625278in}}%
\pgfpathlineto{\pgfqpoint{4.589311in}{0.625278in}}%
\pgfpathlineto{\pgfqpoint{4.590481in}{0.633527in}}%
\pgfpathlineto{\pgfqpoint{4.590652in}{0.625278in}}%
\pgfpathlineto{\pgfqpoint{4.591488in}{0.625278in}}%
\pgfpathlineto{\pgfqpoint{4.592658in}{0.633527in}}%
\pgfpathlineto{\pgfqpoint{4.592828in}{0.625278in}}%
\pgfpathlineto{\pgfqpoint{4.593665in}{0.625278in}}%
\pgfpathlineto{\pgfqpoint{4.594842in}{0.633527in}}%
\pgfpathlineto{\pgfqpoint{4.595176in}{0.625278in}}%
\pgfpathlineto{\pgfqpoint{4.595849in}{0.625278in}}%
\pgfpathlineto{\pgfqpoint{4.597019in}{0.633527in}}%
\pgfpathlineto{\pgfqpoint{4.597189in}{0.625278in}}%
\pgfpathlineto{\pgfqpoint{4.598026in}{0.625278in}}%
\pgfpathlineto{\pgfqpoint{4.599196in}{0.633527in}}%
\pgfpathlineto{\pgfqpoint{4.599366in}{0.625278in}}%
\pgfpathlineto{\pgfqpoint{4.600203in}{0.625278in}}%
\pgfpathlineto{\pgfqpoint{4.601210in}{0.633527in}}%
\pgfpathlineto{\pgfqpoint{4.601210in}{0.625278in}}%
\pgfpathlineto{\pgfqpoint{4.602380in}{0.633527in}}%
\pgfpathlineto{\pgfqpoint{4.602550in}{0.625278in}}%
\pgfpathlineto{\pgfqpoint{4.603224in}{0.625278in}}%
\pgfpathlineto{\pgfqpoint{4.604394in}{0.633527in}}%
\pgfpathlineto{\pgfqpoint{4.604727in}{0.625278in}}%
\pgfpathlineto{\pgfqpoint{4.605401in}{0.625278in}}%
\pgfpathlineto{\pgfqpoint{4.606571in}{0.633527in}}%
\pgfpathlineto{\pgfqpoint{4.606571in}{0.625278in}}%
\pgfpathlineto{\pgfqpoint{4.607578in}{0.625278in}}%
\pgfpathlineto{\pgfqpoint{4.608585in}{0.633527in}}%
\pgfpathlineto{\pgfqpoint{4.608585in}{0.625278in}}%
\pgfpathlineto{\pgfqpoint{4.609762in}{0.633527in}}%
\pgfpathlineto{\pgfqpoint{4.609925in}{0.625278in}}%
\pgfpathlineto{\pgfqpoint{4.610598in}{0.625278in}}%
\pgfpathlineto{\pgfqpoint{4.611768in}{0.633527in}}%
\pgfpathlineto{\pgfqpoint{4.611939in}{0.625278in}}%
\pgfpathlineto{\pgfqpoint{4.612775in}{0.625278in}}%
\pgfpathlineto{\pgfqpoint{4.613953in}{0.633527in}}%
\pgfpathlineto{\pgfqpoint{4.614115in}{0.625278in}}%
\pgfpathlineto{\pgfqpoint{4.614952in}{0.625278in}}%
\pgfpathlineto{\pgfqpoint{4.615959in}{0.633527in}}%
\pgfpathlineto{\pgfqpoint{4.615959in}{0.625278in}}%
\pgfpathlineto{\pgfqpoint{4.617299in}{0.633527in}}%
\pgfpathlineto{\pgfqpoint{4.617299in}{0.625278in}}%
\pgfpathlineto{\pgfqpoint{4.618306in}{0.625278in}}%
\pgfpathlineto{\pgfqpoint{4.619483in}{0.633527in}}%
\pgfpathlineto{\pgfqpoint{4.619483in}{0.625278in}}%
\pgfpathlineto{\pgfqpoint{4.620490in}{0.625278in}}%
\pgfpathlineto{\pgfqpoint{4.621660in}{0.633527in}}%
\pgfpathlineto{\pgfqpoint{4.621831in}{0.625278in}}%
\pgfpathlineto{\pgfqpoint{4.622667in}{0.625278in}}%
\pgfpathlineto{\pgfqpoint{4.623674in}{0.633527in}}%
\pgfpathlineto{\pgfqpoint{4.623674in}{0.625278in}}%
\pgfpathlineto{\pgfqpoint{4.625014in}{0.633527in}}%
\pgfpathlineto{\pgfqpoint{4.625014in}{0.625278in}}%
\pgfpathlineto{\pgfqpoint{4.626021in}{0.625278in}}%
\pgfpathlineto{\pgfqpoint{4.627191in}{0.633527in}}%
\pgfpathlineto{\pgfqpoint{4.627362in}{0.625278in}}%
\pgfpathlineto{\pgfqpoint{4.628028in}{0.625278in}}%
\pgfpathlineto{\pgfqpoint{4.629375in}{0.633527in}}%
\pgfpathlineto{\pgfqpoint{4.629375in}{0.625278in}}%
\pgfpathlineto{\pgfqpoint{4.630375in}{0.625278in}}%
\pgfpathlineto{\pgfqpoint{4.631552in}{0.633527in}}%
\pgfpathlineto{\pgfqpoint{4.631715in}{0.625278in}}%
\pgfpathlineto{\pgfqpoint{4.632559in}{0.625278in}}%
\pgfpathlineto{\pgfqpoint{4.633729in}{0.633527in}}%
\pgfpathlineto{\pgfqpoint{4.633899in}{0.625278in}}%
\pgfpathlineto{\pgfqpoint{4.634758in}{0.650025in}}%
\pgfpathlineto{\pgfqpoint{4.635425in}{0.625278in}}%
\pgfpathlineto{\pgfqpoint{4.635595in}{0.650025in}}%
\pgfpathlineto{\pgfqpoint{4.636750in}{0.658274in}}%
\pgfpathlineto{\pgfqpoint{4.637772in}{0.650025in}}%
\pgfpathlineto{\pgfqpoint{4.639430in}{0.683022in}}%
\pgfpathlineto{\pgfqpoint{4.640622in}{0.625278in}}%
\pgfpathlineto{\pgfqpoint{4.644791in}{0.806759in}}%
\pgfpathlineto{\pgfqpoint{4.644813in}{0.798510in}}%
\pgfpathlineto{\pgfqpoint{4.645465in}{0.757264in}}%
\pgfpathlineto{\pgfqpoint{4.645820in}{0.798510in}}%
\pgfpathlineto{\pgfqpoint{4.646139in}{0.806759in}}%
\pgfpathlineto{\pgfqpoint{4.646487in}{0.773763in}}%
\pgfpathlineto{\pgfqpoint{4.646635in}{0.782012in}}%
\pgfpathlineto{\pgfqpoint{4.646975in}{0.806759in}}%
\pgfpathlineto{\pgfqpoint{4.647664in}{0.773763in}}%
\pgfpathlineto{\pgfqpoint{4.648330in}{0.798510in}}%
\pgfpathlineto{\pgfqpoint{4.648819in}{0.782012in}}%
\pgfpathlineto{\pgfqpoint{4.649826in}{0.757264in}}%
\pgfpathlineto{\pgfqpoint{4.649841in}{0.773763in}}%
\pgfpathlineto{\pgfqpoint{4.650848in}{0.749015in}}%
\pgfpathlineto{\pgfqpoint{4.650996in}{0.765513in}}%
\pgfpathlineto{\pgfqpoint{4.651351in}{0.798510in}}%
\pgfpathlineto{\pgfqpoint{4.652018in}{0.749015in}}%
\pgfpathlineto{\pgfqpoint{4.652669in}{0.732517in}}%
\pgfpathlineto{\pgfqpoint{4.652669in}{0.765513in}}%
\pgfpathlineto{\pgfqpoint{4.653024in}{0.798510in}}%
\pgfpathlineto{\pgfqpoint{4.653698in}{0.749015in}}%
\pgfpathlineto{\pgfqpoint{4.654705in}{0.798510in}}%
\pgfpathlineto{\pgfqpoint{4.655357in}{0.782012in}}%
\pgfpathlineto{\pgfqpoint{4.656379in}{0.773763in}}%
\pgfpathlineto{\pgfqpoint{4.657200in}{0.782012in}}%
\pgfpathlineto{\pgfqpoint{4.657215in}{0.773763in}}%
\pgfpathlineto{\pgfqpoint{4.657534in}{0.757264in}}%
\pgfpathlineto{\pgfqpoint{4.657867in}{0.790261in}}%
\pgfpathlineto{\pgfqpoint{4.658874in}{0.806759in}}%
\pgfpathlineto{\pgfqpoint{4.658896in}{0.798510in}}%
\pgfpathlineto{\pgfqpoint{4.659548in}{0.806759in}}%
\pgfpathlineto{\pgfqpoint{4.659733in}{0.790261in}}%
\pgfpathlineto{\pgfqpoint{4.660555in}{0.757264in}}%
\pgfpathlineto{\pgfqpoint{4.660569in}{0.798510in}}%
\pgfpathlineto{\pgfqpoint{4.661221in}{0.856254in}}%
\pgfpathlineto{\pgfqpoint{4.662080in}{0.823258in}}%
\pgfpathlineto{\pgfqpoint{4.662746in}{0.798510in}}%
\pgfpathlineto{\pgfqpoint{4.662916in}{0.823258in}}%
\pgfpathlineto{\pgfqpoint{4.663901in}{0.856254in}}%
\pgfpathlineto{\pgfqpoint{4.664094in}{0.839756in}}%
\pgfpathlineto{\pgfqpoint{4.664930in}{0.823258in}}%
\pgfpathlineto{\pgfqpoint{4.665079in}{0.839756in}}%
\pgfpathlineto{\pgfqpoint{4.665412in}{0.856254in}}%
\pgfpathlineto{\pgfqpoint{4.665767in}{0.823258in}}%
\pgfpathlineto{\pgfqpoint{4.666100in}{0.823258in}}%
\pgfpathlineto{\pgfqpoint{4.668092in}{0.881002in}}%
\pgfpathlineto{\pgfqpoint{4.668447in}{0.872753in}}%
\pgfpathlineto{\pgfqpoint{4.669099in}{0.913998in}}%
\pgfpathlineto{\pgfqpoint{4.669454in}{0.922247in}}%
\pgfpathlineto{\pgfqpoint{4.669788in}{0.897500in}}%
\pgfpathlineto{\pgfqpoint{4.670128in}{0.897500in}}%
\pgfpathlineto{\pgfqpoint{4.671446in}{0.930497in}}%
\pgfpathlineto{\pgfqpoint{4.671468in}{0.922247in}}%
\pgfpathlineto{\pgfqpoint{4.672120in}{0.905749in}}%
\pgfpathlineto{\pgfqpoint{4.672283in}{0.930497in}}%
\pgfpathlineto{\pgfqpoint{4.672305in}{0.922247in}}%
\pgfpathlineto{\pgfqpoint{4.672623in}{0.930497in}}%
\pgfpathlineto{\pgfqpoint{4.672971in}{0.897500in}}%
\pgfpathlineto{\pgfqpoint{4.673793in}{0.930497in}}%
\pgfpathlineto{\pgfqpoint{4.674126in}{0.905749in}}%
\pgfpathlineto{\pgfqpoint{4.674319in}{0.897500in}}%
\pgfpathlineto{\pgfqpoint{4.674971in}{0.979992in}}%
\pgfpathlineto{\pgfqpoint{4.674985in}{0.971742in}}%
\pgfpathlineto{\pgfqpoint{4.675133in}{0.979992in}}%
\pgfpathlineto{\pgfqpoint{4.675489in}{0.946995in}}%
\pgfpathlineto{\pgfqpoint{4.675637in}{0.955244in}}%
\pgfpathlineto{\pgfqpoint{4.676666in}{0.946995in}}%
\pgfpathlineto{\pgfqpoint{4.678006in}{0.996490in}}%
\pgfpathlineto{\pgfqpoint{4.678317in}{0.979992in}}%
\pgfpathlineto{\pgfqpoint{4.679346in}{0.971742in}}%
\pgfpathlineto{\pgfqpoint{4.680161in}{0.979992in}}%
\pgfpathlineto{\pgfqpoint{4.680183in}{0.971742in}}%
\pgfpathlineto{\pgfqpoint{4.681523in}{0.946995in}}%
\pgfpathlineto{\pgfqpoint{4.682360in}{0.996490in}}%
\pgfpathlineto{\pgfqpoint{4.682863in}{0.963493in}}%
\pgfpathlineto{\pgfqpoint{4.683204in}{0.946995in}}%
\pgfpathlineto{\pgfqpoint{4.683515in}{0.979992in}}%
\pgfpathlineto{\pgfqpoint{4.683537in}{0.971742in}}%
\pgfpathlineto{\pgfqpoint{4.684692in}{1.004739in}}%
\pgfpathlineto{\pgfqpoint{4.684877in}{0.988241in}}%
\pgfpathlineto{\pgfqpoint{4.686217in}{0.946995in}}%
\pgfpathlineto{\pgfqpoint{4.686366in}{0.963493in}}%
\pgfpathlineto{\pgfqpoint{4.686721in}{0.971742in}}%
\pgfpathlineto{\pgfqpoint{4.687054in}{0.946995in}}%
\pgfpathlineto{\pgfqpoint{4.687395in}{0.946995in}}%
\pgfpathlineto{\pgfqpoint{4.689046in}{0.979992in}}%
\pgfpathlineto{\pgfqpoint{4.690075in}{0.946995in}}%
\pgfpathlineto{\pgfqpoint{4.690223in}{0.955244in}}%
\pgfpathlineto{\pgfqpoint{4.690386in}{0.979992in}}%
\pgfpathlineto{\pgfqpoint{4.691082in}{0.946995in}}%
\pgfpathlineto{\pgfqpoint{4.691245in}{0.946995in}}%
\pgfpathlineto{\pgfqpoint{4.691897in}{0.955244in}}%
\pgfpathlineto{\pgfqpoint{4.692082in}{0.938746in}}%
\pgfpathlineto{\pgfqpoint{4.692237in}{0.930497in}}%
\pgfpathlineto{\pgfqpoint{4.692400in}{0.955244in}}%
\pgfpathlineto{\pgfqpoint{4.692733in}{0.955244in}}%
\pgfpathlineto{\pgfqpoint{4.692755in}{0.971742in}}%
\pgfpathlineto{\pgfqpoint{4.693429in}{0.946995in}}%
\pgfpathlineto{\pgfqpoint{4.693762in}{0.946995in}}%
\pgfpathlineto{\pgfqpoint{4.694917in}{0.955244in}}%
\pgfpathlineto{\pgfqpoint{4.695939in}{0.922247in}}%
\pgfpathlineto{\pgfqpoint{4.696776in}{0.897500in}}%
\pgfpathlineto{\pgfqpoint{4.697265in}{0.905749in}}%
\pgfpathlineto{\pgfqpoint{4.697279in}{0.922247in}}%
\pgfpathlineto{\pgfqpoint{4.698286in}{0.922247in}}%
\pgfpathlineto{\pgfqpoint{4.698938in}{0.905749in}}%
\pgfpathlineto{\pgfqpoint{4.699108in}{0.930497in}}%
\pgfpathlineto{\pgfqpoint{4.699123in}{0.922247in}}%
\pgfpathlineto{\pgfqpoint{4.700782in}{0.955244in}}%
\pgfpathlineto{\pgfqpoint{4.701470in}{0.946995in}}%
\pgfpathlineto{\pgfqpoint{4.701640in}{0.971742in}}%
\pgfpathlineto{\pgfqpoint{4.701811in}{0.963493in}}%
\pgfpathlineto{\pgfqpoint{4.701974in}{0.971742in}}%
\pgfpathlineto{\pgfqpoint{4.702958in}{0.955244in}}%
\pgfpathlineto{\pgfqpoint{4.703802in}{1.004739in}}%
\pgfpathlineto{\pgfqpoint{4.703988in}{0.996490in}}%
\pgfpathlineto{\pgfqpoint{4.704469in}{1.004739in}}%
\pgfpathlineto{\pgfqpoint{4.704654in}{0.988241in}}%
\pgfpathlineto{\pgfqpoint{4.704809in}{0.979992in}}%
\pgfpathlineto{\pgfqpoint{4.705143in}{1.012988in}}%
\pgfpathlineto{\pgfqpoint{4.705476in}{1.029487in}}%
\pgfpathlineto{\pgfqpoint{4.706164in}{0.971742in}}%
\pgfpathlineto{\pgfqpoint{4.707675in}{1.021237in}}%
\pgfpathlineto{\pgfqpoint{4.707845in}{1.012988in}}%
\pgfpathlineto{\pgfqpoint{4.708349in}{0.996490in}}%
\pgfpathlineto{\pgfqpoint{4.708660in}{1.029487in}}%
\pgfpathlineto{\pgfqpoint{4.708682in}{1.021237in}}%
\pgfpathlineto{\pgfqpoint{4.710340in}{1.054234in}}%
\pgfpathlineto{\pgfqpoint{4.710859in}{1.021237in}}%
\pgfpathlineto{\pgfqpoint{4.711340in}{1.078981in}}%
\pgfpathlineto{\pgfqpoint{4.711362in}{1.070732in}}%
\pgfpathlineto{\pgfqpoint{4.712036in}{1.120227in}}%
\pgfpathlineto{\pgfqpoint{4.712539in}{1.087231in}}%
\pgfpathlineto{\pgfqpoint{4.713709in}{1.070732in}}%
\pgfpathlineto{\pgfqpoint{4.714864in}{1.128476in}}%
\pgfpathlineto{\pgfqpoint{4.715368in}{1.103729in}}%
\pgfpathlineto{\pgfqpoint{4.715553in}{1.095480in}}%
\pgfpathlineto{\pgfqpoint{4.716390in}{1.144975in}}%
\pgfpathlineto{\pgfqpoint{4.717545in}{1.153224in}}%
\pgfpathlineto{\pgfqpoint{4.717567in}{1.144975in}}%
\pgfpathlineto{\pgfqpoint{4.718218in}{1.202719in}}%
\pgfpathlineto{\pgfqpoint{4.718404in}{1.194470in}}%
\pgfpathlineto{\pgfqpoint{4.719559in}{1.202719in}}%
\pgfpathlineto{\pgfqpoint{4.720580in}{1.186221in}}%
\pgfpathlineto{\pgfqpoint{4.720921in}{1.169722in}}%
\pgfpathlineto{\pgfqpoint{4.721232in}{1.202719in}}%
\pgfpathlineto{\pgfqpoint{4.721254in}{1.194470in}}%
\pgfpathlineto{\pgfqpoint{4.722409in}{1.202719in}}%
\pgfpathlineto{\pgfqpoint{4.723416in}{1.153224in}}%
\pgfpathlineto{\pgfqpoint{4.723431in}{1.169722in}}%
\pgfpathlineto{\pgfqpoint{4.724416in}{1.153224in}}%
\pgfpathlineto{\pgfqpoint{4.724416in}{1.161473in}}%
\pgfpathlineto{\pgfqpoint{4.724919in}{1.202719in}}%
\pgfpathlineto{\pgfqpoint{4.725445in}{1.194470in}}%
\pgfpathlineto{\pgfqpoint{4.725948in}{1.243965in}}%
\pgfpathlineto{\pgfqpoint{4.726785in}{1.210968in}}%
\pgfpathlineto{\pgfqpoint{4.726933in}{1.202719in}}%
\pgfpathlineto{\pgfqpoint{4.726955in}{1.243965in}}%
\pgfpathlineto{\pgfqpoint{4.727451in}{1.235715in}}%
\pgfpathlineto{\pgfqpoint{4.728125in}{1.219217in}}%
\pgfpathlineto{\pgfqpoint{4.728296in}{1.243965in}}%
\pgfpathlineto{\pgfqpoint{4.728777in}{1.252214in}}%
\pgfpathlineto{\pgfqpoint{4.729110in}{1.227466in}}%
\pgfpathlineto{\pgfqpoint{4.730117in}{1.202719in}}%
\pgfpathlineto{\pgfqpoint{4.730139in}{1.219217in}}%
\pgfpathlineto{\pgfqpoint{4.730620in}{1.252214in}}%
\pgfpathlineto{\pgfqpoint{4.731294in}{1.227466in}}%
\pgfpathlineto{\pgfqpoint{4.732316in}{1.210968in}}%
\pgfpathlineto{\pgfqpoint{4.734330in}{1.144975in}}%
\pgfpathlineto{\pgfqpoint{4.734478in}{1.153224in}}%
\pgfpathlineto{\pgfqpoint{4.734833in}{1.120227in}}%
\pgfpathlineto{\pgfqpoint{4.735167in}{1.136726in}}%
\pgfpathlineto{\pgfqpoint{4.736655in}{1.103729in}}%
\pgfpathlineto{\pgfqpoint{4.735337in}{1.144975in}}%
\pgfpathlineto{\pgfqpoint{4.736655in}{1.111978in}}%
\pgfpathlineto{\pgfqpoint{4.737662in}{1.128476in}}%
\pgfpathlineto{\pgfqpoint{4.737684in}{1.120227in}}%
\pgfpathlineto{\pgfqpoint{4.738336in}{1.128476in}}%
\pgfpathlineto{\pgfqpoint{4.738521in}{1.111978in}}%
\pgfpathlineto{\pgfqpoint{4.739024in}{1.095480in}}%
\pgfpathlineto{\pgfqpoint{4.739335in}{1.128476in}}%
\pgfpathlineto{\pgfqpoint{4.739357in}{1.120227in}}%
\pgfpathlineto{\pgfqpoint{4.739676in}{1.153224in}}%
\pgfpathlineto{\pgfqpoint{4.740024in}{1.111978in}}%
\pgfpathlineto{\pgfqpoint{4.741031in}{1.095480in}}%
\pgfpathlineto{\pgfqpoint{4.740675in}{1.128476in}}%
\pgfpathlineto{\pgfqpoint{4.741179in}{1.103729in}}%
\pgfpathlineto{\pgfqpoint{4.741357in}{1.128476in}}%
\pgfpathlineto{\pgfqpoint{4.741875in}{1.095480in}}%
\pgfpathlineto{\pgfqpoint{4.742208in}{1.095480in}}%
\pgfpathlineto{\pgfqpoint{4.742860in}{1.103729in}}%
\pgfpathlineto{\pgfqpoint{4.743045in}{1.087231in}}%
\pgfpathlineto{\pgfqpoint{4.744215in}{1.045985in}}%
\pgfpathlineto{\pgfqpoint{4.744703in}{1.054234in}}%
\pgfpathlineto{\pgfqpoint{4.744718in}{1.070732in}}%
\pgfpathlineto{\pgfqpoint{4.745370in}{1.029487in}}%
\pgfpathlineto{\pgfqpoint{4.745725in}{1.045985in}}%
\pgfpathlineto{\pgfqpoint{4.746214in}{1.054234in}}%
\pgfpathlineto{\pgfqpoint{4.746399in}{1.037736in}}%
\pgfpathlineto{\pgfqpoint{4.746562in}{1.045985in}}%
\pgfpathlineto{\pgfqpoint{4.747569in}{1.021237in}}%
\pgfpathlineto{\pgfqpoint{4.748057in}{1.054234in}}%
\pgfpathlineto{\pgfqpoint{4.748724in}{1.029487in}}%
\pgfpathlineto{\pgfqpoint{4.749916in}{0.963493in}}%
\pgfpathlineto{\pgfqpoint{4.751908in}{0.905749in}}%
\pgfpathlineto{\pgfqpoint{4.750086in}{0.971742in}}%
\pgfpathlineto{\pgfqpoint{4.751908in}{0.913998in}}%
\pgfpathlineto{\pgfqpoint{4.752248in}{0.905749in}}%
\pgfpathlineto{\pgfqpoint{4.752937in}{0.922247in}}%
\pgfpathlineto{\pgfqpoint{4.753773in}{0.897500in}}%
\pgfpathlineto{\pgfqpoint{4.753944in}{0.922247in}}%
\pgfpathlineto{\pgfqpoint{4.754092in}{0.930497in}}%
\pgfpathlineto{\pgfqpoint{4.754447in}{0.897500in}}%
\pgfpathlineto{\pgfqpoint{4.755432in}{0.881002in}}%
\pgfpathlineto{\pgfqpoint{4.754928in}{0.905749in}}%
\pgfpathlineto{\pgfqpoint{4.755447in}{0.897500in}}%
\pgfpathlineto{\pgfqpoint{4.756269in}{0.905749in}}%
\pgfpathlineto{\pgfqpoint{4.756291in}{0.897500in}}%
\pgfpathlineto{\pgfqpoint{4.756602in}{0.881002in}}%
\pgfpathlineto{\pgfqpoint{4.756624in}{0.922247in}}%
\pgfpathlineto{\pgfqpoint{4.756950in}{0.913998in}}%
\pgfpathlineto{\pgfqpoint{4.756957in}{0.922247in}}%
\pgfpathlineto{\pgfqpoint{4.757794in}{0.872753in}}%
\pgfpathlineto{\pgfqpoint{4.757964in}{0.872753in}}%
\pgfpathlineto{\pgfqpoint{4.758460in}{0.881002in}}%
\pgfpathlineto{\pgfqpoint{4.758638in}{0.864503in}}%
\pgfpathlineto{\pgfqpoint{4.759134in}{0.872753in}}%
\pgfpathlineto{\pgfqpoint{4.759808in}{0.848005in}}%
\pgfpathlineto{\pgfqpoint{4.760474in}{0.881002in}}%
\pgfpathlineto{\pgfqpoint{4.760978in}{0.864503in}}%
\pgfpathlineto{\pgfqpoint{4.761800in}{0.831507in}}%
\pgfpathlineto{\pgfqpoint{4.761148in}{0.872753in}}%
\pgfpathlineto{\pgfqpoint{4.762133in}{0.856254in}}%
\pgfpathlineto{\pgfqpoint{4.763147in}{0.881002in}}%
\pgfpathlineto{\pgfqpoint{4.763162in}{0.872753in}}%
\pgfpathlineto{\pgfqpoint{4.763665in}{0.897500in}}%
\pgfpathlineto{\pgfqpoint{4.763999in}{0.872753in}}%
\pgfpathlineto{\pgfqpoint{4.766175in}{0.773763in}}%
\pgfpathlineto{\pgfqpoint{4.766664in}{0.790261in}}%
\pgfpathlineto{\pgfqpoint{4.767182in}{0.782012in}}%
\pgfpathlineto{\pgfqpoint{4.767686in}{0.823258in}}%
\pgfpathlineto{\pgfqpoint{4.767834in}{0.831507in}}%
\pgfpathlineto{\pgfqpoint{4.768189in}{0.798510in}}%
\pgfpathlineto{\pgfqpoint{4.769855in}{0.757264in}}%
\pgfpathlineto{\pgfqpoint{4.769855in}{0.765513in}}%
\pgfpathlineto{\pgfqpoint{4.770862in}{0.831507in}}%
\pgfpathlineto{\pgfqpoint{4.770870in}{0.823258in}}%
\pgfpathlineto{\pgfqpoint{4.771706in}{0.848005in}}%
\pgfpathlineto{\pgfqpoint{4.772040in}{0.831507in}}%
\pgfpathlineto{\pgfqpoint{4.772706in}{0.806759in}}%
\pgfpathlineto{\pgfqpoint{4.773054in}{0.872753in}}%
\pgfpathlineto{\pgfqpoint{4.774705in}{0.831507in}}%
\pgfpathlineto{\pgfqpoint{4.774712in}{0.839756in}}%
\pgfpathlineto{\pgfqpoint{4.775060in}{0.848005in}}%
\pgfpathlineto{\pgfqpoint{4.775712in}{0.806759in}}%
\pgfpathlineto{\pgfqpoint{4.775734in}{0.831507in}}%
\pgfpathlineto{\pgfqpoint{4.775897in}{0.823258in}}%
\pgfpathlineto{\pgfqpoint{4.776053in}{0.864503in}}%
\pgfpathlineto{\pgfqpoint{4.776571in}{0.897500in}}%
\pgfpathlineto{\pgfqpoint{4.777074in}{0.856254in}}%
\pgfpathlineto{\pgfqpoint{4.777911in}{0.872753in}}%
\pgfpathlineto{\pgfqpoint{4.778081in}{0.864503in}}%
\pgfpathlineto{\pgfqpoint{4.778570in}{0.831507in}}%
\pgfpathlineto{\pgfqpoint{4.778918in}{0.872753in}}%
\pgfpathlineto{\pgfqpoint{4.779407in}{0.881002in}}%
\pgfpathlineto{\pgfqpoint{4.779592in}{0.864503in}}%
\pgfpathlineto{\pgfqpoint{4.780584in}{0.831507in}}%
\pgfpathlineto{\pgfqpoint{4.779755in}{0.872753in}}%
\pgfpathlineto{\pgfqpoint{4.780762in}{0.848005in}}%
\pgfpathlineto{\pgfqpoint{4.781258in}{0.856254in}}%
\pgfpathlineto{\pgfqpoint{4.781435in}{0.839756in}}%
\pgfpathlineto{\pgfqpoint{4.782102in}{0.773763in}}%
\pgfpathlineto{\pgfqpoint{4.783109in}{0.798510in}}%
\pgfpathlineto{\pgfqpoint{4.783442in}{0.806759in}}%
\pgfpathlineto{\pgfqpoint{4.783612in}{0.790261in}}%
\pgfpathlineto{\pgfqpoint{4.785626in}{0.650025in}}%
\pgfpathlineto{\pgfqpoint{4.785959in}{0.625278in}}%
\pgfpathlineto{\pgfqpoint{4.786789in}{0.641776in}}%
\pgfpathlineto{\pgfqpoint{4.786796in}{0.650025in}}%
\pgfpathlineto{\pgfqpoint{4.787300in}{0.625278in}}%
\pgfpathlineto{\pgfqpoint{4.787788in}{0.625278in}}%
\pgfpathlineto{\pgfqpoint{4.788795in}{0.633527in}}%
\pgfpathlineto{\pgfqpoint{4.788795in}{0.625278in}}%
\pgfpathlineto{\pgfqpoint{4.790135in}{0.633527in}}%
\pgfpathlineto{\pgfqpoint{4.790320in}{0.625278in}}%
\pgfpathlineto{\pgfqpoint{4.790979in}{0.625278in}}%
\pgfpathlineto{\pgfqpoint{4.792164in}{0.633527in}}%
\pgfpathlineto{\pgfqpoint{4.792497in}{0.625278in}}%
\pgfpathlineto{\pgfqpoint{4.793156in}{0.625278in}}%
\pgfpathlineto{\pgfqpoint{4.794163in}{0.633527in}}%
\pgfpathlineto{\pgfqpoint{4.794163in}{0.625278in}}%
\pgfpathlineto{\pgfqpoint{4.795511in}{0.633527in}}%
\pgfpathlineto{\pgfqpoint{4.795829in}{0.625278in}}%
\pgfpathlineto{\pgfqpoint{4.796347in}{0.625278in}}%
\pgfpathlineto{\pgfqpoint{4.797680in}{0.633527in}}%
\pgfpathlineto{\pgfqpoint{4.797695in}{0.625278in}}%
\pgfpathlineto{\pgfqpoint{4.798702in}{0.625278in}}%
\pgfpathlineto{\pgfqpoint{4.799864in}{0.633527in}}%
\pgfpathlineto{\pgfqpoint{4.800035in}{0.625278in}}%
\pgfpathlineto{\pgfqpoint{4.800871in}{0.625278in}}%
\pgfpathlineto{\pgfqpoint{4.801871in}{0.633527in}}%
\pgfpathlineto{\pgfqpoint{4.801871in}{0.625278in}}%
\pgfpathlineto{\pgfqpoint{4.803048in}{0.633527in}}%
\pgfpathlineto{\pgfqpoint{4.803219in}{0.625278in}}%
\pgfpathlineto{\pgfqpoint{4.803892in}{0.625278in}}%
\pgfpathlineto{\pgfqpoint{4.805070in}{0.633527in}}%
\pgfpathlineto{\pgfqpoint{4.805233in}{0.625278in}}%
\pgfpathlineto{\pgfqpoint{4.806077in}{0.625278in}}%
\pgfpathlineto{\pgfqpoint{4.807239in}{0.633527in}}%
\pgfpathlineto{\pgfqpoint{4.807246in}{0.625278in}}%
\pgfpathlineto{\pgfqpoint{4.808253in}{0.625278in}}%
\pgfpathlineto{\pgfqpoint{4.809246in}{0.633527in}}%
\pgfpathlineto{\pgfqpoint{4.809246in}{0.625278in}}%
\pgfpathlineto{\pgfqpoint{4.810415in}{0.633527in}}%
\pgfpathlineto{\pgfqpoint{4.810749in}{0.625278in}}%
\pgfpathlineto{\pgfqpoint{4.811422in}{0.625278in}}%
\pgfpathlineto{\pgfqpoint{4.812600in}{0.633527in}}%
\pgfpathlineto{\pgfqpoint{4.812763in}{0.625278in}}%
\pgfpathlineto{\pgfqpoint{4.813607in}{0.625278in}}%
\pgfpathlineto{\pgfqpoint{4.814939in}{0.633527in}}%
\pgfpathlineto{\pgfqpoint{4.815110in}{0.625278in}}%
\pgfpathlineto{\pgfqpoint{4.815954in}{0.625278in}}%
\pgfpathlineto{\pgfqpoint{4.817116in}{0.633527in}}%
\pgfpathlineto{\pgfqpoint{4.817301in}{0.625278in}}%
\pgfpathlineto{\pgfqpoint{4.818131in}{0.625278in}}%
\pgfpathlineto{\pgfqpoint{4.819300in}{0.633527in}}%
\pgfpathlineto{\pgfqpoint{4.819308in}{0.625278in}}%
\pgfpathlineto{\pgfqpoint{4.820307in}{0.625278in}}%
\pgfpathlineto{\pgfqpoint{4.821485in}{0.633527in}}%
\pgfpathlineto{\pgfqpoint{4.821811in}{0.625278in}}%
\pgfpathlineto{\pgfqpoint{4.822484in}{0.625278in}}%
\pgfpathlineto{\pgfqpoint{4.823654in}{0.633527in}}%
\pgfpathlineto{\pgfqpoint{4.823654in}{0.625278in}}%
\pgfpathlineto{\pgfqpoint{4.824491in}{0.625278in}}%
\pgfpathlineto{\pgfqpoint{4.825838in}{0.633527in}}%
\pgfpathlineto{\pgfqpoint{4.825853in}{0.625278in}}%
\pgfpathlineto{\pgfqpoint{4.826845in}{0.625278in}}%
\pgfpathlineto{\pgfqpoint{4.827845in}{0.633527in}}%
\pgfpathlineto{\pgfqpoint{4.827845in}{0.625278in}}%
\pgfpathlineto{\pgfqpoint{4.829185in}{0.633527in}}%
\pgfpathlineto{\pgfqpoint{4.829185in}{0.625278in}}%
\pgfpathlineto{\pgfqpoint{4.830192in}{0.625278in}}%
\pgfpathlineto{\pgfqpoint{4.831199in}{0.633527in}}%
\pgfpathlineto{\pgfqpoint{4.831199in}{0.625278in}}%
\pgfpathlineto{\pgfqpoint{4.832376in}{0.633527in}}%
\pgfpathlineto{\pgfqpoint{4.832539in}{0.625278in}}%
\pgfpathlineto{\pgfqpoint{4.833376in}{0.625278in}}%
\pgfpathlineto{\pgfqpoint{4.834553in}{0.633527in}}%
\pgfpathlineto{\pgfqpoint{4.834723in}{0.625278in}}%
\pgfpathlineto{\pgfqpoint{4.835560in}{0.625278in}}%
\pgfpathlineto{\pgfqpoint{4.836730in}{0.633527in}}%
\pgfpathlineto{\pgfqpoint{4.836900in}{0.625278in}}%
\pgfpathlineto{\pgfqpoint{4.837567in}{0.625278in}}%
\pgfpathlineto{\pgfqpoint{4.838914in}{0.633527in}}%
\pgfpathlineto{\pgfqpoint{4.838914in}{0.625278in}}%
\pgfpathlineto{\pgfqpoint{4.839914in}{0.625278in}}%
\pgfpathlineto{\pgfqpoint{4.841091in}{0.633527in}}%
\pgfpathlineto{\pgfqpoint{4.841424in}{0.625278in}}%
\pgfpathlineto{\pgfqpoint{4.842098in}{0.625278in}}%
\pgfpathlineto{\pgfqpoint{4.843268in}{0.633527in}}%
\pgfpathlineto{\pgfqpoint{4.843438in}{0.625278in}}%
\pgfpathlineto{\pgfqpoint{4.844105in}{0.625278in}}%
\pgfpathlineto{\pgfqpoint{4.845282in}{0.633527in}}%
\pgfpathlineto{\pgfqpoint{4.845615in}{0.625278in}}%
\pgfpathlineto{\pgfqpoint{4.846289in}{0.625278in}}%
\pgfpathlineto{\pgfqpoint{4.847459in}{0.633527in}}%
\pgfpathlineto{\pgfqpoint{4.847629in}{0.625278in}}%
\pgfpathlineto{\pgfqpoint{4.848466in}{0.625278in}}%
\pgfpathlineto{\pgfqpoint{4.849635in}{0.633527in}}%
\pgfpathlineto{\pgfqpoint{4.849806in}{0.625278in}}%
\pgfpathlineto{\pgfqpoint{4.850642in}{0.625278in}}%
\pgfpathlineto{\pgfqpoint{4.851820in}{0.633527in}}%
\pgfpathlineto{\pgfqpoint{4.851983in}{0.625278in}}%
\pgfpathlineto{\pgfqpoint{4.852827in}{0.625278in}}%
\pgfpathlineto{\pgfqpoint{4.853997in}{0.633527in}}%
\pgfpathlineto{\pgfqpoint{4.854167in}{0.625278in}}%
\pgfpathlineto{\pgfqpoint{4.854833in}{0.625278in}}%
\pgfpathlineto{\pgfqpoint{4.856173in}{0.633527in}}%
\pgfpathlineto{\pgfqpoint{4.856344in}{0.625278in}}%
\pgfpathlineto{\pgfqpoint{4.857188in}{0.625278in}}%
\pgfpathlineto{\pgfqpoint{4.858187in}{0.633527in}}%
\pgfpathlineto{\pgfqpoint{4.858187in}{0.625278in}}%
\pgfpathlineto{\pgfqpoint{4.859365in}{0.633527in}}%
\pgfpathlineto{\pgfqpoint{4.859527in}{0.625278in}}%
\pgfpathlineto{\pgfqpoint{4.860364in}{0.625278in}}%
\pgfpathlineto{\pgfqpoint{4.861541in}{0.633527in}}%
\pgfpathlineto{\pgfqpoint{4.861712in}{0.625278in}}%
\pgfpathlineto{\pgfqpoint{4.862548in}{0.625278in}}%
\pgfpathlineto{\pgfqpoint{4.863718in}{0.633527in}}%
\pgfpathlineto{\pgfqpoint{4.863889in}{0.625278in}}%
\pgfpathlineto{\pgfqpoint{4.864725in}{0.625278in}}%
\pgfpathlineto{\pgfqpoint{4.865732in}{0.633527in}}%
\pgfpathlineto{\pgfqpoint{4.865732in}{0.625278in}}%
\pgfpathlineto{\pgfqpoint{4.866902in}{0.633527in}}%
\pgfpathlineto{\pgfqpoint{4.867072in}{0.625278in}}%
\pgfpathlineto{\pgfqpoint{4.867746in}{0.625278in}}%
\pgfpathlineto{\pgfqpoint{4.869086in}{0.633527in}}%
\pgfpathlineto{\pgfqpoint{4.869086in}{0.625278in}}%
\pgfpathlineto{\pgfqpoint{4.870093in}{0.625278in}}%
\pgfpathlineto{\pgfqpoint{4.871263in}{0.633527in}}%
\pgfpathlineto{\pgfqpoint{4.871433in}{0.625278in}}%
\pgfpathlineto{\pgfqpoint{4.872270in}{0.625278in}}%
\pgfpathlineto{\pgfqpoint{4.873277in}{0.633527in}}%
\pgfpathlineto{\pgfqpoint{4.873277in}{0.625278in}}%
\pgfpathlineto{\pgfqpoint{4.874447in}{0.633527in}}%
\pgfpathlineto{\pgfqpoint{4.874617in}{0.625278in}}%
\pgfpathlineto{\pgfqpoint{4.875454in}{0.625278in}}%
\pgfpathlineto{\pgfqpoint{4.876631in}{0.633527in}}%
\pgfpathlineto{\pgfqpoint{4.876631in}{0.625278in}}%
\pgfpathlineto{\pgfqpoint{4.877631in}{0.625278in}}%
\pgfpathlineto{\pgfqpoint{4.878808in}{0.633527in}}%
\pgfpathlineto{\pgfqpoint{4.878971in}{0.625278in}}%
\pgfpathlineto{\pgfqpoint{4.879645in}{0.625278in}}%
\pgfpathlineto{\pgfqpoint{4.880822in}{0.633527in}}%
\pgfpathlineto{\pgfqpoint{4.880985in}{0.625278in}}%
\pgfpathlineto{\pgfqpoint{4.881659in}{0.625278in}}%
\pgfpathlineto{\pgfqpoint{4.882828in}{0.633527in}}%
\pgfpathlineto{\pgfqpoint{4.882999in}{0.625278in}}%
\pgfpathlineto{\pgfqpoint{4.883835in}{0.625278in}}%
\pgfpathlineto{\pgfqpoint{4.885013in}{0.633527in}}%
\pgfpathlineto{\pgfqpoint{4.885176in}{0.625278in}}%
\pgfpathlineto{\pgfqpoint{4.886012in}{0.625278in}}%
\pgfpathlineto{\pgfqpoint{4.887190in}{0.633527in}}%
\pgfpathlineto{\pgfqpoint{4.887352in}{0.625278in}}%
\pgfpathlineto{\pgfqpoint{4.888197in}{0.625278in}}%
\pgfpathlineto{\pgfqpoint{4.889366in}{0.633527in}}%
\pgfpathlineto{\pgfqpoint{4.889700in}{0.625278in}}%
\pgfpathlineto{\pgfqpoint{4.890373in}{0.625278in}}%
\pgfpathlineto{\pgfqpoint{4.891714in}{0.633527in}}%
\pgfpathlineto{\pgfqpoint{4.891714in}{0.625278in}}%
\pgfpathlineto{\pgfqpoint{4.892720in}{0.625278in}}%
\pgfpathlineto{\pgfqpoint{4.893890in}{0.633527in}}%
\pgfpathlineto{\pgfqpoint{4.894061in}{0.625278in}}%
\pgfpathlineto{\pgfqpoint{4.894897in}{0.625278in}}%
\pgfpathlineto{\pgfqpoint{4.896237in}{0.633527in}}%
\pgfpathlineto{\pgfqpoint{4.896237in}{0.625278in}}%
\pgfpathlineto{\pgfqpoint{4.897244in}{0.625278in}}%
\pgfpathlineto{\pgfqpoint{4.898251in}{0.633527in}}%
\pgfpathlineto{\pgfqpoint{4.898251in}{0.625278in}}%
\pgfpathlineto{\pgfqpoint{4.899429in}{0.633527in}}%
\pgfpathlineto{\pgfqpoint{4.899592in}{0.625278in}}%
\pgfpathlineto{\pgfqpoint{4.900428in}{0.625278in}}%
\pgfpathlineto{\pgfqpoint{4.901606in}{0.633527in}}%
\pgfpathlineto{\pgfqpoint{4.901776in}{0.625278in}}%
\pgfpathlineto{\pgfqpoint{4.902612in}{0.625278in}}%
\pgfpathlineto{\pgfqpoint{4.903782in}{0.633527in}}%
\pgfpathlineto{\pgfqpoint{4.903842in}{0.625278in}}%
\pgfpathlineto{\pgfqpoint{4.904789in}{0.625278in}}%
\pgfpathlineto{\pgfqpoint{4.905967in}{0.633527in}}%
\pgfpathlineto{\pgfqpoint{4.905967in}{0.625278in}}%
\pgfpathlineto{\pgfqpoint{4.906966in}{0.625278in}}%
\pgfpathlineto{\pgfqpoint{4.908143in}{0.633527in}}%
\pgfpathlineto{\pgfqpoint{4.908306in}{0.625278in}}%
\pgfpathlineto{\pgfqpoint{4.909150in}{0.625278in}}%
\pgfpathlineto{\pgfqpoint{4.910320in}{0.633527in}}%
\pgfpathlineto{\pgfqpoint{4.910491in}{0.625278in}}%
\pgfpathlineto{\pgfqpoint{4.911327in}{0.625278in}}%
\pgfpathlineto{\pgfqpoint{4.912334in}{0.633527in}}%
\pgfpathlineto{\pgfqpoint{4.912334in}{0.625278in}}%
\pgfpathlineto{\pgfqpoint{4.913504in}{0.633527in}}%
\pgfpathlineto{\pgfqpoint{4.913674in}{0.625278in}}%
\pgfpathlineto{\pgfqpoint{4.914348in}{0.625278in}}%
\pgfpathlineto{\pgfqpoint{4.915688in}{0.633527in}}%
\pgfpathlineto{\pgfqpoint{4.915688in}{0.625278in}}%
\pgfpathlineto{\pgfqpoint{4.916688in}{0.625278in}}%
\pgfpathlineto{\pgfqpoint{4.917695in}{0.633527in}}%
\pgfpathlineto{\pgfqpoint{4.917695in}{0.625278in}}%
\pgfpathlineto{\pgfqpoint{4.918872in}{0.633527in}}%
\pgfpathlineto{\pgfqpoint{4.919035in}{0.625278in}}%
\pgfpathlineto{\pgfqpoint{4.919879in}{0.625278in}}%
\pgfpathlineto{\pgfqpoint{4.921049in}{0.633527in}}%
\pgfpathlineto{\pgfqpoint{4.921382in}{0.625278in}}%
\pgfpathlineto{\pgfqpoint{4.922056in}{0.625278in}}%
\pgfpathlineto{\pgfqpoint{4.923226in}{0.633527in}}%
\pgfpathlineto{\pgfqpoint{4.923396in}{0.625278in}}%
\pgfpathlineto{\pgfqpoint{4.924233in}{0.625278in}}%
\pgfpathlineto{\pgfqpoint{4.925410in}{0.633527in}}%
\pgfpathlineto{\pgfqpoint{4.925410in}{0.625278in}}%
\pgfpathlineto{\pgfqpoint{4.926417in}{0.625278in}}%
\pgfpathlineto{\pgfqpoint{4.927757in}{0.633527in}}%
\pgfpathlineto{\pgfqpoint{4.928090in}{0.625278in}}%
\pgfpathlineto{\pgfqpoint{4.928779in}{0.625278in}}%
\pgfpathlineto{\pgfqpoint{4.929934in}{0.633527in}}%
\pgfpathlineto{\pgfqpoint{4.930956in}{0.625278in}}%
\pgfpathlineto{\pgfqpoint{4.931444in}{0.658274in}}%
\pgfpathlineto{\pgfqpoint{4.932111in}{0.633527in}}%
\pgfpathlineto{\pgfqpoint{4.932133in}{0.625278in}}%
\pgfpathlineto{\pgfqpoint{4.932614in}{0.683022in}}%
\pgfpathlineto{\pgfqpoint{4.932970in}{0.674773in}}%
\pgfpathlineto{\pgfqpoint{4.933954in}{0.658274in}}%
\pgfpathlineto{\pgfqpoint{4.933977in}{0.674773in}}%
\pgfpathlineto{\pgfqpoint{4.935132in}{0.683022in}}%
\pgfpathlineto{\pgfqpoint{4.935650in}{0.674773in}}%
\pgfpathlineto{\pgfqpoint{4.935968in}{0.707769in}}%
\pgfpathlineto{\pgfqpoint{4.936153in}{0.699520in}}%
\pgfpathlineto{\pgfqpoint{4.936472in}{0.707769in}}%
\pgfpathlineto{\pgfqpoint{4.936805in}{0.683022in}}%
\pgfpathlineto{\pgfqpoint{4.936990in}{0.691271in}}%
\pgfpathlineto{\pgfqpoint{4.937982in}{0.683022in}}%
\pgfpathlineto{\pgfqpoint{4.937309in}{0.707769in}}%
\pgfpathlineto{\pgfqpoint{4.937982in}{0.691271in}}%
\pgfpathlineto{\pgfqpoint{4.938989in}{0.707769in}}%
\pgfpathlineto{\pgfqpoint{4.938316in}{0.683022in}}%
\pgfpathlineto{\pgfqpoint{4.939004in}{0.699520in}}%
\pgfpathlineto{\pgfqpoint{4.940159in}{0.707769in}}%
\pgfpathlineto{\pgfqpoint{4.941181in}{0.699520in}}%
\pgfpathlineto{\pgfqpoint{4.942173in}{0.707769in}}%
\pgfpathlineto{\pgfqpoint{4.942188in}{0.699520in}}%
\pgfpathlineto{\pgfqpoint{4.942691in}{0.674773in}}%
\pgfpathlineto{\pgfqpoint{4.943010in}{0.707769in}}%
\pgfpathlineto{\pgfqpoint{4.943032in}{0.699520in}}%
\pgfpathlineto{\pgfqpoint{4.943180in}{0.707769in}}%
\pgfpathlineto{\pgfqpoint{4.943528in}{0.666524in}}%
\pgfpathlineto{\pgfqpoint{4.945024in}{0.633527in}}%
\pgfpathlineto{\pgfqpoint{4.945024in}{0.641776in}}%
\pgfpathlineto{\pgfqpoint{4.945357in}{0.633527in}}%
\pgfpathlineto{\pgfqpoint{4.946194in}{0.683022in}}%
\pgfpathlineto{\pgfqpoint{4.947201in}{0.633527in}}%
\pgfpathlineto{\pgfqpoint{4.947386in}{0.650025in}}%
\pgfpathlineto{\pgfqpoint{4.947874in}{0.683022in}}%
\pgfpathlineto{\pgfqpoint{4.948222in}{0.650025in}}%
\pgfpathlineto{\pgfqpoint{4.948541in}{0.633527in}}%
\pgfpathlineto{\pgfqpoint{4.948874in}{0.666524in}}%
\pgfpathlineto{\pgfqpoint{4.949214in}{0.683022in}}%
\pgfpathlineto{\pgfqpoint{4.949562in}{0.650025in}}%
\pgfpathlineto{\pgfqpoint{4.949903in}{0.650025in}}%
\pgfpathlineto{\pgfqpoint{4.950384in}{0.683022in}}%
\pgfpathlineto{\pgfqpoint{4.951058in}{0.658274in}}%
\pgfpathlineto{\pgfqpoint{4.951414in}{0.650025in}}%
\pgfpathlineto{\pgfqpoint{4.951725in}{0.683022in}}%
\pgfpathlineto{\pgfqpoint{4.952080in}{0.666524in}}%
\pgfpathlineto{\pgfqpoint{4.953087in}{0.650025in}}%
\pgfpathlineto{\pgfqpoint{4.953124in}{0.658274in}}%
\pgfpathlineto{\pgfqpoint{4.953420in}{0.699520in}}%
\pgfpathlineto{\pgfqpoint{4.954427in}{0.674773in}}%
\pgfpathlineto{\pgfqpoint{4.955249in}{0.683022in}}%
\pgfpathlineto{\pgfqpoint{4.955264in}{0.674773in}}%
\pgfpathlineto{\pgfqpoint{4.956419in}{0.658274in}}%
\pgfpathlineto{\pgfqpoint{4.955752in}{0.683022in}}%
\pgfpathlineto{\pgfqpoint{4.956419in}{0.666524in}}%
\pgfpathlineto{\pgfqpoint{4.956774in}{0.674773in}}%
\pgfpathlineto{\pgfqpoint{4.957107in}{0.650025in}}%
\pgfpathlineto{\pgfqpoint{4.957448in}{0.650025in}}%
\pgfpathlineto{\pgfqpoint{4.959773in}{0.757264in}}%
\pgfpathlineto{\pgfqpoint{4.957818in}{0.633527in}}%
\pgfpathlineto{\pgfqpoint{4.959958in}{0.740766in}}%
\pgfpathlineto{\pgfqpoint{4.961639in}{0.699520in}}%
\pgfpathlineto{\pgfqpoint{4.962290in}{0.707769in}}%
\pgfpathlineto{\pgfqpoint{4.962475in}{0.691271in}}%
\pgfpathlineto{\pgfqpoint{4.962623in}{0.683022in}}%
\pgfpathlineto{\pgfqpoint{4.962979in}{0.724268in}}%
\pgfpathlineto{\pgfqpoint{4.963482in}{0.749015in}}%
\pgfpathlineto{\pgfqpoint{4.963816in}{0.724268in}}%
\pgfpathlineto{\pgfqpoint{4.964985in}{0.699520in}}%
\pgfpathlineto{\pgfqpoint{4.965134in}{0.707769in}}%
\pgfpathlineto{\pgfqpoint{4.965489in}{0.749015in}}%
\pgfpathlineto{\pgfqpoint{4.966163in}{0.749015in}}%
\pgfpathlineto{\pgfqpoint{4.967147in}{0.732517in}}%
\pgfpathlineto{\pgfqpoint{4.967170in}{0.749015in}}%
\pgfpathlineto{\pgfqpoint{4.967651in}{0.757264in}}%
\pgfpathlineto{\pgfqpoint{4.967984in}{0.732517in}}%
\pgfpathlineto{\pgfqpoint{4.968006in}{0.749015in}}%
\pgfpathlineto{\pgfqpoint{4.968325in}{0.732517in}}%
\pgfpathlineto{\pgfqpoint{4.968340in}{0.773763in}}%
\pgfpathlineto{\pgfqpoint{4.968658in}{0.765513in}}%
\pgfpathlineto{\pgfqpoint{4.969324in}{0.732517in}}%
\pgfpathlineto{\pgfqpoint{4.969680in}{0.773763in}}%
\pgfpathlineto{\pgfqpoint{4.969998in}{0.782012in}}%
\pgfpathlineto{\pgfqpoint{4.970331in}{0.757264in}}%
\pgfpathlineto{\pgfqpoint{4.971360in}{0.749015in}}%
\pgfpathlineto{\pgfqpoint{4.972678in}{0.782012in}}%
\pgfpathlineto{\pgfqpoint{4.972701in}{0.773763in}}%
\pgfpathlineto{\pgfqpoint{4.974019in}{0.732517in}}%
\pgfpathlineto{\pgfqpoint{4.974041in}{0.749015in}}%
\pgfpathlineto{\pgfqpoint{4.975551in}{0.798510in}}%
\pgfpathlineto{\pgfqpoint{4.975714in}{0.790261in}}%
\pgfpathlineto{\pgfqpoint{4.975884in}{0.798510in}}%
\pgfpathlineto{\pgfqpoint{4.976869in}{0.782012in}}%
\pgfpathlineto{\pgfqpoint{4.976891in}{0.798510in}}%
\pgfpathlineto{\pgfqpoint{4.977225in}{0.773763in}}%
\pgfpathlineto{\pgfqpoint{4.977898in}{0.773763in}}%
\pgfpathlineto{\pgfqpoint{4.979401in}{0.823258in}}%
\pgfpathlineto{\pgfqpoint{4.979572in}{0.815008in}}%
\pgfpathlineto{\pgfqpoint{4.981060in}{0.757264in}}%
\pgfpathlineto{\pgfqpoint{4.981245in}{0.773763in}}%
\pgfpathlineto{\pgfqpoint{4.981897in}{0.806759in}}%
\pgfpathlineto{\pgfqpoint{4.982422in}{0.790261in}}%
\pgfpathlineto{\pgfqpoint{4.982904in}{0.782012in}}%
\pgfpathlineto{\pgfqpoint{4.983074in}{0.806759in}}%
\pgfpathlineto{\pgfqpoint{4.983089in}{0.798510in}}%
\pgfpathlineto{\pgfqpoint{4.983577in}{0.806759in}}%
\pgfpathlineto{\pgfqpoint{4.983762in}{0.790261in}}%
\pgfpathlineto{\pgfqpoint{4.984918in}{0.757264in}}%
\pgfpathlineto{\pgfqpoint{4.984940in}{0.773763in}}%
\pgfpathlineto{\pgfqpoint{4.985591in}{0.806759in}}%
\pgfpathlineto{\pgfqpoint{4.986110in}{0.790261in}}%
\pgfpathlineto{\pgfqpoint{4.986258in}{0.782012in}}%
\pgfpathlineto{\pgfqpoint{4.986428in}{0.806759in}}%
\pgfpathlineto{\pgfqpoint{4.986783in}{0.798510in}}%
\pgfpathlineto{\pgfqpoint{4.987953in}{0.823258in}}%
\pgfpathlineto{\pgfqpoint{4.987265in}{0.782012in}}%
\pgfpathlineto{\pgfqpoint{4.988124in}{0.815008in}}%
\pgfpathlineto{\pgfqpoint{4.988286in}{0.823258in}}%
\pgfpathlineto{\pgfqpoint{4.989293in}{0.798510in}}%
\pgfpathlineto{\pgfqpoint{4.992292in}{0.905749in}}%
\pgfpathlineto{\pgfqpoint{4.993484in}{0.848005in}}%
\pgfpathlineto{\pgfqpoint{4.994491in}{0.872753in}}%
\pgfpathlineto{\pgfqpoint{4.994661in}{0.864503in}}%
\pgfpathlineto{\pgfqpoint{4.994972in}{0.881002in}}%
\pgfpathlineto{\pgfqpoint{4.995831in}{0.848005in}}%
\pgfpathlineto{\pgfqpoint{4.996150in}{0.856254in}}%
\pgfpathlineto{\pgfqpoint{4.996483in}{0.831507in}}%
\pgfpathlineto{\pgfqpoint{4.996668in}{0.839756in}}%
\pgfpathlineto{\pgfqpoint{4.996816in}{0.831507in}}%
\pgfpathlineto{\pgfqpoint{4.996838in}{0.872753in}}%
\pgfpathlineto{\pgfqpoint{4.997171in}{0.872753in}}%
\pgfpathlineto{\pgfqpoint{4.997490in}{0.856254in}}%
\pgfpathlineto{\pgfqpoint{4.997823in}{0.889251in}}%
\pgfpathlineto{\pgfqpoint{4.999000in}{0.930497in}}%
\pgfpathlineto{\pgfqpoint{4.998164in}{0.881002in}}%
\pgfpathlineto{\pgfqpoint{4.999015in}{0.922247in}}%
\pgfpathlineto{\pgfqpoint{5.000022in}{0.946995in}}%
\pgfpathlineto{\pgfqpoint{5.000192in}{0.938746in}}%
\pgfpathlineto{\pgfqpoint{5.001703in}{0.897500in}}%
\pgfpathlineto{\pgfqpoint{5.001851in}{0.905749in}}%
\pgfpathlineto{\pgfqpoint{5.002014in}{0.930497in}}%
\pgfpathlineto{\pgfqpoint{5.002873in}{0.922247in}}%
\pgfpathlineto{\pgfqpoint{5.004198in}{0.979992in}}%
\pgfpathlineto{\pgfqpoint{5.004531in}{0.955244in}}%
\pgfpathlineto{\pgfqpoint{5.004887in}{0.946995in}}%
\pgfpathlineto{\pgfqpoint{5.005198in}{0.979992in}}%
\pgfpathlineto{\pgfqpoint{5.005553in}{0.971742in}}%
\pgfpathlineto{\pgfqpoint{5.007212in}{1.004739in}}%
\pgfpathlineto{\pgfqpoint{5.007885in}{0.955244in}}%
\pgfpathlineto{\pgfqpoint{5.008233in}{0.996490in}}%
\pgfpathlineto{\pgfqpoint{5.008885in}{1.029487in}}%
\pgfpathlineto{\pgfqpoint{5.009388in}{1.004739in}}%
\pgfpathlineto{\pgfqpoint{5.009581in}{0.996490in}}%
\pgfpathlineto{\pgfqpoint{5.010232in}{1.054234in}}%
\pgfpathlineto{\pgfqpoint{5.010247in}{1.045985in}}%
\pgfpathlineto{\pgfqpoint{5.010580in}{1.070732in}}%
\pgfpathlineto{\pgfqpoint{5.011084in}{1.045985in}}%
\pgfpathlineto{\pgfqpoint{5.011402in}{1.029487in}}%
\pgfpathlineto{\pgfqpoint{5.011736in}{1.062483in}}%
\pgfpathlineto{\pgfqpoint{5.012239in}{1.078981in}}%
\pgfpathlineto{\pgfqpoint{5.012594in}{1.045985in}}%
\pgfpathlineto{\pgfqpoint{5.012765in}{1.045985in}}%
\pgfpathlineto{\pgfqpoint{5.013935in}{1.070732in}}%
\pgfpathlineto{\pgfqpoint{5.014105in}{1.062483in}}%
\pgfpathlineto{\pgfqpoint{5.015423in}{1.029487in}}%
\pgfpathlineto{\pgfqpoint{5.015445in}{1.045985in}}%
\pgfpathlineto{\pgfqpoint{5.016430in}{1.078981in}}%
\pgfpathlineto{\pgfqpoint{5.016615in}{1.062483in}}%
\pgfpathlineto{\pgfqpoint{5.016785in}{1.070732in}}%
\pgfpathlineto{\pgfqpoint{5.017770in}{1.054234in}}%
\pgfpathlineto{\pgfqpoint{5.017940in}{1.078981in}}%
\pgfpathlineto{\pgfqpoint{5.018799in}{1.062483in}}%
\pgfpathlineto{\pgfqpoint{5.019806in}{1.045985in}}%
\pgfpathlineto{\pgfqpoint{5.019954in}{1.054234in}}%
\pgfpathlineto{\pgfqpoint{5.020117in}{1.078981in}}%
\pgfpathlineto{\pgfqpoint{5.020806in}{1.045985in}}%
\pgfpathlineto{\pgfqpoint{5.020976in}{1.045985in}}%
\pgfpathlineto{\pgfqpoint{5.022464in}{1.078981in}}%
\pgfpathlineto{\pgfqpoint{5.022486in}{1.070732in}}%
\pgfpathlineto{\pgfqpoint{5.023138in}{1.054234in}}%
\pgfpathlineto{\pgfqpoint{5.023308in}{1.078981in}}%
\pgfpathlineto{\pgfqpoint{5.023323in}{1.070732in}}%
\pgfpathlineto{\pgfqpoint{5.024478in}{1.078981in}}%
\pgfpathlineto{\pgfqpoint{5.025670in}{1.045985in}}%
\pgfpathlineto{\pgfqpoint{5.027329in}{1.078981in}}%
\pgfpathlineto{\pgfqpoint{5.028188in}{1.070732in}}%
\pgfpathlineto{\pgfqpoint{5.028351in}{1.095480in}}%
\pgfpathlineto{\pgfqpoint{5.029839in}{1.029487in}}%
\pgfpathlineto{\pgfqpoint{5.030179in}{1.054234in}}%
\pgfpathlineto{\pgfqpoint{5.030194in}{1.070732in}}%
\pgfpathlineto{\pgfqpoint{5.030535in}{1.045985in}}%
\pgfpathlineto{\pgfqpoint{5.031201in}{1.062483in}}%
\pgfpathlineto{\pgfqpoint{5.031520in}{1.078981in}}%
\pgfpathlineto{\pgfqpoint{5.032378in}{1.045985in}}%
\pgfpathlineto{\pgfqpoint{5.033378in}{1.070732in}}%
\pgfpathlineto{\pgfqpoint{5.033548in}{1.062483in}}%
\pgfpathlineto{\pgfqpoint{5.035059in}{1.021237in}}%
\pgfpathlineto{\pgfqpoint{5.034385in}{1.070732in}}%
\pgfpathlineto{\pgfqpoint{5.035096in}{1.029487in}}%
\pgfpathlineto{\pgfqpoint{5.036377in}{1.078981in}}%
\pgfpathlineto{\pgfqpoint{5.036569in}{1.062483in}}%
\pgfpathlineto{\pgfqpoint{5.037073in}{1.045985in}}%
\pgfpathlineto{\pgfqpoint{5.037236in}{1.070732in}}%
\pgfpathlineto{\pgfqpoint{5.037554in}{1.062483in}}%
\pgfpathlineto{\pgfqpoint{5.038561in}{1.078981in}}%
\pgfpathlineto{\pgfqpoint{5.037887in}{1.054234in}}%
\pgfpathlineto{\pgfqpoint{5.038576in}{1.070732in}}%
\pgfpathlineto{\pgfqpoint{5.038894in}{1.078981in}}%
\pgfpathlineto{\pgfqpoint{5.039227in}{1.054234in}}%
\pgfpathlineto{\pgfqpoint{5.039420in}{1.062483in}}%
\pgfpathlineto{\pgfqpoint{5.039901in}{1.054234in}}%
\pgfpathlineto{\pgfqpoint{5.039960in}{1.078981in}}%
\pgfpathlineto{\pgfqpoint{5.040256in}{1.070732in}}%
\pgfpathlineto{\pgfqpoint{5.041071in}{1.078981in}}%
\pgfpathlineto{\pgfqpoint{5.041093in}{1.070732in}}%
\pgfpathlineto{\pgfqpoint{5.042078in}{1.054234in}}%
\pgfpathlineto{\pgfqpoint{5.042100in}{1.070732in}}%
\pgfpathlineto{\pgfqpoint{5.042915in}{1.103729in}}%
\pgfpathlineto{\pgfqpoint{5.043440in}{1.087231in}}%
\pgfpathlineto{\pgfqpoint{5.043588in}{1.078981in}}%
\pgfpathlineto{\pgfqpoint{5.043611in}{1.120227in}}%
\pgfpathlineto{\pgfqpoint{5.043944in}{1.120227in}}%
\pgfpathlineto{\pgfqpoint{5.044610in}{1.144975in}}%
\pgfpathlineto{\pgfqpoint{5.045099in}{1.128476in}}%
\pgfpathlineto{\pgfqpoint{5.045617in}{1.120227in}}%
\pgfpathlineto{\pgfqpoint{5.045935in}{1.153224in}}%
\pgfpathlineto{\pgfqpoint{5.046121in}{1.144975in}}%
\pgfpathlineto{\pgfqpoint{5.047276in}{1.153224in}}%
\pgfpathlineto{\pgfqpoint{5.047298in}{1.144975in}}%
\pgfpathlineto{\pgfqpoint{5.047779in}{1.202719in}}%
\pgfpathlineto{\pgfqpoint{5.048135in}{1.194470in}}%
\pgfpathlineto{\pgfqpoint{5.048616in}{1.202719in}}%
\pgfpathlineto{\pgfqpoint{5.048801in}{1.186221in}}%
\pgfpathlineto{\pgfqpoint{5.049142in}{1.144975in}}%
\pgfpathlineto{\pgfqpoint{5.049978in}{1.169722in}}%
\pgfpathlineto{\pgfqpoint{5.050459in}{1.177971in}}%
\pgfpathlineto{\pgfqpoint{5.050793in}{1.153224in}}%
\pgfpathlineto{\pgfqpoint{5.050815in}{1.169722in}}%
\pgfpathlineto{\pgfqpoint{5.051800in}{1.153224in}}%
\pgfpathlineto{\pgfqpoint{5.051822in}{1.169722in}}%
\pgfpathlineto{\pgfqpoint{5.052473in}{1.177971in}}%
\pgfpathlineto{\pgfqpoint{5.052659in}{1.161473in}}%
\pgfpathlineto{\pgfqpoint{5.054317in}{1.128476in}}%
\pgfpathlineto{\pgfqpoint{5.055176in}{1.169722in}}%
\pgfpathlineto{\pgfqpoint{5.055339in}{1.161473in}}%
\pgfpathlineto{\pgfqpoint{5.056834in}{1.128476in}}%
\pgfpathlineto{\pgfqpoint{5.055509in}{1.169722in}}%
\pgfpathlineto{\pgfqpoint{5.056834in}{1.136726in}}%
\pgfpathlineto{\pgfqpoint{5.057182in}{1.169722in}}%
\pgfpathlineto{\pgfqpoint{5.057168in}{1.128476in}}%
\pgfpathlineto{\pgfqpoint{5.057856in}{1.144975in}}%
\pgfpathlineto{\pgfqpoint{5.058189in}{1.169722in}}%
\pgfpathlineto{\pgfqpoint{5.059174in}{1.128476in}}%
\pgfpathlineto{\pgfqpoint{5.059344in}{1.153224in}}%
\pgfpathlineto{\pgfqpoint{5.060018in}{1.103729in}}%
\pgfpathlineto{\pgfqpoint{5.060203in}{1.136726in}}%
\pgfpathlineto{\pgfqpoint{5.060351in}{1.128476in}}%
\pgfpathlineto{\pgfqpoint{5.060374in}{1.169722in}}%
\pgfpathlineto{\pgfqpoint{5.060870in}{1.161473in}}%
\pgfpathlineto{\pgfqpoint{5.063054in}{1.095480in}}%
\pgfpathlineto{\pgfqpoint{5.063202in}{1.103729in}}%
\pgfpathlineto{\pgfqpoint{5.063557in}{1.070732in}}%
\pgfpathlineto{\pgfqpoint{5.064209in}{1.054234in}}%
\pgfpathlineto{\pgfqpoint{5.064372in}{1.078981in}}%
\pgfpathlineto{\pgfqpoint{5.064394in}{1.070732in}}%
\pgfpathlineto{\pgfqpoint{5.065549in}{1.078981in}}%
\pgfpathlineto{\pgfqpoint{5.065882in}{1.054234in}}%
\pgfpathlineto{\pgfqpoint{5.066571in}{1.095480in}}%
\pgfpathlineto{\pgfqpoint{5.069925in}{0.971742in}}%
\pgfpathlineto{\pgfqpoint{5.070577in}{0.979992in}}%
\pgfpathlineto{\pgfqpoint{5.070762in}{0.963493in}}%
\pgfpathlineto{\pgfqpoint{5.072776in}{0.897500in}}%
\pgfpathlineto{\pgfqpoint{5.073257in}{0.905749in}}%
\pgfpathlineto{\pgfqpoint{5.073442in}{0.889251in}}%
\pgfpathlineto{\pgfqpoint{5.074597in}{0.856254in}}%
\pgfpathlineto{\pgfqpoint{5.073612in}{0.897500in}}%
\pgfpathlineto{\pgfqpoint{5.074619in}{0.872753in}}%
\pgfpathlineto{\pgfqpoint{5.076108in}{0.930497in}}%
\pgfpathlineto{\pgfqpoint{5.076293in}{0.913998in}}%
\pgfpathlineto{\pgfqpoint{5.076463in}{0.897500in}}%
\pgfpathlineto{\pgfqpoint{5.076781in}{0.930497in}}%
\pgfpathlineto{\pgfqpoint{5.077300in}{0.913998in}}%
\pgfpathlineto{\pgfqpoint{5.077300in}{0.922247in}}%
\pgfpathlineto{\pgfqpoint{5.078307in}{0.872753in}}%
\pgfpathlineto{\pgfqpoint{5.079795in}{0.905749in}}%
\pgfpathlineto{\pgfqpoint{5.079817in}{0.897500in}}%
\pgfpathlineto{\pgfqpoint{5.079987in}{0.881002in}}%
\pgfpathlineto{\pgfqpoint{5.080150in}{0.922247in}}%
\pgfpathlineto{\pgfqpoint{5.080632in}{0.930497in}}%
\pgfpathlineto{\pgfqpoint{5.080824in}{0.913998in}}%
\pgfpathlineto{\pgfqpoint{5.081809in}{0.806759in}}%
\pgfpathlineto{\pgfqpoint{5.082497in}{0.823258in}}%
\pgfpathlineto{\pgfqpoint{5.083334in}{0.848005in}}%
\pgfpathlineto{\pgfqpoint{5.083652in}{0.831507in}}%
\pgfpathlineto{\pgfqpoint{5.084489in}{0.806759in}}%
\pgfpathlineto{\pgfqpoint{5.084674in}{0.815008in}}%
\pgfpathlineto{\pgfqpoint{5.086170in}{0.782012in}}%
\pgfpathlineto{\pgfqpoint{5.085518in}{0.823258in}}%
\pgfpathlineto{\pgfqpoint{5.086170in}{0.790261in}}%
\pgfpathlineto{\pgfqpoint{5.087169in}{0.806759in}}%
\pgfpathlineto{\pgfqpoint{5.086503in}{0.782012in}}%
\pgfpathlineto{\pgfqpoint{5.087192in}{0.798510in}}%
\pgfpathlineto{\pgfqpoint{5.088362in}{0.823258in}}%
\pgfpathlineto{\pgfqpoint{5.088532in}{0.815008in}}%
\pgfpathlineto{\pgfqpoint{5.088532in}{0.806759in}}%
\pgfpathlineto{\pgfqpoint{5.089183in}{0.856254in}}%
\pgfpathlineto{\pgfqpoint{5.089539in}{0.823258in}}%
\pgfpathlineto{\pgfqpoint{5.089687in}{0.831507in}}%
\pgfpathlineto{\pgfqpoint{5.090042in}{0.798510in}}%
\pgfpathlineto{\pgfqpoint{5.090205in}{0.798510in}}%
\pgfpathlineto{\pgfqpoint{5.091034in}{0.806759in}}%
\pgfpathlineto{\pgfqpoint{5.091049in}{0.798510in}}%
\pgfpathlineto{\pgfqpoint{5.091886in}{0.773763in}}%
\pgfpathlineto{\pgfqpoint{5.092049in}{0.798510in}}%
\pgfpathlineto{\pgfqpoint{5.093226in}{0.823258in}}%
\pgfpathlineto{\pgfqpoint{5.092700in}{0.782012in}}%
\pgfpathlineto{\pgfqpoint{5.093396in}{0.815008in}}%
\pgfpathlineto{\pgfqpoint{5.093396in}{0.806759in}}%
\pgfpathlineto{\pgfqpoint{5.093715in}{0.831507in}}%
\pgfpathlineto{\pgfqpoint{5.094396in}{0.815008in}}%
\pgfpathlineto{\pgfqpoint{5.095403in}{0.773763in}}%
\pgfpathlineto{\pgfqpoint{5.095573in}{0.798510in}}%
\pgfpathlineto{\pgfqpoint{5.096240in}{0.823258in}}%
\pgfpathlineto{\pgfqpoint{5.096728in}{0.806759in}}%
\pgfpathlineto{\pgfqpoint{5.097084in}{0.798510in}}%
\pgfpathlineto{\pgfqpoint{5.097587in}{0.872753in}}%
\pgfpathlineto{\pgfqpoint{5.099431in}{0.823258in}}%
\pgfpathlineto{\pgfqpoint{5.099431in}{0.831507in}}%
\pgfpathlineto{\pgfqpoint{5.099594in}{0.848005in}}%
\pgfpathlineto{\pgfqpoint{5.100097in}{0.815008in}}%
\pgfpathlineto{\pgfqpoint{5.101104in}{0.798510in}}%
\pgfpathlineto{\pgfqpoint{5.100267in}{0.823258in}}%
\pgfpathlineto{\pgfqpoint{5.101104in}{0.806759in}}%
\pgfpathlineto{\pgfqpoint{5.102615in}{0.848005in}}%
\pgfpathlineto{\pgfqpoint{5.102777in}{0.839756in}}%
\pgfpathlineto{\pgfqpoint{5.103118in}{0.831507in}}%
\pgfpathlineto{\pgfqpoint{5.103436in}{0.856254in}}%
\pgfpathlineto{\pgfqpoint{5.103451in}{0.848005in}}%
\pgfpathlineto{\pgfqpoint{5.104614in}{0.856254in}}%
\pgfpathlineto{\pgfqpoint{5.105465in}{0.848005in}}%
\pgfpathlineto{\pgfqpoint{5.104791in}{0.872753in}}%
\pgfpathlineto{\pgfqpoint{5.105628in}{0.872753in}}%
\pgfpathlineto{\pgfqpoint{5.111477in}{0.633527in}}%
\pgfpathlineto{\pgfqpoint{5.111492in}{0.641776in}}%
\pgfpathlineto{\pgfqpoint{5.111811in}{0.633527in}}%
\pgfpathlineto{\pgfqpoint{5.112507in}{0.674773in}}%
\pgfpathlineto{\pgfqpoint{5.114017in}{0.633527in}}%
\pgfpathlineto{\pgfqpoint{5.114173in}{0.641776in}}%
\pgfpathlineto{\pgfqpoint{5.114180in}{0.658274in}}%
\pgfpathlineto{\pgfqpoint{5.114683in}{0.625278in}}%
\pgfpathlineto{\pgfqpoint{5.115187in}{0.650025in}}%
\pgfpathlineto{\pgfqpoint{5.116364in}{0.625278in}}%
\pgfpathlineto{\pgfqpoint{5.117519in}{0.633527in}}%
\pgfpathlineto{\pgfqpoint{5.117534in}{0.625278in}}%
\pgfpathlineto{\pgfqpoint{5.118371in}{0.625278in}}%
\pgfpathlineto{\pgfqpoint{5.119703in}{0.633527in}}%
\pgfpathlineto{\pgfqpoint{5.119703in}{0.625278in}}%
\pgfpathlineto{\pgfqpoint{5.120718in}{0.625278in}}%
\pgfpathlineto{\pgfqpoint{5.121880in}{0.633527in}}%
\pgfpathlineto{\pgfqpoint{5.122058in}{0.625278in}}%
\pgfpathlineto{\pgfqpoint{5.122887in}{0.625278in}}%
\pgfpathlineto{\pgfqpoint{5.123887in}{0.633527in}}%
\pgfpathlineto{\pgfqpoint{5.123887in}{0.625278in}}%
\pgfpathlineto{\pgfqpoint{5.125234in}{0.633527in}}%
\pgfpathlineto{\pgfqpoint{5.125234in}{0.625278in}}%
\pgfpathlineto{\pgfqpoint{5.126249in}{0.625278in}}%
\pgfpathlineto{\pgfqpoint{5.127419in}{0.633527in}}%
\pgfpathlineto{\pgfqpoint{5.127589in}{0.625278in}}%
\pgfpathlineto{\pgfqpoint{5.128426in}{0.625278in}}%
\pgfpathlineto{\pgfqpoint{5.129603in}{0.633527in}}%
\pgfpathlineto{\pgfqpoint{5.129936in}{0.625278in}}%
\pgfpathlineto{\pgfqpoint{5.130588in}{0.625278in}}%
\pgfpathlineto{\pgfqpoint{5.131602in}{0.633527in}}%
\pgfpathlineto{\pgfqpoint{5.131602in}{0.625278in}}%
\pgfpathlineto{\pgfqpoint{5.132772in}{0.633527in}}%
\pgfpathlineto{\pgfqpoint{5.132935in}{0.625278in}}%
\pgfpathlineto{\pgfqpoint{5.133616in}{0.625278in}}%
\pgfpathlineto{\pgfqpoint{5.134786in}{0.633527in}}%
\pgfpathlineto{\pgfqpoint{5.134956in}{0.625278in}}%
\pgfpathlineto{\pgfqpoint{5.135623in}{0.625278in}}%
\pgfpathlineto{\pgfqpoint{5.136800in}{0.633527in}}%
\pgfpathlineto{\pgfqpoint{5.136955in}{0.625278in}}%
\pgfpathlineto{\pgfqpoint{5.137807in}{0.625278in}}%
\pgfpathlineto{\pgfqpoint{5.138806in}{0.633527in}}%
\pgfpathlineto{\pgfqpoint{5.138806in}{0.625278in}}%
\pgfpathlineto{\pgfqpoint{5.139984in}{0.633527in}}%
\pgfpathlineto{\pgfqpoint{5.140317in}{0.625278in}}%
\pgfpathlineto{\pgfqpoint{5.140991in}{0.625278in}}%
\pgfpathlineto{\pgfqpoint{5.142160in}{0.633527in}}%
\pgfpathlineto{\pgfqpoint{5.142160in}{0.625278in}}%
\pgfpathlineto{\pgfqpoint{5.143167in}{0.625278in}}%
\pgfpathlineto{\pgfqpoint{5.144345in}{0.633527in}}%
\pgfpathlineto{\pgfqpoint{5.144678in}{0.625278in}}%
\pgfpathlineto{\pgfqpoint{5.145344in}{0.625278in}}%
\pgfpathlineto{\pgfqpoint{5.146684in}{0.633527in}}%
\pgfpathlineto{\pgfqpoint{5.146692in}{0.625278in}}%
\pgfpathlineto{\pgfqpoint{5.147684in}{0.625278in}}%
\pgfpathlineto{\pgfqpoint{5.148861in}{0.633527in}}%
\pgfpathlineto{\pgfqpoint{5.149032in}{0.625278in}}%
\pgfpathlineto{\pgfqpoint{5.149868in}{0.625278in}}%
\pgfpathlineto{\pgfqpoint{5.151038in}{0.633527in}}%
\pgfpathlineto{\pgfqpoint{5.151371in}{0.625278in}}%
\pgfpathlineto{\pgfqpoint{5.151875in}{0.625278in}}%
\pgfpathlineto{\pgfqpoint{5.153052in}{0.633527in}}%
\pgfpathlineto{\pgfqpoint{5.153385in}{0.625278in}}%
\pgfpathlineto{\pgfqpoint{5.154059in}{0.625278in}}%
\pgfpathlineto{\pgfqpoint{5.155229in}{0.633527in}}%
\pgfpathlineto{\pgfqpoint{5.155229in}{0.625278in}}%
\pgfpathlineto{\pgfqpoint{5.156236in}{0.625278in}}%
\pgfpathlineto{\pgfqpoint{5.157413in}{0.633527in}}%
\pgfpathlineto{\pgfqpoint{5.157746in}{0.625278in}}%
\pgfpathlineto{\pgfqpoint{5.158420in}{0.625278in}}%
\pgfpathlineto{\pgfqpoint{5.159753in}{0.633527in}}%
\pgfpathlineto{\pgfqpoint{5.159930in}{0.625278in}}%
\pgfpathlineto{\pgfqpoint{5.160767in}{0.625278in}}%
\pgfpathlineto{\pgfqpoint{5.161937in}{0.633527in}}%
\pgfpathlineto{\pgfqpoint{5.161944in}{0.625278in}}%
\pgfpathlineto{\pgfqpoint{5.162951in}{0.625278in}}%
\pgfpathlineto{\pgfqpoint{5.164114in}{0.633527in}}%
\pgfpathlineto{\pgfqpoint{5.164284in}{0.625278in}}%
\pgfpathlineto{\pgfqpoint{5.165121in}{0.625278in}}%
\pgfpathlineto{\pgfqpoint{5.166128in}{0.633527in}}%
\pgfpathlineto{\pgfqpoint{5.166128in}{0.625278in}}%
\pgfpathlineto{\pgfqpoint{5.167298in}{0.633527in}}%
\pgfpathlineto{\pgfqpoint{5.167638in}{0.625278in}}%
\pgfpathlineto{\pgfqpoint{5.168305in}{0.625278in}}%
\pgfpathlineto{\pgfqpoint{5.169482in}{0.633527in}}%
\pgfpathlineto{\pgfqpoint{5.169652in}{0.625278in}}%
\pgfpathlineto{\pgfqpoint{5.170481in}{0.625278in}}%
\pgfpathlineto{\pgfqpoint{5.171659in}{0.633527in}}%
\pgfpathlineto{\pgfqpoint{5.171829in}{0.625278in}}%
\pgfpathlineto{\pgfqpoint{5.172495in}{0.625278in}}%
\pgfusepath{stroke}%
\end{pgfscope}%
\begin{pgfscope}%
\pgfsetrectcap%
\pgfsetmiterjoin%
\pgfsetlinewidth{0.803000pt}%
\definecolor{currentstroke}{rgb}{0.000000,0.000000,0.000000}%
\pgfsetstrokecolor{currentstroke}%
\pgfsetdash{}{0pt}%
\pgfpathmoveto{\pgfqpoint{0.586112in}{0.502778in}}%
\pgfpathlineto{\pgfqpoint{0.586112in}{3.197778in}}%
\pgfusepath{stroke}%
\end{pgfscope}%
\begin{pgfscope}%
\pgfsetrectcap%
\pgfsetmiterjoin%
\pgfsetlinewidth{0.803000pt}%
\definecolor{currentstroke}{rgb}{0.000000,0.000000,0.000000}%
\pgfsetstrokecolor{currentstroke}%
\pgfsetdash{}{0pt}%
\pgfpathmoveto{\pgfqpoint{5.391112in}{0.502778in}}%
\pgfpathlineto{\pgfqpoint{5.391112in}{3.197778in}}%
\pgfusepath{stroke}%
\end{pgfscope}%
\begin{pgfscope}%
\pgfsetrectcap%
\pgfsetmiterjoin%
\pgfsetlinewidth{0.803000pt}%
\definecolor{currentstroke}{rgb}{0.000000,0.000000,0.000000}%
\pgfsetstrokecolor{currentstroke}%
\pgfsetdash{}{0pt}%
\pgfpathmoveto{\pgfqpoint{0.586112in}{0.502778in}}%
\pgfpathlineto{\pgfqpoint{5.391112in}{0.502778in}}%
\pgfusepath{stroke}%
\end{pgfscope}%
\begin{pgfscope}%
\pgfsetrectcap%
\pgfsetmiterjoin%
\pgfsetlinewidth{0.803000pt}%
\definecolor{currentstroke}{rgb}{0.000000,0.000000,0.000000}%
\pgfsetstrokecolor{currentstroke}%
\pgfsetdash{}{0pt}%
\pgfpathmoveto{\pgfqpoint{0.586112in}{3.197778in}}%
\pgfpathlineto{\pgfqpoint{5.391112in}{3.197778in}}%
\pgfusepath{stroke}%
\end{pgfscope}%
\begin{pgfscope}%
\definecolor{textcolor}{rgb}{0.000000,0.000000,0.000000}%
\pgfsetstrokecolor{textcolor}%
\pgfsetfillcolor{textcolor}%
\pgftext[x=2.988612in,y=3.281111in,,base]{\color{textcolor}{\rmfamily\fontsize{12.000000}{14.400000}\selectfont\catcode`\^=\active\def^{\ifmmode\sp\else\^{}\fi}\catcode`\%=\active\def%{\%}Queue Dynamics Comparison}}%
\end{pgfscope}%
\begin{pgfscope}%
\pgfsetbuttcap%
\pgfsetmiterjoin%
\definecolor{currentfill}{rgb}{1.000000,1.000000,1.000000}%
\pgfsetfillcolor{currentfill}%
\pgfsetfillopacity{0.800000}%
\pgfsetlinewidth{1.003750pt}%
\definecolor{currentstroke}{rgb}{0.800000,0.800000,0.800000}%
\pgfsetstrokecolor{currentstroke}%
\pgfsetstrokeopacity{0.800000}%
\pgfsetdash{}{0pt}%
\pgfpathmoveto{\pgfqpoint{4.309445in}{2.496250in}}%
\pgfpathlineto{\pgfqpoint{5.293890in}{2.496250in}}%
\pgfpathquadraticcurveto{\pgfqpoint{5.321667in}{2.496250in}}{\pgfqpoint{5.321667in}{2.524028in}}%
\pgfpathlineto{\pgfqpoint{5.321667in}{3.100556in}}%
\pgfpathquadraticcurveto{\pgfqpoint{5.321667in}{3.128333in}}{\pgfqpoint{5.293890in}{3.128333in}}%
\pgfpathlineto{\pgfqpoint{4.309445in}{3.128333in}}%
\pgfpathquadraticcurveto{\pgfqpoint{4.281667in}{3.128333in}}{\pgfqpoint{4.281667in}{3.100556in}}%
\pgfpathlineto{\pgfqpoint{4.281667in}{2.524028in}}%
\pgfpathquadraticcurveto{\pgfqpoint{4.281667in}{2.496250in}}{\pgfqpoint{4.309445in}{2.496250in}}%
\pgfpathlineto{\pgfqpoint{4.309445in}{2.496250in}}%
\pgfpathclose%
\pgfusepath{stroke,fill}%
\end{pgfscope}%
\begin{pgfscope}%
\pgfsetrectcap%
\pgfsetroundjoin%
\pgfsetlinewidth{1.505625pt}%
\definecolor{currentstroke}{rgb}{0.000000,0.000000,0.000000}%
\pgfsetstrokecolor{currentstroke}%
\pgfsetdash{}{0pt}%
\pgfpathmoveto{\pgfqpoint{4.337223in}{3.024167in}}%
\pgfpathlineto{\pgfqpoint{4.476112in}{3.024167in}}%
\pgfpathlineto{\pgfqpoint{4.615001in}{3.024167in}}%
\pgfusepath{stroke}%
\end{pgfscope}%
\begin{pgfscope}%
\definecolor{textcolor}{rgb}{0.000000,0.000000,0.000000}%
\pgfsetstrokecolor{textcolor}%
\pgfsetfillcolor{textcolor}%
\pgftext[x=4.726112in,y=2.975556in,left,base]{\color{textcolor}{\rmfamily\fontsize{10.000000}{12.000000}\selectfont\catcode`\^=\active\def^{\ifmmode\sp\else\^{}\fi}\catcode`\%=\active\def%{\%}RED}}%
\end{pgfscope}%
\begin{pgfscope}%
\pgfsetbuttcap%
\pgfsetroundjoin%
\pgfsetlinewidth{1.505625pt}%
\definecolor{currentstroke}{rgb}{0.400000,0.400000,0.400000}%
\pgfsetstrokecolor{currentstroke}%
\pgfsetdash{{5.550000pt}{2.400000pt}}{0.000000pt}%
\pgfpathmoveto{\pgfqpoint{4.337223in}{2.827361in}}%
\pgfpathlineto{\pgfqpoint{4.476112in}{2.827361in}}%
\pgfpathlineto{\pgfqpoint{4.615001in}{2.827361in}}%
\pgfusepath{stroke}%
\end{pgfscope}%
\begin{pgfscope}%
\definecolor{textcolor}{rgb}{0.000000,0.000000,0.000000}%
\pgfsetstrokecolor{textcolor}%
\pgfsetfillcolor{textcolor}%
\pgftext[x=4.726112in,y=2.778750in,left,base]{\color{textcolor}{\rmfamily\fontsize{10.000000}{12.000000}\selectfont\catcode`\^=\active\def^{\ifmmode\sp\else\^{}\fi}\catcode`\%=\active\def%{\%}ARED}}%
\end{pgfscope}%
\begin{pgfscope}%
\pgfsetbuttcap%
\pgfsetroundjoin%
\pgfsetlinewidth{1.505625pt}%
\definecolor{currentstroke}{rgb}{0.650000,0.650000,0.650000}%
\pgfsetstrokecolor{currentstroke}%
\pgfsetdash{{1.500000pt}{2.475000pt}}{0.000000pt}%
\pgfpathmoveto{\pgfqpoint{4.337223in}{2.630556in}}%
\pgfpathlineto{\pgfqpoint{4.476112in}{2.630556in}}%
\pgfpathlineto{\pgfqpoint{4.615001in}{2.630556in}}%
\pgfusepath{stroke}%
\end{pgfscope}%
\begin{pgfscope}%
\definecolor{textcolor}{rgb}{0.000000,0.000000,0.000000}%
\pgfsetstrokecolor{textcolor}%
\pgfsetfillcolor{textcolor}%
\pgftext[x=4.726112in,y=2.581944in,left,base]{\color{textcolor}{\rmfamily\fontsize{10.000000}{12.000000}\selectfont\catcode`\^=\active\def^{\ifmmode\sp\else\^{}\fi}\catcode`\%=\active\def%{\%}GENTLE}}%
\end{pgfscope}%
\end{pgfpicture}%
\makeatother%
\endgroup%
}
\caption{Queue length dynamics for RED, ARED, and Gentle RED}
\label{fig:queue}
\end{figure}

Two phases are clearly distinguishable on the plot.

\textbf{Slow-start phase ($t = 0$--$4$~s).} All TCP connections simultaneously enter slow-start and saturate the buffer within seconds. The queue for all algorithms reaches peak values: ARED---up to 100~packets, RED and Gentle RED---up to 84~packets. This burst is identical for all three, since the AQM algorithm has not yet accumulated sufficient statistics to react.

\textbf{Steady-state regime ($t > 5$~s).} After the initial burst dissipates, the algorithms diverge in behavior.

\textbf{RED} (blue curve) maintains the queue in the range of 5--26~packets with relatively regular oscillations. The curve rarely touches zero (empty queue only 13.4\% of time): the forced drop at \texttt{m\_qAvg >= maxTh} reacts quickly and keeps the queue ``in check,'' but at the cost of constant pressure on TCP flows.

\textbf{ARED} (orange curve) demonstrates the largest oscillation amplitude and the most chaotic pattern. Periodically the queue grows to 30--35~packets in steady state and sharply collapses to zero. This is a consequence of $\mathrm{maxP}$ adaptation: the algorithm alternately ``releases'' and ``tightens'' the drop probability, producing wide irregular cycles. p99~$= 49$~packets---the highest value among the three algorithms.

\textbf{Gentle RED} (green curve) returns to zero most frequently (26.6\% of time the queue is empty, median~$= 1$~packet). The smooth probabilistic function in the zone \texttt{[maxTh, 2*maxTh]} allows TCP flows to desynchronize their windows, leading to deeper but shorter queue ``dips.''

Steady-state queue statistics are presented in Table~\ref{tab:queue-stats}.

\begin{table}[H]
\centering
\caption{Queue length statistics in steady-state regime}
\label{tab:queue-stats}
\begin{tabular}{lcccccc}
\toprule
Algorithm & Mean (pkt) & Std.\ dev. & Median & p90 & p99 & Empty (\%) \\
\midrule
RED      & 7.74  & 7.45  & 6 & 18 & 26 & 13.4\% \\
ARED     & 9.57  & 11.94 & 4 & 25 & 49 & 22.3\% \\
Gentle   & 6.40  & 8.86  & 1 & 19 & 34 & 26.6\% \\
\bottomrule
\end{tabular}
\end{table}

The time distribution across EWMA average $\bar{q}$ operating zones is shown in Table~\ref{tab:zones}.

\begin{table}[H]
\centering
\caption{Fraction of time in each queue management zone}
\label{tab:zones}
\begin{tabular}{lcccc}
\toprule
Algorithm & $\bar{q} < \mathrm{min}_{\mathrm{Th}}$ & $[\mathrm{min}_{\mathrm{Th}}, \mathrm{max}_{\mathrm{Th}})$ & $[\mathrm{max}_{\mathrm{Th}}, 2\mathrm{max}_{\mathrm{Th}})$ & $\bar{q} \ge 2\mathrm{max}_{\mathrm{Th}}$ \\
\midrule
RED      & 46.0\% & 19.8\% & 29.3\% & 4.8\%  \\
ARED     & 52.6\% & 8.0\%  & 19.6\% & 19.7\% \\
Gentle   & 60.8\% & 11.8\% & 19.8\% & 7.7\%  \\
\bottomrule
\end{tabular}
\end{table}

RED spends 34.1\% of time in the forced drop zone (\texttt{m\_qAvg >= maxTh}), explaining its high average queue. ARED spends 19.7\% of time in the $\bar{q} \ge 2\cdot\mathrm{max}_{\mathrm{Th}}$ zone where all packets are dropped---hence its anomalously high p99. Gentle RED most frequently stays below minTh (60.8\%) and least frequently enters the hard drop zone (7.7\%).

\subsection{Throughput}

The throughput plot (Fig.~\ref{fig:throughput}) is constructed with moving average smoothing (window of 30~samples), thus displaying data starting from $t \approx 29$~s. The Y-axis is intentionally compressed to the range 1.74--1.86~Mbit/s, allowing examination of differences invisible on the full scale.

\begin{figure}[H]
\centering
\resizebox{\textwidth}{!}{%% Creator: Matplotlib, PGF backend
%%
%% To include the figure in your LaTeX document, write
%%   \input{<filename>.pgf}
%%
%% Make sure the required packages are loaded in your preamble
%%   \usepackage{pgf}
%%
%% Also ensure that all the required font packages are loaded; for instance,
%% the lmodern package is sometimes necessary when using math font.
%%   \usepackage{lmodern}
%%
%% Figures using additional raster images can only be included by \input if
%% they are in the same directory as the main LaTeX file. For loading figures
%% from other directories you can use the `import` package
%%   \usepackage{import}
%%
%% and then include the figures with
%%   \import{<path to file>}{<filename>.pgf}
%%
%% Matplotlib used the following preamble
%%   \def\mathdefault#1{#1}
%%   \everymath=\expandafter{\the\everymath\displaystyle}
%%   \IfFileExists{scrextend.sty}{
%%     \usepackage[fontsize=10.000000pt]{scrextend}
%%   }{
%%     \renewcommand{\normalsize}{\fontsize{10.000000}{12.000000}\selectfont}
%%     \normalsize
%%   }
%%   \usepackage{fontspec}
%%   \setmainfont{TeX Gyre Termes}
%%   \setsansfont{TeX Gyre Heros}
%%   \setmonofont{TeX Gyre Cursor}
%%   \makeatletter\@ifpackageloaded{underscore}{}{\usepackage[strings]{underscore}}\makeatother
%%
\begingroup%
\makeatletter%
\begin{pgfpicture}%
\pgfpathrectangle{\pgfpointorigin}{\pgfqpoint{5.531493in}{3.494944in}}%
\pgfusepath{use as bounding box, clip}%
\begin{pgfscope}%
\pgfsetbuttcap%
\pgfsetmiterjoin%
\definecolor{currentfill}{rgb}{1.000000,1.000000,1.000000}%
\pgfsetfillcolor{currentfill}%
\pgfsetlinewidth{0.000000pt}%
\definecolor{currentstroke}{rgb}{1.000000,1.000000,1.000000}%
\pgfsetstrokecolor{currentstroke}%
\pgfsetdash{}{0pt}%
\pgfpathmoveto{\pgfqpoint{0.000000in}{0.000000in}}%
\pgfpathlineto{\pgfqpoint{5.531493in}{0.000000in}}%
\pgfpathlineto{\pgfqpoint{5.531493in}{3.494944in}}%
\pgfpathlineto{\pgfqpoint{0.000000in}{3.494944in}}%
\pgfpathlineto{\pgfqpoint{0.000000in}{0.000000in}}%
\pgfpathclose%
\pgfusepath{fill}%
\end{pgfscope}%
\begin{pgfscope}%
\pgfsetbuttcap%
\pgfsetmiterjoin%
\definecolor{currentfill}{rgb}{1.000000,1.000000,1.000000}%
\pgfsetfillcolor{currentfill}%
\pgfsetlinewidth{0.000000pt}%
\definecolor{currentstroke}{rgb}{0.000000,0.000000,0.000000}%
\pgfsetstrokecolor{currentstroke}%
\pgfsetstrokeopacity{0.000000}%
\pgfsetdash{}{0pt}%
\pgfpathmoveto{\pgfqpoint{0.624831in}{0.502778in}}%
\pgfpathlineto{\pgfqpoint{5.429831in}{0.502778in}}%
\pgfpathlineto{\pgfqpoint{5.429831in}{3.197778in}}%
\pgfpathlineto{\pgfqpoint{0.624831in}{3.197778in}}%
\pgfpathlineto{\pgfqpoint{0.624831in}{0.502778in}}%
\pgfpathclose%
\pgfusepath{fill}%
\end{pgfscope}%
\begin{pgfscope}%
\pgfpathrectangle{\pgfqpoint{0.624831in}{0.502778in}}{\pgfqpoint{4.805000in}{2.695000in}}%
\pgfusepath{clip}%
\pgfsetrectcap%
\pgfsetroundjoin%
\pgfsetlinewidth{0.803000pt}%
\definecolor{currentstroke}{rgb}{0.690196,0.690196,0.690196}%
\pgfsetstrokecolor{currentstroke}%
\pgfsetdash{}{0pt}%
\pgfpathmoveto{\pgfqpoint{0.843240in}{0.502778in}}%
\pgfpathlineto{\pgfqpoint{0.843240in}{3.197778in}}%
\pgfusepath{stroke}%
\end{pgfscope}%
\begin{pgfscope}%
\pgfsetbuttcap%
\pgfsetroundjoin%
\definecolor{currentfill}{rgb}{0.000000,0.000000,0.000000}%
\pgfsetfillcolor{currentfill}%
\pgfsetlinewidth{0.803000pt}%
\definecolor{currentstroke}{rgb}{0.000000,0.000000,0.000000}%
\pgfsetstrokecolor{currentstroke}%
\pgfsetdash{}{0pt}%
\pgfsys@defobject{currentmarker}{\pgfqpoint{0.000000in}{-0.048611in}}{\pgfqpoint{0.000000in}{0.000000in}}{%
\pgfpathmoveto{\pgfqpoint{0.000000in}{0.000000in}}%
\pgfpathlineto{\pgfqpoint{0.000000in}{-0.048611in}}%
\pgfusepath{stroke,fill}%
}%
\begin{pgfscope}%
\pgfsys@transformshift{0.843240in}{0.502778in}%
\pgfsys@useobject{currentmarker}{}%
\end{pgfscope}%
\end{pgfscope}%
\begin{pgfscope}%
\definecolor{textcolor}{rgb}{0.000000,0.000000,0.000000}%
\pgfsetstrokecolor{textcolor}%
\pgfsetfillcolor{textcolor}%
\pgftext[x=0.843240in,y=0.405556in,,top]{\color{textcolor}{\rmfamily\fontsize{10.000000}{12.000000}\selectfont\catcode`\^=\active\def^{\ifmmode\sp\else\^{}\fi}\catcode`\%=\active\def%{\%}$\mathdefault{30}$}}%
\end{pgfscope}%
\begin{pgfscope}%
\pgfpathrectangle{\pgfqpoint{0.624831in}{0.502778in}}{\pgfqpoint{4.805000in}{2.695000in}}%
\pgfusepath{clip}%
\pgfsetrectcap%
\pgfsetroundjoin%
\pgfsetlinewidth{0.803000pt}%
\definecolor{currentstroke}{rgb}{0.690196,0.690196,0.690196}%
\pgfsetstrokecolor{currentstroke}%
\pgfsetdash{}{0pt}%
\pgfpathmoveto{\pgfqpoint{1.596375in}{0.502778in}}%
\pgfpathlineto{\pgfqpoint{1.596375in}{3.197778in}}%
\pgfusepath{stroke}%
\end{pgfscope}%
\begin{pgfscope}%
\pgfsetbuttcap%
\pgfsetroundjoin%
\definecolor{currentfill}{rgb}{0.000000,0.000000,0.000000}%
\pgfsetfillcolor{currentfill}%
\pgfsetlinewidth{0.803000pt}%
\definecolor{currentstroke}{rgb}{0.000000,0.000000,0.000000}%
\pgfsetstrokecolor{currentstroke}%
\pgfsetdash{}{0pt}%
\pgfsys@defobject{currentmarker}{\pgfqpoint{0.000000in}{-0.048611in}}{\pgfqpoint{0.000000in}{0.000000in}}{%
\pgfpathmoveto{\pgfqpoint{0.000000in}{0.000000in}}%
\pgfpathlineto{\pgfqpoint{0.000000in}{-0.048611in}}%
\pgfusepath{stroke,fill}%
}%
\begin{pgfscope}%
\pgfsys@transformshift{1.596375in}{0.502778in}%
\pgfsys@useobject{currentmarker}{}%
\end{pgfscope}%
\end{pgfscope}%
\begin{pgfscope}%
\definecolor{textcolor}{rgb}{0.000000,0.000000,0.000000}%
\pgfsetstrokecolor{textcolor}%
\pgfsetfillcolor{textcolor}%
\pgftext[x=1.596375in,y=0.405556in,,top]{\color{textcolor}{\rmfamily\fontsize{10.000000}{12.000000}\selectfont\catcode`\^=\active\def^{\ifmmode\sp\else\^{}\fi}\catcode`\%=\active\def%{\%}$\mathdefault{35}$}}%
\end{pgfscope}%
\begin{pgfscope}%
\pgfpathrectangle{\pgfqpoint{0.624831in}{0.502778in}}{\pgfqpoint{4.805000in}{2.695000in}}%
\pgfusepath{clip}%
\pgfsetrectcap%
\pgfsetroundjoin%
\pgfsetlinewidth{0.803000pt}%
\definecolor{currentstroke}{rgb}{0.690196,0.690196,0.690196}%
\pgfsetstrokecolor{currentstroke}%
\pgfsetdash{}{0pt}%
\pgfpathmoveto{\pgfqpoint{2.349510in}{0.502778in}}%
\pgfpathlineto{\pgfqpoint{2.349510in}{3.197778in}}%
\pgfusepath{stroke}%
\end{pgfscope}%
\begin{pgfscope}%
\pgfsetbuttcap%
\pgfsetroundjoin%
\definecolor{currentfill}{rgb}{0.000000,0.000000,0.000000}%
\pgfsetfillcolor{currentfill}%
\pgfsetlinewidth{0.803000pt}%
\definecolor{currentstroke}{rgb}{0.000000,0.000000,0.000000}%
\pgfsetstrokecolor{currentstroke}%
\pgfsetdash{}{0pt}%
\pgfsys@defobject{currentmarker}{\pgfqpoint{0.000000in}{-0.048611in}}{\pgfqpoint{0.000000in}{0.000000in}}{%
\pgfpathmoveto{\pgfqpoint{0.000000in}{0.000000in}}%
\pgfpathlineto{\pgfqpoint{0.000000in}{-0.048611in}}%
\pgfusepath{stroke,fill}%
}%
\begin{pgfscope}%
\pgfsys@transformshift{2.349510in}{0.502778in}%
\pgfsys@useobject{currentmarker}{}%
\end{pgfscope}%
\end{pgfscope}%
\begin{pgfscope}%
\definecolor{textcolor}{rgb}{0.000000,0.000000,0.000000}%
\pgfsetstrokecolor{textcolor}%
\pgfsetfillcolor{textcolor}%
\pgftext[x=2.349510in,y=0.405556in,,top]{\color{textcolor}{\rmfamily\fontsize{10.000000}{12.000000}\selectfont\catcode`\^=\active\def^{\ifmmode\sp\else\^{}\fi}\catcode`\%=\active\def%{\%}$\mathdefault{40}$}}%
\end{pgfscope}%
\begin{pgfscope}%
\pgfpathrectangle{\pgfqpoint{0.624831in}{0.502778in}}{\pgfqpoint{4.805000in}{2.695000in}}%
\pgfusepath{clip}%
\pgfsetrectcap%
\pgfsetroundjoin%
\pgfsetlinewidth{0.803000pt}%
\definecolor{currentstroke}{rgb}{0.690196,0.690196,0.690196}%
\pgfsetstrokecolor{currentstroke}%
\pgfsetdash{}{0pt}%
\pgfpathmoveto{\pgfqpoint{3.102644in}{0.502778in}}%
\pgfpathlineto{\pgfqpoint{3.102644in}{3.197778in}}%
\pgfusepath{stroke}%
\end{pgfscope}%
\begin{pgfscope}%
\pgfsetbuttcap%
\pgfsetroundjoin%
\definecolor{currentfill}{rgb}{0.000000,0.000000,0.000000}%
\pgfsetfillcolor{currentfill}%
\pgfsetlinewidth{0.803000pt}%
\definecolor{currentstroke}{rgb}{0.000000,0.000000,0.000000}%
\pgfsetstrokecolor{currentstroke}%
\pgfsetdash{}{0pt}%
\pgfsys@defobject{currentmarker}{\pgfqpoint{0.000000in}{-0.048611in}}{\pgfqpoint{0.000000in}{0.000000in}}{%
\pgfpathmoveto{\pgfqpoint{0.000000in}{0.000000in}}%
\pgfpathlineto{\pgfqpoint{0.000000in}{-0.048611in}}%
\pgfusepath{stroke,fill}%
}%
\begin{pgfscope}%
\pgfsys@transformshift{3.102644in}{0.502778in}%
\pgfsys@useobject{currentmarker}{}%
\end{pgfscope}%
\end{pgfscope}%
\begin{pgfscope}%
\definecolor{textcolor}{rgb}{0.000000,0.000000,0.000000}%
\pgfsetstrokecolor{textcolor}%
\pgfsetfillcolor{textcolor}%
\pgftext[x=3.102644in,y=0.405556in,,top]{\color{textcolor}{\rmfamily\fontsize{10.000000}{12.000000}\selectfont\catcode`\^=\active\def^{\ifmmode\sp\else\^{}\fi}\catcode`\%=\active\def%{\%}$\mathdefault{45}$}}%
\end{pgfscope}%
\begin{pgfscope}%
\pgfpathrectangle{\pgfqpoint{0.624831in}{0.502778in}}{\pgfqpoint{4.805000in}{2.695000in}}%
\pgfusepath{clip}%
\pgfsetrectcap%
\pgfsetroundjoin%
\pgfsetlinewidth{0.803000pt}%
\definecolor{currentstroke}{rgb}{0.690196,0.690196,0.690196}%
\pgfsetstrokecolor{currentstroke}%
\pgfsetdash{}{0pt}%
\pgfpathmoveto{\pgfqpoint{3.855779in}{0.502778in}}%
\pgfpathlineto{\pgfqpoint{3.855779in}{3.197778in}}%
\pgfusepath{stroke}%
\end{pgfscope}%
\begin{pgfscope}%
\pgfsetbuttcap%
\pgfsetroundjoin%
\definecolor{currentfill}{rgb}{0.000000,0.000000,0.000000}%
\pgfsetfillcolor{currentfill}%
\pgfsetlinewidth{0.803000pt}%
\definecolor{currentstroke}{rgb}{0.000000,0.000000,0.000000}%
\pgfsetstrokecolor{currentstroke}%
\pgfsetdash{}{0pt}%
\pgfsys@defobject{currentmarker}{\pgfqpoint{0.000000in}{-0.048611in}}{\pgfqpoint{0.000000in}{0.000000in}}{%
\pgfpathmoveto{\pgfqpoint{0.000000in}{0.000000in}}%
\pgfpathlineto{\pgfqpoint{0.000000in}{-0.048611in}}%
\pgfusepath{stroke,fill}%
}%
\begin{pgfscope}%
\pgfsys@transformshift{3.855779in}{0.502778in}%
\pgfsys@useobject{currentmarker}{}%
\end{pgfscope}%
\end{pgfscope}%
\begin{pgfscope}%
\definecolor{textcolor}{rgb}{0.000000,0.000000,0.000000}%
\pgfsetstrokecolor{textcolor}%
\pgfsetfillcolor{textcolor}%
\pgftext[x=3.855779in,y=0.405556in,,top]{\color{textcolor}{\rmfamily\fontsize{10.000000}{12.000000}\selectfont\catcode`\^=\active\def^{\ifmmode\sp\else\^{}\fi}\catcode`\%=\active\def%{\%}$\mathdefault{50}$}}%
\end{pgfscope}%
\begin{pgfscope}%
\pgfpathrectangle{\pgfqpoint{0.624831in}{0.502778in}}{\pgfqpoint{4.805000in}{2.695000in}}%
\pgfusepath{clip}%
\pgfsetrectcap%
\pgfsetroundjoin%
\pgfsetlinewidth{0.803000pt}%
\definecolor{currentstroke}{rgb}{0.690196,0.690196,0.690196}%
\pgfsetstrokecolor{currentstroke}%
\pgfsetdash{}{0pt}%
\pgfpathmoveto{\pgfqpoint{4.608914in}{0.502778in}}%
\pgfpathlineto{\pgfqpoint{4.608914in}{3.197778in}}%
\pgfusepath{stroke}%
\end{pgfscope}%
\begin{pgfscope}%
\pgfsetbuttcap%
\pgfsetroundjoin%
\definecolor{currentfill}{rgb}{0.000000,0.000000,0.000000}%
\pgfsetfillcolor{currentfill}%
\pgfsetlinewidth{0.803000pt}%
\definecolor{currentstroke}{rgb}{0.000000,0.000000,0.000000}%
\pgfsetstrokecolor{currentstroke}%
\pgfsetdash{}{0pt}%
\pgfsys@defobject{currentmarker}{\pgfqpoint{0.000000in}{-0.048611in}}{\pgfqpoint{0.000000in}{0.000000in}}{%
\pgfpathmoveto{\pgfqpoint{0.000000in}{0.000000in}}%
\pgfpathlineto{\pgfqpoint{0.000000in}{-0.048611in}}%
\pgfusepath{stroke,fill}%
}%
\begin{pgfscope}%
\pgfsys@transformshift{4.608914in}{0.502778in}%
\pgfsys@useobject{currentmarker}{}%
\end{pgfscope}%
\end{pgfscope}%
\begin{pgfscope}%
\definecolor{textcolor}{rgb}{0.000000,0.000000,0.000000}%
\pgfsetstrokecolor{textcolor}%
\pgfsetfillcolor{textcolor}%
\pgftext[x=4.608914in,y=0.405556in,,top]{\color{textcolor}{\rmfamily\fontsize{10.000000}{12.000000}\selectfont\catcode`\^=\active\def^{\ifmmode\sp\else\^{}\fi}\catcode`\%=\active\def%{\%}$\mathdefault{55}$}}%
\end{pgfscope}%
\begin{pgfscope}%
\pgfpathrectangle{\pgfqpoint{0.624831in}{0.502778in}}{\pgfqpoint{4.805000in}{2.695000in}}%
\pgfusepath{clip}%
\pgfsetrectcap%
\pgfsetroundjoin%
\pgfsetlinewidth{0.803000pt}%
\definecolor{currentstroke}{rgb}{0.690196,0.690196,0.690196}%
\pgfsetstrokecolor{currentstroke}%
\pgfsetdash{}{0pt}%
\pgfpathmoveto{\pgfqpoint{5.362049in}{0.502778in}}%
\pgfpathlineto{\pgfqpoint{5.362049in}{3.197778in}}%
\pgfusepath{stroke}%
\end{pgfscope}%
\begin{pgfscope}%
\pgfsetbuttcap%
\pgfsetroundjoin%
\definecolor{currentfill}{rgb}{0.000000,0.000000,0.000000}%
\pgfsetfillcolor{currentfill}%
\pgfsetlinewidth{0.803000pt}%
\definecolor{currentstroke}{rgb}{0.000000,0.000000,0.000000}%
\pgfsetstrokecolor{currentstroke}%
\pgfsetdash{}{0pt}%
\pgfsys@defobject{currentmarker}{\pgfqpoint{0.000000in}{-0.048611in}}{\pgfqpoint{0.000000in}{0.000000in}}{%
\pgfpathmoveto{\pgfqpoint{0.000000in}{0.000000in}}%
\pgfpathlineto{\pgfqpoint{0.000000in}{-0.048611in}}%
\pgfusepath{stroke,fill}%
}%
\begin{pgfscope}%
\pgfsys@transformshift{5.362049in}{0.502778in}%
\pgfsys@useobject{currentmarker}{}%
\end{pgfscope}%
\end{pgfscope}%
\begin{pgfscope}%
\definecolor{textcolor}{rgb}{0.000000,0.000000,0.000000}%
\pgfsetstrokecolor{textcolor}%
\pgfsetfillcolor{textcolor}%
\pgftext[x=5.362049in,y=0.405556in,,top]{\color{textcolor}{\rmfamily\fontsize{10.000000}{12.000000}\selectfont\catcode`\^=\active\def^{\ifmmode\sp\else\^{}\fi}\catcode`\%=\active\def%{\%}$\mathdefault{60}$}}%
\end{pgfscope}%
\begin{pgfscope}%
\definecolor{textcolor}{rgb}{0.000000,0.000000,0.000000}%
\pgfsetstrokecolor{textcolor}%
\pgfsetfillcolor{textcolor}%
\pgftext[x=3.027331in,y=0.225000in,,top]{\color{textcolor}{\rmfamily\fontsize{10.000000}{12.000000}\selectfont\catcode`\^=\active\def^{\ifmmode\sp\else\^{}\fi}\catcode`\%=\active\def%{\%}Time (s)}}%
\end{pgfscope}%
\begin{pgfscope}%
\pgfpathrectangle{\pgfqpoint{0.624831in}{0.502778in}}{\pgfqpoint{4.805000in}{2.695000in}}%
\pgfusepath{clip}%
\pgfsetrectcap%
\pgfsetroundjoin%
\pgfsetlinewidth{0.803000pt}%
\definecolor{currentstroke}{rgb}{0.690196,0.690196,0.690196}%
\pgfsetstrokecolor{currentstroke}%
\pgfsetdash{}{0pt}%
\pgfpathmoveto{\pgfqpoint{0.624831in}{0.687721in}}%
\pgfpathlineto{\pgfqpoint{5.429831in}{0.687721in}}%
\pgfusepath{stroke}%
\end{pgfscope}%
\begin{pgfscope}%
\pgfsetbuttcap%
\pgfsetroundjoin%
\definecolor{currentfill}{rgb}{0.000000,0.000000,0.000000}%
\pgfsetfillcolor{currentfill}%
\pgfsetlinewidth{0.803000pt}%
\definecolor{currentstroke}{rgb}{0.000000,0.000000,0.000000}%
\pgfsetstrokecolor{currentstroke}%
\pgfsetdash{}{0pt}%
\pgfsys@defobject{currentmarker}{\pgfqpoint{-0.048611in}{0.000000in}}{\pgfqpoint{-0.000000in}{0.000000in}}{%
\pgfpathmoveto{\pgfqpoint{-0.000000in}{0.000000in}}%
\pgfpathlineto{\pgfqpoint{-0.048611in}{0.000000in}}%
\pgfusepath{stroke,fill}%
}%
\begin{pgfscope}%
\pgfsys@transformshift{0.624831in}{0.687721in}%
\pgfsys@useobject{currentmarker}{}%
\end{pgfscope}%
\end{pgfscope}%
\begin{pgfscope}%
\definecolor{textcolor}{rgb}{0.000000,0.000000,0.000000}%
\pgfsetstrokecolor{textcolor}%
\pgfsetfillcolor{textcolor}%
\pgftext[x=0.280694in, y=0.640290in, left, base]{\color{textcolor}{\rmfamily\fontsize{10.000000}{12.000000}\selectfont\catcode`\^=\active\def^{\ifmmode\sp\else\^{}\fi}\catcode`\%=\active\def%{\%}$\mathdefault{1.74}$}}%
\end{pgfscope}%
\begin{pgfscope}%
\pgfpathrectangle{\pgfqpoint{0.624831in}{0.502778in}}{\pgfqpoint{4.805000in}{2.695000in}}%
\pgfusepath{clip}%
\pgfsetrectcap%
\pgfsetroundjoin%
\pgfsetlinewidth{0.803000pt}%
\definecolor{currentstroke}{rgb}{0.690196,0.690196,0.690196}%
\pgfsetstrokecolor{currentstroke}%
\pgfsetdash{}{0pt}%
\pgfpathmoveto{\pgfqpoint{0.624831in}{1.082970in}}%
\pgfpathlineto{\pgfqpoint{5.429831in}{1.082970in}}%
\pgfusepath{stroke}%
\end{pgfscope}%
\begin{pgfscope}%
\pgfsetbuttcap%
\pgfsetroundjoin%
\definecolor{currentfill}{rgb}{0.000000,0.000000,0.000000}%
\pgfsetfillcolor{currentfill}%
\pgfsetlinewidth{0.803000pt}%
\definecolor{currentstroke}{rgb}{0.000000,0.000000,0.000000}%
\pgfsetstrokecolor{currentstroke}%
\pgfsetdash{}{0pt}%
\pgfsys@defobject{currentmarker}{\pgfqpoint{-0.048611in}{0.000000in}}{\pgfqpoint{-0.000000in}{0.000000in}}{%
\pgfpathmoveto{\pgfqpoint{-0.000000in}{0.000000in}}%
\pgfpathlineto{\pgfqpoint{-0.048611in}{0.000000in}}%
\pgfusepath{stroke,fill}%
}%
\begin{pgfscope}%
\pgfsys@transformshift{0.624831in}{1.082970in}%
\pgfsys@useobject{currentmarker}{}%
\end{pgfscope}%
\end{pgfscope}%
\begin{pgfscope}%
\definecolor{textcolor}{rgb}{0.000000,0.000000,0.000000}%
\pgfsetstrokecolor{textcolor}%
\pgfsetfillcolor{textcolor}%
\pgftext[x=0.280694in, y=1.035540in, left, base]{\color{textcolor}{\rmfamily\fontsize{10.000000}{12.000000}\selectfont\catcode`\^=\active\def^{\ifmmode\sp\else\^{}\fi}\catcode`\%=\active\def%{\%}$\mathdefault{1.76}$}}%
\end{pgfscope}%
\begin{pgfscope}%
\pgfpathrectangle{\pgfqpoint{0.624831in}{0.502778in}}{\pgfqpoint{4.805000in}{2.695000in}}%
\pgfusepath{clip}%
\pgfsetrectcap%
\pgfsetroundjoin%
\pgfsetlinewidth{0.803000pt}%
\definecolor{currentstroke}{rgb}{0.690196,0.690196,0.690196}%
\pgfsetstrokecolor{currentstroke}%
\pgfsetdash{}{0pt}%
\pgfpathmoveto{\pgfqpoint{0.624831in}{1.478220in}}%
\pgfpathlineto{\pgfqpoint{5.429831in}{1.478220in}}%
\pgfusepath{stroke}%
\end{pgfscope}%
\begin{pgfscope}%
\pgfsetbuttcap%
\pgfsetroundjoin%
\definecolor{currentfill}{rgb}{0.000000,0.000000,0.000000}%
\pgfsetfillcolor{currentfill}%
\pgfsetlinewidth{0.803000pt}%
\definecolor{currentstroke}{rgb}{0.000000,0.000000,0.000000}%
\pgfsetstrokecolor{currentstroke}%
\pgfsetdash{}{0pt}%
\pgfsys@defobject{currentmarker}{\pgfqpoint{-0.048611in}{0.000000in}}{\pgfqpoint{-0.000000in}{0.000000in}}{%
\pgfpathmoveto{\pgfqpoint{-0.000000in}{0.000000in}}%
\pgfpathlineto{\pgfqpoint{-0.048611in}{0.000000in}}%
\pgfusepath{stroke,fill}%
}%
\begin{pgfscope}%
\pgfsys@transformshift{0.624831in}{1.478220in}%
\pgfsys@useobject{currentmarker}{}%
\end{pgfscope}%
\end{pgfscope}%
\begin{pgfscope}%
\definecolor{textcolor}{rgb}{0.000000,0.000000,0.000000}%
\pgfsetstrokecolor{textcolor}%
\pgfsetfillcolor{textcolor}%
\pgftext[x=0.280694in, y=1.430789in, left, base]{\color{textcolor}{\rmfamily\fontsize{10.000000}{12.000000}\selectfont\catcode`\^=\active\def^{\ifmmode\sp\else\^{}\fi}\catcode`\%=\active\def%{\%}$\mathdefault{1.78}$}}%
\end{pgfscope}%
\begin{pgfscope}%
\pgfpathrectangle{\pgfqpoint{0.624831in}{0.502778in}}{\pgfqpoint{4.805000in}{2.695000in}}%
\pgfusepath{clip}%
\pgfsetrectcap%
\pgfsetroundjoin%
\pgfsetlinewidth{0.803000pt}%
\definecolor{currentstroke}{rgb}{0.690196,0.690196,0.690196}%
\pgfsetstrokecolor{currentstroke}%
\pgfsetdash{}{0pt}%
\pgfpathmoveto{\pgfqpoint{0.624831in}{1.873469in}}%
\pgfpathlineto{\pgfqpoint{5.429831in}{1.873469in}}%
\pgfusepath{stroke}%
\end{pgfscope}%
\begin{pgfscope}%
\pgfsetbuttcap%
\pgfsetroundjoin%
\definecolor{currentfill}{rgb}{0.000000,0.000000,0.000000}%
\pgfsetfillcolor{currentfill}%
\pgfsetlinewidth{0.803000pt}%
\definecolor{currentstroke}{rgb}{0.000000,0.000000,0.000000}%
\pgfsetstrokecolor{currentstroke}%
\pgfsetdash{}{0pt}%
\pgfsys@defobject{currentmarker}{\pgfqpoint{-0.048611in}{0.000000in}}{\pgfqpoint{-0.000000in}{0.000000in}}{%
\pgfpathmoveto{\pgfqpoint{-0.000000in}{0.000000in}}%
\pgfpathlineto{\pgfqpoint{-0.048611in}{0.000000in}}%
\pgfusepath{stroke,fill}%
}%
\begin{pgfscope}%
\pgfsys@transformshift{0.624831in}{1.873469in}%
\pgfsys@useobject{currentmarker}{}%
\end{pgfscope}%
\end{pgfscope}%
\begin{pgfscope}%
\definecolor{textcolor}{rgb}{0.000000,0.000000,0.000000}%
\pgfsetstrokecolor{textcolor}%
\pgfsetfillcolor{textcolor}%
\pgftext[x=0.280694in, y=1.826038in, left, base]{\color{textcolor}{\rmfamily\fontsize{10.000000}{12.000000}\selectfont\catcode`\^=\active\def^{\ifmmode\sp\else\^{}\fi}\catcode`\%=\active\def%{\%}$\mathdefault{1.80}$}}%
\end{pgfscope}%
\begin{pgfscope}%
\pgfpathrectangle{\pgfqpoint{0.624831in}{0.502778in}}{\pgfqpoint{4.805000in}{2.695000in}}%
\pgfusepath{clip}%
\pgfsetrectcap%
\pgfsetroundjoin%
\pgfsetlinewidth{0.803000pt}%
\definecolor{currentstroke}{rgb}{0.690196,0.690196,0.690196}%
\pgfsetstrokecolor{currentstroke}%
\pgfsetdash{}{0pt}%
\pgfpathmoveto{\pgfqpoint{0.624831in}{2.268719in}}%
\pgfpathlineto{\pgfqpoint{5.429831in}{2.268719in}}%
\pgfusepath{stroke}%
\end{pgfscope}%
\begin{pgfscope}%
\pgfsetbuttcap%
\pgfsetroundjoin%
\definecolor{currentfill}{rgb}{0.000000,0.000000,0.000000}%
\pgfsetfillcolor{currentfill}%
\pgfsetlinewidth{0.803000pt}%
\definecolor{currentstroke}{rgb}{0.000000,0.000000,0.000000}%
\pgfsetstrokecolor{currentstroke}%
\pgfsetdash{}{0pt}%
\pgfsys@defobject{currentmarker}{\pgfqpoint{-0.048611in}{0.000000in}}{\pgfqpoint{-0.000000in}{0.000000in}}{%
\pgfpathmoveto{\pgfqpoint{-0.000000in}{0.000000in}}%
\pgfpathlineto{\pgfqpoint{-0.048611in}{0.000000in}}%
\pgfusepath{stroke,fill}%
}%
\begin{pgfscope}%
\pgfsys@transformshift{0.624831in}{2.268719in}%
\pgfsys@useobject{currentmarker}{}%
\end{pgfscope}%
\end{pgfscope}%
\begin{pgfscope}%
\definecolor{textcolor}{rgb}{0.000000,0.000000,0.000000}%
\pgfsetstrokecolor{textcolor}%
\pgfsetfillcolor{textcolor}%
\pgftext[x=0.280694in, y=2.221288in, left, base]{\color{textcolor}{\rmfamily\fontsize{10.000000}{12.000000}\selectfont\catcode`\^=\active\def^{\ifmmode\sp\else\^{}\fi}\catcode`\%=\active\def%{\%}$\mathdefault{1.82}$}}%
\end{pgfscope}%
\begin{pgfscope}%
\pgfpathrectangle{\pgfqpoint{0.624831in}{0.502778in}}{\pgfqpoint{4.805000in}{2.695000in}}%
\pgfusepath{clip}%
\pgfsetrectcap%
\pgfsetroundjoin%
\pgfsetlinewidth{0.803000pt}%
\definecolor{currentstroke}{rgb}{0.690196,0.690196,0.690196}%
\pgfsetstrokecolor{currentstroke}%
\pgfsetdash{}{0pt}%
\pgfpathmoveto{\pgfqpoint{0.624831in}{2.663968in}}%
\pgfpathlineto{\pgfqpoint{5.429831in}{2.663968in}}%
\pgfusepath{stroke}%
\end{pgfscope}%
\begin{pgfscope}%
\pgfsetbuttcap%
\pgfsetroundjoin%
\definecolor{currentfill}{rgb}{0.000000,0.000000,0.000000}%
\pgfsetfillcolor{currentfill}%
\pgfsetlinewidth{0.803000pt}%
\definecolor{currentstroke}{rgb}{0.000000,0.000000,0.000000}%
\pgfsetstrokecolor{currentstroke}%
\pgfsetdash{}{0pt}%
\pgfsys@defobject{currentmarker}{\pgfqpoint{-0.048611in}{0.000000in}}{\pgfqpoint{-0.000000in}{0.000000in}}{%
\pgfpathmoveto{\pgfqpoint{-0.000000in}{0.000000in}}%
\pgfpathlineto{\pgfqpoint{-0.048611in}{0.000000in}}%
\pgfusepath{stroke,fill}%
}%
\begin{pgfscope}%
\pgfsys@transformshift{0.624831in}{2.663968in}%
\pgfsys@useobject{currentmarker}{}%
\end{pgfscope}%
\end{pgfscope}%
\begin{pgfscope}%
\definecolor{textcolor}{rgb}{0.000000,0.000000,0.000000}%
\pgfsetstrokecolor{textcolor}%
\pgfsetfillcolor{textcolor}%
\pgftext[x=0.280694in, y=2.616537in, left, base]{\color{textcolor}{\rmfamily\fontsize{10.000000}{12.000000}\selectfont\catcode`\^=\active\def^{\ifmmode\sp\else\^{}\fi}\catcode`\%=\active\def%{\%}$\mathdefault{1.84}$}}%
\end{pgfscope}%
\begin{pgfscope}%
\pgfpathrectangle{\pgfqpoint{0.624831in}{0.502778in}}{\pgfqpoint{4.805000in}{2.695000in}}%
\pgfusepath{clip}%
\pgfsetrectcap%
\pgfsetroundjoin%
\pgfsetlinewidth{0.803000pt}%
\definecolor{currentstroke}{rgb}{0.690196,0.690196,0.690196}%
\pgfsetstrokecolor{currentstroke}%
\pgfsetdash{}{0pt}%
\pgfpathmoveto{\pgfqpoint{0.624831in}{3.059217in}}%
\pgfpathlineto{\pgfqpoint{5.429831in}{3.059217in}}%
\pgfusepath{stroke}%
\end{pgfscope}%
\begin{pgfscope}%
\pgfsetbuttcap%
\pgfsetroundjoin%
\definecolor{currentfill}{rgb}{0.000000,0.000000,0.000000}%
\pgfsetfillcolor{currentfill}%
\pgfsetlinewidth{0.803000pt}%
\definecolor{currentstroke}{rgb}{0.000000,0.000000,0.000000}%
\pgfsetstrokecolor{currentstroke}%
\pgfsetdash{}{0pt}%
\pgfsys@defobject{currentmarker}{\pgfqpoint{-0.048611in}{0.000000in}}{\pgfqpoint{-0.000000in}{0.000000in}}{%
\pgfpathmoveto{\pgfqpoint{-0.000000in}{0.000000in}}%
\pgfpathlineto{\pgfqpoint{-0.048611in}{0.000000in}}%
\pgfusepath{stroke,fill}%
}%
\begin{pgfscope}%
\pgfsys@transformshift{0.624831in}{3.059217in}%
\pgfsys@useobject{currentmarker}{}%
\end{pgfscope}%
\end{pgfscope}%
\begin{pgfscope}%
\definecolor{textcolor}{rgb}{0.000000,0.000000,0.000000}%
\pgfsetstrokecolor{textcolor}%
\pgfsetfillcolor{textcolor}%
\pgftext[x=0.280694in, y=3.011787in, left, base]{\color{textcolor}{\rmfamily\fontsize{10.000000}{12.000000}\selectfont\catcode`\^=\active\def^{\ifmmode\sp\else\^{}\fi}\catcode`\%=\active\def%{\%}$\mathdefault{1.86}$}}%
\end{pgfscope}%
\begin{pgfscope}%
\definecolor{textcolor}{rgb}{0.000000,0.000000,0.000000}%
\pgfsetstrokecolor{textcolor}%
\pgfsetfillcolor{textcolor}%
\pgftext[x=0.225139in,y=1.850278in,,bottom,rotate=90.000000]{\color{textcolor}{\rmfamily\fontsize{10.000000}{12.000000}\selectfont\catcode`\^=\active\def^{\ifmmode\sp\else\^{}\fi}\catcode`\%=\active\def%{\%}Throughput (Mbps)}}%
\end{pgfscope}%
\begin{pgfscope}%
\pgfpathrectangle{\pgfqpoint{0.624831in}{0.502778in}}{\pgfqpoint{4.805000in}{2.695000in}}%
\pgfusepath{clip}%
\pgfsetrectcap%
\pgfsetroundjoin%
\pgfsetlinewidth{1.505625pt}%
\definecolor{currentstroke}{rgb}{0.121569,0.466667,0.705882}%
\pgfsetstrokecolor{currentstroke}%
\pgfsetdash{}{0pt}%
\pgfpathmoveto{\pgfqpoint{0.843240in}{0.625278in}}%
\pgfpathlineto{\pgfqpoint{0.993867in}{1.807107in}}%
\pgfpathlineto{\pgfqpoint{1.144494in}{2.389928in}}%
\pgfpathlineto{\pgfqpoint{1.295121in}{3.075278in}}%
\pgfpathlineto{\pgfqpoint{1.445748in}{1.842185in}}%
\pgfpathlineto{\pgfqpoint{1.596375in}{2.195650in}}%
\pgfpathlineto{\pgfqpoint{1.747002in}{2.147080in}}%
\pgfpathlineto{\pgfqpoint{1.897629in}{1.839484in}}%
\pgfpathlineto{\pgfqpoint{2.048256in}{2.014870in}}%
\pgfpathlineto{\pgfqpoint{2.198883in}{2.114703in}}%
\pgfpathlineto{\pgfqpoint{2.349510in}{2.157877in}}%
\pgfpathlineto{\pgfqpoint{2.500137in}{1.979791in}}%
\pgfpathlineto{\pgfqpoint{2.650764in}{2.055343in}}%
\pgfpathlineto{\pgfqpoint{2.801391in}{2.074236in}}%
\pgfpathlineto{\pgfqpoint{2.952017in}{2.028367in}}%
\pgfpathlineto{\pgfqpoint{3.102644in}{2.025666in}}%
\pgfpathlineto{\pgfqpoint{3.253271in}{2.055343in}}%
\pgfpathlineto{\pgfqpoint{3.403898in}{1.879958in}}%
\pgfpathlineto{\pgfqpoint{3.554525in}{1.842185in}}%
\pgfpathlineto{\pgfqpoint{3.705152in}{2.130895in}}%
\pgfpathlineto{\pgfqpoint{3.855779in}{1.993282in}}%
\pgfpathlineto{\pgfqpoint{4.006406in}{1.982492in}}%
\pgfpathlineto{\pgfqpoint{4.157033in}{2.120098in}}%
\pgfpathlineto{\pgfqpoint{4.307660in}{2.125493in}}%
\pgfpathlineto{\pgfqpoint{4.458287in}{2.036457in}}%
\pgfpathlineto{\pgfqpoint{4.608914in}{2.090421in}}%
\pgfpathlineto{\pgfqpoint{4.759541in}{2.055343in}}%
\pgfpathlineto{\pgfqpoint{4.910168in}{2.047247in}}%
\pgfpathlineto{\pgfqpoint{5.060795in}{1.995983in}}%
\pgfpathlineto{\pgfqpoint{5.211422in}{2.017570in}}%
\pgfusepath{stroke}%
\end{pgfscope}%
\begin{pgfscope}%
\pgfpathrectangle{\pgfqpoint{0.624831in}{0.502778in}}{\pgfqpoint{4.805000in}{2.695000in}}%
\pgfusepath{clip}%
\pgfsetrectcap%
\pgfsetroundjoin%
\pgfsetlinewidth{1.505625pt}%
\definecolor{currentstroke}{rgb}{1.000000,0.498039,0.054902}%
\pgfsetstrokecolor{currentstroke}%
\pgfsetdash{}{0pt}%
\pgfpathmoveto{\pgfqpoint{0.843240in}{0.692747in}}%
\pgfpathlineto{\pgfqpoint{0.993867in}{1.469840in}}%
\pgfpathlineto{\pgfqpoint{1.144494in}{2.187567in}}%
\pgfpathlineto{\pgfqpoint{1.295121in}{3.015924in}}%
\pgfpathlineto{\pgfqpoint{1.445748in}{2.012182in}}%
\pgfpathlineto{\pgfqpoint{1.596375in}{1.990595in}}%
\pgfpathlineto{\pgfqpoint{1.747002in}{2.568021in}}%
\pgfpathlineto{\pgfqpoint{1.897629in}{2.020278in}}%
\pgfpathlineto{\pgfqpoint{2.048256in}{2.017584in}}%
\pgfpathlineto{\pgfqpoint{2.198883in}{2.017584in}}%
\pgfpathlineto{\pgfqpoint{2.349510in}{1.977110in}}%
\pgfpathlineto{\pgfqpoint{2.500137in}{2.025680in}}%
\pgfpathlineto{\pgfqpoint{2.650764in}{2.025680in}}%
\pgfpathlineto{\pgfqpoint{2.801391in}{2.025680in}}%
\pgfpathlineto{\pgfqpoint{2.952017in}{2.039171in}}%
\pgfpathlineto{\pgfqpoint{3.102644in}{2.017584in}}%
\pgfpathlineto{\pgfqpoint{3.253271in}{1.831401in}}%
\pgfpathlineto{\pgfqpoint{3.403898in}{2.095830in}}%
\pgfpathlineto{\pgfqpoint{3.554525in}{2.152495in}}%
\pgfpathlineto{\pgfqpoint{3.705152in}{2.047267in}}%
\pgfpathlineto{\pgfqpoint{3.855779in}{1.739664in}}%
\pgfpathlineto{\pgfqpoint{4.006406in}{1.680304in}}%
\pgfpathlineto{\pgfqpoint{4.157033in}{2.128214in}}%
\pgfpathlineto{\pgfqpoint{4.307660in}{2.389942in}}%
\pgfpathlineto{\pgfqpoint{4.458287in}{1.985206in}}%
\pgfpathlineto{\pgfqpoint{4.608914in}{1.693795in}}%
\pgfpathlineto{\pgfqpoint{4.759541in}{1.750461in}}%
\pgfpathlineto{\pgfqpoint{4.910168in}{2.114723in}}%
\pgfpathlineto{\pgfqpoint{5.060795in}{2.165993in}}%
\pgfpathlineto{\pgfqpoint{5.211422in}{1.971715in}}%
\pgfusepath{stroke}%
\end{pgfscope}%
\begin{pgfscope}%
\pgfpathrectangle{\pgfqpoint{0.624831in}{0.502778in}}{\pgfqpoint{4.805000in}{2.695000in}}%
\pgfusepath{clip}%
\pgfsetrectcap%
\pgfsetroundjoin%
\pgfsetlinewidth{1.505625pt}%
\definecolor{currentstroke}{rgb}{0.172549,0.627451,0.172549}%
\pgfsetstrokecolor{currentstroke}%
\pgfsetdash{}{0pt}%
\pgfpathmoveto{\pgfqpoint{0.843240in}{0.746702in}}%
\pgfpathlineto{\pgfqpoint{0.993867in}{1.939322in}}%
\pgfpathlineto{\pgfqpoint{1.144494in}{2.049952in}}%
\pgfpathlineto{\pgfqpoint{1.295121in}{3.069883in}}%
\pgfpathlineto{\pgfqpoint{1.445748in}{1.998684in}}%
\pgfpathlineto{\pgfqpoint{1.596375in}{2.473576in}}%
\pgfpathlineto{\pgfqpoint{1.747002in}{1.739651in}}%
\pgfpathlineto{\pgfqpoint{1.897629in}{1.979798in}}%
\pgfpathlineto{\pgfqpoint{2.048256in}{1.987894in}}%
\pgfpathlineto{\pgfqpoint{2.198883in}{2.306287in}}%
\pgfpathlineto{\pgfqpoint{2.349510in}{1.979804in}}%
\pgfpathlineto{\pgfqpoint{2.500137in}{1.415874in}}%
\pgfpathlineto{\pgfqpoint{2.650764in}{1.985198in}}%
\pgfpathlineto{\pgfqpoint{2.801391in}{2.228039in}}%
\pgfpathlineto{\pgfqpoint{2.952017in}{2.112014in}}%
\pgfpathlineto{\pgfqpoint{3.102644in}{1.966305in}}%
\pgfpathlineto{\pgfqpoint{3.253271in}{1.825992in}}%
\pgfpathlineto{\pgfqpoint{3.403898in}{1.993281in}}%
\pgfpathlineto{\pgfqpoint{3.554525in}{2.376429in}}%
\pgfpathlineto{\pgfqpoint{3.705152in}{1.985185in}}%
\pgfpathlineto{\pgfqpoint{3.855779in}{1.610133in}}%
\pgfpathlineto{\pgfqpoint{4.006406in}{1.720763in}}%
\pgfpathlineto{\pgfqpoint{4.157033in}{2.090420in}}%
\pgfpathlineto{\pgfqpoint{4.307660in}{2.206446in}}%
\pgfpathlineto{\pgfqpoint{4.458287in}{1.998683in}}%
\pgfpathlineto{\pgfqpoint{4.608914in}{1.758536in}}%
\pgfpathlineto{\pgfqpoint{4.759541in}{1.990587in}}%
\pgfpathlineto{\pgfqpoint{4.910168in}{2.122798in}}%
\pgfpathlineto{\pgfqpoint{5.060795in}{2.085025in}}%
\pgfpathlineto{\pgfqpoint{5.211422in}{1.812500in}}%
\pgfusepath{stroke}%
\end{pgfscope}%
\begin{pgfscope}%
\pgfsetrectcap%
\pgfsetmiterjoin%
\pgfsetlinewidth{0.803000pt}%
\definecolor{currentstroke}{rgb}{0.000000,0.000000,0.000000}%
\pgfsetstrokecolor{currentstroke}%
\pgfsetdash{}{0pt}%
\pgfpathmoveto{\pgfqpoint{0.624831in}{0.502778in}}%
\pgfpathlineto{\pgfqpoint{0.624831in}{3.197778in}}%
\pgfusepath{stroke}%
\end{pgfscope}%
\begin{pgfscope}%
\pgfsetrectcap%
\pgfsetmiterjoin%
\pgfsetlinewidth{0.803000pt}%
\definecolor{currentstroke}{rgb}{0.000000,0.000000,0.000000}%
\pgfsetstrokecolor{currentstroke}%
\pgfsetdash{}{0pt}%
\pgfpathmoveto{\pgfqpoint{5.429831in}{0.502778in}}%
\pgfpathlineto{\pgfqpoint{5.429831in}{3.197778in}}%
\pgfusepath{stroke}%
\end{pgfscope}%
\begin{pgfscope}%
\pgfsetrectcap%
\pgfsetmiterjoin%
\pgfsetlinewidth{0.803000pt}%
\definecolor{currentstroke}{rgb}{0.000000,0.000000,0.000000}%
\pgfsetstrokecolor{currentstroke}%
\pgfsetdash{}{0pt}%
\pgfpathmoveto{\pgfqpoint{0.624831in}{0.502778in}}%
\pgfpathlineto{\pgfqpoint{5.429831in}{0.502778in}}%
\pgfusepath{stroke}%
\end{pgfscope}%
\begin{pgfscope}%
\pgfsetrectcap%
\pgfsetmiterjoin%
\pgfsetlinewidth{0.803000pt}%
\definecolor{currentstroke}{rgb}{0.000000,0.000000,0.000000}%
\pgfsetstrokecolor{currentstroke}%
\pgfsetdash{}{0pt}%
\pgfpathmoveto{\pgfqpoint{0.624831in}{3.197778in}}%
\pgfpathlineto{\pgfqpoint{5.429831in}{3.197778in}}%
\pgfusepath{stroke}%
\end{pgfscope}%
\begin{pgfscope}%
\definecolor{textcolor}{rgb}{0.000000,0.000000,0.000000}%
\pgfsetstrokecolor{textcolor}%
\pgfsetfillcolor{textcolor}%
\pgftext[x=3.027331in,y=3.281111in,,base]{\color{textcolor}{\rmfamily\fontsize{12.000000}{14.400000}\selectfont\catcode`\^=\active\def^{\ifmmode\sp\else\^{}\fi}\catcode`\%=\active\def%{\%}Throughput Comparison}}%
\end{pgfscope}%
\begin{pgfscope}%
\pgfsetbuttcap%
\pgfsetmiterjoin%
\definecolor{currentfill}{rgb}{1.000000,1.000000,1.000000}%
\pgfsetfillcolor{currentfill}%
\pgfsetfillopacity{0.800000}%
\pgfsetlinewidth{1.003750pt}%
\definecolor{currentstroke}{rgb}{0.800000,0.800000,0.800000}%
\pgfsetstrokecolor{currentstroke}%
\pgfsetstrokeopacity{0.800000}%
\pgfsetdash{}{0pt}%
\pgfpathmoveto{\pgfqpoint{4.348164in}{2.496250in}}%
\pgfpathlineto{\pgfqpoint{5.332609in}{2.496250in}}%
\pgfpathquadraticcurveto{\pgfqpoint{5.360387in}{2.496250in}}{\pgfqpoint{5.360387in}{2.524028in}}%
\pgfpathlineto{\pgfqpoint{5.360387in}{3.100556in}}%
\pgfpathquadraticcurveto{\pgfqpoint{5.360387in}{3.128333in}}{\pgfqpoint{5.332609in}{3.128333in}}%
\pgfpathlineto{\pgfqpoint{4.348164in}{3.128333in}}%
\pgfpathquadraticcurveto{\pgfqpoint{4.320387in}{3.128333in}}{\pgfqpoint{4.320387in}{3.100556in}}%
\pgfpathlineto{\pgfqpoint{4.320387in}{2.524028in}}%
\pgfpathquadraticcurveto{\pgfqpoint{4.320387in}{2.496250in}}{\pgfqpoint{4.348164in}{2.496250in}}%
\pgfpathlineto{\pgfqpoint{4.348164in}{2.496250in}}%
\pgfpathclose%
\pgfusepath{stroke,fill}%
\end{pgfscope}%
\begin{pgfscope}%
\pgfsetrectcap%
\pgfsetroundjoin%
\pgfsetlinewidth{1.505625pt}%
\definecolor{currentstroke}{rgb}{0.121569,0.466667,0.705882}%
\pgfsetstrokecolor{currentstroke}%
\pgfsetdash{}{0pt}%
\pgfpathmoveto{\pgfqpoint{4.375942in}{3.024167in}}%
\pgfpathlineto{\pgfqpoint{4.514831in}{3.024167in}}%
\pgfpathlineto{\pgfqpoint{4.653720in}{3.024167in}}%
\pgfusepath{stroke}%
\end{pgfscope}%
\begin{pgfscope}%
\definecolor{textcolor}{rgb}{0.000000,0.000000,0.000000}%
\pgfsetstrokecolor{textcolor}%
\pgfsetfillcolor{textcolor}%
\pgftext[x=4.764831in,y=2.975556in,left,base]{\color{textcolor}{\rmfamily\fontsize{10.000000}{12.000000}\selectfont\catcode`\^=\active\def^{\ifmmode\sp\else\^{}\fi}\catcode`\%=\active\def%{\%}RED}}%
\end{pgfscope}%
\begin{pgfscope}%
\pgfsetrectcap%
\pgfsetroundjoin%
\pgfsetlinewidth{1.505625pt}%
\definecolor{currentstroke}{rgb}{1.000000,0.498039,0.054902}%
\pgfsetstrokecolor{currentstroke}%
\pgfsetdash{}{0pt}%
\pgfpathmoveto{\pgfqpoint{4.375942in}{2.827361in}}%
\pgfpathlineto{\pgfqpoint{4.514831in}{2.827361in}}%
\pgfpathlineto{\pgfqpoint{4.653720in}{2.827361in}}%
\pgfusepath{stroke}%
\end{pgfscope}%
\begin{pgfscope}%
\definecolor{textcolor}{rgb}{0.000000,0.000000,0.000000}%
\pgfsetstrokecolor{textcolor}%
\pgfsetfillcolor{textcolor}%
\pgftext[x=4.764831in,y=2.778750in,left,base]{\color{textcolor}{\rmfamily\fontsize{10.000000}{12.000000}\selectfont\catcode`\^=\active\def^{\ifmmode\sp\else\^{}\fi}\catcode`\%=\active\def%{\%}ARED}}%
\end{pgfscope}%
\begin{pgfscope}%
\pgfsetrectcap%
\pgfsetroundjoin%
\pgfsetlinewidth{1.505625pt}%
\definecolor{currentstroke}{rgb}{0.172549,0.627451,0.172549}%
\pgfsetstrokecolor{currentstroke}%
\pgfsetdash{}{0pt}%
\pgfpathmoveto{\pgfqpoint{4.375942in}{2.630556in}}%
\pgfpathlineto{\pgfqpoint{4.514831in}{2.630556in}}%
\pgfpathlineto{\pgfqpoint{4.653720in}{2.630556in}}%
\pgfusepath{stroke}%
\end{pgfscope}%
\begin{pgfscope}%
\definecolor{textcolor}{rgb}{0.000000,0.000000,0.000000}%
\pgfsetstrokecolor{textcolor}%
\pgfsetfillcolor{textcolor}%
\pgftext[x=4.764831in,y=2.581944in,left,base]{\color{textcolor}{\rmfamily\fontsize{10.000000}{12.000000}\selectfont\catcode`\^=\active\def^{\ifmmode\sp\else\^{}\fi}\catcode`\%=\active\def%{\%}GENTLE}}%
\end{pgfscope}%
\end{pgfpicture}%
\makeatother%
\endgroup%
}
\caption{Throughput comparison: RED, ARED, and Gentle RED}
\label{fig:throughput}
\end{figure}

All three algorithms stably maintain throughput near \textbf{1.80--1.81~Mbit/s} ($\sim$90\% of the nominal 2~Mbit/s link capacity). This means that the choice of AQM algorithm at these parameters does not affect bandwidth utilization: all three mechanisms equally effectively ``fill'' the link.

The characteristic peak around $t = 33$~s (all curves reach $\sim$1.86~Mbit/s) represents synchronous recovery of TCP flows after the initial wave of drops: congestion windows simultaneously begin growing, creating a short-term excess flow.

\begin{table}[H]
\centering
\caption{Throughput statistics (CV---coefficient of variation)}
\label{tab:throughput}
\begin{tabular}{lccccc}
\toprule
Algorithm & Mean (Mbit/s) & Std.\ dev. & Min. & Max. & CV \\
\midrule
RED      & 1.8047 & 0.3380 & 0.533 & 3.621 & 18.7\% \\
ARED     & 1.8039 & 0.3950 & 0.553 & 3.293 & 21.9\% \\
Gentle   & 1.7998 & 0.4444 & 0.352 & 3.654 & 24.7\% \\
\bottomrule
\end{tabular}
\end{table}

Despite identical means, \textbf{RED demonstrates the lowest variability} (CV~$= 18.7\%$): hard synchronous drops keep all TCP flows in step, making throughput predictable. Gentle RED has the highest CV (24.7\%)---asynchronous window management creates more ``ripple,'' although the mean does not suffer.

\subsection{Packet drops}

The plot (Fig.~\ref{fig:drops}) presents drop moments as a scatter diagram: each point is one dropped packet. Three horizontal bands correspond to the three algorithms.

\begin{figure}[H]
\centering
\resizebox{\textwidth}{!}{%% Creator: Matplotlib, PGF backend
%%
%% To include the figure in your LaTeX document, write
%%   \input{<filename>.pgf}
%%
%% Make sure the required packages are loaded in your preamble
%%   \usepackage{pgf}
%%
%% Also ensure that all the required font packages are loaded; for instance,
%% the lmodern package is sometimes necessary when using math font.
%%   \usepackage{lmodern}
%%
%% Figures using additional raster images can only be included by \input if
%% they are in the same directory as the main LaTeX file. For loading figures
%% from other directories you can use the `import` package
%%   \usepackage{import}
%%
%% and then include the figures with
%%   \import{<path to file>}{<filename>.pgf}
%%
%% Matplotlib used the following preamble
%%   \usepackage{fontspec}
%%   \setmainfont{TeX Gyre Termes}
%%   \setsansfont{TeX Gyre Heros}
%%   \setmonofont{TeX Gyre Cursor}
%%   \usepackage{fontspec}
%%   \makeatletter\@ifpackageloaded{underscore}{}{\usepackage[strings]{underscore}}\makeatother
%%
\begingroup%
\makeatletter%
\begin{pgfpicture}%
\pgfpathrectangle{\pgfpointorigin}{\pgfqpoint{5.618889in}{3.494944in}}%
\pgfusepath{use as bounding box, clip}%
\begin{pgfscope}%
\pgfsetbuttcap%
\pgfsetmiterjoin%
\definecolor{currentfill}{rgb}{1.000000,1.000000,1.000000}%
\pgfsetfillcolor{currentfill}%
\pgfsetlinewidth{0.000000pt}%
\definecolor{currentstroke}{rgb}{1.000000,1.000000,1.000000}%
\pgfsetstrokecolor{currentstroke}%
\pgfsetdash{}{0pt}%
\pgfpathmoveto{\pgfqpoint{0.000000in}{0.000000in}}%
\pgfpathlineto{\pgfqpoint{5.618889in}{0.000000in}}%
\pgfpathlineto{\pgfqpoint{5.618889in}{3.494944in}}%
\pgfpathlineto{\pgfqpoint{0.000000in}{3.494944in}}%
\pgfpathlineto{\pgfqpoint{0.000000in}{0.000000in}}%
\pgfpathclose%
\pgfusepath{fill}%
\end{pgfscope}%
\begin{pgfscope}%
\pgfsetbuttcap%
\pgfsetmiterjoin%
\definecolor{currentfill}{rgb}{1.000000,1.000000,1.000000}%
\pgfsetfillcolor{currentfill}%
\pgfsetlinewidth{0.000000pt}%
\definecolor{currentstroke}{rgb}{0.000000,0.000000,0.000000}%
\pgfsetstrokecolor{currentstroke}%
\pgfsetstrokeopacity{0.000000}%
\pgfsetdash{}{0pt}%
\pgfpathmoveto{\pgfqpoint{0.713889in}{0.502778in}}%
\pgfpathlineto{\pgfqpoint{5.518889in}{0.502778in}}%
\pgfpathlineto{\pgfqpoint{5.518889in}{3.197778in}}%
\pgfpathlineto{\pgfqpoint{0.713889in}{3.197778in}}%
\pgfpathlineto{\pgfqpoint{0.713889in}{0.502778in}}%
\pgfpathclose%
\pgfusepath{fill}%
\end{pgfscope}%
\begin{pgfscope}%
\pgfpathrectangle{\pgfqpoint{0.713889in}{0.502778in}}{\pgfqpoint{4.805000in}{2.695000in}}%
\pgfusepath{clip}%
\pgfsetbuttcap%
\pgfsetroundjoin%
\definecolor{currentfill}{rgb}{0.121569,0.466667,0.705882}%
\pgfsetfillcolor{currentfill}%
\pgfsetlinewidth{1.003750pt}%
\definecolor{currentstroke}{rgb}{0.121569,0.466667,0.705882}%
\pgfsetstrokecolor{currentstroke}%
\pgfsetdash{}{0pt}%
\pgfsys@defobject{currentmarker}{\pgfqpoint{-0.015528in}{-0.015528in}}{\pgfqpoint{0.015528in}{0.015528in}}{%
\pgfpathmoveto{\pgfqpoint{0.000000in}{-0.015528in}}%
\pgfpathcurveto{\pgfqpoint{0.004118in}{-0.015528in}}{\pgfqpoint{0.008068in}{-0.013892in}}{\pgfqpoint{0.010980in}{-0.010980in}}%
\pgfpathcurveto{\pgfqpoint{0.013892in}{-0.008068in}}{\pgfqpoint{0.015528in}{-0.004118in}}{\pgfqpoint{0.015528in}{0.000000in}}%
\pgfpathcurveto{\pgfqpoint{0.015528in}{0.004118in}}{\pgfqpoint{0.013892in}{0.008068in}}{\pgfqpoint{0.010980in}{0.010980in}}%
\pgfpathcurveto{\pgfqpoint{0.008068in}{0.013892in}}{\pgfqpoint{0.004118in}{0.015528in}}{\pgfqpoint{0.000000in}{0.015528in}}%
\pgfpathcurveto{\pgfqpoint{-0.004118in}{0.015528in}}{\pgfqpoint{-0.008068in}{0.013892in}}{\pgfqpoint{-0.010980in}{0.010980in}}%
\pgfpathcurveto{\pgfqpoint{-0.013892in}{0.008068in}}{\pgfqpoint{-0.015528in}{0.004118in}}{\pgfqpoint{-0.015528in}{0.000000in}}%
\pgfpathcurveto{\pgfqpoint{-0.015528in}{-0.004118in}}{\pgfqpoint{-0.013892in}{-0.008068in}}{\pgfqpoint{-0.010980in}{-0.010980in}}%
\pgfpathcurveto{\pgfqpoint{-0.008068in}{-0.013892in}}{\pgfqpoint{-0.004118in}{-0.015528in}}{\pgfqpoint{0.000000in}{-0.015528in}}%
\pgfpathlineto{\pgfqpoint{0.000000in}{-0.015528in}}%
\pgfpathclose%
\pgfusepath{stroke,fill}%
}%
\begin{pgfscope}%
\pgfsys@transformshift{0.932298in}{0.625278in}%
\pgfsys@useobject{currentmarker}{}%
\end{pgfscope}%
\begin{pgfscope}%
\pgfsys@transformshift{0.932298in}{0.625278in}%
\pgfsys@useobject{currentmarker}{}%
\end{pgfscope}%
\begin{pgfscope}%
\pgfsys@transformshift{0.932298in}{0.625278in}%
\pgfsys@useobject{currentmarker}{}%
\end{pgfscope}%
\begin{pgfscope}%
\pgfsys@transformshift{0.932298in}{0.625278in}%
\pgfsys@useobject{currentmarker}{}%
\end{pgfscope}%
\begin{pgfscope}%
\pgfsys@transformshift{0.932466in}{0.625278in}%
\pgfsys@useobject{currentmarker}{}%
\end{pgfscope}%
\begin{pgfscope}%
\pgfsys@transformshift{0.932466in}{0.625278in}%
\pgfsys@useobject{currentmarker}{}%
\end{pgfscope}%
\begin{pgfscope}%
\pgfsys@transformshift{0.932466in}{0.625278in}%
\pgfsys@useobject{currentmarker}{}%
\end{pgfscope}%
\begin{pgfscope}%
\pgfsys@transformshift{0.932466in}{0.625278in}%
\pgfsys@useobject{currentmarker}{}%
\end{pgfscope}%
\begin{pgfscope}%
\pgfsys@transformshift{0.932635in}{0.625278in}%
\pgfsys@useobject{currentmarker}{}%
\end{pgfscope}%
\begin{pgfscope}%
\pgfsys@transformshift{0.932635in}{0.625278in}%
\pgfsys@useobject{currentmarker}{}%
\end{pgfscope}%
\begin{pgfscope}%
\pgfsys@transformshift{0.932635in}{0.625278in}%
\pgfsys@useobject{currentmarker}{}%
\end{pgfscope}%
\begin{pgfscope}%
\pgfsys@transformshift{0.932635in}{0.625278in}%
\pgfsys@useobject{currentmarker}{}%
\end{pgfscope}%
\begin{pgfscope}%
\pgfsys@transformshift{0.932804in}{0.625278in}%
\pgfsys@useobject{currentmarker}{}%
\end{pgfscope}%
\begin{pgfscope}%
\pgfsys@transformshift{0.932804in}{0.625278in}%
\pgfsys@useobject{currentmarker}{}%
\end{pgfscope}%
\begin{pgfscope}%
\pgfsys@transformshift{0.932804in}{0.625278in}%
\pgfsys@useobject{currentmarker}{}%
\end{pgfscope}%
\begin{pgfscope}%
\pgfsys@transformshift{0.932972in}{0.625278in}%
\pgfsys@useobject{currentmarker}{}%
\end{pgfscope}%
\begin{pgfscope}%
\pgfsys@transformshift{0.932972in}{0.625278in}%
\pgfsys@useobject{currentmarker}{}%
\end{pgfscope}%
\begin{pgfscope}%
\pgfsys@transformshift{0.932972in}{0.625278in}%
\pgfsys@useobject{currentmarker}{}%
\end{pgfscope}%
\begin{pgfscope}%
\pgfsys@transformshift{0.933141in}{0.625278in}%
\pgfsys@useobject{currentmarker}{}%
\end{pgfscope}%
\begin{pgfscope}%
\pgfsys@transformshift{0.933141in}{0.625278in}%
\pgfsys@useobject{currentmarker}{}%
\end{pgfscope}%
\begin{pgfscope}%
\pgfsys@transformshift{0.933309in}{0.625278in}%
\pgfsys@useobject{currentmarker}{}%
\end{pgfscope}%
\begin{pgfscope}%
\pgfsys@transformshift{0.933309in}{0.625278in}%
\pgfsys@useobject{currentmarker}{}%
\end{pgfscope}%
\begin{pgfscope}%
\pgfsys@transformshift{0.933478in}{0.625278in}%
\pgfsys@useobject{currentmarker}{}%
\end{pgfscope}%
\begin{pgfscope}%
\pgfsys@transformshift{0.933647in}{0.625278in}%
\pgfsys@useobject{currentmarker}{}%
\end{pgfscope}%
\begin{pgfscope}%
\pgfsys@transformshift{0.934153in}{0.625278in}%
\pgfsys@useobject{currentmarker}{}%
\end{pgfscope}%
\begin{pgfscope}%
\pgfsys@transformshift{0.934490in}{0.625278in}%
\pgfsys@useobject{currentmarker}{}%
\end{pgfscope}%
\begin{pgfscope}%
\pgfsys@transformshift{0.934658in}{0.625278in}%
\pgfsys@useobject{currentmarker}{}%
\end{pgfscope}%
\begin{pgfscope}%
\pgfsys@transformshift{0.934827in}{0.625278in}%
\pgfsys@useobject{currentmarker}{}%
\end{pgfscope}%
\begin{pgfscope}%
\pgfsys@transformshift{0.934827in}{0.625278in}%
\pgfsys@useobject{currentmarker}{}%
\end{pgfscope}%
\begin{pgfscope}%
\pgfsys@transformshift{0.934996in}{0.625278in}%
\pgfsys@useobject{currentmarker}{}%
\end{pgfscope}%
\begin{pgfscope}%
\pgfsys@transformshift{0.935164in}{0.625278in}%
\pgfsys@useobject{currentmarker}{}%
\end{pgfscope}%
\begin{pgfscope}%
\pgfsys@transformshift{0.935164in}{0.625278in}%
\pgfsys@useobject{currentmarker}{}%
\end{pgfscope}%
\begin{pgfscope}%
\pgfsys@transformshift{0.935164in}{0.625278in}%
\pgfsys@useobject{currentmarker}{}%
\end{pgfscope}%
\begin{pgfscope}%
\pgfsys@transformshift{0.935333in}{0.625278in}%
\pgfsys@useobject{currentmarker}{}%
\end{pgfscope}%
\begin{pgfscope}%
\pgfsys@transformshift{0.935501in}{0.625278in}%
\pgfsys@useobject{currentmarker}{}%
\end{pgfscope}%
\begin{pgfscope}%
\pgfsys@transformshift{0.935501in}{0.625278in}%
\pgfsys@useobject{currentmarker}{}%
\end{pgfscope}%
\begin{pgfscope}%
\pgfsys@transformshift{0.935501in}{0.625278in}%
\pgfsys@useobject{currentmarker}{}%
\end{pgfscope}%
\begin{pgfscope}%
\pgfsys@transformshift{0.935671in}{0.625278in}%
\pgfsys@useobject{currentmarker}{}%
\end{pgfscope}%
\begin{pgfscope}%
\pgfsys@transformshift{0.935839in}{0.625278in}%
\pgfsys@useobject{currentmarker}{}%
\end{pgfscope}%
\begin{pgfscope}%
\pgfsys@transformshift{0.935839in}{0.625278in}%
\pgfsys@useobject{currentmarker}{}%
\end{pgfscope}%
\begin{pgfscope}%
\pgfsys@transformshift{0.935839in}{0.625278in}%
\pgfsys@useobject{currentmarker}{}%
\end{pgfscope}%
\begin{pgfscope}%
\pgfsys@transformshift{0.936007in}{0.625278in}%
\pgfsys@useobject{currentmarker}{}%
\end{pgfscope}%
\begin{pgfscope}%
\pgfsys@transformshift{0.936007in}{0.625278in}%
\pgfsys@useobject{currentmarker}{}%
\end{pgfscope}%
\begin{pgfscope}%
\pgfsys@transformshift{0.936176in}{0.625278in}%
\pgfsys@useobject{currentmarker}{}%
\end{pgfscope}%
\begin{pgfscope}%
\pgfsys@transformshift{0.936176in}{0.625278in}%
\pgfsys@useobject{currentmarker}{}%
\end{pgfscope}%
\begin{pgfscope}%
\pgfsys@transformshift{0.936176in}{0.625278in}%
\pgfsys@useobject{currentmarker}{}%
\end{pgfscope}%
\begin{pgfscope}%
\pgfsys@transformshift{0.936345in}{0.625278in}%
\pgfsys@useobject{currentmarker}{}%
\end{pgfscope}%
\begin{pgfscope}%
\pgfsys@transformshift{0.936345in}{0.625278in}%
\pgfsys@useobject{currentmarker}{}%
\end{pgfscope}%
\begin{pgfscope}%
\pgfsys@transformshift{0.936514in}{0.625278in}%
\pgfsys@useobject{currentmarker}{}%
\end{pgfscope}%
\begin{pgfscope}%
\pgfsys@transformshift{0.936514in}{0.625278in}%
\pgfsys@useobject{currentmarker}{}%
\end{pgfscope}%
\begin{pgfscope}%
\pgfsys@transformshift{0.936514in}{0.625278in}%
\pgfsys@useobject{currentmarker}{}%
\end{pgfscope}%
\begin{pgfscope}%
\pgfsys@transformshift{0.936682in}{0.625278in}%
\pgfsys@useobject{currentmarker}{}%
\end{pgfscope}%
\begin{pgfscope}%
\pgfsys@transformshift{0.936682in}{0.625278in}%
\pgfsys@useobject{currentmarker}{}%
\end{pgfscope}%
\begin{pgfscope}%
\pgfsys@transformshift{0.936850in}{0.625278in}%
\pgfsys@useobject{currentmarker}{}%
\end{pgfscope}%
\begin{pgfscope}%
\pgfsys@transformshift{0.936850in}{0.625278in}%
\pgfsys@useobject{currentmarker}{}%
\end{pgfscope}%
\begin{pgfscope}%
\pgfsys@transformshift{0.936850in}{0.625278in}%
\pgfsys@useobject{currentmarker}{}%
\end{pgfscope}%
\begin{pgfscope}%
\pgfsys@transformshift{0.937019in}{0.625278in}%
\pgfsys@useobject{currentmarker}{}%
\end{pgfscope}%
\begin{pgfscope}%
\pgfsys@transformshift{0.937019in}{0.625278in}%
\pgfsys@useobject{currentmarker}{}%
\end{pgfscope}%
\begin{pgfscope}%
\pgfsys@transformshift{0.937188in}{0.625278in}%
\pgfsys@useobject{currentmarker}{}%
\end{pgfscope}%
\begin{pgfscope}%
\pgfsys@transformshift{0.937188in}{0.625278in}%
\pgfsys@useobject{currentmarker}{}%
\end{pgfscope}%
\begin{pgfscope}%
\pgfsys@transformshift{0.937188in}{0.625278in}%
\pgfsys@useobject{currentmarker}{}%
\end{pgfscope}%
\begin{pgfscope}%
\pgfsys@transformshift{0.937188in}{0.625278in}%
\pgfsys@useobject{currentmarker}{}%
\end{pgfscope}%
\begin{pgfscope}%
\pgfsys@transformshift{0.937525in}{0.625278in}%
\pgfsys@useobject{currentmarker}{}%
\end{pgfscope}%
\begin{pgfscope}%
\pgfsys@transformshift{0.937525in}{0.625278in}%
\pgfsys@useobject{currentmarker}{}%
\end{pgfscope}%
\begin{pgfscope}%
\pgfsys@transformshift{0.937525in}{0.625278in}%
\pgfsys@useobject{currentmarker}{}%
\end{pgfscope}%
\begin{pgfscope}%
\pgfsys@transformshift{0.937525in}{0.625278in}%
\pgfsys@useobject{currentmarker}{}%
\end{pgfscope}%
\begin{pgfscope}%
\pgfsys@transformshift{0.937863in}{0.625278in}%
\pgfsys@useobject{currentmarker}{}%
\end{pgfscope}%
\begin{pgfscope}%
\pgfsys@transformshift{0.937863in}{0.625278in}%
\pgfsys@useobject{currentmarker}{}%
\end{pgfscope}%
\begin{pgfscope}%
\pgfsys@transformshift{0.937863in}{0.625278in}%
\pgfsys@useobject{currentmarker}{}%
\end{pgfscope}%
\begin{pgfscope}%
\pgfsys@transformshift{0.937863in}{0.625278in}%
\pgfsys@useobject{currentmarker}{}%
\end{pgfscope}%
\begin{pgfscope}%
\pgfsys@transformshift{0.938200in}{0.625278in}%
\pgfsys@useobject{currentmarker}{}%
\end{pgfscope}%
\begin{pgfscope}%
\pgfsys@transformshift{0.938200in}{0.625278in}%
\pgfsys@useobject{currentmarker}{}%
\end{pgfscope}%
\begin{pgfscope}%
\pgfsys@transformshift{0.938200in}{0.625278in}%
\pgfsys@useobject{currentmarker}{}%
\end{pgfscope}%
\begin{pgfscope}%
\pgfsys@transformshift{0.938200in}{0.625278in}%
\pgfsys@useobject{currentmarker}{}%
\end{pgfscope}%
\begin{pgfscope}%
\pgfsys@transformshift{0.938537in}{0.625278in}%
\pgfsys@useobject{currentmarker}{}%
\end{pgfscope}%
\begin{pgfscope}%
\pgfsys@transformshift{0.938537in}{0.625278in}%
\pgfsys@useobject{currentmarker}{}%
\end{pgfscope}%
\begin{pgfscope}%
\pgfsys@transformshift{0.938537in}{0.625278in}%
\pgfsys@useobject{currentmarker}{}%
\end{pgfscope}%
\begin{pgfscope}%
\pgfsys@transformshift{0.938537in}{0.625278in}%
\pgfsys@useobject{currentmarker}{}%
\end{pgfscope}%
\begin{pgfscope}%
\pgfsys@transformshift{0.938874in}{0.625278in}%
\pgfsys@useobject{currentmarker}{}%
\end{pgfscope}%
\begin{pgfscope}%
\pgfsys@transformshift{0.938874in}{0.625278in}%
\pgfsys@useobject{currentmarker}{}%
\end{pgfscope}%
\begin{pgfscope}%
\pgfsys@transformshift{0.938874in}{0.625278in}%
\pgfsys@useobject{currentmarker}{}%
\end{pgfscope}%
\begin{pgfscope}%
\pgfsys@transformshift{0.938874in}{0.625278in}%
\pgfsys@useobject{currentmarker}{}%
\end{pgfscope}%
\begin{pgfscope}%
\pgfsys@transformshift{0.939211in}{0.625278in}%
\pgfsys@useobject{currentmarker}{}%
\end{pgfscope}%
\begin{pgfscope}%
\pgfsys@transformshift{0.939211in}{0.625278in}%
\pgfsys@useobject{currentmarker}{}%
\end{pgfscope}%
\begin{pgfscope}%
\pgfsys@transformshift{0.939211in}{0.625278in}%
\pgfsys@useobject{currentmarker}{}%
\end{pgfscope}%
\begin{pgfscope}%
\pgfsys@transformshift{0.939211in}{0.625278in}%
\pgfsys@useobject{currentmarker}{}%
\end{pgfscope}%
\begin{pgfscope}%
\pgfsys@transformshift{0.939549in}{0.625278in}%
\pgfsys@useobject{currentmarker}{}%
\end{pgfscope}%
\begin{pgfscope}%
\pgfsys@transformshift{0.939549in}{0.625278in}%
\pgfsys@useobject{currentmarker}{}%
\end{pgfscope}%
\begin{pgfscope}%
\pgfsys@transformshift{0.939549in}{0.625278in}%
\pgfsys@useobject{currentmarker}{}%
\end{pgfscope}%
\begin{pgfscope}%
\pgfsys@transformshift{0.939549in}{0.625278in}%
\pgfsys@useobject{currentmarker}{}%
\end{pgfscope}%
\begin{pgfscope}%
\pgfsys@transformshift{0.939886in}{0.625278in}%
\pgfsys@useobject{currentmarker}{}%
\end{pgfscope}%
\begin{pgfscope}%
\pgfsys@transformshift{0.939886in}{0.625278in}%
\pgfsys@useobject{currentmarker}{}%
\end{pgfscope}%
\begin{pgfscope}%
\pgfsys@transformshift{0.939886in}{0.625278in}%
\pgfsys@useobject{currentmarker}{}%
\end{pgfscope}%
\begin{pgfscope}%
\pgfsys@transformshift{0.940055in}{0.625278in}%
\pgfsys@useobject{currentmarker}{}%
\end{pgfscope}%
\begin{pgfscope}%
\pgfsys@transformshift{0.940223in}{0.625278in}%
\pgfsys@useobject{currentmarker}{}%
\end{pgfscope}%
\begin{pgfscope}%
\pgfsys@transformshift{0.940223in}{0.625278in}%
\pgfsys@useobject{currentmarker}{}%
\end{pgfscope}%
\begin{pgfscope}%
\pgfsys@transformshift{0.940392in}{0.625278in}%
\pgfsys@useobject{currentmarker}{}%
\end{pgfscope}%
\begin{pgfscope}%
\pgfsys@transformshift{0.940560in}{0.625278in}%
\pgfsys@useobject{currentmarker}{}%
\end{pgfscope}%
\begin{pgfscope}%
\pgfsys@transformshift{0.940560in}{0.625278in}%
\pgfsys@useobject{currentmarker}{}%
\end{pgfscope}%
\begin{pgfscope}%
\pgfsys@transformshift{0.940729in}{0.625278in}%
\pgfsys@useobject{currentmarker}{}%
\end{pgfscope}%
\begin{pgfscope}%
\pgfsys@transformshift{0.940898in}{0.625278in}%
\pgfsys@useobject{currentmarker}{}%
\end{pgfscope}%
\begin{pgfscope}%
\pgfsys@transformshift{0.940898in}{0.625278in}%
\pgfsys@useobject{currentmarker}{}%
\end{pgfscope}%
\begin{pgfscope}%
\pgfsys@transformshift{0.941066in}{0.625278in}%
\pgfsys@useobject{currentmarker}{}%
\end{pgfscope}%
\begin{pgfscope}%
\pgfsys@transformshift{0.941235in}{0.625278in}%
\pgfsys@useobject{currentmarker}{}%
\end{pgfscope}%
\begin{pgfscope}%
\pgfsys@transformshift{0.941235in}{0.625278in}%
\pgfsys@useobject{currentmarker}{}%
\end{pgfscope}%
\begin{pgfscope}%
\pgfsys@transformshift{0.941403in}{0.625278in}%
\pgfsys@useobject{currentmarker}{}%
\end{pgfscope}%
\begin{pgfscope}%
\pgfsys@transformshift{0.941573in}{0.625278in}%
\pgfsys@useobject{currentmarker}{}%
\end{pgfscope}%
\begin{pgfscope}%
\pgfsys@transformshift{0.941573in}{0.625278in}%
\pgfsys@useobject{currentmarker}{}%
\end{pgfscope}%
\begin{pgfscope}%
\pgfsys@transformshift{0.941573in}{0.625278in}%
\pgfsys@useobject{currentmarker}{}%
\end{pgfscope}%
\begin{pgfscope}%
\pgfsys@transformshift{0.941741in}{0.625278in}%
\pgfsys@useobject{currentmarker}{}%
\end{pgfscope}%
\begin{pgfscope}%
\pgfsys@transformshift{0.941909in}{0.625278in}%
\pgfsys@useobject{currentmarker}{}%
\end{pgfscope}%
\begin{pgfscope}%
\pgfsys@transformshift{0.941909in}{0.625278in}%
\pgfsys@useobject{currentmarker}{}%
\end{pgfscope}%
\begin{pgfscope}%
\pgfsys@transformshift{0.941909in}{0.625278in}%
\pgfsys@useobject{currentmarker}{}%
\end{pgfscope}%
\begin{pgfscope}%
\pgfsys@transformshift{0.942078in}{0.625278in}%
\pgfsys@useobject{currentmarker}{}%
\end{pgfscope}%
\begin{pgfscope}%
\pgfsys@transformshift{0.942247in}{0.625278in}%
\pgfsys@useobject{currentmarker}{}%
\end{pgfscope}%
\begin{pgfscope}%
\pgfsys@transformshift{0.942247in}{0.625278in}%
\pgfsys@useobject{currentmarker}{}%
\end{pgfscope}%
\begin{pgfscope}%
\pgfsys@transformshift{0.942247in}{0.625278in}%
\pgfsys@useobject{currentmarker}{}%
\end{pgfscope}%
\begin{pgfscope}%
\pgfsys@transformshift{0.942416in}{0.625278in}%
\pgfsys@useobject{currentmarker}{}%
\end{pgfscope}%
\begin{pgfscope}%
\pgfsys@transformshift{0.942584in}{0.625278in}%
\pgfsys@useobject{currentmarker}{}%
\end{pgfscope}%
\begin{pgfscope}%
\pgfsys@transformshift{0.942584in}{0.625278in}%
\pgfsys@useobject{currentmarker}{}%
\end{pgfscope}%
\begin{pgfscope}%
\pgfsys@transformshift{0.942584in}{0.625278in}%
\pgfsys@useobject{currentmarker}{}%
\end{pgfscope}%
\begin{pgfscope}%
\pgfsys@transformshift{0.942752in}{0.625278in}%
\pgfsys@useobject{currentmarker}{}%
\end{pgfscope}%
\begin{pgfscope}%
\pgfsys@transformshift{0.942752in}{0.625278in}%
\pgfsys@useobject{currentmarker}{}%
\end{pgfscope}%
\begin{pgfscope}%
\pgfsys@transformshift{0.942921in}{0.625278in}%
\pgfsys@useobject{currentmarker}{}%
\end{pgfscope}%
\begin{pgfscope}%
\pgfsys@transformshift{0.942921in}{0.625278in}%
\pgfsys@useobject{currentmarker}{}%
\end{pgfscope}%
\begin{pgfscope}%
\pgfsys@transformshift{0.942921in}{0.625278in}%
\pgfsys@useobject{currentmarker}{}%
\end{pgfscope}%
\begin{pgfscope}%
\pgfsys@transformshift{0.943090in}{0.625278in}%
\pgfsys@useobject{currentmarker}{}%
\end{pgfscope}%
\begin{pgfscope}%
\pgfsys@transformshift{0.943090in}{0.625278in}%
\pgfsys@useobject{currentmarker}{}%
\end{pgfscope}%
\begin{pgfscope}%
\pgfsys@transformshift{0.943259in}{0.625278in}%
\pgfsys@useobject{currentmarker}{}%
\end{pgfscope}%
\begin{pgfscope}%
\pgfsys@transformshift{0.943259in}{0.625278in}%
\pgfsys@useobject{currentmarker}{}%
\end{pgfscope}%
\begin{pgfscope}%
\pgfsys@transformshift{0.943259in}{0.625278in}%
\pgfsys@useobject{currentmarker}{}%
\end{pgfscope}%
\begin{pgfscope}%
\pgfsys@transformshift{0.943427in}{0.625278in}%
\pgfsys@useobject{currentmarker}{}%
\end{pgfscope}%
\begin{pgfscope}%
\pgfsys@transformshift{0.943427in}{0.625278in}%
\pgfsys@useobject{currentmarker}{}%
\end{pgfscope}%
\begin{pgfscope}%
\pgfsys@transformshift{0.943595in}{0.625278in}%
\pgfsys@useobject{currentmarker}{}%
\end{pgfscope}%
\begin{pgfscope}%
\pgfsys@transformshift{0.943595in}{0.625278in}%
\pgfsys@useobject{currentmarker}{}%
\end{pgfscope}%
\begin{pgfscope}%
\pgfsys@transformshift{0.943595in}{0.625278in}%
\pgfsys@useobject{currentmarker}{}%
\end{pgfscope}%
\begin{pgfscope}%
\pgfsys@transformshift{0.943765in}{0.625278in}%
\pgfsys@useobject{currentmarker}{}%
\end{pgfscope}%
\begin{pgfscope}%
\pgfsys@transformshift{0.943765in}{0.625278in}%
\pgfsys@useobject{currentmarker}{}%
\end{pgfscope}%
\begin{pgfscope}%
\pgfsys@transformshift{0.943933in}{0.625278in}%
\pgfsys@useobject{currentmarker}{}%
\end{pgfscope}%
\begin{pgfscope}%
\pgfsys@transformshift{0.943933in}{0.625278in}%
\pgfsys@useobject{currentmarker}{}%
\end{pgfscope}%
\begin{pgfscope}%
\pgfsys@transformshift{0.943933in}{0.625278in}%
\pgfsys@useobject{currentmarker}{}%
\end{pgfscope}%
\begin{pgfscope}%
\pgfsys@transformshift{0.943933in}{0.625278in}%
\pgfsys@useobject{currentmarker}{}%
\end{pgfscope}%
\begin{pgfscope}%
\pgfsys@transformshift{0.944270in}{0.625278in}%
\pgfsys@useobject{currentmarker}{}%
\end{pgfscope}%
\begin{pgfscope}%
\pgfsys@transformshift{0.944270in}{0.625278in}%
\pgfsys@useobject{currentmarker}{}%
\end{pgfscope}%
\begin{pgfscope}%
\pgfsys@transformshift{0.944270in}{0.625278in}%
\pgfsys@useobject{currentmarker}{}%
\end{pgfscope}%
\begin{pgfscope}%
\pgfsys@transformshift{0.944608in}{0.625278in}%
\pgfsys@useobject{currentmarker}{}%
\end{pgfscope}%
\begin{pgfscope}%
\pgfsys@transformshift{0.944608in}{0.625278in}%
\pgfsys@useobject{currentmarker}{}%
\end{pgfscope}%
\begin{pgfscope}%
\pgfsys@transformshift{0.944945in}{0.625278in}%
\pgfsys@useobject{currentmarker}{}%
\end{pgfscope}%
\begin{pgfscope}%
\pgfsys@transformshift{0.944945in}{0.625278in}%
\pgfsys@useobject{currentmarker}{}%
\end{pgfscope}%
\begin{pgfscope}%
\pgfsys@transformshift{0.945282in}{0.625278in}%
\pgfsys@useobject{currentmarker}{}%
\end{pgfscope}%
\begin{pgfscope}%
\pgfsys@transformshift{0.945282in}{0.625278in}%
\pgfsys@useobject{currentmarker}{}%
\end{pgfscope}%
\begin{pgfscope}%
\pgfsys@transformshift{0.945619in}{0.625278in}%
\pgfsys@useobject{currentmarker}{}%
\end{pgfscope}%
\begin{pgfscope}%
\pgfsys@transformshift{0.945619in}{0.625278in}%
\pgfsys@useobject{currentmarker}{}%
\end{pgfscope}%
\begin{pgfscope}%
\pgfsys@transformshift{0.945957in}{0.625278in}%
\pgfsys@useobject{currentmarker}{}%
\end{pgfscope}%
\begin{pgfscope}%
\pgfsys@transformshift{0.945957in}{0.625278in}%
\pgfsys@useobject{currentmarker}{}%
\end{pgfscope}%
\begin{pgfscope}%
\pgfsys@transformshift{0.946294in}{0.625278in}%
\pgfsys@useobject{currentmarker}{}%
\end{pgfscope}%
\begin{pgfscope}%
\pgfsys@transformshift{0.946294in}{0.625278in}%
\pgfsys@useobject{currentmarker}{}%
\end{pgfscope}%
\begin{pgfscope}%
\pgfsys@transformshift{0.946631in}{0.625278in}%
\pgfsys@useobject{currentmarker}{}%
\end{pgfscope}%
\begin{pgfscope}%
\pgfsys@transformshift{0.946968in}{0.625278in}%
\pgfsys@useobject{currentmarker}{}%
\end{pgfscope}%
\begin{pgfscope}%
\pgfsys@transformshift{0.947305in}{0.625278in}%
\pgfsys@useobject{currentmarker}{}%
\end{pgfscope}%
\begin{pgfscope}%
\pgfsys@transformshift{0.947643in}{0.625278in}%
\pgfsys@useobject{currentmarker}{}%
\end{pgfscope}%
\begin{pgfscope}%
\pgfsys@transformshift{0.947980in}{0.625278in}%
\pgfsys@useobject{currentmarker}{}%
\end{pgfscope}%
\begin{pgfscope}%
\pgfsys@transformshift{0.948318in}{0.625278in}%
\pgfsys@useobject{currentmarker}{}%
\end{pgfscope}%
\begin{pgfscope}%
\pgfsys@transformshift{0.948318in}{0.625278in}%
\pgfsys@useobject{currentmarker}{}%
\end{pgfscope}%
\begin{pgfscope}%
\pgfsys@transformshift{0.948654in}{0.625278in}%
\pgfsys@useobject{currentmarker}{}%
\end{pgfscope}%
\begin{pgfscope}%
\pgfsys@transformshift{0.948654in}{0.625278in}%
\pgfsys@useobject{currentmarker}{}%
\end{pgfscope}%
\begin{pgfscope}%
\pgfsys@transformshift{0.948992in}{0.625278in}%
\pgfsys@useobject{currentmarker}{}%
\end{pgfscope}%
\begin{pgfscope}%
\pgfsys@transformshift{0.948992in}{0.625278in}%
\pgfsys@useobject{currentmarker}{}%
\end{pgfscope}%
\begin{pgfscope}%
\pgfsys@transformshift{0.949329in}{0.625278in}%
\pgfsys@useobject{currentmarker}{}%
\end{pgfscope}%
\begin{pgfscope}%
\pgfsys@transformshift{0.949329in}{0.625278in}%
\pgfsys@useobject{currentmarker}{}%
\end{pgfscope}%
\begin{pgfscope}%
\pgfsys@transformshift{0.949497in}{0.625278in}%
\pgfsys@useobject{currentmarker}{}%
\end{pgfscope}%
\begin{pgfscope}%
\pgfsys@transformshift{0.949667in}{0.625278in}%
\pgfsys@useobject{currentmarker}{}%
\end{pgfscope}%
\begin{pgfscope}%
\pgfsys@transformshift{0.949667in}{0.625278in}%
\pgfsys@useobject{currentmarker}{}%
\end{pgfscope}%
\begin{pgfscope}%
\pgfsys@transformshift{0.949835in}{0.625278in}%
\pgfsys@useobject{currentmarker}{}%
\end{pgfscope}%
\begin{pgfscope}%
\pgfsys@transformshift{0.950004in}{0.625278in}%
\pgfsys@useobject{currentmarker}{}%
\end{pgfscope}%
\begin{pgfscope}%
\pgfsys@transformshift{0.950004in}{0.625278in}%
\pgfsys@useobject{currentmarker}{}%
\end{pgfscope}%
\begin{pgfscope}%
\pgfsys@transformshift{0.950172in}{0.625278in}%
\pgfsys@useobject{currentmarker}{}%
\end{pgfscope}%
\begin{pgfscope}%
\pgfsys@transformshift{0.950341in}{0.625278in}%
\pgfsys@useobject{currentmarker}{}%
\end{pgfscope}%
\begin{pgfscope}%
\pgfsys@transformshift{0.950341in}{0.625278in}%
\pgfsys@useobject{currentmarker}{}%
\end{pgfscope}%
\begin{pgfscope}%
\pgfsys@transformshift{0.950510in}{0.625278in}%
\pgfsys@useobject{currentmarker}{}%
\end{pgfscope}%
\begin{pgfscope}%
\pgfsys@transformshift{0.950678in}{0.625278in}%
\pgfsys@useobject{currentmarker}{}%
\end{pgfscope}%
\begin{pgfscope}%
\pgfsys@transformshift{0.950678in}{0.625278in}%
\pgfsys@useobject{currentmarker}{}%
\end{pgfscope}%
\begin{pgfscope}%
\pgfsys@transformshift{0.950678in}{0.625278in}%
\pgfsys@useobject{currentmarker}{}%
\end{pgfscope}%
\begin{pgfscope}%
\pgfsys@transformshift{0.951015in}{0.625278in}%
\pgfsys@useobject{currentmarker}{}%
\end{pgfscope}%
\begin{pgfscope}%
\pgfsys@transformshift{0.951015in}{0.625278in}%
\pgfsys@useobject{currentmarker}{}%
\end{pgfscope}%
\begin{pgfscope}%
\pgfsys@transformshift{0.951353in}{0.625278in}%
\pgfsys@useobject{currentmarker}{}%
\end{pgfscope}%
\begin{pgfscope}%
\pgfsys@transformshift{0.951353in}{0.625278in}%
\pgfsys@useobject{currentmarker}{}%
\end{pgfscope}%
\begin{pgfscope}%
\pgfsys@transformshift{0.951690in}{0.625278in}%
\pgfsys@useobject{currentmarker}{}%
\end{pgfscope}%
\begin{pgfscope}%
\pgfsys@transformshift{0.951690in}{0.625278in}%
\pgfsys@useobject{currentmarker}{}%
\end{pgfscope}%
\begin{pgfscope}%
\pgfsys@transformshift{0.952027in}{0.625278in}%
\pgfsys@useobject{currentmarker}{}%
\end{pgfscope}%
\begin{pgfscope}%
\pgfsys@transformshift{0.952027in}{0.625278in}%
\pgfsys@useobject{currentmarker}{}%
\end{pgfscope}%
\begin{pgfscope}%
\pgfsys@transformshift{0.952364in}{0.625278in}%
\pgfsys@useobject{currentmarker}{}%
\end{pgfscope}%
\begin{pgfscope}%
\pgfsys@transformshift{0.952364in}{0.625278in}%
\pgfsys@useobject{currentmarker}{}%
\end{pgfscope}%
\begin{pgfscope}%
\pgfsys@transformshift{0.952702in}{0.625278in}%
\pgfsys@useobject{currentmarker}{}%
\end{pgfscope}%
\begin{pgfscope}%
\pgfsys@transformshift{0.952702in}{0.625278in}%
\pgfsys@useobject{currentmarker}{}%
\end{pgfscope}%
\begin{pgfscope}%
\pgfsys@transformshift{0.953039in}{0.625278in}%
\pgfsys@useobject{currentmarker}{}%
\end{pgfscope}%
\begin{pgfscope}%
\pgfsys@transformshift{0.953039in}{0.625278in}%
\pgfsys@useobject{currentmarker}{}%
\end{pgfscope}%
\begin{pgfscope}%
\pgfsys@transformshift{0.953377in}{0.625278in}%
\pgfsys@useobject{currentmarker}{}%
\end{pgfscope}%
\begin{pgfscope}%
\pgfsys@transformshift{0.953713in}{0.625278in}%
\pgfsys@useobject{currentmarker}{}%
\end{pgfscope}%
\begin{pgfscope}%
\pgfsys@transformshift{0.954051in}{0.625278in}%
\pgfsys@useobject{currentmarker}{}%
\end{pgfscope}%
\begin{pgfscope}%
\pgfsys@transformshift{0.954388in}{0.625278in}%
\pgfsys@useobject{currentmarker}{}%
\end{pgfscope}%
\begin{pgfscope}%
\pgfsys@transformshift{0.954725in}{0.625278in}%
\pgfsys@useobject{currentmarker}{}%
\end{pgfscope}%
\begin{pgfscope}%
\pgfsys@transformshift{0.955063in}{0.625278in}%
\pgfsys@useobject{currentmarker}{}%
\end{pgfscope}%
\begin{pgfscope}%
\pgfsys@transformshift{0.955063in}{0.625278in}%
\pgfsys@useobject{currentmarker}{}%
\end{pgfscope}%
\begin{pgfscope}%
\pgfsys@transformshift{0.955399in}{0.625278in}%
\pgfsys@useobject{currentmarker}{}%
\end{pgfscope}%
\begin{pgfscope}%
\pgfsys@transformshift{0.955399in}{0.625278in}%
\pgfsys@useobject{currentmarker}{}%
\end{pgfscope}%
\begin{pgfscope}%
\pgfsys@transformshift{0.955737in}{0.625278in}%
\pgfsys@useobject{currentmarker}{}%
\end{pgfscope}%
\begin{pgfscope}%
\pgfsys@transformshift{0.955737in}{0.625278in}%
\pgfsys@useobject{currentmarker}{}%
\end{pgfscope}%
\begin{pgfscope}%
\pgfsys@transformshift{0.956074in}{0.625278in}%
\pgfsys@useobject{currentmarker}{}%
\end{pgfscope}%
\begin{pgfscope}%
\pgfsys@transformshift{0.956074in}{0.625278in}%
\pgfsys@useobject{currentmarker}{}%
\end{pgfscope}%
\begin{pgfscope}%
\pgfsys@transformshift{0.956243in}{0.625278in}%
\pgfsys@useobject{currentmarker}{}%
\end{pgfscope}%
\begin{pgfscope}%
\pgfsys@transformshift{0.956412in}{0.625278in}%
\pgfsys@useobject{currentmarker}{}%
\end{pgfscope}%
\begin{pgfscope}%
\pgfsys@transformshift{0.956412in}{0.625278in}%
\pgfsys@useobject{currentmarker}{}%
\end{pgfscope}%
\begin{pgfscope}%
\pgfsys@transformshift{0.956580in}{0.625278in}%
\pgfsys@useobject{currentmarker}{}%
\end{pgfscope}%
\begin{pgfscope}%
\pgfsys@transformshift{0.956749in}{0.625278in}%
\pgfsys@useobject{currentmarker}{}%
\end{pgfscope}%
\begin{pgfscope}%
\pgfsys@transformshift{0.956749in}{0.625278in}%
\pgfsys@useobject{currentmarker}{}%
\end{pgfscope}%
\begin{pgfscope}%
\pgfsys@transformshift{0.956917in}{0.625278in}%
\pgfsys@useobject{currentmarker}{}%
\end{pgfscope}%
\begin{pgfscope}%
\pgfsys@transformshift{0.957086in}{0.625278in}%
\pgfsys@useobject{currentmarker}{}%
\end{pgfscope}%
\begin{pgfscope}%
\pgfsys@transformshift{0.957086in}{0.625278in}%
\pgfsys@useobject{currentmarker}{}%
\end{pgfscope}%
\begin{pgfscope}%
\pgfsys@transformshift{0.957255in}{0.625278in}%
\pgfsys@useobject{currentmarker}{}%
\end{pgfscope}%
\begin{pgfscope}%
\pgfsys@transformshift{0.957423in}{0.625278in}%
\pgfsys@useobject{currentmarker}{}%
\end{pgfscope}%
\begin{pgfscope}%
\pgfsys@transformshift{0.957423in}{0.625278in}%
\pgfsys@useobject{currentmarker}{}%
\end{pgfscope}%
\begin{pgfscope}%
\pgfsys@transformshift{0.957423in}{0.625278in}%
\pgfsys@useobject{currentmarker}{}%
\end{pgfscope}%
\begin{pgfscope}%
\pgfsys@transformshift{0.957761in}{0.625278in}%
\pgfsys@useobject{currentmarker}{}%
\end{pgfscope}%
\begin{pgfscope}%
\pgfsys@transformshift{0.957761in}{0.625278in}%
\pgfsys@useobject{currentmarker}{}%
\end{pgfscope}%
\begin{pgfscope}%
\pgfsys@transformshift{0.958098in}{0.625278in}%
\pgfsys@useobject{currentmarker}{}%
\end{pgfscope}%
\begin{pgfscope}%
\pgfsys@transformshift{0.958098in}{0.625278in}%
\pgfsys@useobject{currentmarker}{}%
\end{pgfscope}%
\begin{pgfscope}%
\pgfsys@transformshift{0.958435in}{0.625278in}%
\pgfsys@useobject{currentmarker}{}%
\end{pgfscope}%
\begin{pgfscope}%
\pgfsys@transformshift{0.958435in}{0.625278in}%
\pgfsys@useobject{currentmarker}{}%
\end{pgfscope}%
\begin{pgfscope}%
\pgfsys@transformshift{0.958772in}{0.625278in}%
\pgfsys@useobject{currentmarker}{}%
\end{pgfscope}%
\begin{pgfscope}%
\pgfsys@transformshift{0.958772in}{0.625278in}%
\pgfsys@useobject{currentmarker}{}%
\end{pgfscope}%
\begin{pgfscope}%
\pgfsys@transformshift{0.959109in}{0.625278in}%
\pgfsys@useobject{currentmarker}{}%
\end{pgfscope}%
\begin{pgfscope}%
\pgfsys@transformshift{0.959109in}{0.625278in}%
\pgfsys@useobject{currentmarker}{}%
\end{pgfscope}%
\begin{pgfscope}%
\pgfsys@transformshift{0.959447in}{0.625278in}%
\pgfsys@useobject{currentmarker}{}%
\end{pgfscope}%
\begin{pgfscope}%
\pgfsys@transformshift{0.959447in}{0.625278in}%
\pgfsys@useobject{currentmarker}{}%
\end{pgfscope}%
\begin{pgfscope}%
\pgfsys@transformshift{0.959784in}{0.625278in}%
\pgfsys@useobject{currentmarker}{}%
\end{pgfscope}%
\begin{pgfscope}%
\pgfsys@transformshift{0.959784in}{0.625278in}%
\pgfsys@useobject{currentmarker}{}%
\end{pgfscope}%
\begin{pgfscope}%
\pgfsys@transformshift{0.960122in}{0.625278in}%
\pgfsys@useobject{currentmarker}{}%
\end{pgfscope}%
\begin{pgfscope}%
\pgfsys@transformshift{0.960458in}{0.625278in}%
\pgfsys@useobject{currentmarker}{}%
\end{pgfscope}%
\begin{pgfscope}%
\pgfsys@transformshift{0.960796in}{0.625278in}%
\pgfsys@useobject{currentmarker}{}%
\end{pgfscope}%
\begin{pgfscope}%
\pgfsys@transformshift{0.961133in}{0.625278in}%
\pgfsys@useobject{currentmarker}{}%
\end{pgfscope}%
\begin{pgfscope}%
\pgfsys@transformshift{0.961471in}{0.625278in}%
\pgfsys@useobject{currentmarker}{}%
\end{pgfscope}%
\begin{pgfscope}%
\pgfsys@transformshift{0.961808in}{0.625278in}%
\pgfsys@useobject{currentmarker}{}%
\end{pgfscope}%
\begin{pgfscope}%
\pgfsys@transformshift{0.961808in}{0.625278in}%
\pgfsys@useobject{currentmarker}{}%
\end{pgfscope}%
\begin{pgfscope}%
\pgfsys@transformshift{0.962145in}{0.625278in}%
\pgfsys@useobject{currentmarker}{}%
\end{pgfscope}%
\begin{pgfscope}%
\pgfsys@transformshift{0.962145in}{0.625278in}%
\pgfsys@useobject{currentmarker}{}%
\end{pgfscope}%
\begin{pgfscope}%
\pgfsys@transformshift{0.962988in}{0.625278in}%
\pgfsys@useobject{currentmarker}{}%
\end{pgfscope}%
\begin{pgfscope}%
\pgfsys@transformshift{0.963494in}{0.625278in}%
\pgfsys@useobject{currentmarker}{}%
\end{pgfscope}%
\begin{pgfscope}%
\pgfsys@transformshift{0.964176in}{0.625278in}%
\pgfsys@useobject{currentmarker}{}%
\end{pgfscope}%
\begin{pgfscope}%
\pgfsys@transformshift{0.964344in}{0.625278in}%
\pgfsys@useobject{currentmarker}{}%
\end{pgfscope}%
\begin{pgfscope}%
\pgfsys@transformshift{0.964681in}{0.625278in}%
\pgfsys@useobject{currentmarker}{}%
\end{pgfscope}%
\begin{pgfscope}%
\pgfsys@transformshift{0.965019in}{0.625278in}%
\pgfsys@useobject{currentmarker}{}%
\end{pgfscope}%
\begin{pgfscope}%
\pgfsys@transformshift{0.965356in}{0.625278in}%
\pgfsys@useobject{currentmarker}{}%
\end{pgfscope}%
\begin{pgfscope}%
\pgfsys@transformshift{0.966038in}{0.625278in}%
\pgfsys@useobject{currentmarker}{}%
\end{pgfscope}%
\begin{pgfscope}%
\pgfsys@transformshift{0.966375in}{0.625278in}%
\pgfsys@useobject{currentmarker}{}%
\end{pgfscope}%
\begin{pgfscope}%
\pgfsys@transformshift{0.966712in}{0.625278in}%
\pgfsys@useobject{currentmarker}{}%
\end{pgfscope}%
\begin{pgfscope}%
\pgfsys@transformshift{0.967049in}{0.625278in}%
\pgfsys@useobject{currentmarker}{}%
\end{pgfscope}%
\begin{pgfscope}%
\pgfsys@transformshift{0.967548in}{0.625278in}%
\pgfsys@useobject{currentmarker}{}%
\end{pgfscope}%
\begin{pgfscope}%
\pgfsys@transformshift{0.967886in}{0.625278in}%
\pgfsys@useobject{currentmarker}{}%
\end{pgfscope}%
\begin{pgfscope}%
\pgfsys@transformshift{0.968061in}{0.625278in}%
\pgfsys@useobject{currentmarker}{}%
\end{pgfscope}%
\begin{pgfscope}%
\pgfsys@transformshift{0.968398in}{0.625278in}%
\pgfsys@useobject{currentmarker}{}%
\end{pgfscope}%
\begin{pgfscope}%
\pgfsys@transformshift{0.968398in}{0.625278in}%
\pgfsys@useobject{currentmarker}{}%
\end{pgfscope}%
\begin{pgfscope}%
\pgfsys@transformshift{0.968736in}{0.625278in}%
\pgfsys@useobject{currentmarker}{}%
\end{pgfscope}%
\begin{pgfscope}%
\pgfsys@transformshift{0.973464in}{0.625278in}%
\pgfsys@useobject{currentmarker}{}%
\end{pgfscope}%
\begin{pgfscope}%
\pgfsys@transformshift{0.973626in}{0.625278in}%
\pgfsys@useobject{currentmarker}{}%
\end{pgfscope}%
\begin{pgfscope}%
\pgfsys@transformshift{0.973802in}{0.625278in}%
\pgfsys@useobject{currentmarker}{}%
\end{pgfscope}%
\begin{pgfscope}%
\pgfsys@transformshift{0.973802in}{0.625278in}%
\pgfsys@useobject{currentmarker}{}%
\end{pgfscope}%
\begin{pgfscope}%
\pgfsys@transformshift{0.973802in}{0.625278in}%
\pgfsys@useobject{currentmarker}{}%
\end{pgfscope}%
\begin{pgfscope}%
\pgfsys@transformshift{0.973963in}{0.625278in}%
\pgfsys@useobject{currentmarker}{}%
\end{pgfscope}%
\begin{pgfscope}%
\pgfsys@transformshift{0.974139in}{0.625278in}%
\pgfsys@useobject{currentmarker}{}%
\end{pgfscope}%
\begin{pgfscope}%
\pgfsys@transformshift{0.974293in}{0.625278in}%
\pgfsys@useobject{currentmarker}{}%
\end{pgfscope}%
\begin{pgfscope}%
\pgfsys@transformshift{0.974300in}{0.625278in}%
\pgfsys@useobject{currentmarker}{}%
\end{pgfscope}%
\begin{pgfscope}%
\pgfsys@transformshift{0.974476in}{0.625278in}%
\pgfsys@useobject{currentmarker}{}%
\end{pgfscope}%
\begin{pgfscope}%
\pgfsys@transformshift{0.974631in}{0.625278in}%
\pgfsys@useobject{currentmarker}{}%
\end{pgfscope}%
\begin{pgfscope}%
\pgfsys@transformshift{0.974638in}{0.625278in}%
\pgfsys@useobject{currentmarker}{}%
\end{pgfscope}%
\begin{pgfscope}%
\pgfsys@transformshift{0.974813in}{0.625278in}%
\pgfsys@useobject{currentmarker}{}%
\end{pgfscope}%
\begin{pgfscope}%
\pgfsys@transformshift{0.974813in}{0.625278in}%
\pgfsys@useobject{currentmarker}{}%
\end{pgfscope}%
\begin{pgfscope}%
\pgfsys@transformshift{0.974967in}{0.625278in}%
\pgfsys@useobject{currentmarker}{}%
\end{pgfscope}%
\begin{pgfscope}%
\pgfsys@transformshift{0.975151in}{0.625278in}%
\pgfsys@useobject{currentmarker}{}%
\end{pgfscope}%
\begin{pgfscope}%
\pgfsys@transformshift{0.975151in}{0.625278in}%
\pgfsys@useobject{currentmarker}{}%
\end{pgfscope}%
\begin{pgfscope}%
\pgfsys@transformshift{0.975151in}{0.625278in}%
\pgfsys@useobject{currentmarker}{}%
\end{pgfscope}%
\begin{pgfscope}%
\pgfsys@transformshift{0.975305in}{0.625278in}%
\pgfsys@useobject{currentmarker}{}%
\end{pgfscope}%
\begin{pgfscope}%
\pgfsys@transformshift{0.975488in}{0.625278in}%
\pgfsys@useobject{currentmarker}{}%
\end{pgfscope}%
\begin{pgfscope}%
\pgfsys@transformshift{0.975488in}{0.625278in}%
\pgfsys@useobject{currentmarker}{}%
\end{pgfscope}%
\begin{pgfscope}%
\pgfsys@transformshift{0.975642in}{0.625278in}%
\pgfsys@useobject{currentmarker}{}%
\end{pgfscope}%
\begin{pgfscope}%
\pgfsys@transformshift{0.975826in}{0.625278in}%
\pgfsys@useobject{currentmarker}{}%
\end{pgfscope}%
\begin{pgfscope}%
\pgfsys@transformshift{0.975826in}{0.625278in}%
\pgfsys@useobject{currentmarker}{}%
\end{pgfscope}%
\begin{pgfscope}%
\pgfsys@transformshift{0.975980in}{0.625278in}%
\pgfsys@useobject{currentmarker}{}%
\end{pgfscope}%
\begin{pgfscope}%
\pgfsys@transformshift{0.976162in}{0.625278in}%
\pgfsys@useobject{currentmarker}{}%
\end{pgfscope}%
\begin{pgfscope}%
\pgfsys@transformshift{0.976162in}{0.625278in}%
\pgfsys@useobject{currentmarker}{}%
\end{pgfscope}%
\begin{pgfscope}%
\pgfsys@transformshift{0.976162in}{0.625278in}%
\pgfsys@useobject{currentmarker}{}%
\end{pgfscope}%
\begin{pgfscope}%
\pgfsys@transformshift{0.976317in}{0.625278in}%
\pgfsys@useobject{currentmarker}{}%
\end{pgfscope}%
\begin{pgfscope}%
\pgfsys@transformshift{0.976500in}{0.625278in}%
\pgfsys@useobject{currentmarker}{}%
\end{pgfscope}%
\begin{pgfscope}%
\pgfsys@transformshift{0.976500in}{0.625278in}%
\pgfsys@useobject{currentmarker}{}%
\end{pgfscope}%
\begin{pgfscope}%
\pgfsys@transformshift{0.976654in}{0.625278in}%
\pgfsys@useobject{currentmarker}{}%
\end{pgfscope}%
\begin{pgfscope}%
\pgfsys@transformshift{0.981896in}{0.625278in}%
\pgfsys@useobject{currentmarker}{}%
\end{pgfscope}%
\begin{pgfscope}%
\pgfsys@transformshift{0.981896in}{0.625278in}%
\pgfsys@useobject{currentmarker}{}%
\end{pgfscope}%
\begin{pgfscope}%
\pgfsys@transformshift{0.981896in}{0.625278in}%
\pgfsys@useobject{currentmarker}{}%
\end{pgfscope}%
\begin{pgfscope}%
\pgfsys@transformshift{0.982050in}{0.625278in}%
\pgfsys@useobject{currentmarker}{}%
\end{pgfscope}%
\begin{pgfscope}%
\pgfsys@transformshift{0.982064in}{0.625278in}%
\pgfsys@useobject{currentmarker}{}%
\end{pgfscope}%
\begin{pgfscope}%
\pgfsys@transformshift{0.982233in}{0.625278in}%
\pgfsys@useobject{currentmarker}{}%
\end{pgfscope}%
\begin{pgfscope}%
\pgfsys@transformshift{0.982233in}{0.625278in}%
\pgfsys@useobject{currentmarker}{}%
\end{pgfscope}%
\begin{pgfscope}%
\pgfsys@transformshift{0.982387in}{0.625278in}%
\pgfsys@useobject{currentmarker}{}%
\end{pgfscope}%
\begin{pgfscope}%
\pgfsys@transformshift{0.982571in}{0.625278in}%
\pgfsys@useobject{currentmarker}{}%
\end{pgfscope}%
\begin{pgfscope}%
\pgfsys@transformshift{0.982571in}{0.625278in}%
\pgfsys@useobject{currentmarker}{}%
\end{pgfscope}%
\begin{pgfscope}%
\pgfsys@transformshift{0.982739in}{0.625278in}%
\pgfsys@useobject{currentmarker}{}%
\end{pgfscope}%
\begin{pgfscope}%
\pgfsys@transformshift{0.982907in}{0.625278in}%
\pgfsys@useobject{currentmarker}{}%
\end{pgfscope}%
\begin{pgfscope}%
\pgfsys@transformshift{0.982907in}{0.625278in}%
\pgfsys@useobject{currentmarker}{}%
\end{pgfscope}%
\begin{pgfscope}%
\pgfsys@transformshift{0.983076in}{0.625278in}%
\pgfsys@useobject{currentmarker}{}%
\end{pgfscope}%
\begin{pgfscope}%
\pgfsys@transformshift{0.983076in}{0.625278in}%
\pgfsys@useobject{currentmarker}{}%
\end{pgfscope}%
\begin{pgfscope}%
\pgfsys@transformshift{0.983245in}{0.625278in}%
\pgfsys@useobject{currentmarker}{}%
\end{pgfscope}%
\begin{pgfscope}%
\pgfsys@transformshift{0.983245in}{0.625278in}%
\pgfsys@useobject{currentmarker}{}%
\end{pgfscope}%
\begin{pgfscope}%
\pgfsys@transformshift{0.983414in}{0.625278in}%
\pgfsys@useobject{currentmarker}{}%
\end{pgfscope}%
\begin{pgfscope}%
\pgfsys@transformshift{0.983414in}{0.625278in}%
\pgfsys@useobject{currentmarker}{}%
\end{pgfscope}%
\begin{pgfscope}%
\pgfsys@transformshift{0.983582in}{0.625278in}%
\pgfsys@useobject{currentmarker}{}%
\end{pgfscope}%
\begin{pgfscope}%
\pgfsys@transformshift{0.983582in}{0.625278in}%
\pgfsys@useobject{currentmarker}{}%
\end{pgfscope}%
\begin{pgfscope}%
\pgfsys@transformshift{0.983750in}{0.625278in}%
\pgfsys@useobject{currentmarker}{}%
\end{pgfscope}%
\begin{pgfscope}%
\pgfsys@transformshift{0.983750in}{0.625278in}%
\pgfsys@useobject{currentmarker}{}%
\end{pgfscope}%
\begin{pgfscope}%
\pgfsys@transformshift{0.984088in}{0.625278in}%
\pgfsys@useobject{currentmarker}{}%
\end{pgfscope}%
\begin{pgfscope}%
\pgfsys@transformshift{0.984088in}{0.625278in}%
\pgfsys@useobject{currentmarker}{}%
\end{pgfscope}%
\begin{pgfscope}%
\pgfsys@transformshift{0.984088in}{0.625278in}%
\pgfsys@useobject{currentmarker}{}%
\end{pgfscope}%
\begin{pgfscope}%
\pgfsys@transformshift{0.984425in}{0.625278in}%
\pgfsys@useobject{currentmarker}{}%
\end{pgfscope}%
\begin{pgfscope}%
\pgfsys@transformshift{0.984425in}{0.625278in}%
\pgfsys@useobject{currentmarker}{}%
\end{pgfscope}%
\begin{pgfscope}%
\pgfsys@transformshift{0.984425in}{0.625278in}%
\pgfsys@useobject{currentmarker}{}%
\end{pgfscope}%
\begin{pgfscope}%
\pgfsys@transformshift{0.984763in}{0.625278in}%
\pgfsys@useobject{currentmarker}{}%
\end{pgfscope}%
\begin{pgfscope}%
\pgfsys@transformshift{0.984763in}{0.625278in}%
\pgfsys@useobject{currentmarker}{}%
\end{pgfscope}%
\begin{pgfscope}%
\pgfsys@transformshift{0.984763in}{0.625278in}%
\pgfsys@useobject{currentmarker}{}%
\end{pgfscope}%
\begin{pgfscope}%
\pgfsys@transformshift{0.984763in}{0.625278in}%
\pgfsys@useobject{currentmarker}{}%
\end{pgfscope}%
\begin{pgfscope}%
\pgfsys@transformshift{0.995049in}{0.625278in}%
\pgfsys@useobject{currentmarker}{}%
\end{pgfscope}%
\begin{pgfscope}%
\pgfsys@transformshift{1.132291in}{0.625278in}%
\pgfsys@useobject{currentmarker}{}%
\end{pgfscope}%
\begin{pgfscope}%
\pgfsys@transformshift{1.132966in}{0.625278in}%
\pgfsys@useobject{currentmarker}{}%
\end{pgfscope}%
\begin{pgfscope}%
\pgfsys@transformshift{1.132966in}{0.625278in}%
\pgfsys@useobject{currentmarker}{}%
\end{pgfscope}%
\begin{pgfscope}%
\pgfsys@transformshift{1.133302in}{0.625278in}%
\pgfsys@useobject{currentmarker}{}%
\end{pgfscope}%
\begin{pgfscope}%
\pgfsys@transformshift{1.133302in}{0.625278in}%
\pgfsys@useobject{currentmarker}{}%
\end{pgfscope}%
\begin{pgfscope}%
\pgfsys@transformshift{1.133302in}{0.625278in}%
\pgfsys@useobject{currentmarker}{}%
\end{pgfscope}%
\begin{pgfscope}%
\pgfsys@transformshift{1.133471in}{0.625278in}%
\pgfsys@useobject{currentmarker}{}%
\end{pgfscope}%
\begin{pgfscope}%
\pgfsys@transformshift{1.133640in}{0.625278in}%
\pgfsys@useobject{currentmarker}{}%
\end{pgfscope}%
\begin{pgfscope}%
\pgfsys@transformshift{1.133640in}{0.625278in}%
\pgfsys@useobject{currentmarker}{}%
\end{pgfscope}%
\begin{pgfscope}%
\pgfsys@transformshift{1.133809in}{0.625278in}%
\pgfsys@useobject{currentmarker}{}%
\end{pgfscope}%
\begin{pgfscope}%
\pgfsys@transformshift{1.133977in}{0.625278in}%
\pgfsys@useobject{currentmarker}{}%
\end{pgfscope}%
\begin{pgfscope}%
\pgfsys@transformshift{1.134315in}{0.625278in}%
\pgfsys@useobject{currentmarker}{}%
\end{pgfscope}%
\begin{pgfscope}%
\pgfsys@transformshift{1.134315in}{0.625278in}%
\pgfsys@useobject{currentmarker}{}%
\end{pgfscope}%
\begin{pgfscope}%
\pgfsys@transformshift{1.134652in}{0.625278in}%
\pgfsys@useobject{currentmarker}{}%
\end{pgfscope}%
\begin{pgfscope}%
\pgfsys@transformshift{1.134652in}{0.625278in}%
\pgfsys@useobject{currentmarker}{}%
\end{pgfscope}%
\begin{pgfscope}%
\pgfsys@transformshift{1.134820in}{0.625278in}%
\pgfsys@useobject{currentmarker}{}%
\end{pgfscope}%
\begin{pgfscope}%
\pgfsys@transformshift{1.134989in}{0.625278in}%
\pgfsys@useobject{currentmarker}{}%
\end{pgfscope}%
\begin{pgfscope}%
\pgfsys@transformshift{1.135158in}{0.625278in}%
\pgfsys@useobject{currentmarker}{}%
\end{pgfscope}%
\begin{pgfscope}%
\pgfsys@transformshift{1.135158in}{0.625278in}%
\pgfsys@useobject{currentmarker}{}%
\end{pgfscope}%
\begin{pgfscope}%
\pgfsys@transformshift{1.135326in}{0.625278in}%
\pgfsys@useobject{currentmarker}{}%
\end{pgfscope}%
\begin{pgfscope}%
\pgfsys@transformshift{1.135326in}{0.625278in}%
\pgfsys@useobject{currentmarker}{}%
\end{pgfscope}%
\begin{pgfscope}%
\pgfsys@transformshift{1.135495in}{0.625278in}%
\pgfsys@useobject{currentmarker}{}%
\end{pgfscope}%
\begin{pgfscope}%
\pgfsys@transformshift{1.135663in}{0.625278in}%
\pgfsys@useobject{currentmarker}{}%
\end{pgfscope}%
\begin{pgfscope}%
\pgfsys@transformshift{1.135663in}{0.625278in}%
\pgfsys@useobject{currentmarker}{}%
\end{pgfscope}%
\begin{pgfscope}%
\pgfsys@transformshift{1.135721in}{0.625278in}%
\pgfsys@useobject{currentmarker}{}%
\end{pgfscope}%
\begin{pgfscope}%
\pgfsys@transformshift{1.136001in}{0.625278in}%
\pgfsys@useobject{currentmarker}{}%
\end{pgfscope}%
\begin{pgfscope}%
\pgfsys@transformshift{1.136001in}{0.625278in}%
\pgfsys@useobject{currentmarker}{}%
\end{pgfscope}%
\begin{pgfscope}%
\pgfsys@transformshift{1.136058in}{0.625278in}%
\pgfsys@useobject{currentmarker}{}%
\end{pgfscope}%
\begin{pgfscope}%
\pgfsys@transformshift{1.136338in}{0.625278in}%
\pgfsys@useobject{currentmarker}{}%
\end{pgfscope}%
\begin{pgfscope}%
\pgfsys@transformshift{1.136338in}{0.625278in}%
\pgfsys@useobject{currentmarker}{}%
\end{pgfscope}%
\begin{pgfscope}%
\pgfsys@transformshift{1.136675in}{0.625278in}%
\pgfsys@useobject{currentmarker}{}%
\end{pgfscope}%
\begin{pgfscope}%
\pgfsys@transformshift{1.136675in}{0.625278in}%
\pgfsys@useobject{currentmarker}{}%
\end{pgfscope}%
\begin{pgfscope}%
\pgfsys@transformshift{1.137012in}{0.625278in}%
\pgfsys@useobject{currentmarker}{}%
\end{pgfscope}%
\begin{pgfscope}%
\pgfsys@transformshift{1.137012in}{0.625278in}%
\pgfsys@useobject{currentmarker}{}%
\end{pgfscope}%
\begin{pgfscope}%
\pgfsys@transformshift{1.137181in}{0.625278in}%
\pgfsys@useobject{currentmarker}{}%
\end{pgfscope}%
\begin{pgfscope}%
\pgfsys@transformshift{1.137350in}{0.625278in}%
\pgfsys@useobject{currentmarker}{}%
\end{pgfscope}%
\begin{pgfscope}%
\pgfsys@transformshift{1.137518in}{0.625278in}%
\pgfsys@useobject{currentmarker}{}%
\end{pgfscope}%
\begin{pgfscope}%
\pgfsys@transformshift{1.137518in}{0.625278in}%
\pgfsys@useobject{currentmarker}{}%
\end{pgfscope}%
\begin{pgfscope}%
\pgfsys@transformshift{1.137687in}{0.625278in}%
\pgfsys@useobject{currentmarker}{}%
\end{pgfscope}%
\begin{pgfscope}%
\pgfsys@transformshift{1.137855in}{0.625278in}%
\pgfsys@useobject{currentmarker}{}%
\end{pgfscope}%
\begin{pgfscope}%
\pgfsys@transformshift{1.137855in}{0.625278in}%
\pgfsys@useobject{currentmarker}{}%
\end{pgfscope}%
\begin{pgfscope}%
\pgfsys@transformshift{1.138024in}{0.625278in}%
\pgfsys@useobject{currentmarker}{}%
\end{pgfscope}%
\begin{pgfscope}%
\pgfsys@transformshift{1.138193in}{0.625278in}%
\pgfsys@useobject{currentmarker}{}%
\end{pgfscope}%
\begin{pgfscope}%
\pgfsys@transformshift{1.138193in}{0.625278in}%
\pgfsys@useobject{currentmarker}{}%
\end{pgfscope}%
\begin{pgfscope}%
\pgfsys@transformshift{1.138530in}{0.625278in}%
\pgfsys@useobject{currentmarker}{}%
\end{pgfscope}%
\begin{pgfscope}%
\pgfsys@transformshift{1.138530in}{0.625278in}%
\pgfsys@useobject{currentmarker}{}%
\end{pgfscope}%
\begin{pgfscope}%
\pgfsys@transformshift{1.138868in}{0.625278in}%
\pgfsys@useobject{currentmarker}{}%
\end{pgfscope}%
\begin{pgfscope}%
\pgfsys@transformshift{1.139036in}{0.625278in}%
\pgfsys@useobject{currentmarker}{}%
\end{pgfscope}%
\begin{pgfscope}%
\pgfsys@transformshift{1.139204in}{0.625278in}%
\pgfsys@useobject{currentmarker}{}%
\end{pgfscope}%
\begin{pgfscope}%
\pgfsys@transformshift{1.139373in}{0.625278in}%
\pgfsys@useobject{currentmarker}{}%
\end{pgfscope}%
\begin{pgfscope}%
\pgfsys@transformshift{1.139373in}{0.625278in}%
\pgfsys@useobject{currentmarker}{}%
\end{pgfscope}%
\begin{pgfscope}%
\pgfsys@transformshift{1.139542in}{0.625278in}%
\pgfsys@useobject{currentmarker}{}%
\end{pgfscope}%
\begin{pgfscope}%
\pgfsys@transformshift{1.139711in}{0.625278in}%
\pgfsys@useobject{currentmarker}{}%
\end{pgfscope}%
\begin{pgfscope}%
\pgfsys@transformshift{1.139711in}{0.625278in}%
\pgfsys@useobject{currentmarker}{}%
\end{pgfscope}%
\begin{pgfscope}%
\pgfsys@transformshift{1.139879in}{0.625278in}%
\pgfsys@useobject{currentmarker}{}%
\end{pgfscope}%
\begin{pgfscope}%
\pgfsys@transformshift{1.140047in}{0.625278in}%
\pgfsys@useobject{currentmarker}{}%
\end{pgfscope}%
\begin{pgfscope}%
\pgfsys@transformshift{1.140217in}{0.625278in}%
\pgfsys@useobject{currentmarker}{}%
\end{pgfscope}%
\begin{pgfscope}%
\pgfsys@transformshift{1.140217in}{0.625278in}%
\pgfsys@useobject{currentmarker}{}%
\end{pgfscope}%
\begin{pgfscope}%
\pgfsys@transformshift{1.140554in}{0.625278in}%
\pgfsys@useobject{currentmarker}{}%
\end{pgfscope}%
\begin{pgfscope}%
\pgfsys@transformshift{1.140554in}{0.625278in}%
\pgfsys@useobject{currentmarker}{}%
\end{pgfscope}%
\begin{pgfscope}%
\pgfsys@transformshift{1.140891in}{0.625278in}%
\pgfsys@useobject{currentmarker}{}%
\end{pgfscope}%
\begin{pgfscope}%
\pgfsys@transformshift{1.142409in}{0.625278in}%
\pgfsys@useobject{currentmarker}{}%
\end{pgfscope}%
\begin{pgfscope}%
\pgfsys@transformshift{1.144095in}{0.625278in}%
\pgfsys@useobject{currentmarker}{}%
\end{pgfscope}%
\begin{pgfscope}%
\pgfsys@transformshift{1.145781in}{0.625278in}%
\pgfsys@useobject{currentmarker}{}%
\end{pgfscope}%
\begin{pgfscope}%
\pgfsys@transformshift{1.146456in}{0.625278in}%
\pgfsys@useobject{currentmarker}{}%
\end{pgfscope}%
\begin{pgfscope}%
\pgfsys@transformshift{1.146624in}{0.625278in}%
\pgfsys@useobject{currentmarker}{}%
\end{pgfscope}%
\begin{pgfscope}%
\pgfsys@transformshift{1.146793in}{0.625278in}%
\pgfsys@useobject{currentmarker}{}%
\end{pgfscope}%
\begin{pgfscope}%
\pgfsys@transformshift{1.146793in}{0.625278in}%
\pgfsys@useobject{currentmarker}{}%
\end{pgfscope}%
\begin{pgfscope}%
\pgfsys@transformshift{1.167211in}{0.625278in}%
\pgfsys@useobject{currentmarker}{}%
\end{pgfscope}%
\begin{pgfscope}%
\pgfsys@transformshift{1.171419in}{0.625278in}%
\pgfsys@useobject{currentmarker}{}%
\end{pgfscope}%
\begin{pgfscope}%
\pgfsys@transformshift{1.171595in}{0.625278in}%
\pgfsys@useobject{currentmarker}{}%
\end{pgfscope}%
\begin{pgfscope}%
\pgfsys@transformshift{1.171757in}{0.625278in}%
\pgfsys@useobject{currentmarker}{}%
\end{pgfscope}%
\begin{pgfscope}%
\pgfsys@transformshift{1.171926in}{0.625278in}%
\pgfsys@useobject{currentmarker}{}%
\end{pgfscope}%
\begin{pgfscope}%
\pgfsys@transformshift{1.172094in}{0.625278in}%
\pgfsys@useobject{currentmarker}{}%
\end{pgfscope}%
\begin{pgfscope}%
\pgfsys@transformshift{1.172263in}{0.625278in}%
\pgfsys@useobject{currentmarker}{}%
\end{pgfscope}%
\begin{pgfscope}%
\pgfsys@transformshift{1.172432in}{0.625278in}%
\pgfsys@useobject{currentmarker}{}%
\end{pgfscope}%
\begin{pgfscope}%
\pgfsys@transformshift{1.172600in}{0.625278in}%
\pgfsys@useobject{currentmarker}{}%
\end{pgfscope}%
\begin{pgfscope}%
\pgfsys@transformshift{1.172607in}{0.625278in}%
\pgfsys@useobject{currentmarker}{}%
\end{pgfscope}%
\begin{pgfscope}%
\pgfsys@transformshift{1.172762in}{0.625278in}%
\pgfsys@useobject{currentmarker}{}%
\end{pgfscope}%
\begin{pgfscope}%
\pgfsys@transformshift{1.172769in}{0.625278in}%
\pgfsys@useobject{currentmarker}{}%
\end{pgfscope}%
\begin{pgfscope}%
\pgfsys@transformshift{1.173099in}{0.625278in}%
\pgfsys@useobject{currentmarker}{}%
\end{pgfscope}%
\begin{pgfscope}%
\pgfsys@transformshift{1.173106in}{0.625278in}%
\pgfsys@useobject{currentmarker}{}%
\end{pgfscope}%
\begin{pgfscope}%
\pgfsys@transformshift{1.173113in}{0.625278in}%
\pgfsys@useobject{currentmarker}{}%
\end{pgfscope}%
\begin{pgfscope}%
\pgfsys@transformshift{1.173443in}{0.625278in}%
\pgfsys@useobject{currentmarker}{}%
\end{pgfscope}%
\begin{pgfscope}%
\pgfsys@transformshift{1.173443in}{0.625278in}%
\pgfsys@useobject{currentmarker}{}%
\end{pgfscope}%
\begin{pgfscope}%
\pgfsys@transformshift{1.173781in}{0.625278in}%
\pgfsys@useobject{currentmarker}{}%
\end{pgfscope}%
\begin{pgfscope}%
\pgfsys@transformshift{1.173942in}{0.625278in}%
\pgfsys@useobject{currentmarker}{}%
\end{pgfscope}%
\begin{pgfscope}%
\pgfsys@transformshift{1.174000in}{0.625278in}%
\pgfsys@useobject{currentmarker}{}%
\end{pgfscope}%
\begin{pgfscope}%
\pgfsys@transformshift{1.174118in}{0.625278in}%
\pgfsys@useobject{currentmarker}{}%
\end{pgfscope}%
\begin{pgfscope}%
\pgfsys@transformshift{1.174280in}{0.625278in}%
\pgfsys@useobject{currentmarker}{}%
\end{pgfscope}%
\begin{pgfscope}%
\pgfsys@transformshift{1.174455in}{0.625278in}%
\pgfsys@useobject{currentmarker}{}%
\end{pgfscope}%
\begin{pgfscope}%
\pgfsys@transformshift{1.174624in}{0.625278in}%
\pgfsys@useobject{currentmarker}{}%
\end{pgfscope}%
\begin{pgfscope}%
\pgfsys@transformshift{1.174792in}{0.625278in}%
\pgfsys@useobject{currentmarker}{}%
\end{pgfscope}%
\begin{pgfscope}%
\pgfsys@transformshift{1.174954in}{0.625278in}%
\pgfsys@useobject{currentmarker}{}%
\end{pgfscope}%
\begin{pgfscope}%
\pgfsys@transformshift{1.174961in}{0.625278in}%
\pgfsys@useobject{currentmarker}{}%
\end{pgfscope}%
\begin{pgfscope}%
\pgfsys@transformshift{1.175298in}{0.625278in}%
\pgfsys@useobject{currentmarker}{}%
\end{pgfscope}%
\begin{pgfscope}%
\pgfsys@transformshift{1.188950in}{0.625278in}%
\pgfsys@useobject{currentmarker}{}%
\end{pgfscope}%
\begin{pgfscope}%
\pgfsys@transformshift{1.224025in}{0.625278in}%
\pgfsys@useobject{currentmarker}{}%
\end{pgfscope}%
\begin{pgfscope}%
\pgfsys@transformshift{1.246452in}{0.625278in}%
\pgfsys@useobject{currentmarker}{}%
\end{pgfscope}%
\begin{pgfscope}%
\pgfsys@transformshift{1.298390in}{0.625278in}%
\pgfsys@useobject{currentmarker}{}%
\end{pgfscope}%
\begin{pgfscope}%
\pgfsys@transformshift{1.349316in}{0.625278in}%
\pgfsys@useobject{currentmarker}{}%
\end{pgfscope}%
\begin{pgfscope}%
\pgfsys@transformshift{1.362300in}{0.625278in}%
\pgfsys@useobject{currentmarker}{}%
\end{pgfscope}%
\begin{pgfscope}%
\pgfsys@transformshift{1.372586in}{0.625278in}%
\pgfsys@useobject{currentmarker}{}%
\end{pgfscope}%
\begin{pgfscope}%
\pgfsys@transformshift{1.375284in}{0.625278in}%
\pgfsys@useobject{currentmarker}{}%
\end{pgfscope}%
\begin{pgfscope}%
\pgfsys@transformshift{1.375453in}{0.625278in}%
\pgfsys@useobject{currentmarker}{}%
\end{pgfscope}%
\begin{pgfscope}%
\pgfsys@transformshift{1.375621in}{0.625278in}%
\pgfsys@useobject{currentmarker}{}%
\end{pgfscope}%
\begin{pgfscope}%
\pgfsys@transformshift{1.375790in}{0.625278in}%
\pgfsys@useobject{currentmarker}{}%
\end{pgfscope}%
\begin{pgfscope}%
\pgfsys@transformshift{1.375959in}{0.625278in}%
\pgfsys@useobject{currentmarker}{}%
\end{pgfscope}%
\begin{pgfscope}%
\pgfsys@transformshift{1.376127in}{0.625278in}%
\pgfsys@useobject{currentmarker}{}%
\end{pgfscope}%
\begin{pgfscope}%
\pgfsys@transformshift{1.376296in}{0.625278in}%
\pgfsys@useobject{currentmarker}{}%
\end{pgfscope}%
\begin{pgfscope}%
\pgfsys@transformshift{1.376464in}{0.625278in}%
\pgfsys@useobject{currentmarker}{}%
\end{pgfscope}%
\begin{pgfscope}%
\pgfsys@transformshift{1.376634in}{0.625278in}%
\pgfsys@useobject{currentmarker}{}%
\end{pgfscope}%
\begin{pgfscope}%
\pgfsys@transformshift{1.376802in}{0.625278in}%
\pgfsys@useobject{currentmarker}{}%
\end{pgfscope}%
\begin{pgfscope}%
\pgfsys@transformshift{1.376970in}{0.625278in}%
\pgfsys@useobject{currentmarker}{}%
\end{pgfscope}%
\begin{pgfscope}%
\pgfsys@transformshift{1.377308in}{0.625278in}%
\pgfsys@useobject{currentmarker}{}%
\end{pgfscope}%
\begin{pgfscope}%
\pgfsys@transformshift{1.377477in}{0.625278in}%
\pgfsys@useobject{currentmarker}{}%
\end{pgfscope}%
\begin{pgfscope}%
\pgfsys@transformshift{1.377645in}{0.625278in}%
\pgfsys@useobject{currentmarker}{}%
\end{pgfscope}%
\begin{pgfscope}%
\pgfsys@transformshift{1.377813in}{0.625278in}%
\pgfsys@useobject{currentmarker}{}%
\end{pgfscope}%
\begin{pgfscope}%
\pgfsys@transformshift{1.377813in}{0.625278in}%
\pgfsys@useobject{currentmarker}{}%
\end{pgfscope}%
\begin{pgfscope}%
\pgfsys@transformshift{1.377982in}{0.625278in}%
\pgfsys@useobject{currentmarker}{}%
\end{pgfscope}%
\begin{pgfscope}%
\pgfsys@transformshift{1.378151in}{0.625278in}%
\pgfsys@useobject{currentmarker}{}%
\end{pgfscope}%
\begin{pgfscope}%
\pgfsys@transformshift{1.378320in}{0.625278in}%
\pgfsys@useobject{currentmarker}{}%
\end{pgfscope}%
\begin{pgfscope}%
\pgfsys@transformshift{1.378488in}{0.625278in}%
\pgfsys@useobject{currentmarker}{}%
\end{pgfscope}%
\begin{pgfscope}%
\pgfsys@transformshift{1.378656in}{0.625278in}%
\pgfsys@useobject{currentmarker}{}%
\end{pgfscope}%
\begin{pgfscope}%
\pgfsys@transformshift{1.423175in}{0.625278in}%
\pgfsys@useobject{currentmarker}{}%
\end{pgfscope}%
\begin{pgfscope}%
\pgfsys@transformshift{1.448806in}{0.625278in}%
\pgfsys@useobject{currentmarker}{}%
\end{pgfscope}%
\begin{pgfscope}%
\pgfsys@transformshift{1.467186in}{0.625278in}%
\pgfsys@useobject{currentmarker}{}%
\end{pgfscope}%
\begin{pgfscope}%
\pgfsys@transformshift{1.474944in}{0.625278in}%
\pgfsys@useobject{currentmarker}{}%
\end{pgfscope}%
\begin{pgfscope}%
\pgfsys@transformshift{1.475112in}{0.625278in}%
\pgfsys@useobject{currentmarker}{}%
\end{pgfscope}%
\begin{pgfscope}%
\pgfsys@transformshift{1.475281in}{0.625278in}%
\pgfsys@useobject{currentmarker}{}%
\end{pgfscope}%
\begin{pgfscope}%
\pgfsys@transformshift{1.475449in}{0.625278in}%
\pgfsys@useobject{currentmarker}{}%
\end{pgfscope}%
\begin{pgfscope}%
\pgfsys@transformshift{1.475618in}{0.625278in}%
\pgfsys@useobject{currentmarker}{}%
\end{pgfscope}%
\begin{pgfscope}%
\pgfsys@transformshift{1.475787in}{0.625278in}%
\pgfsys@useobject{currentmarker}{}%
\end{pgfscope}%
\begin{pgfscope}%
\pgfsys@transformshift{1.475955in}{0.625278in}%
\pgfsys@useobject{currentmarker}{}%
\end{pgfscope}%
\begin{pgfscope}%
\pgfsys@transformshift{1.476124in}{0.625278in}%
\pgfsys@useobject{currentmarker}{}%
\end{pgfscope}%
\begin{pgfscope}%
\pgfsys@transformshift{1.476293in}{0.625278in}%
\pgfsys@useobject{currentmarker}{}%
\end{pgfscope}%
\begin{pgfscope}%
\pgfsys@transformshift{1.476461in}{0.625278in}%
\pgfsys@useobject{currentmarker}{}%
\end{pgfscope}%
\begin{pgfscope}%
\pgfsys@transformshift{1.476630in}{0.625278in}%
\pgfsys@useobject{currentmarker}{}%
\end{pgfscope}%
\begin{pgfscope}%
\pgfsys@transformshift{1.476798in}{0.625278in}%
\pgfsys@useobject{currentmarker}{}%
\end{pgfscope}%
\begin{pgfscope}%
\pgfsys@transformshift{1.476967in}{0.625278in}%
\pgfsys@useobject{currentmarker}{}%
\end{pgfscope}%
\begin{pgfscope}%
\pgfsys@transformshift{1.477136in}{0.625278in}%
\pgfsys@useobject{currentmarker}{}%
\end{pgfscope}%
\begin{pgfscope}%
\pgfsys@transformshift{1.477304in}{0.625278in}%
\pgfsys@useobject{currentmarker}{}%
\end{pgfscope}%
\begin{pgfscope}%
\pgfsys@transformshift{1.477473in}{0.625278in}%
\pgfsys@useobject{currentmarker}{}%
\end{pgfscope}%
\begin{pgfscope}%
\pgfsys@transformshift{1.504959in}{0.625278in}%
\pgfsys@useobject{currentmarker}{}%
\end{pgfscope}%
\begin{pgfscope}%
\pgfsys@transformshift{1.529242in}{0.625278in}%
\pgfsys@useobject{currentmarker}{}%
\end{pgfscope}%
\begin{pgfscope}%
\pgfsys@transformshift{1.545598in}{0.625278in}%
\pgfsys@useobject{currentmarker}{}%
\end{pgfscope}%
\begin{pgfscope}%
\pgfsys@transformshift{1.554873in}{0.625278in}%
\pgfsys@useobject{currentmarker}{}%
\end{pgfscope}%
\begin{pgfscope}%
\pgfsys@transformshift{1.555042in}{0.625278in}%
\pgfsys@useobject{currentmarker}{}%
\end{pgfscope}%
\begin{pgfscope}%
\pgfsys@transformshift{1.555210in}{0.625278in}%
\pgfsys@useobject{currentmarker}{}%
\end{pgfscope}%
\begin{pgfscope}%
\pgfsys@transformshift{1.555548in}{0.625278in}%
\pgfsys@useobject{currentmarker}{}%
\end{pgfscope}%
\begin{pgfscope}%
\pgfsys@transformshift{1.555885in}{0.625278in}%
\pgfsys@useobject{currentmarker}{}%
\end{pgfscope}%
\begin{pgfscope}%
\pgfsys@transformshift{1.555885in}{0.625278in}%
\pgfsys@useobject{currentmarker}{}%
\end{pgfscope}%
\begin{pgfscope}%
\pgfsys@transformshift{1.556223in}{0.625278in}%
\pgfsys@useobject{currentmarker}{}%
\end{pgfscope}%
\begin{pgfscope}%
\pgfsys@transformshift{1.556559in}{0.625278in}%
\pgfsys@useobject{currentmarker}{}%
\end{pgfscope}%
\begin{pgfscope}%
\pgfsys@transformshift{1.556728in}{0.625278in}%
\pgfsys@useobject{currentmarker}{}%
\end{pgfscope}%
\begin{pgfscope}%
\pgfsys@transformshift{1.556897in}{0.625278in}%
\pgfsys@useobject{currentmarker}{}%
\end{pgfscope}%
\begin{pgfscope}%
\pgfsys@transformshift{1.556897in}{0.625278in}%
\pgfsys@useobject{currentmarker}{}%
\end{pgfscope}%
\begin{pgfscope}%
\pgfsys@transformshift{1.557066in}{0.625278in}%
\pgfsys@useobject{currentmarker}{}%
\end{pgfscope}%
\begin{pgfscope}%
\pgfsys@transformshift{1.557066in}{0.625278in}%
\pgfsys@useobject{currentmarker}{}%
\end{pgfscope}%
\begin{pgfscope}%
\pgfsys@transformshift{1.557234in}{0.625278in}%
\pgfsys@useobject{currentmarker}{}%
\end{pgfscope}%
\begin{pgfscope}%
\pgfsys@transformshift{1.557234in}{0.625278in}%
\pgfsys@useobject{currentmarker}{}%
\end{pgfscope}%
\begin{pgfscope}%
\pgfsys@transformshift{1.557402in}{0.625278in}%
\pgfsys@useobject{currentmarker}{}%
\end{pgfscope}%
\begin{pgfscope}%
\pgfsys@transformshift{1.557572in}{0.625278in}%
\pgfsys@useobject{currentmarker}{}%
\end{pgfscope}%
\begin{pgfscope}%
\pgfsys@transformshift{1.557572in}{0.625278in}%
\pgfsys@useobject{currentmarker}{}%
\end{pgfscope}%
\begin{pgfscope}%
\pgfsys@transformshift{1.581183in}{0.625278in}%
\pgfsys@useobject{currentmarker}{}%
\end{pgfscope}%
\begin{pgfscope}%
\pgfsys@transformshift{1.594337in}{0.625278in}%
\pgfsys@useobject{currentmarker}{}%
\end{pgfscope}%
\begin{pgfscope}%
\pgfsys@transformshift{1.753521in}{0.625278in}%
\pgfsys@useobject{currentmarker}{}%
\end{pgfscope}%
\begin{pgfscope}%
\pgfsys@transformshift{1.768186in}{0.625278in}%
\pgfsys@useobject{currentmarker}{}%
\end{pgfscope}%
\begin{pgfscope}%
\pgfsys@transformshift{1.772238in}{0.625278in}%
\pgfsys@useobject{currentmarker}{}%
\end{pgfscope}%
\begin{pgfscope}%
\pgfsys@transformshift{1.772402in}{0.625278in}%
\pgfsys@useobject{currentmarker}{}%
\end{pgfscope}%
\begin{pgfscope}%
\pgfsys@transformshift{1.772573in}{0.625278in}%
\pgfsys@useobject{currentmarker}{}%
\end{pgfscope}%
\begin{pgfscope}%
\pgfsys@transformshift{1.772745in}{0.625278in}%
\pgfsys@useobject{currentmarker}{}%
\end{pgfscope}%
\begin{pgfscope}%
\pgfsys@transformshift{1.772909in}{0.625278in}%
\pgfsys@useobject{currentmarker}{}%
\end{pgfscope}%
\begin{pgfscope}%
\pgfsys@transformshift{1.773080in}{0.625278in}%
\pgfsys@useobject{currentmarker}{}%
\end{pgfscope}%
\begin{pgfscope}%
\pgfsys@transformshift{1.773080in}{0.625278in}%
\pgfsys@useobject{currentmarker}{}%
\end{pgfscope}%
\begin{pgfscope}%
\pgfsys@transformshift{1.773244in}{0.625278in}%
\pgfsys@useobject{currentmarker}{}%
\end{pgfscope}%
\begin{pgfscope}%
\pgfsys@transformshift{1.773415in}{0.625278in}%
\pgfsys@useobject{currentmarker}{}%
\end{pgfscope}%
\begin{pgfscope}%
\pgfsys@transformshift{1.773415in}{0.625278in}%
\pgfsys@useobject{currentmarker}{}%
\end{pgfscope}%
\begin{pgfscope}%
\pgfsys@transformshift{1.773586in}{0.625278in}%
\pgfsys@useobject{currentmarker}{}%
\end{pgfscope}%
\begin{pgfscope}%
\pgfsys@transformshift{1.773750in}{0.625278in}%
\pgfsys@useobject{currentmarker}{}%
\end{pgfscope}%
\begin{pgfscope}%
\pgfsys@transformshift{1.773750in}{0.625278in}%
\pgfsys@useobject{currentmarker}{}%
\end{pgfscope}%
\begin{pgfscope}%
\pgfsys@transformshift{1.773921in}{0.625278in}%
\pgfsys@useobject{currentmarker}{}%
\end{pgfscope}%
\begin{pgfscope}%
\pgfsys@transformshift{1.774093in}{0.625278in}%
\pgfsys@useobject{currentmarker}{}%
\end{pgfscope}%
\begin{pgfscope}%
\pgfsys@transformshift{1.774257in}{0.625278in}%
\pgfsys@useobject{currentmarker}{}%
\end{pgfscope}%
\begin{pgfscope}%
\pgfsys@transformshift{1.774428in}{0.625278in}%
\pgfsys@useobject{currentmarker}{}%
\end{pgfscope}%
\begin{pgfscope}%
\pgfsys@transformshift{1.774599in}{0.625278in}%
\pgfsys@useobject{currentmarker}{}%
\end{pgfscope}%
\begin{pgfscope}%
\pgfsys@transformshift{1.774763in}{0.625278in}%
\pgfsys@useobject{currentmarker}{}%
\end{pgfscope}%
\begin{pgfscope}%
\pgfsys@transformshift{1.774934in}{0.625278in}%
\pgfsys@useobject{currentmarker}{}%
\end{pgfscope}%
\begin{pgfscope}%
\pgfsys@transformshift{1.775098in}{0.625278in}%
\pgfsys@useobject{currentmarker}{}%
\end{pgfscope}%
\begin{pgfscope}%
\pgfsys@transformshift{1.775441in}{0.625278in}%
\pgfsys@useobject{currentmarker}{}%
\end{pgfscope}%
\begin{pgfscope}%
\pgfsys@transformshift{1.775605in}{0.625278in}%
\pgfsys@useobject{currentmarker}{}%
\end{pgfscope}%
\begin{pgfscope}%
\pgfsys@transformshift{1.775776in}{0.625278in}%
\pgfsys@useobject{currentmarker}{}%
\end{pgfscope}%
\begin{pgfscope}%
\pgfsys@transformshift{1.775947in}{0.625278in}%
\pgfsys@useobject{currentmarker}{}%
\end{pgfscope}%
\begin{pgfscope}%
\pgfsys@transformshift{1.776283in}{0.625278in}%
\pgfsys@useobject{currentmarker}{}%
\end{pgfscope}%
\begin{pgfscope}%
\pgfsys@transformshift{1.806299in}{0.625278in}%
\pgfsys@useobject{currentmarker}{}%
\end{pgfscope}%
\begin{pgfscope}%
\pgfsys@transformshift{1.849469in}{0.625278in}%
\pgfsys@useobject{currentmarker}{}%
\end{pgfscope}%
\begin{pgfscope}%
\pgfsys@transformshift{1.864812in}{0.625278in}%
\pgfsys@useobject{currentmarker}{}%
\end{pgfscope}%
\begin{pgfscope}%
\pgfsys@transformshift{1.867345in}{0.625278in}%
\pgfsys@useobject{currentmarker}{}%
\end{pgfscope}%
\begin{pgfscope}%
\pgfsys@transformshift{1.867509in}{0.625278in}%
\pgfsys@useobject{currentmarker}{}%
\end{pgfscope}%
\begin{pgfscope}%
\pgfsys@transformshift{1.867680in}{0.625278in}%
\pgfsys@useobject{currentmarker}{}%
\end{pgfscope}%
\begin{pgfscope}%
\pgfsys@transformshift{1.868015in}{0.625278in}%
\pgfsys@useobject{currentmarker}{}%
\end{pgfscope}%
\begin{pgfscope}%
\pgfsys@transformshift{1.868015in}{0.625278in}%
\pgfsys@useobject{currentmarker}{}%
\end{pgfscope}%
\begin{pgfscope}%
\pgfsys@transformshift{1.868350in}{0.625278in}%
\pgfsys@useobject{currentmarker}{}%
\end{pgfscope}%
\begin{pgfscope}%
\pgfsys@transformshift{1.868350in}{0.625278in}%
\pgfsys@useobject{currentmarker}{}%
\end{pgfscope}%
\begin{pgfscope}%
\pgfsys@transformshift{1.868350in}{0.625278in}%
\pgfsys@useobject{currentmarker}{}%
\end{pgfscope}%
\begin{pgfscope}%
\pgfsys@transformshift{1.868693in}{0.625278in}%
\pgfsys@useobject{currentmarker}{}%
\end{pgfscope}%
\begin{pgfscope}%
\pgfsys@transformshift{1.868693in}{0.625278in}%
\pgfsys@useobject{currentmarker}{}%
\end{pgfscope}%
\begin{pgfscope}%
\pgfsys@transformshift{1.869028in}{0.625278in}%
\pgfsys@useobject{currentmarker}{}%
\end{pgfscope}%
\begin{pgfscope}%
\pgfsys@transformshift{1.869199in}{0.625278in}%
\pgfsys@useobject{currentmarker}{}%
\end{pgfscope}%
\begin{pgfscope}%
\pgfsys@transformshift{1.869363in}{0.625278in}%
\pgfsys@useobject{currentmarker}{}%
\end{pgfscope}%
\begin{pgfscope}%
\pgfsys@transformshift{1.869535in}{0.625278in}%
\pgfsys@useobject{currentmarker}{}%
\end{pgfscope}%
\begin{pgfscope}%
\pgfsys@transformshift{1.869535in}{0.625278in}%
\pgfsys@useobject{currentmarker}{}%
\end{pgfscope}%
\begin{pgfscope}%
\pgfsys@transformshift{1.869706in}{0.625278in}%
\pgfsys@useobject{currentmarker}{}%
\end{pgfscope}%
\begin{pgfscope}%
\pgfsys@transformshift{1.875940in}{0.625278in}%
\pgfsys@useobject{currentmarker}{}%
\end{pgfscope}%
\begin{pgfscope}%
\pgfsys@transformshift{1.888766in}{0.625278in}%
\pgfsys@useobject{currentmarker}{}%
\end{pgfscope}%
\begin{pgfscope}%
\pgfsys@transformshift{1.897026in}{0.625278in}%
\pgfsys@useobject{currentmarker}{}%
\end{pgfscope}%
\begin{pgfscope}%
\pgfsys@transformshift{1.897696in}{0.625278in}%
\pgfsys@useobject{currentmarker}{}%
\end{pgfscope}%
\begin{pgfscope}%
\pgfsys@transformshift{1.897704in}{0.625278in}%
\pgfsys@useobject{currentmarker}{}%
\end{pgfscope}%
\begin{pgfscope}%
\pgfsys@transformshift{1.898039in}{0.625278in}%
\pgfsys@useobject{currentmarker}{}%
\end{pgfscope}%
\begin{pgfscope}%
\pgfsys@transformshift{1.898047in}{0.625278in}%
\pgfsys@useobject{currentmarker}{}%
\end{pgfscope}%
\begin{pgfscope}%
\pgfsys@transformshift{1.898091in}{0.625278in}%
\pgfsys@useobject{currentmarker}{}%
\end{pgfscope}%
\begin{pgfscope}%
\pgfsys@transformshift{1.898374in}{0.625278in}%
\pgfsys@useobject{currentmarker}{}%
\end{pgfscope}%
\begin{pgfscope}%
\pgfsys@transformshift{1.898382in}{0.625278in}%
\pgfsys@useobject{currentmarker}{}%
\end{pgfscope}%
\begin{pgfscope}%
\pgfsys@transformshift{1.898546in}{0.625278in}%
\pgfsys@useobject{currentmarker}{}%
\end{pgfscope}%
\begin{pgfscope}%
\pgfsys@transformshift{1.898717in}{0.625278in}%
\pgfsys@useobject{currentmarker}{}%
\end{pgfscope}%
\begin{pgfscope}%
\pgfsys@transformshift{1.898724in}{0.625278in}%
\pgfsys@useobject{currentmarker}{}%
\end{pgfscope}%
\begin{pgfscope}%
\pgfsys@transformshift{1.899052in}{0.625278in}%
\pgfsys@useobject{currentmarker}{}%
\end{pgfscope}%
\begin{pgfscope}%
\pgfsys@transformshift{1.899052in}{0.625278in}%
\pgfsys@useobject{currentmarker}{}%
\end{pgfscope}%
\begin{pgfscope}%
\pgfsys@transformshift{1.899059in}{0.625278in}%
\pgfsys@useobject{currentmarker}{}%
\end{pgfscope}%
\begin{pgfscope}%
\pgfsys@transformshift{1.899387in}{0.625278in}%
\pgfsys@useobject{currentmarker}{}%
\end{pgfscope}%
\begin{pgfscope}%
\pgfsys@transformshift{1.899395in}{0.625278in}%
\pgfsys@useobject{currentmarker}{}%
\end{pgfscope}%
\begin{pgfscope}%
\pgfsys@transformshift{1.899730in}{0.625278in}%
\pgfsys@useobject{currentmarker}{}%
\end{pgfscope}%
\begin{pgfscope}%
\pgfsys@transformshift{1.899886in}{0.625278in}%
\pgfsys@useobject{currentmarker}{}%
\end{pgfscope}%
\begin{pgfscope}%
\pgfsys@transformshift{1.900072in}{0.625278in}%
\pgfsys@useobject{currentmarker}{}%
\end{pgfscope}%
\begin{pgfscope}%
\pgfsys@transformshift{1.905621in}{0.625278in}%
\pgfsys@useobject{currentmarker}{}%
\end{pgfscope}%
\begin{pgfscope}%
\pgfsys@transformshift{2.149119in}{0.625278in}%
\pgfsys@useobject{currentmarker}{}%
\end{pgfscope}%
\begin{pgfscope}%
\pgfsys@transformshift{2.159405in}{0.625278in}%
\pgfsys@useobject{currentmarker}{}%
\end{pgfscope}%
\begin{pgfscope}%
\pgfsys@transformshift{2.159577in}{0.625278in}%
\pgfsys@useobject{currentmarker}{}%
\end{pgfscope}%
\begin{pgfscope}%
\pgfsys@transformshift{2.159912in}{0.625278in}%
\pgfsys@useobject{currentmarker}{}%
\end{pgfscope}%
\begin{pgfscope}%
\pgfsys@transformshift{2.160083in}{0.625278in}%
\pgfsys@useobject{currentmarker}{}%
\end{pgfscope}%
\begin{pgfscope}%
\pgfsys@transformshift{2.160247in}{0.625278in}%
\pgfsys@useobject{currentmarker}{}%
\end{pgfscope}%
\begin{pgfscope}%
\pgfsys@transformshift{2.160418in}{0.625278in}%
\pgfsys@useobject{currentmarker}{}%
\end{pgfscope}%
\begin{pgfscope}%
\pgfsys@transformshift{2.160418in}{0.625278in}%
\pgfsys@useobject{currentmarker}{}%
\end{pgfscope}%
\begin{pgfscope}%
\pgfsys@transformshift{2.160754in}{0.625278in}%
\pgfsys@useobject{currentmarker}{}%
\end{pgfscope}%
\begin{pgfscope}%
\pgfsys@transformshift{2.160754in}{0.625278in}%
\pgfsys@useobject{currentmarker}{}%
\end{pgfscope}%
\begin{pgfscope}%
\pgfsys@transformshift{2.161089in}{0.625278in}%
\pgfsys@useobject{currentmarker}{}%
\end{pgfscope}%
\begin{pgfscope}%
\pgfsys@transformshift{2.161148in}{0.625278in}%
\pgfsys@useobject{currentmarker}{}%
\end{pgfscope}%
\begin{pgfscope}%
\pgfsys@transformshift{2.161431in}{0.625278in}%
\pgfsys@useobject{currentmarker}{}%
\end{pgfscope}%
\begin{pgfscope}%
\pgfsys@transformshift{2.161431in}{0.625278in}%
\pgfsys@useobject{currentmarker}{}%
\end{pgfscope}%
\begin{pgfscope}%
\pgfsys@transformshift{2.161595in}{0.625278in}%
\pgfsys@useobject{currentmarker}{}%
\end{pgfscope}%
\begin{pgfscope}%
\pgfsys@transformshift{2.161767in}{0.625278in}%
\pgfsys@useobject{currentmarker}{}%
\end{pgfscope}%
\begin{pgfscope}%
\pgfsys@transformshift{2.161767in}{0.625278in}%
\pgfsys@useobject{currentmarker}{}%
\end{pgfscope}%
\begin{pgfscope}%
\pgfsys@transformshift{2.161938in}{0.625278in}%
\pgfsys@useobject{currentmarker}{}%
\end{pgfscope}%
\begin{pgfscope}%
\pgfsys@transformshift{2.162102in}{0.625278in}%
\pgfsys@useobject{currentmarker}{}%
\end{pgfscope}%
\begin{pgfscope}%
\pgfsys@transformshift{2.162333in}{0.625278in}%
\pgfsys@useobject{currentmarker}{}%
\end{pgfscope}%
\begin{pgfscope}%
\pgfsys@transformshift{2.162444in}{0.625278in}%
\pgfsys@useobject{currentmarker}{}%
\end{pgfscope}%
\begin{pgfscope}%
\pgfsys@transformshift{2.162780in}{0.625278in}%
\pgfsys@useobject{currentmarker}{}%
\end{pgfscope}%
\begin{pgfscope}%
\pgfsys@transformshift{2.185884in}{0.625278in}%
\pgfsys@useobject{currentmarker}{}%
\end{pgfscope}%
\begin{pgfscope}%
\pgfsys@transformshift{2.193816in}{0.625278in}%
\pgfsys@useobject{currentmarker}{}%
\end{pgfscope}%
\begin{pgfscope}%
\pgfsys@transformshift{2.194144in}{0.625278in}%
\pgfsys@useobject{currentmarker}{}%
\end{pgfscope}%
\begin{pgfscope}%
\pgfsys@transformshift{2.194152in}{0.625278in}%
\pgfsys@useobject{currentmarker}{}%
\end{pgfscope}%
\begin{pgfscope}%
\pgfsys@transformshift{2.194152in}{0.625278in}%
\pgfsys@useobject{currentmarker}{}%
\end{pgfscope}%
\begin{pgfscope}%
\pgfsys@transformshift{2.194479in}{0.625278in}%
\pgfsys@useobject{currentmarker}{}%
\end{pgfscope}%
\begin{pgfscope}%
\pgfsys@transformshift{2.194487in}{0.625278in}%
\pgfsys@useobject{currentmarker}{}%
\end{pgfscope}%
\begin{pgfscope}%
\pgfsys@transformshift{2.194487in}{0.625278in}%
\pgfsys@useobject{currentmarker}{}%
\end{pgfscope}%
\begin{pgfscope}%
\pgfsys@transformshift{2.194658in}{0.625278in}%
\pgfsys@useobject{currentmarker}{}%
\end{pgfscope}%
\begin{pgfscope}%
\pgfsys@transformshift{2.194814in}{0.625278in}%
\pgfsys@useobject{currentmarker}{}%
\end{pgfscope}%
\begin{pgfscope}%
\pgfsys@transformshift{2.194993in}{0.625278in}%
\pgfsys@useobject{currentmarker}{}%
\end{pgfscope}%
\begin{pgfscope}%
\pgfsys@transformshift{2.195045in}{0.625278in}%
\pgfsys@useobject{currentmarker}{}%
\end{pgfscope}%
\begin{pgfscope}%
\pgfsys@transformshift{2.195157in}{0.625278in}%
\pgfsys@useobject{currentmarker}{}%
\end{pgfscope}%
\begin{pgfscope}%
\pgfsys@transformshift{2.195164in}{0.625278in}%
\pgfsys@useobject{currentmarker}{}%
\end{pgfscope}%
\begin{pgfscope}%
\pgfsys@transformshift{2.195328in}{0.625278in}%
\pgfsys@useobject{currentmarker}{}%
\end{pgfscope}%
\begin{pgfscope}%
\pgfsys@transformshift{2.195671in}{0.625278in}%
\pgfsys@useobject{currentmarker}{}%
\end{pgfscope}%
\begin{pgfscope}%
\pgfsys@transformshift{2.195671in}{0.625278in}%
\pgfsys@useobject{currentmarker}{}%
\end{pgfscope}%
\begin{pgfscope}%
\pgfsys@transformshift{2.195827in}{0.625278in}%
\pgfsys@useobject{currentmarker}{}%
\end{pgfscope}%
\begin{pgfscope}%
\pgfsys@transformshift{2.196006in}{0.625278in}%
\pgfsys@useobject{currentmarker}{}%
\end{pgfscope}%
\begin{pgfscope}%
\pgfsys@transformshift{2.196170in}{0.625278in}%
\pgfsys@useobject{currentmarker}{}%
\end{pgfscope}%
\begin{pgfscope}%
\pgfsys@transformshift{2.196177in}{0.625278in}%
\pgfsys@useobject{currentmarker}{}%
\end{pgfscope}%
\begin{pgfscope}%
\pgfsys@transformshift{2.196513in}{0.625278in}%
\pgfsys@useobject{currentmarker}{}%
\end{pgfscope}%
\begin{pgfscope}%
\pgfsys@transformshift{2.196840in}{0.625278in}%
\pgfsys@useobject{currentmarker}{}%
\end{pgfscope}%
\begin{pgfscope}%
\pgfsys@transformshift{2.197176in}{0.625278in}%
\pgfsys@useobject{currentmarker}{}%
\end{pgfscope}%
\begin{pgfscope}%
\pgfsys@transformshift{2.396163in}{0.625278in}%
\pgfsys@useobject{currentmarker}{}%
\end{pgfscope}%
\begin{pgfscope}%
\pgfsys@transformshift{2.408639in}{0.625278in}%
\pgfsys@useobject{currentmarker}{}%
\end{pgfscope}%
\begin{pgfscope}%
\pgfsys@transformshift{2.422127in}{0.625278in}%
\pgfsys@useobject{currentmarker}{}%
\end{pgfscope}%
\begin{pgfscope}%
\pgfsys@transformshift{2.422299in}{0.625278in}%
\pgfsys@useobject{currentmarker}{}%
\end{pgfscope}%
\begin{pgfscope}%
\pgfsys@transformshift{2.422470in}{0.625278in}%
\pgfsys@useobject{currentmarker}{}%
\end{pgfscope}%
\begin{pgfscope}%
\pgfsys@transformshift{2.422634in}{0.625278in}%
\pgfsys@useobject{currentmarker}{}%
\end{pgfscope}%
\begin{pgfscope}%
\pgfsys@transformshift{2.422634in}{0.625278in}%
\pgfsys@useobject{currentmarker}{}%
\end{pgfscope}%
\begin{pgfscope}%
\pgfsys@transformshift{2.422813in}{0.625278in}%
\pgfsys@useobject{currentmarker}{}%
\end{pgfscope}%
\begin{pgfscope}%
\pgfsys@transformshift{2.422969in}{0.625278in}%
\pgfsys@useobject{currentmarker}{}%
\end{pgfscope}%
\begin{pgfscope}%
\pgfsys@transformshift{2.423312in}{0.625278in}%
\pgfsys@useobject{currentmarker}{}%
\end{pgfscope}%
\begin{pgfscope}%
\pgfsys@transformshift{2.423647in}{0.625278in}%
\pgfsys@useobject{currentmarker}{}%
\end{pgfscope}%
\begin{pgfscope}%
\pgfsys@transformshift{2.423654in}{0.625278in}%
\pgfsys@useobject{currentmarker}{}%
\end{pgfscope}%
\begin{pgfscope}%
\pgfsys@transformshift{2.423818in}{0.625278in}%
\pgfsys@useobject{currentmarker}{}%
\end{pgfscope}%
\begin{pgfscope}%
\pgfsys@transformshift{2.424153in}{0.625278in}%
\pgfsys@useobject{currentmarker}{}%
\end{pgfscope}%
\begin{pgfscope}%
\pgfsys@transformshift{2.424489in}{0.625278in}%
\pgfsys@useobject{currentmarker}{}%
\end{pgfscope}%
\begin{pgfscope}%
\pgfsys@transformshift{2.424660in}{0.625278in}%
\pgfsys@useobject{currentmarker}{}%
\end{pgfscope}%
\begin{pgfscope}%
\pgfsys@transformshift{2.424824in}{0.625278in}%
\pgfsys@useobject{currentmarker}{}%
\end{pgfscope}%
\begin{pgfscope}%
\pgfsys@transformshift{2.424824in}{0.625278in}%
\pgfsys@useobject{currentmarker}{}%
\end{pgfscope}%
\begin{pgfscope}%
\pgfsys@transformshift{2.424995in}{0.625278in}%
\pgfsys@useobject{currentmarker}{}%
\end{pgfscope}%
\begin{pgfscope}%
\pgfsys@transformshift{2.425166in}{0.625278in}%
\pgfsys@useobject{currentmarker}{}%
\end{pgfscope}%
\begin{pgfscope}%
\pgfsys@transformshift{2.425330in}{0.625278in}%
\pgfsys@useobject{currentmarker}{}%
\end{pgfscope}%
\begin{pgfscope}%
\pgfsys@transformshift{2.425501in}{0.625278in}%
\pgfsys@useobject{currentmarker}{}%
\end{pgfscope}%
\begin{pgfscope}%
\pgfsys@transformshift{2.425673in}{0.625278in}%
\pgfsys@useobject{currentmarker}{}%
\end{pgfscope}%
\begin{pgfscope}%
\pgfsys@transformshift{2.425837in}{0.625278in}%
\pgfsys@useobject{currentmarker}{}%
\end{pgfscope}%
\begin{pgfscope}%
\pgfsys@transformshift{2.464970in}{0.625278in}%
\pgfsys@useobject{currentmarker}{}%
\end{pgfscope}%
\begin{pgfscope}%
\pgfsys@transformshift{2.494472in}{0.625278in}%
\pgfsys@useobject{currentmarker}{}%
\end{pgfscope}%
\begin{pgfscope}%
\pgfsys@transformshift{2.506948in}{0.625278in}%
\pgfsys@useobject{currentmarker}{}%
\end{pgfscope}%
\begin{pgfscope}%
\pgfsys@transformshift{2.524489in}{0.625278in}%
\pgfsys@useobject{currentmarker}{}%
\end{pgfscope}%
\begin{pgfscope}%
\pgfsys@transformshift{2.531401in}{0.625278in}%
\pgfsys@useobject{currentmarker}{}%
\end{pgfscope}%
\begin{pgfscope}%
\pgfsys@transformshift{2.531572in}{0.625278in}%
\pgfsys@useobject{currentmarker}{}%
\end{pgfscope}%
\begin{pgfscope}%
\pgfsys@transformshift{2.531736in}{0.625278in}%
\pgfsys@useobject{currentmarker}{}%
\end{pgfscope}%
\begin{pgfscope}%
\pgfsys@transformshift{2.531915in}{0.625278in}%
\pgfsys@useobject{currentmarker}{}%
\end{pgfscope}%
\begin{pgfscope}%
\pgfsys@transformshift{2.532071in}{0.625278in}%
\pgfsys@useobject{currentmarker}{}%
\end{pgfscope}%
\begin{pgfscope}%
\pgfsys@transformshift{2.532250in}{0.625278in}%
\pgfsys@useobject{currentmarker}{}%
\end{pgfscope}%
\begin{pgfscope}%
\pgfsys@transformshift{2.532414in}{0.625278in}%
\pgfsys@useobject{currentmarker}{}%
\end{pgfscope}%
\begin{pgfscope}%
\pgfsys@transformshift{2.532414in}{0.625278in}%
\pgfsys@useobject{currentmarker}{}%
\end{pgfscope}%
\begin{pgfscope}%
\pgfsys@transformshift{2.532585in}{0.625278in}%
\pgfsys@useobject{currentmarker}{}%
\end{pgfscope}%
\begin{pgfscope}%
\pgfsys@transformshift{2.532749in}{0.625278in}%
\pgfsys@useobject{currentmarker}{}%
\end{pgfscope}%
\begin{pgfscope}%
\pgfsys@transformshift{2.532920in}{0.625278in}%
\pgfsys@useobject{currentmarker}{}%
\end{pgfscope}%
\begin{pgfscope}%
\pgfsys@transformshift{2.532928in}{0.625278in}%
\pgfsys@useobject{currentmarker}{}%
\end{pgfscope}%
\begin{pgfscope}%
\pgfsys@transformshift{2.533084in}{0.625278in}%
\pgfsys@useobject{currentmarker}{}%
\end{pgfscope}%
\begin{pgfscope}%
\pgfsys@transformshift{2.533084in}{0.625278in}%
\pgfsys@useobject{currentmarker}{}%
\end{pgfscope}%
\begin{pgfscope}%
\pgfsys@transformshift{2.533255in}{0.625278in}%
\pgfsys@useobject{currentmarker}{}%
\end{pgfscope}%
\begin{pgfscope}%
\pgfsys@transformshift{2.533427in}{0.625278in}%
\pgfsys@useobject{currentmarker}{}%
\end{pgfscope}%
\begin{pgfscope}%
\pgfsys@transformshift{2.533427in}{0.625278in}%
\pgfsys@useobject{currentmarker}{}%
\end{pgfscope}%
\begin{pgfscope}%
\pgfsys@transformshift{2.533598in}{0.625278in}%
\pgfsys@useobject{currentmarker}{}%
\end{pgfscope}%
\begin{pgfscope}%
\pgfsys@transformshift{2.533762in}{0.625278in}%
\pgfsys@useobject{currentmarker}{}%
\end{pgfscope}%
\begin{pgfscope}%
\pgfsys@transformshift{2.535110in}{0.625278in}%
\pgfsys@useobject{currentmarker}{}%
\end{pgfscope}%
\begin{pgfscope}%
\pgfsys@transformshift{2.546744in}{0.625278in}%
\pgfsys@useobject{currentmarker}{}%
\end{pgfscope}%
\begin{pgfscope}%
\pgfsys@transformshift{2.562773in}{0.625278in}%
\pgfsys@useobject{currentmarker}{}%
\end{pgfscope}%
\begin{pgfscope}%
\pgfsys@transformshift{2.566482in}{0.625278in}%
\pgfsys@useobject{currentmarker}{}%
\end{pgfscope}%
\begin{pgfscope}%
\pgfsys@transformshift{2.567160in}{0.625278in}%
\pgfsys@useobject{currentmarker}{}%
\end{pgfscope}%
\begin{pgfscope}%
\pgfsys@transformshift{2.567495in}{0.625278in}%
\pgfsys@useobject{currentmarker}{}%
\end{pgfscope}%
\begin{pgfscope}%
\pgfsys@transformshift{2.567651in}{0.625278in}%
\pgfsys@useobject{currentmarker}{}%
\end{pgfscope}%
\begin{pgfscope}%
\pgfsys@transformshift{2.567830in}{0.625278in}%
\pgfsys@useobject{currentmarker}{}%
\end{pgfscope}%
\begin{pgfscope}%
\pgfsys@transformshift{2.568001in}{0.625278in}%
\pgfsys@useobject{currentmarker}{}%
\end{pgfscope}%
\begin{pgfscope}%
\pgfsys@transformshift{2.572887in}{0.625278in}%
\pgfsys@useobject{currentmarker}{}%
\end{pgfscope}%
\begin{pgfscope}%
\pgfsys@transformshift{2.596997in}{0.625278in}%
\pgfsys@useobject{currentmarker}{}%
\end{pgfscope}%
\begin{pgfscope}%
\pgfsys@transformshift{2.609309in}{0.625278in}%
\pgfsys@useobject{currentmarker}{}%
\end{pgfscope}%
\begin{pgfscope}%
\pgfsys@transformshift{2.609481in}{0.625278in}%
\pgfsys@useobject{currentmarker}{}%
\end{pgfscope}%
\begin{pgfscope}%
\pgfsys@transformshift{2.609644in}{0.625278in}%
\pgfsys@useobject{currentmarker}{}%
\end{pgfscope}%
\begin{pgfscope}%
\pgfsys@transformshift{2.609652in}{0.625278in}%
\pgfsys@useobject{currentmarker}{}%
\end{pgfscope}%
\begin{pgfscope}%
\pgfsys@transformshift{2.609987in}{0.625278in}%
\pgfsys@useobject{currentmarker}{}%
\end{pgfscope}%
\begin{pgfscope}%
\pgfsys@transformshift{2.610039in}{0.625278in}%
\pgfsys@useobject{currentmarker}{}%
\end{pgfscope}%
\begin{pgfscope}%
\pgfsys@transformshift{2.610322in}{0.625278in}%
\pgfsys@useobject{currentmarker}{}%
\end{pgfscope}%
\begin{pgfscope}%
\pgfsys@transformshift{2.610374in}{0.625278in}%
\pgfsys@useobject{currentmarker}{}%
\end{pgfscope}%
\begin{pgfscope}%
\pgfsys@transformshift{2.610665in}{0.625278in}%
\pgfsys@useobject{currentmarker}{}%
\end{pgfscope}%
\begin{pgfscope}%
\pgfsys@transformshift{2.610821in}{0.625278in}%
\pgfsys@useobject{currentmarker}{}%
\end{pgfscope}%
\begin{pgfscope}%
\pgfsys@transformshift{2.610993in}{0.625278in}%
\pgfsys@useobject{currentmarker}{}%
\end{pgfscope}%
\begin{pgfscope}%
\pgfsys@transformshift{2.611164in}{0.625278in}%
\pgfsys@useobject{currentmarker}{}%
\end{pgfscope}%
\begin{pgfscope}%
\pgfsys@transformshift{2.611499in}{0.625278in}%
\pgfsys@useobject{currentmarker}{}%
\end{pgfscope}%
\begin{pgfscope}%
\pgfsys@transformshift{2.611670in}{0.625278in}%
\pgfsys@useobject{currentmarker}{}%
\end{pgfscope}%
\begin{pgfscope}%
\pgfsys@transformshift{2.612341in}{0.625278in}%
\pgfsys@useobject{currentmarker}{}%
\end{pgfscope}%
\begin{pgfscope}%
\pgfsys@transformshift{2.650625in}{0.625278in}%
\pgfsys@useobject{currentmarker}{}%
\end{pgfscope}%
\begin{pgfscope}%
\pgfsys@transformshift{2.663264in}{0.625278in}%
\pgfsys@useobject{currentmarker}{}%
\end{pgfscope}%
\begin{pgfscope}%
\pgfsys@transformshift{2.666303in}{0.625278in}%
\pgfsys@useobject{currentmarker}{}%
\end{pgfscope}%
\begin{pgfscope}%
\pgfsys@transformshift{2.666311in}{0.625278in}%
\pgfsys@useobject{currentmarker}{}%
\end{pgfscope}%
\begin{pgfscope}%
\pgfsys@transformshift{2.666475in}{0.625278in}%
\pgfsys@useobject{currentmarker}{}%
\end{pgfscope}%
\begin{pgfscope}%
\pgfsys@transformshift{2.666638in}{0.625278in}%
\pgfsys@useobject{currentmarker}{}%
\end{pgfscope}%
\begin{pgfscope}%
\pgfsys@transformshift{2.666810in}{0.625278in}%
\pgfsys@useobject{currentmarker}{}%
\end{pgfscope}%
\begin{pgfscope}%
\pgfsys@transformshift{2.666974in}{0.625278in}%
\pgfsys@useobject{currentmarker}{}%
\end{pgfscope}%
\begin{pgfscope}%
\pgfsys@transformshift{2.667145in}{0.625278in}%
\pgfsys@useobject{currentmarker}{}%
\end{pgfscope}%
\begin{pgfscope}%
\pgfsys@transformshift{2.667316in}{0.625278in}%
\pgfsys@useobject{currentmarker}{}%
\end{pgfscope}%
\begin{pgfscope}%
\pgfsys@transformshift{2.667480in}{0.625278in}%
\pgfsys@useobject{currentmarker}{}%
\end{pgfscope}%
\begin{pgfscope}%
\pgfsys@transformshift{2.667480in}{0.625278in}%
\pgfsys@useobject{currentmarker}{}%
\end{pgfscope}%
\begin{pgfscope}%
\pgfsys@transformshift{2.667823in}{0.625278in}%
\pgfsys@useobject{currentmarker}{}%
\end{pgfscope}%
\begin{pgfscope}%
\pgfsys@transformshift{2.667823in}{0.625278in}%
\pgfsys@useobject{currentmarker}{}%
\end{pgfscope}%
\begin{pgfscope}%
\pgfsys@transformshift{2.667830in}{0.625278in}%
\pgfsys@useobject{currentmarker}{}%
\end{pgfscope}%
\begin{pgfscope}%
\pgfsys@transformshift{2.668158in}{0.625278in}%
\pgfsys@useobject{currentmarker}{}%
\end{pgfscope}%
\begin{pgfscope}%
\pgfsys@transformshift{2.668337in}{0.625278in}%
\pgfsys@useobject{currentmarker}{}%
\end{pgfscope}%
\begin{pgfscope}%
\pgfsys@transformshift{2.668493in}{0.625278in}%
\pgfsys@useobject{currentmarker}{}%
\end{pgfscope}%
\begin{pgfscope}%
\pgfsys@transformshift{2.668828in}{0.625278in}%
\pgfsys@useobject{currentmarker}{}%
\end{pgfscope}%
\begin{pgfscope}%
\pgfsys@transformshift{2.669007in}{0.625278in}%
\pgfsys@useobject{currentmarker}{}%
\end{pgfscope}%
\begin{pgfscope}%
\pgfsys@transformshift{2.669171in}{0.625278in}%
\pgfsys@useobject{currentmarker}{}%
\end{pgfscope}%
\begin{pgfscope}%
\pgfsys@transformshift{2.669171in}{0.625278in}%
\pgfsys@useobject{currentmarker}{}%
\end{pgfscope}%
\begin{pgfscope}%
\pgfsys@transformshift{2.669506in}{0.625278in}%
\pgfsys@useobject{currentmarker}{}%
\end{pgfscope}%
\begin{pgfscope}%
\pgfsys@transformshift{2.695985in}{0.625278in}%
\pgfsys@useobject{currentmarker}{}%
\end{pgfscope}%
\begin{pgfscope}%
\pgfsys@transformshift{2.706107in}{0.625278in}%
\pgfsys@useobject{currentmarker}{}%
\end{pgfscope}%
\begin{pgfscope}%
\pgfsys@transformshift{2.706278in}{0.625278in}%
\pgfsys@useobject{currentmarker}{}%
\end{pgfscope}%
\begin{pgfscope}%
\pgfsys@transformshift{2.706442in}{0.625278in}%
\pgfsys@useobject{currentmarker}{}%
\end{pgfscope}%
\begin{pgfscope}%
\pgfsys@transformshift{2.706613in}{0.625278in}%
\pgfsys@useobject{currentmarker}{}%
\end{pgfscope}%
\begin{pgfscope}%
\pgfsys@transformshift{2.706777in}{0.625278in}%
\pgfsys@useobject{currentmarker}{}%
\end{pgfscope}%
\begin{pgfscope}%
\pgfsys@transformshift{2.706948in}{0.625278in}%
\pgfsys@useobject{currentmarker}{}%
\end{pgfscope}%
\begin{pgfscope}%
\pgfsys@transformshift{2.707120in}{0.625278in}%
\pgfsys@useobject{currentmarker}{}%
\end{pgfscope}%
\begin{pgfscope}%
\pgfsys@transformshift{2.707284in}{0.625278in}%
\pgfsys@useobject{currentmarker}{}%
\end{pgfscope}%
\begin{pgfscope}%
\pgfsys@transformshift{2.707455in}{0.625278in}%
\pgfsys@useobject{currentmarker}{}%
\end{pgfscope}%
\begin{pgfscope}%
\pgfsys@transformshift{2.707626in}{0.625278in}%
\pgfsys@useobject{currentmarker}{}%
\end{pgfscope}%
\begin{pgfscope}%
\pgfsys@transformshift{2.707790in}{0.625278in}%
\pgfsys@useobject{currentmarker}{}%
\end{pgfscope}%
\begin{pgfscope}%
\pgfsys@transformshift{2.707961in}{0.625278in}%
\pgfsys@useobject{currentmarker}{}%
\end{pgfscope}%
\begin{pgfscope}%
\pgfsys@transformshift{2.708133in}{0.625278in}%
\pgfsys@useobject{currentmarker}{}%
\end{pgfscope}%
\begin{pgfscope}%
\pgfsys@transformshift{2.708297in}{0.625278in}%
\pgfsys@useobject{currentmarker}{}%
\end{pgfscope}%
\begin{pgfscope}%
\pgfsys@transformshift{2.708468in}{0.625278in}%
\pgfsys@useobject{currentmarker}{}%
\end{pgfscope}%
\begin{pgfscope}%
\pgfsys@transformshift{2.708632in}{0.625278in}%
\pgfsys@useobject{currentmarker}{}%
\end{pgfscope}%
\begin{pgfscope}%
\pgfsys@transformshift{2.708803in}{0.625278in}%
\pgfsys@useobject{currentmarker}{}%
\end{pgfscope}%
\begin{pgfscope}%
\pgfsys@transformshift{2.708974in}{0.625278in}%
\pgfsys@useobject{currentmarker}{}%
\end{pgfscope}%
\begin{pgfscope}%
\pgfsys@transformshift{2.708974in}{0.625278in}%
\pgfsys@useobject{currentmarker}{}%
\end{pgfscope}%
\begin{pgfscope}%
\pgfsys@transformshift{2.709138in}{0.625278in}%
\pgfsys@useobject{currentmarker}{}%
\end{pgfscope}%
\begin{pgfscope}%
\pgfsys@transformshift{2.709309in}{0.625278in}%
\pgfsys@useobject{currentmarker}{}%
\end{pgfscope}%
\begin{pgfscope}%
\pgfsys@transformshift{2.709481in}{0.625278in}%
\pgfsys@useobject{currentmarker}{}%
\end{pgfscope}%
\begin{pgfscope}%
\pgfsys@transformshift{2.709481in}{0.625278in}%
\pgfsys@useobject{currentmarker}{}%
\end{pgfscope}%
\begin{pgfscope}%
\pgfsys@transformshift{2.709816in}{0.625278in}%
\pgfsys@useobject{currentmarker}{}%
\end{pgfscope}%
\begin{pgfscope}%
\pgfsys@transformshift{2.710151in}{0.625278in}%
\pgfsys@useobject{currentmarker}{}%
\end{pgfscope}%
\begin{pgfscope}%
\pgfsys@transformshift{2.878437in}{0.625278in}%
\pgfsys@useobject{currentmarker}{}%
\end{pgfscope}%
\begin{pgfscope}%
\pgfsys@transformshift{2.891255in}{0.625278in}%
\pgfsys@useobject{currentmarker}{}%
\end{pgfscope}%
\begin{pgfscope}%
\pgfsys@transformshift{2.900193in}{0.625278in}%
\pgfsys@useobject{currentmarker}{}%
\end{pgfscope}%
\begin{pgfscope}%
\pgfsys@transformshift{2.900193in}{0.625278in}%
\pgfsys@useobject{currentmarker}{}%
\end{pgfscope}%
\begin{pgfscope}%
\pgfsys@transformshift{2.900528in}{0.625278in}%
\pgfsys@useobject{currentmarker}{}%
\end{pgfscope}%
\begin{pgfscope}%
\pgfsys@transformshift{2.900528in}{0.625278in}%
\pgfsys@useobject{currentmarker}{}%
\end{pgfscope}%
\begin{pgfscope}%
\pgfsys@transformshift{2.900692in}{0.625278in}%
\pgfsys@useobject{currentmarker}{}%
\end{pgfscope}%
\begin{pgfscope}%
\pgfsys@transformshift{2.900863in}{0.625278in}%
\pgfsys@useobject{currentmarker}{}%
\end{pgfscope}%
\begin{pgfscope}%
\pgfsys@transformshift{2.901035in}{0.625278in}%
\pgfsys@useobject{currentmarker}{}%
\end{pgfscope}%
\begin{pgfscope}%
\pgfsys@transformshift{2.901035in}{0.625278in}%
\pgfsys@useobject{currentmarker}{}%
\end{pgfscope}%
\begin{pgfscope}%
\pgfsys@transformshift{2.901370in}{0.625278in}%
\pgfsys@useobject{currentmarker}{}%
\end{pgfscope}%
\begin{pgfscope}%
\pgfsys@transformshift{2.901705in}{0.625278in}%
\pgfsys@useobject{currentmarker}{}%
\end{pgfscope}%
\begin{pgfscope}%
\pgfsys@transformshift{2.901713in}{0.625278in}%
\pgfsys@useobject{currentmarker}{}%
\end{pgfscope}%
\begin{pgfscope}%
\pgfsys@transformshift{2.902048in}{0.625278in}%
\pgfsys@useobject{currentmarker}{}%
\end{pgfscope}%
\begin{pgfscope}%
\pgfsys@transformshift{2.902055in}{0.625278in}%
\pgfsys@useobject{currentmarker}{}%
\end{pgfscope}%
\begin{pgfscope}%
\pgfsys@transformshift{2.902212in}{0.625278in}%
\pgfsys@useobject{currentmarker}{}%
\end{pgfscope}%
\begin{pgfscope}%
\pgfsys@transformshift{2.902554in}{0.625278in}%
\pgfsys@useobject{currentmarker}{}%
\end{pgfscope}%
\begin{pgfscope}%
\pgfsys@transformshift{2.902554in}{0.625278in}%
\pgfsys@useobject{currentmarker}{}%
\end{pgfscope}%
\begin{pgfscope}%
\pgfsys@transformshift{2.902718in}{0.625278in}%
\pgfsys@useobject{currentmarker}{}%
\end{pgfscope}%
\begin{pgfscope}%
\pgfsys@transformshift{2.902889in}{0.625278in}%
\pgfsys@useobject{currentmarker}{}%
\end{pgfscope}%
\begin{pgfscope}%
\pgfsys@transformshift{2.903053in}{0.625278in}%
\pgfsys@useobject{currentmarker}{}%
\end{pgfscope}%
\begin{pgfscope}%
\pgfsys@transformshift{2.903225in}{0.625278in}%
\pgfsys@useobject{currentmarker}{}%
\end{pgfscope}%
\begin{pgfscope}%
\pgfsys@transformshift{2.903560in}{0.625278in}%
\pgfsys@useobject{currentmarker}{}%
\end{pgfscope}%
\begin{pgfscope}%
\pgfsys@transformshift{2.940332in}{0.625278in}%
\pgfsys@useobject{currentmarker}{}%
\end{pgfscope}%
\begin{pgfscope}%
\pgfsys@transformshift{2.954491in}{0.625278in}%
\pgfsys@useobject{currentmarker}{}%
\end{pgfscope}%
\begin{pgfscope}%
\pgfsys@transformshift{2.973879in}{0.625278in}%
\pgfsys@useobject{currentmarker}{}%
\end{pgfscope}%
\begin{pgfscope}%
\pgfsys@transformshift{2.986198in}{0.625278in}%
\pgfsys@useobject{currentmarker}{}%
\end{pgfscope}%
\begin{pgfscope}%
\pgfsys@transformshift{2.986362in}{0.625278in}%
\pgfsys@useobject{currentmarker}{}%
\end{pgfscope}%
\begin{pgfscope}%
\pgfsys@transformshift{2.986526in}{0.625278in}%
\pgfsys@useobject{currentmarker}{}%
\end{pgfscope}%
\begin{pgfscope}%
\pgfsys@transformshift{2.986533in}{0.625278in}%
\pgfsys@useobject{currentmarker}{}%
\end{pgfscope}%
\begin{pgfscope}%
\pgfsys@transformshift{2.986868in}{0.625278in}%
\pgfsys@useobject{currentmarker}{}%
\end{pgfscope}%
\begin{pgfscope}%
\pgfsys@transformshift{2.987040in}{0.625278in}%
\pgfsys@useobject{currentmarker}{}%
\end{pgfscope}%
\begin{pgfscope}%
\pgfsys@transformshift{2.987204in}{0.625278in}%
\pgfsys@useobject{currentmarker}{}%
\end{pgfscope}%
\begin{pgfscope}%
\pgfsys@transformshift{2.987367in}{0.625278in}%
\pgfsys@useobject{currentmarker}{}%
\end{pgfscope}%
\begin{pgfscope}%
\pgfsys@transformshift{2.987375in}{0.625278in}%
\pgfsys@useobject{currentmarker}{}%
\end{pgfscope}%
\begin{pgfscope}%
\pgfsys@transformshift{2.987539in}{0.625278in}%
\pgfsys@useobject{currentmarker}{}%
\end{pgfscope}%
\begin{pgfscope}%
\pgfsys@transformshift{2.987710in}{0.625278in}%
\pgfsys@useobject{currentmarker}{}%
\end{pgfscope}%
\begin{pgfscope}%
\pgfsys@transformshift{2.987874in}{0.625278in}%
\pgfsys@useobject{currentmarker}{}%
\end{pgfscope}%
\begin{pgfscope}%
\pgfsys@transformshift{2.988045in}{0.625278in}%
\pgfsys@useobject{currentmarker}{}%
\end{pgfscope}%
\begin{pgfscope}%
\pgfsys@transformshift{2.988045in}{0.625278in}%
\pgfsys@useobject{currentmarker}{}%
\end{pgfscope}%
\begin{pgfscope}%
\pgfsys@transformshift{2.988216in}{0.625278in}%
\pgfsys@useobject{currentmarker}{}%
\end{pgfscope}%
\begin{pgfscope}%
\pgfsys@transformshift{2.988380in}{0.625278in}%
\pgfsys@useobject{currentmarker}{}%
\end{pgfscope}%
\begin{pgfscope}%
\pgfsys@transformshift{2.989565in}{0.625278in}%
\pgfsys@useobject{currentmarker}{}%
\end{pgfscope}%
\begin{pgfscope}%
\pgfsys@transformshift{3.020594in}{0.625278in}%
\pgfsys@useobject{currentmarker}{}%
\end{pgfscope}%
\begin{pgfscope}%
\pgfsys@transformshift{3.027342in}{0.625278in}%
\pgfsys@useobject{currentmarker}{}%
\end{pgfscope}%
\begin{pgfscope}%
\pgfsys@transformshift{3.027350in}{0.625278in}%
\pgfsys@useobject{currentmarker}{}%
\end{pgfscope}%
\begin{pgfscope}%
\pgfsys@transformshift{3.027685in}{0.625278in}%
\pgfsys@useobject{currentmarker}{}%
\end{pgfscope}%
\begin{pgfscope}%
\pgfsys@transformshift{3.027849in}{0.625278in}%
\pgfsys@useobject{currentmarker}{}%
\end{pgfscope}%
\begin{pgfscope}%
\pgfsys@transformshift{3.028027in}{0.625278in}%
\pgfsys@useobject{currentmarker}{}%
\end{pgfscope}%
\begin{pgfscope}%
\pgfsys@transformshift{3.028184in}{0.625278in}%
\pgfsys@useobject{currentmarker}{}%
\end{pgfscope}%
\begin{pgfscope}%
\pgfsys@transformshift{3.028363in}{0.625278in}%
\pgfsys@useobject{currentmarker}{}%
\end{pgfscope}%
\begin{pgfscope}%
\pgfsys@transformshift{3.028519in}{0.625278in}%
\pgfsys@useobject{currentmarker}{}%
\end{pgfscope}%
\begin{pgfscope}%
\pgfsys@transformshift{3.028690in}{0.625278in}%
\pgfsys@useobject{currentmarker}{}%
\end{pgfscope}%
\begin{pgfscope}%
\pgfsys@transformshift{3.028854in}{0.625278in}%
\pgfsys@useobject{currentmarker}{}%
\end{pgfscope}%
\begin{pgfscope}%
\pgfsys@transformshift{3.028869in}{0.625278in}%
\pgfsys@useobject{currentmarker}{}%
\end{pgfscope}%
\begin{pgfscope}%
\pgfsys@transformshift{3.029539in}{0.625278in}%
\pgfsys@useobject{currentmarker}{}%
\end{pgfscope}%
\begin{pgfscope}%
\pgfsys@transformshift{3.029703in}{0.625278in}%
\pgfsys@useobject{currentmarker}{}%
\end{pgfscope}%
\begin{pgfscope}%
\pgfsys@transformshift{3.029867in}{0.625278in}%
\pgfsys@useobject{currentmarker}{}%
\end{pgfscope}%
\begin{pgfscope}%
\pgfsys@transformshift{3.030038in}{0.625278in}%
\pgfsys@useobject{currentmarker}{}%
\end{pgfscope}%
\begin{pgfscope}%
\pgfsys@transformshift{3.030381in}{0.625278in}%
\pgfsys@useobject{currentmarker}{}%
\end{pgfscope}%
\begin{pgfscope}%
\pgfsys@transformshift{3.030537in}{0.625278in}%
\pgfsys@useobject{currentmarker}{}%
\end{pgfscope}%
\begin{pgfscope}%
\pgfsys@transformshift{3.030716in}{0.625278in}%
\pgfsys@useobject{currentmarker}{}%
\end{pgfscope}%
\begin{pgfscope}%
\pgfsys@transformshift{3.030880in}{0.625278in}%
\pgfsys@useobject{currentmarker}{}%
\end{pgfscope}%
\begin{pgfscope}%
\pgfsys@transformshift{3.058036in}{0.625278in}%
\pgfsys@useobject{currentmarker}{}%
\end{pgfscope}%
\begin{pgfscope}%
\pgfsys@transformshift{3.067309in}{0.625278in}%
\pgfsys@useobject{currentmarker}{}%
\end{pgfscope}%
\begin{pgfscope}%
\pgfsys@transformshift{3.067473in}{0.625278in}%
\pgfsys@useobject{currentmarker}{}%
\end{pgfscope}%
\begin{pgfscope}%
\pgfsys@transformshift{3.067645in}{0.625278in}%
\pgfsys@useobject{currentmarker}{}%
\end{pgfscope}%
\begin{pgfscope}%
\pgfsys@transformshift{3.067816in}{0.625278in}%
\pgfsys@useobject{currentmarker}{}%
\end{pgfscope}%
\begin{pgfscope}%
\pgfsys@transformshift{3.067816in}{0.625278in}%
\pgfsys@useobject{currentmarker}{}%
\end{pgfscope}%
\begin{pgfscope}%
\pgfsys@transformshift{3.067972in}{0.625278in}%
\pgfsys@useobject{currentmarker}{}%
\end{pgfscope}%
\begin{pgfscope}%
\pgfsys@transformshift{3.068151in}{0.625278in}%
\pgfsys@useobject{currentmarker}{}%
\end{pgfscope}%
\begin{pgfscope}%
\pgfsys@transformshift{3.068315in}{0.625278in}%
\pgfsys@useobject{currentmarker}{}%
\end{pgfscope}%
\begin{pgfscope}%
\pgfsys@transformshift{3.068486in}{0.625278in}%
\pgfsys@useobject{currentmarker}{}%
\end{pgfscope}%
\begin{pgfscope}%
\pgfsys@transformshift{3.068650in}{0.625278in}%
\pgfsys@useobject{currentmarker}{}%
\end{pgfscope}%
\begin{pgfscope}%
\pgfsys@transformshift{3.068829in}{0.625278in}%
\pgfsys@useobject{currentmarker}{}%
\end{pgfscope}%
\begin{pgfscope}%
\pgfsys@transformshift{3.068985in}{0.625278in}%
\pgfsys@useobject{currentmarker}{}%
\end{pgfscope}%
\begin{pgfscope}%
\pgfsys@transformshift{3.069157in}{0.625278in}%
\pgfsys@useobject{currentmarker}{}%
\end{pgfscope}%
\begin{pgfscope}%
\pgfsys@transformshift{3.069164in}{0.625278in}%
\pgfsys@useobject{currentmarker}{}%
\end{pgfscope}%
\begin{pgfscope}%
\pgfsys@transformshift{3.069492in}{0.625278in}%
\pgfsys@useobject{currentmarker}{}%
\end{pgfscope}%
\begin{pgfscope}%
\pgfsys@transformshift{3.069499in}{0.625278in}%
\pgfsys@useobject{currentmarker}{}%
\end{pgfscope}%
\begin{pgfscope}%
\pgfsys@transformshift{3.069827in}{0.625278in}%
\pgfsys@useobject{currentmarker}{}%
\end{pgfscope}%
\begin{pgfscope}%
\pgfsys@transformshift{3.069834in}{0.625278in}%
\pgfsys@useobject{currentmarker}{}%
\end{pgfscope}%
\begin{pgfscope}%
\pgfsys@transformshift{3.239975in}{0.625278in}%
\pgfsys@useobject{currentmarker}{}%
\end{pgfscope}%
\begin{pgfscope}%
\pgfsys@transformshift{3.247564in}{0.625278in}%
\pgfsys@useobject{currentmarker}{}%
\end{pgfscope}%
\begin{pgfscope}%
\pgfsys@transformshift{3.247564in}{0.625278in}%
\pgfsys@useobject{currentmarker}{}%
\end{pgfscope}%
\begin{pgfscope}%
\pgfsys@transformshift{3.247900in}{0.625278in}%
\pgfsys@useobject{currentmarker}{}%
\end{pgfscope}%
\begin{pgfscope}%
\pgfsys@transformshift{3.247900in}{0.625278in}%
\pgfsys@useobject{currentmarker}{}%
\end{pgfscope}%
\begin{pgfscope}%
\pgfsys@transformshift{3.248242in}{0.625278in}%
\pgfsys@useobject{currentmarker}{}%
\end{pgfscope}%
\begin{pgfscope}%
\pgfsys@transformshift{3.248242in}{0.625278in}%
\pgfsys@useobject{currentmarker}{}%
\end{pgfscope}%
\begin{pgfscope}%
\pgfsys@transformshift{3.248294in}{0.625278in}%
\pgfsys@useobject{currentmarker}{}%
\end{pgfscope}%
\begin{pgfscope}%
\pgfsys@transformshift{3.248577in}{0.625278in}%
\pgfsys@useobject{currentmarker}{}%
\end{pgfscope}%
\begin{pgfscope}%
\pgfsys@transformshift{3.248577in}{0.625278in}%
\pgfsys@useobject{currentmarker}{}%
\end{pgfscope}%
\begin{pgfscope}%
\pgfsys@transformshift{3.248913in}{0.625278in}%
\pgfsys@useobject{currentmarker}{}%
\end{pgfscope}%
\begin{pgfscope}%
\pgfsys@transformshift{3.248913in}{0.625278in}%
\pgfsys@useobject{currentmarker}{}%
\end{pgfscope}%
\begin{pgfscope}%
\pgfsys@transformshift{3.249248in}{0.625278in}%
\pgfsys@useobject{currentmarker}{}%
\end{pgfscope}%
\begin{pgfscope}%
\pgfsys@transformshift{3.249248in}{0.625278in}%
\pgfsys@useobject{currentmarker}{}%
\end{pgfscope}%
\begin{pgfscope}%
\pgfsys@transformshift{3.249590in}{0.625278in}%
\pgfsys@useobject{currentmarker}{}%
\end{pgfscope}%
\begin{pgfscope}%
\pgfsys@transformshift{3.249590in}{0.625278in}%
\pgfsys@useobject{currentmarker}{}%
\end{pgfscope}%
\begin{pgfscope}%
\pgfsys@transformshift{3.249926in}{0.625278in}%
\pgfsys@useobject{currentmarker}{}%
\end{pgfscope}%
\begin{pgfscope}%
\pgfsys@transformshift{3.249926in}{0.625278in}%
\pgfsys@useobject{currentmarker}{}%
\end{pgfscope}%
\begin{pgfscope}%
\pgfsys@transformshift{3.250261in}{0.625278in}%
\pgfsys@useobject{currentmarker}{}%
\end{pgfscope}%
\begin{pgfscope}%
\pgfsys@transformshift{3.250261in}{0.625278in}%
\pgfsys@useobject{currentmarker}{}%
\end{pgfscope}%
\begin{pgfscope}%
\pgfsys@transformshift{3.250596in}{0.625278in}%
\pgfsys@useobject{currentmarker}{}%
\end{pgfscope}%
\begin{pgfscope}%
\pgfsys@transformshift{3.250596in}{0.625278in}%
\pgfsys@useobject{currentmarker}{}%
\end{pgfscope}%
\begin{pgfscope}%
\pgfsys@transformshift{3.250938in}{0.625278in}%
\pgfsys@useobject{currentmarker}{}%
\end{pgfscope}%
\begin{pgfscope}%
\pgfsys@transformshift{3.250938in}{0.625278in}%
\pgfsys@useobject{currentmarker}{}%
\end{pgfscope}%
\begin{pgfscope}%
\pgfsys@transformshift{3.250938in}{0.625278in}%
\pgfsys@useobject{currentmarker}{}%
\end{pgfscope}%
\begin{pgfscope}%
\pgfsys@transformshift{3.251274in}{0.625278in}%
\pgfsys@useobject{currentmarker}{}%
\end{pgfscope}%
\begin{pgfscope}%
\pgfsys@transformshift{3.251274in}{0.625278in}%
\pgfsys@useobject{currentmarker}{}%
\end{pgfscope}%
\begin{pgfscope}%
\pgfsys@transformshift{3.251609in}{0.625278in}%
\pgfsys@useobject{currentmarker}{}%
\end{pgfscope}%
\begin{pgfscope}%
\pgfsys@transformshift{3.251951in}{0.625278in}%
\pgfsys@useobject{currentmarker}{}%
\end{pgfscope}%
\begin{pgfscope}%
\pgfsys@transformshift{3.266960in}{0.625278in}%
\pgfsys@useobject{currentmarker}{}%
\end{pgfscope}%
\begin{pgfscope}%
\pgfsys@transformshift{3.277082in}{0.625278in}%
\pgfsys@useobject{currentmarker}{}%
\end{pgfscope}%
\begin{pgfscope}%
\pgfsys@transformshift{3.286683in}{0.625278in}%
\pgfsys@useobject{currentmarker}{}%
\end{pgfscope}%
\begin{pgfscope}%
\pgfsys@transformshift{3.287025in}{0.625278in}%
\pgfsys@useobject{currentmarker}{}%
\end{pgfscope}%
\begin{pgfscope}%
\pgfsys@transformshift{3.287204in}{0.625278in}%
\pgfsys@useobject{currentmarker}{}%
\end{pgfscope}%
\begin{pgfscope}%
\pgfsys@transformshift{3.287360in}{0.625278in}%
\pgfsys@useobject{currentmarker}{}%
\end{pgfscope}%
\begin{pgfscope}%
\pgfsys@transformshift{3.287547in}{0.625278in}%
\pgfsys@useobject{currentmarker}{}%
\end{pgfscope}%
\begin{pgfscope}%
\pgfsys@transformshift{3.287696in}{0.625278in}%
\pgfsys@useobject{currentmarker}{}%
\end{pgfscope}%
\begin{pgfscope}%
\pgfsys@transformshift{3.288038in}{0.625278in}%
\pgfsys@useobject{currentmarker}{}%
\end{pgfscope}%
\begin{pgfscope}%
\pgfsys@transformshift{3.288046in}{0.625278in}%
\pgfsys@useobject{currentmarker}{}%
\end{pgfscope}%
\begin{pgfscope}%
\pgfsys@transformshift{3.294451in}{0.625278in}%
\pgfsys@useobject{currentmarker}{}%
\end{pgfscope}%
\begin{pgfscope}%
\pgfsys@transformshift{3.329182in}{0.625278in}%
\pgfsys@useobject{currentmarker}{}%
\end{pgfscope}%
\begin{pgfscope}%
\pgfsys@transformshift{3.477574in}{0.625278in}%
\pgfsys@useobject{currentmarker}{}%
\end{pgfscope}%
\begin{pgfscope}%
\pgfsys@transformshift{3.489208in}{0.625278in}%
\pgfsys@useobject{currentmarker}{}%
\end{pgfscope}%
\begin{pgfscope}%
\pgfsys@transformshift{3.492753in}{0.625278in}%
\pgfsys@useobject{currentmarker}{}%
\end{pgfscope}%
\begin{pgfscope}%
\pgfsys@transformshift{3.492753in}{0.625278in}%
\pgfsys@useobject{currentmarker}{}%
\end{pgfscope}%
\begin{pgfscope}%
\pgfsys@transformshift{3.493088in}{0.625278in}%
\pgfsys@useobject{currentmarker}{}%
\end{pgfscope}%
\begin{pgfscope}%
\pgfsys@transformshift{3.493088in}{0.625278in}%
\pgfsys@useobject{currentmarker}{}%
\end{pgfscope}%
\begin{pgfscope}%
\pgfsys@transformshift{3.493424in}{0.625278in}%
\pgfsys@useobject{currentmarker}{}%
\end{pgfscope}%
\begin{pgfscope}%
\pgfsys@transformshift{3.493424in}{0.625278in}%
\pgfsys@useobject{currentmarker}{}%
\end{pgfscope}%
\begin{pgfscope}%
\pgfsys@transformshift{3.493759in}{0.625278in}%
\pgfsys@useobject{currentmarker}{}%
\end{pgfscope}%
\begin{pgfscope}%
\pgfsys@transformshift{3.493759in}{0.625278in}%
\pgfsys@useobject{currentmarker}{}%
\end{pgfscope}%
\begin{pgfscope}%
\pgfsys@transformshift{3.494101in}{0.625278in}%
\pgfsys@useobject{currentmarker}{}%
\end{pgfscope}%
\begin{pgfscope}%
\pgfsys@transformshift{3.494101in}{0.625278in}%
\pgfsys@useobject{currentmarker}{}%
\end{pgfscope}%
\begin{pgfscope}%
\pgfsys@transformshift{3.494437in}{0.625278in}%
\pgfsys@useobject{currentmarker}{}%
\end{pgfscope}%
\begin{pgfscope}%
\pgfsys@transformshift{3.494437in}{0.625278in}%
\pgfsys@useobject{currentmarker}{}%
\end{pgfscope}%
\begin{pgfscope}%
\pgfsys@transformshift{3.494772in}{0.625278in}%
\pgfsys@useobject{currentmarker}{}%
\end{pgfscope}%
\begin{pgfscope}%
\pgfsys@transformshift{3.495107in}{0.625278in}%
\pgfsys@useobject{currentmarker}{}%
\end{pgfscope}%
\begin{pgfscope}%
\pgfsys@transformshift{3.495278in}{0.625278in}%
\pgfsys@useobject{currentmarker}{}%
\end{pgfscope}%
\begin{pgfscope}%
\pgfsys@transformshift{3.495449in}{0.625278in}%
\pgfsys@useobject{currentmarker}{}%
\end{pgfscope}%
\begin{pgfscope}%
\pgfsys@transformshift{3.495613in}{0.625278in}%
\pgfsys@useobject{currentmarker}{}%
\end{pgfscope}%
\begin{pgfscope}%
\pgfsys@transformshift{3.495785in}{0.625278in}%
\pgfsys@useobject{currentmarker}{}%
\end{pgfscope}%
\begin{pgfscope}%
\pgfsys@transformshift{3.496120in}{0.625278in}%
\pgfsys@useobject{currentmarker}{}%
\end{pgfscope}%
\begin{pgfscope}%
\pgfsys@transformshift{3.496120in}{0.625278in}%
\pgfsys@useobject{currentmarker}{}%
\end{pgfscope}%
\begin{pgfscope}%
\pgfsys@transformshift{3.496291in}{0.625278in}%
\pgfsys@useobject{currentmarker}{}%
\end{pgfscope}%
\begin{pgfscope}%
\pgfsys@transformshift{3.496462in}{0.625278in}%
\pgfsys@useobject{currentmarker}{}%
\end{pgfscope}%
\begin{pgfscope}%
\pgfsys@transformshift{3.496462in}{0.625278in}%
\pgfsys@useobject{currentmarker}{}%
\end{pgfscope}%
\begin{pgfscope}%
\pgfsys@transformshift{3.496462in}{0.625278in}%
\pgfsys@useobject{currentmarker}{}%
\end{pgfscope}%
\begin{pgfscope}%
\pgfsys@transformshift{3.496626in}{0.625278in}%
\pgfsys@useobject{currentmarker}{}%
\end{pgfscope}%
\begin{pgfscope}%
\pgfsys@transformshift{3.496798in}{0.625278in}%
\pgfsys@useobject{currentmarker}{}%
\end{pgfscope}%
\begin{pgfscope}%
\pgfsys@transformshift{3.515351in}{0.625278in}%
\pgfsys@useobject{currentmarker}{}%
\end{pgfscope}%
\begin{pgfscope}%
\pgfsys@transformshift{3.527827in}{0.625278in}%
\pgfsys@useobject{currentmarker}{}%
\end{pgfscope}%
\begin{pgfscope}%
\pgfsys@transformshift{3.536087in}{0.625278in}%
\pgfsys@useobject{currentmarker}{}%
\end{pgfscope}%
\begin{pgfscope}%
\pgfsys@transformshift{3.536095in}{0.625278in}%
\pgfsys@useobject{currentmarker}{}%
\end{pgfscope}%
\begin{pgfscope}%
\pgfsys@transformshift{3.536266in}{0.625278in}%
\pgfsys@useobject{currentmarker}{}%
\end{pgfscope}%
\begin{pgfscope}%
\pgfsys@transformshift{3.536430in}{0.625278in}%
\pgfsys@useobject{currentmarker}{}%
\end{pgfscope}%
\begin{pgfscope}%
\pgfsys@transformshift{3.536594in}{0.625278in}%
\pgfsys@useobject{currentmarker}{}%
\end{pgfscope}%
\begin{pgfscope}%
\pgfsys@transformshift{3.536772in}{0.625278in}%
\pgfsys@useobject{currentmarker}{}%
\end{pgfscope}%
\begin{pgfscope}%
\pgfsys@transformshift{3.536929in}{0.625278in}%
\pgfsys@useobject{currentmarker}{}%
\end{pgfscope}%
\begin{pgfscope}%
\pgfsys@transformshift{3.536936in}{0.625278in}%
\pgfsys@useobject{currentmarker}{}%
\end{pgfscope}%
\begin{pgfscope}%
\pgfsys@transformshift{3.537100in}{0.625278in}%
\pgfsys@useobject{currentmarker}{}%
\end{pgfscope}%
\begin{pgfscope}%
\pgfsys@transformshift{3.537271in}{0.625278in}%
\pgfsys@useobject{currentmarker}{}%
\end{pgfscope}%
\begin{pgfscope}%
\pgfsys@transformshift{3.537435in}{0.625278in}%
\pgfsys@useobject{currentmarker}{}%
\end{pgfscope}%
\begin{pgfscope}%
\pgfsys@transformshift{3.537443in}{0.625278in}%
\pgfsys@useobject{currentmarker}{}%
\end{pgfscope}%
\begin{pgfscope}%
\pgfsys@transformshift{3.537778in}{0.625278in}%
\pgfsys@useobject{currentmarker}{}%
\end{pgfscope}%
\begin{pgfscope}%
\pgfsys@transformshift{3.538284in}{0.625278in}%
\pgfsys@useobject{currentmarker}{}%
\end{pgfscope}%
\begin{pgfscope}%
\pgfsys@transformshift{3.538448in}{0.625278in}%
\pgfsys@useobject{currentmarker}{}%
\end{pgfscope}%
\begin{pgfscope}%
\pgfsys@transformshift{3.538783in}{0.625278in}%
\pgfsys@useobject{currentmarker}{}%
\end{pgfscope}%
\begin{pgfscope}%
\pgfsys@transformshift{3.538791in}{0.625278in}%
\pgfsys@useobject{currentmarker}{}%
\end{pgfscope}%
\begin{pgfscope}%
\pgfsys@transformshift{3.574542in}{0.625278in}%
\pgfsys@useobject{currentmarker}{}%
\end{pgfscope}%
\begin{pgfscope}%
\pgfsys@transformshift{3.731186in}{0.625278in}%
\pgfsys@useobject{currentmarker}{}%
\end{pgfscope}%
\begin{pgfscope}%
\pgfsys@transformshift{3.743163in}{0.625278in}%
\pgfsys@useobject{currentmarker}{}%
\end{pgfscope}%
\begin{pgfscope}%
\pgfsys@transformshift{3.751930in}{0.625278in}%
\pgfsys@useobject{currentmarker}{}%
\end{pgfscope}%
\begin{pgfscope}%
\pgfsys@transformshift{3.752101in}{0.625278in}%
\pgfsys@useobject{currentmarker}{}%
\end{pgfscope}%
\begin{pgfscope}%
\pgfsys@transformshift{3.752265in}{0.625278in}%
\pgfsys@useobject{currentmarker}{}%
\end{pgfscope}%
\begin{pgfscope}%
\pgfsys@transformshift{3.752265in}{0.625278in}%
\pgfsys@useobject{currentmarker}{}%
\end{pgfscope}%
\begin{pgfscope}%
\pgfsys@transformshift{3.752608in}{0.625278in}%
\pgfsys@useobject{currentmarker}{}%
\end{pgfscope}%
\begin{pgfscope}%
\pgfsys@transformshift{3.752772in}{0.625278in}%
\pgfsys@useobject{currentmarker}{}%
\end{pgfscope}%
\begin{pgfscope}%
\pgfsys@transformshift{3.753114in}{0.625278in}%
\pgfsys@useobject{currentmarker}{}%
\end{pgfscope}%
\begin{pgfscope}%
\pgfsys@transformshift{3.753114in}{0.625278in}%
\pgfsys@useobject{currentmarker}{}%
\end{pgfscope}%
\begin{pgfscope}%
\pgfsys@transformshift{3.753449in}{0.625278in}%
\pgfsys@useobject{currentmarker}{}%
\end{pgfscope}%
\begin{pgfscope}%
\pgfsys@transformshift{3.753621in}{0.625278in}%
\pgfsys@useobject{currentmarker}{}%
\end{pgfscope}%
\begin{pgfscope}%
\pgfsys@transformshift{3.753673in}{0.625278in}%
\pgfsys@useobject{currentmarker}{}%
\end{pgfscope}%
\begin{pgfscope}%
\pgfsys@transformshift{3.753784in}{0.625278in}%
\pgfsys@useobject{currentmarker}{}%
\end{pgfscope}%
\begin{pgfscope}%
\pgfsys@transformshift{3.753956in}{0.625278in}%
\pgfsys@useobject{currentmarker}{}%
\end{pgfscope}%
\begin{pgfscope}%
\pgfsys@transformshift{3.754127in}{0.625278in}%
\pgfsys@useobject{currentmarker}{}%
\end{pgfscope}%
\begin{pgfscope}%
\pgfsys@transformshift{3.754291in}{0.625278in}%
\pgfsys@useobject{currentmarker}{}%
\end{pgfscope}%
\begin{pgfscope}%
\pgfsys@transformshift{3.754462in}{0.625278in}%
\pgfsys@useobject{currentmarker}{}%
\end{pgfscope}%
\begin{pgfscope}%
\pgfsys@transformshift{3.754626in}{0.625278in}%
\pgfsys@useobject{currentmarker}{}%
\end{pgfscope}%
\begin{pgfscope}%
\pgfsys@transformshift{3.754626in}{0.625278in}%
\pgfsys@useobject{currentmarker}{}%
\end{pgfscope}%
\begin{pgfscope}%
\pgfsys@transformshift{3.754969in}{0.625278in}%
\pgfsys@useobject{currentmarker}{}%
\end{pgfscope}%
\begin{pgfscope}%
\pgfsys@transformshift{3.754969in}{0.625278in}%
\pgfsys@useobject{currentmarker}{}%
\end{pgfscope}%
\begin{pgfscope}%
\pgfsys@transformshift{3.755304in}{0.625278in}%
\pgfsys@useobject{currentmarker}{}%
\end{pgfscope}%
\begin{pgfscope}%
\pgfsys@transformshift{3.785656in}{0.625278in}%
\pgfsys@useobject{currentmarker}{}%
\end{pgfscope}%
\begin{pgfscope}%
\pgfsys@transformshift{3.807747in}{0.625278in}%
\pgfsys@useobject{currentmarker}{}%
\end{pgfscope}%
\begin{pgfscope}%
\pgfsys@transformshift{3.820223in}{0.625278in}%
\pgfsys@useobject{currentmarker}{}%
\end{pgfscope}%
\begin{pgfscope}%
\pgfsys@transformshift{3.838940in}{0.625278in}%
\pgfsys@useobject{currentmarker}{}%
\end{pgfscope}%
\begin{pgfscope}%
\pgfsys@transformshift{3.846873in}{0.625278in}%
\pgfsys@useobject{currentmarker}{}%
\end{pgfscope}%
\begin{pgfscope}%
\pgfsys@transformshift{3.847037in}{0.625278in}%
\pgfsys@useobject{currentmarker}{}%
\end{pgfscope}%
\begin{pgfscope}%
\pgfsys@transformshift{3.847208in}{0.625278in}%
\pgfsys@useobject{currentmarker}{}%
\end{pgfscope}%
\begin{pgfscope}%
\pgfsys@transformshift{3.847215in}{0.625278in}%
\pgfsys@useobject{currentmarker}{}%
\end{pgfscope}%
\begin{pgfscope}%
\pgfsys@transformshift{3.847372in}{0.625278in}%
\pgfsys@useobject{currentmarker}{}%
\end{pgfscope}%
\begin{pgfscope}%
\pgfsys@transformshift{3.847543in}{0.625278in}%
\pgfsys@useobject{currentmarker}{}%
\end{pgfscope}%
\begin{pgfscope}%
\pgfsys@transformshift{3.847550in}{0.625278in}%
\pgfsys@useobject{currentmarker}{}%
\end{pgfscope}%
\begin{pgfscope}%
\pgfsys@transformshift{3.847878in}{0.625278in}%
\pgfsys@useobject{currentmarker}{}%
\end{pgfscope}%
\begin{pgfscope}%
\pgfsys@transformshift{3.848049in}{0.625278in}%
\pgfsys@useobject{currentmarker}{}%
\end{pgfscope}%
\begin{pgfscope}%
\pgfsys@transformshift{3.848221in}{0.625278in}%
\pgfsys@useobject{currentmarker}{}%
\end{pgfscope}%
\begin{pgfscope}%
\pgfsys@transformshift{3.848385in}{0.625278in}%
\pgfsys@useobject{currentmarker}{}%
\end{pgfscope}%
\begin{pgfscope}%
\pgfsys@transformshift{3.848392in}{0.625278in}%
\pgfsys@useobject{currentmarker}{}%
\end{pgfscope}%
\begin{pgfscope}%
\pgfsys@transformshift{3.848556in}{0.625278in}%
\pgfsys@useobject{currentmarker}{}%
\end{pgfscope}%
\begin{pgfscope}%
\pgfsys@transformshift{3.848891in}{0.625278in}%
\pgfsys@useobject{currentmarker}{}%
\end{pgfscope}%
\begin{pgfscope}%
\pgfsys@transformshift{3.849226in}{0.625278in}%
\pgfsys@useobject{currentmarker}{}%
\end{pgfscope}%
\begin{pgfscope}%
\pgfsys@transformshift{3.849398in}{0.625278in}%
\pgfsys@useobject{currentmarker}{}%
\end{pgfscope}%
\begin{pgfscope}%
\pgfsys@transformshift{3.849569in}{0.625278in}%
\pgfsys@useobject{currentmarker}{}%
\end{pgfscope}%
\begin{pgfscope}%
\pgfsys@transformshift{3.849569in}{0.625278in}%
\pgfsys@useobject{currentmarker}{}%
\end{pgfscope}%
\begin{pgfscope}%
\pgfsys@transformshift{3.849733in}{0.625278in}%
\pgfsys@useobject{currentmarker}{}%
\end{pgfscope}%
\begin{pgfscope}%
\pgfsys@transformshift{3.962045in}{0.625278in}%
\pgfsys@useobject{currentmarker}{}%
\end{pgfscope}%
\begin{pgfscope}%
\pgfsys@transformshift{3.982781in}{0.625278in}%
\pgfsys@useobject{currentmarker}{}%
\end{pgfscope}%
\begin{pgfscope}%
\pgfsys@transformshift{3.998973in}{0.625278in}%
\pgfsys@useobject{currentmarker}{}%
\end{pgfscope}%
\begin{pgfscope}%
\pgfsys@transformshift{4.010272in}{0.625278in}%
\pgfsys@useobject{currentmarker}{}%
\end{pgfscope}%
\begin{pgfscope}%
\pgfsys@transformshift{4.010436in}{0.625278in}%
\pgfsys@useobject{currentmarker}{}%
\end{pgfscope}%
\begin{pgfscope}%
\pgfsys@transformshift{4.010607in}{0.625278in}%
\pgfsys@useobject{currentmarker}{}%
\end{pgfscope}%
\begin{pgfscope}%
\pgfsys@transformshift{4.010779in}{0.625278in}%
\pgfsys@useobject{currentmarker}{}%
\end{pgfscope}%
\begin{pgfscope}%
\pgfsys@transformshift{4.010943in}{0.625278in}%
\pgfsys@useobject{currentmarker}{}%
\end{pgfscope}%
\begin{pgfscope}%
\pgfsys@transformshift{4.010943in}{0.625278in}%
\pgfsys@useobject{currentmarker}{}%
\end{pgfscope}%
\begin{pgfscope}%
\pgfsys@transformshift{4.011114in}{0.625278in}%
\pgfsys@useobject{currentmarker}{}%
\end{pgfscope}%
\begin{pgfscope}%
\pgfsys@transformshift{4.011285in}{0.625278in}%
\pgfsys@useobject{currentmarker}{}%
\end{pgfscope}%
\begin{pgfscope}%
\pgfsys@transformshift{4.011620in}{0.625278in}%
\pgfsys@useobject{currentmarker}{}%
\end{pgfscope}%
\begin{pgfscope}%
\pgfsys@transformshift{4.011956in}{0.625278in}%
\pgfsys@useobject{currentmarker}{}%
\end{pgfscope}%
\begin{pgfscope}%
\pgfsys@transformshift{4.011956in}{0.625278in}%
\pgfsys@useobject{currentmarker}{}%
\end{pgfscope}%
\begin{pgfscope}%
\pgfsys@transformshift{4.012127in}{0.625278in}%
\pgfsys@useobject{currentmarker}{}%
\end{pgfscope}%
\begin{pgfscope}%
\pgfsys@transformshift{4.014145in}{0.625278in}%
\pgfsys@useobject{currentmarker}{}%
\end{pgfscope}%
\begin{pgfscope}%
\pgfsys@transformshift{4.033213in}{0.625278in}%
\pgfsys@useobject{currentmarker}{}%
\end{pgfscope}%
\begin{pgfscope}%
\pgfsys@transformshift{4.052601in}{0.625278in}%
\pgfsys@useobject{currentmarker}{}%
\end{pgfscope}%
\begin{pgfscope}%
\pgfsys@transformshift{4.070804in}{0.625278in}%
\pgfsys@useobject{currentmarker}{}%
\end{pgfscope}%
\begin{pgfscope}%
\pgfsys@transformshift{4.081090in}{0.625278in}%
\pgfsys@useobject{currentmarker}{}%
\end{pgfscope}%
\begin{pgfscope}%
\pgfsys@transformshift{4.093909in}{0.625278in}%
\pgfsys@useobject{currentmarker}{}%
\end{pgfscope}%
\begin{pgfscope}%
\pgfsys@transformshift{4.102176in}{0.625278in}%
\pgfsys@useobject{currentmarker}{}%
\end{pgfscope}%
\begin{pgfscope}%
\pgfsys@transformshift{4.106221in}{0.625278in}%
\pgfsys@useobject{currentmarker}{}%
\end{pgfscope}%
\begin{pgfscope}%
\pgfsys@transformshift{4.106384in}{0.625278in}%
\pgfsys@useobject{currentmarker}{}%
\end{pgfscope}%
\begin{pgfscope}%
\pgfsys@transformshift{4.106392in}{0.625278in}%
\pgfsys@useobject{currentmarker}{}%
\end{pgfscope}%
\begin{pgfscope}%
\pgfsys@transformshift{4.106727in}{0.625278in}%
\pgfsys@useobject{currentmarker}{}%
\end{pgfscope}%
\begin{pgfscope}%
\pgfsys@transformshift{4.106727in}{0.625278in}%
\pgfsys@useobject{currentmarker}{}%
\end{pgfscope}%
\begin{pgfscope}%
\pgfsys@transformshift{4.106735in}{0.625278in}%
\pgfsys@useobject{currentmarker}{}%
\end{pgfscope}%
\begin{pgfscope}%
\pgfsys@transformshift{4.107062in}{0.625278in}%
\pgfsys@useobject{currentmarker}{}%
\end{pgfscope}%
\begin{pgfscope}%
\pgfsys@transformshift{4.107062in}{0.625278in}%
\pgfsys@useobject{currentmarker}{}%
\end{pgfscope}%
\begin{pgfscope}%
\pgfsys@transformshift{4.107070in}{0.625278in}%
\pgfsys@useobject{currentmarker}{}%
\end{pgfscope}%
\begin{pgfscope}%
\pgfsys@transformshift{4.107397in}{0.625278in}%
\pgfsys@useobject{currentmarker}{}%
\end{pgfscope}%
\begin{pgfscope}%
\pgfsys@transformshift{4.107397in}{0.625278in}%
\pgfsys@useobject{currentmarker}{}%
\end{pgfscope}%
\begin{pgfscope}%
\pgfsys@transformshift{4.107569in}{0.625278in}%
\pgfsys@useobject{currentmarker}{}%
\end{pgfscope}%
\begin{pgfscope}%
\pgfsys@transformshift{4.107740in}{0.625278in}%
\pgfsys@useobject{currentmarker}{}%
\end{pgfscope}%
\begin{pgfscope}%
\pgfsys@transformshift{4.107740in}{0.625278in}%
\pgfsys@useobject{currentmarker}{}%
\end{pgfscope}%
\begin{pgfscope}%
\pgfsys@transformshift{4.107904in}{0.625278in}%
\pgfsys@useobject{currentmarker}{}%
\end{pgfscope}%
\begin{pgfscope}%
\pgfsys@transformshift{4.108246in}{0.625278in}%
\pgfsys@useobject{currentmarker}{}%
\end{pgfscope}%
\begin{pgfscope}%
\pgfsys@transformshift{4.108410in}{0.625278in}%
\pgfsys@useobject{currentmarker}{}%
\end{pgfscope}%
\begin{pgfscope}%
\pgfsys@transformshift{4.108582in}{0.625278in}%
\pgfsys@useobject{currentmarker}{}%
\end{pgfscope}%
\begin{pgfscope}%
\pgfsys@transformshift{4.108746in}{0.625278in}%
\pgfsys@useobject{currentmarker}{}%
\end{pgfscope}%
\begin{pgfscope}%
\pgfsys@transformshift{4.108917in}{0.625278in}%
\pgfsys@useobject{currentmarker}{}%
\end{pgfscope}%
\begin{pgfscope}%
\pgfsys@transformshift{4.109252in}{0.625278in}%
\pgfsys@useobject{currentmarker}{}%
\end{pgfscope}%
\begin{pgfscope}%
\pgfsys@transformshift{4.109423in}{0.625278in}%
\pgfsys@useobject{currentmarker}{}%
\end{pgfscope}%
\begin{pgfscope}%
\pgfsys@transformshift{4.109595in}{0.625278in}%
\pgfsys@useobject{currentmarker}{}%
\end{pgfscope}%
\begin{pgfscope}%
\pgfsys@transformshift{4.109937in}{0.625278in}%
\pgfsys@useobject{currentmarker}{}%
\end{pgfscope}%
\begin{pgfscope}%
\pgfsys@transformshift{4.258493in}{0.625278in}%
\pgfsys@useobject{currentmarker}{}%
\end{pgfscope}%
\begin{pgfscope}%
\pgfsys@transformshift{4.267430in}{0.625278in}%
\pgfsys@useobject{currentmarker}{}%
\end{pgfscope}%
\begin{pgfscope}%
\pgfsys@transformshift{4.267430in}{0.625278in}%
\pgfsys@useobject{currentmarker}{}%
\end{pgfscope}%
\begin{pgfscope}%
\pgfsys@transformshift{4.267430in}{0.625278in}%
\pgfsys@useobject{currentmarker}{}%
\end{pgfscope}%
\begin{pgfscope}%
\pgfsys@transformshift{4.267766in}{0.625278in}%
\pgfsys@useobject{currentmarker}{}%
\end{pgfscope}%
\begin{pgfscope}%
\pgfsys@transformshift{4.267766in}{0.625278in}%
\pgfsys@useobject{currentmarker}{}%
\end{pgfscope}%
\begin{pgfscope}%
\pgfsys@transformshift{4.267937in}{0.625278in}%
\pgfsys@useobject{currentmarker}{}%
\end{pgfscope}%
\begin{pgfscope}%
\pgfsys@transformshift{4.268101in}{0.625278in}%
\pgfsys@useobject{currentmarker}{}%
\end{pgfscope}%
\begin{pgfscope}%
\pgfsys@transformshift{4.268272in}{0.625278in}%
\pgfsys@useobject{currentmarker}{}%
\end{pgfscope}%
\begin{pgfscope}%
\pgfsys@transformshift{4.268443in}{0.625278in}%
\pgfsys@useobject{currentmarker}{}%
\end{pgfscope}%
\begin{pgfscope}%
\pgfsys@transformshift{4.268779in}{0.625278in}%
\pgfsys@useobject{currentmarker}{}%
\end{pgfscope}%
\begin{pgfscope}%
\pgfsys@transformshift{4.269114in}{0.625278in}%
\pgfsys@useobject{currentmarker}{}%
\end{pgfscope}%
\begin{pgfscope}%
\pgfsys@transformshift{4.269285in}{0.625278in}%
\pgfsys@useobject{currentmarker}{}%
\end{pgfscope}%
\begin{pgfscope}%
\pgfsys@transformshift{4.269449in}{0.625278in}%
\pgfsys@useobject{currentmarker}{}%
\end{pgfscope}%
\begin{pgfscope}%
\pgfsys@transformshift{4.269449in}{0.625278in}%
\pgfsys@useobject{currentmarker}{}%
\end{pgfscope}%
\begin{pgfscope}%
\pgfsys@transformshift{4.269620in}{0.625278in}%
\pgfsys@useobject{currentmarker}{}%
\end{pgfscope}%
\begin{pgfscope}%
\pgfsys@transformshift{4.269620in}{0.625278in}%
\pgfsys@useobject{currentmarker}{}%
\end{pgfscope}%
\begin{pgfscope}%
\pgfsys@transformshift{4.269792in}{0.625278in}%
\pgfsys@useobject{currentmarker}{}%
\end{pgfscope}%
\begin{pgfscope}%
\pgfsys@transformshift{4.269792in}{0.625278in}%
\pgfsys@useobject{currentmarker}{}%
\end{pgfscope}%
\begin{pgfscope}%
\pgfsys@transformshift{4.269955in}{0.625278in}%
\pgfsys@useobject{currentmarker}{}%
\end{pgfscope}%
\begin{pgfscope}%
\pgfsys@transformshift{4.270127in}{0.625278in}%
\pgfsys@useobject{currentmarker}{}%
\end{pgfscope}%
\begin{pgfscope}%
\pgfsys@transformshift{4.270462in}{0.625278in}%
\pgfsys@useobject{currentmarker}{}%
\end{pgfscope}%
\begin{pgfscope}%
\pgfsys@transformshift{4.270693in}{0.625278in}%
\pgfsys@useobject{currentmarker}{}%
\end{pgfscope}%
\begin{pgfscope}%
\pgfsys@transformshift{4.270797in}{0.625278in}%
\pgfsys@useobject{currentmarker}{}%
\end{pgfscope}%
\begin{pgfscope}%
\pgfsys@transformshift{4.270797in}{0.625278in}%
\pgfsys@useobject{currentmarker}{}%
\end{pgfscope}%
\begin{pgfscope}%
\pgfsys@transformshift{4.270968in}{0.625278in}%
\pgfsys@useobject{currentmarker}{}%
\end{pgfscope}%
\begin{pgfscope}%
\pgfsys@transformshift{4.271140in}{0.625278in}%
\pgfsys@useobject{currentmarker}{}%
\end{pgfscope}%
\begin{pgfscope}%
\pgfsys@transformshift{4.271140in}{0.625278in}%
\pgfsys@useobject{currentmarker}{}%
\end{pgfscope}%
\begin{pgfscope}%
\pgfsys@transformshift{4.271475in}{0.625278in}%
\pgfsys@useobject{currentmarker}{}%
\end{pgfscope}%
\begin{pgfscope}%
\pgfsys@transformshift{4.284457in}{0.625278in}%
\pgfsys@useobject{currentmarker}{}%
\end{pgfscope}%
\begin{pgfscope}%
\pgfsys@transformshift{4.295592in}{0.625278in}%
\pgfsys@useobject{currentmarker}{}%
\end{pgfscope}%
\begin{pgfscope}%
\pgfsys@transformshift{4.295600in}{0.625278in}%
\pgfsys@useobject{currentmarker}{}%
\end{pgfscope}%
\begin{pgfscope}%
\pgfsys@transformshift{4.295764in}{0.625278in}%
\pgfsys@useobject{currentmarker}{}%
\end{pgfscope}%
\begin{pgfscope}%
\pgfsys@transformshift{4.295927in}{0.625278in}%
\pgfsys@useobject{currentmarker}{}%
\end{pgfscope}%
\begin{pgfscope}%
\pgfsys@transformshift{4.296099in}{0.625278in}%
\pgfsys@useobject{currentmarker}{}%
\end{pgfscope}%
\begin{pgfscope}%
\pgfsys@transformshift{4.296099in}{0.625278in}%
\pgfsys@useobject{currentmarker}{}%
\end{pgfscope}%
\begin{pgfscope}%
\pgfsys@transformshift{4.296263in}{0.625278in}%
\pgfsys@useobject{currentmarker}{}%
\end{pgfscope}%
\begin{pgfscope}%
\pgfsys@transformshift{4.296277in}{0.625278in}%
\pgfsys@useobject{currentmarker}{}%
\end{pgfscope}%
\begin{pgfscope}%
\pgfsys@transformshift{4.296441in}{0.625278in}%
\pgfsys@useobject{currentmarker}{}%
\end{pgfscope}%
\begin{pgfscope}%
\pgfsys@transformshift{4.296598in}{0.625278in}%
\pgfsys@useobject{currentmarker}{}%
\end{pgfscope}%
\begin{pgfscope}%
\pgfsys@transformshift{4.296784in}{0.625278in}%
\pgfsys@useobject{currentmarker}{}%
\end{pgfscope}%
\begin{pgfscope}%
\pgfsys@transformshift{4.296940in}{0.625278in}%
\pgfsys@useobject{currentmarker}{}%
\end{pgfscope}%
\begin{pgfscope}%
\pgfsys@transformshift{4.296948in}{0.625278in}%
\pgfsys@useobject{currentmarker}{}%
\end{pgfscope}%
\begin{pgfscope}%
\pgfsys@transformshift{4.297112in}{0.625278in}%
\pgfsys@useobject{currentmarker}{}%
\end{pgfscope}%
\begin{pgfscope}%
\pgfsys@transformshift{4.297283in}{0.625278in}%
\pgfsys@useobject{currentmarker}{}%
\end{pgfscope}%
\begin{pgfscope}%
\pgfsys@transformshift{4.297454in}{0.625278in}%
\pgfsys@useobject{currentmarker}{}%
\end{pgfscope}%
\begin{pgfscope}%
\pgfsys@transformshift{4.297789in}{0.625278in}%
\pgfsys@useobject{currentmarker}{}%
\end{pgfscope}%
\begin{pgfscope}%
\pgfsys@transformshift{4.297789in}{0.625278in}%
\pgfsys@useobject{currentmarker}{}%
\end{pgfscope}%
\begin{pgfscope}%
\pgfsys@transformshift{4.298125in}{0.625278in}%
\pgfsys@useobject{currentmarker}{}%
\end{pgfscope}%
\begin{pgfscope}%
\pgfsys@transformshift{4.298125in}{0.625278in}%
\pgfsys@useobject{currentmarker}{}%
\end{pgfscope}%
\begin{pgfscope}%
\pgfsys@transformshift{4.312626in}{0.625278in}%
\pgfsys@useobject{currentmarker}{}%
\end{pgfscope}%
\begin{pgfscope}%
\pgfsys@transformshift{4.316000in}{0.625278in}%
\pgfsys@useobject{currentmarker}{}%
\end{pgfscope}%
\begin{pgfscope}%
\pgfsys@transformshift{4.316164in}{0.625278in}%
\pgfsys@useobject{currentmarker}{}%
\end{pgfscope}%
\begin{pgfscope}%
\pgfsys@transformshift{4.316336in}{0.625278in}%
\pgfsys@useobject{currentmarker}{}%
\end{pgfscope}%
\begin{pgfscope}%
\pgfsys@transformshift{4.316499in}{0.625278in}%
\pgfsys@useobject{currentmarker}{}%
\end{pgfscope}%
\begin{pgfscope}%
\pgfsys@transformshift{4.316507in}{0.625278in}%
\pgfsys@useobject{currentmarker}{}%
\end{pgfscope}%
\begin{pgfscope}%
\pgfsys@transformshift{4.316678in}{0.625278in}%
\pgfsys@useobject{currentmarker}{}%
\end{pgfscope}%
\begin{pgfscope}%
\pgfsys@transformshift{4.316835in}{0.625278in}%
\pgfsys@useobject{currentmarker}{}%
\end{pgfscope}%
\begin{pgfscope}%
\pgfsys@transformshift{4.317013in}{0.625278in}%
\pgfsys@useobject{currentmarker}{}%
\end{pgfscope}%
\begin{pgfscope}%
\pgfsys@transformshift{4.317013in}{0.625278in}%
\pgfsys@useobject{currentmarker}{}%
\end{pgfscope}%
\begin{pgfscope}%
\pgfsys@transformshift{4.320551in}{0.625278in}%
\pgfsys@useobject{currentmarker}{}%
\end{pgfscope}%
\begin{pgfscope}%
\pgfsys@transformshift{4.512954in}{0.625278in}%
\pgfsys@useobject{currentmarker}{}%
\end{pgfscope}%
\begin{pgfscope}%
\pgfsys@transformshift{4.527456in}{0.625278in}%
\pgfsys@useobject{currentmarker}{}%
\end{pgfscope}%
\begin{pgfscope}%
\pgfsys@transformshift{4.527955in}{0.625278in}%
\pgfsys@useobject{currentmarker}{}%
\end{pgfscope}%
\begin{pgfscope}%
\pgfsys@transformshift{4.528126in}{0.625278in}%
\pgfsys@useobject{currentmarker}{}%
\end{pgfscope}%
\begin{pgfscope}%
\pgfsys@transformshift{4.528298in}{0.625278in}%
\pgfsys@useobject{currentmarker}{}%
\end{pgfscope}%
\begin{pgfscope}%
\pgfsys@transformshift{4.528462in}{0.625278in}%
\pgfsys@useobject{currentmarker}{}%
\end{pgfscope}%
\begin{pgfscope}%
\pgfsys@transformshift{4.528633in}{0.625278in}%
\pgfsys@useobject{currentmarker}{}%
\end{pgfscope}%
\begin{pgfscope}%
\pgfsys@transformshift{4.528804in}{0.625278in}%
\pgfsys@useobject{currentmarker}{}%
\end{pgfscope}%
\begin{pgfscope}%
\pgfsys@transformshift{4.528968in}{0.625278in}%
\pgfsys@useobject{currentmarker}{}%
\end{pgfscope}%
\begin{pgfscope}%
\pgfsys@transformshift{4.529139in}{0.625278in}%
\pgfsys@useobject{currentmarker}{}%
\end{pgfscope}%
\begin{pgfscope}%
\pgfsys@transformshift{4.529311in}{0.625278in}%
\pgfsys@useobject{currentmarker}{}%
\end{pgfscope}%
\begin{pgfscope}%
\pgfsys@transformshift{4.529475in}{0.625278in}%
\pgfsys@useobject{currentmarker}{}%
\end{pgfscope}%
\begin{pgfscope}%
\pgfsys@transformshift{4.529475in}{0.625278in}%
\pgfsys@useobject{currentmarker}{}%
\end{pgfscope}%
\begin{pgfscope}%
\pgfsys@transformshift{4.529646in}{0.625278in}%
\pgfsys@useobject{currentmarker}{}%
\end{pgfscope}%
\begin{pgfscope}%
\pgfsys@transformshift{4.529817in}{0.625278in}%
\pgfsys@useobject{currentmarker}{}%
\end{pgfscope}%
\begin{pgfscope}%
\pgfsys@transformshift{4.529981in}{0.625278in}%
\pgfsys@useobject{currentmarker}{}%
\end{pgfscope}%
\begin{pgfscope}%
\pgfsys@transformshift{4.530316in}{0.625278in}%
\pgfsys@useobject{currentmarker}{}%
\end{pgfscope}%
\begin{pgfscope}%
\pgfsys@transformshift{4.530316in}{0.625278in}%
\pgfsys@useobject{currentmarker}{}%
\end{pgfscope}%
\begin{pgfscope}%
\pgfsys@transformshift{4.530488in}{0.625278in}%
\pgfsys@useobject{currentmarker}{}%
\end{pgfscope}%
\begin{pgfscope}%
\pgfsys@transformshift{4.530659in}{0.625278in}%
\pgfsys@useobject{currentmarker}{}%
\end{pgfscope}%
\begin{pgfscope}%
\pgfsys@transformshift{4.530823in}{0.625278in}%
\pgfsys@useobject{currentmarker}{}%
\end{pgfscope}%
\begin{pgfscope}%
\pgfsys@transformshift{4.531165in}{0.625278in}%
\pgfsys@useobject{currentmarker}{}%
\end{pgfscope}%
\begin{pgfscope}%
\pgfsys@transformshift{4.531165in}{0.625278in}%
\pgfsys@useobject{currentmarker}{}%
\end{pgfscope}%
\begin{pgfscope}%
\pgfsys@transformshift{4.555119in}{0.625278in}%
\pgfsys@useobject{currentmarker}{}%
\end{pgfscope}%
\begin{pgfscope}%
\pgfsys@transformshift{4.568109in}{0.625278in}%
\pgfsys@useobject{currentmarker}{}%
\end{pgfscope}%
\begin{pgfscope}%
\pgfsys@transformshift{4.581583in}{0.625278in}%
\pgfsys@useobject{currentmarker}{}%
\end{pgfscope}%
\begin{pgfscope}%
\pgfsys@transformshift{4.581754in}{0.625278in}%
\pgfsys@useobject{currentmarker}{}%
\end{pgfscope}%
\begin{pgfscope}%
\pgfsys@transformshift{4.581918in}{0.625278in}%
\pgfsys@useobject{currentmarker}{}%
\end{pgfscope}%
\begin{pgfscope}%
\pgfsys@transformshift{4.582089in}{0.625278in}%
\pgfsys@useobject{currentmarker}{}%
\end{pgfscope}%
\begin{pgfscope}%
\pgfsys@transformshift{4.582096in}{0.625278in}%
\pgfsys@useobject{currentmarker}{}%
\end{pgfscope}%
\begin{pgfscope}%
\pgfsys@transformshift{4.582260in}{0.625278in}%
\pgfsys@useobject{currentmarker}{}%
\end{pgfscope}%
\begin{pgfscope}%
\pgfsys@transformshift{4.582595in}{0.625278in}%
\pgfsys@useobject{currentmarker}{}%
\end{pgfscope}%
\begin{pgfscope}%
\pgfsys@transformshift{4.582603in}{0.625278in}%
\pgfsys@useobject{currentmarker}{}%
\end{pgfscope}%
\begin{pgfscope}%
\pgfsys@transformshift{4.582938in}{0.625278in}%
\pgfsys@useobject{currentmarker}{}%
\end{pgfscope}%
\begin{pgfscope}%
\pgfsys@transformshift{4.583102in}{0.625278in}%
\pgfsys@useobject{currentmarker}{}%
\end{pgfscope}%
\begin{pgfscope}%
\pgfsys@transformshift{4.583273in}{0.625278in}%
\pgfsys@useobject{currentmarker}{}%
\end{pgfscope}%
\begin{pgfscope}%
\pgfsys@transformshift{4.583437in}{0.625278in}%
\pgfsys@useobject{currentmarker}{}%
\end{pgfscope}%
\begin{pgfscope}%
\pgfsys@transformshift{4.583608in}{0.625278in}%
\pgfsys@useobject{currentmarker}{}%
\end{pgfscope}%
\begin{pgfscope}%
\pgfsys@transformshift{4.583616in}{0.625278in}%
\pgfsys@useobject{currentmarker}{}%
\end{pgfscope}%
\begin{pgfscope}%
\pgfsys@transformshift{4.583944in}{0.625278in}%
\pgfsys@useobject{currentmarker}{}%
\end{pgfscope}%
\begin{pgfscope}%
\pgfsys@transformshift{4.584279in}{0.625278in}%
\pgfsys@useobject{currentmarker}{}%
\end{pgfscope}%
\begin{pgfscope}%
\pgfsys@transformshift{4.617498in}{0.625278in}%
\pgfsys@useobject{currentmarker}{}%
\end{pgfscope}%
\begin{pgfscope}%
\pgfsys@transformshift{4.625937in}{0.625278in}%
\pgfsys@useobject{currentmarker}{}%
\end{pgfscope}%
\begin{pgfscope}%
\pgfsys@transformshift{4.628462in}{0.625278in}%
\pgfsys@useobject{currentmarker}{}%
\end{pgfscope}%
\begin{pgfscope}%
\pgfsys@transformshift{4.628797in}{0.625278in}%
\pgfsys@useobject{currentmarker}{}%
\end{pgfscope}%
\begin{pgfscope}%
\pgfsys@transformshift{4.628797in}{0.625278in}%
\pgfsys@useobject{currentmarker}{}%
\end{pgfscope}%
\begin{pgfscope}%
\pgfsys@transformshift{4.628968in}{0.625278in}%
\pgfsys@useobject{currentmarker}{}%
\end{pgfscope}%
\begin{pgfscope}%
\pgfsys@transformshift{4.629132in}{0.625278in}%
\pgfsys@useobject{currentmarker}{}%
\end{pgfscope}%
\begin{pgfscope}%
\pgfsys@transformshift{4.629132in}{0.625278in}%
\pgfsys@useobject{currentmarker}{}%
\end{pgfscope}%
\begin{pgfscope}%
\pgfsys@transformshift{4.629475in}{0.625278in}%
\pgfsys@useobject{currentmarker}{}%
\end{pgfscope}%
\begin{pgfscope}%
\pgfsys@transformshift{4.629810in}{0.625278in}%
\pgfsys@useobject{currentmarker}{}%
\end{pgfscope}%
\begin{pgfscope}%
\pgfsys@transformshift{4.630145in}{0.625278in}%
\pgfsys@useobject{currentmarker}{}%
\end{pgfscope}%
\begin{pgfscope}%
\pgfsys@transformshift{4.630145in}{0.625278in}%
\pgfsys@useobject{currentmarker}{}%
\end{pgfscope}%
\begin{pgfscope}%
\pgfsys@transformshift{4.630488in}{0.625278in}%
\pgfsys@useobject{currentmarker}{}%
\end{pgfscope}%
\begin{pgfscope}%
\pgfsys@transformshift{4.630488in}{0.625278in}%
\pgfsys@useobject{currentmarker}{}%
\end{pgfscope}%
\begin{pgfscope}%
\pgfsys@transformshift{4.630652in}{0.625278in}%
\pgfsys@useobject{currentmarker}{}%
\end{pgfscope}%
\begin{pgfscope}%
\pgfsys@transformshift{4.630823in}{0.625278in}%
\pgfsys@useobject{currentmarker}{}%
\end{pgfscope}%
\begin{pgfscope}%
\pgfsys@transformshift{4.630823in}{0.625278in}%
\pgfsys@useobject{currentmarker}{}%
\end{pgfscope}%
\begin{pgfscope}%
\pgfsys@transformshift{4.630882in}{0.625278in}%
\pgfsys@useobject{currentmarker}{}%
\end{pgfscope}%
\begin{pgfscope}%
\pgfsys@transformshift{4.630994in}{0.625278in}%
\pgfsys@useobject{currentmarker}{}%
\end{pgfscope}%
\begin{pgfscope}%
\pgfsys@transformshift{4.631329in}{0.625278in}%
\pgfsys@useobject{currentmarker}{}%
\end{pgfscope}%
\begin{pgfscope}%
\pgfsys@transformshift{4.631337in}{0.625278in}%
\pgfsys@useobject{currentmarker}{}%
\end{pgfscope}%
\begin{pgfscope}%
\pgfsys@transformshift{4.631665in}{0.625278in}%
\pgfsys@useobject{currentmarker}{}%
\end{pgfscope}%
\begin{pgfscope}%
\pgfsys@transformshift{4.631672in}{0.625278in}%
\pgfsys@useobject{currentmarker}{}%
\end{pgfscope}%
\begin{pgfscope}%
\pgfsys@transformshift{4.668600in}{0.625278in}%
\pgfsys@useobject{currentmarker}{}%
\end{pgfscope}%
\begin{pgfscope}%
\pgfsys@transformshift{4.783765in}{0.625278in}%
\pgfsys@useobject{currentmarker}{}%
\end{pgfscope}%
\begin{pgfscope}%
\pgfsys@transformshift{4.801641in}{0.625278in}%
\pgfsys@useobject{currentmarker}{}%
\end{pgfscope}%
\begin{pgfscope}%
\pgfsys@transformshift{4.807376in}{0.625278in}%
\pgfsys@useobject{currentmarker}{}%
\end{pgfscope}%
\begin{pgfscope}%
\pgfsys@transformshift{4.807547in}{0.625278in}%
\pgfsys@useobject{currentmarker}{}%
\end{pgfscope}%
\begin{pgfscope}%
\pgfsys@transformshift{4.807555in}{0.625278in}%
\pgfsys@useobject{currentmarker}{}%
\end{pgfscope}%
\begin{pgfscope}%
\pgfsys@transformshift{4.807883in}{0.625278in}%
\pgfsys@useobject{currentmarker}{}%
\end{pgfscope}%
\begin{pgfscope}%
\pgfsys@transformshift{4.807883in}{0.625278in}%
\pgfsys@useobject{currentmarker}{}%
\end{pgfscope}%
\begin{pgfscope}%
\pgfsys@transformshift{4.807890in}{0.625278in}%
\pgfsys@useobject{currentmarker}{}%
\end{pgfscope}%
\begin{pgfscope}%
\pgfsys@transformshift{4.808218in}{0.625278in}%
\pgfsys@useobject{currentmarker}{}%
\end{pgfscope}%
\begin{pgfscope}%
\pgfsys@transformshift{4.808225in}{0.625278in}%
\pgfsys@useobject{currentmarker}{}%
\end{pgfscope}%
\begin{pgfscope}%
\pgfsys@transformshift{4.808389in}{0.625278in}%
\pgfsys@useobject{currentmarker}{}%
\end{pgfscope}%
\begin{pgfscope}%
\pgfsys@transformshift{4.808553in}{0.625278in}%
\pgfsys@useobject{currentmarker}{}%
\end{pgfscope}%
\begin{pgfscope}%
\pgfsys@transformshift{4.808724in}{0.625278in}%
\pgfsys@useobject{currentmarker}{}%
\end{pgfscope}%
\begin{pgfscope}%
\pgfsys@transformshift{4.808724in}{0.625278in}%
\pgfsys@useobject{currentmarker}{}%
\end{pgfscope}%
\begin{pgfscope}%
\pgfsys@transformshift{4.808896in}{0.625278in}%
\pgfsys@useobject{currentmarker}{}%
\end{pgfscope}%
\begin{pgfscope}%
\pgfsys@transformshift{4.809059in}{0.625278in}%
\pgfsys@useobject{currentmarker}{}%
\end{pgfscope}%
\begin{pgfscope}%
\pgfsys@transformshift{4.809566in}{0.625278in}%
\pgfsys@useobject{currentmarker}{}%
\end{pgfscope}%
\begin{pgfscope}%
\pgfsys@transformshift{4.809737in}{0.625278in}%
\pgfsys@useobject{currentmarker}{}%
\end{pgfscope}%
\begin{pgfscope}%
\pgfsys@transformshift{4.809909in}{0.625278in}%
\pgfsys@useobject{currentmarker}{}%
\end{pgfscope}%
\begin{pgfscope}%
\pgfsys@transformshift{4.810072in}{0.625278in}%
\pgfsys@useobject{currentmarker}{}%
\end{pgfscope}%
\begin{pgfscope}%
\pgfsys@transformshift{4.810244in}{0.625278in}%
\pgfsys@useobject{currentmarker}{}%
\end{pgfscope}%
\begin{pgfscope}%
\pgfsys@transformshift{4.810244in}{0.625278in}%
\pgfsys@useobject{currentmarker}{}%
\end{pgfscope}%
\begin{pgfscope}%
\pgfsys@transformshift{4.810579in}{0.625278in}%
\pgfsys@useobject{currentmarker}{}%
\end{pgfscope}%
\begin{pgfscope}%
\pgfsys@transformshift{4.835210in}{0.625278in}%
\pgfsys@useobject{currentmarker}{}%
\end{pgfscope}%
\begin{pgfscope}%
\pgfsys@transformshift{4.846666in}{0.625278in}%
\pgfsys@useobject{currentmarker}{}%
\end{pgfscope}%
\begin{pgfscope}%
\pgfsys@transformshift{4.883594in}{0.625278in}%
\pgfsys@useobject{currentmarker}{}%
\end{pgfscope}%
\begin{pgfscope}%
\pgfsys@transformshift{4.901298in}{0.625278in}%
\pgfsys@useobject{currentmarker}{}%
\end{pgfscope}%
\begin{pgfscope}%
\pgfsys@transformshift{4.910743in}{0.625278in}%
\pgfsys@useobject{currentmarker}{}%
\end{pgfscope}%
\begin{pgfscope}%
\pgfsys@transformshift{4.919852in}{0.625278in}%
\pgfsys@useobject{currentmarker}{}%
\end{pgfscope}%
\begin{pgfscope}%
\pgfsys@transformshift{4.922548in}{0.625278in}%
\pgfsys@useobject{currentmarker}{}%
\end{pgfscope}%
\begin{pgfscope}%
\pgfsys@transformshift{4.922720in}{0.625278in}%
\pgfsys@useobject{currentmarker}{}%
\end{pgfscope}%
\begin{pgfscope}%
\pgfsys@transformshift{4.922884in}{0.625278in}%
\pgfsys@useobject{currentmarker}{}%
\end{pgfscope}%
\begin{pgfscope}%
\pgfsys@transformshift{4.923055in}{0.625278in}%
\pgfsys@useobject{currentmarker}{}%
\end{pgfscope}%
\begin{pgfscope}%
\pgfsys@transformshift{4.923226in}{0.625278in}%
\pgfsys@useobject{currentmarker}{}%
\end{pgfscope}%
\begin{pgfscope}%
\pgfsys@transformshift{4.923390in}{0.625278in}%
\pgfsys@useobject{currentmarker}{}%
\end{pgfscope}%
\begin{pgfscope}%
\pgfsys@transformshift{4.923561in}{0.625278in}%
\pgfsys@useobject{currentmarker}{}%
\end{pgfscope}%
\begin{pgfscope}%
\pgfsys@transformshift{4.923733in}{0.625278in}%
\pgfsys@useobject{currentmarker}{}%
\end{pgfscope}%
\begin{pgfscope}%
\pgfsys@transformshift{4.923896in}{0.625278in}%
\pgfsys@useobject{currentmarker}{}%
\end{pgfscope}%
\begin{pgfscope}%
\pgfsys@transformshift{4.924068in}{0.625278in}%
\pgfsys@useobject{currentmarker}{}%
\end{pgfscope}%
\begin{pgfscope}%
\pgfsys@transformshift{4.924232in}{0.625278in}%
\pgfsys@useobject{currentmarker}{}%
\end{pgfscope}%
\begin{pgfscope}%
\pgfsys@transformshift{4.924239in}{0.625278in}%
\pgfsys@useobject{currentmarker}{}%
\end{pgfscope}%
\begin{pgfscope}%
\pgfsys@transformshift{4.924574in}{0.625278in}%
\pgfsys@useobject{currentmarker}{}%
\end{pgfscope}%
\begin{pgfscope}%
\pgfsys@transformshift{4.924574in}{0.625278in}%
\pgfsys@useobject{currentmarker}{}%
\end{pgfscope}%
\begin{pgfscope}%
\pgfsys@transformshift{4.924746in}{0.625278in}%
\pgfsys@useobject{currentmarker}{}%
\end{pgfscope}%
\begin{pgfscope}%
\pgfsys@transformshift{4.924909in}{0.625278in}%
\pgfsys@useobject{currentmarker}{}%
\end{pgfscope}%
\begin{pgfscope}%
\pgfsys@transformshift{4.928283in}{0.625278in}%
\pgfsys@useobject{currentmarker}{}%
\end{pgfscope}%
\begin{pgfscope}%
\pgfsys@transformshift{4.939925in}{0.625278in}%
\pgfsys@useobject{currentmarker}{}%
\end{pgfscope}%
\begin{pgfscope}%
\pgfsys@transformshift{4.955939in}{0.625278in}%
\pgfsys@useobject{currentmarker}{}%
\end{pgfscope}%
\begin{pgfscope}%
\pgfsys@transformshift{4.968593in}{0.625278in}%
\pgfsys@useobject{currentmarker}{}%
\end{pgfscope}%
\begin{pgfscope}%
\pgfsys@transformshift{5.068243in}{0.625278in}%
\pgfsys@useobject{currentmarker}{}%
\end{pgfscope}%
\begin{pgfscope}%
\pgfsys@transformshift{5.087803in}{0.625278in}%
\pgfsys@useobject{currentmarker}{}%
\end{pgfscope}%
\begin{pgfscope}%
\pgfsys@transformshift{5.097076in}{0.625278in}%
\pgfsys@useobject{currentmarker}{}%
\end{pgfscope}%
\begin{pgfscope}%
\pgfsys@transformshift{5.097076in}{0.625278in}%
\pgfsys@useobject{currentmarker}{}%
\end{pgfscope}%
\begin{pgfscope}%
\pgfsys@transformshift{5.097418in}{0.625278in}%
\pgfsys@useobject{currentmarker}{}%
\end{pgfscope}%
\begin{pgfscope}%
\pgfsys@transformshift{5.097418in}{0.625278in}%
\pgfsys@useobject{currentmarker}{}%
\end{pgfscope}%
\begin{pgfscope}%
\pgfsys@transformshift{5.097418in}{0.625278in}%
\pgfsys@useobject{currentmarker}{}%
\end{pgfscope}%
\begin{pgfscope}%
\pgfsys@transformshift{5.097753in}{0.625278in}%
\pgfsys@useobject{currentmarker}{}%
\end{pgfscope}%
\begin{pgfscope}%
\pgfsys@transformshift{5.097753in}{0.625278in}%
\pgfsys@useobject{currentmarker}{}%
\end{pgfscope}%
\begin{pgfscope}%
\pgfsys@transformshift{5.098089in}{0.625278in}%
\pgfsys@useobject{currentmarker}{}%
\end{pgfscope}%
\begin{pgfscope}%
\pgfsys@transformshift{5.098431in}{0.625278in}%
\pgfsys@useobject{currentmarker}{}%
\end{pgfscope}%
\begin{pgfscope}%
\pgfsys@transformshift{5.098595in}{0.625278in}%
\pgfsys@useobject{currentmarker}{}%
\end{pgfscope}%
\begin{pgfscope}%
\pgfsys@transformshift{5.098766in}{0.625278in}%
\pgfsys@useobject{currentmarker}{}%
\end{pgfscope}%
\begin{pgfscope}%
\pgfsys@transformshift{5.098938in}{0.625278in}%
\pgfsys@useobject{currentmarker}{}%
\end{pgfscope}%
\begin{pgfscope}%
\pgfsys@transformshift{5.098938in}{0.625278in}%
\pgfsys@useobject{currentmarker}{}%
\end{pgfscope}%
\begin{pgfscope}%
\pgfsys@transformshift{5.099102in}{0.625278in}%
\pgfsys@useobject{currentmarker}{}%
\end{pgfscope}%
\begin{pgfscope}%
\pgfsys@transformshift{5.099273in}{0.625278in}%
\pgfsys@useobject{currentmarker}{}%
\end{pgfscope}%
\begin{pgfscope}%
\pgfsys@transformshift{5.099608in}{0.625278in}%
\pgfsys@useobject{currentmarker}{}%
\end{pgfscope}%
\begin{pgfscope}%
\pgfsys@transformshift{5.099608in}{0.625278in}%
\pgfsys@useobject{currentmarker}{}%
\end{pgfscope}%
\begin{pgfscope}%
\pgfsys@transformshift{5.099943in}{0.625278in}%
\pgfsys@useobject{currentmarker}{}%
\end{pgfscope}%
\begin{pgfscope}%
\pgfsys@transformshift{5.099943in}{0.625278in}%
\pgfsys@useobject{currentmarker}{}%
\end{pgfscope}%
\begin{pgfscope}%
\pgfsys@transformshift{5.116642in}{0.625278in}%
\pgfsys@useobject{currentmarker}{}%
\end{pgfscope}%
\begin{pgfscope}%
\pgfsys@transformshift{5.129289in}{0.625278in}%
\pgfsys@useobject{currentmarker}{}%
\end{pgfscope}%
\begin{pgfscope}%
\pgfsys@transformshift{5.129297in}{0.625278in}%
\pgfsys@useobject{currentmarker}{}%
\end{pgfscope}%
\begin{pgfscope}%
\pgfsys@transformshift{5.129624in}{0.625278in}%
\pgfsys@useobject{currentmarker}{}%
\end{pgfscope}%
\begin{pgfscope}%
\pgfsys@transformshift{5.129632in}{0.625278in}%
\pgfsys@useobject{currentmarker}{}%
\end{pgfscope}%
\begin{pgfscope}%
\pgfsys@transformshift{5.129960in}{0.625278in}%
\pgfsys@useobject{currentmarker}{}%
\end{pgfscope}%
\begin{pgfscope}%
\pgfsys@transformshift{5.129960in}{0.625278in}%
\pgfsys@useobject{currentmarker}{}%
\end{pgfscope}%
\begin{pgfscope}%
\pgfsys@transformshift{5.129967in}{0.625278in}%
\pgfsys@useobject{currentmarker}{}%
\end{pgfscope}%
\begin{pgfscope}%
\pgfsys@transformshift{5.130146in}{0.625278in}%
\pgfsys@useobject{currentmarker}{}%
\end{pgfscope}%
\begin{pgfscope}%
\pgfsys@transformshift{5.130302in}{0.625278in}%
\pgfsys@useobject{currentmarker}{}%
\end{pgfscope}%
\begin{pgfscope}%
\pgfsys@transformshift{5.130302in}{0.625278in}%
\pgfsys@useobject{currentmarker}{}%
\end{pgfscope}%
\begin{pgfscope}%
\pgfsys@transformshift{5.130474in}{0.625278in}%
\pgfsys@useobject{currentmarker}{}%
\end{pgfscope}%
\begin{pgfscope}%
\pgfsys@transformshift{5.130637in}{0.625278in}%
\pgfsys@useobject{currentmarker}{}%
\end{pgfscope}%
\begin{pgfscope}%
\pgfsys@transformshift{5.130645in}{0.625278in}%
\pgfsys@useobject{currentmarker}{}%
\end{pgfscope}%
\begin{pgfscope}%
\pgfsys@transformshift{5.130809in}{0.625278in}%
\pgfsys@useobject{currentmarker}{}%
\end{pgfscope}%
\begin{pgfscope}%
\pgfsys@transformshift{5.131151in}{0.625278in}%
\pgfsys@useobject{currentmarker}{}%
\end{pgfscope}%
\begin{pgfscope}%
\pgfsys@transformshift{5.131308in}{0.625278in}%
\pgfsys@useobject{currentmarker}{}%
\end{pgfscope}%
\begin{pgfscope}%
\pgfsys@transformshift{5.131487in}{0.625278in}%
\pgfsys@useobject{currentmarker}{}%
\end{pgfscope}%
\begin{pgfscope}%
\pgfsys@transformshift{5.131650in}{0.625278in}%
\pgfsys@useobject{currentmarker}{}%
\end{pgfscope}%
\begin{pgfscope}%
\pgfsys@transformshift{5.131822in}{0.625278in}%
\pgfsys@useobject{currentmarker}{}%
\end{pgfscope}%
\begin{pgfscope}%
\pgfsys@transformshift{5.131986in}{0.625278in}%
\pgfsys@useobject{currentmarker}{}%
\end{pgfscope}%
\begin{pgfscope}%
\pgfsys@transformshift{5.132045in}{0.625278in}%
\pgfsys@useobject{currentmarker}{}%
\end{pgfscope}%
\begin{pgfscope}%
\pgfsys@transformshift{5.132321in}{0.625278in}%
\pgfsys@useobject{currentmarker}{}%
\end{pgfscope}%
\begin{pgfscope}%
\pgfsys@transformshift{5.132328in}{0.625278in}%
\pgfsys@useobject{currentmarker}{}%
\end{pgfscope}%
\begin{pgfscope}%
\pgfsys@transformshift{5.186961in}{0.625278in}%
\pgfsys@useobject{currentmarker}{}%
\end{pgfscope}%
\begin{pgfscope}%
\pgfsys@transformshift{5.202640in}{0.625278in}%
\pgfsys@useobject{currentmarker}{}%
\end{pgfscope}%
\begin{pgfscope}%
\pgfsys@transformshift{5.209224in}{0.625278in}%
\pgfsys@useobject{currentmarker}{}%
\end{pgfscope}%
\begin{pgfscope}%
\pgfsys@transformshift{5.209388in}{0.625278in}%
\pgfsys@useobject{currentmarker}{}%
\end{pgfscope}%
\begin{pgfscope}%
\pgfsys@transformshift{5.209552in}{0.625278in}%
\pgfsys@useobject{currentmarker}{}%
\end{pgfscope}%
\begin{pgfscope}%
\pgfsys@transformshift{5.209559in}{0.625278in}%
\pgfsys@useobject{currentmarker}{}%
\end{pgfscope}%
\begin{pgfscope}%
\pgfsys@transformshift{5.209723in}{0.625278in}%
\pgfsys@useobject{currentmarker}{}%
\end{pgfscope}%
\begin{pgfscope}%
\pgfsys@transformshift{5.209894in}{0.625278in}%
\pgfsys@useobject{currentmarker}{}%
\end{pgfscope}%
\begin{pgfscope}%
\pgfsys@transformshift{5.209902in}{0.625278in}%
\pgfsys@useobject{currentmarker}{}%
\end{pgfscope}%
\begin{pgfscope}%
\pgfsys@transformshift{5.210058in}{0.625278in}%
\pgfsys@useobject{currentmarker}{}%
\end{pgfscope}%
\begin{pgfscope}%
\pgfsys@transformshift{5.210058in}{0.625278in}%
\pgfsys@useobject{currentmarker}{}%
\end{pgfscope}%
\begin{pgfscope}%
\pgfsys@transformshift{5.210401in}{0.625278in}%
\pgfsys@useobject{currentmarker}{}%
\end{pgfscope}%
\begin{pgfscope}%
\pgfsys@transformshift{5.210401in}{0.625278in}%
\pgfsys@useobject{currentmarker}{}%
\end{pgfscope}%
\begin{pgfscope}%
\pgfsys@transformshift{5.210572in}{0.625278in}%
\pgfsys@useobject{currentmarker}{}%
\end{pgfscope}%
\begin{pgfscope}%
\pgfsys@transformshift{5.210736in}{0.625278in}%
\pgfsys@useobject{currentmarker}{}%
\end{pgfscope}%
\begin{pgfscope}%
\pgfsys@transformshift{5.210907in}{0.625278in}%
\pgfsys@useobject{currentmarker}{}%
\end{pgfscope}%
\begin{pgfscope}%
\pgfsys@transformshift{5.211071in}{0.625278in}%
\pgfsys@useobject{currentmarker}{}%
\end{pgfscope}%
\begin{pgfscope}%
\pgfsys@transformshift{5.211250in}{0.625278in}%
\pgfsys@useobject{currentmarker}{}%
\end{pgfscope}%
\begin{pgfscope}%
\pgfsys@transformshift{5.211406in}{0.625278in}%
\pgfsys@useobject{currentmarker}{}%
\end{pgfscope}%
\begin{pgfscope}%
\pgfsys@transformshift{5.211585in}{0.625278in}%
\pgfsys@useobject{currentmarker}{}%
\end{pgfscope}%
\begin{pgfscope}%
\pgfsys@transformshift{5.211749in}{0.625278in}%
\pgfsys@useobject{currentmarker}{}%
\end{pgfscope}%
\begin{pgfscope}%
\pgfsys@transformshift{5.211920in}{0.625278in}%
\pgfsys@useobject{currentmarker}{}%
\end{pgfscope}%
\begin{pgfscope}%
\pgfsys@transformshift{5.223718in}{0.625278in}%
\pgfsys@useobject{currentmarker}{}%
\end{pgfscope}%
\begin{pgfscope}%
\pgfsys@transformshift{5.238905in}{0.625278in}%
\pgfsys@useobject{currentmarker}{}%
\end{pgfscope}%
\begin{pgfscope}%
\pgfsys@transformshift{5.259477in}{0.625278in}%
\pgfsys@useobject{currentmarker}{}%
\end{pgfscope}%
\end{pgfscope}%
\begin{pgfscope}%
\pgfpathrectangle{\pgfqpoint{0.713889in}{0.502778in}}{\pgfqpoint{4.805000in}{2.695000in}}%
\pgfusepath{clip}%
\pgfsetbuttcap%
\pgfsetroundjoin%
\definecolor{currentfill}{rgb}{1.000000,0.498039,0.054902}%
\pgfsetfillcolor{currentfill}%
\pgfsetlinewidth{1.003750pt}%
\definecolor{currentstroke}{rgb}{1.000000,0.498039,0.054902}%
\pgfsetstrokecolor{currentstroke}%
\pgfsetdash{}{0pt}%
\pgfsys@defobject{currentmarker}{\pgfqpoint{-0.015528in}{-0.015528in}}{\pgfqpoint{0.015528in}{0.015528in}}{%
\pgfpathmoveto{\pgfqpoint{0.000000in}{-0.015528in}}%
\pgfpathcurveto{\pgfqpoint{0.004118in}{-0.015528in}}{\pgfqpoint{0.008068in}{-0.013892in}}{\pgfqpoint{0.010980in}{-0.010980in}}%
\pgfpathcurveto{\pgfqpoint{0.013892in}{-0.008068in}}{\pgfqpoint{0.015528in}{-0.004118in}}{\pgfqpoint{0.015528in}{0.000000in}}%
\pgfpathcurveto{\pgfqpoint{0.015528in}{0.004118in}}{\pgfqpoint{0.013892in}{0.008068in}}{\pgfqpoint{0.010980in}{0.010980in}}%
\pgfpathcurveto{\pgfqpoint{0.008068in}{0.013892in}}{\pgfqpoint{0.004118in}{0.015528in}}{\pgfqpoint{0.000000in}{0.015528in}}%
\pgfpathcurveto{\pgfqpoint{-0.004118in}{0.015528in}}{\pgfqpoint{-0.008068in}{0.013892in}}{\pgfqpoint{-0.010980in}{0.010980in}}%
\pgfpathcurveto{\pgfqpoint{-0.013892in}{0.008068in}}{\pgfqpoint{-0.015528in}{0.004118in}}{\pgfqpoint{-0.015528in}{0.000000in}}%
\pgfpathcurveto{\pgfqpoint{-0.015528in}{-0.004118in}}{\pgfqpoint{-0.013892in}{-0.008068in}}{\pgfqpoint{-0.010980in}{-0.010980in}}%
\pgfpathcurveto{\pgfqpoint{-0.008068in}{-0.013892in}}{\pgfqpoint{-0.004118in}{-0.015528in}}{\pgfqpoint{0.000000in}{-0.015528in}}%
\pgfpathlineto{\pgfqpoint{0.000000in}{-0.015528in}}%
\pgfpathclose%
\pgfusepath{stroke,fill}%
}%
\begin{pgfscope}%
\pgfsys@transformshift{0.935501in}{1.850278in}%
\pgfsys@useobject{currentmarker}{}%
\end{pgfscope}%
\begin{pgfscope}%
\pgfsys@transformshift{0.936345in}{1.850278in}%
\pgfsys@useobject{currentmarker}{}%
\end{pgfscope}%
\begin{pgfscope}%
\pgfsys@transformshift{0.936850in}{1.850278in}%
\pgfsys@useobject{currentmarker}{}%
\end{pgfscope}%
\begin{pgfscope}%
\pgfsys@transformshift{0.937188in}{1.850278in}%
\pgfsys@useobject{currentmarker}{}%
\end{pgfscope}%
\begin{pgfscope}%
\pgfsys@transformshift{0.937188in}{1.850278in}%
\pgfsys@useobject{currentmarker}{}%
\end{pgfscope}%
\begin{pgfscope}%
\pgfsys@transformshift{0.937525in}{1.850278in}%
\pgfsys@useobject{currentmarker}{}%
\end{pgfscope}%
\begin{pgfscope}%
\pgfsys@transformshift{0.937525in}{1.850278in}%
\pgfsys@useobject{currentmarker}{}%
\end{pgfscope}%
\begin{pgfscope}%
\pgfsys@transformshift{0.937863in}{1.850278in}%
\pgfsys@useobject{currentmarker}{}%
\end{pgfscope}%
\begin{pgfscope}%
\pgfsys@transformshift{0.937863in}{1.850278in}%
\pgfsys@useobject{currentmarker}{}%
\end{pgfscope}%
\begin{pgfscope}%
\pgfsys@transformshift{0.938200in}{1.850278in}%
\pgfsys@useobject{currentmarker}{}%
\end{pgfscope}%
\begin{pgfscope}%
\pgfsys@transformshift{0.938200in}{1.850278in}%
\pgfsys@useobject{currentmarker}{}%
\end{pgfscope}%
\begin{pgfscope}%
\pgfsys@transformshift{0.938537in}{1.850278in}%
\pgfsys@useobject{currentmarker}{}%
\end{pgfscope}%
\begin{pgfscope}%
\pgfsys@transformshift{0.938537in}{1.850278in}%
\pgfsys@useobject{currentmarker}{}%
\end{pgfscope}%
\begin{pgfscope}%
\pgfsys@transformshift{0.938874in}{1.850278in}%
\pgfsys@useobject{currentmarker}{}%
\end{pgfscope}%
\begin{pgfscope}%
\pgfsys@transformshift{0.938874in}{1.850278in}%
\pgfsys@useobject{currentmarker}{}%
\end{pgfscope}%
\begin{pgfscope}%
\pgfsys@transformshift{0.939211in}{1.850278in}%
\pgfsys@useobject{currentmarker}{}%
\end{pgfscope}%
\begin{pgfscope}%
\pgfsys@transformshift{0.939211in}{1.850278in}%
\pgfsys@useobject{currentmarker}{}%
\end{pgfscope}%
\begin{pgfscope}%
\pgfsys@transformshift{0.939549in}{1.850278in}%
\pgfsys@useobject{currentmarker}{}%
\end{pgfscope}%
\begin{pgfscope}%
\pgfsys@transformshift{0.939549in}{1.850278in}%
\pgfsys@useobject{currentmarker}{}%
\end{pgfscope}%
\begin{pgfscope}%
\pgfsys@transformshift{0.939886in}{1.850278in}%
\pgfsys@useobject{currentmarker}{}%
\end{pgfscope}%
\begin{pgfscope}%
\pgfsys@transformshift{0.940055in}{1.850278in}%
\pgfsys@useobject{currentmarker}{}%
\end{pgfscope}%
\begin{pgfscope}%
\pgfsys@transformshift{0.940223in}{1.850278in}%
\pgfsys@useobject{currentmarker}{}%
\end{pgfscope}%
\begin{pgfscope}%
\pgfsys@transformshift{0.940560in}{1.850278in}%
\pgfsys@useobject{currentmarker}{}%
\end{pgfscope}%
\begin{pgfscope}%
\pgfsys@transformshift{0.940729in}{1.850278in}%
\pgfsys@useobject{currentmarker}{}%
\end{pgfscope}%
\begin{pgfscope}%
\pgfsys@transformshift{0.940898in}{1.850278in}%
\pgfsys@useobject{currentmarker}{}%
\end{pgfscope}%
\begin{pgfscope}%
\pgfsys@transformshift{0.941235in}{1.850278in}%
\pgfsys@useobject{currentmarker}{}%
\end{pgfscope}%
\begin{pgfscope}%
\pgfsys@transformshift{0.941403in}{1.850278in}%
\pgfsys@useobject{currentmarker}{}%
\end{pgfscope}%
\begin{pgfscope}%
\pgfsys@transformshift{0.941573in}{1.850278in}%
\pgfsys@useobject{currentmarker}{}%
\end{pgfscope}%
\begin{pgfscope}%
\pgfsys@transformshift{0.941741in}{1.850278in}%
\pgfsys@useobject{currentmarker}{}%
\end{pgfscope}%
\begin{pgfscope}%
\pgfsys@transformshift{0.941909in}{1.850278in}%
\pgfsys@useobject{currentmarker}{}%
\end{pgfscope}%
\begin{pgfscope}%
\pgfsys@transformshift{0.942078in}{1.850278in}%
\pgfsys@useobject{currentmarker}{}%
\end{pgfscope}%
\begin{pgfscope}%
\pgfsys@transformshift{0.942247in}{1.850278in}%
\pgfsys@useobject{currentmarker}{}%
\end{pgfscope}%
\begin{pgfscope}%
\pgfsys@transformshift{0.942416in}{1.850278in}%
\pgfsys@useobject{currentmarker}{}%
\end{pgfscope}%
\begin{pgfscope}%
\pgfsys@transformshift{0.942584in}{1.850278in}%
\pgfsys@useobject{currentmarker}{}%
\end{pgfscope}%
\begin{pgfscope}%
\pgfsys@transformshift{0.942752in}{1.850278in}%
\pgfsys@useobject{currentmarker}{}%
\end{pgfscope}%
\begin{pgfscope}%
\pgfsys@transformshift{0.942921in}{1.850278in}%
\pgfsys@useobject{currentmarker}{}%
\end{pgfscope}%
\begin{pgfscope}%
\pgfsys@transformshift{0.942921in}{1.850278in}%
\pgfsys@useobject{currentmarker}{}%
\end{pgfscope}%
\begin{pgfscope}%
\pgfsys@transformshift{0.943090in}{1.850278in}%
\pgfsys@useobject{currentmarker}{}%
\end{pgfscope}%
\begin{pgfscope}%
\pgfsys@transformshift{0.943259in}{1.850278in}%
\pgfsys@useobject{currentmarker}{}%
\end{pgfscope}%
\begin{pgfscope}%
\pgfsys@transformshift{0.943427in}{1.850278in}%
\pgfsys@useobject{currentmarker}{}%
\end{pgfscope}%
\begin{pgfscope}%
\pgfsys@transformshift{0.943595in}{1.850278in}%
\pgfsys@useobject{currentmarker}{}%
\end{pgfscope}%
\begin{pgfscope}%
\pgfsys@transformshift{0.943595in}{1.850278in}%
\pgfsys@useobject{currentmarker}{}%
\end{pgfscope}%
\begin{pgfscope}%
\pgfsys@transformshift{0.943765in}{1.850278in}%
\pgfsys@useobject{currentmarker}{}%
\end{pgfscope}%
\begin{pgfscope}%
\pgfsys@transformshift{0.943933in}{1.850278in}%
\pgfsys@useobject{currentmarker}{}%
\end{pgfscope}%
\begin{pgfscope}%
\pgfsys@transformshift{0.943933in}{1.850278in}%
\pgfsys@useobject{currentmarker}{}%
\end{pgfscope}%
\begin{pgfscope}%
\pgfsys@transformshift{0.943933in}{1.850278in}%
\pgfsys@useobject{currentmarker}{}%
\end{pgfscope}%
\begin{pgfscope}%
\pgfsys@transformshift{0.943933in}{1.850278in}%
\pgfsys@useobject{currentmarker}{}%
\end{pgfscope}%
\begin{pgfscope}%
\pgfsys@transformshift{0.944270in}{1.850278in}%
\pgfsys@useobject{currentmarker}{}%
\end{pgfscope}%
\begin{pgfscope}%
\pgfsys@transformshift{0.944270in}{1.850278in}%
\pgfsys@useobject{currentmarker}{}%
\end{pgfscope}%
\begin{pgfscope}%
\pgfsys@transformshift{0.944270in}{1.850278in}%
\pgfsys@useobject{currentmarker}{}%
\end{pgfscope}%
\begin{pgfscope}%
\pgfsys@transformshift{0.944270in}{1.850278in}%
\pgfsys@useobject{currentmarker}{}%
\end{pgfscope}%
\begin{pgfscope}%
\pgfsys@transformshift{0.944608in}{1.850278in}%
\pgfsys@useobject{currentmarker}{}%
\end{pgfscope}%
\begin{pgfscope}%
\pgfsys@transformshift{0.944608in}{1.850278in}%
\pgfsys@useobject{currentmarker}{}%
\end{pgfscope}%
\begin{pgfscope}%
\pgfsys@transformshift{0.944608in}{1.850278in}%
\pgfsys@useobject{currentmarker}{}%
\end{pgfscope}%
\begin{pgfscope}%
\pgfsys@transformshift{0.944608in}{1.850278in}%
\pgfsys@useobject{currentmarker}{}%
\end{pgfscope}%
\begin{pgfscope}%
\pgfsys@transformshift{0.944945in}{1.850278in}%
\pgfsys@useobject{currentmarker}{}%
\end{pgfscope}%
\begin{pgfscope}%
\pgfsys@transformshift{0.944945in}{1.850278in}%
\pgfsys@useobject{currentmarker}{}%
\end{pgfscope}%
\begin{pgfscope}%
\pgfsys@transformshift{0.944945in}{1.850278in}%
\pgfsys@useobject{currentmarker}{}%
\end{pgfscope}%
\begin{pgfscope}%
\pgfsys@transformshift{0.944945in}{1.850278in}%
\pgfsys@useobject{currentmarker}{}%
\end{pgfscope}%
\begin{pgfscope}%
\pgfsys@transformshift{0.945282in}{1.850278in}%
\pgfsys@useobject{currentmarker}{}%
\end{pgfscope}%
\begin{pgfscope}%
\pgfsys@transformshift{0.945282in}{1.850278in}%
\pgfsys@useobject{currentmarker}{}%
\end{pgfscope}%
\begin{pgfscope}%
\pgfsys@transformshift{0.945282in}{1.850278in}%
\pgfsys@useobject{currentmarker}{}%
\end{pgfscope}%
\begin{pgfscope}%
\pgfsys@transformshift{0.945282in}{1.850278in}%
\pgfsys@useobject{currentmarker}{}%
\end{pgfscope}%
\begin{pgfscope}%
\pgfsys@transformshift{0.945619in}{1.850278in}%
\pgfsys@useobject{currentmarker}{}%
\end{pgfscope}%
\begin{pgfscope}%
\pgfsys@transformshift{0.945619in}{1.850278in}%
\pgfsys@useobject{currentmarker}{}%
\end{pgfscope}%
\begin{pgfscope}%
\pgfsys@transformshift{0.945619in}{1.850278in}%
\pgfsys@useobject{currentmarker}{}%
\end{pgfscope}%
\begin{pgfscope}%
\pgfsys@transformshift{0.945619in}{1.850278in}%
\pgfsys@useobject{currentmarker}{}%
\end{pgfscope}%
\begin{pgfscope}%
\pgfsys@transformshift{0.945957in}{1.850278in}%
\pgfsys@useobject{currentmarker}{}%
\end{pgfscope}%
\begin{pgfscope}%
\pgfsys@transformshift{0.945957in}{1.850278in}%
\pgfsys@useobject{currentmarker}{}%
\end{pgfscope}%
\begin{pgfscope}%
\pgfsys@transformshift{0.945957in}{1.850278in}%
\pgfsys@useobject{currentmarker}{}%
\end{pgfscope}%
\begin{pgfscope}%
\pgfsys@transformshift{0.945957in}{1.850278in}%
\pgfsys@useobject{currentmarker}{}%
\end{pgfscope}%
\begin{pgfscope}%
\pgfsys@transformshift{0.946294in}{1.850278in}%
\pgfsys@useobject{currentmarker}{}%
\end{pgfscope}%
\begin{pgfscope}%
\pgfsys@transformshift{0.946294in}{1.850278in}%
\pgfsys@useobject{currentmarker}{}%
\end{pgfscope}%
\begin{pgfscope}%
\pgfsys@transformshift{0.946294in}{1.850278in}%
\pgfsys@useobject{currentmarker}{}%
\end{pgfscope}%
\begin{pgfscope}%
\pgfsys@transformshift{0.946294in}{1.850278in}%
\pgfsys@useobject{currentmarker}{}%
\end{pgfscope}%
\begin{pgfscope}%
\pgfsys@transformshift{0.946631in}{1.850278in}%
\pgfsys@useobject{currentmarker}{}%
\end{pgfscope}%
\begin{pgfscope}%
\pgfsys@transformshift{0.946631in}{1.850278in}%
\pgfsys@useobject{currentmarker}{}%
\end{pgfscope}%
\begin{pgfscope}%
\pgfsys@transformshift{0.946631in}{1.850278in}%
\pgfsys@useobject{currentmarker}{}%
\end{pgfscope}%
\begin{pgfscope}%
\pgfsys@transformshift{0.946800in}{1.850278in}%
\pgfsys@useobject{currentmarker}{}%
\end{pgfscope}%
\begin{pgfscope}%
\pgfsys@transformshift{0.946968in}{1.850278in}%
\pgfsys@useobject{currentmarker}{}%
\end{pgfscope}%
\begin{pgfscope}%
\pgfsys@transformshift{0.946968in}{1.850278in}%
\pgfsys@useobject{currentmarker}{}%
\end{pgfscope}%
\begin{pgfscope}%
\pgfsys@transformshift{0.947137in}{1.850278in}%
\pgfsys@useobject{currentmarker}{}%
\end{pgfscope}%
\begin{pgfscope}%
\pgfsys@transformshift{0.947305in}{1.850278in}%
\pgfsys@useobject{currentmarker}{}%
\end{pgfscope}%
\begin{pgfscope}%
\pgfsys@transformshift{0.947305in}{1.850278in}%
\pgfsys@useobject{currentmarker}{}%
\end{pgfscope}%
\begin{pgfscope}%
\pgfsys@transformshift{0.947475in}{1.850278in}%
\pgfsys@useobject{currentmarker}{}%
\end{pgfscope}%
\begin{pgfscope}%
\pgfsys@transformshift{0.947643in}{1.850278in}%
\pgfsys@useobject{currentmarker}{}%
\end{pgfscope}%
\begin{pgfscope}%
\pgfsys@transformshift{0.947643in}{1.850278in}%
\pgfsys@useobject{currentmarker}{}%
\end{pgfscope}%
\begin{pgfscope}%
\pgfsys@transformshift{0.947811in}{1.850278in}%
\pgfsys@useobject{currentmarker}{}%
\end{pgfscope}%
\begin{pgfscope}%
\pgfsys@transformshift{0.947980in}{1.850278in}%
\pgfsys@useobject{currentmarker}{}%
\end{pgfscope}%
\begin{pgfscope}%
\pgfsys@transformshift{0.947980in}{1.850278in}%
\pgfsys@useobject{currentmarker}{}%
\end{pgfscope}%
\begin{pgfscope}%
\pgfsys@transformshift{0.948149in}{1.850278in}%
\pgfsys@useobject{currentmarker}{}%
\end{pgfscope}%
\begin{pgfscope}%
\pgfsys@transformshift{0.948318in}{1.850278in}%
\pgfsys@useobject{currentmarker}{}%
\end{pgfscope}%
\begin{pgfscope}%
\pgfsys@transformshift{0.948318in}{1.850278in}%
\pgfsys@useobject{currentmarker}{}%
\end{pgfscope}%
\begin{pgfscope}%
\pgfsys@transformshift{0.948318in}{1.850278in}%
\pgfsys@useobject{currentmarker}{}%
\end{pgfscope}%
\begin{pgfscope}%
\pgfsys@transformshift{0.948486in}{1.850278in}%
\pgfsys@useobject{currentmarker}{}%
\end{pgfscope}%
\begin{pgfscope}%
\pgfsys@transformshift{0.948654in}{1.850278in}%
\pgfsys@useobject{currentmarker}{}%
\end{pgfscope}%
\begin{pgfscope}%
\pgfsys@transformshift{0.948654in}{1.850278in}%
\pgfsys@useobject{currentmarker}{}%
\end{pgfscope}%
\begin{pgfscope}%
\pgfsys@transformshift{0.948654in}{1.850278in}%
\pgfsys@useobject{currentmarker}{}%
\end{pgfscope}%
\begin{pgfscope}%
\pgfsys@transformshift{0.948823in}{1.850278in}%
\pgfsys@useobject{currentmarker}{}%
\end{pgfscope}%
\begin{pgfscope}%
\pgfsys@transformshift{0.948992in}{1.850278in}%
\pgfsys@useobject{currentmarker}{}%
\end{pgfscope}%
\begin{pgfscope}%
\pgfsys@transformshift{0.948992in}{1.850278in}%
\pgfsys@useobject{currentmarker}{}%
\end{pgfscope}%
\begin{pgfscope}%
\pgfsys@transformshift{0.948992in}{1.850278in}%
\pgfsys@useobject{currentmarker}{}%
\end{pgfscope}%
\begin{pgfscope}%
\pgfsys@transformshift{0.949161in}{1.850278in}%
\pgfsys@useobject{currentmarker}{}%
\end{pgfscope}%
\begin{pgfscope}%
\pgfsys@transformshift{0.949329in}{1.850278in}%
\pgfsys@useobject{currentmarker}{}%
\end{pgfscope}%
\begin{pgfscope}%
\pgfsys@transformshift{0.949329in}{1.850278in}%
\pgfsys@useobject{currentmarker}{}%
\end{pgfscope}%
\begin{pgfscope}%
\pgfsys@transformshift{0.949329in}{1.850278in}%
\pgfsys@useobject{currentmarker}{}%
\end{pgfscope}%
\begin{pgfscope}%
\pgfsys@transformshift{0.949497in}{1.850278in}%
\pgfsys@useobject{currentmarker}{}%
\end{pgfscope}%
\begin{pgfscope}%
\pgfsys@transformshift{0.949497in}{1.850278in}%
\pgfsys@useobject{currentmarker}{}%
\end{pgfscope}%
\begin{pgfscope}%
\pgfsys@transformshift{0.949667in}{1.850278in}%
\pgfsys@useobject{currentmarker}{}%
\end{pgfscope}%
\begin{pgfscope}%
\pgfsys@transformshift{0.949667in}{1.850278in}%
\pgfsys@useobject{currentmarker}{}%
\end{pgfscope}%
\begin{pgfscope}%
\pgfsys@transformshift{0.949667in}{1.850278in}%
\pgfsys@useobject{currentmarker}{}%
\end{pgfscope}%
\begin{pgfscope}%
\pgfsys@transformshift{0.949835in}{1.850278in}%
\pgfsys@useobject{currentmarker}{}%
\end{pgfscope}%
\begin{pgfscope}%
\pgfsys@transformshift{0.949835in}{1.850278in}%
\pgfsys@useobject{currentmarker}{}%
\end{pgfscope}%
\begin{pgfscope}%
\pgfsys@transformshift{0.950004in}{1.850278in}%
\pgfsys@useobject{currentmarker}{}%
\end{pgfscope}%
\begin{pgfscope}%
\pgfsys@transformshift{0.950004in}{1.850278in}%
\pgfsys@useobject{currentmarker}{}%
\end{pgfscope}%
\begin{pgfscope}%
\pgfsys@transformshift{0.950004in}{1.850278in}%
\pgfsys@useobject{currentmarker}{}%
\end{pgfscope}%
\begin{pgfscope}%
\pgfsys@transformshift{0.950172in}{1.850278in}%
\pgfsys@useobject{currentmarker}{}%
\end{pgfscope}%
\begin{pgfscope}%
\pgfsys@transformshift{0.950172in}{1.850278in}%
\pgfsys@useobject{currentmarker}{}%
\end{pgfscope}%
\begin{pgfscope}%
\pgfsys@transformshift{0.950341in}{1.850278in}%
\pgfsys@useobject{currentmarker}{}%
\end{pgfscope}%
\begin{pgfscope}%
\pgfsys@transformshift{0.950341in}{1.850278in}%
\pgfsys@useobject{currentmarker}{}%
\end{pgfscope}%
\begin{pgfscope}%
\pgfsys@transformshift{0.950341in}{1.850278in}%
\pgfsys@useobject{currentmarker}{}%
\end{pgfscope}%
\begin{pgfscope}%
\pgfsys@transformshift{0.950510in}{1.850278in}%
\pgfsys@useobject{currentmarker}{}%
\end{pgfscope}%
\begin{pgfscope}%
\pgfsys@transformshift{0.950510in}{1.850278in}%
\pgfsys@useobject{currentmarker}{}%
\end{pgfscope}%
\begin{pgfscope}%
\pgfsys@transformshift{0.950678in}{1.850278in}%
\pgfsys@useobject{currentmarker}{}%
\end{pgfscope}%
\begin{pgfscope}%
\pgfsys@transformshift{0.950678in}{1.850278in}%
\pgfsys@useobject{currentmarker}{}%
\end{pgfscope}%
\begin{pgfscope}%
\pgfsys@transformshift{0.950678in}{1.850278in}%
\pgfsys@useobject{currentmarker}{}%
\end{pgfscope}%
\begin{pgfscope}%
\pgfsys@transformshift{0.950678in}{1.850278in}%
\pgfsys@useobject{currentmarker}{}%
\end{pgfscope}%
\begin{pgfscope}%
\pgfsys@transformshift{0.951015in}{1.850278in}%
\pgfsys@useobject{currentmarker}{}%
\end{pgfscope}%
\begin{pgfscope}%
\pgfsys@transformshift{0.951015in}{1.850278in}%
\pgfsys@useobject{currentmarker}{}%
\end{pgfscope}%
\begin{pgfscope}%
\pgfsys@transformshift{0.951015in}{1.850278in}%
\pgfsys@useobject{currentmarker}{}%
\end{pgfscope}%
\begin{pgfscope}%
\pgfsys@transformshift{0.951015in}{1.850278in}%
\pgfsys@useobject{currentmarker}{}%
\end{pgfscope}%
\begin{pgfscope}%
\pgfsys@transformshift{0.951353in}{1.850278in}%
\pgfsys@useobject{currentmarker}{}%
\end{pgfscope}%
\begin{pgfscope}%
\pgfsys@transformshift{0.951353in}{1.850278in}%
\pgfsys@useobject{currentmarker}{}%
\end{pgfscope}%
\begin{pgfscope}%
\pgfsys@transformshift{0.951353in}{1.850278in}%
\pgfsys@useobject{currentmarker}{}%
\end{pgfscope}%
\begin{pgfscope}%
\pgfsys@transformshift{0.951353in}{1.850278in}%
\pgfsys@useobject{currentmarker}{}%
\end{pgfscope}%
\begin{pgfscope}%
\pgfsys@transformshift{0.951690in}{1.850278in}%
\pgfsys@useobject{currentmarker}{}%
\end{pgfscope}%
\begin{pgfscope}%
\pgfsys@transformshift{0.951690in}{1.850278in}%
\pgfsys@useobject{currentmarker}{}%
\end{pgfscope}%
\begin{pgfscope}%
\pgfsys@transformshift{0.951690in}{1.850278in}%
\pgfsys@useobject{currentmarker}{}%
\end{pgfscope}%
\begin{pgfscope}%
\pgfsys@transformshift{0.951690in}{1.850278in}%
\pgfsys@useobject{currentmarker}{}%
\end{pgfscope}%
\begin{pgfscope}%
\pgfsys@transformshift{0.952027in}{1.850278in}%
\pgfsys@useobject{currentmarker}{}%
\end{pgfscope}%
\begin{pgfscope}%
\pgfsys@transformshift{0.952027in}{1.850278in}%
\pgfsys@useobject{currentmarker}{}%
\end{pgfscope}%
\begin{pgfscope}%
\pgfsys@transformshift{0.952027in}{1.850278in}%
\pgfsys@useobject{currentmarker}{}%
\end{pgfscope}%
\begin{pgfscope}%
\pgfsys@transformshift{0.952027in}{1.850278in}%
\pgfsys@useobject{currentmarker}{}%
\end{pgfscope}%
\begin{pgfscope}%
\pgfsys@transformshift{0.952364in}{1.850278in}%
\pgfsys@useobject{currentmarker}{}%
\end{pgfscope}%
\begin{pgfscope}%
\pgfsys@transformshift{0.952364in}{1.850278in}%
\pgfsys@useobject{currentmarker}{}%
\end{pgfscope}%
\begin{pgfscope}%
\pgfsys@transformshift{0.952364in}{1.850278in}%
\pgfsys@useobject{currentmarker}{}%
\end{pgfscope}%
\begin{pgfscope}%
\pgfsys@transformshift{0.952364in}{1.850278in}%
\pgfsys@useobject{currentmarker}{}%
\end{pgfscope}%
\begin{pgfscope}%
\pgfsys@transformshift{0.952702in}{1.850278in}%
\pgfsys@useobject{currentmarker}{}%
\end{pgfscope}%
\begin{pgfscope}%
\pgfsys@transformshift{0.952702in}{1.850278in}%
\pgfsys@useobject{currentmarker}{}%
\end{pgfscope}%
\begin{pgfscope}%
\pgfsys@transformshift{0.952702in}{1.850278in}%
\pgfsys@useobject{currentmarker}{}%
\end{pgfscope}%
\begin{pgfscope}%
\pgfsys@transformshift{0.952702in}{1.850278in}%
\pgfsys@useobject{currentmarker}{}%
\end{pgfscope}%
\begin{pgfscope}%
\pgfsys@transformshift{0.953039in}{1.850278in}%
\pgfsys@useobject{currentmarker}{}%
\end{pgfscope}%
\begin{pgfscope}%
\pgfsys@transformshift{0.953039in}{1.850278in}%
\pgfsys@useobject{currentmarker}{}%
\end{pgfscope}%
\begin{pgfscope}%
\pgfsys@transformshift{0.953039in}{1.850278in}%
\pgfsys@useobject{currentmarker}{}%
\end{pgfscope}%
\begin{pgfscope}%
\pgfsys@transformshift{0.953039in}{1.850278in}%
\pgfsys@useobject{currentmarker}{}%
\end{pgfscope}%
\begin{pgfscope}%
\pgfsys@transformshift{0.953377in}{1.850278in}%
\pgfsys@useobject{currentmarker}{}%
\end{pgfscope}%
\begin{pgfscope}%
\pgfsys@transformshift{0.953377in}{1.850278in}%
\pgfsys@useobject{currentmarker}{}%
\end{pgfscope}%
\begin{pgfscope}%
\pgfsys@transformshift{0.953377in}{1.850278in}%
\pgfsys@useobject{currentmarker}{}%
\end{pgfscope}%
\begin{pgfscope}%
\pgfsys@transformshift{0.953545in}{1.850278in}%
\pgfsys@useobject{currentmarker}{}%
\end{pgfscope}%
\begin{pgfscope}%
\pgfsys@transformshift{0.953713in}{1.850278in}%
\pgfsys@useobject{currentmarker}{}%
\end{pgfscope}%
\begin{pgfscope}%
\pgfsys@transformshift{0.953713in}{1.850278in}%
\pgfsys@useobject{currentmarker}{}%
\end{pgfscope}%
\begin{pgfscope}%
\pgfsys@transformshift{0.953882in}{1.850278in}%
\pgfsys@useobject{currentmarker}{}%
\end{pgfscope}%
\begin{pgfscope}%
\pgfsys@transformshift{0.954051in}{1.850278in}%
\pgfsys@useobject{currentmarker}{}%
\end{pgfscope}%
\begin{pgfscope}%
\pgfsys@transformshift{0.954051in}{1.850278in}%
\pgfsys@useobject{currentmarker}{}%
\end{pgfscope}%
\begin{pgfscope}%
\pgfsys@transformshift{0.954220in}{1.850278in}%
\pgfsys@useobject{currentmarker}{}%
\end{pgfscope}%
\begin{pgfscope}%
\pgfsys@transformshift{0.954388in}{1.850278in}%
\pgfsys@useobject{currentmarker}{}%
\end{pgfscope}%
\begin{pgfscope}%
\pgfsys@transformshift{0.954388in}{1.850278in}%
\pgfsys@useobject{currentmarker}{}%
\end{pgfscope}%
\begin{pgfscope}%
\pgfsys@transformshift{0.954556in}{1.850278in}%
\pgfsys@useobject{currentmarker}{}%
\end{pgfscope}%
\begin{pgfscope}%
\pgfsys@transformshift{0.954725in}{1.850278in}%
\pgfsys@useobject{currentmarker}{}%
\end{pgfscope}%
\begin{pgfscope}%
\pgfsys@transformshift{0.954725in}{1.850278in}%
\pgfsys@useobject{currentmarker}{}%
\end{pgfscope}%
\begin{pgfscope}%
\pgfsys@transformshift{0.954894in}{1.850278in}%
\pgfsys@useobject{currentmarker}{}%
\end{pgfscope}%
\begin{pgfscope}%
\pgfsys@transformshift{0.955063in}{1.850278in}%
\pgfsys@useobject{currentmarker}{}%
\end{pgfscope}%
\begin{pgfscope}%
\pgfsys@transformshift{0.955063in}{1.850278in}%
\pgfsys@useobject{currentmarker}{}%
\end{pgfscope}%
\begin{pgfscope}%
\pgfsys@transformshift{0.955063in}{1.850278in}%
\pgfsys@useobject{currentmarker}{}%
\end{pgfscope}%
\begin{pgfscope}%
\pgfsys@transformshift{0.955231in}{1.850278in}%
\pgfsys@useobject{currentmarker}{}%
\end{pgfscope}%
\begin{pgfscope}%
\pgfsys@transformshift{0.955399in}{1.850278in}%
\pgfsys@useobject{currentmarker}{}%
\end{pgfscope}%
\begin{pgfscope}%
\pgfsys@transformshift{0.955399in}{1.850278in}%
\pgfsys@useobject{currentmarker}{}%
\end{pgfscope}%
\begin{pgfscope}%
\pgfsys@transformshift{0.955399in}{1.850278in}%
\pgfsys@useobject{currentmarker}{}%
\end{pgfscope}%
\begin{pgfscope}%
\pgfsys@transformshift{0.955569in}{1.850278in}%
\pgfsys@useobject{currentmarker}{}%
\end{pgfscope}%
\begin{pgfscope}%
\pgfsys@transformshift{0.955737in}{1.850278in}%
\pgfsys@useobject{currentmarker}{}%
\end{pgfscope}%
\begin{pgfscope}%
\pgfsys@transformshift{0.955737in}{1.850278in}%
\pgfsys@useobject{currentmarker}{}%
\end{pgfscope}%
\begin{pgfscope}%
\pgfsys@transformshift{0.955737in}{1.850278in}%
\pgfsys@useobject{currentmarker}{}%
\end{pgfscope}%
\begin{pgfscope}%
\pgfsys@transformshift{0.955906in}{1.850278in}%
\pgfsys@useobject{currentmarker}{}%
\end{pgfscope}%
\begin{pgfscope}%
\pgfsys@transformshift{0.956074in}{1.850278in}%
\pgfsys@useobject{currentmarker}{}%
\end{pgfscope}%
\begin{pgfscope}%
\pgfsys@transformshift{0.956074in}{1.850278in}%
\pgfsys@useobject{currentmarker}{}%
\end{pgfscope}%
\begin{pgfscope}%
\pgfsys@transformshift{0.956074in}{1.850278in}%
\pgfsys@useobject{currentmarker}{}%
\end{pgfscope}%
\begin{pgfscope}%
\pgfsys@transformshift{0.956243in}{1.850278in}%
\pgfsys@useobject{currentmarker}{}%
\end{pgfscope}%
\begin{pgfscope}%
\pgfsys@transformshift{0.956243in}{1.850278in}%
\pgfsys@useobject{currentmarker}{}%
\end{pgfscope}%
\begin{pgfscope}%
\pgfsys@transformshift{0.956412in}{1.850278in}%
\pgfsys@useobject{currentmarker}{}%
\end{pgfscope}%
\begin{pgfscope}%
\pgfsys@transformshift{0.956412in}{1.850278in}%
\pgfsys@useobject{currentmarker}{}%
\end{pgfscope}%
\begin{pgfscope}%
\pgfsys@transformshift{0.956412in}{1.850278in}%
\pgfsys@useobject{currentmarker}{}%
\end{pgfscope}%
\begin{pgfscope}%
\pgfsys@transformshift{0.956580in}{1.850278in}%
\pgfsys@useobject{currentmarker}{}%
\end{pgfscope}%
\begin{pgfscope}%
\pgfsys@transformshift{0.956580in}{1.850278in}%
\pgfsys@useobject{currentmarker}{}%
\end{pgfscope}%
\begin{pgfscope}%
\pgfsys@transformshift{0.956749in}{1.850278in}%
\pgfsys@useobject{currentmarker}{}%
\end{pgfscope}%
\begin{pgfscope}%
\pgfsys@transformshift{0.956749in}{1.850278in}%
\pgfsys@useobject{currentmarker}{}%
\end{pgfscope}%
\begin{pgfscope}%
\pgfsys@transformshift{0.956749in}{1.850278in}%
\pgfsys@useobject{currentmarker}{}%
\end{pgfscope}%
\begin{pgfscope}%
\pgfsys@transformshift{0.956917in}{1.850278in}%
\pgfsys@useobject{currentmarker}{}%
\end{pgfscope}%
\begin{pgfscope}%
\pgfsys@transformshift{0.956917in}{1.850278in}%
\pgfsys@useobject{currentmarker}{}%
\end{pgfscope}%
\begin{pgfscope}%
\pgfsys@transformshift{0.957086in}{1.850278in}%
\pgfsys@useobject{currentmarker}{}%
\end{pgfscope}%
\begin{pgfscope}%
\pgfsys@transformshift{0.957086in}{1.850278in}%
\pgfsys@useobject{currentmarker}{}%
\end{pgfscope}%
\begin{pgfscope}%
\pgfsys@transformshift{0.957086in}{1.850278in}%
\pgfsys@useobject{currentmarker}{}%
\end{pgfscope}%
\begin{pgfscope}%
\pgfsys@transformshift{0.957255in}{1.850278in}%
\pgfsys@useobject{currentmarker}{}%
\end{pgfscope}%
\begin{pgfscope}%
\pgfsys@transformshift{0.957255in}{1.850278in}%
\pgfsys@useobject{currentmarker}{}%
\end{pgfscope}%
\begin{pgfscope}%
\pgfsys@transformshift{0.957423in}{1.850278in}%
\pgfsys@useobject{currentmarker}{}%
\end{pgfscope}%
\begin{pgfscope}%
\pgfsys@transformshift{0.957423in}{1.850278in}%
\pgfsys@useobject{currentmarker}{}%
\end{pgfscope}%
\begin{pgfscope}%
\pgfsys@transformshift{0.957423in}{1.850278in}%
\pgfsys@useobject{currentmarker}{}%
\end{pgfscope}%
\begin{pgfscope}%
\pgfsys@transformshift{0.957423in}{1.850278in}%
\pgfsys@useobject{currentmarker}{}%
\end{pgfscope}%
\begin{pgfscope}%
\pgfsys@transformshift{0.957761in}{1.850278in}%
\pgfsys@useobject{currentmarker}{}%
\end{pgfscope}%
\begin{pgfscope}%
\pgfsys@transformshift{0.957761in}{1.850278in}%
\pgfsys@useobject{currentmarker}{}%
\end{pgfscope}%
\begin{pgfscope}%
\pgfsys@transformshift{0.957761in}{1.850278in}%
\pgfsys@useobject{currentmarker}{}%
\end{pgfscope}%
\begin{pgfscope}%
\pgfsys@transformshift{0.957761in}{1.850278in}%
\pgfsys@useobject{currentmarker}{}%
\end{pgfscope}%
\begin{pgfscope}%
\pgfsys@transformshift{0.958098in}{1.850278in}%
\pgfsys@useobject{currentmarker}{}%
\end{pgfscope}%
\begin{pgfscope}%
\pgfsys@transformshift{0.958098in}{1.850278in}%
\pgfsys@useobject{currentmarker}{}%
\end{pgfscope}%
\begin{pgfscope}%
\pgfsys@transformshift{0.958098in}{1.850278in}%
\pgfsys@useobject{currentmarker}{}%
\end{pgfscope}%
\begin{pgfscope}%
\pgfsys@transformshift{0.958098in}{1.850278in}%
\pgfsys@useobject{currentmarker}{}%
\end{pgfscope}%
\begin{pgfscope}%
\pgfsys@transformshift{0.958435in}{1.850278in}%
\pgfsys@useobject{currentmarker}{}%
\end{pgfscope}%
\begin{pgfscope}%
\pgfsys@transformshift{0.958435in}{1.850278in}%
\pgfsys@useobject{currentmarker}{}%
\end{pgfscope}%
\begin{pgfscope}%
\pgfsys@transformshift{0.958435in}{1.850278in}%
\pgfsys@useobject{currentmarker}{}%
\end{pgfscope}%
\begin{pgfscope}%
\pgfsys@transformshift{0.958435in}{1.850278in}%
\pgfsys@useobject{currentmarker}{}%
\end{pgfscope}%
\begin{pgfscope}%
\pgfsys@transformshift{0.958772in}{1.850278in}%
\pgfsys@useobject{currentmarker}{}%
\end{pgfscope}%
\begin{pgfscope}%
\pgfsys@transformshift{0.958772in}{1.850278in}%
\pgfsys@useobject{currentmarker}{}%
\end{pgfscope}%
\begin{pgfscope}%
\pgfsys@transformshift{0.958772in}{1.850278in}%
\pgfsys@useobject{currentmarker}{}%
\end{pgfscope}%
\begin{pgfscope}%
\pgfsys@transformshift{0.958772in}{1.850278in}%
\pgfsys@useobject{currentmarker}{}%
\end{pgfscope}%
\begin{pgfscope}%
\pgfsys@transformshift{0.959109in}{1.850278in}%
\pgfsys@useobject{currentmarker}{}%
\end{pgfscope}%
\begin{pgfscope}%
\pgfsys@transformshift{0.959109in}{1.850278in}%
\pgfsys@useobject{currentmarker}{}%
\end{pgfscope}%
\begin{pgfscope}%
\pgfsys@transformshift{0.959109in}{1.850278in}%
\pgfsys@useobject{currentmarker}{}%
\end{pgfscope}%
\begin{pgfscope}%
\pgfsys@transformshift{0.959109in}{1.850278in}%
\pgfsys@useobject{currentmarker}{}%
\end{pgfscope}%
\begin{pgfscope}%
\pgfsys@transformshift{0.959447in}{1.850278in}%
\pgfsys@useobject{currentmarker}{}%
\end{pgfscope}%
\begin{pgfscope}%
\pgfsys@transformshift{0.959447in}{1.850278in}%
\pgfsys@useobject{currentmarker}{}%
\end{pgfscope}%
\begin{pgfscope}%
\pgfsys@transformshift{0.959447in}{1.850278in}%
\pgfsys@useobject{currentmarker}{}%
\end{pgfscope}%
\begin{pgfscope}%
\pgfsys@transformshift{0.959447in}{1.850278in}%
\pgfsys@useobject{currentmarker}{}%
\end{pgfscope}%
\begin{pgfscope}%
\pgfsys@transformshift{0.959784in}{1.850278in}%
\pgfsys@useobject{currentmarker}{}%
\end{pgfscope}%
\begin{pgfscope}%
\pgfsys@transformshift{0.959784in}{1.850278in}%
\pgfsys@useobject{currentmarker}{}%
\end{pgfscope}%
\begin{pgfscope}%
\pgfsys@transformshift{0.959784in}{1.850278in}%
\pgfsys@useobject{currentmarker}{}%
\end{pgfscope}%
\begin{pgfscope}%
\pgfsys@transformshift{0.959784in}{1.850278in}%
\pgfsys@useobject{currentmarker}{}%
\end{pgfscope}%
\begin{pgfscope}%
\pgfsys@transformshift{0.960122in}{1.850278in}%
\pgfsys@useobject{currentmarker}{}%
\end{pgfscope}%
\begin{pgfscope}%
\pgfsys@transformshift{0.960122in}{1.850278in}%
\pgfsys@useobject{currentmarker}{}%
\end{pgfscope}%
\begin{pgfscope}%
\pgfsys@transformshift{0.960122in}{1.850278in}%
\pgfsys@useobject{currentmarker}{}%
\end{pgfscope}%
\begin{pgfscope}%
\pgfsys@transformshift{0.960458in}{1.850278in}%
\pgfsys@useobject{currentmarker}{}%
\end{pgfscope}%
\begin{pgfscope}%
\pgfsys@transformshift{0.960796in}{1.850278in}%
\pgfsys@useobject{currentmarker}{}%
\end{pgfscope}%
\begin{pgfscope}%
\pgfsys@transformshift{0.961133in}{1.850278in}%
\pgfsys@useobject{currentmarker}{}%
\end{pgfscope}%
\begin{pgfscope}%
\pgfsys@transformshift{0.961471in}{1.850278in}%
\pgfsys@useobject{currentmarker}{}%
\end{pgfscope}%
\begin{pgfscope}%
\pgfsys@transformshift{0.961808in}{1.850278in}%
\pgfsys@useobject{currentmarker}{}%
\end{pgfscope}%
\begin{pgfscope}%
\pgfsys@transformshift{0.961808in}{1.850278in}%
\pgfsys@useobject{currentmarker}{}%
\end{pgfscope}%
\begin{pgfscope}%
\pgfsys@transformshift{0.962145in}{1.850278in}%
\pgfsys@useobject{currentmarker}{}%
\end{pgfscope}%
\begin{pgfscope}%
\pgfsys@transformshift{0.962145in}{1.850278in}%
\pgfsys@useobject{currentmarker}{}%
\end{pgfscope}%
\begin{pgfscope}%
\pgfsys@transformshift{0.962988in}{1.850278in}%
\pgfsys@useobject{currentmarker}{}%
\end{pgfscope}%
\begin{pgfscope}%
\pgfsys@transformshift{0.963494in}{1.850278in}%
\pgfsys@useobject{currentmarker}{}%
\end{pgfscope}%
\begin{pgfscope}%
\pgfsys@transformshift{0.975143in}{1.850278in}%
\pgfsys@useobject{currentmarker}{}%
\end{pgfscope}%
\begin{pgfscope}%
\pgfsys@transformshift{0.975474in}{1.850278in}%
\pgfsys@useobject{currentmarker}{}%
\end{pgfscope}%
\begin{pgfscope}%
\pgfsys@transformshift{0.975811in}{1.850278in}%
\pgfsys@useobject{currentmarker}{}%
\end{pgfscope}%
\begin{pgfscope}%
\pgfsys@transformshift{0.975818in}{1.850278in}%
\pgfsys@useobject{currentmarker}{}%
\end{pgfscope}%
\begin{pgfscope}%
\pgfsys@transformshift{0.976324in}{1.850278in}%
\pgfsys@useobject{currentmarker}{}%
\end{pgfscope}%
\begin{pgfscope}%
\pgfsys@transformshift{0.977490in}{1.850278in}%
\pgfsys@useobject{currentmarker}{}%
\end{pgfscope}%
\begin{pgfscope}%
\pgfsys@transformshift{0.977828in}{1.850278in}%
\pgfsys@useobject{currentmarker}{}%
\end{pgfscope}%
\begin{pgfscope}%
\pgfsys@transformshift{0.977842in}{1.850278in}%
\pgfsys@useobject{currentmarker}{}%
\end{pgfscope}%
\begin{pgfscope}%
\pgfsys@transformshift{0.978348in}{1.850278in}%
\pgfsys@useobject{currentmarker}{}%
\end{pgfscope}%
\begin{pgfscope}%
\pgfsys@transformshift{0.979191in}{1.850278in}%
\pgfsys@useobject{currentmarker}{}%
\end{pgfscope}%
\begin{pgfscope}%
\pgfsys@transformshift{0.979528in}{1.850278in}%
\pgfsys@useobject{currentmarker}{}%
\end{pgfscope}%
\begin{pgfscope}%
\pgfsys@transformshift{0.979535in}{1.850278in}%
\pgfsys@useobject{currentmarker}{}%
\end{pgfscope}%
\begin{pgfscope}%
\pgfsys@transformshift{0.980378in}{1.850278in}%
\pgfsys@useobject{currentmarker}{}%
\end{pgfscope}%
\begin{pgfscope}%
\pgfsys@transformshift{0.981221in}{1.850278in}%
\pgfsys@useobject{currentmarker}{}%
\end{pgfscope}%
\begin{pgfscope}%
\pgfsys@transformshift{0.981558in}{1.850278in}%
\pgfsys@useobject{currentmarker}{}%
\end{pgfscope}%
\begin{pgfscope}%
\pgfsys@transformshift{0.981713in}{1.850278in}%
\pgfsys@useobject{currentmarker}{}%
\end{pgfscope}%
\begin{pgfscope}%
\pgfsys@transformshift{0.981728in}{1.850278in}%
\pgfsys@useobject{currentmarker}{}%
\end{pgfscope}%
\begin{pgfscope}%
\pgfsys@transformshift{0.981896in}{1.850278in}%
\pgfsys@useobject{currentmarker}{}%
\end{pgfscope}%
\begin{pgfscope}%
\pgfsys@transformshift{0.982064in}{1.850278in}%
\pgfsys@useobject{currentmarker}{}%
\end{pgfscope}%
\begin{pgfscope}%
\pgfsys@transformshift{0.982233in}{1.850278in}%
\pgfsys@useobject{currentmarker}{}%
\end{pgfscope}%
\begin{pgfscope}%
\pgfsys@transformshift{0.982571in}{1.850278in}%
\pgfsys@useobject{currentmarker}{}%
\end{pgfscope}%
\begin{pgfscope}%
\pgfsys@transformshift{0.982739in}{1.850278in}%
\pgfsys@useobject{currentmarker}{}%
\end{pgfscope}%
\begin{pgfscope}%
\pgfsys@transformshift{0.982907in}{1.850278in}%
\pgfsys@useobject{currentmarker}{}%
\end{pgfscope}%
\begin{pgfscope}%
\pgfsys@transformshift{0.983076in}{1.850278in}%
\pgfsys@useobject{currentmarker}{}%
\end{pgfscope}%
\begin{pgfscope}%
\pgfsys@transformshift{0.983245in}{1.850278in}%
\pgfsys@useobject{currentmarker}{}%
\end{pgfscope}%
\begin{pgfscope}%
\pgfsys@transformshift{0.983414in}{1.850278in}%
\pgfsys@useobject{currentmarker}{}%
\end{pgfscope}%
\begin{pgfscope}%
\pgfsys@transformshift{0.983905in}{1.850278in}%
\pgfsys@useobject{currentmarker}{}%
\end{pgfscope}%
\begin{pgfscope}%
\pgfsys@transformshift{0.983920in}{1.850278in}%
\pgfsys@useobject{currentmarker}{}%
\end{pgfscope}%
\begin{pgfscope}%
\pgfsys@transformshift{0.984257in}{1.850278in}%
\pgfsys@useobject{currentmarker}{}%
\end{pgfscope}%
\begin{pgfscope}%
\pgfsys@transformshift{0.984594in}{1.850278in}%
\pgfsys@useobject{currentmarker}{}%
\end{pgfscope}%
\begin{pgfscope}%
\pgfsys@transformshift{0.984931in}{1.850278in}%
\pgfsys@useobject{currentmarker}{}%
\end{pgfscope}%
\begin{pgfscope}%
\pgfsys@transformshift{0.985268in}{1.850278in}%
\pgfsys@useobject{currentmarker}{}%
\end{pgfscope}%
\begin{pgfscope}%
\pgfsys@transformshift{0.985268in}{1.850278in}%
\pgfsys@useobject{currentmarker}{}%
\end{pgfscope}%
\begin{pgfscope}%
\pgfsys@transformshift{0.985268in}{1.850278in}%
\pgfsys@useobject{currentmarker}{}%
\end{pgfscope}%
\begin{pgfscope}%
\pgfsys@transformshift{0.985606in}{1.850278in}%
\pgfsys@useobject{currentmarker}{}%
\end{pgfscope}%
\begin{pgfscope}%
\pgfsys@transformshift{0.985606in}{1.850278in}%
\pgfsys@useobject{currentmarker}{}%
\end{pgfscope}%
\begin{pgfscope}%
\pgfsys@transformshift{0.985606in}{1.850278in}%
\pgfsys@useobject{currentmarker}{}%
\end{pgfscope}%
\begin{pgfscope}%
\pgfsys@transformshift{0.985943in}{1.850278in}%
\pgfsys@useobject{currentmarker}{}%
\end{pgfscope}%
\begin{pgfscope}%
\pgfsys@transformshift{0.985943in}{1.850278in}%
\pgfsys@useobject{currentmarker}{}%
\end{pgfscope}%
\begin{pgfscope}%
\pgfsys@transformshift{0.986280in}{1.850278in}%
\pgfsys@useobject{currentmarker}{}%
\end{pgfscope}%
\begin{pgfscope}%
\pgfsys@transformshift{0.986280in}{1.850278in}%
\pgfsys@useobject{currentmarker}{}%
\end{pgfscope}%
\begin{pgfscope}%
\pgfsys@transformshift{0.986280in}{1.850278in}%
\pgfsys@useobject{currentmarker}{}%
\end{pgfscope}%
\begin{pgfscope}%
\pgfsys@transformshift{0.986617in}{1.850278in}%
\pgfsys@useobject{currentmarker}{}%
\end{pgfscope}%
\begin{pgfscope}%
\pgfsys@transformshift{0.986617in}{1.850278in}%
\pgfsys@useobject{currentmarker}{}%
\end{pgfscope}%
\begin{pgfscope}%
\pgfsys@transformshift{0.986617in}{1.850278in}%
\pgfsys@useobject{currentmarker}{}%
\end{pgfscope}%
\begin{pgfscope}%
\pgfsys@transformshift{0.986955in}{1.850278in}%
\pgfsys@useobject{currentmarker}{}%
\end{pgfscope}%
\begin{pgfscope}%
\pgfsys@transformshift{0.986955in}{1.850278in}%
\pgfsys@useobject{currentmarker}{}%
\end{pgfscope}%
\begin{pgfscope}%
\pgfsys@transformshift{0.986955in}{1.850278in}%
\pgfsys@useobject{currentmarker}{}%
\end{pgfscope}%
\begin{pgfscope}%
\pgfsys@transformshift{0.987292in}{1.850278in}%
\pgfsys@useobject{currentmarker}{}%
\end{pgfscope}%
\begin{pgfscope}%
\pgfsys@transformshift{0.987630in}{1.850278in}%
\pgfsys@useobject{currentmarker}{}%
\end{pgfscope}%
\begin{pgfscope}%
\pgfsys@transformshift{0.987966in}{1.850278in}%
\pgfsys@useobject{currentmarker}{}%
\end{pgfscope}%
\begin{pgfscope}%
\pgfsys@transformshift{0.988135in}{1.850278in}%
\pgfsys@useobject{currentmarker}{}%
\end{pgfscope}%
\begin{pgfscope}%
\pgfsys@transformshift{0.988304in}{1.850278in}%
\pgfsys@useobject{currentmarker}{}%
\end{pgfscope}%
\begin{pgfscope}%
\pgfsys@transformshift{0.988641in}{1.850278in}%
\pgfsys@useobject{currentmarker}{}%
\end{pgfscope}%
\begin{pgfscope}%
\pgfsys@transformshift{0.988964in}{1.850278in}%
\pgfsys@useobject{currentmarker}{}%
\end{pgfscope}%
\begin{pgfscope}%
\pgfsys@transformshift{0.989147in}{1.850278in}%
\pgfsys@useobject{currentmarker}{}%
\end{pgfscope}%
\begin{pgfscope}%
\pgfsys@transformshift{0.989316in}{1.850278in}%
\pgfsys@useobject{currentmarker}{}%
\end{pgfscope}%
\begin{pgfscope}%
\pgfsys@transformshift{0.989652in}{1.850278in}%
\pgfsys@useobject{currentmarker}{}%
\end{pgfscope}%
\begin{pgfscope}%
\pgfsys@transformshift{0.989822in}{1.850278in}%
\pgfsys@useobject{currentmarker}{}%
\end{pgfscope}%
\begin{pgfscope}%
\pgfsys@transformshift{0.989990in}{1.850278in}%
\pgfsys@useobject{currentmarker}{}%
\end{pgfscope}%
\begin{pgfscope}%
\pgfsys@transformshift{0.990327in}{1.850278in}%
\pgfsys@useobject{currentmarker}{}%
\end{pgfscope}%
\begin{pgfscope}%
\pgfsys@transformshift{0.990496in}{1.850278in}%
\pgfsys@useobject{currentmarker}{}%
\end{pgfscope}%
\begin{pgfscope}%
\pgfsys@transformshift{0.990826in}{1.850278in}%
\pgfsys@useobject{currentmarker}{}%
\end{pgfscope}%
\begin{pgfscope}%
\pgfsys@transformshift{0.991170in}{1.850278in}%
\pgfsys@useobject{currentmarker}{}%
\end{pgfscope}%
\begin{pgfscope}%
\pgfsys@transformshift{0.991508in}{1.850278in}%
\pgfsys@useobject{currentmarker}{}%
\end{pgfscope}%
\begin{pgfscope}%
\pgfsys@transformshift{0.991830in}{1.850278in}%
\pgfsys@useobject{currentmarker}{}%
\end{pgfscope}%
\begin{pgfscope}%
\pgfsys@transformshift{0.992168in}{1.850278in}%
\pgfsys@useobject{currentmarker}{}%
\end{pgfscope}%
\begin{pgfscope}%
\pgfsys@transformshift{0.992505in}{1.850278in}%
\pgfsys@useobject{currentmarker}{}%
\end{pgfscope}%
\begin{pgfscope}%
\pgfsys@transformshift{0.992843in}{1.850278in}%
\pgfsys@useobject{currentmarker}{}%
\end{pgfscope}%
\begin{pgfscope}%
\pgfsys@transformshift{0.993180in}{1.850278in}%
\pgfsys@useobject{currentmarker}{}%
\end{pgfscope}%
\begin{pgfscope}%
\pgfsys@transformshift{0.993362in}{1.850278in}%
\pgfsys@useobject{currentmarker}{}%
\end{pgfscope}%
\begin{pgfscope}%
\pgfsys@transformshift{0.993868in}{1.850278in}%
\pgfsys@useobject{currentmarker}{}%
\end{pgfscope}%
\begin{pgfscope}%
\pgfsys@transformshift{0.994199in}{1.850278in}%
\pgfsys@useobject{currentmarker}{}%
\end{pgfscope}%
\begin{pgfscope}%
\pgfsys@transformshift{0.994375in}{1.850278in}%
\pgfsys@useobject{currentmarker}{}%
\end{pgfscope}%
\begin{pgfscope}%
\pgfsys@transformshift{0.994711in}{1.850278in}%
\pgfsys@useobject{currentmarker}{}%
\end{pgfscope}%
\begin{pgfscope}%
\pgfsys@transformshift{0.995049in}{1.850278in}%
\pgfsys@useobject{currentmarker}{}%
\end{pgfscope}%
\begin{pgfscope}%
\pgfsys@transformshift{0.995386in}{1.850278in}%
\pgfsys@useobject{currentmarker}{}%
\end{pgfscope}%
\begin{pgfscope}%
\pgfsys@transformshift{0.995724in}{1.850278in}%
\pgfsys@useobject{currentmarker}{}%
\end{pgfscope}%
\begin{pgfscope}%
\pgfsys@transformshift{0.996061in}{1.850278in}%
\pgfsys@useobject{currentmarker}{}%
\end{pgfscope}%
\begin{pgfscope}%
\pgfsys@transformshift{0.996567in}{1.850278in}%
\pgfsys@useobject{currentmarker}{}%
\end{pgfscope}%
\begin{pgfscope}%
\pgfsys@transformshift{0.996735in}{1.850278in}%
\pgfsys@useobject{currentmarker}{}%
\end{pgfscope}%
\begin{pgfscope}%
\pgfsys@transformshift{0.996904in}{1.850278in}%
\pgfsys@useobject{currentmarker}{}%
\end{pgfscope}%
\begin{pgfscope}%
\pgfsys@transformshift{0.997072in}{1.850278in}%
\pgfsys@useobject{currentmarker}{}%
\end{pgfscope}%
\begin{pgfscope}%
\pgfsys@transformshift{0.997410in}{1.850278in}%
\pgfsys@useobject{currentmarker}{}%
\end{pgfscope}%
\begin{pgfscope}%
\pgfsys@transformshift{0.997747in}{1.850278in}%
\pgfsys@useobject{currentmarker}{}%
\end{pgfscope}%
\begin{pgfscope}%
\pgfsys@transformshift{0.998084in}{1.850278in}%
\pgfsys@useobject{currentmarker}{}%
\end{pgfscope}%
\begin{pgfscope}%
\pgfsys@transformshift{0.998421in}{1.850278in}%
\pgfsys@useobject{currentmarker}{}%
\end{pgfscope}%
\begin{pgfscope}%
\pgfsys@transformshift{0.998759in}{1.850278in}%
\pgfsys@useobject{currentmarker}{}%
\end{pgfscope}%
\begin{pgfscope}%
\pgfsys@transformshift{0.998920in}{1.850278in}%
\pgfsys@useobject{currentmarker}{}%
\end{pgfscope}%
\begin{pgfscope}%
\pgfsys@transformshift{0.999096in}{1.850278in}%
\pgfsys@useobject{currentmarker}{}%
\end{pgfscope}%
\begin{pgfscope}%
\pgfsys@transformshift{0.999594in}{1.850278in}%
\pgfsys@useobject{currentmarker}{}%
\end{pgfscope}%
\begin{pgfscope}%
\pgfsys@transformshift{0.999939in}{1.850278in}%
\pgfsys@useobject{currentmarker}{}%
\end{pgfscope}%
\begin{pgfscope}%
\pgfsys@transformshift{1.000108in}{1.850278in}%
\pgfsys@useobject{currentmarker}{}%
\end{pgfscope}%
\begin{pgfscope}%
\pgfsys@transformshift{1.000277in}{1.850278in}%
\pgfsys@useobject{currentmarker}{}%
\end{pgfscope}%
\begin{pgfscope}%
\pgfsys@transformshift{1.000782in}{1.850278in}%
\pgfsys@useobject{currentmarker}{}%
\end{pgfscope}%
\begin{pgfscope}%
\pgfsys@transformshift{1.001626in}{1.850278in}%
\pgfsys@useobject{currentmarker}{}%
\end{pgfscope}%
\begin{pgfscope}%
\pgfsys@transformshift{1.003473in}{1.850278in}%
\pgfsys@useobject{currentmarker}{}%
\end{pgfscope}%
\begin{pgfscope}%
\pgfsys@transformshift{1.003810in}{1.850278in}%
\pgfsys@useobject{currentmarker}{}%
\end{pgfscope}%
\begin{pgfscope}%
\pgfsys@transformshift{1.004822in}{1.850278in}%
\pgfsys@useobject{currentmarker}{}%
\end{pgfscope}%
\begin{pgfscope}%
\pgfsys@transformshift{1.006003in}{1.850278in}%
\pgfsys@useobject{currentmarker}{}%
\end{pgfscope}%
\begin{pgfscope}%
\pgfsys@transformshift{1.132139in}{1.850278in}%
\pgfsys@useobject{currentmarker}{}%
\end{pgfscope}%
\begin{pgfscope}%
\pgfsys@transformshift{1.137366in}{1.850278in}%
\pgfsys@useobject{currentmarker}{}%
\end{pgfscope}%
\begin{pgfscope}%
\pgfsys@transformshift{1.141413in}{1.850278in}%
\pgfsys@useobject{currentmarker}{}%
\end{pgfscope}%
\begin{pgfscope}%
\pgfsys@transformshift{1.143605in}{1.850278in}%
\pgfsys@useobject{currentmarker}{}%
\end{pgfscope}%
\begin{pgfscope}%
\pgfsys@transformshift{1.144954in}{1.850278in}%
\pgfsys@useobject{currentmarker}{}%
\end{pgfscope}%
\begin{pgfscope}%
\pgfsys@transformshift{1.146304in}{1.850278in}%
\pgfsys@useobject{currentmarker}{}%
\end{pgfscope}%
\begin{pgfscope}%
\pgfsys@transformshift{1.147484in}{1.850278in}%
\pgfsys@useobject{currentmarker}{}%
\end{pgfscope}%
\begin{pgfscope}%
\pgfsys@transformshift{1.148159in}{1.850278in}%
\pgfsys@useobject{currentmarker}{}%
\end{pgfscope}%
\begin{pgfscope}%
\pgfsys@transformshift{1.148833in}{1.850278in}%
\pgfsys@useobject{currentmarker}{}%
\end{pgfscope}%
\begin{pgfscope}%
\pgfsys@transformshift{1.149677in}{1.850278in}%
\pgfsys@useobject{currentmarker}{}%
\end{pgfscope}%
\begin{pgfscope}%
\pgfsys@transformshift{1.150182in}{1.850278in}%
\pgfsys@useobject{currentmarker}{}%
\end{pgfscope}%
\begin{pgfscope}%
\pgfsys@transformshift{1.150520in}{1.850278in}%
\pgfsys@useobject{currentmarker}{}%
\end{pgfscope}%
\begin{pgfscope}%
\pgfsys@transformshift{1.151699in}{1.850278in}%
\pgfsys@useobject{currentmarker}{}%
\end{pgfscope}%
\begin{pgfscope}%
\pgfsys@transformshift{1.152037in}{1.850278in}%
\pgfsys@useobject{currentmarker}{}%
\end{pgfscope}%
\begin{pgfscope}%
\pgfsys@transformshift{1.152374in}{1.850278in}%
\pgfsys@useobject{currentmarker}{}%
\end{pgfscope}%
\begin{pgfscope}%
\pgfsys@transformshift{1.153049in}{1.850278in}%
\pgfsys@useobject{currentmarker}{}%
\end{pgfscope}%
\begin{pgfscope}%
\pgfsys@transformshift{1.153386in}{1.850278in}%
\pgfsys@useobject{currentmarker}{}%
\end{pgfscope}%
\begin{pgfscope}%
\pgfsys@transformshift{1.153892in}{1.850278in}%
\pgfsys@useobject{currentmarker}{}%
\end{pgfscope}%
\begin{pgfscope}%
\pgfsys@transformshift{1.154398in}{1.850278in}%
\pgfsys@useobject{currentmarker}{}%
\end{pgfscope}%
\begin{pgfscope}%
\pgfsys@transformshift{1.154904in}{1.850278in}%
\pgfsys@useobject{currentmarker}{}%
\end{pgfscope}%
\begin{pgfscope}%
\pgfsys@transformshift{1.155241in}{1.850278in}%
\pgfsys@useobject{currentmarker}{}%
\end{pgfscope}%
\begin{pgfscope}%
\pgfsys@transformshift{1.155409in}{1.850278in}%
\pgfsys@useobject{currentmarker}{}%
\end{pgfscope}%
\begin{pgfscope}%
\pgfsys@transformshift{1.156084in}{1.850278in}%
\pgfsys@useobject{currentmarker}{}%
\end{pgfscope}%
\begin{pgfscope}%
\pgfsys@transformshift{1.156422in}{1.850278in}%
\pgfsys@useobject{currentmarker}{}%
\end{pgfscope}%
\begin{pgfscope}%
\pgfsys@transformshift{1.156927in}{1.850278in}%
\pgfsys@useobject{currentmarker}{}%
\end{pgfscope}%
\begin{pgfscope}%
\pgfsys@transformshift{1.157265in}{1.850278in}%
\pgfsys@useobject{currentmarker}{}%
\end{pgfscope}%
\begin{pgfscope}%
\pgfsys@transformshift{1.157601in}{1.850278in}%
\pgfsys@useobject{currentmarker}{}%
\end{pgfscope}%
\begin{pgfscope}%
\pgfsys@transformshift{1.158276in}{1.850278in}%
\pgfsys@useobject{currentmarker}{}%
\end{pgfscope}%
\begin{pgfscope}%
\pgfsys@transformshift{1.158614in}{1.850278in}%
\pgfsys@useobject{currentmarker}{}%
\end{pgfscope}%
\begin{pgfscope}%
\pgfsys@transformshift{1.158951in}{1.850278in}%
\pgfsys@useobject{currentmarker}{}%
\end{pgfscope}%
\begin{pgfscope}%
\pgfsys@transformshift{1.159127in}{1.850278in}%
\pgfsys@useobject{currentmarker}{}%
\end{pgfscope}%
\begin{pgfscope}%
\pgfsys@transformshift{1.159457in}{1.850278in}%
\pgfsys@useobject{currentmarker}{}%
\end{pgfscope}%
\begin{pgfscope}%
\pgfsys@transformshift{1.159794in}{1.850278in}%
\pgfsys@useobject{currentmarker}{}%
\end{pgfscope}%
\begin{pgfscope}%
\pgfsys@transformshift{1.160138in}{1.850278in}%
\pgfsys@useobject{currentmarker}{}%
\end{pgfscope}%
\begin{pgfscope}%
\pgfsys@transformshift{1.160637in}{1.850278in}%
\pgfsys@useobject{currentmarker}{}%
\end{pgfscope}%
\begin{pgfscope}%
\pgfsys@transformshift{1.160974in}{1.850278in}%
\pgfsys@useobject{currentmarker}{}%
\end{pgfscope}%
\begin{pgfscope}%
\pgfsys@transformshift{1.161150in}{1.850278in}%
\pgfsys@useobject{currentmarker}{}%
\end{pgfscope}%
\begin{pgfscope}%
\pgfsys@transformshift{1.161487in}{1.850278in}%
\pgfsys@useobject{currentmarker}{}%
\end{pgfscope}%
\begin{pgfscope}%
\pgfsys@transformshift{1.161986in}{1.850278in}%
\pgfsys@useobject{currentmarker}{}%
\end{pgfscope}%
\begin{pgfscope}%
\pgfsys@transformshift{1.162324in}{1.850278in}%
\pgfsys@useobject{currentmarker}{}%
\end{pgfscope}%
\begin{pgfscope}%
\pgfsys@transformshift{1.162492in}{1.850278in}%
\pgfsys@useobject{currentmarker}{}%
\end{pgfscope}%
\begin{pgfscope}%
\pgfsys@transformshift{1.162829in}{1.850278in}%
\pgfsys@useobject{currentmarker}{}%
\end{pgfscope}%
\begin{pgfscope}%
\pgfsys@transformshift{1.163167in}{1.850278in}%
\pgfsys@useobject{currentmarker}{}%
\end{pgfscope}%
\begin{pgfscope}%
\pgfsys@transformshift{1.163335in}{1.850278in}%
\pgfsys@useobject{currentmarker}{}%
\end{pgfscope}%
\begin{pgfscope}%
\pgfsys@transformshift{1.163673in}{1.850278in}%
\pgfsys@useobject{currentmarker}{}%
\end{pgfscope}%
\begin{pgfscope}%
\pgfsys@transformshift{1.164178in}{1.850278in}%
\pgfsys@useobject{currentmarker}{}%
\end{pgfscope}%
\begin{pgfscope}%
\pgfsys@transformshift{1.164347in}{1.850278in}%
\pgfsys@useobject{currentmarker}{}%
\end{pgfscope}%
\begin{pgfscope}%
\pgfsys@transformshift{1.164684in}{1.850278in}%
\pgfsys@useobject{currentmarker}{}%
\end{pgfscope}%
\begin{pgfscope}%
\pgfsys@transformshift{1.165021in}{1.850278in}%
\pgfsys@useobject{currentmarker}{}%
\end{pgfscope}%
\begin{pgfscope}%
\pgfsys@transformshift{1.165197in}{1.850278in}%
\pgfsys@useobject{currentmarker}{}%
\end{pgfscope}%
\begin{pgfscope}%
\pgfsys@transformshift{1.165359in}{1.850278in}%
\pgfsys@useobject{currentmarker}{}%
\end{pgfscope}%
\begin{pgfscope}%
\pgfsys@transformshift{1.165696in}{1.850278in}%
\pgfsys@useobject{currentmarker}{}%
\end{pgfscope}%
\begin{pgfscope}%
\pgfsys@transformshift{1.166033in}{1.850278in}%
\pgfsys@useobject{currentmarker}{}%
\end{pgfscope}%
\begin{pgfscope}%
\pgfsys@transformshift{1.166202in}{1.850278in}%
\pgfsys@useobject{currentmarker}{}%
\end{pgfscope}%
\begin{pgfscope}%
\pgfsys@transformshift{1.166708in}{1.850278in}%
\pgfsys@useobject{currentmarker}{}%
\end{pgfscope}%
\begin{pgfscope}%
\pgfsys@transformshift{1.167213in}{1.850278in}%
\pgfsys@useobject{currentmarker}{}%
\end{pgfscope}%
\begin{pgfscope}%
\pgfsys@transformshift{1.167551in}{1.850278in}%
\pgfsys@useobject{currentmarker}{}%
\end{pgfscope}%
\begin{pgfscope}%
\pgfsys@transformshift{1.167888in}{1.850278in}%
\pgfsys@useobject{currentmarker}{}%
\end{pgfscope}%
\begin{pgfscope}%
\pgfsys@transformshift{1.168226in}{1.850278in}%
\pgfsys@useobject{currentmarker}{}%
\end{pgfscope}%
\begin{pgfscope}%
\pgfsys@transformshift{1.168226in}{1.850278in}%
\pgfsys@useobject{currentmarker}{}%
\end{pgfscope}%
\begin{pgfscope}%
\pgfsys@transformshift{1.168562in}{1.850278in}%
\pgfsys@useobject{currentmarker}{}%
\end{pgfscope}%
\begin{pgfscope}%
\pgfsys@transformshift{1.168562in}{1.850278in}%
\pgfsys@useobject{currentmarker}{}%
\end{pgfscope}%
\begin{pgfscope}%
\pgfsys@transformshift{1.169069in}{1.850278in}%
\pgfsys@useobject{currentmarker}{}%
\end{pgfscope}%
\begin{pgfscope}%
\pgfsys@transformshift{1.169405in}{1.850278in}%
\pgfsys@useobject{currentmarker}{}%
\end{pgfscope}%
\begin{pgfscope}%
\pgfsys@transformshift{1.169575in}{1.850278in}%
\pgfsys@useobject{currentmarker}{}%
\end{pgfscope}%
\begin{pgfscope}%
\pgfsys@transformshift{1.169743in}{1.850278in}%
\pgfsys@useobject{currentmarker}{}%
\end{pgfscope}%
\begin{pgfscope}%
\pgfsys@transformshift{1.169912in}{1.850278in}%
\pgfsys@useobject{currentmarker}{}%
\end{pgfscope}%
\begin{pgfscope}%
\pgfsys@transformshift{1.170418in}{1.850278in}%
\pgfsys@useobject{currentmarker}{}%
\end{pgfscope}%
\begin{pgfscope}%
\pgfsys@transformshift{1.170762in}{1.850278in}%
\pgfsys@useobject{currentmarker}{}%
\end{pgfscope}%
\begin{pgfscope}%
\pgfsys@transformshift{1.171099in}{1.850278in}%
\pgfsys@useobject{currentmarker}{}%
\end{pgfscope}%
\begin{pgfscope}%
\pgfsys@transformshift{1.171429in}{1.850278in}%
\pgfsys@useobject{currentmarker}{}%
\end{pgfscope}%
\begin{pgfscope}%
\pgfsys@transformshift{1.171774in}{1.850278in}%
\pgfsys@useobject{currentmarker}{}%
\end{pgfscope}%
\begin{pgfscope}%
\pgfsys@transformshift{1.172272in}{1.850278in}%
\pgfsys@useobject{currentmarker}{}%
\end{pgfscope}%
\begin{pgfscope}%
\pgfsys@transformshift{1.172448in}{1.850278in}%
\pgfsys@useobject{currentmarker}{}%
\end{pgfscope}%
\begin{pgfscope}%
\pgfsys@transformshift{1.172785in}{1.850278in}%
\pgfsys@useobject{currentmarker}{}%
\end{pgfscope}%
\begin{pgfscope}%
\pgfsys@transformshift{1.173123in}{1.850278in}%
\pgfsys@useobject{currentmarker}{}%
\end{pgfscope}%
\begin{pgfscope}%
\pgfsys@transformshift{1.173460in}{1.850278in}%
\pgfsys@useobject{currentmarker}{}%
\end{pgfscope}%
\begin{pgfscope}%
\pgfsys@transformshift{1.173510in}{1.850278in}%
\pgfsys@useobject{currentmarker}{}%
\end{pgfscope}%
\begin{pgfscope}%
\pgfsys@transformshift{1.173797in}{1.850278in}%
\pgfsys@useobject{currentmarker}{}%
\end{pgfscope}%
\begin{pgfscope}%
\pgfsys@transformshift{1.173959in}{1.850278in}%
\pgfsys@useobject{currentmarker}{}%
\end{pgfscope}%
\begin{pgfscope}%
\pgfsys@transformshift{1.174134in}{1.850278in}%
\pgfsys@useobject{currentmarker}{}%
\end{pgfscope}%
\begin{pgfscope}%
\pgfsys@transformshift{1.174472in}{1.850278in}%
\pgfsys@useobject{currentmarker}{}%
\end{pgfscope}%
\begin{pgfscope}%
\pgfsys@transformshift{1.174809in}{1.850278in}%
\pgfsys@useobject{currentmarker}{}%
\end{pgfscope}%
\begin{pgfscope}%
\pgfsys@transformshift{1.175147in}{1.850278in}%
\pgfsys@useobject{currentmarker}{}%
\end{pgfscope}%
\begin{pgfscope}%
\pgfsys@transformshift{1.175315in}{1.850278in}%
\pgfsys@useobject{currentmarker}{}%
\end{pgfscope}%
\begin{pgfscope}%
\pgfsys@transformshift{1.175491in}{1.850278in}%
\pgfsys@useobject{currentmarker}{}%
\end{pgfscope}%
\begin{pgfscope}%
\pgfsys@transformshift{1.175821in}{1.850278in}%
\pgfsys@useobject{currentmarker}{}%
\end{pgfscope}%
\begin{pgfscope}%
\pgfsys@transformshift{1.176158in}{1.850278in}%
\pgfsys@useobject{currentmarker}{}%
\end{pgfscope}%
\begin{pgfscope}%
\pgfsys@transformshift{1.176502in}{1.850278in}%
\pgfsys@useobject{currentmarker}{}%
\end{pgfscope}%
\begin{pgfscope}%
\pgfsys@transformshift{1.176840in}{1.850278in}%
\pgfsys@useobject{currentmarker}{}%
\end{pgfscope}%
\begin{pgfscope}%
\pgfsys@transformshift{1.177177in}{1.850278in}%
\pgfsys@useobject{currentmarker}{}%
\end{pgfscope}%
\begin{pgfscope}%
\pgfsys@transformshift{1.177514in}{1.850278in}%
\pgfsys@useobject{currentmarker}{}%
\end{pgfscope}%
\begin{pgfscope}%
\pgfsys@transformshift{1.178343in}{1.850278in}%
\pgfsys@useobject{currentmarker}{}%
\end{pgfscope}%
\begin{pgfscope}%
\pgfsys@transformshift{1.178863in}{1.850278in}%
\pgfsys@useobject{currentmarker}{}%
\end{pgfscope}%
\begin{pgfscope}%
\pgfsys@transformshift{1.179861in}{1.850278in}%
\pgfsys@useobject{currentmarker}{}%
\end{pgfscope}%
\begin{pgfscope}%
\pgfsys@transformshift{1.180036in}{1.850278in}%
\pgfsys@useobject{currentmarker}{}%
\end{pgfscope}%
\begin{pgfscope}%
\pgfsys@transformshift{1.180374in}{1.850278in}%
\pgfsys@useobject{currentmarker}{}%
\end{pgfscope}%
\begin{pgfscope}%
\pgfsys@transformshift{1.181773in}{1.850278in}%
\pgfsys@useobject{currentmarker}{}%
\end{pgfscope}%
\begin{pgfscope}%
\pgfsys@transformshift{1.182228in}{1.850278in}%
\pgfsys@useobject{currentmarker}{}%
\end{pgfscope}%
\begin{pgfscope}%
\pgfsys@transformshift{1.182735in}{1.850278in}%
\pgfsys@useobject{currentmarker}{}%
\end{pgfscope}%
\begin{pgfscope}%
\pgfsys@transformshift{1.183071in}{1.850278in}%
\pgfsys@useobject{currentmarker}{}%
\end{pgfscope}%
\begin{pgfscope}%
\pgfsys@transformshift{1.183416in}{1.850278in}%
\pgfsys@useobject{currentmarker}{}%
\end{pgfscope}%
\begin{pgfscope}%
\pgfsys@transformshift{1.183746in}{1.850278in}%
\pgfsys@useobject{currentmarker}{}%
\end{pgfscope}%
\begin{pgfscope}%
\pgfsys@transformshift{1.183930in}{1.850278in}%
\pgfsys@useobject{currentmarker}{}%
\end{pgfscope}%
\begin{pgfscope}%
\pgfsys@transformshift{1.184435in}{1.850278in}%
\pgfsys@useobject{currentmarker}{}%
\end{pgfscope}%
\begin{pgfscope}%
\pgfsys@transformshift{1.184773in}{1.850278in}%
\pgfsys@useobject{currentmarker}{}%
\end{pgfscope}%
\begin{pgfscope}%
\pgfsys@transformshift{1.185109in}{1.850278in}%
\pgfsys@useobject{currentmarker}{}%
\end{pgfscope}%
\begin{pgfscope}%
\pgfsys@transformshift{1.185770in}{1.850278in}%
\pgfsys@useobject{currentmarker}{}%
\end{pgfscope}%
\begin{pgfscope}%
\pgfsys@transformshift{1.186122in}{1.850278in}%
\pgfsys@useobject{currentmarker}{}%
\end{pgfscope}%
\begin{pgfscope}%
\pgfsys@transformshift{1.192178in}{1.850278in}%
\pgfsys@useobject{currentmarker}{}%
\end{pgfscope}%
\begin{pgfscope}%
\pgfsys@transformshift{1.192515in}{1.850278in}%
\pgfsys@useobject{currentmarker}{}%
\end{pgfscope}%
\begin{pgfscope}%
\pgfsys@transformshift{1.193197in}{1.850278in}%
\pgfsys@useobject{currentmarker}{}%
\end{pgfscope}%
\begin{pgfscope}%
\pgfsys@transformshift{1.195564in}{1.850278in}%
\pgfsys@useobject{currentmarker}{}%
\end{pgfscope}%
\begin{pgfscope}%
\pgfsys@transformshift{1.196408in}{1.850278in}%
\pgfsys@useobject{currentmarker}{}%
\end{pgfscope}%
\begin{pgfscope}%
\pgfsys@transformshift{1.198263in}{1.850278in}%
\pgfsys@useobject{currentmarker}{}%
\end{pgfscope}%
\begin{pgfscope}%
\pgfsys@transformshift{1.199443in}{1.850278in}%
\pgfsys@useobject{currentmarker}{}%
\end{pgfscope}%
\begin{pgfscope}%
\pgfsys@transformshift{1.200953in}{1.850278in}%
\pgfsys@useobject{currentmarker}{}%
\end{pgfscope}%
\begin{pgfscope}%
\pgfsys@transformshift{1.202310in}{1.850278in}%
\pgfsys@useobject{currentmarker}{}%
\end{pgfscope}%
\begin{pgfscope}%
\pgfsys@transformshift{1.203658in}{1.850278in}%
\pgfsys@useobject{currentmarker}{}%
\end{pgfscope}%
\begin{pgfscope}%
\pgfsys@transformshift{1.204502in}{1.850278in}%
\pgfsys@useobject{currentmarker}{}%
\end{pgfscope}%
\begin{pgfscope}%
\pgfsys@transformshift{1.206188in}{1.850278in}%
\pgfsys@useobject{currentmarker}{}%
\end{pgfscope}%
\begin{pgfscope}%
\pgfsys@transformshift{1.457086in}{1.850278in}%
\pgfsys@useobject{currentmarker}{}%
\end{pgfscope}%
\begin{pgfscope}%
\pgfsys@transformshift{1.471082in}{1.850278in}%
\pgfsys@useobject{currentmarker}{}%
\end{pgfscope}%
\begin{pgfscope}%
\pgfsys@transformshift{1.481368in}{1.850278in}%
\pgfsys@useobject{currentmarker}{}%
\end{pgfscope}%
\begin{pgfscope}%
\pgfsys@transformshift{1.488956in}{1.850278in}%
\pgfsys@useobject{currentmarker}{}%
\end{pgfscope}%
\begin{pgfscope}%
\pgfsys@transformshift{1.494697in}{1.850278in}%
\pgfsys@useobject{currentmarker}{}%
\end{pgfscope}%
\begin{pgfscope}%
\pgfsys@transformshift{1.499756in}{1.850278in}%
\pgfsys@useobject{currentmarker}{}%
\end{pgfscope}%
\begin{pgfscope}%
\pgfsys@transformshift{1.505145in}{1.850278in}%
\pgfsys@useobject{currentmarker}{}%
\end{pgfscope}%
\begin{pgfscope}%
\pgfsys@transformshift{1.511047in}{1.850278in}%
\pgfsys@useobject{currentmarker}{}%
\end{pgfscope}%
\begin{pgfscope}%
\pgfsys@transformshift{1.514250in}{1.850278in}%
\pgfsys@useobject{currentmarker}{}%
\end{pgfscope}%
\begin{pgfscope}%
\pgfsys@transformshift{1.516787in}{1.850278in}%
\pgfsys@useobject{currentmarker}{}%
\end{pgfscope}%
\begin{pgfscope}%
\pgfsys@transformshift{1.519478in}{1.850278in}%
\pgfsys@useobject{currentmarker}{}%
\end{pgfscope}%
\begin{pgfscope}%
\pgfsys@transformshift{1.521678in}{1.850278in}%
\pgfsys@useobject{currentmarker}{}%
\end{pgfscope}%
\begin{pgfscope}%
\pgfsys@transformshift{1.523694in}{1.850278in}%
\pgfsys@useobject{currentmarker}{}%
\end{pgfscope}%
\begin{pgfscope}%
\pgfsys@transformshift{1.525387in}{1.850278in}%
\pgfsys@useobject{currentmarker}{}%
\end{pgfscope}%
\begin{pgfscope}%
\pgfsys@transformshift{1.526737in}{1.850278in}%
\pgfsys@useobject{currentmarker}{}%
\end{pgfscope}%
\begin{pgfscope}%
\pgfsys@transformshift{1.528415in}{1.850278in}%
\pgfsys@useobject{currentmarker}{}%
\end{pgfscope}%
\begin{pgfscope}%
\pgfsys@transformshift{1.530109in}{1.850278in}%
\pgfsys@useobject{currentmarker}{}%
\end{pgfscope}%
\begin{pgfscope}%
\pgfsys@transformshift{1.531458in}{1.850278in}%
\pgfsys@useobject{currentmarker}{}%
\end{pgfscope}%
\begin{pgfscope}%
\pgfsys@transformshift{1.533144in}{1.850278in}%
\pgfsys@useobject{currentmarker}{}%
\end{pgfscope}%
\begin{pgfscope}%
\pgfsys@transformshift{1.534156in}{1.850278in}%
\pgfsys@useobject{currentmarker}{}%
\end{pgfscope}%
\begin{pgfscope}%
\pgfsys@transformshift{1.534831in}{1.850278in}%
\pgfsys@useobject{currentmarker}{}%
\end{pgfscope}%
\begin{pgfscope}%
\pgfsys@transformshift{1.535666in}{1.850278in}%
\pgfsys@useobject{currentmarker}{}%
\end{pgfscope}%
\begin{pgfscope}%
\pgfsys@transformshift{1.536678in}{1.850278in}%
\pgfsys@useobject{currentmarker}{}%
\end{pgfscope}%
\begin{pgfscope}%
\pgfsys@transformshift{1.537866in}{1.850278in}%
\pgfsys@useobject{currentmarker}{}%
\end{pgfscope}%
\begin{pgfscope}%
\pgfsys@transformshift{1.538877in}{1.850278in}%
\pgfsys@useobject{currentmarker}{}%
\end{pgfscope}%
\begin{pgfscope}%
\pgfsys@transformshift{1.540388in}{1.850278in}%
\pgfsys@useobject{currentmarker}{}%
\end{pgfscope}%
\begin{pgfscope}%
\pgfsys@transformshift{1.541231in}{1.850278in}%
\pgfsys@useobject{currentmarker}{}%
\end{pgfscope}%
\begin{pgfscope}%
\pgfsys@transformshift{1.542763in}{1.850278in}%
\pgfsys@useobject{currentmarker}{}%
\end{pgfscope}%
\begin{pgfscope}%
\pgfsys@transformshift{1.543936in}{1.850278in}%
\pgfsys@useobject{currentmarker}{}%
\end{pgfscope}%
\begin{pgfscope}%
\pgfsys@transformshift{1.545110in}{1.850278in}%
\pgfsys@useobject{currentmarker}{}%
\end{pgfscope}%
\begin{pgfscope}%
\pgfsys@transformshift{1.545960in}{1.850278in}%
\pgfsys@useobject{currentmarker}{}%
\end{pgfscope}%
\begin{pgfscope}%
\pgfsys@transformshift{1.547646in}{1.850278in}%
\pgfsys@useobject{currentmarker}{}%
\end{pgfscope}%
\begin{pgfscope}%
\pgfsys@transformshift{1.549157in}{1.850278in}%
\pgfsys@useobject{currentmarker}{}%
\end{pgfscope}%
\begin{pgfscope}%
\pgfsys@transformshift{1.550843in}{1.850278in}%
\pgfsys@useobject{currentmarker}{}%
\end{pgfscope}%
\begin{pgfscope}%
\pgfsys@transformshift{1.552537in}{1.850278in}%
\pgfsys@useobject{currentmarker}{}%
\end{pgfscope}%
\begin{pgfscope}%
\pgfsys@transformshift{1.554553in}{1.850278in}%
\pgfsys@useobject{currentmarker}{}%
\end{pgfscope}%
\begin{pgfscope}%
\pgfsys@transformshift{1.556583in}{1.850278in}%
\pgfsys@useobject{currentmarker}{}%
\end{pgfscope}%
\begin{pgfscope}%
\pgfsys@transformshift{1.558101in}{1.850278in}%
\pgfsys@useobject{currentmarker}{}%
\end{pgfscope}%
\begin{pgfscope}%
\pgfsys@transformshift{1.561312in}{1.850278in}%
\pgfsys@useobject{currentmarker}{}%
\end{pgfscope}%
\begin{pgfscope}%
\pgfsys@transformshift{1.563666in}{1.850278in}%
\pgfsys@useobject{currentmarker}{}%
\end{pgfscope}%
\begin{pgfscope}%
\pgfsys@transformshift{1.813397in}{1.850278in}%
\pgfsys@useobject{currentmarker}{}%
\end{pgfscope}%
\begin{pgfscope}%
\pgfsys@transformshift{1.824867in}{1.850278in}%
\pgfsys@useobject{currentmarker}{}%
\end{pgfscope}%
\begin{pgfscope}%
\pgfsys@transformshift{1.831437in}{1.850278in}%
\pgfsys@useobject{currentmarker}{}%
\end{pgfscope}%
\begin{pgfscope}%
\pgfsys@transformshift{1.841053in}{1.850278in}%
\pgfsys@useobject{currentmarker}{}%
\end{pgfscope}%
\begin{pgfscope}%
\pgfsys@transformshift{1.846445in}{1.850278in}%
\pgfsys@useobject{currentmarker}{}%
\end{pgfscope}%
\begin{pgfscope}%
\pgfsys@transformshift{1.852352in}{1.850278in}%
\pgfsys@useobject{currentmarker}{}%
\end{pgfscope}%
\begin{pgfscope}%
\pgfsys@transformshift{1.858429in}{1.850278in}%
\pgfsys@useobject{currentmarker}{}%
\end{pgfscope}%
\begin{pgfscope}%
\pgfsys@transformshift{1.864321in}{1.850278in}%
\pgfsys@useobject{currentmarker}{}%
\end{pgfscope}%
\begin{pgfscope}%
\pgfsys@transformshift{1.868373in}{1.850278in}%
\pgfsys@useobject{currentmarker}{}%
\end{pgfscope}%
\begin{pgfscope}%
\pgfsys@transformshift{1.873773in}{1.850278in}%
\pgfsys@useobject{currentmarker}{}%
\end{pgfscope}%
\begin{pgfscope}%
\pgfsys@transformshift{1.876305in}{1.850278in}%
\pgfsys@useobject{currentmarker}{}%
\end{pgfscope}%
\begin{pgfscope}%
\pgfsys@transformshift{1.878495in}{1.850278in}%
\pgfsys@useobject{currentmarker}{}%
\end{pgfscope}%
\begin{pgfscope}%
\pgfsys@transformshift{1.881184in}{1.850278in}%
\pgfsys@useobject{currentmarker}{}%
\end{pgfscope}%
\begin{pgfscope}%
\pgfsys@transformshift{1.882711in}{1.850278in}%
\pgfsys@useobject{currentmarker}{}%
\end{pgfscope}%
\begin{pgfscope}%
\pgfsys@transformshift{1.884900in}{1.850278in}%
\pgfsys@useobject{currentmarker}{}%
\end{pgfscope}%
\begin{pgfscope}%
\pgfsys@transformshift{1.886748in}{1.850278in}%
\pgfsys@useobject{currentmarker}{}%
\end{pgfscope}%
\begin{pgfscope}%
\pgfsys@transformshift{1.888438in}{1.850278in}%
\pgfsys@useobject{currentmarker}{}%
\end{pgfscope}%
\begin{pgfscope}%
\pgfsys@transformshift{1.890293in}{1.850278in}%
\pgfsys@useobject{currentmarker}{}%
\end{pgfscope}%
\begin{pgfscope}%
\pgfsys@transformshift{1.891641in}{1.850278in}%
\pgfsys@useobject{currentmarker}{}%
\end{pgfscope}%
\begin{pgfscope}%
\pgfsys@transformshift{1.893324in}{1.850278in}%
\pgfsys@useobject{currentmarker}{}%
\end{pgfscope}%
\begin{pgfscope}%
\pgfsys@transformshift{1.896192in}{1.850278in}%
\pgfsys@useobject{currentmarker}{}%
\end{pgfscope}%
\begin{pgfscope}%
\pgfsys@transformshift{1.897540in}{1.850278in}%
\pgfsys@useobject{currentmarker}{}%
\end{pgfscope}%
\begin{pgfscope}%
\pgfsys@transformshift{1.899737in}{1.850278in}%
\pgfsys@useobject{currentmarker}{}%
\end{pgfscope}%
\begin{pgfscope}%
\pgfsys@transformshift{1.901763in}{1.850278in}%
\pgfsys@useobject{currentmarker}{}%
\end{pgfscope}%
\begin{pgfscope}%
\pgfsys@transformshift{1.903454in}{1.850278in}%
\pgfsys@useobject{currentmarker}{}%
\end{pgfscope}%
\begin{pgfscope}%
\pgfsys@transformshift{1.905815in}{1.850278in}%
\pgfsys@useobject{currentmarker}{}%
\end{pgfscope}%
\begin{pgfscope}%
\pgfsys@transformshift{1.908176in}{1.850278in}%
\pgfsys@useobject{currentmarker}{}%
\end{pgfscope}%
\begin{pgfscope}%
\pgfsys@transformshift{1.910701in}{1.850278in}%
\pgfsys@useobject{currentmarker}{}%
\end{pgfscope}%
\begin{pgfscope}%
\pgfsys@transformshift{1.913740in}{1.850278in}%
\pgfsys@useobject{currentmarker}{}%
\end{pgfscope}%
\begin{pgfscope}%
\pgfsys@transformshift{1.917963in}{1.850278in}%
\pgfsys@useobject{currentmarker}{}%
\end{pgfscope}%
\begin{pgfscope}%
\pgfsys@transformshift{1.922350in}{1.850278in}%
\pgfsys@useobject{currentmarker}{}%
\end{pgfscope}%
\begin{pgfscope}%
\pgfsys@transformshift{1.929091in}{1.850278in}%
\pgfsys@useobject{currentmarker}{}%
\end{pgfscope}%
\begin{pgfscope}%
\pgfsys@transformshift{1.932614in}{1.850278in}%
\pgfsys@useobject{currentmarker}{}%
\end{pgfscope}%
\begin{pgfscope}%
\pgfsys@transformshift{1.934982in}{1.850278in}%
\pgfsys@useobject{currentmarker}{}%
\end{pgfscope}%
\begin{pgfscope}%
\pgfsys@transformshift{2.052009in}{1.850278in}%
\pgfsys@useobject{currentmarker}{}%
\end{pgfscope}%
\begin{pgfscope}%
\pgfsys@transformshift{2.064820in}{1.850278in}%
\pgfsys@useobject{currentmarker}{}%
\end{pgfscope}%
\begin{pgfscope}%
\pgfsys@transformshift{2.071062in}{1.850278in}%
\pgfsys@useobject{currentmarker}{}%
\end{pgfscope}%
\begin{pgfscope}%
\pgfsys@transformshift{2.076454in}{1.850278in}%
\pgfsys@useobject{currentmarker}{}%
\end{pgfscope}%
\begin{pgfscope}%
\pgfsys@transformshift{2.080007in}{1.850278in}%
\pgfsys@useobject{currentmarker}{}%
\end{pgfscope}%
\begin{pgfscope}%
\pgfsys@transformshift{2.085057in}{1.850278in}%
\pgfsys@useobject{currentmarker}{}%
\end{pgfscope}%
\begin{pgfscope}%
\pgfsys@transformshift{2.087753in}{1.850278in}%
\pgfsys@useobject{currentmarker}{}%
\end{pgfscope}%
\begin{pgfscope}%
\pgfsys@transformshift{2.093153in}{1.850278in}%
\pgfsys@useobject{currentmarker}{}%
\end{pgfscope}%
\begin{pgfscope}%
\pgfsys@transformshift{2.095343in}{1.850278in}%
\pgfsys@useobject{currentmarker}{}%
\end{pgfscope}%
\begin{pgfscope}%
\pgfsys@transformshift{2.096922in}{1.850278in}%
\pgfsys@useobject{currentmarker}{}%
\end{pgfscope}%
\begin{pgfscope}%
\pgfsys@transformshift{2.097875in}{1.850278in}%
\pgfsys@useobject{currentmarker}{}%
\end{pgfscope}%
\begin{pgfscope}%
\pgfsys@transformshift{2.099224in}{1.850278in}%
\pgfsys@useobject{currentmarker}{}%
\end{pgfscope}%
\begin{pgfscope}%
\pgfsys@transformshift{2.100237in}{1.850278in}%
\pgfsys@useobject{currentmarker}{}%
\end{pgfscope}%
\begin{pgfscope}%
\pgfsys@transformshift{2.101242in}{1.850278in}%
\pgfsys@useobject{currentmarker}{}%
\end{pgfscope}%
\begin{pgfscope}%
\pgfsys@transformshift{2.102426in}{1.850278in}%
\pgfsys@useobject{currentmarker}{}%
\end{pgfscope}%
\begin{pgfscope}%
\pgfsys@transformshift{2.103454in}{1.850278in}%
\pgfsys@useobject{currentmarker}{}%
\end{pgfscope}%
\begin{pgfscope}%
\pgfsys@transformshift{2.104452in}{1.850278in}%
\pgfsys@useobject{currentmarker}{}%
\end{pgfscope}%
\begin{pgfscope}%
\pgfsys@transformshift{2.105458in}{1.850278in}%
\pgfsys@useobject{currentmarker}{}%
\end{pgfscope}%
\begin{pgfscope}%
\pgfsys@transformshift{2.106650in}{1.850278in}%
\pgfsys@useobject{currentmarker}{}%
\end{pgfscope}%
\begin{pgfscope}%
\pgfsys@transformshift{2.107484in}{1.850278in}%
\pgfsys@useobject{currentmarker}{}%
\end{pgfscope}%
\begin{pgfscope}%
\pgfsys@transformshift{2.108005in}{1.850278in}%
\pgfsys@useobject{currentmarker}{}%
\end{pgfscope}%
\begin{pgfscope}%
\pgfsys@transformshift{2.109018in}{1.850278in}%
\pgfsys@useobject{currentmarker}{}%
\end{pgfscope}%
\begin{pgfscope}%
\pgfsys@transformshift{2.110180in}{1.850278in}%
\pgfsys@useobject{currentmarker}{}%
\end{pgfscope}%
\begin{pgfscope}%
\pgfsys@transformshift{2.110858in}{1.850278in}%
\pgfsys@useobject{currentmarker}{}%
\end{pgfscope}%
\begin{pgfscope}%
\pgfsys@transformshift{2.111550in}{1.850278in}%
\pgfsys@useobject{currentmarker}{}%
\end{pgfscope}%
\begin{pgfscope}%
\pgfsys@transformshift{2.112042in}{1.850278in}%
\pgfsys@useobject{currentmarker}{}%
\end{pgfscope}%
\begin{pgfscope}%
\pgfsys@transformshift{2.113070in}{1.850278in}%
\pgfsys@useobject{currentmarker}{}%
\end{pgfscope}%
\begin{pgfscope}%
\pgfsys@transformshift{2.113912in}{1.850278in}%
\pgfsys@useobject{currentmarker}{}%
\end{pgfscope}%
\begin{pgfscope}%
\pgfsys@transformshift{2.114403in}{1.850278in}%
\pgfsys@useobject{currentmarker}{}%
\end{pgfscope}%
\begin{pgfscope}%
\pgfsys@transformshift{2.115252in}{1.850278in}%
\pgfsys@useobject{currentmarker}{}%
\end{pgfscope}%
\begin{pgfscope}%
\pgfsys@transformshift{2.116444in}{1.850278in}%
\pgfsys@useobject{currentmarker}{}%
\end{pgfscope}%
\begin{pgfscope}%
\pgfsys@transformshift{2.117435in}{1.850278in}%
\pgfsys@useobject{currentmarker}{}%
\end{pgfscope}%
\begin{pgfscope}%
\pgfsys@transformshift{2.118470in}{1.850278in}%
\pgfsys@useobject{currentmarker}{}%
\end{pgfscope}%
\begin{pgfscope}%
\pgfsys@transformshift{2.119289in}{1.850278in}%
\pgfsys@useobject{currentmarker}{}%
\end{pgfscope}%
\begin{pgfscope}%
\pgfsys@transformshift{2.120473in}{1.850278in}%
\pgfsys@useobject{currentmarker}{}%
\end{pgfscope}%
\begin{pgfscope}%
\pgfsys@transformshift{2.121658in}{1.850278in}%
\pgfsys@useobject{currentmarker}{}%
\end{pgfscope}%
\begin{pgfscope}%
\pgfsys@transformshift{2.122321in}{1.850278in}%
\pgfsys@useobject{currentmarker}{}%
\end{pgfscope}%
\begin{pgfscope}%
\pgfsys@transformshift{2.123334in}{1.850278in}%
\pgfsys@useobject{currentmarker}{}%
\end{pgfscope}%
\begin{pgfscope}%
\pgfsys@transformshift{2.124183in}{1.850278in}%
\pgfsys@useobject{currentmarker}{}%
\end{pgfscope}%
\begin{pgfscope}%
\pgfsys@transformshift{2.125196in}{1.850278in}%
\pgfsys@useobject{currentmarker}{}%
\end{pgfscope}%
\begin{pgfscope}%
\pgfsys@transformshift{2.126395in}{1.850278in}%
\pgfsys@useobject{currentmarker}{}%
\end{pgfscope}%
\begin{pgfscope}%
\pgfsys@transformshift{2.144606in}{1.850278in}%
\pgfsys@useobject{currentmarker}{}%
\end{pgfscope}%
\begin{pgfscope}%
\pgfsys@transformshift{2.148807in}{1.850278in}%
\pgfsys@useobject{currentmarker}{}%
\end{pgfscope}%
\begin{pgfscope}%
\pgfsys@transformshift{2.154542in}{1.850278in}%
\pgfsys@useobject{currentmarker}{}%
\end{pgfscope}%
\begin{pgfscope}%
\pgfsys@transformshift{2.393310in}{1.850278in}%
\pgfsys@useobject{currentmarker}{}%
\end{pgfscope}%
\begin{pgfscope}%
\pgfsys@transformshift{2.409331in}{1.850278in}%
\pgfsys@useobject{currentmarker}{}%
\end{pgfscope}%
\begin{pgfscope}%
\pgfsys@transformshift{2.419625in}{1.850278in}%
\pgfsys@useobject{currentmarker}{}%
\end{pgfscope}%
\begin{pgfscope}%
\pgfsys@transformshift{2.425181in}{1.850278in}%
\pgfsys@useobject{currentmarker}{}%
\end{pgfscope}%
\begin{pgfscope}%
\pgfsys@transformshift{2.430239in}{1.850278in}%
\pgfsys@useobject{currentmarker}{}%
\end{pgfscope}%
\begin{pgfscope}%
\pgfsys@transformshift{2.434626in}{1.850278in}%
\pgfsys@useobject{currentmarker}{}%
\end{pgfscope}%
\begin{pgfscope}%
\pgfsys@transformshift{2.438164in}{1.850278in}%
\pgfsys@useobject{currentmarker}{}%
\end{pgfscope}%
\begin{pgfscope}%
\pgfsys@transformshift{2.442551in}{1.850278in}%
\pgfsys@useobject{currentmarker}{}%
\end{pgfscope}%
\begin{pgfscope}%
\pgfsys@transformshift{2.446766in}{1.850278in}%
\pgfsys@useobject{currentmarker}{}%
\end{pgfscope}%
\begin{pgfscope}%
\pgfsys@transformshift{2.453008in}{1.850278in}%
\pgfsys@useobject{currentmarker}{}%
\end{pgfscope}%
\begin{pgfscope}%
\pgfsys@transformshift{2.455868in}{1.850278in}%
\pgfsys@useobject{currentmarker}{}%
\end{pgfscope}%
\begin{pgfscope}%
\pgfsys@transformshift{2.458572in}{1.850278in}%
\pgfsys@useobject{currentmarker}{}%
\end{pgfscope}%
\begin{pgfscope}%
\pgfsys@transformshift{2.459756in}{1.850278in}%
\pgfsys@useobject{currentmarker}{}%
\end{pgfscope}%
\begin{pgfscope}%
\pgfsys@transformshift{2.461774in}{1.850278in}%
\pgfsys@useobject{currentmarker}{}%
\end{pgfscope}%
\begin{pgfscope}%
\pgfsys@transformshift{2.462452in}{1.850278in}%
\pgfsys@useobject{currentmarker}{}%
\end{pgfscope}%
\begin{pgfscope}%
\pgfsys@transformshift{2.463629in}{1.850278in}%
\pgfsys@useobject{currentmarker}{}%
\end{pgfscope}%
\begin{pgfscope}%
\pgfsys@transformshift{2.464985in}{1.850278in}%
\pgfsys@useobject{currentmarker}{}%
\end{pgfscope}%
\begin{pgfscope}%
\pgfsys@transformshift{2.466497in}{1.850278in}%
\pgfsys@useobject{currentmarker}{}%
\end{pgfscope}%
\begin{pgfscope}%
\pgfsys@transformshift{2.467167in}{1.850278in}%
\pgfsys@useobject{currentmarker}{}%
\end{pgfscope}%
\begin{pgfscope}%
\pgfsys@transformshift{2.468359in}{1.850278in}%
\pgfsys@useobject{currentmarker}{}%
\end{pgfscope}%
\begin{pgfscope}%
\pgfsys@transformshift{2.469193in}{1.850278in}%
\pgfsys@useobject{currentmarker}{}%
\end{pgfscope}%
\begin{pgfscope}%
\pgfsys@transformshift{2.469707in}{1.850278in}%
\pgfsys@useobject{currentmarker}{}%
\end{pgfscope}%
\begin{pgfscope}%
\pgfsys@transformshift{2.470548in}{1.850278in}%
\pgfsys@useobject{currentmarker}{}%
\end{pgfscope}%
\begin{pgfscope}%
\pgfsys@transformshift{2.471234in}{1.850278in}%
\pgfsys@useobject{currentmarker}{}%
\end{pgfscope}%
\begin{pgfscope}%
\pgfsys@transformshift{2.471904in}{1.850278in}%
\pgfsys@useobject{currentmarker}{}%
\end{pgfscope}%
\begin{pgfscope}%
\pgfsys@transformshift{2.472917in}{1.850278in}%
\pgfsys@useobject{currentmarker}{}%
\end{pgfscope}%
\begin{pgfscope}%
\pgfsys@transformshift{2.473759in}{1.850278in}%
\pgfsys@useobject{currentmarker}{}%
\end{pgfscope}%
\begin{pgfscope}%
\pgfsys@transformshift{2.474600in}{1.850278in}%
\pgfsys@useobject{currentmarker}{}%
\end{pgfscope}%
\begin{pgfscope}%
\pgfsys@transformshift{2.475435in}{1.850278in}%
\pgfsys@useobject{currentmarker}{}%
\end{pgfscope}%
\begin{pgfscope}%
\pgfsys@transformshift{2.476276in}{1.850278in}%
\pgfsys@useobject{currentmarker}{}%
\end{pgfscope}%
\begin{pgfscope}%
\pgfsys@transformshift{2.477133in}{1.850278in}%
\pgfsys@useobject{currentmarker}{}%
\end{pgfscope}%
\begin{pgfscope}%
\pgfsys@transformshift{2.478302in}{1.850278in}%
\pgfsys@useobject{currentmarker}{}%
\end{pgfscope}%
\begin{pgfscope}%
\pgfsys@transformshift{2.479315in}{1.850278in}%
\pgfsys@useobject{currentmarker}{}%
\end{pgfscope}%
\begin{pgfscope}%
\pgfsys@transformshift{2.479993in}{1.850278in}%
\pgfsys@useobject{currentmarker}{}%
\end{pgfscope}%
\begin{pgfscope}%
\pgfsys@transformshift{2.480835in}{1.850278in}%
\pgfsys@useobject{currentmarker}{}%
\end{pgfscope}%
\begin{pgfscope}%
\pgfsys@transformshift{2.481341in}{1.850278in}%
\pgfsys@useobject{currentmarker}{}%
\end{pgfscope}%
\begin{pgfscope}%
\pgfsys@transformshift{2.482354in}{1.850278in}%
\pgfsys@useobject{currentmarker}{}%
\end{pgfscope}%
\begin{pgfscope}%
\pgfsys@transformshift{2.483367in}{1.850278in}%
\pgfsys@useobject{currentmarker}{}%
\end{pgfscope}%
\begin{pgfscope}%
\pgfsys@transformshift{2.484380in}{1.850278in}%
\pgfsys@useobject{currentmarker}{}%
\end{pgfscope}%
\begin{pgfscope}%
\pgfsys@transformshift{2.485385in}{1.850278in}%
\pgfsys@useobject{currentmarker}{}%
\end{pgfscope}%
\begin{pgfscope}%
\pgfsys@transformshift{2.486734in}{1.850278in}%
\pgfsys@useobject{currentmarker}{}%
\end{pgfscope}%
\begin{pgfscope}%
\pgfsys@transformshift{2.488417in}{1.850278in}%
\pgfsys@useobject{currentmarker}{}%
\end{pgfscope}%
\begin{pgfscope}%
\pgfsys@transformshift{2.489266in}{1.850278in}%
\pgfsys@useobject{currentmarker}{}%
\end{pgfscope}%
\begin{pgfscope}%
\pgfsys@transformshift{2.490450in}{1.850278in}%
\pgfsys@useobject{currentmarker}{}%
\end{pgfscope}%
\begin{pgfscope}%
\pgfsys@transformshift{2.494159in}{1.850278in}%
\pgfsys@useobject{currentmarker}{}%
\end{pgfscope}%
\begin{pgfscope}%
\pgfsys@transformshift{2.495508in}{1.850278in}%
\pgfsys@useobject{currentmarker}{}%
\end{pgfscope}%
\begin{pgfscope}%
\pgfsys@transformshift{2.497191in}{1.850278in}%
\pgfsys@useobject{currentmarker}{}%
\end{pgfscope}%
\begin{pgfscope}%
\pgfsys@transformshift{2.499217in}{1.850278in}%
\pgfsys@useobject{currentmarker}{}%
\end{pgfscope}%
\begin{pgfscope}%
\pgfsys@transformshift{2.500729in}{1.850278in}%
\pgfsys@useobject{currentmarker}{}%
\end{pgfscope}%
\begin{pgfscope}%
\pgfsys@transformshift{2.503604in}{1.850278in}%
\pgfsys@useobject{currentmarker}{}%
\end{pgfscope}%
\begin{pgfscope}%
\pgfsys@transformshift{2.764121in}{1.850278in}%
\pgfsys@useobject{currentmarker}{}%
\end{pgfscope}%
\begin{pgfscope}%
\pgfsys@transformshift{2.778965in}{1.850278in}%
\pgfsys@useobject{currentmarker}{}%
\end{pgfscope}%
\begin{pgfscope}%
\pgfsys@transformshift{2.788745in}{1.850278in}%
\pgfsys@useobject{currentmarker}{}%
\end{pgfscope}%
\begin{pgfscope}%
\pgfsys@transformshift{2.795828in}{1.850278in}%
\pgfsys@useobject{currentmarker}{}%
\end{pgfscope}%
\begin{pgfscope}%
\pgfsys@transformshift{2.803582in}{1.850278in}%
\pgfsys@useobject{currentmarker}{}%
\end{pgfscope}%
\begin{pgfscope}%
\pgfsys@transformshift{2.808818in}{1.850278in}%
\pgfsys@useobject{currentmarker}{}%
\end{pgfscope}%
\begin{pgfscope}%
\pgfsys@transformshift{2.813875in}{1.850278in}%
\pgfsys@useobject{currentmarker}{}%
\end{pgfscope}%
\begin{pgfscope}%
\pgfsys@transformshift{2.817078in}{1.850278in}%
\pgfsys@useobject{currentmarker}{}%
\end{pgfscope}%
\begin{pgfscope}%
\pgfsys@transformshift{2.821465in}{1.850278in}%
\pgfsys@useobject{currentmarker}{}%
\end{pgfscope}%
\begin{pgfscope}%
\pgfsys@transformshift{2.824660in}{1.850278in}%
\pgfsys@useobject{currentmarker}{}%
\end{pgfscope}%
\begin{pgfscope}%
\pgfsys@transformshift{2.829554in}{1.850278in}%
\pgfsys@useobject{currentmarker}{}%
\end{pgfscope}%
\begin{pgfscope}%
\pgfsys@transformshift{2.831915in}{1.850278in}%
\pgfsys@useobject{currentmarker}{}%
\end{pgfscope}%
\begin{pgfscope}%
\pgfsys@transformshift{2.833606in}{1.850278in}%
\pgfsys@useobject{currentmarker}{}%
\end{pgfscope}%
\begin{pgfscope}%
\pgfsys@transformshift{2.836637in}{1.850278in}%
\pgfsys@useobject{currentmarker}{}%
\end{pgfscope}%
\begin{pgfscope}%
\pgfsys@transformshift{2.839170in}{1.850278in}%
\pgfsys@useobject{currentmarker}{}%
\end{pgfscope}%
\begin{pgfscope}%
\pgfsys@transformshift{2.841359in}{1.850278in}%
\pgfsys@useobject{currentmarker}{}%
\end{pgfscope}%
\begin{pgfscope}%
\pgfsys@transformshift{2.842708in}{1.850278in}%
\pgfsys@useobject{currentmarker}{}%
\end{pgfscope}%
\begin{pgfscope}%
\pgfsys@transformshift{2.844391in}{1.850278in}%
\pgfsys@useobject{currentmarker}{}%
\end{pgfscope}%
\begin{pgfscope}%
\pgfsys@transformshift{2.846245in}{1.850278in}%
\pgfsys@useobject{currentmarker}{}%
\end{pgfscope}%
\begin{pgfscope}%
\pgfsys@transformshift{2.848785in}{1.850278in}%
\pgfsys@useobject{currentmarker}{}%
\end{pgfscope}%
\begin{pgfscope}%
\pgfsys@transformshift{2.850968in}{1.850278in}%
\pgfsys@useobject{currentmarker}{}%
\end{pgfscope}%
\begin{pgfscope}%
\pgfsys@transformshift{2.852651in}{1.850278in}%
\pgfsys@useobject{currentmarker}{}%
\end{pgfscope}%
\begin{pgfscope}%
\pgfsys@transformshift{2.854007in}{1.850278in}%
\pgfsys@useobject{currentmarker}{}%
\end{pgfscope}%
\begin{pgfscope}%
\pgfsys@transformshift{2.855861in}{1.850278in}%
\pgfsys@useobject{currentmarker}{}%
\end{pgfscope}%
\begin{pgfscope}%
\pgfsys@transformshift{2.857209in}{1.850278in}%
\pgfsys@useobject{currentmarker}{}%
\end{pgfscope}%
\begin{pgfscope}%
\pgfsys@transformshift{2.859570in}{1.850278in}%
\pgfsys@useobject{currentmarker}{}%
\end{pgfscope}%
\begin{pgfscope}%
\pgfsys@transformshift{2.862103in}{1.850278in}%
\pgfsys@useobject{currentmarker}{}%
\end{pgfscope}%
\begin{pgfscope}%
\pgfsys@transformshift{2.864807in}{1.850278in}%
\pgfsys@useobject{currentmarker}{}%
\end{pgfscope}%
\begin{pgfscope}%
\pgfsys@transformshift{2.868009in}{1.850278in}%
\pgfsys@useobject{currentmarker}{}%
\end{pgfscope}%
\begin{pgfscope}%
\pgfsys@transformshift{2.872218in}{1.850278in}%
\pgfsys@useobject{currentmarker}{}%
\end{pgfscope}%
\begin{pgfscope}%
\pgfsys@transformshift{2.876947in}{1.850278in}%
\pgfsys@useobject{currentmarker}{}%
\end{pgfscope}%
\begin{pgfscope}%
\pgfsys@transformshift{2.886228in}{1.850278in}%
\pgfsys@useobject{currentmarker}{}%
\end{pgfscope}%
\begin{pgfscope}%
\pgfsys@transformshift{3.013697in}{1.850278in}%
\pgfsys@useobject{currentmarker}{}%
\end{pgfscope}%
\begin{pgfscope}%
\pgfsys@transformshift{3.022799in}{1.850278in}%
\pgfsys@useobject{currentmarker}{}%
\end{pgfscope}%
\begin{pgfscope}%
\pgfsys@transformshift{3.033085in}{1.850278in}%
\pgfsys@useobject{currentmarker}{}%
\end{pgfscope}%
\begin{pgfscope}%
\pgfsys@transformshift{3.040004in}{1.850278in}%
\pgfsys@useobject{currentmarker}{}%
\end{pgfscope}%
\begin{pgfscope}%
\pgfsys@transformshift{3.043036in}{1.850278in}%
\pgfsys@useobject{currentmarker}{}%
\end{pgfscope}%
\begin{pgfscope}%
\pgfsys@transformshift{3.047586in}{1.850278in}%
\pgfsys@useobject{currentmarker}{}%
\end{pgfscope}%
\begin{pgfscope}%
\pgfsys@transformshift{3.051296in}{1.850278in}%
\pgfsys@useobject{currentmarker}{}%
\end{pgfscope}%
\begin{pgfscope}%
\pgfsys@transformshift{3.056860in}{1.850278in}%
\pgfsys@useobject{currentmarker}{}%
\end{pgfscope}%
\begin{pgfscope}%
\pgfsys@transformshift{3.060241in}{1.850278in}%
\pgfsys@useobject{currentmarker}{}%
\end{pgfscope}%
\begin{pgfscope}%
\pgfsys@transformshift{3.062595in}{1.850278in}%
\pgfsys@useobject{currentmarker}{}%
\end{pgfscope}%
\begin{pgfscope}%
\pgfsys@transformshift{3.064285in}{1.850278in}%
\pgfsys@useobject{currentmarker}{}%
\end{pgfscope}%
\begin{pgfscope}%
\pgfsys@transformshift{3.065805in}{1.850278in}%
\pgfsys@useobject{currentmarker}{}%
\end{pgfscope}%
\begin{pgfscope}%
\pgfsys@transformshift{3.066639in}{1.850278in}%
\pgfsys@useobject{currentmarker}{}%
\end{pgfscope}%
\begin{pgfscope}%
\pgfsys@transformshift{3.067823in}{1.850278in}%
\pgfsys@useobject{currentmarker}{}%
\end{pgfscope}%
\begin{pgfscope}%
\pgfsys@transformshift{3.069000in}{1.850278in}%
\pgfsys@useobject{currentmarker}{}%
\end{pgfscope}%
\begin{pgfscope}%
\pgfsys@transformshift{3.070184in}{1.850278in}%
\pgfsys@useobject{currentmarker}{}%
\end{pgfscope}%
\begin{pgfscope}%
\pgfsys@transformshift{3.071250in}{1.850278in}%
\pgfsys@useobject{currentmarker}{}%
\end{pgfscope}%
\begin{pgfscope}%
\pgfsys@transformshift{3.072553in}{1.850278in}%
\pgfsys@useobject{currentmarker}{}%
\end{pgfscope}%
\begin{pgfscope}%
\pgfsys@transformshift{3.073894in}{1.850278in}%
\pgfsys@useobject{currentmarker}{}%
\end{pgfscope}%
\begin{pgfscope}%
\pgfsys@transformshift{3.074735in}{1.850278in}%
\pgfsys@useobject{currentmarker}{}%
\end{pgfscope}%
\begin{pgfscope}%
\pgfsys@transformshift{3.076083in}{1.850278in}%
\pgfsys@useobject{currentmarker}{}%
\end{pgfscope}%
\begin{pgfscope}%
\pgfsys@transformshift{3.076933in}{1.850278in}%
\pgfsys@useobject{currentmarker}{}%
\end{pgfscope}%
\begin{pgfscope}%
\pgfsys@transformshift{3.077610in}{1.850278in}%
\pgfsys@useobject{currentmarker}{}%
\end{pgfscope}%
\begin{pgfscope}%
\pgfsys@transformshift{3.078281in}{1.850278in}%
\pgfsys@useobject{currentmarker}{}%
\end{pgfscope}%
\begin{pgfscope}%
\pgfsys@transformshift{3.079122in}{1.850278in}%
\pgfsys@useobject{currentmarker}{}%
\end{pgfscope}%
\begin{pgfscope}%
\pgfsys@transformshift{3.079971in}{1.850278in}%
\pgfsys@useobject{currentmarker}{}%
\end{pgfscope}%
\begin{pgfscope}%
\pgfsys@transformshift{3.080470in}{1.850278in}%
\pgfsys@useobject{currentmarker}{}%
\end{pgfscope}%
\begin{pgfscope}%
\pgfsys@transformshift{3.081483in}{1.850278in}%
\pgfsys@useobject{currentmarker}{}%
\end{pgfscope}%
\begin{pgfscope}%
\pgfsys@transformshift{3.082154in}{1.850278in}%
\pgfsys@useobject{currentmarker}{}%
\end{pgfscope}%
\begin{pgfscope}%
\pgfsys@transformshift{3.082832in}{1.850278in}%
\pgfsys@useobject{currentmarker}{}%
\end{pgfscope}%
\begin{pgfscope}%
\pgfsys@transformshift{3.084180in}{1.850278in}%
\pgfsys@useobject{currentmarker}{}%
\end{pgfscope}%
\begin{pgfscope}%
\pgfsys@transformshift{3.084857in}{1.850278in}%
\pgfsys@useobject{currentmarker}{}%
\end{pgfscope}%
\begin{pgfscope}%
\pgfsys@transformshift{3.085863in}{1.850278in}%
\pgfsys@useobject{currentmarker}{}%
\end{pgfscope}%
\begin{pgfscope}%
\pgfsys@transformshift{3.086712in}{1.850278in}%
\pgfsys@useobject{currentmarker}{}%
\end{pgfscope}%
\begin{pgfscope}%
\pgfsys@transformshift{3.087226in}{1.850278in}%
\pgfsys@useobject{currentmarker}{}%
\end{pgfscope}%
\begin{pgfscope}%
\pgfsys@transformshift{3.088403in}{1.850278in}%
\pgfsys@useobject{currentmarker}{}%
\end{pgfscope}%
\begin{pgfscope}%
\pgfsys@transformshift{3.089081in}{1.850278in}%
\pgfsys@useobject{currentmarker}{}%
\end{pgfscope}%
\begin{pgfscope}%
\pgfsys@transformshift{3.089922in}{1.850278in}%
\pgfsys@useobject{currentmarker}{}%
\end{pgfscope}%
\begin{pgfscope}%
\pgfsys@transformshift{3.090764in}{1.850278in}%
\pgfsys@useobject{currentmarker}{}%
\end{pgfscope}%
\begin{pgfscope}%
\pgfsys@transformshift{3.091598in}{1.850278in}%
\pgfsys@useobject{currentmarker}{}%
\end{pgfscope}%
\begin{pgfscope}%
\pgfsys@transformshift{3.092276in}{1.850278in}%
\pgfsys@useobject{currentmarker}{}%
\end{pgfscope}%
\begin{pgfscope}%
\pgfsys@transformshift{3.093132in}{1.850278in}%
\pgfsys@useobject{currentmarker}{}%
\end{pgfscope}%
\begin{pgfscope}%
\pgfsys@transformshift{3.093639in}{1.850278in}%
\pgfsys@useobject{currentmarker}{}%
\end{pgfscope}%
\begin{pgfscope}%
\pgfsys@transformshift{3.093974in}{1.850278in}%
\pgfsys@useobject{currentmarker}{}%
\end{pgfscope}%
\begin{pgfscope}%
\pgfsys@transformshift{3.094652in}{1.850278in}%
\pgfsys@useobject{currentmarker}{}%
\end{pgfscope}%
\begin{pgfscope}%
\pgfsys@transformshift{3.095307in}{1.850278in}%
\pgfsys@useobject{currentmarker}{}%
\end{pgfscope}%
\begin{pgfscope}%
\pgfsys@transformshift{3.095829in}{1.850278in}%
\pgfsys@useobject{currentmarker}{}%
\end{pgfscope}%
\begin{pgfscope}%
\pgfsys@transformshift{3.096320in}{1.850278in}%
\pgfsys@useobject{currentmarker}{}%
\end{pgfscope}%
\begin{pgfscope}%
\pgfsys@transformshift{3.096842in}{1.850278in}%
\pgfsys@useobject{currentmarker}{}%
\end{pgfscope}%
\begin{pgfscope}%
\pgfsys@transformshift{3.097855in}{1.850278in}%
\pgfsys@useobject{currentmarker}{}%
\end{pgfscope}%
\begin{pgfscope}%
\pgfsys@transformshift{3.098525in}{1.850278in}%
\pgfsys@useobject{currentmarker}{}%
\end{pgfscope}%
\begin{pgfscope}%
\pgfsys@transformshift{3.099717in}{1.850278in}%
\pgfsys@useobject{currentmarker}{}%
\end{pgfscope}%
\begin{pgfscope}%
\pgfsys@transformshift{3.100730in}{1.850278in}%
\pgfsys@useobject{currentmarker}{}%
\end{pgfscope}%
\begin{pgfscope}%
\pgfsys@transformshift{3.101720in}{1.850278in}%
\pgfsys@useobject{currentmarker}{}%
\end{pgfscope}%
\begin{pgfscope}%
\pgfsys@transformshift{3.103083in}{1.850278in}%
\pgfsys@useobject{currentmarker}{}%
\end{pgfscope}%
\begin{pgfscope}%
\pgfsys@transformshift{3.104268in}{1.850278in}%
\pgfsys@useobject{currentmarker}{}%
\end{pgfscope}%
\begin{pgfscope}%
\pgfsys@transformshift{3.105452in}{1.850278in}%
\pgfsys@useobject{currentmarker}{}%
\end{pgfscope}%
\begin{pgfscope}%
\pgfsys@transformshift{3.106942in}{1.850278in}%
\pgfsys@useobject{currentmarker}{}%
\end{pgfscope}%
\begin{pgfscope}%
\pgfsys@transformshift{3.107813in}{1.850278in}%
\pgfsys@useobject{currentmarker}{}%
\end{pgfscope}%
\begin{pgfscope}%
\pgfsys@transformshift{3.109161in}{1.850278in}%
\pgfsys@useobject{currentmarker}{}%
\end{pgfscope}%
\begin{pgfscope}%
\pgfsys@transformshift{3.110844in}{1.850278in}%
\pgfsys@useobject{currentmarker}{}%
\end{pgfscope}%
\begin{pgfscope}%
\pgfsys@transformshift{3.111857in}{1.850278in}%
\pgfsys@useobject{currentmarker}{}%
\end{pgfscope}%
\begin{pgfscope}%
\pgfsys@transformshift{3.113362in}{1.850278in}%
\pgfsys@useobject{currentmarker}{}%
\end{pgfscope}%
\begin{pgfscope}%
\pgfsys@transformshift{3.331052in}{1.850278in}%
\pgfsys@useobject{currentmarker}{}%
\end{pgfscope}%
\begin{pgfscope}%
\pgfsys@transformshift{3.340161in}{1.850278in}%
\pgfsys@useobject{currentmarker}{}%
\end{pgfscope}%
\begin{pgfscope}%
\pgfsys@transformshift{3.347743in}{1.850278in}%
\pgfsys@useobject{currentmarker}{}%
\end{pgfscope}%
\begin{pgfscope}%
\pgfsys@transformshift{3.356517in}{1.850278in}%
\pgfsys@useobject{currentmarker}{}%
\end{pgfscope}%
\begin{pgfscope}%
\pgfsys@transformshift{3.361403in}{1.850278in}%
\pgfsys@useobject{currentmarker}{}%
\end{pgfscope}%
\begin{pgfscope}%
\pgfsys@transformshift{3.366967in}{1.850278in}%
\pgfsys@useobject{currentmarker}{}%
\end{pgfscope}%
\begin{pgfscope}%
\pgfsys@transformshift{3.372539in}{1.850278in}%
\pgfsys@useobject{currentmarker}{}%
\end{pgfscope}%
\begin{pgfscope}%
\pgfsys@transformshift{3.377097in}{1.850278in}%
\pgfsys@useobject{currentmarker}{}%
\end{pgfscope}%
\begin{pgfscope}%
\pgfsys@transformshift{3.380970in}{1.850278in}%
\pgfsys@useobject{currentmarker}{}%
\end{pgfscope}%
\begin{pgfscope}%
\pgfsys@transformshift{3.383160in}{1.850278in}%
\pgfsys@useobject{currentmarker}{}%
\end{pgfscope}%
\begin{pgfscope}%
\pgfsys@transformshift{3.385521in}{1.850278in}%
\pgfsys@useobject{currentmarker}{}%
\end{pgfscope}%
\begin{pgfscope}%
\pgfsys@transformshift{3.388724in}{1.850278in}%
\pgfsys@useobject{currentmarker}{}%
\end{pgfscope}%
\begin{pgfscope}%
\pgfsys@transformshift{3.390750in}{1.850278in}%
\pgfsys@useobject{currentmarker}{}%
\end{pgfscope}%
\begin{pgfscope}%
\pgfsys@transformshift{3.392448in}{1.850278in}%
\pgfsys@useobject{currentmarker}{}%
\end{pgfscope}%
\begin{pgfscope}%
\pgfsys@transformshift{3.393788in}{1.850278in}%
\pgfsys@useobject{currentmarker}{}%
\end{pgfscope}%
\begin{pgfscope}%
\pgfsys@transformshift{3.395807in}{1.850278in}%
\pgfsys@useobject{currentmarker}{}%
\end{pgfscope}%
\begin{pgfscope}%
\pgfsys@transformshift{3.396991in}{1.850278in}%
\pgfsys@useobject{currentmarker}{}%
\end{pgfscope}%
\begin{pgfscope}%
\pgfsys@transformshift{3.398511in}{1.850278in}%
\pgfsys@useobject{currentmarker}{}%
\end{pgfscope}%
\begin{pgfscope}%
\pgfsys@transformshift{3.399516in}{1.850278in}%
\pgfsys@useobject{currentmarker}{}%
\end{pgfscope}%
\begin{pgfscope}%
\pgfsys@transformshift{3.401207in}{1.850278in}%
\pgfsys@useobject{currentmarker}{}%
\end{pgfscope}%
\begin{pgfscope}%
\pgfsys@transformshift{3.402726in}{1.850278in}%
\pgfsys@useobject{currentmarker}{}%
\end{pgfscope}%
\begin{pgfscope}%
\pgfsys@transformshift{3.404238in}{1.850278in}%
\pgfsys@useobject{currentmarker}{}%
\end{pgfscope}%
\begin{pgfscope}%
\pgfsys@transformshift{3.405765in}{1.850278in}%
\pgfsys@useobject{currentmarker}{}%
\end{pgfscope}%
\begin{pgfscope}%
\pgfsys@transformshift{3.407620in}{1.850278in}%
\pgfsys@useobject{currentmarker}{}%
\end{pgfscope}%
\begin{pgfscope}%
\pgfsys@transformshift{3.409296in}{1.850278in}%
\pgfsys@useobject{currentmarker}{}%
\end{pgfscope}%
\begin{pgfscope}%
\pgfsys@transformshift{3.411322in}{1.850278in}%
\pgfsys@useobject{currentmarker}{}%
\end{pgfscope}%
\begin{pgfscope}%
\pgfsys@transformshift{3.413012in}{1.850278in}%
\pgfsys@useobject{currentmarker}{}%
\end{pgfscope}%
\begin{pgfscope}%
\pgfsys@transformshift{3.414867in}{1.850278in}%
\pgfsys@useobject{currentmarker}{}%
\end{pgfscope}%
\begin{pgfscope}%
\pgfsys@transformshift{3.416394in}{1.850278in}%
\pgfsys@useobject{currentmarker}{}%
\end{pgfscope}%
\begin{pgfscope}%
\pgfsys@transformshift{3.418911in}{1.850278in}%
\pgfsys@useobject{currentmarker}{}%
\end{pgfscope}%
\begin{pgfscope}%
\pgfsys@transformshift{3.422129in}{1.850278in}%
\pgfsys@useobject{currentmarker}{}%
\end{pgfscope}%
\begin{pgfscope}%
\pgfsys@transformshift{3.426836in}{1.850278in}%
\pgfsys@useobject{currentmarker}{}%
\end{pgfscope}%
\begin{pgfscope}%
\pgfsys@transformshift{3.431231in}{1.850278in}%
\pgfsys@useobject{currentmarker}{}%
\end{pgfscope}%
\begin{pgfscope}%
\pgfsys@transformshift{3.434426in}{1.850278in}%
\pgfsys@useobject{currentmarker}{}%
\end{pgfscope}%
\begin{pgfscope}%
\pgfsys@transformshift{3.688039in}{1.850278in}%
\pgfsys@useobject{currentmarker}{}%
\end{pgfscope}%
\begin{pgfscope}%
\pgfsys@transformshift{3.701192in}{1.850278in}%
\pgfsys@useobject{currentmarker}{}%
\end{pgfscope}%
\begin{pgfscope}%
\pgfsys@transformshift{3.712149in}{1.850278in}%
\pgfsys@useobject{currentmarker}{}%
\end{pgfscope}%
\begin{pgfscope}%
\pgfsys@transformshift{3.719232in}{1.850278in}%
\pgfsys@useobject{currentmarker}{}%
\end{pgfscope}%
\begin{pgfscope}%
\pgfsys@transformshift{3.726651in}{1.850278in}%
\pgfsys@useobject{currentmarker}{}%
\end{pgfscope}%
\begin{pgfscope}%
\pgfsys@transformshift{3.732386in}{1.850278in}%
\pgfsys@useobject{currentmarker}{}%
\end{pgfscope}%
\begin{pgfscope}%
\pgfsys@transformshift{3.738456in}{1.850278in}%
\pgfsys@useobject{currentmarker}{}%
\end{pgfscope}%
\begin{pgfscope}%
\pgfsys@transformshift{3.742679in}{1.850278in}%
\pgfsys@useobject{currentmarker}{}%
\end{pgfscope}%
\begin{pgfscope}%
\pgfsys@transformshift{3.745539in}{1.850278in}%
\pgfsys@useobject{currentmarker}{}%
\end{pgfscope}%
\begin{pgfscope}%
\pgfsys@transformshift{3.748913in}{1.850278in}%
\pgfsys@useobject{currentmarker}{}%
\end{pgfscope}%
\begin{pgfscope}%
\pgfsys@transformshift{3.751610in}{1.850278in}%
\pgfsys@useobject{currentmarker}{}%
\end{pgfscope}%
\begin{pgfscope}%
\pgfsys@transformshift{3.753300in}{1.850278in}%
\pgfsys@useobject{currentmarker}{}%
\end{pgfscope}%
\begin{pgfscope}%
\pgfsys@transformshift{3.755319in}{1.850278in}%
\pgfsys@useobject{currentmarker}{}%
\end{pgfscope}%
\begin{pgfscope}%
\pgfsys@transformshift{3.757352in}{1.850278in}%
\pgfsys@useobject{currentmarker}{}%
\end{pgfscope}%
\begin{pgfscope}%
\pgfsys@transformshift{3.758872in}{1.850278in}%
\pgfsys@useobject{currentmarker}{}%
\end{pgfscope}%
\begin{pgfscope}%
\pgfsys@transformshift{3.760547in}{1.850278in}%
\pgfsys@useobject{currentmarker}{}%
\end{pgfscope}%
\begin{pgfscope}%
\pgfsys@transformshift{3.762067in}{1.850278in}%
\pgfsys@useobject{currentmarker}{}%
\end{pgfscope}%
\begin{pgfscope}%
\pgfsys@transformshift{3.763087in}{1.850278in}%
\pgfsys@useobject{currentmarker}{}%
\end{pgfscope}%
\begin{pgfscope}%
\pgfsys@transformshift{3.763922in}{1.850278in}%
\pgfsys@useobject{currentmarker}{}%
\end{pgfscope}%
\begin{pgfscope}%
\pgfsys@transformshift{3.764942in}{1.850278in}%
\pgfsys@useobject{currentmarker}{}%
\end{pgfscope}%
\begin{pgfscope}%
\pgfsys@transformshift{3.766290in}{1.850278in}%
\pgfsys@useobject{currentmarker}{}%
\end{pgfscope}%
\begin{pgfscope}%
\pgfsys@transformshift{3.767303in}{1.850278in}%
\pgfsys@useobject{currentmarker}{}%
\end{pgfscope}%
\begin{pgfscope}%
\pgfsys@transformshift{3.768651in}{1.850278in}%
\pgfsys@useobject{currentmarker}{}%
\end{pgfscope}%
\begin{pgfscope}%
\pgfsys@transformshift{3.769999in}{1.850278in}%
\pgfsys@useobject{currentmarker}{}%
\end{pgfscope}%
\begin{pgfscope}%
\pgfsys@transformshift{3.771176in}{1.850278in}%
\pgfsys@useobject{currentmarker}{}%
\end{pgfscope}%
\begin{pgfscope}%
\pgfsys@transformshift{3.772688in}{1.850278in}%
\pgfsys@useobject{currentmarker}{}%
\end{pgfscope}%
\begin{pgfscope}%
\pgfsys@transformshift{3.774036in}{1.850278in}%
\pgfsys@useobject{currentmarker}{}%
\end{pgfscope}%
\begin{pgfscope}%
\pgfsys@transformshift{3.775392in}{1.850278in}%
\pgfsys@useobject{currentmarker}{}%
\end{pgfscope}%
\begin{pgfscope}%
\pgfsys@transformshift{3.776911in}{1.850278in}%
\pgfsys@useobject{currentmarker}{}%
\end{pgfscope}%
\begin{pgfscope}%
\pgfsys@transformshift{3.778103in}{1.850278in}%
\pgfsys@useobject{currentmarker}{}%
\end{pgfscope}%
\begin{pgfscope}%
\pgfsys@transformshift{3.779771in}{1.850278in}%
\pgfsys@useobject{currentmarker}{}%
\end{pgfscope}%
\begin{pgfscope}%
\pgfsys@transformshift{3.780784in}{1.850278in}%
\pgfsys@useobject{currentmarker}{}%
\end{pgfscope}%
\begin{pgfscope}%
\pgfsys@transformshift{3.782155in}{1.850278in}%
\pgfsys@useobject{currentmarker}{}%
\end{pgfscope}%
\begin{pgfscope}%
\pgfsys@transformshift{3.783309in}{1.850278in}%
\pgfsys@useobject{currentmarker}{}%
\end{pgfscope}%
\begin{pgfscope}%
\pgfsys@transformshift{3.784166in}{1.850278in}%
\pgfsys@useobject{currentmarker}{}%
\end{pgfscope}%
\begin{pgfscope}%
\pgfsys@transformshift{3.785521in}{1.850278in}%
\pgfsys@useobject{currentmarker}{}%
\end{pgfscope}%
\begin{pgfscope}%
\pgfsys@transformshift{3.786534in}{1.850278in}%
\pgfsys@useobject{currentmarker}{}%
\end{pgfscope}%
\begin{pgfscope}%
\pgfsys@transformshift{3.787719in}{1.850278in}%
\pgfsys@useobject{currentmarker}{}%
\end{pgfscope}%
\begin{pgfscope}%
\pgfsys@transformshift{3.789044in}{1.850278in}%
\pgfsys@useobject{currentmarker}{}%
\end{pgfscope}%
\begin{pgfscope}%
\pgfsys@transformshift{3.791249in}{1.850278in}%
\pgfsys@useobject{currentmarker}{}%
\end{pgfscope}%
\begin{pgfscope}%
\pgfsys@transformshift{3.792597in}{1.850278in}%
\pgfsys@useobject{currentmarker}{}%
\end{pgfscope}%
\begin{pgfscope}%
\pgfsys@transformshift{3.795115in}{1.850278in}%
\pgfsys@useobject{currentmarker}{}%
\end{pgfscope}%
\begin{pgfscope}%
\pgfsys@transformshift{3.797818in}{1.850278in}%
\pgfsys@useobject{currentmarker}{}%
\end{pgfscope}%
\begin{pgfscope}%
\pgfsys@transformshift{4.056325in}{1.850278in}%
\pgfsys@useobject{currentmarker}{}%
\end{pgfscope}%
\begin{pgfscope}%
\pgfsys@transformshift{4.070484in}{1.850278in}%
\pgfsys@useobject{currentmarker}{}%
\end{pgfscope}%
\begin{pgfscope}%
\pgfsys@transformshift{4.077738in}{1.850278in}%
\pgfsys@useobject{currentmarker}{}%
\end{pgfscope}%
\begin{pgfscope}%
\pgfsys@transformshift{4.085999in}{1.850278in}%
\pgfsys@useobject{currentmarker}{}%
\end{pgfscope}%
\begin{pgfscope}%
\pgfsys@transformshift{4.092918in}{1.850278in}%
\pgfsys@useobject{currentmarker}{}%
\end{pgfscope}%
\begin{pgfscope}%
\pgfsys@transformshift{4.098482in}{1.850278in}%
\pgfsys@useobject{currentmarker}{}%
\end{pgfscope}%
\begin{pgfscope}%
\pgfsys@transformshift{4.105066in}{1.850278in}%
\pgfsys@useobject{currentmarker}{}%
\end{pgfscope}%
\begin{pgfscope}%
\pgfsys@transformshift{4.107584in}{1.850278in}%
\pgfsys@useobject{currentmarker}{}%
\end{pgfscope}%
\begin{pgfscope}%
\pgfsys@transformshift{4.109952in}{1.850278in}%
\pgfsys@useobject{currentmarker}{}%
\end{pgfscope}%
\begin{pgfscope}%
\pgfsys@transformshift{4.111643in}{1.850278in}%
\pgfsys@useobject{currentmarker}{}%
\end{pgfscope}%
\begin{pgfscope}%
\pgfsys@transformshift{4.113997in}{1.850278in}%
\pgfsys@useobject{currentmarker}{}%
\end{pgfscope}%
\begin{pgfscope}%
\pgfsys@transformshift{4.115509in}{1.850278in}%
\pgfsys@useobject{currentmarker}{}%
\end{pgfscope}%
\begin{pgfscope}%
\pgfsys@transformshift{4.117207in}{1.850278in}%
\pgfsys@useobject{currentmarker}{}%
\end{pgfscope}%
\begin{pgfscope}%
\pgfsys@transformshift{4.119218in}{1.850278in}%
\pgfsys@useobject{currentmarker}{}%
\end{pgfscope}%
\begin{pgfscope}%
\pgfsys@transformshift{4.120909in}{1.850278in}%
\pgfsys@useobject{currentmarker}{}%
\end{pgfscope}%
\begin{pgfscope}%
\pgfsys@transformshift{4.121921in}{1.850278in}%
\pgfsys@useobject{currentmarker}{}%
\end{pgfscope}%
\begin{pgfscope}%
\pgfsys@transformshift{4.123612in}{1.850278in}%
\pgfsys@useobject{currentmarker}{}%
\end{pgfscope}%
\begin{pgfscope}%
\pgfsys@transformshift{4.124789in}{1.850278in}%
\pgfsys@useobject{currentmarker}{}%
\end{pgfscope}%
\begin{pgfscope}%
\pgfsys@transformshift{4.126979in}{1.850278in}%
\pgfsys@useobject{currentmarker}{}%
\end{pgfscope}%
\begin{pgfscope}%
\pgfsys@transformshift{4.128156in}{1.850278in}%
\pgfsys@useobject{currentmarker}{}%
\end{pgfscope}%
\begin{pgfscope}%
\pgfsys@transformshift{4.129504in}{1.850278in}%
\pgfsys@useobject{currentmarker}{}%
\end{pgfscope}%
\begin{pgfscope}%
\pgfsys@transformshift{4.130189in}{1.850278in}%
\pgfsys@useobject{currentmarker}{}%
\end{pgfscope}%
\begin{pgfscope}%
\pgfsys@transformshift{4.131708in}{1.850278in}%
\pgfsys@useobject{currentmarker}{}%
\end{pgfscope}%
\begin{pgfscope}%
\pgfsys@transformshift{4.133906in}{1.850278in}%
\pgfsys@useobject{currentmarker}{}%
\end{pgfscope}%
\begin{pgfscope}%
\pgfsys@transformshift{4.135753in}{1.850278in}%
\pgfsys@useobject{currentmarker}{}%
\end{pgfscope}%
\begin{pgfscope}%
\pgfsys@transformshift{4.136758in}{1.850278in}%
\pgfsys@useobject{currentmarker}{}%
\end{pgfscope}%
\begin{pgfscope}%
\pgfsys@transformshift{4.138285in}{1.850278in}%
\pgfsys@useobject{currentmarker}{}%
\end{pgfscope}%
\begin{pgfscope}%
\pgfsys@transformshift{4.139805in}{1.850278in}%
\pgfsys@useobject{currentmarker}{}%
\end{pgfscope}%
\begin{pgfscope}%
\pgfsys@transformshift{4.141838in}{1.850278in}%
\pgfsys@useobject{currentmarker}{}%
\end{pgfscope}%
\begin{pgfscope}%
\pgfsys@transformshift{4.143171in}{1.850278in}%
\pgfsys@useobject{currentmarker}{}%
\end{pgfscope}%
\begin{pgfscope}%
\pgfsys@transformshift{4.144691in}{1.850278in}%
\pgfsys@useobject{currentmarker}{}%
\end{pgfscope}%
\begin{pgfscope}%
\pgfsys@transformshift{4.146218in}{1.850278in}%
\pgfsys@useobject{currentmarker}{}%
\end{pgfscope}%
\begin{pgfscope}%
\pgfsys@transformshift{4.148393in}{1.850278in}%
\pgfsys@useobject{currentmarker}{}%
\end{pgfscope}%
\begin{pgfscope}%
\pgfsys@transformshift{4.150098in}{1.850278in}%
\pgfsys@useobject{currentmarker}{}%
\end{pgfscope}%
\begin{pgfscope}%
\pgfsys@transformshift{4.150932in}{1.850278in}%
\pgfsys@useobject{currentmarker}{}%
\end{pgfscope}%
\begin{pgfscope}%
\pgfsys@transformshift{4.152109in}{1.850278in}%
\pgfsys@useobject{currentmarker}{}%
\end{pgfscope}%
\begin{pgfscope}%
\pgfsys@transformshift{4.153636in}{1.850278in}%
\pgfsys@useobject{currentmarker}{}%
\end{pgfscope}%
\begin{pgfscope}%
\pgfsys@transformshift{4.156154in}{1.850278in}%
\pgfsys@useobject{currentmarker}{}%
\end{pgfscope}%
\begin{pgfscope}%
\pgfsys@transformshift{4.158522in}{1.850278in}%
\pgfsys@useobject{currentmarker}{}%
\end{pgfscope}%
\begin{pgfscope}%
\pgfsys@transformshift{4.161055in}{1.850278in}%
\pgfsys@useobject{currentmarker}{}%
\end{pgfscope}%
\begin{pgfscope}%
\pgfsys@transformshift{4.163416in}{1.850278in}%
\pgfsys@useobject{currentmarker}{}%
\end{pgfscope}%
\begin{pgfscope}%
\pgfsys@transformshift{4.166618in}{1.850278in}%
\pgfsys@useobject{currentmarker}{}%
\end{pgfscope}%
\begin{pgfscope}%
\pgfsys@transformshift{4.172011in}{1.850278in}%
\pgfsys@useobject{currentmarker}{}%
\end{pgfscope}%
\begin{pgfscope}%
\pgfsys@transformshift{4.174208in}{1.850278in}%
\pgfsys@useobject{currentmarker}{}%
\end{pgfscope}%
\begin{pgfscope}%
\pgfsys@transformshift{4.178923in}{1.850278in}%
\pgfsys@useobject{currentmarker}{}%
\end{pgfscope}%
\begin{pgfscope}%
\pgfsys@transformshift{4.187682in}{1.850278in}%
\pgfsys@useobject{currentmarker}{}%
\end{pgfscope}%
\begin{pgfscope}%
\pgfsys@transformshift{4.426801in}{1.850278in}%
\pgfsys@useobject{currentmarker}{}%
\end{pgfscope}%
\begin{pgfscope}%
\pgfsys@transformshift{4.442479in}{1.850278in}%
\pgfsys@useobject{currentmarker}{}%
\end{pgfscope}%
\begin{pgfscope}%
\pgfsys@transformshift{4.449898in}{1.850278in}%
\pgfsys@useobject{currentmarker}{}%
\end{pgfscope}%
\begin{pgfscope}%
\pgfsys@transformshift{4.459848in}{1.850278in}%
\pgfsys@useobject{currentmarker}{}%
\end{pgfscope}%
\begin{pgfscope}%
\pgfsys@transformshift{4.465584in}{1.850278in}%
\pgfsys@useobject{currentmarker}{}%
\end{pgfscope}%
\begin{pgfscope}%
\pgfsys@transformshift{4.470984in}{1.850278in}%
\pgfsys@useobject{currentmarker}{}%
\end{pgfscope}%
\begin{pgfscope}%
\pgfsys@transformshift{4.476205in}{1.850278in}%
\pgfsys@useobject{currentmarker}{}%
\end{pgfscope}%
\begin{pgfscope}%
\pgfsys@transformshift{4.480927in}{1.850278in}%
\pgfsys@useobject{currentmarker}{}%
\end{pgfscope}%
\begin{pgfscope}%
\pgfsys@transformshift{4.485649in}{1.850278in}%
\pgfsys@useobject{currentmarker}{}%
\end{pgfscope}%
\begin{pgfscope}%
\pgfsys@transformshift{4.488524in}{1.850278in}%
\pgfsys@useobject{currentmarker}{}%
\end{pgfscope}%
\begin{pgfscope}%
\pgfsys@transformshift{4.491049in}{1.850278in}%
\pgfsys@useobject{currentmarker}{}%
\end{pgfscope}%
\begin{pgfscope}%
\pgfsys@transformshift{4.493075in}{1.850278in}%
\pgfsys@useobject{currentmarker}{}%
\end{pgfscope}%
\begin{pgfscope}%
\pgfsys@transformshift{4.494587in}{1.850278in}%
\pgfsys@useobject{currentmarker}{}%
\end{pgfscope}%
\begin{pgfscope}%
\pgfsys@transformshift{4.495779in}{1.850278in}%
\pgfsys@useobject{currentmarker}{}%
\end{pgfscope}%
\begin{pgfscope}%
\pgfsys@transformshift{4.496784in}{1.850278in}%
\pgfsys@useobject{currentmarker}{}%
\end{pgfscope}%
\begin{pgfscope}%
\pgfsys@transformshift{4.498125in}{1.850278in}%
\pgfsys@useobject{currentmarker}{}%
\end{pgfscope}%
\begin{pgfscope}%
\pgfsys@transformshift{4.499816in}{1.850278in}%
\pgfsys@useobject{currentmarker}{}%
\end{pgfscope}%
\begin{pgfscope}%
\pgfsys@transformshift{4.501000in}{1.850278in}%
\pgfsys@useobject{currentmarker}{}%
\end{pgfscope}%
\begin{pgfscope}%
\pgfsys@transformshift{4.502013in}{1.850278in}%
\pgfsys@useobject{currentmarker}{}%
\end{pgfscope}%
\begin{pgfscope}%
\pgfsys@transformshift{4.502847in}{1.850278in}%
\pgfsys@useobject{currentmarker}{}%
\end{pgfscope}%
\begin{pgfscope}%
\pgfsys@transformshift{4.503532in}{1.850278in}%
\pgfsys@useobject{currentmarker}{}%
\end{pgfscope}%
\begin{pgfscope}%
\pgfsys@transformshift{4.504702in}{1.850278in}%
\pgfsys@useobject{currentmarker}{}%
\end{pgfscope}%
\begin{pgfscope}%
\pgfsys@transformshift{4.506057in}{1.850278in}%
\pgfsys@useobject{currentmarker}{}%
\end{pgfscope}%
\begin{pgfscope}%
\pgfsys@transformshift{4.506914in}{1.850278in}%
\pgfsys@useobject{currentmarker}{}%
\end{pgfscope}%
\begin{pgfscope}%
\pgfsys@transformshift{4.507912in}{1.850278in}%
\pgfsys@useobject{currentmarker}{}%
\end{pgfscope}%
\begin{pgfscope}%
\pgfsys@transformshift{4.508754in}{1.850278in}%
\pgfsys@useobject{currentmarker}{}%
\end{pgfscope}%
\begin{pgfscope}%
\pgfsys@transformshift{4.509595in}{1.850278in}%
\pgfsys@useobject{currentmarker}{}%
\end{pgfscope}%
\begin{pgfscope}%
\pgfsys@transformshift{4.510772in}{1.850278in}%
\pgfsys@useobject{currentmarker}{}%
\end{pgfscope}%
\begin{pgfscope}%
\pgfsys@transformshift{4.511621in}{1.850278in}%
\pgfsys@useobject{currentmarker}{}%
\end{pgfscope}%
\begin{pgfscope}%
\pgfsys@transformshift{4.512478in}{1.850278in}%
\pgfsys@useobject{currentmarker}{}%
\end{pgfscope}%
\begin{pgfscope}%
\pgfsys@transformshift{4.513483in}{1.850278in}%
\pgfsys@useobject{currentmarker}{}%
\end{pgfscope}%
\begin{pgfscope}%
\pgfsys@transformshift{4.514489in}{1.850278in}%
\pgfsys@useobject{currentmarker}{}%
\end{pgfscope}%
\begin{pgfscope}%
\pgfsys@transformshift{4.515673in}{1.850278in}%
\pgfsys@useobject{currentmarker}{}%
\end{pgfscope}%
\begin{pgfscope}%
\pgfsys@transformshift{4.516507in}{1.850278in}%
\pgfsys@useobject{currentmarker}{}%
\end{pgfscope}%
\begin{pgfscope}%
\pgfsys@transformshift{4.517699in}{1.850278in}%
\pgfsys@useobject{currentmarker}{}%
\end{pgfscope}%
\begin{pgfscope}%
\pgfsys@transformshift{4.518205in}{1.850278in}%
\pgfsys@useobject{currentmarker}{}%
\end{pgfscope}%
\begin{pgfscope}%
\pgfsys@transformshift{4.519554in}{1.850278in}%
\pgfsys@useobject{currentmarker}{}%
\end{pgfscope}%
\begin{pgfscope}%
\pgfsys@transformshift{4.520902in}{1.850278in}%
\pgfsys@useobject{currentmarker}{}%
\end{pgfscope}%
\begin{pgfscope}%
\pgfsys@transformshift{4.522071in}{1.850278in}%
\pgfsys@useobject{currentmarker}{}%
\end{pgfscope}%
\begin{pgfscope}%
\pgfsys@transformshift{4.522920in}{1.850278in}%
\pgfsys@useobject{currentmarker}{}%
\end{pgfscope}%
\begin{pgfscope}%
\pgfsys@transformshift{4.524440in}{1.850278in}%
\pgfsys@useobject{currentmarker}{}%
\end{pgfscope}%
\begin{pgfscope}%
\pgfsys@transformshift{4.525274in}{1.850278in}%
\pgfsys@useobject{currentmarker}{}%
\end{pgfscope}%
\begin{pgfscope}%
\pgfsys@transformshift{4.526480in}{1.850278in}%
\pgfsys@useobject{currentmarker}{}%
\end{pgfscope}%
\begin{pgfscope}%
\pgfsys@transformshift{4.527635in}{1.850278in}%
\pgfsys@useobject{currentmarker}{}%
\end{pgfscope}%
\begin{pgfscope}%
\pgfsys@transformshift{4.528998in}{1.850278in}%
\pgfsys@useobject{currentmarker}{}%
\end{pgfscope}%
\begin{pgfscope}%
\pgfsys@transformshift{4.530011in}{1.850278in}%
\pgfsys@useobject{currentmarker}{}%
\end{pgfscope}%
\begin{pgfscope}%
\pgfsys@transformshift{4.530689in}{1.850278in}%
\pgfsys@useobject{currentmarker}{}%
\end{pgfscope}%
\begin{pgfscope}%
\pgfsys@transformshift{4.532379in}{1.850278in}%
\pgfsys@useobject{currentmarker}{}%
\end{pgfscope}%
\begin{pgfscope}%
\pgfsys@transformshift{4.533392in}{1.850278in}%
\pgfsys@useobject{currentmarker}{}%
\end{pgfscope}%
\begin{pgfscope}%
\pgfsys@transformshift{4.534405in}{1.850278in}%
\pgfsys@useobject{currentmarker}{}%
\end{pgfscope}%
\begin{pgfscope}%
\pgfsys@transformshift{4.535917in}{1.850278in}%
\pgfsys@useobject{currentmarker}{}%
\end{pgfscope}%
\begin{pgfscope}%
\pgfsys@transformshift{4.538115in}{1.850278in}%
\pgfsys@useobject{currentmarker}{}%
\end{pgfscope}%
\begin{pgfscope}%
\pgfsys@transformshift{4.540453in}{1.850278in}%
\pgfsys@useobject{currentmarker}{}%
\end{pgfscope}%
\begin{pgfscope}%
\pgfsys@transformshift{4.541995in}{1.850278in}%
\pgfsys@useobject{currentmarker}{}%
\end{pgfscope}%
\begin{pgfscope}%
\pgfsys@transformshift{4.543671in}{1.850278in}%
\pgfsys@useobject{currentmarker}{}%
\end{pgfscope}%
\begin{pgfscope}%
\pgfsys@transformshift{4.546375in}{1.850278in}%
\pgfsys@useobject{currentmarker}{}%
\end{pgfscope}%
\begin{pgfscope}%
\pgfsys@transformshift{4.550412in}{1.850278in}%
\pgfsys@useobject{currentmarker}{}%
\end{pgfscope}%
\begin{pgfscope}%
\pgfsys@transformshift{4.557502in}{1.850278in}%
\pgfsys@useobject{currentmarker}{}%
\end{pgfscope}%
\begin{pgfscope}%
\pgfsys@transformshift{4.560869in}{1.850278in}%
\pgfsys@useobject{currentmarker}{}%
\end{pgfscope}%
\begin{pgfscope}%
\pgfsys@transformshift{4.813804in}{1.850278in}%
\pgfsys@useobject{currentmarker}{}%
\end{pgfscope}%
\begin{pgfscope}%
\pgfsys@transformshift{4.825602in}{1.850278in}%
\pgfsys@useobject{currentmarker}{}%
\end{pgfscope}%
\begin{pgfscope}%
\pgfsys@transformshift{4.836394in}{1.850278in}%
\pgfsys@useobject{currentmarker}{}%
\end{pgfscope}%
\begin{pgfscope}%
\pgfsys@transformshift{4.840953in}{1.850278in}%
\pgfsys@useobject{currentmarker}{}%
\end{pgfscope}%
\begin{pgfscope}%
\pgfsys@transformshift{4.847023in}{1.850278in}%
\pgfsys@useobject{currentmarker}{}%
\end{pgfscope}%
\begin{pgfscope}%
\pgfsys@transformshift{4.853429in}{1.850278in}%
\pgfsys@useobject{currentmarker}{}%
\end{pgfscope}%
\begin{pgfscope}%
\pgfsys@transformshift{4.857138in}{1.850278in}%
\pgfsys@useobject{currentmarker}{}%
\end{pgfscope}%
\begin{pgfscope}%
\pgfsys@transformshift{4.860676in}{1.850278in}%
\pgfsys@useobject{currentmarker}{}%
\end{pgfscope}%
\begin{pgfscope}%
\pgfsys@transformshift{4.865398in}{1.850278in}%
\pgfsys@useobject{currentmarker}{}%
\end{pgfscope}%
\begin{pgfscope}%
\pgfsys@transformshift{4.870291in}{1.850278in}%
\pgfsys@useobject{currentmarker}{}%
\end{pgfscope}%
\begin{pgfscope}%
\pgfsys@transformshift{4.873666in}{1.850278in}%
\pgfsys@useobject{currentmarker}{}%
\end{pgfscope}%
\begin{pgfscope}%
\pgfsys@transformshift{4.876533in}{1.850278in}%
\pgfsys@useobject{currentmarker}{}%
\end{pgfscope}%
\begin{pgfscope}%
\pgfsys@transformshift{4.878894in}{1.850278in}%
\pgfsys@useobject{currentmarker}{}%
\end{pgfscope}%
\begin{pgfscope}%
\pgfsys@transformshift{4.880406in}{1.850278in}%
\pgfsys@useobject{currentmarker}{}%
\end{pgfscope}%
\begin{pgfscope}%
\pgfsys@transformshift{4.881419in}{1.850278in}%
\pgfsys@useobject{currentmarker}{}%
\end{pgfscope}%
\begin{pgfscope}%
\pgfsys@transformshift{4.882268in}{1.850278in}%
\pgfsys@useobject{currentmarker}{}%
\end{pgfscope}%
\begin{pgfscope}%
\pgfsys@transformshift{4.883281in}{1.850278in}%
\pgfsys@useobject{currentmarker}{}%
\end{pgfscope}%
\begin{pgfscope}%
\pgfsys@transformshift{4.884793in}{1.850278in}%
\pgfsys@useobject{currentmarker}{}%
\end{pgfscope}%
\begin{pgfscope}%
\pgfsys@transformshift{4.885471in}{1.850278in}%
\pgfsys@useobject{currentmarker}{}%
\end{pgfscope}%
\begin{pgfscope}%
\pgfsys@transformshift{4.886476in}{1.850278in}%
\pgfsys@useobject{currentmarker}{}%
\end{pgfscope}%
\begin{pgfscope}%
\pgfsys@transformshift{4.887333in}{1.850278in}%
\pgfsys@useobject{currentmarker}{}%
\end{pgfscope}%
\begin{pgfscope}%
\pgfsys@transformshift{4.888674in}{1.850278in}%
\pgfsys@useobject{currentmarker}{}%
\end{pgfscope}%
\begin{pgfscope}%
\pgfsys@transformshift{4.889515in}{1.850278in}%
\pgfsys@useobject{currentmarker}{}%
\end{pgfscope}%
\begin{pgfscope}%
\pgfsys@transformshift{4.890528in}{1.850278in}%
\pgfsys@useobject{currentmarker}{}%
\end{pgfscope}%
\begin{pgfscope}%
\pgfsys@transformshift{4.891377in}{1.850278in}%
\pgfsys@useobject{currentmarker}{}%
\end{pgfscope}%
\begin{pgfscope}%
\pgfsys@transformshift{4.892390in}{1.850278in}%
\pgfsys@useobject{currentmarker}{}%
\end{pgfscope}%
\begin{pgfscope}%
\pgfsys@transformshift{4.893560in}{1.850278in}%
\pgfsys@useobject{currentmarker}{}%
\end{pgfscope}%
\begin{pgfscope}%
\pgfsys@transformshift{4.894908in}{1.850278in}%
\pgfsys@useobject{currentmarker}{}%
\end{pgfscope}%
\begin{pgfscope}%
\pgfsys@transformshift{4.895928in}{1.850278in}%
\pgfsys@useobject{currentmarker}{}%
\end{pgfscope}%
\begin{pgfscope}%
\pgfsys@transformshift{4.896941in}{1.850278in}%
\pgfsys@useobject{currentmarker}{}%
\end{pgfscope}%
\begin{pgfscope}%
\pgfsys@transformshift{4.897947in}{1.850278in}%
\pgfsys@useobject{currentmarker}{}%
\end{pgfscope}%
\begin{pgfscope}%
\pgfsys@transformshift{4.898960in}{1.850278in}%
\pgfsys@useobject{currentmarker}{}%
\end{pgfscope}%
\begin{pgfscope}%
\pgfsys@transformshift{4.900315in}{1.850278in}%
\pgfsys@useobject{currentmarker}{}%
\end{pgfscope}%
\begin{pgfscope}%
\pgfsys@transformshift{4.901492in}{1.850278in}%
\pgfsys@useobject{currentmarker}{}%
\end{pgfscope}%
\begin{pgfscope}%
\pgfsys@transformshift{4.902684in}{1.850278in}%
\pgfsys@useobject{currentmarker}{}%
\end{pgfscope}%
\begin{pgfscope}%
\pgfsys@transformshift{4.904017in}{1.850278in}%
\pgfsys@useobject{currentmarker}{}%
\end{pgfscope}%
\begin{pgfscope}%
\pgfsys@transformshift{4.905544in}{1.850278in}%
\pgfsys@useobject{currentmarker}{}%
\end{pgfscope}%
\begin{pgfscope}%
\pgfsys@transformshift{4.906728in}{1.850278in}%
\pgfsys@useobject{currentmarker}{}%
\end{pgfscope}%
\begin{pgfscope}%
\pgfsys@transformshift{4.907577in}{1.850278in}%
\pgfsys@useobject{currentmarker}{}%
\end{pgfscope}%
\begin{pgfscope}%
\pgfsys@transformshift{4.908903in}{1.850278in}%
\pgfsys@useobject{currentmarker}{}%
\end{pgfscope}%
\begin{pgfscope}%
\pgfsys@transformshift{4.910594in}{1.850278in}%
\pgfsys@useobject{currentmarker}{}%
\end{pgfscope}%
\begin{pgfscope}%
\pgfsys@transformshift{4.912463in}{1.850278in}%
\pgfsys@useobject{currentmarker}{}%
\end{pgfscope}%
\begin{pgfscope}%
\pgfsys@transformshift{4.913298in}{1.850278in}%
\pgfsys@useobject{currentmarker}{}%
\end{pgfscope}%
\begin{pgfscope}%
\pgfsys@transformshift{4.914154in}{1.850278in}%
\pgfsys@useobject{currentmarker}{}%
\end{pgfscope}%
\begin{pgfscope}%
\pgfsys@transformshift{4.915651in}{1.850278in}%
\pgfsys@useobject{currentmarker}{}%
\end{pgfscope}%
\begin{pgfscope}%
\pgfsys@transformshift{4.917186in}{1.850278in}%
\pgfsys@useobject{currentmarker}{}%
\end{pgfscope}%
\begin{pgfscope}%
\pgfsys@transformshift{4.918869in}{1.850278in}%
\pgfsys@useobject{currentmarker}{}%
\end{pgfscope}%
\begin{pgfscope}%
\pgfsys@transformshift{4.922407in}{1.850278in}%
\pgfsys@useobject{currentmarker}{}%
\end{pgfscope}%
\begin{pgfscope}%
\pgfsys@transformshift{4.924760in}{1.850278in}%
\pgfsys@useobject{currentmarker}{}%
\end{pgfscope}%
\begin{pgfscope}%
\pgfsys@transformshift{4.929498in}{1.850278in}%
\pgfsys@useobject{currentmarker}{}%
\end{pgfscope}%
\begin{pgfscope}%
\pgfsys@transformshift{4.934890in}{1.850278in}%
\pgfsys@useobject{currentmarker}{}%
\end{pgfscope}%
\begin{pgfscope}%
\pgfsys@transformshift{4.939106in}{1.850278in}%
\pgfsys@useobject{currentmarker}{}%
\end{pgfscope}%
\begin{pgfscope}%
\pgfsys@transformshift{5.191862in}{1.850278in}%
\pgfsys@useobject{currentmarker}{}%
\end{pgfscope}%
\begin{pgfscope}%
\pgfsys@transformshift{5.201984in}{1.850278in}%
\pgfsys@useobject{currentmarker}{}%
\end{pgfscope}%
\begin{pgfscope}%
\pgfsys@transformshift{5.209231in}{1.850278in}%
\pgfsys@useobject{currentmarker}{}%
\end{pgfscope}%
\begin{pgfscope}%
\pgfsys@transformshift{5.216821in}{1.850278in}%
\pgfsys@useobject{currentmarker}{}%
\end{pgfscope}%
\begin{pgfscope}%
\pgfsys@transformshift{5.221715in}{1.850278in}%
\pgfsys@useobject{currentmarker}{}%
\end{pgfscope}%
\begin{pgfscope}%
\pgfsys@transformshift{5.227792in}{1.850278in}%
\pgfsys@useobject{currentmarker}{}%
\end{pgfscope}%
\begin{pgfscope}%
\pgfsys@transformshift{5.233021in}{1.850278in}%
\pgfsys@useobject{currentmarker}{}%
\end{pgfscope}%
\begin{pgfscope}%
\pgfsys@transformshift{5.237564in}{1.850278in}%
\pgfsys@useobject{currentmarker}{}%
\end{pgfscope}%
\begin{pgfscope}%
\pgfsys@transformshift{5.241951in}{1.850278in}%
\pgfsys@useobject{currentmarker}{}%
\end{pgfscope}%
\begin{pgfscope}%
\pgfsys@transformshift{5.244141in}{1.850278in}%
\pgfsys@useobject{currentmarker}{}%
\end{pgfscope}%
\begin{pgfscope}%
\pgfsys@transformshift{5.246331in}{1.850278in}%
\pgfsys@useobject{currentmarker}{}%
\end{pgfscope}%
\begin{pgfscope}%
\pgfsys@transformshift{5.249027in}{1.850278in}%
\pgfsys@useobject{currentmarker}{}%
\end{pgfscope}%
\begin{pgfscope}%
\pgfsys@transformshift{5.251225in}{1.850278in}%
\pgfsys@useobject{currentmarker}{}%
\end{pgfscope}%
\begin{pgfscope}%
\pgfsys@transformshift{5.253250in}{1.850278in}%
\pgfsys@useobject{currentmarker}{}%
\end{pgfscope}%
\begin{pgfscope}%
\pgfsys@transformshift{5.254427in}{1.850278in}%
\pgfsys@useobject{currentmarker}{}%
\end{pgfscope}%
\begin{pgfscope}%
\pgfsys@transformshift{5.255775in}{1.850278in}%
\pgfsys@useobject{currentmarker}{}%
\end{pgfscope}%
\begin{pgfscope}%
\pgfsys@transformshift{5.257302in}{1.850278in}%
\pgfsys@useobject{currentmarker}{}%
\end{pgfscope}%
\begin{pgfscope}%
\pgfsys@transformshift{5.259313in}{1.850278in}%
\pgfsys@useobject{currentmarker}{}%
\end{pgfscope}%
\begin{pgfscope}%
\pgfsys@transformshift{5.260669in}{1.850278in}%
\pgfsys@useobject{currentmarker}{}%
\end{pgfscope}%
\begin{pgfscope}%
\pgfsys@transformshift{5.261853in}{1.850278in}%
\pgfsys@useobject{currentmarker}{}%
\end{pgfscope}%
\begin{pgfscope}%
\pgfsys@transformshift{5.262695in}{1.850278in}%
\pgfsys@useobject{currentmarker}{}%
\end{pgfscope}%
\begin{pgfscope}%
\pgfsys@transformshift{5.264378in}{1.850278in}%
\pgfsys@useobject{currentmarker}{}%
\end{pgfscope}%
\begin{pgfscope}%
\pgfsys@transformshift{5.266233in}{1.850278in}%
\pgfsys@useobject{currentmarker}{}%
\end{pgfscope}%
\begin{pgfscope}%
\pgfsys@transformshift{5.267916in}{1.850278in}%
\pgfsys@useobject{currentmarker}{}%
\end{pgfscope}%
\begin{pgfscope}%
\pgfsys@transformshift{5.269436in}{1.850278in}%
\pgfsys@useobject{currentmarker}{}%
\end{pgfscope}%
\begin{pgfscope}%
\pgfsys@transformshift{5.271804in}{1.850278in}%
\pgfsys@useobject{currentmarker}{}%
\end{pgfscope}%
\begin{pgfscope}%
\pgfsys@transformshift{5.272981in}{1.850278in}%
\pgfsys@useobject{currentmarker}{}%
\end{pgfscope}%
\begin{pgfscope}%
\pgfsys@transformshift{5.274329in}{1.850278in}%
\pgfsys@useobject{currentmarker}{}%
\end{pgfscope}%
\begin{pgfscope}%
\pgfsys@transformshift{5.276176in}{1.850278in}%
\pgfsys@useobject{currentmarker}{}%
\end{pgfscope}%
\begin{pgfscope}%
\pgfsys@transformshift{5.278038in}{1.850278in}%
\pgfsys@useobject{currentmarker}{}%
\end{pgfscope}%
\begin{pgfscope}%
\pgfsys@transformshift{5.279893in}{1.850278in}%
\pgfsys@useobject{currentmarker}{}%
\end{pgfscope}%
\begin{pgfscope}%
\pgfsys@transformshift{5.282425in}{1.850278in}%
\pgfsys@useobject{currentmarker}{}%
\end{pgfscope}%
\begin{pgfscope}%
\pgfsys@transformshift{5.285464in}{1.850278in}%
\pgfsys@useobject{currentmarker}{}%
\end{pgfscope}%
\begin{pgfscope}%
\pgfsys@transformshift{5.287825in}{1.850278in}%
\pgfsys@useobject{currentmarker}{}%
\end{pgfscope}%
\begin{pgfscope}%
\pgfsys@transformshift{5.293397in}{1.850278in}%
\pgfsys@useobject{currentmarker}{}%
\end{pgfscope}%
\begin{pgfscope}%
\pgfsys@transformshift{5.297270in}{1.850278in}%
\pgfsys@useobject{currentmarker}{}%
\end{pgfscope}%
\begin{pgfscope}%
\pgfsys@transformshift{5.300480in}{1.850278in}%
\pgfsys@useobject{currentmarker}{}%
\end{pgfscope}%
\end{pgfscope}%
\begin{pgfscope}%
\pgfpathrectangle{\pgfqpoint{0.713889in}{0.502778in}}{\pgfqpoint{4.805000in}{2.695000in}}%
\pgfusepath{clip}%
\pgfsetbuttcap%
\pgfsetroundjoin%
\definecolor{currentfill}{rgb}{0.172549,0.627451,0.172549}%
\pgfsetfillcolor{currentfill}%
\pgfsetlinewidth{1.003750pt}%
\definecolor{currentstroke}{rgb}{0.172549,0.627451,0.172549}%
\pgfsetstrokecolor{currentstroke}%
\pgfsetdash{}{0pt}%
\pgfsys@defobject{currentmarker}{\pgfqpoint{-0.015528in}{-0.015528in}}{\pgfqpoint{0.015528in}{0.015528in}}{%
\pgfpathmoveto{\pgfqpoint{0.000000in}{-0.015528in}}%
\pgfpathcurveto{\pgfqpoint{0.004118in}{-0.015528in}}{\pgfqpoint{0.008068in}{-0.013892in}}{\pgfqpoint{0.010980in}{-0.010980in}}%
\pgfpathcurveto{\pgfqpoint{0.013892in}{-0.008068in}}{\pgfqpoint{0.015528in}{-0.004118in}}{\pgfqpoint{0.015528in}{0.000000in}}%
\pgfpathcurveto{\pgfqpoint{0.015528in}{0.004118in}}{\pgfqpoint{0.013892in}{0.008068in}}{\pgfqpoint{0.010980in}{0.010980in}}%
\pgfpathcurveto{\pgfqpoint{0.008068in}{0.013892in}}{\pgfqpoint{0.004118in}{0.015528in}}{\pgfqpoint{0.000000in}{0.015528in}}%
\pgfpathcurveto{\pgfqpoint{-0.004118in}{0.015528in}}{\pgfqpoint{-0.008068in}{0.013892in}}{\pgfqpoint{-0.010980in}{0.010980in}}%
\pgfpathcurveto{\pgfqpoint{-0.013892in}{0.008068in}}{\pgfqpoint{-0.015528in}{0.004118in}}{\pgfqpoint{-0.015528in}{0.000000in}}%
\pgfpathcurveto{\pgfqpoint{-0.015528in}{-0.004118in}}{\pgfqpoint{-0.013892in}{-0.008068in}}{\pgfqpoint{-0.010980in}{-0.010980in}}%
\pgfpathcurveto{\pgfqpoint{-0.008068in}{-0.013892in}}{\pgfqpoint{-0.004118in}{-0.015528in}}{\pgfqpoint{0.000000in}{-0.015528in}}%
\pgfpathlineto{\pgfqpoint{0.000000in}{-0.015528in}}%
\pgfpathclose%
\pgfusepath{stroke,fill}%
}%
\begin{pgfscope}%
\pgfsys@transformshift{0.932298in}{3.075278in}%
\pgfsys@useobject{currentmarker}{}%
\end{pgfscope}%
\begin{pgfscope}%
\pgfsys@transformshift{0.932635in}{3.075278in}%
\pgfsys@useobject{currentmarker}{}%
\end{pgfscope}%
\begin{pgfscope}%
\pgfsys@transformshift{0.932972in}{3.075278in}%
\pgfsys@useobject{currentmarker}{}%
\end{pgfscope}%
\begin{pgfscope}%
\pgfsys@transformshift{0.933309in}{3.075278in}%
\pgfsys@useobject{currentmarker}{}%
\end{pgfscope}%
\begin{pgfscope}%
\pgfsys@transformshift{0.934658in}{3.075278in}%
\pgfsys@useobject{currentmarker}{}%
\end{pgfscope}%
\begin{pgfscope}%
\pgfsys@transformshift{0.934996in}{3.075278in}%
\pgfsys@useobject{currentmarker}{}%
\end{pgfscope}%
\begin{pgfscope}%
\pgfsys@transformshift{0.935164in}{3.075278in}%
\pgfsys@useobject{currentmarker}{}%
\end{pgfscope}%
\begin{pgfscope}%
\pgfsys@transformshift{0.935501in}{3.075278in}%
\pgfsys@useobject{currentmarker}{}%
\end{pgfscope}%
\begin{pgfscope}%
\pgfsys@transformshift{0.935839in}{3.075278in}%
\pgfsys@useobject{currentmarker}{}%
\end{pgfscope}%
\begin{pgfscope}%
\pgfsys@transformshift{0.936007in}{3.075278in}%
\pgfsys@useobject{currentmarker}{}%
\end{pgfscope}%
\begin{pgfscope}%
\pgfsys@transformshift{0.936176in}{3.075278in}%
\pgfsys@useobject{currentmarker}{}%
\end{pgfscope}%
\begin{pgfscope}%
\pgfsys@transformshift{0.936176in}{3.075278in}%
\pgfsys@useobject{currentmarker}{}%
\end{pgfscope}%
\begin{pgfscope}%
\pgfsys@transformshift{0.936345in}{3.075278in}%
\pgfsys@useobject{currentmarker}{}%
\end{pgfscope}%
\begin{pgfscope}%
\pgfsys@transformshift{0.936514in}{3.075278in}%
\pgfsys@useobject{currentmarker}{}%
\end{pgfscope}%
\begin{pgfscope}%
\pgfsys@transformshift{0.936682in}{3.075278in}%
\pgfsys@useobject{currentmarker}{}%
\end{pgfscope}%
\begin{pgfscope}%
\pgfsys@transformshift{0.936850in}{3.075278in}%
\pgfsys@useobject{currentmarker}{}%
\end{pgfscope}%
\begin{pgfscope}%
\pgfsys@transformshift{0.936850in}{3.075278in}%
\pgfsys@useobject{currentmarker}{}%
\end{pgfscope}%
\begin{pgfscope}%
\pgfsys@transformshift{0.937019in}{3.075278in}%
\pgfsys@useobject{currentmarker}{}%
\end{pgfscope}%
\begin{pgfscope}%
\pgfsys@transformshift{0.937188in}{3.075278in}%
\pgfsys@useobject{currentmarker}{}%
\end{pgfscope}%
\begin{pgfscope}%
\pgfsys@transformshift{0.937188in}{3.075278in}%
\pgfsys@useobject{currentmarker}{}%
\end{pgfscope}%
\begin{pgfscope}%
\pgfsys@transformshift{0.937525in}{3.075278in}%
\pgfsys@useobject{currentmarker}{}%
\end{pgfscope}%
\begin{pgfscope}%
\pgfsys@transformshift{0.937525in}{3.075278in}%
\pgfsys@useobject{currentmarker}{}%
\end{pgfscope}%
\begin{pgfscope}%
\pgfsys@transformshift{0.937863in}{3.075278in}%
\pgfsys@useobject{currentmarker}{}%
\end{pgfscope}%
\begin{pgfscope}%
\pgfsys@transformshift{0.937863in}{3.075278in}%
\pgfsys@useobject{currentmarker}{}%
\end{pgfscope}%
\begin{pgfscope}%
\pgfsys@transformshift{0.938200in}{3.075278in}%
\pgfsys@useobject{currentmarker}{}%
\end{pgfscope}%
\begin{pgfscope}%
\pgfsys@transformshift{0.938200in}{3.075278in}%
\pgfsys@useobject{currentmarker}{}%
\end{pgfscope}%
\begin{pgfscope}%
\pgfsys@transformshift{0.938537in}{3.075278in}%
\pgfsys@useobject{currentmarker}{}%
\end{pgfscope}%
\begin{pgfscope}%
\pgfsys@transformshift{0.938537in}{3.075278in}%
\pgfsys@useobject{currentmarker}{}%
\end{pgfscope}%
\begin{pgfscope}%
\pgfsys@transformshift{0.938537in}{3.075278in}%
\pgfsys@useobject{currentmarker}{}%
\end{pgfscope}%
\begin{pgfscope}%
\pgfsys@transformshift{0.938874in}{3.075278in}%
\pgfsys@useobject{currentmarker}{}%
\end{pgfscope}%
\begin{pgfscope}%
\pgfsys@transformshift{0.938874in}{3.075278in}%
\pgfsys@useobject{currentmarker}{}%
\end{pgfscope}%
\begin{pgfscope}%
\pgfsys@transformshift{0.938874in}{3.075278in}%
\pgfsys@useobject{currentmarker}{}%
\end{pgfscope}%
\begin{pgfscope}%
\pgfsys@transformshift{0.938874in}{3.075278in}%
\pgfsys@useobject{currentmarker}{}%
\end{pgfscope}%
\begin{pgfscope}%
\pgfsys@transformshift{0.939211in}{3.075278in}%
\pgfsys@useobject{currentmarker}{}%
\end{pgfscope}%
\begin{pgfscope}%
\pgfsys@transformshift{0.939211in}{3.075278in}%
\pgfsys@useobject{currentmarker}{}%
\end{pgfscope}%
\begin{pgfscope}%
\pgfsys@transformshift{0.939211in}{3.075278in}%
\pgfsys@useobject{currentmarker}{}%
\end{pgfscope}%
\begin{pgfscope}%
\pgfsys@transformshift{0.939211in}{3.075278in}%
\pgfsys@useobject{currentmarker}{}%
\end{pgfscope}%
\begin{pgfscope}%
\pgfsys@transformshift{0.939549in}{3.075278in}%
\pgfsys@useobject{currentmarker}{}%
\end{pgfscope}%
\begin{pgfscope}%
\pgfsys@transformshift{0.939549in}{3.075278in}%
\pgfsys@useobject{currentmarker}{}%
\end{pgfscope}%
\begin{pgfscope}%
\pgfsys@transformshift{0.939549in}{3.075278in}%
\pgfsys@useobject{currentmarker}{}%
\end{pgfscope}%
\begin{pgfscope}%
\pgfsys@transformshift{0.939549in}{3.075278in}%
\pgfsys@useobject{currentmarker}{}%
\end{pgfscope}%
\begin{pgfscope}%
\pgfsys@transformshift{0.939886in}{3.075278in}%
\pgfsys@useobject{currentmarker}{}%
\end{pgfscope}%
\begin{pgfscope}%
\pgfsys@transformshift{0.939886in}{3.075278in}%
\pgfsys@useobject{currentmarker}{}%
\end{pgfscope}%
\begin{pgfscope}%
\pgfsys@transformshift{0.939886in}{3.075278in}%
\pgfsys@useobject{currentmarker}{}%
\end{pgfscope}%
\begin{pgfscope}%
\pgfsys@transformshift{0.940055in}{3.075278in}%
\pgfsys@useobject{currentmarker}{}%
\end{pgfscope}%
\begin{pgfscope}%
\pgfsys@transformshift{0.940223in}{3.075278in}%
\pgfsys@useobject{currentmarker}{}%
\end{pgfscope}%
\begin{pgfscope}%
\pgfsys@transformshift{0.940223in}{3.075278in}%
\pgfsys@useobject{currentmarker}{}%
\end{pgfscope}%
\begin{pgfscope}%
\pgfsys@transformshift{0.940392in}{3.075278in}%
\pgfsys@useobject{currentmarker}{}%
\end{pgfscope}%
\begin{pgfscope}%
\pgfsys@transformshift{0.940560in}{3.075278in}%
\pgfsys@useobject{currentmarker}{}%
\end{pgfscope}%
\begin{pgfscope}%
\pgfsys@transformshift{0.940560in}{3.075278in}%
\pgfsys@useobject{currentmarker}{}%
\end{pgfscope}%
\begin{pgfscope}%
\pgfsys@transformshift{0.940729in}{3.075278in}%
\pgfsys@useobject{currentmarker}{}%
\end{pgfscope}%
\begin{pgfscope}%
\pgfsys@transformshift{0.940898in}{3.075278in}%
\pgfsys@useobject{currentmarker}{}%
\end{pgfscope}%
\begin{pgfscope}%
\pgfsys@transformshift{0.940898in}{3.075278in}%
\pgfsys@useobject{currentmarker}{}%
\end{pgfscope}%
\begin{pgfscope}%
\pgfsys@transformshift{0.941066in}{3.075278in}%
\pgfsys@useobject{currentmarker}{}%
\end{pgfscope}%
\begin{pgfscope}%
\pgfsys@transformshift{0.941235in}{3.075278in}%
\pgfsys@useobject{currentmarker}{}%
\end{pgfscope}%
\begin{pgfscope}%
\pgfsys@transformshift{0.941235in}{3.075278in}%
\pgfsys@useobject{currentmarker}{}%
\end{pgfscope}%
\begin{pgfscope}%
\pgfsys@transformshift{0.941403in}{3.075278in}%
\pgfsys@useobject{currentmarker}{}%
\end{pgfscope}%
\begin{pgfscope}%
\pgfsys@transformshift{0.941573in}{3.075278in}%
\pgfsys@useobject{currentmarker}{}%
\end{pgfscope}%
\begin{pgfscope}%
\pgfsys@transformshift{0.941573in}{3.075278in}%
\pgfsys@useobject{currentmarker}{}%
\end{pgfscope}%
\begin{pgfscope}%
\pgfsys@transformshift{0.941573in}{3.075278in}%
\pgfsys@useobject{currentmarker}{}%
\end{pgfscope}%
\begin{pgfscope}%
\pgfsys@transformshift{0.941741in}{3.075278in}%
\pgfsys@useobject{currentmarker}{}%
\end{pgfscope}%
\begin{pgfscope}%
\pgfsys@transformshift{0.941909in}{3.075278in}%
\pgfsys@useobject{currentmarker}{}%
\end{pgfscope}%
\begin{pgfscope}%
\pgfsys@transformshift{0.941909in}{3.075278in}%
\pgfsys@useobject{currentmarker}{}%
\end{pgfscope}%
\begin{pgfscope}%
\pgfsys@transformshift{0.941909in}{3.075278in}%
\pgfsys@useobject{currentmarker}{}%
\end{pgfscope}%
\begin{pgfscope}%
\pgfsys@transformshift{0.942078in}{3.075278in}%
\pgfsys@useobject{currentmarker}{}%
\end{pgfscope}%
\begin{pgfscope}%
\pgfsys@transformshift{0.942247in}{3.075278in}%
\pgfsys@useobject{currentmarker}{}%
\end{pgfscope}%
\begin{pgfscope}%
\pgfsys@transformshift{0.942247in}{3.075278in}%
\pgfsys@useobject{currentmarker}{}%
\end{pgfscope}%
\begin{pgfscope}%
\pgfsys@transformshift{0.942247in}{3.075278in}%
\pgfsys@useobject{currentmarker}{}%
\end{pgfscope}%
\begin{pgfscope}%
\pgfsys@transformshift{0.942416in}{3.075278in}%
\pgfsys@useobject{currentmarker}{}%
\end{pgfscope}%
\begin{pgfscope}%
\pgfsys@transformshift{0.942584in}{3.075278in}%
\pgfsys@useobject{currentmarker}{}%
\end{pgfscope}%
\begin{pgfscope}%
\pgfsys@transformshift{0.942584in}{3.075278in}%
\pgfsys@useobject{currentmarker}{}%
\end{pgfscope}%
\begin{pgfscope}%
\pgfsys@transformshift{0.942584in}{3.075278in}%
\pgfsys@useobject{currentmarker}{}%
\end{pgfscope}%
\begin{pgfscope}%
\pgfsys@transformshift{0.942752in}{3.075278in}%
\pgfsys@useobject{currentmarker}{}%
\end{pgfscope}%
\begin{pgfscope}%
\pgfsys@transformshift{0.942752in}{3.075278in}%
\pgfsys@useobject{currentmarker}{}%
\end{pgfscope}%
\begin{pgfscope}%
\pgfsys@transformshift{0.942921in}{3.075278in}%
\pgfsys@useobject{currentmarker}{}%
\end{pgfscope}%
\begin{pgfscope}%
\pgfsys@transformshift{0.942921in}{3.075278in}%
\pgfsys@useobject{currentmarker}{}%
\end{pgfscope}%
\begin{pgfscope}%
\pgfsys@transformshift{0.942921in}{3.075278in}%
\pgfsys@useobject{currentmarker}{}%
\end{pgfscope}%
\begin{pgfscope}%
\pgfsys@transformshift{0.943090in}{3.075278in}%
\pgfsys@useobject{currentmarker}{}%
\end{pgfscope}%
\begin{pgfscope}%
\pgfsys@transformshift{0.943090in}{3.075278in}%
\pgfsys@useobject{currentmarker}{}%
\end{pgfscope}%
\begin{pgfscope}%
\pgfsys@transformshift{0.943259in}{3.075278in}%
\pgfsys@useobject{currentmarker}{}%
\end{pgfscope}%
\begin{pgfscope}%
\pgfsys@transformshift{0.943259in}{3.075278in}%
\pgfsys@useobject{currentmarker}{}%
\end{pgfscope}%
\begin{pgfscope}%
\pgfsys@transformshift{0.943259in}{3.075278in}%
\pgfsys@useobject{currentmarker}{}%
\end{pgfscope}%
\begin{pgfscope}%
\pgfsys@transformshift{0.943427in}{3.075278in}%
\pgfsys@useobject{currentmarker}{}%
\end{pgfscope}%
\begin{pgfscope}%
\pgfsys@transformshift{0.943427in}{3.075278in}%
\pgfsys@useobject{currentmarker}{}%
\end{pgfscope}%
\begin{pgfscope}%
\pgfsys@transformshift{0.943595in}{3.075278in}%
\pgfsys@useobject{currentmarker}{}%
\end{pgfscope}%
\begin{pgfscope}%
\pgfsys@transformshift{0.943595in}{3.075278in}%
\pgfsys@useobject{currentmarker}{}%
\end{pgfscope}%
\begin{pgfscope}%
\pgfsys@transformshift{0.943595in}{3.075278in}%
\pgfsys@useobject{currentmarker}{}%
\end{pgfscope}%
\begin{pgfscope}%
\pgfsys@transformshift{0.943765in}{3.075278in}%
\pgfsys@useobject{currentmarker}{}%
\end{pgfscope}%
\begin{pgfscope}%
\pgfsys@transformshift{0.943765in}{3.075278in}%
\pgfsys@useobject{currentmarker}{}%
\end{pgfscope}%
\begin{pgfscope}%
\pgfsys@transformshift{0.943933in}{3.075278in}%
\pgfsys@useobject{currentmarker}{}%
\end{pgfscope}%
\begin{pgfscope}%
\pgfsys@transformshift{0.943933in}{3.075278in}%
\pgfsys@useobject{currentmarker}{}%
\end{pgfscope}%
\begin{pgfscope}%
\pgfsys@transformshift{0.943933in}{3.075278in}%
\pgfsys@useobject{currentmarker}{}%
\end{pgfscope}%
\begin{pgfscope}%
\pgfsys@transformshift{0.943933in}{3.075278in}%
\pgfsys@useobject{currentmarker}{}%
\end{pgfscope}%
\begin{pgfscope}%
\pgfsys@transformshift{0.944270in}{3.075278in}%
\pgfsys@useobject{currentmarker}{}%
\end{pgfscope}%
\begin{pgfscope}%
\pgfsys@transformshift{0.944270in}{3.075278in}%
\pgfsys@useobject{currentmarker}{}%
\end{pgfscope}%
\begin{pgfscope}%
\pgfsys@transformshift{0.944270in}{3.075278in}%
\pgfsys@useobject{currentmarker}{}%
\end{pgfscope}%
\begin{pgfscope}%
\pgfsys@transformshift{0.944270in}{3.075278in}%
\pgfsys@useobject{currentmarker}{}%
\end{pgfscope}%
\begin{pgfscope}%
\pgfsys@transformshift{0.944608in}{3.075278in}%
\pgfsys@useobject{currentmarker}{}%
\end{pgfscope}%
\begin{pgfscope}%
\pgfsys@transformshift{0.944608in}{3.075278in}%
\pgfsys@useobject{currentmarker}{}%
\end{pgfscope}%
\begin{pgfscope}%
\pgfsys@transformshift{0.944608in}{3.075278in}%
\pgfsys@useobject{currentmarker}{}%
\end{pgfscope}%
\begin{pgfscope}%
\pgfsys@transformshift{0.944608in}{3.075278in}%
\pgfsys@useobject{currentmarker}{}%
\end{pgfscope}%
\begin{pgfscope}%
\pgfsys@transformshift{0.944945in}{3.075278in}%
\pgfsys@useobject{currentmarker}{}%
\end{pgfscope}%
\begin{pgfscope}%
\pgfsys@transformshift{0.944945in}{3.075278in}%
\pgfsys@useobject{currentmarker}{}%
\end{pgfscope}%
\begin{pgfscope}%
\pgfsys@transformshift{0.944945in}{3.075278in}%
\pgfsys@useobject{currentmarker}{}%
\end{pgfscope}%
\begin{pgfscope}%
\pgfsys@transformshift{0.944945in}{3.075278in}%
\pgfsys@useobject{currentmarker}{}%
\end{pgfscope}%
\begin{pgfscope}%
\pgfsys@transformshift{0.945282in}{3.075278in}%
\pgfsys@useobject{currentmarker}{}%
\end{pgfscope}%
\begin{pgfscope}%
\pgfsys@transformshift{0.945282in}{3.075278in}%
\pgfsys@useobject{currentmarker}{}%
\end{pgfscope}%
\begin{pgfscope}%
\pgfsys@transformshift{0.945282in}{3.075278in}%
\pgfsys@useobject{currentmarker}{}%
\end{pgfscope}%
\begin{pgfscope}%
\pgfsys@transformshift{0.945282in}{3.075278in}%
\pgfsys@useobject{currentmarker}{}%
\end{pgfscope}%
\begin{pgfscope}%
\pgfsys@transformshift{0.945619in}{3.075278in}%
\pgfsys@useobject{currentmarker}{}%
\end{pgfscope}%
\begin{pgfscope}%
\pgfsys@transformshift{0.945619in}{3.075278in}%
\pgfsys@useobject{currentmarker}{}%
\end{pgfscope}%
\begin{pgfscope}%
\pgfsys@transformshift{0.945619in}{3.075278in}%
\pgfsys@useobject{currentmarker}{}%
\end{pgfscope}%
\begin{pgfscope}%
\pgfsys@transformshift{0.945619in}{3.075278in}%
\pgfsys@useobject{currentmarker}{}%
\end{pgfscope}%
\begin{pgfscope}%
\pgfsys@transformshift{0.945957in}{3.075278in}%
\pgfsys@useobject{currentmarker}{}%
\end{pgfscope}%
\begin{pgfscope}%
\pgfsys@transformshift{0.945957in}{3.075278in}%
\pgfsys@useobject{currentmarker}{}%
\end{pgfscope}%
\begin{pgfscope}%
\pgfsys@transformshift{0.945957in}{3.075278in}%
\pgfsys@useobject{currentmarker}{}%
\end{pgfscope}%
\begin{pgfscope}%
\pgfsys@transformshift{0.945957in}{3.075278in}%
\pgfsys@useobject{currentmarker}{}%
\end{pgfscope}%
\begin{pgfscope}%
\pgfsys@transformshift{0.946294in}{3.075278in}%
\pgfsys@useobject{currentmarker}{}%
\end{pgfscope}%
\begin{pgfscope}%
\pgfsys@transformshift{0.946294in}{3.075278in}%
\pgfsys@useobject{currentmarker}{}%
\end{pgfscope}%
\begin{pgfscope}%
\pgfsys@transformshift{0.946294in}{3.075278in}%
\pgfsys@useobject{currentmarker}{}%
\end{pgfscope}%
\begin{pgfscope}%
\pgfsys@transformshift{0.946294in}{3.075278in}%
\pgfsys@useobject{currentmarker}{}%
\end{pgfscope}%
\begin{pgfscope}%
\pgfsys@transformshift{0.946631in}{3.075278in}%
\pgfsys@useobject{currentmarker}{}%
\end{pgfscope}%
\begin{pgfscope}%
\pgfsys@transformshift{0.946631in}{3.075278in}%
\pgfsys@useobject{currentmarker}{}%
\end{pgfscope}%
\begin{pgfscope}%
\pgfsys@transformshift{0.946631in}{3.075278in}%
\pgfsys@useobject{currentmarker}{}%
\end{pgfscope}%
\begin{pgfscope}%
\pgfsys@transformshift{0.946800in}{3.075278in}%
\pgfsys@useobject{currentmarker}{}%
\end{pgfscope}%
\begin{pgfscope}%
\pgfsys@transformshift{0.946968in}{3.075278in}%
\pgfsys@useobject{currentmarker}{}%
\end{pgfscope}%
\begin{pgfscope}%
\pgfsys@transformshift{0.946968in}{3.075278in}%
\pgfsys@useobject{currentmarker}{}%
\end{pgfscope}%
\begin{pgfscope}%
\pgfsys@transformshift{0.947137in}{3.075278in}%
\pgfsys@useobject{currentmarker}{}%
\end{pgfscope}%
\begin{pgfscope}%
\pgfsys@transformshift{0.947305in}{3.075278in}%
\pgfsys@useobject{currentmarker}{}%
\end{pgfscope}%
\begin{pgfscope}%
\pgfsys@transformshift{0.947305in}{3.075278in}%
\pgfsys@useobject{currentmarker}{}%
\end{pgfscope}%
\begin{pgfscope}%
\pgfsys@transformshift{0.947475in}{3.075278in}%
\pgfsys@useobject{currentmarker}{}%
\end{pgfscope}%
\begin{pgfscope}%
\pgfsys@transformshift{0.947643in}{3.075278in}%
\pgfsys@useobject{currentmarker}{}%
\end{pgfscope}%
\begin{pgfscope}%
\pgfsys@transformshift{0.947643in}{3.075278in}%
\pgfsys@useobject{currentmarker}{}%
\end{pgfscope}%
\begin{pgfscope}%
\pgfsys@transformshift{0.947811in}{3.075278in}%
\pgfsys@useobject{currentmarker}{}%
\end{pgfscope}%
\begin{pgfscope}%
\pgfsys@transformshift{0.947980in}{3.075278in}%
\pgfsys@useobject{currentmarker}{}%
\end{pgfscope}%
\begin{pgfscope}%
\pgfsys@transformshift{0.947980in}{3.075278in}%
\pgfsys@useobject{currentmarker}{}%
\end{pgfscope}%
\begin{pgfscope}%
\pgfsys@transformshift{0.948149in}{3.075278in}%
\pgfsys@useobject{currentmarker}{}%
\end{pgfscope}%
\begin{pgfscope}%
\pgfsys@transformshift{0.948318in}{3.075278in}%
\pgfsys@useobject{currentmarker}{}%
\end{pgfscope}%
\begin{pgfscope}%
\pgfsys@transformshift{0.948318in}{3.075278in}%
\pgfsys@useobject{currentmarker}{}%
\end{pgfscope}%
\begin{pgfscope}%
\pgfsys@transformshift{0.948318in}{3.075278in}%
\pgfsys@useobject{currentmarker}{}%
\end{pgfscope}%
\begin{pgfscope}%
\pgfsys@transformshift{0.948486in}{3.075278in}%
\pgfsys@useobject{currentmarker}{}%
\end{pgfscope}%
\begin{pgfscope}%
\pgfsys@transformshift{0.948654in}{3.075278in}%
\pgfsys@useobject{currentmarker}{}%
\end{pgfscope}%
\begin{pgfscope}%
\pgfsys@transformshift{0.948654in}{3.075278in}%
\pgfsys@useobject{currentmarker}{}%
\end{pgfscope}%
\begin{pgfscope}%
\pgfsys@transformshift{0.948654in}{3.075278in}%
\pgfsys@useobject{currentmarker}{}%
\end{pgfscope}%
\begin{pgfscope}%
\pgfsys@transformshift{0.948823in}{3.075278in}%
\pgfsys@useobject{currentmarker}{}%
\end{pgfscope}%
\begin{pgfscope}%
\pgfsys@transformshift{0.948992in}{3.075278in}%
\pgfsys@useobject{currentmarker}{}%
\end{pgfscope}%
\begin{pgfscope}%
\pgfsys@transformshift{0.948992in}{3.075278in}%
\pgfsys@useobject{currentmarker}{}%
\end{pgfscope}%
\begin{pgfscope}%
\pgfsys@transformshift{0.948992in}{3.075278in}%
\pgfsys@useobject{currentmarker}{}%
\end{pgfscope}%
\begin{pgfscope}%
\pgfsys@transformshift{0.949161in}{3.075278in}%
\pgfsys@useobject{currentmarker}{}%
\end{pgfscope}%
\begin{pgfscope}%
\pgfsys@transformshift{0.949329in}{3.075278in}%
\pgfsys@useobject{currentmarker}{}%
\end{pgfscope}%
\begin{pgfscope}%
\pgfsys@transformshift{0.949329in}{3.075278in}%
\pgfsys@useobject{currentmarker}{}%
\end{pgfscope}%
\begin{pgfscope}%
\pgfsys@transformshift{0.949329in}{3.075278in}%
\pgfsys@useobject{currentmarker}{}%
\end{pgfscope}%
\begin{pgfscope}%
\pgfsys@transformshift{0.949497in}{3.075278in}%
\pgfsys@useobject{currentmarker}{}%
\end{pgfscope}%
\begin{pgfscope}%
\pgfsys@transformshift{0.949497in}{3.075278in}%
\pgfsys@useobject{currentmarker}{}%
\end{pgfscope}%
\begin{pgfscope}%
\pgfsys@transformshift{0.949667in}{3.075278in}%
\pgfsys@useobject{currentmarker}{}%
\end{pgfscope}%
\begin{pgfscope}%
\pgfsys@transformshift{0.949667in}{3.075278in}%
\pgfsys@useobject{currentmarker}{}%
\end{pgfscope}%
\begin{pgfscope}%
\pgfsys@transformshift{0.949667in}{3.075278in}%
\pgfsys@useobject{currentmarker}{}%
\end{pgfscope}%
\begin{pgfscope}%
\pgfsys@transformshift{0.949835in}{3.075278in}%
\pgfsys@useobject{currentmarker}{}%
\end{pgfscope}%
\begin{pgfscope}%
\pgfsys@transformshift{0.949835in}{3.075278in}%
\pgfsys@useobject{currentmarker}{}%
\end{pgfscope}%
\begin{pgfscope}%
\pgfsys@transformshift{0.950004in}{3.075278in}%
\pgfsys@useobject{currentmarker}{}%
\end{pgfscope}%
\begin{pgfscope}%
\pgfsys@transformshift{0.950004in}{3.075278in}%
\pgfsys@useobject{currentmarker}{}%
\end{pgfscope}%
\begin{pgfscope}%
\pgfsys@transformshift{0.950004in}{3.075278in}%
\pgfsys@useobject{currentmarker}{}%
\end{pgfscope}%
\begin{pgfscope}%
\pgfsys@transformshift{0.950172in}{3.075278in}%
\pgfsys@useobject{currentmarker}{}%
\end{pgfscope}%
\begin{pgfscope}%
\pgfsys@transformshift{0.950172in}{3.075278in}%
\pgfsys@useobject{currentmarker}{}%
\end{pgfscope}%
\begin{pgfscope}%
\pgfsys@transformshift{0.950341in}{3.075278in}%
\pgfsys@useobject{currentmarker}{}%
\end{pgfscope}%
\begin{pgfscope}%
\pgfsys@transformshift{0.950341in}{3.075278in}%
\pgfsys@useobject{currentmarker}{}%
\end{pgfscope}%
\begin{pgfscope}%
\pgfsys@transformshift{0.950341in}{3.075278in}%
\pgfsys@useobject{currentmarker}{}%
\end{pgfscope}%
\begin{pgfscope}%
\pgfsys@transformshift{0.950510in}{3.075278in}%
\pgfsys@useobject{currentmarker}{}%
\end{pgfscope}%
\begin{pgfscope}%
\pgfsys@transformshift{0.950510in}{3.075278in}%
\pgfsys@useobject{currentmarker}{}%
\end{pgfscope}%
\begin{pgfscope}%
\pgfsys@transformshift{0.950678in}{3.075278in}%
\pgfsys@useobject{currentmarker}{}%
\end{pgfscope}%
\begin{pgfscope}%
\pgfsys@transformshift{0.950678in}{3.075278in}%
\pgfsys@useobject{currentmarker}{}%
\end{pgfscope}%
\begin{pgfscope}%
\pgfsys@transformshift{0.950678in}{3.075278in}%
\pgfsys@useobject{currentmarker}{}%
\end{pgfscope}%
\begin{pgfscope}%
\pgfsys@transformshift{0.950678in}{3.075278in}%
\pgfsys@useobject{currentmarker}{}%
\end{pgfscope}%
\begin{pgfscope}%
\pgfsys@transformshift{0.951015in}{3.075278in}%
\pgfsys@useobject{currentmarker}{}%
\end{pgfscope}%
\begin{pgfscope}%
\pgfsys@transformshift{0.951015in}{3.075278in}%
\pgfsys@useobject{currentmarker}{}%
\end{pgfscope}%
\begin{pgfscope}%
\pgfsys@transformshift{0.951015in}{3.075278in}%
\pgfsys@useobject{currentmarker}{}%
\end{pgfscope}%
\begin{pgfscope}%
\pgfsys@transformshift{0.951015in}{3.075278in}%
\pgfsys@useobject{currentmarker}{}%
\end{pgfscope}%
\begin{pgfscope}%
\pgfsys@transformshift{0.951353in}{3.075278in}%
\pgfsys@useobject{currentmarker}{}%
\end{pgfscope}%
\begin{pgfscope}%
\pgfsys@transformshift{0.951353in}{3.075278in}%
\pgfsys@useobject{currentmarker}{}%
\end{pgfscope}%
\begin{pgfscope}%
\pgfsys@transformshift{0.951353in}{3.075278in}%
\pgfsys@useobject{currentmarker}{}%
\end{pgfscope}%
\begin{pgfscope}%
\pgfsys@transformshift{0.951353in}{3.075278in}%
\pgfsys@useobject{currentmarker}{}%
\end{pgfscope}%
\begin{pgfscope}%
\pgfsys@transformshift{0.951690in}{3.075278in}%
\pgfsys@useobject{currentmarker}{}%
\end{pgfscope}%
\begin{pgfscope}%
\pgfsys@transformshift{0.951690in}{3.075278in}%
\pgfsys@useobject{currentmarker}{}%
\end{pgfscope}%
\begin{pgfscope}%
\pgfsys@transformshift{0.951690in}{3.075278in}%
\pgfsys@useobject{currentmarker}{}%
\end{pgfscope}%
\begin{pgfscope}%
\pgfsys@transformshift{0.951690in}{3.075278in}%
\pgfsys@useobject{currentmarker}{}%
\end{pgfscope}%
\begin{pgfscope}%
\pgfsys@transformshift{0.952027in}{3.075278in}%
\pgfsys@useobject{currentmarker}{}%
\end{pgfscope}%
\begin{pgfscope}%
\pgfsys@transformshift{0.952027in}{3.075278in}%
\pgfsys@useobject{currentmarker}{}%
\end{pgfscope}%
\begin{pgfscope}%
\pgfsys@transformshift{0.952027in}{3.075278in}%
\pgfsys@useobject{currentmarker}{}%
\end{pgfscope}%
\begin{pgfscope}%
\pgfsys@transformshift{0.952027in}{3.075278in}%
\pgfsys@useobject{currentmarker}{}%
\end{pgfscope}%
\begin{pgfscope}%
\pgfsys@transformshift{0.952364in}{3.075278in}%
\pgfsys@useobject{currentmarker}{}%
\end{pgfscope}%
\begin{pgfscope}%
\pgfsys@transformshift{0.952364in}{3.075278in}%
\pgfsys@useobject{currentmarker}{}%
\end{pgfscope}%
\begin{pgfscope}%
\pgfsys@transformshift{0.952364in}{3.075278in}%
\pgfsys@useobject{currentmarker}{}%
\end{pgfscope}%
\begin{pgfscope}%
\pgfsys@transformshift{0.952364in}{3.075278in}%
\pgfsys@useobject{currentmarker}{}%
\end{pgfscope}%
\begin{pgfscope}%
\pgfsys@transformshift{0.952702in}{3.075278in}%
\pgfsys@useobject{currentmarker}{}%
\end{pgfscope}%
\begin{pgfscope}%
\pgfsys@transformshift{0.952702in}{3.075278in}%
\pgfsys@useobject{currentmarker}{}%
\end{pgfscope}%
\begin{pgfscope}%
\pgfsys@transformshift{0.953039in}{3.075278in}%
\pgfsys@useobject{currentmarker}{}%
\end{pgfscope}%
\begin{pgfscope}%
\pgfsys@transformshift{0.953039in}{3.075278in}%
\pgfsys@useobject{currentmarker}{}%
\end{pgfscope}%
\begin{pgfscope}%
\pgfsys@transformshift{0.953377in}{3.075278in}%
\pgfsys@useobject{currentmarker}{}%
\end{pgfscope}%
\begin{pgfscope}%
\pgfsys@transformshift{0.953713in}{3.075278in}%
\pgfsys@useobject{currentmarker}{}%
\end{pgfscope}%
\begin{pgfscope}%
\pgfsys@transformshift{0.954051in}{3.075278in}%
\pgfsys@useobject{currentmarker}{}%
\end{pgfscope}%
\begin{pgfscope}%
\pgfsys@transformshift{0.954388in}{3.075278in}%
\pgfsys@useobject{currentmarker}{}%
\end{pgfscope}%
\begin{pgfscope}%
\pgfsys@transformshift{0.954725in}{3.075278in}%
\pgfsys@useobject{currentmarker}{}%
\end{pgfscope}%
\begin{pgfscope}%
\pgfsys@transformshift{0.955063in}{3.075278in}%
\pgfsys@useobject{currentmarker}{}%
\end{pgfscope}%
\begin{pgfscope}%
\pgfsys@transformshift{0.955063in}{3.075278in}%
\pgfsys@useobject{currentmarker}{}%
\end{pgfscope}%
\begin{pgfscope}%
\pgfsys@transformshift{0.955399in}{3.075278in}%
\pgfsys@useobject{currentmarker}{}%
\end{pgfscope}%
\begin{pgfscope}%
\pgfsys@transformshift{0.955399in}{3.075278in}%
\pgfsys@useobject{currentmarker}{}%
\end{pgfscope}%
\begin{pgfscope}%
\pgfsys@transformshift{0.955737in}{3.075278in}%
\pgfsys@useobject{currentmarker}{}%
\end{pgfscope}%
\begin{pgfscope}%
\pgfsys@transformshift{0.955737in}{3.075278in}%
\pgfsys@useobject{currentmarker}{}%
\end{pgfscope}%
\begin{pgfscope}%
\pgfsys@transformshift{0.956074in}{3.075278in}%
\pgfsys@useobject{currentmarker}{}%
\end{pgfscope}%
\begin{pgfscope}%
\pgfsys@transformshift{0.956074in}{3.075278in}%
\pgfsys@useobject{currentmarker}{}%
\end{pgfscope}%
\begin{pgfscope}%
\pgfsys@transformshift{0.956243in}{3.075278in}%
\pgfsys@useobject{currentmarker}{}%
\end{pgfscope}%
\begin{pgfscope}%
\pgfsys@transformshift{0.956412in}{3.075278in}%
\pgfsys@useobject{currentmarker}{}%
\end{pgfscope}%
\begin{pgfscope}%
\pgfsys@transformshift{0.956412in}{3.075278in}%
\pgfsys@useobject{currentmarker}{}%
\end{pgfscope}%
\begin{pgfscope}%
\pgfsys@transformshift{0.956580in}{3.075278in}%
\pgfsys@useobject{currentmarker}{}%
\end{pgfscope}%
\begin{pgfscope}%
\pgfsys@transformshift{0.956749in}{3.075278in}%
\pgfsys@useobject{currentmarker}{}%
\end{pgfscope}%
\begin{pgfscope}%
\pgfsys@transformshift{0.956749in}{3.075278in}%
\pgfsys@useobject{currentmarker}{}%
\end{pgfscope}%
\begin{pgfscope}%
\pgfsys@transformshift{0.956917in}{3.075278in}%
\pgfsys@useobject{currentmarker}{}%
\end{pgfscope}%
\begin{pgfscope}%
\pgfsys@transformshift{0.957086in}{3.075278in}%
\pgfsys@useobject{currentmarker}{}%
\end{pgfscope}%
\begin{pgfscope}%
\pgfsys@transformshift{0.957086in}{3.075278in}%
\pgfsys@useobject{currentmarker}{}%
\end{pgfscope}%
\begin{pgfscope}%
\pgfsys@transformshift{0.957255in}{3.075278in}%
\pgfsys@useobject{currentmarker}{}%
\end{pgfscope}%
\begin{pgfscope}%
\pgfsys@transformshift{0.957423in}{3.075278in}%
\pgfsys@useobject{currentmarker}{}%
\end{pgfscope}%
\begin{pgfscope}%
\pgfsys@transformshift{0.957423in}{3.075278in}%
\pgfsys@useobject{currentmarker}{}%
\end{pgfscope}%
\begin{pgfscope}%
\pgfsys@transformshift{0.957423in}{3.075278in}%
\pgfsys@useobject{currentmarker}{}%
\end{pgfscope}%
\begin{pgfscope}%
\pgfsys@transformshift{0.957761in}{3.075278in}%
\pgfsys@useobject{currentmarker}{}%
\end{pgfscope}%
\begin{pgfscope}%
\pgfsys@transformshift{0.957761in}{3.075278in}%
\pgfsys@useobject{currentmarker}{}%
\end{pgfscope}%
\begin{pgfscope}%
\pgfsys@transformshift{0.958098in}{3.075278in}%
\pgfsys@useobject{currentmarker}{}%
\end{pgfscope}%
\begin{pgfscope}%
\pgfsys@transformshift{0.958098in}{3.075278in}%
\pgfsys@useobject{currentmarker}{}%
\end{pgfscope}%
\begin{pgfscope}%
\pgfsys@transformshift{0.958435in}{3.075278in}%
\pgfsys@useobject{currentmarker}{}%
\end{pgfscope}%
\begin{pgfscope}%
\pgfsys@transformshift{0.958435in}{3.075278in}%
\pgfsys@useobject{currentmarker}{}%
\end{pgfscope}%
\begin{pgfscope}%
\pgfsys@transformshift{0.958772in}{3.075278in}%
\pgfsys@useobject{currentmarker}{}%
\end{pgfscope}%
\begin{pgfscope}%
\pgfsys@transformshift{0.958772in}{3.075278in}%
\pgfsys@useobject{currentmarker}{}%
\end{pgfscope}%
\begin{pgfscope}%
\pgfsys@transformshift{0.959109in}{3.075278in}%
\pgfsys@useobject{currentmarker}{}%
\end{pgfscope}%
\begin{pgfscope}%
\pgfsys@transformshift{0.959109in}{3.075278in}%
\pgfsys@useobject{currentmarker}{}%
\end{pgfscope}%
\begin{pgfscope}%
\pgfsys@transformshift{0.959447in}{3.075278in}%
\pgfsys@useobject{currentmarker}{}%
\end{pgfscope}%
\begin{pgfscope}%
\pgfsys@transformshift{0.959447in}{3.075278in}%
\pgfsys@useobject{currentmarker}{}%
\end{pgfscope}%
\begin{pgfscope}%
\pgfsys@transformshift{0.959784in}{3.075278in}%
\pgfsys@useobject{currentmarker}{}%
\end{pgfscope}%
\begin{pgfscope}%
\pgfsys@transformshift{0.959784in}{3.075278in}%
\pgfsys@useobject{currentmarker}{}%
\end{pgfscope}%
\begin{pgfscope}%
\pgfsys@transformshift{0.960122in}{3.075278in}%
\pgfsys@useobject{currentmarker}{}%
\end{pgfscope}%
\begin{pgfscope}%
\pgfsys@transformshift{0.960458in}{3.075278in}%
\pgfsys@useobject{currentmarker}{}%
\end{pgfscope}%
\begin{pgfscope}%
\pgfsys@transformshift{0.960796in}{3.075278in}%
\pgfsys@useobject{currentmarker}{}%
\end{pgfscope}%
\begin{pgfscope}%
\pgfsys@transformshift{0.961133in}{3.075278in}%
\pgfsys@useobject{currentmarker}{}%
\end{pgfscope}%
\begin{pgfscope}%
\pgfsys@transformshift{0.961471in}{3.075278in}%
\pgfsys@useobject{currentmarker}{}%
\end{pgfscope}%
\begin{pgfscope}%
\pgfsys@transformshift{0.961808in}{3.075278in}%
\pgfsys@useobject{currentmarker}{}%
\end{pgfscope}%
\begin{pgfscope}%
\pgfsys@transformshift{0.961808in}{3.075278in}%
\pgfsys@useobject{currentmarker}{}%
\end{pgfscope}%
\begin{pgfscope}%
\pgfsys@transformshift{0.962145in}{3.075278in}%
\pgfsys@useobject{currentmarker}{}%
\end{pgfscope}%
\begin{pgfscope}%
\pgfsys@transformshift{0.962145in}{3.075278in}%
\pgfsys@useobject{currentmarker}{}%
\end{pgfscope}%
\begin{pgfscope}%
\pgfsys@transformshift{0.962988in}{3.075278in}%
\pgfsys@useobject{currentmarker}{}%
\end{pgfscope}%
\begin{pgfscope}%
\pgfsys@transformshift{0.963494in}{3.075278in}%
\pgfsys@useobject{currentmarker}{}%
\end{pgfscope}%
\begin{pgfscope}%
\pgfsys@transformshift{0.968553in}{3.075278in}%
\pgfsys@useobject{currentmarker}{}%
\end{pgfscope}%
\begin{pgfscope}%
\pgfsys@transformshift{0.968736in}{3.075278in}%
\pgfsys@useobject{currentmarker}{}%
\end{pgfscope}%
\begin{pgfscope}%
\pgfsys@transformshift{0.969565in}{3.075278in}%
\pgfsys@useobject{currentmarker}{}%
\end{pgfscope}%
\begin{pgfscope}%
\pgfsys@transformshift{0.971595in}{3.075278in}%
\pgfsys@useobject{currentmarker}{}%
\end{pgfscope}%
\begin{pgfscope}%
\pgfsys@transformshift{0.971940in}{3.075278in}%
\pgfsys@useobject{currentmarker}{}%
\end{pgfscope}%
\begin{pgfscope}%
\pgfsys@transformshift{0.972951in}{3.075278in}%
\pgfsys@useobject{currentmarker}{}%
\end{pgfscope}%
\begin{pgfscope}%
\pgfsys@transformshift{0.974308in}{3.075278in}%
\pgfsys@useobject{currentmarker}{}%
\end{pgfscope}%
\begin{pgfscope}%
\pgfsys@transformshift{0.974645in}{3.075278in}%
\pgfsys@useobject{currentmarker}{}%
\end{pgfscope}%
\begin{pgfscope}%
\pgfsys@transformshift{0.974813in}{3.075278in}%
\pgfsys@useobject{currentmarker}{}%
\end{pgfscope}%
\begin{pgfscope}%
\pgfsys@transformshift{0.974982in}{3.075278in}%
\pgfsys@useobject{currentmarker}{}%
\end{pgfscope}%
\begin{pgfscope}%
\pgfsys@transformshift{0.975143in}{3.075278in}%
\pgfsys@useobject{currentmarker}{}%
\end{pgfscope}%
\begin{pgfscope}%
\pgfsys@transformshift{0.975319in}{3.075278in}%
\pgfsys@useobject{currentmarker}{}%
\end{pgfscope}%
\begin{pgfscope}%
\pgfsys@transformshift{0.975488in}{3.075278in}%
\pgfsys@useobject{currentmarker}{}%
\end{pgfscope}%
\begin{pgfscope}%
\pgfsys@transformshift{0.975994in}{3.075278in}%
\pgfsys@useobject{currentmarker}{}%
\end{pgfscope}%
\begin{pgfscope}%
\pgfsys@transformshift{0.976162in}{3.075278in}%
\pgfsys@useobject{currentmarker}{}%
\end{pgfscope}%
\begin{pgfscope}%
\pgfsys@transformshift{0.976500in}{3.075278in}%
\pgfsys@useobject{currentmarker}{}%
\end{pgfscope}%
\begin{pgfscope}%
\pgfsys@transformshift{0.977174in}{3.075278in}%
\pgfsys@useobject{currentmarker}{}%
\end{pgfscope}%
\begin{pgfscope}%
\pgfsys@transformshift{0.977343in}{3.075278in}%
\pgfsys@useobject{currentmarker}{}%
\end{pgfscope}%
\begin{pgfscope}%
\pgfsys@transformshift{0.977512in}{3.075278in}%
\pgfsys@useobject{currentmarker}{}%
\end{pgfscope}%
\begin{pgfscope}%
\pgfsys@transformshift{0.977512in}{3.075278in}%
\pgfsys@useobject{currentmarker}{}%
\end{pgfscope}%
\begin{pgfscope}%
\pgfsys@transformshift{0.977680in}{3.075278in}%
\pgfsys@useobject{currentmarker}{}%
\end{pgfscope}%
\begin{pgfscope}%
\pgfsys@transformshift{0.977848in}{3.075278in}%
\pgfsys@useobject{currentmarker}{}%
\end{pgfscope}%
\begin{pgfscope}%
\pgfsys@transformshift{0.977848in}{3.075278in}%
\pgfsys@useobject{currentmarker}{}%
\end{pgfscope}%
\begin{pgfscope}%
\pgfsys@transformshift{0.978186in}{3.075278in}%
\pgfsys@useobject{currentmarker}{}%
\end{pgfscope}%
\begin{pgfscope}%
\pgfsys@transformshift{0.978186in}{3.075278in}%
\pgfsys@useobject{currentmarker}{}%
\end{pgfscope}%
\begin{pgfscope}%
\pgfsys@transformshift{0.978355in}{3.075278in}%
\pgfsys@useobject{currentmarker}{}%
\end{pgfscope}%
\begin{pgfscope}%
\pgfsys@transformshift{0.978523in}{3.075278in}%
\pgfsys@useobject{currentmarker}{}%
\end{pgfscope}%
\begin{pgfscope}%
\pgfsys@transformshift{0.978523in}{3.075278in}%
\pgfsys@useobject{currentmarker}{}%
\end{pgfscope}%
\begin{pgfscope}%
\pgfsys@transformshift{0.978861in}{3.075278in}%
\pgfsys@useobject{currentmarker}{}%
\end{pgfscope}%
\begin{pgfscope}%
\pgfsys@transformshift{0.978861in}{3.075278in}%
\pgfsys@useobject{currentmarker}{}%
\end{pgfscope}%
\begin{pgfscope}%
\pgfsys@transformshift{0.978861in}{3.075278in}%
\pgfsys@useobject{currentmarker}{}%
\end{pgfscope}%
\begin{pgfscope}%
\pgfsys@transformshift{0.979198in}{3.075278in}%
\pgfsys@useobject{currentmarker}{}%
\end{pgfscope}%
\begin{pgfscope}%
\pgfsys@transformshift{0.979198in}{3.075278in}%
\pgfsys@useobject{currentmarker}{}%
\end{pgfscope}%
\begin{pgfscope}%
\pgfsys@transformshift{0.979198in}{3.075278in}%
\pgfsys@useobject{currentmarker}{}%
\end{pgfscope}%
\begin{pgfscope}%
\pgfsys@transformshift{0.979535in}{3.075278in}%
\pgfsys@useobject{currentmarker}{}%
\end{pgfscope}%
\begin{pgfscope}%
\pgfsys@transformshift{0.979535in}{3.075278in}%
\pgfsys@useobject{currentmarker}{}%
\end{pgfscope}%
\begin{pgfscope}%
\pgfsys@transformshift{0.979704in}{3.075278in}%
\pgfsys@useobject{currentmarker}{}%
\end{pgfscope}%
\begin{pgfscope}%
\pgfsys@transformshift{0.979872in}{3.075278in}%
\pgfsys@useobject{currentmarker}{}%
\end{pgfscope}%
\begin{pgfscope}%
\pgfsys@transformshift{0.980041in}{3.075278in}%
\pgfsys@useobject{currentmarker}{}%
\end{pgfscope}%
\begin{pgfscope}%
\pgfsys@transformshift{0.980210in}{3.075278in}%
\pgfsys@useobject{currentmarker}{}%
\end{pgfscope}%
\begin{pgfscope}%
\pgfsys@transformshift{0.980378in}{3.075278in}%
\pgfsys@useobject{currentmarker}{}%
\end{pgfscope}%
\begin{pgfscope}%
\pgfsys@transformshift{0.980547in}{3.075278in}%
\pgfsys@useobject{currentmarker}{}%
\end{pgfscope}%
\begin{pgfscope}%
\pgfsys@transformshift{0.980884in}{3.075278in}%
\pgfsys@useobject{currentmarker}{}%
\end{pgfscope}%
\begin{pgfscope}%
\pgfsys@transformshift{0.981053in}{3.075278in}%
\pgfsys@useobject{currentmarker}{}%
\end{pgfscope}%
\begin{pgfscope}%
\pgfsys@transformshift{0.981390in}{3.075278in}%
\pgfsys@useobject{currentmarker}{}%
\end{pgfscope}%
\begin{pgfscope}%
\pgfsys@transformshift{0.981728in}{3.075278in}%
\pgfsys@useobject{currentmarker}{}%
\end{pgfscope}%
\begin{pgfscope}%
\pgfsys@transformshift{0.981896in}{3.075278in}%
\pgfsys@useobject{currentmarker}{}%
\end{pgfscope}%
\begin{pgfscope}%
\pgfsys@transformshift{0.981896in}{3.075278in}%
\pgfsys@useobject{currentmarker}{}%
\end{pgfscope}%
\begin{pgfscope}%
\pgfsys@transformshift{0.982064in}{3.075278in}%
\pgfsys@useobject{currentmarker}{}%
\end{pgfscope}%
\begin{pgfscope}%
\pgfsys@transformshift{0.982233in}{3.075278in}%
\pgfsys@useobject{currentmarker}{}%
\end{pgfscope}%
\begin{pgfscope}%
\pgfsys@transformshift{0.982571in}{3.075278in}%
\pgfsys@useobject{currentmarker}{}%
\end{pgfscope}%
\begin{pgfscope}%
\pgfsys@transformshift{0.982907in}{3.075278in}%
\pgfsys@useobject{currentmarker}{}%
\end{pgfscope}%
\begin{pgfscope}%
\pgfsys@transformshift{0.983076in}{3.075278in}%
\pgfsys@useobject{currentmarker}{}%
\end{pgfscope}%
\begin{pgfscope}%
\pgfsys@transformshift{0.983414in}{3.075278in}%
\pgfsys@useobject{currentmarker}{}%
\end{pgfscope}%
\begin{pgfscope}%
\pgfsys@transformshift{0.983582in}{3.075278in}%
\pgfsys@useobject{currentmarker}{}%
\end{pgfscope}%
\begin{pgfscope}%
\pgfsys@transformshift{0.983920in}{3.075278in}%
\pgfsys@useobject{currentmarker}{}%
\end{pgfscope}%
\begin{pgfscope}%
\pgfsys@transformshift{0.983920in}{3.075278in}%
\pgfsys@useobject{currentmarker}{}%
\end{pgfscope}%
\begin{pgfscope}%
\pgfsys@transformshift{0.984257in}{3.075278in}%
\pgfsys@useobject{currentmarker}{}%
\end{pgfscope}%
\begin{pgfscope}%
\pgfsys@transformshift{0.984425in}{3.075278in}%
\pgfsys@useobject{currentmarker}{}%
\end{pgfscope}%
\begin{pgfscope}%
\pgfsys@transformshift{0.984763in}{3.075278in}%
\pgfsys@useobject{currentmarker}{}%
\end{pgfscope}%
\begin{pgfscope}%
\pgfsys@transformshift{0.984931in}{3.075278in}%
\pgfsys@useobject{currentmarker}{}%
\end{pgfscope}%
\begin{pgfscope}%
\pgfsys@transformshift{0.985100in}{3.075278in}%
\pgfsys@useobject{currentmarker}{}%
\end{pgfscope}%
\begin{pgfscope}%
\pgfsys@transformshift{0.985268in}{3.075278in}%
\pgfsys@useobject{currentmarker}{}%
\end{pgfscope}%
\begin{pgfscope}%
\pgfsys@transformshift{0.985437in}{3.075278in}%
\pgfsys@useobject{currentmarker}{}%
\end{pgfscope}%
\begin{pgfscope}%
\pgfsys@transformshift{0.985774in}{3.075278in}%
\pgfsys@useobject{currentmarker}{}%
\end{pgfscope}%
\begin{pgfscope}%
\pgfsys@transformshift{0.985943in}{3.075278in}%
\pgfsys@useobject{currentmarker}{}%
\end{pgfscope}%
\begin{pgfscope}%
\pgfsys@transformshift{0.986280in}{3.075278in}%
\pgfsys@useobject{currentmarker}{}%
\end{pgfscope}%
\begin{pgfscope}%
\pgfsys@transformshift{0.986449in}{3.075278in}%
\pgfsys@useobject{currentmarker}{}%
\end{pgfscope}%
\begin{pgfscope}%
\pgfsys@transformshift{0.986786in}{3.075278in}%
\pgfsys@useobject{currentmarker}{}%
\end{pgfscope}%
\begin{pgfscope}%
\pgfsys@transformshift{0.986786in}{3.075278in}%
\pgfsys@useobject{currentmarker}{}%
\end{pgfscope}%
\begin{pgfscope}%
\pgfsys@transformshift{0.987460in}{3.075278in}%
\pgfsys@useobject{currentmarker}{}%
\end{pgfscope}%
\begin{pgfscope}%
\pgfsys@transformshift{0.987460in}{3.075278in}%
\pgfsys@useobject{currentmarker}{}%
\end{pgfscope}%
\begin{pgfscope}%
\pgfsys@transformshift{0.987798in}{3.075278in}%
\pgfsys@useobject{currentmarker}{}%
\end{pgfscope}%
\begin{pgfscope}%
\pgfsys@transformshift{0.988135in}{3.075278in}%
\pgfsys@useobject{currentmarker}{}%
\end{pgfscope}%
\begin{pgfscope}%
\pgfsys@transformshift{0.988135in}{3.075278in}%
\pgfsys@useobject{currentmarker}{}%
\end{pgfscope}%
\begin{pgfscope}%
\pgfsys@transformshift{0.988473in}{3.075278in}%
\pgfsys@useobject{currentmarker}{}%
\end{pgfscope}%
\begin{pgfscope}%
\pgfsys@transformshift{0.988809in}{3.075278in}%
\pgfsys@useobject{currentmarker}{}%
\end{pgfscope}%
\begin{pgfscope}%
\pgfsys@transformshift{0.988978in}{3.075278in}%
\pgfsys@useobject{currentmarker}{}%
\end{pgfscope}%
\begin{pgfscope}%
\pgfsys@transformshift{0.989147in}{3.075278in}%
\pgfsys@useobject{currentmarker}{}%
\end{pgfscope}%
\begin{pgfscope}%
\pgfsys@transformshift{0.989484in}{3.075278in}%
\pgfsys@useobject{currentmarker}{}%
\end{pgfscope}%
\begin{pgfscope}%
\pgfsys@transformshift{0.989484in}{3.075278in}%
\pgfsys@useobject{currentmarker}{}%
\end{pgfscope}%
\begin{pgfscope}%
\pgfsys@transformshift{0.989822in}{3.075278in}%
\pgfsys@useobject{currentmarker}{}%
\end{pgfscope}%
\begin{pgfscope}%
\pgfsys@transformshift{0.990159in}{3.075278in}%
\pgfsys@useobject{currentmarker}{}%
\end{pgfscope}%
\begin{pgfscope}%
\pgfsys@transformshift{0.990327in}{3.075278in}%
\pgfsys@useobject{currentmarker}{}%
\end{pgfscope}%
\begin{pgfscope}%
\pgfsys@transformshift{0.990665in}{3.075278in}%
\pgfsys@useobject{currentmarker}{}%
\end{pgfscope}%
\begin{pgfscope}%
\pgfsys@transformshift{0.990833in}{3.075278in}%
\pgfsys@useobject{currentmarker}{}%
\end{pgfscope}%
\begin{pgfscope}%
\pgfsys@transformshift{0.991170in}{3.075278in}%
\pgfsys@useobject{currentmarker}{}%
\end{pgfscope}%
\begin{pgfscope}%
\pgfsys@transformshift{0.991508in}{3.075278in}%
\pgfsys@useobject{currentmarker}{}%
\end{pgfscope}%
\begin{pgfscope}%
\pgfsys@transformshift{0.991676in}{3.075278in}%
\pgfsys@useobject{currentmarker}{}%
\end{pgfscope}%
\begin{pgfscope}%
\pgfsys@transformshift{0.992014in}{3.075278in}%
\pgfsys@useobject{currentmarker}{}%
\end{pgfscope}%
\begin{pgfscope}%
\pgfsys@transformshift{0.992182in}{3.075278in}%
\pgfsys@useobject{currentmarker}{}%
\end{pgfscope}%
\begin{pgfscope}%
\pgfsys@transformshift{0.992688in}{3.075278in}%
\pgfsys@useobject{currentmarker}{}%
\end{pgfscope}%
\begin{pgfscope}%
\pgfsys@transformshift{0.993025in}{3.075278in}%
\pgfsys@useobject{currentmarker}{}%
\end{pgfscope}%
\begin{pgfscope}%
\pgfsys@transformshift{0.993362in}{3.075278in}%
\pgfsys@useobject{currentmarker}{}%
\end{pgfscope}%
\begin{pgfscope}%
\pgfsys@transformshift{0.993700in}{3.075278in}%
\pgfsys@useobject{currentmarker}{}%
\end{pgfscope}%
\begin{pgfscope}%
\pgfsys@transformshift{0.994206in}{3.075278in}%
\pgfsys@useobject{currentmarker}{}%
\end{pgfscope}%
\begin{pgfscope}%
\pgfsys@transformshift{0.994711in}{3.075278in}%
\pgfsys@useobject{currentmarker}{}%
\end{pgfscope}%
\begin{pgfscope}%
\pgfsys@transformshift{0.995386in}{3.075278in}%
\pgfsys@useobject{currentmarker}{}%
\end{pgfscope}%
\begin{pgfscope}%
\pgfsys@transformshift{0.996229in}{3.075278in}%
\pgfsys@useobject{currentmarker}{}%
\end{pgfscope}%
\begin{pgfscope}%
\pgfsys@transformshift{0.996735in}{3.075278in}%
\pgfsys@useobject{currentmarker}{}%
\end{pgfscope}%
\begin{pgfscope}%
\pgfsys@transformshift{0.997578in}{3.075278in}%
\pgfsys@useobject{currentmarker}{}%
\end{pgfscope}%
\begin{pgfscope}%
\pgfsys@transformshift{0.998084in}{3.075278in}%
\pgfsys@useobject{currentmarker}{}%
\end{pgfscope}%
\begin{pgfscope}%
\pgfsys@transformshift{0.998421in}{3.075278in}%
\pgfsys@useobject{currentmarker}{}%
\end{pgfscope}%
\begin{pgfscope}%
\pgfsys@transformshift{0.998759in}{3.075278in}%
\pgfsys@useobject{currentmarker}{}%
\end{pgfscope}%
\begin{pgfscope}%
\pgfsys@transformshift{1.126221in}{3.075278in}%
\pgfsys@useobject{currentmarker}{}%
\end{pgfscope}%
\begin{pgfscope}%
\pgfsys@transformshift{1.127738in}{3.075278in}%
\pgfsys@useobject{currentmarker}{}%
\end{pgfscope}%
\begin{pgfscope}%
\pgfsys@transformshift{1.128638in}{3.075278in}%
\pgfsys@useobject{currentmarker}{}%
\end{pgfscope}%
\begin{pgfscope}%
\pgfsys@transformshift{1.129930in}{3.075278in}%
\pgfsys@useobject{currentmarker}{}%
\end{pgfscope}%
\begin{pgfscope}%
\pgfsys@transformshift{1.130605in}{3.075278in}%
\pgfsys@useobject{currentmarker}{}%
\end{pgfscope}%
\begin{pgfscope}%
\pgfsys@transformshift{1.131110in}{3.075278in}%
\pgfsys@useobject{currentmarker}{}%
\end{pgfscope}%
\begin{pgfscope}%
\pgfsys@transformshift{1.131785in}{3.075278in}%
\pgfsys@useobject{currentmarker}{}%
\end{pgfscope}%
\begin{pgfscope}%
\pgfsys@transformshift{1.132291in}{3.075278in}%
\pgfsys@useobject{currentmarker}{}%
\end{pgfscope}%
\begin{pgfscope}%
\pgfsys@transformshift{1.133134in}{3.075278in}%
\pgfsys@useobject{currentmarker}{}%
\end{pgfscope}%
\begin{pgfscope}%
\pgfsys@transformshift{1.133809in}{3.075278in}%
\pgfsys@useobject{currentmarker}{}%
\end{pgfscope}%
\begin{pgfscope}%
\pgfsys@transformshift{1.134483in}{3.075278in}%
\pgfsys@useobject{currentmarker}{}%
\end{pgfscope}%
\begin{pgfscope}%
\pgfsys@transformshift{1.134820in}{3.075278in}%
\pgfsys@useobject{currentmarker}{}%
\end{pgfscope}%
\begin{pgfscope}%
\pgfsys@transformshift{1.135158in}{3.075278in}%
\pgfsys@useobject{currentmarker}{}%
\end{pgfscope}%
\begin{pgfscope}%
\pgfsys@transformshift{1.136169in}{3.075278in}%
\pgfsys@useobject{currentmarker}{}%
\end{pgfscope}%
\begin{pgfscope}%
\pgfsys@transformshift{1.136338in}{3.075278in}%
\pgfsys@useobject{currentmarker}{}%
\end{pgfscope}%
\begin{pgfscope}%
\pgfsys@transformshift{1.136675in}{3.075278in}%
\pgfsys@useobject{currentmarker}{}%
\end{pgfscope}%
\begin{pgfscope}%
\pgfsys@transformshift{1.136844in}{3.075278in}%
\pgfsys@useobject{currentmarker}{}%
\end{pgfscope}%
\begin{pgfscope}%
\pgfsys@transformshift{1.137181in}{3.075278in}%
\pgfsys@useobject{currentmarker}{}%
\end{pgfscope}%
\begin{pgfscope}%
\pgfsys@transformshift{1.137855in}{3.075278in}%
\pgfsys@useobject{currentmarker}{}%
\end{pgfscope}%
\begin{pgfscope}%
\pgfsys@transformshift{1.138193in}{3.075278in}%
\pgfsys@useobject{currentmarker}{}%
\end{pgfscope}%
\begin{pgfscope}%
\pgfsys@transformshift{1.138193in}{3.075278in}%
\pgfsys@useobject{currentmarker}{}%
\end{pgfscope}%
\begin{pgfscope}%
\pgfsys@transformshift{1.138530in}{3.075278in}%
\pgfsys@useobject{currentmarker}{}%
\end{pgfscope}%
\begin{pgfscope}%
\pgfsys@transformshift{1.138537in}{3.075278in}%
\pgfsys@useobject{currentmarker}{}%
\end{pgfscope}%
\begin{pgfscope}%
\pgfsys@transformshift{1.138874in}{3.075278in}%
\pgfsys@useobject{currentmarker}{}%
\end{pgfscope}%
\begin{pgfscope}%
\pgfsys@transformshift{1.139212in}{3.075278in}%
\pgfsys@useobject{currentmarker}{}%
\end{pgfscope}%
\begin{pgfscope}%
\pgfsys@transformshift{1.139542in}{3.075278in}%
\pgfsys@useobject{currentmarker}{}%
\end{pgfscope}%
\begin{pgfscope}%
\pgfsys@transformshift{1.139879in}{3.075278in}%
\pgfsys@useobject{currentmarker}{}%
\end{pgfscope}%
\begin{pgfscope}%
\pgfsys@transformshift{1.139887in}{3.075278in}%
\pgfsys@useobject{currentmarker}{}%
\end{pgfscope}%
\begin{pgfscope}%
\pgfsys@transformshift{1.140217in}{3.075278in}%
\pgfsys@useobject{currentmarker}{}%
\end{pgfscope}%
\begin{pgfscope}%
\pgfsys@transformshift{1.140385in}{3.075278in}%
\pgfsys@useobject{currentmarker}{}%
\end{pgfscope}%
\begin{pgfscope}%
\pgfsys@transformshift{1.140561in}{3.075278in}%
\pgfsys@useobject{currentmarker}{}%
\end{pgfscope}%
\begin{pgfscope}%
\pgfsys@transformshift{1.140891in}{3.075278in}%
\pgfsys@useobject{currentmarker}{}%
\end{pgfscope}%
\begin{pgfscope}%
\pgfsys@transformshift{1.141228in}{3.075278in}%
\pgfsys@useobject{currentmarker}{}%
\end{pgfscope}%
\begin{pgfscope}%
\pgfsys@transformshift{1.141565in}{3.075278in}%
\pgfsys@useobject{currentmarker}{}%
\end{pgfscope}%
\begin{pgfscope}%
\pgfsys@transformshift{1.141903in}{3.075278in}%
\pgfsys@useobject{currentmarker}{}%
\end{pgfscope}%
\begin{pgfscope}%
\pgfsys@transformshift{1.142071in}{3.075278in}%
\pgfsys@useobject{currentmarker}{}%
\end{pgfscope}%
\begin{pgfscope}%
\pgfsys@transformshift{1.142247in}{3.075278in}%
\pgfsys@useobject{currentmarker}{}%
\end{pgfscope}%
\begin{pgfscope}%
\pgfsys@transformshift{1.142577in}{3.075278in}%
\pgfsys@useobject{currentmarker}{}%
\end{pgfscope}%
\begin{pgfscope}%
\pgfsys@transformshift{1.142584in}{3.075278in}%
\pgfsys@useobject{currentmarker}{}%
\end{pgfscope}%
\begin{pgfscope}%
\pgfsys@transformshift{1.142922in}{3.075278in}%
\pgfsys@useobject{currentmarker}{}%
\end{pgfscope}%
\begin{pgfscope}%
\pgfsys@transformshift{1.143589in}{3.075278in}%
\pgfsys@useobject{currentmarker}{}%
\end{pgfscope}%
\begin{pgfscope}%
\pgfsys@transformshift{1.143757in}{3.075278in}%
\pgfsys@useobject{currentmarker}{}%
\end{pgfscope}%
\begin{pgfscope}%
\pgfsys@transformshift{1.143933in}{3.075278in}%
\pgfsys@useobject{currentmarker}{}%
\end{pgfscope}%
\begin{pgfscope}%
\pgfsys@transformshift{1.144271in}{3.075278in}%
\pgfsys@useobject{currentmarker}{}%
\end{pgfscope}%
\begin{pgfscope}%
\pgfsys@transformshift{1.144432in}{3.075278in}%
\pgfsys@useobject{currentmarker}{}%
\end{pgfscope}%
\begin{pgfscope}%
\pgfsys@transformshift{1.144770in}{3.075278in}%
\pgfsys@useobject{currentmarker}{}%
\end{pgfscope}%
\begin{pgfscope}%
\pgfsys@transformshift{1.145106in}{3.075278in}%
\pgfsys@useobject{currentmarker}{}%
\end{pgfscope}%
\begin{pgfscope}%
\pgfsys@transformshift{1.145444in}{3.075278in}%
\pgfsys@useobject{currentmarker}{}%
\end{pgfscope}%
\begin{pgfscope}%
\pgfsys@transformshift{1.145781in}{3.075278in}%
\pgfsys@useobject{currentmarker}{}%
\end{pgfscope}%
\begin{pgfscope}%
\pgfsys@transformshift{1.145949in}{3.075278in}%
\pgfsys@useobject{currentmarker}{}%
\end{pgfscope}%
\begin{pgfscope}%
\pgfsys@transformshift{1.146456in}{3.075278in}%
\pgfsys@useobject{currentmarker}{}%
\end{pgfscope}%
\begin{pgfscope}%
\pgfsys@transformshift{1.146793in}{3.075278in}%
\pgfsys@useobject{currentmarker}{}%
\end{pgfscope}%
\begin{pgfscope}%
\pgfsys@transformshift{1.146962in}{3.075278in}%
\pgfsys@useobject{currentmarker}{}%
\end{pgfscope}%
\begin{pgfscope}%
\pgfsys@transformshift{1.147299in}{3.075278in}%
\pgfsys@useobject{currentmarker}{}%
\end{pgfscope}%
\begin{pgfscope}%
\pgfsys@transformshift{1.147636in}{3.075278in}%
\pgfsys@useobject{currentmarker}{}%
\end{pgfscope}%
\begin{pgfscope}%
\pgfsys@transformshift{1.147805in}{3.075278in}%
\pgfsys@useobject{currentmarker}{}%
\end{pgfscope}%
\begin{pgfscope}%
\pgfsys@transformshift{1.147973in}{3.075278in}%
\pgfsys@useobject{currentmarker}{}%
\end{pgfscope}%
\begin{pgfscope}%
\pgfsys@transformshift{1.148311in}{3.075278in}%
\pgfsys@useobject{currentmarker}{}%
\end{pgfscope}%
\begin{pgfscope}%
\pgfsys@transformshift{1.148985in}{3.075278in}%
\pgfsys@useobject{currentmarker}{}%
\end{pgfscope}%
\begin{pgfscope}%
\pgfsys@transformshift{1.149322in}{3.075278in}%
\pgfsys@useobject{currentmarker}{}%
\end{pgfscope}%
\begin{pgfscope}%
\pgfsys@transformshift{1.149659in}{3.075278in}%
\pgfsys@useobject{currentmarker}{}%
\end{pgfscope}%
\begin{pgfscope}%
\pgfsys@transformshift{1.149828in}{3.075278in}%
\pgfsys@useobject{currentmarker}{}%
\end{pgfscope}%
\begin{pgfscope}%
\pgfsys@transformshift{1.150165in}{3.075278in}%
\pgfsys@useobject{currentmarker}{}%
\end{pgfscope}%
\begin{pgfscope}%
\pgfsys@transformshift{1.150672in}{3.075278in}%
\pgfsys@useobject{currentmarker}{}%
\end{pgfscope}%
\begin{pgfscope}%
\pgfsys@transformshift{1.151008in}{3.075278in}%
\pgfsys@useobject{currentmarker}{}%
\end{pgfscope}%
\begin{pgfscope}%
\pgfsys@transformshift{1.151346in}{3.075278in}%
\pgfsys@useobject{currentmarker}{}%
\end{pgfscope}%
\begin{pgfscope}%
\pgfsys@transformshift{1.151683in}{3.075278in}%
\pgfsys@useobject{currentmarker}{}%
\end{pgfscope}%
\begin{pgfscope}%
\pgfsys@transformshift{1.151909in}{3.075278in}%
\pgfsys@useobject{currentmarker}{}%
\end{pgfscope}%
\begin{pgfscope}%
\pgfsys@transformshift{1.152020in}{3.075278in}%
\pgfsys@useobject{currentmarker}{}%
\end{pgfscope}%
\begin{pgfscope}%
\pgfsys@transformshift{1.152189in}{3.075278in}%
\pgfsys@useobject{currentmarker}{}%
\end{pgfscope}%
\begin{pgfscope}%
\pgfsys@transformshift{1.152358in}{3.075278in}%
\pgfsys@useobject{currentmarker}{}%
\end{pgfscope}%
\begin{pgfscope}%
\pgfsys@transformshift{1.152695in}{3.075278in}%
\pgfsys@useobject{currentmarker}{}%
\end{pgfscope}%
\begin{pgfscope}%
\pgfsys@transformshift{1.153039in}{3.075278in}%
\pgfsys@useobject{currentmarker}{}%
\end{pgfscope}%
\begin{pgfscope}%
\pgfsys@transformshift{1.153377in}{3.075278in}%
\pgfsys@useobject{currentmarker}{}%
\end{pgfscope}%
\begin{pgfscope}%
\pgfsys@transformshift{1.154051in}{3.075278in}%
\pgfsys@useobject{currentmarker}{}%
\end{pgfscope}%
\begin{pgfscope}%
\pgfsys@transformshift{1.155730in}{3.075278in}%
\pgfsys@useobject{currentmarker}{}%
\end{pgfscope}%
\begin{pgfscope}%
\pgfsys@transformshift{1.156742in}{3.075278in}%
\pgfsys@useobject{currentmarker}{}%
\end{pgfscope}%
\begin{pgfscope}%
\pgfsys@transformshift{1.158597in}{3.075278in}%
\pgfsys@useobject{currentmarker}{}%
\end{pgfscope}%
\begin{pgfscope}%
\pgfsys@transformshift{1.159953in}{3.075278in}%
\pgfsys@useobject{currentmarker}{}%
\end{pgfscope}%
\begin{pgfscope}%
\pgfsys@transformshift{1.160115in}{3.075278in}%
\pgfsys@useobject{currentmarker}{}%
\end{pgfscope}%
\begin{pgfscope}%
\pgfsys@transformshift{1.160459in}{3.075278in}%
\pgfsys@useobject{currentmarker}{}%
\end{pgfscope}%
\begin{pgfscope}%
\pgfsys@transformshift{1.160965in}{3.075278in}%
\pgfsys@useobject{currentmarker}{}%
\end{pgfscope}%
\begin{pgfscope}%
\pgfsys@transformshift{1.161302in}{3.075278in}%
\pgfsys@useobject{currentmarker}{}%
\end{pgfscope}%
\begin{pgfscope}%
\pgfsys@transformshift{1.162160in}{3.075278in}%
\pgfsys@useobject{currentmarker}{}%
\end{pgfscope}%
\begin{pgfscope}%
\pgfsys@transformshift{1.163509in}{3.075278in}%
\pgfsys@useobject{currentmarker}{}%
\end{pgfscope}%
\begin{pgfscope}%
\pgfsys@transformshift{1.164162in}{3.075278in}%
\pgfsys@useobject{currentmarker}{}%
\end{pgfscope}%
\begin{pgfscope}%
\pgfsys@transformshift{1.164499in}{3.075278in}%
\pgfsys@useobject{currentmarker}{}%
\end{pgfscope}%
\begin{pgfscope}%
\pgfsys@transformshift{1.164836in}{3.075278in}%
\pgfsys@useobject{currentmarker}{}%
\end{pgfscope}%
\begin{pgfscope}%
\pgfsys@transformshift{1.165173in}{3.075278in}%
\pgfsys@useobject{currentmarker}{}%
\end{pgfscope}%
\begin{pgfscope}%
\pgfsys@transformshift{1.165525in}{3.075278in}%
\pgfsys@useobject{currentmarker}{}%
\end{pgfscope}%
\begin{pgfscope}%
\pgfsys@transformshift{1.165869in}{3.075278in}%
\pgfsys@useobject{currentmarker}{}%
\end{pgfscope}%
\begin{pgfscope}%
\pgfsys@transformshift{1.166375in}{3.075278in}%
\pgfsys@useobject{currentmarker}{}%
\end{pgfscope}%
\begin{pgfscope}%
\pgfsys@transformshift{1.167050in}{3.075278in}%
\pgfsys@useobject{currentmarker}{}%
\end{pgfscope}%
\begin{pgfscope}%
\pgfsys@transformshift{1.167879in}{3.075278in}%
\pgfsys@useobject{currentmarker}{}%
\end{pgfscope}%
\begin{pgfscope}%
\pgfsys@transformshift{1.168216in}{3.075278in}%
\pgfsys@useobject{currentmarker}{}%
\end{pgfscope}%
\begin{pgfscope}%
\pgfsys@transformshift{1.168553in}{3.075278in}%
\pgfsys@useobject{currentmarker}{}%
\end{pgfscope}%
\begin{pgfscope}%
\pgfsys@transformshift{1.168736in}{3.075278in}%
\pgfsys@useobject{currentmarker}{}%
\end{pgfscope}%
\begin{pgfscope}%
\pgfsys@transformshift{1.169073in}{3.075278in}%
\pgfsys@useobject{currentmarker}{}%
\end{pgfscope}%
\begin{pgfscope}%
\pgfsys@transformshift{1.169221in}{3.075278in}%
\pgfsys@useobject{currentmarker}{}%
\end{pgfscope}%
\begin{pgfscope}%
\pgfsys@transformshift{1.169579in}{3.075278in}%
\pgfsys@useobject{currentmarker}{}%
\end{pgfscope}%
\begin{pgfscope}%
\pgfsys@transformshift{1.170078in}{3.075278in}%
\pgfsys@useobject{currentmarker}{}%
\end{pgfscope}%
\begin{pgfscope}%
\pgfsys@transformshift{1.170760in}{3.075278in}%
\pgfsys@useobject{currentmarker}{}%
\end{pgfscope}%
\begin{pgfscope}%
\pgfsys@transformshift{1.171251in}{3.075278in}%
\pgfsys@useobject{currentmarker}{}%
\end{pgfscope}%
\begin{pgfscope}%
\pgfsys@transformshift{1.171757in}{3.075278in}%
\pgfsys@useobject{currentmarker}{}%
\end{pgfscope}%
\begin{pgfscope}%
\pgfsys@transformshift{1.172094in}{3.075278in}%
\pgfsys@useobject{currentmarker}{}%
\end{pgfscope}%
\begin{pgfscope}%
\pgfsys@transformshift{1.172614in}{3.075278in}%
\pgfsys@useobject{currentmarker}{}%
\end{pgfscope}%
\begin{pgfscope}%
\pgfsys@transformshift{1.173289in}{3.075278in}%
\pgfsys@useobject{currentmarker}{}%
\end{pgfscope}%
\begin{pgfscope}%
\pgfsys@transformshift{1.173964in}{3.075278in}%
\pgfsys@useobject{currentmarker}{}%
\end{pgfscope}%
\begin{pgfscope}%
\pgfsys@transformshift{1.174624in}{3.075278in}%
\pgfsys@useobject{currentmarker}{}%
\end{pgfscope}%
\begin{pgfscope}%
\pgfsys@transformshift{1.175481in}{3.075278in}%
\pgfsys@useobject{currentmarker}{}%
\end{pgfscope}%
\begin{pgfscope}%
\pgfsys@transformshift{1.175818in}{3.075278in}%
\pgfsys@useobject{currentmarker}{}%
\end{pgfscope}%
\begin{pgfscope}%
\pgfsys@transformshift{1.176324in}{3.075278in}%
\pgfsys@useobject{currentmarker}{}%
\end{pgfscope}%
\begin{pgfscope}%
\pgfsys@transformshift{1.176830in}{3.075278in}%
\pgfsys@useobject{currentmarker}{}%
\end{pgfscope}%
\begin{pgfscope}%
\pgfsys@transformshift{1.178010in}{3.075278in}%
\pgfsys@useobject{currentmarker}{}%
\end{pgfscope}%
\begin{pgfscope}%
\pgfsys@transformshift{1.178854in}{3.075278in}%
\pgfsys@useobject{currentmarker}{}%
\end{pgfscope}%
\begin{pgfscope}%
\pgfsys@transformshift{1.179866in}{3.075278in}%
\pgfsys@useobject{currentmarker}{}%
\end{pgfscope}%
\begin{pgfscope}%
\pgfsys@transformshift{1.406987in}{3.075278in}%
\pgfsys@useobject{currentmarker}{}%
\end{pgfscope}%
\begin{pgfscope}%
\pgfsys@transformshift{1.423343in}{3.075278in}%
\pgfsys@useobject{currentmarker}{}%
\end{pgfscope}%
\begin{pgfscope}%
\pgfsys@transformshift{1.429077in}{3.075278in}%
\pgfsys@useobject{currentmarker}{}%
\end{pgfscope}%
\begin{pgfscope}%
\pgfsys@transformshift{1.431606in}{3.075278in}%
\pgfsys@useobject{currentmarker}{}%
\end{pgfscope}%
\begin{pgfscope}%
\pgfsys@transformshift{1.434135in}{3.075278in}%
\pgfsys@useobject{currentmarker}{}%
\end{pgfscope}%
\begin{pgfscope}%
\pgfsys@transformshift{1.435990in}{3.075278in}%
\pgfsys@useobject{currentmarker}{}%
\end{pgfscope}%
\begin{pgfscope}%
\pgfsys@transformshift{1.438351in}{3.075278in}%
\pgfsys@useobject{currentmarker}{}%
\end{pgfscope}%
\begin{pgfscope}%
\pgfsys@transformshift{1.440375in}{3.075278in}%
\pgfsys@useobject{currentmarker}{}%
\end{pgfscope}%
\begin{pgfscope}%
\pgfsys@transformshift{1.441386in}{3.075278in}%
\pgfsys@useobject{currentmarker}{}%
\end{pgfscope}%
\begin{pgfscope}%
\pgfsys@transformshift{1.442567in}{3.075278in}%
\pgfsys@useobject{currentmarker}{}%
\end{pgfscope}%
\begin{pgfscope}%
\pgfsys@transformshift{1.444085in}{3.075278in}%
\pgfsys@useobject{currentmarker}{}%
\end{pgfscope}%
\begin{pgfscope}%
\pgfsys@transformshift{1.445771in}{3.075278in}%
\pgfsys@useobject{currentmarker}{}%
\end{pgfscope}%
\begin{pgfscope}%
\pgfsys@transformshift{1.446284in}{3.075278in}%
\pgfsys@useobject{currentmarker}{}%
\end{pgfscope}%
\begin{pgfscope}%
\pgfsys@transformshift{1.447633in}{3.075278in}%
\pgfsys@useobject{currentmarker}{}%
\end{pgfscope}%
\begin{pgfscope}%
\pgfsys@transformshift{1.448307in}{3.075278in}%
\pgfsys@useobject{currentmarker}{}%
\end{pgfscope}%
\begin{pgfscope}%
\pgfsys@transformshift{1.449144in}{3.075278in}%
\pgfsys@useobject{currentmarker}{}%
\end{pgfscope}%
\begin{pgfscope}%
\pgfsys@transformshift{1.449818in}{3.075278in}%
\pgfsys@useobject{currentmarker}{}%
\end{pgfscope}%
\begin{pgfscope}%
\pgfsys@transformshift{1.450998in}{3.075278in}%
\pgfsys@useobject{currentmarker}{}%
\end{pgfscope}%
\begin{pgfscope}%
\pgfsys@transformshift{1.451841in}{3.075278in}%
\pgfsys@useobject{currentmarker}{}%
\end{pgfscope}%
\begin{pgfscope}%
\pgfsys@transformshift{1.453528in}{3.075278in}%
\pgfsys@useobject{currentmarker}{}%
\end{pgfscope}%
\begin{pgfscope}%
\pgfsys@transformshift{1.454203in}{3.075278in}%
\pgfsys@useobject{currentmarker}{}%
\end{pgfscope}%
\begin{pgfscope}%
\pgfsys@transformshift{1.455221in}{3.075278in}%
\pgfsys@useobject{currentmarker}{}%
\end{pgfscope}%
\begin{pgfscope}%
\pgfsys@transformshift{1.456570in}{3.075278in}%
\pgfsys@useobject{currentmarker}{}%
\end{pgfscope}%
\begin{pgfscope}%
\pgfsys@transformshift{1.457575in}{3.075278in}%
\pgfsys@useobject{currentmarker}{}%
\end{pgfscope}%
\begin{pgfscope}%
\pgfsys@transformshift{1.458587in}{3.075278in}%
\pgfsys@useobject{currentmarker}{}%
\end{pgfscope}%
\begin{pgfscope}%
\pgfsys@transformshift{1.459613in}{3.075278in}%
\pgfsys@useobject{currentmarker}{}%
\end{pgfscope}%
\begin{pgfscope}%
\pgfsys@transformshift{1.460779in}{3.075278in}%
\pgfsys@useobject{currentmarker}{}%
\end{pgfscope}%
\begin{pgfscope}%
\pgfsys@transformshift{1.461622in}{3.075278in}%
\pgfsys@useobject{currentmarker}{}%
\end{pgfscope}%
\begin{pgfscope}%
\pgfsys@transformshift{1.462634in}{3.075278in}%
\pgfsys@useobject{currentmarker}{}%
\end{pgfscope}%
\begin{pgfscope}%
\pgfsys@transformshift{1.463484in}{3.075278in}%
\pgfsys@useobject{currentmarker}{}%
\end{pgfscope}%
\begin{pgfscope}%
\pgfsys@transformshift{1.464840in}{3.075278in}%
\pgfsys@useobject{currentmarker}{}%
\end{pgfscope}%
\begin{pgfscope}%
\pgfsys@transformshift{1.466021in}{3.075278in}%
\pgfsys@useobject{currentmarker}{}%
\end{pgfscope}%
\begin{pgfscope}%
\pgfsys@transformshift{1.467018in}{3.075278in}%
\pgfsys@useobject{currentmarker}{}%
\end{pgfscope}%
\begin{pgfscope}%
\pgfsys@transformshift{1.468199in}{3.075278in}%
\pgfsys@useobject{currentmarker}{}%
\end{pgfscope}%
\begin{pgfscope}%
\pgfsys@transformshift{1.469386in}{3.075278in}%
\pgfsys@useobject{currentmarker}{}%
\end{pgfscope}%
\begin{pgfscope}%
\pgfsys@transformshift{1.470910in}{3.075278in}%
\pgfsys@useobject{currentmarker}{}%
\end{pgfscope}%
\begin{pgfscope}%
\pgfsys@transformshift{1.471739in}{3.075278in}%
\pgfsys@useobject{currentmarker}{}%
\end{pgfscope}%
\begin{pgfscope}%
\pgfsys@transformshift{1.472752in}{3.075278in}%
\pgfsys@useobject{currentmarker}{}%
\end{pgfscope}%
\begin{pgfscope}%
\pgfsys@transformshift{1.473609in}{3.075278in}%
\pgfsys@useobject{currentmarker}{}%
\end{pgfscope}%
\begin{pgfscope}%
\pgfsys@transformshift{1.474789in}{3.075278in}%
\pgfsys@useobject{currentmarker}{}%
\end{pgfscope}%
\begin{pgfscope}%
\pgfsys@transformshift{1.475963in}{3.075278in}%
\pgfsys@useobject{currentmarker}{}%
\end{pgfscope}%
\begin{pgfscope}%
\pgfsys@transformshift{1.476798in}{3.075278in}%
\pgfsys@useobject{currentmarker}{}%
\end{pgfscope}%
\begin{pgfscope}%
\pgfsys@transformshift{1.478316in}{3.075278in}%
\pgfsys@useobject{currentmarker}{}%
\end{pgfscope}%
\begin{pgfscope}%
\pgfsys@transformshift{1.479665in}{3.075278in}%
\pgfsys@useobject{currentmarker}{}%
\end{pgfscope}%
\begin{pgfscope}%
\pgfsys@transformshift{1.481535in}{3.075278in}%
\pgfsys@useobject{currentmarker}{}%
\end{pgfscope}%
\begin{pgfscope}%
\pgfsys@transformshift{1.482722in}{3.075278in}%
\pgfsys@useobject{currentmarker}{}%
\end{pgfscope}%
\begin{pgfscope}%
\pgfsys@transformshift{1.720801in}{3.075278in}%
\pgfsys@useobject{currentmarker}{}%
\end{pgfscope}%
\begin{pgfscope}%
\pgfsys@transformshift{1.735809in}{3.075278in}%
\pgfsys@useobject{currentmarker}{}%
\end{pgfscope}%
\begin{pgfscope}%
\pgfsys@transformshift{1.741037in}{3.075278in}%
\pgfsys@useobject{currentmarker}{}%
\end{pgfscope}%
\begin{pgfscope}%
\pgfsys@transformshift{1.744240in}{3.075278in}%
\pgfsys@useobject{currentmarker}{}%
\end{pgfscope}%
\begin{pgfscope}%
\pgfsys@transformshift{1.745931in}{3.075278in}%
\pgfsys@useobject{currentmarker}{}%
\end{pgfscope}%
\begin{pgfscope}%
\pgfsys@transformshift{1.747108in}{3.075278in}%
\pgfsys@useobject{currentmarker}{}%
\end{pgfscope}%
\begin{pgfscope}%
\pgfsys@transformshift{1.748635in}{3.075278in}%
\pgfsys@useobject{currentmarker}{}%
\end{pgfscope}%
\begin{pgfscope}%
\pgfsys@transformshift{1.750653in}{3.075278in}%
\pgfsys@useobject{currentmarker}{}%
\end{pgfscope}%
\begin{pgfscope}%
\pgfsys@transformshift{1.752165in}{3.075278in}%
\pgfsys@useobject{currentmarker}{}%
\end{pgfscope}%
\begin{pgfscope}%
\pgfsys@transformshift{1.753521in}{3.075278in}%
\pgfsys@useobject{currentmarker}{}%
\end{pgfscope}%
\begin{pgfscope}%
\pgfsys@transformshift{1.754362in}{3.075278in}%
\pgfsys@useobject{currentmarker}{}%
\end{pgfscope}%
\begin{pgfscope}%
\pgfsys@transformshift{1.755711in}{3.075278in}%
\pgfsys@useobject{currentmarker}{}%
\end{pgfscope}%
\begin{pgfscope}%
\pgfsys@transformshift{1.757230in}{3.075278in}%
\pgfsys@useobject{currentmarker}{}%
\end{pgfscope}%
\begin{pgfscope}%
\pgfsys@transformshift{1.757960in}{3.075278in}%
\pgfsys@useobject{currentmarker}{}%
\end{pgfscope}%
\begin{pgfscope}%
\pgfsys@transformshift{1.758921in}{3.075278in}%
\pgfsys@useobject{currentmarker}{}%
\end{pgfscope}%
\begin{pgfscope}%
\pgfsys@transformshift{1.759934in}{3.075278in}%
\pgfsys@useobject{currentmarker}{}%
\end{pgfscope}%
\begin{pgfscope}%
\pgfsys@transformshift{1.760939in}{3.075278in}%
\pgfsys@useobject{currentmarker}{}%
\end{pgfscope}%
\begin{pgfscope}%
\pgfsys@transformshift{1.761952in}{3.075278in}%
\pgfsys@useobject{currentmarker}{}%
\end{pgfscope}%
\begin{pgfscope}%
\pgfsys@transformshift{1.762794in}{3.075278in}%
\pgfsys@useobject{currentmarker}{}%
\end{pgfscope}%
\begin{pgfscope}%
\pgfsys@transformshift{1.763807in}{3.075278in}%
\pgfsys@useobject{currentmarker}{}%
\end{pgfscope}%
\begin{pgfscope}%
\pgfsys@transformshift{1.764991in}{3.075278in}%
\pgfsys@useobject{currentmarker}{}%
\end{pgfscope}%
\begin{pgfscope}%
\pgfsys@transformshift{1.766168in}{3.075278in}%
\pgfsys@useobject{currentmarker}{}%
\end{pgfscope}%
\begin{pgfscope}%
\pgfsys@transformshift{1.766838in}{3.075278in}%
\pgfsys@useobject{currentmarker}{}%
\end{pgfscope}%
\begin{pgfscope}%
\pgfsys@transformshift{1.768358in}{3.075278in}%
\pgfsys@useobject{currentmarker}{}%
\end{pgfscope}%
\begin{pgfscope}%
\pgfsys@transformshift{1.769371in}{3.075278in}%
\pgfsys@useobject{currentmarker}{}%
\end{pgfscope}%
\begin{pgfscope}%
\pgfsys@transformshift{1.770726in}{3.075278in}%
\pgfsys@useobject{currentmarker}{}%
\end{pgfscope}%
\begin{pgfscope}%
\pgfsys@transformshift{1.772238in}{3.075278in}%
\pgfsys@useobject{currentmarker}{}%
\end{pgfscope}%
\begin{pgfscope}%
\pgfsys@transformshift{1.773080in}{3.075278in}%
\pgfsys@useobject{currentmarker}{}%
\end{pgfscope}%
\begin{pgfscope}%
\pgfsys@transformshift{1.774763in}{3.075278in}%
\pgfsys@useobject{currentmarker}{}%
\end{pgfscope}%
\begin{pgfscope}%
\pgfsys@transformshift{1.775955in}{3.075278in}%
\pgfsys@useobject{currentmarker}{}%
\end{pgfscope}%
\begin{pgfscope}%
\pgfsys@transformshift{1.776789in}{3.075278in}%
\pgfsys@useobject{currentmarker}{}%
\end{pgfscope}%
\begin{pgfscope}%
\pgfsys@transformshift{1.777809in}{3.075278in}%
\pgfsys@useobject{currentmarker}{}%
\end{pgfscope}%
\begin{pgfscope}%
\pgfsys@transformshift{1.779508in}{3.075278in}%
\pgfsys@useobject{currentmarker}{}%
\end{pgfscope}%
\begin{pgfscope}%
\pgfsys@transformshift{1.781184in}{3.075278in}%
\pgfsys@useobject{currentmarker}{}%
\end{pgfscope}%
\begin{pgfscope}%
\pgfsys@transformshift{1.782696in}{3.075278in}%
\pgfsys@useobject{currentmarker}{}%
\end{pgfscope}%
\begin{pgfscope}%
\pgfsys@transformshift{1.784721in}{3.075278in}%
\pgfsys@useobject{currentmarker}{}%
\end{pgfscope}%
\begin{pgfscope}%
\pgfsys@transformshift{1.786911in}{3.075278in}%
\pgfsys@useobject{currentmarker}{}%
\end{pgfscope}%
\begin{pgfscope}%
\pgfsys@transformshift{1.788438in}{3.075278in}%
\pgfsys@useobject{currentmarker}{}%
\end{pgfscope}%
\begin{pgfscope}%
\pgfsys@transformshift{1.790285in}{3.075278in}%
\pgfsys@useobject{currentmarker}{}%
\end{pgfscope}%
\begin{pgfscope}%
\pgfsys@transformshift{1.792982in}{3.075278in}%
\pgfsys@useobject{currentmarker}{}%
\end{pgfscope}%
\begin{pgfscope}%
\pgfsys@transformshift{1.794330in}{3.075278in}%
\pgfsys@useobject{currentmarker}{}%
\end{pgfscope}%
\begin{pgfscope}%
\pgfsys@transformshift{1.797197in}{3.075278in}%
\pgfsys@useobject{currentmarker}{}%
\end{pgfscope}%
\begin{pgfscope}%
\pgfsys@transformshift{1.799573in}{3.075278in}%
\pgfsys@useobject{currentmarker}{}%
\end{pgfscope}%
\begin{pgfscope}%
\pgfsys@transformshift{1.802605in}{3.075278in}%
\pgfsys@useobject{currentmarker}{}%
\end{pgfscope}%
\begin{pgfscope}%
\pgfsys@transformshift{2.075933in}{3.075278in}%
\pgfsys@useobject{currentmarker}{}%
\end{pgfscope}%
\begin{pgfscope}%
\pgfsys@transformshift{2.086897in}{3.075278in}%
\pgfsys@useobject{currentmarker}{}%
\end{pgfscope}%
\begin{pgfscope}%
\pgfsys@transformshift{2.090271in}{3.075278in}%
\pgfsys@useobject{currentmarker}{}%
\end{pgfscope}%
\begin{pgfscope}%
\pgfsys@transformshift{2.093138in}{3.075278in}%
\pgfsys@useobject{currentmarker}{}%
\end{pgfscope}%
\begin{pgfscope}%
\pgfsys@transformshift{2.095328in}{3.075278in}%
\pgfsys@useobject{currentmarker}{}%
\end{pgfscope}%
\begin{pgfscope}%
\pgfsys@transformshift{2.095999in}{3.075278in}%
\pgfsys@useobject{currentmarker}{}%
\end{pgfscope}%
\begin{pgfscope}%
\pgfsys@transformshift{2.097689in}{3.075278in}%
\pgfsys@useobject{currentmarker}{}%
\end{pgfscope}%
\begin{pgfscope}%
\pgfsys@transformshift{2.099037in}{3.075278in}%
\pgfsys@useobject{currentmarker}{}%
\end{pgfscope}%
\begin{pgfscope}%
\pgfsys@transformshift{2.099879in}{3.075278in}%
\pgfsys@useobject{currentmarker}{}%
\end{pgfscope}%
\begin{pgfscope}%
\pgfsys@transformshift{2.101227in}{3.075278in}%
\pgfsys@useobject{currentmarker}{}%
\end{pgfscope}%
\begin{pgfscope}%
\pgfsys@transformshift{2.101905in}{3.075278in}%
\pgfsys@useobject{currentmarker}{}%
\end{pgfscope}%
\begin{pgfscope}%
\pgfsys@transformshift{2.102918in}{3.075278in}%
\pgfsys@useobject{currentmarker}{}%
\end{pgfscope}%
\begin{pgfscope}%
\pgfsys@transformshift{2.103923in}{3.075278in}%
\pgfsys@useobject{currentmarker}{}%
\end{pgfscope}%
\begin{pgfscope}%
\pgfsys@transformshift{2.105108in}{3.075278in}%
\pgfsys@useobject{currentmarker}{}%
\end{pgfscope}%
\begin{pgfscope}%
\pgfsys@transformshift{2.106121in}{3.075278in}%
\pgfsys@useobject{currentmarker}{}%
\end{pgfscope}%
\begin{pgfscope}%
\pgfsys@transformshift{2.107134in}{3.075278in}%
\pgfsys@useobject{currentmarker}{}%
\end{pgfscope}%
\begin{pgfscope}%
\pgfsys@transformshift{2.107804in}{3.075278in}%
\pgfsys@useobject{currentmarker}{}%
\end{pgfscope}%
\begin{pgfscope}%
\pgfsys@transformshift{2.108817in}{3.075278in}%
\pgfsys@useobject{currentmarker}{}%
\end{pgfscope}%
\begin{pgfscope}%
\pgfsys@transformshift{2.109659in}{3.075278in}%
\pgfsys@useobject{currentmarker}{}%
\end{pgfscope}%
\begin{pgfscope}%
\pgfsys@transformshift{2.111007in}{3.075278in}%
\pgfsys@useobject{currentmarker}{}%
\end{pgfscope}%
\begin{pgfscope}%
\pgfsys@transformshift{2.111856in}{3.075278in}%
\pgfsys@useobject{currentmarker}{}%
\end{pgfscope}%
\begin{pgfscope}%
\pgfsys@transformshift{2.112861in}{3.075278in}%
\pgfsys@useobject{currentmarker}{}%
\end{pgfscope}%
\begin{pgfscope}%
\pgfsys@transformshift{2.114053in}{3.075278in}%
\pgfsys@useobject{currentmarker}{}%
\end{pgfscope}%
\begin{pgfscope}%
\pgfsys@transformshift{2.114716in}{3.075278in}%
\pgfsys@useobject{currentmarker}{}%
\end{pgfscope}%
\begin{pgfscope}%
\pgfsys@transformshift{2.115565in}{3.075278in}%
\pgfsys@useobject{currentmarker}{}%
\end{pgfscope}%
\begin{pgfscope}%
\pgfsys@transformshift{2.116414in}{3.075278in}%
\pgfsys@useobject{currentmarker}{}%
\end{pgfscope}%
\begin{pgfscope}%
\pgfsys@transformshift{2.117427in}{3.075278in}%
\pgfsys@useobject{currentmarker}{}%
\end{pgfscope}%
\begin{pgfscope}%
\pgfsys@transformshift{2.118090in}{3.075278in}%
\pgfsys@useobject{currentmarker}{}%
\end{pgfscope}%
\begin{pgfscope}%
\pgfsys@transformshift{2.118775in}{3.075278in}%
\pgfsys@useobject{currentmarker}{}%
\end{pgfscope}%
\begin{pgfscope}%
\pgfsys@transformshift{2.119952in}{3.075278in}%
\pgfsys@useobject{currentmarker}{}%
\end{pgfscope}%
\begin{pgfscope}%
\pgfsys@transformshift{2.120630in}{3.075278in}%
\pgfsys@useobject{currentmarker}{}%
\end{pgfscope}%
\begin{pgfscope}%
\pgfsys@transformshift{2.121814in}{3.075278in}%
\pgfsys@useobject{currentmarker}{}%
\end{pgfscope}%
\begin{pgfscope}%
\pgfsys@transformshift{2.122827in}{3.075278in}%
\pgfsys@useobject{currentmarker}{}%
\end{pgfscope}%
\begin{pgfscope}%
\pgfsys@transformshift{2.123840in}{3.075278in}%
\pgfsys@useobject{currentmarker}{}%
\end{pgfscope}%
\begin{pgfscope}%
\pgfsys@transformshift{2.125345in}{3.075278in}%
\pgfsys@useobject{currentmarker}{}%
\end{pgfscope}%
\begin{pgfscope}%
\pgfsys@transformshift{2.126864in}{3.075278in}%
\pgfsys@useobject{currentmarker}{}%
\end{pgfscope}%
\begin{pgfscope}%
\pgfsys@transformshift{2.127713in}{3.075278in}%
\pgfsys@useobject{currentmarker}{}%
\end{pgfscope}%
\begin{pgfscope}%
\pgfsys@transformshift{2.129061in}{3.075278in}%
\pgfsys@useobject{currentmarker}{}%
\end{pgfscope}%
\begin{pgfscope}%
\pgfsys@transformshift{2.130409in}{3.075278in}%
\pgfsys@useobject{currentmarker}{}%
\end{pgfscope}%
\begin{pgfscope}%
\pgfsys@transformshift{2.131251in}{3.075278in}%
\pgfsys@useobject{currentmarker}{}%
\end{pgfscope}%
\begin{pgfscope}%
\pgfsys@transformshift{2.132264in}{3.075278in}%
\pgfsys@useobject{currentmarker}{}%
\end{pgfscope}%
\begin{pgfscope}%
\pgfsys@transformshift{2.133158in}{3.075278in}%
\pgfsys@useobject{currentmarker}{}%
\end{pgfscope}%
\begin{pgfscope}%
\pgfsys@transformshift{2.135310in}{3.075278in}%
\pgfsys@useobject{currentmarker}{}%
\end{pgfscope}%
\begin{pgfscope}%
\pgfsys@transformshift{2.136487in}{3.075278in}%
\pgfsys@useobject{currentmarker}{}%
\end{pgfscope}%
\begin{pgfscope}%
\pgfsys@transformshift{2.137657in}{3.075278in}%
\pgfsys@useobject{currentmarker}{}%
\end{pgfscope}%
\begin{pgfscope}%
\pgfsys@transformshift{2.139340in}{3.075278in}%
\pgfsys@useobject{currentmarker}{}%
\end{pgfscope}%
\begin{pgfscope}%
\pgfsys@transformshift{2.141217in}{3.075278in}%
\pgfsys@useobject{currentmarker}{}%
\end{pgfscope}%
\begin{pgfscope}%
\pgfsys@transformshift{2.143064in}{3.075278in}%
\pgfsys@useobject{currentmarker}{}%
\end{pgfscope}%
\begin{pgfscope}%
\pgfsys@transformshift{2.144740in}{3.075278in}%
\pgfsys@useobject{currentmarker}{}%
\end{pgfscope}%
\begin{pgfscope}%
\pgfsys@transformshift{2.423476in}{3.075278in}%
\pgfsys@useobject{currentmarker}{}%
\end{pgfscope}%
\begin{pgfscope}%
\pgfsys@transformshift{2.437136in}{3.075278in}%
\pgfsys@useobject{currentmarker}{}%
\end{pgfscope}%
\begin{pgfscope}%
\pgfsys@transformshift{2.439832in}{3.075278in}%
\pgfsys@useobject{currentmarker}{}%
\end{pgfscope}%
\begin{pgfscope}%
\pgfsys@transformshift{2.441858in}{3.075278in}%
\pgfsys@useobject{currentmarker}{}%
\end{pgfscope}%
\begin{pgfscope}%
\pgfsys@transformshift{2.444725in}{3.075278in}%
\pgfsys@useobject{currentmarker}{}%
\end{pgfscope}%
\begin{pgfscope}%
\pgfsys@transformshift{2.446580in}{3.075278in}%
\pgfsys@useobject{currentmarker}{}%
\end{pgfscope}%
\begin{pgfscope}%
\pgfsys@transformshift{2.447928in}{3.075278in}%
\pgfsys@useobject{currentmarker}{}%
\end{pgfscope}%
\begin{pgfscope}%
\pgfsys@transformshift{2.449611in}{3.075278in}%
\pgfsys@useobject{currentmarker}{}%
\end{pgfscope}%
\begin{pgfscope}%
\pgfsys@transformshift{2.451131in}{3.075278in}%
\pgfsys@useobject{currentmarker}{}%
\end{pgfscope}%
\begin{pgfscope}%
\pgfsys@transformshift{2.452144in}{3.075278in}%
\pgfsys@useobject{currentmarker}{}%
\end{pgfscope}%
\begin{pgfscope}%
\pgfsys@transformshift{2.453157in}{3.075278in}%
\pgfsys@useobject{currentmarker}{}%
\end{pgfscope}%
\begin{pgfscope}%
\pgfsys@transformshift{2.454334in}{3.075278in}%
\pgfsys@useobject{currentmarker}{}%
\end{pgfscope}%
\begin{pgfscope}%
\pgfsys@transformshift{2.455183in}{3.075278in}%
\pgfsys@useobject{currentmarker}{}%
\end{pgfscope}%
\begin{pgfscope}%
\pgfsys@transformshift{2.456360in}{3.075278in}%
\pgfsys@useobject{currentmarker}{}%
\end{pgfscope}%
\begin{pgfscope}%
\pgfsys@transformshift{2.457879in}{3.075278in}%
\pgfsys@useobject{currentmarker}{}%
\end{pgfscope}%
\begin{pgfscope}%
\pgfsys@transformshift{2.458549in}{3.075278in}%
\pgfsys@useobject{currentmarker}{}%
\end{pgfscope}%
\begin{pgfscope}%
\pgfsys@transformshift{2.459227in}{3.075278in}%
\pgfsys@useobject{currentmarker}{}%
\end{pgfscope}%
\begin{pgfscope}%
\pgfsys@transformshift{2.460069in}{3.075278in}%
\pgfsys@useobject{currentmarker}{}%
\end{pgfscope}%
\begin{pgfscope}%
\pgfsys@transformshift{2.461253in}{3.075278in}%
\pgfsys@useobject{currentmarker}{}%
\end{pgfscope}%
\begin{pgfscope}%
\pgfsys@transformshift{2.462430in}{3.075278in}%
\pgfsys@useobject{currentmarker}{}%
\end{pgfscope}%
\begin{pgfscope}%
\pgfsys@transformshift{2.463614in}{3.075278in}%
\pgfsys@useobject{currentmarker}{}%
\end{pgfscope}%
\begin{pgfscope}%
\pgfsys@transformshift{2.464299in}{3.075278in}%
\pgfsys@useobject{currentmarker}{}%
\end{pgfscope}%
\begin{pgfscope}%
\pgfsys@transformshift{2.465633in}{3.075278in}%
\pgfsys@useobject{currentmarker}{}%
\end{pgfscope}%
\begin{pgfscope}%
\pgfsys@transformshift{2.467003in}{3.075278in}%
\pgfsys@useobject{currentmarker}{}%
\end{pgfscope}%
\begin{pgfscope}%
\pgfsys@transformshift{2.468001in}{3.075278in}%
\pgfsys@useobject{currentmarker}{}%
\end{pgfscope}%
\begin{pgfscope}%
\pgfsys@transformshift{2.468679in}{3.075278in}%
\pgfsys@useobject{currentmarker}{}%
\end{pgfscope}%
\begin{pgfscope}%
\pgfsys@transformshift{2.470355in}{3.075278in}%
\pgfsys@useobject{currentmarker}{}%
\end{pgfscope}%
\begin{pgfscope}%
\pgfsys@transformshift{2.471539in}{3.075278in}%
\pgfsys@useobject{currentmarker}{}%
\end{pgfscope}%
\begin{pgfscope}%
\pgfsys@transformshift{2.472552in}{3.075278in}%
\pgfsys@useobject{currentmarker}{}%
\end{pgfscope}%
\begin{pgfscope}%
\pgfsys@transformshift{2.473908in}{3.075278in}%
\pgfsys@useobject{currentmarker}{}%
\end{pgfscope}%
\begin{pgfscope}%
\pgfsys@transformshift{2.475248in}{3.075278in}%
\pgfsys@useobject{currentmarker}{}%
\end{pgfscope}%
\begin{pgfscope}%
\pgfsys@transformshift{2.476433in}{3.075278in}%
\pgfsys@useobject{currentmarker}{}%
\end{pgfscope}%
\begin{pgfscope}%
\pgfsys@transformshift{2.477781in}{3.075278in}%
\pgfsys@useobject{currentmarker}{}%
\end{pgfscope}%
\begin{pgfscope}%
\pgfsys@transformshift{2.479129in}{3.075278in}%
\pgfsys@useobject{currentmarker}{}%
\end{pgfscope}%
\begin{pgfscope}%
\pgfsys@transformshift{2.480142in}{3.075278in}%
\pgfsys@useobject{currentmarker}{}%
\end{pgfscope}%
\begin{pgfscope}%
\pgfsys@transformshift{2.480983in}{3.075278in}%
\pgfsys@useobject{currentmarker}{}%
\end{pgfscope}%
\begin{pgfscope}%
\pgfsys@transformshift{2.482495in}{3.075278in}%
\pgfsys@useobject{currentmarker}{}%
\end{pgfscope}%
\begin{pgfscope}%
\pgfsys@transformshift{2.483680in}{3.075278in}%
\pgfsys@useobject{currentmarker}{}%
\end{pgfscope}%
\begin{pgfscope}%
\pgfsys@transformshift{2.484521in}{3.075278in}%
\pgfsys@useobject{currentmarker}{}%
\end{pgfscope}%
\begin{pgfscope}%
\pgfsys@transformshift{2.485534in}{3.075278in}%
\pgfsys@useobject{currentmarker}{}%
\end{pgfscope}%
\begin{pgfscope}%
\pgfsys@transformshift{2.486398in}{3.075278in}%
\pgfsys@useobject{currentmarker}{}%
\end{pgfscope}%
\begin{pgfscope}%
\pgfsys@transformshift{2.486890in}{3.075278in}%
\pgfsys@useobject{currentmarker}{}%
\end{pgfscope}%
\begin{pgfscope}%
\pgfsys@transformshift{2.488074in}{3.075278in}%
\pgfsys@useobject{currentmarker}{}%
\end{pgfscope}%
\begin{pgfscope}%
\pgfsys@transformshift{2.489266in}{3.075278in}%
\pgfsys@useobject{currentmarker}{}%
\end{pgfscope}%
\begin{pgfscope}%
\pgfsys@transformshift{2.490279in}{3.075278in}%
\pgfsys@useobject{currentmarker}{}%
\end{pgfscope}%
\begin{pgfscope}%
\pgfsys@transformshift{2.491277in}{3.075278in}%
\pgfsys@useobject{currentmarker}{}%
\end{pgfscope}%
\begin{pgfscope}%
\pgfsys@transformshift{2.492119in}{3.075278in}%
\pgfsys@useobject{currentmarker}{}%
\end{pgfscope}%
\begin{pgfscope}%
\pgfsys@transformshift{2.493482in}{3.075278in}%
\pgfsys@useobject{currentmarker}{}%
\end{pgfscope}%
\begin{pgfscope}%
\pgfsys@transformshift{2.494495in}{3.075278in}%
\pgfsys@useobject{currentmarker}{}%
\end{pgfscope}%
\begin{pgfscope}%
\pgfsys@transformshift{2.496178in}{3.075278in}%
\pgfsys@useobject{currentmarker}{}%
\end{pgfscope}%
\begin{pgfscope}%
\pgfsys@transformshift{2.789572in}{3.075278in}%
\pgfsys@useobject{currentmarker}{}%
\end{pgfscope}%
\begin{pgfscope}%
\pgfsys@transformshift{2.804580in}{3.075278in}%
\pgfsys@useobject{currentmarker}{}%
\end{pgfscope}%
\begin{pgfscope}%
\pgfsys@transformshift{2.808453in}{3.075278in}%
\pgfsys@useobject{currentmarker}{}%
\end{pgfscope}%
\begin{pgfscope}%
\pgfsys@transformshift{2.810479in}{3.075278in}%
\pgfsys@useobject{currentmarker}{}%
\end{pgfscope}%
\begin{pgfscope}%
\pgfsys@transformshift{2.812840in}{3.075278in}%
\pgfsys@useobject{currentmarker}{}%
\end{pgfscope}%
\begin{pgfscope}%
\pgfsys@transformshift{2.814188in}{3.075278in}%
\pgfsys@useobject{currentmarker}{}%
\end{pgfscope}%
\begin{pgfscope}%
\pgfsys@transformshift{2.815536in}{3.075278in}%
\pgfsys@useobject{currentmarker}{}%
\end{pgfscope}%
\begin{pgfscope}%
\pgfsys@transformshift{2.817227in}{3.075278in}%
\pgfsys@useobject{currentmarker}{}%
\end{pgfscope}%
\begin{pgfscope}%
\pgfsys@transformshift{2.818240in}{3.075278in}%
\pgfsys@useobject{currentmarker}{}%
\end{pgfscope}%
\begin{pgfscope}%
\pgfsys@transformshift{2.819417in}{3.075278in}%
\pgfsys@useobject{currentmarker}{}%
\end{pgfscope}%
\begin{pgfscope}%
\pgfsys@transformshift{2.820430in}{3.075278in}%
\pgfsys@useobject{currentmarker}{}%
\end{pgfscope}%
\begin{pgfscope}%
\pgfsys@transformshift{2.821778in}{3.075278in}%
\pgfsys@useobject{currentmarker}{}%
\end{pgfscope}%
\begin{pgfscope}%
\pgfsys@transformshift{2.822791in}{3.075278in}%
\pgfsys@useobject{currentmarker}{}%
\end{pgfscope}%
\begin{pgfscope}%
\pgfsys@transformshift{2.823804in}{3.075278in}%
\pgfsys@useobject{currentmarker}{}%
\end{pgfscope}%
\begin{pgfscope}%
\pgfsys@transformshift{2.824809in}{3.075278in}%
\pgfsys@useobject{currentmarker}{}%
\end{pgfscope}%
\begin{pgfscope}%
\pgfsys@transformshift{2.825316in}{3.075278in}%
\pgfsys@useobject{currentmarker}{}%
\end{pgfscope}%
\begin{pgfscope}%
\pgfsys@transformshift{2.826835in}{3.075278in}%
\pgfsys@useobject{currentmarker}{}%
\end{pgfscope}%
\begin{pgfscope}%
\pgfsys@transformshift{2.827677in}{3.075278in}%
\pgfsys@useobject{currentmarker}{}%
\end{pgfscope}%
\begin{pgfscope}%
\pgfsys@transformshift{2.828690in}{3.075278in}%
\pgfsys@useobject{currentmarker}{}%
\end{pgfscope}%
\begin{pgfscope}%
\pgfsys@transformshift{2.830038in}{3.075278in}%
\pgfsys@useobject{currentmarker}{}%
\end{pgfscope}%
\begin{pgfscope}%
\pgfsys@transformshift{2.830887in}{3.075278in}%
\pgfsys@useobject{currentmarker}{}%
\end{pgfscope}%
\begin{pgfscope}%
\pgfsys@transformshift{2.831558in}{3.075278in}%
\pgfsys@useobject{currentmarker}{}%
\end{pgfscope}%
\begin{pgfscope}%
\pgfsys@transformshift{2.831900in}{3.075278in}%
\pgfsys@useobject{currentmarker}{}%
\end{pgfscope}%
\begin{pgfscope}%
\pgfsys@transformshift{2.833241in}{3.075278in}%
\pgfsys@useobject{currentmarker}{}%
\end{pgfscope}%
\begin{pgfscope}%
\pgfsys@transformshift{2.833755in}{3.075278in}%
\pgfsys@useobject{currentmarker}{}%
\end{pgfscope}%
\begin{pgfscope}%
\pgfsys@transformshift{2.834596in}{3.075278in}%
\pgfsys@useobject{currentmarker}{}%
\end{pgfscope}%
\begin{pgfscope}%
\pgfsys@transformshift{2.835446in}{3.075278in}%
\pgfsys@useobject{currentmarker}{}%
\end{pgfscope}%
\begin{pgfscope}%
\pgfsys@transformshift{2.836615in}{3.075278in}%
\pgfsys@useobject{currentmarker}{}%
\end{pgfscope}%
\begin{pgfscope}%
\pgfsys@transformshift{2.837457in}{3.075278in}%
\pgfsys@useobject{currentmarker}{}%
\end{pgfscope}%
\begin{pgfscope}%
\pgfsys@transformshift{2.838306in}{3.075278in}%
\pgfsys@useobject{currentmarker}{}%
\end{pgfscope}%
\begin{pgfscope}%
\pgfsys@transformshift{2.839147in}{3.075278in}%
\pgfsys@useobject{currentmarker}{}%
\end{pgfscope}%
\begin{pgfscope}%
\pgfsys@transformshift{2.840495in}{3.075278in}%
\pgfsys@useobject{currentmarker}{}%
\end{pgfscope}%
\begin{pgfscope}%
\pgfsys@transformshift{2.841508in}{3.075278in}%
\pgfsys@useobject{currentmarker}{}%
\end{pgfscope}%
\begin{pgfscope}%
\pgfsys@transformshift{2.842529in}{3.075278in}%
\pgfsys@useobject{currentmarker}{}%
\end{pgfscope}%
\begin{pgfscope}%
\pgfsys@transformshift{2.843534in}{3.075278in}%
\pgfsys@useobject{currentmarker}{}%
\end{pgfscope}%
\begin{pgfscope}%
\pgfsys@transformshift{2.844212in}{3.075278in}%
\pgfsys@useobject{currentmarker}{}%
\end{pgfscope}%
\begin{pgfscope}%
\pgfsys@transformshift{2.845560in}{3.075278in}%
\pgfsys@useobject{currentmarker}{}%
\end{pgfscope}%
\begin{pgfscope}%
\pgfsys@transformshift{2.846908in}{3.075278in}%
\pgfsys@useobject{currentmarker}{}%
\end{pgfscope}%
\begin{pgfscope}%
\pgfsys@transformshift{2.848100in}{3.075278in}%
\pgfsys@useobject{currentmarker}{}%
\end{pgfscope}%
\begin{pgfscope}%
\pgfsys@transformshift{2.849612in}{3.075278in}%
\pgfsys@useobject{currentmarker}{}%
\end{pgfscope}%
\begin{pgfscope}%
\pgfsys@transformshift{2.850796in}{3.075278in}%
\pgfsys@useobject{currentmarker}{}%
\end{pgfscope}%
\begin{pgfscope}%
\pgfsys@transformshift{2.852822in}{3.075278in}%
\pgfsys@useobject{currentmarker}{}%
\end{pgfscope}%
\begin{pgfscope}%
\pgfsys@transformshift{2.854163in}{3.075278in}%
\pgfsys@useobject{currentmarker}{}%
\end{pgfscope}%
\begin{pgfscope}%
\pgfsys@transformshift{2.857045in}{3.075278in}%
\pgfsys@useobject{currentmarker}{}%
\end{pgfscope}%
\begin{pgfscope}%
\pgfsys@transformshift{2.858900in}{3.075278in}%
\pgfsys@useobject{currentmarker}{}%
\end{pgfscope}%
\begin{pgfscope}%
\pgfsys@transformshift{2.861246in}{3.075278in}%
\pgfsys@useobject{currentmarker}{}%
\end{pgfscope}%
\begin{pgfscope}%
\pgfsys@transformshift{3.130031in}{3.075278in}%
\pgfsys@useobject{currentmarker}{}%
\end{pgfscope}%
\begin{pgfscope}%
\pgfsys@transformshift{3.141837in}{3.075278in}%
\pgfsys@useobject{currentmarker}{}%
\end{pgfscope}%
\begin{pgfscope}%
\pgfsys@transformshift{3.146894in}{3.075278in}%
\pgfsys@useobject{currentmarker}{}%
\end{pgfscope}%
\begin{pgfscope}%
\pgfsys@transformshift{3.150097in}{3.075278in}%
\pgfsys@useobject{currentmarker}{}%
\end{pgfscope}%
\begin{pgfscope}%
\pgfsys@transformshift{3.151445in}{3.075278in}%
\pgfsys@useobject{currentmarker}{}%
\end{pgfscope}%
\begin{pgfscope}%
\pgfsys@transformshift{3.152972in}{3.075278in}%
\pgfsys@useobject{currentmarker}{}%
\end{pgfscope}%
\begin{pgfscope}%
\pgfsys@transformshift{3.154491in}{3.075278in}%
\pgfsys@useobject{currentmarker}{}%
\end{pgfscope}%
\begin{pgfscope}%
\pgfsys@transformshift{3.155996in}{3.075278in}%
\pgfsys@useobject{currentmarker}{}%
\end{pgfscope}%
\begin{pgfscope}%
\pgfsys@transformshift{3.157009in}{3.075278in}%
\pgfsys@useobject{currentmarker}{}%
\end{pgfscope}%
\begin{pgfscope}%
\pgfsys@transformshift{3.158357in}{3.075278in}%
\pgfsys@useobject{currentmarker}{}%
\end{pgfscope}%
\begin{pgfscope}%
\pgfsys@transformshift{3.159548in}{3.075278in}%
\pgfsys@useobject{currentmarker}{}%
\end{pgfscope}%
\begin{pgfscope}%
\pgfsys@transformshift{3.160897in}{3.075278in}%
\pgfsys@useobject{currentmarker}{}%
\end{pgfscope}%
\begin{pgfscope}%
\pgfsys@transformshift{3.161910in}{3.075278in}%
\pgfsys@useobject{currentmarker}{}%
\end{pgfscope}%
\begin{pgfscope}%
\pgfsys@transformshift{3.162744in}{3.075278in}%
\pgfsys@useobject{currentmarker}{}%
\end{pgfscope}%
\begin{pgfscope}%
\pgfsys@transformshift{3.163928in}{3.075278in}%
\pgfsys@useobject{currentmarker}{}%
\end{pgfscope}%
\begin{pgfscope}%
\pgfsys@transformshift{3.164770in}{3.075278in}%
\pgfsys@useobject{currentmarker}{}%
\end{pgfscope}%
\begin{pgfscope}%
\pgfsys@transformshift{3.166282in}{3.075278in}%
\pgfsys@useobject{currentmarker}{}%
\end{pgfscope}%
\begin{pgfscope}%
\pgfsys@transformshift{3.166960in}{3.075278in}%
\pgfsys@useobject{currentmarker}{}%
\end{pgfscope}%
\begin{pgfscope}%
\pgfsys@transformshift{3.167809in}{3.075278in}%
\pgfsys@useobject{currentmarker}{}%
\end{pgfscope}%
\begin{pgfscope}%
\pgfsys@transformshift{3.169157in}{3.075278in}%
\pgfsys@useobject{currentmarker}{}%
\end{pgfscope}%
\begin{pgfscope}%
\pgfsys@transformshift{3.169998in}{3.075278in}%
\pgfsys@useobject{currentmarker}{}%
\end{pgfscope}%
\begin{pgfscope}%
\pgfsys@transformshift{3.171175in}{3.075278in}%
\pgfsys@useobject{currentmarker}{}%
\end{pgfscope}%
\begin{pgfscope}%
\pgfsys@transformshift{3.172024in}{3.075278in}%
\pgfsys@useobject{currentmarker}{}%
\end{pgfscope}%
\begin{pgfscope}%
\pgfsys@transformshift{3.173030in}{3.075278in}%
\pgfsys@useobject{currentmarker}{}%
\end{pgfscope}%
\begin{pgfscope}%
\pgfsys@transformshift{3.173708in}{3.075278in}%
\pgfsys@useobject{currentmarker}{}%
\end{pgfscope}%
\begin{pgfscope}%
\pgfsys@transformshift{3.174713in}{3.075278in}%
\pgfsys@useobject{currentmarker}{}%
\end{pgfscope}%
\begin{pgfscope}%
\pgfsys@transformshift{3.175391in}{3.075278in}%
\pgfsys@useobject{currentmarker}{}%
\end{pgfscope}%
\begin{pgfscope}%
\pgfsys@transformshift{3.176069in}{3.075278in}%
\pgfsys@useobject{currentmarker}{}%
\end{pgfscope}%
\begin{pgfscope}%
\pgfsys@transformshift{3.177417in}{3.075278in}%
\pgfsys@useobject{currentmarker}{}%
\end{pgfscope}%
\begin{pgfscope}%
\pgfsys@transformshift{3.178087in}{3.075278in}%
\pgfsys@useobject{currentmarker}{}%
\end{pgfscope}%
\begin{pgfscope}%
\pgfsys@transformshift{3.178772in}{3.075278in}%
\pgfsys@useobject{currentmarker}{}%
\end{pgfscope}%
\begin{pgfscope}%
\pgfsys@transformshift{3.179778in}{3.075278in}%
\pgfsys@useobject{currentmarker}{}%
\end{pgfscope}%
\begin{pgfscope}%
\pgfsys@transformshift{3.181126in}{3.075278in}%
\pgfsys@useobject{currentmarker}{}%
\end{pgfscope}%
\begin{pgfscope}%
\pgfsys@transformshift{3.182139in}{3.075278in}%
\pgfsys@useobject{currentmarker}{}%
\end{pgfscope}%
\begin{pgfscope}%
\pgfsys@transformshift{3.182809in}{3.075278in}%
\pgfsys@useobject{currentmarker}{}%
\end{pgfscope}%
\begin{pgfscope}%
\pgfsys@transformshift{3.184329in}{3.075278in}%
\pgfsys@useobject{currentmarker}{}%
\end{pgfscope}%
\begin{pgfscope}%
\pgfsys@transformshift{3.185684in}{3.075278in}%
\pgfsys@useobject{currentmarker}{}%
\end{pgfscope}%
\begin{pgfscope}%
\pgfsys@transformshift{3.186533in}{3.075278in}%
\pgfsys@useobject{currentmarker}{}%
\end{pgfscope}%
\begin{pgfscope}%
\pgfsys@transformshift{3.188209in}{3.075278in}%
\pgfsys@useobject{currentmarker}{}%
\end{pgfscope}%
\begin{pgfscope}%
\pgfsys@transformshift{3.189058in}{3.075278in}%
\pgfsys@useobject{currentmarker}{}%
\end{pgfscope}%
\begin{pgfscope}%
\pgfsys@transformshift{3.191256in}{3.075278in}%
\pgfsys@useobject{currentmarker}{}%
\end{pgfscope}%
\begin{pgfscope}%
\pgfsys@transformshift{3.199851in}{3.075278in}%
\pgfsys@useobject{currentmarker}{}%
\end{pgfscope}%
\begin{pgfscope}%
\pgfsys@transformshift{3.205757in}{3.075278in}%
\pgfsys@useobject{currentmarker}{}%
\end{pgfscope}%
\begin{pgfscope}%
\pgfsys@transformshift{3.452786in}{3.075278in}%
\pgfsys@useobject{currentmarker}{}%
\end{pgfscope}%
\begin{pgfscope}%
\pgfsys@transformshift{3.467459in}{3.075278in}%
\pgfsys@useobject{currentmarker}{}%
\end{pgfscope}%
\begin{pgfscope}%
\pgfsys@transformshift{3.470319in}{3.075278in}%
\pgfsys@useobject{currentmarker}{}%
\end{pgfscope}%
\begin{pgfscope}%
\pgfsys@transformshift{3.472516in}{3.075278in}%
\pgfsys@useobject{currentmarker}{}%
\end{pgfscope}%
\begin{pgfscope}%
\pgfsys@transformshift{3.474877in}{3.075278in}%
\pgfsys@useobject{currentmarker}{}%
\end{pgfscope}%
\begin{pgfscope}%
\pgfsys@transformshift{3.476732in}{3.075278in}%
\pgfsys@useobject{currentmarker}{}%
\end{pgfscope}%
\begin{pgfscope}%
\pgfsys@transformshift{3.478922in}{3.075278in}%
\pgfsys@useobject{currentmarker}{}%
\end{pgfscope}%
\begin{pgfscope}%
\pgfsys@transformshift{3.480270in}{3.075278in}%
\pgfsys@useobject{currentmarker}{}%
\end{pgfscope}%
\begin{pgfscope}%
\pgfsys@transformshift{3.482132in}{3.075278in}%
\pgfsys@useobject{currentmarker}{}%
\end{pgfscope}%
\begin{pgfscope}%
\pgfsys@transformshift{3.483473in}{3.075278in}%
\pgfsys@useobject{currentmarker}{}%
\end{pgfscope}%
\begin{pgfscope}%
\pgfsys@transformshift{3.484151in}{3.075278in}%
\pgfsys@useobject{currentmarker}{}%
\end{pgfscope}%
\begin{pgfscope}%
\pgfsys@transformshift{3.484992in}{3.075278in}%
\pgfsys@useobject{currentmarker}{}%
\end{pgfscope}%
\begin{pgfscope}%
\pgfsys@transformshift{3.485834in}{3.075278in}%
\pgfsys@useobject{currentmarker}{}%
\end{pgfscope}%
\begin{pgfscope}%
\pgfsys@transformshift{3.486847in}{3.075278in}%
\pgfsys@useobject{currentmarker}{}%
\end{pgfscope}%
\begin{pgfscope}%
\pgfsys@transformshift{3.488366in}{3.075278in}%
\pgfsys@useobject{currentmarker}{}%
\end{pgfscope}%
\begin{pgfscope}%
\pgfsys@transformshift{3.489208in}{3.075278in}%
\pgfsys@useobject{currentmarker}{}%
\end{pgfscope}%
\begin{pgfscope}%
\pgfsys@transformshift{3.490050in}{3.075278in}%
\pgfsys@useobject{currentmarker}{}%
\end{pgfscope}%
\begin{pgfscope}%
\pgfsys@transformshift{3.490891in}{3.075278in}%
\pgfsys@useobject{currentmarker}{}%
\end{pgfscope}%
\begin{pgfscope}%
\pgfsys@transformshift{3.491912in}{3.075278in}%
\pgfsys@useobject{currentmarker}{}%
\end{pgfscope}%
\begin{pgfscope}%
\pgfsys@transformshift{3.492925in}{3.075278in}%
\pgfsys@useobject{currentmarker}{}%
\end{pgfscope}%
\begin{pgfscope}%
\pgfsys@transformshift{3.493759in}{3.075278in}%
\pgfsys@useobject{currentmarker}{}%
\end{pgfscope}%
\begin{pgfscope}%
\pgfsys@transformshift{3.494437in}{3.075278in}%
\pgfsys@useobject{currentmarker}{}%
\end{pgfscope}%
\begin{pgfscope}%
\pgfsys@transformshift{3.495449in}{3.075278in}%
\pgfsys@useobject{currentmarker}{}%
\end{pgfscope}%
\begin{pgfscope}%
\pgfsys@transformshift{3.496299in}{3.075278in}%
\pgfsys@useobject{currentmarker}{}%
\end{pgfscope}%
\begin{pgfscope}%
\pgfsys@transformshift{3.497319in}{3.075278in}%
\pgfsys@useobject{currentmarker}{}%
\end{pgfscope}%
\begin{pgfscope}%
\pgfsys@transformshift{3.497974in}{3.075278in}%
\pgfsys@useobject{currentmarker}{}%
\end{pgfscope}%
\begin{pgfscope}%
\pgfsys@transformshift{3.499323in}{3.075278in}%
\pgfsys@useobject{currentmarker}{}%
\end{pgfscope}%
\begin{pgfscope}%
\pgfsys@transformshift{3.500507in}{3.075278in}%
\pgfsys@useobject{currentmarker}{}%
\end{pgfscope}%
\begin{pgfscope}%
\pgfsys@transformshift{3.501356in}{3.075278in}%
\pgfsys@useobject{currentmarker}{}%
\end{pgfscope}%
\begin{pgfscope}%
\pgfsys@transformshift{3.502034in}{3.075278in}%
\pgfsys@useobject{currentmarker}{}%
\end{pgfscope}%
\begin{pgfscope}%
\pgfsys@transformshift{3.503047in}{3.075278in}%
\pgfsys@useobject{currentmarker}{}%
\end{pgfscope}%
\begin{pgfscope}%
\pgfsys@transformshift{3.504216in}{3.075278in}%
\pgfsys@useobject{currentmarker}{}%
\end{pgfscope}%
\begin{pgfscope}%
\pgfsys@transformshift{3.505743in}{3.075278in}%
\pgfsys@useobject{currentmarker}{}%
\end{pgfscope}%
\begin{pgfscope}%
\pgfsys@transformshift{3.507084in}{3.075278in}%
\pgfsys@useobject{currentmarker}{}%
\end{pgfscope}%
\begin{pgfscope}%
\pgfsys@transformshift{3.507754in}{3.075278in}%
\pgfsys@useobject{currentmarker}{}%
\end{pgfscope}%
\begin{pgfscope}%
\pgfsys@transformshift{3.508938in}{3.075278in}%
\pgfsys@useobject{currentmarker}{}%
\end{pgfscope}%
\begin{pgfscope}%
\pgfsys@transformshift{3.509959in}{3.075278in}%
\pgfsys@useobject{currentmarker}{}%
\end{pgfscope}%
\begin{pgfscope}%
\pgfsys@transformshift{3.510629in}{3.075278in}%
\pgfsys@useobject{currentmarker}{}%
\end{pgfscope}%
\begin{pgfscope}%
\pgfsys@transformshift{3.511821in}{3.075278in}%
\pgfsys@useobject{currentmarker}{}%
\end{pgfscope}%
\begin{pgfscope}%
\pgfsys@transformshift{3.513169in}{3.075278in}%
\pgfsys@useobject{currentmarker}{}%
\end{pgfscope}%
\begin{pgfscope}%
\pgfsys@transformshift{3.515031in}{3.075278in}%
\pgfsys@useobject{currentmarker}{}%
\end{pgfscope}%
\begin{pgfscope}%
\pgfsys@transformshift{3.516550in}{3.075278in}%
\pgfsys@useobject{currentmarker}{}%
\end{pgfscope}%
\begin{pgfscope}%
\pgfsys@transformshift{3.779920in}{3.075278in}%
\pgfsys@useobject{currentmarker}{}%
\end{pgfscope}%
\begin{pgfscope}%
\pgfsys@transformshift{3.791391in}{3.075278in}%
\pgfsys@useobject{currentmarker}{}%
\end{pgfscope}%
\begin{pgfscope}%
\pgfsys@transformshift{3.795606in}{3.075278in}%
\pgfsys@useobject{currentmarker}{}%
\end{pgfscope}%
\begin{pgfscope}%
\pgfsys@transformshift{3.797461in}{3.075278in}%
\pgfsys@useobject{currentmarker}{}%
\end{pgfscope}%
\begin{pgfscope}%
\pgfsys@transformshift{3.799993in}{3.075278in}%
\pgfsys@useobject{currentmarker}{}%
\end{pgfscope}%
\begin{pgfscope}%
\pgfsys@transformshift{3.802012in}{3.075278in}%
\pgfsys@useobject{currentmarker}{}%
\end{pgfscope}%
\begin{pgfscope}%
\pgfsys@transformshift{3.803703in}{3.075278in}%
\pgfsys@useobject{currentmarker}{}%
\end{pgfscope}%
\begin{pgfscope}%
\pgfsys@transformshift{3.804879in}{3.075278in}%
\pgfsys@useobject{currentmarker}{}%
\end{pgfscope}%
\begin{pgfscope}%
\pgfsys@transformshift{3.806228in}{3.075278in}%
\pgfsys@useobject{currentmarker}{}%
\end{pgfscope}%
\begin{pgfscope}%
\pgfsys@transformshift{3.807412in}{3.075278in}%
\pgfsys@useobject{currentmarker}{}%
\end{pgfscope}%
\begin{pgfscope}%
\pgfsys@transformshift{3.808589in}{3.075278in}%
\pgfsys@useobject{currentmarker}{}%
\end{pgfscope}%
\begin{pgfscope}%
\pgfsys@transformshift{3.809937in}{3.075278in}%
\pgfsys@useobject{currentmarker}{}%
\end{pgfscope}%
\begin{pgfscope}%
\pgfsys@transformshift{3.811121in}{3.075278in}%
\pgfsys@useobject{currentmarker}{}%
\end{pgfscope}%
\begin{pgfscope}%
\pgfsys@transformshift{3.812134in}{3.075278in}%
\pgfsys@useobject{currentmarker}{}%
\end{pgfscope}%
\begin{pgfscope}%
\pgfsys@transformshift{3.813482in}{3.075278in}%
\pgfsys@useobject{currentmarker}{}%
\end{pgfscope}%
\begin{pgfscope}%
\pgfsys@transformshift{3.814495in}{3.075278in}%
\pgfsys@useobject{currentmarker}{}%
\end{pgfscope}%
\begin{pgfscope}%
\pgfsys@transformshift{3.815672in}{3.075278in}%
\pgfsys@useobject{currentmarker}{}%
\end{pgfscope}%
\begin{pgfscope}%
\pgfsys@transformshift{3.816685in}{3.075278in}%
\pgfsys@useobject{currentmarker}{}%
\end{pgfscope}%
\begin{pgfscope}%
\pgfsys@transformshift{3.817415in}{3.075278in}%
\pgfsys@useobject{currentmarker}{}%
\end{pgfscope}%
\begin{pgfscope}%
\pgfsys@transformshift{3.818547in}{3.075278in}%
\pgfsys@useobject{currentmarker}{}%
\end{pgfscope}%
\begin{pgfscope}%
\pgfsys@transformshift{3.819225in}{3.075278in}%
\pgfsys@useobject{currentmarker}{}%
\end{pgfscope}%
\begin{pgfscope}%
\pgfsys@transformshift{3.820402in}{3.075278in}%
\pgfsys@useobject{currentmarker}{}%
\end{pgfscope}%
\begin{pgfscope}%
\pgfsys@transformshift{3.821415in}{3.075278in}%
\pgfsys@useobject{currentmarker}{}%
\end{pgfscope}%
\begin{pgfscope}%
\pgfsys@transformshift{3.822428in}{3.075278in}%
\pgfsys@useobject{currentmarker}{}%
\end{pgfscope}%
\begin{pgfscope}%
\pgfsys@transformshift{3.823433in}{3.075278in}%
\pgfsys@useobject{currentmarker}{}%
\end{pgfscope}%
\begin{pgfscope}%
\pgfsys@transformshift{3.824446in}{3.075278in}%
\pgfsys@useobject{currentmarker}{}%
\end{pgfscope}%
\begin{pgfscope}%
\pgfsys@transformshift{3.825295in}{3.075278in}%
\pgfsys@useobject{currentmarker}{}%
\end{pgfscope}%
\begin{pgfscope}%
\pgfsys@transformshift{3.826301in}{3.075278in}%
\pgfsys@useobject{currentmarker}{}%
\end{pgfscope}%
\begin{pgfscope}%
\pgfsys@transformshift{3.827500in}{3.075278in}%
\pgfsys@useobject{currentmarker}{}%
\end{pgfscope}%
\begin{pgfscope}%
\pgfsys@transformshift{3.828319in}{3.075278in}%
\pgfsys@useobject{currentmarker}{}%
\end{pgfscope}%
\begin{pgfscope}%
\pgfsys@transformshift{3.829332in}{3.075278in}%
\pgfsys@useobject{currentmarker}{}%
\end{pgfscope}%
\begin{pgfscope}%
\pgfsys@transformshift{3.830010in}{3.075278in}%
\pgfsys@useobject{currentmarker}{}%
\end{pgfscope}%
\begin{pgfscope}%
\pgfsys@transformshift{3.830859in}{3.075278in}%
\pgfsys@useobject{currentmarker}{}%
\end{pgfscope}%
\begin{pgfscope}%
\pgfsys@transformshift{3.831529in}{3.075278in}%
\pgfsys@useobject{currentmarker}{}%
\end{pgfscope}%
\begin{pgfscope}%
\pgfsys@transformshift{3.832535in}{3.075278in}%
\pgfsys@useobject{currentmarker}{}%
\end{pgfscope}%
\begin{pgfscope}%
\pgfsys@transformshift{3.833213in}{3.075278in}%
\pgfsys@useobject{currentmarker}{}%
\end{pgfscope}%
\begin{pgfscope}%
\pgfsys@transformshift{3.834397in}{3.075278in}%
\pgfsys@useobject{currentmarker}{}%
\end{pgfscope}%
\begin{pgfscope}%
\pgfsys@transformshift{3.834911in}{3.075278in}%
\pgfsys@useobject{currentmarker}{}%
\end{pgfscope}%
\begin{pgfscope}%
\pgfsys@transformshift{3.835752in}{3.075278in}%
\pgfsys@useobject{currentmarker}{}%
\end{pgfscope}%
\begin{pgfscope}%
\pgfsys@transformshift{3.836587in}{3.075278in}%
\pgfsys@useobject{currentmarker}{}%
\end{pgfscope}%
\begin{pgfscope}%
\pgfsys@transformshift{3.837101in}{3.075278in}%
\pgfsys@useobject{currentmarker}{}%
\end{pgfscope}%
\begin{pgfscope}%
\pgfsys@transformshift{3.838277in}{3.075278in}%
\pgfsys@useobject{currentmarker}{}%
\end{pgfscope}%
\begin{pgfscope}%
\pgfsys@transformshift{3.839626in}{3.075278in}%
\pgfsys@useobject{currentmarker}{}%
\end{pgfscope}%
\begin{pgfscope}%
\pgfsys@transformshift{3.840467in}{3.075278in}%
\pgfsys@useobject{currentmarker}{}%
\end{pgfscope}%
\begin{pgfscope}%
\pgfsys@transformshift{3.841659in}{3.075278in}%
\pgfsys@useobject{currentmarker}{}%
\end{pgfscope}%
\begin{pgfscope}%
\pgfsys@transformshift{3.843000in}{3.075278in}%
\pgfsys@useobject{currentmarker}{}%
\end{pgfscope}%
\begin{pgfscope}%
\pgfsys@transformshift{3.844363in}{3.075278in}%
\pgfsys@useobject{currentmarker}{}%
\end{pgfscope}%
\begin{pgfscope}%
\pgfsys@transformshift{3.845361in}{3.075278in}%
\pgfsys@useobject{currentmarker}{}%
\end{pgfscope}%
\begin{pgfscope}%
\pgfsys@transformshift{4.154113in}{3.075278in}%
\pgfsys@useobject{currentmarker}{}%
\end{pgfscope}%
\begin{pgfscope}%
\pgfsys@transformshift{4.163051in}{3.075278in}%
\pgfsys@useobject{currentmarker}{}%
\end{pgfscope}%
\begin{pgfscope}%
\pgfsys@transformshift{4.165911in}{3.075278in}%
\pgfsys@useobject{currentmarker}{}%
\end{pgfscope}%
\begin{pgfscope}%
\pgfsys@transformshift{4.168108in}{3.075278in}%
\pgfsys@useobject{currentmarker}{}%
\end{pgfscope}%
\begin{pgfscope}%
\pgfsys@transformshift{4.169620in}{3.075278in}%
\pgfsys@useobject{currentmarker}{}%
\end{pgfscope}%
\begin{pgfscope}%
\pgfsys@transformshift{4.171139in}{3.075278in}%
\pgfsys@useobject{currentmarker}{}%
\end{pgfscope}%
\begin{pgfscope}%
\pgfsys@transformshift{4.171981in}{3.075278in}%
\pgfsys@useobject{currentmarker}{}%
\end{pgfscope}%
\begin{pgfscope}%
\pgfsys@transformshift{4.172830in}{3.075278in}%
\pgfsys@useobject{currentmarker}{}%
\end{pgfscope}%
\begin{pgfscope}%
\pgfsys@transformshift{4.173836in}{3.075278in}%
\pgfsys@useobject{currentmarker}{}%
\end{pgfscope}%
\begin{pgfscope}%
\pgfsys@transformshift{4.174514in}{3.075278in}%
\pgfsys@useobject{currentmarker}{}%
\end{pgfscope}%
\begin{pgfscope}%
\pgfsys@transformshift{4.175355in}{3.075278in}%
\pgfsys@useobject{currentmarker}{}%
\end{pgfscope}%
\begin{pgfscope}%
\pgfsys@transformshift{4.176539in}{3.075278in}%
\pgfsys@useobject{currentmarker}{}%
\end{pgfscope}%
\begin{pgfscope}%
\pgfsys@transformshift{4.177716in}{3.075278in}%
\pgfsys@useobject{currentmarker}{}%
\end{pgfscope}%
\begin{pgfscope}%
\pgfsys@transformshift{4.178223in}{3.075278in}%
\pgfsys@useobject{currentmarker}{}%
\end{pgfscope}%
\begin{pgfscope}%
\pgfsys@transformshift{4.178901in}{3.075278in}%
\pgfsys@useobject{currentmarker}{}%
\end{pgfscope}%
\begin{pgfscope}%
\pgfsys@transformshift{4.179750in}{3.075278in}%
\pgfsys@useobject{currentmarker}{}%
\end{pgfscope}%
\begin{pgfscope}%
\pgfsys@transformshift{4.180413in}{3.075278in}%
\pgfsys@useobject{currentmarker}{}%
\end{pgfscope}%
\begin{pgfscope}%
\pgfsys@transformshift{4.181262in}{3.075278in}%
\pgfsys@useobject{currentmarker}{}%
\end{pgfscope}%
\begin{pgfscope}%
\pgfsys@transformshift{4.181768in}{3.075278in}%
\pgfsys@useobject{currentmarker}{}%
\end{pgfscope}%
\begin{pgfscope}%
\pgfsys@transformshift{4.182774in}{3.075278in}%
\pgfsys@useobject{currentmarker}{}%
\end{pgfscope}%
\begin{pgfscope}%
\pgfsys@transformshift{4.183280in}{3.075278in}%
\pgfsys@useobject{currentmarker}{}%
\end{pgfscope}%
\begin{pgfscope}%
\pgfsys@transformshift{4.184293in}{3.075278in}%
\pgfsys@useobject{currentmarker}{}%
\end{pgfscope}%
\begin{pgfscope}%
\pgfsys@transformshift{4.185135in}{3.075278in}%
\pgfsys@useobject{currentmarker}{}%
\end{pgfscope}%
\begin{pgfscope}%
\pgfsys@transformshift{4.185984in}{3.075278in}%
\pgfsys@useobject{currentmarker}{}%
\end{pgfscope}%
\begin{pgfscope}%
\pgfsys@transformshift{4.186654in}{3.075278in}%
\pgfsys@useobject{currentmarker}{}%
\end{pgfscope}%
\begin{pgfscope}%
\pgfsys@transformshift{4.187667in}{3.075278in}%
\pgfsys@useobject{currentmarker}{}%
\end{pgfscope}%
\begin{pgfscope}%
\pgfsys@transformshift{4.188680in}{3.075278in}%
\pgfsys@useobject{currentmarker}{}%
\end{pgfscope}%
\begin{pgfscope}%
\pgfsys@transformshift{4.189350in}{3.075278in}%
\pgfsys@useobject{currentmarker}{}%
\end{pgfscope}%
\begin{pgfscope}%
\pgfsys@transformshift{4.190028in}{3.075278in}%
\pgfsys@useobject{currentmarker}{}%
\end{pgfscope}%
\begin{pgfscope}%
\pgfsys@transformshift{4.190699in}{3.075278in}%
\pgfsys@useobject{currentmarker}{}%
\end{pgfscope}%
\begin{pgfscope}%
\pgfsys@transformshift{4.191205in}{3.075278in}%
\pgfsys@useobject{currentmarker}{}%
\end{pgfscope}%
\begin{pgfscope}%
\pgfsys@transformshift{4.192054in}{3.075278in}%
\pgfsys@useobject{currentmarker}{}%
\end{pgfscope}%
\begin{pgfscope}%
\pgfsys@transformshift{4.192732in}{3.075278in}%
\pgfsys@useobject{currentmarker}{}%
\end{pgfscope}%
\begin{pgfscope}%
\pgfsys@transformshift{4.193231in}{3.075278in}%
\pgfsys@useobject{currentmarker}{}%
\end{pgfscope}%
\begin{pgfscope}%
\pgfsys@transformshift{4.193745in}{3.075278in}%
\pgfsys@useobject{currentmarker}{}%
\end{pgfscope}%
\begin{pgfscope}%
\pgfsys@transformshift{4.194244in}{3.075278in}%
\pgfsys@useobject{currentmarker}{}%
\end{pgfscope}%
\begin{pgfscope}%
\pgfsys@transformshift{4.194922in}{3.075278in}%
\pgfsys@useobject{currentmarker}{}%
\end{pgfscope}%
\begin{pgfscope}%
\pgfsys@transformshift{4.195935in}{3.075278in}%
\pgfsys@useobject{currentmarker}{}%
\end{pgfscope}%
\begin{pgfscope}%
\pgfsys@transformshift{4.196612in}{3.075278in}%
\pgfsys@useobject{currentmarker}{}%
\end{pgfscope}%
\begin{pgfscope}%
\pgfsys@transformshift{4.197454in}{3.075278in}%
\pgfsys@useobject{currentmarker}{}%
\end{pgfscope}%
\begin{pgfscope}%
\pgfsys@transformshift{4.198132in}{3.075278in}%
\pgfsys@useobject{currentmarker}{}%
\end{pgfscope}%
\begin{pgfscope}%
\pgfsys@transformshift{4.198467in}{3.075278in}%
\pgfsys@useobject{currentmarker}{}%
\end{pgfscope}%
\begin{pgfscope}%
\pgfsys@transformshift{4.198974in}{3.075278in}%
\pgfsys@useobject{currentmarker}{}%
\end{pgfscope}%
\begin{pgfscope}%
\pgfsys@transformshift{4.199987in}{3.075278in}%
\pgfsys@useobject{currentmarker}{}%
\end{pgfscope}%
\begin{pgfscope}%
\pgfsys@transformshift{4.200493in}{3.075278in}%
\pgfsys@useobject{currentmarker}{}%
\end{pgfscope}%
\begin{pgfscope}%
\pgfsys@transformshift{4.200992in}{3.075278in}%
\pgfsys@useobject{currentmarker}{}%
\end{pgfscope}%
\begin{pgfscope}%
\pgfsys@transformshift{4.202683in}{3.075278in}%
\pgfsys@useobject{currentmarker}{}%
\end{pgfscope}%
\begin{pgfscope}%
\pgfsys@transformshift{4.203696in}{3.075278in}%
\pgfsys@useobject{currentmarker}{}%
\end{pgfscope}%
\begin{pgfscope}%
\pgfsys@transformshift{4.204709in}{3.075278in}%
\pgfsys@useobject{currentmarker}{}%
\end{pgfscope}%
\begin{pgfscope}%
\pgfsys@transformshift{4.205714in}{3.075278in}%
\pgfsys@useobject{currentmarker}{}%
\end{pgfscope}%
\begin{pgfscope}%
\pgfsys@transformshift{4.206906in}{3.075278in}%
\pgfsys@useobject{currentmarker}{}%
\end{pgfscope}%
\begin{pgfscope}%
\pgfsys@transformshift{4.207748in}{3.075278in}%
\pgfsys@useobject{currentmarker}{}%
\end{pgfscope}%
\begin{pgfscope}%
\pgfsys@transformshift{4.208589in}{3.075278in}%
\pgfsys@useobject{currentmarker}{}%
\end{pgfscope}%
\begin{pgfscope}%
\pgfsys@transformshift{4.210615in}{3.075278in}%
\pgfsys@useobject{currentmarker}{}%
\end{pgfscope}%
\begin{pgfscope}%
\pgfsys@transformshift{4.211792in}{3.075278in}%
\pgfsys@useobject{currentmarker}{}%
\end{pgfscope}%
\begin{pgfscope}%
\pgfsys@transformshift{4.213818in}{3.075278in}%
\pgfsys@useobject{currentmarker}{}%
\end{pgfscope}%
\begin{pgfscope}%
\pgfsys@transformshift{4.215851in}{3.075278in}%
\pgfsys@useobject{currentmarker}{}%
\end{pgfscope}%
\begin{pgfscope}%
\pgfsys@transformshift{4.218205in}{3.075278in}%
\pgfsys@useobject{currentmarker}{}%
\end{pgfscope}%
\begin{pgfscope}%
\pgfsys@transformshift{4.475177in}{3.075278in}%
\pgfsys@useobject{currentmarker}{}%
\end{pgfscope}%
\begin{pgfscope}%
\pgfsys@transformshift{4.487318in}{3.075278in}%
\pgfsys@useobject{currentmarker}{}%
\end{pgfscope}%
\begin{pgfscope}%
\pgfsys@transformshift{4.496255in}{3.075278in}%
\pgfsys@useobject{currentmarker}{}%
\end{pgfscope}%
\begin{pgfscope}%
\pgfsys@transformshift{4.499458in}{3.075278in}%
\pgfsys@useobject{currentmarker}{}%
\end{pgfscope}%
\begin{pgfscope}%
\pgfsys@transformshift{4.501484in}{3.075278in}%
\pgfsys@useobject{currentmarker}{}%
\end{pgfscope}%
\begin{pgfscope}%
\pgfsys@transformshift{4.503674in}{3.075278in}%
\pgfsys@useobject{currentmarker}{}%
\end{pgfscope}%
\begin{pgfscope}%
\pgfsys@transformshift{4.505022in}{3.075278in}%
\pgfsys@useobject{currentmarker}{}%
\end{pgfscope}%
\begin{pgfscope}%
\pgfsys@transformshift{4.506035in}{3.075278in}%
\pgfsys@useobject{currentmarker}{}%
\end{pgfscope}%
\begin{pgfscope}%
\pgfsys@transformshift{4.507554in}{3.075278in}%
\pgfsys@useobject{currentmarker}{}%
\end{pgfscope}%
\begin{pgfscope}%
\pgfsys@transformshift{4.508903in}{3.075278in}%
\pgfsys@useobject{currentmarker}{}%
\end{pgfscope}%
\begin{pgfscope}%
\pgfsys@transformshift{4.509744in}{3.075278in}%
\pgfsys@useobject{currentmarker}{}%
\end{pgfscope}%
\begin{pgfscope}%
\pgfsys@transformshift{4.510928in}{3.075278in}%
\pgfsys@useobject{currentmarker}{}%
\end{pgfscope}%
\begin{pgfscope}%
\pgfsys@transformshift{4.511599in}{3.075278in}%
\pgfsys@useobject{currentmarker}{}%
\end{pgfscope}%
\begin{pgfscope}%
\pgfsys@transformshift{4.512619in}{3.075278in}%
\pgfsys@useobject{currentmarker}{}%
\end{pgfscope}%
\begin{pgfscope}%
\pgfsys@transformshift{4.513796in}{3.075278in}%
\pgfsys@useobject{currentmarker}{}%
\end{pgfscope}%
\begin{pgfscope}%
\pgfsys@transformshift{4.514310in}{3.075278in}%
\pgfsys@useobject{currentmarker}{}%
\end{pgfscope}%
\begin{pgfscope}%
\pgfsys@transformshift{4.515487in}{3.075278in}%
\pgfsys@useobject{currentmarker}{}%
\end{pgfscope}%
\begin{pgfscope}%
\pgfsys@transformshift{4.516328in}{3.075278in}%
\pgfsys@useobject{currentmarker}{}%
\end{pgfscope}%
\begin{pgfscope}%
\pgfsys@transformshift{4.517170in}{3.075278in}%
\pgfsys@useobject{currentmarker}{}%
\end{pgfscope}%
\begin{pgfscope}%
\pgfsys@transformshift{4.518012in}{3.075278in}%
\pgfsys@useobject{currentmarker}{}%
\end{pgfscope}%
\begin{pgfscope}%
\pgfsys@transformshift{4.518682in}{3.075278in}%
\pgfsys@useobject{currentmarker}{}%
\end{pgfscope}%
\begin{pgfscope}%
\pgfsys@transformshift{4.519695in}{3.075278in}%
\pgfsys@useobject{currentmarker}{}%
\end{pgfscope}%
\begin{pgfscope}%
\pgfsys@transformshift{4.520380in}{3.075278in}%
\pgfsys@useobject{currentmarker}{}%
\end{pgfscope}%
\begin{pgfscope}%
\pgfsys@transformshift{4.521386in}{3.075278in}%
\pgfsys@useobject{currentmarker}{}%
\end{pgfscope}%
\begin{pgfscope}%
\pgfsys@transformshift{4.522563in}{3.075278in}%
\pgfsys@useobject{currentmarker}{}%
\end{pgfscope}%
\begin{pgfscope}%
\pgfsys@transformshift{4.523419in}{3.075278in}%
\pgfsys@useobject{currentmarker}{}%
\end{pgfscope}%
\begin{pgfscope}%
\pgfsys@transformshift{4.523911in}{3.075278in}%
\pgfsys@useobject{currentmarker}{}%
\end{pgfscope}%
\begin{pgfscope}%
\pgfsys@transformshift{4.524767in}{3.075278in}%
\pgfsys@useobject{currentmarker}{}%
\end{pgfscope}%
\begin{pgfscope}%
\pgfsys@transformshift{4.525937in}{3.075278in}%
\pgfsys@useobject{currentmarker}{}%
\end{pgfscope}%
\begin{pgfscope}%
\pgfsys@transformshift{4.526443in}{3.075278in}%
\pgfsys@useobject{currentmarker}{}%
\end{pgfscope}%
\begin{pgfscope}%
\pgfsys@transformshift{4.527456in}{3.075278in}%
\pgfsys@useobject{currentmarker}{}%
\end{pgfscope}%
\begin{pgfscope}%
\pgfsys@transformshift{4.528305in}{3.075278in}%
\pgfsys@useobject{currentmarker}{}%
\end{pgfscope}%
\begin{pgfscope}%
\pgfsys@transformshift{4.528976in}{3.075278in}%
\pgfsys@useobject{currentmarker}{}%
\end{pgfscope}%
\begin{pgfscope}%
\pgfsys@transformshift{4.529996in}{3.075278in}%
\pgfsys@useobject{currentmarker}{}%
\end{pgfscope}%
\begin{pgfscope}%
\pgfsys@transformshift{4.533035in}{3.075278in}%
\pgfsys@useobject{currentmarker}{}%
\end{pgfscope}%
\begin{pgfscope}%
\pgfsys@transformshift{4.534539in}{3.075278in}%
\pgfsys@useobject{currentmarker}{}%
\end{pgfscope}%
\begin{pgfscope}%
\pgfsys@transformshift{4.535716in}{3.075278in}%
\pgfsys@useobject{currentmarker}{}%
\end{pgfscope}%
\begin{pgfscope}%
\pgfsys@transformshift{4.536744in}{3.075278in}%
\pgfsys@useobject{currentmarker}{}%
\end{pgfscope}%
\begin{pgfscope}%
\pgfsys@transformshift{4.538591in}{3.075278in}%
\pgfsys@useobject{currentmarker}{}%
\end{pgfscope}%
\begin{pgfscope}%
\pgfsys@transformshift{4.539954in}{3.075278in}%
\pgfsys@useobject{currentmarker}{}%
\end{pgfscope}%
\begin{pgfscope}%
\pgfsys@transformshift{4.542814in}{3.075278in}%
\pgfsys@useobject{currentmarker}{}%
\end{pgfscope}%
\begin{pgfscope}%
\pgfsys@transformshift{4.545183in}{3.075278in}%
\pgfsys@useobject{currentmarker}{}%
\end{pgfscope}%
\begin{pgfscope}%
\pgfsys@transformshift{4.548036in}{3.075278in}%
\pgfsys@useobject{currentmarker}{}%
\end{pgfscope}%
\begin{pgfscope}%
\pgfsys@transformshift{4.845146in}{3.075278in}%
\pgfsys@useobject{currentmarker}{}%
\end{pgfscope}%
\begin{pgfscope}%
\pgfsys@transformshift{4.857458in}{3.075278in}%
\pgfsys@useobject{currentmarker}{}%
\end{pgfscope}%
\begin{pgfscope}%
\pgfsys@transformshift{4.861503in}{3.075278in}%
\pgfsys@useobject{currentmarker}{}%
\end{pgfscope}%
\begin{pgfscope}%
\pgfsys@transformshift{4.863357in}{3.075278in}%
\pgfsys@useobject{currentmarker}{}%
\end{pgfscope}%
\begin{pgfscope}%
\pgfsys@transformshift{4.865048in}{3.075278in}%
\pgfsys@useobject{currentmarker}{}%
\end{pgfscope}%
\begin{pgfscope}%
\pgfsys@transformshift{4.867238in}{3.075278in}%
\pgfsys@useobject{currentmarker}{}%
\end{pgfscope}%
\begin{pgfscope}%
\pgfsys@transformshift{4.868757in}{3.075278in}%
\pgfsys@useobject{currentmarker}{}%
\end{pgfscope}%
\begin{pgfscope}%
\pgfsys@transformshift{4.869599in}{3.075278in}%
\pgfsys@useobject{currentmarker}{}%
\end{pgfscope}%
\begin{pgfscope}%
\pgfsys@transformshift{4.870448in}{3.075278in}%
\pgfsys@useobject{currentmarker}{}%
\end{pgfscope}%
\begin{pgfscope}%
\pgfsys@transformshift{4.871796in}{3.075278in}%
\pgfsys@useobject{currentmarker}{}%
\end{pgfscope}%
\begin{pgfscope}%
\pgfsys@transformshift{4.872809in}{3.075278in}%
\pgfsys@useobject{currentmarker}{}%
\end{pgfscope}%
\begin{pgfscope}%
\pgfsys@transformshift{4.873814in}{3.075278in}%
\pgfsys@useobject{currentmarker}{}%
\end{pgfscope}%
\begin{pgfscope}%
\pgfsys@transformshift{4.874500in}{3.075278in}%
\pgfsys@useobject{currentmarker}{}%
\end{pgfscope}%
\begin{pgfscope}%
\pgfsys@transformshift{4.875557in}{3.075278in}%
\pgfsys@useobject{currentmarker}{}%
\end{pgfscope}%
\begin{pgfscope}%
\pgfsys@transformshift{4.876347in}{3.075278in}%
\pgfsys@useobject{currentmarker}{}%
\end{pgfscope}%
\begin{pgfscope}%
\pgfsys@transformshift{4.877017in}{3.075278in}%
\pgfsys@useobject{currentmarker}{}%
\end{pgfscope}%
\begin{pgfscope}%
\pgfsys@transformshift{4.878030in}{3.075278in}%
\pgfsys@useobject{currentmarker}{}%
\end{pgfscope}%
\begin{pgfscope}%
\pgfsys@transformshift{4.879214in}{3.075278in}%
\pgfsys@useobject{currentmarker}{}%
\end{pgfscope}%
\begin{pgfscope}%
\pgfsys@transformshift{4.880056in}{3.075278in}%
\pgfsys@useobject{currentmarker}{}%
\end{pgfscope}%
\begin{pgfscope}%
\pgfsys@transformshift{4.881077in}{3.075278in}%
\pgfsys@useobject{currentmarker}{}%
\end{pgfscope}%
\begin{pgfscope}%
\pgfsys@transformshift{4.882253in}{3.075278in}%
\pgfsys@useobject{currentmarker}{}%
\end{pgfscope}%
\begin{pgfscope}%
\pgfsys@transformshift{4.883430in}{3.075278in}%
\pgfsys@useobject{currentmarker}{}%
\end{pgfscope}%
\begin{pgfscope}%
\pgfsys@transformshift{4.884100in}{3.075278in}%
\pgfsys@useobject{currentmarker}{}%
\end{pgfscope}%
\begin{pgfscope}%
\pgfsys@transformshift{4.885113in}{3.075278in}%
\pgfsys@useobject{currentmarker}{}%
\end{pgfscope}%
\begin{pgfscope}%
\pgfsys@transformshift{4.886134in}{3.075278in}%
\pgfsys@useobject{currentmarker}{}%
\end{pgfscope}%
\begin{pgfscope}%
\pgfsys@transformshift{4.887475in}{3.075278in}%
\pgfsys@useobject{currentmarker}{}%
\end{pgfscope}%
\begin{pgfscope}%
\pgfsys@transformshift{4.888651in}{3.075278in}%
\pgfsys@useobject{currentmarker}{}%
\end{pgfscope}%
\begin{pgfscope}%
\pgfsys@transformshift{4.889009in}{3.075278in}%
\pgfsys@useobject{currentmarker}{}%
\end{pgfscope}%
\begin{pgfscope}%
\pgfsys@transformshift{4.890350in}{3.075278in}%
\pgfsys@useobject{currentmarker}{}%
\end{pgfscope}%
\begin{pgfscope}%
\pgfsys@transformshift{4.891020in}{3.075278in}%
\pgfsys@useobject{currentmarker}{}%
\end{pgfscope}%
\begin{pgfscope}%
\pgfsys@transformshift{4.892025in}{3.075278in}%
\pgfsys@useobject{currentmarker}{}%
\end{pgfscope}%
\begin{pgfscope}%
\pgfsys@transformshift{4.893217in}{3.075278in}%
\pgfsys@useobject{currentmarker}{}%
\end{pgfscope}%
\begin{pgfscope}%
\pgfsys@transformshift{4.894900in}{3.075278in}%
\pgfsys@useobject{currentmarker}{}%
\end{pgfscope}%
\begin{pgfscope}%
\pgfsys@transformshift{4.895757in}{3.075278in}%
\pgfsys@useobject{currentmarker}{}%
\end{pgfscope}%
\begin{pgfscope}%
\pgfsys@transformshift{4.896926in}{3.075278in}%
\pgfsys@useobject{currentmarker}{}%
\end{pgfscope}%
\begin{pgfscope}%
\pgfsys@transformshift{4.898438in}{3.075278in}%
\pgfsys@useobject{currentmarker}{}%
\end{pgfscope}%
\begin{pgfscope}%
\pgfsys@transformshift{4.899451in}{3.075278in}%
\pgfsys@useobject{currentmarker}{}%
\end{pgfscope}%
\begin{pgfscope}%
\pgfsys@transformshift{4.900978in}{3.075278in}%
\pgfsys@useobject{currentmarker}{}%
\end{pgfscope}%
\begin{pgfscope}%
\pgfsys@transformshift{4.902654in}{3.075278in}%
\pgfsys@useobject{currentmarker}{}%
\end{pgfscope}%
\begin{pgfscope}%
\pgfsys@transformshift{4.904337in}{3.075278in}%
\pgfsys@useobject{currentmarker}{}%
\end{pgfscope}%
\begin{pgfscope}%
\pgfsys@transformshift{4.905358in}{3.075278in}%
\pgfsys@useobject{currentmarker}{}%
\end{pgfscope}%
\begin{pgfscope}%
\pgfsys@transformshift{4.907056in}{3.075278in}%
\pgfsys@useobject{currentmarker}{}%
\end{pgfscope}%
\begin{pgfscope}%
\pgfsys@transformshift{4.908896in}{3.075278in}%
\pgfsys@useobject{currentmarker}{}%
\end{pgfscope}%
\begin{pgfscope}%
\pgfsys@transformshift{4.910743in}{3.075278in}%
\pgfsys@useobject{currentmarker}{}%
\end{pgfscope}%
\begin{pgfscope}%
\pgfsys@transformshift{4.912955in}{3.075278in}%
\pgfsys@useobject{currentmarker}{}%
\end{pgfscope}%
\begin{pgfscope}%
\pgfsys@transformshift{5.165033in}{3.075278in}%
\pgfsys@useobject{currentmarker}{}%
\end{pgfscope}%
\begin{pgfscope}%
\pgfsys@transformshift{5.179535in}{3.075278in}%
\pgfsys@useobject{currentmarker}{}%
\end{pgfscope}%
\begin{pgfscope}%
\pgfsys@transformshift{5.183252in}{3.075278in}%
\pgfsys@useobject{currentmarker}{}%
\end{pgfscope}%
\begin{pgfscope}%
\pgfsys@transformshift{5.185948in}{3.075278in}%
\pgfsys@useobject{currentmarker}{}%
\end{pgfscope}%
\begin{pgfscope}%
\pgfsys@transformshift{5.188145in}{3.075278in}%
\pgfsys@useobject{currentmarker}{}%
\end{pgfscope}%
\begin{pgfscope}%
\pgfsys@transformshift{5.190328in}{3.075278in}%
\pgfsys@useobject{currentmarker}{}%
\end{pgfscope}%
\begin{pgfscope}%
\pgfsys@transformshift{5.191683in}{3.075278in}%
\pgfsys@useobject{currentmarker}{}%
\end{pgfscope}%
\begin{pgfscope}%
\pgfsys@transformshift{5.192689in}{3.075278in}%
\pgfsys@useobject{currentmarker}{}%
\end{pgfscope}%
\begin{pgfscope}%
\pgfsys@transformshift{5.194208in}{3.075278in}%
\pgfsys@useobject{currentmarker}{}%
\end{pgfscope}%
\begin{pgfscope}%
\pgfsys@transformshift{5.195221in}{3.075278in}%
\pgfsys@useobject{currentmarker}{}%
\end{pgfscope}%
\begin{pgfscope}%
\pgfsys@transformshift{5.196234in}{3.075278in}%
\pgfsys@useobject{currentmarker}{}%
\end{pgfscope}%
\begin{pgfscope}%
\pgfsys@transformshift{5.197582in}{3.075278in}%
\pgfsys@useobject{currentmarker}{}%
\end{pgfscope}%
\begin{pgfscope}%
\pgfsys@transformshift{5.199266in}{3.075278in}%
\pgfsys@useobject{currentmarker}{}%
\end{pgfscope}%
\begin{pgfscope}%
\pgfsys@transformshift{5.199943in}{3.075278in}%
\pgfsys@useobject{currentmarker}{}%
\end{pgfscope}%
\begin{pgfscope}%
\pgfsys@transformshift{5.200956in}{3.075278in}%
\pgfsys@useobject{currentmarker}{}%
\end{pgfscope}%
\begin{pgfscope}%
\pgfsys@transformshift{5.202304in}{3.075278in}%
\pgfsys@useobject{currentmarker}{}%
\end{pgfscope}%
\begin{pgfscope}%
\pgfsys@transformshift{5.203146in}{3.075278in}%
\pgfsys@useobject{currentmarker}{}%
\end{pgfscope}%
\begin{pgfscope}%
\pgfsys@transformshift{5.204159in}{3.075278in}%
\pgfsys@useobject{currentmarker}{}%
\end{pgfscope}%
\begin{pgfscope}%
\pgfsys@transformshift{5.205001in}{3.075278in}%
\pgfsys@useobject{currentmarker}{}%
\end{pgfscope}%
\begin{pgfscope}%
\pgfsys@transformshift{5.205679in}{3.075278in}%
\pgfsys@useobject{currentmarker}{}%
\end{pgfscope}%
\begin{pgfscope}%
\pgfsys@transformshift{5.206364in}{3.075278in}%
\pgfsys@useobject{currentmarker}{}%
\end{pgfscope}%
\begin{pgfscope}%
\pgfsys@transformshift{5.207034in}{3.075278in}%
\pgfsys@useobject{currentmarker}{}%
\end{pgfscope}%
\begin{pgfscope}%
\pgfsys@transformshift{5.207704in}{3.075278in}%
\pgfsys@useobject{currentmarker}{}%
\end{pgfscope}%
\begin{pgfscope}%
\pgfsys@transformshift{5.208881in}{3.075278in}%
\pgfsys@useobject{currentmarker}{}%
\end{pgfscope}%
\begin{pgfscope}%
\pgfsys@transformshift{5.209552in}{3.075278in}%
\pgfsys@useobject{currentmarker}{}%
\end{pgfscope}%
\begin{pgfscope}%
\pgfsys@transformshift{5.210565in}{3.075278in}%
\pgfsys@useobject{currentmarker}{}%
\end{pgfscope}%
\begin{pgfscope}%
\pgfsys@transformshift{5.211913in}{3.075278in}%
\pgfsys@useobject{currentmarker}{}%
\end{pgfscope}%
\begin{pgfscope}%
\pgfsys@transformshift{5.213097in}{3.075278in}%
\pgfsys@useobject{currentmarker}{}%
\end{pgfscope}%
\begin{pgfscope}%
\pgfsys@transformshift{5.214281in}{3.075278in}%
\pgfsys@useobject{currentmarker}{}%
\end{pgfscope}%
\begin{pgfscope}%
\pgfsys@transformshift{5.215622in}{3.075278in}%
\pgfsys@useobject{currentmarker}{}%
\end{pgfscope}%
\begin{pgfscope}%
\pgfsys@transformshift{5.216635in}{3.075278in}%
\pgfsys@useobject{currentmarker}{}%
\end{pgfscope}%
\begin{pgfscope}%
\pgfsys@transformshift{5.217990in}{3.075278in}%
\pgfsys@useobject{currentmarker}{}%
\end{pgfscope}%
\begin{pgfscope}%
\pgfsys@transformshift{5.219175in}{3.075278in}%
\pgfsys@useobject{currentmarker}{}%
\end{pgfscope}%
\begin{pgfscope}%
\pgfsys@transformshift{5.220687in}{3.075278in}%
\pgfsys@useobject{currentmarker}{}%
\end{pgfscope}%
\begin{pgfscope}%
\pgfsys@transformshift{5.222035in}{3.075278in}%
\pgfsys@useobject{currentmarker}{}%
\end{pgfscope}%
\begin{pgfscope}%
\pgfsys@transformshift{5.223390in}{3.075278in}%
\pgfsys@useobject{currentmarker}{}%
\end{pgfscope}%
\begin{pgfscope}%
\pgfsys@transformshift{5.225409in}{3.075278in}%
\pgfsys@useobject{currentmarker}{}%
\end{pgfscope}%
\begin{pgfscope}%
\pgfsys@transformshift{5.226765in}{3.075278in}%
\pgfsys@useobject{currentmarker}{}%
\end{pgfscope}%
\begin{pgfscope}%
\pgfsys@transformshift{5.227770in}{3.075278in}%
\pgfsys@useobject{currentmarker}{}%
\end{pgfscope}%
\begin{pgfscope}%
\pgfsys@transformshift{5.229304in}{3.075278in}%
\pgfsys@useobject{currentmarker}{}%
\end{pgfscope}%
\begin{pgfscope}%
\pgfsys@transformshift{5.231330in}{3.075278in}%
\pgfsys@useobject{currentmarker}{}%
\end{pgfscope}%
\begin{pgfscope}%
\pgfsys@transformshift{5.233505in}{3.075278in}%
\pgfsys@useobject{currentmarker}{}%
\end{pgfscope}%
\begin{pgfscope}%
\pgfsys@transformshift{5.235360in}{3.075278in}%
\pgfsys@useobject{currentmarker}{}%
\end{pgfscope}%
\begin{pgfscope}%
\pgfsys@transformshift{5.238242in}{3.075278in}%
\pgfsys@useobject{currentmarker}{}%
\end{pgfscope}%
\end{pgfscope}%
\begin{pgfscope}%
\pgfpathrectangle{\pgfqpoint{0.713889in}{0.502778in}}{\pgfqpoint{4.805000in}{2.695000in}}%
\pgfusepath{clip}%
\pgfsetrectcap%
\pgfsetroundjoin%
\pgfsetlinewidth{0.803000pt}%
\definecolor{currentstroke}{rgb}{0.690196,0.690196,0.690196}%
\pgfsetstrokecolor{currentstroke}%
\pgfsetdash{}{0pt}%
\pgfpathmoveto{\pgfqpoint{0.835189in}{0.502778in}}%
\pgfpathlineto{\pgfqpoint{0.835189in}{3.197778in}}%
\pgfusepath{stroke}%
\end{pgfscope}%
\begin{pgfscope}%
\pgfsetbuttcap%
\pgfsetroundjoin%
\definecolor{currentfill}{rgb}{0.000000,0.000000,0.000000}%
\pgfsetfillcolor{currentfill}%
\pgfsetlinewidth{0.803000pt}%
\definecolor{currentstroke}{rgb}{0.000000,0.000000,0.000000}%
\pgfsetstrokecolor{currentstroke}%
\pgfsetdash{}{0pt}%
\pgfsys@defobject{currentmarker}{\pgfqpoint{0.000000in}{-0.048611in}}{\pgfqpoint{0.000000in}{0.000000in}}{%
\pgfpathmoveto{\pgfqpoint{0.000000in}{0.000000in}}%
\pgfpathlineto{\pgfqpoint{0.000000in}{-0.048611in}}%
\pgfusepath{stroke,fill}%
}%
\begin{pgfscope}%
\pgfsys@transformshift{0.835189in}{0.502778in}%
\pgfsys@useobject{currentmarker}{}%
\end{pgfscope}%
\end{pgfscope}%
\begin{pgfscope}%
\definecolor{textcolor}{rgb}{0.000000,0.000000,0.000000}%
\pgfsetstrokecolor{textcolor}%
\pgfsetfillcolor{textcolor}%
\pgftext[x=0.835189in,y=0.405556in,,top]{\color{textcolor}\rmfamily\fontsize{10.000000}{12.000000}\selectfont \(\displaystyle {0}\)}%
\end{pgfscope}%
\begin{pgfscope}%
\pgfpathrectangle{\pgfqpoint{0.713889in}{0.502778in}}{\pgfqpoint{4.805000in}{2.695000in}}%
\pgfusepath{clip}%
\pgfsetrectcap%
\pgfsetroundjoin%
\pgfsetlinewidth{0.803000pt}%
\definecolor{currentstroke}{rgb}{0.690196,0.690196,0.690196}%
\pgfsetstrokecolor{currentstroke}%
\pgfsetdash{}{0pt}%
\pgfpathmoveto{\pgfqpoint{1.580014in}{0.502778in}}%
\pgfpathlineto{\pgfqpoint{1.580014in}{3.197778in}}%
\pgfusepath{stroke}%
\end{pgfscope}%
\begin{pgfscope}%
\pgfsetbuttcap%
\pgfsetroundjoin%
\definecolor{currentfill}{rgb}{0.000000,0.000000,0.000000}%
\pgfsetfillcolor{currentfill}%
\pgfsetlinewidth{0.803000pt}%
\definecolor{currentstroke}{rgb}{0.000000,0.000000,0.000000}%
\pgfsetstrokecolor{currentstroke}%
\pgfsetdash{}{0pt}%
\pgfsys@defobject{currentmarker}{\pgfqpoint{0.000000in}{-0.048611in}}{\pgfqpoint{0.000000in}{0.000000in}}{%
\pgfpathmoveto{\pgfqpoint{0.000000in}{0.000000in}}%
\pgfpathlineto{\pgfqpoint{0.000000in}{-0.048611in}}%
\pgfusepath{stroke,fill}%
}%
\begin{pgfscope}%
\pgfsys@transformshift{1.580014in}{0.502778in}%
\pgfsys@useobject{currentmarker}{}%
\end{pgfscope}%
\end{pgfscope}%
\begin{pgfscope}%
\definecolor{textcolor}{rgb}{0.000000,0.000000,0.000000}%
\pgfsetstrokecolor{textcolor}%
\pgfsetfillcolor{textcolor}%
\pgftext[x=1.580014in,y=0.405556in,,top]{\color{textcolor}\rmfamily\fontsize{10.000000}{12.000000}\selectfont \(\displaystyle {10}\)}%
\end{pgfscope}%
\begin{pgfscope}%
\pgfpathrectangle{\pgfqpoint{0.713889in}{0.502778in}}{\pgfqpoint{4.805000in}{2.695000in}}%
\pgfusepath{clip}%
\pgfsetrectcap%
\pgfsetroundjoin%
\pgfsetlinewidth{0.803000pt}%
\definecolor{currentstroke}{rgb}{0.690196,0.690196,0.690196}%
\pgfsetstrokecolor{currentstroke}%
\pgfsetdash{}{0pt}%
\pgfpathmoveto{\pgfqpoint{2.324838in}{0.502778in}}%
\pgfpathlineto{\pgfqpoint{2.324838in}{3.197778in}}%
\pgfusepath{stroke}%
\end{pgfscope}%
\begin{pgfscope}%
\pgfsetbuttcap%
\pgfsetroundjoin%
\definecolor{currentfill}{rgb}{0.000000,0.000000,0.000000}%
\pgfsetfillcolor{currentfill}%
\pgfsetlinewidth{0.803000pt}%
\definecolor{currentstroke}{rgb}{0.000000,0.000000,0.000000}%
\pgfsetstrokecolor{currentstroke}%
\pgfsetdash{}{0pt}%
\pgfsys@defobject{currentmarker}{\pgfqpoint{0.000000in}{-0.048611in}}{\pgfqpoint{0.000000in}{0.000000in}}{%
\pgfpathmoveto{\pgfqpoint{0.000000in}{0.000000in}}%
\pgfpathlineto{\pgfqpoint{0.000000in}{-0.048611in}}%
\pgfusepath{stroke,fill}%
}%
\begin{pgfscope}%
\pgfsys@transformshift{2.324838in}{0.502778in}%
\pgfsys@useobject{currentmarker}{}%
\end{pgfscope}%
\end{pgfscope}%
\begin{pgfscope}%
\definecolor{textcolor}{rgb}{0.000000,0.000000,0.000000}%
\pgfsetstrokecolor{textcolor}%
\pgfsetfillcolor{textcolor}%
\pgftext[x=2.324838in,y=0.405556in,,top]{\color{textcolor}\rmfamily\fontsize{10.000000}{12.000000}\selectfont \(\displaystyle {20}\)}%
\end{pgfscope}%
\begin{pgfscope}%
\pgfpathrectangle{\pgfqpoint{0.713889in}{0.502778in}}{\pgfqpoint{4.805000in}{2.695000in}}%
\pgfusepath{clip}%
\pgfsetrectcap%
\pgfsetroundjoin%
\pgfsetlinewidth{0.803000pt}%
\definecolor{currentstroke}{rgb}{0.690196,0.690196,0.690196}%
\pgfsetstrokecolor{currentstroke}%
\pgfsetdash{}{0pt}%
\pgfpathmoveto{\pgfqpoint{3.069663in}{0.502778in}}%
\pgfpathlineto{\pgfqpoint{3.069663in}{3.197778in}}%
\pgfusepath{stroke}%
\end{pgfscope}%
\begin{pgfscope}%
\pgfsetbuttcap%
\pgfsetroundjoin%
\definecolor{currentfill}{rgb}{0.000000,0.000000,0.000000}%
\pgfsetfillcolor{currentfill}%
\pgfsetlinewidth{0.803000pt}%
\definecolor{currentstroke}{rgb}{0.000000,0.000000,0.000000}%
\pgfsetstrokecolor{currentstroke}%
\pgfsetdash{}{0pt}%
\pgfsys@defobject{currentmarker}{\pgfqpoint{0.000000in}{-0.048611in}}{\pgfqpoint{0.000000in}{0.000000in}}{%
\pgfpathmoveto{\pgfqpoint{0.000000in}{0.000000in}}%
\pgfpathlineto{\pgfqpoint{0.000000in}{-0.048611in}}%
\pgfusepath{stroke,fill}%
}%
\begin{pgfscope}%
\pgfsys@transformshift{3.069663in}{0.502778in}%
\pgfsys@useobject{currentmarker}{}%
\end{pgfscope}%
\end{pgfscope}%
\begin{pgfscope}%
\definecolor{textcolor}{rgb}{0.000000,0.000000,0.000000}%
\pgfsetstrokecolor{textcolor}%
\pgfsetfillcolor{textcolor}%
\pgftext[x=3.069663in,y=0.405556in,,top]{\color{textcolor}\rmfamily\fontsize{10.000000}{12.000000}\selectfont \(\displaystyle {30}\)}%
\end{pgfscope}%
\begin{pgfscope}%
\pgfpathrectangle{\pgfqpoint{0.713889in}{0.502778in}}{\pgfqpoint{4.805000in}{2.695000in}}%
\pgfusepath{clip}%
\pgfsetrectcap%
\pgfsetroundjoin%
\pgfsetlinewidth{0.803000pt}%
\definecolor{currentstroke}{rgb}{0.690196,0.690196,0.690196}%
\pgfsetstrokecolor{currentstroke}%
\pgfsetdash{}{0pt}%
\pgfpathmoveto{\pgfqpoint{3.814488in}{0.502778in}}%
\pgfpathlineto{\pgfqpoint{3.814488in}{3.197778in}}%
\pgfusepath{stroke}%
\end{pgfscope}%
\begin{pgfscope}%
\pgfsetbuttcap%
\pgfsetroundjoin%
\definecolor{currentfill}{rgb}{0.000000,0.000000,0.000000}%
\pgfsetfillcolor{currentfill}%
\pgfsetlinewidth{0.803000pt}%
\definecolor{currentstroke}{rgb}{0.000000,0.000000,0.000000}%
\pgfsetstrokecolor{currentstroke}%
\pgfsetdash{}{0pt}%
\pgfsys@defobject{currentmarker}{\pgfqpoint{0.000000in}{-0.048611in}}{\pgfqpoint{0.000000in}{0.000000in}}{%
\pgfpathmoveto{\pgfqpoint{0.000000in}{0.000000in}}%
\pgfpathlineto{\pgfqpoint{0.000000in}{-0.048611in}}%
\pgfusepath{stroke,fill}%
}%
\begin{pgfscope}%
\pgfsys@transformshift{3.814488in}{0.502778in}%
\pgfsys@useobject{currentmarker}{}%
\end{pgfscope}%
\end{pgfscope}%
\begin{pgfscope}%
\definecolor{textcolor}{rgb}{0.000000,0.000000,0.000000}%
\pgfsetstrokecolor{textcolor}%
\pgfsetfillcolor{textcolor}%
\pgftext[x=3.814488in,y=0.405556in,,top]{\color{textcolor}\rmfamily\fontsize{10.000000}{12.000000}\selectfont \(\displaystyle {40}\)}%
\end{pgfscope}%
\begin{pgfscope}%
\pgfpathrectangle{\pgfqpoint{0.713889in}{0.502778in}}{\pgfqpoint{4.805000in}{2.695000in}}%
\pgfusepath{clip}%
\pgfsetrectcap%
\pgfsetroundjoin%
\pgfsetlinewidth{0.803000pt}%
\definecolor{currentstroke}{rgb}{0.690196,0.690196,0.690196}%
\pgfsetstrokecolor{currentstroke}%
\pgfsetdash{}{0pt}%
\pgfpathmoveto{\pgfqpoint{4.559312in}{0.502778in}}%
\pgfpathlineto{\pgfqpoint{4.559312in}{3.197778in}}%
\pgfusepath{stroke}%
\end{pgfscope}%
\begin{pgfscope}%
\pgfsetbuttcap%
\pgfsetroundjoin%
\definecolor{currentfill}{rgb}{0.000000,0.000000,0.000000}%
\pgfsetfillcolor{currentfill}%
\pgfsetlinewidth{0.803000pt}%
\definecolor{currentstroke}{rgb}{0.000000,0.000000,0.000000}%
\pgfsetstrokecolor{currentstroke}%
\pgfsetdash{}{0pt}%
\pgfsys@defobject{currentmarker}{\pgfqpoint{0.000000in}{-0.048611in}}{\pgfqpoint{0.000000in}{0.000000in}}{%
\pgfpathmoveto{\pgfqpoint{0.000000in}{0.000000in}}%
\pgfpathlineto{\pgfqpoint{0.000000in}{-0.048611in}}%
\pgfusepath{stroke,fill}%
}%
\begin{pgfscope}%
\pgfsys@transformshift{4.559312in}{0.502778in}%
\pgfsys@useobject{currentmarker}{}%
\end{pgfscope}%
\end{pgfscope}%
\begin{pgfscope}%
\definecolor{textcolor}{rgb}{0.000000,0.000000,0.000000}%
\pgfsetstrokecolor{textcolor}%
\pgfsetfillcolor{textcolor}%
\pgftext[x=4.559312in,y=0.405556in,,top]{\color{textcolor}\rmfamily\fontsize{10.000000}{12.000000}\selectfont \(\displaystyle {50}\)}%
\end{pgfscope}%
\begin{pgfscope}%
\pgfpathrectangle{\pgfqpoint{0.713889in}{0.502778in}}{\pgfqpoint{4.805000in}{2.695000in}}%
\pgfusepath{clip}%
\pgfsetrectcap%
\pgfsetroundjoin%
\pgfsetlinewidth{0.803000pt}%
\definecolor{currentstroke}{rgb}{0.690196,0.690196,0.690196}%
\pgfsetstrokecolor{currentstroke}%
\pgfsetdash{}{0pt}%
\pgfpathmoveto{\pgfqpoint{5.304137in}{0.502778in}}%
\pgfpathlineto{\pgfqpoint{5.304137in}{3.197778in}}%
\pgfusepath{stroke}%
\end{pgfscope}%
\begin{pgfscope}%
\pgfsetbuttcap%
\pgfsetroundjoin%
\definecolor{currentfill}{rgb}{0.000000,0.000000,0.000000}%
\pgfsetfillcolor{currentfill}%
\pgfsetlinewidth{0.803000pt}%
\definecolor{currentstroke}{rgb}{0.000000,0.000000,0.000000}%
\pgfsetstrokecolor{currentstroke}%
\pgfsetdash{}{0pt}%
\pgfsys@defobject{currentmarker}{\pgfqpoint{0.000000in}{-0.048611in}}{\pgfqpoint{0.000000in}{0.000000in}}{%
\pgfpathmoveto{\pgfqpoint{0.000000in}{0.000000in}}%
\pgfpathlineto{\pgfqpoint{0.000000in}{-0.048611in}}%
\pgfusepath{stroke,fill}%
}%
\begin{pgfscope}%
\pgfsys@transformshift{5.304137in}{0.502778in}%
\pgfsys@useobject{currentmarker}{}%
\end{pgfscope}%
\end{pgfscope}%
\begin{pgfscope}%
\definecolor{textcolor}{rgb}{0.000000,0.000000,0.000000}%
\pgfsetstrokecolor{textcolor}%
\pgfsetfillcolor{textcolor}%
\pgftext[x=5.304137in,y=0.405556in,,top]{\color{textcolor}\rmfamily\fontsize{10.000000}{12.000000}\selectfont \(\displaystyle {60}\)}%
\end{pgfscope}%
\begin{pgfscope}%
\definecolor{textcolor}{rgb}{0.000000,0.000000,0.000000}%
\pgfsetstrokecolor{textcolor}%
\pgfsetfillcolor{textcolor}%
\pgftext[x=3.116389in,y=0.225000in,,top]{\color{textcolor}\rmfamily\fontsize{10.000000}{12.000000}\selectfont Time (s)}%
\end{pgfscope}%
\begin{pgfscope}%
\pgfpathrectangle{\pgfqpoint{0.713889in}{0.502778in}}{\pgfqpoint{4.805000in}{2.695000in}}%
\pgfusepath{clip}%
\pgfsetrectcap%
\pgfsetroundjoin%
\pgfsetlinewidth{0.803000pt}%
\definecolor{currentstroke}{rgb}{0.690196,0.690196,0.690196}%
\pgfsetstrokecolor{currentstroke}%
\pgfsetdash{}{0pt}%
\pgfpathmoveto{\pgfqpoint{0.713889in}{0.625278in}}%
\pgfpathlineto{\pgfqpoint{5.518889in}{0.625278in}}%
\pgfusepath{stroke}%
\end{pgfscope}%
\begin{pgfscope}%
\pgfsetbuttcap%
\pgfsetroundjoin%
\definecolor{currentfill}{rgb}{0.000000,0.000000,0.000000}%
\pgfsetfillcolor{currentfill}%
\pgfsetlinewidth{0.803000pt}%
\definecolor{currentstroke}{rgb}{0.000000,0.000000,0.000000}%
\pgfsetstrokecolor{currentstroke}%
\pgfsetdash{}{0pt}%
\pgfsys@defobject{currentmarker}{\pgfqpoint{-0.048611in}{0.000000in}}{\pgfqpoint{-0.000000in}{0.000000in}}{%
\pgfpathmoveto{\pgfqpoint{-0.000000in}{0.000000in}}%
\pgfpathlineto{\pgfqpoint{-0.048611in}{0.000000in}}%
\pgfusepath{stroke,fill}%
}%
\begin{pgfscope}%
\pgfsys@transformshift{0.713889in}{0.625278in}%
\pgfsys@useobject{currentmarker}{}%
\end{pgfscope}%
\end{pgfscope}%
\begin{pgfscope}%
\definecolor{textcolor}{rgb}{0.000000,0.000000,0.000000}%
\pgfsetstrokecolor{textcolor}%
\pgfsetfillcolor{textcolor}%
\pgftext[x=0.439306in, y=0.577847in, left, base]{\color{textcolor}\rmfamily\fontsize{10.000000}{12.000000}\selectfont red}%
\end{pgfscope}%
\begin{pgfscope}%
\pgfpathrectangle{\pgfqpoint{0.713889in}{0.502778in}}{\pgfqpoint{4.805000in}{2.695000in}}%
\pgfusepath{clip}%
\pgfsetrectcap%
\pgfsetroundjoin%
\pgfsetlinewidth{0.803000pt}%
\definecolor{currentstroke}{rgb}{0.690196,0.690196,0.690196}%
\pgfsetstrokecolor{currentstroke}%
\pgfsetdash{}{0pt}%
\pgfpathmoveto{\pgfqpoint{0.713889in}{1.850278in}}%
\pgfpathlineto{\pgfqpoint{5.518889in}{1.850278in}}%
\pgfusepath{stroke}%
\end{pgfscope}%
\begin{pgfscope}%
\pgfsetbuttcap%
\pgfsetroundjoin%
\definecolor{currentfill}{rgb}{0.000000,0.000000,0.000000}%
\pgfsetfillcolor{currentfill}%
\pgfsetlinewidth{0.803000pt}%
\definecolor{currentstroke}{rgb}{0.000000,0.000000,0.000000}%
\pgfsetstrokecolor{currentstroke}%
\pgfsetdash{}{0pt}%
\pgfsys@defobject{currentmarker}{\pgfqpoint{-0.048611in}{0.000000in}}{\pgfqpoint{-0.000000in}{0.000000in}}{%
\pgfpathmoveto{\pgfqpoint{-0.000000in}{0.000000in}}%
\pgfpathlineto{\pgfqpoint{-0.048611in}{0.000000in}}%
\pgfusepath{stroke,fill}%
}%
\begin{pgfscope}%
\pgfsys@transformshift{0.713889in}{1.850278in}%
\pgfsys@useobject{currentmarker}{}%
\end{pgfscope}%
\end{pgfscope}%
\begin{pgfscope}%
\definecolor{textcolor}{rgb}{0.000000,0.000000,0.000000}%
\pgfsetstrokecolor{textcolor}%
\pgfsetfillcolor{textcolor}%
\pgftext[x=0.377639in, y=1.802847in, left, base]{\color{textcolor}\rmfamily\fontsize{10.000000}{12.000000}\selectfont ared}%
\end{pgfscope}%
\begin{pgfscope}%
\pgfpathrectangle{\pgfqpoint{0.713889in}{0.502778in}}{\pgfqpoint{4.805000in}{2.695000in}}%
\pgfusepath{clip}%
\pgfsetrectcap%
\pgfsetroundjoin%
\pgfsetlinewidth{0.803000pt}%
\definecolor{currentstroke}{rgb}{0.690196,0.690196,0.690196}%
\pgfsetstrokecolor{currentstroke}%
\pgfsetdash{}{0pt}%
\pgfpathmoveto{\pgfqpoint{0.713889in}{3.075278in}}%
\pgfpathlineto{\pgfqpoint{5.518889in}{3.075278in}}%
\pgfusepath{stroke}%
\end{pgfscope}%
\begin{pgfscope}%
\pgfsetbuttcap%
\pgfsetroundjoin%
\definecolor{currentfill}{rgb}{0.000000,0.000000,0.000000}%
\pgfsetfillcolor{currentfill}%
\pgfsetlinewidth{0.803000pt}%
\definecolor{currentstroke}{rgb}{0.000000,0.000000,0.000000}%
\pgfsetstrokecolor{currentstroke}%
\pgfsetdash{}{0pt}%
\pgfsys@defobject{currentmarker}{\pgfqpoint{-0.048611in}{0.000000in}}{\pgfqpoint{-0.000000in}{0.000000in}}{%
\pgfpathmoveto{\pgfqpoint{-0.000000in}{0.000000in}}%
\pgfpathlineto{\pgfqpoint{-0.048611in}{0.000000in}}%
\pgfusepath{stroke,fill}%
}%
\begin{pgfscope}%
\pgfsys@transformshift{0.713889in}{3.075278in}%
\pgfsys@useobject{currentmarker}{}%
\end{pgfscope}%
\end{pgfscope}%
\begin{pgfscope}%
\definecolor{textcolor}{rgb}{0.000000,0.000000,0.000000}%
\pgfsetstrokecolor{textcolor}%
\pgfsetfillcolor{textcolor}%
\pgftext[x=0.280694in, y=3.027847in, left, base]{\color{textcolor}\rmfamily\fontsize{10.000000}{12.000000}\selectfont gentle}%
\end{pgfscope}%
\begin{pgfscope}%
\definecolor{textcolor}{rgb}{0.000000,0.000000,0.000000}%
\pgfsetstrokecolor{textcolor}%
\pgfsetfillcolor{textcolor}%
\pgftext[x=0.225139in,y=1.850278in,,bottom,rotate=90.000000]{\color{textcolor}\rmfamily\fontsize{10.000000}{12.000000}\selectfont Algorithm}%
\end{pgfscope}%
\begin{pgfscope}%
\pgfsetrectcap%
\pgfsetmiterjoin%
\pgfsetlinewidth{0.803000pt}%
\definecolor{currentstroke}{rgb}{0.000000,0.000000,0.000000}%
\pgfsetstrokecolor{currentstroke}%
\pgfsetdash{}{0pt}%
\pgfpathmoveto{\pgfqpoint{0.713889in}{0.502778in}}%
\pgfpathlineto{\pgfqpoint{0.713889in}{3.197778in}}%
\pgfusepath{stroke}%
\end{pgfscope}%
\begin{pgfscope}%
\pgfsetrectcap%
\pgfsetmiterjoin%
\pgfsetlinewidth{0.803000pt}%
\definecolor{currentstroke}{rgb}{0.000000,0.000000,0.000000}%
\pgfsetstrokecolor{currentstroke}%
\pgfsetdash{}{0pt}%
\pgfpathmoveto{\pgfqpoint{5.518889in}{0.502778in}}%
\pgfpathlineto{\pgfqpoint{5.518889in}{3.197778in}}%
\pgfusepath{stroke}%
\end{pgfscope}%
\begin{pgfscope}%
\pgfsetrectcap%
\pgfsetmiterjoin%
\pgfsetlinewidth{0.803000pt}%
\definecolor{currentstroke}{rgb}{0.000000,0.000000,0.000000}%
\pgfsetstrokecolor{currentstroke}%
\pgfsetdash{}{0pt}%
\pgfpathmoveto{\pgfqpoint{0.713889in}{0.502778in}}%
\pgfpathlineto{\pgfqpoint{5.518889in}{0.502778in}}%
\pgfusepath{stroke}%
\end{pgfscope}%
\begin{pgfscope}%
\pgfsetrectcap%
\pgfsetmiterjoin%
\pgfsetlinewidth{0.803000pt}%
\definecolor{currentstroke}{rgb}{0.000000,0.000000,0.000000}%
\pgfsetstrokecolor{currentstroke}%
\pgfsetdash{}{0pt}%
\pgfpathmoveto{\pgfqpoint{0.713889in}{3.197778in}}%
\pgfpathlineto{\pgfqpoint{5.518889in}{3.197778in}}%
\pgfusepath{stroke}%
\end{pgfscope}%
\begin{pgfscope}%
\definecolor{textcolor}{rgb}{0.000000,0.000000,0.000000}%
\pgfsetstrokecolor{textcolor}%
\pgfsetfillcolor{textcolor}%
\pgftext[x=3.116389in,y=3.281111in,,base]{\color{textcolor}\rmfamily\fontsize{12.000000}{14.400000}\selectfont Packet Drops Comparison}%
\end{pgfscope}%
\end{pgfpicture}%
\makeatother%
\endgroup%
}
\caption{Packet drop moments for each algorithm}
\label{fig:drops}
\end{figure}

\textbf{RED} (bottom band, blue)---the densest and most continuous band: points fill virtually the entire simulation time, with only rare small gaps. RED actively drops in 47 out of 58~seconds with a median inter-drop interval of only \textbf{2.3~ms}. The forced drop at \texttt{m\_qAvg >= maxTh} creates a continuous stream of dropped packets each time the average queue exceeds the threshold.

\textbf{ARED} (middle band, orange)---clearly visible drop clusters separated by noticeable pauses. The algorithm is active in 33 out of 58~seconds. The $\mathrm{maxP}$ adaptation periodically reduces drop probability to zero, giving flows a ``break.'' However, when drops occur, their density within a cluster is high (maximum burst---276 drops per second).

\textbf{Gentle RED} (top band, green)---the most pronounced cluster structure with the longest pauses between clusters. The algorithm is active in only 27 out of 58~seconds. Within each cluster, drops are dense (maximum burst---305/s), but between them---extended ``quiet'' intervals.

Drop statistics summary is presented in Table~\ref{tab:drops}.

\begin{table}[H]
\centering
\caption{Packet drop statistics}
\label{tab:drops}
\begin{tabular}{lccccc}
\toprule
Algorithm & Total drops & Rate (drops/s) & Median interval & Active seconds & Max.\ burst \\
\midrule
RED      & 1240 & 21.34 & 2.3 ms  & 47 & 321 \\
ARED     & 1003 & 17.11 & 11.3 ms & 33 & 276 \\
Gentle   & 1026 & 17.75 & 9.1 ms  & 27 & 305 \\
\bottomrule
\end{tabular}
\end{table}

RED generates \textbf{19--21\% more drops} than the competitors. The median inter-drop interval for RED (2.3~ms) is 4--5~times smaller than for ARED and Gentle---RED practically continuously drops packets during congestion periods, while ARED and Gentle operate in ``pulses.''

\section{Discussion}
\label{sec:discussion}

The final comparison of algorithms across all key metrics is presented in Table~\ref{tab:summary}.

\begin{table}[H]
\centering
\caption{Summary comparison of RED, ARED, and Gentle RED algorithms}
\label{tab:summary}
\begin{tabular}{lccc}
\toprule
Metric & RED & ARED & Gentle RED \\
\midrule
Mean throughput & 1.8047 Mbit/s & 1.8039 Mbit/s & 1.7998 Mbit/s \\
Throughput CV & \textbf{18.7\%} & 21.9\% & 24.7\% \\
Mean queue & 7.74 pkt & 9.57 pkt & \textbf{6.40 pkt} \\
Queue std.\ dev. & 7.45 & 11.94 & 8.86 \\
p99 queue & 26 pkt & 49 pkt & 34 pkt \\
Empty queue (\%) & 13.4\% & 22.3\% & \textbf{26.6\%} \\
Total drops & 1240 & \textbf{1003} & 1026 \\
Drop pattern & continuous & clustered & pulsed \\
Median inter-drop interval & 2.3 ms & 11.3 ms & 9.1 ms \\
\bottomrule
\end{tabular}
\end{table}

\textbf{RED} provides the most predictable and stable throughput (CV~$= 18.7\%$) but achieves this at the cost of the highest number of drops (1240) and the highest average queue. The continuous nature of drops synchronizes TCP flows, creating rigid, ``quantized'' congestion management dynamics.

\textbf{ARED} minimizes the number of drops (1003, $-19\%$ vs.\ RED) and least frequently disturbs flows. However, $\mathrm{maxP}$ adaptation creates significant queue instability (std.~$= 11.94$, p99~$= 49$~packets) with periods of excessive buffer growth---potentially undesirable from a delay perspective.

\textbf{Gentle RED} achieves the lowest average queue (6.40~packets) and the highest fraction of empty buffer (26.6\%). The number of drops is comparable to ARED (1026), but they are distributed in the most compact pulses with long pauses between them. The cost is the highest throughput variability (CV~$= 24.7\%$) due to TCP flow desynchronization.

\section{Conclusion}
\label{sec:conclusion}

The conducted experiment showed that all three AQM algorithms---RED, ARED, and Gentle RED---provide practically the same average throughput ($\sim$1.80~Mbit/s, $\sim$90\% of link capacity), but fundamentally differ in the nature of buffer management and the load on the TCP congestion control mechanism.

The key difference lies in the trade-off between predictability and aggressiveness of queue management. RED acts most aggressively: forced drop upon exceeding $\mathrm{max}_{\mathrm{Th}}$ keeps the queue ``in check'' (mean 7.74~packets) and ensures minimum throughput variability (CV~$= 18.7\%$), but at the cost of the highest number of drops (1240) and continuous pressure on TCP flows. ARED, conversely, sacrifices queue stability for reduced drops: $\mathrm{maxP}$ adaptation decreases drops by 19\% compared to RED, but produces wide oscillations (std.~$= 11.94$, p99~$= 49$~packets), unacceptable for delay-sensitive applications. Gentle RED occupies an intermediate position, achieving the lowest average queue (6.40~packets) due to a smoother drop function that desynchronizes TCP flows and reduces buffer occupancy---but increasing throughput variability (CV~$= 24.7\%$).

The obtained results confirm that AQM algorithm selection is determined by specific scenario requirements: there is no single ``best'' solution---each algorithm is optimal for its target metric.

\textbf{Application recommendations:}
\begin{itemize}
  \item When \textbf{minimum buffer delay} is priority---Gentle RED (lowest average queue 6.40~pkt).
  \item When \textbf{minimum drop count} is priority---ARED (19\% fewer drops than RED).
  \item When \textbf{predictable throughput} is priority---RED (lowest CV 18.7\%).
\end{itemize}

\printbibliography

}% end \selectlanguage{english}

\end{document}
